% Options for packages loaded elsewhere
\PassOptionsToPackage{unicode}{hyperref}
\PassOptionsToPackage{hyphens}{url}
\documentclass[
  british,
  a4paper,
  openany]{book}
\usepackage{xcolor}
\usepackage[margin=1in]{geometry}
\usepackage{amsmath,amssymb}
\setcounter{secnumdepth}{5}
\usepackage{iftex}
\ifPDFTeX
  \usepackage[T1]{fontenc}
  \usepackage[utf8]{inputenc}
  \usepackage{textcomp} % provide euro and other symbols
\else % if luatex or xetex
  \usepackage{unicode-math} % this also loads fontspec
  \defaultfontfeatures{Scale=MatchLowercase}
  \defaultfontfeatures[\rmfamily]{Ligatures=TeX,Scale=1}
\fi
\usepackage{lmodern}
\ifPDFTeX\else
  % xetex/luatex font selection
  \setmainfont[]{Noto Serif}
  \setmonofont[]{Noto Sans Mono}
\fi
% Use upquote if available, for straight quotes in verbatim environments
\IfFileExists{upquote.sty}{\usepackage{upquote}}{}
\IfFileExists{microtype.sty}{% use microtype if available
  \usepackage[]{microtype}
  \UseMicrotypeSet[protrusion]{basicmath} % disable protrusion for tt fonts
}{}
\makeatletter
\@ifundefined{KOMAClassName}{% if non-KOMA class
  \IfFileExists{parskip.sty}{%
    \usepackage{parskip}
  }{% else
    \setlength{\parindent}{0pt}
    \setlength{\parskip}{6pt plus 2pt minus 1pt}}
}{% if KOMA class
  \KOMAoptions{parskip=half}}
\makeatother
\usepackage{longtable,booktabs,array}
\newcounter{none} % for unnumbered tables
\usepackage{calc} % for calculating minipage widths
% Correct order of tables after \paragraph or \subparagraph
\usepackage{etoolbox}
\makeatletter
\patchcmd\longtable{\par}{\if@noskipsec\mbox{}\fi\par}{}{}
\makeatother
% Allow footnotes in longtable head/foot
\IfFileExists{footnotehyper.sty}{\usepackage{footnotehyper}}{\usepackage{footnote}}
\makesavenoteenv{longtable}
% definitions for citeproc citations
\NewDocumentCommand\citeproctext{}{}
\NewDocumentCommand\citeproc{mm}{%
  \begingroup\def\citeproctext{#2}\cite{#1}\endgroup}
\makeatletter
 % allow citations to break across lines
 \let\@cite@ofmt\@firstofone
 % avoid brackets around text for \cite:
 \def\@biblabel#1{}
 \def\@cite#1#2{{#1\if@tempswa , #2\fi}}
\makeatother
\newlength{\cslhangindent}
\setlength{\cslhangindent}{1.5em}
\newlength{\csllabelwidth}
\setlength{\csllabelwidth}{3em}
\newenvironment{CSLReferences}[2] % #1 hanging-indent, #2 entry-spacing
 {\begin{list}{}{%
  \setlength{\itemindent}{0pt}
  \setlength{\leftmargin}{0pt}
  \setlength{\parsep}{0pt}
  % turn on hanging indent if param 1 is 1
  \ifodd #1
   \setlength{\leftmargin}{\cslhangindent}
   \setlength{\itemindent}{-1\cslhangindent}
  \fi
  % set entry spacing
  \setlength{\itemsep}{#2\baselineskip}}}
 {\end{list}}
\usepackage{calc}
\newcommand{\CSLBlock}[1]{\hfill\break\parbox[t]{\linewidth}{\strut\ignorespaces#1\strut}}
\newcommand{\CSLLeftMargin}[1]{\parbox[t]{\csllabelwidth}{\strut#1\strut}}
\newcommand{\CSLRightInline}[1]{\parbox[t]{\linewidth - \csllabelwidth}{\strut#1\strut}}
\newcommand{\CSLIndent}[1]{\hspace{\cslhangindent}#1}
\ifLuaTeX
\usepackage[bidi=basic,shorthands=off]{babel}
\else
\usepackage[bidi=default,shorthands=off]{babel}
\fi
\ifPDFTeX
\else
\babelfont{rm}[]{Noto Serif}
\fi
\ifLuaTeX
  \usepackage{selnolig} % disable illegal ligatures
\fi
\setlength{\emergencystretch}{3em} % prevent overfull lines
\providecommand{\tightlist}{%
  \setlength{\itemsep}{0pt}\setlength{\parskip}{0pt}}
\usepackage{array}
\usepackage{tabularx}
\usepackage{morewrites}
\usepackage{imakeidx}
\indexsetup{level=\section*,noclearpage}
\setcounter{tocdepth}{1}
\setcounter{secnumdepth}{1}
\renewcommand{\sectionmark}[1]{\markright{\thesection.\ #1}}
\newif\ifivfirstlangheader
\ivfirstlangheadertrue
\providecommand{\ivlangheader}[2]{%
  \ifivfirstlangheader
    \global\ivfirstlangheaderfalse
  \else
    \addvspace{0.9\baselineskip}%
  \fi
  \nopagebreak[3]\textbf{#1}%
  \if\relax\detokenize{#2}\relax\else\ \textit{(#2)}\fi
  \par\nobreak\vspace{0.2\baselineskip}\nopagebreak[3]%
}
% Index Verborum form entry (all languages): italicize the form (including any
% single reconstruction asterisk) plus an optional roman variety label.
% A blank second argument renders identically to \emph{#1}.
\providecommand{\iventry}[2]{%
  \emph{#1}%
  \if\relax\detokenize{#2}\relax\else\nobreakspace\textnormal{(#2)}\fi
}
% Predicted (counterfactual) form marker: superscript dagger immediately before the form
\providecommand{\Pred}[1]{\textsuperscript{\dag}\emph{#1}}
% Reconstructed lexical form: asterisk + italic (asterisk supplied by macro)
\providecommand{\Recon}[1]{\emph{*#1}}
% U+2192 RIGHTWARDS ARROW — Noto Serif lacks this glyph; remap to math-mode arrow
{\catcode`→=\active \global\def→{\ensuremath{\rightarrow}}}
\catcode`→=\active
% Preserve orthographic angle notation with the existing Unicode math font.
{\catcode`⟨=\active \global\def⟨{\ensuremath{\langle}}}
\catcode`⟨=\active
{\catcode`⟩=\active \global\def⟩{\ensuremath{\rangle}}}
\catcode`⟩=\active
\providecommand{\ivlangheader}[2]{\textbf{#1}\if\relax\detokenize{#2}\relax\else\ \textit{(#2)}\fi}
\providecommand{\iventry}[2]{\emph{#1}\if\relax\detokenize{#2}\relax\else\nobreakspace\textnormal{(#2)}\fi}
\makeindex[name=iv,title={},columns=3]
\usepackage{fvextra}
\providecommand{\texorpdfstring}[2]{#1}
\DeclareRobustCommand{\CAPRRuleHeading}[2]{%
  #1\texorpdfstring{\\(#2)}{ (#2)}%
}
\DeclareRobustCommand{\CAPRHeadingBreak}{%
  \texorpdfstring{\\}{ }%
}
\DefineVerbatimEnvironment{ReaderFacingFoma}{Verbatim}{
  breaklines=true,
  breakanywhere=true,
  fontsize=\small,
  frame=single,
  framesep=2mm
}
% Predicted (counterfactual) form marker: superscript dagger immediately before the form
\providecommand{\Pred}[1]{\textsuperscript{\dag}\emph{#1}}
% Reconstructed lexical form: asterisk + italic (asterisk supplied by macro)
\providecommand{\Recon}[1]{\emph{*#1}}
% U+2192 RIGHTWARDS ARROW — Noto Serif lacks this glyph; remap to math-mode arrow
{\catcode`→=\active \global\def→{\ensuremath{\rightarrow}}}
\catcode`→=\active
% Preserve orthographic angle notation with the existing Unicode math font.
{\catcode`⟨=\active \global\def⟨{\ensuremath{\langle}}}
\catcode`⟨=\active
{\catcode`⟩=\active \global\def⟩{\ensuremath{\rangle}}}
\catcode`⟩=\active
\usepackage{bookmark}
\IfFileExists{xurl.sty}{\usepackage{xurl}}{} % add URL line breaks if available
\urlstyle{same}
\hypersetup{
  pdflang={en-GB},
  hidelinks,
  pdfcreator={LaTeX via pandoc}}

\author{}
\date{}

\begin{document}
\frontmatter

{
\setcounter{tocdepth}{0}
\tableofcontents
}
\mainmatter
\mainmatter

\chapter*{Introduction}\label{introduction}
\addcontentsline{toc}{chapter}{Introduction}

\section*{From sound law to
derivation}\label{from-sound-law-to-derivation}
\addcontentsline{toc}{section}{From sound law to derivation}

Historical linguists ordinarily test an etymology by carrying a
reconstructed form through the sound changes that separate it from its
alleged reflex. Most such derivations remain implicit. A scholar knows
that Proto-Germanic \emph{*p} yields Old English \emph{f}, that West
Germanic \emph{*z} became \emph{r} under the appropriate conditions, and
that one change must have preceded another because the reverse order
produces the wrong form. For a single word this mental arithmetic
presents little difficulty. Across hundreds of words and scores of
interacting changes it becomes treacherous. Each step may look familiar
while the derivation as a whole is false.

In computer-assisted phonological reconstruction (CAPR), I represent
each sound law as a finite-state transducer and compose the transducers
in the order argued for below. The resulting cascade applies every
change without discretion: the same rule, with the same environment,
reaches every eligible form. A proposed reconstruction either yields the
Old English comparator or it does not. Moving a rule may repair one
derivation and destroy another. Both results deserve attention, for a
failure locates the point at which phonology ceases to suffice and
philology must decide among analogy, borrowing, dialect mixture,
uncertain attestation, or a faulty reconstruction.

No new principle is involved. The Neogrammarians already demanded
exceptionless sound laws and an account of apparent exceptions
(\citeproc{ref-OsthoffBrugmann1881}{Osthoff and Brugmann 1878-\/-1910}).
Their formula was the \emph{Ausnahmslosigkeit der Lautgesetze}.
Formalization merely enforces that demand mechanically: every eligible
form bears the stated consequences. The computer cannot decide which
reconstruction is historically defensible, but it prevents the
investigator from forgetting what that reconstruction entails.

I apply this method to the development of Proto-Germanic and early West
Germanic forms into Old English. Germanic makes a severe test. Its
historical grammar rests on two centuries of philological labor, while
Old English offers abundant but orthographically and dialectally varied
testimony (\citeproc{ref-Campbell1959}{Campbell 1959};
\citeproc{ref-HoggGrammar1992}{Hogg 1992a};
\citeproc{ref-RingeTaylor2014}{Ringe and Taylor 2014};
\citeproc{ref-Fulk2018}{Fulk 2018}). A computational account cannot
plead scarcity of evidence. It must reproduce familiar developments,
identify the evidence for their order, and explain why its input
sometimes differs from the headword printed in an etymological
dictionary.

\section*{The formal claim}\label{the-formal-claim}
\addcontentsline{toc}{section}{The formal claim}

A sound change defines a relation between strings. For each etymology I
supply a reconstructed input to an ordered series of such relations and
compare the result with an Old English form. Kaplan and Kay demonstrated
that the familiar rewrite rules of phonology admit a finite-state
interpretation (\citeproc{ref-KaplanKay1994}{Kaplan and Kay 1994}); Foma
gives that interpretation executable form
(\citeproc{ref-Hulden2009}{Hulden 2009}). A toy rule deleting final
\emph{*z} may be written:

\noindent\begin{minipage}{\linewidth}
\begin{ReaderFacingFoma}
define ToyFinalZLoss [{*z} -> 0 || _ .#.];
\end{ReaderFacingFoma}
\end{minipage}

The notation states a target, an output, and an environment. Actual
rules also refer to segment classes, stress, syllable weight, word
boundaries, and the output of earlier changes. A loose prose formulation
can conceal choices that code must resolve. Does the rule affect every
vowel or only short vowels? Does a following consonant block it? Does it
apply before or after the loss of a final syllable? A transducer forces
each question into the open.

Order gives the cascade its historical content. An early change may
create the environment for a later one and thus feed it; alternatively,
it may remove that environment and bleed it. A word affected by both
changes can therefore establish their relative chronology. Other changes
commute across the available lexicon: reversing them alters no output in
the active dataset. Such a result does not prove simultaneity or
historical indifference. It shows only that the present evidence fails
to order them. Throughout this book I distinguish chronology compelled
by the tested forms from placement adopted on the authority of the
historical grammars.

Backward reconstruction requires a further restriction. An unrestricted
inverse transducer will propose formally possible strings that no
Germanic language could have inherited. I therefore restrict backward
reconstruction with an inventory and a statement of permissible
ancestral forms. Reconstruction always combines correspondences with a
theory of what could count as a word in the \emph{Grundsprache}; this
restriction states that theory rather than leaving it tacit.

\section*{Inputs, targets, and
success}\label{inputs-targets-and-success}
\addcontentsline{toc}{section}{Inputs, targets, and success}

I compare a selected earlier Germanic form with a selected Old English
target. Neither selection is innocent. Dictionaries cite lexemes, but
sound change operates on word-forms. The ancestor of an Old English
plural, preterite, or oblique case may differ from the reconstructed
lemma in precisely the material on which a later sound law acts. Kroonen
and Orel provide indispensable lexical reconstructions, while the
grammars often supply the paradigm history needed to choose the actual
input (\citeproc{ref-Orel2003}{Orel 2003};
\citeproc{ref-Kroonen2013}{Kroonen 2013};
\citeproc{ref-RingeTaylor2014}{Ringe and Taylor 2014}).

I therefore distinguish the citation reconstruction from the transducer
input. The former identifies the etymon; the latter represents the
paradigm cell or remodeled stem whose history is at issue. Where the two
differ, the lexical entry states the difference and argues for it. This
prevents a convenient input from masquerading as a received
reconstruction. An unmotivated alteration made solely to secure the
desired output would empty the exercise of historical meaning.

The distinction is concrete in Old English
{\emph{sċuldrum}}\index[iv]{02oe@\ivlangheader{Old English}{West Saxon normalization unmarked}!sculdrum@\iventry{sċuldrum}{}}
`shoulders, dative plural'. Its inflectional history requires an input
of the type
\Recon{skúldramiz}\index[iv]{03pgmc@\ivlangheader{Proto-Germanic}{}!skuldramiz@\iventry{*skúldramiz}{}}
`shoulders, reconstructed dative-plural input'; the singular headword
represents a different paradigm cell.

The target also requires judgment. Old English spelling varies by date,
dialect, manuscript, and editorial practice. The cascade may produce an
internal phonological symbol that an orthographic transducer then maps
to a normalized written form. A string match at this final stage cannot
by itself establish an etymology; conversely, a superficial spelling
mismatch need not disprove one. I treat phonological development and
orthographic normalization as separate operations so that exact
computation does not confuse notation with history.

Within these limits a successful derivation has three senses. It
succeeds formally when the output string matches the target. It succeeds
philologically when the chosen input and comparator are the proper forms
to compare. It succeeds historically when the proposed path agrees with
the wider Germanic evidence. The first kind of success is cheap. The
argument of the book concerns the conjunction of all three.

\section*{The evidence of failure}\label{the-evidence-of-failure}
\addcontentsline{toc}{section}{The evidence of failure}

Regular sound change makes irregularity legible. If an inherited form
refuses to pass through an otherwise successful cascade, the mismatch
demands a name. Analogy may have replaced the expected reflex with a
form drawn from another paradigm cell. Borrowing may have introduced the
word after the relevant changes. A dialectal form may lie outside the
modeled West Saxon path. The target may be late, corrupt, or normalized
beyond what the manuscript evidence warrants. Finally, either the
reconstruction or the rule may be wrong.

These possibilities should not be suppressed by narrow, lexeme-specific
``sound laws.'' A system that derives every target by multiplying
exceptions has only encoded its answers. I instead distinguish several
classes of non-regular result: attested variants, early and late
analogy, reconstructed Old English comparators, known but unmodelled
phonological and morphological developments, and unexplained exceptions.
The categories are claims, not housekeeping labels. They identify where
regular phonology ends and what additional history the evidence
requires.

This treatment follows the original Burmish CAPR work, in which
resistant forms often disclosed loans or mistaken cognate assignments.
Old English shifts the balance toward morphology and analogy, but the
methodological advantage remains the same. Failure concentrates inquiry.
It tells us which assumption---input, target, environment, order, or
lexical history---must bear the explanation.

\section*{Evidence for relative
chronology}\label{evidence-for-relative-chronology}
\addcontentsline{toc}{section}{Evidence for relative chronology}

The chronology chapters combine three kinds of evidence. First come the
statements of the standard historical grammars. These establish the
received description and often the broad order of developments
(\citeproc{ref-Campbell1959}{Campbell 1959};
\citeproc{ref-HoggGrammar1992}{Hogg 1992a};
\citeproc{ref-RingeTaylor2014}{Ringe and Taylor 2014};
\citeproc{ref-Fulk2018}{Fulk 2018}). Second come individual witness
words. A derivation that succeeds under one order and fails under the
reverse order supplies direct lexical evidence for that relation. Third
come exhaustive order tests across the active dataset. These reveal
whether an apparently decisive relation is local, whether other words
contradict it, and how far a rule can move without disturbing any
output.

Suppose that a consonant change creates the environment for a later
vowel change. Under the received order both apply and the Old English
target emerges. Reverse them and the vowel change misses its
environment. The word then supports the priority of the consonant
change. The converse case is equally informative: if an early rule
creates a segment that a later rule would wrongly alter, the creating
rule must follow the other. When reversal changes no output, the lexical
evidence leaves the order open even if philological considerations still
favor one placement.

A note on notation: a superscript dagger placed immediately before an
italicized form marks a concrete predicted output generated by a
counterfactual manipulation of the rule order (for example, moving a
change earlier or later than argued). This dagger indicates that the
marked form is a model prediction, not an independently attested
historical form; reconstructions remain starred in the usual way.

Sims-Williams argues for mechanizing precisely this kind of reasoning
(\citeproc{ref-SimsWilliams2018}{Sims-Williams 2018}). Computation does
not replace the historical argument; it makes the extent of that
argument measurable. A traditional chronology may prove correct but less
tightly constrained than its customary presentation suggests. A relation
described as local may in fact rest only on a broad terminus. Such
negative results are salutary. They separate what the data demonstrate
from what a convenient exposition merely presupposes.

\section*{Rules and words}\label{rules-and-words}
\addcontentsline{toc}{section}{Rules and words}

The book accordingly moves twice through the same history. Part I begins
with the ordered rules. For each development it states the historical
problem, gives the rule in formal notation, and examines the evidence
for placement. The code appears because it is part of the claim, but the
code-name is never an explanation. \texttt{OEIUmlaut}, for example, is
only a rule name; its linguistic content lies in the stated environment,
the historical discussion, and the words whose derivations depend upon
it.

Part II begins with the words. Each entry identifies the reconstruction,
the selected input, the Old English comparator, and the derivational
class. A trace shows which rules altered the form; the accompanying
prose explains the philological decisions that the trace cannot make.
Short entries record ordinary derivations. Longer ones treat doubtful
reconstructions, variant attestations, paradigm-cell selection, or
analogy. The asymmetry is intentional: equal database rows do not
warrant equal historical discussion.

The reader can thus move in either direction. A chronology chapter names
the lexical witnesses that constrain a rule; their entries display the
complete derivations. A lexical entry invokes a change; the
corresponding chapter explains its formulation and place in the cascade.
The index verborum provides a third route through the material,
gathering reconstructed and attested forms by language.

\section*{Reproducibility and
disagreement}\label{reproducibility-and-disagreement}
\addcontentsline{toc}{section}{Reproducibility and disagreement}

An executable derivation identifies the exact point of disagreement. One
reader may accept a sound change but reject its environment; another may
accept the rule and dispute its order; a third may object to the
selected paradigm cell or to the normalization of the Old English
target. Each objection addresses a recorded decision. Given the same
inputs, rules, and order, the stated outputs follow. Reproducibility
here concerns those consequences, not the surrender of philological
judgment to an algorithm.

Prose can sound settled while concealing several incompatible
derivations. Code is less accommodating. A rule that is too broad
damages words it was never designed to explain; one that is too narrow
leaves legitimate reflexes untouched. A reconstruction plausible in
isolation may fail after three later changes. The resulting
embarrassment is a virtue of the method. Applying every rule to the
whole lexicon subjects local analyses to a global test.

The converse danger is false precision. A deterministic cascade may
tempt the reader to mistake exact strings for exact history. Every
result still depends on choices about segmentation, symbol inventory,
reconstruction, morphology, dialect, chronology, and orthography. I
therefore cite the philological sources, display the selected forms,
preserve unresolved exceptions, and distinguish tested order from
source-based order. Those choices remain accountable to philology.

\section*{The argument}\label{the-argument}
\addcontentsline{toc}{section}{The argument}

A book must advance an argument. Mine is that traditional rule-based
reconstruction becomes clearer when every proposed derivation can be
executed, every rule must face the whole lexicon, and every failure is
reported. The Germanic-to-Old-English case shows both the power and the
boundary of that claim. Much of the history admits a coherent ordered
account. The residue does not disappear: it resolves into morphology,
analogy, variation, borrowing, imperfect attestation, and a small number
of genuine problems.

The sound-change chapters ask how tightly the lexical evidence
constrains the cascade. The lexical chapters ask whether individual
etymologies survive that cascade without philological sleight of hand.
Neither inquiry suffices alone. A formal system that reaches the targets
by hiding doubtful inputs is bad historical linguistics; elegant prose
that never reproduces its derivations is an incomplete formal account.
Executable derivations make those two standards answer to one another.

\part{Sound changes, formalization, and relative chronology}

\chapter*{Sound-change overview}\label{sound-change-overview}
\addcontentsline{toc}{chapter}{Sound-change overview}

\section*{Scope and orientation}\label{scope-and-orientation}
\addcontentsline{toc}{section}{Scope and orientation}

The sequence follows the executable reconstruction from its early
Germanic developments through the Old English written surface.

Rhotacism, brightening, breaking, umlaut, and apocope alternate with
narrowly conditioned changes whose relative order rests on particular
witness words.

The evidence ranges from broadly attested sound laws to lexical
constraints that establish only one chronological boundary.

\section*{Numbering note}\label{numbering-note}
\addcontentsline{toc}{section}{Numbering note}

SC numbers remain stable legacy identifiers, not chronological
positions. Each published entry has one authoritative source. The book
automatically follows the individual stages of the canonical executable
reconstruction; neither filenames nor chapter membership can override
that order.

This is the order adopted by the executable reconstruction, not a claim
that every pairwise historical chronology has been independently proved.
Adjudication records distinguish demonstrated chronology, stage-entailed
relations and computational serialization. Chapter boundaries and titles
remain editorial; they divide the sequence without rearranging its
entries.

Explicitly retained technical explanations are published at their
executable positions without becoming adjudicated historical sound laws.

SC038, SC062, and SC084 mark technical or prosodic stages rather than
sound changes; SC113 annotates stress-foot context and is not a
published law. SC077 is unused.

\newpage

\chapter{From Proto-Germanic to Proto-Northwest
Germanic}\label{from-proto-germanic-to-proto-northwest-germanic}

\section{Historical interval}\label{historical-interval}

This chapter covers developments that took place within the
Proto-Germanic period, from the inherited consonant system to the first
changes that separate the Northwest Germanic line from the rest of the
Germanic family.

The reconstruction labelled Proto-Germanic here is the common ancestor
of Gothic, North Germanic, and West Germanic, reconstructed through the
classical comparative method from attested descendant languages. The
label Proto-Northwest Germanic designates the hypothetical node linking
the ancestors of North Germanic (Old Norse and its relatives) and West
Germanic (Old English, Old High German, Old Saxon, and Old Frisian among
others), to the exclusion of Gothic and the other East Germanic
varieties.

\section{What this chapter contains}\label{what-this-chapter-contains}

Chapter 1 contains one reader-facing sound-change section: the
Proto-Germanic loss of a nasal before \emph{*x}
(\texttt{SC103\ PGmcNasalLossBeforeX}), with compensatory lengthening
and nasalization of the preceding vowel. Its results are shared by every
daughter language, Gothic included, so the change belongs to
Proto-Germanic itself and precedes every Northwest Germanic and West
Germanic development treated in the chapters that follow
(\citeproc{ref-Campbell1959}{Campbell 1959, 44}, §119;
\citeproc{ref-Fulk2018}{Fulk 2018, 55}, §4.1;
\citeproc{ref-Ringe2017}{Ringe 2017, 149--50}, §3.2.7).

Another historically Proto-Germanic rule, the positional allophony of
Proto-Germanic \emph{*b} (\texttt{SC049\ PGmcBAllophony}), executes late
in the computational cascade because the alternation interacts with
environments shaped by intermediate rule applications; its section
therefore appears in Chapter 4, where the divergence between cascade
placement and historical stage is noted explicitly.

One other historically Proto-Germanic change, Gm-simplification
(\texttt{SC002\ PGmcGmSimplification}), is documented in the book-entry
plan and its literature dossier confirms the source base is narrow (two
lexical families:
\Recon{draugma-}\index[iv]{03pgmc@\ivlangheader{Proto-Germanic}{}!draugma@\iventry{*draugma-}{}}
`dream' and
\Recon{taugma-}\index[iv]{03pgmc@\ivlangheader{Proto-Germanic}{}!taugma@\iventry{*taugma-}{}}
`team'; (\citeproc{ref-Kroonen2013}{Kroonen 2013, 101, 511})). A
reader-facing section for SC002 awaits a stronger explanatory source
base and is not yet assembled in the reader-facing sequence.

A generative counterpart to the handbook accounts is Barrack's formal
cascade from Proto-Germanic to West Saxon, which states the nasal-loss
development as its Proto-Germanic rule 119
(\citeproc{ref-Barrack1975}{Barrack 1975, 42--43}). Barrack's rule
numbers are anchors to Campbell's paragraph numbering, which he names as
his principal data source (\citeproc{ref-Barrack1975}{Barrack 1975,
21}), and his ordering statements are formal, not chronological:
ordering between rules ``only implies a logical relationship---in no
sense a chronological relationship'' (\citeproc{ref-Barrack1975}{Barrack
1975, 34}). His cascade therefore corroborates the handbook conditioning
for this change without constituting independent dating evidence, and it
keeps the Proto-Germanic law distinct from the later Ingvaeonic
nasal-spirant development, which he states separately as rule 121
(\citeproc{ref-Barrack1975}{Barrack 1975, 44}).

\section{Scope and genealogical
context}\label{scope-and-genealogical-context}

Changes in this chapter are pan-Germanic in scope: they apply to
ancestral forms that feed both the North Germanic and West Germanic
descendants, or they represent internal Proto-Germanic processes visible
across the Germanic family.

The boundary between Proto-Germanic and Proto-Northwest Germanic is not
sharp in the textbook literature. Ringe and Taylor treat many of the
traditionally ``Proto-Germanic'' changes as part of a shared innovation
package that is diagnostically older than North--West Germanic
divergence but not necessarily earlier than the separation of the East
Germanic line (\citeproc{ref-RingeTaylor2014}{Ringe and Taylor 2014,
1--30}). For book purposes, the distinction matters primarily because it
separates the features inherited uniformly from all Germanic from those
shared selectively by North and West Germanic to the exclusion of
Gothic.

\section{A note on the rule names}\label{a-note-on-the-rule-names}

The CAPR rules implemented in this chapter carry names beginning with
\texttt{PGmc}. Those names are intended as stable internal identifiers,
not as claims about the precise historical stage of every rule so
labelled. A rule named \texttt{PGmcX} may in some cases be a later
development that affects only the West Germanic or Northwest Germanic
branch. Each entry's source-supported historical analysis takes priority
over both its identifier and its chapter placement; the latter follows
executable order, not an independent historical-stage map.

\section{\texorpdfstring{Proto-Germanic loss of a nasal before
\emph{*x}}{Proto-Germanic loss of a nasal before }}\label{proto-germanic-loss-of-a-nasal-before}

\subsection{Historical discussion}\label{historical-discussion}

The oldest change treated in this book is common to the whole family. In
the group \Recon{-nx-} the nasal consonant was lost, the preceding vowel
was lengthened in compensation, and that lengthened vowel was nasalized.
Its results are shared by every daughter language: Gothic \emph{þeihan}
`thrive', \emph{brāhta} `brought', \emph{þūhta} `seemed' stand beside
Old High German \emph{dīhan} `thrive', \emph{brāhta}, \emph{fūht} `damp'
and Old English \emph{þēon} `thrive', \emph{brōhte} `brought',
\emph{þūhte} `seemed', \emph{fūht} `damp'
(\citeproc{ref-Campbell1959}{Campbell 1959, 44}, §119;
\citeproc{ref-Fulk2018}{Fulk 2018, 55}, §4.1;
\citeproc{ref-Ringe2017}{Ringe 2017, 149--50}, §3.2.7). Because no
daughter keeps the nasal, the change belongs to Proto-Germanic itself,
and it precedes every Northwest Germanic, West Germanic, North Sea
Germanic and Anglo-Frisian development described in the chapters that
follow.

Only \Recon{a}, \Recon{i} and \Recon{u} occur in this position. Germanic
had already raised \Recon{e} to \Recon{i} and \Recon{o} to \Recon{u}
before a nasal followed by a consonant, so the mid vowels are absent
from the input (\citeproc{ref-Fulk2018}{Fulk 2018, 55}, §4.1).

The vowel that the change creates is long and nasalized, and it is not
yet rounded. Fulk emphasizes that the lengthened vowels remained
nasalized for a considerable time, well past the close of the Northwest
Germanic period, since the nasalized low vowel produced in this way went
on to develop to \emph{ō} in Anglo-Frisian and did not fall together
with the Old English \emph{ā} that came from \Recon{ai}
(\citeproc{ref-Fulk2018}{Fulk 2018, 55}, §4.1). The comparative material
makes the same point directly: Gothic, Old Norse, Old High German and
Old Saxon all reflect the Proto-Germanic nasalized low vowel as
unrounded \emph{ā}, and only Old English and Old Frisian show \emph{ō}
(\citeproc{ref-Campbell1959}{Campbell 1959, 44}, §119;
\citeproc{ref-Fulk2018}{Fulk 2018, 72}, §4.11). Ringe describes
Proto-Germanic \emph{*hanhaną} as \emph{*[xą̄xaną]} and observes that its
low vowel was rounded, along with the other nasalized low vowels, in the
northernmost West Germanic dialects (\citeproc{ref-Ringe2017}{Ringe
2017, 149--50}, §3.2.7). That rounding is a separate and much later
change, and it is treated in the chapter on the rounding of the long
nasalized low vowel.

Barrack's formal cascade states this development as his Proto-Germanic
rule 119 (\citeproc{ref-Barrack1975}{Barrack 1975, 42--43}) and, in
agreement with the comparative argument above, keeps it distinct from
the later North Sea law, which he states separately as Ingvaeonic rule
121 and explicitly orders after 119 (\citeproc{ref-Barrack1975}{Barrack
1975, 44}). His rule numbers are anchors to Campbell's paragraph
numbering (\citeproc{ref-Barrack1975}{Barrack 1975, 21}), and his
orderings are formal rather than chronological
(\citeproc{ref-Barrack1975}{Barrack 1975, 34}), so the agreement
concerns the identity and conditioning of the two laws, not an
independent dating.

\subsection{\texorpdfstring{SC103. Proto-Germanic nasal loss before
\emph{*x}
(\texttt{PGmcNasalLossBeforeX})}{SC103. Proto-Germanic nasal loss before  (PGmcNasalLossBeforeX)}}\label{rule-PGmcNasalLossBeforeX}

\begin{ReaderFacingFoma}
define PGmcNasalLossBeforeX Ctx([
    [
    {*a} -> {*ã} || _ EnglishStarNasal {*x},
    {*i} -> {*ī} || _ EnglishStarNasal {*x},
    {*u} -> {*ū} || _ EnglishStarNasal {*x},
    {*á} -> {*ã} || _ EnglishStarNasal {*x},
    {*í} -> {*ī} || _ EnglishStarNasal {*x},
    {*ú} -> {*ū} || _ EnglishStarNasal {*x}
] .o. [
    EnglishStarNasal -> 0 || _ {*x}
]
]);
\end{ReaderFacingFoma}

The rule performs the three parts of the change together: it lengthens
the vowel, it marks the lengthened low vowel as nasalized, and it then
removes the conditioning nasal. The low vowel is written \Recon{ą̄} and
the high vowels are written \Recon{ī} and \Recon{ū}, because the
nasality of the high vowels has no further consequence: they go on to
develop exactly as the inherited long \Recon{ī} and \Recon{ū} do
(\citeproc{ref-Campbell1959}{Campbell 1959, 47}, §121). The nasality of
the low vowel is carried forward because its later fate depends on it.
In the executable definitions, \texttt{ã} is the single-symbol encoding
of long nasalized \emph{*ą̄}, not a different or short vowel. The
historical notation in the discussion retains its explicit length mark.

The corpus witnesses both branches of the rule. The high branch is
\Recon{fúnxstiz} `fist', which becomes \Recon{fū́xsti} `fist' and, after
the loss of \Recon{x} before the cluster in
\hyperref[rule-PNWGmcPreconsonantalXLoss]{SC028
PNWGmcPreconsonantalXLoss}, gives Old English \emph{fȳst} `fist'. The
long vowel of Old High German \emph{fūst} `fist', Dutch \emph{vuist} and
German \emph{Faust} shows that the word belongs here and to no later law
(\citeproc{ref-Kroonen2013}{Kroonen 2013, 160}). The rule also supplies
the \Recon{xst} cluster on which the loss of \Recon{x} before a
consonant operates, so the two stand in a feeding relation.

The low branch is \Recon{θánxtē} `thought', the preterite of the verb
`to think', whose principal parts are reconstructed as \Recon{þankijaną}
`to think', \Recon{þanhtē} `thought', \Recon{þanhtaz} `thought (past
participle)' with the nasal still standing before the fricative
(\citeproc{ref-Ringe2017}{Ringe 2017, 281};
\citeproc{ref-Ringe2017}{2017, 136}). Here the rule yields \Recon{θą̄xtē}
`thought', and the nasalized low vowel is later rounded by
\hyperref[rule-EAFNasalizedLowRounding]{SC104 EAFNasalizedLowRounding}
to give Old English \emph{þōhte} `thought'. This is the form the
handbooks themselves cite for the Anglo-Frisian rounding of the vowel
produced here (\citeproc{ref-Fulk2018}{Fulk 2018, 55}, §4.1). It does
not enter the loss of \Recon{x} before a consonant, in which respect it
differs from `fist', so the fricative survives to the surface and the
word displays the vowel history alone.

\newpage

\chapter{From Proto-Northwest Germanic to Proto-West
Germanic}\label{from-proto-northwest-germanic-to-proto-west-germanic}

\section{Historical interval}\label{historical-interval-1}

This chapter covers the sound changes that took place in the
proto-language shared by the West Germanic languages --- Old English,
Old Frisian, Old Saxon, Old High German, and Old Dutch --- before the
individual languages diverged. The starting reconstruction is
Proto-Northwest Germanic (PNWGmc), the hypothetical common ancestor of
North Germanic and West Germanic together; the ending reconstruction is
Proto-West Germanic (PWGmc), the immediate common ancestor of the West
Germanic languages specifically.

\section{Scope and internal
diversity}\label{scope-and-internal-diversity}

Changes in this chapter are not all equally pan-West-Germanic in scope.
They may be grouped broadly as follows:

Northwest Germanic innovations (shared by both North and West Germanic):
innovations in the unstressed vowel system, certain final-syllable vowel
changes, and selected consonant cluster simplifications. Their
historical scope must be established from comparative evidence, not from
executable identifiers.

Proto-West-Germanic innovations (shared within West Germanic but not in
North Germanic): the cluster of morphological and phonological changes
that distinguish Old English, Old High German, Old Saxon, and Old
Frisian from Old Norse. They include early apocope rules, certain
consonant assimilations, and the West Germanic gemination of consonants
before \texttt{*j}.

\section{Major changes}\label{major-changes}

The chapter opens with the root-noun nominative \texttt{*-z} loss
(SC096), the generalization of endingless nominatives through the
athematic consonant stems, complete before Proto-West Germanic: none of
the West Germanic daughters shows any ending in this class
(\citeproc{ref-RingeTaylor2014}{Ringe and Taylor 2014, 118}). It is the
earliest of the three historically distinct final-\texttt{*z}
developments; the other two (SC020 and SC097) open Chapter 3.

Barrack's formal cascade covers the same final-syllable territory with a
family of rules keyed to Campbell's paragraphs: rhotacism and the loss
and retention of final \emph{*z} appear as rules 404a--e
(\citeproc{ref-Barrack1975}{Barrack 1975, 101--2}), and the
final-syllable reductions as the 331 series, including final \emph{*a}
syncope as West Germanic rule 331.4b (\citeproc{ref-Barrack1975}{Barrack
1975, 78}). He grades the evidence candidly --- the material for 404b,
which deletes the rhotic before a voiced dental with compensatory
lengthening to derive \emph{mēd} `reward', he calls ``very sporadic'',
and he refers the rival development that yields \emph{meord} to his rule
139zc (\citeproc{ref-Barrack1975}{Barrack 1975, 101--2}) --- and his
orderings among these rules are formal rather than chronological
(\citeproc{ref-Barrack1975}{Barrack 1975, 34}). The crosswalk in the
Barrack dossier records where these statements agree with, and where
they subdivide differently from, the SC096/SC020/SC097 split adopted
here.

The unstressed \emph{*ai > *ē} development
(\hyperref[rule-PNWGmcUnstressedAiMonophthongization]{SC014
PNWGmcUnstressedAiMonophthongization}) is one of the widespread
Northwest Germanic vowel shifts. Unstressed dative-singular and
strong-adjective plural endings provide examples of monophthongization
and merger with inherited unstressed \emph{*ē}. Ringe and Taylor
distinguish this broad correspondence from the more questionable claim
of a single completed Proto-Northwest Germanic event
(\citeproc{ref-RingeTaylor2014}{Ringe and Taylor 2014, 25--26}). The
stressed monophthongization is treated separately under
\hyperref[rule-EAFAiMonophthongization]{SC004 EAFAiMonophthongization}.

The later \hyperref[rule-PNWGmcILowering]{SC015 PNWGmcILowering}
concerns stressed first-syllable \emph{*i}, not further unstressed-vowel
reduction. Its following-vowel hypothesis has traditional authority, but
its comparative regularity and precise consonantal conditioning remain
disputed. The section compares vowel-conditioned and
consonant-conditioned accounts and explains the provisional restrictions
of the current implementation (\citeproc{ref-Campbell1959}{Campbell
1959, 43}, §§114--116; \citeproc{ref-Fulk2018}{Fulk 2018, 56--57};
\citeproc{ref-Hill2017}{Hill 2017, 149--58}).

The West Germanic consonant changes of this chapter --- j-gemination
(SC010), early i-apocope (SC006), coronal-w assimilation (SC008), and
related rules --- represent the most productive phonological territory
for the CAPR derivations. They feed a large proportion of the
distinctive consonant clusters of Old English. Handbooks vary in exactly
how they group and name these changes
(\citeproc{ref-Campbell1959}{Campbell 1959, §§ 404, 406};
\citeproc{ref-HoggGrammar1992}{Hogg 1992a, 71}, §4.11 (held 2011
reissue)).

Coronal assimilation also supplies the geminate in \emph{fēower} `four'
and \emph{ēow} `you', before inherited short-vowel-plus-geminate-glide
reanalysis. \hyperref[rule-OEWWSimplification]{SC031 OEWWSimplification}
now represents that earlier development, not an unrestricted late
deletion. The resulting diphthong-plus-glide sequence is inherited by
later English realization; j-created and singleton paths remain distinct
(\citeproc{ref-RingeTaylor2014}{Ringe and Taylor 2014, 41--42, 65--66};
\citeproc{ref-Campbell1959}{Campbell 1959, 45--47}).

The quoted definitions use \texttt{Ctx(...)} to carry selected sentence
context outside the segmental operation. It is computational transport,
not a reconstructed segment or an additional sound law. Lexical accent,
sentence stress and phonological-word finality are separate: the early
high-vowel-loss account requires a heavy syllable in a
sentence-unstressed, phonologically final word, not simply the absence
of an acute. Ordinary evaluation selects strong-final citation context;
the explicit weak-final selection for \emph{ēow} is discussed with
\hyperref[rule-PWGmcUnstressedWordFinalIApocope]{SC098
PWGmcUnstressedWordFinalIApocope} (\citeproc{ref-RingeTaylor2014}{Ringe
and Taylor 2014, 55, 57--58}).

The nasal spirant corridor (SC026--SC027), treated in Chapter 3 at its
cascade position, illustrates a type of change common in historical
grammars of the ``Ingvaeonic'' or ``North Sea Germanic'' area: nasals
disappear before voiceless fricatives, with compensatory vowel
lengthening (\citeproc{ref-Campbell1959}{Campbell 1959, §§ 462--463};
\citeproc{ref-HoggGrammar1992}{Hogg 1992a, 56}, §3.14 (held 2011
reissue)). The CAPR model splits this into two ordered steps to make the
vowel effect computationally tractable; the book prose explains that
split against the handbook tradition, which typically presents the
change as a single process.

Chapters in this part of the book follow the executable reconstruction's
adopted serialization, not an independently demonstrated total
historical chronology. Several rules that carry \texttt{PWGmc} labels
--- final bare-\texttt{*a} loss (SC041), surviving bimoric \texttt{*ō}
unrounding (SC042), and Sievers-law syncope (SC050) --- execute later in
the cascade and are therefore presented in Chapter 4, where their
individual sections discuss their historical stage labels. Conversely,
one rule with a West Saxon label, the palatal-glide rule (SC016), is an
orthographic rule of the written surface: it executes after the Old
English orthography stage and is presented in Chapter 5.

\section{A note on source terminology and
subgrouping}\label{a-note-on-source-terminology-and-subgrouping}

The literature uses several partly overlapping stage labels for this
period:

\begin{itemize}
\tightlist
\item
  Northwest Germanic: the node uniting North and West Germanic.
\item
  Proto-West Germanic: the node uniting only the West Germanic
  languages.
\item
  North Sea Germanic or Ingvaeonic: a proposed subgroup within West
  Germanic covering Old English, Old Frisian, and Old Saxon (and
  sometimes Old Low Franconian), sharing certain innovations over a
  broader area.
\item
  Anglo-Frisian: a narrower proposed subgroup linking only Old English
  and Old Frisian.
\end{itemize}

These labels are not always used consistently across sources. Ringe and
Taylor are cautious about reconstructing a discrete Proto-West-Germanic
node (\citeproc{ref-RingeTaylor2014}{Ringe and Taylor 2014, 50--55}).
Campbell notes that many of the ``West Germanic'' shared features could
alternatively be treated as parallel developments rather than common
inheritance (\citeproc{ref-Campbell1959}{Campbell 1959, §§ 1--5}).

CAPR uses \texttt{PWGmc} and \texttt{NWGmc} as organizing labels for
this chapter without claiming to have settled all questions about West
Germanic subgrouping. Changes that appear in the literature under
``Ingvaeonic'' labels but affect the Old English--to-Proto-Germanic
derivation chain are treated here as late expressions of the same West
Germanic developmental period unless existing CAPR dossier research
specifically argues for Anglo-Frisian or English-specific placement.

\section{Rule names}\label{rule-names}

Most CAPR rules in this chapter carry names beginning with
\texttt{NWGmc} or \texttt{PWGmc}; the reformulated inherited-glide rule
retains its older \texttt{OE} identifier. These names are stable
internal identifiers. A name beginning with \texttt{NWGmc} does not
guarantee that the change is exclusive to Northwest Germanic, and a name
beginning with \texttt{PWGmc} does not guarantee that it is absent from
North Germanic. The historical analysis in each sound-change section
takes priority over the name prefix.

\section{Early Northwest Germanic vowel
developments}\label{early-northwest-germanic-vowel-developments}

These two developments affect different parts of the vowel system.
\hyperref[rule-PNWGmcUnstressedAiMonophthongization]{SC014
PNWGmcUnstressedAiMonophthongization} monophthongizes unstressed
\emph{*ai}; \hyperref[rule-PNWGmcILowering]{SC015 PNWGmcILowering}
represents the disputed lowering of stressed short \emph{*i}. This
thematic comparison does not make their separated rule sections adjacent
or identify them as successive stages of unstressed-vowel reduction.

\section{\texorpdfstring{Root-noun nominative \emph{*-z}
loss}{Root-noun nominative  loss}}\label{root-noun-nominative-loss}

\subsection{Historical discussion}\label{historical-discussion-1}

The athematic consonant stems --- the ``root nouns'' of the handbooks
--- attached the nominative-singular marker directly to a
consonant-final root, and the Proto-Germanic outcome of that collision
is genuinely uncertain. Ringe gives the consonant-stem nominative ending
as zero, \emph{*-z}, or possibly \emph{*-s}, and states plainly that the
distribution is unrecoverable for monosyllabic stems
(\citeproc{ref-Ringe2017}{Ringe 2017, 306}, §4.3.4). His own paradigm
tables carry the uncertainty into print: the nominative of `foot'
appears as ``fōts? (fōs?)'', while `mouse' is plain \emph{*mūs}, its
expected extra sibilant already absorbed by degemination
(\citeproc{ref-Ringe2017}{Ringe 2017, 149, 313}). Ringe and Taylor
repeat the same three-way agnosticism --- the root-noun nominative
``either ended in \emph{*-s} or \emph{*-z}, or was endingless''
(\citeproc{ref-RingeTaylor2014}{Ringe and Taylor 2014, 28}, §2.3.1).

The dictionary traditions encode this situation in different notations,
and the differences are conventions of citation rather than competing
claims of fact. Orel prints morphologically explicit nominatives with
final \emph{-z} across the whole class: \emph{*bōkz} `book',
\emph{*ǥansz} `goose', \emph{*lūsz} `louse'
(\citeproc{ref-Orel2003}{Orel 2003, 52, 126, 252}). Kroonen cites the
same words as bare stems or endingless forms, \emph{*bōk-} and
\emph{*gans-} (\citeproc{ref-Kroonen2013}{Kroonen 2013, 71--72,
168--69}), and Kluge/Seebold print a third variant with voiceless
\emph{-s}, as in \emph{*bōks} `Buch'
(\citeproc{ref-KlugeSeebold2011}{Kluge 2011, 158}). Bammesberger shows
why the marker is nonetheless real: for voiced-final stems the overt
nominative \emph{*-z} is positively reconstructible --- \emph{*burgz}
(Gothic \emph{baúrgs}), \emph{*frijōndz} (Gothic \emph{frijonds}) ---
even though \emph{*fōt-z} is ``phonotaktisch kaum denkbar'', and in West
Germanic the ending simply ``fiel es ab''
(\citeproc{ref-Bammesberger1990}{Bammesberger 1990, 190--92}, §8.2.3.1).
CAPR retains Orel's morphologically explicit forms as its inputs
precisely because they record the inflectional marker whose fate this
rule describes.

The three focal words are only superficially parallel. The root-final
\emph{s} of `louse' is itself an extension of \emph{*luw-} on the model
of `mouse' (\citeproc{ref-Bammesberger1990}{Bammesberger 1990, 195},
§8.3). For `goose', Szemerényi's lengthening would give a
Proto-Indo-European nominative \emph{*ǵʰanss} \textgreater{}
\emph{*ǵʰān}, after which the Proto-Germanic nominative was rebuilt with
a final voiced sibilant, \emph{*ganz}, reanalyzed within the paradigm
(\citeproc{ref-Bammesberger1990}{Bammesberger 1990, 196}, §8.3;
\citeproc{ref-Kroonen2013}{Kroonen 2013, 168--69}). `Book' preserves a
plain obstruent-final root. What unites them is not a single phonetic
history but membership in a paradigm class that generalized the
endingless nominative.

The branch evidence dates and localizes that generalization. Gothic
keeps its sibilant throughout the class (\emph{baúrgs}, \emph{nahts},
\emph{reiks}) (\citeproc{ref-Fulk2018}{Fulk 2018, 165--66},
§§7.26--7.27). Old Norse redistributes the ending morphologically:
masculine root nouns keep \emph{-r} (\emph{fótr}), feminines are
endingless (\emph{nótt}, \emph{geit}), while the vocalic-stem feminines
\emph{kýr}, \emph{sýr}, \emph{ær} retain \emph{-r} from \emph{*-z} with
R-umlaut (\citeproc{ref-Fulk2018}{Fulk 2018, 167}, §7.28;
\citeproc{ref-Bammesberger1990}{Bammesberger 1990, 192--93}). West
Germanic alone is uniform: there was ``no ending in PWGmc, as none of
the daughters exhibits any'' (\citeproc{ref-RingeTaylor2014}{Ringe and
Taylor 2014, 118}, §3.4). Fulk supplies the mechanism: Szemerényi's law
removed the nominative sibilant after sonorant-final stems, and
endinglessness then spread analogically through the class
(\citeproc{ref-Fulk2018}{Fulk 2018, 143}, §7.2). The development is
therefore best understood as a morphological generalization enacted
differently in each branch --- absolute in West Germanic, consonant- and
gender-conditioned in North Germanic, absent in Gothic --- rather than
as one exceptionless sound law.

This change is distinct from the two later final-\emph{*z} developments.
It was complete before Proto-West Germanic, whereas the loss of
\emph{*-z} in unstressed syllables
(\hyperref[rule-EAFFinalZDeletion]{SC020 EAFFinalZDeletion}) is itself a
Proto-West Germanic change: polysyllabic consonant-stem nominatives such
as \emph{*fadurz} `father' --- and, in this corpus, \emph{*frijōndz}
`friend', \emph{*melukz} `milk', and \emph{*mēnōþz} `month' --- kept
their ending into Proto-West Germanic and lost it there in an unstressed
syllable (\citeproc{ref-RingeTaylor2014}{Ringe and Taylor 2014, 44--45},
§3.1.1). The still later northern loss of \emph{*-z} in stressed
monosyllables (\hyperref[rule-MonosyllabicFinalZLoss]{SC097
MonosyllabicFinalZLoss}) affects vowel-final monosyllables like
\emph{*hwaz} and does not touch the consonant-final root nouns at all,
whose ending was gone long before.

\subsection{\texorpdfstring{SC096. Root-noun nominative \emph{*-z} loss
(\texttt{RootNounNomZLoss})}{SC096. Root-noun nominative  loss (RootNounNomZLoss)}}\label{rule-RootNounNomZLoss}

\begin{ReaderFacingFoma}
define RootNounNomZLoss Ctx([
    [{*z} -> 0 || .#. [EnglishStarConsonant | EnglishPalatalConsonant]* EnglishStarVocalic+ [EnglishStarConsonant | EnglishPalatalConsonant]+ _ .#.]
]);
\end{ReaderFacingFoma}

The rule deletes word-final \emph{*z} after a consonant in a
monosyllable. Three claims of different kinds meet here and must be kept
apart. The historical claim is morphological: the nominative-singular
ending was lost in the athematic root-noun class, so that this one
paradigm cell came to lack its marker --- a development complete before
Proto-West Germanic (\citeproc{ref-RingeTaylor2014}{Ringe and Taylor
2014, 118}, §3.4). The lexical claim belongs to the dictionaries: Orel's
citation forms, which supply the corpus inputs, print that marker
explicitly as \emph{-z} in \emph{*bōkz}, \emph{*flauxz}, \emph{*ǥansz},
and \emph{*lūsz} (\citeproc{ref-Orel2003}{Orel 2003, 52, 105, 126,
252}). The executable statement is neither of these but a computational
proxy for them: `delete word-final \emph{*z} after a consonant in a
monosyllable'. It is not proposed as a Proto-Germanic sound law; it
earns its place only because every form the corpus submits to the
morphological development is a consonant-final monosyllable, so the
narrow phonological statement covers the class exactly. Four corpus
derivations witness the rule, each yielding its expected Old English
outcome: PGmc \Recon{bōkz} `book' yields OE \emph{bōc} `book',
\Recon{gánsz} `goose' yields \emph{gōs} `goose', \Recon{lūsz} `louse'
yields \emph{lūs} `louse', and \Recon{fláuxz} `flea' yields \emph{flēah}
`flea'. Should the corpus ever acquire a consonant-final monosyllable in
\emph{*-z} that is not a root-noun nominative, the proxy and the
morphology would come apart, and the rule would need to be re-scoped;
the project's regression tests pin the firing population to exactly
these four words so that any fifth firing forces that adjudication
rather than passing silently.

The rule applies at the head of the English line, before
\hyperref[rule-PWGmcIjContraction]{SC009 PWGmcIjContraction}. That
ordering is fixed by the identity of the process rather than by a wrong
form: contraction turns the polysyllabic \Recon{fríjōndz} `friend' into
a monosyllable, and if the root-noun rule applied after contraction it
would capture \emph{*friundz} --- yet the ending of `friend' survived
into Proto-West Germanic and fell in an unstressed syllable, the change
described under \hyperref[rule-EAFFinalZDeletion]{SC020
EAFFinalZDeletion} (\citeproc{ref-RingeTaylor2014}{Ringe and Taylor
2014, 44--45}, §3.1.1). Because both rules delete the same segment,
moving this rule later changes no Old English output; the early
placement keeps the derivation of `friend' aligned with the historical
account rather than with an accident of the cascade.

Negative controls behave as the morphology predicts. Stressed
monosyllables whose \emph{*z} follows a vowel --- the domain of the
later northern change --- do not meet the post-consonantal environment.
Medial \emph{*z} is untouched and remains available for rhotacism
(\hyperref[rule-EAFRhotacism]{SC003 EAFRhotacism}): PGmc \Recon{déuzą}
`deer' still yields OE \emph{dēor} `deer' with its rhotacized medial
consonant.

\section{Unstressed ai
monophthongization}\label{unstressed-ai-monophthongization}

\subsection{\texorpdfstring{Historical discussion of unstressed
\emph{*ai}
monophthongization}{Historical discussion of unstressed  monophthongization}}\label{historical-discussion-of-unstressed-monophthongization}

Ringe and Taylor describe the widespread monophthongization of
unstressed \emph{*ai} to a long mid vowel, largely merging with
inherited unstressed \emph{*ē}. Their examples include final
dative-singular and nonfinal inflectional syllables. They nevertheless
question a single completed Proto-Northwest Germanic merger: Runic
evidence suggests that some unstressed diphthongs survived later than
others (\citeproc{ref-RingeTaylor2014}{Ringe and Taylor 2014, 25--26}).
The stressed development of \emph{*ái} to \emph{*ā} is treated
separately as \hyperref[rule-EAFAiMonophthongization]{SC004
EAFAiMonophthongization}.

\subsection{\texorpdfstring{\CAPRRuleHeading{SC014. Monophthongization of unstressed \emph{*ai}}{PNWGmcUnstressedAiMonophthongization}}{}}\label{rule-PNWGmcUnstressedAiMonophthongization}

\noindent\begin{minipage}{\linewidth}
\begin{ReaderFacingFoma}
define PNWGmcUnstressedAiMonophthongization Ctx([
    [
    {*ai} -> {*ē}
]
]);
\end{ReaderFacingFoma}
\end{minipage}

The dative-singular endings span \Recon{spánnai} `span' and meed
\Recon{mízdai} `meed' carry this change; both give a final \emph{*ē}. If
\hyperref[rule-PNWGmcUnstressedAiMonophthongization]{SC014
PNWGmcUnstressedAiMonophthongization} is delayed until after
\hyperref[rule-OEUnstressedLongVowelShortening]{SC072
OEUnstressedLongVowelShortening}, the \emph{*ē} is no longer present for
shortening, so PGmc \Recon{spánnai} `span' yields \Pred{spannē} rather
than expected OE \emph{spanne} `span'. This shows that
\hyperref[rule-PNWGmcUnstressedAiMonophthongization]{SC014
PNWGmcUnstressedAiMonophthongization} must come before
\hyperref[rule-OEUnstressedLongVowelShortening]{SC072
OEUnstressedLongVowelShortening} in the modeled sequence.

\section{\texorpdfstring{Unstressed \emph{*a}-raising before final
\emph{*m}}{Unstressed -raising before final }}\label{unstressed--raising-before-final}

\subsection{Historical discussion}\label{historical-discussion-2}

Campbell notes that unstressed \emph{u} is especially well preserved
before \emph{m}, with dat.pl. \emph{-um} and related endings as the
clearest evidence (\citeproc{ref-Campbell1959}{Campbell 1959, 156},
§373). Fulk likewise includes the development of early unstressed
\emph{*o} to \emph{u} before \emph{m} among the similarities shared by
North and West Germanic (\citeproc{ref-Fulk2018}{Fulk 2018, 16}, §5.2).

I restrict the change to unstressed vowels in inflectional material
because the strongest evidence concerns noninitial unstressed material
before final \emph{*m}. Final \emph{*m} conditions the raising.

\subsection{\texorpdfstring{SC005. Unstressed \emph{*a}-raising before
final \emph{*m}
(\texttt{PNWGmcAToUBeforeM})}{SC005. Unstressed -raising before final  (PNWGmcAToUBeforeM)}}\label{rule-PNWGmcAToUBeforeM}

\begin{ReaderFacingFoma}
define PNWGmcAToUBeforeM Ctx([
    [
    {*a} -> {*u} || EnglishStarVocalic EnglishStarConsonant+ _ {*m} ({*i})? ({*z})? .#.
]
]);
\end{ReaderFacingFoma}

Here the witness word and the comparative evidence serve different
purposes. If raising is delayed until after
\hyperref[rule-PNWGmcULowering]{SC017 PNWGmcULowering}, PGmc
\Recon{skúldramiz} `shoulders' yields \Pred{sċoldrum} rather than
expected OE \emph{sċuldrum} `shoulders'; earlier placements converge on
the expected output. The scope of the change is established by
inflectional evidence across multiple paradigm types: a-stem dative
plural ON
{\emph{dǫgum}}\index[iv]{09on@\ivlangheader{Old Norse}{}!dogum@\iventry{dǫgum}{}}
`days', OE
{\emph{dagum}}\index[iv]{02oe@\ivlangheader{Old English}{West Saxon normalization unmarked}!dagum@\iventry{dagum}{}}
`days', OS
{\emph{dagun}}\index[iv]{13os@\ivlangheader{Old Saxon}{}!dagun@\iventry{dagun}{}}
`days', OHG
{\emph{tagum}}\index[iv]{10ohg@\ivlangheader{Old High German}{}!tagum@\iventry{tagum}{}}
`days', beside Gothic
{\emph{dagam}}\index[iv]{12goth@\ivlangheader{Gothic}{}!dagam@\iventry{dagam}{}}
`days'; strong-adjective dative singular ON
{\emph{góðum}}\index[iv]{09on@\ivlangheader{Old Norse}{}!godum@\iventry{góðum}{}}
`good', OE
{\emph{gōdum}}\index[iv]{02oe@\ivlangheader{Old English}{West Saxon normalization unmarked}!godum@\iventry{gōdum}{}}
`good', OS
{\emph{gōdum}}\index[iv]{13os@\ivlangheader{Old Saxon}{}!godum@\iventry{gōdum}{}}
`good', beside Gothic
{\emph{godamma}}\index[iv]{12goth@\ivlangheader{Gothic}{}!godamma@\iventry{godamma}{}}
`good' (OS also shows variant forms gōdumu and -un); and first-plural
present ON
{\emph{berum}}\index[iv]{09on@\ivlangheader{Old Norse}{}!berum@\iventry{berum}{}}
`we carry', OHG
{\emph{berumēs}}\index[iv]{10ohg@\ivlangheader{Old High German}{}!berumes@\iventry{berumēs}{}}
`we carry', beside Gothic
{\emph{baíram}}\index[iv]{12goth@\ivlangheader{Gothic}{}!bairam@\iventry{baíram}{}}
`we carry'. Across these sets, North/West Germanic shows unstressed
\emph{-um} where Gothic preserves \emph{-am}. The derivation of
\emph{sċuldrum} `shoulders' supplies a CAPR ordering witness for the
relative chronology, but the cognate set for `shoulder' does not
contribute comparative evidence for the rule's historical scope.

\section{\texorpdfstring{Northwest Germanic lowering of long
\emph{ē}}{Northwest Germanic lowering of long }}\label{northwest-germanic-lowering-of-long}

\subsection{Historical discussion}\label{historical-discussion-3}

Proto-Germanic \emph{*ē₁} (the long mid vowel of PIE origin, as against
the later \emph{*ē₂}) split East Germanic from the rest of the family.
Gothic keeps a mid vowel, written ⟨e⟩, in \emph{gadēþs} `deed',
\emph{slēpan} `to sleep', \emph{mēna} `moon', and \emph{jēr} `year',
while Norse and all of West Germanic show a low vowel, as in Old Norse
\emph{ráða} `to consider, to advise', \emph{láta} `to let go, to allow',
\emph{ár} `year', \emph{nál} `needle', Old High German \emph{tāt}
`deed', \emph{slāfan} `to sleep', \emph{jār} `year', \emph{māno} `moon',
Old Saxon \emph{dād} `deed', \emph{slāpan} `to sleep', \emph{jār}. Ringe
and Taylor assemble the comparative set and reconstruct a Northwest
Germanic sound change \emph{*ē₁} \textgreater{} \emph{*ā}
(\citeproc{ref-RingeTaylor2014}{Ringe and Taylor 2014, 11--13}). The
change is directly dated by the earliest epigraphy. The Early Runic
accusative \emph{mākija} `sword' (beside Gothic \emph{mēkeis} `sword')
already shows ⟨a⟩ in the second half of the second century AD
(\citeproc{ref-RingeTaylor2014}{Ringe and Taylor 2014, 12}), and the
a-rune likewise writes the stressed reflex in the Opedal stone's
\emph{swestar} `sister' (\citeproc{ref-Stiles2017}{Stiles 2017, 4}).
Loanwords borrowed into Sami from early Proto-Scandinavian point the
same way. Lid showed that the root vowel of Sami \emph{mānno} `moon',
Proto-Germanic \emph{*ē} as in Gothic \emph{mēna}, ``had gone over to
\emph{ā} before the word was borrowed into Sami'', and adduced the
parallel Sami \emph{saððo} `bran' beside Proto-Nordic \emph{*sāðō} from
\emph{*sēðō-}, Old Norse \emph{sáð} `bran' (\citeproc{ref-Lid1952}{Lid
1952, 238}). This is Sami borrowing from Proto-Scandinavian, so it dates
and localizes the \emph{ā} stage in the North; its application to West
Germanic rests on the daughter correspondences above.

The change affected stressed syllables only. Ringe and Taylor restrict
it explicitly: unstressed \emph{*ē} did not lower but was eventually
shortened, as in \emph{*fadēr} \textgreater{} \emph{*fader} `father'
(\citeproc{ref-RingeTaylor2014}{Ringe and Taylor 2014, 13}, and p.~147);
in the runic material the unstressed final vowel of \emph{faþiR}
`father' has already merged with \emph{*-i}
(\citeproc{ref-Stiles2017}{Stiles 2017, 4}). Within stressed syllables,
however, the lowering was unconditioned: it applied before nasals
exactly as elsewhere, and Ringe and Taylor accordingly print the
intermediates \emph{*mānō} `moon', \emph{*mānōþ-} `month', and
\emph{*spānuz} `spoon' (\citeproc{ref-RingeTaylor2014}{Ringe and Taylor
2014, 11}). The later, dialectally narrower fates of this \emph{*ā} ---
fronting in the North Sea area when oral, rounding when nasalized ---
are separate sound changes treated in their own chapters.

The reconstruction of the intermediate value is disputed. The objection
is an old one: Bennett observed that the received account obliges Gothic
to move \emph{ē} to \emph{ǣ} and back to \emph{ē}, and the
non-West-Saxon dialects and Old Frisian to run through
\emph{ē > ǣ > ā > ǣ > ē}, ``with no apparent agreement among the
languages and no discernible phonological trend'', and that the
Proto-Germanic \emph{ǣ} is on that account preserved directly nowhere
(\citeproc{ref-Bennett1950}{Bennett 1950, 232--33, 235}). In its place
he proposed that the backing to \emph{ā} began in the north of the
Germanic homeland and spread southward, never reaching Gothic in the
east or Anglian and Frisian in the west, so that West Saxon \emph{ǣ} is
a retention of the intermediate stage rather than a re-fronting
(\citeproc{ref-Bennett1950}{Bennett 1950, 234--35}). Fulk argues in the
same direction on modern evidence: the Northwest Germanic vowel was a
low front \emph{*ǣ}, retained unchanged in Anglo-Frisian, with the
backing to \emph{ā} of Norse and inner West Germanic a separate areal
development spreading northward from Upper German territory; on that
reading the runic ⟨a⟩ spellings write a front {[}æː{]} for which the
futhark had no better grapheme (\citeproc{ref-Fulk2018}{Fulk 2018,
60--61}, §4.6). Campbell is deliberately noncommittal, finding the
\emph{*ā} stage ``tempting to assume, though not definitely
demonstrable'' (\citeproc{ref-Campbell1959}{Campbell 1959, 50--51},
§§128--129). Ringe and Taylor answer with the lengthened place-adverbs
\emph{þǣr} `there' and \emph{hwǣr} `where', whose front vowels are most
naturally the output of fronting applied to a back \emph{*ā}
(\citeproc{ref-RingeTaylor2014}{Ringe and Taylor 2014, 13--14}). The
present model adopts the two-step reconstruction while recording that
the alternative remains live.

\subsection{\texorpdfstring{SC024. Lowering of stressed long \emph{ē}
(\texttt{PNWGmcLongELowering})}{SC024. Lowering of stressed long  (PNWGmcLongELowering)}}\label{rule-PNWGmcLongELowering}

\noindent\begin{minipage}{\linewidth}
\begin{ReaderFacingFoma}
define PNWGmcLongELowering Ctx([
    [
    {*ḗ} -> {*ā}
]
]);
\end{ReaderFacingFoma}
\end{minipage}

The rule reads the stressed tier \emph{*ḗ} only, in keeping with the
stress restriction; unstressed \emph{*ē}, as in \emph{*fadēr} `father',
is left for the unstressed-shortening rules of the Old English stage and
never lowers. No segmental environment is imposed: nasal forms such as
\Recon{mḗnōþz} `month' and \Recon{spḗnuz} `spoon' pass through
\emph{*mānōþ-} and \emph{*spānuz} on their way to \emph{mōnaþ} `month'
and \emph{spōn} `spoon', exactly as reconstructed by Ringe and Taylor
(\citeproc{ref-RingeTaylor2014}{Ringe and Taylor 2014, 11}).

The rule stands near the head of the cascade, before the genuinely West
Germanic innovations such as early \emph{i}-apocope,
\emph{*ij}-contraction, and \emph{j}-gemination --- the placement
follows from the dating: the second-century runic evidence puts the
lowering among the earliest Northwest Germanic developments, well before
the changes that separate West Germanic from Norse. Its output \emph{*ā}
is consumed much later in the cascade by
\hyperref[rule-EAFHiatusWInsertion]{SC102 EAFHiatusWInsertion},
\hyperref[rule-EAFLongANasalRounding]{SC025 EAFLongANasalRounding}, and
\hyperref[rule-EAFLongAFronting]{SC101 EAFLongAFronting}: if the
lowering is instead displaced after those rules, \Recon{mḗnōθz} `month'
yields \Pred{mānaþ} rather than OE \emph{mōnaþ} `month', \Recon{skḗpą}
`sheep' yields \Pred{sċāp} rather than \emph{sċēap} `sheep', and
\Recon{sḗaną} `to sow' never acquires the hiatus-filling \emph{w} of
\emph{sāwan} `to sow'. Within the cascade no earlier boundary has been
demonstrated; the second-century runic attestation supplies the absolute
dating.

\section{Early i-apocope}\label{early-i-apocope}

\subsection{Historical discussion}\label{historical-discussion-4}

Sievers/Brunner treats the early loss of final \emph{*i} after
unstressed syllables as established by the fact that these endings no
longer trigger later i-umlaut in Old English, and Ringe and Taylor make
the same point through the pathway to \emph{geoguþ} `youth'
(\citeproc{ref-SieversBrunner1965}{Brunner 1965, §§ 145--146};
\citeproc{ref-RingeTaylor2014}{Ringe and Taylor 2014, 141}). Campbell's
\emph{dugup} `troop' and \emph{geogup} `youth' examples belong to the
same pattern (\citeproc{ref-Campbell1959}{Campbell 1959, § 332}).

The ending vowel disappears in a weak suffixal environment early enough
to block later umlaut. This anti-umlaut chronology distinguishes the
change from later final-vowel losses.

\subsection{\texorpdfstring{SC006. Early i-apocope
(\texttt{PWGmcEarlyIApocope})}{SC006. Early i-apocope (PWGmcEarlyIApocope)}}\label{rule-PWGmcEarlyIApocope}

\begin{ReaderFacingFoma}
define PWGmcEarlyIApocope Ctx([
    [
    {*i} -> 0 || PGmcStarStressedVowel PGmcStarConsonant+ PGmcStarVocalic PGmcStarConsonant+ _ .#.,
    {*i} -> 0 || PGmcStarStressedVowel PGmcStarConsonant+ PGmcStarVocalic PGmcStarConsonant+ _ {*z} .#.
]
]);
\end{ReaderFacingFoma}

The absence of umlaut in \emph{geoguþ} `youth' provides the historical
argument for early deletion. The ordered derivation supplies a different
test: if apocope is delayed until after
\hyperref[rule-OEAwLongDiphthong]{SC034 OEAwLongDiphthong}, PGmc
\Recon{skáwōθi} `shows' yields \Pred{sċēaweþ} rather than expected OE
\emph{sċēawaþ} `shows'.

Early i-apocope must therefore precede the long-diphthong development.
Moving it earlier within the tested range leaves every output unchanged;
its early date rests on the anti-umlaut evidence, not on a lower
boundary supplied by the witness words.

\section{\texorpdfstring{Final \emph{*ō}-lowering before
\emph{*r}}{Final -lowering before }}\label{final--lowering-before}

\subsection{Historical discussion}\label{historical-discussion-5}

Ringe and Taylor separate two points here: a broader shortening of
vowels before word-final \emph{*r} in unstressed syllables (for which
kinship \emph{*r}-stems such as PGmc \emph{*fadér} \textgreater{} PWGmc
\emph{*fader} are a key diagnostic), and the specific
\emph{*ō}-before-\emph{*r} development needed here
(\citeproc{ref-RingeTaylor2014}{Ringe and Taylor 2014, 58--59}). The
direct lexical witnesses for \emph{*ō} in that environment are two
independent etyma: PGmc
\Recon{fedwōr}\index[iv]{03pgmc@\ivlangheader{Proto-Germanic}{}!fedwor@\iventry{*fedwōr}{}}
`four' with WGmc reflexes OE
{\emph{fēower}}\index[iv]{02oe@\ivlangheader{Old English}{West Saxon normalization unmarked}!feower@\iventry{fēower}{}}
`four', OFris
{\emph{fiuwer}}\index[iv]{11ofris@\ivlangheader{Old Frisian}{}!fiuwer@\iventry{fiuwer}{}}
`four', OS
{\emph{fiuwar}}\index[iv]{13os@\ivlangheader{Old Saxon}{}!fiuwar@\iventry{fiuwar}{}}
`four'; and PGmc
\Recon{watōr}\index[iv]{03pgmc@\ivlangheader{Proto-Germanic}{}!wator@\iventry{*watōr}{}}
`water' with OE
{\emph{wæter}}\index[iv]{02oe@\ivlangheader{Old English}{West Saxon normalization unmarked}!waeter@\iventry{wæter}{}}
`water'.

The rule is historically secure but narrow: final or pre-final \emph{*ō}
before word-final \emph{*r}. The clearest evidence remains concentrated
in the \texttt{four} and \texttt{water} material. No broader environment
for \emph{*ō} is attested.

\subsection{\texorpdfstring{\CAPRRuleHeading{SC007. Lowering of final bimoric \emph{*ō} before \emph{*r}}{PWGmcFinalOrLowering}}{}}\label{rule-PWGmcFinalOrLowering}

\noindent\begin{minipage}{\linewidth}
\begin{ReaderFacingFoma}
define PWGmcFinalOrLowering Ctx([
    [
    {*ō} -> {*a} || _ {*r} .#.
]
]);
\end{ReaderFacingFoma}
\end{minipage}

OE \emph{wæter} `water' reveals why lowering must precede
\hyperref[rule-EAFBrightening]{SC043 EAFBrightening}. If
\hyperref[rule-PWGmcFinalOrLowering]{SC007 PWGmcFinalOrLowering} is
delayed until afterwards, PGmc \Recon{wátōr} `water' yields \Pred{water}
rather than expected OE \emph{wæter} `water': brightening can affect the
vowel only after lowering has created its input. Moving the change
earlier within the tested range alters no output.

The witness thus supplies a terminus ante quem at brightening but no
earlier boundary. Comparative support for the \emph{*ō}-before-\emph{*r}
rule comes from the two lexical witnesses
\Recon{fedwōr}\index[iv]{03pgmc@\ivlangheader{Proto-Germanic}{}!fedwor@\iventry{*fedwōr}{}}
`four' and
\Recon{watōr}\index[iv]{03pgmc@\ivlangheader{Proto-Germanic}{}!wator@\iventry{*watōr}{}}
`water' and their WGmc reflexes; kinship \emph{*r}-stems belong to the
broader pre-\emph{*r} shortening context, not to a direct \emph{*ō}
\textgreater{} \emph{*a} control. Within CAPR,
{\emph{wæter}}\index[iv]{02oe@\ivlangheader{Old English}{West Saxon normalization unmarked}!waeter@\iventry{wæter}{}}
`water' is the form that establishes ordering before brightening. No
broader lowering of \emph{*ō} is attested.

\section{Coronal-w assimilation}\label{coronal-w-assimilation}

\subsection{Historical discussion}\label{historical-discussion-6}

Ringe and Taylor treat the assimilation of \emph{*dw} and \emph{*zw} to
\emph{*ww} as a shared Proto-West-Germanic innovation supported by one
example of each input cluster (\citeproc{ref-RingeTaylor2014}{Ringe and
Taylor 2014, 56--57}; \citeproc{ref-Stiles1985}{Stiles 1985, 89--94}).
The \emph{*dw} example is the numeral `four': PGmc \emph{*feðwor}
(Gothic \emph{fidwor}) → WGmc \emph{*fewwar} → OE \emph{fēower}, Old
Frisian \emph{fiuwer}, Old Saxon \emph{fiuwar}. The \emph{*zw} example
is the second-person plural pronoun, where two oblique case forms show
the change: acc./dat.~PGmc \emph{*izwiz} (Gothic \emph{izwis}) → OE
\emph{eow}, Old Frisian \emph{iu}, Old Saxon \emph{iu}, OHG \emph{iu};
and gen.~Ringe and Taylor's PGmc \emph{*izweraz} (Gothic \emph{izwara})
→ OE \emph{eower}, OHG \emph{iuwer}
(\citeproc{ref-RingeTaylor2014}{Ringe and Taylor 2014, 56}). Stiles
discusses the same pronominal material using his own reconstruction
conventions and explicitly treats Gothic \emph{izwara} among the
relevant comparanda (\citeproc{ref-Stiles1985}{Stiles 1985, 89--94}).
These two case forms belong to a single pronominal paradigm, not to two
independent etyma.

The historical support rests on a small witness set. Both coronal inputs
assimilate before \emph{*w}, but the evidence for each cluster is
confined: the numeral alone supplies the \emph{*dw} instance, and the
oblique case forms of the second-person plural pronoun supply the
\emph{*zw} instance. Both clusters are now witnessed in the corpus:
`four' for \emph{*dw}, and `you' --- selected in its dat.(-acc.) plural
cell \emph{*izwiz} --- for \emph{*zw}, deriving through \emph{*iwwi},
apocopated \emph{*iww}, to OE \emph{ēow} `you'
(\citeproc{ref-RingeTaylor2014}{Ringe and Taylor 2014, 41--42}, §3.1.1;
\citeproc{ref-Fulk2018}{Fulk 2018, § 8.3}, pp.~204--205).

\subsection{\texorpdfstring{\CAPRRuleHeading{SC008. Assimilation of coronal consonants before \emph{*w}}{PWGmcCoronalWAssimilation}}{}}\label{rule-PWGmcCoronalWAssimilation}

\noindent\begin{minipage}{\linewidth}
\begin{ReaderFacingFoma}
define PWGmcCoronalWAssimilation Ctx([
    [
    {*d} -> {*w} || _ {*w},
    {*z} -> {*w} || _ {*w}
]
]);
\end{ReaderFacingFoma}
\end{minipage}

OE \emph{fēower} `four' exposes a feeding relation: coronal assimilation
creates the \emph{*ww} input of the early reanalysis in
\hyperref[rule-OEWWSimplification]{SC031 OEWWSimplification}, yielding
the diphthongal intermediate before later English realization. Ringe and
Taylor explicitly give the coronal-assimilation and glide-reanalysis
sequence (\citeproc{ref-RingeTaylor2014}{Ringe and Taylor 2014,
41--42}). The earlier displacement result belonged to the superseded
late \emph{*ww}-deletion proxy and does not by itself predict the result
of displacing the reformulated early operation.

The numeral fixes that relative order. The pronoun now fixes a second
one: assimilation must precede rhotacism
(\hyperref[rule-EAFRhotacism]{SC003 EAFRhotacism}). The \emph{*z} of
\emph{*izwiz} stands between vowel and \emph{*w}; had rhotacism applied
first, it would have produced \Pred{irwiz}, from which OE \emph{ēow}
`you' can never be derived. The executable cascade composes the
assimilation well before rhotacism, and the corpus derivation of
\emph{ēow} `you' fails if the two are reversed. `Four' remains the sole
\emph{*dw} witness and the sole source of the coronal-assimilation →
\emph{ww}-simplification ordering constraint. The earlier boundary of
the assimilation remains undetermined.

\section{Inherited short-vowel-plus-geminate-glide
reanalysis}\label{inherited-short-vowel-plus-geminate-glide-reanalysis}

\subsection{Historical discussion}\label{historical-discussion-7}

Inherited geminate glides after short vowels underwent reanalysis in
Proto-West Germanic: the first glide became the offglide of a diphthong,
while the second remained consonantal. This is not simply deletion of
one \emph{w}, and it is not the later English realization of the new
diphthong. Ringe and Taylor distinguish those histories explicitly; Fulk
specifies the short-vowel environment
((\citeproc{ref-RingeTaylor2014}{Ringe and Taylor 2014, 41--42, 65--66,
171--75}; \citeproc{ref-Fulk2018}{Fulk 2018, 117})).

Coronal assimilation feeds the change in \emph{fēower} `four' and
\emph{ēow} `you': \hyperref[rule-PWGmcCoronalWAssimilation]{SC008
PWGmcCoronalWAssimilation} supplies the geminate that becomes
\emph{eu+w} or \emph{iu+w}. Inherited geminates also supply the
histories of \emph{ċēowan} `chew', \emph{dēaw} `dew' and \emph{hēawan}
`hew'. The resulting diphthongs enter the separate English realization
described under \hyperref[rule-OEDiphthongLeveling]{SC032
OEDiphthongLeveling} ((\citeproc{ref-RingeTaylor2014}{Ringe and Taylor
2014, 41--42, 65--66, 172--75}; \citeproc{ref-Campbell1959}{Campbell
1959, 45--47})).

Homorganic \emph{uww} requires a separate quantity decision. We adopt
Ringe and Taylor's preferred long-\emph{u} account for \emph{sċūwa}
`shadow', but retain their qualification of its evidence rather than
treating the quantity as independently certain. Hypothetical inherited
long-vowel-plus- \emph{ww} sequences are not additional reconstructed
witnesses: the held description specifies inherited geminates after
short vowels. The parallel inherited \emph{jj} development is
independently governed and is not changed here
((\citeproc{ref-RingeTaylor2014}{Ringe and Taylor 2014, 65--66};
\citeproc{ref-Fulk2018}{Fulk 2018, 117})).

\subsection{\texorpdfstring{SC031. Inherited short-Vww reanalysis
(\texttt{OEWWSimplification})}{SC031. Inherited short-Vww reanalysis (OEWWSimplification)}}\label{rule-OEWWSimplification}

\begin{ReaderFacingFoma}
define OEWWSimplification Ctx([
    [
    {*a} {*w} {*w} -> {*au} {*w},
    {*á} {*w} {*w} -> {*áu} {*w},
    {*e} {*w} {*w} -> {*eu} {*w},
    {*é} {*w} {*w} -> {*éu} {*w},
    {*i} {*w} {*w} -> {*iu} {*w},
    {*í} {*w} {*w} -> {*íu} {*w},
    {*u} {*w} {*w} -> {*ū} {*w},
    {*ú} {*w} {*w} -> {*ū} {*w}
]
]);
\end{ReaderFacingFoma}

The transport wrapper carries an independently selected sentence context
outside the segmental rewrite; it adds no reconstructed sound and does
not condition the reanalysis. Acute and unmarked spellings are notation
variants, not alternatives in sentence stress. The retained executable
name likewise does not determine the historical stage.

The modeled position after coronal assimilation and before j-gemination
is a representative serialization that keeps inherited/coronal-created
geminates distinct from the subsequently created j-path. It is not a
newly proved strict chronology between the whole events. \emph{Hīeġ}
`hay' and reconstructed West Saxon strew retain the separately
adjudicated \hyperref[rule-OEAwwjResolution]{SC029 OEAwwjResolution}
path. \emph{Hīew} `hue' has distinct j-created reanalysis at
\hyperref[rule-OEEwLongDiphthong]{SC033 OEEwLongDiphthong} and long-io
realization at \hyperref[rule-OEJGlideIO]{SC110 OEJGlideIO}. The
retained \hyperref[rule-OEJWWSimplification]{SC106 OEJWWSimplification}
compatibility operation is no longer fed by hue
((\citeproc{ref-Campbell1959}{Campbell 1959, 45--47};
\citeproc{ref-RingeTaylor2014}{Ringe and Taylor 2014, 173, 250})).

For you, early \emph{iu+w} now reaches
\hyperref[rule-PWGmcUnstressedWordFinalIApocope]{SC098
PWGmcUnstressedWordFinalIApocope}. Its weak-final context, not the
former presence of literal \emph{ww}, licenses apocope before mutation.
Earlier source-backed reanalysis thus coexists with genuine prosodic
conditioning rather than an output-restoring segment proxy
((\citeproc{ref-RingeTaylor2014}{Ringe and Taylor 2014, 41--42, 55,
57--58})).

\section{\texorpdfstring{\emph{ij}-contraction in
\emph{friend}}{-contraction in }}\label{contraction-in}

\subsection{Historical discussion}\label{historical-discussion-8}

Ringe and Taylor describe a change of \emph{*ijo} to \emph{*iu} in the
ancestor of \emph{friend}, with the pathway PGmc \emph{*frijōnd-}
(Gothic \emph{frijonds}) → PWGmc \emph{*friund} → OE \emph{frēond}, Old
Frisian \emph{frīund}, Old Saxon \emph{friund}, Old High German
\emph{friunt} (\citeproc{ref-RingeTaylor2014}{Ringe and Taylor 2014,
62}). The same source immediately warns that the \emph{*ijo} sequence is
unique enough that wider generalization is inadvisable
(\citeproc{ref-RingeTaylor2014}{Ringe and Taylor 2014, 62}). Luick
(printed p.~118) notes that \emph{iu} generalised within several
\emph{j}-stem paradigms through a related but differently conditioned
loss of \emph{j}, but does not supply a second example of the exact
stressed \emph{*ijo} sequence (\citeproc{ref-Luick1914}{Luick
1914-\/-1940, 118}).

The change concerns a rare sequence attested only in the
\emph{*frijōnd-} etymon and cannot safely be generalized into a broadly
productive rule.

\subsection{\texorpdfstring{SC009. \emph{ij}-contraction in
\emph{friend}
(\texttt{PWGmcIjContraction})}{SC009. -contraction in  (PWGmcIjContraction)}}\label{rule-PWGmcIjContraction}

\noindent\begin{minipage}{\linewidth}
\begin{ReaderFacingFoma}
define PWGmcIjContraction Ctx([
    [
    {*i} {*j} {*ō} -> {*iu} || _ EnglishStarConsonant,
    # Stressed variant (stress preserved):
    {*í} {*j} {*ō} -> {*íu} || _ EnglishStarConsonant
]
]);
\end{ReaderFacingFoma}
\end{minipage}

Only the \emph{*frijōnd-} etymon tests this contraction. If the rare
\emph{*ijō} sequence survives until after
\hyperref[rule-OEDiphthongLeveling]{SC032 OEDiphthongLeveling}, PGmc
\Recon{fríjōndz} `friend' yields \Pred{friund} rather than expected OE
\emph{frēond} `friend'; moving contraction earlier within the tested
range changes no output.

That single contrast places \hyperref[rule-PWGmcIjContraction]{SC009
PWGmcIjContraction} before diphthong leveling but gives no lower
boundary. It cannot establish a productive sound law beyond the
\emph{*frijōnd-} etymon, precisely the reservation made by Ringe and
Taylor.

\section{West Germanic j-gemination}\label{west-germanic-j-gemination}

\subsection{Historical discussion}\label{historical-discussion-9}

Fulk treats West Germanic consonant gemination before \texttt{*j} after
a short vowel as a regular development and illustrates it with forms
such as OE \emph{settan} `set' and \emph{lecgan} `lay'
(\citeproc{ref-Fulk2018}{Fulk 2018, 127}, §6.15).

The change applies specifically after a short vowel before \emph{*j},
not to geminate consonants generally.

\subsection{\texorpdfstring{SC010. West Germanic j-gemination
(\texttt{PWGmcJGemination})}{SC010. West Germanic j-gemination (PWGmcJGemination)}}\label{rule-PWGmcJGemination}

\begin{ReaderFacingFoma}
define PWGmcJGemination Ctx([
    [
    {*p} -> {*p} {*p} || EnglishStarShortVowel _ {*j},
    {*b} -> {*b} {*b} || EnglishStarShortVowel _ {*j},
    {*t} -> {*t} {*t} || EnglishStarShortVowel _ {*j},
    {*d} -> {*d} {*d} || EnglishStarShortVowel _ {*j},
    {*k} -> {*k} {*k} || EnglishStarShortVowel _ {*j},
    {*g} -> {*g} {*g} || EnglishStarShortVowel _ {*j},
    {*f} -> {*f} {*f} || EnglishStarShortVowel _ {*j},
    {*s} -> {*s} {*s} || EnglishStarShortVowel _ {*j},
    {*m} -> {*m} {*m} || EnglishStarShortVowel _ {*j},
    {*n} -> {*n} {*n} || EnglishStarShortVowel _ {*j},
    {*l} -> {*l} {*l} || EnglishStarShortVowel _ {*j},
    {*ŋ} -> {*ŋ} {*ŋ} || EnglishStarShortVowel _ {*j},
    {*x} -> {*x} {*x} || EnglishStarShortVowel _ {*j},
    # *w is an ordinary member of the same law. Campbell §407 p. 167:
    # "every consonant except r being affected after short syllables...
    # The forms in which w is doubled are dealt with in §120.2."
    # Ringe & Taylor p. 53: "it appears that that cluster too underwent
    # gemination in PWGmc." The geminate this creates is the input that
    # SC029 later resolves. Fulk §4.10 n. 1 dissents; see
    # audits/sc010-w-gemination-and-hay-depth-adjudication.md.
    {*w} -> {*w} {*w} || EnglishStarShortVowel _ {*j}
]
]);
\end{ReaderFacingFoma}

OE \emph{nett} `net' fixes the order because the syllabic-\emph{j}
development would remove the glide that conditions gemination. If
\hyperref[rule-PWGmcSyllabicJ]{SC011 PWGmcSyllabicJ} precedes
\hyperref[rule-PWGmcJGemination]{SC010 PWGmcJGemination}, PGmc
\Recon{nátją} `net' yields \Pred{nete} rather than expected OE
\emph{nett} `net'. Earlier movement of gemination changes no output.

The chronology is phonologically transparent: the consonant must
geminate before \emph{*j} ceases to be consonantal. The witness
establishes no earlier boundary.

\section{Syllabic j after final-vowel
loss}\label{syllabic-j-after-final-vowel-loss}

\subsection{Historical discussion}\label{historical-discussion-10}

Ringe and Taylor state directly that after final unstressed \texttt{*a}
and \texttt{*ą} were lost, postconsonantal \texttt{*j} became syllabic
\texttt{*i}, with outcomes behind OE \emph{here} `army' and \emph{rice}
`kingdom' (\citeproc{ref-RingeTaylor2014}{Ringe and Taylor 2014, 46}).

The sources establish the development, although the lexical evidence
supplies little independent support for its position. Its scope is
postconsonantal j, not high-vowel vocalization generally.

\subsection{\texorpdfstring{SC011. Syllabic \emph{*j} after final-vowel
loss
(\texttt{PWGmcSyllabicJ})}{SC011. Syllabic  after final-vowel loss (PWGmcSyllabicJ)}}\label{rule-PWGmcSyllabicJ}

\noindent\begin{minipage}{\linewidth}
\begin{ReaderFacingFoma}
define PWGmcSyllabicJ Ctx([
    [
    {*j} {*a} -> {*i} || EnglishStarShortVowel EnglishStarConsonant _ .#.,
    {*j} {*ą} -> {*i} || EnglishStarShortVowel EnglishStarConsonant _ .#.
]
]);
\end{ReaderFacingFoma}
\end{minipage}

The same PGmc \Recon{nátją} `net' witness supplies the only firm
boundary. Placing \hyperref[rule-PWGmcSyllabicJ]{SC011 PWGmcSyllabicJ}
before \hyperref[rule-PWGmcJGemination]{SC010 PWGmcJGemination} yields
\Pred{nete} rather than expected OE \emph{nett} `net'; moving it later
changes no output.

Comparative evidence establishes postconsonantal \emph{*j} to syllabic
\emph{*i} after final unstressed \emph{*a} or \emph{*ą} loss, with
\emph{here} `army' and \emph{rice} `kingdom' as outcomes. The lexicon
adds only that vocalization followed gemination, not where it falls
among subsequent changes.

\section{\texorpdfstring{Simplification of the cluster \emph{*xs} before
a
consonant}{Simplification of the cluster  before a consonant}}\label{simplification-of-the-cluster-before-a-consonant}

\subsection{Historical discussion}\label{historical-discussion-11}

The change treated here is narrow. When \emph{*x} stood before \emph{*s}
and that \emph{*s} was itself followed by a further consonant, the
\emph{*x} was lost and the cluster was reduced. Campbell states the rule
in exactly these terms, that when a consonant follows, \emph{*xs}
becomes \emph{s}, and he illustrates it with \emph{wæstm} `growth' and
\emph{wæsma} `growth' beside \emph{weaxan} `to grow', and with
Northumbrian \emph{sesta} `sixth' beside West Saxon \emph{siexta}
(\citeproc{ref-Campbell1959}{Campbell 1959, 170}, §417). Brunner gives
the same statement and adds that the following consonant may be \emph{j}
as well as a true obstruent, citing \emph{nēosian} `to visit',
\emph{þīsl} `pole', \emph{wæsma} and \emph{wæstm}
(\citeproc{ref-SieversBrunner1965}{Brunner 1965, 184}, §221.2). Bülbring
formulates it as a pre-English change and is careful to say that it held
as a rule, with exceptions (\citeproc{ref-Bulbring1902}{Bülbring 1902,
215}, §527).

The cluster that the change requires is a specific one, and two
neighbouring clusters show that the restriction is real. Where
\emph{*xs} stood with no consonant after it, the \emph{*x} was not lost
at all. It survived long enough to cause breaking and then hardened to
the sound written \emph{x}, that is {[}ks{]}, as in \emph{fox} `fox',
\emph{siex} `six', \emph{weaxan}, \emph{oxa} `ox' and \emph{fleax}
`flax' (\citeproc{ref-Campbell1959}{Campbell 1959, 170}, §416). Where
\emph{*x} stood before a single consonant it was likewise kept. Campbell
observes that once \emph{*xs} had become {[}ks{]}, the only group in
which \emph{*x} still stood before a voiceless consonant in earliest Old
English was \emph{*xt}, and that this group remained, as in
\emph{feohtan} `to fight', \emph{miht} `might', \emph{niht} `night' and
\emph{sōhte} `sought' (\citeproc{ref-Campbell1959}{Campbell 1959, 186},
§464). The contrast is attested outside English as well. Old High German
and Old Saxon \emph{lastar} `reproach' comes from \emph{*laxstra-} with
the \emph{*x} lost before \emph{*st}, while Old English \emph{leahtor}
comes from \emph{*laxtra-} with the \emph{*x} kept before \emph{*t}
(\citeproc{ref-Campbell1959}{Campbell 1959, 170}, §417).

The change is not Proto-Germanic. Gothic keeps the \emph{*h} of this
cluster in \emph{bi-niuhsjan} `to spy out' and \emph{saihsta} `sixth',
so the loss must be later than the separation of Gothic
(\citeproc{ref-RingeTaylor2014}{Ringe and Taylor 2014, 157--58}).
Campbell reports the loss from the whole West Germanic area and from
North Germanic as well, citing Old Norse \emph{ísl} `axle' and
\emph{nýsa} `to investigate', Old Saxon \emph{weslon} `to exchange',
\emph{wastum} `growth' and \emph{niustan}, and Old High German
\emph{niusen} `to try' (\citeproc{ref-Campbell1959}{Campbell 1959, 170},
§417). Ringe and Taylor set out the same comparative material and reach
a more guarded conclusion. They derive Proto-Germanic \emph{*niuhsijaną}
through Proto-West Germanic \emph{*niusjan} to Old English \emph{nēosan}
`to seek out, to visit', and \emph{*sehstō} to Northumbrian
\emph{sesta}, and they take the Old Saxon agreement in \emph{wastum},
\emph{thisla} and \emph{niusian} to show a shared northern West Germanic
change. Against that they weigh the competition between \emph{þixl}
`pole' and \emph{þīsl} `pole' in early Mercian, the survival of
\emph{*x} in \emph{eaxl} `shoulder' from \emph{*ahslu}, and the
retention in Old High German \emph{sehsto} and \emph{dihsala}. Their
conclusion is that the \emph{*h} was lost, possibly variably and
possibly only in some dialects, when two or more consonants followed,
and that the loss may have been in part a parallel development in the
diverging Northwest Germanic dialects
(\citeproc{ref-RingeTaylor2014}{Ringe and Taylor 2014, 157--58}). The
existence of the change is therefore secure while its exact date and
extent are not, and the rating given here reflects that division.

Two chronological anchors are available. Ringe and Taylor place the loss
after the Proto-West Germanic syncope of \emph{*-CijV-}, since it is
that syncope which brings the \emph{*s} and the \emph{*j} of
\emph{*niuhsjan} together (\citeproc{ref-RingeTaylor2014}{Ringe and
Taylor 2014, 157}). They place it before breaking, observing that the
undiphthongized vowels of \emph{wæstm} `growth' and \emph{þīsl} `pole'
can be accounted for only by supposing that these \emph{*h} were lost
before breaking took place (\citeproc{ref-RingeTaylor2014}{Ringe and
Taylor 2014, 158}). The rule is stated between those two points.

One witness in this collection undergoes the change. Proto-Germanic
\emph{*funxstiz} `fist' reaches the rule as \emph{*fū́xstiz}, the long
vowel having been produced by the Proto-Germanic loss of a nasal before
\emph{*x} described in \hyperref[rule-PGmcNasalLossBeforeX]{SC103
PGmcNasalLossBeforeX}. The cluster \emph{*xst} is then reduced to
\emph{*st}, and the word continues to Old English \emph{fȳst} `fist'.
The comparative set for this particular word is West Germanic
throughout, with Old Frisian \emph{fēst} `fist', Old Saxon and Old High
German \emph{fūst} `fist', Dutch \emph{vuist} and German \emph{Faust}
(\citeproc{ref-Kroonen2013}{Kroonen 2013, 160};
\citeproc{ref-Orel2003}{Orel 2003, 157}). Neither Gothic nor Old Norse
preserves a reflex of it, so the word bears on the domain of the change
rather than on its date.

The word \emph{thought} is the control. Proto-Germanic \emph{*θánxtē}
also passes through \hyperref[rule-PGmcNasalLossBeforeX]{SC103
PGmcNasalLossBeforeX}, which leaves \emph{*θą̄xtē}, but the \emph{*x}
there stands before a single \emph{*t} and not before \emph{*s}, so the
present rule does not touch it and Old English has \emph{þōhte}
`thought' with its \emph{h} intact. The pair \emph{fist} and
\emph{thought} reproduces within this collection the same contrast that
\emph{lastar} and \emph{leahtor} show across the West Germanic
languages.

The relation between the two rules should not be overstated. The earlier
rule does supply the cluster that this one simplifies, but the order of
the two is established by their stages and not by any word in this
collection. If the present rule were stated first, \emph{fist} would
still reach \emph{fȳst} `fist', because removing the \emph{*x} from
\emph{*funxstiz} leaves a nasal before \emph{s}, and that nasal is
removed with compensatory lengthening by the North Sea Germanic
nasal-spirant law. The outcome is overdetermined, and the chronology
rests on Gothic \emph{bi-niuhsjan} `to spy out' and \emph{saihsta}
instead.

\subsection{\texorpdfstring{\CAPRRuleHeading{SC028. Simplification of \emph{*xs} before a consonant}{PNWGmcPreconsonantalXLoss}}{}}\label{rule-PNWGmcPreconsonantalXLoss}

\noindent\begin{minipage}{\linewidth}
\begin{ReaderFacingFoma}
define PNWGmcPreconsonantalXLoss Ctx([
    [
    {*x} -> 0 || _ {*s} EnglishStarConsonant
]
]);
\end{ReaderFacingFoma}
\end{minipage}

The \emph{*s} in the structural description carries the whole weight of
the rule. Without it the rule would delete \emph{*x} before any two
consonants, and it would then wrongly remove the first element of a
geminate \emph{*xx} before \emph{*j}, where Old English in fact has
\emph{hliehhan} `to laugh' with the geminate written \emph{hh}
(\citeproc{ref-Campbell1959}{Campbell 1959, 186}, §464). The rule fires
on \emph{fist} and on no other word in this collection, and it leaves
\emph{fox}, \emph{six}, \emph{wax}, \emph{flax} and \emph{ox} untouched,
as it must, along with \emph{thought}, \emph{fight}, \emph{night},
\emph{light}, \emph{might}, \emph{knight}, \emph{fright} and
\emph{wight}.

The implementation name retains an older description of the rule and
will be brought into line in a separate pass.

\section{\texorpdfstring{\emph{lþ}-voicing}{-voicing}}\label{voicing}

\subsection{Historical discussion}\label{historical-discussion-12}

Ringe and Taylor treat word-internal \emph{*lþ} \textgreater{}
\emph{*ld} as a regular sound change in northern West Germanic and
illustrate it with forms such as \emph{fealdan} `fold', \emph{beald}
`bold', \emph{wuldor} `glory', and \emph{gylden} `golden'
(\citeproc{ref-RingeTaylor2014}{Ringe and Taylor 2014, 170--71}).
Campbell gives a similar West-Germanic-facing formulation with examples
such as \emph{fealdan}, \emph{wuldor}, \emph{beald}, \emph{gold} `gold',
and \emph{feld} `field' (\citeproc{ref-Campbell1959}{Campbell 1959,
169}, §414).

The comparative evidence supports \emph{lþ > ld} most clearly in
northern West Germanic, not as an unqualified pan-PWGmc development.

\subsection{\texorpdfstring{SC012. Northern West Germanic
\emph{lþ}-voicing
(\texttt{EAFLThVoicing})}{SC012. Northern West Germanic -voicing (EAFLThVoicing)}}\label{rule-EAFLThVoicing}

\noindent\begin{minipage}{\linewidth}
\begin{ReaderFacingFoma}
define EAFLThVoicing Ctx([
    [
    {*θ} -> {*d} || {*l} _
]
]);
\end{ReaderFacingFoma}
\end{minipage}

The \texttt{field}, \texttt{fold}, \texttt{gold}, and \texttt{wold}
families preserve \emph{*lþ} to \emph{*ld}, but none dates the change
against a neighboring rule. Every output remains unchanged when the
voicing is moved in either direction.

Comparative reconstruction therefore establishes northern West Germanic
\emph{lþ > ld}, but the witness forms fix no date. Neither a pan-PWGmc
attribution nor an exact local placement follows from the evidence
presented here.

\section{Dental hardening}\label{dental-hardening}

\subsection{Historical discussion}\label{historical-discussion-13}

Ringe and Taylor state directly that in PWGmc voiced dental fricative
\texttt{*ð} became stop \texttt{*d} in all positions
(\citeproc{ref-RingeTaylor2014}{Ringe and Taylor 2014, 43}).

The change is systemic across early West Germanic and extends beyond any
one lexical family.

\subsection{\texorpdfstring{SC013. Dental hardening
(\texttt{PWGmcDentalHardening})}{SC013. Dental hardening (PWGmcDentalHardening)}}\label{rule-PWGmcDentalHardening}

\noindent\begin{minipage}{\linewidth}
\begin{ReaderFacingFoma}
define PWGmcDentalHardening Ctx([
    [
    {*ð} -> {*d}
]
]);
\end{ReaderFacingFoma}
\end{minipage}

Dental hardening has systemic scope: voiced fricative \emph{*ð} became
stop \emph{*d} throughout early West Germanic. Moving
\hyperref[rule-PWGmcDentalHardening]{SC013 PWGmcDentalHardening} earlier
or later changes no output.

Comparative evidence establishes the sound law; the present lexicon
leaves its exact position approximate.

\section{\texorpdfstring{Endingless \emph{aw}: the regular straw
branch}{Endingless : the regular straw branch}}\label{endingless-the-regular-straw-branch}

\subsection{Historical discussion}\label{historical-discussion-14}

The inherited nominative/accusative \emph{*stráwą} loses its final
ending and vocalizes resulting final \emph{w}, giving \emph{*strau}.
Ordinary English diphthong development then produces reconstructed
\emph{*strēa}. Attested \emph{strēaw} `straw' instead combines that long
diphthong with \emph{w} restored from inflected forms; the long nucleus
also spreads into those forms. The regular free simplex and short
obliques are explicitly unattested. The compound \emph{strēaberige}
`strawberry' preserves the endingless branch, but is not substituted for
an attested free noun (\citeproc{ref-Campbell1959}{Campbell 1959,
46--47, 232--33}; \citeproc{ref-SieversBrunner1965}{Brunner 1965, 49,
110, 204--5}; \citeproc{ref-RingeTaylor2014}{Ringe and Taylor 2014,
172--73, 387}).

The model keeps attested \emph{strēaw} `straw' as its comparison and
produces regular reconstructed \emph{*strēa}: this is an explicit
mismatch. The fragment below combines final-vowel loss and ensuing
vocalization. Its narrow matcher does not make general final-vowel loss
\emph{w}-conditioned.

\subsection{\texorpdfstring{SC111. Endingless \emph{aw} fragment
(\texttt{PWGmcEndinglessAw})}{SC111. Endingless  fragment (PWGmcEndinglessAw)}}\label{rule-PWGmcEndinglessAw}

\noindent\begin{minipage}{\linewidth}
\begin{ReaderFacingFoma}
define PWGmcEndinglessAw [
    {*a} {*w} {*ą} -> {*au} || _ .#. ,
    {*á} {*w} {*ą} -> {*áu} || _ .#.
];
\end{ReaderFacingFoma}
\end{minipage}

The ordinary \emph{au} product feeds
\hyperref[rule-OEAuBrightening]{SC030 OEAuBrightening} and
\hyperref[rule-OEDiphthongLeveling]{SC032 OEDiphthongLeveling}, not the
retained show proxy \hyperref[rule-OEAwLongDiphthong]{SC034
OEAwLongDiphthong}. The early holding site is compatible with the
source's West Germanic history, not evidence for strict relative order
against unrelated dental hardening (\citeproc{ref-RingeTaylor2014}{Ringe
and Taylor 2014, 172--73}).

\section{Consonant-conditioned lowering of stressed short
i}\label{consonant-conditioned-lowering-of-stressed-short-i}

\subsection{\texorpdfstring{\CAPRRuleHeading{SC015. Consonant-conditioned lowering of stressed short \emph{*i}}{PNWGmcILowering}}{}}\label{rule-PNWGmcILowering}

\subsubsection{A disputed Germanic
correspondence}\label{a-disputed-germanic-correspondence}

Old English \emph{nest} `nest' and \emph{wer} `man' have long occupied a
disproportionate place in the history of Germanic vowels. Their mid
vowels appear to continue earlier short \emph{*i}, yet many words with
apparently comparable endings retain a high vowel. The question is not
whether lowering occurred somewhere, but which conditions explain its
distribution, and whether the relevant developments belonged to a common
ancestor or separate daughter histories. Campbell identifies only these
nouns and isolated \emph{spec} `bacon', beside \emph{spic} `bacon', as
clear English evidence for the traditional formulation
(\citeproc{ref-Campbell1959}{Campbell 1959, 43--44}).

The account adopted here lowers stressed short \emph{*i} immediately
before inherited \emph{*r}, tautosyllabic \emph{*z}, or \emph{*st}. It
follows Hill's consonant-conditioned proposal, not a general
assimilation to the following vowel. This is a working reconstruction
with unequal evidential support for its components: the rhotic
environment is strongest; the sibilant environments are more difficult.
Its value lies in relating positive correspondences to independently
important negative evidence, not in making every inherited formation
unproblematic (\citeproc{ref-Hill2017}{Hill 2017, 153--58}).

\subsubsection{The vowel-conditioned
tradition}\label{the-vowel-conditioned-tradition}

Streitberg states a common-Germanic lowering before a following
\emph{*a} or \emph{*ā}, blocked by intervening \emph{*j} or a nasal
followed by another consonant. The latter qualification matters: a
single nasal is not automatically a blocker. His presentation already
acknowledges extensive analogical disturbance. Hirt likewise connects
lowering with a following low or mid vowel, but qualifies its
geographical unity and admits unknown conditions. Neither author offers
an uncontested exceptionless distribution
(\citeproc{ref-Streitberg1896}{Streitberg 1896, 56--57};
\citeproc{ref-Hirt1931}{Hirt 1931-\/-1934, 45--46}).

Luick is more guarded for English: lowering appears in only some
formations and seems dependent on additional conditions not yet
recovered. Campbell preserves the short statement while separating the
much more regular lowering of \emph{*u} in
\hyperref[rule-PNWGmcULowering]{SC017 PNWGmcULowering}. The apparent
parallel between the two vowels is historically important, but parallel
notation does not establish identical conditions or chronology. Fulk
retains a serious role for vowel-conditioned lowering and subsequent
levelling, particularly the elimination of lowered vowels from paradigms
(\citeproc{ref-Luick1914}{Luick 1914-\/-1940, 108--9};
\citeproc{ref-Campbell1959}{Campbell 1959, 43--44};
\citeproc{ref-Fulk2018}{Fulk 2018, 56--59}).

There was consequently no single traditional consensus subsequently
overturned by one discovery. Lloyd's survey distinguishes authors who
accepted a common change, those who restricted it to dialectal
prehistories, and those who merely registered the forms or denied the
law. That history cautions against treating every reconstructed
\emph{*e} as an endorsement of the same vowel-harmony mechanism. The
inference from a correspondence to a particular conditioner must be
argued independently (\citeproc{ref-Lloyd1966}{Lloyd 1966, 738--39}).

\subsubsection{Why the following vowel is
insufficient}\label{why-the-following-vowel-is-insufficient}

Class I strong-verb participles supply a major difficulty. Short
\emph{*i} survives before a non-high participial vowel where the
parallel class II formation commonly shows lowered \emph{*u}. The
contrast cannot be removed simply by noting that participles belong to
paradigms: a proposed restoration must identify the source of the
restored vowel and explain its distribution. Lloyd challenges precisely
those morphological premises, rather than denying that analogy exists
(\citeproc{ref-Lloyd1966}{Lloyd 1966, 739--40};
\citeproc{ref-Hill2017}{Hill 2017, 150--52}).

The uninflected adverbs are harder still. OHG \emph{hina} `hence',
\emph{hinana} `from here', and \emph{nidana} `from below' retain high
vowels before non-high continuations without an ordinary case paradigm
from which \emph{*i} could be restored. Fish vocabulary supplies another
independent class: PGmc \emph{*fiskaz} `fish' and \emph{*wiskō} `wisp'
preserve \emph{*i} before \emph{*sk}. Changing their inflection after
the event does not itself explain the inherited root vocalism
(\citeproc{ref-Lloyd1966}{Lloyd 1966, 741--42};
\citeproc{ref-Hill2017}{Hill 2017, 152, 154--55}).

\subsubsection{Hock's defence of
lowering}\label{hocks-defence-of-lowering}

Hock's detailed reply must nevertheless be distinguished from a mere
assumption of analogy. He accepts an inherited contrast between
\emph{*e} and \emph{*i} before \emph{*u} in Norse and English, and in
certain monosyllables, while still defending lowering before low vowels.
Complete phonemic merger and conditioned lowering are therefore
different claims: rejecting the former does not refute the latter. His
later raising before \emph{*u} is geographically restricted to the
continental languages, not projected indiscriminately into their common
ancestor (\citeproc{ref-Hock1973}{Hock 1973, 328--29, 345--46}).

For class I participles, Hock identifies the \emph{i}-timbre of the
present, rather than simply the preterite plural, as the restoring
model; he also appeals to differentiation from other strong classes.
Hill directly challenges this account on quantity: the proposed present
model has long \emph{*ī}, whereas the participles require short
\emph{*i}. Selective restoration of timbre without quantity needs its
own historical justification. Cercignani and Fulk nevertheless continue
to allow present-stem influence, so the dispute is not settled by citing
agreement on phonemic distinctness (\citeproc{ref-Hock1973}{Hock 1973,
340--41}; \citeproc{ref-Hill2017}{Hill 2017, 151--52};
\citeproc{ref-Cercignani1980}{Cercignani 1980, 130};
\citeproc{ref-Fulk2018}{Fulk 2018, 58--59}).

Hock also confronts the adverbs, proposing different formative histories
and subsequent levelling, and adducing OE
{\emph{heonan}}\index[iv]{02oe@\ivlangheader{Old English}{West Saxon normalization unmarked}!heonan@\iventry{heonan}{}}
`hence' and
{\emph{neoðan}}\index[iv]{02oe@\ivlangheader{Old English}{West Saxon normalization unmarked}!neothan@\iventry{neoðan}{}}
`from below' alongside the high-vowel forms. Those are comparative
hypotheses, not evidence that every relevant suffix or restoration has
been established. In particular, Ringe and Taylor note that Old English
\emph{eo} for expected \emph{io} can result from a much later merger;
Hill likewise distinguishes later Norse developments. The surface vowels
alone cannot date a shared ancestral lowering
(\citeproc{ref-Hock1973}{Hock 1973, 342};
\citeproc{ref-RingeTaylor2014}{Ringe and Taylor 2014, 35}, n.~17;
\citeproc{ref-Hill2017}{Hill 2017, 152}, n.~24).

For `fish', Hock proposes an original i-stem transferred to the a-stems
only after lowered vowels had become phonemic. That is a specific
alternative history, but the stem transfer and its date require
independent support; Cercignani explicitly questions this explanation.
Hock's final appeal to nominative/vocative syncretism likewise supplies
a proposed mechanism for nominal doublets, not an independently fixed
chronology. His own defence and the later objections are part of the
case for evaluating particular histories rather than treating analogy as
either universally available or universally excluded
(\citeproc{ref-Hock1973}{Hock 1973, 343--45}, p.~349, n.~29;
\citeproc{ref-Cercignani1979}{Cercignani 1979, 77--78};
\citeproc{ref-Cercignani1980}{1980, 131}).

Cercignani's defence of distinct \emph{*i} and \emph{*e} phonemes
addresses a related but different proposition: phonemic opposition does
not exclude a conditioned lowering. His criticism concerns the extension
of one distribution and date across Germanic. In particular, the Norse
restriction after preceding \emph{*k/g} cannot simply be transferred to
West Germanic. Ringe and Taylor also separate the southern West Germanic
lowering before labial and velar obstruents from the widely shared nest
and man correspondences. An independent continental development
therefore cannot be assumed to establish a general ancestral rule
(\citeproc{ref-Cercignani1979}{Cercignani 1979, 75--76, 82};
\citeproc{ref-Cercignani1980}{1980, 129--31};
\citeproc{ref-RingeTaylor2014}{Ringe and Taylor 2014, 33--36}).

\subsubsection{Hill's consonantal
environments}\label{hills-consonantal-environments}

Hill redirects attention to the immediately following consonant. Before
tautosyllabic \emph{*z}, the principal comparisons include PGmc
\emph{*iz} `he', \emph{*hiz} `this here', and \emph{*mizdōn} `reward'.
The last belongs with OE \emph{meord} `reward' and \emph{mēd} `reward',
whose consonants and quantities subsequently diverge. Their mid vowels
need not have been induced by a later low vowel: the monosyllabic
pronouns are especially significant because they remove that presumed
determinant (\citeproc{ref-Hill2017}{Hill 2017, 153}).

Before \emph{*st}, \emph{*nistan} `nest' supplies the central example.
Its \emph{*e} is compatible with both accounts when a non-high ending is
reconstructed, so it does not by itself choose the conditioner. The
contrast with retained \emph{*sk} is more informative than nest alone.
The comparative formation also supports an earlier high vowel: Kroonen
relates the noun to the downward particle and the zero grade of `sit'.
Nevertheless, one exceptionally clear noun is a thin foundation for the
full productivity of an \emph{*st} law
(\citeproc{ref-Kroonen2013}{Kroonen 2013, 391};
\citeproc{ref-Hill2017}{Hill 2017, 153--55}).

Before \emph{*r}, early \emph{*wiraz} `man' is compared with OE
\emph{wer}, ON \emph{verr} `man', and Gothic \emph{wair} `man'. Lloyd
himself allows a special rhotic influence despite rejecting general
a-umlaut. Howell's account of Gothic breaking adds phonetic
plausibility: rhotics, especially coda rhotics, commonly lower
neighbouring vowels. But recurrence of a natural phonetic process is not
proof of a single inherited event. Gothic could have lowered
independently, and its breaking can obscure an older distinction between
high and mid inputs (\citeproc{ref-Lloyd1966}{Lloyd 1966, 744}, n.~9;
\citeproc{ref-Hill2017}{Hill 2017, 156};
\citeproc{ref-Howell2024}{Howell 2024, 81--82, 90--91}).

Gothic consequently has a double evidential role. Its distribution can
illuminate the phonetic environment of lowering, while its own breaking
makes some comparisons historically ambiguous. A Gothic mid-vowel reflex
before a rhotic is compatible with an earlier lowered vowel, but it is
also compatible with an inherited high vowel affected by Gothic
breaking. The correspondence cannot choose between those histories
without additional chronological evidence. Hill's proposal gains its
force from combining the Gothic facts with West Germanic quantity and
consonantal correspondences, not from treating Gothic spelling as a
transparent record of one ancestral vowel. Conversely, the possibility
of a Gothic innovation does not compel an independent innovation in
every other daughter. A common development remains an explanatory
hypothesis where its conditions and relative chronology can be defended.
The inherited rhotic, the later rhotacized sibilant and the short versus
long vowel must therefore be compared separately, rather than counted as
interchangeable instances of the same environment
(\citeproc{ref-Hill2017}{Hill 2017, 149, 156--58};
\citeproc{ref-Howell2024}{Howell 2024, 81--82}).

Hill further proposes that later i-mutation conceals some earlier
lowering. His lowered Northwest Germanic \emph{*nẹstijana} `to nest' and
\emph{*lẹstiz} `skill' can return to high vowels before a following
high-front trigger. Here the later raising has a specific phonological
cause; it is not an unspecified repair assigned separately to every
exception. The same account must nevertheless establish that the trigger
existed in the relevant historical formation. Identical final high
vowels alone cannot demonstrate the two-stage history
(\citeproc{ref-Hill2017}{Hill 2017, 153}).

\subsubsection{Lowering and the length of
`here'}\label{lowering-and-the-length-of-here}

The most consequential chronological proposal concerns OE \emph{hēr}
`here'. Hill begins with early \emph{*hir}, parallel to the deictic
adverbs `where' and `there', and argues for the sequence \emph{*hir}
\textgreater{} \emph{*her} \textgreater{} \emph{*hēr}. His closer
mid-vowel notation distinguishes the lowered result from an automatic
identification with inherited \emph{*ē₁}. If monosyllabic lengthening
had occurred first, an additional lowering of long \emph{*ī} would be
needed, without the same independent support
(\citeproc{ref-Hill2017}{Hill 2017, 148--49, 156--58}).

On Hill's interpretation, lowering before \emph{*r} must therefore
precede the lengthening. His broader argument treats West Germanic
monosyllabic lengthening and aspects of Gothic breaking as potentially
partly Proto-Germanic, making an early shared chronology a serious
hypothesis. The conclusion is conditional, however: the source of the
long vowel, the reconstruction of the short precursor and the identity
of the lengthening must all be accepted. Ringe and Taylor discuss
alternative ablaut and later-development interpretations. The modern
correspondence does not uniquely force Hill's sequence
(\citeproc{ref-Hill2017}{Hill 2017, 149, 156--58};
\citeproc{ref-RingeTaylor2014}{Ringe and Taylor 2014, 36--37}).

Nor does the argument date every consonantal component. Even if
\emph{*hir > her} is established before lengthening, neither \emph{*z}
nor \emph{*st} participates in that derivation. Their absolute dates
require their own evidence. Likewise inherited \emph{*r} is distinct
from \emph{*r} subsequently created by rhotacism: a later surface rhotic
cannot simply be projected backwards as the conditioner of an earlier
law (\citeproc{ref-Hill2017}{Hill 2017, 153--58};
\citeproc{ref-Crist2002}{Crist 2002, 3--6}).

\subsubsection{Twine: a genuine difficulty for the sibilant
condition}\label{twine-a-genuine-difficulty-for-the-sibilant-condition}

OE \emph{twīn} `twine' and High German \emph{zwirn-} present the
sharpest difficulty. Orel reconstructs masculine \emph{*twiznaz};
Kluge--Seebold gives \emph{*twizna-} and distinguishes the early verbal
and nominal attestations. The modern long vowel is not evidence for
inherited long \emph{*ī}. A short-vowel formation containing \emph{*zn}
is supported independently of the lowering hypothesis
(\citeproc{ref-Orel2003}{Orel 2003, 414};
\citeproc{ref-KlugeSeebold2011}{Kluge 2011, 1021}).

Two familiar histories preserve its timbre: \emph{*twizna- > twirna-} by
rhotacism, and \emph{*twizna- > twīna-} through loss and compensatory
lengthening. De Vaan explicitly connects those treatments with
assimilation as another outcome of \emph{*zn}; Campbell groups
\emph{twīn} `twine' with lengthening after sibilant loss. If early
lowering instead produces \emph{*twezna-}, neither rhotacism nor loss
restores the high vowel: \emph{*twerna-} and \emph{*twēna-} remain
different predictions. The conflict is not merely orthographical
(\citeproc{ref-Campbell1959}{Campbell 1959, 166}, §404;
\citeproc{ref-deVaan2016}{Vaan 2016}, unpaginated, paragraphs ``Verwante
vormen'' and ``Zowel'').

Kilday offers a specific alternative: the English word may be borrowed,
with a linen-related semantic history and possible verbal borrowing
before nominal back-formation. That proposal is worth considering
because it has an argument independent of the present rule. It is not
thereby proved. Moreover, removing English from the inherited evidence
would not explain away the High German high vowel. Crist's
lowering-before-loss-before-rhotacism chronology makes the
incompatibility particularly clear when applied to the same inherited
\emph{*izn} sequence (\citeproc{ref-Kilday2024}{Kilday 2024, 1, 4--5};
\citeproc{ref-Crist2002}{Crist 2002, 3--6}).

Twine must therefore remain counterevidence, not be assigned a
convenient exceptional exclusion. A syllable-sensitive account might
distinguish some sequences, but the proposed syllabification would need
independent historical support. The fact that a consonant follows
\emph{*z} is not itself proof that this sibilant had precisely the
tautosyllabic status required by Hill's phonetic generalization
(\citeproc{ref-Hill2017}{Hill 2017, 153}; \citeproc{ref-Orel2003}{Orel
2003, 414}).

\subsubsection{Numeral adverbs: conditional rather than
discharged}\label{numeral-adverbs-conditional-rather-than-discharged}

OHG \emph{zwiro} `twice', \emph{zwiror} `twice', and \emph{driror}
`thrice' raise a related question. Their prehistory is less secure than
that of twine. Kroonen gives \emph{*twis} and alternatives
\emph{*twiswar ~ twizwar}; Orel also preserves different sibilant
formations. Fulk starts from \emph{*twis-wōz}. These cannot be reduced
to one certain bare \emph{*twiz} simply because a later form contains a
rhotic (\citeproc{ref-Hirt1931}{Hirt 1931-\/-1934, 116};
\citeproc{ref-Kroonen2013}{Kroonen 2013, 530};
\citeproc{ref-Orel2003}{Orel 2003, 414}; \citeproc{ref-Fulk2018}{Fulk
2018, 275}).

The etymological discussion of \emph{driror} explains final \emph{-or}
through comparative-adverbial remodelling and considers a
\emph{*wa}-formation. That explains neither the root vowel nor
automatically the first rhotic; its own \emph{*þrizwa-/twizwa-} retains
a potentially relevant sibilant. A history with \emph{*sw}, one with
intervocalic \emph{*z}, and one with coda \emph{*z} before \emph{*w}
test different conditions. These adverbs are thus unresolved conditional
controls, not three independent refutations and not three exceptions
already disposed of by analogy
(\citeproc{ref-LloydLuhrSpringer1998}{Lloyd et al. 1998, cols 784,
s.v}.~driror; \citeproc{ref-Noreen1923}{Noreen 1923, 309};
\citeproc{ref-Campbell1959}{Campbell 1959, 287}).

\subsubsection{The Norse participle and the disputed leather
word}\label{the-norse-participle-and-the-disputed-leather-word}

The participle of ON \emph{bíða} `await' is \emph{beðinn} `awaited',
also cited as \emph{beðenn} `awaited', not an attested \emph{*biðinn}
`awaited'. Noreen considers a lowered ending and influence from the
participle of `bid'; Seebold allows regular lowering but prefers the
latter connection. Fulk treats the lowering interpretation seriously and
denies that the analogy has been demonstrated. This isolated low-vowel
participle is therefore evidence the consonantal account must confront,
but it does not establish a general inherited vowel-conditioned law for
all class I participles (\citeproc{ref-Noreen1923}{Noreen 1923, 326,
332}; \citeproc{ref-Seebold1970}{Seebold 1970, 94--95};
\citeproc{ref-Fulk2018}{Fulk 2018, 58--59, 286}).

OHG \emph{irh} `chamois leather' potentially contradicts lowering before
\emph{*r}. Hill favours an existing loanword interpretation; Kroonen
reconstructs \emph{*irha-} `he-goat' within a wider comparative family.
The alternatives concern lexical ancestry, not merely the operation of a
vowel law. If the relevant formation is inherited, its high vowel needs
an explanation. Borrowing cannot be declared certain solely because it
would remove a counterexample (\citeproc{ref-Hill2017}{Hill 2017, 156},
n.~30; \citeproc{ref-Kroonen2013}{Kroonen 2013, 270}).

\subsubsection{Learn and the limits of paradigm
reconstruction}\label{learn-and-the-limits-of-paradigm-reconstruction}

OE \emph{leornian} `learn' alternates with \emph{liornian} `learn'; OHG
\emph{lernēn} `learn' likewise stands beside \emph{lirnēn} `learn'. The
comparative base is \emph{*lizn-}, but a single uncontroversial
inherited paradigm has not been recovered. Jasanoff reconstructs a class
III history and explicitly considers a lost third-singular
\emph{*lizniþ} as a possible source of high-vowel forms. Hill instead
emphasizes earlier lowering followed by mutation in appropriately
extended stems. Both arguments depend on morphology not recoverable from
the infinitive alone (\citeproc{ref-Kroonen2013}{Kroonen 2013, 340};
\citeproc{ref-Orel2003}{Orel 2003, 247};
\citeproc{ref-Jasanoff1973}{Jasanoff 1973, 865, 868};
\citeproc{ref-Hill2017}{Hill 2017, 153}, n.~25).

The inherited high vowel and the later \emph{io} spelling must be
distinguished. Breaking can represent the continuation of a high vowel,
but that high vowel could itself have resulted from mutation after an
earlier lowering. Conversely, the existence of \emph{leorn-} does not
establish every reconstructed mutation trigger or levelling direction.
Traditional dialect descriptions illuminate the distribution, not a
simple opposition between conservative Northumbrian and innovative West
Saxon: early West Saxon also has \emph{liorn-}. Later weak-class
remodelling further complicates the inherited cells
(\citeproc{ref-Campbell1959}{Campbell 1959, 48, 59, 125};
\citeproc{ref-Jasanoff1973}{Jasanoff 1973, 865, 868};
\citeproc{ref-TarsiLaker2026}{Tarsi and Laker 2026, 42--43}).

CAPR selects \emph{leornian}, imperative \emph{leorna} `learn!', and
third-singular \emph{leornaþ} `learns'. These have independent textual
support: Ælfric's Treatise supplies an infinitive after \emph{scealt},
Deuteronomy has the singular command, and Bede has \emph{leornað}
`learns' with a singular `whoever' subject. The dental spelling in the
last is normalized without changing its vowel or cell. These are
selections among attested forms, not three independent etymologies or
proof that every \emph{liorn-} history has been explained
(\citeproc{ref-Crawford1922}{Crawford 1922, 71, 351};
\citeproc{ref-Miller1891}{Miller 1891, 380}).

\subsubsection{Adoption and its chronological
limits}\label{adoption-and-its-chronological-limits}

The evidential asymmetry is itself informative. A law whose only
argument were the absence of exceptions in a small English sample would
have little comparative force. Here the positive correspondences belong
to several different formations: pronominal stems, a verbal family, a
reward noun and the inherited nest word. Their consonantal overlap is
more specific than the mere presence of a non-high vowel somewhere to
the right. Conversely, the class I participles and the uninflected
adverbs test the old following-vowel condition without automatically
entering any of the proposed consonantal environments. The same examples
thus have different diagnostic value under the two hypotheses
(\citeproc{ref-Lloyd1966}{Lloyd 1966, 739--45};
\citeproc{ref-Hill2017}{Hill 2017, 153--56}).

Later raising must be argued with corresponding precision. A visible
high-front trigger can make a low intermediate compatible with a later
high reflex; it does not prove that lowering happened. Once the trigger
has been lost or transferred across a paradigm, morphological evidence
must establish its former presence. That is why a mutation-masked
derivative can be a useful compatibility control while the proposed
learn paradigm remains uncertain. Neither should be treated as an
unlimited mechanism for restoring the high vowel of every
counterexample. The distinction preserves testable content in the
consonantal hypothesis rather than immunizing it against comparison
(\citeproc{ref-Hill2017}{Hill 2017, 153};
\citeproc{ref-Jasanoff1973}{Jasanoff 1973, 865, 868}).

The unified rule is adopted because it relates important positive and
negative correspondences without requiring general vowel-conditioned
lowering and unspecified restoration in each resisting class. Its
strongest historical foundation remains the rhotic component. The reward
and pronoun families support the sibilant component, but twine remains a
real difficulty; the productivity of the \emph{st} component is less
securely demonstrated beyond nest. Explanatory economy is not permission
to erase those differences (\citeproc{ref-Hill2017}{Hill 2017, 153--58};
\citeproc{ref-Lloyd1966}{Lloyd 1966, 739--45}).

\begin{ReaderFacingFoma}
define PNWGmcILowering Ctx([
    {*i} -> {*e} ||
        .#. EnglishStarConsonant* _
        [{*z} [EnglishStarConsonant | .#.]
         | {*s} {*t}
         | {*r}],
    {*í} -> {*é} ||
        .#. EnglishStarConsonant* _
        [{*z} [EnglishStarConsonant | .#.]
         | {*s} {*t}
         | {*r}]
]);
\end{ReaderFacingFoma}

The definition affects the first short-vowel nucleus, with acute and
unmarked forms treated in parallel; it does not lower long \emph{*ī}.
Its following consonants, not the next vowel or the presence of an
earlier velar, determine application. Preconsonantal or word-final
\emph{*z} is an explicit approximation to tautosyllabicity. It must not
be mistaken for independent evidence about every \emph{*zn} sequence.
The Northwest Germanic placement is a qualified working reconstruction;
the possible earlier date of the rhotic component remains separate.

The selected reward dative \emph{*mízdai} is not the nominative
\emph{*mizdōn}; their shared root is the relevant lowering witness. In
\emph{nest}, the inherited formation and \emph{st} are central. Selected
thistle lowers before \emph{st} and later returns to a high vowel
through mutation, illustrating a concealed intermediate without proving
all proposed concealed histories. Masculine PGmc \emph{*weþruz} `wether'
already has \emph{*e} and therefore supplies no \emph{i}-lowering
evidence (\citeproc{ref-Hill2017}{Hill 2017, 153};
\citeproc{ref-Kroonen2013}{Kroonen 2013, 370, 391, 584};
\citeproc{ref-Orel2003}{Orel 2003, 459--60}).

Finally, \emph{world} contains the inherited `man' element: Ringe and
Taylor's \emph{*weraldiz} already has \emph{*e}. An earlier-vowel
compound representation cannot independently prove lowering before
\hyperref[rule-OEInterStressRaising]{SC036 OEInterStressRaising}. The
consonantal account has no following-vowel trigger for interstress
raising to remove. The compound is compatible with the selected history,
not a unique chronological discriminator. Neither it nor the conditional
\emph{hēr} `here' argument establishes a total order for all three
lowering environments (\citeproc{ref-RingeTaylor2014}{Ringe and Taylor
2014, 326}; \citeproc{ref-Orel2003}{Orel 2003, 462};
\citeproc{ref-Hill2017}{Hill 2017, 156--58}).

\section{Northwest Germanic
u-lowering}\label{northwest-germanic-u-lowering}

\subsection{Historical discussion}\label{historical-discussion-15}

Northwest Germanic lowered \emph{*u} to \emph{*o} when the following
syllable contained a non-high vowel. Campbell describes the change and
lists \emph{ġeoc} `yoke' among its regular outcomes
(\citeproc{ref-Campbell1959}{Campbell 1959, 42--43}, §115); Fulk gives
the same word as a standard example --- ``OIcel. ok, OE geoc, OHG joh
beside juh and OS juk'' --- and notes the paradigmatic alternation
between lowered and unlowered stems that the conditioning produced
(\citeproc{ref-Fulk2018}{Fulk 2018, 56}, §4.3). A word-initial \emph{*j}
does not block the change: the blocking effect of \emph{j} concerns only
a consonantal \emph{j} standing between the target vowel and the
conditioning vowel, as in the class I weak verbs of the \emph{cnyssan}
`strike' type (\citeproc{ref-Fulk2018}{Fulk 2018, 56}, §4.3). Ringe and
Taylor accordingly reconstruct the Proto-West Germanic paradigm of
`yoke' with the lowering applied (\citeproc{ref-RingeTaylor2014}{Ringe
and Taylor 2014, 129}).

Regular corpus witnesses include
{\emph{ġeoc}}\index[iv]{02oe@\ivlangheader{Old English}{West Saxon normalization unmarked}!geoc@\iventry{ġeoc}{}}
`yoke', \emph{nosu} `nose', \emph{sċofl} `shovel', and \emph{sorg}
`sorrow'. Where the following syllable kept a high vowel the lowering
did not apply, as in \emph{ġeoguþ} `youth', whose root \emph{u} survived
(\citeproc{ref-SieversBrunner1965}{Brunner 1965, 64--65}, §92.1).

Regular lowering does not imply that every attested inherited word
preserves the lowered vowel. The masculine noun
{\emph{rust}}\index[iv]{02oe@\ivlangheader{Old English}{West Saxon normalization unmarked}!rust@\iventry{rust}{}}
`rust' is a lexical residue: selected PGmc \Recon{rústaz} `rust'
regularly predicts \Recon{rost} `rust'. Orel's alternative \Recon{rustō}
`rust' gives the same lowered root; his masculine alternative is
independently supported by Ringe (\citeproc{ref-Orel2003}{Orel 2003,
309}; \citeproc{ref-Ringe2017}{Ringe 2017, 247--48}). Since
\Recon{frústą} `frost' regularly gives
{\emph{forst}}\index[iv]{02oe@\ivlangheader{Old English}{West Saxon normalization unmarked}!forst@\iventry{forst}{}}
`frost', a general \emph{-st-} exemption is unwarranted. The retained
\emph{u} remains unexplained, alongside the other retained-\emph{u}
cases such as \emph{bucc} `buck', \emph{fugol} `bird', \emph{wulf}
`wolf' and \emph{wulle} `wool'; their consonantal environments do not
establish a general replacement law
(\citeproc{ref-Campbell1959}{Campbell 1959, 43}, §115;
\citeproc{ref-RingeTaylor2014}{Ringe and Taylor 2014, 27--34}). The
lowering law is retained; its non-high conditioning class includes
nasalized long \emph{*ǭ}, not just oral \emph{*ō}. The independently
reconstructed accusative \emph{*wullǭ} therefore lowers just as the
nominative \emph{*wullō} does under the preferred chronology
(\citeproc{ref-Ringe2017}{Ringe 2017, 242--43}). This does not merge
either ending with trimoric \emph{*ô}. The \hyperref[rust-oe-rust]{rust
lexical discussion} distinguishes the selected masculine citation,
Orel's alternative and diagnostic paradigm constructions.

\subsubsection{The second root of
rainbow}\label{the-second-root-of-rainbow}

The independent noun \emph{boga} `bow' supplies an important comparison.
Its short \emph{o} continues inherited \emph{u} lowered before a
following non-high vowel; Kroonen identifies the noun as a masculine
n-stem derived from `bend'. Colman explicitly gives this early
explanation of \emph{boga}, while Ringe and Taylor include its ancestor
among the lowering examples (\citeproc{ref-Kroonen2013}{Kroonen 2013,
82}; \citeproc{ref-Colman2011}{Colman 2011, 95--96};
\citeproc{ref-RingeTaylor2014}{Ringe and Taylor 2014, 27--29}).

CAPR extends the same assimilation to the subordinate-stressed root in
\emph{*régna-bùgô} `rainbow'. The reconstruction retains primary accent
on `rain' and subordinate accent on `bow'; these are reconstructed
prosodic properties, not manuscript accent marks. The proposed
development begins \emph{*régna-bùgô > *régna-bògô > *régnabogô},
followed by palatalization of the first velar and the ordinary ending
developments to \emph{reġnboga} `rainbow'. The early non-high ending
still conditions lowering. The resulting \emph{ò} is decoded by the
existing technical stress operation; the first-member vowel survives
until \hyperref[rule-OENonHighSyn]{SC114 OENonHighSyn}, giving
\emph{*réʝnbogô}. Positional realization of underlying \emph{b} is
computed after that loss, without a new late fricative-to-stop sound
law. Neither early lowering nor syncope requires an optional rewrite.

The general early change is source-backed; its extension to this precise
compound and its reconstructed secondary accent is CAPR's phonological
inference. The authorities do not print the complete accented derivation
(\citeproc{ref-Campbell1959}{Campbell 1959, 43}, §§115--116;
\citeproc{ref-HoggGrammar2011}{Hogg 2011, 54--55}, §§3.10--3.11;
\citeproc{ref-Fulk2018}{Fulk 2018, 56--57}, §4.3). Analogy from the
independent noun is conceivable but is not independently demonstrated
here. The significant consequence is chronological: \emph{rainbow}
already has \emph{o} before
\hyperref[rule-OEMedUnstressedULowering]{SC040
OEMedUnstressedULowering}. It therefore cannot establish that secondary
stress had to disappear before that late lowering.

Colman's moneyer-name evidence remains a genuine qualification.
\emph{Buga} beside \emph{Boga}, including coin spellings with both
vowels, cannot be dismissed by exempting names from sound change.
Transmission to die-cutters and the interpretation of individual
spellings complicate phonological and dialect attribution. The ordinary
noun's regular early history does not by itself resolve the name's
retained vowel (\citeproc{ref-Colman2011}{Colman 2011, 95--98}).

\subsection{\texorpdfstring{\CAPRRuleHeading{SC017. Lowering of \emph{*u} before following non-high vowels}{PNWGmcULowering}}{}}\label{rule-PNWGmcULowering}

\begin{ReaderFacingFoma}
define PNWGmcULowering [Ctx([
    [
    {*u} -> {*o}
        || .#. EnglishStarConsonant* _
           [EnglishStarConsonantNoJ - EnglishStarNasal]
           EnglishStarConsonantNoJ* EnglishStarNonHighVowel,
    {*ú} -> {*ó}
        || .#. EnglishStarConsonant* _
           [EnglishStarConsonantNoJ - EnglishStarNasal]
           EnglishStarConsonantNoJ* EnglishStarNonHighVowel
    ]
]) .o. Ctx([
    {*ù} -> {*ò}
        || _
           [EnglishStarConsonantNoJ - EnglishStarNasal]
           EnglishStarConsonantNoJ* EnglishStarNonHighVowel
])];
\end{ReaderFacingFoma}

The plain and acute clauses lower the first nucleus; the grave clause
recognizes subordinate root stress without lowering ordinary unstressed
medial \emph{u}. Each uses a consonantal interval without \emph{*j}; the
immediately following consonant must not be nasal. This literal guard is
not a claim that every nasal anywhere in the interval blocks the
historical process. Fulk distinguishes tautosyllabic and heterosyllabic
nasal environments (\citeproc{ref-Fulk2018}{Fulk 2018, 56}).

The lowering feeds the much later West Saxon palatal-glide spelling
(\hyperref[rule-OEWsPalatalGlide]{SC016 OEWsPalatalGlide}): the \emph{o}
that the scribes wrote in \emph{ġeoc} `yoke' is the output of this
change, so PGmc \Recon{júką} `yoke' passes through \emph{*joką} on its
way to the attested spelling (\citeproc{ref-Fulk2018}{Fulk 2018, 56},
§4.3; \citeproc{ref-RingeTaylor2014}{Ringe and Taylor 2014, 129}). If
lowering is placed after \hyperref[rule-PNWGmcFinalLongORaising]{SC019
PNWGmcFinalLongORaising}, selected \Recon{núsō} `nose' yields
\Pred{nusu} rather than expected \emph{nosu}, PGmc \Recon{skúflō}
`shovel' yields \Pred{sċufl} rather than expected \emph{sċofl} `shovel',
and PGmc \Recon{súrgō} `sorrow' yields \Pred{surg} rather than expected
\emph{sorg}. These are conditional computational consequences, not
independent proof that those paradigms underwent no levelling. CAPR
provisionally prefers lowering before final raising because it regularly
accounts for these low-root forms without additional, individually
unproved analogical generalizations. Nose already involves earlier stem
remodelling, and the retained high vowels of wool and love remain a
genuine cost of that preference. The main comparison of Fulk, Ringe and
Taylor, and Stiles, including the Runic evidence and the plausible but
unproved nominative-generalization account of wool, is in
\hyperref[rule-PNWGmcFinalLongORaising]{SC019 PNWGmcFinalLongORaising}.

\section{\texorpdfstring{Stressed monosyllable
\emph{*ō}-raising}{Stressed monosyllable -raising}}\label{stressed-monosyllable--raising}

\subsection{Historical discussion}\label{historical-discussion-16}

Campbell treats the development of final accented \emph{ō} to \emph{ū}
in stressed monosyllables directly, with the familiar outcomes behind
\emph{cū} `cow', \emph{hū} `how', \emph{tū} `two', and \emph{bū} `both'
(\citeproc{ref-Campbell1959}{Campbell 1959, 47}, §122).

The change is historically secure, but the tested forms determine no
close relative position for it. Its input is final \emph{*ō} in a
stressed monosyllable.

\subsection{\texorpdfstring{\CAPRRuleHeading{SC018. Raising of final stressed monosyllabic \emph{*ō}}{PNWGmcStressedMonosyllableORaising}}{}}\label{rule-PNWGmcStressedMonosyllableORaising}

\noindent\begin{minipage}{\linewidth}
\begin{ReaderFacingFoma}
define PNWGmcStressedMonosyllableORaising Ctx([
    [
    {*ō} -> {*ū} || .#. [EnglishStarConsonant | EnglishPalatalConsonant]* _ .#.
]
]);
\end{ReaderFacingFoma}
\end{minipage}

Campbell's \emph{cū} `cow', \emph{hū} `how', and \emph{tū} `two'
establish final stressed monosyllabic \emph{*ō} \textgreater{}
\emph{*ū}.

Reversing \hyperref[rule-PNWGmcStressedMonosyllableORaising]{SC018
PNWGmcStressedMonosyllableORaising} with neighboring changes leaves
every output unchanged. The sound change is secure, but its exact
position in the early history of long vowels rests on the handbooks.

\section{\texorpdfstring{Raising of final unstressed long
\emph{*ō}}{Raising of final unstressed long }}\label{raising-of-final-unstressed-long}

\subsection{Historical discussion}\label{historical-discussion-17}

Final unstressed oral bimoric \emph{*ō} became short \emph{*u} in
Northwest Germanic. The qualification matters: oral \emph{*ō}, nasalized
\emph{*ǭ} and trimoric \emph{*ô} were not interchangeable endings. Ringe
and Taylor distinguish the first two developments explicitly; Fulk also
distinguishes the quantities when reconstructing the endings
(\citeproc{ref-RingeTaylor2014}{Ringe and Taylor 2014, 15--16, 58--60};
\citeproc{ref-Fulk2018}{Fulk 2018, 82--86, 155--56}).

The nominative singular \emph{*wullō} `wool' and accusative singular
\emph{*wullǭ} are independently reconstructed by Ringe. Their distinct
endings must survive the comparison even when a later paradigm levels
the root vowel (\citeproc{ref-Ringe2017}{Ringe 2017, 242--43}).
Likewise, an oral vowel followed by \emph{*z} was not yet word-final.
These distinctions explain both the scope of raising and its interaction
with later loss.

\subsection{\texorpdfstring{\CAPRRuleHeading{SC019. Raising of final unstressed long \emph{*ō}}{PNWGmcFinalLongORaising}}{}}\label{rule-PNWGmcFinalLongORaising}

\noindent\begin{minipage}{\linewidth}
\begin{ReaderFacingFoma}
define PNWGmcFinalLongORaising Ctx([
    [
    {*ō} -> {*u}
        || EnglishStarVocalic
           [EnglishStarConsonant | EnglishPalatalConsonant]+ _ .#.
    ]
]);
\end{ReaderFacingFoma}
\end{minipage}

The preceding vowel and consonants select a polysyllabic input. The
distinct stressed-monosyllable development is retained separately in
\hyperref[rule-PNWGmcStressedMonosyllableORaising]{SC018
PNWGmcStressedMonosyllableORaising}; it is not absorbed into this
unstressed ending law.

\subsubsection{Lowering before raising: the preferred
account}\label{lowering-before-raising-the-preferred-account}

Fulk argues that \emph{scolu} `troop, shoal', from \emph{*skulō},
requires lowering of the root before final \emph{*ō} became \emph{*u}.
If raising had already removed the non-high conditioner, an unchanged
nominative would retain root \emph{*u} (\citeproc{ref-Fulk2018}{Fulk
2018, 56}). CAPR provisionally follows this order:
\hyperref[rule-PNWGmcULowering]{SC017 PNWGmcULowering}, then final oral
raising. It regularly gives the low roots of \emph{nosu} `nose',
\emph{sċofl} `shovel' and \emph{sorg} `sorrow', rather than requiring a
later root-vowel generalization for each selected form.

This is an argument from explanatory economy, not from the number of
matching outputs. Fulk supplies the direct \emph{scolu} argument; the
comparison of its explanatory costs with the other selected paradigms is
CAPR's assessment. The witnesses are not equally independent.
\emph{Nosu} already depends on an earlier remodeled zero-grade stem
beside the comparative \emph{*nasō-}; it is not an untouched inherited
nominative (\citeproc{ref-Kroonen2013}{Kroonen 2013, 383};
\citeproc{ref-Campbell1959}{Campbell 1959, 43}, §115). For \emph{sċofl}
`shovel', Luick also infers an earlier high-root doublet,
\emph{*scufel}, beside \emph{scofel} `shovel'; that inference makes an
alternating paradigm credible but does not establish its direction of
levelling (\citeproc{ref-Luick1914}{Luick 1914-\/-1940, 105--6}, §78).
The comparative formations of shovel and sorrow support the selected
short-vowel stems, not a directly attested early English levelling
operation (\citeproc{ref-Orel2003}{Orel 2003, 347};
\citeproc{ref-Kroonen2013}{Kroonen 2013, 493}). The preferred sequence
avoids those additional operations, but it does not thereby explain
every retained high vowel.

\subsubsection{Raising before lowering: a coherent
alternative}\label{raising-before-lowering-a-coherent-alternative}

Stiles expressly places final \emph{*ō} raising before a-umlaut of
\emph{*u} (\citeproc{ref-Stiles2012}{Stiles 2012, 45}). Ringe and Taylor
explain how \emph{wull} `wool' and \emph{lufu} `love' could preserve a
nominative root \emph{*u} by generalizing it through a paradigm whose
other cells had lowered \emph{*o}; they also leave open genuine
exceptions (\citeproc{ref-RingeTaylor2014}{Ringe and Taylor 2014, 32}).
With the nasal accusative correctly recognized, the alternative has the
following regular consequences. The table gives reconstructed sound-law
predictions, not unattested forms promoted to attestations.

{\def\LTcaptype{none} % do not increment counter
\begin{longtable}[]{@{}
  >{\raggedright\arraybackslash}p{(\linewidth - 4\tabcolsep) * \real{0.3333}}
  >{\raggedright\arraybackslash}p{(\linewidth - 4\tabcolsep) * \real{0.3333}}
  >{\raggedright\arraybackslash}p{(\linewidth - 4\tabcolsep) * \real{0.3333}}@{}}
\toprule\noalign{}
\begin{minipage}[b]{\linewidth}\raggedright
Paradigm
\end{minipage} & \begin{minipage}[b]{\linewidth}\raggedright
Lowering before raising
\end{minipage} & \begin{minipage}[b]{\linewidth}\raggedright
Raising before lowering
\end{minipage} \\
\midrule\noalign{}
\endhead
\bottomrule\noalign{}
\endlastfoot
Wool: nominative / accusative & \emph{*woll} / \emph{*wolle} &
\emph{*wull} / \emph{*wolle} \\
Love: nominative / accusative & \emph{*lofu} / \emph{*lofe} &
\emph{*lufu} / \emph{*lofe} \\
Shoal: nominative / accusative & \emph{*sċolu} / \emph{*sċole} &
\emph{*sċulu} / \emph{*sċole} \\
\end{longtable}
}

The inherited ending opposition follows Ringe's reconstructed paradigm
and the independently described ō-stem inflections; it is not a special
ending invented to rescue wool (\citeproc{ref-Ringe2017}{Ringe 2017,
242--43}; \citeproc{ref-Fulk2018}{Fulk 2018, 154--56}). Under earlier
raising, nominative-to-oblique generalization can account for attested
\emph{wull} and \emph{lufu/lufe} `love'. Conversely, generalization of
the lowered oblique root can account for \emph{scolu} without replacing
its nominative ending. The same possibility exists for nose, shovel and
sorrow. Neither particular direction is directly witnessed in the
reconstructed early English paradigms.

There is independent evidence that levelling occurred: Ringe and Taylor
discuss early \emph{Holtijaz} and the daughter and ox paradigms, while
Stiles distinguishes conditioned Runic forms from analogical
generalizations (\citeproc{ref-RingeTaylor2014}{Ringe and Taylor 2014,
33--34}; \citeproc{ref-Stiles2012}{Stiles 2012, 47--48}). That
establishes the mechanism's credibility, not its occurrence in every
inconvenient lexeme. Continental \emph{skola} is not an untouched
nominative control either: Fulk describes accusative-ending
generalization into the nominative in Old Saxon and Old High German
(\citeproc{ref-Fulk2018}{Fulk 2018, 154--55}). Thus his \emph{scolu}
argument is strong only conditional on the inherited nominative root not
having been remodeled. It does not refute Ringe and Taylor's
phonologically coherent wool account.

There is a further lexical caution: Fulk elsewhere calls \emph{scolu}
`host' a loan from Latin \emph{schola} `school', whereas his chronology
argument expressly starts from inherited \emph{*skulō}
(\citeproc{ref-Fulk2018}{Fulk 2018, 47}, n.~1, 56). These analyses
cannot silently supply independent confirmation of one another. The
comparison here evaluates his inherited analysis.

The balance nevertheless favors retaining lowering first as a working
account: several independently selected low-root forms follow without
additional, individually unproved levelling. Wool and love remain its
real explanatory cost. Ringe and Taylor's alternative is plausible but
unproved, not a demonstrated error. Analogy belongs to the separately
argued history of a paradigm, not to the exceptionless sound laws.

\subsubsection{What the Runic evidence
dates}\label{what-the-runic-evidence-dates}

Fulk cites \emph{horna} `horn' on the Gallehus horn, about 400, for
lowering and \emph{waritu} `I write' on the Järsberg stone, about 450,
for final oral raising (\citeproc{ref-Fulk2018}{Fulk 2018, 56, 82}).
These are separate positive chronological bounds. An occurrence of
raising around 450 does not establish that it began after 400. Nor does
a written final \emph{o} in another inflection establish retention of
oral bimoric \emph{*ō}: nasalization, quantity and a following consonant
must first be identified. Stiles's treatment of the Runic distribution
also requires distinguishing conditioned realization, phonemicization
and analogical extension (\citeproc{ref-Stiles2012}{Stiles 2012, 47--48,
54--60}). No new general revaluation law follows from those
distinctions.

Fulk's note cites Ringe and Taylor on lowering and final \emph{*-az}
loss, not as an explicit rebuttal of their wool explanation
(\citeproc{ref-Fulk2018}{Fulk 2018, 57}, n.~1;
\citeproc{ref-RingeTaylor2014}{Ringe and Taylor 2014, 27--34}). Stiles
predates their book. Luick discusses alternation and rejects a proposed
labial protection law using regular low-root counterexamples; Campbell
and Bülbring list retained-high-vowel cases
(\citeproc{ref-Luick1914}{Luick 1914-\/-1940, 104--6}, §§77--79;
\citeproc{ref-Campbell1959}{Campbell 1959, 43}, §115;
\citeproc{ref-Bulbring1902}{Bülbring 1902, 45--46}, §116). None of these
checked passages supplies our exact comparative economy argument. No
specific later criticism or defence was verified in the bounded
literature search; this is not a claim of novelty or of its absence from
all scholarship. The Campbell and Bülbring passages remain text-checked
rather than newly certified against original images.

\subsubsection{Final consonants and later ending
reduction}\label{final-consonants-and-later-ending-reduction}

In \emph{*rástōz} `rest', \emph{*z} protects oral \emph{*ō} from this
raising. After \hyperref[rule-EAFFinalZDeletion]{SC020
EAFFinalZDeletion}, the newly final long vowel can undergo
\hyperref[rule-PWGmcSurvivingBimoricOUnrounding]{SC042
PWGmcSurvivingBimoricOUnrounding} and later fronting and shortening,
yielding \emph{ræste} `rest'. Moving raising after final consonant loss
wrongly exposes that vowel and gives \emph{*rast} instead. This relation
is retained separately; qualifying the lowering/raising comparison does
not weaken it.

Nasal \emph{*ǭ} does not raise here. Its later reduction yields the
accusative ending in \emph{*wolle}; the current serialization carries
the nasal symbol to that reduction rather than claiming that nasality
historically lasted until then. The genitive alternatives must also
remain distinct: Ringe and Taylor reconstruct bimoric \emph{*-ōz}, where
Fulk argues for trimoric \emph{*-ôz}
(\citeproc{ref-RingeTaylor2014}{Ringe and Taylor 2014, 59, 299};
\citeproc{ref-Fulk2018}{Fulk 2018, 155}). The trimoric genitive's
complete pathway remains unmodeled; it is not silently replaced by a
bimoric genitive or covered by the present nasal-vowel repair.

\newpage

\chapter{From Proto-West Germanic to
Anglo-Frisian}\label{from-proto-west-germanic-to-anglo-frisian}

\section{Historical interval}\label{historical-interval-2}

This chapter covers the sound changes that occurred after the Proto-West
Germanic period and before, or during the emergence of, the specifically
English line. The starting reconstruction is Proto-West Germanic; the
end point is the Anglo-Frisian ancestor required by the model's
historical tree. Its existence and the identification of its innovations
are separate questions. Its inventory must be reconstructed from
conditioned laws and daughter chronologies, not from the traditional
subgroup label alone.

\section{A necessary terminological
caution}\label{a-necessary-terminological-caution}

The title of this chapter uses ``Anglo-Frisian'' as an organizing
historical concept. That choice requires an explicit qualification.

The scholarly literature uses several overlapping terms for this
developmental period:

\begin{itemize}
\tightlist
\item
  North Sea Germanic and Ingvaeonic: labels used by some scholars for a
  proposed subgroup comprising Old English, Old Frisian, and Old Saxon
  (or more narrowly, Old English and Old Frisian only). The innovations
  associated with this label --- especially the nasal spirant changes
  and certain vowel developments --- are sometimes described as
  diffusion rather than shared inheritance
  (\citeproc{ref-Bremmer2009}{Bremmer 2009, 24--27, 126--28}).
\item
  Anglo-Frisian: a label used specifically for the Old English / Old
  Frisian branch, or for innovations shared between the two languages.
  Its use presupposes a tighter relationship between English and Frisian
  than between either and Old Saxon.
\item
  Proto-Anglo-Frisian (PAF): the strongest interpretation, positing a
  discrete reconstructed common ancestor for Old English and Old Frisian
  specifically. This is the position of Ringe and Taylor, who
  reconstruct a PAF stage between Proto-West Germanic and
  Proto-Old-English (\citeproc{ref-RingeTaylor2014}{Ringe and Taylor
  2014, 54--68}).
\end{itemize}

CAPR requires a discrete Anglo-Frisian ancestral node. This does not
guarantee that every traditionally associated innovation is inherited
from it, or exclusive to English and Frisian. A broader innovation can
already be present in its input; similar daughter outcomes can also
result from independent events. Authors' diffusion interpretations
remain accurately attributed, but they cannot resolve an incompatible
chronology within CAPR's ordered tree
(\citeproc{ref-Versloot2017}{Versloot 2017, 302--19};
\citeproc{ref-Bremmer2008}{Bremmer 2008, 284, 292--93}).

\section{A worked historical tree
test}\label{a-worked-historical-tree-test}

If a supposed shared innovation obligatorily follows a daughter-only
change in Frisian which English never underwent, it cannot occupy the
common stem under those premises. Either the ordering or scope premise
must change, the innovation happened independently, or the apparently
identical events are historically different. The required ancestor is
not removed.

Campbell reconstructs Frisian au contraction before ordinary fronting;
Goblirsch instead proposes a shared diphthong followed by Frisian stress
shift and contraction (\citeproc{ref-Campbell1939}{Campbell 1939, 91,
104}; \citeproc{ref-Goblirsch1991}{Goblirsch 1991, 17, 20--21}). The
daughter-only event is a predecessor in the first account and a follower
in the second. The second is compatible with sharing, but does not
attest the reconstructed Frisian intermediate. Surface agreement,
compatibility with inheritance, compelled parallelism under premises,
and unresolved history must therefore remain distinct judgments.

\section{The working account and why it is
used}\label{the-working-account-and-why-it-is-used}

The English reconstruction adopted as a working scaffold is the
conventional account developed by Campbell and critically refined by
Ringe and Taylor. It distinguishes inherited long vowels from the long
vowels produced by diphthong contraction, and it treats fronting,
breaking, restoration and mutation as conditioned phonological
developments. That scaffold is useful because it specifies contrasting
inputs and intermediate states, not because one author's authority
decides the history (\citeproc{ref-Campbell1959}{Campbell 1959, 52--53,
60--72}; \citeproc{ref-RingeTaylor2014}{Ringe and Taylor 2014, 170--73,
215--37}).

The comparative specialists test particular load-bearing premises.
Campbell's argument about English and Frisian fronting, Versloot's runic
interpretations and Goblirsch's proposed lost Frisian diphthong do not
form one automatically compatible package
(\citeproc{ref-Campbell1939}{Campbell 1939, 90--91, 104};
\citeproc{ref-Versloot2017}{Versloot 2017, 295--97, 318};
\citeproc{ref-Goblirsch1991}{Goblirsch 1991, 17--21}). Where an
alternative changes the inherited inventory or the identity of an event,
it must be evaluated as that alternative, not inserted into the
conventional chain without its premises.

CAPR implements regular sound laws. Prosodic and phonological conditions
may distinguish environments, but lexical diffusion, grammatical
categories and individually selected exceptions cannot repair a
contradiction. Paradigm remodeling and borrowing may affect which
historical form a lexical entry represents; they are not thereby
conditions on a sound law. This methodological choice is distinct from
faithfully reporting an author's different theoretical account.

Barrack's West Saxon cascade supplies a formal generative restatement of
parts of this interval. He states the Ingvaeonic nasal-spirant
development as rule 121, locating it in the Old English/Old Frisian area
with Old Saxon retaining free variation, with \emph{ōþer} `other' among
his witnesses (\citeproc{ref-Barrack1975}{Barrack 1975, 44}); the
fronting-and-rounding complex appears as rules 130a--c
(\citeproc{ref-Barrack1975}{Barrack 1975, 47--48}); and a loss of
\emph{*i} after \emph{*w} is marked ``apparently Anglo-Frisian'' as rule
406 (\citeproc{ref-Barrack1975}{Barrack 1975, 104}). Two cautions from
the Barrack dossier apply throughout. First, his changes are modelled as
windows of free variation with overlapping phases --- his worked example
is precisely the retraction/breaking overlap of rules 143a--b yielding
attested \emph{æll}/\emph{all}/\emph{eall} variation
(\citeproc{ref-Barrack1975}{Barrack 1975, 21--23, 53--54}) --- so his
rule boundaries do not map one-to-one onto CAPR's categorical laws.
Second, his orderings are formal, not chronological
(\citeproc{ref-Barrack1975}{Barrack 1975, 34}), and cannot by themselves
settle the stem-versus-daughter placement questions this chapter has
just outlined.

\section{Long vowels, diphthongs and the ancestral
inventory}\label{long-vowels-diphthongs-and-the-ancestral-inventory}

The inherited-long contrast is central. On the conventional
reconstruction, inherited oral long \emph{*ā} has already acquired a
front quality when English stressed \emph{*ai} completes its development
to a new back \emph{*ā}. Otherwise the new vowel should share the
inherited vowel's fronting. Ringe and Taylor formulate this as a
relation between temporal milestones: inherited-long fronting must be
well under way before contraction completes. They allow overlap; a
serial implementation is not proof that the changes occupied disjoint
periods (\citeproc{ref-Campbell1939}{Campbell 1939, 90};
\citeproc{ref-RingeTaylor2014}{Ringe and Taylor 2014, 170}).

The inherited vowel itself is disputed. Fulk's retained-front
reconstruction does not require precisely the same earlier backing and
re-fronting as Ringe and Taylor's account. That is a disagreement about
the input inventory, not a small adjustment to the position of an
otherwise identical rule (\citeproc{ref-Fulk2018}{Fulk 2018, 60--61};
\citeproc{ref-RingeTaylor2014}{Ringe and Taylor 2014, 10--13}). Nielsen
also proposes a different causal relation between contraction and
long-vowel restructuring, allowing co-occurrence rather than imposing
the conventional strict sequence (\citeproc{ref-Nielsen2001}{Nielsen
2001, 514--16, 521}).

Luick's proposed fronted-ai path is a further alternative, not an
observed intermediate series. His conjecture is
\emph{*æe > *æə > *æa > *ā}; the middle offglide is schwa. Campbell
rejects that explanation, but the rejection is not an independent
observation that such a phonetic path is impossible
(\citeproc{ref-Luick1914}{Luick 1914-\/-1940, 132--33};
\citeproc{ref-Campbell1939}{Campbell 1939, 90}, n.~3). The
inherited-long and secondary-long histories must consequently remain
separate even where later mutation makes their final vowels converge.

Nasal environments require another distinction. Nasalization, loss of a
nasal with compensatory lengthening, rounding and eventual merger are
not one event. The nasalized input can escape an oral-fronting law
before its later rounded outcome is established. A shared final vowel
therefore does not place every substep at the same ancestral node
(\citeproc{ref-Bremmer2009}{Bremmer 2009, 24--27};
\citeproc{ref-RingeTaylor2014}{Ringe and Taylor 2014, 142--46}).

The following table specifies the relevant inventory questions at a
conservative ancestral cut. It is a working reconstruction, not a newly
adjudicated full phoneme inventory. Earlier West Germanic developments
can be inherited at this cut without being exclusively Anglo-Frisian
innovations.

{\def\LTcaptype{none} % do not increment counter
\begin{longtable}[]{@{}
  >{\raggedright\arraybackslash}p{(\linewidth - 6\tabcolsep) * \real{0.2500}}
  >{\raggedright\arraybackslash}p{(\linewidth - 6\tabcolsep) * \real{0.2500}}
  >{\raggedright\arraybackslash}p{(\linewidth - 6\tabcolsep) * \real{0.2500}}
  >{\raggedright\arraybackslash}p{(\linewidth - 6\tabcolsep) * \real{0.2500}}@{}}
\toprule\noalign{}
\begin{minipage}[b]{\linewidth}\raggedright
Input class
\end{minipage} & \begin{minipage}[b]{\linewidth}\raggedright
Conservative working ancestral state
\end{minipage} & \begin{minipage}[b]{\linewidth}\raggedright
English daughter consequence
\end{minipage} & \begin{minipage}[b]{\linewidth}\raggedright
Remaining premise
\end{minipage} \\
\midrule\noalign{}
\endhead
\bottomrule\noalign{}
\endlastfoot
Inherited oral long vowel & Front long vowel under the conventional
account & Subsequent dialectal and palatal treatment & Earlier back
versus retained-front input \\
Nasal long vowel & Distinguished from the oral long vowel & Rounding and
merger must be dated separately & Which nasal substeps precede the
cut \\
Stressed ai & Retained diphthong before completed English contraction &
Secondary back long vowel, eligible for later mutation & Runic
interpretation and any limited ancestral onset \\
Ordinary short non-nasal a & Unfronted at this conservative cut &
Ordinary English fronting & Daughter contraction before fronting \\
Inherited au & Not identified with completed English or Frisian outcome
& English front element and later offglide development & Identity with
ordinary short fronting \\
Supported inherited Vww classes & Adopted West Germanic a/e/i-ww
reanalysis; separate homorganic uww quantity & Later English diphthong
realization, not inherited English surface vowels & Qualified long-u
shadow account; j-created classes remain separate \\
Velars and palatal tendencies & No blanket completed assibilation is
stipulated & Class-specific daughter histories & Articulation, cutoff
and merger are distinct \\
\end{longtable}
}

The evidence underlying these qualifications is component-specific
(\citeproc{ref-RingeTaylor2014}{Ringe and Taylor 2014, 41--42, 65--66,
170--73}; \citeproc{ref-Versloot2017}{Versloot 2017, 295--97, 318};
\citeproc{ref-Laker2007}{Laker 2007, 167--84}). The table does not
assign inherited-short fronting to the stem merely because the node is
called Anglo-Frisian.

The inherited-glide component is now implemented rather than left as a
possible future decomposition. \hyperref[rule-OEWWSimplification]{SC031
OEWWSimplification} belongs to the earlier West Germanic account
discussed in Chapter 2: short nonhomorganic \emph{*Vww} gives
\emph{*Vuw}, with retained consonantal \emph{*w}; homorganic \emph{*uww}
separately gives long \emph{*ūw}. The latter follows Ringe and Taylor's
preferred but qualified shadow analysis. These earlier products can be
inherited at the Anglo-Frisian node without placing the later English
long diphthongs there. The selected sentence context for early apocope
is independent of the segmental reconstruction and of lexical accent
(\citeproc{ref-RingeTaylor2014}{Ringe and Taylor 2014, 41--42, 55,
57--58, 65--66, 171--75}; \citeproc{ref-Fulk2018}{Fulk 2018, 117}).

\section{Alternative cuts and their tree
consequences}\label{alternative-cuts-and-their-tree-consequences}

A more expansive common-stem reconstruction is possible only with
different premises. In Luick's account, participation of the diphthong
nucleus in fronting can coexist with a conjectured later Frisian
development. In Goblirsch's account, a common fronted au diphthong
precedes Frisian stress shift and contraction. Both require historical
intermediates whose existence must be argued, not obtained from the name
of the subgroup (\citeproc{ref-Luick1914}{Luick 1914-\/-1940, 130--33};
\citeproc{ref-Goblirsch1991}{Goblirsch 1991, 17, 20--21}).

{\def\LTcaptype{none} % do not increment counter
\begin{longtable}[]{@{}
  >{\raggedright\arraybackslash}p{(\linewidth - 4\tabcolsep) * \real{0.3333}}
  >{\raggedright\arraybackslash}p{(\linewidth - 4\tabcolsep) * \real{0.3333}}
  >{\raggedright\arraybackslash}p{(\linewidth - 4\tabcolsep) * \real{0.3333}}@{}}
\toprule\noalign{}
\begin{minipage}[b]{\linewidth}\raggedright
Proposed common-stem content
\end{minipage} & \begin{minipage}[b]{\linewidth}\raggedright
Required premise
\end{minipage} & \begin{minipage}[b]{\linewidth}\raggedright
Conditional consequence
\end{minipage} \\
\midrule\noalign{}
\endhead
\bottomrule\noalign{}
\endlastfoot
Inherited-long restructuring, but retained ai and short a & Completed
English contraction is later & Compatible conservative cut \\
Ordinary short fronting & Its secure predecessors can also occur before
the split & Excluded if daughter-only English contraction obligatorily
precedes it \\
Shared fronted-au development & Frisian stress shift/contraction follows
that development & Compatible in Goblirsch's account, not proof of
inheritance \\
Shared consonantal palatalization & Corresponding consonant classes and
necessary predecessors are genuinely identical & Unresolved; apparent
similarities do not settle all layers \\
\end{longtable}
}

Campbell's English chain makes the second row especially consequential:
if completed contraction is English-only and ordinary fronting must
follow it, the latter cannot be inherited from the common stem. The
corresponding English and Frisian frontings are then parallel under
those premises. Changing a disputed premise can change that conclusion,
but cannot remove the required ancestor
(\citeproc{ref-Campbell1939}{Campbell 1939, 90--91};
\citeproc{ref-Versloot2017}{Versloot 2017, 295--97, 318}).

\section{Palatalization is not a single subgroup
character}\label{palatalization-is-not-a-single-subgroup-character}

A palatal articulation can precede phonemic contrast, and the productive
cutoff of a conditioning process can precede its eventual assibilated
reflex. Fricative-g merger with j is different again. Geminate and
postnasal stops, singleton fricatives, k, h and sk must not be treated
as one dated change simply because later spellings look palatal
(\citeproc{ref-Hogg1979}{Hogg 1979, 100--111};
\citeproc{ref-Laker2007}{Laker 2007, 167--68, 175--84}).

The shared-palatalization hypothesis must therefore identify which layer
is inherited. A daughter-specific predecessor can exclude a completed
component from the stem without excluding every earlier articulatory
tendency. Conversely, removing an argument against sharing establishes
compatibility, not that one common event actually occurred
(\citeproc{ref-Laker2007}{Laker 2007, 175--84}). The vowel
diphthongization following initial palatals is another event; Luick
explicitly distinguishes the similar Frisian and English outcomes
(\citeproc{ref-Luick1914}{Luick 1914-\/-1940, 162--63}).

The English implementation now distinguishes singleton palatal
fricatives from gg/ng stops and adopts postmutation fricative merger as
a working account (\citeproc{ref-RingeTaylor2014}{Ringe and Taylor 2014,
203--4}; \citeproc{ref-Fulk2018}{Fulk 2018, 130--32}). Hogg's objection
remains explicit (\citeproc{ref-Hogg1979}{Hogg 1979, 102--11}). This
daughter account does not select a complete ancestral inventory, exclude
every shared early articulatory tendency, or date stop affrication from
an eventual written reflex.

\section{What the runes date}\label{what-the-runes-date}

An inscription dates a written object. A sound-change terminus
additionally requires the reading, etymology, linguistic affiliation and
relation of the grapheme to phonetic or phonemic structure. The
interpretation of the Caistor inscription as retaining ai and the
interpretation of Frisian retention are consequential but conditional
(\citeproc{ref-Versloot2017}{Versloot 2017, 295--97, 318}). Waxenberger
likewise distinguishes early allophonic mutation from later
phonemicization and lacks immediate runic evidence for every short-vowel
split (\citeproc{ref-Waxenberger2019}{Waxenberger 2019, 63--74}).

No one calendar date is therefore attached here to every change
described as Anglo-Frisian. The English discussion that follows starts
from the working inventory and makes its necessary daughter developments
explicit. Canonical classification, preferred reconstruction and an
unresolved component are kept distinct until individual adjudication.

\section{Major changes and their historical
basis}\label{major-changes-and-their-historical-basis}

\subsection{West Germanic rhotacism
(SC003)}\label{west-germanic-rhotacism-sc003}

The medial change of \texttt{*z} to \texttt{*r} in environments such as
\Recon{déuzaz}\index[iv]{03pgmc@\ivlangheader{Proto-Germanic}{}!deuzaz@\iventry{*déuzaz}{}}
`deer',
\Recon{xúrdaz}\index[iv]{03pgmc@\ivlangheader{Proto-Germanic}{}!xurdaz@\iventry{*xúrdaz}{}}
`hoard', and
\Recon{líznōjaną}\index[iv]{03pgmc@\ivlangheader{Proto-Germanic}{}!liznojana@\iventry{*líznōjaną}{}}
`learn' is historically a post-Proto-West-Germanic development. Ringe
and Taylor argue that rhotacism was not inherited from Proto-Northwest
Germanic and was not uniform within West Germanic
(\citeproc{ref-RingeTaylor2014}{Ringe and Taylor 2014, 52, 98, 102}).
Crist separates this change explicitly from the deletion of word-final
\texttt{*z} and argues that rhotacism must follow the deletion rules
(\citeproc{ref-Crist2001}{Crist 2001, 104--6};
\citeproc{ref-Crist2002}{2002, 1, 4}). Bammesberger gives the standard
Old English--facing summary: \texttt{*z} yielded \texttt{*r} in
intervocalic position but was generally lost in final position
(\citeproc{ref-Bammesberger1992}{Bammesberger 1992, 39}).

The executable identifier \texttt{EAFRhotacism} is retained for
stability, not as independent evidence of historical stage or scope. The
preceding source discussion motivates the post-Proto-West-Germanic
account; the chapter's placement follows the adopted executable
sequence.

\subsection{\texorpdfstring{Word-final \texttt{*z} deletion (SC020) and
the three final-\texttt{*z}
developments}{Word-final *z deletion (SC020) and the three final-*z developments}}\label{word-final-z-deletion-sc020-and-the-three-final-z-developments}

The loss of word-final \texttt{*z} is not one process but three
historically distinct developments, and this chapter contains two of
them. The central one, SC020, is the Proto-West Germanic loss of
\texttt{*z} in unstressed syllables, seen in forms such as
\Recon{rástōz}\index[iv]{03pgmc@\ivlangheader{Proto-Germanic}{}!rastoz@\iventry{*rástōz}{}}
`rest (nom.sg.)' and stated for the whole branch by Ringe and Taylor
(\citeproc{ref-RingeTaylor2014}{Ringe and Taylor 2014, 44--45}); Crist's
analysis distinguishes it both from the earlier NWGmc changes and from
the later narrower Ingvaeonic deletion rules
(\citeproc{ref-Crist2002}{Crist 2002, 1, 4}). Earlier still, the
consonant-stem root nouns had generalized endingless nominatives before
Proto-West Germanic (SC096, Chapter 2). Later, and only in the north,
\texttt{*z} was lost in stressed monosyllables with compensatory
lengthening (SC097, this chapter); the southern dialects instead
retained and rhotacized it. The standard handbooks confirm the West
Germanic deletion in general terms: Bammesberger gives a clean statement
that Germanic \texttt{*z} is generally lost in final position
(\citeproc{ref-Bammesberger1992}{Bammesberger 1992, 39}); the three-way
division refines those summaries rather than contradicting them.

The CAPR rule for the unstressed loss is still named
\texttt{EAFFinalZDeletion}, an identifier that predates the restaging of
the change to Proto-West Germanic; the name is retained as a stable
identifier pending a global renaming pass, and the historical stage
recorded in the staging metadata takes priority over the name prefix.
SC020 and rhotacism jointly account for the fate of the remaining
\texttt{*z}, but their sections are separated by intervening executable
stages. Their thematic connection does not determine their publication
positions.

\subsection{Completed English ai contraction (SC004, treated in Chapter
4)}\label{completed-english-ai-contraction-sc004-treated-in-chapter-4}

The English outcome of stressed/root \emph{*ai} is \emph{*ā}. Versloot
distinguishes an early velar-conditioned Frisian contraction from later
non-velar treatment, with mutation and delabialization intervening
(\citeproc{ref-Versloot2017}{Versloot 2017, 302--9, 316}). His wave
account is an author's interpretation, not CAPR's own genealogy. CAPR
now places completed English contraction on the daughter branch, while
preserving the distinct possibility of a narrower conditioned ancestral
onset (\citeproc{ref-Campbell1939}{Campbell 1939, 90--91};
\citeproc{ref-Versloot2017}{Versloot 2017, 295--97, 318}). This local
decision does not select every component of the ancestral inventory.

The current \hyperref[rule-EAFAiMonophthongization]{SC004
EAFAiMonophthongization} models English contraction after
\hyperref[rule-EAFLongAFronting]{SC101 EAFLongAFronting}. The new
ai-derived \emph{*ā} escaped the earlier fronting of inherited oral
\emph{*ā}; Campbell states that inference and Ringe and Taylor retain it
(\citeproc{ref-Campbell1959}{Campbell 1959, 52--53}, §132;
\citeproc{ref-RingeTaylor2014}{Ringe and Taylor 2014, 169--70}). Soul
provides another implementation boundary at interstress raising, not the
sole historical argument or proof of inherited contraction.

Fulk offers a materially different reconstruction. He allows fronting of
the nucleus before contraction and prefers a surviving \emph{*æi}, with
back-vowel restoration and subsequent English generalization of an
unfronted paradigmatic alternant yielding \emph{*ā}. Paradigms with
consistent \emph{*j} conditioning need not participate in that
restoration (\citeproc{ref-Fulk1998}{Fulk 1998, 143--45}). This is not
Campbell's completed-contraction argument in different terminology: it
preserves a different intermediate and invokes paradigmatic levelling.
Ringe and Taylor's warning that diphthong nuclei need not behave like
the plain short vowel likewise qualifies, rather than independently
proves, Campbell's ordering (\citeproc{ref-RingeTaylor2014}{Ringe and
Taylor 2014, 170--75}). CAPR retains the conventional regular
phonological scaffold; adopting Fulk's alternative would require a
separate assessment of its conditioned alternations and the evidence for
levelling.

Fulk's reconstruction demonstrates that several objections to shared
Anglo-Frisian developments are not decisive under his premises, not that
those developments have been proved inherited. His final displayed
sequence is specifically Northumbrian, with qualifications for other
English dialects, and is not a complete West Saxon chronology
(\citeproc{ref-Fulk1998}{Fulk 1998, 153}). He explicitly calls his
reconstruction biased toward the Anglo-Frisian hypothesis, finds no
compelling reason for placing palatalization so early apart from that
interpretation, and states that the hypothesis is not proved
(\citeproc{ref-Fulk1998}{Fulk 1998, 154}). Those reservations are part
of his conclusion, not objections supplied only by later critics.

The unstressed development \texttt{*ai\ \textgreater{}\ *ē} (in final
and nonfinal syllables) is a separate and earlier Proto-Northwest
Germanic change (SC014), discussed in Chapter 2; its corpus witnesses
are the dative-singular endings of \texttt{span} (\Recon{spánnai} `span'
\textgreater{} \emph{spanne} `span') and \texttt{meed} (\Recon{mízdai}
`meed' \textgreater{} \emph{meorde} `meed').

\subsection{Anglo-Frisian brightening (SC043, treated in Chapter
4)}\label{anglo-frisian-brightening-sc043-treated-in-chapter-4}

The fronting of low \emph{*a} to \emph{*æ} outside nasal environments is
traditionally called Anglo-Frisian brightening. The name does not
establish one inherited event. It executes later in the cascade than the
changes of this chapter, so its full section appears in Chapter 4; it is
introduced here because it anchors the ``Anglo-Frisian'' label that
names this period. Campbell gives the classical statement: ``By a very
early change Prim. Gmc. \texttt{a\ \textgreater{}\ æ} in OE and OFris.
when not followed by a nasal consonant''
(\citeproc{ref-Campbell1959}{Campbell 1959, 52}, §131). The traditional
name also appears in modern comparative treatments; it does not
independently establish a shared event (\citeproc{ref-Fulk2018}{Fulk
2018, 72--73}).

The change is notable for what follows it: OE Breaking presupposes the
fronted input; OE a-Restoration partially undoes it in back-vowel
environments. The three-change sequence is the conventional English
working account, but its components and conditioning remain individually
arguable. Versloot proposes restricted early fronting and late
e-breaking; final forms can coincide under front-and-restore and
never-front histories (\citeproc{ref-Versloot2025}{Versloot 2026, 123,
125--28, 131--35}).

Campbell notes that English and Frisian may not simply reflect one
undifferentiated shared prehistoric event, and Ringe and Taylor leave
open whether the wider spread of fronted outcomes happened mainly on the
continent or in Britain (\citeproc{ref-RingeTaylor2014}{Ringe and Taylor
2014, 60--62}). CAPR's implementation treats the stressed English
component separately from the retained unstressed and final-vowel
representations.

The completed ordinary stressed component is now characterized on the
English daughter under Campbell's conventional working chain
(\citeproc{ref-Campbell1939}{Campbell 1939, 90--91};
\citeproc{ref-Campbell1959}{1959, 52--53}). Ringe and Taylor's
diphthong-nucleus objection remains explicit
(\citeproc{ref-RingeTaylor2014}{Ringe and Taylor 2014, 170--75});
earlier restricted ancestral tendencies are not excluded. English
plain-a/au process identity is distinct from inherited stem identity,
and their current serialization is retained. The unstressed and
final-vowel contributions are not assigned the same date from this
argument. This local adoption does not select the complete ancestral
inventory.

\section{Cascade vs.~historical order in this
chapter}\label{cascade-vs.-historical-order-in-this-chapter}

The sections follow executable serialization so that derivations are
inspectable. That order is not an independently demonstrated total
historical chronology. The broader final-z loss, later northern loss and
rhotacism are distinguished by their actual environments
(\citeproc{ref-RingeTaylor2014}{Ringe and Taylor 2014, 44--45};
\citeproc{ref-Crist2002}{Crist 2002, 1, 4}).

Earlier-stage proxies remain in this editorial corridor. The cluster
change represented by \hyperref[rule-PNWGmcMnDissimilation]{SC022
PNWGmcMnDissimilation} and final-n loss represented by
\hyperref[rule-PNWGmcNStemNLoss]{SC023 PNWGmcNStemNLoss} are classified
as Common Germanic/Proto-Germanic despite retained name prefixes
(\citeproc{ref-Fulk2018}{Fulk 2018, 121}; \citeproc{ref-Ringe2017}{Ringe
2017, 101--3}). The retired unstressed-o raising is not an active
member. Book order, computational position and historical stage must
remain separate.

\section{\texorpdfstring{West Germanic final
\emph{*z}-deletion}{West Germanic final -deletion}}\label{west-germanic-final--deletion}

\subsection{Historical discussion}\label{historical-discussion-18}

Word-final \emph{*z} in unstressed syllables was lost in Proto-West
Germanic. Ringe and Taylor state the change for the whole branch and
illustrate it with the nominative plural \emph{*dagōz} \textgreater{}
\emph{*dagō} and the consonant-stem nominative \emph{*fadurz}
\textgreater{} \emph{*fadur}, noting that the ending is lost after
consonants as well as after vowels (\citeproc{ref-RingeTaylor2014}{Ringe
and Taylor 2014, 44--45}, §3.1.1). Crist's handout formulates the same
development and its Ingvaeonic sequels (\citeproc{ref-Crist2002}{Crist
2002, 2}, §§5--6). The change is pan-West-Germanic, not specifically
Ingvaeonic: every West Germanic daughter shows the loss, and the
Frienstedt comb inscription \emph{kaba} \textless{} \emph{*kambaz}
`comb' (c.~250--300 CE) supplies early epigraphic confirmation
(\citeproc{ref-Fulk2018}{Fulk 2018, 25}, n.~1).

The conditioning segment is specifically the voiced sibilant \emph{*z},
never \emph{*s}: Ringe and Taylor's near-minimal pair of nominative
singular \emph{*dagaz} \textgreater{} \emph{*dag} beside genitive
singular \emph{*dagas}, which keeps its sibilant into Old English
\emph{dæġes}, shows that the change reads the Verner voicing distinction
(\citeproc{ref-RingeTaylor2014}{Ringe and Taylor 2014, 212}, §6.1).
Where the handbooks disagree about whether a given ending had \emph{*-s}
or \emph{*-z} --- as for the nominative plural \emph{*-ōz} --- the
disagreement matters directly to whether this rule applies
(\citeproc{ref-RingeTaylor2014}{Ringe and Taylor 2014, 115--16},
§4.2.1).

This is the middle of three historically distinct final-\emph{*z}
developments, and Ringe and Taylor explicitly separate it from the later
loss in stressed monosyllables, citing Crist's demonstration that they
are two changes (\citeproc{ref-RingeTaylor2014}{Ringe and Taylor 2014,
44--45}, §3.1.1). Earlier, the consonant-stem (root-noun) nominatives of
monosyllables had already generalized endinglessness before Proto-West
Germanic (\hyperref[rule-RootNounNomZLoss]{SC096 RootNounNomZLoss}), so
forms like \emph{*bōkz} `book' never reach this rule with their marker
intact. Later, and only in the north, \emph{*z} was lost in stressed
monosyllables with compensatory lengthening
(\hyperref[rule-MonosyllabicFinalZLoss]{SC097 MonosyllabicFinalZLoss});
the present rule leaves stressed monosyllables untouched. Older accounts
that grouped all of these under one loss of final \emph{*z}, such as
Campbell's, are superseded by this three-way division
(\citeproc{ref-Campbell1959}{Campbell 1959, 166}). Barrack's cascade
stands closer to the older grouping: his rules 404a--e treat rhotacism
and the losses and retentions of \emph{*z}/R as one West Germanic family
(\citeproc{ref-Barrack1975}{Barrack 1975, 101--2}), with the deletion of
R after a vowel as rule 404c (\citeproc{ref-Barrack1975}{Barrack 1975,
102}), and he states final \emph{*a} syncope separately as West Germanic
rule 331.4b (\citeproc{ref-Barrack1975}{Barrack 1975, 78}). His
subdivision thus cuts the material differently from the three-way
division adopted here; since his rule numbers are Campbell paragraph
anchors and his orderings formal rather than chronological
(\citeproc{ref-Barrack1975}{Barrack 1975, 21, 34}), the difference is
one of analysis, not of additional evidence against the
Ringe--Taylor--Crist sequence.

At the boundary with Chapter 2's Northwest Germanic sequence, the
derivation of \emph{ræste} `rest' shows that final \emph{*ō}-raising
(\hyperref[rule-PNWGmcFinalLongORaising]{SC019 PNWGmcFinalLongORaising})
must precede this rule: raising applies to \emph{*-ō} but not to
\emph{*-ōz}, whose final vowel is still sheltered by the sibilant when
raising runs (\citeproc{ref-RingeTaylor2014}{Ringe and Taylor 2014,
15--16, 24}). On the later side, Ringe and Taylor order the loss of
\emph{*z} before the loss of word-final bare \emph{*-a}, since
\emph{*dagaz} first becomes \emph{*daga} and only then \emph{*dag}
(\citeproc{ref-RingeTaylor2014}{Ringe and Taylor 2014, 45--46}, §3.1.2).

\subsection{\texorpdfstring{SC020. West Germanic final
\emph{*z}-deletion
(\texttt{EAFFinalZDeletion})}{SC020. West Germanic final -deletion (EAFFinalZDeletion)}}\label{rule-EAFFinalZDeletion}

\begin{ReaderFacingFoma}
define EAFFinalZDeletion Ctx([
    [{*z} -> 0 ||
    .#. ?* EnglishStarVocalic [EnglishStarConsonant | EnglishPalatalConsonant]+ [EnglishStarVocalic | {*ô}] ?* _ .#.,
    .#. [EnglishStarConsonant | EnglishPalatalConsonant]* EnglishStarVocalic+ [EnglishStarConsonant | EnglishPalatalConsonant]+ _ .#.]
]);
\end{ReaderFacingFoma}

The rule deletes word-final \emph{*z} in unstressed syllables, stated
through two environments. The first clause covers polysyllables, where
the final syllable of these corpus forms is unstressed: this is the
ordinary case, with 110 corpus derivations, such as PGmc \Recon{bárdaz}
`beard' on its way to OE \emph{beard} `beard' and \Recon{rástōz} `rest'
on its way to \emph{ræste} `rest'. The second clause covers
post-consonantal \emph{*z} in monosyllables. By the time this rule runs,
\hyperref[rule-RootNounNomZLoss]{SC096 RootNounNomZLoss} has already
removed the genuine root-noun nominative endings, so the only form
reaching the second clause is \Recon{fríjōndz} `friend', contracted to
monosyllabic \emph{*fríundz} by \hyperref[rule-PWGmcIjContraction]{SC009
PWGmcIjContraction}; its ending, like that of \emph{*fadurz}, stood in
an unstressed syllable when the Proto-West Germanic change applied and
so belongs here rather than to the root-noun development
(\citeproc{ref-RingeTaylor2014}{Ringe and Taylor 2014, 44--45}, §3.1.1).
Stressed monosyllables ending in vowel plus \emph{*z} meet neither
clause and are left for \hyperref[rule-MonosyllabicFinalZLoss]{SC097
MonosyllabicFinalZLoss}.

The polysyllabic environment counts trimoric vowels as nuclei as well as
bimoric vowels. Ringe and Taylor's contrast between the trimoric
nominative plural ending in \emph{*gebō̄z} and the bimoric genitive
singular ending in \emph{*gebōz} requires loss of \emph{*z} in both,
without merging their quantities: the former leads to OE \emph{ġiefa}
`gifts', the latter to \emph{ġiefe} `of a gift'
(\citeproc{ref-RingeTaylor2014}{Ringe and Taylor 2014, 59}). The extra
trimoric symbol in the first environment represents this quantity
distinction; it does not extend the rule to stressed monosyllables.

The chronology of word-final \emph{*z}-loss is unusually well delimited:
\emph{ræste} `rest' supplies its early boundary, while later weak
syllables supply its late boundary.

Before \hyperref[rule-PNWGmcFinalLongORaising]{SC019
PNWGmcFinalLongORaising}, PGmc \Recon{rástōz} `rest' yields \Pred{rast}
rather than expected OE \emph{ræste} `rest'. After
\hyperref[rule-OEMedUnstressedULowering]{SC040
OEMedUnstressedULowering}, PGmc \Recon{bébruz} `beaver' yields
\Pred{befro} rather than expected \emph{befer} `beaver', PGmc
\Recon{kwéðuz} `cud' yields \Pred{cwedo} rather than expected
\emph{cwedu} `cud', and PGmc \Recon{félθuz} `field' yields \Pred{feldo}
rather than expected \emph{feld} `field', alongside eight other newly
failing rows. Final \emph{z}-loss therefore follows
\hyperref[rule-PNWGmcFinalLongORaising]{SC019 PNWGmcFinalLongORaising}
and precedes \hyperref[rule-OEMedUnstressedULowering]{SC040
OEMedUnstressedULowering}.

The \Recon{rástōz} `rest' derivation fixes the local relation to
\hyperref[rule-PNWGmcFinalLongORaising]{SC019 PNWGmcFinalLongORaising}.
The distant boundary at \hyperref[rule-OEMedUnstressedULowering]{SC040
OEMedUnstressedULowering} shows only that word-final \emph{*z}-loss
precedes the later weak-syllable sequence; its placement within that
wider interval follows the handbook chronology after final
\emph{*ō}-raising.

\section{Early apocope in unstressed
words}\label{early-apocope-in-unstressed-words}

\subsection{Historical discussion}\label{historical-discussion-19}

Alongside the regular loss of word-final short high vowels in third
syllables (\hyperref[rule-PWGmcEarlyIApocope]{SC006
PWGmcEarlyIApocope}), Ringe and Taylor identify a second, earlier
apocope: ``Short high vowels were also lost after heavy syllables in
unstressed words'' (\citeproc{ref-RingeTaylor2014}{Ringe and Taylor
2014, 57--58}, §3.1.4). The two laws must not be conflated. Fully
stressed disyllables kept their final \emph{*-i} long enough to cause
i-umlaut --- OE \emph{ġiest} `guest' \textless{} \emph{*gastiz} and
\emph{fȳr} `fire' \textless{} \emph{*fūri} require exactly that survival
(\citeproc{ref-RingeTaylor2014}{Ringe and Taylor 2014, 55}, §3.1.4) ---
whereas words that carried no sentence stress lost the vowel already in
Proto-West Germanic.

The conditioning is prosodic. Like Verner's law, the change is governed
by accent: it applied in words unstressed in the sentence, and its
apparent exceptions are systematic, not sporadic. Forms such as OE
\emph{ymbe} `around' and OHG \emph{umbi} kept their final vowel because,
as Ringe and Taylor observe, proclitics ``were not phonologically
word-final'' and so stood outside the environment altogether
(\citeproc{ref-RingeTaylor2014}{Ringe and Taylor 2014, 57--58}, §3.1.4).
Sentence-level accent placement therefore decides which sandhi variant
each daughter language continues, and doublets across the family reflect
the stressed and unstressed sentence forms of the same word --- regular
sandhi, not lexical diffusion.

The diagnostic witness is the second-person plural pronoun. Ringe and
Taylor print the Proto-West Germanic form as a doublet: PGmc
\emph{*izwiz} (Gothic \emph{izwis}) → \emph{*iwwi} (by
\hyperref[rule-PWGmcCoronalWAssimilation]{SC008
PWGmcCoronalWAssimilation} and the loss of final \emph{*z},
\hyperref[rule-EAFFinalZDeletion]{SC020 EAFFinalZDeletion}) → PWGmc
\emph{*iuwi} \textasciitilde{} \emph{*iuw}
(\citeproc{ref-RingeTaylor2014}{Ringe and Taylor 2014, 41--42}, §3.1.1).
Old English continues the apocopated, unstressed variant, and Ringe and
Taylor's proof is the vocalism itself: ``OE iow `you (dat. pl.)'
definitely does {[}show early apocope{]} (since it does not exhibit
i-umlaut)'' (\citeproc{ref-RingeTaylor2014}{Ringe and Taylor 2014,
57--58}, §3.1.4). Had the \emph{*-i} survived, i-umlaut
(\hyperref[rule-OEIUmlaut]{SC055 OEIUmlaut}) would have fronted the
diphthong; West Saxon \emph{ēow} `you' beside early West Saxon and
Northumbrian \emph{īow} `you' shows the normal unumlauted development
(\citeproc{ref-Campbell1959}{Campbell 1959, § 702}, p.~283).

\subsection{\texorpdfstring{\CAPRRuleHeading{SC098. Early apocope in unstressed words}{PWGmcUnstressedWordFinalIApocope}}{}}\label{rule-PWGmcUnstressedWordFinalIApocope}

\begin{ReaderFacingFoma}
define PWGmcUnstressedWordFinalIApocope [
    [ [{*i}|{*u}] -> 0 || .#. {*ᵘ} ?* [
        [EnglishStarLongVowel | EnglishStarLongDiphthong] EnglishStarConsonant*
        | EnglishStarShortVowel EnglishStarConsonant EnglishStarConsonant+
    ] _ .#. ]
    .o. [{*ᵘ} -> 0]
    .o. [
        [{*ᶜ}:0 [EnglishStarAlphabet - {*ᶜ}]* 0:{*ᶜ}]
        | [EnglishStarAlphabet - {*ᶜ}]*
    ]
];
\end{ReaderFacingFoma}

Sentence stress and phonological finality are now explicit,
independently selected context. They do not alter the PGmc
reconstruction, the selected segmental input or its historical stage.
The computational marks in the rule distinguish weak-final context from
weak-nonfinal context; they are not reconstructed phonemes and are
absent from displayed forms. Ordinary evaluation explicitly chooses
strong-final citation context. You selects the independently
source-discussed weak-final variant, not an unstressed classification
inferred from a pronoun label or missing acute.

The worked history is \emph{*izwiz} → \emph{*iwwiz} by assimilation →
\emph{*iuwiz} by \hyperref[rule-OEWWSimplification]{SC031
OEWWSimplification} → \emph{*iuwi} by final \emph{z}-loss → weak-final
\emph{*iuw} by this rule → OE \emph{ēow}. The operation requires a heavy
syllable and a phonologically final short high vowel: long
vowel/diphthong weight or a short vowel followed by a closing consonant
cluster, not literal \emph{ww}. Its condition therefore survives glide
reanalysis. Loss of final \emph{z} still feeds it, and removal of
\emph{i} bleeds \hyperref[rule-OEIUmlaut]{SC055 OEIUmlaut}. Its PWGmc
placement precedes the later northern loss of stressed-monosyllabic
final \emph{z} (\hyperref[rule-MonosyllabicFinalZLoss]{SC097
MonosyllabicFinalZLoss}) ((\citeproc{ref-RingeTaylor2014}{Ringe and
Taylor 2014, 41--42, 55, 57--58})).

The staged strong contexts retain the high vowel in the histories of
\emph{ġiest} `guest' and \emph{fȳr} `fire'. A weak-final and context
yields \emph{and} `and', whereas the proposed proclitic context retains
the vowel in \emph{ymbe} through the later apocope corridor. The
right-boundary annotation is removed only after
\hyperref[rule-OEHighVowelApocope]{SC063 OEHighVowelApocope} has had the
opportunity to apply. These controls implement the defended
exceptionless account; they do not claim direct observation of
prehistoric sentence stress or lexical optionality
((\citeproc{ref-RingeTaylor2014}{Ringe and Taylor 2014, 55, 57--58})).

The corresponding English realization of \emph{iu+w} belongs to
\hyperref[rule-OEDiphthongLeveling]{SC032 OEDiphthongLeveling}, not a
late vocalization of retained \emph{ww}. The model's retained-i
counterfactual produces an umlauted result after intervening w-loss; its
exact terminal spelling is a computational prediction, not an attested
strong OE form. The attested weak-final target and the source's mutation
argument must not be replaced by that counterfactual
((\citeproc{ref-RingeTaylor2014}{Ringe and Taylor 2014, 41--42, 57--58,
173--75}; \citeproc{ref-Campbell1959}{Campbell 1959, 283})).

\section{\texorpdfstring{Northern monosyllabic final
\emph{*z}-loss}{Northern monosyllabic final -loss}}\label{northern-monosyllabic-final--loss}

\subsection{Historical discussion}\label{historical-discussion-20}

Long after the Proto-West Germanic loss of final \emph{*z} in unstressed
syllables (\hyperref[rule-EAFFinalZDeletion]{SC020 EAFFinalZDeletion}),
the northern West Germanic dialects lost word-final \emph{*z} in
stressed monosyllables as well, with compensatory lengthening of a short
nucleus. Ringe and Taylor's witness set is OE \emph{mā} `more'
\textless{} \emph{*maiz}, the pronouns \emph{wē} `we', \emph{ġē} `you',
\emph{mē} `me', \emph{þē} `thee', \emph{hē} `he', \emph{hwā} `who'
\textless{} \emph{*hwaz}, and --- hedged in their own print with
question marks --- \emph{cū} `cow' (\citeproc{ref-RingeTaylor2014}{Ringe
and Taylor 2014, 86}, §3.3.1). The southern dialects retained the
sibilant and rhotacized it: Old High German \emph{mir}, \emph{wir},
\emph{mēr}, \emph{er} answer the Old English endingless forms, which is
why Fulk counts this loss among the diagnostic Ingvaeonic features
(\citeproc{ref-Fulk2018}{Fulk 2018, 18}, n.~6). Ringe and Taylor
explicitly treat this as a change separate from the Proto-West Germanic
unstressed loss, citing Crist's demonstration that the two must be
distinguished (\citeproc{ref-RingeTaylor2014}{Ringe and Taylor 2014,
44--45}, §3.1.1).

The scholarship disagrees about the exact conditioning, and the
disagreement is worth recording. An older account, represented by
Campbell and going back to Luick, derived the endingless pronouns from
unaccented sentence variants rather than from a regular sound change
(\citeproc{ref-Campbell1959}{Campbell 1959, 166};
\citeproc{ref-Luick1914}{Luick 1914-\/-1940, 819}). Ringe and Taylor
reject that analysis because \emph{mā} `more' and \emph{cū} `cow' are
not plausibly unaccented words (\citeproc{ref-RingeTaylor2014}{Ringe and
Taylor 2014, 86}, §3.3.1). Crist formulates an Ingvaeonic rule in which
\emph{*z} is lost after front vowels, with compensatory lengthening,
covering preconsonantal cases as well; his data contain no word-final
back-vowel monosyllables, so forms like \emph{*hwaz} and the ancestor of
\emph{cū} fall outside what his statement can decide --- a documented
gap rather than a refutation (\citeproc{ref-Crist2002}{Crist 2002, 1,
4}, §§1, 10). Kilday narrows the preconsonantal subcase to Old Saxon and
Old Frisian while accepting the word-final monosyllabic loss for Old
English, contrasting regular \emph{meord} `reward' with the
loanword-influenced \emph{mēd} `reward'
(\citeproc{ref-Kilday2024}{Kilday 2024, 1--3}). CAPR adopts Ringe and
Taylor's quality-neutral formulation because it alone generates the
back-vowel witnesses, while noting that the front-vowel forms are
compatible with both analyses.

Apparent counterexamples are analogical, not phonological: OE
\emph{dēor} `deer', \emph{ār} `oar', and \emph{gār} `spear' show final
\emph{-r} from levelling out of inflected forms where the sibilant was
word-internal and regularly rhotacized, not from retention of word-final
\emph{*z} (\citeproc{ref-RingeTaylor2014}{Ringe and Taylor 2014, 86},
§3.3.1, n.~24). The change precedes rhotacism
(\hyperref[rule-EAFRhotacism]{SC003 EAFRhotacism}), which Ringe and
Taylor place last in this sequence of northern developments
(\citeproc{ref-RingeTaylor2014}{Ringe and Taylor 2014, 87}, §3.3.1).

The corpus now witnesses this change directly. The interrogative pronoun
`who' is selected in its nominative singular masculine cell, PGmc
\emph{*hwaz} (Gothic \emph{hwas}), precisely the form Ringe and Taylor
cite for this loss (\citeproc{ref-RingeTaylor2014}{Ringe and Taylor
2014, 86}, §3.3.1): the rule lengthens the short nucleus and deletes the
sibilant, giving \emph{*hwā}, whence OE \emph{hwā} `who'. The resulting
back vowel never undergoes Anglo-Frisian brightening --- ``\emph{*hwǣ}
does not exist'', as Campbell puts it
(\citeproc{ref-Campbell1959}{Campbell 1959, § 125}, p.~49;
\citeproc{ref-SieversBrunner1965}{Brunner 1965, § 137} Anm. 1,
p.~129) --- so the derivation ends with the attested form. Other members
of Ringe and Taylor's witness set remain outside the corpus because it
selects oblique or plural cells for them --- `cow' and `meed', for
instance, enter the cascade in inflected forms whose \emph{*z}, where
present, is word-internal.

\subsection{\texorpdfstring{\CAPRRuleHeading{SC097. Northern monosyllabic final \emph{*z}-loss}{MonosyllabicFinalZLoss}}{}}\label{rule-MonosyllabicFinalZLoss}

\begin{ReaderFacingFoma}
define MonosyllabicFinalZLoss [
    {*a} -> {*ā}, {*á} -> {*ā},
    {*e} -> {*ē}, {*é} -> {*ḗ},
    {*i} -> {*ī}, {*í} -> {*ḯ},
    {*o} -> {*ō}, {*ó} -> {*ō},
    {*u} -> {*ū}, {*ú} -> {*ū}
        || .#. [EnglishStarConsonant | EnglishPalatalConsonant]*
            _ {*z} .#.
] .o. [
    {*z} -> 0 ||
        .#. [EnglishStarConsonant | EnglishPalatalConsonant]*
            EnglishStarVocalic+ _ .#.
];
\end{ReaderFacingFoma}

The rule first lengthens a short nucleus standing immediately before
word-final \emph{*z} in a monosyllable, then deletes the \emph{*z} after
any vowel in a monosyllable. The principal synthetic controls run the
change on Ringe and Taylor's own witnesses, fed to the rule in their
chronologically correct intermediate shapes. Short-nucleus inputs show
both halves of the change at once: \emph{*hwaz} yields \emph{*hwā} and
\emph{*hiz} yields \emph{*hī}, each with loss of the sibilant and
compensatory lengthening (\citeproc{ref-RingeTaylor2014}{Ringe and
Taylor 2014, 86}, §3.3.1). A form whose nucleus is already bimoric skips
the lengthening step and simply loses the sibilant: \emph{*maiz} yields
\emph{*mai} at this stage, with its diphthong intact; the attested OE
\emph{mā} `more' arises only later, when the stressed monophthongization
(\hyperref[rule-EAFAiMonophthongization]{SC004 EAFAiMonophthongization})
takes \emph{*ai} to \emph{*ā}. The word `cow', whose history Ringe and
Taylor themselves print with question marks and whose analysis remains
disputed, is deliberately not used as a principal control; a long-vowel
input of that shape (\emph{*kūz} yielding \emph{*kū}) merely repeats
what \emph{*maiz} already demonstrates, and the word's evidentiary
weight is discussed in the historical dossier rather than leaned on
here.

The corpus derivation of \emph{hwā} `who' now fixes this rule's position
empirically as well as philologically. It stands after
\hyperref[rule-EAFFinalZDeletion]{SC020 EAFFinalZDeletion}, since the
two losses are historically distinct changes with the unstressed loss
earlier (\citeproc{ref-RingeTaylor2014}{Ringe and Taylor 2014, 44--45},
§3.1.1). Rhotacism (\hyperref[rule-EAFRhotacism]{SC003 EAFRhotacism})
follows in the executable cascade as in the historical account: Ringe
and Taylor place rhotacism after this loss, at the end of the sequence
of \emph{*z}-losses (\citeproc{ref-RingeTaylor2014}{Ringe and Taylor
2014, 87}, §3.3.1), and the cascade composes it immediately after this
rule, so a sibilant removed here can never surface as \emph{-r} --- were
the order reversed, \emph{*hwaz} would rhotacize to \emph{*hwar} and
`who' could never be derived; the negative controls below confirm this.
Consonant-final monosyllables are untouched: their nominative
\emph{*-z}, where it ever existed, was eliminated before Proto-West
Germanic under \hyperref[rule-RootNounNomZLoss]{SC096 RootNounNomZLoss}.
Word-internal \emph{*z}, as in PGmc \Recon{déuzą} `deer' on its way to
OE \emph{dēor} `deer', does not meet the environment of this rule and
duly rhotacizes.

\section{\texorpdfstring{Monosyllabic lengthening before final
\emph{r}}{Monosyllabic lengthening before final }}\label{monosyllabic-lengthening-before-final}

The model admits Old English \emph{þǣr} `there' and \emph{hwǣr} `where'
from source-supported short-\emph{a} antecedents. Their quantity and the
date of the lengthening remain qualified: the forms do not establish a
Proto-Germanic date, and they do not determine the order of
\hyperref[rule-PNWGmcILowering]{SC015 PNWGmcILowering}.

\subsection{\texorpdfstring{\CAPRRuleHeading{SC112. Stressed monosyllabic lengthening before inherited final \emph{r}}{PNWGmcMonosyllabicFinalRLengthening}}{}}\label{rule-PNWGmcMonosyllabicFinalRLengthening}

\begin{ReaderFacingFoma}
define PNWGmcMonosyllabicFinalRLengthening [
    [
        {*á} -> {*ā},
        {*é} -> {*ē},
        {*í} -> {*ī},
        {*ó} -> {*ō},
        {*ú} -> {*ū}
        || .#. EnglishStarConsonant* _ {*r} .#.
    ]
];
\end{ReaderFacingFoma}

The early working placement preserves inherited final \emph{r} before
rhotacism, so secondary \emph{r} does not trigger the rule. The guard
also excludes disyllabic forms. This is a general formalization of the
two admitted witnesses, not a lexical exception list or a claim of
strict relative chronology against \hyperref[rule-PNWGmcILowering]{SC015
PNWGmcILowering} or \hyperref[rule-PNWGmcLongELowering]{SC024
PNWGmcLongELowering}.

\section{West Germanic rhotacism}\label{west-germanic-rhotacism}

\subsection{Historical discussion}\label{historical-discussion-21}

Bammesberger states that Germanic \emph{*z} yielded \emph{*r} in
intervocalic position in Old English, while final \emph{*z} was
generally lost (\citeproc{ref-Bammesberger1992}{Bammesberger 1992, 39}).
Ringe and Taylor argue that this merger of \emph{*z} with \emph{*r} was
independent in Norse and West Germanic and belongs after the
Proto-West-Germanic stage (\citeproc{ref-RingeTaylor2014}{Ringe and
Taylor 2014, 52, 98, 102}). Crist likewise places rhotacism after
earlier West Germanic \emph{*z}-deletion rules and rejects treating it
as an inherited Proto-Northwest-Germanic innovation
(\citeproc{ref-Crist2001}{Crist 2001, 104--6};
\citeproc{ref-Crist2002}{2002, 1, 4}).

The internal identifier \hyperref[rule-EAFRhotacism]{SC003 EAFRhotacism}
places the change in CAPR's Early Anglo-Frisian corridor, the
operational post-Proto-West-Germanic stage on the English line;
historically the change is a West Germanic rhotacism, later than
Proto-Germanic. It is also distinct from
\hyperref[rule-EAFFinalZDeletion]{SC020 EAFFinalZDeletion}, which
removes final \emph{*z} before the surviving medial consonant becomes
\emph{*r}.

Barrack states the same complex as his 404 family: \emph{*z} becomes a
rhotic R as Proto-Germanic rule 404a (\citeproc{ref-Barrack1975}{Barrack
1975, 101}), and R becomes \emph{r} as West Germanic rule 404e
(\citeproc{ref-Barrack1975}{Barrack 1975, 102}), with the deletions of R
stated between them. His two-step treatment, in which an intermediate
rhotic persists while the deletions apply, is structurally parallel to
the ordering argument from Crist adopted here, though he assigns the
first step to Proto-Germanic where Ringe and Taylor argue for
independent Norse and West Germanic mergers. He also registers the
material directly relevant to this rule's \emph{meorde} witness: his
rule 404b, which he grades as resting on ``very sporadic'' evidence,
deletes the rhotic before a voiced dental with compensatory lengthening
and so derives \emph{mēd} `reward', while the alternate development that
yields \emph{meord} is referred to his rule 139zc
(\citeproc{ref-Barrack1975}{Barrack 1975, 101--2}). His orderings are
formal, not chronological (\citeproc{ref-Barrack1975}{Barrack 1975,
34}).

\subsection{\texorpdfstring{SC003. West Germanic rhotacism
(\texttt{EAFRhotacism})}{SC003. West Germanic rhotacism (EAFRhotacism)}}\label{rule-EAFRhotacism}

\noindent\begin{minipage}{\linewidth}
\begin{ReaderFacingFoma}
define EAFRhotacism [
    {*z} -> {*r} || EnglishStarVocalic _ ?
];
\end{ReaderFacingFoma}
\end{minipage}

Breaking supplies the decisive upper boundary. If rhotacism is delayed
until after \hyperref[rule-OEBreaking]{SC044 OEBreaking}, PGmc
\Recon{mízdai} `meed' yields \Pred{merde} rather than expected OE
\emph{meorde} `meed'. Moving rhotacism earlier within the tested range
changes no output.

The lexical evidence thus supplies a terminus ante quem but no terminus
post quem. The lower boundary rests on the historical analyses cited
above: Ringe and Taylor put rhotacism at the end of the sequence of
\emph{*z}-losses --- ``first \emph{*z} was lost in a variety of
environments \ldots, then all surviving \emph{*z} became \emph{*r}''
(\citeproc{ref-RingeTaylor2014}{Ringe and Taylor 2014, 87}, §3.3.1) ---
and Crist observes that the deletions distinguish \emph{*z} from
\emph{*r} and so must precede the merger (\citeproc{ref-Crist2002}{Crist
2002, 2--3}). The rule is accordingly ordered after the three
final-\emph{*z} losses (\hyperref[rule-RootNounNomZLoss]{SC096
RootNounNomZLoss}, \hyperref[rule-EAFFinalZDeletion]{SC020
EAFFinalZDeletion}, and \hyperref[rule-MonosyllabicFinalZLoss]{SC097
MonosyllabicFinalZLoss}) and before breaking, so that the derivation
follows the reconstructed chronology.

\section{\texorpdfstring{\emph{mn}-dissimilation}{-dissimilation}}\label{dissimilation}

\subsection{Historical discussion}\label{historical-discussion-22}

In the inherited \emph{n}-stem paradigm the zero-grade oblique cells
brought \emph{m} and \emph{n} into direct contact, and in that adjacent
cluster the labial nasal dissimilated to a labial spirant: \emph{mn}
\textgreater{} \emph{βn} (surfacing as \emph{fn}). Old Norse preserves
the older paradigmatic distribution, with the labial confined to the
oblique cluster (\emph{himinn} `heaven' beside dative \emph{hifni}); Old
English and Old Saxon generalized it. Fulk treats the cluster change
among developments common to Germanic, while warning that its surface
results are irregular and that reverse \emph{bn} \textgreater{}
\emph{mn} is later well attested in Northwest Germanic
(\citeproc{ref-Fulk2018}{Fulk 2018, 121}, §6.11). The relevant
\emph{heofon} `heaven' and \emph{mōnaþ} `month' material is discussed by
Campbell (\citeproc{ref-Campbell1959}{Campbell 1959, 189, 195}, §§470,
484).

The underlying cluster change is therefore late Proto-Germanic / Common
Germanic, not securely a pan-Northwest-Germanic innovation;
\texttt{PNWGmc} remains a stable executable identifier only. The lexical
evidence does not constrain a positive local cascade position.

\subsection{\texorpdfstring{SC022. Dissimilation of adjacent \emph{mn}
(\texttt{PNWGmcMnDissimilation})}{SC022. Dissimilation of adjacent  (PNWGmcMnDissimilation)}}\label{rule-PNWGmcMnDissimilation}

\noindent\begin{minipage}{\linewidth}
\begin{ReaderFacingFoma}
define PNWGmcMnDissimilation [
    {*m} -> {*β}
        || EnglishStarVocalic _ {*n}
];
\end{ReaderFacingFoma}
\end{minipage}

The rule fires only where \emph{m} stands directly before \emph{n}. It
supplies the labial of \emph{stefn} `stem, trunk' from the
\emph{mn}-cluster of the \emph{stamn}-family, and it is the historical
change behind the labial of \emph{heofon} `heaven', which was
generalized from the oblique cluster into the vowel-bearing stem before
the Old English vocalic changes. (An earlier cross-syllable formulation
that labialized an intervocalic \emph{m} before a later nasal has been
retired: it simulated paradigm levelling rather than a sound law.)

Moving \hyperref[rule-PNWGmcMnDissimilation]{SC022
PNWGmcMnDissimilation} earlier or later leaves every output unchanged.
Its executable place in this holding zone is therefore
editorial/computational, while its historical classification rests on
the handbook account of \emph{mn}-dissimilation.

\section{\texorpdfstring{Word-final
\emph{n}-loss}{Word-final -loss}}\label{word-final--loss}

\subsection{Historical discussion}\label{historical-discussion-23}

The change isolated here is far older than its position in the cascade
suggests: it is the general (pre-)Proto-Germanic loss of word-final
\emph{*n}, with nasalization of the preceding vowel, in polysyllables.
Ringe's proof set for the law spans the whole grammar --- nouns such as
\emph{*yugón} \textgreater{} \emph{*juką} `yoke', pronouns such as
\emph{*tón} \textgreater{} \emph{*þanǭ}, and even the verb form
\emph{*dedǭ} `I did' --- so it is general phonology, not a fact about
any one declension (\citeproc{ref-Ringe2017}{Ringe 2017, 101--3}).
Gothic \emph{tuggo} shares the weak nominative-singular outcome, and the
nasalized reflex \emph{*-ǭ} remained contrastive into Proto-West
Germanic before yielding OE \emph{-e}
(\citeproc{ref-RingeTaylor2014}{Ringe and Taylor 2014, 54--55, 58--59}).

Within the present corpus the change surfaces in exactly one shape: the
weak nouns are cited in the stem form \emph{*-ōn-}, and this rule
carries them to the Proto-Germanic nominative singular in \emph{*-ǭ}, as
in \emph{*túngōn} \textgreater{} \emph{*túngǭ} \textgreater{}
\emph{tunge} `tongue', alongside \emph{eorþe} `earth', \emph{heorte}
`heart', \emph{nǣdre} `adder', and thirteen further weak nouns. The
masculine weak nominative singular in trimoric \emph{*-ô} never had a
final \emph{*-n} to lose, and Proto-Germanic \emph{*sebun} `seven',
\emph{*nigun} `nine', and \emph{*tehun} `ten' kept their \emph{-n} by
lexical analogy among the numerals (\citeproc{ref-Ringe2017}{Ringe 2017,
103}); the rule's narrow \emph{*-ōn} environment leaves all of these
correctly untouched.

\subsection{\texorpdfstring{SC023. Loss of word-final \emph{*n} after
\emph{*ō}
(\texttt{PNWGmcNStemNLoss})}{SC023. Loss of word-final  after  (PNWGmcNStemNLoss)}}\label{rule-PNWGmcNStemNLoss}

\noindent\begin{minipage}{\linewidth}
\begin{ReaderFacingFoma}
define PNWGmcNStemNLoss [
    {*ō} {*n} -> {*ǭ} || _ .#.
];
\end{ReaderFacingFoma}
\end{minipage}

The verb \emph{dōn} `do' supplies the negative, counterfeeding witness
for the chronology. PGmc \Recon{dōną} `do' passes this rule untouched;
only \hyperref[rule-OEHeavySyllableNasalApocope]{SC047
OEHeavySyllableNasalApocope} later strips the final \Recon{ą} and
creates a new word-final \Recon{-ōn}. That secondary \Recon{-n} survives
into \emph{dōn} precisely because the old loss was no longer active: if
\hyperref[rule-PNWGmcNStemNLoss]{SC023 PNWGmcNStemNLoss} is displaced
after the apocope, it consumes the new nasal and the derivation
collapses entirely (\emph{+?}).

The retained \emph{-n} of \emph{dōn} `do' therefore supplies a terminus
ante quem for the loss --- it must be dead before the apocope --- while
the seventeen weak nouns above are its positive witnesses; the lower
boundary remains unattested within the cascade, as befits a change
already complete in Proto-Germanic.

\section{Nasal-spirant lengthening}\label{nasal-spirant-lengthening}

\subsection{Historical discussion}\label{historical-discussion-24}

Germanic lost nasal consonants before voiceless fricatives twice, in two
changes that are easily confused because their outcomes look alike. Both
replace a sequence of vowel, nasal and fricative with a long nasalized
vowel and the fricative. They differ in date, in geography, and in which
fricatives they affect.

The earlier change is common Germanic and is treated in the opening
chapter of this book, on the Proto-Germanic loss of a nasal before
\Recon{x}. Its results are shared by every daughter language, and only
\Recon{a}, \Recon{i} and \Recon{u} occur in its input
(\citeproc{ref-Campbell1959}{Campbell 1959, 44}, §119;
\citeproc{ref-Fulk2018}{Fulk 2018, 55}, §4.1;
\citeproc{ref-Ringe2017}{Ringe 2017, 149--50}, §3.2.7).

The later change belongs to the dialects bordering the North Sea, that
is to Old English, Old Frisian and Old Saxon, the group traditionally
called Ingvaeonic. Here the groups \Recon{mf}, \Recon{ns} and \Recon{nþ}
likewise reject the nasal with compensatory lengthening and nasalization
(\citeproc{ref-Campbell1959}{Campbell 1959, 47}, §121;
\citeproc{ref-Fulk2018}{Fulk 2018, 72}, §4.11;
\citeproc{ref-SieversBrunner1965}{Brunner 1965, 176}, §186.1;
\citeproc{ref-Luick1914}{Luick 1914-\/-1940, 276}, §301.1). Ringe and
Taylor call it the most obvious phonological innovation of the northern
dialects and list some thirty examples, among them \Recon{gans} `goose'
and \Recon{jugunþi} `youth' (\citeproc{ref-RingeTaylor2014}{Ringe and
Taylor 2014, 139--41}). Campbell describes it as a later change similar
to the common Germanic one, and Fulk as comparable to it; neither treats
them as the same law.

The two are told apart by what the languages outside the North Sea area
show. The common Germanic change left no nasal anywhere, so Old High
German has \emph{fūht} `damp' and \emph{fūst} `fist' exactly as Old
English does. The later change was confined to the north, so the
cognates of its witnesses keep the nasal: Old High German \emph{fimf},
\emph{gans}, \emph{ander}, \emph{jugund} answer Old English \emph{fīf}
`five', \emph{gōs} `goose', \emph{ōþer} `other', \emph{ġeoguþ} `youth'.
On this test \Recon{funhsti-} `fist', whose long vowel is shared by Old
High German \emph{fūst}, Dutch \emph{vuist} and German \emph{Faust},
belongs to the earlier change and not to the North Sea law at all
(\citeproc{ref-Kroonen2013}{Kroonen 2013, 160}).

Where the vowel was \Recon{a}, the long nasalized vowel that this law
produced was afterwards rounded to \emph{ō} in Anglo-Frisian, which is
why Old English has \emph{gōs} `goose', \emph{tōþ} `tooth' and
\emph{ōþer} `other'. That rounding is a third change again, and it is
the same rounding that gives \emph{fōn} `seize' and \emph{þōhte}
`thought' from the common Germanic law and \emph{mōna} `moon' and
\emph{spōn} `chip' from inherited long \emph{ā} before a surviving
nasal; Campbell states that it reached all three sources at one and the
same time (\citeproc{ref-Campbell1959}{Campbell 1959, 50}, §128 n.~1;
\citeproc{ref-SieversBrunner1965}{Brunner 1965, 33}, §26;
\citeproc{ref-Fulk2018}{Fulk 2018, 72}, §4.11). It is treated in the
chapter on the rounding of the long nasalized low vowel. That the vowel
was rounded and did not simply merge shows that its nasality survived
the law that created it: as Fulk observes, it did not fall together with
the \emph{ā} that came from \Recon{ai} (\citeproc{ref-Fulk2018}{Fulk
2018, 55}, §4.1). Ringe and Taylor take the nasalization to have
remained subphonemic until it was lost separately in each daughter
(\citeproc{ref-RingeTaylor2014}{Ringe and Taylor 2014, 141}). Old Saxon
shares the loss of the nasal and the nasalization, and rounds only
variably, which is why the law itself is described as North Sea Germanic
and the rounding as Anglo-Frisian (\citeproc{ref-Campbell1959}{Campbell
1959, 47}, §121; \citeproc{ref-Fulk2018}{Fulk 2018, 72}, §4.11).

The North Sea Germanic law is a single connected sound change: in every
handbook account the nasal is lost \emph{with} compensatory lengthening,
and the lengthening is the compensation for the loss. It is stated here
as two rules only because the vowel must be adjusted while the
conditioning nasal is still present, before the nasal can be removed.
The order of \hyperref[rule-EAFNasalSpirantLengthening]{SC026
EAFNasalSpirantLengthening} before
\hyperref[rule-EAFNasalSpirantLoss]{SC027 EAFNasalSpirantLoss} is a
requirement of the statement, not evidence for two successive historical
stages.

Barrack's cascade concurs on both the geography and the two-law
distinction. He states the North Sea development as Ingvaeonic rule 121,
locating it in the Old English and Old Frisian area with Old Saxon
retaining free variation, and cites \emph{ōþer} among his witnesses
(\citeproc{ref-Barrack1975}{Barrack 1975, 44}); he orders it explicitly
after his Proto-Germanic rule 119, the common Germanic nasal loss before
the velar fricative (\citeproc{ref-Barrack1975}{Barrack 1975, 42--44}).
Two differences of framework apply: he models the change as a window of
free variation rather than a categorical law
(\citeproc{ref-Barrack1975}{Barrack 1975, 21--23}), and his orderings
are formal rather than chronological (\citeproc{ref-Barrack1975}{Barrack
1975, 34}).

\subsection{\texorpdfstring{\CAPRRuleHeading{SC026. North Sea Germanic nasal-spirant law, first step}{EAFNasalSpirantLengthening}}{}}\label{rule-EAFNasalSpirantLengthening}

\begin{ReaderFacingFoma}
define EAFNasalSpirantLengthening [
    {*a} -> {*ã} || _ EnglishStarNasal EnglishStarNSGmcSpirant,
    {*i} -> {*ī} || _ EnglishStarNasal EnglishStarNSGmcSpirant,
    {*u} -> {*ū} || _ EnglishStarNasal EnglishStarNSGmcSpirant,
    # Stressed variants (*á -> *ã: stress absorbed into length, no stressed
    # long low tier exists):
    {*á} -> {*ã} || _ EnglishStarNasal EnglishStarNSGmcSpirant,
    {*í} -> {*ī} || _ EnglishStarNasal EnglishStarNSGmcSpirant,
    {*ú} -> {*ū} || _ EnglishStarNasal EnglishStarNSGmcSpirant
];
\end{ReaderFacingFoma}

The environment is nasal plus \Recon{f}, \Recon{þ} or \Recon{s}. The
fricative \Recon{x} is excluded: nasal loss before \Recon{x} is the
earlier change stated as \hyperref[rule-PGmcNasalLossBeforeX]{SC103
PGmcNasalLossBeforeX}. Campbell names the groups \Recon{mf}, \Recon{ns}
and \Recon{nþ}, Fulk says the change affects \Recon{mf}, \Recon{ns} and
\Recon{nþ}, Sievers and Brunner name the fricatives \Recon{f}, \Recon{þ}
and \Recon{s}, and no example in Ringe and Taylor's list contains
\Recon{x} (\citeproc{ref-Campbell1959}{Campbell 1959, 47}, §121;
\citeproc{ref-Fulk2018}{Fulk 2018, 72}, §4.11;
\citeproc{ref-SieversBrunner1965}{Brunner 1965, 176}, §186.1;
\citeproc{ref-RingeTaylor2014}{Ringe and Taylor 2014, 139--41}). As in
the earlier change, only \Recon{a}, \Recon{i} and \Recon{u} occur. The
outcome of \Recon{a} is the long nasalized \Recon{ą̄}, which
\hyperref[rule-EAFNasalizedLowRounding]{SC104 EAFNasalizedLowRounding}
later rounds; the outcomes of \Recon{i} and \Recon{u} are written
\Recon{ī} and \Recon{ū}, their nasality having no further consequence
(\citeproc{ref-Campbell1959}{Campbell 1959, 47}, §121).

Two witnesses apply in the present corpus. PGmc \Recon{gánsz} `goose'
becomes \Recon{gą̄ns} `goose', and PGmc \Recon{júgunθ} `youth' becomes
\Recon{júgūnθ} `youth'. In \emph{ġeoguþ} `youth' the syllable carrying
the lengthened vowel is unstressed, and the length is given up again by
the later shortening of unstressed syllables; Sievers and Brunner note
the same course in \emph{beraþ} `they carry' from \Recon{beranþi} `they
carry' through \Recon{berōþ} `they carry'
(\citeproc{ref-SieversBrunner1965}{Brunner 1965, 176}, §186.1 Anm. 3;
\citeproc{ref-Luick1914}{Luick 1914-\/-1940, 276}, §301.1).

If lengthening is stated after the loss of the nasal, it can no longer
read its conditioning environment. The vowel must therefore be adjusted
before that nasal is removed. This is an executable dependency within
the statement of one historical change, not a date for either component.
The earlier experiment's youth output also depended on the retired early
glide-insertion architecture; it is not a current historical boundary.

\section{Nasal-spirant loss}\label{nasal-spirant-loss}

\subsection{\texorpdfstring{\CAPRRuleHeading{SC027. North Sea Germanic nasal-spirant law, second step}{EAFNasalSpirantLoss}}{}}\label{rule-EAFNasalSpirantLoss}

\noindent\begin{minipage}{\linewidth}
\begin{ReaderFacingFoma}
define EAFNasalSpirantLoss [
    EnglishStarNasal -> 0 || _ EnglishStarNSGmcSpirant
];
\end{ReaderFacingFoma}
\end{minipage}

The nasal is removed in the environment that conditioned the
lengthening, giving \Recon{gą̄s} `goose' and \Recon{júgūθ} `youth'. The
rule completes the statement of the single change begun in
\hyperref[rule-EAFNasalSpirantLengthening]{SC026
EAFNasalSpirantLengthening}; the two are not independent sound laws. The
converse test, stating the loss first, merely reproduces the same two
wrong forms. Nothing in the material fixes a later boundary for the
loss.

\section{\texorpdfstring{The hiatus-breaking \emph{w} of the verba
pura}{The hiatus-breaking  of the verba pura}}\label{the-hiatus-breaking-of-the-verba-pura}

\subsection{Historical discussion}\label{historical-discussion-25}

The small class of Germanic strong verbs whose roots ended in a vowel,
the \emph{verba pura}, reached Northwest Germanic with a morphologically
expected hiatus. Ringe and Taylor reconstruct Proto-Germanic
\emph{*sēaną} `to sow' (Gothic \emph{saian}), which the Northwest
Germanic lowering of \emph{*ē₁} carried to \emph{*sāaną}; the West
Germanic languages then ``exhibit innovative consonants that
eliminated'' the hiatus (\citeproc{ref-RingeTaylor2014}{Ringe and Taylor
2014, 12}). The repair differs by branch, and the difference dates and
localizes the change. Old English and Old Frisian inserted \emph{w}, as
in Old English \emph{sāwan} `to sow' and Old Frisian \emph{sāwinge}
`sowing' beside \emph{grōwinge} `growth', while Old Saxon and Old High
German used \emph{j} (\emph{sāian} `to sow', \emph{sāen} `to sow',
\emph{sājen} `to sow') (\citeproc{ref-RingeTaylor2014}{Ringe and Taylor
2014, 12, 151}). The insertion of \emph{w} is therefore an Anglo-Frisian
development and the consonant of \emph{sāwan} is not inherited; it must
not be projected back into the protoform.

The origin of that \emph{w} is neither inherited nor morphological.
Þórhallsdóttir, in the standard treatment of the problem, rejects in
turn the derivation of the West Germanic \emph{w} from a
Proto-Indo-European \emph{u}-perfect
(\citeproc{ref-Thorhallsdottir1993}{Þórhallsdóttir 1993, 115--17}), from
a \emph{*ue/o}-present
(\citeproc{ref-Thorhallsdottir1993}{Þórhallsdóttir 1993, 117--19}), and
from the transfer of a preterite \emph{w} into the present
(\citeproc{ref-Thorhallsdottir1993}{Þórhallsdóttir 1993, 121--22}). The
\emph{w} ``rather has an inner-Germanic phonological explanation''
(\citeproc{ref-Thorhallsdottir1993}{Þórhallsdóttir 1993, 136}). It began
as an ordinary phonetic glide filling the hiatus, and it began in one
narrow environment, before \emph{*u}. The cells in question, as Bremer
had already seen, are the first singular and first plural of the present
indicative together with the plural of the reduplicated preterite in
\emph{*-um}, \emph{*-uþ}, \emph{*-un}
(\citeproc{ref-Thorhallsdottir1993}{Þórhallsdóttir 1993, 120}).
Þórhallsdóttir sets out the pre-Old English present indicative with the
hiatus still intact, first singular \emph{*sā'u}, second \emph{*sā'is},
third \emph{*sā'iþ}, first plural \emph{*sā'um}, second plural
\emph{*sā'iþ}, third plural \emph{*sā'anþ}, and observes that ``the
glide \emph{w} was a natural hiatus filler before the \emph{u} of the
endings of the 1st person singular and plural''
(\citeproc{ref-Thorhallsdottir1993}{Þórhallsdóttir 1993, 127}). The
preterite stem \emph{*seuwun} arose by the same insertion before the
\emph{*u} of the plural endings
(\citeproc{ref-Thorhallsdottir1993}{Þórhallsdóttir 1993, 126}).

From those few cells the \emph{w} was carried through the whole paradigm
by analogy. Three forces converge. Levelling proceeded out of the
\emph{*u}-cells themselves; the preterite stem \emph{*seuw},
\emph{*seuwun} ``would have assisted in extending the new \emph{w} to
all forms of the present paradigm'' and to the preterite participle; and
verbs with an etymological \emph{*w}, above all \emph{flōwan} `to flow'
from the root \emph{*pleu-}, exerted their own attraction
(\citeproc{ref-Thorhallsdottir1993}{Þórhallsdóttir 1993, 127}). The
result, complete already in pre-Old English, is the uniform pattern
\emph{*sāwan}, \emph{*seuw}, \emph{*seuwun}, \emph{*sāwan-}
(\citeproc{ref-Thorhallsdottir1993}{Þórhallsdóttir 1993, 136}). Old
Frisian inherits the same generalized type, so that its \emph{w} and the
Old English \emph{w} are one Anglo-Frisian event
(\citeproc{ref-Thorhallsdottir1993}{Þórhallsdóttir 1993, 130, 134}). The
scattered \emph{w}-forms of Old High German are a separate and later
inner-German development and are not to be identified with it
(\citeproc{ref-Thorhallsdottir1993}{Þórhallsdóttir 1993, 119, 123}).

This layering matters for how the rule below should be read. Old English
citation forms are infinitives, and the infinitive \emph{*sā'an} is
precisely one of the cells that never received \emph{w} by the
phonological rule, so the \emph{w} of \emph{sāwan} `to sow' is
analogical. The rule therefore models the outcome of the generalization,
and its environment, any vowel after \emph{ā}, is a citation-form proxy
for a paradigm-wide result. It is not a reconstruction of the
conditioning of the original insertion.

The chronology is fixed on both sides by Ringe and Taylor. The \emph{w}
must postdate the West Germanic loss of intervocalic \emph{*w}, or the
new glide would itself have been swept away
(\citeproc{ref-RingeTaylor2014}{Ringe and Taylor 2014, 151}, n.~9); and
it ``must have occurred early enough to prevent fronting of \emph{*ā}''
in pre-Old English (\citeproc{ref-RingeTaylor2014}{Ringe and Taylor
2014, 151}). The whole reason \emph{sāwan} `to sow', \emph{cnāwan} `to
know', \emph{blāwan} `to blow', and \emph{māwan} `to mow' keep their
back vowel is that the \emph{w} was already in place when the North Sea
Germanic fronting applied. Since it is the generalized \emph{w} that
stands in the infinitive, this later boundary constrains the
generalization itself. The rule accordingly stands before the fronting
in the cascade, even though on the present corpus the two orders happen
to produce the same outputs, since the fronting rule as implemented does
not touch a prevocalic \emph{ā} in any case; the ordering encodes the
historical chronology, and no corpus-internal contrast demonstrates it.
That chronology rests on Ringe and Taylor alone, for Þórhallsdóttir's
chapter nowhere discusses the fronting of \emph{*ā} and supplies no
independent evidence for the order.

\subsection{\texorpdfstring{\CAPRRuleHeading{SC102. Generalized hiatus-breaking \emph{w} after long \emph{ā}}{EAFHiatusWInsertion}}{}}\label{rule-EAFHiatusWInsertion}

\noindent\begin{minipage}{\linewidth}
\begin{ReaderFacingFoma}
define EAFHiatusWInsertion [
    [..] -> {*w} || {*ā} _ EnglishStarVocalic
];
\end{ReaderFacingFoma}
\end{minipage}

The rule is fed by \hyperref[rule-PNWGmcLongELowering]{SC024
PNWGmcLongELowering}, which creates the \emph{ā}-initial hiatus it
repairs; \Recon{sḗaną} `to sow' passes through \emph{*sāaną} to
\emph{*sāwaną} on its way to \emph{sāwan} `to sow'. Displaced before the
lowering, the rule can never apply, since the root vowel is still
\emph{*ē} and no hiatus after \emph{ā} exists, and the derivation loses
its consonant altogether.

Its output in turn feeds the blocking environment of
\hyperref[rule-EAFLongAFronting]{SC101 EAFLongAFronting}: the inserted
\emph{w} is precisely what shields the \emph{ā} of \emph{sāwan} `to sow'
from fronting, as treated in that chapter. The corpus carries
\emph{sāwan} as the diagnostic witness of the class; \emph{cnāwan} `to
know', \emph{blāwan} `to blow', \emph{māwan} `to mow', \emph{wāwan} `to
blow (of wind)', and \emph{þrāwan} `to turn' instantiate the same
derivation (\citeproc{ref-RingeTaylor2014}{Ringe and Taylor 2014, 151}).

\section{\texorpdfstring{North Sea Germanic nasalization of long
\emph{ā} before a
nasal}{North Sea Germanic nasalization of long  before a nasal}}\label{north-sea-germanic-nasalization-of-long-before-a-nasal}

\subsection{Historical discussion}\label{historical-discussion-26}

In the dialects along the North Sea coast the long low vowel \emph{*ā}
--- centrally the vowel produced from \emph{*ē₁} by the Northwest
Germanic lowering --- was nasalized when a nasal consonant followed and
survived. The nasalized vowel was afterwards rounded in Anglo-Frisian,
which is why Old English has \emph{mōna} `moon', \emph{mōnaþ} `month'
and \emph{spōn} `spoon' against Old High German \emph{māno} `moon',
\emph{mānōd} `month', \emph{spān} `spoon' and Old Norse \emph{máni}
`moon', \emph{mánaðr} `month', \emph{spánn} `spoon'. Campbell describes
the split of Germanic \emph{ǣ¹} before nasals in these terms and
identifies the vowel that the rounding operated on as a nasalized and
unrounded \Recon{ą̄} {[}Campbell (\citeproc{ref-Campbell1959}{1959}),
p.~50, §127; p.~50, §128 n.~1{]}. Fulk places the rounded outcome of
\emph{*ǣ} before nasals among the Anglo-Frisian changes and derives it
from an earlier nasalized vowel (\citeproc{ref-Fulk2018}{Fulk 2018,
72--73}, §4.12).

The nasalization and the rounding have different geographies, and
separating them resolves an apparent disagreement in the handbooks.
Ringe and Taylor state that stressed low vowels were nasalized in the
northern West Germanic dialects, Old Saxon among them
(\citeproc{ref-RingeTaylor2014}{Ringe and Taylor 2014, 142}, §5.1.2);
Old Saxon accordingly shows \emph{ōdar} `other' and \emph{sōd} `true'
beside unrounded \emph{quān} `wife' and \emph{sāno} `immediately'
(\citeproc{ref-RingeTaylor2014}{Ringe and Taylor 2014, 150--51}). What
Old Saxon shares is the nasalization; what it shares only in part is the
rounding, and for the vowel inherited from Proto-Germanic it does not
share the rounding at all (\citeproc{ref-Campbell1959}{Campbell 1959,
44}, §119; \citeproc{ref-Fulk2018}{Fulk 2018, 72}, §4.11). The
nasalization stated here is therefore North Sea Germanic, and the
rounding treated in the chapter on the long nasalized low vowel is
Anglo-Frisian.

The nasalization is the conditioned counterpart of the fronting of oral
\emph{*ā} treated in the fronting chapter: a following nasal gives
nasalization and eventual rounding, and its absence gives fronting.
Together the two exhaust the fate of the old low vowel in this area. The
comparative material does not order the two branches against each other,
and no such ordering is claimed here.

\subsection{\texorpdfstring{\CAPRRuleHeading{SC025. Nasalization of long \emph{ā} before nasals}{EAFLongANasalRounding}}{}}\label{rule-EAFLongANasalRounding}

\noindent\begin{minipage}{\linewidth}
\begin{ReaderFacingFoma}
define EAFLongANasalRounding [
    {*ā} -> {*ã} || _ EnglishStarNasal
];
\end{ReaderFacingFoma}
\end{minipage}

The rule consumes the \emph{*ā} created by
\hyperref[rule-PNWGmcLongELowering]{SC024 PNWGmcLongELowering}:
displacing the lowering after this rule leaves the nasalization without
an input, and \Recon{mḗnōθz} `month' surfaces as \Pred{mānaþ} in place
of OE \emph{mōnaþ} `month', \Recon{spḗnuz} `spoon' as \Pred{spān} in
place of \emph{spōn} `spoon'. Its output is consumed in turn by
\hyperref[rule-EAFNasalizedLowRounding]{SC104 EAFNasalizedLowRounding},
which supplies the rounded vowel that the two words actually show.

Equally important is what the rule must precede. The monophthongization
of \emph{*ai} in \hyperref[rule-EAFAiMonophthongization]{SC004
EAFAiMonophthongization} creates a new long \emph{ā}, and that vowel was
never nasalized and never rounded before nasals: \emph{stān} `stone' and
\emph{hām} `home' keep \emph{ā}. Stating the monophthongization before
this rule makes the new vowel eligible for nasalization and hence for
rounding, and the cascade then wrongly yields \Pred{stōn} and
\Pred{hōm}. This is the same chronological inference Campbell draws for
the fronting: the treatments of the old low vowel were complete, or at
least under way, before \emph{ai}-monophthongization supplied a new one
(\citeproc{ref-Campbell1959}{Campbell 1959, 52--53}, §132;
\citeproc{ref-RingeTaylor2014}{Ringe and Taylor 2014, 169--70}).

\section{Anglo-Frisian rounding of the long nasalized low
vowel}\label{anglo-frisian-rounding-of-the-long-nasalized-low-vowel}

\subsection{Historical discussion}\label{historical-discussion-27}

Three separate developments described in earlier chapters end in the
same sound: a long, nasalized, low vowel. The oldest is the
Proto-Germanic loss of a nasal before \Recon{x}; the second is the North
Sea Germanic nasal-spirant law; the third is the nasalization of the
inherited long \emph{ā} before a nasal that survived. In Old English and
Old Frisian all three surface as \emph{ō}, and Campbell states the
unification without qualification: the nasalized vowel became identical
with the \emph{ō} inherited from Proto-Germanic already in prehistoric
Old English, and the same change affected the nasalized vowel of the
Proto-Germanic law and the nasalized vowel of the Ingvaeonic law at one
and the same time (\citeproc{ref-Campbell1959}{Campbell 1959, 50}, §128
n.~1). One later sound change therefore accounts for all three, and the
appearance of three independent roads to \emph{ō} is an illusion created
by looking only at the Old English surface.

That the intermediate vowel was nasalized and unrounded is shown by a
form in which it was shortened very early. Campbell's example is
\emph{samcucu} `half alive', where the shortened reflex is \emph{a} and
never \emph{o}; the vowel that the earlier laws produced was therefore a
nasalized \Recon{ą̄}, and the rounding is a distinct and later event
(\citeproc{ref-Campbell1959}{Campbell 1959, 50}, §128 n.~1). Sievers and
Brunner describe the same nasalized long vowel and its rounded Old
English outcome {[}Brunner (\citeproc{ref-SieversBrunner1965}{1965}),
p.~33, §26; p.~58, §64{]}.

The geography is Anglo-Frisian. Old Saxon nasalized its stressed low
vowels along with the rest of the northern West Germanic area
(\citeproc{ref-RingeTaylor2014}{Ringe and Taylor 2014, 142}, §5.1.2),
and it shares the loss of the nasal in both of the earlier laws; the
categorical, systematic rounding, however, it does not share. For the
Proto-Germanic nasalized low vowel Old Saxon retains \emph{ā}
consistently, and for the vowel created by the nasal-spirant law it has
\emph{ā} or \emph{ō} according to word and dialect, while Old English
and Old Frisian have \emph{ō} throughout
(\citeproc{ref-Campbell1959}{Campbell 1959, 44}, §119;
\citeproc{ref-Fulk2018}{Fulk 2018, 72}, §4.11). That partial and
lexically variable Old Saxon rounding is comparative evidence bearing on
the innovation rather than participation in it. Ringe treats the
rounding as a parallel development of the diverging northern dialects
and locates it in the northernmost of them, that is in Anglo-Frisian
(\citeproc{ref-Ringe2017}{Ringe 2017, 149--50}, §3.2.7;
\citeproc{ref-RingeTaylor2014}{Ringe and Taylor 2014, 142}, §5.1.2).
Luick groups the whole set of changes among those peculiar to the
Anglo-Frisian dialect group and notes that the nasalized vowel later
gave up its nasality (\citeproc{ref-Luick1914}{Luick 1914-\/-1940, 276},
§301.1). The rounding is accordingly Anglo-Frisian, and the nasalization
that feeds it is the wider North Sea Germanic property.

Ringe and Taylor confirm that a single rounding covers all the sources:
the rounding affected the nasalized low vowels of the nasal-spirant law,
the nasalized low vowels of the accompanying list, and the reflexes of
Proto-Germanic \emph{*/anh/} alike (\citeproc{ref-RingeTaylor2014}{Ringe
and Taylor 2014, 142}, §5.1.2). The condition on the input is nasality
and nothing else, which is why the long \emph{ā} that Old English later
won from \Recon{ai} escapes: that vowel was oral, and it came into being
after the rounding had run its course. Fulk makes this the central
argument for the long survival of the nasality, since the nasalized
vowel developed to \emph{ō} and did not fall together with Old English
\emph{ā} from \Recon{ai} (\citeproc{ref-Fulk2018}{Fulk 2018, 55}, §4.1).

\subsection{\texorpdfstring{\CAPRRuleHeading{SC104. Rounding of the long nasalized low vowel}{EAFNasalizedLowRounding}}{}}\label{rule-EAFNasalizedLowRounding}

\noindent\begin{minipage}{\linewidth}
\begin{ReaderFacingFoma}
define EAFNasalizedLowRounding [
    {*ã} -> {*ō}
];
\end{ReaderFacingFoma}
\end{minipage}

The rule is unconditioned, since the vowel it operates on exists only
where one of the three earlier changes created it. Its inputs arrive
from \hyperref[rule-PGmcNasalLossBeforeX]{SC103 PGmcNasalLossBeforeX},
from \hyperref[rule-EAFNasalSpirantLengthening]{SC026
EAFNasalSpirantLengthening} and from
\hyperref[rule-EAFLongANasalRounding]{SC025 EAFLongANasalRounding}, and
each of those three stands in a feeding relation to it: stated before
any one of them, the rule leaves that source's nasalized vowel untouched
and the cascade returns no form at all for its witnesses.

Four lexemes in the present corpus reach Old English through this rule,
and between them they witness all three sources. \Recon{gánsz} `goose'
arrives as \Recon{gą̄s} `goose' from the North Sea Germanic law and gives
\emph{gōs} `goose'; \Recon{mḗnōθz} `month' arrives as \Recon{mą̄nōþ}
`month' and gives \emph{mōnaþ} `month'; \Recon{spḗnuz} `spoon' arrives
as \Recon{spą̄nu} `spoon' and gives \emph{spōn} `spoon'. The
Proto-Germanic law contributes \Recon{θánxtē} `thought', the preterite
of the verb `to think', which arrives as \Recon{θą̄xtē} `thought' and
gives \emph{þōhte} `thought'. That is the very form cited for this
rounding, beside Old Frisian \emph{thochte}, and it is the evidence that
the vowel did not fall together with the \emph{ā} of \emph{stān} `stone'
(\citeproc{ref-Fulk2018}{Fulk 2018, 55}, §4.1;
\citeproc{ref-Campbell1959}{Campbell 1959, 44}, §119). The other firing
of the Proto-Germanic law, the high vowel of \emph{fȳst} `fist', does
not reach this rule at all.

The counterpart is what the rule leaves alone. \Recon{stáinaz} `stone'
and \Recon{xáimaz} `home' acquire their long \emph{ā} from
\hyperref[rule-EAFAiMonophthongization]{SC004 EAFAiMonophthongization},
which follows this rule, and their vowel was never nasalized; they
surface as \emph{stān} `stone' and \emph{hām} `home'. Placing the
monophthongization before the nasalization and the rounding makes that
vowel eligible and the cascade then yields \Pred{stōn} and \Pred{hōm},
which is the chronological inference Campbell draws for the treatments
of the old low vowel generally (\citeproc{ref-Campbell1959}{Campbell
1959, 52--53}, §132; \citeproc{ref-RingeTaylor2014}{Ringe and Taylor
2014, 169--70}).

\section{\texorpdfstring{North Sea Germanic fronting of long
\emph{ā}}{North Sea Germanic fronting of long }}\label{north-sea-germanic-fronting-of-long}

\subsection{Historical discussion}\label{historical-discussion-28}

Long after the Northwest Germanic lowering of \emph{*ē₁} to \emph{*ā},
the dialects of the North Sea coast fronted the surviving oral \emph{*ā}
to a low front vowel: West Saxon \emph{ǣ} in \emph{dǣd} `deed',
\emph{slǣpan} `to sleep', \emph{lǣtan} `to let', \emph{rǣdan} `to read',
and \emph{mǣl} `meal', against Anglian, Kentish, and Old Frisian
\emph{ē} (\emph{dēd} `deed', \emph{slēpa} `to sleep', \emph{jēr} `year')
(\citeproc{ref-RingeTaylor2014}{Ringe and Taylor 2014, 146--50};
\citeproc{ref-Campbell1959}{Campbell 1959, 50--51}, §128). Old Saxon and
Old High German keep the back vowel (Old Saxon \emph{dād} `deed', Old
High German \emph{tāt} `deed', \emph{slāfan} `to sleep', \emph{lāzan}
`to let'), but sporadic Old Saxon spellings in ⟨e⟩ show that the
fronting lapped unevenly into Old Saxon territory
(\citeproc{ref-RingeTaylor2014}{Ringe and Taylor 2014, 150}). The change
is therefore a North Sea Germanic development rather than an exclusively
Anglo-Frisian one, though only Old English and Old Frisian carry it
through systematically.

The fronting affected stressed, non-nasalized \emph{*ā}; nasalized
\emph{*ą̄} was instead rounded, as treated in the rounding chapter. Ringe
and Taylor further establish a conditioning by a following \emph{*w}.
Before \emph{*w} plus a back or non-high vowel the \emph{*ā} was
retained: the clearest witnesses are the \emph{verba pura} ---
\emph{sāwan} `to sow', \emph{cnāwan} `to know', \emph{blāwan} `to blow',
\emph{māwan} `to mow', \emph{þrāwan} `to turn' --- together with
\emph{clāwu} `claw' (\citeproc{ref-RingeTaylor2014}{Ringe and Taylor
2014, 151}); the simplest hypothesis, which Ringe and Taylor adopt from
Hogg, is that fronting never occurred in that environment
(\citeproc{ref-RingeTaylor2014}{Ringe and Taylor 2014, 151};
\citeproc{ref-HoggGrammar1992}{Hogg 1992a, 81}). Before \emph{*w} plus a
high front vocalic, by contrast, the fronting did apply:
\emph{*lēwijaną} `to betray' (Gothic \emph{lēwjan} `to betray', Old High
German \emph{gilāen} `to betray') yields West Saxon \emph{lǣwan} `to
betray', and the same environment appears in \emph{eltǣwe} `entire' and
\emph{brǣw} `eyelid' (\citeproc{ref-RingeTaylor2014}{Ringe and Taylor
2014, 150}). Both sides of the condition are encoded in the rule below
and witnessed in the corpus. Bennett states the retaining environment
independently from the West Saxon evidence alone, and slightly more
broadly: West Saxon ``shows \emph{ǣ} as a regular isolative development
of IE \emph{ē} but has \emph{ā} before \emph{w} or \emph{g} plus a back
vowel'', with the paradigmatic alternation \emph{wǣg} `wave' beside
plural \emph{wāgas} `waves' (\citeproc{ref-Bennett1950}{Bennett 1950,
235}, n.~6). The rule below encodes the \emph{w} environment, for which
the corpus supplies witnesses on both sides; the parallel \emph{g}
environment is left for separate treatment.

Whether this fronting restored a front vowel that had earlier been
backed, or whether --- as Fulk argues --- the North Sea dialects simply
retained an old front \emph{*ǣ} that was never backed at all, is the
same dispute recorded in the lowering chapter
(\citeproc{ref-Fulk2018}{Fulk 2018, 60--61}, §4.6;
\citeproc{ref-Campbell1959}{Campbell 1959, 50--51}, §§128--129). On the
retention analysis this chapter's change dissolves into the non-event of
staying put; the present model follows Ringe and Taylor's two-step
reconstruction, on the strength of the runic evidence for an early
{[}aː{]} and the place-adverbs \emph{þǣr} `there' and \emph{hwǣr}
`where' (\citeproc{ref-RingeTaylor2014}{Ringe and Taylor 2014, 13--14}).

\subsection{\texorpdfstring{\CAPRRuleHeading{SC101. Fronting of long \emph{ā} outside nasal and \emph{w} environments}{EAFLongAFronting}}{}}\label{rule-EAFLongAFronting}

\noindent\begin{minipage}{\linewidth}
\begin{ReaderFacingFoma}
define EAFLongAFronting [
    {*ā} -> {*ǣ} || _ [EnglishStarConsonant - EnglishStarNasal - {*w}],
    {*ā} -> {*ǣ} || _ {*w} EnglishIUmlautTrigger
];
\end{ReaderFacingFoma}
\end{minipage}

The first clause fronts \emph{*ā} before any oral consonant other than
\emph{*w}; the second admits the fronting before \emph{*w} exactly when
a high front vocalic follows. The two corpus witnesses of the \emph{w}
condition form a minimal contrast. \Recon{sḗaną} `to sow', whose
hiatus-filling \emph{w} is supplied by
\hyperref[rule-EAFHiatusWInsertion]{SC102 EAFHiatusWInsertion}, reaches
this rule as \emph{*sāwaną} and is left unfronted, surfacing as
\emph{sāwan} `to sow'; were the \emph{w}-block removed, the cascade
would deliver \Pred{sǣwan} instead. \Recon{lḗwijaną} `to betray', whose
inherited \emph{w} is followed by \emph{*i}, is fronted by the second
clause to \emph{*lǣwijaną} at this rule's own stage --- before, and
independently of, the much later i-umlaut --- and surfaces as
\emph{lǣwan} `to betray'. The rule consumes the \emph{*ā} created by
\hyperref[rule-PNWGmcLongELowering]{SC024 PNWGmcLongELowering};
displaced before the lowering, it has nothing to front, and
\Recon{skḗpą} `sheep' surfaces as \Pred{sċāp} rather than OE
\emph{sċēap} `sheep', \Recon{jḗrą} `year' as \Pred{ġār} rather than
\emph{ġēar} `year', \Recon{slḗpaną} `to sleep' as \Pred{slāpan} rather
than \emph{slǣpan} `to sleep'.

Two later boundaries carry real historical content. First, the fronting
must precede the completion of
\hyperref[rule-EAFAiMonophthongization]{SC004 EAFAiMonophthongization}:
the \emph{ā} that arose from \emph{*ai} was never fronted ---
\emph{stān} `stone', \emph{hām} `home', \emph{lāþ} `hostile', \emph{rāp}
`rope', \emph{tācn} `token', \emph{gāst} `spirit' all keep the back
vowel. Campbell draws exactly this chronological inference
(\citeproc{ref-Campbell1959}{Campbell 1959, 52--53}, §132), and Ringe
and Taylor endorse it as cogent (\citeproc{ref-RingeTaylor2014}{Ringe
and Taylor 2014, 169--70}); in the present cascade the inference is
enforced by rule order, and displacing the fronting after the
monophthongization wrongly yields \Pred{lǣþ}, \Pred{rǣp}, \Pred{tǣcn},
\Pred{sǣwol}, and \Pred{ġēast}. Historically the two changes may well
have overlapped in time; the discrete ordering is the grammar's way of
stating that inherited \emph{ā} had been fronted before the new \emph{ā}
arose.

Second, the fronting must precede
\hyperref[rule-OEWsPalatalDiphthongization]{SC056
OEWsPalatalDiphthongization}, which operated on the already-fronted
vowel: \emph{ǣ} \textgreater{} \emph{ēa} after the palatals, as in
\emph{sċēap} `sheep' and \emph{ġēar} `year'
(\citeproc{ref-Campbell1959}{Campbell 1959, 69--70}, §185;
\citeproc{ref-RingeTaylor2014}{Ringe and Taylor 2014, 215--16}, §6.5.1).
Displaced after the diphthongization, the cascade yields \Pred{sċǣp} and
\Pred{ġǣr} instead.

\newpage

\chapter{From Anglo-Frisian to Old
English}\label{from-anglo-frisian-to-old-english}

\section{Historical interval}\label{historical-interval-3}

This chapter follows the English daughter from the required
Anglo-Frisian ancestor through prehistoric English to attested Old
English. It includes changes earlier than the manuscript period, not
merely changes within it. West Saxon is the principal target; each
comparison still requires its actual dialect, paradigm cell and
evidential status.

\section{Scope and dialect variation}\label{scope-and-dialect-variation}

Not every change in this chapter has pan-Old-English scope. Some changes
--- most notably ordinary West Saxon palatal diphthongization and parts
of the back-mutation history --- are specifically West Saxon or more
broadly southern Old English phenomena. The diphthong corridor also
represents earlier West Germanic feeders; a West Saxon final outcome
does not date every feeder to West Saxon
(\citeproc{ref-RingeTaylor2014}{Ringe and Taylor 2014, 41--42, 65--66,
171--73, 215--17}). (The West Saxon palatal-glide spellings, SC016,
belong to the written surface of Old English and are treated in Chapter
5.)

The CAPR derivations target West Saxon Old English citation forms as the
default comparator. Changes that belong to other dialects, or that are
absent from West Saxon, may appear in lexical entries as comparanda
rather than as derivational steps.

The distinction between West Saxon and Anglian comparanda requires the
actual text and dialect, not merely a regional label. Hogg discusses the
limits of the traditional divisions and the evidence for the principal
textual varieties (\citeproc{ref-HoggGrammar1992}{Hogg 1992a, 3--8},
§§1.5--1.12).

\section{Chapter structure}\label{chapter-structure}

The changes in this chapter fall into several natural historical
subgroups, though the boundaries between them are not always sharp:

Prehistoric English contraction and fronting: Completed stressed
\emph{*ai > *ā} is now adopted on the English daughter in
\hyperref[rule-EAFAiMonophthongization]{SC004 EAFAiMonophthongization}.
The new long vowel is distinguished from inherited oral long vowels; the
source milestones allow overlap with their earlier restructuring
(\citeproc{ref-Campbell1939}{Campbell 1939, 90--91};
\citeproc{ref-RingeTaylor2014}{Ringe and Taylor 2014, 170--71}). This
placement leaves possible conditioned ancestral onset and the runic
interpretation qualifications explicit
(\citeproc{ref-Versloot2017}{Versloot 2017, 295--97, 318}).

Early Old English changes linked to the Anglo-Frisian inheritance:
Changes that feed directly on, or are closely related to, Anglo-Frisian
brightening (SC043), whose section follows the earlier vowel
developments at its adopted executable position in this chapter. The
\emph{*awj} resolution, \emph{*au} fronting and diphthong operations
currently execute before ordinary-a brightening. They cannot be
described as consuming its output. Their histories combine earlier glide
reanalysis with English diphthong realization
(\citeproc{ref-Campbell1959}{Campbell 1959, 44--47};
\citeproc{ref-RingeTaylor2014}{Ringe and Taylor 2014, 65--66, 171--73}).
The conventional breaking/restoration account instead presupposes
fronted \emph{*æ}. English plain-a/au process identity and inherited
stem identity remain separate questions.

A coordinated formalization of ordinary \emph{*a} fronting and
\emph{*au} fronting with its completion has now been tested over the
selected lexical material. Its twenty au paths wait until the ordinary
fronting corridor and then converge with the current derivations before
breaking. This establishes computational compatibility, not historical
event identity. The production serialization remains unchanged, and the
unstressed and surviving-long-final fronting components are not dated
from this test. The completed ordinary stressed component is now
independently characterized on the English daughter. Its English-episode
interpretation remains explicitly dependent on its contraction and
nucleus premises (\citeproc{ref-Campbell1939}{Campbell 1939, 90--91};
\citeproc{ref-Campbell1959}{1959, 52};
\citeproc{ref-RingeTaylor2014}{Ringe and Taylor 2014, 170--75}).

The earlier unstressed contribution is not the complete unstressed law:
\hyperref[rule-OEUnstressedFrontingEarly]{SC070
OEUnstressedFrontingEarly} separately implements fronting before
heterosyllabic nasals, whereas the coda-nasal history protects other
endings (\citeproc{ref-Campbell1959}{Campbell 1959, 140--41}). The
retained final-vowel helper connects unrounding's carried quantity to
later shortening and merger; it does not establish an independent
historical long-a fronting (\citeproc{ref-RingeTaylor2014}{Ringe and
Taylor 2014, 58--59, 299--300}).

Old English consonantal changes: Velar palatalization (SC052),
palatalization of \texttt{*sk} (SC051), j-cluster coalescence (SC057),
and related changes produce the characteristically Old English consonant
phonemes. Hogg discusses these as OE consonant changes with
class-specific conditions and chronology (\citeproc{ref-Hogg1979}{Hogg
1979, 90--111}).

Old English i-umlaut and its context: The i-umlaut changes vowels under
a following high front vocoid; paradigm alternations are evidence for
that phonological conditioning, not grammatical conditions on the law
(\citeproc{ref-Campbell1959}{Campbell 1959, 69--72}, §§190--197;
\citeproc{ref-RingeTaylor2014}{Ringe and Taylor 2014, 222}). Ordinary
West Saxon palatal diphthongization precedes mutation in the handbook
account, while the later treatment of some mutation products after
\emph{sċ} follows it (\citeproc{ref-Fulk2018}{Fulk 2018, 74};
\citeproc{ref-Campbell1959}{Campbell 1959, 68--69};
\citeproc{ref-RingeTaylor2014}{Ringe and Taylor 2014, 215--17, 235}).
The adopted cascade now places the ordinary component before mutation
and retains the late approximation afterward. Gift enters with the
earlier-raised vowel, following Ringe's explicitly PGmc reconstruction
(\citeproc{ref-Ringe2017}{Ringe 2017, 135}, pp.~151--153). Its lexical
section discusses the published e-alternatives and the confidence in
this choice.

Syncope, apocope and weak-vowel quality: Mutation precedes the ordinary
high-vowel losses: a subsequently lost \emph{*i} has already fronted the
root, as in \emph{fȳlþe} `filth'. The former conventional working
sequence distinguished ordinary nonhigh syncope (SC114), high-vowel
apocope (SC063) and heavy/open high-vowel syncope (SC065). The adopted
quantitative account instead gives SC063 one final/medial operation,
preserving heavy--light--light \emph{hēafudu} `heads' and distinguishing
the following-heavy \emph{nētna} `of livestock'. Early medial shortening
belongs to SC072 through its separate component SC117; SC065 remains
only a trace anchor (\citeproc{ref-Fulk2010}{Fulk 2010, 129--34,
140--42}; \citeproc{ref-Goering2016}{Goering 2016, 186--88};
\citeproc{ref-Goering2023}{2023, 67--68}). These are prehistoric Old
English changes, not simply late literary reductions toward Middle
English. The former dental and heavy pre-\emph{*l} fragments implement
the same high-vowel event and are no longer separate active deletions.
Later light-syllable pre-\emph{*l} loss (SC066) retains its own
consonant and openness conditions (\citeproc{ref-HoggGrammar2011}{Hogg
2011, 220--27}, §§6.14--6.25; \citeproc{ref-SieversBrunner1965}{Brunner
1965, 132--33}, §159; \citeproc{ref-Luick1914}{Luick 1914-\/-1940,
310--11}, §336).

The trisyllabic analysis remains disputed: Fulk's conservative Mercian
evidence, Ringe and Taylor's syncope-before-apocope account, and Hogg's
tentative phonological alternative are discussed with SC063. Early weak
front-vowel unrounding and later lowering are distinct quality processes
(SC074), linked to transmitted mutation under SC055, not additional
syncopes. The early unrounding support component is printed at SC115;
its separately positioned stress decoder is technical. Ordinary nonhigh
loss now also removes compound \emph{*a}; the former early proxy (SC037)
is an identity trace anchor. Positional realization of inherited
\emph{*b} is computed afterward (SC049) without postnasal hardening of
existing fricatives or a newly dated lenition law
(\citeproc{ref-HoggGrammar2011}{Hogg 2011, 29--30}, §§2.53--2.54,
p.~219, §6.10; \citeproc{ref-Campbell1959}{Campbell 1959, 179--81},
§§444--448).

Across this chapter the densest independent parallel is Barrack's West
Saxon cascade, whose focus --- West Saxon to the beginning of the
eleventh century (\citeproc{ref-Barrack1975}{Barrack 1975, 19}) ---
coincides with the present target. His statements most relevant here
include velar palatalization as rule 170a
(\citeproc{ref-Barrack1975}{Barrack 1975, 57}); back umlaut as rules 210
and 220 (\citeproc{ref-Barrack1975}{Barrack 1975, 63--66}); the
\emph{weor-} \textgreater{} \emph{wur-} development as rule 320, his
showcase for progressively widening environments
(\citeproc{ref-Barrack1975}{Barrack 1975, 75--76, 23--25}); high-vowel
apocope and syncope of originally long vowels as rules 345 and 358
(\citeproc{ref-Barrack1975}{Barrack 1975, 85--91}); and, among the late
consonant adjustments, \emph{sk}-metathesis in
\emph{fiscas}/\emph{fixas} `fishes' as rule 440b
(\citeproc{ref-Barrack1975}{Barrack 1975, 113}), the \emph{bringþ}-type
dissimilation as 477a (\citeproc{ref-Barrack1975}{Barrack 1975, 123}),
stop fortition of the \emph{fint}-type as 481.1
(\citeproc{ref-Barrack1975}{Barrack 1975, 124--25}), and the nasal
assimilation of \emph{efn} \textgreater{} \emph{emn} `even' as 484d
(\citeproc{ref-Barrack1975}{Barrack 1975, 127}). On r-metathesis (rule
459.1) he records that the relative order of breaking and metathesis
varies by dialect in the \emph{beornan}/\emph{birnan} `to burn' family
(\citeproc{ref-Barrack1975}{Barrack 1975, 115--17}), a
conditioning-relevant observation held for separate adjudication in the
Barrack dossier. His orderings are formal, not chronological
(\citeproc{ref-Barrack1975}{Barrack 1975, 34}), and are cited here as
corroborating statements of conditioning, never as dating evidence.

\section{The English partial chains}\label{the-english-partial-chains}

The working historical scaffold has several connected chains, not one
uniquely demonstrated total order. Their evidential strength differs. An
explicit handbook derivation, a contrast between original and secondary
vowels, and the survival of a modeled output are not interchangeable
arguments (\citeproc{ref-RingeTaylor2014}{Ringe and Taylor 2014,
170--73, 215--37}; \citeproc{ref-Hogg1979}{Hogg 1979, 100--110}).

{\def\LTcaptype{none} % do not increment counter
\begin{longtable}[]{@{}
  >{\raggedright\arraybackslash}p{(\linewidth - 4\tabcolsep) * \real{0.3333}}
  >{\raggedright\arraybackslash}p{(\linewidth - 4\tabcolsep) * \real{0.3333}}
  >{\raggedright\arraybackslash}p{(\linewidth - 4\tabcolsep) * \real{0.3333}}@{}}
\toprule\noalign{}
\begin{minipage}[b]{\linewidth}\raggedright
Historical relation
\end{minipage} & \begin{minipage}[b]{\linewidth}\raggedright
Evidence and interpretation
\end{minipage} & \begin{minipage}[b]{\linewidth}\raggedright
Limitation
\end{minipage} \\
\midrule\noalign{}
\endhead
\bottomrule\noalign{}
\endlastfoot
Coronal assimilation before inherited ww reanalysis & Four and the
second-person pronoun have assimilated inputs to vocalization & Does not
date every j-created geminate \\
Inherited-long restructuring before completed ai contraction & New
ai-derived long a escapes inherited-long fronting & Temporal overlap is
possible \\
Completed ai contraction before ordinary short fronting & Conventional
Campbell chain & Depends on the diphthong-nucleus premise \\
Ordinary fronting before breaking/restoration & Conventional English
reconstruction & Never-fronted and restored outputs can coincide \\
Breaking before ordinary palatal diphthongization before mutation &
Broken inputs and worked handbook vowel histories & Does not date every
consonantal palatal layer \\
Mutation before the later sc treatment & Ai-derived mutation products
such as sheath & Exact later conditioner remains unresolved \\
Initial productive palatal cutoff before mutation-created fronts &
Unrounded mutation vowels are the stronger controls & Medial fricative
merger is a separate problem \\
\end{longtable}
}

The first three relations are supported and qualified by the earlier
vowel accounts (\citeproc{ref-RingeTaylor2014}{Ringe and Taylor 2014,
41--42, 65--66, 170--73}; \citeproc{ref-Campbell1939}{Campbell 1939,
90--91}). The palatal and mutation relations require their own evidence
(\citeproc{ref-Luick1914}{Luick 1914-\/-1940, 162--63};
\citeproc{ref-Fulk2018}{Fulk 2018, 74};
\citeproc{ref-RingeTaylor2014}{Ringe and Taylor 2014, 215--17, 235};
\citeproc{ref-Hogg1979}{Hogg 1979, 100--110}). The table states the
defended historical targets; it does not claim that every present
executable proxy already realizes them faithfully.

\section{Early glide history and the English
outcome}\label{early-glide-history-and-the-english-outcome}

The diphthong corridor includes inherited ww, assimilated ww and
geminates created later before j. These inputs must be distinguished
before assigning one date to their long-vowel outputs. Ringe and Taylor
describe inherited glide reanalysis separately from the English
realization of its products (\citeproc{ref-RingeTaylor2014}{Ringe and
Taylor 2014, 41--42, 65--66, 171--73}). The model now separates the
earlier reanalysis from English realization.

{\def\LTcaptype{none} % do not increment counter
\begin{longtable}[]{@{}
  >{\raggedright\arraybackslash}p{(\linewidth - 6\tabcolsep) * \real{0.2500}}
  >{\raggedright\arraybackslash}p{(\linewidth - 6\tabcolsep) * \real{0.2500}}
  >{\raggedright\arraybackslash}p{(\linewidth - 6\tabcolsep) * \real{0.2500}}
  >{\raggedright\arraybackslash}p{(\linewidth - 6\tabcolsep) * \real{0.2500}}@{}}
\toprule\noalign{}
\begin{minipage}[b]{\linewidth}\raggedright
Witness
\end{minipage} & \begin{minipage}[b]{\linewidth}\raggedright
Adopted earlier checkpoint
\end{minipage} & \begin{minipage}[b]{\linewidth}\raggedright
Later English outcome
\end{minipage} & \begin{minipage}[b]{\linewidth}\raggedright
What it tests
\end{minipage} \\
\midrule\noalign{}
\endhead
\bottomrule\noalign{}
\endlastfoot
Four & Assimilated ww is reanalyzed as eu plus w & fēower `four' &
Assimilation feeder and retained glide \\
Hew & Inherited a plus ww gives au plus w & hēawan `hew' & Reanalysis
distinct from English realization \\
You & Vocalized iu plus w with a surviving final i before apocope & ēow
`you' & Prosody/apocope compatibility, not merely quantity \\
Hue & Later j-gemination creates a distinct glide input & hīew `hue' &
Negative control for the earlier inherited-ww subset \\
Hay & Later awj history remains separate & hīeġ `hay' & Gemination and
secondary glide resolution \\
\end{longtable}
}

These are existing lexical witnesses, not new corpus admissions. The
current spelling differences also require their actual dialectal
interpretation (\citeproc{ref-Campbell1959}{Campbell 1959, 44--47};
\citeproc{ref-RingeTaylor2014}{Ringe and Taylor 2014, 41--42, 57--58,
171--73}). The reconstructed West Saxon strew target is a computational
control, not an additional attested form.

The completed decomposition replaces the old consonantal-ww apocope
proxy with the historical condition: a short final high vowel after a
heavy syllable in a sentence-unstressed, phonologically final word. You
selects that weak-final context independently of its unchanged PGmc
reconstruction. Thus \emph{*iuwi} loses the final vowel to give
\emph{*iuw} before mutation. A proclitic context retains the vowel, as
in \emph{ymbe} `around'; weak-final \emph{and} `and' and stressed
\emph{ġiest} `guest' / \emph{fȳr} `fire' check the positive and negative
conditions (\citeproc{ref-RingeTaylor2014}{Ringe and Taylor 2014,
41--42, 55, 57--58}; \citeproc{ref-Campbell1959}{Campbell 1959, 283}).

The computational retained-i counterfactual for you yields a predicted
\Pred{īei}, not a source-backed strong Old English spelling: the
existing w-loss before i intervenes before mutation. It tests the
represented alternative context, not an additional attestation. Unmarked
evaluation explicitly selects strong-final citation context; absence of
an acute is never used to infer sentence stress. All selected final
outputs remain unchanged.

The later realization paths are now explicit. Chew, dew, four, hew and
you complete their earlier products through
\hyperref[rule-OEDiphthongLeveling]{SC032 OEDiphthongLeveling}.
\hyperref[rule-OEEwLongDiphthong]{SC033 OEEwLongDiphthong} supplies the
separate j-created offglide reanalysis for \emph{hīew} `hue'. The
selected input is inherited PGmc \emph{*xíwją}; its path is
\emph{*xíwwją} to \emph{*xíuwją}, then long \emph{*xīowją} through
\hyperref[rule-OEJGlideIO]{SC110 OEJGlideIO}, before mutation gives the
West Saxon long \emph{īe} reflex. Ringe and Taylor explicitly give the
inherited i and long-io chain while leaving its phonemic analysis
unclear (\citeproc{ref-RingeTaylor2014}{Ringe and Taylor 2014, 53,
250}). Orel's and Kroonen's e reconstructions are real alternatives, and
Fulk disputes traditional geminate dismantling; these qualifications are
discussed in the lexical account (\citeproc{ref-Orel2003}{Orel 2003,
171--72}; \citeproc{ref-Kroonen2013}{Kroonen 2013, 224};
\citeproc{ref-Fulk2018}{Fulk 2018, 59, 71--72}). Knee instead selects
the regular dative \emph{cneowe} `knee (dat.sg.)': its short pre-ending
diphthong comes from \hyperref[rule-OEBreaking]{SC044 OEBreaking}, not
singleton lengthening. The long endingless \emph{cnēo} `knee' and
analogically restored final \emph{w} in \emph{cnēow} `knee' remain a
separate paradigm comparison; later long obliques must not be confused
with the selected short cell (\citeproc{ref-Campbell1959}{Campbell 1959,
232--33}; \citeproc{ref-HoggGrammar2011}{Hogg 2011, 21--22, 86};
\citeproc{ref-RingeTaylor2014}{Ringe and Taylor 2014, 187--88, 387}).
\hyperref[rule-PWGmcEndinglessAw]{SC111 PWGmcEndinglessAw} now supplies
straw's regular endingless \emph{*strau}, feeding ordinary English
realization to reconstructed \emph{*strēa}. The comparison remains
attested \emph{strēaw} `straw': its restored final \emph{w} belongs to
paradigm levelling, so this regular derivation explicitly mismatches the
target. The sources mark the free regular simplex and short obliques as
unattested; the compound \emph{strēaberige} `strawberry' is not
substituted for the selected noun (\citeproc{ref-Campbell1959}{Campbell
1959, 232--33}; \citeproc{ref-SieversBrunner1965}{Brunner 1965, 49,
110}; \citeproc{ref-RingeTaylor2014}{Ringe and Taylor 2014, 172--73,
387}). \hyperref[rule-OEAwLongDiphthong]{SC034 OEAwLongDiphthong}
retains only the three show cells as a disclosed proxy. Dictionary
singleton stems and the inherited-geminate handbook analysis remain
distinct, unresolved premises (\citeproc{ref-Orel2003}{Orel 2003, 337};
\citeproc{ref-Kroonen2013}{Kroonen 2013, 442};
\citeproc{ref-Campbell1959}{Campbell 1959, 45};
\citeproc{ref-SieversBrunner1965}{Brunner 1965, 49}). The unchanged
\hyperref[rule-OEJWWSimplification]{SC106 OEJWWSimplification} remains a
technical compatibility operation, but hue now reaches it with singleton
w and no longer feeds it. The before-j placement of early reanalysis is
a representative serialization, not a new strict historical chronology
claim (\citeproc{ref-Campbell1959}{Campbell 1959, 45--47};
\citeproc{ref-RingeTaylor2014}{Ringe and Taylor 2014, 65--66, 171--75}).

\section{Ordinary fronting, breaking and
restoration}\label{ordinary-fronting-breaking-and-restoration}

The preferred conventional account identifies ordinary non-nasal short-a
fronting and fronting of au's first element as one English episode. The
later offglide development is distinct. Early spellings record the
fronted element, but do not measure the simultaneity of two laws
(\citeproc{ref-Campbell1939}{Campbell 1939, 91};
\citeproc{ref-Campbell1959}{1959, 52--53};
\citeproc{ref-Luick1914}{Luick 1914-\/-1940, 130--31}). English episode
identity is not inherited Anglo-Frisian event identity. Nor does it
identify the stressed-short process with every unstressed or long-final
clause in the present model.

Breaking then has vowel, quantity, consonantal and dialectal conditions.
Restoration in a following back-vowel environment can remove a
previously fronted vowel. A final restored vowel alone cannot
distinguish that history from original blocking of fronting
(\citeproc{ref-Hogg1979}{Hogg 1979, 90--97};
\citeproc{ref-Versloot2025}{Versloot 2026, 104, 123}). The conventional
account is retained as the working baseline because of the combined
argument, not because the more elaborate intermediate path is inherently
preferable.

Versloot's late Anglian alternatives must be compared in their exact
domains. His mutation-created e argument and the unbroken neighboring
environments are relevant to a restricted e-breaking component, not to a
wholesale postmutation relocation of every West Saxon breaking law
(\citeproc{ref-Versloot2025}{Versloot 2026, 126--28, 131--35}). A
different dialectal target cannot be silently substituted for a West
Saxon lexical target. Similarly, Frisian closed-syllable breaking is not
English rC breaking merely because both are called breaking
(\citeproc{ref-Bremmer2009}{Bremmer 2009, 33--35, 37}).

Fulk likewise distinguishes the later Frisian process from English
breaking. His Northumbrian restoration-before-breaking account differs
from the usual English ordering illustrated by \emph{slēan} `slay'; it
cannot replace the West Saxon sequence merely because the final vowels
coincide (\citeproc{ref-Fulk1998}{Fulk 1998, 150--53}). Chapter 3
compares his broader reconstruction and its explicit limits on shared
inheritance.

\section{Palatal consonants and the key--day
problem}\label{palatal-consonants-and-the-keyday-problem}

The useful chronological contrast is between original conditioning
fronts and front vowels created by mutation. Rounded mutation products
alone may support a weaker conclusion, a cutoff before later unrounding.
Unrounded mutation-created fronts more directly test whether initial
velar palatalization was still productive (\citeproc{ref-Hogg1979}{Hogg
1979, 100--103}; \citeproc{ref-Laker2007}{Laker 2007, 167--68}). This
cutoff does not by itself date assibilation or every earlier
articulatory tendency.

Hogg's key and day paradigms identify the disputed merger premise. The
initial key consonant survives a mutation-created unrounded front vowel;
the day paradigm does not show the mutation that would follow if its new
medial palatal element were already equivalent to inherited j. These
statements concern particular historical cells, not arbitrary exceptions
to a sound law (\citeproc{ref-Hogg1979}{Hogg 1979, 102--10}).

{\def\LTcaptype{none} % do not increment counter
\begin{longtable}[]{@{}
  >{\raggedright\arraybackslash}p{(\linewidth - 4\tabcolsep) * \real{0.3333}}
  >{\raggedright\arraybackslash}p{(\linewidth - 4\tabcolsep) * \real{0.3333}}
  >{\raggedright\arraybackslash}p{(\linewidth - 4\tabcolsep) * \real{0.3333}}@{}}
\toprule\noalign{}
\begin{minipage}[b]{\linewidth}\raggedright
Diagnostic
\end{minipage} & \begin{minipage}[b]{\linewidth}\raggedright
Required distinction
\end{minipage} & \begin{minipage}[b]{\linewidth}\raggedright
What is not established
\end{minipage} \\
\midrule\noalign{}
\endhead
\bottomrule\noalign{}
\endlastfoot
Key oblique vowel & Original fronts versus mutation-created fronts &
Date of all palatal articulation \\
Day oblique vowel & Inherited j versus a palatal fricative at mutation &
Unanimous agreement on merger timing \\
Geminate/postnasal g & Stop-class history versus singleton fricative &
One universal g-to-j change \\
h & Breaking trigger versus subsequent palatal/weakening history &
Identity with voiced-fricative chronology \\
sk & Own consonantal law versus later sc vowel treatment & A single
combined palatalization event \\
\end{longtable}
}

Hogg considers several resolutions without endorsing one as demonstrated
(\citeproc{ref-Hogg1979}{Hogg 1979, 103--11}). Ringe and Taylor,
however, explicitly adopt merger after mutation
(\citeproc{ref-RingeTaylor2014}{Ringe and Taylor 2014, 204}). CAPR now
implements that regular working account, preserving Hogg's phonetic
objection rather than claiming consensus. Singleton g is a fricative,
including initially; gg/ng are stops (\citeproc{ref-Fulk2018}{Fulk 2018,
130--32}).

At Hogg's pre-palatal, pre-OE checkpoint, key has \emph{*kājæ} and day
\emph{*dæɣæ}. The inherited j in the former causes mutation, producing
\emph{*kǣjæ} without reactivating initial-k palatalization. In the
latter, the newly palatal fricative remains ʝ and the vowel stays æ:
\emph{*dæʝæ}. Native realization then gives \emph{cǣġe} `key' and
\emph{dæġe} `day' (\citeproc{ref-Hogg1979}{Hogg 1979, 105};
\citeproc{ref-RingeTaylor2014}{Ringe and Taylor 2014, 204}). Premature
merger instead supplies a j trigger and wrongly raises the day vowel.
These oblique cells are non-corpus diagnostics, not new selected PGmc
reconstructions; the existing day nominative cannot substitute for them.

The complete key suffix exposed a second defect: unrestricted
inherited-j normalization would erase its retained glide. The repaired
\hyperref[rule-OEIntervocalicJVocalization]{SC082
OEIntervocalicJVocalization} excludes non-high long front monophthongs.
Its residual weak-suffix normalization is a telescoped representation,
not a universal historical VjV vocalization law
(\citeproc{ref-HoggGrammar2011}{Hogg 2011, 283--86};
\citeproc{ref-RingeTaylor2014}{Ringe and Taylor 2014, 228}).

The fricative merger has its own visible
\hyperref[rule-OEPalatalFricativeMerger]{SC109
OEPalatalFricativeMerger}. Its late execution after suffix raising is a
tested computational holding zone, not a precisely proved historical
date. Likewise stop ʧ/ʤ are eventual-reflex proxies, not early affricate
claims. Cluster coalescence and a dotted written reflex cannot identify
all these classes as one historical sound
(\citeproc{ref-RingeTaylor2014}{Ringe and Taylor 2014, 203--4};
\citeproc{ref-Laker2007}{Laker 2007, 167--68};
\citeproc{ref-Fulk2018}{Fulk 2018, 131--32}).

\section{Three different mutation and diphthongization
arguments}\label{three-different-mutation-and-diphthongization-arguments}

Guest, sheath and gift do not establish the same chronological relation.
The guest derivation passes through ordinary palatal diphthongization
before mutation in both the handbook account and the adopted model. The
former model raised a simple vowel before diphthongizing it. Both paths
reach \emph{ġiest} `guest', so their final agreement does not select the
earlier path (\citeproc{ref-RingeTaylor2014}{Ringe and Taylor 2014,
216}).

The following worked comparison abstracts from consonant notation and
retains the vocalic trigger where it matters. For guest, the source
starts with PGmc \emph{*gastiz} `guest', continued as PWGmc
\emph{*gasti} `guest'. The vowel history is \emph{*a > *æ > *ea > *ie};
final loss of the trigger follows the mutation it caused
(\citeproc{ref-RingeTaylor2014}{Ringe and Taylor 2014, 216, 287}).

{\def\LTcaptype{none} % do not increment counter
\begin{longtable}[]{@{}
  >{\raggedright\arraybackslash}p{(\linewidth - 4\tabcolsep) * \real{0.3333}}
  >{\raggedright\arraybackslash}p{(\linewidth - 4\tabcolsep) * \real{0.3333}}
  >{\raggedright\arraybackslash}p{(\linewidth - 4\tabcolsep) * \real{0.3333}}@{}}
\toprule\noalign{}
\begin{minipage}[b]{\linewidth}\raggedright
Checkpoint
\end{minipage} & \begin{minipage}[b]{\linewidth}\raggedright
Handbook and adopted guest path
\end{minipage} & \begin{minipage}[b]{\linewidth}\raggedright
Superseded modeled guest path
\end{minipage} \\
\midrule\noalign{}
\endhead
\bottomrule\noalign{}
\endlastfoot
Ordinary fronting & Short a becomes æ & Short a becomes æ \\
After palatal conditioning is established & Original front vowel is
available to ordinary PD & Original front vowel is available, but PD
remains later \\
First disputed vowel operation & æ becomes ea by ordinary PD & æ becomes
e at mutation \\
Second disputed vowel operation & ea becomes ie at mutation & e becomes
ie in the bundled PD proxy \\
Written result & ġiest `guest' & ġiest `guest' \\
\end{longtable}
}

The agreement in the last row is precisely why the intermediate rows
must be tested. This table does not treat the implementation's
palatal-consonant symbol as proof of a uniquely dated phonetic merger
(\citeproc{ref-Hogg1979}{Hogg 1979, 103--10}).

{\def\LTcaptype{none} % do not increment counter
\begin{longtable}[]{@{}
  >{\raggedright\arraybackslash}p{(\linewidth - 4\tabcolsep) * \real{0.3333}}
  >{\raggedright\arraybackslash}p{(\linewidth - 4\tabcolsep) * \real{0.3333}}
  >{\raggedright\arraybackslash}p{(\linewidth - 4\tabcolsep) * \real{0.3333}}@{}}
\toprule\noalign{}
\begin{minipage}[b]{\linewidth}\raggedright
Witness
\end{minipage} & \begin{minipage}[b]{\linewidth}\raggedright
Source-supported distinction
\end{minipage} & \begin{minipage}[b]{\linewidth}\raggedright
Present model consequence
\end{minipage} \\
\midrule\noalign{}
\endhead
\bottomrule\noalign{}
\endlastfoot
Guest & Ordinary palatal diphthongization feeds mutation & Adopted æ
\textgreater{} ea \textgreater{} ie replaces the output-equivalent old
path \\
Sheath & Ai-derived ā mutates; a later sc layer can affect its product &
Its unchanged late clause is separately retained as an approximation \\
Gift & Earlier inherited e-to-i raising is distinct from English
mutation & Selected PGmc i is outside ordinary diphthongization \\
Sheep/year & Ordinary palatal diphthongization without a mutation
trigger & Controls for the ordinary component \\
Cow/lung & Mutation changes the productive consonantal environment &
Controls for cutoff, not diphthongal mutation \\
\end{longtable}
}

For sheath, \emph{sċēaþ} `sheath' beside \emph{sċǣþ} `sheath' belongs to
the separately described later group. Reporting both outcomes does not
authorize lexical optionality in the law; the exact phonological layer
and conditioner still require a defended specification
(\citeproc{ref-Campbell1959}{Campbell 1959, 68--69};
\citeproc{ref-RingeTaylor2014}{Ringe and Taylor 2014, 235}). The
incremental repair preserves the literal late clause as a visible
approximation, not a lexical exception or a newly established
conditioner.

The sheath comparison instead follows the ai-derived long vowel:
\emph{*ai > *ā > *ǣ}, with the final step supplied by mutation. The
separately discussed later sc treatment can then give \emph{*ēa}. The
vowel available to the earlier ordinary process and the vowel created
later by mutation are therefore different historical inputs, which the
adopted model now separates (\citeproc{ref-Campbell1959}{Campbell 1959,
68--69}; \citeproc{ref-RingeTaylor2014}{Ringe and Taylor 2014, 235}).

For gift, Ringe reconstructs PGmc \emph{*giftiz} `gift', with the Old
English plural \emph{ġifta} `wedding', and describes the earlier raising
separately (\citeproc{ref-Ringe2017}{Ringe 2017, 135}, pp.~151--153).
Ringe and Taylor explicitly distinguish that early change from the much
later English vowel history (\citeproc{ref-RingeTaylor2014}{Ringe and
Taylor 2014, 220}). The selected input now follows that earlier-raised
reconstruction. The choice is lexical and source-led, not a new raising
rule justified by output fit. Orel and Kluge--Seebold print e-vowel
alternatives (\citeproc{ref-Orel2003}{Orel 2003, 130};
\citeproc{ref-KlugeSeebold2011}{Kluge 2011, 359}); the lexical section
compares their evidence with the e-grade presentations and Ringe's
qualified PGmc dating rather than claiming unanimous agreement.

\section{Quantity, reduction and the written
surface}\label{quantity-reduction-and-the-written-surface}

Later apocope and syncope depend on prosody and syllable structure;
earlier input conventions and surviving triggers must be inspected
before their effects are used as dates for other changes. Morphological
cells provide evidence for such environments, not grammatical conditions
on phonological rules (\citeproc{ref-RingeTaylor2014}{Ringe and Taylor
2014, 57--59, 299--300}). Inherited quantity, compensatory lengthening,
diphthong quantity and later shortening must likewise remain separate
dimensions.

The transition to Chapter 5 is a transition to the written surface, not
proof that an orthographic change is a newly dated sound change. A
single spelling can represent different historical consonant inputs, and
a changed spelling need not change the phonetic derivation. The source
discussions and worked lexical entries distinguish implemented laws from
working reconstructions and localized unresolved components.

\section{Executable serialization and unresolved
interfaces}\label{executable-serialization-and-unresolved-interfaces}

Sections follow executable order to make derivations inspectable, not
because that serialization proves a unique total history. Final
agreement can conceal different intermediate paths. Some proxies execute
here with earlier historical classifications; renaming identifiers would
not resolve that substantive distinction:

\begin{itemize}
\tightlist
\item
  SC041 (PWGmc Final Bare-\texttt{*a} Loss) and SC042 (Surviving Bimoric
  \texttt{*ō} Unrounding) carry Proto-West Germanic labels but execute
  in this stretch of the cascade.
\item
  SC064 (NWGmc \texttt{*-n} Stem \texttt{*n} Loss) carries a Northwest
  Germanic label but executes after OE High-Vowel Apocope (SC063).
\item
  SC049 (PGmc B Allophony) carries a Proto-Germanic label but executes
  here.
\end{itemize}

\section{Sources}\label{sources}

Campbell's \emph{Old English Grammar} is the primary source for the
dating and scope of individual changes in this chapter
(\citeproc{ref-Campbell1959}{Campbell 1959, 52--72}). Hogg provides
critical reassessments of palatal conditioning and its relative
chronology (\citeproc{ref-Hogg1979}{Hogg 1979, 90--112}). Ringe and
Taylor supply the most detailed relative-chronology analysis for the
earlier portion of the chapter, through back-mutation
(\citeproc{ref-RingeTaylor2014}{Ringe and Taylor 2014, 169--73,
215--37}). Fulk's \emph{Comparative Grammar} provides additional
coverage for phonological conditioning (\citeproc{ref-Fulk2018}{Fulk
2018, 58--61, 72--74}). For individual changes, source-specific
citations appear in the relevant sound-change sections.

\section{Awj resolution and the English brightening of
au}\label{awj-resolution-and-the-english-brightening-of-au}

Two changes stand between the Proto-Germanic diphthong \emph{*au} and
its Old English reflex \emph{ēa}. The earlier one repairs a West
Germanic gemination and restores a diphthong that the gemination had
obscured. The later one fronts the first element of every \emph{*au},
whether inherited or newly created, and belongs with Anglo-Frisian
brightening.

The two are separate developments with separate domains. The first
concerns a handful of words in which \emph{*w} stood before \emph{*j}.
The second concerns the whole Old English history of \emph{*au}, for
which \emph{lēaf} `leaf', \emph{strēam} `stream' and \emph{brēad}
`bread' are ordinary witnesses, and to which the first change merely
adds two more inputs.

\section{Post-velar w-loss and loss of w before final
i}\label{post-velar-w-loss-and-loss-of-w-before-final-i}

The first rule is a narrow loss of \emph{*w} after velars in the
\emph{*ngw} sequence. Ringe and Taylor derive PGmc \Recon{singwan}
`sing' to Old English \emph{singan} `sing'
(\citeproc{ref-RingeTaylor2014}{Ringe and Taylor 2014, 214}, §6.4.2).
This comparative evidence establishes the change, although no lexical
evidence fixes its order relative to a neighboring rule.

The second rule is historically more legible. Campbell notes the
recurring loss of \emph{*w} before \emph{*i} in unstressed position
(\citeproc{ref-Campbell1959}{Campbell 1959, 167}, §406). Ringe and
Taylor trace the development of \emph{sǣ} `sea' from earlier
\emph{*saiwi-} / \emph{*sawi-} (\citeproc{ref-RingeTaylor2014}{Ringe and
Taylor 2014, 257}, §6.7.1), and Luick gives the same trajectory in his
own historical grammar (\citeproc{ref-Luick1914}{Luick 1914-\/-1940,
173}, §187). The first rule is restricted to the \emph{*ngw} sequence;
the second has a specific lexical witness and defined earlier and later
limits.

\section{I-umlaut and West Saxon palatal
diphthongization}\label{i-umlaut-and-west-saxon-palatal-diphthongization}

Luick gives the change its traditional scale:

\begin{quote}
Der wichtigste Fall von palataler Beeinflussung \ldots{} war die
Veränderung der urenglischen Vokale durch i oder j der Folgesilbe.

(\citeproc{ref-Luick1914}{Luick 1914-\/-1940, 166--67}, §182)
\end{quote}

Campbell gives the most compact classical formulation in English when he
writes that ``the process known as i-umlaut or i-mutation operates on
practically all the sounds which it could theoretically affect in OE''
(\citeproc{ref-Campbell1959}{Campbell 1959, 69}, §190). He immediately
defines the core conditioning environment as a following \texttt{i} or
\texttt{j}, and he goes on to trace the consequences across much of the
vowel system, including forms such as \emph{giest} `guest',
\emph{giefan} `give', \emph{hierde} `shepherd', and \emph{ieldra}
`older' (\citeproc{ref-Campbell1959}{Campbell 1959, 69--72},
§§190--197).

Hogg continues in the same vein: ``we come now to a change which is
almost as uncontroversial as it is important''
(\citeproc{ref-HoggPhonology1992}{Hogg 1992b, 113}). His examples, such
as \emph{bryd} `bride', \emph{trymman} `strengthen', \emph{bedd} `bed',
\emph{ciest} `chest', and \emph{wiersa} `worse', likewise emphasize that
the change is a broad redistribution of vowel quality across the Old
English vowel system (\citeproc{ref-HoggPhonology1992}{Hogg 1992b,
113--14}).

The narrower palatal-diphthongal material is described differently.
Ringe and Taylor treat West-Saxon diphthongization after initial
palatals as a distinct process (\citeproc{ref-RingeTaylor2014}{Ringe and
Taylor 2014, 215}, §6.5.1), and Fulk is even more explicit about its
chronological delicacy when he calls it ``diphthongization by initial
palatal consonants (which precedes front umlaut but not breaking)''
(\citeproc{ref-Fulk2018}{Fulk 2018, 74}, §4.13). Ringe and Taylor's
examples such as \emph{gieldan} `pay', \emph{scield} `shield', and
\emph{scieppan} `create' show that this narrower process is triggered by
already palatal consonants and leads to specifically West-Saxon
diphthongal outputs (\citeproc{ref-RingeTaylor2014}{Ringe and Taylor
2014, 215--16}, §6.5.1).

Luick, Campbell, and Hogg treat i-umlaut as a system-wide change. Ringe
and Taylor and Fulk distinguish from it a narrower West-Saxon process
affecting words after initial palatals. The two changes act in different
environments and produce different lexical consequences.

Two further distinctions are essential to their relative chronology.
First, the raising of inherited Germanic \emph{*e} to \emph{*i} before a
following high front vocoid is much earlier than Old English i-umlaut.
Ringe reconstructs PGmc
\Recon{giftiz}\index[iv]{03pgmc@\ivlangheader{Proto-Germanic}{}!giftiz@\iventry{*giftiz}{}}
`gift', with OE plural \emph{ġifta} `wedding', and discusses that
earlier raising separately (\citeproc{ref-Ringe2017}{Ringe 2017, 135},
pp.~151--153). Ringe and Taylor likewise caution that raising of
inherited \emph{*e} occurred hundreds of years before the Old English
changes; a later repetition is rare and doubtful
(\citeproc{ref-RingeTaylor2014}{Ringe and Taylor 2014, 220}). The
selected gift input now has the earlier-raised vowel. Its lack of
ordinary diphthongization therefore does not date the general West Saxon
process. Published e-reconstructions remain genuine alternatives
(\citeproc{ref-Orel2003}{Orel 2003, 130};
\citeproc{ref-KlugeSeebold2011}{Kluge 2011, 359}); the lexical
discussion compares their evidence and explains the attributed working
choice.

Second, the ordinary diphthongization of inherited/fronted non-high
vowels must be distinguished from the later treatment of some mutation
products after \emph{sċ}. Campbell separates that small group from his
general front-vowel diphthongization: the group includes \emph{sċēaþ}
`sheath' beside \emph{sċǣþ} `sheath'
(\citeproc{ref-Campbell1959}{Campbell 1959, 68--69}, §§184--185). Ringe
and Taylor explicitly place the treatment of the ai-derived mutation
product after mutation, giving the same two sheath outcomes
(\citeproc{ref-RingeTaylor2014}{Ringe and Taylor 2014, 235}). Their
ordinary West Saxon derivations, including \emph{ġiest} `guest' and
\emph{ċytel} `kettle', instead pass through palatal diphthongization
before mutation (\citeproc{ref-RingeTaylor2014}{Ringe and Taylor 2014,
215--17, 222}).

The historically preferred working account therefore has ordinary
palatal diphthongization before Old English mutation, with the
separately evidenced later \emph{sċ} treatment after it. The adopted
cascade separates these operations: ordinary diphthongization precedes
mutation, and the unchanged late approximation follows it. The sources'
alternate spellings do not by themselves supply an exceptionless
conditioning law for the later group; its full conditioner remains
unresolved.

\section{English stressed ai
contraction}\label{english-stressed-ai-contraction}

\subsection{Historical discussion}\label{historical-discussion-29}

Inherited stressed \emph{*ái} yields \emph{*ā} in English. Ringe and
Taylor discuss this development separately from the earlier unstressed
contraction. The distinction between inherited long vowels and new
\emph{*ā} requires inherited-long fronting to have been well under way
before contraction completed; it does not prove that the two
developments could not overlap (\citeproc{ref-RingeTaylor2014}{Ringe and
Taylor 2014, 170--71}). Campbell gives the stricter conventional
sequence and places contraction before ordinary short-vowel fronting
(\citeproc{ref-Campbell1959}{Campbell 1959, 52--53}). Later i-mutation
of the new long vowel is a separate conditioned development, not that
ordinary fronting (\citeproc{ref-Campbell1959}{Campbell 1959, 69}).

Versloot proposes a wave account of the regional outcomes, but CAPR does
not adopt diffusion as a solution to the comparative tree problem. His
readings of early English and Frisian inscriptions are relevant evidence
independently of that mechanism, conditional on their provenance,
etymology and phonetic interpretation
(\citeproc{ref-Versloot2017}{Versloot 2017, 295--97, 318}). The required
Anglo-Frisian ancestral node is retained. CAPR places completed English
contraction on the daughter branch, while leaving a possible earlier
conditioned onset on the common stem distinct. This is the defended
working placement under the cited runic and comparative premises, not
proof against every ancestral onset
(\citeproc{ref-Campbell1939}{Campbell 1939, 90--91};
\citeproc{ref-Versloot2017}{Versloot 2017, 295--97, 318}). The
inherited-long relation permits temporal overlap; the retained
executable order does not assert a strict order between entire events.

The live selected-corpus census has twenty-three applications, all
carrying stressed \emph{*ái}. Loam's selected \emph{*láimą} `loam' is
explicitly a pre-Old-English model input, not an independent
Proto-Germanic witness. The raw corpus's additional roe reconstruction
has no attested target and is excluded from that census. The unstressed
development \emph{*ai > *ē} is the separate earlier change
\hyperref[rule-PNWGmcUnstressedAiMonophthongization]{SC014
PNWGmcUnstressedAiMonophthongization}.

\subsection{\texorpdfstring{SC004. English stressed ai contraction
(\texttt{EAFAiMonophthongization})}{SC004. English stressed ai contraction (EAFAiMonophthongization)}}\label{rule-EAFAiMonophthongization}

\noindent\begin{minipage}{\linewidth}
\begin{ReaderFacingFoma}
define EAFAiMonophthongization Ctx([
    {*ái} -> {*ā}
]);
\end{ReaderFacingFoma}
\end{minipage}

The soul form fixes the relation to interstress raising. If the
monophthongization is delayed until after that change, PGmc
\Recon{sáiwalō} `soul' yields \Pred{sāwel} rather than expected OE
\emph{sāwol} `soul'. An earlier placement changes no output. This shows
that \hyperref[rule-EAFAiMonophthongization]{SC004
EAFAiMonophthongization} must come before
\hyperref[rule-OEInterStressRaising]{SC036 OEInterStressRaising} in the
modeled sequence.

The unstressed development \emph{*ai > *ē} in final and nonfinal
syllables is a separate and earlier change; see
\hyperref[rule-PNWGmcUnstressedAiMonophthongization]{SC014
PNWGmcUnstressedAiMonophthongization}.

\section{Resolution of awj}\label{resolution-of-awj}

\subsection{\texorpdfstring{Historical discussion of the resolution of
\emph{*awj}}{Historical discussion of the resolution of }}\label{historical-discussion-of-the-resolution-of}

Old English \emph{hīeġ} `hay' and \emph{strīeġan} `strew' go back to
forms in which \emph{*w} preceded \emph{*j}. West Germanic doubled every
consonant except \emph{*r} before \emph{*j} after a short syllable, and
\emph{*w} was an ordinary member of that law
(\citeproc{ref-Campbell1959}{Campbell 1959, 167}, §407). Proto-Germanic
\emph{*hawja-} therefore appears as West Germanic \emph{*hauuj}
(\citeproc{ref-Campbell1959}{Campbell 1959, 46}, §120.2). Campbell
writes the sequence as \emph{auj} \textgreater{} \emph{auuj}
\textgreater{} \emph{auj}, with the diphthong restored before the Old
English developments begin.

Ringe and Taylor reach the same result and explain why it is possible
(\citeproc{ref-RingeTaylor2014}{Ringe and Taylor 2014, 53}, §3.1.3).
They find the gemination of \emph{*wj} clearest where the preceding
vowel was \emph{*i}, as in Proto-Germanic \emph{*niwjaz} `new' and
\emph{*siwjaną} `sew', which give Old Saxon and Old High German
\emph{niuwi} and \emph{siuwen}. Gemination was reversible, since it
merged nothing and altered no underlying form, so the sequence Northwest
Germanic \emph{*awj} to West Germanic \emph{*[aw'w']} to pre-Old English
\emph{*[auj]} can have run its course and then undone itself. Their
derivations give Proto-Germanic \emph{*hawja} through \emph{*hauj-} to
\emph{hīeġ} `hay', and Proto-Germanic \emph{*strawjaną} through
\emph{*straujan} to Anglian \emph{strēgan} `strew'
(\citeproc{ref-RingeTaylor2014}{Ringe and Taylor 2014, 173}).

The strongest comparative argument for the doubled stage comes from
paradigms in which some cells had \emph{*j} in the ending and others had
\emph{*i}. Only the first group could double, and the two outcomes then
sat side by side. Old High German has \emph{hewi} `hay' beside
\emph{houwi}, and Old English itself preserves both in one word,
\emph{glīg} `mirth' from the undoubled nominative beside \emph{glīowes}
`of mirth' in the genitive (\citeproc{ref-RingeTaylor2014}{Ringe and
Taylor 2014, 53}, §3.1.3). The continental forms \emph{houwi} and
\emph{gistrouwen} `bestrew', and Old Saxon \emph{hoi} `hay', point to
the same West Germanic stage, which English alone went on to resolve
(\citeproc{ref-RingeTaylor2014}{Ringe and Taylor 2014, 173}).

The resolution did not treat every doubled \emph{*w} alike, and the
difference is what allows the doubling to be seen apart from it.
Campbell sets the two vowel types side by side: \emph{auj} becomes
\emph{auuj} and then \emph{auj}, while \emph{iuj} becomes \emph{iuuj}
and then \emph{iuj}. Thereafter they part company, since the \emph{u} of
\emph{auuj} is generally lost while the \emph{j} of \emph{iuuj} is lost
(\citeproc{ref-Campbell1959}{Campbell 1959, 46}, §120.2). Old English
\emph{hīeġ} `hay' and \emph{hīew} `form, hue' begin from shapes that
differ in a single vowel and end with different survivors. \emph{Hīeġ}
keeps its \emph{*j}, written \emph{ġ}, while \emph{hīew} keeps its
\emph{*w}. Campbell's other examples of the second type are \emph{nīowe}
`new', \emph{nīewe} `new' and \emph{glīow} `mirth', \emph{glīw} `mirth'.

One qualification belongs in the record. Fulk holds that \emph{*w} was
never consonantal here, so that Proto-Germanic already had
\emph{*straujaną} with its diphthong in place and there is no gemination
to undo (\citeproc{ref-Fulk2018}{Fulk 2018, 73}, §4.10 n.~1). The
account followed here is the handbook one, which the comparative
material supports: the doubled stage that the continental and
paradigm-internal forms point to is precisely the one Fulk denies ever
existed. Both accounts agree that \emph{*auj} is what enters Old
English, and the disagreement concerns whether a discrete change took
place.

\subsection{\texorpdfstring{SC029. Resolution of \emph{*awj} to
\emph{*auj}
(\texttt{OEAwwjResolution})}{SC029. Resolution of  to  (OEAwwjResolution)}}\label{rule-OEAwwjResolution}

\noindent\begin{minipage}{\linewidth}
\begin{ReaderFacingFoma}
define OEAwwjResolution Ctx([
    {*á} {*w} {*w} {*j} -> {*áu} {*j},
    {*a} {*w} {*w} {*j} -> {*au} {*j}
]);
\end{ReaderFacingFoma}
\end{minipage}

The change is confined to the \emph{*a} type. PGmc \Recon{xáwją} `hay'
is doubled to \emph{*xáwwją} by the West Germanic law and then resolved
to \emph{*xáują}, and PGmc \Recon{stráwjaną} `strew' is doubled to
\emph{*stráwwjaną} and resolved to \emph{*stráujaną}, yielding
\emph{hīeġ} `hay' and \emph{strīeġan} `strew' once the later diphthong
changes and \emph{i}-umlaut have applied. PGmc \Recon{xíwją} `form, hue'
is doubled by the same law, and there the doubling holds: it passes
through this change untouched and surfaces as \emph{hīew} `form, hue'.

The resolution has to precede \hyperref[rule-OEAuBrightening]{SC030
OEAuBrightening}, since the fronting needs a diphthong to work on. Ringe
and Taylor state the same dependence when they observe that these new
instances of \emph{*au} went on to share the ordinary development
(\citeproc{ref-RingeTaylor2014}{Ringe and Taylor 2014, 173}).

\section{English brightening of au}\label{english-brightening-of-au}

\subsection{\texorpdfstring{Historical discussion of the brightening of
\emph{*au}}{Historical discussion of the brightening of }}\label{historical-discussion-of-the-brightening-of}

The fronting of \emph{*au} is the application to a diphthong of the same
process that fronts plain \emph{*a}, the process usually called
Anglo-Frisian brightening. Campbell arrives at the point while
establishing the order of the early vowel changes, remarking that the
normal development of Proto-Germanic \emph{*au} to Old English \emph{ēa}
shows that the change of \emph{a} to \emph{æ} would affect the first
element of a diphthong (\citeproc{ref-Campbell1959}{Campbell 1959, 52},
§132). His ordered list accordingly places West Germanic \emph{a}
\textgreater{} Old English \emph{æ} and West Germanic \emph{*au}
\textgreater{} Old English \emph{*æu} in one and the same step. Fulk
puts it in a single sentence, saying that this fronting of \emph{a}
applied also to the diphthong \emph{au} in Old English
(\citeproc{ref-Fulk2018}{Fulk 2018, 73}, §4.12).

The intermediate stage is directly attested. Ringe and Taylor describe
\emph{*au} as first tensed and fronted to \emph{*æu}, a spelling still
found occasionally in eighth-century documents, with the offglide
unrounded and lowered only later (\citeproc{ref-RingeTaylor2014}{Ringe
and Taylor 2014, 172}). Early spellings such as \emph{Eadbald} with an
initial \emph{aeo} preserve the rounded offglide, and the rounding
survives in late Northumbrian (\citeproc{ref-Fulk2018}{Fulk 2018, 73},
§4.12). The outcome of the whole sequence is the \emph{ēa} of
\emph{dēaþ} `death', \emph{ēage} `eye' and \emph{lēaf} `leaf'
(\citeproc{ref-Campbell1959}{Campbell 1959, 53}, §135).

The geographical reach of the diphthongal fronting is narrower than that
of the plain one. Old Frisian shows no such fronting and has \emph{ā},
so that Old English \emph{ēac} `also', \emph{ēage} `eye' and \emph{bēam}
`tree' stand against Old Frisian \emph{āk} `also', \emph{āge} `eye' and
\emph{bām} `tree' (\citeproc{ref-Fulk2018}{Fulk 2018, 73}, §4.12;
\citeproc{ref-RingeTaylor2014}{Ringe and Taylor 2014, 172}). A later
change supplies independent confirmation. Old English \emph{gēac}
`cuckoo' has a palatalized initial, which requires a front vowel to have
followed it, while Old Frisian \emph{gāk} `cuckoo' has none
(\citeproc{ref-Fulk2018}{Fulk 2018, 73}, §4.12). Brightening of plain
\emph{*a} is shared with Frisian; brightening of the first element of
\emph{*au} is English.

\subsection{\texorpdfstring{SC030. Brightening of \emph{*au} to
\emph{*æu}
(\texttt{OEAuBrightening})}{SC030. Brightening of  to  (OEAuBrightening)}}\label{rule-OEAuBrightening}

\noindent\begin{minipage}{\linewidth}
\begin{ReaderFacingFoma}
define OEAuBrightening Ctx([
    {*au} -> {*aeu},
    {*áu} -> {*áeu}
]);
\end{ReaderFacingFoma}
\end{minipage}

Most of the words that pass through the change carry \emph{*au}
inherited straight from Proto-Germanic. PGmc \Recon{láubą} `leaf' gives
\emph{lēaf} `leaf', PGmc \Recon{stráumaz} `stream' gives \emph{strēam}
`stream', PGmc \Recon{bráudą} `bread' gives \emph{brēad} `bread', and
PGmc \Recon{dráugmaz} `dream' gives \emph{drēam} `dream'. Where a
following \emph{*j} or \emph{*i} survives long enough to cause
\emph{i}-umlaut, the \emph{ēa} appears in West Saxon as \emph{īe}, as in
\emph{ġelīefan} `believe' from PGmc \Recon{galáubijaną} `believe' and
\emph{nīed} `need' from PGmc \Recon{náudiz} `need'
(\citeproc{ref-Campbell1959}{Campbell 1959, 46}, §120.2). The two words
supplied by \hyperref[rule-OEAwwjResolution]{SC029 OEAwwjResolution},
\emph{hīeġ} `hay' and \emph{strīeġan} `strew', join this second group.

The proportions matter for what the rule is. This is the general Old
English treatment of \emph{*au}, to which the resolution of \emph{*awj}
contributes two further inputs; the history of \emph{hīeġ} `hay' and
\emph{strīeġan} `strew' does not define it.

The fronted \emph{*æu} has no independent life. It is taken up at once
by the \hyperref[rule-OEDiphthongLeveling]{SC032 OEDiphthongLeveling},
which lowers the offglide and delivers \emph{ēa}. Ringe and Taylor give
that order explicitly when they place the tensing and fronting first and
the unrounding and lowering later (\citeproc{ref-RingeTaylor2014}{Ringe
and Taylor 2014, 172}).

\section{Short pre-w diphthongs and the j-created glide
chain}\label{short-pre-w-diphthongs-and-the-j-created-glide-chain}

\subsection{Historical discussion}\label{historical-discussion-30}

Ordinary diphthongization before singleton \emph{w} must be
distinguished from long diphthong formation in endingless forms and from
the independently created \emph{wwj} sequences. Their similar Old
English spellings do not establish one sound law assigning long quantity
before every retained vowel ending.

For \emph{cneowe} `knee (dat.sg.)', the regular pre-ending stem has
short \emph{eo}. Hogg explicitly identifies the dative as a
short-diphthong example, and Campbell and Luick distinguish its genitive
and dative from the long endingless form
(\citeproc{ref-HoggGrammar2011}{Hogg 2011, 21--22, 86}, §§2.33, 5.22;
\citeproc{ref-Campbell1959}{Campbell 1959, 232--33}, §584;
\citeproc{ref-Luick1914}{Luick 1914-\/-1940, 139}, §134). Ringe and
Taylor likewise distinguish short \emph{cneow-} before syllabic endings
from long endingless \emph{cnēo} `knee'. The final \emph{w} of
\emph{cnēow} `knee' is generalized from the obliques, not retained by an
unrestricted prevocalic lengthening law
(\citeproc{ref-RingeTaylor2014}{Ringe and Taylor 2014, 187--88, 387},
§§6.2.4, 7.2.4).

The selected comparison is therefore the dative \emph{cneowe}, from the
constructed PGmc dative \emph{*knéwai}, while citation \emph{*knéwą}
remains the lexeme-level reconstruction. Fulk prefers the \emph{*-ai}
analysis for the West Germanic dative line but discusses alternatives;
the complete selected input is a paradigm construction, not a quotation
of a whole-word reconstruction (\citeproc{ref-Fulk2018}{Fulk 2018, 147},
§7.8).

Its short diphthong is supplied by \hyperref[rule-OEBreaking]{SC044
OEBreaking}, with the source's exclusion before a following high front
vocalic. Long quantity could later spread from endingless forms into
obliques, but that leveling is not the selected short dative's sound-law
history. Campbell describes the extension; Ringe and Taylor qualify the
knee quantity inference because decisive verse evidence is elusive
(\citeproc{ref-Campbell1959}{Campbell 1959, 233};
\citeproc{ref-RingeTaylor2014}{Ringe and Taylor 2014, 387}).

Earlier inherited-glide reanalysis is independently represented by
\hyperref[rule-OEWWSimplification]{SC031 OEWWSimplification}. Chew, dew,
four, hew and you complete its products through
\hyperref[rule-OEDiphthongLeveling]{SC032 OEDiphthongLeveling}. The
operation below neither repeats that inherited event nor implements the
regular singleton history (\citeproc{ref-RingeTaylor2014}{Ringe and
Taylor 2014, 41--42, 65--66, 171--75}).

\subsection{\texorpdfstring{SC033. J-created palatalized-glide
reanalysis
(\texttt{OEEwLongDiphthong})}{SC033. J-created palatalized-glide reanalysis (OEEwLongDiphthong)}}\label{rule-OEEwLongDiphthong}

\noindent\begin{minipage}{\linewidth}
\begin{ReaderFacingFoma}
define OEEwLongDiphthong Ctx([
    {*i} {*w} {*w} -> {*iu} {*w} || _ {*j},
    {*í} {*w} {*w} -> {*íu} {*w} || _ {*j}
]);
\end{ReaderFacingFoma}
\end{minipage}

\emph{Hīew} `hue' enters with inherited PGmc \emph{*xíwją}, following
Ringe and Taylor's explicit \emph{*hiwją}. Their palatalized geminate
develops through \emph{*iuw} before the separately realized long
diphthong. The conditioner is palatalized \emph{w}, carried here as
\emph{ww+j} and subsequently \emph{w+j}; this does not assert a second
independent phoneme where the authors leave the phonemic analysis
unclear (\citeproc{ref-RingeTaylor2014}{Ringe and Taylor 2014, 53,
250}).

Orel's \emph{*xewjan} and Kroonen's \emph{*heuja-} remain competing e
reconstructions, not earlier stages invented to reconcile the authors.
Their derivational explanations differ. The realized-i choice uses the
inherited raising account, not a later hue-specific law
(\citeproc{ref-Orel2003}{Orel 2003, 171--72};
\citeproc{ref-Kroonen2013}{Kroonen 2013, 224};
\citeproc{ref-Ringe2017}{Ringe 2017, 151--53};
\citeproc{ref-Fulk2018}{Fulk 2018, 59}). Campbell supports the
contrasting \emph{iwj} outcome, but Fulk questions traditional geminate
dismantling and its consonantal premises. The detailed lexical account
preserves those objections (\citeproc{ref-Campbell1959}{Campbell 1959,
46}; \citeproc{ref-Fulk2018}{Fulk 2018, 71--72, 126}).

The stable component now supplies offglide reanalysis, not long
quantity. Singleton \emph{iwj}, nonpalatal \emph{iww}, unraised
\emph{ewwj} and the separate low-vowel \emph{awj} history are excluded.
Knee remains a short singleton-w comparison. No exact historical date or
independent confidence is fabricated for this support encoding.

\section{Technical support: j-created glide
compatibility}\label{technical-support-j-created-glide-compatibility}

\subsection{\texorpdfstring{\CAPRRuleHeading{SC106. Retained j-created glide compatibility}{OEJWWSimplification}}{}}\label{rule-OEJWWSimplification}

\noindent\begin{minipage}{\linewidth}
\begin{ReaderFacingFoma}
define OEJWWSimplification Ctx([
    {*w} {*w} -> {*w} || _ {*j}
]);
\end{ReaderFacingFoma}
\end{minipage}

The unchanged technical operation remains composed for compatibility.
Hue's reanalysis already leaves a singleton glide, so it no longer feeds
this simplification. Absence of a current application does not establish
universal redundancy or a newly dated historical event. It neither
restores nominative \emph{w} in knee nor duplicates earlier inherited
short-Vww reanalysis.

\section{Technical support: palatalized-glide long-io
realization}\label{technical-support-palatalized-glide-long-io-realization}

\subsection{\texorpdfstring{SC110. Palatalized-glide long-io realization
(\texttt{OEJGlideIO})}{SC110. Palatalized-glide long-io realization (OEJGlideIO)}}\label{rule-OEJGlideIO}

\noindent\begin{minipage}{\linewidth}
\begin{ReaderFacingFoma}
define OEJGlideIO Ctx([
    [{*iu}|{*íu}] -> {*īo} || _ {*w} {*j}
]);
\end{ReaderFacingFoma}
\end{minipage}

The palatalized-glide product realizes long \emph{īo}, not the ordinary
long \emph{ēo} supplied by \hyperref[rule-OEDiphthongLeveling]{SC032
OEDiphthongLeveling}. Ringe and Taylor explicitly supply pre-OE long
\emph{īow} and early West Saxon long \emph{īew}; the retained semivowel
still conditions mutation (\citeproc{ref-RingeTaylor2014}{Ringe and
Taylor 2014, 250}). Existing \hyperref[rule-OEIUmlaut]{SC055 OEIUmlaut}
supplies the latter change. Ordinary inherited \emph{iu} without this
palatal conditioner remains outside the component, as does short
\emph{io}. This distinct realization does not adjudicate all the
inherited atomic and split-symbol clauses in
\hyperref[rule-OEDiphthongLeveling]{SC032 OEDiphthongLeveling}.

\section{Quantity-distinct English diphthong
realization}\label{quantity-distinct-english-diphthong-realization}

\subsection{Historical discussion}\label{historical-discussion-31}

The English histories of inherited diphthongs and their secondary
equivalents must be distinguished from both smoothing and paradigm
leveling. Ringe and Taylor identify a common tensing or lengthening of
the syllabic nuclei, followed by further developments of the offglides;
they place that tensing before breaking. This does not make all the
eventual English spellings the products of a single simultaneous change
(\citeproc{ref-RingeTaylor2014}{Ringe and Taylor 2014, 170--75}).

The fronted continuation of \emph{*au} gives long \emph{ēa}, including
in \emph{brēad} `bread', \emph{dēaw} `dew' and \emph{hēawan} `hew'.
First-element fronting is distinct from subsequent offglide unrounding
and lowering. The latter was not universal: some Northumbrian varieties
retained a rounded offglide. Thus this outcome is not exclusively West
Saxon, but neither is it an exceptionless description of every English
dialect (\citeproc{ref-RingeTaylor2014}{Ringe and Taylor 2014, 172--73};
\citeproc{ref-Campbell1959}{Campbell 1959, 53--54, 116--18}).

The long \emph{ēo} outcomes in \emph{ċēowan} `chew' and \emph{fēower}
`four' belong to the \emph{*eu} history. Earlier inherited or
assimilation-created geminate-glide reanalysis supplies the diphthong
while retaining the following consonantal glide. Dew and hew similarly
acquire \emph{*au} through \hyperref[rule-OEWWSimplification]{SC031
OEWWSimplification}, then reach the fronted-au path. Their early
reanalysis is not the later English realization
(\citeproc{ref-RingeTaylor2014}{Ringe and Taylor 2014, 41--42, 65--66,
172--74}).

The \emph{*iu} history is different again. Early \emph{īo} is not simply
another spelling for the initial \emph{*eu} development. Subsequent
io/eo convergence has dialect-specific histories: Northumbrian retains
io, whereas West Saxon and parts of Mercian show convergence. The long
\emph{ēo} of \emph{frēond} `friend' and \emph{ēow} `you' is the selected
endpoint here, not evidence that the merger preceded every intervening
English change. Friend's contraction-created input and you's earlier
glide reanalysis are also distinct origins
(\citeproc{ref-RingeTaylor2014}{Ringe and Taylor 2014, 41--42, 174--75};
\citeproc{ref-Campbell1959}{Campbell 1959, 124--26}).

\subsection{\texorpdfstring{\CAPRRuleHeading{SC032. Quantity-distinct diphthong realization}{OEDiphthongLeveling}}{}}\label{rule-OEDiphthongLeveling}

\begin{ReaderFacingFoma}
define OEDiphthongLeveling Ctx([
    {*aeu} -> {*ēa},
    {*áeu} -> {*ēa},
    {*eu} -> {*ēo},
    {*éu} -> {*ēo},
    {*iu} -> {*ēo},
    {*íu} -> {*ēo},
    # Split symbols from palatal glide (stress preserved where present):
    {*e} {*u} -> {*eo},
    {*é} {*u} -> {*éo},
    {*i} {*u} -> {*eo}
]);
\end{ReaderFacingFoma}

The retained definition is a compound realization proxy. Its atomic
\emph{*aeu}, \emph{*eu} and \emph{*iu} symbols belong to the heavy,
long-output tier; an acute mark does not itself indicate length. Atomic
\emph{*aeu} records a fronted-au intermediate, not a claim that a
three-vowel phonetic sequence is demonstrated. The direct \emph{*iu} to
\emph{ēo} mapping packages the selected later endpoint. The split-symbol
e+u, é+u and i+u clauses instead produce short \emph{eo}, retaining
stress only in the é+u clause. Campbell distinguishes short breaking and
back-umlaut diphthongs from their long equivalents, but his account does
not specify these computational token boundaries
(\citeproc{ref-Campbell1959}{Campbell 1959, 116--18}).

All 32 current applications use atomic long-tier inputs. The unaccented
atomic alternatives and all three split-symbol clauses remain formal
compatibility paths without a selected-word application. Their retention
does not establish an independent palatal-glide origin, and no
additional stressed-i clause is inferred.

The short dative \emph{cneowe} `knee' remains outside this long-output
operation; its pre-w breaking is supplied by
\hyperref[rule-OEBreaking]{SC044 OEBreaking}. The separately
palatalized-glide history of \emph{hīew} `hue' reaches long \emph{īo}
through \hyperref[rule-OEJGlideIO]{SC110 OEJGlideIO} before the ordinary
iu matcher. It then receives long \emph{īe} through mutation, not
ordinary \emph{ēo} here. These contrasts preserve quantity and
conditioning within the adopted selected cells
(\citeproc{ref-Campbell1959}{Campbell 1959, 232--33};
\citeproc{ref-RingeTaylor2014}{Ringe and Taylor 2014, 187--88, 250,
387}).

Fronting must supply its product before its offglide continuation:
\hyperref[rule-OEAuBrightening]{SC030 OEAuBrightening} feeds this
component, not every clause in the definition. Ringe and Taylor
explicitly distinguish the two developments
(\citeproc{ref-RingeTaylor2014}{Ringe and Taylor 2014, 172}). Earlier
displacement tests that left \emph{*aeu} unrealized produced \emph{+?},
a computational rejection with no attested linguistic status. They
therefore cannot independently establish the historical interval.

The input of \emph{hēafod} `head' must be realized into a recognized
vowel before \hyperref[rule-OEMedUnstressedULowering]{SC040
OEMedUnstressedULowering} can lower its remaining medial vowel in the
model. That internal recognition requirement is not proof that
historical offglide change conditioned the lowering. Campbell records
both \emph{hēafud} and \emph{hēafod}, with varying preservation of
protected unstressed u (\citeproc{ref-Campbell1959}{Campbell 1959,
155--56}). No historical order for the entire nine-clause bundle follows
from that one computational dependency.

\section{\texorpdfstring{Endingless \emph{aw} and the retained show
proxy}{Endingless  and the retained show proxy}}\label{endingless-and-the-retained-show-proxy}

\subsection{Historical discussion}\label{historical-discussion-32}

The attested long nucleus does not establish one general law changing
singleton \emph{aw} to \emph{ēaw} before every following vowel. For
straw, the inherited nominative/accusative \emph{*stráwą} loses its
ending and vocalizes final \emph{w}, giving \emph{*strau}; the ordinary
English development then gives reconstructed \emph{*strēa}. Attested
\emph{strēaw} `straw' combines that long diphthong with \emph{w}
restored from inflected forms. Conversely, the long diphthong spreads
into those forms. The regular free simplex and short obliques are
explicitly unattested, although the endingless pattern survives in the
compound \emph{strēaberige} `strawberry'
(\citeproc{ref-Campbell1959}{Campbell 1959, 46--47, 232--33};
\citeproc{ref-SieversBrunner1965}{Brunner 1965, 49, 110, 204--5};
\citeproc{ref-RingeTaylor2014}{Ringe and Taylor 2014, 172--73, 387}).

The model retains attested \emph{strēaw} `straw' as its comparison
target and produces reconstructed \emph{*strēa} through the regular
cascade. This is an explicit mismatch, not an attested reconstructed
endpoint or a sound law for analogical restoration.
\hyperref[rule-PWGmcEndinglessAw]{SC111 PWGmcEndinglessAw} supplies the
narrow endingless fragment; it does not imply that general final-vowel
loss was conditioned by an adjacent \emph{w}.

The inherited geminates of \emph{dēaw} `dew' and \emph{hēawan} `hew'
instead undergo early \hyperref[rule-OEWWSimplification]{SC031
OEWWSimplification}, producing \emph{*au} plus retained \emph{*w},
followed by English fronting and
\hyperref[rule-OEDiphthongLeveling]{SC032 OEDiphthongLeveling}. Campbell
distinguishes the glide histories; Ringe and Taylor distinguish earlier
reanalysis from later English realization
(\citeproc{ref-Campbell1959}{Campbell 1959, 45--47, 53--54};
\citeproc{ref-RingeTaylor2014}{Ringe and Taylor 2014, 65--66, 171--75}).
These retained-glide histories differ from endingless straw.

For \emph{sċēawian} `show', Orel and Kroonen print singleton dictionary
inputs, while Campbell and Brunner include the verb among inherited
geminate examples. The class-II suffix places a vowel between root
\emph{w} and \emph{j}; this excludes direct root \emph{aw+j} gemination
but does not resolve the different root histories
(\citeproc{ref-Orel2003}{Orel 2003, 337};
\citeproc{ref-Kroonen2013}{Kroonen 2013, 442};
\citeproc{ref-Campbell1959}{Campbell 1959, 45};
\citeproc{ref-SieversBrunner1965}{Brunner 1965, 49}).
\hyperref[rule-OEAwLongDiphthong]{SC034 OEAwLongDiphthong} therefore
remains a disclosed support proxy for the three selected show cells, not
an established general singleton sound law or a newly dated West Saxon
event.

\subsection{\texorpdfstring{SC034. Retained show long-diphthong proxy
(\texttt{OEAwLongDiphthong})}{SC034. Retained show long-diphthong proxy (OEAwLongDiphthong)}}\label{rule-OEAwLongDiphthong}

\noindent\begin{minipage}{\linewidth}
\begin{ReaderFacingFoma}
define OEAwLongDiphthong Ctx([
    {*a} {*w} -> {*ēa} {*w} || _ [EnglishStarVocalic | {*ô}],
    # Stressed variant (prosodic tier):
    {*á} {*w} -> {*ḗa} {*w} || _ [EnglishStarVocalic | {*ô}]
]);
\end{ReaderFacingFoma}
\end{minipage}

The current operation changes the show infinitive, imperative singular
and present third singular. Straw already has its regular long nucleus
without \emph{w} and is unchanged here. Dew and hew no longer
demonstrate a local simplification-then-lengthening chain here: their
earlier reanalysis has already removed that input shape. The older
displacement results for those geminates describe the superseded
representation, not independent dates for the adopted early event. The
two proxy clauses remain unchanged; their successful show finals do not
decide the unresolved singleton-versus-geminate source premise. The old
straw displacement observation measured dependence on a following vowel
in the superseded proxy, not independent historical order
(\citeproc{ref-Campbell1959}{Campbell 1959, 45--47};
\citeproc{ref-RingeTaylor2014}{Ringe and Taylor 2014, 171--75}).

\section{Former compound-syncope trace
anchor}\label{former-compound-syncope-trace-anchor}

The former early deletion is no longer active. Ordinary nonhigh loss
belongs to \hyperref[rule-OENonHighSyn]{SC114 OENonHighSyn}, and
weight-conditioned high-vowel loss to
\hyperref[rule-OEMedialSyncope]{SC065 OEMedialSyncope}. The comparative
compound discussion accompanies those later operations. This identity
preserves earlier cross-references without representing an independently
dated early sound law.

\subsection{\texorpdfstring{\CAPRRuleHeading{SC037. Former compound-syncope trace anchor}{OECompoundLinkingSyncope}}{}}\label{rule-OECompoundLinkingSyncope}

\noindent\begin{minipage}{\linewidth}
\begin{ReaderFacingFoma}
define OECompoundLinkingSyncope ?*;
\end{ReaderFacingFoma}
\end{minipage}

\section{Weak medial vocalism}\label{weak-medial-vocalism}

\subsection{Historical discussion}\label{historical-discussion-33}

\subsubsection{The inherited compound `age of
men'}\label{the-inherited-compound-age-of-men}

Old English \emph{weorold} `world' continues the compound `age of men'.
Orel's \emph{*wira-alđiz} and Ringe and Taylor's \emph{*weraldiz}
support an inherited i-stem starting point. CAPR selects the compact
nominative \emph{*wír-àldiz}, not an early ō-stem replacement. Kroonen
supports the two components, but neither their exact junction nor the
compound's accentuation follows automatically from their separate
dictionary entries (\citeproc{ref-Orel2003}{Orel 2003, 462};
\citeproc{ref-Kroonen2013}{Kroonen 2013, 20, 588};
\citeproc{ref-RingeTaylor2014}{Ringe and Taylor 2014, 34, 202, 326}).

Original \emph{i} beside \emph{e} remains a comparative uncertainty. The
selected \emph{i} undergoes \hyperref[rule-PNWGmcILowering]{SC015
PNWGmcILowering} before \emph{r}; an originally \emph{e}-vowelled
reconstruction would not need that step. Early short-\emph{i} apocope
removes the final vowel of \emph{-iz}, and subsequent final-\emph{z}
loss leaves \emph{*wér-àld}. The compact stem presupposes resolution or
reduction of the earlier compound junction. Its formation is not derived
by treating a literal \emph{a+a} sequence as an ordinary
consonant-separated linking vowel (\citeproc{ref-Orel2003}{Orel 2003,
462}; \citeproc{ref-RingeTaylor2014}{Ringe and Taylor 2014, 34, 326}).

This nominative does not require early ō-stem analogy. That does not
mean its paradigm never underwent remodelling. Ringe and Taylor's i-stem
versus ō-stem comparison and Kluge--Seebold's \emph{*wira-aldō} remain
relevant; late Norse \emph{verǫld} `world' is probably borrowed and is
not an independent inherited vote. Campbell's frequent endingless
accusative \emph{weorold} supports survival of conservative inflection
without requiring an entirely unchanged paradigm
(\citeproc{ref-KlugeSeebold2011}{Kluge 2011, 981--82};
\citeproc{ref-RingeTaylor2014}{Ringe and Taylor 2014, 202, 326};
\citeproc{ref-Campbell1959}{Campbell 1959, 244}, §610(1) n.2).

\subsubsection{Earlier secondary accent and its
weakening}\label{earlier-secondary-accent-and-its-weakening}

An earlier compound could carry secondary prominence on the second root.
Kaluza's older account, presented in Stanley's \emph{Metricalities},
allows subsidiary stress in \emph{weorold} and many other disyllables.
This is a metrical theory, not direct written evidence of pronunciation
(\citeproc{ref-Stanley2022}{Stanley 2022, 34}). Goering instead argues
that early Old English \emph{weruld} belonged to a single initial
trimoraic foot, not a light initial foot followed by a separately
stressed medial foot. Initial alliteration and vowel development
establish initial stress; apocope alone cannot choose between the
footings. The weak vowel's behaviour and resolution of \emph{worulde} in
\emph{Beowulf} 3068a and 2343a provide the stronger discriminants
(\citeproc{ref-Goering2023}{Goering 2023, 79--82}, §4.5.1.2).

The adopted relative claim is loss of independent second-element
prominence before the weak-vowel development. Goering's reconstructed
early Old English system, perhaps around 600 and compatible with the
seventh-/eighth-century poetic horizon, provides an approximate upper
horizon for that analysis, not a precisely dated deaccentuation law
(\citeproc{ref-Goering2023}{Goering 2023, 84--85}, §4.6). The grave
therefore represents an earlier accent; it does not assert that the
vowel still bore independent stress when it rounded. Hogg's warning
against an unexplained ``tertiary stress'' remains important: a
descriptive label cannot supply the missing phonological argument
(\citeproc{ref-HoggGrammar2011}{Hogg 2011, 215--18}, §§6.4--6.7).

\subsubsection{Weak rounding and the selected high-vowel
path}\label{weak-rounding-and-the-selected-high-vowel-path}

Campbell describes weak retraction before \emph{l/r} plus consonant,
rounding to \emph{o}, and further \emph{u} in some forms. Hogg
distinguishes weak retraction from breaking under stronger prominence.
Bülbring proposes following-\emph{u} assimilation for \emph{world}.
Luick, however, particularly associates weak rounding with labials and
following \emph{l}, and expressly finds the following-\emph{u} account
insufficient. His uncertainty over the boundary between \emph{o} and
\emph{u} outcomes does not exclude a high-vowel lineage
(\citeproc{ref-Campbell1959}{Campbell 1959, 142}, §338;
\citeproc{ref-Bulbring1902}{Bülbring 1902, 147--49}, §§366--367;
\citeproc{ref-Luick1914}{Luick 1914-\/-1940, 125--26}, §113;
\citeproc{ref-HoggGrammar2011}{Hogg 2011, 215--18}, §§6.4--6.7).

CAPR adopts weak \emph{a > u} before \emph{l} plus consonant in a
light-initial domain: an initial primary-stressed short vowel followed
by one consonant and the weak medial vowel. This bounded prosodic
generalization combines Luick's environment with Goering's LH analysis;
neither author supplies this exact formal rule. It excludes
heavy-initial shapes, primary medial accent and vowels outside the
specified cluster, without identifying words or grammatical classes.
Retained-prominence compounds such as \emph{ānfeald} `onefold' caution
against unrestricted post-accent-stripping raising; weak \emph{twifold}
`twofold' supplies a relevant rounded comparison. The rule is a
qualified reconstruction of the chosen lineage, not a claim that every
historically weak vowel became high
(\citeproc{ref-Campbell1959}{Campbell 1959, 142}, §338;
\citeproc{ref-Bulbring1902}{Bülbring 1902, 148--49}, §367;
\citeproc{ref-HoggGrammar2011}{Hogg 2011, 217--18}, §§6.6--6.7;
\citeproc{ref-Goering2023}{Goering 2023, 79--82}).

One rewrite compresses possible gradual rounding and raising; it does
not posit an independently dated intermediate \emph{o} law. The
resulting medial \emph{u} feeds initial back mutation in
\hyperref[rule-OEBackMutation]{SC059 OEBackMutation}.
\hyperref[rule-OEMedUnstressedULowering]{SC040 OEMedUnstressedULowering}
later supplies medial \emph{o}, giving the selected West Saxon
\emph{weorold}. The inherited nominative has already lost final short
\emph{i}; no later final high vowel needs to be removed. The initial and
medial dialect differences are separate: West Saxon
\emph{weorold/worold}, also with medial \emph{u}, Mercian
\emph{weoruld}, Northumbrian \emph{woruld} and Kentish \emph{wiarald} do
not require one universal \emph{a > u > o} chain. Subsequent syncopated
\emph{world} concerns loss of the medial vowel, not another raising
witness (\citeproc{ref-RingeTaylor2014}{Ringe and Taylor 2014, 202,
326}; \citeproc{ref-Campbell1959}{Campbell 1959, 142}, §338, p.~160,
§391; \citeproc{ref-Luick1914}{Luick 1914-\/-1940, 126}, §113).

\subsubsection{Placement and its limits}\label{placement-and-its-limits}

The selected weak-domain raising precedes early unstressed fronting and
breaking. A post-fronting account, compressing weak \emph{æ} retraction,
rounding and raising, also derives \emph{weorold} and is closer to
Hogg's stated sequence. Output alone cannot choose the dates. The
adopted serialization avoids requiring an otherwise unobserved fronted
intermediate in this weak domain. Versloot's alternative
conditional-fronting/later-breaking chronology challenges compulsory
fronting before \emph{lC}; its main argument concerns other domains,
especially stressed \emph{a}, and is not evidence that he proposed weak
medial \emph{a > u} in \emph{world} (\citeproc{ref-HoggGrammar2011}{Hogg
2011, 217--18}, §§6.6--6.7; \citeproc{ref-Versloot2025}{Versloot 2026,
111--14, 122--23}).

The earlier grave survives until
\hyperref[rule-OECompoundLinkingSyncope]{SC037 OECompoundLinkingSyncope}
has read the relevant accent information; technical grave removal then
precedes weak raising. This represents annotation lifetime and a
compatible relative history, not proof that every linking syncope
preceded every compound's deaccentuation. The chronology of \emph{world}
is independent of the early second-root lowering in \emph{reġnboga}
`rainbow'. That word still requires the grave through linking syncope,
but its \emph{o} already arose under
\hyperref[rule-PNWGmcULowering]{SC017 PNWGmcULowering}; it does not date
grave loss before late unstressed lowering. \emph{World} retains the
distinct requirement that the weak vowel lack independent prominence
before raising. The back-mutation implementation also admits lower
back-vowel triggers through shared contexts; the selected high-\emph{u}
history succeeds under the narrower intended \emph{u}-triggered
interpretation as well. Neither that separate defect nor the later
vowel-merger laws are changed here.

\subsubsection{The distinct history of
soul}\label{the-distinct-history-of-soul}

Ringe and Taylor derive \emph{*saiwalō} through \emph{*saiwalu},
\emph{*sawælu}, \emph{*sawlu} and \emph{sāwl} `soul', then
\emph{sāwul/sāwol} `soul'. Contraction of the stressed diphthong, loss
of the medial nonhigh vowel, loss of final \emph{u}, and epenthesis are
separate developments. Their final \emph{u} is described as restored
after earlier regular loss (\citeproc{ref-RingeTaylor2014}{Ringe and
Taylor 2014, 265, 330}). Campbell's \emph{sāwl} `soul' and Luick's
epenthetic \emph{sāwul} `soul' support this distinction; the oblique
\emph{sāwle} `soul' also matters. The expanded \emph{sāwol} `soul' does
not license suppressing the unexpanded forms
(\citeproc{ref-Campbell1959}{Campbell 1959, 143}, §341;
\citeproc{ref-Luick1914}{Luick 1914-\/-1940, 287, 294}, §§309, 317;
\citeproc{ref-SieversBrunner1965}{Brunner 1965, 132}, §159).

Fulk's competing metrical account of the early syllabic history must be
evaluated independently of later spellings (\citeproc{ref-Fulk2010}{Fulk
2010, 133--39}). CAPR selects the inherited accusative \emph{*sáiwalǭ},
yielding \emph{sāwle} by \hyperref[rule-OENonHighSyn]{SC114
OENonHighSyn}, without medial raising, epenthesis or analogy inside that
span. Luke12:20 provides an unambiguous accusative; the nearby
\emph{eala sawel} is vocative, not nominative
(\citeproc{ref-Skeat1874}{Skeat 1874, 130}). The wider nominative
history remains distinct and unresolved.

The former continuous medial-raising proxy is retired. It was never the
source-led syncope/apocope/epenthesis history, and its temporary success
did not demonstrate one common raising event for \emph{world} and
\emph{soul}.

\subsection{\texorpdfstring{SC036. Weak medial vocalism
(\texttt{OEInterStressRaising})}{SC036. Weak medial vocalism (OEInterStressRaising)}}\label{rule-OEInterStressRaising}

\noindent\begin{minipage}{\linewidth}
\begin{ReaderFacingFoma}
define OEInterStressRaising OELHWeakRaising;
\end{ReaderFacingFoma}
\end{minipage}

The primary network now names only the adopted weak-LH operation. It
does not include the distinct nonhigh syncope process. Confidence C
expresses the qualified domain and chronology, independently of the Old
English scope and secure attestations.

\subsection{\texorpdfstring{SC036. Light-initial weak raising
(\texttt{OELHWeakRaising})}{SC036. Light-initial weak raising (OELHWeakRaising)}}\label{rule-OELHWeakRaising}

\noindent\begin{minipage}{\linewidth}
\begin{ReaderFacingFoma}
define OELHWeakRaising Ctx([
    {*a} -> {*u} || .#. EnglishStarConsonant* [{*á}|{*é}|{*í}|{*ó}|{*ú}]
        EnglishStarConsonant _ {*l} EnglishStarConsonant
]);
\end{ReaderFacingFoma}
\end{minipage}

The boundary and short acute nucleus restrict the initial domain; one
intervening consonant distinguishes the light initial syllable. Only
plain short \emph{a} is rewritten, after removal of the earlier grave.

\section{Combinative u-umlaut in
wi-forms}\label{combinative-u-umlaut-in-wi-forms}

\subsection{Historical discussion}\label{historical-discussion-34}

The history of \emph{wuduwe} `widow' orders these two changes within the
same low-stress vocalic development. Campbell discusses both the
\emph{w}-conditioned \emph{u} forms and the later \emph{weorold} `world'
/ \emph{weoruld} `world' alternation, while Ringe and Taylor give the
same connection comparatively in \emph{*widuwon-}, \Recon{weraldu}
`world', and \Recon{jugunþi} `youth'
(\citeproc{ref-Campbell1959}{Campbell 1959, 92}, §218;
\citeproc{ref-Campbell1959}{Campbell 1959, 140}, §332;
\citeproc{ref-Campbell1959}{Campbell 1959, 141--42}, §§338--339;
\citeproc{ref-RingeTaylor2014}{Ringe and Taylor 2014, 267};
\citeproc{ref-RingeTaylor2014}{Ringe and Taylor 2014, 322}, §6.9.4).

\hyperref[rule-OEWICombinativeUUmlaut]{SC039 OEWICombinativeUUmlaut}
feeds the vowel sequence subsequently reshaped by
\hyperref[rule-OEMedUnstressedULowering]{SC040
OEMedUnstressedULowering}. Initial \emph{w} conditions the first change.

\subsection{\texorpdfstring{SC039. Combinative \emph{*u}-umlaut in
\emph{wi}-forms
(\texttt{OEWICombinativeUUmlaut})}{SC039. Combinative -umlaut in -forms (OEWICombinativeUUmlaut)}}\label{rule-OEWICombinativeUUmlaut}

\noindent\begin{minipage}{\linewidth}
\begin{ReaderFacingFoma}
define OEWICombinativeUUmlaut Ctx([
    {*í} -> {*ú} || .#. {*w} _ EnglishStarConsonant [{*u} | {*o}]
]);
\end{ReaderFacingFoma}
\end{minipage}

The \emph{wuduwe} `widow' derivation answers one narrow question about
\emph{wi}-forms. If \hyperref[rule-OEWICombinativeUUmlaut]{SC039
OEWICombinativeUUmlaut} is delayed until after
\hyperref[rule-OEMedUnstressedULowering]{SC040
OEMedUnstressedULowering}, PGmc \Recon{wíduwōn} `widow' yields
\Pred{wudowe} rather than expected OE \emph{wuduwe}. This one-sided
witness requires combinative u-umlaut to precede medial lowering and
supplies no lower boundary.

\subsection{Historical limits of the represented
account}\label{historical-limits-of-the-represented-account}

The phonetic relevance of preceding \emph{w} and a following back vowel
is independently supported by \emph{wudu} `wood', \emph{wuduwe} `widow'
and West Saxon \emph{wucu} `week'. Unlike ordinary West Saxon back
umlaut, this retraction can occur across a dental or velar. These
comparisons do not, however, establish that every inherited short
\emph{i} in that environment acquired \emph{u}
(\citeproc{ref-Campbell1959}{Campbell 1959, 92--93}, §§218--221;
\citeproc{ref-HoggGrammar2011}{Hogg 2011, 158--63}, §§5.108--5.112;
\citeproc{ref-SieversBrunner1965}{Brunner 1965, 93--94}, §114;
\citeproc{ref-RingeTaylor2014}{Ringe and Taylor 2014, 322--24};
\citeproc{ref-Fulk2018}{Fulk 2018, 68--69}, §4.8).

The independently reconstructed \emph{*witun} `they know' supplies a
substantive limit: the present rule predicts a root \emph{u}, whereas
the cited West Saxon forms \emph{witon} and \emph{wioton} retain
\emph{i} or show ordinary back umlaut. This is a historical comparison,
not a rejected artificial word. The general conditioning is therefore
deferred; the selected widow derivation is retained without claiming a
complete account of all these reflexes. No lexical exception,
grammatical condition or optional sound law is introduced. Brunner's
specific optative-based explanation of \emph{witon} `let us' cannot
simply be transferred to `they know'.

Barrack distinguishes combinative retraction from ordinary back umlaut
and expressly leaves their relative order in West Saxon uncertain. His
early-glossary argument does not by itself establish an earlier
historical date for the combinative change; Hogg explicitly questions
that inference. Bülbring's account places a two-step development around
700, but the present representation does not establish those separate
steps or a new ordering against ordinary back umlaut
(\citeproc{ref-Barrack1975}{Barrack 1975, 63--65}, rules 210, 212, 218,
220; \citeproc{ref-Bulbring1902}{Bülbring 1902, 106--7}, §264;
\citeproc{ref-HoggGrammar2011}{Hogg 2011, 162--63}, §5.112).

\section{\texorpdfstring{Final bare-\emph{a}
loss}{Final bare- loss}}\label{final-bare--loss}

\subsection{Historical discussion}\label{historical-discussion-35}

The loss of final short low vowels belongs to the erosion of Germanic
final syllables, but not every part of that erosion is one event.
Campbell distinguishes the early loss of vowels which had already been
absolutely final in Indo-European from the later loss of protected
\emph{*a}; the latter includes the endings represented by nominative
\emph{*stainaz} and accusative \emph{*stainam}. His subsequent reference
to final loss introduces a separate discussion of prehistoric English
medial syncope, not a new Old English date for the final loss
{[}Campbell (\citeproc{ref-Campbell1959}{1959}), pp.~137--138, §331,
especially §331.4 and n.~4; p.~143, §341{]}.

Ringe and Taylor reconstruct West Germanic loss of unstressed short oral
\emph{*a} and nasal \emph{*ą}, word-finally or before final \emph{*z}.
Fulk likewise distinguishes this loss from subsequent reductions and
states that it was independent of the weight of the preceding syllable
(\citeproc{ref-RingeTaylor2014}{Ringe and Taylor 2014, 44--47}, §3.1.2;
\citeproc{ref-Fulk2018}{Fulk 2018, 91}, §5.6). The conditioning is
phonological and prosodic, not membership in a declension. Retention
before final \emph{*s}, contrasted with loss in \emph{*-az}, is
particularly important: a rule deleting a vowel before any final
consonant would be too broad.

Loss of \emph{*z} before apocope is the economical West Germanic
analysis, not a logical consequence of an endingless descendant. Old
Norse demonstrates the possibility of the opposite order. The West
Germanic survival of vowels before \emph{*s} and \emph{*r} favors the
\emph{*z}-first account, while the possible Matronae-name
counterexamples have uncertain linguistic affiliation and different
syllabic positions (\citeproc{ref-RingeTaylor2014}{Ringe and Taylor
2014, 46}). The Frienstedt comb inscription \emph{kaba} `comb',
interpreted as \emph{kamba} and dated approximately 250--300 CE,
provides a more concrete local witness: on that interpretation,
\emph{*z} has gone while \emph{*a} remains. Fulk expressly cautions
against extending that date to completed \emph{*z}-loss throughout West
Germanic (\citeproc{ref-Fulk2018}{Fulk 2018, 25}, n.~1). Neither
observation supplies an exact universal date for apocope; the runic and
Latinized evidence remains qualified
(\citeproc{ref-RingeTaylor2014}{Ringe and Taylor 2014, 45}, n.~4).

The broader apocope also stranded glides, producing syllabic outcomes,
and preceded West Germanic consonant gemination before \emph{*j}. Ringe
and Taylor distinguish those early outcomes from later paradigmatic
leveling (\citeproc{ref-RingeTaylor2014}{Ringe and Taylor 2014,
46--47}). This is a historical constraint on the broader account, not a
reason to date apocope from the late position of the following
executable fragment. In particular, a successful geminated ja-stem
output does not by itself establish an uninterrupted regular history of
its selected nominative.

\subsection{\texorpdfstring{SC041. Loss of final bare \emph{*a}
(\texttt{PWGmcFinalBareALoss})}{SC041. Loss of final bare  (PWGmcFinalBareALoss)}}\label{rule-PWGmcFinalBareALoss}

\noindent\begin{minipage}{\linewidth}
\begin{ReaderFacingFoma}
define PWGmcFinalBareALoss Ctx([
    {*a} -> 0 || _ .#.
]);
\end{ReaderFacingFoma}
\end{minipage}

This is a narrower proxy than the historical change. It deletes an
unaccented, short, oral final \emph{*a}; it does not delete nasal
\emph{*ą}, long \emph{*ā}, stressed \emph{*á}, or a vowel still
protected by a consonant. Nasal endings have a separate executable
treatment, including \hyperref[rule-OEHeavySyllableNasalApocope]{SC047
OEHeavySyllableNasalApocope}; that division of labor does not establish
an Old English date for their historical loss.

The adjudicated population contains seventy live applications, all with
selected inputs ending in \emph{*-az}. That agreement is a consequence
of the inputs and preceding sound changes, not a grammatical condition
on the rule. The inherited oblique used for \emph{fyrhte} `fright' is
among them, so they are not simply seventy nominatives of one
declension. The selected \emph{*stáinaz} gives \emph{stān} `stone'; an
accusative discussed in the comparative literature is not silently
substituted for this selected nominative. Nasal-final \emph{horn}
`horn', the infinitive \emph{bacan} `bake', and consonant-protected
\emph{bannes} `ban' do not undergo this particular operation.

\hyperref[rule-EAFFinalZDeletion]{SC020 EAFFinalZDeletion} exposes the
final vowel in \emph{*bárdaz} and \emph{*kámbaz}, giving the immediately
pre-loss forms \emph{*bárda} and \emph{*kámba}. The live rule removes
their final vowel and eventually gives \emph{beard} `beard' and
\emph{camb} `comb'. Moving it before the exposing loss instead gives
\Pred{bearda} and \Pred{camba}. This is feeding across a broad
computational interval, not evidence that the two historical changes
were adjacent.

The relationship to \hyperref[rule-OEARestoration]{SC046 OEARestoration}
is different. In the live order, apocope removes a potential back-vowel
trigger before restoration can use it: loss therefore bleeds
restoration. For \emph{cræft} `craft' and \emph{dæġ} `day', delaying
loss leaves a trigger present and gives \Pred{craft} and \Pred{dag}.
This does not mean that restoration requires an environment created by
apocope; it is the environment destroyed by apocope that matters. The
source account of restoration independently requires a surviving back
vowel after the relevant consonantal material
(\citeproc{ref-RingeTaylor2014}{Ringe and Taylor 2014, 189--90},
§6.3.1).

The selected fright derivation imposes another, explicitly computational
requirement: loss exposes final \emph{*n} for
\hyperref[rule-NWGmcInStemNLoss]{SC064 NWGmcInStemNLoss}. That
visibility dependency is not independent proof of adjacent historical
dates. The present position after medial developments is accordingly
retained as an implementation accommodation, not as the historical
chronology of West Germanic apocope. The unchanged operation correctly
removes the exposed short oral vowel in its current firing population;
it does not, by itself, certify the earlier history of nasal endings,
stranded glides, or every selected paradigm cell.

\section{Stem-final n-loss after long
i}\label{stem-final-n-loss-after-long-i}

\subsection{Historical discussion}\label{historical-discussion-36}

The inherited feminine \emph{in}-stem represented by Gothic
\emph{faurhtei}, OE \emph{fyrhtu}, and oblique OE \emph{fyrhte} supplies
the relevant evidence (\citeproc{ref-Orel2003}{Orel 2003, 120};
\citeproc{ref-RingeTaylor2014}{Ringe and Taylor 2014, 380--81};
\citeproc{ref-Campbell1959}{Campbell 1959, 236}, §589.7). Final
\emph{*n} after unstressed \emph{*ī} was lost after Proto-West-Germanic
differentiation; conservative southern Old High German prevents treating
it as an undifferentiated Northwest Germanic law
(\citeproc{ref-RingeTaylor2014}{Ringe and Taylor 2014, 87--88};
\citeproc{ref-Fulk2018}{Fulk 2018, 173}, §7.34). This early loss must be
distinguished from the later weak-syllable developments described by
Hogg, Ringe and Taylor, and Fulk (\citeproc{ref-HoggPhonology1992}{Hogg
1992b, 121}; \citeproc{ref-RingeTaylor2014}{Ringe and Taylor 2014,
264--303}, §§6.7.3--6.8.4; \citeproc{ref-Fulk2018}{Fulk 2018, 91},
§5.6).

\subsection{\texorpdfstring{SC064. Loss of stem-final \emph{*n} after
long \emph{*ī}
(\texttt{NWGmcInStemNLoss})}{SC064. Loss of stem-final  after long  (NWGmcInStemNLoss)}}\label{rule-NWGmcInStemNLoss}

\noindent\begin{minipage}{\linewidth}
\begin{ReaderFacingFoma}
define NWGmcInStemNLoss Ctx([{*n} -> 0 || {*ī} _ .#.]);
\end{ReaderFacingFoma}
\end{minipage}

Only word-final \emph{*n} after unstressed long \emph{*ī} is at issue.
The conditioning is phonological, not membership in the feminine
\emph{īn}-stems. Besides their obliques, Campbell cites \emph{strenge}
`strength', the loanword \emph{pyle} `pillow' and early past-subjunctive
plural endings (\citeproc{ref-Campbell1959}{Campbell 1959, 189}, §473;
\citeproc{ref-SieversBrunner1965}{Brunner 1965, 225--26}, §280).

CAPR models the oblique OE form through the Proto-Germanic genitive
singular \emph{*fúrxtīnaz} `fright', following the project convention of
using an appropriate non-nominative paradigm cell when the nominative
does not supply the required derivation. Final ending loss exposes
\emph{*fúrxtīn}; early loss of \emph{*n}, mutation and later shortening
yield \emph{fyrhte} `fright'. This selected encoding imposes two
implementation boundaries. Before
\hyperref[rule-PWGmcFinalBareALoss]{SC041 PWGmcFinalBareALoss}, PGmc
\emph{*fúrxtīnaz} `fright' yields \Pred{fyrhten} rather than expected OE
\emph{fyrhte} `fright'. After
\hyperref[rule-OEUnstressedLongVowelShortening]{SC072
OEUnstressedLongVowelShortening}, PGmc \emph{*fúrxtīnaz} again yields
\Pred{fyrhten} rather than expected \emph{fyrhte} `fright'. The former
boundary exposes the segment; the latter keeps its long-vowel
environment visible. Neither dates the loss to a late post-apocope
period.

Campbell attributes retained \emph{n} in \emph{nēten} `livestock' to the
inflected forms (\citeproc{ref-Campbell1959}{Campbell 1959, 189}, §473).
Goering's later \emph{*nētīn} therefore does not require immunity to the
early law: a later paradigmatic state and uninterrupted inherited
retention are different claims. An exposed early final \emph{*-īn} would
regularly lose its consonant; the later \emph{*nētīn} cannot be
presented as an exemption or as uninterrupted phonological survival. Its
complete ancestral derivative and intervening paradigmatic history are
separate questions. Final \emph{*ī} remains long until later shortening,
while medial \emph{*ī} has already shortened in the inflected forms
(\citeproc{ref-Goering2016}{Goering 2016, 186--88};
\citeproc{ref-Goering2023}{2023, 67--68}). No word-specific protection
or added accent is required. The exact date of paradigmatic retention or
restoration is not fixed by these comparisons.

The selected plural avoids that singular problem. Fulk prints
\emph{*nautīnō}, the ancestor of the Mercian plural, while the Hatton
\emph{Pastoral Care} attests West Saxon \emph{nietenu} `animals,
livestock' (\citeproc{ref-Fulk2018}{Fulk 2018, 92}, §5.6;
\citeproc{ref-Sweet1871}{Sweet 1871, 109}, lines 8--11;
\citeproc{ref-RingeTaylor2014}{Ringe and Taylor 2014, 298--99}). Its
medial nasal is followed by the plural vowel throughout this early loss
and therefore survives regularly. The full plural reaches \emph{nīetenu}
`animals' without either an exception to this law or a transponent with
restored final \emph{n}.

\section{\texorpdfstring{Surviving bimoric \emph{*ō}
unrounding}{Surviving bimoric  unrounding}}\label{surviving-bimoric-unrounding}

\subsection{Historical discussion}\label{historical-discussion-37}

The development has independent comparative support, not merely a
successful derivation of one Old English word. Ringe and Taylor
explicitly reconstruct surviving bimoric \emph{*ō} as West Germanic
\emph{*ā} word-finally and before final \emph{*r}; trimoric vowels
instead give \emph{*ō}, while the two quantities merge as \emph{*ō} in
other unstressed environments (\citeproc{ref-RingeTaylor2014}{Ringe and
Taylor 2014, 58--60}). Quantity, finality and the remaining consonants
are therefore essential to their account. A phonological law, not
membership in a declension, distinguishes the outcomes.

Their comparative examples include the histories of \emph{wæter} `water'
and \emph{fēower} `four' before final \emph{*r}, and final-vowel
histories continued by \emph{ġiefe} `gift', \emph{tunge} `tongue' and
\emph{heorte} `heart'. The different quantity underlying \emph{nama}
`name' supplies a contrast, not an exception. They place unrounding
before shortening before final \emph{*r} and argue for a West Germanic
development partly from the retention of rounded vowels in Early Runic
(\citeproc{ref-RingeTaylor2014}{Ringe and Taylor 2014, 59--60}). These
are the authors' comparative and runic arguments, not additional
applications of the operation printed below or new independent readings
of inscriptions.

The quantity analysis is not uniform across handbooks. Hogg discusses
nasal-associated low-vowel outcomes and the survival and shortening of
long vowels, but includes the genitive singular among the endings
considered in his trimoric-vowel discussion and explicitly qualifies the
security of that classification {[}Hogg
(\citeproc{ref-HoggGrammar2011}{2011}), p.~64, §3.32; pp.~227--228,
§§6.26--28, especially §6.27 n.~2{]}. Ringe and Taylor instead
explicitly assign the relevant genitive a bimoric input and reconstruct
long \emph{*ā} followed by long \emph{*ǣ} on its way to final \emph{-e}
(\citeproc{ref-RingeTaylor2014}{Ringe and Taylor 2014, 59, 299--300}).
The length is consequently not an invention required only to obtain a
particular output, although the detailed prehistoric quantity remains a
matter for source criticism.

Campbell's discussion of fronting and restoration, Hogg's separate
historical chapter, and Ringe and Taylor's English restoration account
explain the surrounding English developments, rather than independently
establishing the West Germanic unrounding law
(\citeproc{ref-Campbell1959}{Campbell 1959, 52, 60--61}, §§131,
157--158; \citeproc{ref-HoggPhonology1992}{Hogg 1992b, 102, 105};
\citeproc{ref-RingeTaylor2014}{Ringe and Taylor 2014, 189--90}). The
grammar cited for long-vowel quantity above is Hogg's monograph in its
2011 pagination, not that 1992 chapter.

\subsection{\texorpdfstring{\CAPRRuleHeading{SC042. Unrounding of the surviving bimoric \emph{*ō}}{PWGmcSurvivingBimoricOUnrounding}}{}}\label{rule-PWGmcSurvivingBimoricOUnrounding}

\begin{ReaderFacingFoma}
define PWGmcSurvivingBimoricOUnrounding Ctx([
    {*ō} -> {*ā} || EnglishStarVocalic [EnglishStarConsonant | EnglishPalatalConsonant]+ _ .#.
]);
\end{ReaderFacingFoma}

This operation represents only the surviving oral-final fragment of the
broader development. Its preceding-nucleus and intervening-consonant
guard selects a final vowel in a polysyllabic shape; it does not
independently reconstruct historical stress. Before-final-\emph{*r}
histories and nasal-vowel histories have other executable treatments. At
this point the selected water and four inputs already contain \emph{*a}
before \emph{*r}, whereas the tongue and heart inputs retain the
distinct nasal token \emph{*ǭ}. Name retains a separate trimoric marker.
Their nonapplication here is therefore not a counterexample to the
broader comparative account.

The complete current population has one live application: \emph{ræste}
`rest', from the selected Proto-Germanic genitive singular
\emph{*rástōz}, not the citation nominative \emph{*rastō}. The selected
gift input is instead \emph{*gíftiz}, so the literature's genitive
\emph{*gebōz} must not silently replace it as another corpus witness.
Neither a lexical exception nor a grammatical condition is introduced.

After \hyperref[rule-EAFFinalZDeletion]{SC020 EAFFinalZDeletion} exposes
\emph{*rástō}, unrounding gives \emph{*rástā}.
\hyperref[rule-EAFBrightening]{SC043 EAFBrightening} fronts the root,
and the final-vowel support
\hyperref[rule-EAFBrighteningLongFinal]{SC108 EAFBrighteningLongFinal}
gives \emph{*ræstǣ}; subsequent shortening and merger yield \emph{ræste}
`rest'. The final support carries quantity in the model and is not
another independently dated sound law. The West Germanic long-vowel
reconstruction has a source basis, but this exact serialized chain
remains a computational representation of the selected history.

For the selected \emph{*rástōz}, moving unrounding before z-loss or
after final-vowel support gives \Pred{rasta}, rather than \emph{ræste}
`rest'. In the early displacement the vowel is still protected by
\emph{*z}; in the late displacement its unrounded output arrives too
late for final-vowel fronting. In both, a surviving back vowel permits
\hyperref[rule-OEARestoration]{SC046 OEARestoration} to retract the
root, and later shortening leaves final \emph{-a}. By contrast, placing
unrounding after ordinary stressed fronting but before final-vowel
support preserves all current outputs. Thus z-loss feeds the
exposed-final-vowel operation, which feeds the final support; it is not
the stressed fronting operation that needs the unrounded final vowel.

The historical West Germanic development precedes the English fronting
history in the adopted account. That does not establish an independently
witnessed local relation to ordinary stressed fronting: the single
selected rest derivation distinguishes exposure of a final vowel and the
use of its unrounded output by final-vowel support. The rule, its order
and the selected forms are retained; its historical interpretation is
reformulated to distinguish this narrow executable contribution from the
independently argued broader law.

\section{Technical support: retained unstressed
fronting}\label{technical-support-retained-unstressed-fronting}

\subsection{\texorpdfstring{\CAPRRuleHeading{SC107. Retained unstressed contribution}{EAFBrighteningUnstressed}}{}}\label{rule-EAFBrighteningUnstressed}

\noindent\begin{minipage}{\linewidth}
\begin{ReaderFacingFoma}
define EAFBrighteningUnstressed Ctx([
    {*a} -> {*æ} || _ [EnglishStarConsonant - EnglishStarNasal]
]);
\end{ReaderFacingFoma}
\end{minipage}

This earlier contribution is preserved without identifying it with the
entire unstressed history. The later syllabic rule also fronts surviving
unstressed vowels before heterosyllabic nasals; an exception before
every nasal would therefore be too broad
(\citeproc{ref-Campbell1959}{Campbell 1959, 140--41}).

\section{Ordinary English fronting}\label{ordinary-english-fronting}

\subsection{Historical discussion}\label{historical-discussion-38}

Anglo-Frisian Brightening or First Fronting turns low \emph{*a} into
fronted \emph{*æ}-type outcomes outside nasal environments. Later Old
English developments presuppose this fronted stage even where they
partly conceal it. Campbell gives the classical statement of the change,
Hogg supplies the standard modern labels, and Ringe and Taylor establish
its local chronology with breaking and restoration
(\citeproc{ref-Campbell1959}{Campbell 1959, 52}, §131;
\citeproc{ref-HoggPhonology1992}{Hogg 1992b, 102, 105};
\citeproc{ref-RingeTaylor2014}{Ringe and Taylor 2014, 157--58, 189--90};
\citeproc{ref-Fulk2018}{Fulk 2018, 73--74}, §§4.12--4.13).

Brightening creates the input to \hyperref[rule-OEBreaking]{SC044
OEBreaking}, while \hyperref[rule-OEARestoration]{SC046 OEARestoration}
later partly reverses its outcome before back vowels.

The traditional name does not itself establish one inherited event.
Under Campbell's premise that the daughter contractions precede ordinary
fronting, comparable English and Frisian outcomes require separate
daughter frontings in the strict tree. Chapter 3 develops that
conditional argument without removing the ancestral node
(\citeproc{ref-Campbell1939}{Campbell 1939, 90--91}). The model adopts
the completed ordinary English stressed component on the English
daughter under this conventional working account. This is not a claim
that every author proves contraction before ordinary fronting: Ringe and
Taylor explicitly question whether diphthong nuclei must behave like the
plain short vowel (\citeproc{ref-RingeTaylor2014}{Ringe and Taylor 2014,
170--75}). Earlier restricted ancestral fronting remains possible.

Fulk's surviving fronted \emph{*æi} and proposed paradigmatic levelling
provide a competing account, discussed in Chapter 3, rather than direct
confirmation of contraction before ordinary fronting
(\citeproc{ref-Fulk1998}{Fulk 1998, 143--45}). No shared fronting is
demonstrated merely by making that alternative compatible with the
daughter outcomes.

The separately retained unstressed and final-vowel components are not
dated by that argument. The first supplies an earlier nonnasal
contribution; \hyperref[rule-OEUnstressedFrontingEarly]{SC070
OEUnstressedFrontingEarly} separately implements the broader unstressed
domain, including the contrast between coda and heterosyllabic nasals
(\citeproc{ref-Campbell1959}{Campbell 1959, 140--41}, §§333--334). The
second carries the model's preserved final-vowel quantity into later
shortening and merger. Its long intermediate is a representation choice,
not independent evidence for a historical long-vowel fronting
(\citeproc{ref-RingeTaylor2014}{Ringe and Taylor 2014, 58--59,
299--300}).

\subsection{\texorpdfstring{\CAPRRuleHeading{SC043. Fronting of low \emph{*a} outside nasal environments}{EAFBrightening}}{}}\label{rule-EAFBrightening}

\noindent\begin{minipage}{\linewidth}
\begin{ReaderFacingFoma}
define EAFBrightening Ctx([
    EAFBrighteningStressed
]);
\end{ReaderFacingFoma}
\end{minipage}

\subsubsection{Stressed component}\label{stressed-component}

\noindent\begin{minipage}{\linewidth}
\begin{ReaderFacingFoma}
define EAFBrighteningStressed [
    {*á} -> {*æ} || _ [EnglishStarConsonant - EnglishStarNasal | .#.]
];
\end{ReaderFacingFoma}
\end{minipage}

Slay requires breaking to receive the fronted root vowel: delaying the
stressed member until after \hyperref[rule-OEBreaking]{SC044 OEBreaking}
gives \emph{sleaan | slēaan}, rather than OE \emph{slēan} `slay'. Rest
involves two different changes, root fronting and the separately
retained final-vowel representation. Its final-vowel dependency must not
be used to date this stressed rule.

\section{Technical support: retained final-vowel
representation}\label{technical-support-retained-final-vowel-representation}

\subsection{\texorpdfstring{\CAPRRuleHeading{SC108. Retained final-vowel representation}{EAFBrighteningLongFinal}}{}}\label{rule-EAFBrighteningLongFinal}

\begin{ReaderFacingFoma}
define EAFBrighteningLongFinal Ctx([
    {*ā} -> {*ǣ} || EnglishStarVocalic [EnglishStarConsonant | EnglishPalatalConsonant]+ _ .#.
]);
\end{ReaderFacingFoma}

In \emph{ræste} `rest', this helper receives the length-preserved
outcome of \hyperref[rule-PWGmcSurvivingBimoricOUnrounding]{SC042
PWGmcSurvivingBimoricOUnrounding}. The preceding-nucleus guard excludes
stressed monosyllabic \emph{hwā} `who'. Later shortening and merger
complete the ending; the carried long vowel does not establish a
separate dated long-a law (\citeproc{ref-RingeTaylor2014}{Ringe and
Taylor 2014, 58--59, 299--300}; \citeproc{ref-Campbell1959}{Campbell
1959, 49}, §125).

\section{Breaking}\label{breaking}

\subsection{Historical discussion}\label{historical-discussion-39}

Breaking creates \emph{eo}-type outputs before \emph{h}, \emph{rC}, and
\emph{lC}; velar-fricative palatalization then operates in that reshaped
environment. Campbell, Ringe and Taylor, and Fulk place breaking after
brightening. The following fricative palatalization is more narrowly
conditioned (\citeproc{ref-Campbell1959}{Campbell 1959, 54, 166}, §§139,
405--406; \citeproc{ref-RingeTaylor2014}{Ringe and Taylor 2014, 168--69,
213--14}, §§6.2.1--6.2.3, 6.4.1--6.4.2; \citeproc{ref-Fulk2018}{Fulk
2018, 73--74}, §4.13).

Breaking has the fuller handbook treatment. The \emph{feoh} `cattle' and
\emph{feohtan} `fight' type derivations preserve its velar-fricative
trigger; they do not subsequently undergo velar-fricative
palatalization.

Fulk cautions that breaking before \emph{h} does not necessarily date
the articulation of \emph{k/g}: the consonantal histories must be
distinguished. His English and later Frisian breaking processes are also
not one inherited event (\citeproc{ref-Fulk1998}{Fulk 1998, 145--46,
152}).

\subsection{\texorpdfstring{SC044. Breaking before \emph{h}, \emph{rC},
\emph{lC}, and conditioned \emph{w}
(\texttt{OEBreaking})}{SC044. Breaking before , , , and conditioned  (OEBreaking)}}\label{rule-OEBreaking}

\noindent\begin{minipage}{\linewidth}
\begin{ReaderFacingFoma}
define OEBreaking FootCtx(OEBreakingA
    .o. OEBreakingE
    .o. OEBreakingI);
\end{ReaderFacingFoma}
\end{minipage}

Short \emph{e} and \emph{i} also develop short \emph{eo} and \emph{io}
before singleton \emph{w}, except when a high front vowel or \emph{j}
follows it. This is the regular pre-ending history of \emph{cneowe}
`knee (dat.sg.)', not long-diphthong promotion
(\citeproc{ref-Luick1914}{Luick 1914-\/-1940, 139}, §134;
\citeproc{ref-HoggGrammar2011}{Hogg 2011, 86}, §5.22;
\citeproc{ref-RingeTaylor2014}{Ringe and Taylor 2014, 187--88}, §6.2.4).
The shared w conditioner makes that exclusion explicit:

\noindent\begin{minipage}{\linewidth}
\begin{ReaderFacingFoma}
define EnglishBreakingWContext [
    {*w} [[EnglishStarVocalic | EnglishStarConsonant]
        - [{*i} | {*ì} | {*í} | {*ī} | {*ḯ} | {*j}]] |
    {*w} .#.
];
\end{ReaderFacingFoma}
\end{minipage}

Root stress notation does not determine vowel length. The short knee
dative is the selected comparison; long endingless \emph{cnēo} `knee'
and the restored \emph{w} of \emph{cnēow} `knee' are separate paradigm
histories, discussed with \hyperref[rule-OEEwLongDiphthong]{SC033
OEEwLongDiphthong}.

Breaking must encounter the vowel created by brightening and must
precede the rule that would otherwise palatalize its velar trigger in
\emph{feoh} `cattle' and \emph{feohtan} `fight'. Before
\hyperref[rule-EAFBrightening]{SC043 EAFBrightening}, PGmc
\Recon{sláxaną} `slay' yields \emph{sleaan | slēaan} rather than
expected OE \emph{slēan} `slay'. After
\hyperref[rule-OEVelarFricativePalatalization]{SC045
OEVelarFricativePalatalization}, PGmc \Recon{féxu} `cattle' yields
\Pred{fehu} rather than expected OE \emph{feoh}, and PGmc
\Recon{féxtaną} `fight' yields \Pred{fehtan} rather than expected
\emph{feohtan}. The fronting relation feeds breaking. The fee/fight
relation instead protects breaking's velar trigger from premature
palatalization: those forms are displacement negatives, not live
\hyperref[rule-OEVelarFricativePalatalization]{SC045
OEVelarFricativePalatalization} applications.

\subsection{\texorpdfstring{SC044. Stress-sensitive short-i breaking
(\texttt{OEBreakingI})}{SC044. Stress-sensitive short-i breaking (OEBreakingI)}}\label{rule-OEBreakingI}

\begin{ReaderFacingFoma}
define OEBreakingI [
    {*i} -> {*io} || .#. OEAnyConsonant* _ EnglishBreakingRContext,
    {*i} -> {*io} || .#. OEAnyConsonant* _ EnglishBreakingHContext,
    {*i} -> {*io} || .#. OEAnyConsonant* _ EnglishBreakingWContext
] .o. [
    {*í} -> {*ío} || _ EnglishBreakingRContext,
    {*í} -> {*ío} || _ EnglishBreakingHContext,
    {*í} -> {*ío} || _ EnglishBreakingWContext,
    {*ì} -> {*io} || _ EnglishBreakingRContext,
    {*ì} -> {*io} || _ EnglishBreakingHContext,
    {*ì} -> {*io} || _ EnglishBreakingWContext
];
\end{ReaderFacingFoma}

Breaking is not an ordinary unstressed-vowel development. \emph{fulwiht}
`baptism' and \emph{nawiht} `nothing' supply independent negative
evidence; the latter has separately explained stressed byforms. A plain
vowel can nevertheless represent the stressed initial nucleus after an
earlier accent-consuming change. This component therefore admits that
nucleus, explicit primary \emph{í}, and preserved subordinate \emph{ì},
but excludes plain medial \emph{i}. Positive pretonic foot marking
retains its existing scope (\citeproc{ref-HoggGrammar2011}{Hogg 2011,
217--18}, §§6.7--6.8; \citeproc{ref-Campbell1959}{Campbell 1959,
54--56}, §§139--144; \citeproc{ref-RingeTaylor2014}{Ringe and Taylor
2014, 187--88}).

In the conditional \emph{seleweard} history the retained medial \emph{i}
causes mutation without breaking and later lowers to \emph{e}. This is a
general stress distinction, not suppression of a diphthong by lexical
identity. The second root's subordinate marker survives until its
phonological consumers; its technical decoding is not a date for
obligatory historical destressing (\citeproc{ref-Russom2001}{Russom
2001, 39--64}).

\section{Velar-fricative
palatalization}\label{velar-fricative-palatalization}

\subsection{\texorpdfstring{\CAPRRuleHeading{SC045. Palatalization of velar fricatives beside front vowels}{OEVelarFricativePalatalization}}{}}\label{rule-OEVelarFricativePalatalization}

\noindent\begin{minipage}{\linewidth}
\begin{ReaderFacingFoma}
define OEVelarFricativePalatalization Ctx([
    {*x} -> {*ç} || _ EnglishStarFrontVowel,
    {*x} -> {*ç} || EnglishStarFrontVowel _,
    {*x} -> {*ç} || _ {*j}
]
    .o. EnglishStarAlphabet*);
\end{ReaderFacingFoma}
\end{minipage}

The live population concerns voiceless x, as in \emph{hēafod} `head' and
\emph{heofon} `heaven', not voiced-fricative merger. In \emph{feoh}
`cattle' and \emph{feohtan} `fight', breaking removes the original
front-vowel context: they do not change at this rule in the live
derivation. Moving the rule before \hyperref[rule-OEBreaking]{SC044
OEBreaking} instead consumes its velar trigger, yielding \Pred{fehu} and
\Pred{fehtan}. This is counterbleeding protection, not direct feeding.

The distant \emph{six} `six' displacement test supplies only a broader
constraint: after \hyperref[rule-OEWsPalatalUmlaut]{SC060
OEWsPalatalUmlaut}, PGmc \Recon{séxs} `six' yields \Pred{sihs} rather
than expected OE \emph{six}. Neither test dates voiced g. Its
articulation and disputed merger are separately treated under
\hyperref[rule-OEVelarPalatalization]{SC052 OEVelarPalatalization} and
\hyperref[rule-OEPalatalFricativeMerger]{SC109 OEPalatalFricativeMerger}
(\citeproc{ref-RingeTaylor2014}{Ringe and Taylor 2014, 203--4};
\citeproc{ref-Fulk2018}{Fulk 2018, 130--32}).

\section{A-restoration}\label{a-restoration}

\subsection{Historical discussion of
A-restoration}\label{historical-discussion-of-a-restoration}

Campbell's restoration of \emph{a} before following back vowels and
Ringe and Taylor's later retraction describe the same post-brightening
development (\citeproc{ref-Campbell1959}{Campbell 1959, 60--61},
§§157--159; \citeproc{ref-RingeTaylor2014}{Ringe and Taylor 2014,
189--90}, §6.3.1; \citeproc{ref-Fulk2018}{Fulk 2018, 74}, §4.13). Some
outcomes of Anglo-Frisian fronting survive only in environments where
restoration does not return them to back \emph{a}.

\hyperref[rule-OEARestoration]{SC046 OEARestoration} has firmer handbook
support than the two following nasal rules.

\subsection{\texorpdfstring{\CAPRRuleHeading{SC046. Restoration of \emph{*a} before following back vowels}{OEARestoration}}{}}\label{rule-OEARestoration}

\noindent\begin{minipage}{\linewidth}
\begin{ReaderFacingFoma}
define OEARestoration FootCtx((
    {*æ} -> {*a} || _
        OEARestorationIntervening OEARestorationTriggerVowel
        - OEARestorationIntervening OEARestorationWeakTailVowel
));
\end{ReaderFacingFoma}
\end{minipage}

Restoration must receive fronted \emph{*æ} and return \emph{*a} before
the nasal-tail changes. Before \hyperref[rule-EAFBrightening]{SC043
EAFBrightening}, PGmc \Recon{bákaną} `bake' yields \Pred{bæcan} rather
than expected OE \emph{bacan} `bake', and PGmc \Recon{fáraną} `fare'
yields \Pred{færan} rather than expected \emph{faran} `fare'. After
\hyperref[rule-OESecondaryNasalization]{SC048 OESecondaryNasalization},
\Recon{bákaną} again yields \Pred{bæcan} instead of \emph{bacan}, while
PGmc \Recon{wádaną} `wade' yields \Pred{wædan} instead of \emph{wadan}
`wade'. These independent witness pairs place restoration after
brightening and before secondary nasalization.

\section{Heavy-syllable nasal
apocope}\label{heavy-syllable-nasal-apocope}

\subsection{Historical discussion of heavy-syllable nasal loss and
secondary
nasalization}\label{historical-discussion-of-heavy-syllable-nasal-loss-and-secondary-nasalization}

Heavy-syllable nasal apocope removes the final nasalized vowel;
secondary nasalization then marks the preceding \emph{a} before final
\emph{n}. The handbooks do not isolate both developments under equally
prominent labels. Campbell describes later nasal loss and the
back-mutation environment; Ringe and Taylor provide the later relation
to back mutation (\citeproc{ref-Campbell1959}{Campbell 1959, 86, 166},
§§205--206, 403; \citeproc{ref-RingeTaylor2014}{Ringe and Taylor 2014,
319}, §6.9.4).

The reciprocal failure set fixes the order: apocope removes the ending
before secondary nasalization acts on the remaining structure.
Restoration receives the fuller historical treatment in the handbooks.

\subsection{\texorpdfstring{\CAPRRuleHeading{SC047. Heavy-syllable nasal apocope of final \emph{*ą}}{OEHeavySyllableNasalApocope}}{}}\label{rule-OEHeavySyllableNasalApocope}

\noindent\begin{minipage}{\linewidth}
\begin{ReaderFacingFoma}
define OEHeavySyllableNasalApocope Ctx([
    {*ą} -> 0 || OEAnyConsonant _ .#.
]);
\end{ReaderFacingFoma}
\end{minipage}

The old lower-bound test depended on the superseded straw proxy: moving
apocope before \hyperref[rule-OEAwLongDiphthong]{SC034
OEAwLongDiphthong} made PGmc \Recon{stráwą} `straw' yield \Pred{stræw}
rather than attested OE \emph{strēaw} `straw'. The adopted regular
endingless branch now reaches reconstructed \emph{*strēa} through
\hyperref[rule-PWGmcEndinglessAw]{SC111 PWGmcEndinglessAw}, before this
apocope. The attested target retains analogically restored \emph{w} and
explicitly mismatches the regular output. The old failure therefore does
not independently place nasal apocope after a universal singleton
long-diphthong law (\citeproc{ref-Campbell1959}{Campbell 1959, 46--47,
232--33}; \citeproc{ref-RingeTaylor2014}{Ringe and Taylor 2014, 172--73,
387}).

The separate recorded upper-bound test is unchanged: after
\hyperref[rule-OESecondaryNasalization]{SC048 OESecondaryNasalization},
PGmc \Recon{bákaną} `bake' yields \Pred{bacen} rather than expected OE
\emph{bacan}, and PGmc \Recon{bíndaną} `bind' yields \Pred{binden}
rather than expected \emph{bindan}, alongside a broad \emph{-en} failure
set. That secondary-nasalization dependency is not newly adjudicated
here.

\section{Secondary nasalization}\label{secondary-nasalization}

\subsection{\texorpdfstring{\CAPRRuleHeading{SC048. Secondary nasalization before final \emph{*n}}{OESecondaryNasalization}}{}}\label{rule-OESecondaryNasalization}

\noindent\begin{minipage}{\linewidth}
\begin{ReaderFacingFoma}
define OESecondaryNasalization Ctx([
    {*a} -> {*ą} || _ {*n} .#.
]);
\end{ReaderFacingFoma}
\end{minipage}

The broad \emph{-an}/\emph{-en} split fixes the lower boundary of final
\emph{*a} nasalization before \emph{n}. Before
\hyperref[rule-OEHeavySyllableNasalApocope]{SC047
OEHeavySyllableNasalApocope}, PGmc \Recon{bákaną} `bake' yields
\Pred{bacen} rather than expected OE \emph{bacan} `bake', and PGmc
\Recon{bíndaną} `bind' yields \Pred{binden} rather than expected
\emph{bindan} `bind'. The upper boundary comes from back mutation. After
\hyperref[rule-OEBackMutation]{SC059 OEBackMutation}, PGmc
\Recon{stélaną} `steal' yields \Pred{steolan} rather than expected OE
\emph{stelan} `steal', and PGmc \Recon{wébaną} `weave' yields
\Pred{weofan} rather than expected \emph{wefan} `weave'. Reciprocal
nasal-tail failures place secondary nasalization after apocope, and the
later mutation witnesses place it before back mutation;
\hyperref[rule-OEARestoration]{SC046 OEARestoration} retains the
clearest independent historical support.

\section{Sievers-law syncope}\label{sievers-law-syncope}

\subsection{Historical discussion}\label{historical-discussion-40}

Sievers' Law concerns a prosodic and morphological adjustment in heavy
stems. It is a distributional rule distinct from b-allophony
(\hyperref[rule-PGmcBAllophony]{SC049 PGmcBAllophony}). Adamczyk treats
the Old English reflexes of the law as historical evidence from weak
verbs and related formations (\citeproc{ref-Adamczyk2001}{Adamczyk 2001,
61--72}). Fulk gives the compact comparative summary through familiar
forms such as \emph{biddan} `ask', \emph{sellan} `give', and
\emph{nerian} `save' (\citeproc{ref-Fulk2018}{Fulk 2018, 127}, §6.15).

Sievers-law syncope is narrow in scope, but its relation to the
following palatalization is lexically secure. Its earlier limit is less
sharply defined than that of the preceding allophony rule.

\subsection{\texorpdfstring{SC050. Sievers-law syncope
(\texttt{SieversLawSyncope})}{SC050. Sievers-law syncope (SieversLawSyncope)}}\label{rule-SieversLawSyncope}

\noindent\begin{minipage}{\linewidth}
\begin{ReaderFacingFoma}
define SieversLawSyncope Ctx([
    {*i} -> 0 || [EnglishStarConsonant | EnglishPalatalConsonant] _ {*j}
]);
\end{ReaderFacingFoma}
\end{minipage}

The Sievers-law reduction \emph{*-CijV-*} \textgreater{} \emph{*-CjV-*},
including loss of \emph{*i} before \emph{*j}, must precede
palatalization. If \hyperref[rule-SieversLawSyncope]{SC050
SieversLawSyncope} follows \hyperref[rule-OEVelarPalatalization]{SC052
OEVelarPalatalization}, PGmc \Recon{strákkijaną} `stretch' yields
\Pred{strecċan} rather than expected OE \emph{streċċan} `stretch';
earlier placement creates no comparably precise error. The single
cluster witness therefore places syncope before velar palatalization.

\section{\texorpdfstring{Palatalization of \emph{*sk} to
\emph{*sc}}{Palatalization of  to }}\label{palatalization-of-to}

\subsection{Historical discussion}\label{historical-discussion-41}

The palatalization of \emph{*sk} to Old English \emph{*sc} is one of the
recognizable early cluster changes in the larger palatalization zone.
Campbell distinguishes the cluster from plain velars when he remarks
that \emph{*sk} is especially prone to palatalization and assibilation
(\citeproc{ref-Campbell1959}{Campbell 1959, 278}, §440). Barrack treats
the cluster in the same zone as his 440 series: 440a summarizes the
developments (\citeproc{ref-Barrack1975}{Barrack 1975, 112--13}) and
440c states the palatalization of \emph{*sk} itself
(\citeproc{ref-Barrack1975}{Barrack 1975, 113}), kept beside, but
distinct from, plain velar palatalization (his rule 170a)
(\citeproc{ref-Barrack1975}{Barrack 1975, 57}), in agreement with the
structural placement adopted here. He also registers the typologically
separate Late West Saxon metathesis of the cluster in
\emph{fiscas}/\emph{fixas} `fishes' as rule 440b
(\citeproc{ref-Barrack1975}{Barrack 1975, 113}), a written-period
development that presupposes the earlier palatalization boundary rather
than contradicting it. Hogg gives the same change a clearer structural
place by treating \emph{*sk} beside the palatalization of plain velars
and before the later West Saxon diphthongal developments
(\citeproc{ref-HoggPhonology1992}{Hogg 1992b, 106--8, 113}). Ringe and
Taylor make the same sequence explicit when they distinguish the earlier
palatalization of velars and \emph{*sk} from the later diphthongization
after already palatal consonants (\citeproc{ref-RingeTaylor2014}{Ringe
and Taylor 2014, 213--16}, §§6.4.1, 6.5.1).

Luick places the cluster change within a broader early movement toward
palatal articulation, while still allowing later vowel consequences to
form a different chapter of the history (\citeproc{ref-Luick1914}{Luick
1914-\/-1940, 157}, §168). Fulk's summary is the most concise warning
against overextension: Old English \emph{*sc} is palatal except in the
well-known back-vowel environments that preserve harder outcomes
(\citeproc{ref-Fulk2018}{Fulk 2018, 28}). The result is a historically
clear rule, but not an identity between the cluster change and the later
umlautal developments.

\subsection{\texorpdfstring{SC051. Palatalization of \emph{*sk} to
\emph{*sc}
(\texttt{OESkPalatalization})}{SC051. Palatalization of  to  (OESkPalatalization)}}\label{rule-OESkPalatalization}

\begin{ReaderFacingFoma}
define OESkPalatalization Ctx([
    {*s} {*k} -> {*ʃ} || .#. _
] .o. [
    {*s} {*k} -> {*ʃ} || EnglishStarFrontVowel _ (EnglishStarConsonant | .#.)
] .o. [
    {*s} {*k} -> {*ʃ} || (EnglishStarConsonant | .#.) _ EnglishStarFrontVowel
] .o. [
    {*s} {*k} -> {*ʃ} || _ {*j}
] .o. [
    {*s} {*k} -> {*ʃ} || {*j} _
]);
\end{ReaderFacingFoma}

The non-fronted vowels of \emph{flasce} `flask' and \emph{wascan} `wash'
fix the lower boundary of \emph{*sk} \textgreater{} \emph{*sc}. Before
\hyperref[rule-OEARestoration]{SC046 OEARestoration}, the forms are
fronted too soon, yielding \emph{flæsce} `flask' and \emph{wæscan}
`wash' rather than expected OE \emph{flasce} and \emph{wascan}. This
places \hyperref[rule-OESkPalatalization]{SC051 OESkPalatalization}
after restoration.

Five witnesses establish the upper boundary collectively. The palatal
cluster must already underlie \emph{sċeaft} `shaft', \emph{sċēar}
`shear', \emph{sċēaþ} `sheath', \emph{sċēap} `sheep', and \emph{sċield}
`shield' before \hyperref[rule-OEWsPalatalDiphthongization]{SC056
OEWsPalatalDiphthongization}. The \emph{*sċea-* 'sea'}/\emph{*sċie-*}
set therefore places cluster palatalization before the West Saxon vowel
change. The cluster change occupies the same palatalization zone as
\hyperref[rule-OEVelarPalatalization]{SC052 OEVelarPalatalization} while
remaining distinct from plain-velar palatalization and the later vowel
changes.

\section{Class-distinct velar palatalization and fricative
merger}\label{class-distinct-velar-palatalization-and-fricative-merger}

\subsection{Historical discussion}\label{historical-discussion-42}

Singleton g, including initial g, was a fricative when palatalization
occurred; gg and postnasal g were stops. Palatal articulation,
productive cutoff, merger and later affrication are therefore separate
questions (\citeproc{ref-Fulk2018}{Fulk 2018, 130--32};
\citeproc{ref-Laker2007}{Laker 2007, 166--68, 175--84}). The dotted
spelling in \emph{dæġ} `day' cannot establish an early affricate or
identify every palatal consonant with inherited j.

The adopted working account keeps a palatal fricative ʝ distinct from j
during mutation, then merges it with j. Ringe and Taylor explicitly
propose this chronology (\citeproc{ref-RingeTaylor2014}{Ringe and Taylor
2014, 203--4}). It is not unanimous: Hogg objects to the phonetic
plausibility of delayed merger and leaves the key/day contradiction
unresolved (\citeproc{ref-Hogg1979}{Hogg 1979, 102--11}). We prefer the
regular class-distinct account because it explains the contrasting
mutation triggers without lexical or grammatical conditions. The
objection remains substantive, not silently superseded by a matching
final spelling.

Fulk's earlier comparison treats reversible palatal articulation and
later assibilation separately, but does not establish the particular
delayed singleton-fricative merger used here
(\citeproc{ref-Fulk1998}{Fulk 1998, 145--49}). That merger remains the
attributed Ringe--Taylor working account, with Hogg's objection
retained.

Eligibility also requires attribution. Ringe and Taylor's medial-g table
is narrower than Fulk's before-front-vowel description
(\citeproc{ref-RingeTaylor2014}{Ringe and Taylor 2014, 203--4};
\citeproc{ref-Fulk2018}{Fulk 2018, 130--32}). The present fricative
implementation retains its broader domain; this merger decision is not a
covert adoption of the narrower table.

Barrack's cascade distributes the velar material over several rules:
palatalization proper as rule 170a (\citeproc{ref-Barrack1975}{Barrack
1975, 57}), the palatal/velar alternations of the written period as rule
428, with \emph{streccan} `stretch' among the cited forms
(\citeproc{ref-Barrack1975}{Barrack 1975, 109}), and the back-and-forth
adjustments between palatals and velars as rules 432--435, including the
\emph{hwilic} `which' type and \emph{ortgeard}/\emph{feccan}
(\citeproc{ref-Barrack1975}{Barrack 1975, 110--12}). His decomposition
is finer than the class-distinct account adopted here and is keyed to
Campbell's paragraphs rather than to the fricative/stop class
distinction; it corroborates the breadth of the palatalization zone
without bearing on the disputed merger chronology, since his orderings
are formal, not chronological (\citeproc{ref-Barrack1975}{Barrack 1975,
34}).

\subsection{\texorpdfstring{SC052. Singleton-fricative articulation
(\texttt{OEGFricativePalatal})}{SC052. Singleton-fricative articulation (OEGFricativePalatal)}}\label{rule-OEGFricativePalatal}

\begin{ReaderFacingFoma}
define OEGFricativePalatal [
    [{*g}|{*ɣ}] -> {*ʝ} ||
        [.#.|[EnglishStarAlphabet - [{*g}|{*n}|{*ŋ}]]]
        _ EnglishStarFrontVowel,
    [{*g}|{*ɣ}] -> {*ʝ} || EnglishStarFrontVowel _ .#.,
    [{*g}|{*ɣ}] -> {*ʝ} || EnglishStarFrontVowel
        _ [EnglishStarConsonant - [{*j}|{*g}]]
] .o. [
    [{*g}|{*ɣ}] -> {*ʝ} ||
        [.#.|[EnglishStarAlphabet - [{*g}|{*n}|{*ŋ}]]] _ {*j}
];
\end{ReaderFacingFoma}

The guards keep gg/ng in their stop paths. The value ʝ is a real
phonetic class, not a tag recording a consonant's ancestry. It is
transparent to a following mutation trigger, but is not itself j at that
date. The source comparison normalizes older fricative notation to
modern ɣ/ʝ, never to the digit 3 (\citeproc{ref-Hogg1979}{Hogg 1979,
105}; \citeproc{ref-RingeTaylor2014}{Ringe and Taylor 2014, 204}).

\subsection{\texorpdfstring{\CAPRRuleHeading{SC052. K articulation and eventual reflex}{OEVelarPalatalizationKFront}}{}}\label{rule-OEVelarPalatalizationKFront}

\begin{ReaderFacingFoma}
define OEVelarPalatalizationKFront [
    {*k} -> {*ʧ} || .#. _ EnglishStarFrontVowel,
    {*k} -> {*ʧ} || _ [{*i} | {*ī}],
    {*k} -> {*ʧ} || _ {*ḯ},
    {*k} -> {*ʧ} || [{*i} | {*ī}] _ EnglishStarFrontVowel,
    {*k} -> {*ʧ} || {*ḯ} _ EnglishStarFrontVowel,
    {*k} -> {*ʧ} || [{*i} | {*ī}] _ .#.,
    {*k} -> {*ʧ} || {*ḯ} _ .#.
] .o. [
    {*k} {*k} -> {*ʧ} {*ʧ} || _ {*j}
] .o. [
    {*k} -> {*ʧ} || _ {*j}
];
\end{ReaderFacingFoma}

Here ʧ is a telescoped eventual-reflex proxy, not a claim that early
palatal stops were already affricates. Ringe and Taylor explicitly
separate the initial palatal stop from subsequent affrication and its
syncope dependencies (\citeproc{ref-RingeTaylor2014}{Ringe and Taylor
2014, 203--4}). The \emph{weccan} `wake', \emph{licgan} `lie' and
\emph{lecgan} `lay' examples concern inherited j-clusters, not
independent proof of every plain-velar conditioner
(\citeproc{ref-RingeTaylor2014}{Ringe and Taylor 2014, 213--14}).

The deterministic kk-before-j path protects \emph{streċċan} `stretch'.
The original-front versus secondary-front distinction is more directly
tested by unrounded key than by rounded \emph{cȳ} `cow'; rounding alone
can give only a cutoff before later unrounding
(\citeproc{ref-Hogg1979}{Hogg 1979, 100--105};
\citeproc{ref-Laker2007}{Laker 2007, 167--68}).

\subsection{\texorpdfstring{SC052. Retained stop paths
(\texttt{OEVelarPalatalizationStops})}{SC052. Retained stop paths (OEVelarPalatalizationStops)}}\label{rule-OEVelarPalatalizationStops}

\begin{ReaderFacingFoma}
define OEVelarPalatalizationStops [
    OEVelarPalatalizationKFront
] .o. [
    {*g} -> {*ʤ} || _ EnglishStarFrontVowel,
    {*g} -> {*ʤ} || EnglishStarFrontVowel _ .#.,
    {*g} -> {*ʤ} || EnglishStarFrontVowel _ EnglishStarFrontVowel,
    {*g} -> {*ʤ} || EnglishStarFrontVowel _ [EnglishStarConsonant - {*j}],,
    {*g} {*g} -> {*ʤ} {*ʤ} || _ {*j}
] .o. [
    {*g} -> {*ʤ} || _ {*j}
];
\end{ReaderFacingFoma}

Singleton fricatives have already left this domain. The remaining gg/ng
paths use ʤ as an eventual-reflex proxy; the initial singleton fricative
is no longer represented as an early affricate. \emph{Wicg} `horse' is a
geminate-class illustration, \emph{senġan} `singe' a postnasal control,
and \emph{lungen} `lung' the unchanged back-vowel negative
(\citeproc{ref-RingeTaylor2014}{Ringe and Taylor 2014, 203--4, 213--14};
\citeproc{ref-Fulk2018}{Fulk 2018, 131--32}).

The double comma separates the geminate-before-\emph{j} rule's context
from the singleton rules' front-vowel contexts. It corrects rewrite-rule
scoping, not the proposed historical environment, date or
consonant-class distinction. The separate geminate rule must replace the
whole cluster, not permit a competing singleton covering.

\subsection{\texorpdfstring{SC052. Combined articulation
(\texttt{OEVelarPalatalization})}{SC052. Combined articulation (OEVelarPalatalization)}}\label{rule-OEVelarPalatalization}

\noindent\begin{minipage}{\linewidth}
\begin{ReaderFacingFoma}
define OEVelarPalatalization Ctx([
    OEGFricativePalatal .o. OEVelarPalatalizationStops
]);
\end{ReaderFacingFoma}
\end{minipage}

The combined operation does not give every class one sharply dated
palatalization/affrication/merger event. It supplies the class
distinctions needed by ordinary palatal diphthongization and mutation.
The common-stem onset remains a separate question
(\citeproc{ref-Luick1914}{Luick 1914-\/-1940, 835--41};
\citeproc{ref-Fulk2018}{Fulk 2018, 130--32}).

\subsubsection{Key and day at the mutation
checkpoint}\label{key-and-day-at-the-mutation-checkpoint}

Hogg's oblique comparison starts here from pre-palatal, pre-OE
\emph{*kājæ} and \emph{*dæɣæ}; neither is silently fed through a PGmc
prefix (\citeproc{ref-Hogg1979}{Hogg 1979, 105}). The source target is
\emph{cǣġe} `key', with velar initial k, versus \emph{dæġe} `day'.

{\def\LTcaptype{none} % do not increment counter
\begin{longtable}[]{@{}
  >{\raggedright\arraybackslash}p{(\linewidth - 4\tabcolsep) * \real{0.3333}}
  >{\raggedright\arraybackslash}p{(\linewidth - 4\tabcolsep) * \real{0.3333}}
  >{\raggedright\arraybackslash}p{(\linewidth - 4\tabcolsep) * \real{0.3333}}@{}}
\toprule\noalign{}
\begin{minipage}[b]{\linewidth}\raggedright
Checkpoint
\end{minipage} & \begin{minipage}[b]{\linewidth}\raggedright
Key
\end{minipage} & \begin{minipage}[b]{\linewidth}\raggedright
Day under the adopted account
\end{minipage} \\
\midrule\noalign{}
\endhead
\bottomrule\noalign{}
\endlastfoot
Before palatalization & \emph{*kājæ} & \emph{*dæɣæ} \\
Before mutation & Inherited j remains & Palatal ʝ is distinct from j \\
After mutation & \emph{*kǣjæ}, without new initial palatalization &
\emph{*dæʝæ}, without æ-raising \\
After merger and native realization & \emph{cǣġe} & \emph{dæġe} \\
\end{longtable}
}

Premature merger instead predicts pre-OE \Pred{dejæ} under these
premises. The experiment tests that source-based contrast; the existing
day nominative is not the oblique discriminator
(\citeproc{ref-Hogg1979}{Hogg 1979, 105--10};
\citeproc{ref-RingeTaylor2014}{Ringe and Taylor 2014, 204}).

The complete key suffix additionally requires
\hyperref[rule-OEIntervocalicJVocalization]{SC082
OEIntervocalicJVocalization} to retain j after its non-high long front
vowel. A correct mutation checkpoint followed by model-only \Pred{cǣie}
would not constitute a correct native derivation
(\citeproc{ref-Hogg1979}{Hogg 1979, 105}).

\section{Pretonic front-vowel
raising}\label{pretonic-front-vowel-raising}

\subsection{Inheritance and
accentuation}\label{inheritance-and-accentuation}

Old English \emph{ġelīefan} `believe' continues a formation with
inherited \emph{*ga-}. Gothic \emph{ga-} preserves the older vowel,
while Old Saxon and Old High German \emph{gi-} and Old English
\emph{ge-} have reduced reflexes. Ringe and Taylor reconstruct
\emph{*ga-} for Proto-Germanic and Proto-West-Germanic; similar reduced
forms do not establish a single shared reduction event. Continental
\emph{gasokun} on the Pforzen buckle, externally dated 585--610,
supplies neither an English palatal trigger nor a date for English
raising (\citeproc{ref-RingeTaylor2014}{Ringe and Taylor 2014, 206};
\citeproc{ref-RuneSPforzen}{Runische Schriftlichkeit in den germanischen
Sprachen (RuneS), n.d.}, find 201, inscription and external dating;
\citeproc{ref-Berg2014}{Berg 2014, 60--61}).

The inherited accent is disputed. Bammesberger distinguishes verbal
constructions from genuine compounds with nominal second members,
reconstructing initial stress for the latter, not for every nominal
formation. His strongest phonological argument uses the dental contrast
between Gothic \emph{gabaurþs} `birth' and West Germanic cognates:
starting from a voiced dental, Thurneysen's Law would require an
unstressed root and hence a stressed prefix. His account also requires
particular reconstructions of the formation and inherited accent.
Quinlin instead analyses inherited \emph{*ga-} as a clitic without
primary word stress and contests those premises. An alternative account
of inherited dental alternation further prevents the Gothic consonant
from uniquely determining prefix stress. These are competing historical
arguments, not observations of Old English metrical stress
(\citeproc{ref-Bammesberger1981}{Bammesberger 1981, 377--82, 384},
especially p.~381 and nn.~9--10; \citeproc{ref-Quinlin1991}{Quinlin
1991, 150--57}; \citeproc{ref-Ringe2017}{Ringe 2017, 305--6}).

The often-cited \emph{ġeatwe} `equipment' is not a secure exception with
accented inherited \emph{*ga-}. The traditional prefixed analysis
connects Old Norse \emph{gǫtvar} `equipment'; Bammesberger retains it.
Quinlin separates \emph{ġeatwe} from \emph{getāwe} `deeds' and proposes
an unprefixed stem related to Gothic \emph{gatwō} `street'. That
alternative requires a semantic bridge which is not independently
proved. Campbell himself qualifies the relationship of the associated
forms, and Orel retains the equipment word under a different
reconstructed base of unknown origin. The disagreement remains
unresolved: neither prefixed category membership nor exceptional
inherited accent can be assumed. Accordingly this word cannot securely
date the prefix's reduction relative to palatal diphthongization
(\citeproc{ref-Campbell1959}{Campbell 1959, 31}, §74 n.~5, and p.~149,
§358 n.~6; \citeproc{ref-Bammesberger1981}{Bammesberger 1981, 384}
n.~17; \citeproc{ref-Quinlin1991}{Quinlin 1991, 155--56};
\citeproc{ref-Orel2003}{Orel 2003, 403}).

Ordinary Old English \emph{ge-} and \emph{be-} are weak, with the
following root independently stressable. Alliteration and metrical
analysis support this distinction; they do not prove an identical
accentual system at every prehistoric stage. Grammatical structure can
help reconstruct stress without itself conditioning a sound law
(\citeproc{ref-HoggPhonology1992}{Hogg 1992b, 99--100};
\citeproc{ref-Minkova2008}{Minkova 2008, 22--23};
\citeproc{ref-MinkovaZhou2022}{Minkova and Zhou 2022, 535--36}).

\subsection{Manuscript distribution}\label{manuscript-distribution}

Pheifer's Épinal--Erfurt edition provides the central aggregate
evidence. His initial-prefix counts are:

{\def\LTcaptype{none} % do not increment counter
\begin{longtable}[]{@{}lrrrr@{}}
\toprule\noalign{}
Witness & \emph{gi-} & \emph{ge-} & Total \emph{gi-/ge-} & \emph{bi-}
with \emph{i} \\
\midrule\noalign{}
\endhead
\bottomrule\noalign{}
\endlastfoot
Épinal & 40 & 13 & 53 & 11 \\
Erfurt & 31 & 34 & 65 & 11 \\
\end{longtable}
}

The proportions of \emph{gi-} are 75.47\% and 47.69\%, respectively.
These are descriptive proportions in related transmitted texts, not
independent samples measuring the rate of a sound change. The
denominators differ, and the edition does not enumerate the complete
token lists or explicitly state how every corrected prefix enters the
totals. Its separate medial/final-vowel exclusions cannot simply be
transferred to this sample (\citeproc{ref-Pheifer1974}{Pheifer 1974,
lxxv}, §63).

The larger distribution complements, rather than replaces, the paired
readings. Nine Épinal--Erfurt alignments identified by Toller include
differences in both directions; their corresponding Corpus readings have
\emph{ge-}. Corpus nevertheless also has \emph{gibrec} `phlegm',
glossing Latin \emph{umecta}. A written correction from \emph{i} to
\emph{e} does not itself date phonological lowering. These readings
overlap the larger edition's material and must not be counted again as
independent confirmations (\citeproc{ref-Toller1921}{Toller 1921, 284,
s.v}.~ge-; \citeproc{ref-Sweet1885}{Sweet 1885, 48--49, 107, 602}).

Difficult forms require lexical and editorial discrimination. Épinal
\emph{gybyrdid}, beside Erfurt \emph{gebyrdid}, means `striped(?)' in
Pheifer's interpretation, not `born'; it cannot witness the accent of
\emph{ġebyrd} `birth'. Sweet and Pheifer both print Épinal
\emph{gihiodum}, beside Erfurt \emph{gaeadun} and Corpus \emph{geeodun}.
The manuscript confirms \emph{h}; Pheifer interprets it as a possible
hiatus marker. Pheifer also compares Épinal \emph{saegesetu} with Erfurt
\emph{saegaesetu} `maritime regions', where the disputed vowel is
internal. Neither exceptional \emph{ga-} spelling establishes by itself
an unraised short vowel after an independently established palatal
fricative: early prosody, vowel value and transmission remain material
questions (\citeproc{ref-Pheifer1974}{Pheifer 1974, xxiii, lxxv,
lxxxiii}, §§4, 63, 72, and pp.~6, 13, 39, 74, 108;
\citeproc{ref-Sweet1885}{Sweet 1885, 27--28, 40, 50}).

The revised critical edition confirms the \emph{h} in \emph{gi hiodum}
and distinguishes manuscript readings from its reconstructed archetype.
It also prints Corpus \emph{saegeseotu}, rather than Pheifer's
\emph{saegesetu}: real editorial differences must remain distinct from
modern transcription errors. Neither this root-vowel difference nor
editorial word division decides the prefix's early phonetic history
(\citeproc{ref-HerrenEtAl2026}{Herren et al. 2026}, letter A, PDF p.~20,
Ep A192, and letter P, PDF p.~3, Cp P599).

The dates of manuscripts are not the dates of all their linguistic
material. Épinal is placed around 700. Herren and Sauer date the
continental Erfurt copy to the second or third decade of the ninth
century; Pheifer argues for a common ancestor with Épinal and an
intervening English copying layer. Corpus is a related expanded
recension, not an unrelated later sample; ninth-century dating does not
give every gloss a new phonological date. Dialect, exemplar, English
revision and continental copying must remain distinct
(\citeproc{ref-LucasLucas2012}{Lucas and Lucas 2012, 67};
\citeproc{ref-Lucas2016}{Lucas 2016, 11};
\citeproc{ref-HerrenSauer2016}{Herren and Sauer 2016, 129--32};
\citeproc{ref-Pheifer1974}{Pheifer 1974, xli}, §24;
\citeproc{ref-Pheifer1987}{Pheifer 1987, 34--37}).

Seiler's analysis of inherited spelling systems and source batches
reinforces this caution: an aggregate need not represent one scribal
phonology. Her evidence chiefly concerns wyn and thorn, not a new
\emph{gi-/ge-} census or a reclassification of Pheifer's tokens
(\citeproc{ref-Seiler2017}{Seiler 2017, 159--64, 167--68}).

Seiler's direct comparison further distinguishes independent copying
from independent linguistic evidence. Épinal and Erfurt descend from a
shared textual ancestor, with an English intermediary in Erfurt's
transmission; the continental copy can preserve inherited orthography as
well as revise it. Positional spelling conventions and source batches
therefore need not express the phonology of one copyist
(\citeproc{ref-Seiler2019}{Seiler 2019, 527--28, 535--37}). Her analysis
chiefly concerns the representation of /w/ and /þ/, not prefix vowels.
It neither supplies a new \emph{gi-/ge-} census nor corrects Pheifer's
figures or the individually checked readings
(\citeproc{ref-Seiler2019}{Seiler 2019, 529--31, 537--38}). Erfurt's
independent act of copying must not be counted as an independent origin
for every spelling.

Early Northumbrian \emph{gistiga} `climb into/onto' and later
Northumbrian \emph{giher} `hear' and \emph{gitrivadon} `we trust'
prevent a universal early-\emph{gi-}/late-\emph{ge-} spelling calendar.
Samuels's Lindisfarne--Rushworth comparison likewise documents late
\emph{gi-} forms, but its principal contrast concerns the presence or
absence of the whole prefix, not a vowel-frequency census or proof of
influence from \emph{bi-}. Hogg explicitly allows spelling confusion
and, in late Northumbrian, a raised allophone. The attested spellings
are more secure than a unique phonemic interpretation of each vowel
(\citeproc{ref-Sweet1885}{Sweet 1885, 125};
\citeproc{ref-Sweet1871}{1871, 3};
\citeproc{ref-Stevenson1840}{Stevenson 1840, 7}, collects 2--3;
\citeproc{ref-Samuels1949}{Samuels 1949, 109--11};
\citeproc{ref-HoggGrammar2011}{Hogg 2011, 238--39}, §§6.53--6.54).

\subsection{Regular reconstruction and prosodic
domain}\label{regular-reconstruction-and-prosodic-domain}

The preferred reconstruction remains ordinary unstressed fronting,
palatalization, raising and later lowering:
\emph{*ga- > *gæ- > *gi- > *ge-}. Ringe and Taylor consider regular
raising of short pretonic \emph{*æ} after a palatal fricative as one
possibility, alongside influence from \emph{bi-}. The preceding high
front articulation supplies a plausible phonetic motivation, not direct
measurement of the intermediate vowel. Fronting and later weak-vowel
merger are independently supported (\citeproc{ref-Campbell1959}{Campbell
1959, 140--41}, §§333--334, and pp.~153--155, §§369--370;
\citeproc{ref-RingeTaylor2014}{Ringe and Taylor 2014, 334--35}).

Fronting and merger alone give late \emph{ge-} but do not explain
genuine early \emph{gi-} without an additional account. A raised
realization rather than an obligatory separate high-vowel stage remains
possible; Hogg's posttonic raising before following palatals and late
Northumbrian allophony, however, do not independently establish that
early pathway after a preceding fricative. Influence from \emph{bi-} is
another published possibility, not a demonstrated replacement. Pheifer's
uniform \emph{bi-} and variable \emph{gi-} are not matched for stress or
lexical distribution; he considers stressed \emph{bi-} and an ablaut
variant corresponding to continental \emph{ga-} as possible
explanations. No positive redistribution or substitution has been
established. Equally, spelling alone does not demonstrate an obligatory
separate high-vowel phonemic stage. Regular raising is the more
economical working reconstruction within the adopted regular account,
not an experimentally observed mechanism
(\citeproc{ref-RingeTaylor2014}{Ringe and Taylor 2014, 335};
\citeproc{ref-Berg2014}{Berg 2014, 60--61};
\citeproc{ref-Pheifer1974}{Pheifer 1974, lxxv}, §63;
\citeproc{ref-HoggGrammar1992}{Hogg 1992a, §§ 6.51--6.54} (held 2011
reprint, pp.~238--239)).

Weakness does not require word-initial position. Minkova's
\emph{mísgewìeld} `error' has the grouping \emph{(mís-ge)-(wìeld)}, with
weak internal \emph{ge-} incorporated into the first prosodic word. In
\emph{geédnìwad} `restored', \emph{ge-(éd)-(nì.wad)} has a weak outer
prefix before an independently stressed inner element and hiatus. The
form \emph{únforcùð} `dishonest' supplies a comparable multi-domain
structure, not the same palatal raising environment. These analyses are
grounded in Old English evidence; they cannot simply be projected back
as Proto-Germanic accent. Minkova also notes the lack of equivalent
evidence for extending the analysis to finite double-prefixed verbs
(\citeproc{ref-Minkova2008}{Minkova 2008, 38--41}).

Manuscript grouping can similarly join weak material to a preceding
unit. Prosodic words, stress feet and written words therefore need not
coincide. The present derivations cover one initial short weak syllable
before an independently headed domain, not every attested compound
structure. Internal \emph{modgidanc} `thought' and \emph{saegaesetu}
deserve historical domain analysis; they are not irrelevant because the
formula is word-initial, but neither yet supplies a securely
reconstructed early same-environment test. No independent
non-\emph{*ga-} positive or genuine same-environment unraised negative
has been established (\citeproc{ref-Sweet1885}{Sweet 1885, 149};
\citeproc{ref-Pheifer1974}{Pheifer 1974, lxxv, 39};
\citeproc{ref-RieggerBogelWalkden2024}{Riegger et al. 2024, 2065--68}).

The lexical concentration of the evidence is consequential: a rule
without a prefix label is not thereby independently demonstrated across
the vocabulary. Its regular phonological conditioning remains a testable
reconstruction of this prefix family's history, rather than a securely
established generalization from several unrelated families.

\subsection{\texorpdfstring{\CAPRRuleHeading{SC035. Raising of pretonic \emph{*æ} after a palatal fricative}{OEPrefixAReduction}}{}}\label{rule-OEPrefixAReduction}

\noindent\begin{minipage}{\linewidth}
\begin{ReaderFacingFoma}
define OEPrefixAReduction SentenceCtx([
    {*æ} -> {*i} || .#. {*ᶠ} OEAnyConsonant* {*ʝ} _
]);
\end{ReaderFacingFoma}
\end{minipage}

The conditioning is shortness, pretonic weakness and an immediately
preceding palatal fricative, not prefix identity. The superscript
\texttt{*ᶠ} is supplied by an earlier stress-foot-context operation: it
marks one initial short unaccented vowel separated by one or more
consonants from an explicitly acute-accented following vowel in the
reconstructed input. It preserves that positively encoded headed
structure when later rules remove the input accents. This is
nonsegmental notation for reconstructed prosody, not a manuscript sign,
a sound or an unrestricted parser of Old English stress. Hiatus, long
weak vowels, multiple pretonics and internal weak syllables require
their own historical analysis; the annotation does not parse them.

The complementary lack of weak-vowel diphthongization follows from
\hyperref[rule-OEWsPalatalDiphthongization]{SC056
OEWsPalatalDiphthongization}: the ordinary process requires a stressed
vowel at the beginning of a headed foot. Both \emph{sċeaft} `shaft' and
\emph{ġesċeaft} `creation' diphthongize the stressed root after the
sh-sound, but the latter's weak first vowel does not diphthongize. The
model also derives \emph{ġestīgan} `climb into/onto'. These comparisons
support the stress distinction without making contested \emph{ġeatwe}
`equipment' its sole diagnostic (\citeproc{ref-HoggGrammar1992}{Hogg
1992a, § 5.48} n.~1 (held 2011 reprint, p.~104);
\citeproc{ref-RingeTaylor2014}{Ringe and Taylor 2014, 335}).

For \emph{ġelīefan} `believe',
\hyperref[rule-EAFBrighteningUnstressed]{SC107 EAFBrighteningUnstressed}
supplies the front vowel and \hyperref[rule-OEVelarPalatalization]{SC052
OEVelarPalatalization} the fricative trigger. The rule raises the weak
vowel; \hyperref[rule-OEUnstressedAEMerger]{SC073 OEUnstressedAEMerger}
leaves that high vowel intact; \hyperref[rule-OEPrefixIReduction]{SC076
OEPrefixIReduction} later lowers it to \emph{e}. These are feeding
requirements of the adopted reconstruction, not independently dated
events inferred from a final match. Hogg places weak \emph{i} lowering
marginally after the earlier front-vowel merger; Pheifer's transmitted
spelling counts alone do not date that interval
(\citeproc{ref-RingeTaylor2014}{Ringe and Taylor 2014, 245, 267, 335};
\citeproc{ref-HoggGrammar1992}{Hogg 1992a, §§ 6.50--6.54} (held 2011
reprint, pp.~237--239)).

The staged account is retained with qualified confidence: it explains
early high-vowel evidence through regular phonology, while the bounded
implementation neither claims universal compound coverage nor treats one
successful derivation as independent proof of historical generality.

\section{Post-velar w-loss}\label{post-velar-w-loss}

\subsection{\texorpdfstring{SC053. Loss of \emph{*w} after velars
(\texttt{OEPostVelarWLoss})}{SC053. Loss of  after velars (OEPostVelarWLoss)}}\label{rule-OEPostVelarWLoss}

\noindent\begin{minipage}{\linewidth}
\begin{ReaderFacingFoma}
define OEPostVelarWLoss Ctx([
    {*w} -> 0 || {*n} {*g} _
]);
\end{ReaderFacingFoma}
\end{minipage}

The comparative development \texttt{*singwan\ \textgreater{}\ singan}
establishes narrow post-velar \emph{*w}-loss in the \emph{*ngw}
sequence, yielding \emph{singan} `sing'. Moving
\hyperref[rule-OEPostVelarWLoss]{SC053 OEPostVelarWLoss} earlier or
later leaves every output unchanged. Its pre-umlaut position therefore
rests on comparative evidence, while the present lexicon supplies no
neighboring boundary.

\section{West Saxon palatal
diphthongization}\label{west-saxon-palatal-diphthongization}

\subsection{\texorpdfstring{\CAPRRuleHeading{SC056. West Saxon palatal diphthongization}{OEWsPalatalDiphthongization}}{}}\label{rule-OEWsPalatalDiphthongization}

\begin{ReaderFacingFoma}
define OEWsPalatalDiphthongization FootCtx([
    {*æ} -> {*ea} || .#. [{*ʧ} | {*ʤ} | {*ʝ} | {*ʃ} | {*j}] _ [EnglishStarConsonant | EnglishPalatalConsonant | .#.],
    {*ǣ} -> {*ēa} || .#. [{*ʧ} | {*ʤ} | {*ʝ} | {*ʃ} | {*j}] _ [EnglishStarConsonant | EnglishPalatalConsonant | .#.],
    {*e} -> {*ie} || .#. [{*ʧ} | {*ʤ} | {*ʝ} | {*ʃ} | {*j}] _ [EnglishStarConsonant | EnglishPalatalConsonant | .#.],
    {*ē} -> {*īe} || .#. [{*ʧ} | {*ʤ} | {*ʝ} | {*ʃ} | {*j}] _ [EnglishStarConsonant | EnglishPalatalConsonant | .#.],
    {*é} -> {*íe} || .#. [{*ʧ} | {*ʤ} | {*ʝ} | {*ʃ} | {*j}] _ [EnglishStarConsonant | EnglishPalatalConsonant | .#.],
    {*ḗ} -> {*īe} || .#. [{*ʧ} | {*ʤ} | {*ʝ} | {*ʃ} | {*j}] _ [EnglishStarConsonant | EnglishPalatalConsonant | .#.]
]);
\end{ReaderFacingFoma}

West Saxon \emph{gieldan} `pay', \emph{scield} `shield', and
\emph{scieppan} `create' show diphthongization after an already palatal
consonant (\citeproc{ref-RingeTaylor2014}{Ringe and Taylor 2014,
215--16}, §6.5.1). Their dialectal and phonological restriction
separates this development from system-wide i-umlaut.

Hogg's \emph{giefan} `give' and \emph{sceap} `sheep' belong to the same
palatal-consonant environment (\citeproc{ref-HoggPhonology1992}{Hogg
1992b, 112}). Fulk likewise assigns this diphthongization a place before
front mutation and distinguishes the two processes
(\citeproc{ref-Fulk2018}{Fulk 2018, 74}, §4.13).

The relevant beginning is that of a headed stress foot, not necessarily
the word. Hogg's \emph{sċeaft} `shaft' and \emph{ġesċeaft} `creation'
both have diphthongization in the stressed root after the sh-sound. The
weak first vowel of \emph{ġesċeaft} `creation' does not diphthongize
after the initial j-sound. The context operation in the printed rule
preserves an explicitly identified pretonic syllable and applies the
ordinary rule to the following headed domain. It does not suppress
diphthongization throughout a prefixed word or infer weakness from a
missing late accent (\citeproc{ref-HoggGrammar1992}{Hogg 1992a, § 5.48} n.~1 (held 2011 reprint, p.~104)).

This foot-based reformulation retains ordinary diphthongization before
mutation; it does not reopen the separate later treatment of mutation
products. Its complementary pretonic vocalism is discussed under
\hyperref[rule-OEPrefixAReduction]{SC035 OEPrefixAReduction}.

The ordinary live applications comprise \emph{ġiefan} `give',
\emph{ġiest} `guest', \emph{sċeaft} `shaft', \emph{sċieran} `shear',
\emph{sċēap} `sheep', \emph{sċield} `shield' and \emph{ġēar} `year'.
Gift is unchanged at this rule because its selected vowel already
reflects earlier raising. Sheath changes from the mutation product
\emph{*ǣ} to \emph{*ēa}; it supplies evidence for the later \emph{sċ}
layer, not the date of the ordinary treatment represented by the other
examples.

The adopted ordinary process precedes mutation and preserves the split-
diphthong and high-i negative controls. The reported absence of a later
displacement failure supplies no positive historical terminus.

\section{Old English i-umlaut}\label{old-english-i-umlaut}

\subsection{\texorpdfstring{SC055. Fronting under i-umlaut
(\texttt{OEIUmlautFronting})}{SC055. Fronting under i-umlaut (OEIUmlautFronting)}}\label{rule-OEIUmlautFronting}

\begin{ReaderFacingFoma}
define OEIUmlautFronting [
    {*a} -> {*æ} || _ EnglishIUmlautThroughU EnglishIUmlautTrigger,
    {*ā} -> {*ǣ} || _ EnglishIUmlautThroughU EnglishIUmlautTrigger,
    {*e} -> {*i} || _ EnglishIUmlautThroughU EnglishIUmlautTrigger,
    {*o} -> {*e} || _ EnglishIUmlautThroughU EnglishIUmlautTrigger,
    {*ō} -> {*ē} || _ EnglishIUmlautThroughU EnglishIUmlautTrigger,
    {*ù} -> {*ỳ} || _ EnglishIUmlautThroughU EnglishIUmlautTrigger,
    {*u} -> {*y} || _ EnglishIUmlautThroughU EnglishIUmlautTrigger,
    {*ū} -> {*ȳ} || _ EnglishIUmlautThroughU EnglishIUmlautTrigger,
    {*á} -> {*æ} || _ EnglishIUmlautThroughU EnglishIUmlautTrigger,
    {*é} -> {*i} || _ EnglishIUmlautThroughU EnglishIUmlautTrigger,
    {*ó} -> {*e} || _ EnglishIUmlautThroughU EnglishIUmlautTrigger,
    {*ú} -> {*y} || _ EnglishIUmlautThroughU EnglishIUmlautTrigger
];
\end{ReaderFacingFoma}

The breadth of i-umlaut appears in lexical classes that share only a
following high front vocoid. The forms \emph{fylgan} `follow',
\emph{gylden} `golden', \emph{wyrm} `worm', and \emph{giest} `guest'
exemplify the same \texttt{i}- or \texttt{j}-conditioned fronting across
different vowels (\citeproc{ref-RingeTaylor2014}{Ringe and Taylor 2014,
222}, §6.6.1; \citeproc{ref-Campbell1959}{Campbell 1959, 69--72},
§§190--191).

The selected cow and lung inputs are negative controls on the present
palatalization rule. Moving the composite umlaut rule before that rule
yields \Pred{ċȳ} and \Pred{lunġen} instead of \emph{cȳ} `cow' and
\emph{lungen} `lung'. This constrains the current rule's productive
domain; it does not assign every consonantal layer the same date. In
particular the selected \emph{cȳ} is a dative-singular input, not
evidence from an assumed generic plural.

The stored gift and sheath displacement failures constrain the current
serialization, not a universal historical upper boundary. The former
gift input left an earlier raising to this late rule; sheath represents
the later \emph{sċ} treatment of a mutation product. Moving the entire
palatal-diphthongization bundle earlier strands sheath and diphthongizes
the model's still-unraised gift vowel. Neither failure overturns the
independently supported ordinary
palatal-diphthongization-before-mutation account.

Short medial \emph{u} can undergo mutation while transmitting it to the
preceding stressed monophthong. Hogg explicitly describes this iterative
development, and Campbell calls it double umlaut. The fronting context
therefore includes intervening plain medial \emph{u}, not a
secondary-stressed \emph{ù}. The latter keeps its stress as the internal
carrier \emph{ỳ} through both high-syncope components and weak-vowel
unrounding. The ordinary raising and diphthong components keep their
existing contexts; newly fronted \emph{æ} across a mutation-produced
\emph{y} is not indiscriminately raised again
(\citeproc{ref-HoggGrammar2011}{Hogg 2011, 120--21}, §5.76, p.~219,
§6.10; \citeproc{ref-Campbell1959}{Campbell 1959, 80--81}, §203).

\begin{ReaderFacingFoma}
define EnglishIUmlautThroughU [EnglishIUmlautIntervening [{*u} EnglishIUmlautIntervening]*];
\end{ReaderFacingFoma}

The independent outmost/furthest reconstruction and the separate
nonhigh-vowel and subordinate-stress controls are discussed with
\hyperref[rule-OEMedialSyncope]{SC065 OEMedialSyncope}; retained
weak-vowel unrounding is the separate early component
\hyperref[rule-OEHighLinkSyncope]{SC115 OEHighLinkSyncope}. This coupled
correction preserves the previously selected mutation and breaking
histories.

\subsection{\texorpdfstring{SC055. Raising under i-umlaut
(\texttt{OEIUmlautRaising})}{SC055. Raising under i-umlaut (OEIUmlautRaising)}}\label{rule-OEIUmlautRaising}

\noindent\begin{minipage}{\linewidth}
\begin{ReaderFacingFoma}
define OEIUmlautRaising [
    {*æ} -> {*e} || _ EnglishIUmlautIntervening EnglishIUmlautTrigger
];
\end{ReaderFacingFoma}
\end{minipage}

Raising of umlauted \texttt{æ} to \texttt{e} continues the same
assimilatory event as fronting and therefore shares the chronology of
general i-umlaut.

The four displacement controls concern the composite rule. They are not
independent tests of this raising component: for example, the cow vowel
changes in the fronting component, not in this \emph{*æ}-raising clause.
Component chronology must follow the source-supported assimilatory
event, with ordinary palatal diphthongization and the later \emph{sċ}
extension kept distinct.

The sources do not describe umlaut as simple fronting alone. Campbell
notes that the low front vowel changes again before \texttt{m} and
\texttt{n} in most dialects (\citeproc{ref-Campbell1959}{Campbell 1959,
69}, §190), and Hogg likewise treats short front vowels as part of the
same assimilatory system (\citeproc{ref-HoggPhonology1992}{Hogg 1992b,
113}).

\subsection{\texorpdfstring{SC055. Diphthongal outcomes under i-umlaut
(\texttt{OEIUmlautDiphthong})}{SC055. Diphthongal outcomes under i-umlaut (OEIUmlautDiphthong)}}\label{rule-OEIUmlautDiphthong}

\begin{ReaderFacingFoma}
define OEIUmlautDiphthong [
    {*ea} -> {*ie} || _ EnglishIUmlautIntervening EnglishIUmlautTrigger,
    {*ēa} -> {*īe} || _ EnglishIUmlautIntervening EnglishIUmlautTrigger,
    {*io} -> {*ie} || _ EnglishIUmlautIntervening EnglishIUmlautTrigger,
    {*īo} -> {*īe} || _ EnglishIUmlautIntervening EnglishIUmlautTrigger,
    {*eo} -> {*ie} || _ EnglishIUmlautIntervening EnglishIUmlautTrigger,
    {*ēo} -> {*īe} || _ EnglishIUmlautIntervening EnglishIUmlautTrigger,
    {*éa} -> {*íe} || _ EnglishIUmlautIntervening EnglishIUmlautTrigger,
    {*éo} -> {*íe} || _ EnglishIUmlautIntervening EnglishIUmlautTrigger,
    {*ío} -> {*íe} || _ EnglishIUmlautIntervening EnglishIUmlautTrigger
];
\end{ReaderFacingFoma}

Diphthongal outcomes belong to the same system-wide assimilation as
simple-vowel fronting and raising. All three therefore belong to a
single historical event.

The relevant examples are the recurring West-Saxon \texttt{ie} forms
cited in the handbooks, including \emph{giest} `guest', \emph{giefan}
`give', and \emph{hierde} `shepherd' in Campbell and \emph{ciest}
`chest' in Hogg (\citeproc{ref-Campbell1959}{Campbell 1959, 69--72,
78--80}, §§190--191, 248--251; \citeproc{ref-HoggPhonology1992}{Hogg
1992b, 113--14}). These diphthongal outcomes form a distinct part of the
general umlautal development alongside simple fronting.

Those composite displacement controls do not isolate diphthongal
mutation: neither the cow input nor the gift input has a diphthong here.
In the conventional history, diphthongs produced by ordinary palatal
diphthongization can be mutated, as in the source derivation of
\emph{ġiest} `guest' (\citeproc{ref-RingeTaylor2014}{Ringe and Taylor
2014, 216}). The adopted derivation now follows that ea-to-ie source
path. The former simple-vowel route converged on the same final form,
showing why final agreement alone could not establish the relative
chronology.

\subsection{\texorpdfstring{SC055. The composite i-umlaut rule
(\texttt{OEIUmlaut})}{SC055. The composite i-umlaut rule (OEIUmlaut)}}\label{rule-OEIUmlaut}

\noindent\begin{minipage}{\linewidth}
\begin{ReaderFacingFoma}
define OEIUmlaut FootCtx(OEIUmlautFronting
    .o. OEIUmlautRaising
    .o. OEIUmlautDiphthong);
\end{ReaderFacingFoma}
\end{minipage}

The literature presents fronting, raising, and diphthongal mutation as
effects of one historical development. They consequently occupy a single
place in the Old English chronology.

The cow/lung displacement results require the present productive
palatalization rule not to consume those mutation-created environments.
The historical interpretation still depends on the relevant consonantal
layer and inherited versus secondary front-vowel conditions.

There is no single source-supported upper boundary at the entire West
Saxon bundle. Ordinary palatal diphthongization precedes mutation in the
handbook account; the later \emph{sċ} treatment of some mutation
products follows it. The former gift input and sheath's later layer
explained the old bundled serialization, not a date for every historical
palatal process.

\section{Noninitial w-loss before unstressed
i}\label{noninitial-w-loss-before-unstressed-i}

\subsection{\texorpdfstring{SC054. Noninitial \emph{*w}-loss before
unstressed \emph{*i/ī}
(\texttt{OEWLossBeforeI})}{SC054. Noninitial -loss before unstressed  (OEWLossBeforeI)}}\label{rule-OEWLossBeforeI}

\noindent\begin{minipage}{\linewidth}
\begin{ReaderFacingFoma}
define OEWLossBeforeI Ctx([
    {*w} -> 0 || EnglishStarVocalic EnglishStarConsonant* _ [{*i}|{*ī}]
]);
\end{ReaderFacingFoma}
\end{minipage}

The history of \emph{sǣ} `sea' exemplifies noninitial \emph{*w}-loss
before unstressed \emph{*i}. Campbell describes the loss, Ringe and
Taylor derive the form from \emph{*saiwi-}/\emph{*sawi-}, and Luick
gives the parallel trajectory (\citeproc{ref-Campbell1959}{Campbell
1959, 166--67}, §406; \citeproc{ref-RingeTaylor2014}{Ringe and Taylor
2014, 258--60}, §6.7.1; \citeproc{ref-Luick1914}{Luick 1914-\/-1940,
173}, §187). The environment is neither exclusively postvocalic nor
exclusively final: Ringe and Taylor also give weak-verb forms in
\emph{*-Cwi-} and \emph{*pulwīn > *pylwī > *pylī}, yielding \emph{pyle}
`pillow'. The preceding nucleus distinguishes noninitial \emph{w} from
an initial cluster; the unaccented target excludes stressed high vowels.
No lexical exemption or loanword-specific clause is used.

Ringe and Taylor independently place this loss after i-umlaut and before
general syncope: the histories of \emph{eowu} `ewe' and \emph{gierede}
`prepared' supply the chronology, rather than successful output alone
(\citeproc{ref-RingeTaylor2014}{Ringe and Taylor 2014, 258--60}). The
early final-n loss in \hyperref[rule-NWGmcInStemNLoss]{SC064
NWGmcInStemNLoss} and this distinct glide loss together explain
\emph{pyle}.

The former sea-only displacement supplied two distant implementation
limits: before \hyperref[rule-EAFFinalZDeletion]{SC020
EAFFinalZDeletion} or after \hyperref[rule-OEHighVowelApocope]{SC063
OEHighVowelApocope}, that older final-only rule yielded \Pred{sǣw}
rather than expected OE \emph{sǣ} `sea'. Those receipts remain
historical evidence for the former implementation, not current tests of
the adopted noninitial environment. In particular, they do not
independently demonstrate that the present rule requires final
\emph{z}-deletion to expose its target. Its placement after mutation and
before general syncope rests on the cited histories, not on promoting
the former rule's displacement bounds.

\section{Technical support: late palatal
diphthongization}\label{technical-support-late-palatal-diphthongization}

\subsection{\texorpdfstring{\CAPRRuleHeading{SC105. Retained late palatal-diphthongization approximation}{OELatePalatalDiphthong}}{}}\label{rule-OELatePalatalDiphthong}

\begin{ReaderFacingFoma}
define OELatePalatalDiphthong Ctx([
    {*ǣ} -> {*ēa} || .#. [{*ʧ} | {*ʤ} | {*ʝ} | {*ʃ} | {*j}] _ [EnglishStarConsonant | EnglishPalatalConsonant | .#.]
]);
\end{ReaderFacingFoma}

This support operation retains exactly the former postmutation
long-vowel clause. Its sole current application is \emph{sċēaþ}
`sheath'. The four initial triggers are an unchanged approximation, not
a source-proven complete historical conditioner. Campbell and
Ringe--Taylor distinguish the later class
(\citeproc{ref-Campbell1959}{Campbell 1959, 68--69};
\citeproc{ref-RingeTaylor2014}{Ringe and Taylor 2014, 235}); the
incremental implementation does not claim to have resolved it.

\section{J-cluster coalescence}\label{j-cluster-coalescence}

\subsection{Historical discussion}\label{historical-discussion-43}

Only a small lexical group reveals the coalescence of velars with
\emph{*j}. Plain-velar and \emph{*sk} palatalization must already have
run before \emph{*gj} and \emph{*kj} acquire their later outcomes.
Campbell, Ringe and Taylor, and Fulk discuss the palatalized and fronted
outcomes in \emph{bīeġan} `bend' and \emph{sēċan} `seek' without
assigning this later cluster adjustment the status of a major sound law
(\citeproc{ref-Campbell1959}{Campbell 1959, 89, 107--8}, §§170,
248--251; \citeproc{ref-RingeTaylor2014}{Ringe and Taylor 2014,
213--51}, §§6.4.1, 6.5.1, 6.6.1--6.6.4; \citeproc{ref-Fulk2018}{Fulk
2018, 65, 75}, §§4.7, 4.13).

\subsection{\texorpdfstring{SC057. Coalescence of velar + \emph{*j}
clusters
(\texttt{OEJClusterCoalescence})}{SC057. Coalescence of velar +  clusters (OEJClusterCoalescence)}}\label{rule-OEJClusterCoalescence}

\noindent\begin{minipage}{\linewidth}
\begin{ReaderFacingFoma}
define OEJClusterCoalescence Ctx((
    # *bj coalescence removed: BAllophony (*b→*β) + JLossAfterHeavy handle
    # non-geminated *bj; for geminated *bbj, *j is consumed by JLossAfterHeavy.
    [{*g} {*j} -> {*ʤ}]
    .o. [{*k} {*j} -> {*ʧ}]
));
\end{ReaderFacingFoma}
\end{minipage}

The forms \emph{bīeġan} `bend' and \emph{sēċan} `seek' determine the
earlier boundary. If coalescence precedes
\hyperref[rule-OEVelarPalatalization]{SC052 OEVelarPalatalization}, the
developments behind \emph{bīeġan} `bend' and \emph{sēċan} `seek' are
lost. Related forms such as \emph{fylġan} `follow', \emph{heċġ} `hedge',
and \emph{sengan} `singe' fail in the same broader palatalization zone.
PGmc \Recon{báugijaną} `bow' yields \Pred{bēaġan} rather than expected
OE \emph{bīeġan}, and PGmc \Recon{sōkijaną} `seek' yields \Pred{sōċan}
rather than expected \emph{sēċan}. This constrains this computational
cluster consumer. It does not independently date all singleton-fricative
articulation or merger: the latter is distinct from this rule and is now
explicit under \hyperref[rule-OEPalatalFricativeMerger]{SC109
OEPalatalFricativeMerger} (\citeproc{ref-RingeTaylor2014}{Ringe and
Taylor 2014, 204}; \citeproc{ref-Laker2007}{Laker 2007, 167--68}).
Nothing in the present lexicon supplies a terminus ante quem for the
cluster consumer.

\section{Back mutation}\label{back-mutation}

\subsection{Historical discussion}\label{historical-discussion-44}

West Saxon \emph{giefan} `give' and \emph{wefan} `weave' stand against
non-West-Saxon \emph{geofad} `gave' and \emph{weofan} `weave'. Ringe and
Taylor use this contrast to define the dialectal profile of back
mutation (\citeproc{ref-RingeTaylor2014}{Ringe and Taylor 2014, 319},
§6.9.4). Campbell's treatment of diphthongization before following back
vowels includes \emph{heofon} `heaven'
(\citeproc{ref-Campbell1959}{Campbell 1959, 86}, §207), while Hogg draws
the instructive comparison with breaking
(\citeproc{ref-HoggPhonology1992}{Hogg 1992b, 113}). Fulk accordingly
separates back mutation from the earlier umlautal changes
(\citeproc{ref-Fulk2018}{Fulk 2018, 69}, §4.8). Barrack likewise states
back mutation apart from the earlier mutations, as his rules 210 and 212
with a continuation at 220 (\citeproc{ref-Barrack1975}{Barrack 1975,
63--66}); his residual-ordering discussion uses a back umlaut form,
\emph{heolap}, as the worked example of a rule pair whose formal order
must be reversed for a residue (210 against his rule 355.5)
(\citeproc{ref-Barrack1975}{Barrack 1975, 36}), which illustrates his
insistence that such orderings are logical rather than chronological
relations (\citeproc{ref-Barrack1975}{Barrack 1975, 34}).

\subsection{\texorpdfstring{SC059. Back mutation before labials and
liquids
(\texttt{OEBackMutation})}{SC059. Back mutation before labials and liquids (OEBackMutation)}}\label{rule-OEBackMutation}

\begin{ReaderFacingFoma}
define OEBackMutation Ctx([
    {*e} -> {*eo} || _ [EnglishStarLabial | EnglishStarLiquid] {*u},
    {*æ} -> {*ea} || _ [EnglishStarLabial | EnglishStarLiquid] EnglishBackMutationTrigger,
    # Stressed variant (stress preserved on first element):
    {*é} -> {*éo} || _ [EnglishStarLabial | EnglishStarLiquid] {*u}
]);
\end{ReaderFacingFoma}

Three witness forms bracket the chronology. If back mutation precedes
\hyperref[rule-OESecondaryNasalization]{SC048 OESecondaryNasalization},
forms such as \Recon{gébaną} `give' produce \emph{ġeofan} `give'; the
expected form is \emph{ġiefan} `give'. \Recon{stélaną} `steal' likewise
produces \emph{steolan} `steal'; the expected form is \emph{stelan}
`steal'. At the other edge, delaying back mutation until after
\hyperref[rule-OEWeakTailReduction]{SC078 OEWeakTailReduction} makes
\Recon{wébaną} `weave' yield \emph{weofan} `weave'; the expected form is
\emph{wefan} `weave'. Thus back mutation follows secondary nasalization
but precedes the weak-tail reductions.

\section{West Saxon palatal umlaut}\label{west-saxon-palatal-umlaut}

\subsection{Historical discussion}\label{historical-discussion-45}

The reflexes \emph{miht} `might' and \emph{niht} `night' place West
Saxon palatal umlaut after the principal umlautal developments. Campbell
and Ringe and Taylor describe the forms themselves; Fulk supplies the
broader chronology of palatal-vowel change
(\citeproc{ref-Campbell1959}{Campbell 1959, 107--8}, §§248--251;
\citeproc{ref-RingeTaylor2014}{Ringe and Taylor 2014, 215--51}, §§6.5.1,
6.6.1--6.6.4; \citeproc{ref-Fulk2018}{Fulk 2018, 65, 75}, §§4.7, 4.13).

\subsection{\texorpdfstring{\CAPRRuleHeading{SC060. West Saxon palatal umlaut before \emph{*h}-clusters}{OEWsPalatalUmlaut}}{}}\label{rule-OEWsPalatalUmlaut}

\begin{ReaderFacingFoma}
define OEWsPalatalUmlaut Ctx([
    {*eo} -> {*i} || _ OEHCluster .#.,
    {*io} -> {*i} || _ OEHCluster .#.,
    {*ie} -> {*i} || _ OEHCluster .#.,
    {*eo} -> {*i} || _ OEHCluster EnglishStarFrontVowel,
    {*io} -> {*i} || _ OEHCluster EnglishStarFrontVowel,
    {*ie} -> {*i} || _ OEHCluster EnglishStarFrontVowel,
    # Stressed variants (stress absorbed: result is short i):
    {*éo} -> {*i} || _ OEHCluster .#.,
    {*ío} -> {*i} || _ OEHCluster .#.,
    {*íe} -> {*i} || _ OEHCluster .#.,
    {*éo} -> {*i} || _ OEHCluster EnglishStarFrontVowel,
    {*ío} -> {*i} || _ OEHCluster EnglishStarFrontVowel,
    {*íe} -> {*i} || _ OEHCluster EnglishStarFrontVowel
]);
\end{ReaderFacingFoma}

The change to \emph{*i} before \emph{*h}-clusters can be ordered only on
its earlier side. If palatal umlaut precedes
\hyperref[rule-OEIUmlaut]{SC055 OEIUmlaut}, the forms behind \emph{miht}
`might' and \emph{niht} `night' remain at the overdeveloped stage
\Pred{mieht} and \Pred{nieht} rather than expected OE \emph{miht} and
\emph{niht}. Consequently, i-umlaut precedes palatal umlaut. Reordering
the latter against any tested later change leaves both witness forms
unchanged.

\section{Weak-tail nasal loss}\label{weak-tail-nasal-loss}

\subsection{Historical discussion}\label{historical-discussion-46}

The pathway from \Recon{dōną} `do' to \emph{dōn} `do' supplies the sole
lexical thread through this reduction. Campbell, Hogg, and Fulk place
such weak-tail losses among apocope and related late reductions
(\citeproc{ref-Campbell1959}{Campbell 1959, 144--45}, §§345--349;
\citeproc{ref-HoggPhonology1992}{Hogg 1992b, 121};
\citeproc{ref-Fulk2018}{Fulk 2018, 91}, §5.6). The witness, however,
ties the change to a much older development. Its immediate neighbors
remain untested.

\subsection{\texorpdfstring{\CAPRRuleHeading{SC061. Reduction of final nasal weak-tail endings}{OEWeakTailNasalLoss}}{}}\label{rule-OEWeakTailNasalLoss}

\noindent\begin{minipage}{\linewidth}
\begin{ReaderFacingFoma}
define OEWeakTailNasalLoss Ctx([
    {*n} {*ą} -> {*n} || _ .#.,
    {*m} {*ą} -> {*m} || _ .#.
]);
\end{ReaderFacingFoma}
\end{minipage}

Final weak-tail \emph{*-ną} and \emph{*-mą} accordingly yield plain
\emph{*-n} and \emph{*-m}.

Only \emph{dōn} `do' constrains the relative order. Placing this loss
before the older n-stem loss makes the derivation record no output
instead of expected OE \emph{dōn} `do'. The older loss must therefore
precede weak-tail nasal loss. Nothing in the current lexicon
distinguishes among its possible later positions, and one witness cannot
establish a wider historical development.

\section{Ordinary nonhigh-vowel
syncope}\label{ordinary-nonhigh-vowel-syncope}

\subsection{Historical discussion}\label{historical-discussion-47}

Old English nonhigh-vowel syncope must be distinguished from the later
high-vowel losses. Hogg describes loss of a short nonhigh vowel in an
open medial syllable, after both light and heavy roots, with restoration
and paradigm interactions considered separately. His account places the
process after umlaut and before unstressed \emph{i} lowers
(\citeproc{ref-HoggGrammar1992}{Hogg 1992a, 220--22}, §§6.14--6.18).
Luick's longer-word cases warn against reducing the domain to exactly
three syllables (\citeproc{ref-Luick1914}{Luick 1914-\/-1940, 279--80},
§303). These are phonological and prosodic conditions, not permissions
for particular words or grammatical classes.

The implemented core deletes plain short \emph{a, æ, e} from an open
posttonic syllable. Primary-root weight and final-syllable weight are
not additional deletion conditions. A following singleton consonant and
vowel identify the bounded open-medial domain; this does not silently
claim complete syllabification of arbitrary consonant clusters. A
positively identified primary root protects its own vowel, and the
existing stress-domain machinery admits the supported single pretonic
syllable. Wider multiprefix histories are not inferred from this
restricted account.

\subsubsection{Selected applications and contrasting
forms}\label{selected-applications-and-contrasting-forms}

The strong masculine nominative plural participle \emph{*fúnðanai} gives
\emph{fundne} `found', attested in \emph{Juliana}335. The inherited
ending becomes long \emph{ē}; medial \emph{a} then disappears before its
still-heavy final syllable, and later shortening gives final \emph{e}.
The cell is supported independently by the edition's syntax and
glossary, while the reconstructed whole word combines participial
formation and adjectival inflection (\citeproc{ref-Strunk1904}{Strunk
1904, 15, 82}; \citeproc{ref-Fulk2018}{Fulk 2018, 209}, §9.2, p.~284,
§12.30). The genuine longer \emph{fundene} `found' and competing
inherited participial suffixes remain evidence for the wider paradigm,
not exceptions manufactured by changing the sound law.

The inherited accusative \emph{*sáiwalǭ} gives \emph{sāwle} `soul'.
After stressed-diphthong monophthongization, medial \emph{a} is lost
before the quantity-bearing inherited ending; its later decoding and
shortening supply final \emph{e}. The ending notation does not imply
unchanged late phonetic nasal rounding. Luke12:20 establishes the
accusative, distinct from the neighbouring dative and vocative
(\citeproc{ref-Skeat1874}{Skeat 1874, 130}). No raising proxy or
epenthesis is required in this selected span. Fulk's competing metrical
account and the syncope/apocope/epenthesis history proposed for
nominatives remain independent questions (\citeproc{ref-Fulk2010}{Fulk
2010, 133--39}; \citeproc{ref-RingeTaylor2014}{Ringe and Taylor 2014,
265, 330}).

Citation \emph{hamer} `hammer' is a nonapplication: earlier final-vowel
loss leaves the nonhigh vowel in the last syllable. By contrast, the
longer genitive input predicts \emph{hamres}, not attested
\emph{hameres} `of a hammer'; that wider history remains unresolved
(\citeproc{ref-WrightWuelcker1884}{Wright 1884, vol. 1}, p.~272, line
36; \citeproc{ref-Toller1921}{Toller 1921, 505}). The inherited neuter
base \emph{*wéstą} gives \emph{west} `westwards' through existing
nasal-vowel apocope, with no medial syncope target
(\citeproc{ref-Orel2003}{Orel 2003, 459}; \citeproc{ref-Fulk2018}{Fulk
2018, 146}, §7.8).

\subsubsection{Relative chronology and
limits}\label{relative-chronology-and-limits}

The executable position immediately before
\hyperref[rule-OEHighVowelApocope]{SC063 OEHighVowelApocope} preserves
visibility of the inherited final quantity. That serialization is a
model-local representation choice, not a new independently demonstrated
historical ordering inferred from output success. The source-supported
chronology and the notation's lifetime must be kept separate. The rules
\hyperref[rule-OEMedialSyncope]{SC065 OEMedialSyncope} and
\hyperref[rule-OELAdjacentSyncope]{SC066 OELAdjacentSyncope} now
distinguish ordinary heavy/open high-vowel loss from later light
pre-\emph{*l} loss. \emph{Weorold} `world' retains its established
selected history. \emph{Eorþcunde} `earthly' and \emph{reġnboga}
`rainbow' now lose the linking \emph{*a} here rather than through the
former early proxy.

\subsection{Compound formation and comparative linking-vowel
evidence}\label{compound-formation-and-comparative-linking-vowel-evidence}

Compound members with weakened force often lose or reshape their linking
vowels, and Campbell treats that broad pattern through reduced second
elements, connecting vowels, and obscured compounds
(\citeproc{ref-Campbell1959}{Campbell 1959, 148--49}, §§356--357;
\citeproc{ref-Campbell1959}{Campbell 1959, 153}, §367;
\citeproc{ref-Campbell1959}{Campbell 1959, 159}, §§386--387).

Those discussions concern several distinct developments: modification
and shortening within obscured second roots, connecting-vowel insertion,
and weak-vowel quality changes do not by themselves demonstrate deletion
of a linking vowel. For syncope, the relevant distinction is between
short nonhigh vowels, whose loss does not require a heavy first
syllable, and high linking \emph{i/u}, whose regular heavy-root loss
supplies the weight condition modelled under
\hyperref[rule-OEMedialSyncope]{SC065 OEMedialSyncope}. That condition
is not an exhaustive account of all historical compound losses. The
handbooks compare retention in light-root \emph{Winifrið/Winefrið} and
\emph{duruweard} `doorkeeper' with loss in \emph{dælnimend} `partaker'
and \emph{handclāð} `towel' (\citeproc{ref-Campbell1959}{Campbell 1959,
143}, §341; \citeproc{ref-Campbell1959}{Campbell 1959, 145--46},
§§348--353; \citeproc{ref-HoggGrammar2011}{Hogg 2011, 220--23},
§§6.14--6.19; \citeproc{ref-RingeTaylor2014}{Ringe and Taylor 2014,
266--67}; \citeproc{ref-SieversBrunner1965}{Brunner 1965, 137--38},
§§167--168). These are comparison forms, not additional selected
derivations. Variant first stems and the \emph{hild-/hilde-} material
require source-critical analysis rather than an automatic rule
conditioned by the identity or weight of the second member
(\citeproc{ref-RingeTaylor2014}{Ringe and Taylor 2014, 303}).

The handbook comparisons do not yet provide a selected complete history
testing the newly restricted high-vowel clause. In particular, Orel
reconstructs Proto-Germanic
\Recon{đura-warđaz}\index[iv]{03pgmc@\ivlangheader{Proto-Germanic}{}!durawardaz@\iventry{*đura-warđaz}{}}
`doorkeeper', with an \emph{a}-linker, and compares Gothic
{\emph{daura-wards}}\index[iv]{12goth@\ivlangheader{Gothic}{}!daurawards@\iventry{daura-wards}{}}
`doorkeeper'; Kroonen derives the Old English simplex
{\emph{duru}}\index[iv]{02oe@\ivlangheader{Old English}{West Saxon normalization unmarked}!duru@\iventry{duru}{}}
`door' from
\Recon{durō-}\index[iv]{03pgmc@\ivlangheader{Proto-Germanic}{}!duro@\iventry{*durō-}{}}.
Neither fact establishes an inherited short linking \emph{u} in
{\emph{duruweard}}\index[iv]{02oe@\ivlangheader{Old English}{West Saxon normalization unmarked}!duruweard@\iventry{duruweard}{}}
`doorkeeper', nor an independently demonstrated change supplying it
before syncope. Its attested vowel cannot simply be substituted into the
earlier compound reconstruction. The retained name vowels also need
their own formation histories (\citeproc{ref-Orel2003}{Orel 2003, 79};
\citeproc{ref-Kroonen2013}{Kroonen 2013, 110};
\citeproc{ref-SieversBrunner1965}{Brunner 1965, 137--38}, §168).

The older names cannot simply be dismissed as exceptions to a universal
heavy-root requirement. Sievers and Brunner, Campbell, Luick and Hogg
explicitly recognize high-vowel loss after light first roots, including
{\emph{Sigfrið}}\index[iv]{02oe@\ivlangheader{Old English}{West Saxon normalization unmarked}!sigfrith@\iventry{Sigfrið}{}}
and
{\emph{Hyglāc}}\index[iv]{02oe@\ivlangheader{Old English}{West Saxon normalization unmarked}!hyglac@\iventry{Hyglāc}{}}.
Luick distinguishes chronological layers and says that the circumstances
of the change are not yet understood. Hogg also discusses the second
member's weight. These observations require a bounded reconsideration of
an exhaustive weight-only account; they do not license lexical
optionality, a rule restricted to personal names, or an unsupported
appeal to analogy. Nor does the manuscript date of a name independently
date its formation or syncope (\citeproc{ref-SieversBrunner1965}{Brunner
1965, 137--38}, §168; \citeproc{ref-Campbell1959}{Campbell 1959,
145--46}, §§348--349; \citeproc{ref-Luick1914}{Luick 1914-\/-1940,
283--84}, §305; \citeproc{ref-HoggGrammar2011}{Hogg 2011, 224--25},
§6.21).

{\emph{seleweard}}\index[iv]{02oe@\ivlangheader{Old English}{West Saxon normalization unmarked}!seleweard@\iventry{seleweard}{}}
`hall-warden' provides a stronger retained-vowel candidate. Orel prints
\Recon{sali-warđaz}\index[iv]{03pgmc@\ivlangheader{Proto-Germanic}{}!saliwardaz@\iventry{*sali-warđaz}{}}
and
\Recon{sali-warđuz}\index[iv]{03pgmc@\ivlangheader{Proto-Germanic}{}!saliwarduz@\iventry{*sali-warđuz}{}},
with Old Norse
{\emph{sal-vǫrðr}}\index[iv]{09on@\ivlangheader{Old Norse}{}!salvordr@\iventry{sal-vǫrðr}{}}
`hall-warden'. Both reconstructions have first-member \emph{i}; their
alternatives concern the second member's inflection, not an \emph{a/i}
linker choice. The first member's secondary i-stem history is
independently supported. Nevertheless, these transparent formations do
not discriminate inheritance of this particular compound from later
independent formation: Orel expressly reports Carr's parallel-formation
assessment. Beowulf attests the Old English accusative singular, but
does not independently place the compound before the relevant syncope.
Neither an ending nor a high linker has to be invented; the unresolved
premise is formation chronology (\citeproc{ref-Orel2003}{Orel 2003, 315,
448}; \citeproc{ref-Kroonen2013}{Kroonen 2013, 424};
\citeproc{ref-ClarkHall1960}{Clark Hall 1960, 260}, Beowulf 667).

The canonical grammar now recognizes inherited high-vowel first members
and the second member's dental cluster without sacrificing older
accepted analyses. Both printed hypotheses, and the ordinary a-stem
accusative \emph{*sáli-wàrðą}, regularly produce \emph{seleweard}.
Previously the retained plain medial \emph{i} incorrectly underwent
breaking, producing \Pred{seleoweard}. Restoring the former early
high-vowel deletion instead produces \Pred{sealweard}; removing the
later weight condition produces \Pred{selweard}. These are three
distinct wrong analyses, not three attestations. Correct output
establishes a conditional derivation, not the date or inheritance of the
compound.

No new lexical selection is made without that independent formation
premise. Complete selected weight controls and independently
reconstructed later-stage controls now test the high-vowel core
described under \hyperref[rule-OEMedialSyncope]{SC065 OEMedialSyncope};
they do not supply a complete inherited-compound discriminator for the
former early rule.
{\emph{reġnboga}}\index[iv]{02oe@\ivlangheader{Old English}{West Saxon normalization unmarked}!regnboga@\iventry{reġnboga}{}}
`rainbow' and
{\emph{eorþcunde}}\index[iv]{02oe@\ivlangheader{Old English}{West Saxon normalization unmarked}!eorthcunde@\iventry{eorþcunde}{}}
`earthly' remain genuine selected \emph{a}-deletion witnesses, not
evidence for the new high-vowel restriction. The source-proposed
\emph{ūtumist} history supplies a distinct transmitted-mutation and
weak-vowel-loss control; its final root alone does not independently
order mutation against loss of \emph{u}
(\citeproc{ref-RingeTaylor2014}{Ringe and Taylor 2014, 267}).

Barrack's weak-vowel fronting and quality rules 333a/b and his rounding
rule 335 belong to this wider discussion, but are not equivalent
linking-vowel deletion laws. Their formal ordering is not independent
evidence for dating syncope or removal of secondary stress
(\citeproc{ref-Barrack1975}{Barrack 1975, 82--84}).

The former \hyperref[rule-OECompoundLinkingSyncope]{SC037
OECompoundLinkingSyncope} proxy captured this pattern in \emph{reġnboga}
`rainbow'. Its second-root vowel has already lowered under subordinate
stress: \emph{*régna-bùgô > *régna-bògô}. The grave survives that early
change and formerly licensed early deletion of the linking vowel, giving
\emph{*régnbògô}; ordinary accent stripping formerly gave
\emph{*régnbogô}. The full account and its reconstruction qualifications
belong to \hyperref[rule-PNWGmcULowering]{SC017 PNWGmcULowering}, not to
late unstressed lowering. Its immediate annotation dependency does not
date historical nonhigh syncope. The handbooks place ordinary nonhigh
syncope after breaking, and high-vowel loss normally after mutation. The
early \emph{a} fragment is retained as a selected-input holding proxy in
the earlier model. It is now an identity: ordinary nonhigh deletion
applies here, independently of a surviving grave, and inherited
\emph{*b} receives positional realization afterward under
\hyperref[rule-PGmcBAllophony]{SC049 PGmcBAllophony}. The high-vowel
component follows mutation (\citeproc{ref-SieversBrunner1965}{Brunner
1965, 137--38}, §167 and notes; \citeproc{ref-HoggGrammar2011}{Hogg
2011, 220--24}, §§6.14--6.20; \citeproc{ref-RingeTaylor2014}{Ringe and
Taylor 2014, 266--67}).

The earlier proxy's apparent dependency on grave removal did not date a
universal loss of compound secondary stress. The relevant junctures also
relate to \hyperref[rule-OEPrefixAReduction]{SC035 OEPrefixAReduction}
and \hyperref[rule-OEInterStressRaising]{SC036 OEInterStressRaising}.
Russom's metrical review challenges obligatory positional destressing;
modern stress marks are annotations, not independent prosodic evidence
for every historical name (\citeproc{ref-Russom2001}{Russom 2001,
39--64}).

\subsection{\texorpdfstring{SC114. Ordinary nonhigh syncope
(\texttt{OENonHighSyn})}{SC114. Ordinary nonhigh syncope (OENonHighSyn)}}\label{rule-OENonHighSyn}

\noindent\begin{minipage}{\linewidth}
\begin{ReaderFacingFoma}
define OENonHighSyn FootCtx([
    [{*a}|{*æ}|{*e}] -> 0 ||
        .#. OEAnyConsonant* EnglishStarVocalic
        [OEAnyConsonant+ EnglishStarVocalic]* OEAnyConsonant+
        _ OEAnyConsonant EnglishStarVocalic
        [OEAnyConsonant|EnglishStarVocalic|{*ô}]* .#.
]);
\end{ReaderFacingFoma}
\end{minipage}

\section{B allophony}\label{b-allophony}

\subsection{Historical discussion}\label{historical-discussion-48}

The positional alternation of Germanic \emph{*b} is a Proto-Germanic
distributional feature. Hogg states the Old English distribution
clearly: /b/ is a stop initially, after nasals, and in gemination, while
the same segment is otherwise realized as a voiced bilabial fricative
(\citeproc{ref-HoggGrammar2011}{Hogg 2011, 29--30}, §§2.53--2.54). Ringe
and Taylor support the broader West Germanic background by treating
Proto-West-Germanic \emph{*b} as a segment whose stop and fricative
values depend on position (\citeproc{ref-RingeTaylor2014}{Ringe and
Taylor 2014, 121}), and Luick's spelling evidence shows the same labial
fricative pattern in Old English (\citeproc{ref-Luick1914}{Luick
1914-\/-1940, 107}).

The distribution is narrow, but later changes presuppose the
stop-fricative alternation. CAPR implements the rule at a late cascade
position for computational reasons: the alternation must interact with
consonant environments shaped by intermediate rule applications. The
inherited distribution is Proto-Germanic; the position at which CAPR
realizes it is not a new Proto-Germanic event.

Hogg distinguishes the early bilabial fricative, still an allophone of
inherited /b/, from its later split and merger with /f/. Campbell's
spelling discussion likewise distinguishes inherited labial realization
from secondary voicing of /f/ (\citeproc{ref-HoggGrammar2011}{Hogg 2011,
29--30}, §§2.53--2.54; \citeproc{ref-Campbell1959}{Campbell 1959,
179--81}, §§444--448). CAPR retains underlying \emph{*b} until ordinary
nonhigh syncope has established its conditioning environment. This
removes the need for early compound deletion; it does not insert a later
\emph{*β > *b} sound law.

The distinction matters: already represented \emph{*β} and inherited
\emph{*f} remain fricatives after nasals, whereas underlying \emph{*b}
has stop realization there. Initial, postnasal and geminate stops are
preserved; intervocalic \emph{*b} receives fricative realization.
Syncope-created velar clusters do not license analogous hardening: Ringe
and Taylor contrast \emph{syngian} `sin' with inherited stop clusters
(\citeproc{ref-RingeTaylor2014}{Ringe and Taylor 2014, 4}). The chosen
serialization precedes the labial split in the model, not a newly
demonstrated strict historical boundary.

\subsection{\texorpdfstring{\CAPRRuleHeading{SC049. Distribution of \emph{*b} after vowels and liquids}{PGmcBAllophony}}{}}\label{rule-PGmcBAllophony}

\noindent\begin{minipage}{\linewidth}
\begin{ReaderFacingFoma}
define PGmcBAllophony Ctx([
    {*b} -> {*β} || PGmcStarVocalic _,
    {*b} -> {*β} || [{*l} | {*r}] _
] .o. [
    {*β} -> {*b} || _ {*b}
]);
\end{ReaderFacingFoma}
\end{minipage}

Realization before \hyperref[rule-OENonHighSyn]{SC114 OENonHighSyn}
gives \Pred{reġnfoga} instead of \emph{reġnboga} `rainbow'; realization
afterward preserves the selected output. This discriminates the
representations, not an independently dated new lenition law or
compound-specific exemption.

\section{Early medial realization of unstressed
shortening}\label{early-medial-realization-of-unstressed-shortening}

\subsection{\texorpdfstring{\CAPRRuleHeading{SC117. Medial component of SC072 unstressed shortening}{OEMedLongShort}}{}}\label{rule-OEMedLongShort}

\begin{ReaderFacingFoma}
define OEMedLongShort FootCtx([
    {*ī} -> {*i} || EnglishStarVocalic OEAnyConsonant+ _ OEAnyConsonant+ [EnglishStarVocalic|{*ô}],
    {*ū} -> {*u} || EnglishStarVocalic OEAnyConsonant+ _ OEAnyConsonant+ [EnglishStarVocalic|{*ô}],
    {*ȳ} -> {*y} || EnglishStarVocalic OEAnyConsonant+ _ OEAnyConsonant+ [EnglishStarVocalic|{*ô}],
    {*ā} -> {*a} || EnglishStarVocalic OEAnyConsonant+ _ OEAnyConsonant+ [EnglishStarVocalic|{*ô}],
    {*ē} -> {*e} || EnglishStarVocalic OEAnyConsonant+ _ OEAnyConsonant+ [EnglishStarVocalic|{*ô}],
    {*ǣ} -> {*æ} || EnglishStarVocalic OEAnyConsonant+ _ OEAnyConsonant+ [EnglishStarVocalic|{*ô}],
    {*ō} -> {*o} || EnglishStarVocalic OEAnyConsonant+ _ OEAnyConsonant+ [EnglishStarVocalic|{*ô}]
]);
\end{ReaderFacingFoma}

This separately visible component belongs to
\hyperref[rule-OEUnstressedLongVowelShortening]{SC072
OEUnstressedLongVowelShortening}. It realizes the medial quantity
required before \hyperref[rule-OEHighVowelApocope]{SC063
OEHighVowelApocope}, without prematurely shortening final long vowels.
Goering distinguishes later \emph{*nētīn}, yielding \emph{nēten}
`livestock', from the medially shortened inflected forms that yield
\emph{nētenu} `livestock' and \emph{nētna} `of livestock'
(\citeproc{ref-Goering2016}{Goering 2016, 186--88};
\citeproc{ref-Goering2023}{2023, 67--68}). The same quantity distinction
is required in the earlier comparison of unstressed vowels
(\citeproc{ref-Fulk2018}{Fulk 2018, 90--96}, §§5.6--5.7). The later
singular's final consonant is paradigmatically preserved or restored,
not a survival through early final-n loss; the source does not fix the
restoration's date (\citeproc{ref-Campbell1959}{Campbell 1959, 189},
§473).

The root and transmitted prominence are protected. Final long vowels
remain for the later component of shortening; this early component is
not an extra high-vowel deletion or a morphological condition. The
separately traced identifier records implementation ownership, not a new
independently dated historical sound change.

The selected oblique \emph{fyrhte} `fright' supplies a final-quantity
control. Its inherited \emph{*fúrxtīnaz} exposes final \emph{*ī} after
early ending and final-\emph{*n} loss; that vowel remains long here
(\citeproc{ref-RingeTaylor2014}{Ringe and Taylor 2014, 87--88,
380--81}). Premature shortening makes it a target of
\hyperref[rule-OEHighVowelApocope]{SC063 OEHighVowelApocope}, wrongly
giving \Pred{fyrht}. The later component, not this medial realization,
licenses the regular final vowel of \emph{fyrhte}.

The inherited neuter plural \emph{*nautīnō} provides the complementary
positive dependency: West Saxon \emph{nīetenu} `animals, livestock'
requires medial shortening before quantitative deletion
(\citeproc{ref-Fulk2018}{Fulk 2018, 92}, §5.6;
\citeproc{ref-Sweet1871}{Sweet 1871, 109}, lines 8--11;
\citeproc{ref-Goering2016}{Goering 2016, 186--88}). Shortening produces
HLL \emph{*nīetinu}, retaining both high vowels; omission wrongly gives
\Pred{nīeten}. The selected \emph{lȳtelu} `little' already has a short
medial vowel and is unchanged by the same omission. This full inherited
plural protects the dependency without relying on a singular with
restored final \emph{n} or inventing another sound law.

\section{Quantitative high-vowel
deletion}\label{quantitative-high-vowel-deletion}

\subsection{Historical discussion}\label{historical-discussion-49}

Final and medial short high vowels undergo one quantitative deletion.
Campbell, Hogg, Ringe and Taylor, and Fulk agree that the relevant loss
follows mutation, but disagree over trisyllables and surrounding syncope
(\citeproc{ref-Campbell1959}{Campbell 1959, 144--46}, §§345--353;
\citeproc{ref-HoggPhonology1992}{Hogg 1992b, 121};
\citeproc{ref-RingeTaylor2014}{Ringe and Taylor 2014, 264--93}, §§6.7.3,
6.8.1; \citeproc{ref-Fulk2018}{Fulk 2018, 91}, §5.6). Barrack separately
presents final deletion as rule 345 and medial syncope as rule 351, with
progressive simplification of the former environment
(\citeproc{ref-Barrack1975}{Barrack 1975, 23--25, 85--87}). Those
descriptions remain comparative evidence, not proof of one particular
prehistoric serialization.

The adopted account follows Fulk's conservative Mercian evidence and
Goering's quantitative analysis. A heavy first syllable followed by two
light syllables preserves both high vowels: \emph{hēafudu} `heads',
\emph{lȳtelu} `little' and \emph{īdelu} `idle'. A heavy following
syllable instead licenses medial loss, as in \emph{nētna} `of
livestock'. The distinction is independently motivated by quantity and
metrical grouping, not inferred from the selected corpus's match count
(\citeproc{ref-Fulk2010}{Fulk 2010, 129--34, 140--42};
\citeproc{ref-Goering2016}{Goering 2016, 186--88};
\citeproc{ref-Goering2023}{2023, 67--68}). Early medial shortening,
realized in \hyperref[rule-OEMedLongShort]{SC117 OEMedLongShort}, gives
short medial \emph{*i} in the inflected livestock forms while final
\emph{*ī} remains long in the later compared state \emph{*nētīn}. Its
final consonant is not an uninterrupted survival through early final-n
loss: Campbell attributes it to the inflected forms, without fixing the
restoration's precise date (\citeproc{ref-Campbell1959}{Campbell 1959,
189}, §473). Later shortening yields \emph{nēten} `livestock'; the
earlier final-n law does not restart here.

The alternatives remain substantive. Conventional apocope followed by
syncope gives \emph{*hēafdu}, whereas Ringe and Taylor's
syncope-before-apocope account gives \emph{*hēafd}, with restored
attested plurals attributed to analogy. Hogg also discusses a tentative
one-syllable-reduction account (\citeproc{ref-HoggGrammar2011}{Hogg
2011, 225--27}, §§6.22--6.25; \citeproc{ref-RingeTaylor2014}{Ringe and
Taylor 2014, 289--93}). Output equivalence on selected words does not
erase these contrasting histories.

The conservative neuter plural \emph{lȳtelu} is particularly
discriminating. The Vespasian Psalter has \emph{netenu lytelu 7 micel}
`small and great animals', glossing Latin \emph{animalia pusilla et
magna}; agreement fixes the neuter plural, while Fulk independently
establishes the West Mercian quantity (\citeproc{ref-Sweet1885}{Sweet
1885, 336}, Psalm 103:25; \citeproc{ref-Fulk2010}{Fulk 2010, 133--34};
\citeproc{ref-Goering2016}{Goering 2016, 176--77}). Orel's comparative
\emph{*lūtilaz} and the inherited strong neuter plural ending \emph{*-ō}
support \emph{*lūtilō}; Fulk also prints this complete ancestral plural
explicitly {[}Orel (\citeproc{ref-Orel2003}{2003}), p.~248; Fulk
(\citeproc{ref-Fulk2018}{2018}), p.~92, §5.6; p.~209, §9.2{]}. This
selected plural preserves both weak high vowels: indiscriminate
heavy-root syncope gives \Pred{lȳtlu}, and final apocope in HLL gives
\Pred{lȳtel}. Neither error is exposed by the selected singular
\emph{hēafod} `head'. The comparative root alternation remains
unexplained; retention is not a grammatical exception
(\citeproc{ref-Kroonen2013}{Kroonen 2013, 339}).

The selected West Saxon \emph{nīetenu} `animals, livestock' supplies a
complementary test. The Hatton \emph{Pastoral Care} has plural
\emph{nietenu}; Fulk prints the complete ancestral \emph{*nautīnō}, and
the dialectal comparison distinguishes West Saxon \emph{nīetenu} from
Mercian \emph{nētenu} `animals' (\citeproc{ref-Sweet1871}{Sweet 1871,
109}, lines 8--11; \citeproc{ref-Fulk2018}{Fulk 2018, 92}, §5.6;
\citeproc{ref-RingeTaylor2014}{Ringe and Taylor 2014, 298--99};
\citeproc{ref-Goering2016}{Goering 2016, 186}). The medial nasal is
never final in this plural and needs no restoration or exception. After
mutation, early medial shortening makes \emph{*nīetinu} HLL: both high
vowels survive. Omitting that shortening instead gives \Pred{nīeten},
while the short-medial \emph{lȳtelu} is unaffected. Indiscriminate
heavy-root medial deletion gives \Pred{nīetnu}. These independently
reconstructed plurals protect different aspects of the quantitative
account.

This loss is distinct from West Germanic weak-word deletion in
\hyperref[rule-PWGmcUnstressedWordFinalIApocope]{SC098
PWGmcUnstressedWordFinalIApocope}. The separately selected sentence
context remains effective; its boundary annotation is not a segment of
the reconstruction (\citeproc{ref-RingeTaylor2014}{Ringe and Taylor
2014, 57--58}). Later light-syllable pre-l loss in
\hyperref[rule-OELAdjacentSyncope]{SC066 OELAdjacentSyncope}, weak-vowel
unrounding, and technical stress decoding remain separate operations.

\subsection{\texorpdfstring{\CAPRRuleHeading{SC063. Quantitative deletion of short unstressed high vowels}{OEHighVowelApocope}}{}}\label{rule-OEHighVowelApocope}

\noindent\begin{minipage}{\linewidth}
\begin{ReaderFacingFoma}
define OEHighVowelApocope FootCtx([
    [{*i}|{*u}|{*y}|{*ų}] -> 0 ||
        .#. OEAnyConsonant* HvI HvP* _
        [.#. | OEAnyConsonant [HvL | HvS OEAnyConsonant [.#.|OEAnyConsonant]]]
] .o. [{*ᶜ} -> 0]);
\end{ReaderFacingFoma}
\end{minipage}

The auxiliary expressions distinguish short and long nuclei, heavy
syllables, the initial quantitative group and following permitted
groups. Deletion targets a short unstressed high vowel at a group
boundary; a medial target requires a following singleton onset and heavy
syllable. The first root and separately transmitted prominence remain
protected. The singleton-onset representation bounds this
implementation; it is not a general parser for arbitrary historical
syllabification.

Mutation must precede loss: \emph{cȳ} `cow' and \emph{brȳd} `bride'
require the high vowel to perform its umlautal work. The selected
\emph{fȳlþe} `filth' and \emph{spinle} `spindle' retain their
independently supported accusative cells; their nonhigh final vowels
produce the regular medial deletion. The comparative nominatives
\emph{*fȳleþu} and \emph{*spinnelu} retain their light medial syllables
in this account. Their attested alternatives and the genuine comparative
gemination problem remain separate from the selected accusatives
(\citeproc{ref-Fulk2010}{Fulk 2010, 140--42};
\citeproc{ref-Goering2016}{Goering 2016, 186--88};
\citeproc{ref-RingeTaylor2014}{Ringe and Taylor 2014, 289--93}).

No special final-u loss after \emph{*h} is retained. \emph{Feoh}
`cattle, wealth' continues an independently supported,
already-endingless a-stem transponent \emph{*feh}, not an inherited
u-stem forced through an additional deletion
(\citeproc{ref-RingeTaylor2014}{Ringe and Taylor 2014, 177}, n.~4).

\section{Former syncope anchor and comparative
evidence}\label{former-syncope-anchor-and-comparative-evidence}

\subsection{Historical discussion}\label{historical-discussion-50}

The conventional heavy/open account discussed below is retained as
comparative evidence. The adopted quantitative operation now belongs to
\hyperref[rule-OEHighVowelApocope]{SC063 OEHighVowelApocope}, which
preserves heavy--light--light forms such as \emph{hēafudu} `heads'. This
former medial component is an identity trace anchor, not a second
productive deletion (\citeproc{ref-Fulk2010}{Fulk 2010, 129--34};
\citeproc{ref-Goering2016}{Goering 2016, 186--88}).

Old English lost unstressed medial high vowels in open syllables after
heavy syllables. The handbooks state the law generally, as part of the
prehistoric reduction of unstressed vowels: Ringe and Taylor describe a
weight-conditioned loss of medial \emph{*i} and \emph{*u}, and Campbell
surveys related later developments. Hogg explicitly distinguishes
prehistoric losses from changes after the earliest texts
(\citeproc{ref-RingeTaylor2014}{Ringe and Taylor 2014, 264--67}, §6.7.3;
\citeproc{ref-Campbell1959}{Campbell 1959, 157--59}, §§388--393;
\citeproc{ref-HoggPhonology1992}{Hogg 1992b, 121};
\citeproc{ref-Fulk2018}{Fulk 2018, 91}, §5.6). The loss is English, not
common West Germanic: the continental cognates cited below keep their
medial vowels.

The clearest witnesses are the feminine abstracts in \emph{*-iþō}, which
uniformly show syncope: \emph{strengþu} `strength', \emph{hǣlþ}
`health', \emph{fȳlþ} `filth', \emph{þīefþ} `theft', \emph{wiermþu}
`warmth' (\citeproc{ref-Campbell1959}{Campbell 1959, 236}, §589.6;
\citeproc{ref-Bammesberger1990}{Bammesberger 1990, 117--18}, §4.5). The
corpus now tests the type through accusative \emph{fȳlþe}, from
\emph{*fūliþǭ}, a composite of Kroonen's base \emph{*fūla-} with the
abstract suffix (\citeproc{ref-Kroonen2013}{Kroonen 2013, 158}). The
medial vowel survives in Old Saxon
{fulitha}\index[iv]{13os@\ivlangheader{Old Saxon}{}!fulitha@\iventry{fulitha}{}}
`filth' and in Old High German
{fülida}\index[iv]{10ohg@\ivlangheader{Old High German}{}!fulida@\iventry{fülida}{}}
`filth' (\citeproc{ref-Toller1921}{Toller 1921, 276}).

The same word fixes a relative chronology. Had syncope preceded i-umlaut
(\hyperref[rule-OEIUmlaut]{SC055 OEIUmlaut}), the deleted \emph{*i}
could not have fronted the root vowel, and the selected result would
have been \emph{*fūlþe} rather than attested \emph{fȳlþe}; umlaut must
therefore precede this syncope. Campbell derives the endingless
nominative singular of these abstracts analogically after the \emph{lār}
type (\citeproc{ref-Campbell1959}{Campbell 1959, 236}, §589.6), whereas
the former cascade reached the same form through excess apocope.
Corrected heavy--light retention no longer licenses that endingless
regular nominative. Fulk independently supplies paradigm evidence for
analogical syncope in the nominatives of the abstract class
(\citeproc{ref-Fulk2010}{Fulk 2010, 140--42}). The selected accusative
\emph{fȳlþe} is attested as the object in \emph{Ðā fylde ādōn of þām
Godes temple}; ordinary feminine accusative inflection supplies
\emph{*-ǭ}, and its nonhigh ending avoids the disputed nominative
history (\citeproc{ref-Toller1921}{Toller 1921, 276, Hml.S25, 378}). The
earlier selection and the competing accounts remain comparison history,
not a second selected cell.

In Barrack's cascade the corresponding territory is covered by his
syncope rule 351 (\citeproc{ref-Barrack1975}{Barrack 1975, 87}), with
further medial losses at rules 388 and 390, the latter citing
\emph{meolc} `milk' (\citeproc{ref-Barrack1975}{Barrack 1975, 96}); his
rules are broader than the dental-conditioned statement formerly used
here and are not coextensive with it. Rule 347 concerns final verbal
territory, not an equivalent general medial-syncope law
(\citeproc{ref-Barrack1975}{Barrack 1975, 86}).

\subsection{Independent historical stage
controls}\label{independent-historical-stage-controls}

The regular high-vowel core loses a short unstressed high vowel in an
open medial syllable after a heavy syllable. Light-root retention and
the earlier light-root contractions must remain distinguished; the
latter are not explained away as optionality, grammatical conditioning
or unsupported analogy (\citeproc{ref-SieversBrunner1965}{Brunner 1965,
137--38}, §§167--168; \citeproc{ref-Campbell1959}{Campbell 1959,
145--46}, §§348--349; \citeproc{ref-HoggGrammar2011}{Hogg 2011,
222--25}, §§6.18--6.21).

Ringe and Taylor reconstruct pre-Old-English
\Recon{ūtumist}\index[iv]{08preoe@\ivlangheader{Pre-Old English / prehistoric English}{}!utumist@\iventry{*ūtumist}{}},
then
\Recon{ȳtymist}\index[iv]{08preoe@\ivlangheader{Pre-Old English / prehistoric English}{}!ytymist@\iventry{*ȳtymist}{}},
and attested
{\emph{ȳtmest}}\index[iv]{02oe@\ivlangheader{Old English}{West Saxon normalization unmarked}!ytmest@\iventry{ȳtmest}{}}
`furthest'. In their proposed history, mutation precedes the loss of the
medial high vowel; the final root alone does not establish that
ordering, since earlier loss of \emph{u} would leave suffixal \emph{i}
as a direct trigger. Skipping that loss, while preserving subsequent
unrounding and lowering, gives \Pred{ȳtemest}. That variant is also
attested: Ringe and Taylor explicitly analyse it by analogy with
\emph{yfemest} and similar forms. It fails the specified regular history
to \emph{ȳtmest}, not the entire Old English lexicon
(\citeproc{ref-RingeTaylor2014}{Ringe and Taylor 2014, 267}). The
transmitted mutation is independently described by Hogg and Campbell; it
is not inferred from the need to make this derivation succeed
(\citeproc{ref-HoggGrammar2011}{Hogg 2011, 120--21}, §5.76, p.~219,
§6.10; \citeproc{ref-Campbell1959}{Campbell 1959, 80--81}, §203).

Hogg's
\Recon{gaduling}\index[iv]{08preoe@\ivlangheader{Pre-Old English / prehistoric English}{}!gaduling@\iventry{*gaduling}{}}
gives
{\emph{gædeling}}\index[iv]{02oe@\ivlangheader{Old English}{West Saxon normalization unmarked}!gaedeling@\iventry{gædeling}{}}
`companion'. Its light first syllable preserves the medial high vowel
through syncope; mutation-produced \emph{y} then unrounds and lowers.
Unrestricted deletion predicts \Pred{gædling}; omitting weak unrounding
leaves \Pred{gædyling}. This supplies an independent later-stage
light-root retention control without presuming a new inherited compound
(\citeproc{ref-HoggGrammar2011}{Hogg 2011, 120--21}, §5.76, p.~219,
§6.10).

The bailiff reconstruction provides a second, strictly stage-level
control:
\Recon{skuldihātijā}\index[iv]{08preoe@\ivlangheader{Pre-Old English / prehistoric English}{}!skuldihatija@\iventry{*skuldihātijā}{}}
becomes
\Recon{sċyldihǣtjā}\index[iv]{08preoe@\ivlangheader{Pre-Old English / prehistoric English}{}!scyldihatja@\iventry{*sċyldihǣtjā}{}},
with medial \emph{i} lost in
{\emph{sċyldhǣta}}\index[iv]{02oe@\ivlangheader{Old English}{West Saxon normalization unmarked}!scyldhaeta@\iventry{sċyldhǣta}{}}
`bailiff' (\citeproc{ref-RingeTaylor2014}{Ringe and Taylor 2014, 267}).
The syncope component correctly deletes that \emph{i}. Its complete
suffix derivation is not certified: subordinate prominence on the long
second root is not fully represented, and the current model shortens its
vowel. This limitation must not be concealed by substituting a
convenient short reconstruction or adding a lexical repair.

The inherited
\Recon{wiđuwōn}\index[iv]{03pgmc@\ivlangheader{Proto-Germanic}{}!widuwon@\iventry{*wiđuwōn}{}}
and Old English
{\emph{widuwe}}\index[iv]{02oe@\ivlangheader{Old English}{West Saxon normalization unmarked}!widuwe@\iventry{widuwe}{}}
`widow' independently corroborate the retained medial vowel
(\citeproc{ref-Orel2003}{Orel 2003, 462}). The selected root-vowel
variant
{\emph{wuduwe}}\index[iv]{02oe@\ivlangheader{Old English}{West Saxon normalization unmarked}!wuduwe@\iventry{wuduwe}{}}
retains that same \emph{u}. Removing the weight restriction predicts
\Pred{wudwe}. This is a complete selected negative control for the
general high-vowel component, not a complete discriminator of the former
early compound-only formulation. Its root-vowel history remains with
\hyperref[rule-OEWICombinativeUUmlaut]{SC039 OEWICombinativeUUmlaut}.

\subsection{\texorpdfstring{SC065. Former general-syncope trace anchor
(\texttt{OEMedialSyncope})}{SC065. Former general-syncope trace anchor (OEMedialSyncope)}}\label{rule-OEMedialSyncope}

\noindent\begin{minipage}{\linewidth}
\begin{ReaderFacingFoma}
define OEMedialSyncope ?*;
\end{ReaderFacingFoma}
\end{minipage}

The preceding consolidation replaced the dental enumeration and the
heavy part of \hyperref[rule-OELAdjacentSyncope]{SC066
OELAdjacentSyncope} with a single conventional general operation. Under
that account, short unstressed \emph{*i/u/y} disappears after a heavy
first syllable when a following singleton consonant and vowel establish
an open medial syllable. The retained quantity carrier \emph{*ô} is
included locally as a following vowel; it was omitted from the shared
ordinary-vowel class. This is a bounded onset domain, not a claim to
complete historical syllabification of arbitrary clusters. Light-root
\emph{wuduwe} `widow', final-closed \emph{cetel} `kettle' and
\emph{engel} `angel' environments do not meet the deletion condition.

Earlier controlled experiments found that an unrestricted extension to
\emph{*u} syncopated \emph{hēafod} `head', dropping the dental
restriction destroyed \emph{hierfest} `harvest' and \emph{sċilling}
`shilling', and the open-syllable formulation was vacuous after the old
excess apocope, while moving it earlier reduced the former \emph{sāwol}
`soul' selection. These measurements remain the history of the narrow
proxy, not arguments against the corrected general law. Independently
selected vocalic accusatives now give \emph{fȳlþe} `filth' and
\emph{spinle} `spindle'; the corrected apocope also gives conventional
\emph{hēafdu} `heads'. The adopted quantitative operation likewise
requires neither dental nor lateral identity, but does not repeat the
conventional heavy/open loss here. Mutation must still precede deletion;
syncope feeds the subsequent preconsonantal degemination in
\emph{spinle}.

The selected \emph{spinle} `spindle' is the strong feminine accusative
\emph{*spénnilǭ}, not the former endingless nominative. Clark Hall gives
\emph{spinel}, feminine, with genitive \emph{spinle}; Orel reconstructs
\emph{*spennilō(n)} (\citeproc{ref-ClarkHall1960}{Clark Hall 1960, 272};
\citeproc{ref-Orel2003}{Orel 2003, 364}). Toller's tool enumeration
reads \emph{He sceal habban fela towtōla spinle, reól, gearnwindan}
(\citeproc{ref-Toller1921}{Toller 1921, 288, s.v}.~gearn-winde, Angl. ix
263, 11). The quantifier \emph{fela} governs genitive \emph{towtōla};
the subsequently enumerated implements are interpreted as objects of
\emph{habban}, with singular strong feminine \emph{spinle}. This is a
syntactic interpretation of the quotation, not a dictionary's explicit
case label. Ordinary accusative \emph{*-ǭ} supplies a complete
reconstructed cell. The nominative and the competing
syncope-before-apocope history remain comparison evidence rather than
another selected derivation.

\section{Early weak front-vowel
unrounding}\label{early-weak-front-vowel-unrounding}

\subsection{Historical discussion}\label{historical-discussion-51}

Mutation-produced weak \emph{*y} unrounds to \emph{*i}; retained
\emph{*i} later lowers to \emph{*e}. These are quality changes, not
syncope. Hogg places the unrounding at a very early prehistoric stage
and illustrates transmitted mutation with \emph{*gaduling > gædeling}
(\citeproc{ref-HoggGrammar2011}{Hogg 2011, 219}, §6.10 and note 1). The
early component belongs to the same weak-front-vowel history whose later
lowering is \hyperref[rule-OEMedUnstressedILowering1]{SC074
OEMedUnstressedILowering1}. Its stable identifier is retained for
scholarly cross-references; it no longer implements high-vowel deletion.

\subsection{\texorpdfstring{\CAPRRuleHeading{SC115. Early weak-front-vowel unrounding component}{OEHighLinkSyncope}}{}}\label{rule-OEHighLinkSyncope}

\noindent\begin{minipage}{\linewidth}
\begin{ReaderFacingFoma}
define OEHighLinkSyncope FootCtx([
    {*y} -> {*i} || EnglishStarVocalic OEAnyConsonant+ _
]);
\end{ReaderFacingFoma}
\end{minipage}

The component excludes initial-root \emph{*y} and still-marked
subordinate \emph{*ỳ}. It is computed after
\hyperref[rule-OEMedialSyncope]{SC065 OEMedialSyncope} and before
\hyperref[rule-OELAdjacentSyncope]{SC066 OELAdjacentSyncope}. Unrounding
preserves height and commutes with the bounded general high-vowel loss;
the serialization does not date unrounding after syncope. The \emph{d/l}
protection in the later light component permits \emph{gædeling}
`companion' on either side of that component, but this particular
comparison does not establish universal commutation.

The separate \hyperref[rule-OEHighStressDecode]{SC116
OEHighStressDecode} helper runs after the light pre-\emph{*l} component,
not at this unrounding position. It preserves secondary-vowel protection
until both deletion components have finished. Neither technical decoding
nor this component's serialized location dates disappearance of
historical secondary stress.

\section{L-adjacent syncope}\label{l-adjacent-syncope}

\subsection{Historical discussion}\label{historical-discussion-52}

Vowel loss creates the clusters upon which later assimilation and
degemination operate. Hogg and Ringe and Taylor describe this
dependence, while Brunner's \emph{netle} `nettle' beside later
\emph{netele} `nettle' supplies a concrete lexical type
(\citeproc{ref-HoggPhonology1992}{Hogg 1992b, 121};
\citeproc{ref-RingeTaylor2014}{Ringe and Taylor 2014, 264--96},
§§6.7.3--6.8.2; \citeproc{ref-SieversBrunner1965}{Brunner 1965,
131--33}, §§158--159). Fulk places this syncope after i-umlaut
(\citeproc{ref-Fulk2018}{Fulk 2018, 91}, §5.6).

The three relations are not equally secure. Lexical evidence orders
syncope and degemination; the intervening dental assimilation has no
independent ordering witness.

\subsection{\texorpdfstring{SC066. Later light-syllable pre-\emph{*l}
syncope
(\texttt{OELAdjacentSyncope})}{SC066. Later light-syllable pre- syncope (OELAdjacentSyncope)}}\label{rule-OELAdjacentSyncope}

\noindent\begin{minipage}{\linewidth}
\begin{ReaderFacingFoma}
define OELAdjacentSyncope Ctx([
    {*i} -> 0 ||
        [EnglishStarShortVowel|EnglishStarShortDiphthong]
        [OEAnyConsonant - [{*d}|{*θ}|{*f}]] _ {*l} EnglishStarVocalic
]);
\end{ReaderFacingFoma}
\end{minipage}

The loss of medial \emph{*i} before \emph{*l} is late enough to preserve
earlier umlaut, as \emph{netle} `nettle' demonstrates.

The former statement combined two historically distinct strata. Its
heavy part now belongs to the single quantitative deletion in
\hyperref[rule-OEHighVowelApocope]{SC063 OEHighVowelApocope}, not an
independently productive second syncope. The selected accusative
\emph{*spénnilǭ} `spindle' and the genitive \emph{*θístilas} `thistle'
lose their \emph{*i} after the heavy sequences \emph{-nn-} and
\emph{-st-} (\citeproc{ref-RingeTaylor2014}{Ringe and Taylor 2014,
270--72}, §6.7.3; \citeproc{ref-Campbell1959}{Campbell 1959, 158},
§392). The nonhigh, quantity-bearing endings distinguish these selected
histories from the HLL nominative comparisons;
\hyperref[rule-OEMedialSyncope]{SC065 OEMedialSyncope} remains an
identity anchor. After a light syllable, however, the handbooks describe
a later and more particular development: unstressed \emph{*i} next to
\emph{l} is lost even where the preceding syllable is light, as in
\emph{netle} `nettle' from the diminutive \Recon{nátilōn} `nettle'
(\citeproc{ref-Kroonen2013}{Kroonen 2013, 384};
\citeproc{ref-RingeTaylor2014}{Ringe and Taylor 2014, 272}, §6.7.4).
Luick, following Weyhe, restricts this light-syllable deletion to
trisyllabic forms whose middle syllable was open, dating it before the
literary period after dentals and to the eighth century after gutturals
(\citeproc{ref-Luick1914}{Luick 1914-\/-1940, § 336}); Brunner
likewise treats it as regular in the oldest sources after light
syllables (\citeproc{ref-SieversBrunner1965}{Brunner 1965, 132--33},
§159c); Campbell describes medial vowels as dropped after short
syllables when a group of consonant plus \emph{l} or \emph{r} would
result (\citeproc{ref-Campbell1959}{Campbell 1959, 157--58},
§§388--390).

The open-syllable condition carries the explanatory weight. \emph{Cetel}
`kettle', an early Latin loan (\Recon{katilaz} `kettle', from
\emph{catillus}), famously fails to syncopate
(\citeproc{ref-Orel2003}{Orel 2003, 211};
\citeproc{ref-RingeTaylor2014}{Ringe and Taylor 2014, 272}). It is not a
lexical exception. Its nominative \emph{*i} stands, after apocope, in a
word-final closed syllable with a short nucleus, but heavy by closure
--- a position the adopted light/open formulation does not reach. The
same structural point explains the nominative--oblique split that
Campbell records for the \emph{netel}/\emph{spinel} class, where the
nominative retains its vowel while the syllabically open oblique cases
syncopate (\citeproc{ref-Campbell1959}{Campbell 1959, 235--36}, §589).
The current component contains only the later light-syllable loss.

The selected \emph{sċytel} `dart' supplies a complete inherited
closed-final control. Orel reconstructs \emph{*skutilaz}, beside Old
Norse \emph{skutill}; the Corpus Glossary prints \emph{scytel}, with
retained medial \emph{e} (\citeproc{ref-Orel2003}{Orel 2003, 346--47};
\citeproc{ref-Sweet1885}{Sweet 1885, 79}, Corpus 1325). The
corresponding Épinal and Erfurt glosses have \emph{scytil}. Their Latin
\emph{momentum} denotes the moving beam or lever of a balance, not an
unqualified dart or a balance pointer {[}Pheifer
(\citeproc{ref-Pheifer1974}{1974}), p.~34, gloss 632; p.~101, commentary
632{]}. The palatal dot is editorial, and the Corpus form is not
relabelled West Saxon. After apocope and mutation, \emph{*sċytil} has
the closed final syllable \emph{-til}. Its vowel survives this law and
later lowers to \emph{e}. Removing the following-vowel restriction from
the adopted operation wrongly gives \Pred{sċytl}. This tests the
open-medial domain, without forcing an unverified West Saxon spelling of
\emph{cetel} into the selected population.

Brunner's original passage excludes the environments between \emph{d}
and \emph{l}, and between inherited Germanic \emph{þ/f} and \emph{l};
the \emph{f} descended from earlier \emph{bh} is explicitly
distinguished (\citeproc{ref-SieversBrunner1965}{Brunner 1965, 132},
§159c, §192.2). CAPR excludes the corresponding still-distinct
\emph{*d/θ/f} segments, not words or grammatical classes.
Source-supported \emph{*gaduling > gædeling} remains a retention control
even when mutation-produced weak \emph{*y} has already unrounded to
\emph{*i}.

The sources disagree about the residue of variation. Weyhe, Luick, and
Brunner treat the law as categorical within its consonant-class and
dialect conditions, explaining late West Saxon \emph{netele} as
anaptyxis, a secondary vowel reinserted in the cluster
(\citeproc{ref-Luick1914}{Luick 1914-\/-1940, § 336};
\citeproc{ref-SieversBrunner1965}{Brunner 1965, 133}, §159 Anm. 3).
Ringe and Taylor instead read pairs such as \emph{crypel} beside dative
\emph{cryple} as a change arrested before completion, leaving stable
variation (\citeproc{ref-RingeTaylor2014}{Ringe and Taylor 2014, 272}).
The selected positive \emph{netle} does not discriminate between these
readings of open-medial variation; both accounts are therefore recorded
without a forced choice. The closed-final \emph{sċytel} control protects
the distinct open-versus-closed domain of the adopted categorical law.

Barrack reaches the same material through his framework of progressive
environment simplification. His rule for medial high-vowel deletion in
open syllables after long accented syllables corresponds to the
heavy-syllable stratum, and in West Saxon he drops the length condition,
generalizing the deletion to light syllables
(\citeproc{ref-Barrack1975}{Barrack 1975, 87}). A separate rule deletes
an unstressed vowel before a liquid or nasal followed by a vowel, and he
reports that West Saxon simplified the environment further by dropping
the following-vowel condition entirely
(\citeproc{ref-Barrack1975}{Barrack 1975, 96}). The present law retains
the following-vowel condition for the light stratum: dropping it, as
Barrack's West Saxon formulation would, predicts syncope in the
\emph{cetel} nominative. His neighbouring rule deleting a vowel after a
liquid, which derives \emph{meolc} `milk' from \Recon{meluk} `milk'
(\citeproc{ref-Barrack1975}{Barrack 1975, 96}), is likewise not adopted:
West Saxon \emph{meoloc} retains the vowel, and medial \emph{*u} next to
\emph{l} is in general kept far longer than \emph{*i}
(\citeproc{ref-Luick1914}{Luick 1914-\/-1940, § 336.2};
\citeproc{ref-RingeTaylor2014}{Ringe and Taylor 2014, 273}).

In the preserved experiments on the former combined statement, placement
before i-umlaut makes PGmc \Recon{nátilōn} `nettle' yield \Pred{nætle}
rather than expected OE \emph{netle} `nettle', and PGmc \Recon{spénnilō}
`spindle' yields \Pred{spenl} rather than expected \emph{spinl}
`spindle'. Placed after preconsonantal degemination, PGmc
\Recon{spénnilō} yields \Pred{spinnl} rather than expected \emph{spinl}.
Those earlier experiments remain evidence for the former combined
statement. The spindle feeding relation now belongs to quantitative
\hyperref[rule-OEHighVowelApocope]{SC063 OEHighVowelApocope}, not to
this light-only component. The nettle mutation relation remains with
\hyperref[rule-OELAdjacentSyncope]{SC066 OELAdjacentSyncope}; no
independent chronology is invented for the intervening dental
assimilation.

\section{Secondary high-vowel annotation
decoding}\label{secondary-high-vowel-annotation-decoding}

This independently traced support operation decodes retained subordinate
high-vowel marks only after their phonological consumers have finished.
It is a computational identity, not an additional historical sound
change.

\subsection{\texorpdfstring{\CAPRRuleHeading{SC116. Technical secondary high-marker decoding}{OEHighStressDecode}}{}}\label{rule-OEHighStressDecode}

\noindent\begin{minipage}{\linewidth}
\begin{ReaderFacingFoma}
define OEHighStressDecode Ctx([{*ì} -> {*i}, {*ù} -> {*u}, {*ỳ} -> {*y}]);
\end{ReaderFacingFoma}
\end{minipage}

The decoder follows \hyperref[rule-OELAdjacentSyncope]{SC066
OELAdjacentSyncope} and precedes later weak-vowel lowering. Marked
subordinate vowels are protected from
\hyperref[rule-OEMedialSyncope]{SC065 OEMedialSyncope}, early weak-vowel
unrounding at \hyperref[rule-OEHighLinkSyncope]{SC115
OEHighLinkSyncope}, and the light-syllable deletion component before
these marks are removed. No historical date for obligatory loss of
subordinate stress follows from this serialization.

\section{Dental assimilation}\label{dental-assimilation}

\subsection{\texorpdfstring{\CAPRRuleHeading{SC067. Dental assimilation in newly formed clusters}{OEDentalAssimilation}}{}}\label{rule-OEDentalAssimilation}

\noindent\begin{minipage}{\linewidth}
\begin{ReaderFacingFoma}
define OEDentalAssimilation Ctx([
    {*θ} -> 0 || {*t} _
]);
\end{ReaderFacingFoma}
\end{minipage}

Loss of \emph{*θ} after \emph{*t} resolves a dental cluster produced by
syncope (\citeproc{ref-HoggPhonology1992}{Hogg 1992b, 121};
\citeproc{ref-RingeTaylor2014}{Ringe and Taylor 2014, 279--96}, §§6.7.5,
6.8.2). No witness distinguishes its position: moving dental
assimilation across every tested neighbor leaves the outputs unchanged.
I nevertheless place it after syncope, which supplies its input, and
before the more general cluster simplification described in the
handbooks. This order is phonologically motivated, not established by a
lexical contrast.

\section{Preconsonantal degemination}\label{preconsonantal-degemination}

\subsection{\texorpdfstring{\CAPRRuleHeading{SC068. Preconsonantal degemination before sonorants}{OEPreconsonantalDegemination}}{}}\label{rule-OEPreconsonantalDegemination}

\begin{ReaderFacingFoma}
define OEPreconsonantalDegemination Ctx(OEPreconsonantalDegemTT .o. OEPreconsonantalDegemNN);
\end{ReaderFacingFoma}

Preconsonantal \emph{*tt} and \emph{*nn} simplify only after syncope has
created a following sonorant cluster, as in \emph{spinle} `spindle'
(\citeproc{ref-RingeTaylor2014}{Ringe and Taylor 2014, 279--96},
§§6.7.5, 6.8.2).

The historical change is a general one: Ringe and Taylor state simply
that geminates were degeminated next to another consonant
(\citeproc{ref-RingeTaylor2014}{Ringe and Taylor 2014, 296}, §6.8.2).
Barrack describes the same territory twice over, once as a degemination
of long consonants after a vowel before another consonant, and once as a
late and general simplification of triple clusters in both directions,
which he orders after his syncope rules
(\citeproc{ref-Barrack1975}{Barrack 1975, 115, 122}). The present law
states the narrower slice of this process that the witnesses attest:
\emph{*tt} before a nasal in \emph{botm} `bottom' and \emph{*nn} before
\emph{l} in \emph{spinle} `spindle'. The two clauses are segmental cases
of one change, not two historical developments, and no witness in the
corpus licenses the broader formulation.

The two witnesses reach the law by different routes. For \emph{spinle}
the geminate cluster exists only because SC065 Medial high-vowel syncope
has created it; SC066 Light-syllable syncope before \emph{l} does not
delete this heavy-root vowel. For \emph{botm} the cluster \emph{*ttm} is
already present in the stem. Kroonen reconstructs the word with an
inherited geminate, \Recon{buttman-} beside \Recon{budman-}, deriving OE
\emph{botm} directly from the geminated variant
(\citeproc{ref-Kroonen2013}{Kroonen 2013, 82}). Ringe and Taylor instead
derive \emph{botm} from \Recon{buþm-} with a later West Saxon
development of \emph{þm} to \emph{tm}, while observing that the
prehistory of this word cannot be recovered in detail
(\citeproc{ref-RingeTaylor2014}{Ringe and Taylor 2014, 332, 354}). The
derivation here follows the geminate reconstruction; the alternative
route is recorded as an open question of reconstruction selection, not
of the sound law itself.

Placed before SC065 Medial high-vowel syncope, the selected accusative
\emph{*spénnilǭ} `spindle' yields \Pred{spinnle} rather than expected OE
\emph{spinle}. Syncope must therefore create the cluster before
degemination simplifies it. The earlier nominative comparison
\emph{*spénnilō} gave \Pred{spinnl} against \emph{spinl} under the
former narrower syncope clauses; that comparison does not establish a
separate heavy-root \emph{l}-adjacent law. \emph{Botm} does not order
the pair, since its cluster does not depend on syncope. Reordering
degemination against any tested later change leaves the witness
unchanged, so no terminus ante quem is known.

\section{Early o-shortening}\label{early-o-shortening}

\subsection{Historical discussion}\label{historical-discussion-53}

After the principal palatal and umlautal changes, unstressed vowels
undergo shortening, fronting, merger, and sometimes complete loss.
Campbell describes the early shortening of unaccented long vowels, while
Hogg, Ringe and Taylor, and Fulk relate it to apocope, syncope, and the
later reductions (\citeproc{ref-Campbell1959}{Campbell 1959, 148}, §355;
\citeproc{ref-HoggPhonology1992}{Hogg 1992b, 121};
\citeproc{ref-RingeTaylor2014}{Ringe and Taylor 2014, 298--314},
§§6.8.3--6.9.3; \citeproc{ref-Fulk2018}{Fulk 2018, 90--96}, §§5.6--5.7).

Early o-shortening has only a distant earlier boundary. The rules that
follow, especially \hyperref[rule-OEUnstressedFrontingEarly]{SC070
OEUnstressedFrontingEarly} and
\hyperref[rule-OEUnstressedLongVowelShortening]{SC072
OEUnstressedLongVowelShortening}, have more closely defined relations.

\subsection{\texorpdfstring{\CAPRRuleHeading{SC069. Early shortening of unstressed \emph{*ō} before nasals}{OEEarlyOShortening}}{}}\label{rule-OEEarlyOShortening}

\begin{ReaderFacingFoma}
define OEEarlyOShortening Ctx([
    # EARLY: *ō → *a before nasals (runs BEFORE fronting)
    # This *a will undergo fronting: *a → *æ → e (Campbell §333, §369)
    {*ō} -> {*a} || EnglishStarVocalic [EnglishStarConsonant | EnglishPalatalConsonant]+ _ EnglishStarNasal
]);
\end{ReaderFacingFoma}

The rule shortens unstressed long \emph{*ō} before a following nasal.
Because this shortening happens early, the resulting \emph{*a} can still
participate in the later fronting and merger that shape many weak final
syllables.

Moving the rule before \hyperref[rule-PNWGmcNStemNLoss]{SC023
PNWGmcNStemNLoss}, PGmc \Recon{nḗdrōn} `adder' yields \Pred{nǣdran}
rather than expected OE \emph{nǣdre} `adder', PGmc \Recon{érθōn} `earth'
yields \Pred{eorþan} rather than expected \emph{eorþe} `earth', and PGmc
\Recon{fláskōn} `flask' yields \Pred{flascan} rather than expected
\emph{flasce} `flask'. The same earlier shift also disrupts forms such
as \emph{heorte} `heart' and \emph{līne} `line'. This broad set of
failures requires \hyperref[rule-OEEarlyOShortening]{SC069
OEEarlyOShortening} to follow \hyperref[rule-PNWGmcNStemNLoss]{SC023
PNWGmcNStemNLoss}.

If the rule is moved later within the tested sequence, no output differs
from the expected one. The lexical evidence therefore does not identify
a corresponding later constraint. The sources place early
\emph{*ō}-shortening before the later weak-tail changes without fixing a
closer local order.

\section{Early unstressed fronting}\label{early-unstressed-fronting}

\subsection{Historical discussion}\label{historical-discussion-54}

Campbell distinguishes the shortening of unaccented long vowels, while
Hogg, Ringe and Taylor, and Fulk place fronting and shortening within a
later history of syncope and final-vowel adjustment
(\citeproc{ref-Campbell1959}{Campbell 1959, 148}, §355;
\citeproc{ref-HoggPhonology1992}{Hogg 1992b, 121};
\citeproc{ref-RingeTaylor2014}{Ringe and Taylor 2014, 298--314},
§§6.8.3--6.9.3; \citeproc{ref-Fulk2018}{Fulk 2018, 90--96}, §§5.6--5.7).
Earlier unstressed fronting precedes later o-shortening.

\hyperref[rule-OEUnstressedFrontingEarly]{SC070
OEUnstressedFrontingEarly} has both an earlier and a later lexical
breakpoint. \hyperref[rule-OELateOShortening]{SC071 OELateOShortening}
confirms their reciprocal order, but no lexical evidence fixes its later
boundary.

\subsection{\texorpdfstring{\CAPRRuleHeading{SC070. Early fronting of unstressed \emph{*a}}{OEUnstressedFrontingEarly}}{}}\label{rule-OEUnstressedFrontingEarly}

\noindent\begin{minipage}{\linewidth}
\begin{ReaderFacingFoma}
define OEUnstressedFrontingEarly FootCtx(OEUnstressedAFronting);
\end{ReaderFacingFoma}
\end{minipage}

The rule fronts unstressed \emph{*a} to \emph{*æ} after the earlier
shortening has created a frontable vowel but before the later shortening
of unstressed \emph{*ō}. It produces endings such as OE \emph{-en} in
\emph{lungen} `lungs'.

If the rule is moved before \hyperref[rule-OEVelarPalatalization]{SC052
OEVelarPalatalization}, PGmc \Recon{lúnganjō} `lungs' yields
\Pred{lunġen} rather than expected OE \emph{lungen} `lungs'. If the rule
is delayed until after \hyperref[rule-OELateOShortening]{SC071
OELateOShortening}, PGmc \Recon{búrōθi} `bears' yields \Pred{boreþ}
rather than expected OE \emph{boraþ} `bears', and PGmc \Recon{mḗnōθz}
`month' yields \Pred{mōneþ} rather than expected \emph{mōnaþ} `month'.
The witness forms require
\hyperref[rule-OEUnstressedFrontingEarly]{SC070
OEUnstressedFrontingEarly} to follow
\hyperref[rule-OEVelarPalatalization]{SC052 OEVelarPalatalization} and
precede \hyperref[rule-OELateOShortening]{SC071 OELateOShortening}.

The relation to \hyperref[rule-OELateOShortening]{SC071
OELateOShortening} is local. The earlier boundary at
\hyperref[rule-OEVelarPalatalization]{SC052 OEVelarPalatalization}
places fronting after the older palatal developments.

\section{Late o-shortening}\label{late-o-shortening}

\subsection{\texorpdfstring{SC071. Later shortening of unstressed
\emph{*ō}
(\texttt{OELateOShortening})}{SC071. Later shortening of unstressed  (OELateOShortening)}}\label{rule-OELateOShortening}

The following rule handles the later shortening stage.

\begin{ReaderFacingFoma}
define OELateOShortening Ctx([
    # LATE: *ō → *o elsewhere (runs AFTER fronting; Campbell §355: the late
    # product is "of too late origin to become æ").
    # Stausland Johnsen 2015:28-30: unstressed long vowels shorten regularly in
    # pre-OE AFTER the loss of short unstressed vowels (the *bursti-/*burstīz
    # argument). The shortened vowel is a plain short *o — the ONLY source of
    # unstressed *o at this stage (the earliest OE texts have just two
    # unstressed back vowels, u and a: Stausland Johnsen 2015:21; classical-OE
    # unstressed o arises only from the later 9th-c. lowering, SC040). Its
    # split — u in medial syllables, a in final syllables — is carried out by
    # OEMedUnstressedORaising (SC099) and OEFinalUnstressedOLowering (SC100)
    # immediately below, replacing the old one-step *ō → *a compression that
    # hard-coded the final-syllable outcome for all positions (van Helten's
    # refuted conditioning aside, medial *ō would have surfaced wrongly as a:
    # **wundade).
    {*ō} -> {*o} || EnglishStarVocalic [EnglishStarConsonant | EnglishPalatalConsonant]+ _ [EnglishStarConsonant | EnglishPalatalConsonant]*
]);
\end{ReaderFacingFoma}

The rule shortens the remaining unstressed long \emph{*ō} after
fronting. The shortened vowel is then resolved by the following
medial/final distribution, not directly as \emph{a}
(\citeproc{ref-StauslandJohnsen2015}{Stausland Johnsen 2015, 28--31}).

Moving the rule before \hyperref[rule-OEUnstressedFrontingEarly]{SC070
OEUnstressedFrontingEarly} makes PGmc \Recon{búrōθi} `bears' yield
\Pred{boreþ} rather than expected OE \emph{boraþ} `bears'. The contrast
requires \hyperref[rule-OELateOShortening]{SC071 OELateOShortening} to
follow \hyperref[rule-OEUnstressedFrontingEarly]{SC070
OEUnstressedFrontingEarly}.

\section{Medial raising of shortened unstressed
o}\label{medial-raising-of-shortened-unstressed-o}

\subsection{\texorpdfstring{\CAPRRuleHeading{SC099. Medial raising of shortened unstressed \emph{*o}}{OEMedUnstressedORaising}}{}}\label{rule-OEMedUnstressedORaising}

\begin{ReaderFacingFoma}
define OEMedUnstressedORaising Ctx([
    {*o} -> {*u} || EnglishStarVocalic [EnglishStarConsonant | EnglishPalatalConsonant]+ _ [EnglishStarConsonant | EnglishPalatalConsonant]* EnglishStarVocalic
]);
\end{ReaderFacingFoma}

After \hyperref[rule-OELateOShortening]{SC071 OELateOShortening}, the
shortened vowel gives \emph{u} in an unstressed medial syllable. The
rule encodes Stausland Johnsen's statistically supported account of West
Saxon ō-verb pasts, not a general rule for inherited short \emph{*o} or
for nominal morphology (\citeproc{ref-StauslandJohnsen2015}{Stausland
Johnsen 2015, 28--31, 36}). His diagnostic derivation is PGmc
\Recon{wúndōdē} `wounded' \textgreater{} \Pred{wundode} \textgreater{}
OE
{\emph{wundude}}\index[iv]{02oe@\ivlangheader{Old English}{West Saxon normalization unmarked}!wundude@\iventry{wundude}{}}
`wounded' (\citeproc{ref-StauslandJohnsen2015}{Stausland Johnsen 2015,
28--29}).

\section{Final lowering of shortened unstressed
o}\label{final-lowering-of-shortened-unstressed-o}

\subsection{\texorpdfstring{\CAPRRuleHeading{SC100. Final lowering of shortened unstressed \emph{*o}}{OEFinalUnstressedOLowering}}{}}\label{rule-OEFinalUnstressedOLowering}

\begin{ReaderFacingFoma}
define OEFinalUnstressedOLowering Ctx([
    {*o} -> {*a} || EnglishStarVocalic [EnglishStarConsonant | EnglishPalatalConsonant]+ _ [EnglishStarConsonant | EnglishPalatalConsonant]* .#.
]);
\end{ReaderFacingFoma}

In a final syllable the same shortened vowel gives \emph{a}. Thus the
existing month control continues PGmc \Recon{mḗnōθz} `month' through
shortened \emph{*o} to OE
{\emph{mōnaþ}}\index[iv]{02oe@\ivlangheader{Old English}{West Saxon normalization unmarked}!monath@\iventry{mōnaþ}{}}
`month', while \Recon{wúndōdē} `wounded' takes
\hyperref[rule-OEMedUnstressedORaising]{SC099 OEMedUnstressedORaising}
instead. The medial/final contrast and its chronology after long-vowel
shortening are Stausland Johnsen's analysis
(\citeproc{ref-StauslandJohnsen2015}{Stausland Johnsen 2015, 28--31}).

\section{Unstressed long-vowel
shortening}\label{unstressed-long-vowel-shortening}

\subsection{Historical discussion}\label{historical-discussion-55}

Campbell describes the shortening of unaccented long vowels, and Ringe
and Taylor place the later component among the last prehistoric Old
English changes before the merger of unstressed \emph{*æ} with \emph{*e}
(\citeproc{ref-Campbell1959}{Campbell 1959, 148}, §355;
\citeproc{ref-HoggPhonology1992}{Hogg 1992b, 121};
\citeproc{ref-RingeTaylor2014}{Ringe and Taylor 2014, 298--314},
§§6.8.3--6.9.3; \citeproc{ref-Fulk2018}{Fulk 2018, 90--96}, §§5.6--5.7).

\hyperref[rule-OEUnstressedLongVowelShortening]{SC072
OEUnstressedLongVowelShortening} and
\hyperref[rule-OEUnstressedAEMerger]{SC073 OEUnstressedAEMerger} have a
reciprocal ordering relation. \hyperref[rule-NWGmcInStemNLoss]{SC064
NWGmcInStemNLoss} supplies the earlier boundary of shortening, and
\hyperref[rule-OEHLoss]{SC085 OEHLoss} the later boundary of the merger.

\subsection{\texorpdfstring{\CAPRRuleHeading{SC072. Shortening of unstressed long vowels}{OEUnstressedLongVowelShortening}}{}}\label{rule-OEUnstressedLongVowelShortening}

\noindent\begin{minipage}{\linewidth}
\begin{ReaderFacingFoma}
define OEUnstressedLongVowelShortening FootCtx(OEUnstressedLongVowelShortening1
    .o. OEUnstressedLongVowelShortening2
    .o. OEUnstressedLongVowelShortening3
    .o. OEUnstressedLongVowelShortening5
    .o. OEUnstressedLongVowelShortening6
    .o. OEUnstressedLongVowelShortening7
    .o. OEUnstressedLongVowelShortening8);
\end{ReaderFacingFoma}
\end{minipage}

An earlier medial component is separately realized in
\hyperref[rule-OEMedLongShort]{SC117 OEMedLongShort}, before
\hyperref[rule-OEHighVowelApocope]{SC063 OEHighVowelApocope}. Goering's
contrast between final \emph{*nētīn} and medially shortened inflected
forms requires that distinction (\citeproc{ref-Goering2016}{Goering
2016, 186--88}; \citeproc{ref-Goering2023}{2023, 67--68}). The rule
below shortens the remaining unstressed long vowels before weak final
syllables reach their later forms. A small group of lexical witnesses
fixes its chronology.

If the rule is moved before \hyperref[rule-NWGmcInStemNLoss]{SC064
NWGmcInStemNLoss}, PGmc \Recon{fúrxtīnaz} `fright' yields \Pred{fyrhten}
rather than expected OE \emph{fyrhte} `fright'. If the rule is delayed
until after \hyperref[rule-OEUnstressedAEMerger]{SC073
OEUnstressedAEMerger}, PGmc \Recon{nḗdrōn} `adder' yields \Pred{nǣdræ}
rather than expected OE \emph{nǣdre} `adder', and PGmc \Recon{fádēr}
`father' yields \Pred{fædær} rather than expected \emph{fæder} `father'.
These outputs require
\hyperref[rule-OEUnstressedLongVowelShortening]{SC072
OEUnstressedLongVowelShortening} to follow
\hyperref[rule-NWGmcInStemNLoss]{SC064 NWGmcInStemNLoss} and precede
\hyperref[rule-OEUnstressedAEMerger]{SC073 OEUnstressedAEMerger}.

The late component therefore follows the earlier weak-tail preparation
and immediately precedes the merger. The old fright displacement records
loss of the long-i environment in the executable; it does not
independently date the early final-n law to this late corridor.

\section{Unstressed ae-merger}\label{unstressed-ae-merger}

\subsection{\texorpdfstring{SC073. Merger of unstressed \emph{*æ} with
\emph{*e}
(\texttt{OEUnstressedAEMerger})}{SC073. Merger of unstressed  with  (OEUnstressedAEMerger)}}\label{rule-OEUnstressedAEMerger}

The following rule handles the merger stage.

\noindent\begin{minipage}{\linewidth}
\begin{ReaderFacingFoma}
define OEUnstressedAEMerger SentenceCtx([
    [{*æ} -> {*e} || .#. {*ᶠ} OEAnyConsonant* _]
    .o. FootCtx(OEWeakTailReduction3)
]);
\end{ReaderFacingFoma}
\end{minipage}

The rule merges unstressed \emph{*æ} with \emph{*e} after shortening has
produced the weak final vowels, yielding the ordinary OE \emph{-e}
spellings.

Pretonic vowels also participate. A positively identified weak initial
syllable is eligible without the preceding-vowel condition appropriate
to post-tonic endings. The superscript \texttt{*ᶠ} records this prosodic
information, not a segment or a grammatical prefix. Stressed roots are
not thereby included in the merger. Later lowering of pretonic
\emph{*i}, including the early \emph{gi-} of \emph{gistīga} `climb
into/onto', belongs to \hyperref[rule-OEPrefixIReduction]{SC076
OEPrefixIReduction}; Hogg places it marginally after the earlier merger
(\citeproc{ref-Campbell1959}{Campbell 1959, 153--55}, §§369--370;
\citeproc{ref-HoggGrammar1992}{Hogg 1992a, §§ 6.50--6.54} (held 2011
reprint, pp.~237--239); \citeproc{ref-RingeTaylor2014}{Ringe and Taylor
2014, 334--35}).

Its earlier and later relations are both concrete. If the rule is moved
before \hyperref[rule-OEUnstressedLongVowelShortening]{SC072
OEUnstressedLongVowelShortening}, PGmc \Recon{nḗdrōn} `adder' yields
\Pred{nǣdræ} rather than expected OE \emph{nǣdre} `adder', and PGmc
\Recon{fádēr} `father' yields \Pred{fædær} rather than expected
\emph{fæder} `father'. If the rule is delayed until after
\hyperref[rule-OEHLoss]{SC085 OEHLoss}, PGmc \Recon{táixōn} `toe' yields
\Pred{tāæ} rather than expected OE \emph{tā} `toe'. These failures show
that \hyperref[rule-OEUnstressedLongVowelShortening]{SC072
OEUnstressedLongVowelShortening} must come before
\hyperref[rule-OEUnstressedAEMerger]{SC073 OEUnstressedAEMerger}, and
that \hyperref[rule-OEUnstressedAEMerger]{SC073 OEUnstressedAEMerger}
must come before \hyperref[rule-OEHLoss]{SC085 OEHLoss}.

The lexical evidence fixes the local order after
\hyperref[rule-OEUnstressedLongVowelShortening]{SC072
OEUnstressedLongVowelShortening} and places the merger before the later
h-loss and contraction.

\section{Medial weak front-vowel quality
changes}\label{medial-weak-front-vowel-quality-changes}

\subsection{Historical discussion}\label{historical-discussion-56}

Hogg and Ringe and Taylor treat the late weakening and merger of
unstressed vowels as a continuing history
(\citeproc{ref-HoggPhonology1992}{Hogg 1992b, 121};
\citeproc{ref-RingeTaylor2014}{Ringe and Taylor 2014, 327--32},
§§6.9.5--6.9.6). \hyperref[rule-OEMedUnstressedILowering1]{SC074
OEMedUnstressedILowering1} lowers medial unstressed \emph{i};
\hyperref[rule-OEMedUnstressedILowering]{SC075 OEMedUnstressedILowering}
preserves \emph{i} before \emph{*ng} in words of the \emph{sċilling}
`shilling' type.

Two stages must be distinguished. Mutation-produced weak \emph{*y}
unrounds to \emph{*i}; retained \emph{*i} later lowers to \emph{*e}.
Hogg dates the unrounding to a very early prehistoric stage and gives
\emph{*gaduling > gædeling}, with root mutation transmitted through the
weak vowel. Initial-root \emph{*y} and still marked subordinate
\emph{*ỳ} are not weak-vowel applications
(\citeproc{ref-HoggGrammar2011}{Hogg 2011, 219}, §6.10 and note 1).

Épinal--Erfurt \emph{scytil} `balance beam', glossing Latin
\emph{momentum}, retains medial unstressed \emph{i}; Corpus
\emph{scytel} `balance beam' has \emph{e}
(\citeproc{ref-Pheifer1974}{Pheifer 1974, 34, 101}, item 632;
\citeproc{ref-Sweet1885}{Sweet 1885, 79}, Corpus 1325;
\citeproc{ref-Christiansen2018}{Christiansen 2018, 583}, note 3, and
pp.~599--600, note 52). After ending loss and the appropriate earlier
changes, inherited \emph{*skútilaz} retains its \emph{i} in final closed
\emph{*-til}: \hyperref[rule-OELAdjacentSyncope]{SC066
OELAdjacentSyncope} does not delete that vowel. It subsequently lowers
to \emph{e} under \hyperref[rule-OEMedUnstressedILowering1]{SC074
OEMedUnstressedILowering1}, giving normalized \emph{sċytel} `dart;
balance beam'. Pheifer's discussion already recognizes both retained
\emph{i} and \emph{e} spellings in the early glossaries and variation
attributable to their common exemplar
(\citeproc{ref-Pheifer1974}{Pheifer 1974, lxxv}, §63). The contrast
illustrates retained and lowered vowel qualities, not a precisely dated
interval between surviving manuscripts; their exemplars need not have
the same dates or orthographic histories. Nor does the Corpus form by
itself establish a specifically West Saxon reflex. This medial
\emph{i/e} evidence is separate from Christiansen's etymological
argument about stressed root \emph{y} against an alleged connection with
the \emph{scit-} `excrement' family
(\citeproc{ref-Christiansen2018}{Christiansen 2018, 593--98}).

The unrounding component is computed after general high-vowel syncope
and before the later light pre-\emph{*l} loss. It is separate from both
deletion and later lowering. Unrounding preserves vowel height and
commutes with the bounded general high-vowel deletion; this adjacency
does not prove an unrounding-after-syncope chronology. With the
source-supported \emph{d/l} retention condition, \emph{gædeling}
`companion' also survives on either side of the light pre-\emph{*l}
component. That particular result does not establish universal
commutation.

The early component is printed at its actual executable location under
\hyperref[rule-OEHighLinkSyncope]{SC115 OEHighLinkSyncope}, a support
component of this vowel-quality history, not a second historical syncope
law. The separate technical decoder runs after the light pre-\emph{*l}
component and before ordinary lowering. It does not turn protected
subordinate \emph{*ỳ} into an application of weak unrounding.

General lowering precedes the restricted restoration before \emph{*ng}.
The evidence is narrower than that for
\hyperref[rule-OEUnstressedLongVowelShortening]{SC072
OEUnstressedLongVowelShortening} and
\hyperref[rule-OEUnstressedAEMerger]{SC073 OEUnstressedAEMerger}.

\subsection{\texorpdfstring{\CAPRRuleHeading{SC074. First medial unstressed-\emph{i} lowering}{OEMedUnstressedILowering1}}{}}\label{rule-OEMedUnstressedILowering1}

\begin{ReaderFacingFoma}
define OEMedUnstressedILowering1 FootCtx([
    {*i} -> {*e} || EnglishStarVocalic [EnglishStarConsonant | EnglishPalatalConsonant]+ _
]);
\end{ReaderFacingFoma}

The rule lowers medial unstressed \emph{*i} to \emph{*e} after a
preceding vocalic syllable. The resulting \emph{e}-outcome is reversed
before \emph{*ng}.

If the rule is moved before
\hyperref[rule-OEUnstressedLongVowelShortening]{SC072
OEUnstressedLongVowelShortening}, PGmc \Recon{fúrxtīnaz} `fright' yields
\Pred{fyrhti} rather than expected OE \emph{fyrhte} `fright'. If it is
delayed until after \hyperref[rule-OEMedUnstressedILowering]{SC075
OEMedUnstressedILowering}, PGmc \Recon{skíllingaz} `shilling' yields
\Pred{sċilleng} rather than expected \emph{sċilling} `shilling'. The
derivations require \hyperref[rule-OEMedUnstressedILowering1]{SC074
OEMedUnstressedILowering1} to follow
\hyperref[rule-OEUnstressedLongVowelShortening]{SC072
OEUnstressedLongVowelShortening} and precede
\hyperref[rule-OEMedUnstressedILowering]{SC075
OEMedUnstressedILowering}.

The evidence is narrow on each side. The rule follows unstressed
long-vowel shortening and precedes the more specific \emph{*ng}
preservation.

\section{Preservation of i before ng}\label{preservation-of-i-before-ng}

\subsection{\texorpdfstring{\CAPRRuleHeading{SC075. Preservation of medial unstressed \emph{*i} before \emph{*ng}}{OEMedUnstressedILowering}}{}}\label{rule-OEMedUnstressedILowering}

The following rule reverses the lowering before \emph{*ng}.

\noindent\begin{minipage}{\linewidth}
\begin{ReaderFacingFoma}
define OEMedUnstressedILowering Ctx([
    {*e} -> {*i} || _ {*n} {*g}
]);
\end{ReaderFacingFoma}
\end{minipage}

The rule restores \emph{*i} before \emph{*ng}, preventing the broader
lowering from producing the wrong medial vowel in forms such as
\emph{sċilling} `shilling'.

Moving the rule before \hyperref[rule-OEMedUnstressedILowering1]{SC074
OEMedUnstressedILowering1} makes PGmc \Recon{skíllingaz} `shilling'
yield \Pred{sċilleng} rather than expected OE \emph{sċilling}
`shilling'. On this evidence, I take
\hyperref[rule-OEMedUnstressedILowering]{SC075 OEMedUnstressedILowering}
to follow \hyperref[rule-OEMedUnstressedILowering1]{SC074
OEMedUnstressedILowering1}. Moving it later within the tested range
creates no equally sharp failure.

\section{Medial unstressed-u
lowering}\label{medial-unstressed-u-lowering}

\subsection{\texorpdfstring{\CAPRRuleHeading{SC040. Lowering of medial unstressed \emph{*u}}{OEMedUnstressedULowering}}{}}\label{rule-OEMedUnstressedULowering}

\noindent\begin{minipage}{\linewidth}
\begin{ReaderFacingFoma}
define OEMedUnstressedULowering [MPFoot .o. MPULower .o. MPDecode];
\end{ReaderFacingFoma}
\end{minipage}

Medial unstressed \emph{*u} lowers to \emph{*o} in the later Old English
period, while accented \emph{*u} preserves the following unstressed
vowel (\citeproc{ref-Campbell1959}{Campbell 1959, 155--56}, §373). This
lowering follows the earlier shortening and resolution of unstressed
long vowels, rather than preceding that history. The eighth-century
unstressed back-vowel system and its later development distinguish the
two periods (\citeproc{ref-StauslandJohnsen2015}{Stausland Johnsen 2015,
21--22, 28--31}).

The inherited ordering evidence from \emph{wuduwe} `widow' remains
relevant: delaying \hyperref[rule-OEWICombinativeUUmlaut]{SC039
OEWICombinativeUUmlaut} until after lowering yields \Pred{wudowe} rather
than the expected form. In the adopted sequence, \emph{wuduwe} retains
its medial vowel after accented \emph{*u}, whereas \emph{weorold}
`world' undergoes lowering.

\subsubsection{Medial prominence and the initial
foot}\label{medial-prominence-and-the-initial-foot}

The retained \emph{u} of \emph{eorþcunde} `earthly' is not an exemption
for a suffix or an inflection. Its heavy first syllable is followed by a
nonfinal heavy syllable which can head a medial foot. In contrast, an
initial light syllable is accommodated with a following syllable,
including a heavy one, in the initial foot. These distinct
configurations are supported by the quantity-sensitive analyses of
McCully and Hogg and Goering (\citeproc{ref-McCullyHogg1990}{McCully and
Hogg 1990, 323--33}, especially pp.~330--333;
\citeproc{ref-Goering2023}{Goering 2023, 79--84}, §§4.5.1.1--4.5.2).

The relevant short diphthong \emph{*éo} counts as one short nucleus, not
as two separate vowels. In the inflected control \emph{*wéoruldes},
syllabification gives light \emph{*weo}, heavy \emph{*rul} and final
\emph{*des}. The first two belong to the initial light--heavy foot; the
medial \emph{u} is not a separate head and lowers. The shorter selected
\emph{weorold} has the same initial accommodation. In
\emph{eorþ.cun.de}, the initial \emph{*eorþ} is heavy and the subsequent
nonfinal \emph{*cun} is independently heavy: its \emph{u} is protected.
The genitive control \emph{eorþcundes} has the same contrast. This
explanation depends on the phonological representation at lowering, not
on the spelling of a comparative citation form.

The independent comparisons \emph{æþelingas} `princes' versus
\emph{cyningas} `kings', and \emph{rīdende} `riding' versus
\emph{wesende} `being', test initial-foot accommodation with different
vowel classes. Inflected \emph{godcundne} `divine' and \emph{godcundre}
`divine' supply medial retention comparisons, not permission to exempt
the whole formation (\citeproc{ref-McCullyHogg1990}{McCully and Hogg
1990, 330--33}; \citeproc{ref-HoggGrammar2011}{Hogg 2011, 49--50},
§2.89). Final uninflected \emph{godcund} remains outside this nonfinal
medial-foot computation; the complete paradigm is not thereby derived.

\subsubsection{Temporary computation, not another inherited-accent
loss}\label{temporary-computation-not-another-inherited-accent-loss}

The implementation syllabifies the current string, accommodates the
initial foot and marks eligible later closed short-vowel heads. It uses
the existing short-vowel and short-diphthong inventories. Temporary
graves on \emph{a e i o u} distinguish computed prominence from the
inherited compound accent, already removed by inherited secondary-accent
deletion. They are decoded immediately after this lowering. The existing
pretonic-prefix representation is accommodated without inserting a new
initial boundary or relocating inherited accent deletion.

Segmental conditioning is retained: the preceding \emph{u}-class nuclei
and following \emph{m} block lowering, with the same consonant classes.
A temporarily marked \emph{ù} also belongs to that preceding exclusion;
changing its technical marking must not make it license lowering of a
subsequent \emph{u}. This is preservation of the original condition, not
a new harmony law.

Hogg cautions against using ``tertiary stress'' merely to rename a
desired retention, and distinguishes labial-nasal preservation from
ordinary \emph{-un > -on} (\citeproc{ref-Hogg1996}{Hogg 1996, 8--11},
especially p.~11 n.~7; \citeproc{ref-HoggGrammar2011}{Hogg 2011,
239--41}, §§6.55--6.58). The computation consequently uses weight and
domain rather than lexical identity or a blanket dental-nasal exception.
Nor does metrical prominence by itself certify speech-stress history
(\citeproc{ref-MinkovaZhou2023}{Minkova and Zhou 2023, 86--90}). This is
a bounded prosodic projection for the present sound law, not a complete
cyclic Old English stress grammar.

\subsubsection{Existing chronological
distinctions}\label{existing-chronological-distinctions}

\emph{Reġnboga} `rainbow' is not a further live example: its second-root
\emph{o} arose while subordinate stress was still represented, under
\hyperref[rule-PNWGmcULowering]{SC017 PNWGmcULowering}. Its
linking-vowel history therefore does not require secondary-accent loss
before the present late lowering.

The reciprocal early-lowering experiment is no longer equivalent: the
intervening vowel-resolution history restores medial \emph{*u}, so that
movement does not now produce the same \emph{widow} failure.

The formerly reported \emph{ġeoguþ} `youth' boundary does not establish
that lowering preceded
\hyperref[rule-OEUnstressedLongVowelShortening]{SC072
OEUnstressedLongVowelShortening}. That experiment inserted a front glide
vowel early in the derivation; delayed lowering then encountered an
altered preceding nucleus and produced \Pred{ġeogoþ}. With the adopted
late \hyperref[rule-OEWsPalatalGlide]{SC016 OEWsPalatalGlide} spelling,
the root retains phonological \emph{*u} through lowering and preserves
the medial vowel. Moving lowering across shortening in this current
architecture leaves \emph{ġeoguþ} unchanged. The spelling evidence is
therefore not an independent date for the lowering; the later placement
rests on the historical vowel chronology, not a reversal of that
obsolete experiment.

\section{Pretonic i-lowering}\label{pretonic-i-lowering}

\subsection{Historical discussion}\label{historical-discussion-57}

Late weak-tail reduction affects unstressed prefixes as well as
inflectional endings and medial vowels. Fulk's discussion of prefix
vowels accounts for OE \emph{*be-} and \emph{*ne-}
(\citeproc{ref-Fulk2018}{Fulk 2018, 97}, §5.7). Hogg explicitly includes
\emph{be-} and \emph{ge-} in the lowering of unstressed \emph{*i},
marginally later than the earlier front-vowel merger. Ringe and Taylor
supply early \emph{gistīga} `climb into/onto' beside later
\emph{gestīga} `climb into/onto', and place \emph{bi-} and \emph{ni}
alongside this development (\citeproc{ref-HoggGrammar1992}{Hogg 1992a,
§§ 6.50--6.54} (held 2011 reprint, pp.~237--239);
\citeproc{ref-RingeTaylor2014}{Ringe and Taylor 2014, 334--35}).

The limited relative chronology follows the handbook reconstruction, not
the older proxy's displacement results or a simple count of manuscript
spellings. Pheifer reports Épinal \emph{gi-/ge-} as 40/13 and Erfurt as
31/34; both related witnesses already mix the spellings. Their different
denominators and transmitted linguistic layers do not independently date
a universal merger. The Durham Ritual still writes \emph{gi-} in the
tenth century. Different varieties, possible raised allophony and
transmitted orthography must be distinguished from successive
reconstructed vowels; the accentual status of contested \emph{ġeatwe}
`equipment' cannot supply a secure prefix-specific date either
(\citeproc{ref-Pheifer1974}{Pheifer 1974, lxxv}, §63;
\citeproc{ref-Toller1921}{Toller 1921, 284};
\citeproc{ref-Stevenson1840}{Stevenson 1840, 7};
\citeproc{ref-HoggGrammar2011}{Hogg 2011, 238--39}, §§6.53--6.54;
\citeproc{ref-Quinlin1991}{Quinlin 1991, 155--56}).

\subsection{\texorpdfstring{SC076. Pretonic i-lowering
(\texttt{OEPrefixIReduction})}{SC076. Pretonic i-lowering (OEPrefixIReduction)}}\label{rule-OEPrefixIReduction}

\noindent\begin{minipage}{\linewidth}
\begin{ReaderFacingFoma}
define OEPrefixIReduction SentenceCtx([
    [{*i} -> {*e} || .#. {*ᶠ} OEAnyConsonant* _]
    .o. [{*ᶠ} -> 0]
]);
\end{ReaderFacingFoma}
\end{minipage}

The rule lowers short \emph{*i} in a positively identified weak pretonic
syllable, not specifically in morphemes beginning with a particular
consonant. Thus the same conditioning treats inherited weak \emph{*bi-}
and the raised vowel in \emph{*gi-}; stressed roots remain outside its
domain. The superscript \texttt{*ᶠ} is a nonsegmental prosodic
annotation, removed once it is no longer needed. The rule therefore
produces \emph{*e} directly, not a protected intermediate breve-e.

This bounded implementation requires a following headed syllable. The
separately unstressed negative particle \emph{ni} `not' \textgreater{}
\emph{ne} `not' supports the wider history of weak-vowel lowering, but
is not made into a newly encoded selected input or a claim that the
present pretonic condition covers every weak particle
(\citeproc{ref-RingeTaylor2014}{Ringe and Taylor 2014, 335}).

The earlier displacement experiment found no final-output boundary for
the old consonant-identity proxy. That negative result is not a test of
the reformulated prosodic rule and does not date it. The relative
placement after the earlier merger follows the handbook account; its
precise interval remains approximate.

\section{Weak-tail reduction}\label{weak-tail-reduction}

\subsection{Historical discussion}\label{historical-discussion-58}

Campbell, Hogg, Ringe and Taylor, and Fulk describe a late history in
which apocope, shortening, contraction, and further weak-tail reductions
reshape final syllables (\citeproc{ref-Campbell1959}{Campbell 1959,
148}, §355; \citeproc{ref-HoggPhonology1992}{Hogg 1992b, 121};
\citeproc{ref-RingeTaylor2014}{Ringe and Taylor 2014, 298--314},
§§6.8.3--6.9.3; \citeproc{ref-Fulk2018}{Fulk 2018, 90--91}, §5.6).
Lexical failures place the remaining weak-tail reduction after
unstressed fronting and before contraction.

\subsection{\texorpdfstring{\CAPRRuleHeading{SC078. Reduction of remaining weak-tail vowels}{OEWeakTailReduction}}{}}\label{rule-OEWeakTailReduction}

\noindent\begin{minipage}{\linewidth}
\begin{ReaderFacingFoma}
define OEWeakTailReduction OEWeakTailReduction1;
\end{ReaderFacingFoma}
\end{minipage}

The rule reduces the remaining weak-tail vowels, preventing a broad
class of \emph{-en} and extra-vowel outcomes.

I place the change after \hyperref[rule-OEUnstressedFrontingEarly]{SC070
OEUnstressedFrontingEarly} and before
\hyperref[rule-OEContraction]{SC086 OEContraction}. Moving it before
\hyperref[rule-OEUnstressedFrontingEarly]{SC070
OEUnstressedFrontingEarly}, PGmc \Recon{bákaną} `bake' yields
\Pred{bacen} rather than expected OE \emph{bacan} `bake', and PGmc
\Recon{bíndaną} `bind' yields \Pred{binden} rather than expected
\emph{bindan} `bind', alongside a much wider set of comparable
\emph{-en} failures. If the rule is delayed until after
\hyperref[rule-OEContraction]{SC086 OEContraction}, PGmc
\Recon{fléuxaną} `flee' yields \Pred{flēoan} rather than expected OE
\emph{flēon} `flee', and PGmc \Recon{sláxaną} `slay' yields
\Pred{sleaan} rather than expected \emph{slēan} `slay'.

The earlier boundary spans a wide interval and does not establish a
close neighboring relation. The later boundary is narrower:
\hyperref[rule-OEWeakTailReduction]{SC078 OEWeakTailReduction} precedes
\hyperref[rule-OEContraction]{SC086 OEContraction}.

\section{J-loss after heavy
syllables}\label{j-loss-after-heavy-syllables}

\subsection{Historical discussion}\label{historical-discussion-59}

After \hyperref[rule-OEJLossAfterHeavy]{SC079 OEJLossAfterHeavy} removes
\emph{*j} in heavy environments, forms such as \emph{lungen} `lungs'
acquire a final geminate.
\hyperref[rule-OEFinalGeminateSimplification]{SC080
OEFinalGeminateSimplification} then removes the second nasal.

\hyperref[rule-OEJLossAfterHeavy]{SC079 OEJLossAfterHeavy} has a broad
earlier boundary at \hyperref[rule-OEIUmlaut]{SC055 OEIUmlaut}.
\hyperref[rule-OEFinalGeminateSimplification]{SC080
OEFinalGeminateSimplification} is fixed only by the final \emph{nn}
outcome in the following derivation.

\subsection{\texorpdfstring{SC079. Loss of \emph{*j} after heavy
syllables
(\texttt{OEJLossAfterHeavy})}{SC079. Loss of  after heavy syllables (OEJLossAfterHeavy)}}\label{rule-OEJLossAfterHeavy}

\begin{ReaderFacingFoma}
define OEJLossAfterHeavy [
    {*j} -> 0 || (EnglishStarLongVowel | EnglishStarDiphthong) [EnglishStarConsonantNoR | EnglishPalatalConsonant] _,
    {*j} -> 0 || EnglishStarShortVowel [EnglishStarConsonant | EnglishPalatalConsonant] [EnglishStarConsonantNoR | EnglishPalatalConsonant] _
];
\end{ReaderFacingFoma}

The rule removes \emph{*j} after the relevant heavy-syllable
configurations, after the earlier umlaut-sensitive vocalism has
developed. The affected glide is \emph{*j}.

If the rule is moved before \hyperref[rule-OEIUmlaut]{SC055 OEIUmlaut},
PGmc \Recon{galáubijaną} `believe' yields \Pred{ġelēafan} rather than
expected OE \emph{ġelīefan} `believe', PGmc \Recon{báugijaną} `bow'
yields \Pred{bēaġan} rather than expected \emph{bīeġan} `bow', and PGmc
\Recon{fúlgijaną} `follow' yields \Pred{fulġan} rather than expected
\emph{fylġan} `follow'. If it is delayed until after
\hyperref[rule-OEFinalGeminateSimplification]{SC080
OEFinalGeminateSimplification}, PGmc \Recon{lúnganjō} `lungs' yields
\Pred{lungenn} rather than expected OE \emph{lungen} `lungs'. I
accordingly take \hyperref[rule-OEJLossAfterHeavy]{SC079
OEJLossAfterHeavy} to follow \hyperref[rule-OEIUmlaut]{SC055 OEIUmlaut}
and precede \hyperref[rule-OEFinalGeminateSimplification]{SC080
OEFinalGeminateSimplification}.

The earlier boundary is broad, but the relation to final geminate
simplification is local.

\section{Final geminate
simplification}\label{final-geminate-simplification}

\subsection{\texorpdfstring{\CAPRRuleHeading{SC080. Simplification of final geminates}{OEFinalGeminateSimplification}}{}}\label{rule-OEFinalGeminateSimplification}

The following rule handles the final simplification directly.

\noindent\begin{minipage}{\linewidth}
\begin{ReaderFacingFoma}
define OEFinalGeminateSimplification [
    {*n} -> 0 || {*n} _ .#.
];
\end{ReaderFacingFoma}
\end{minipage}

The rule removes the extra final nasal in forms where the preceding
derivation has already created a final geminate.

Moving the rule before \hyperref[rule-OEJLossAfterHeavy]{SC079
OEJLossAfterHeavy} makes PGmc \Recon{lúnganjō} `lungs' yield
\Pred{lungenn} rather than expected OE \emph{lungen} `lungs'. These
failures require \hyperref[rule-OEFinalGeminateSimplification]{SC080
OEFinalGeminateSimplification} to follow
\hyperref[rule-OEJLossAfterHeavy]{SC079 OEJLossAfterHeavy}. Moving it
later within the tested range before \hyperref[rule-OERMetathesis]{SC087
OERMetathesis} creates no new failure.

\section{J-strengthening after front
diphthongs}\label{j-strengthening-after-front-diphthongs}

\subsection{Historical discussion}\label{historical-discussion-60}

\hyperref[rule-OEJStrengtheningAfterFrontDiphthong]{SC081
OEJStrengtheningAfterFrontDiphthong} preserves a consonantal outcome
after front diphthongs.
\hyperref[rule-OEIntervocalicJVocalization]{SC082
OEIntervocalicJVocalization} then normalizes part of the inherited
intervocalic \emph{*j} domain, and
\hyperref[rule-OEUnstressedEIContraction]{SC083
OEUnstressedEIContraction} removes the resulting \emph{ei}-like sequence
in weak verbal endings.

The output of each rule conditions the next.
\hyperref[rule-OEIntervocalicJVocalization]{SC082
OEIntervocalicJVocalization} has local lexical evidence on both sides;
\hyperref[rule-OEJStrengtheningAfterFrontDiphthong]{SC081
OEJStrengtheningAfterFrontDiphthong} has a distant earlier boundary, and
\hyperref[rule-OEUnstressedEIContraction]{SC083
OEUnstressedEIContraction} has no tested later boundary.

\subsection{\texorpdfstring{\CAPRRuleHeading{SC081. Strengthening of \emph{*j} after front diphthongs}{OEJStrengtheningAfterFrontDiphthong}}{}}\label{rule-OEJStrengtheningAfterFrontDiphthong}

\noindent\begin{minipage}{\linewidth}
\begin{ReaderFacingFoma}
define OEJStrengtheningAfterFrontDiphthong [
    {*j} -> {*ʒ} || [{*ēa}|{*ḗa}|{*íe}|{*īe}|{*éa}] _ EnglishStarVocalic
];
\end{ReaderFacingFoma}
\end{minipage}

After the relevant front diphthongs, \emph{*j} first strengthened to a
consonantal outcome; otherwise it would have vocalized too early.

If the rule is moved before \hyperref[rule-OEIUmlaut]{SC055 OEIUmlaut},
PGmc \Recon{stráwjaną} `strew' yields \Pred{strēaġan} rather than
expected OE \emph{strīeġan} `strew'. If it is delayed until after
\hyperref[rule-OEIntervocalicJVocalization]{SC082
OEIntervocalicJVocalization}, the same PGmc form yields \Pred{strīeian}
rather than \emph{strīeġan}. The order test requires
\hyperref[rule-OEJStrengtheningAfterFrontDiphthong]{SC081
OEJStrengtheningAfterFrontDiphthong} to follow
\hyperref[rule-OEIUmlaut]{SC055 OEIUmlaut} and precede
\hyperref[rule-OEIntervocalicJVocalization]{SC082
OEIntervocalicJVocalization}.

The earlier constraint reaches back to \hyperref[rule-OEIUmlaut]{SC055
OEIUmlaut} and therefore defines a wide interval. The \emph{strīeġan}
`strew' derivation fixes the local relation to
\hyperref[rule-OEIntervocalicJVocalization]{SC082
OEIntervocalicJVocalization}.

\section{Intervocalic j-vocalization}\label{intervocalic-j-vocalization}

\subsection{\texorpdfstring{\CAPRRuleHeading{SC082. Bounded inherited-j normalization}{OEIntervocalicJVocalization}}{}}\label{rule-OEIntervocalicJVocalization}

\noindent\begin{minipage}{\linewidth}
\begin{ReaderFacingFoma}
define OEIntervocalicJVocalization [
    {*j} -> {*i} ||
        [EnglishStarVocalic - [{*ǣ}|{*ē}|{*ḗ}]] _ EnglishStarVocalic
];
\end{ReaderFacingFoma}
\end{minipage}

The former unrestricted VjV matcher wrongly gave model-only \Pred{cǣie}
instead of \emph{cǣġe} `key'. Hogg's pre-OE \emph{*kǣjæ} retains the
glide (\citeproc{ref-Hogg1979}{Hogg 1979, 105}). The bounded rule
therefore excludes preceding non-high long front monophthongs.

This is not a generic physical vocalization law. Hogg distinguishes
high-vowel coalescence from spellings that merely represent a
consonantal glide and from unstressed alternations
(\citeproc{ref-HoggGrammar2011}{Hogg 2011, 283--86}, §§7.69--7.76). The
remaining weak-suffix e+j to e+i to i path telescopes later
raising/contraction, rather than asserting that its e+i intermediate is
a source reconstruction (\citeproc{ref-RingeTaylor2014}{Ringe and Taylor
2014, 228}). \hyperref[rule-OEUnstressedEIContraction]{SC083
OEUnstressedEIContraction} consumes that modeled sequence. The full
residual proxy domain is not claimed to be one independently established
historical event.

The separately represented fricative ʝ passes through this zone without
being treated as inherited j. Its later merger belongs to
\hyperref[rule-OEPalatalFricativeMerger]{SC109
OEPalatalFricativeMerger}.

Moving the rule before
\hyperref[rule-OEJStrengtheningAfterFrontDiphthong]{SC081
OEJStrengtheningAfterFrontDiphthong} makes PGmc \Recon{stráwjaną}
`strew' yield \Pred{strīeian} rather than expected OE \emph{strīeġan}
`strew'. Delaying it until after
\hyperref[rule-OEUnstressedEIContraction]{SC083
OEUnstressedEIContraction} makes PGmc \Recon{búrōjaną} `bore' yield
\Pred{boreian} rather than expected OE \emph{borian} `bore', PGmc
\Recon{xándlōjaną} `handle' yield \Pred{handleian} rather than expected
\emph{handlian} `handle', and PGmc \Recon{mákōjaną} `make' yield
\Pred{maceian} rather than expected \emph{macian} `make'. The witness
forms require \hyperref[rule-OEIntervocalicJVocalization]{SC082
OEIntervocalicJVocalization} to follow
\hyperref[rule-OEJStrengtheningAfterFrontDiphthong]{SC081
OEJStrengtheningAfterFrontDiphthong} and precede
\hyperref[rule-OEUnstressedEIContraction]{SC083
OEUnstressedEIContraction}.

\hyperref[rule-OEIntervocalicJVocalization]{SC082
OEIntervocalicJVocalization} is therefore ordered between strengthening
and contraction.

\section{Unstressed ei-contraction}\label{unstressed-ei-contraction}

\subsection{\texorpdfstring{SC083. Contraction of unstressed \emph{ei}
(\texttt{OEUnstressedEIContraction})}{SC083. Contraction of unstressed  (OEUnstressedEIContraction)}}\label{rule-OEUnstressedEIContraction}

The final rule removes the extra unstressed \emph{e} before \emph{i}.

\begin{ReaderFacingFoma}
define OEUnstressedEIContraction [
    {*e} -> 0 || EnglishStarVocalic [EnglishStarConsonant | EnglishPalatalConsonant]+ _ {*i}
];
\end{ReaderFacingFoma}

The rule contracts the unstressed \emph{ei}-like sequence that the
preceding vocalization would otherwise leave behind in forms such as
\emph{borian} `bore' and \emph{liccian} `lick'.

Moving the rule before \hyperref[rule-OEIntervocalicJVocalization]{SC082
OEIntervocalicJVocalization} makes PGmc \Recon{búrōjaną} `bore' yield
\Pred{boreian} rather than expected OE \emph{borian} `bore', and PGmc
\Recon{líkkōjaną} `lick' yield \Pred{licceian} rather than expected
\emph{liccian} `lick'. The contrast requires
\hyperref[rule-OEUnstressedEIContraction]{SC083
OEUnstressedEIContraction} to follow
\hyperref[rule-OEIntervocalicJVocalization]{SC082
OEIntervocalicJVocalization}. Moving it later within the tested range
before \hyperref[rule-OERMetathesis]{SC087 OERMetathesis} creates no new
failure.

\section{Intervocalic h-loss}\label{intervocalic-h-loss}

\subsection{Historical discussion}\label{historical-discussion-61}

When \hyperref[rule-OEHLoss]{SC085 OEHLoss} removes intervocalic
\emph{*h}, it creates hiatus. \hyperref[rule-OEContraction]{SC086
OEContraction} immediately resolves the resulting vowel sequence.

Ringe and Taylor describe this late sequence of \emph{h}-loss and
contraction (\citeproc{ref-RingeTaylor2014}{Ringe and Taylor 2014,
305--14}, §§6.9.1--6.9.3). Fulk places the contracted verbs in a broader
Germanic context (\citeproc{ref-Fulk2018}{Fulk 2018, 270}, §12.21), and
Luick describes the corresponding West Germanic contractions
(\citeproc{ref-Luick1914}{Luick 1914-\/-1940, 165}).

\subsection{\texorpdfstring{SC085. Loss of intervocalic \emph{*h}
(\texttt{OEHLoss})}{SC085. Loss of intervocalic  (OEHLoss)}}\label{rule-OEHLoss}

\noindent\begin{minipage}{\linewidth}
\begin{ReaderFacingFoma}
define OEHLoss [
    {*x} -> 0 || EnglishStarVocalic _ EnglishStarVocalic
];
\end{ReaderFacingFoma}
\end{minipage}

The rule removes intervocalic \emph{*h}, creating the hiatus that later
contraction must resolve.

If the rule is moved before \hyperref[rule-OEUnstressedAEMerger]{SC073
OEUnstressedAEMerger}, PGmc \Recon{táixōn} `toe' yields \Pred{tāæ}
rather than expected OE \emph{tā} `toe'. If it is delayed until after
\hyperref[rule-OEContraction]{SC086 OEContraction}, PGmc
\Recon{fléuxaną} `flee' yields \Pred{flēoan} rather than expected OE
\emph{flēon} `flee', PGmc \Recon{sláxaną} `slay' yields \Pred{sleaan}
rather than expected \emph{slēan} `slay', PGmc \Recon{téxun} `draw'
yields \Pred{teoon} rather than expected \emph{tēon} `draw', and PGmc
\Recon{táixōn} yields \Pred{tāe} rather than expected \emph{tā}. These
outputs require \hyperref[rule-OEHLoss]{SC085 OEHLoss} to follow
\hyperref[rule-OEUnstressedAEMerger]{SC073 OEUnstressedAEMerger} and
precede \hyperref[rule-OEContraction]{SC086 OEContraction}.

The earlier boundary rests on one witness; the four later witnesses
establish the immediate relation to contraction.

\section{Contraction after h-loss}\label{contraction-after-h-loss}

\subsection{\texorpdfstring{SC086. Contraction of the resulting hiatus
(\texttt{OEContraction})}{SC086. Contraction of the resulting hiatus (OEContraction)}}\label{rule-OEContraction}

The following rule contracts the hiatus left by
\hyperref[rule-OEHLoss]{SC085 OEHLoss}.

\begin{ReaderFacingFoma}
define OEContraction [
    {*a} {*a} -> {*ā},
    {*e} {*e} -> {*ē},
    {*i} {*i} -> {*ī},
    {*o} {*o} -> {*ō},
    {*u} {*u} -> {*ū},
    {*ea} {*a} -> {*ēa},
    {*ēa} {*a} -> {*ēa},
    {*eo} {*a} -> {*ēo},
    {*ēo} {*a} -> {*ēo},
    {*eo} {*o} -> {*ēo},
    {*ēo} {*o} -> {*ēo},
    {*éo} {*o} -> {*ḗo},
    {*ḗo} {*o} -> {*ḗo},
    {*ā} {*a} -> {*ā},
    {*ā} {*e} -> {*ā},
    {*ē} {*a} -> {*ē},
    {*ē} {*e} -> {*ē},
    {*ḗ} {*a} -> {*ḗ},
    {*ḗ} {*e} -> {*ḗ},
    {*ī} {*a} -> {*ī},
    {*ī} {*e} -> {*ī},
    {*ḯ} {*a} -> {*ḯ},
    {*ḯ} {*e} -> {*ḯ},
    {*ō} {*a} -> {*ō},
    {*ō} {*e} -> {*ō},
    {*ū} {*a} -> {*ū},
    {*ū} {*e} -> {*ū}
];
\end{ReaderFacingFoma}

The rule contracts the vowel sequences created after \emph{h}-loss,
producing \emph{flēon} `flee', \emph{slēan} `slay', and \emph{tēon}
`draw'.

Moving contraction before \hyperref[rule-OEHLoss]{SC085 OEHLoss} makes
PGmc \Recon{fléuxaną} `flee' yield \Pred{flēoan} rather than expected OE
\emph{flēon} `flee', PGmc \Recon{sláxaną} `slay' yield \Pred{sleaan}
rather than expected \emph{slēan} `slay', PGmc \Recon{téxun} `draw'
yield \Pred{teoon} rather than expected \emph{tēon} `draw', and PGmc
\Recon{táixōn} `toe' yield \Pred{tāe} rather than expected \emph{tā}
`toe'. The derivations require \hyperref[rule-OEContraction]{SC086
OEContraction} to follow \hyperref[rule-OEHLoss]{SC085 OEHLoss}. Moving
it later within the tested range before
\hyperref[rule-OERMetathesis]{SC087 OERMetathesis} creates no new
failure. The more distant \hyperref[rule-OEWeakTailReduction]{SC078
OEWeakTailReduction} relation establishes only that weak-tail reduction
precedes contraction.

\section{R-metathesis}\label{r-metathesis}

\subsection{Historical discussion}\label{historical-discussion-62}

Sievers-Brunner describes r-metathesis in forms such as \emph{berstan}
`burst', \emph{forst} `frost', and \emph{cærse} `cress'
(\citeproc{ref-SieversBrunner1965}{Brunner 1965, 146--47}, §179). Luick
likewise treats it as a later rearrangement whose interaction with
breaking remains variable (\citeproc{ref-Luick1914}{Luick 1914-\/-1940,
201}).

For the selected West Saxon \emph{berstan}, the evidence places breaking
before metathesis; this is not a chronology for every dialectal
metathesis. It does not establish an ordering relation between
\hyperref[rule-OEContraction]{SC086 OEContraction} and
\hyperref[rule-OERMetathesis]{SC087 OERMetathesis}.

Barrack devotes his longest metathesis discussion to this change as rule
459.1, whose feature matrix moves \emph{*r} across a short vowel before
\emph{s} or \emph{n} and whose worked examples include the \emph{cyrps}
and \emph{-firiþ}/\emph{-ferþ} types; a companion rule 459.3 treats the
regressive movement of \emph{r} back across an accented short vowel
before \emph{ht}, as in late West Saxon \emph{worhte} beside
\emph{wrohte} (\citeproc{ref-Barrack1975}{Barrack 1975, 115--18}). His
account agrees that in West Saxon the \emph{beornan}/\emph{birnan} `to
burn' family shows metathesis operating on the output of earlier
developments, which he schematizes as 157 followed by 170a followed by
459.1; but he records that the relative order of breaking and metathesis
varies by dialect, separating the West Saxon from the Anglian treatment
(\citeproc{ref-Barrack1975}{Barrack 1975, 115--17}). Because his
orderings are formal rather than chronological
(\citeproc{ref-Barrack1975}{Barrack 1975, 34}), that dialect-order
observation is conditioning evidence to be weighed in its own right, not
a dating, and does not alter the West Saxon ordering adopted here.

Barrack describes generalized environments of 459.1 as ``a late stage of
free variation'' (\citeproc{ref-Barrack1975}{Barrack 1975, 116}), while
grading the evidence for 459.3 as meagre
(\citeproc{ref-Barrack1975}{Barrack 1975, 118}). These are descriptions
of the distribution of forms, not an explanation by optional sound
change. Campbell records byforms such as \emph{forst} beside
\emph{frost}, \emph{gærs} `grass' beside \emph{græs}, \emph{cerse}
beside \emph{cresse}, and \emph{rinnan} beside \emph{iernan},
\emph{eornan} (\citeproc{ref-Campbell1959}{Campbell 1959, 184}, §459).
Regular sound change remains exceptionless within its phonological,
chronological and dialectal domain. Distinct dialect histories,
transmission of earlier and later forms, borrowing or independently
evidenced analogy could account for coexistence; the byforms alone do
not decide among them.

The categorical operation below correctly handles the selected
\emph{berstan} and \emph{forst}. Controlled comparisons extending only
the vowel inventory to include \emph{*æ} and \emph{*y}, while retaining
\emph{st}, alter no selected output. With that same inventory, removing
the \emph{t} requirement changes \emph{græs} `grass' to attested
\emph{gærs}. Thus the selected forms discriminate these implementations,
although neither spelling establishes which conditioning was
historically productive. The earlier zero-difference comparison before
\emph{s} versus \emph{st} depended on an enumeration excluding
\emph{*æ}. A further extension to nasal environments also alters
selected \emph{brandes} `firebrand (gen.)', \emph{drenċ} `drink' and
\emph{strenġ} `string'; a regressive pre-\emph{ht} operation changes
selected \emph{fyrhte} `fright'. These are reasons not to adopt those
operations unchanged across the present selections, not proof that
broader regular metathesis was historically impossible.

The cress family supplies independent evidence beyond \emph{st}:
Campbell and Hogg describe \emph{cresse} through reconstructed
\emph{*cersse} to \emph{cerse}, with simplification of the geminate
(\citeproc{ref-Campbell1959}{Campbell 1959, 190}, §476;
\citeproc{ref-HoggGrammar2011}{Hogg 2011, 288}, §7.80). Ringe and Taylor
give \emph{cærse} from reconstructed \emph{*cræsse}
(\citeproc{ref-RingeTaylor2014}{Ringe and Taylor 2014, 263}). Campbell's
Middle English \emph{kers} evidence also places the metathesis after
assibilation (\citeproc{ref-Campbell1959}{Campbell 1959, 184}, §459
n.~1). These forms merit discussion even without selecting a new corpus
target: their inherited stem and exact attested cell and dialect must be
established separately. The retained operation represents the changes
presently required by the selected forms; it does not delimit the whole
historical phenomenon.

\subsection{\texorpdfstring{\CAPRRuleHeading{SC087. Metathesis of \emph{*r} with a following short vowel}{OERMetathesis}}{}}\label{rule-OERMetathesis}

\begin{ReaderFacingFoma}
define OERMetathesis [
    {*r} {*e} -> {*e} {*r} || EnglishStarConsonant _ {*s} {*t},
    {*r} {*u} -> {*u} {*r} || EnglishStarConsonant _ {*s} {*t},
    {*r} {*i} -> {*i} {*r} || EnglishStarConsonant _ {*s} {*t},
    {*r} {*o} -> {*o} {*r} || EnglishStarConsonant _ {*s} {*t},
    {*r} {*a} -> {*a} {*r} || EnglishStarConsonant _ {*s} {*t},
    {*r} {*é} -> {*é} {*r} || EnglishStarConsonant _ {*s} {*t},
    {*r} {*ó} -> {*ó} {*r} || EnglishStarConsonant _ {*s} {*t},
    {*r} {*á} -> {*á} {*r} || EnglishStarConsonant _ {*s} {*t}
];
\end{ReaderFacingFoma}

The rule moves \emph{*r} across a following short vowel in the relevant
late clusters, producing forms such as \emph{berstan} `burst' where an
earlier order would still show a broken vowel sequence.

Moving the rule before \hyperref[rule-OEBreaking]{SC044 OEBreaking}
makes PGmc \Recon{bréstaną} `burst' yield \Pred{beorstan} rather than
expected OE \emph{berstan} `burst'. On this evidence, I take
\hyperref[rule-OERMetathesis]{SC087 OERMetathesis} to follow
\hyperref[rule-OEBreaking]{SC044 OEBreaking}. Moving it later within the
tested sequence alters no output.

The lexical evidence fixes the earlier relation but does not identify a
corresponding later constraint. For the selected West Saxon history, the
sources support a late rearrangement after breaking without placing it
immediately beside contraction.

\section{Late unstressed suffix
raising}\label{late-unstressed-suffix-raising}

\subsection{Historical discussion}\label{historical-discussion-63}

The late history of unstressed \emph{*-ag-} differs from prehistoric
i-mutation. Ringe and Taylor give the sequence through fronting and
raising to \emph{-ig-}, explicitly placing its final stage long after
mutation. Their examples include \emph{huniġ} `honey'; a velar
intermediate can survive before a back vowel
(\citeproc{ref-RingeTaylor2014}{Ringe and Taylor 2014, 334--35},
§6.9.6). These are the printed pages, not the PDF sheet labels.

\subsection{\texorpdfstring{SC089. Late suffix consumer
(\texttt{OELateUnstressedAgSuffix})}{SC089. Late suffix consumer (OELateUnstressedAgSuffix)}}\label{rule-OELateUnstressedAgSuffix}

\begin{ReaderFacingFoma}
define OELateUnstressedAgSuffix (
    [{*a} -> {*e} ||
        EnglishStarVocalic [EnglishStarConsonant | EnglishPalatalConsonant]+ _ {*g}]
    .o.
    [{*g} -> {*ʝ} ||
        EnglishStarVocalic [EnglishStarConsonant | EnglishPalatalConsonant]+ {*e} _ .#.]
    .o.
    [{*e} -> {*i} ||
        EnglishStarVocalic [EnglishStarConsonant | EnglishPalatalConsonant]+ _ [{*ʝ}|{*ʤ}]]
);
\end{ReaderFacingFoma}

The singleton consonant is a palatal fricative, not an early affricate
(\citeproc{ref-Fulk2018}{Fulk 2018, 130--32}). The raising consumer
therefore recognizes ʝ alongside retained stop-reflex compatibility.
This does not make ʝ a prehistoric mutation trigger: later raising and
earlier i-mutation have different conditioning histories.

The present late applications are \emph{huniġ} `honey' and \emph{wīþiġ}
`withy'. The word-final guard remains narrow; this implementation is not
an exhaustive model of every inflected suffix form.
\hyperref[rule-OEPalatalFricativeMerger]{SC109 OEPalatalFricativeMerger}
follows here as a disclosed serialization, not proof of a uniquely dated
historical merger immediately after this raising
(\citeproc{ref-RingeTaylor2014}{Ringe and Taylor 2014, 204},
pp.~334--335).

\section{Technical support: postmutation fricative
merger}\label{technical-support-postmutation-fricative-merger}

\subsection{\texorpdfstring{\CAPRRuleHeading{SC109. Postmutation merger serialization}{OEPalatalFricativeMerger}}{}}\label{rule-OEPalatalFricativeMerger}

\noindent\begin{minipage}{\linewidth}
\begin{ReaderFacingFoma}
define OEPalatalFricativeMerger [
    {*ʝ} -> {*j}
];
\end{ReaderFacingFoma}
\end{minipage}

The historical commitment is merger after mutation
(\citeproc{ref-RingeTaylor2014}{Ringe and Taylor 2014, 204}). Its
execution after \hyperref[rule-OELateUnstressedAgSuffix]{SC089
OELateUnstressedAgSuffix} is a computational holding zone, not proof of
an exact historical date. It keeps the fricative out of inherited-j
normalization while preserving late palatal raising, syllable weight,
reduction and native rendering.
\hyperref[rule-OEJClusterCoalescence]{SC057 OEJClusterCoalescence}
remains a different cluster process. Stop affrication dates and the
precise disputed eligibility domain are not settled by this merger.

\newpage

\chapter{Old English orthography and the written
surface}\label{old-english-orthography-and-the-written-surface}

\section{Historical interval}\label{historical-interval-4}

This short chapter stands apart from the derivational chapters that
precede it. The earlier chapters primarily discuss changes to the spoken
language, alongside explicitly identified technical support operations.
The material treated here belongs instead to the written surface of Old
English --- scribal conventions that determined how the results of the
completed phonological history were committed to parchment.

In the executable model these conventions apply after every phonological
rule, in the final written-surface block. That serialization renders the
results of the phonological reconstruction; it does not independently
date the origin of each spelling practice relative to every phonological
change.

\section{Scope}\label{scope}

The one rule treated here is the West Saxon palatal-glide spelling
(SC016), by which back vowels following word-initial {[}j{]} --- spelled
\emph{g} --- came to be written with a preceding front glide letter, as
in \emph{geoc} `yoke' for spoken {[}jok{]} and \emph{geoguþ} `youth' for
a form whose root vowel remained {[}u{]}. Ringe and Taylor state the
modern assessment directly: the \emph{eo} of \emph{geoc} is a spelling
convention, and the word was pronounced {[}jok{]}
(\citeproc{ref-RingeTaylor2014}{Ringe and Taylor 2014, 5}). Hogg reaches
the same verdict for the back-vowel cases generally
(\citeproc{ref-HoggPhonology1992}{Hogg 1992b, 113}). The older handbooks
--- Campbell, Brunner, Bülbring, Luick --- analysed the same spellings
as rising diphthongs; the section below presents both views
(\citeproc{ref-Campbell1959}{Campbell 1959, 17}, § 44;
\citeproc{ref-SieversBrunner1965}{Brunner 1965, 64--65}, § 92).

\section{Sources}\label{sources-1}

Ringe and Taylor provide the modern phonological interpretation
(\citeproc{ref-RingeTaylor2014}{Ringe and Taylor 2014, 5}). Campbell
(\citeproc{ref-Campbell1959}{Campbell 1959, 17, 64--67}, §§ 44,
170--176), Brunner (\citeproc{ref-SieversBrunner1965}{Brunner 1965,
64--65}, § 92), and Bülbring (\citeproc{ref-Bulbring1902}{Bülbring 1902,
120}, §§ 298--299) document the distribution of the spellings; Hogg
supplies the critical reassessment
(\citeproc{ref-HoggPhonology1992}{Hogg 1992b, 113}).

Barrack's West Saxon cascade registers the same written-surface facts at
two points: the \emph{iung}/\emph{giong}/\emph{geong} `young' spellings
under his rule 171b (\citeproc{ref-Barrack1975}{Barrack 1975, 58}), and
an Early West Saxon spelling variation, with \emph{Geoworpa} and
\emph{Iugurtha} among the cited forms, as a Campbell-footnote rule 430n3
(\citeproc{ref-Barrack1975}{Barrack 1975, 110}). He treats these within
his free-variation framework rather than as a single orthographic
convention; that framing is a rival description of the same evidence,
not support for a spoken rising diphthong.

\section{West Saxon palatal-glide spelling before back
vowels}\label{west-saxon-palatal-glide-spelling-before-back-vowels}

\subsection{Historical discussion}\label{historical-discussion-64}

West Saxon spellings such as \emph{ġeoc} `yoke', \emph{ġeong} `young',
and \emph{ġeoguþ} `youth' write a front glide letter between a
word-initial palatal and a following back vowel. Campbell describes the
phenomenon as the development of rising diphthongs when ``palatal glides
developed before back vowels'' and cites \emph{ġeoc} directly
(\citeproc{ref-Campbell1959}{Campbell 1959, 17}, §44); Brunner separates
the \emph{u}-cases (\emph{ġeong}, \emph{ġeoguþ}) from the \emph{o}-cases
(\emph{ġioc} `yoke', \emph{ġeoc})
(\citeproc{ref-SieversBrunner1965}{Brunner 1965, 64--65}, §92.1);
Bülbring likewise treats \emph{iuguð} `youth' and \emph{iuc} under
\emph{ju} but derives \emph{ġioc}, \emph{ġeoc} from West Germanic
\emph{*jok} (\citeproc{ref-Bulbring1902}{Bülbring 1902, 120},
§§298--299); and Luick groups all of these under his ``schwebende
Diphthonge'' after palatal onsets (\citeproc{ref-Luick1914}{Luick
1914-\/-1940, 158--59}, §169).

The phonological interpretation of these spellings is disputed. The
older handbook tradition --- Campbell, Brunner, Bülbring, Luick ---
reads them as genuine rising diphthongs. The modern assessment is
orthographic: Ringe and Taylor state flatly that \emph{ġeoc} ``is
/jok/'', the digraph being a spelling convention that became universal
after word-initial /j/ (\citeproc{ref-RingeTaylor2014}{Ringe and Taylor
2014, 5}), and Hogg concludes that the back-vowel cases were ``never
anything more than an orthographic variation'', judging Campbell's
arguments to the contrary ``insubstantial''
(\citeproc{ref-HoggPhonology1992}{Hogg 1992b, 113};
\citeproc{ref-Campbell1959}{Campbell 1959, 66--67}, §176). This model
follows Ringe and Taylor and Hogg: the rule is a spelling convention
applied to the finished phonology, and it therefore stands at the end of
the derivation, in the written-surface stage of the cascade.

Its position also settles a relative chronology. The \emph{o} of
\emph{ġeoc} `yoke' is itself the product of Northwest Germanic
u-lowering (\hyperref[rule-PNWGmcULowering]{SC017 PNWGmcULowering}):
Fulk lists \emph{ġeoc} as a regular lowering example beside OIcel
\emph{ok} and OHG \emph{joh} (\citeproc{ref-Fulk2018}{Fulk 2018, 56},
§4.3), and Campbell gives \emph{ġeoc} among the regular \emph{u}
\textgreater{} \emph{o} words (\citeproc{ref-Campbell1959}{Campbell
1959, 43}, §115). The lowering therefore feeds the spelling: first
\emph{*juk-} became \emph{*jok-} in Northwest Germanic, and only much
later did West Saxon scribes write the result as \emph{ġeoc}. Where
lowering did not apply, as in \emph{ġeoguþ} `youth', whose root \emph{u}
was protected by the high vowel of the following syllable, the same
convention wrote the retained \emph{u} with the same digraph
(\citeproc{ref-SieversBrunner1965}{Brunner 1965, 64--65}, §92.1).

Barrack registers the same spellings at two points of his cascade: the
\emph{iung}/\emph{giong}/\emph{geong} `young' alternation under his rule
171b (\citeproc{ref-Barrack1975}{Barrack 1975, 58}), and an Early West
Saxon variation, with \emph{Geoworpa} and \emph{Iugurtha} among the
cited forms, as the Campbell-footnote rule 430n3
(\citeproc{ref-Barrack1975}{Barrack 1975, 110}). Writing before Hogg's
reassessment, he handles the alternation inside his free-variation
framework rather than as a convention of the written surface
(\citeproc{ref-Barrack1975}{Barrack 1975, 21--23}); his treatment is
thus a formal restatement of the older rising-diphthong reading, and the
orthographic interpretation adopted here follows Ringe and Taylor and
Hogg against it.

\subsection{\texorpdfstring{\CAPRRuleHeading{SC016. West Saxon palatal-glide spelling before back vowels}{OEWsPalatalGlide}}{}}\label{rule-OEWsPalatalGlide}

\noindent\begin{minipage}{\linewidth}
\begin{ReaderFacingFoma}
define OEWsPalatalGlide [
    {*ó} -> {*éo} || .#. ġ _ ,
    {*ú} -> {*éo} || .#. ġ _ ,
    {*o} -> {*eo} || .#. ġ _ ,
    {*u} -> {*eo} || .#. ġ _
];
\end{ReaderFacingFoma}
\end{minipage}

The rule rewrites a back vowel after word-initial \emph{ġ} as the
digraph spelling, covering both the lowered \emph{o}-cases (\emph{ġeoc}
`yoke') and the retained \emph{u}-cases (\emph{ġeoguþ} `youth'). Because
it is a convention of the written language, it applies after every
phonological change; in particular it follows
\hyperref[rule-PNWGmcULowering]{SC017 PNWGmcULowering}, which supplies
the \emph{o} of \emph{ġeoc}. If the spelling rule were placed before
u-lowering, the derivation would have to treat an Old English scribal
practice as a Northwest Germanic sound change, an ordering that no
source supports. The witnesses \emph{ġeoc} and \emph{ġeoguþ} between
them fix both faces of the rule: one shows the convention applied to
lowered \emph{o}, the other to unlowered \emph{u}. The handbook domain
is broader (it also includes \emph{a}/\emph{ā}/\emph{ō} contexts after
word-initial palatals), but this executable rule is intentionally
complete for the currently selected corpus witnesses rather than a
maximal dialectal enumeration.

\newpage

\part{Lexical derivations}

\chapter*{Word-by-word derivations}\label{word-by-word-derivations}
\addcontentsline{toc}{chapter}{Word-by-word derivations}

\section*{Introduction}\label{introduction-1}
\addcontentsline{toc}{section}{Introduction}

The catalogue groups words by the kind of historical explanation they
require. Regular derivations establish the reach of the sound laws. The
remaining classes mark, without concealment, the places where
attestation, morphology, analogy, or an unresolved mismatch intervenes.

The catalogue proceeds from words; the sound-change account proceeds
from rules. Together they allow each proposed sound law to be tested
against its lexical witnesses and each etymology against the complete
ordered history. SC-numbered rules cited in these entries are discussed
in the companion sound-change volume, or in the sound-change chapters of
the combined book.

\section*{Data and sources}\label{data-and-sources}
\addcontentsline{toc}{section}{Data and sources}

This volume assembles the lexical corpus from the aligned Germanic
dataset and the compact derivation traces that accompany each entry.
Comparative dictionaries, Old English dictionaries, and historical
grammars are cited in the prose where they bear on particular lexical
arguments.

The result is a lexical catalogue rather than a separate report on
citation method or trace machinery.

\section*{Transducer and derivation
method}\label{transducer-and-derivation-method}
\addcontentsline{toc}{section}{Transducer and derivation method}

Four objects must be distinguished in every derivation: the citation
reconstruction, the selected input, the transducer outcome, and the Old
English target. The summary identifies them where they differ; the boxed
trace then divides the changes into Earlier Germanic and Old English
stages.

\section*{Derivation classes}\label{derivation-classes}
\addcontentsline{toc}{section}{Derivation classes}

The lexical catalogue is ordered by seven derivation classes and, within
each class, alphabetically by headword. Counts in this alpha are:

\begin{itemize}
\tightlist
\item
  Regular derivations: 97
\item
  Attested variants: 6
\item
  Early analogy: 36
\item
  Late analogy: 28
\item
  Reconstructed Old English comparators: 3
\item
  Known but unmodelled developments: 2
\item
  Unexplained or deliberately unmodelled exceptions: 5
\end{itemize}

\clearpage

\chapter{Regular derivations}\label{regular-derivations}

In these entries the selected Germanic input yields the Old English
reflex by regular sound change. They establish the standard against
which the following analogical and exceptional histories must be judged.

\section{\texorpdfstring{adder --- OE
\emph{nǣdre}}{adder --- OE nǣdre}}\label{adder-oe-nux1e3dre}

\index[iv]{02oe@\ivlangheader{Old English}{West Saxon normalization unmarked}!naedre@\iventry{nǣdre}{}}
\index[iv]{03pgmc@\ivlangheader{Proto-Germanic}{}!nedron@\iventry{*nḗdrōn}{}}

Derivation: \emph{*nḗdrōn} \textgreater{} \emph{nǣdre} (regular).

\subsection{Derivation trace}\label{derivation-trace}

Proto input: \emph{*nḗdrōn}

\begingroup
\setlength{\fboxsep}{6pt}

\noindent\fbox{%
\begin{minipage}{0.97\linewidth}
\small
\begin{tabular}{@{}>{\raggedright\arraybackslash}p{0.485\linewidth}@{\hspace{0.030\linewidth}}>{\raggedright\arraybackslash}p{0.485\linewidth}@{}}
\begin{minipage}[t]{\linewidth}
\centering\textbf{Earlier Germanic changes}\par
\vspace{0.35em}
\raggedright
\centering\textbf{Northwest and West Germanic}\par
\raggedright
\vspace{0.2em}
\begin{tabularx}{\linewidth}{@{}>{\raggedright\arraybackslash}X@{\hspace{0.55em}}>{\raggedright\arraybackslash}p{0.14\linewidth}@{\hspace{0.25em}}}
\mbox{PNWGmc Long E Lowering} & \emph{*nādrōn} \\
\mbox{PNWGmc N Stem N Loss} & \emph{*nādrǭ} \\
\mbox{EAF Long A Fronting} & \emph{*nǣdrǭ} \\
\end{tabularx}
\vspace{0.6em}
\centering\textbf{Early Anglo-Frisian}\par
\raggedright
\vspace{0.2em}
\raggedright [no change]\par
\end{minipage}
&
\begin{minipage}[t]{\linewidth}
\centering\textbf{Old English changes}\par
\vspace{0.35em}
\raggedright
\begin{tabularx}{\linewidth}{@{}>{\raggedright\arraybackslash}X@{\hspace{0.55em}}>{\raggedright\arraybackslash}p{0.14\linewidth}@{\hspace{0.25em}}}
OE Unstressed Long Vowel Shortening & \emph{*nǣdræ} \\
\mbox{OE Unstressed AE Merger} & \emph{*nǣdre} \\
\end{tabularx}
\end{minipage}
\\
\end{tabular}
\end{minipage}%
} \endgroup

Old English form: \emph{nǣdre}

\subsection{Reconstruction and comparative
evidence}\label{reconstruction-and-comparative-evidence}

Kroonen distinguishes the masculine snake word
{\emph{*nadra-}}\index[iv]{03pgmc@\ivlangheader{Proto-Germanic}{}!nadra@\iventry{*nadra-}{}}
from a feminine ablauting formation
{\emph{*nēdrōn-}}\index[iv]{03pgmc@\ivlangheader{Proto-Germanic}{}!nedron@\iventry{*nēdrōn-}{}},
and gives Old English
{\emph{nǣdre}}\index[iv]{02oe@\ivlangheader{Old English}{West Saxon normalization unmarked}!naedre@\iventry{nǣdre}{}}
`adder',
{\emph{næddre}}\index[iv]{02oe@\ivlangheader{Old English}{West Saxon normalization unmarked}!naeddre@\iventry{næddre}{}}
`adder' under the latter (\citeproc{ref-Kroonen2013}{Kroonen 2013, 381,
386}). Orel's feminine entry gives
\Recon{nēdrōn}\index[iv]{03pgmc@\ivlangheader{Proto-Germanic}{}!nedron@\iventry{*nēdrōn}{}}
`adder' \textasciitilde{}
\Recon{nadrōn}\index[iv]{03pgmc@\ivlangheader{Proto-Germanic}{}!nadron@\iventry{*nadrōn}{}}
`adder' and relates them to the masculine formation. These comparison
spellings use plain \emph{d} for Orel's printed barred \emph{đ}; the
quantity and the two alternatives are retained
(\citeproc{ref-Orel2003}{Orel 2003, 286}).

The derivational input retains the long-\emph{e} feminine alternative.
The lexical comparison supports that formation without erasing Orel's
alternative or the distinction between his notation and CAPR's
derivational input.

\subsection{Reconstruction alternatives and their
limits}\label{reconstruction-alternatives-and-their-limits}

Orel's short-\emph{a} feminine is an admitted alternative within his own
account, not proof that he and Kroonen advocate exclusive histories. The
unresolved premise is the reflex distribution and formation history
licensing that alternative beside the selected long-\emph{ē} feminine
(\citeproc{ref-Orel2003}{Orel 2003, 286};
\citeproc{ref-Kroonen2013}{Kroonen 2013, 381, 386}). Kroonen's uncertain
sewing or `winder' connection and masculine comparanda cannot substitute
for that feminine argument. CAPR's acute marks its input stress; it is
not independently quoted in these headwords.

Fulk's probable shortening accompanying pre-\emph{r} gemination concerns
the later Old English variant. He calls the motivation obscure and
considers analogical extension; this does not explain inherited \emph{ē}
versus \emph{a} or supply another complete ancestral adder word
(\citeproc{ref-Fulk2018}{Fulk 2018, 132}). Regular production of
\emph{nǣdre} therefore supports the selected path without settling the
alternative formation.

\subsection{Old English evidence}\label{old-english-evidence}

The Old English word is securely represented by
{\emph{nǣdre}}\index[iv]{02oe@\ivlangheader{Old English}{West Saxon normalization unmarked}!naedre@\iventry{nǣdre}{}}
`adder', with
{\emph{næddre}}\index[iv]{02oe@\ivlangheader{Old English}{West Saxon normalization unmarked}!naeddre@\iventry{næddre}{}}
`adder' as a secondary variant. Clark Hall cross-references
{\emph{næddre}}\index[iv]{02oe@\ivlangheader{Old English}{West Saxon normalization unmarked}!naeddre@\iventry{næddre}{}}
to
{\emph{nædre}}\index[iv]{02oe@\ivlangheader{Old English}{West Saxon normalization unmarked}!naedre@\iventry{nædre}{}}
`adder', and Fulk treats
{\emph{næddre}}\index[iv]{02oe@\ivlangheader{Old English}{West Saxon normalization unmarked}!naeddre@\iventry{næddre}{}}
as the later geminated form beside the older base
(\citeproc{ref-ClarkHall1960}{Clark Hall 1960, 225};
\citeproc{ref-Fulk2018}{Fulk 2018, 132}).

\subsection{Development to Old
English}\label{development-to-old-english}

From
\Recon{nḗdrōn}\index[iv]{03pgmc@\ivlangheader{Proto-Germanic}{}!nedron@\iventry{*nḗdrōn}{}}
`adder', the stressed long mid vowel develops to Old English
{\emph{nǣdre}}\index[iv]{02oe@\ivlangheader{Old English}{West Saxon normalization unmarked}!naedre@\iventry{nǣdre}{}}
`adder', and the weak feminine ending remains as final \emph{-e}, giving
{\emph{nǣdre}}\index[iv]{02oe@\ivlangheader{Old English}{West Saxon normalization unmarked}!naedre@\iventry{nǣdre}{}}.
The doubled consonant of
{\emph{næddre}}\index[iv]{02oe@\ivlangheader{Old English}{West Saxon normalization unmarked}!naeddre@\iventry{næddre}{}}
`adder' is secondary and does not alter the inherited base form.

\section{\texorpdfstring{bake --- OE
\emph{bacan}}{bake --- OE bacan}}\label{bake-oe-bacan}

\index[iv]{02oe@\ivlangheader{Old English}{West Saxon normalization unmarked}!bacan@\iventry{bacan}{}}
\index[iv]{03pgmc@\ivlangheader{Proto-Germanic}{}!bakana@\iventry{*bákaną}{}}
\index[iv]{10ohg@\ivlangheader{Old High German}{}!backan@\iventry{backan}{}}
\index[iv]{10ohg@\ivlangheader{Old High German}{}!bahhan@\iventry{bahhan}{}}

Derivation: \emph{*bákaną} \textgreater{} \emph{bacan} (regular).

\subsection{Derivation trace}\label{derivation-trace-1}

Proto input: \emph{*bákaną}

\begingroup
\setlength{\fboxsep}{6pt}

\noindent\fbox{%
\begin{minipage}{0.97\linewidth}
\small
\begin{tabular}{@{}>{\raggedright\arraybackslash}p{0.359\linewidth}@{\hspace{0.030\linewidth}}>{\raggedright\arraybackslash}p{0.611\linewidth}@{}}
\begin{minipage}[t]{\linewidth}
\centering\textbf{Earlier Germanic changes}\par
\vspace{0.35em}
\raggedright
\centering\textbf{Northwest and West Germanic}\par
\raggedright
\vspace{0.2em}
\raggedright [no change]\par
\vspace{0.6em}
\centering\textbf{Early Anglo-Frisian}\par
\raggedright
\vspace{0.2em}
\raggedright [no change]\par
\end{minipage}
&
\begin{minipage}[t]{\linewidth}
\centering\textbf{Old English changes}\par
\vspace{0.35em}
\raggedright
\begin{tabularx}{\linewidth}{@{}>{\raggedright\arraybackslash}X@{\hspace{0.55em}}>{\raggedright\arraybackslash}p{0.14\linewidth}@{\hspace{0.25em}}}
\mbox{EAF Brightening} & \emph{*bækaną} \\
\mbox{OE A Restoration} & \emph{*bakaną} \\
OE Heavy Syllable Nasal Apocope & \emph{*bakan} \\
\mbox{OE Secondary Nasalization} & \emph{*bakąn} \\
\mbox{OE Weak Tail Reduction} & \emph{*bakan} \\
\end{tabularx}
\end{minipage}
\\
\end{tabular}
\end{minipage}%
} \endgroup

Old English form: \emph{bacan}

\subsection{Reconstruction and comparative
evidence}\label{reconstruction-and-comparative-evidence-1}

Orel reconstructs the verb as
\Recon{bakanan}\index[iv]{03pgmc@\ivlangheader{Proto-Germanic}{}!bakanan@\iventry{*bakanan}{}}
`bake' and cites Old English
{\emph{bacan}}\index[iv]{02oe@\ivlangheader{Old English}{West Saxon normalization unmarked}!bacan@\iventry{bacan}{}}
`bake' beside Old High German \emph{backan, bahhan}
(\citeproc{ref-Orel2003}{Orel 2003}). Campbell gives
{\emph{bacan}}\index[iv]{02oe@\ivlangheader{Old English}{West Saxon normalization unmarked}!bacan@\iventry{bacan}{}}
as one of the standard examples of Old English A-restoration before a
single consonant, and Ringe and Taylor state the same development from
\Recon{bakan}\index[iv]{05pwgmc@\ivlangheader{Proto-West Germanic}{}!bakan@\iventry{*bakan}{}}
`bake' to Old English
{\emph{bacan}}\index[iv]{02oe@\ivlangheader{Old English}{West Saxon normalization unmarked}!bacan@\iventry{bacan}{}}
(\citeproc{ref-Campbell1959}{Campbell 1959, 61};
\citeproc{ref-RingeTaylor2014}{Ringe and Taylor 2014}).

\subsection{Singleton and geminate present
alternatives}\label{singleton-and-geminate-present-alternatives}

The selected infinitive keeps singleton \emph{k}. Kroonen explicitly
distinguishes \emph{*bakan-} beside \emph{*bakkan-}, considering
contamination with, or back-formation from, the iterative
\emph{*bak(k)ōn-}; a present \emph{n}-suffix explanation is less likely
(\citeproc{ref-Kroonen2013}{Kroonen 2013, 49}). Orel prints singleton
\emph{*bakanan} and separately the long-\emph{ō} preterite
(\citeproc{ref-Orel2003}{Orel 2003, 33}). The inspected passages do not
establish Orel's reason for excluding the geminate from this citation.
This is a qualified present-formation question, not proof of a mutually
exclusive reconstruction debate.

Later finite-cell front vowels and the inherited preterite grade are not
rival infinitive vowels. Ringe and Taylor distinguish West Saxon
infinitive \emph{bacan} `bake' from second singular \emph{becst} `you
bake', preserving the umlauted vowel in syncopated forms while \emph{a}
spread elsewhere (\citeproc{ref-RingeTaylor2014}{Ringe and Taylor 2014,
232--33}). Regular infinitive production does not resolve the earlier
singleton/geminate distribution.

\subsection{Old English evidence}\label{old-english-evidence-1}

Bosworth-Toller and Clark Hall both record
{\emph{bacan}}\index[iv]{02oe@\ivlangheader{Old English}{West Saxon normalization unmarked}!bacan@\iventry{bacan}{}}
`bake' as the ordinary Old English verb
(\citeproc{ref-BosworthToller1898}{Bosworth and Toller 1898, 72};
\citeproc{ref-ClarkHall1960}{Clark Hall 1960}). The target in this entry
is therefore the attested infinitive headword itself, not a selected
oblique or finite paradigm cell.

\subsection{Development to Old
English}\label{development-to-old-english-1}

From
\Recon{bákaną}\index[iv]{03pgmc@\ivlangheader{Proto-Germanic}{}!bakana@\iventry{*bákaną}{}}
`bake', Anglo-Frisian brightening first gives \Recon{bækaną} `bake'.
A-restoration then returns the stem vowel to \emph{a} before single
\emph{k} plus the back-vocalic infinitive suffix, and later apocope and
weak-tail reduction yield
{\emph{bacan}}\index[iv]{02oe@\ivlangheader{Old English}{West Saxon normalization unmarked}!bacan@\iventry{bacan}{}}
`bake' (\citeproc{ref-Campbell1959}{Campbell 1959, 61};
\citeproc{ref-RingeTaylor2014}{Ringe and Taylor 2014}). The development
is therefore straightforward:
\Recon{bákaną}\index[iv]{03pgmc@\ivlangheader{Proto-Germanic}{}!bakana@\iventry{*bákaną}{}}
\textgreater{}
{\emph{bacan}}\index[iv]{02oe@\ivlangheader{Old English}{West Saxon normalization unmarked}!bacan@\iventry{bacan}{}}.

\section{\texorpdfstring{beech --- OE
\emph{bōc}}{beech --- OE bōc}}\label{beech-oe-bux14dc}

\index[iv]{02oe@\ivlangheader{Old English}{West Saxon normalization unmarked}!boc@\iventry{boc}{}}
\index[iv]{02oe@\ivlangheader{Old English}{West Saxon normalization unmarked}!boc@\iventry{bōc}{}}
\index[iv]{03pgmc@\ivlangheader{Proto-Germanic}{}!boko@\iventry{*bōkō}{}}

Derivation: \emph{*bōkō} \textgreater{} \emph{bōc} (regular).

\subsection{Derivation trace}\label{derivation-trace-2}

Proto input: \emph{*bōkō}

\begingroup
\setlength{\fboxsep}{6pt}

\noindent\fbox{%
\begin{minipage}{0.97\linewidth}
\small
\begin{tabular}{@{}>{\raggedright\arraybackslash}p{0.611\linewidth}@{\hspace{0.030\linewidth}}>{\raggedright\arraybackslash}p{0.359\linewidth}@{}}
\begin{minipage}[t]{\linewidth}
\centering\textbf{Earlier Germanic changes}\par
\vspace{0.35em}
\raggedright
\centering\textbf{Northwest and West Germanic}\par
\raggedright
\vspace{0.2em}
\begin{tabularx}{\linewidth}{@{}>{\raggedright\arraybackslash}X@{\hspace{0.55em}}>{\raggedright\arraybackslash}p{0.14\linewidth}@{\hspace{0.25em}}}
\mbox{PNWGmc Final Long O Raising} & \emph{*bōku} \\
\end{tabularx}
\vspace{0.6em}
\centering\textbf{Early Anglo-Frisian}\par
\raggedright
\vspace{0.2em}
\raggedright [no change]\par
\end{minipage}
&
\begin{minipage}[t]{\linewidth}
\centering\textbf{Old English changes}\par
\vspace{0.35em}
\raggedright
\begin{tabularx}{\linewidth}{@{}>{\raggedright\arraybackslash}X@{\hspace{0.55em}}>{\raggedright\arraybackslash}p{0.14\linewidth}@{\hspace{0.25em}}}
\mbox{OE High Vowel Apocope} & \emph{*bōk} \\
\end{tabularx}
\end{minipage}
\\
\end{tabular}
\end{minipage}%
} \endgroup

Old English form: \emph{bōc}

\subsection{Reconstruction and comparative
evidence}\label{reconstruction-and-comparative-evidence-2}

Kroonen gives the beech noun as \emph{*bōk(j)ō-} and cites Old English
\emph{boc} `beech',
{\emph{bēce}}\index[iv]{02oe@\ivlangheader{Old English}{West Saxon normalization unmarked}!bece@\iventry{bēce}{}}
`beech' among its reflexes (\citeproc{ref-Kroonen2013}{Kroonen 2013}).
The derivational input \Recon{bōkō} `beech' is the nominative-singular
shape of that family, which is the relevant comparison form here.

\subsection{Old English evidence}\label{old-english-evidence-2}

Kroonen's Old English evidence already separates the paradigm material:
\emph{boc} `beech' as the nominative form and
{\emph{bēce}}\index[iv]{02oe@\ivlangheader{Old English}{West Saxon normalization unmarked}!bece@\iventry{bēce}{}}
`beech' as an oblique form (\citeproc{ref-Kroonen2013}{Kroonen 2013}).
The relevant comparator is therefore \emph{bōc} `beech';
{\emph{bēċe}}\index[iv]{02oe@\ivlangheader{Old English}{West Saxon normalization unmarked}!bece@\iventry{bēċe}{}}
`beech' remains related paradigm evidence rather than the form chosen
for this comparison.

\subsection{Development to Old
English}\label{development-to-old-english-2}

With nominative input \Recon{bōkō} `beech', the development is compact.
Northwest Germanic final long \emph{ō} raises to \emph{u}, and later
high-vowel apocope leaves \emph{bōc} `beech'. The regular comparison is
therefore \emph{*bōkō} \textgreater{} \emph{bōc}.

\section{\texorpdfstring{begin --- OE
\emph{beġinnan}}{begin --- OE beġinnan}}\label{begin-oe-beux121innan}

\index[iv]{02oe@\ivlangheader{Old English}{West Saxon normalization unmarked}!beginnan@\iventry{beġinnan}{}}
\index[iv]{03pgmc@\ivlangheader{Proto-Germanic}{}!biginnana@\iventry{*bigínnaną}{}}

Derivation: \emph{*bigínnaną} \textgreater{} \emph{beġinnan} (regular).

\subsection{Derivation trace}\label{derivation-trace-3}

Proto input: \emph{*bigínnaną}

\begingroup
\setlength{\fboxsep}{6pt}

\noindent\fbox{%
\begin{minipage}{0.97\linewidth}
\footnotesize
\begin{tabular}{@{}>{\raggedright\arraybackslash}p{0.339\linewidth}@{\hspace{0.030\linewidth}}>{\raggedright\arraybackslash}p{0.630\linewidth}@{}}
\begin{minipage}[t]{\linewidth}
\centering\textbf{Earlier Germanic changes}\par
\vspace{0.35em}
\raggedright
\centering\textbf{Northwest and West Germanic}\par
\raggedright
\vspace{0.2em}
\raggedright [no change]\par
\vspace{0.6em}
\centering\textbf{Early Anglo-Frisian}\par
\raggedright
\vspace{0.2em}
\raggedright [no change]\par
\end{minipage}
&
\begin{minipage}[t]{\linewidth}
\centering\textbf{Old English changes}\par
\vspace{0.35em}
\raggedright
\begin{tabularx}{\linewidth}{@{}>{\raggedright\arraybackslash}X@{\hspace{0.55em}}>{\raggedright\arraybackslash}p{0.22\linewidth}@{\hspace{0.25em}}}
OE Heavy Syllable Nasal Apocope & \emph{*bigínnan} \\
\mbox{OE Secondary Nasalization} & \emph{*bigínnąn} \\
\mbox{OE Velar Palatalization} & \emph{*biʝínnąn} \\
\mbox{OE Prefix I Reduction} & \emph{*beʝínnąn} \\
\mbox{OE Weak Tail Reduction} & \emph{*beʝínnan} \\
\mbox{OE Palatal Fricative Merger} & \emph{*bejínnan} \\
\end{tabularx}
\end{minipage}
\\
\end{tabular}
\end{minipage}%
} \endgroup

Old English form: \emph{beġinnan}

\subsection{Reconstruction and comparative
evidence}\label{reconstruction-and-comparative-evidence-3}

The verb is modeled here as inherited
\Recon{bigínnaną}\index[iv]{03pgmc@\ivlangheader{Proto-Germanic}{}!biginnana@\iventry{*bigínnaną}{}}
`begin'. Ringe and Taylor state that intervocalic \emph{*g} is
palatalized between front vowels in Old English
(\citeproc{ref-RingeTaylor2014}{Ringe and Taylor 2014}), and Campbell
lists
{\emph{ginnan}}\index[iv]{02oe@\ivlangheader{Old English}{West Saxon normalization unmarked}!ginnan@\iventry{ginnan}{}}
`begin' among familiar examples of palatal \emph{g} in this verb family
(\citeproc{ref-Campbell1959}{Campbell 1959, 174}).

\subsection{Old English evidence}\label{old-english-evidence-3}

Bosworth-Toller and Clark Hall lemmatize the verb as
{\emph{be-ginnan}}\index[iv]{02oe@\ivlangheader{Old English}{West Saxon normalization unmarked}!beginnan@\iventry{be-ginnan}{}}
`begin' /
{\emph{beginnan}}\index[iv]{02oe@\ivlangheader{Old English}{West Saxon normalization unmarked}!beginnan@\iventry{beginnan}{}}
`begin' (\citeproc{ref-BosworthToller1898}{Bosworth and Toller 1898,
84}; \citeproc{ref-ClarkHall1960}{Clark Hall 1960}). Those
plain-\emph{g} dictionary spellings support the same verb that appears
here in normalized form as
{\emph{beġinnan}}\index[iv]{02oe@\ivlangheader{Old English}{West Saxon normalization unmarked}!beginnan@\iventry{beġinnan}{}}
`begin'.

\subsection{Development note}\label{development-note}

The prefix deserves separate notice. Ringe and Taylor explicitly cite
\emph{bi- \textgreater{} be-} as an Old English unstressed-prefix
development (\citeproc{ref-RingeTaylor2014}{Ringe and Taylor 2014}).

\subsection{Development to Old
English}\label{development-to-old-english-3}

From
\Recon{bigínnaną}\index[iv]{03pgmc@\ivlangheader{Proto-Germanic}{}!biginnana@\iventry{*bigínnaną}{}}
`begin', heavy-syllable nasal apocope yields \Recon{bigínnan} `begin'.
Intervocalic \emph{*g} between front vowels then palatalizes to
\emph{ġ}, and the unstressed prefix reduces \emph{bi-} to \emph{be-},
giving
{\emph{beġinnan}}\index[iv]{02oe@\ivlangheader{Old English}{West Saxon normalization unmarked}!beginnan@\iventry{beġinnan}{}}
`begin'.

\section{\texorpdfstring{bier --- OE
\emph{bǣr}}{bier --- OE bǣr}}\label{bier-oe-bux1e3r}

\index[iv]{02oe@\ivlangheader{Old English}{West Saxon normalization unmarked}!baer@\iventry{bǣr}{}}
\index[iv]{03pgmc@\ivlangheader{Proto-Germanic}{}!bero@\iventry{*bḗrō}{}}

Derivation: \emph{*bḗrō} \textgreater{} \emph{bǣr} (regular).

\subsection{Derivation trace}\label{derivation-trace-4}

Proto input: \emph{*bḗrō}

\begingroup
\setlength{\fboxsep}{6pt}

\noindent\fbox{%
\begin{minipage}{0.97\linewidth}
\small
\begin{tabular}{@{}>{\raggedright\arraybackslash}p{0.611\linewidth}@{\hspace{0.030\linewidth}}>{\raggedright\arraybackslash}p{0.359\linewidth}@{}}
\begin{minipage}[t]{\linewidth}
\centering\textbf{Earlier Germanic changes}\par
\vspace{0.35em}
\raggedright
\centering\textbf{Northwest and West Germanic}\par
\raggedright
\vspace{0.2em}
\begin{tabularx}{\linewidth}{@{}>{\raggedright\arraybackslash}X@{\hspace{0.55em}}>{\raggedright\arraybackslash}p{0.14\linewidth}@{\hspace{0.25em}}}
\mbox{PNWGmc Long E Lowering} & \emph{*bārō} \\
\mbox{PNWGmc Final Long O Raising} & \emph{*bāru} \\
\mbox{EAF Long A Fronting} & \emph{*bǣru} \\
\end{tabularx}
\vspace{0.6em}
\centering\textbf{Early Anglo-Frisian}\par
\raggedright
\vspace{0.2em}
\raggedright [no change]\par
\end{minipage}
&
\begin{minipage}[t]{\linewidth}
\centering\textbf{Old English changes}\par
\vspace{0.35em}
\raggedright
\begin{tabularx}{\linewidth}{@{}>{\raggedright\arraybackslash}X@{\hspace{0.55em}}>{\raggedright\arraybackslash}p{0.14\linewidth}@{\hspace{0.25em}}}
\mbox{OE High Vowel Apocope} & \emph{*bǣr} \\
\end{tabularx}
\end{minipage}
\\
\end{tabular}
\end{minipage}%
} \endgroup

Old English form: \emph{bǣr}

\subsection{Reconstruction and comparative
evidence}\label{reconstruction-and-comparative-evidence-4}

Kroonen reconstructs the noun as
{\emph{*bērō-}}\index[iv]{03pgmc@\ivlangheader{Proto-Germanic}{}!bero@\iventry{*bērō-}{}}
f.~`bier' and cites Old English
{\emph{bar}}\index[iv]{02oe@\ivlangheader{Old English}{West Saxon normalization unmarked}!bar@\iventry{bar}{}}
`bier',
{\emph{bær}}\index[iv]{02oe@\ivlangheader{Old English}{West Saxon normalization unmarked}!baer@\iventry{bær}{}}
`bier' among the reflexes (\citeproc{ref-Kroonen2013}{Kroonen 2013,
717}). The derivational input
\Recon{bḗrō}\index[iv]{03pgmc@\ivlangheader{Proto-Germanic}{}!bero@\iventry{*bḗrō}{}}
`bier' is the same lexeme in the accent notation used here.

\subsection{Old English evidence}\label{old-english-evidence-4}

Clark Hall and Bosworth-Toller lemmatize the noun as
{\emph{bær}}\index[iv]{02oe@\ivlangheader{Old English}{West Saxon normalization unmarked}!baer@\iventry{bær}{}}
`bier', and Kroonen also records
{\emph{bar}}\index[iv]{02oe@\ivlangheader{Old English}{West Saxon normalization unmarked}!bar@\iventry{bar}{}}
`bier' beside it (\citeproc{ref-ClarkHall1960}{Clark Hall 1960};
\citeproc{ref-BosworthToller1898}{Bosworth and Toller 1898, 73};
\citeproc{ref-Kroonen2013}{Kroonen 2013, 717}). The target
{\emph{bǣr}}\index[iv]{02oe@\ivlangheader{Old English}{West Saxon normalization unmarked}!baer@\iventry{bǣr}{}}
`bier' is therefore a normalized long-vowel spelling of the same noun.

\subsection{Source note}\label{source-note}

Lexicographic spellings vary between
{\emph{bær}}\index[iv]{02oe@\ivlangheader{Old English}{West Saxon normalization unmarked}!baer@\iventry{bær}{}}
`bier' and
{\emph{bar}}\index[iv]{02oe@\ivlangheader{Old English}{West Saxon normalization unmarked}!bar@\iventry{bar}{}}
`bier'. The normalized target
{\emph{bǣr}}\index[iv]{02oe@\ivlangheader{Old English}{West Saxon normalization unmarked}!baer@\iventry{bǣr}{}}
`bier' represents the same long vowel
(\citeproc{ref-ClarkHall1960}{Clark Hall 1960};
\citeproc{ref-BosworthToller1898}{Bosworth and Toller 1898, 73};
\citeproc{ref-Kroonen2013}{Kroonen 2013, 717}).

\subsection{Development to Old
English}\label{development-to-old-english-4}

From
\Recon{bḗrō}\index[iv]{03pgmc@\ivlangheader{Proto-Germanic}{}!bero@\iventry{*bḗrō}{}}
`bier', Northwest Germanic final long \emph{ō} raises to \emph{u}, long
\emph{ē} lowers to \emph{ǣ}, and high-vowel apocope yields
{\emph{bǣr}}\index[iv]{02oe@\ivlangheader{Old English}{West Saxon normalization unmarked}!baer@\iventry{bǣr}{}}
`bier'. The resulting noun matches the normalized Old English target.

\section{\texorpdfstring{birth --- OE
\emph{byrd}}{birth --- OE byrd}}\label{birth-oe-byrd}

\index[iv]{02oe@\ivlangheader{Old English}{West Saxon normalization unmarked}!byrd@\iventry{byrd}{}}
\index[iv]{02oe@\ivlangheader{Old English}{West Saxon normalization unmarked}!gebyrd@\iventry{ge-byrd}{}}
\index[iv]{02oe@\ivlangheader{Old English}{West Saxon normalization unmarked}!gebyrd@\iventry{gebyrd}{}}
\index[iv]{02oe@\ivlangheader{Old English}{West Saxon normalization unmarked}!gebyrdu@\iventry{gebyrdu}{}}
\index[iv]{03pgmc@\ivlangheader{Proto-Germanic}{}!burdiz@\iventry{*búrdiz}{}}

Derivation: \emph{*búrdiz} \textgreater{} \emph{byrd} (regular).

\subsection{Derivation trace}\label{derivation-trace-5}

Proto input: \emph{*búrdiz}

\begingroup
\setlength{\fboxsep}{6pt}

\noindent\fbox{%
\begin{minipage}{0.97\linewidth}
\small
\begin{tabular}{@{}>{\raggedright\arraybackslash}p{0.485\linewidth}@{\hspace{0.030\linewidth}}>{\raggedright\arraybackslash}p{0.485\linewidth}@{}}
\begin{minipage}[t]{\linewidth}
\centering\textbf{Earlier Germanic changes}\par
\vspace{0.35em}
\raggedright
\centering\textbf{Northwest and West Germanic}\par
\raggedright
\vspace{0.2em}
\begin{tabularx}{\linewidth}{@{}>{\raggedright\arraybackslash}X@{\hspace{0.55em}}>{\raggedright\arraybackslash}p{0.14\linewidth}@{\hspace{0.25em}}}
\mbox{EAF Final Z Deletion} & \emph{*búrdi} \\
\end{tabularx}
\vspace{0.6em}
\centering\textbf{Early Anglo-Frisian}\par
\raggedright
\vspace{0.2em}
\raggedright [no change]\par
\end{minipage}
&
\begin{minipage}[t]{\linewidth}
\centering\textbf{Old English changes}\par
\vspace{0.35em}
\raggedright
\begin{tabularx}{\linewidth}{@{}>{\raggedright\arraybackslash}X@{\hspace{0.55em}}>{\raggedright\arraybackslash}p{0.14\linewidth}@{\hspace{0.25em}}}
\mbox{OE I Umlaut} & \emph{*byrdi} \\
\mbox{OE High Vowel Apocope} & \emph{*byrd} \\
\end{tabularx}
\end{minipage}
\\
\end{tabular}
\end{minipage}%
} \endgroup

Old English form: \emph{byrd}

\subsection{Reconstruction and comparative
evidence}\label{reconstruction-and-comparative-evidence-5}

Kroonen cites the noun under stem-level \emph{*burdi-} and gives Old
English \emph{(ge-)byrd} among the reflexes
(\citeproc{ref-Kroonen2013}{Kroonen 2013}). The derivational input
\Recon{búrdiz} `birth' is the nominative-style form that stands behind
that stem label.

\subsection{Old English evidence}\label{old-english-evidence-5}

Clark Hall and Bosworth-Toller both attest simplex
{\emph{byrd}}\index[iv]{02oe@\ivlangheader{Old English}{West Saxon normalization unmarked}!byrd@\iventry{byrd}{}}
`birth' as an Old English noun meaning `birth'
(\citeproc{ref-ClarkHall1960}{Clark Hall 1960};
\citeproc{ref-BosworthToller1898}{Bosworth and Toller 1898, 125}). The
prefixed form \emph{gebyrd} `birth' is also well established in the
tradition: Kroonen lists \emph{(ge-)byrd}, Bosworth-Toller has a
separate \emph{ge-byrd} `birth' entry, and Campbell cites \emph{gebyrd}
and \emph{gebyrdu} in his grammatical discussion
(\citeproc{ref-Kroonen2013}{Kroonen 2013};
\citeproc{ref-BosworthToller1898}{Bosworth and Toller 1898, 125};
\citeproc{ref-Campbell1959}{Campbell 1959}).

\subsection{Form note}\label{form-note}

The relevant comparator here is the simplex noun \emph{byrd}. The
prefixed forms remain related attested material within the same lexical
family, and Kastovsky's discussion of deverbal feminines provides the
broader derivational setting (\citeproc{ref-Kastovsky1992}{Kastovsky
1992, 358}).

\subsection{Development to Old
English}\label{development-to-old-english-5}

From \Recon{búrdiz} `birth', loss of final \emph{z} gives \Recon{búrdi}
`birth'. I-umlaut fronts \emph{u} to \emph{y}, and high-vowel apocope
then yields \emph{byrd}. The result is the ordinary simplex Old English
noun.

\section{\texorpdfstring{bone --- OE
\emph{bān}}{bone --- OE bān}}\label{bone-oe-bux101n}

\index[iv]{02oe@\ivlangheader{Old English}{West Saxon normalization unmarked}!ban@\iventry{bān}{}}
\index[iv]{03pgmc@\ivlangheader{Proto-Germanic}{}!baina@\iventry{*báiną}{}}

Derivation: \emph{*báiną} \textgreater{} \emph{bān} (regular).

\subsection{Derivation trace}\label{derivation-trace-6}

Proto input: \emph{*báiną}

\begingroup
\setlength{\fboxsep}{6pt}

\noindent\fbox{%
\begin{minipage}{0.97\linewidth}
\small
\begin{tabular}{@{}>{\raggedright\arraybackslash}p{0.485\linewidth}@{\hspace{0.030\linewidth}}>{\raggedright\arraybackslash}p{0.485\linewidth}@{}}
\begin{minipage}[t]{\linewidth}
\centering\textbf{Earlier Germanic changes}\par
\vspace{0.35em}
\raggedright
\centering\textbf{Northwest and West Germanic}\par
\raggedright
\vspace{0.2em}
\raggedright [no change]\par
\vspace{0.6em}
\centering\textbf{Early Anglo-Frisian}\par
\raggedright
\vspace{0.2em}
\begin{tabularx}{\linewidth}{@{}>{\raggedright\arraybackslash}X@{\hspace{0.55em}}>{\raggedright\arraybackslash}p{0.14\linewidth}@{\hspace{0.25em}}}
\mbox{EAF Ai Monophthongization} & \emph{*bāną} \\
\end{tabularx}
\end{minipage}
&
\begin{minipage}[t]{\linewidth}
\centering\textbf{Old English changes}\par
\vspace{0.35em}
\raggedright
\begin{tabularx}{\linewidth}{@{}>{\raggedright\arraybackslash}X@{\hspace{0.55em}}>{\raggedright\arraybackslash}p{0.14\linewidth}@{\hspace{0.25em}}}
OE Heavy Syllable Nasal Apocope & \emph{*bān} \\
\end{tabularx}
\end{minipage}
\\
\end{tabular}
\end{minipage}%
} \endgroup

Old English form: \emph{bān}

\subsection{Reconstruction and comparative
evidence}\label{reconstruction-and-comparative-evidence-6}

Kroonen cites the noun as
{\emph{*baina-}}\index[iv]{03pgmc@\ivlangheader{Proto-Germanic}{}!baina@\iventry{*baina-}{}},
and Orel gives the same lexeme under
\Recon{bainan}\index[iv]{03pgmc@\ivlangheader{Proto-Germanic}{}!bainan@\iventry{*bainan}{}}
`bone' (\citeproc{ref-Kroonen2013}{Kroonen 2013};
\citeproc{ref-Orel2003}{Orel 2003}). Both are comparative headword
conventions for the same neuter noun whose Old English reflex is
{\emph{bān}}\index[iv]{02oe@\ivlangheader{Old English}{West Saxon normalization unmarked}!ban@\iventry{bān}{}}
`bone'.

\subsection{Old English evidence}\label{old-english-evidence-6}

Clark Hall and Bosworth-Toller record
{\emph{bān}}\index[iv]{02oe@\ivlangheader{Old English}{West Saxon normalization unmarked}!ban@\iventry{bān}{}}
`bone' as the ordinary Old English noun
(\citeproc{ref-ClarkHall1960}{Clark Hall 1960};
\citeproc{ref-BosworthToller1898}{Bosworth and Toller 1898}). Bright's
glossary also distinguishes citation-form
{\emph{bān}}\index[iv]{02oe@\ivlangheader{Old English}{West Saxon normalization unmarked}!ban@\iventry{bān}{}}
`bone' from oblique
{\emph{bāne}}\index[iv]{02oe@\ivlangheader{Old English}{West Saxon normalization unmarked}!bane@\iventry{bāne}{}}
`bone' (\citeproc{ref-BrightCassidyRingler1971}{Bright 1971}).

\subsection{Source note}\label{source-note-1}

The comparative headwords
{\emph{*baina-}}\index[iv]{03pgmc@\ivlangheader{Proto-Germanic}{}!baina@\iventry{*baina-}{}}
and
\Recon{bainan}\index[iv]{03pgmc@\ivlangheader{Proto-Germanic}{}!bainan@\iventry{*bainan}{}}
`bone' provide lexeme background. The relevant comparison form here is
the nominative-accusative singular
\Recon{báiną}\index[iv]{03pgmc@\ivlangheader{Proto-Germanic}{}!baina@\iventry{*báiną}{}}
`bone'.

\subsection{Development to Old
English}\label{development-to-old-english-6}

West Germanic monophthongization turns stressed \emph{*ai} into
\emph{ā}, giving \Recon{bāną} `bone'; heavy-syllable nasal apocope then
yields
{\emph{bān}}\index[iv]{02oe@\ivlangheader{Old English}{West Saxon normalization unmarked}!ban@\iventry{bān}{}}
`bone'. The resulting form matches the attested Old English citation
noun.

\section{\texorpdfstring{book --- OE
\emph{bōc}}{book --- OE bōc}}\label{book-oe-bux14dc}

\index[iv]{02oe@\ivlangheader{Old English}{West Saxon normalization unmarked}!boc@\iventry{bōc}{}}
\index[iv]{03pgmc@\ivlangheader{Proto-Germanic}{}!bokz@\iventry{*bōkz}{}}

Derivation: \emph{*bōkz} \textgreater{} \emph{bōc} (regular).

\subsection{Derivation trace}\label{derivation-trace-7}

Proto input: \emph{*bōkz}

\begingroup
\setlength{\fboxsep}{6pt}

\noindent\fbox{%
\begin{minipage}{0.97\linewidth}
\small
\begin{tabular}{@{}>{\raggedright\arraybackslash}p{0.485\linewidth}@{\hspace{0.030\linewidth}}>{\raggedright\arraybackslash}p{0.485\linewidth}@{}}
\begin{minipage}[t]{\linewidth}
\centering\textbf{Earlier Germanic changes}\par
\vspace{0.35em}
\raggedright
\centering\textbf{Northwest and West Germanic}\par
\raggedright
\vspace{0.2em}
\raggedright [no change]\par
\vspace{0.6em}
\centering\textbf{Early Anglo-Frisian}\par
\raggedright
\vspace{0.2em}
\raggedright [no change]\par
\end{minipage}
&
\begin{minipage}[t]{\linewidth}
\centering\textbf{Old English changes}\par
\vspace{0.35em}
\raggedright
\raggedright [no change]\par
\end{minipage}
\\
\end{tabular}
\end{minipage}%
} \endgroup

Old English form: \emph{bōc}

\subsection{Reconstruction and comparative
evidence}\label{reconstruction-and-comparative-evidence-7}

The noun is an athematic consonant stem (root noun), and the sources
print its nominative singular in three different notations that encode
the same morphology. Orel's citation form is \emph{*bōkz} `book', with
the explicit nominative marker \emph{*-z} that he writes across the
whole root-noun class (\citeproc{ref-Orel2003}{Orel 2003, 52}). Kroonen
cites the bare stem
{\emph{*bōk-}}\index[iv]{03pgmc@\ivlangheader{Proto-Germanic}{}!bok@\iventry{*bōk-}{}}
`book, beech' and treats Gothic \emph{bōka} `letter' as a separate
ō-stem formation (\citeproc{ref-Kroonen2013}{Kroonen 2013, 71--72}).
Kluge/Seebold print a third variant with voiceless sibilant,
{\emph{*bōks}}\index[iv]{03pgmc@\ivlangheader{Proto-Germanic}{}!boks@\iventry{*bōks}{}}
`book' (\citeproc{ref-KlugeSeebold2011}{Kluge 2011, 158}). These are
citation conventions, not competing facts: Ringe gives the
consonant-stem nominative ending as zero, \emph{*-z}, or possibly
\emph{*-s}, and states that the distribution is unrecoverable for
monosyllabic stems (\citeproc{ref-Ringe2017}{Ringe 2017, 306}). No
scholar is simply ``wrong'' here; the notations differ in whether they
make the inflectional marker visible.

The derivational input
\Recon{bōkz}\index[iv]{03pgmc@\ivlangheader{Proto-Germanic}{}!bokz@\iventry{*bōkz}{}}
`book' follows Orel's morphologically explicit convention: it records
the nominative marker whose historical elimination the derivation then
models. The word is distinct from the ō-stem `beech' word, whose
derivational input here is the separate nominative \Recon{bōkō} `beech'.

\subsection{Old English evidence}\label{old-english-evidence-7}

Old English
{\emph{bōc}}\index[iv]{02oe@\ivlangheader{Old English}{West Saxon normalization unmarked}!boc@\iventry{bōc}{}}
`book' is the endingless nominative-accusative singular of a root noun,
with the mutated plural \emph{bēċ} `books' preserving the class's
characteristic i-umlaut alternation (\citeproc{ref-Fulk2018}{Fulk 2018,
165--66}). Comparative root-noun evidence shows the branch-by-branch
fate of the ending: Gothic keeps its sibilant in this class
(\emph{baúrgs} `city', \emph{nahts} `night'), while Old Norse \emph{bók}
`book' keeps a reflex of \emph{*-z} as \emph{-r} only in the masculine
root nouns and drops it in feminines (\citeproc{ref-Fulk2018}{Fulk 2018,
165--67}). West Germanic alone is uniform: none of its daughters shows
any ending in the root-noun nominative
(\citeproc{ref-RingeTaylor2014}{Ringe and Taylor 2014, 118}).

\subsection{Development to Old
English}\label{development-to-old-english-7}

The nominative marker of the consonant stems was eliminated before
Proto-West Germanic: endingless nominatives arose by Szemerényi's law
after sonorant-final stems and spread analogically through the class
(\citeproc{ref-Fulk2018}{Fulk 2018, 143}), so that no West Germanic
daughter inherits an ending here (\citeproc{ref-RingeTaylor2014}{Ringe
and Taylor 2014, 118}). From the derivational input \Recon{bōkz} `book'
the loss of the marker gives \Recon{bōk} `book', which is already the
Old English form: \emph{bōc} `book' differs only in spelling. This early
morphological loss is a different development from the later West
Germanic loss of final \emph{*z} in unstressed syllables (as in
\Recon{bárdaz} `beard') and from the still later northern loss in
stressed monosyllables (as in the pronouns); the root-noun ending was
gone before either of those changes applied.

\section{\texorpdfstring{both --- OE
\emph{bū}}{both --- OE bū}}\label{both-oe-bux16b}

\index[iv]{02oe@\ivlangheader{Old English}{West Saxon normalization unmarked}!bu@\iventry{bū}{}}
\index[iv]{02oe@\ivlangheader{Old English}{West Saxon normalization unmarked}!cu@\iventry{cū}{}}
\index[iv]{02oe@\ivlangheader{Old English}{West Saxon normalization unmarked}!hu@\iventry{hū}{}}
\index[iv]{02oe@\ivlangheader{Old English}{West Saxon normalization unmarked}!tu@\iventry{tū}{}}
\index[iv]{03pgmc@\ivlangheader{Proto-Germanic}{}!bo@\iventry{*bō}{}}

Derivation: \emph{*bō} \textgreater{} \emph{bū} (regular).

\subsection{Derivation trace}\label{derivation-trace-8}

Proto input: \emph{*bō}

\begingroup
\setlength{\fboxsep}{6pt}

\noindent\fbox{%
\begin{minipage}{0.97\linewidth}
\small
\begin{tabular}{@{}>{\raggedright\arraybackslash}p{0.611\linewidth}@{\hspace{0.030\linewidth}}>{\raggedright\arraybackslash}p{0.359\linewidth}@{}}
\begin{minipage}[t]{\linewidth}
\centering\textbf{Earlier Germanic changes}\par
\vspace{0.35em}
\raggedright
\centering\textbf{Northwest and West Germanic}\par
\raggedright
\vspace{0.2em}
\begin{tabularx}{\linewidth}{@{}>{\raggedright\arraybackslash}X@{\hspace{0.55em}}>{\raggedright\arraybackslash}p{0.14\linewidth}@{\hspace{0.25em}}}
PNWGmc Stressed Monosyllable O Raising & \emph{*bū} \\
\end{tabularx}
\vspace{0.6em}
\centering\textbf{Early Anglo-Frisian}\par
\raggedright
\vspace{0.2em}
\raggedright [no change]\par
\end{minipage}
&
\begin{minipage}[t]{\linewidth}
\centering\textbf{Old English changes}\par
\vspace{0.35em}
\raggedright
\raggedright [no change]\par
\end{minipage}
\\
\end{tabular}
\end{minipage}%
} \endgroup

Old English form: \emph{bū}

\subsection{Reconstruction and comparative
evidence}\label{reconstruction-and-comparative-evidence-8}

Kroonen treats the Germanic numeral under \emph{*ba-} and gives the
inherited paradigm \Recon{bai} `both', \Recon{bans} `both (pl.)',
\Recon{bōz} `both'/ \Recon{bōns} `both', \Recon{bō} `both', with Old
English
{\emph{bēġen}}\index[iv]{02oe@\ivlangheader{Old English}{West Saxon normalization unmarked}!begen@\iventry{bēġen}{}}
`both',
{\emph{bā}}\index[iv]{02oe@\ivlangheader{Old English}{West Saxon normalization unmarked}!ba@\iventry{bā}{}}
`both', and neuter
{\emph{bū}}\index[iv]{02oe@\ivlangheader{Old English}{West Saxon normalization unmarked}!bu@\iventry{bū}{}}
`both' (\citeproc{ref-Kroonen2013}{Kroonen 2013, 47}). For the present
entry, the relevant inherited form is the unextended neuter dual
{\emph{*bō}}\index[iv]{03pgmc@\ivlangheader{Proto-Germanic}{}!bo@\iventry{*bō}{}}
`both'.

The older explanation of
{\emph{bēġen}}\index[iv]{02oe@\ivlangheader{Old English}{West Saxon normalization unmarked}!begen@\iventry{bēġen}{}}
`both' derives it from \emph{*bō-jen-}, and Orel still gives OE
{\emph{bezen}}\index[iv]{02oe@\ivlangheader{Old English}{West Saxon normalization unmarked}!bezen@\iventry{bezen}{}}
`both' (\textless{} \Recon{bō-jenō} `both') beside ON
{\emph{báðir}}\index[iv]{09on@\ivlangheader{Old Norse}{}!badir@\iventry{báðir}{}}
`both', OFris
{\emph{bēthe}}\index[iv]{11ofris@\ivlangheader{Old Frisian}{}!bethe@\iventry{bēthe}{}}
`both', OS
{\emph{be-thia}}\index[iv]{13os@\ivlangheader{Old Saxon}{}!bethia@\iventry{be-thia}{}}
`both', and OHG
{\emph{bēde}}\index[iv]{10ohg@\ivlangheader{Old High German}{}!bede@\iventry{bēde}{}}
`both' (\citeproc{ref-Orel2003}{Orel 2003, 65}). Fulk reports that
explanation cautiously and notes Seebold's preference for a
\emph{*bō-þ-} analysis instead (\citeproc{ref-Fulk2018}{Fulk 2018, § 10.1}). That debate concerns
{\emph{bēġen}}\index[iv]{02oe@\ivlangheader{Old English}{West Saxon normalization unmarked}!begen@\iventry{bēġen}{}}
and the extended forms behind Modern English
{\emph{both}}\index[iv]{22modeng@\ivlangheader{Modern English}{}!both@\iventry{both}{}}
`both', German
{\emph{beide}}\index[iv]{17german@\ivlangheader{German}{}!beide@\iventry{beide}{}}
`both', and Dutch
{\emph{beide}}\index[iv]{16dutch@\ivlangheader{Dutch}{}!beide@\iventry{beide}{}}
`both'; it does not displace the inherited neuter
{\emph{*bō}}\index[iv]{03pgmc@\ivlangheader{Proto-Germanic}{}!bo@\iventry{*bō}{}}
`both' \textgreater{}
{\emph{bū}}\index[iv]{02oe@\ivlangheader{Old English}{West Saxon normalization unmarked}!bu@\iventry{bū}{}}
`both' treated here.

\subsection{Old English evidence}\label{old-english-evidence-8}

The Old English dual paradigm is well established. Brunner gives
masculine
{\emph{bēġen}}\index[iv]{02oe@\ivlangheader{Old English}{West Saxon normalization unmarked}!begen@\iventry{bēġen}{}}
`both', feminine
{\emph{bā}}\index[iv]{02oe@\ivlangheader{Old English}{West Saxon normalization unmarked}!ba@\iventry{bā}{}}
`both', and neuter
{\emph{bū}}\index[iv]{02oe@\ivlangheader{Old English}{West Saxon normalization unmarked}!bu@\iventry{bū}{}}
`both' beside \emph{bā}, with compounds such as \emph{bā} \emph{twā}
`both (two)', \emph{bū} \emph{tū} `both (two, neut.)', and \emph{bām}
`both' \emph{twām} `two' (\citeproc{ref-SieversBrunner1965}{Brunner
1965, § 324} Anm. 2). Campbell and Fulk present the same basic
pattern: masculine
{\emph{bēġen}}\index[iv]{02oe@\ivlangheader{Old English}{West Saxon normalization unmarked}!begen@\iventry{bēġen}{}}
`both', feminine \emph{bā}, neuter \emph{bā}, \emph{bū}, genitive
\emph{bēġra} `both', \emph{bēġ(e)a}, and dative
{\emph{bǣm}}\index[iv]{02oe@\ivlangheader{Old English}{West Saxon normalization unmarked}!baem@\iventry{bǣm}{}}
`both' (\citeproc{ref-Campbell1959}{Campbell 1959, § 683};
\citeproc{ref-Fulk2018}{Fulk 2018, § 10.1}).

{\emph{bū}}\index[iv]{02oe@\ivlangheader{Old English}{West Saxon normalization unmarked}!bu@\iventry{bū}{}}
`both' is therefore an attested neuter dual form, not a reconstruction.
It is the cleanest target for this entry because
{\emph{bēġen}}\index[iv]{02oe@\ivlangheader{Old English}{West Saxon normalization unmarked}!begen@\iventry{bēġen}{}}
`both' belongs to the historically more contested \emph{*bō-jen-} /
analogical zone, while \emph{bā} `both' remains a partner form within
the dual paradigm rather than the most straightforward monosyllabic
comparison.

\subsection{Development to Old
English}\label{development-to-old-english-8}

{\emph{*bō}}\index[iv]{03pgmc@\ivlangheader{Proto-Germanic}{}!bo@\iventry{*bō}{}}
`both' is a stressed monosyllabic form. Campbell cites \emph{cū} `cow',
\emph{hū} `how', \emph{tū} `two', and
{\emph{bū}}\index[iv]{02oe@\ivlangheader{Old English}{West Saxon normalization unmarked}!bu@\iventry{bū}{}}
`both' as examples of final accented \emph{ō} \textgreater{} \emph{ū} in
the West Germanic stage leading to Old English
(\citeproc{ref-Campbell1959}{Campbell 1959, § 122}). Brunner states
the same development more directly: Auslautendes \emph{ō} erscheint als
û in \emph{bū} \ldots{} cu \ldots{} \emph{hū}, \emph{tū}
(\citeproc{ref-SieversBrunner1965}{Brunner 1965, § 69}).

The development is therefore straightforward:
{\emph{*bō}}\index[iv]{03pgmc@\ivlangheader{Proto-Germanic}{}!bo@\iventry{*bō}{}}
`both' \textgreater{}
{\emph{bū}}\index[iv]{02oe@\ivlangheader{Old English}{West Saxon normalization unmarked}!bu@\iventry{bū}{}}
`both'.

\subsection{Form comparison}\label{form-comparison}

The comparison below sets the relevant forms side by side. It separates
the inherited OE target from the other forms that belong to the same
broader lexical history.

{\def\LTcaptype{none} % do not increment counter
\begin{longtable}[]{@{}
  >{\raggedright\arraybackslash}p{(\linewidth - 6\tabcolsep) * \real{0.2500}}
  >{\raggedright\arraybackslash}p{(\linewidth - 6\tabcolsep) * \real{0.2500}}
  >{\raggedright\arraybackslash}p{(\linewidth - 6\tabcolsep) * \real{0.2500}}
  >{\raggedright\arraybackslash}p{(\linewidth - 6\tabcolsep) * \real{0.2500}}@{}}
\toprule\noalign{}
\begin{minipage}[b]{\linewidth}\raggedright
Form
\end{minipage} & \begin{minipage}[b]{\linewidth}\raggedright
Source / stage
\end{minipage} & \begin{minipage}[b]{\linewidth}\raggedright
Status
\end{minipage} & \begin{minipage}[b]{\linewidth}\raggedright
Relevance to this entry
\end{minipage} \\
\midrule\noalign{}
\endhead
\bottomrule\noalign{}
\endlastfoot
{\emph{*bō}}\index[iv]{03pgmc@\ivlangheader{Proto-Germanic}{}!bo@\iventry{*bō}{}}
\textgreater{}
{\emph{bū}}\index[iv]{02oe@\ivlangheader{Old English}{West Saxon normalization unmarked}!bu@\iventry{bū}{}}
& PGmc neuter dual \textgreater{} OE neuter dual & selected regular
comparison & main line of the entry \\
{\emph{bēġen}}\index[iv]{02oe@\ivlangheader{Old English}{West Saxon normalization unmarked}!begen@\iventry{bēġen}{}}
& OE masculine dual & attested, but historically contested and at least
partly analogical in Kroonen & real OE evidence, not the Old English
form here \\
{\emph{bā}}\index[iv]{02oe@\ivlangheader{Old English}{West Saxon normalization unmarked}!ba@\iventry{bā}{}}
& OE feminine dual; also neuter variant & attested partner form & part
of the OE paradigm, but not the chosen monosyllabic comparator \\
{\emph{báðir}}\index[iv]{09on@\ivlangheader{Old Norse}{}!badir@\iventry{báðir}{}},
German
{\emph{beide}}\index[iv]{17german@\ivlangheader{German}{}!beide@\iventry{beide}{}},
Dutch
{\emph{beide}}\index[iv]{16dutch@\ivlangheader{Dutch}{}!beide@\iventry{beide}{}},
Modern English
{\emph{both}}\index[iv]{22modeng@\ivlangheader{Modern English}{}!both@\iventry{both}{}}
& Norse, continental West Germanic, Modern English extended forms &
related but different formation & useful background, not the direct
continuation of OE \emph{bū} \\
\end{longtable}
}

\section{\texorpdfstring{bow --- OE
\emph{bīeġan}}{bow --- OE bīeġan}}\label{bow-oe-bux12beux121an}

\index[iv]{02oe@\ivlangheader{Old English}{West Saxon normalization unmarked}!biegan@\iventry{bīeġan}{}}
\index[iv]{02oe@\ivlangheader{Old English}{West Saxon normalization unmarked}!bigan@\iventry{bigan}{}}
\index[iv]{03pgmc@\ivlangheader{Proto-Germanic}{}!baugijana@\iventry{*báugijaną}{}}

Derivation: \emph{*báugijaną} \textgreater{} \emph{bīeġan} (regular).

\subsection{Derivation trace}\label{derivation-trace-9}

Proto input: \emph{*báugijaną}

\begingroup
\setlength{\fboxsep}{6pt}

\noindent\fbox{%
\begin{minipage}{0.97\linewidth}
\footnotesize
\begin{tabular}{@{}>{\raggedright\arraybackslash}p{0.339\linewidth}@{\hspace{0.030\linewidth}}>{\raggedright\arraybackslash}p{0.630\linewidth}@{}}
\begin{minipage}[t]{\linewidth}
\centering\textbf{Earlier Germanic changes}\par
\vspace{0.35em}
\raggedright
\centering\textbf{Northwest and West Germanic}\par
\raggedright
\vspace{0.2em}
\raggedright [no change]\par
\vspace{0.6em}
\centering\textbf{Early Anglo-Frisian}\par
\raggedright
\vspace{0.2em}
\raggedright [no change]\par
\end{minipage}
&
\begin{minipage}[t]{\linewidth}
\centering\textbf{Old English changes}\par
\vspace{0.35em}
\raggedright
\begin{tabularx}{\linewidth}{@{}>{\raggedright\arraybackslash}X@{\hspace{0.55em}}>{\raggedright\arraybackslash}p{0.26\linewidth}@{\hspace{0.25em}}}
\mbox{OE Au Brightening} & \emph{*báeugijaną} \\
\mbox{OE Diphthong Leveling} & \emph{*bēagijaną} \\
OE Heavy Syllable Nasal Apocope & \emph{*bēagijan} \\
\mbox{OE Secondary Nasalization} & \emph{*bēagijąn} \\
\mbox{Sievers Law Syncope} & \emph{*bēagjąn} \\
\mbox{OE Velar Palatalization} & \emph{*bēaʝjąn} \\
\mbox{OE I Umlaut} & \emph{*bīeʝjąn} \\
\mbox{OE Weak Tail Reduction} & \emph{*bīeʝjan} \\
\mbox{OE J Loss After Heavy} & \emph{*bīeʝan} \\
\mbox{OE Palatal Fricative Merger} & \emph{*bīejan} \\
\end{tabularx}
\end{minipage}
\\
\end{tabular}
\end{minipage}%
} \endgroup

Old English form: \emph{bīeġan}

\subsection{Reconstruction and comparative
evidence}\label{reconstruction-and-comparative-evidence-9}

Kroonen reconstructs the weak verb as \emph{*baugjan-} `to (make) bend'
and cites Old English
{\emph{biegan}}\index[iv]{02oe@\ivlangheader{Old English}{West Saxon normalization unmarked}!biegan@\iventry{biegan}{}}
`bow, bend' among its reflexes (\citeproc{ref-Kroonen2013}{Kroonen
2013}). Ringe and Taylor give the northwest Germanic preform
\Recon{baugijana} `bow', with later \Recon{béagjan} `bow' behind West
Saxon Old English
{\emph{biegan}}\index[iv]{02oe@\ivlangheader{Old English}{West Saxon normalization unmarked}!biegan@\iventry{biegan}{}}
`bow, bend' (\citeproc{ref-RingeTaylor2014}{Ringe and Taylor 2014}). The
entry therefore concerns the weak causative member of the bend-family,
alongside the related strong verb and noun.

\subsection{Old English evidence}\label{old-english-evidence-9}

Clark Hall lemmatizes \emph{biegan} `bow, bend' as an Old English weak
verb, and Bosworth-Toller records \emph{bigan} `bow, bend' as an
attested form (\citeproc{ref-ClarkHall1960}{Clark Hall 1960};
\citeproc{ref-BosworthToller1898}{Bosworth and Toller 1898, 102}). The
form \emph{bīeġan} `bow, bend' used here is a normalized spelling of
that attested Old English weak verb.

\subsection{Development to Old
English}\label{development-to-old-english-9}

From \Recon{báugijaną} `bow', the stem reaches pre-Old-English
\Recon{bēagjan} `bow', after which palatalization of \emph{*gj} and
i-umlaut yield West Saxon \emph{biegan} `bow'; Campbell lists
\emph{biegan} `bow' among the regular \emph{ie} outcomes of \emph{*éa}
under i-umlaut (\citeproc{ref-RingeTaylor2014}{Ringe and Taylor 2014};
\citeproc{ref-Campbell1959}{Campbell 1959, 80}). The development is
therefore straightforward: \emph{*báugijaną} \textgreater{}
\emph{bīeġan} `bow'.

\section{\texorpdfstring{breeches --- OE
\emph{brēċ}}{breeches --- OE brēċ}}\label{breeches-oe-brux113ux10b}

\index[iv]{02oe@\ivlangheader{Old English}{West Saxon normalization unmarked}!brec@\iventry{bréc}{}}
\index[iv]{02oe@\ivlangheader{Old English}{West Saxon normalization unmarked}!brec@\iventry{brēc}{}}
\index[iv]{02oe@\ivlangheader{Old English}{West Saxon normalization unmarked}!brec@\iventry{brēċ}{}}
\index[iv]{02oe@\ivlangheader{Old English}{West Saxon normalization unmarked}!broc@\iventry{broc}{}}
\index[iv]{02oe@\ivlangheader{Old English}{West Saxon normalization unmarked}!broc@\iventry{brōc}{}}
\index[iv]{03pgmc@\ivlangheader{Proto-Germanic}{}!brokiz@\iventry{*brōkiz}{}}

Derivation: \emph{*brōkiz} \textgreater{} \emph{brēċ} (regular).

\subsection{Derivation trace}\label{derivation-trace-10}

Proto input: \emph{*brōkiz}

\begingroup
\setlength{\fboxsep}{6pt}

\noindent\fbox{%
\begin{minipage}{0.97\linewidth}
\small
\begin{tabular}{@{}>{\raggedright\arraybackslash}p{0.388\linewidth}@{\hspace{0.030\linewidth}}>{\raggedright\arraybackslash}p{0.582\linewidth}@{}}
\begin{minipage}[t]{\linewidth}
\centering\textbf{Earlier Germanic changes}\par
\vspace{0.35em}
\raggedright
\centering\textbf{Northwest and West Germanic}\par
\raggedright
\vspace{0.2em}
\begin{tabularx}{\linewidth}{@{}>{\raggedright\arraybackslash}X@{\hspace{0.55em}}>{\raggedright\arraybackslash}p{0.14\linewidth}@{\hspace{0.25em}}}
\mbox{EAF Final Z Deletion} & \emph{*brōki} \\
\end{tabularx}
\vspace{0.6em}
\centering\textbf{Early Anglo-Frisian}\par
\raggedright
\vspace{0.2em}
\raggedright [no change]\par
\end{minipage}
&
\begin{minipage}[t]{\linewidth}
\centering\textbf{Old English changes}\par
\vspace{0.35em}
\raggedright
\begin{tabularx}{\linewidth}{@{}>{\raggedright\arraybackslash}X@{\hspace{0.55em}}>{\raggedright\arraybackslash}p{0.14\linewidth}@{\hspace{0.25em}}}
\mbox{OE Velar Palatalization} & \emph{*brōʧi} \\
\mbox{OE I Umlaut} & \emph{*brēʧi} \\
\mbox{OE High Vowel Apocope} & \emph{*brēʧ} \\
\end{tabularx}
\end{minipage}
\\
\end{tabular}
\end{minipage}%
} \endgroup

Old English form: \emph{brēċ}

\subsection{Reconstruction and comparative
evidence}\label{reconstruction-and-comparative-evidence-10}

Kroonen cites the noun under \emph{*brōk-}, with Old English \emph{brōc}
`breeches' and plural `breeches' among its reflexes
(\citeproc{ref-Kroonen2013}{Kroonen 2013}). Ringe and Taylor give the
plural development directly as PNWGmc
\Recon{brōkiz}\index[iv]{04pnwgmc@\ivlangheader{Proto-Northwest Germanic}{}!brokiz@\iventry{*brōkiz}{}}
`breeches' \textgreater{} \Recon{breeci} `breeches' \textgreater{} OE
\emph{bréc} `breeches' (\citeproc{ref-RingeTaylor2014}{Ringe and Taylor
2014}). The deeper verbal base belongs to the noun's etymological
background, while the derivational input here is the plural noun form
\Recon{brōkiz}\index[iv]{04pnwgmc@\ivlangheader{Proto-Northwest Germanic}{}!brokiz@\iventry{*brōkiz}{}}
`breeches'.

\subsection{Old English evidence}\label{old-english-evidence-10}

Bright notes \emph{brōc} `breeches' with plural \emph{brēc} `breeches',
and Clark Hall gives \emph{brēc} fp. breeches while also listing
\emph{broc} `breeches' as a feminine noun probably represented chiefly
in the plural (\citeproc{ref-BrightCassidyRingler1971}{Bright 1971};
\citeproc{ref-ClarkHall1960}{Clark Hall 1960, 64}). I write \emph{brēċ}
`breeches' for the long vowel and palatal consonant; the attested plural
is \emph{brēc}.

\subsection{Development to Old
English}\label{development-to-old-english-10}

After loss of final \emph{-z}, the stem ends in \emph{-ki}, so the velar
palatalizes and \emph{ō} `breeches' undergoes i-umlaut to \emph{ē}
`breeches'; final high-vowel apocope then yields \emph{brēċ} `breeches'
(\citeproc{ref-RingeTaylor2014}{Ringe and Taylor 2014}). The development
is therefore regular: \Recon{brōkiz} `breeches' \textgreater{}
\emph{brēċ} `breeches'.

\section{\texorpdfstring{calf --- OE
\emph{ċealf}}{calf --- OE ċealf}}\label{calf-oe-ux10bealf}

\index[iv]{02oe@\ivlangheader{Old English}{West Saxon normalization unmarked}!calfur@\iventry{calfur}{}}
\index[iv]{02oe@\ivlangheader{Old English}{West Saxon normalization unmarked}!cealf@\iventry{ċealf}{}}
\index[iv]{03pgmc@\ivlangheader{Proto-Germanic}{}!kalbaz@\iventry{*kálbaz}{}}

Derivation: \emph{*kálbaz} \textgreater{} \emph{ċealf} (regular).

\subsection{Derivation trace}\label{derivation-trace-11}

Proto input: \emph{*kálbaz}

\begingroup
\setlength{\fboxsep}{6pt}

\noindent\fbox{%
\begin{minipage}{0.97\linewidth}
\small
\begin{tabular}{@{}>{\raggedright\arraybackslash}p{0.388\linewidth}@{\hspace{0.030\linewidth}}>{\raggedright\arraybackslash}p{0.582\linewidth}@{}}
\begin{minipage}[t]{\linewidth}
\centering\textbf{Earlier Germanic changes}\par
\vspace{0.35em}
\raggedright
\centering\textbf{Northwest and West Germanic}\par
\raggedright
\vspace{0.2em}
\begin{tabularx}{\linewidth}{@{}>{\raggedright\arraybackslash}X@{\hspace{0.55em}}>{\raggedright\arraybackslash}p{0.14\linewidth}@{\hspace{0.25em}}}
\mbox{EAF Final Z Deletion} & \emph{*kálba} \\
\end{tabularx}
\vspace{0.6em}
\centering\textbf{Early Anglo-Frisian}\par
\raggedright
\vspace{0.2em}
\raggedright [no change]\par
\end{minipage}
&
\begin{minipage}[t]{\linewidth}
\centering\textbf{Old English changes}\par
\vspace{0.35em}
\raggedright
\begin{tabularx}{\linewidth}{@{}>{\raggedright\arraybackslash}X@{\hspace{0.55em}}>{\raggedright\arraybackslash}p{0.14\linewidth}@{\hspace{0.25em}}}
\mbox{PWGmc Final Bare A Loss} & \emph{*kálb} \\
\mbox{EAF Brightening} & \emph{*kælb} \\
\mbox{OE Breaking} & \emph{*kealb} \\
\mbox{OE Velar Palatalization} & \emph{*ʧealb} \\
\mbox{PGmc B Allophony} & \emph{*ʧealβ} \\
\end{tabularx}
\end{minipage}
\\
\end{tabular}
\end{minipage}%
} \endgroup

Old English form: \emph{ċealf}

\subsection{Reconstruction and comparative
evidence}\label{reconstruction-and-comparative-evidence-11}

Kroonen treats the noun under \emph{*kalbiz-} and notes an older s-stem
\Recon{kalbaz} `calf' with plural \Recon{kalbizō} `calves', while Orel
cites \Recon{kalbaz} `calf' as the citation form and Ringe and Taylor
derive West Saxon \emph{Cealf} from \emph{*kalbaz}, \emph{*kalbiz-}
(\citeproc{ref-Kroonen2013}{Kroonen 2013}; \citeproc{ref-Orel2003}{Orel
2003, 248}; \citeproc{ref-RingeTaylor2014}{Ringe and Taylor 2014, 220}).
The derivational input here is the singular \Recon{kálbaz} `calf', since
the entry concerns the citation-form noun.

\subsection{Old English evidence}\label{old-english-evidence-11}

Clark Hall gives
{\emph{cealf}}\index[iv]{02oe@\ivlangheader{Old English}{West Saxon normalization unmarked}!cealf@\iventry{cealf}{}}
`calf' I. (æ, e) nm. (nap.
{\emph{cealfru}}\index[iv]{02oe@\ivlangheader{Old English}{West Saxon normalization unmarked}!cealfru@\iventry{cealfru}{}}
`calves'), and Bosworth-Toller likewise records \emph{Caelf} /
\emph{Cealf} `calf' beside plural forms such as \emph{calfur} and
{\emph{cealfru}}\index[iv]{02oe@\ivlangheader{Old English}{West Saxon normalization unmarked}!cealfru@\iventry{cealfru}{}}
`calves' (\citeproc{ref-ClarkHall1960}{Clark Hall 1960};
\citeproc{ref-BosworthToller1898}{Bosworth and Toller 1898, 131}).
Campbell and Fulk show the same singular-plus-\emph{-r-} plural pattern
(\citeproc{ref-Campbell1959}{Campbell 1959};
\citeproc{ref-Fulk2018}{Fulk 2018, 193}). I write \emph{ċealf} `calf'
for the palatalized initial; the attested dictionary headword is
{\emph{cealf}}\index[iv]{02oe@\ivlangheader{Old English}{West Saxon normalization unmarked}!cealf@\iventry{cealf}{}}
`calf'.

\subsection{Development to Old
English}\label{development-to-old-english-11}

After loss of final \emph{-z} and bare \emph{-a}, Anglo-Frisian
brightening gives \Recon{kælb} `calf', and breaking before \emph{l} plus
consonant yields \Recon{kealb} `calf'. Ringe and Taylor's account of the
lexeme and their rule for initial \emph{k} in front-vocalic environments
support the West Saxon palatalized onset represented here as \emph{ċ-},
so \Recon{kálbaz} `calf' develops regularly to \emph{ċealf} `calf'
(\citeproc{ref-RingeTaylor2014}{Ringe and Taylor 2014, 220}).

\section{\texorpdfstring{coat --- OE
\emph{rocc}}{coat --- OE rocc}}\label{coat-oe-rocc}

\index[iv]{02oe@\ivlangheader{Old English}{West Saxon normalization unmarked}!rocc@\iventry{rocc}{}}
\index[iv]{02oe@\ivlangheader{Old English}{West Saxon normalization unmarked}!stanrocc@\iventry{stānrocc}{}}
\index[iv]{03pgmc@\ivlangheader{Proto-Germanic}{}!rukkaz@\iventry{*rúkkaz}{}}

Derivation: \emph{*rúkkaz} \textgreater{} \emph{rocc} (regular).

\subsection{Derivation trace}\label{derivation-trace-12}

Proto input: \emph{*rúkkaz}

\begingroup
\setlength{\fboxsep}{6pt}

\noindent\fbox{%
\begin{minipage}{0.97\linewidth}
\small
\begin{tabular}{@{}>{\raggedright\arraybackslash}p{0.485\linewidth}@{\hspace{0.030\linewidth}}>{\raggedright\arraybackslash}p{0.485\linewidth}@{}}
\begin{minipage}[t]{\linewidth}
\centering\textbf{Earlier Germanic changes}\par
\vspace{0.35em}
\raggedright
\centering\textbf{Northwest and West Germanic}\par
\raggedright
\vspace{0.2em}
\begin{tabularx}{\linewidth}{@{}>{\raggedright\arraybackslash}X@{\hspace{0.55em}}>{\raggedright\arraybackslash}p{0.14\linewidth}@{\hspace{0.25em}}}
\mbox{PNWGmc U Lowering} & \emph{*rókkaz} \\
\mbox{EAF Final Z Deletion} & \emph{*rókka} \\
\end{tabularx}
\vspace{0.6em}
\centering\textbf{Early Anglo-Frisian}\par
\raggedright
\vspace{0.2em}
\raggedright [no change]\par
\end{minipage}
&
\begin{minipage}[t]{\linewidth}
\centering\textbf{Old English changes}\par
\vspace{0.35em}
\raggedright
\begin{tabularx}{\linewidth}{@{}>{\raggedright\arraybackslash}X@{\hspace{0.55em}}>{\raggedright\arraybackslash}p{0.14\linewidth}@{\hspace{0.25em}}}
\mbox{PWGmc Final Bare A Loss} & \emph{*rókk} \\
\end{tabularx}
\end{minipage}
\\
\end{tabular}
\end{minipage}%
} \endgroup

Old English form: \emph{rocc}

\subsection{Reconstruction and comparative
evidence}\label{reconstruction-and-comparative-evidence-12}

Orel cites a masculine \Recon{rukkaz} `coat' for the garment word, while
Kroonen gives \emph{*hrukka-}. Both treat this as the garment lexeme and
not as the separate stone word (\citeproc{ref-Orel2003}{Orel 2003};
\citeproc{ref-Kroonen2013}{Kroonen 2013, 290}).

\subsection{Old English evidence}\label{old-english-evidence-12}

Clark Hall and Bosworth-Toller record \emph{rocc} `coat' as an
over-garment or tunic and preserve compounds such as \emph{bisceoprocc}
and \emph{breóstrocc} `breast-coat' (\citeproc{ref-ClarkHall1960}{Clark
Hall 1960}; \citeproc{ref-BosworthToller1898}{Bosworth and Toller
1898}).

\subsection{Development to Old
English}\label{development-to-old-english-12}

With \Recon{rúkkaz} `coat' as the derivational input, Northwest Germanic
u-lowering and later loss of final \emph{-a} yield \emph{rocc} as a
regular outcome.

\subsection{Source note}\label{source-note-2}

This entry concerns the garment noun only. The stone word seen in
\emph{stānrocc} `stone-coat' belongs to a different lexical history
(\citeproc{ref-ClarkHall1960}{Clark Hall 1960};
\citeproc{ref-BosworthToller1898}{Bosworth and Toller 1898}).

\section{\texorpdfstring{corn --- OE
\emph{corn}}{corn --- OE corn}}\label{corn-oe-corn}

\index[iv]{02oe@\ivlangheader{Old English}{West Saxon normalization unmarked}!corn@\iventry{corn}{}}
\index[iv]{02oe@\ivlangheader{Old English}{West Saxon normalization unmarked}!cornes@\iventry{cornes}{}}
\index[iv]{03pgmc@\ivlangheader{Proto-Germanic}{}!kurna@\iventry{*kúrną}{}}

Derivation: \emph{*kúrną} \textgreater{} \emph{corn} (regular).

\subsection{Derivation trace}\label{derivation-trace-13}

Proto input: \emph{*kúrną}

\begingroup
\setlength{\fboxsep}{6pt}

\noindent\fbox{%
\begin{minipage}{0.97\linewidth}
\small
\begin{tabular}{@{}>{\raggedright\arraybackslash}p{0.359\linewidth}@{\hspace{0.030\linewidth}}>{\raggedright\arraybackslash}p{0.611\linewidth}@{}}
\begin{minipage}[t]{\linewidth}
\centering\textbf{Earlier Germanic changes}\par
\vspace{0.35em}
\raggedright
\centering\textbf{Northwest and West Germanic}\par
\raggedright
\vspace{0.2em}
\begin{tabularx}{\linewidth}{@{}>{\raggedright\arraybackslash}X@{\hspace{0.55em}}>{\raggedright\arraybackslash}p{0.14\linewidth}@{\hspace{0.25em}}}
\mbox{PNWGmc U Lowering} & \emph{*kórną} \\
\end{tabularx}
\vspace{0.6em}
\centering\textbf{Early Anglo-Frisian}\par
\raggedright
\vspace{0.2em}
\raggedright [no change]\par
\end{minipage}
&
\begin{minipage}[t]{\linewidth}
\centering\textbf{Old English changes}\par
\vspace{0.35em}
\raggedright
\begin{tabularx}{\linewidth}{@{}>{\raggedright\arraybackslash}X@{\hspace{0.55em}}>{\raggedright\arraybackslash}p{0.14\linewidth}@{\hspace{0.25em}}}
OE Heavy Syllable Nasal Apocope & \emph{*kórn} \\
\end{tabularx}
\end{minipage}
\\
\end{tabular}
\end{minipage}%
} \endgroup

Old English form: \emph{corn}

\subsection{Reconstruction and comparative
evidence}\label{reconstruction-and-comparative-evidence-13}

Kroonen cites the noun as \emph{*kurna-}, and Orel gives the citation
form \Recon{kurnan} `corn', both with Old English
{\emph{corn}}\index[iv]{02oe@\ivlangheader{Old English}{West Saxon normalization unmarked}!corn@\iventry{corn}{}}
`corn, grain' among the reflexes (\citeproc{ref-Kroonen2013}{Kroonen
2013}; \citeproc{ref-Orel2003}{Orel 2003, 264}). The singular form
\Recon{kúrną} `corn' is the nominative-accusative singular appropriate
to the citation noun.

\subsection{Old English evidence}\label{old-english-evidence-13}

Clark Hall gives \emph{corn n.~`corn,' grain}, Bright's glossary lists
\emph{corn, n.} with genitive singular \emph{cornes}, and
Bosworth-Toller treats
{\emph{corn}}\index[iv]{02oe@\ivlangheader{Old English}{West Saxon normalization unmarked}!corn@\iventry{corn}{}}
`corn, grain' as an ordinary noun headword
(\citeproc{ref-ClarkHall1960}{Clark Hall 1960};
\citeproc{ref-BrightCassidyRingler1971}{Bright 1971, 347};
\citeproc{ref-BosworthToller1898}{Bosworth and Toller 1898, 144}). The
target is therefore an attested citation form, while forms such as
\emph{cornes} simply provide paradigm background.

\subsection{Development to Old
English}\label{development-to-old-english-13}

With northwest Germanic lowering, \Recon{kúrną} `corn' becomes
\Recon{kórną} `corn', and later loss of final nasal after a heavy
syllable yields \Recon{kórn} `corn', whence
{\emph{corn}}\index[iv]{02oe@\ivlangheader{Old English}{West Saxon normalization unmarked}!corn@\iventry{corn}{}}
`corn, grain'. The oblique form \Recon{kurnăn} `corn' belongs to
comparative background rather than to the derivational input of this
entry.

\section{\texorpdfstring{dart --- OE
\emph{sċytel}}{dart --- OE sċytel}}\label{dart-oe-sux10bytel}

\index[iv]{02oe@\ivlangheader{Old English}{West Saxon normalization unmarked}!scytel@\iventry{sċytel}{}}
\index[iv]{03pgmc@\ivlangheader{Proto-Germanic}{}!skutilaz@\iventry{*skutilaz}{}}
\index[iv]{03pgmc@\ivlangheader{Proto-Germanic}{}!skutilaz@\iventry{*skútilaz}{}}

Derivation: citation reconstruction \emph{*skutilaz}; form followed here
\emph{*skútilaz} \textgreater{} \emph{sċytel} (regular).

\subsection{Derivation trace}\label{derivation-trace-14}

Proto input: \emph{*skútilaz}

\begingroup
\setlength{\fboxsep}{6pt}

\noindent\fbox{%
\begin{minipage}{0.97\linewidth}
\small
\begin{tabular}{@{}>{\raggedright\arraybackslash}p{0.359\linewidth}@{\hspace{0.030\linewidth}}>{\raggedright\arraybackslash}p{0.611\linewidth}@{}}
\begin{minipage}[t]{\linewidth}
\centering\textbf{Earlier Germanic changes}\par
\vspace{0.35em}
\raggedright
\centering\textbf{Northwest and West Germanic}\par
\raggedright
\vspace{0.2em}
\begin{tabularx}{\linewidth}{@{}>{\raggedright\arraybackslash}X@{\hspace{0.55em}}>{\raggedright\arraybackslash}p{0.24\linewidth}@{\hspace{0.25em}}}
\mbox{EAF Final Z Deletion} & \emph{*skútila} \\
\end{tabularx}
\vspace{0.6em}
\centering\textbf{Early Anglo-Frisian}\par
\raggedright
\vspace{0.2em}
\raggedright [no change]\par
\end{minipage}
&
\begin{minipage}[t]{\linewidth}
\centering\textbf{Old English changes}\par
\vspace{0.35em}
\raggedright
\begin{tabularx}{\linewidth}{@{}>{\raggedright\arraybackslash}X@{\hspace{0.55em}}>{\raggedright\arraybackslash}p{0.14\linewidth}@{\hspace{0.25em}}}
\mbox{PWGmc Final Bare A Loss} & \emph{*skútil} \\
\mbox{OE Sk Palatalization} & \emph{*ʃútil} \\
\mbox{OE I Umlaut} & \emph{*ʃytil} \\
OE Med Unstressed I Lowering1 & \emph{*ʃytel} \\
\end{tabularx}
\end{minipage}
\\
\end{tabular}
\end{minipage}%
} \endgroup

Old English form: \emph{sċytel}

\subsection{Old English evidence and selected
cell}\label{old-english-evidence-and-selected-cell}

The Corpus Glossary has
{\emph{scytel}}\index[iv]{02oe@\ivlangheader{Old English}{West Saxon normalization unmarked}!scytel@\iventry{scytel}{}}
`balance beam', glossing Latin \emph{momentum}
(\citeproc{ref-Sweet1885}{Sweet 1885, 79}, Corpus 1325). The
corresponding Épinal and Erfurt item632 has
{\emph{scytil}}\index[iv]{02oe@\ivlangheader{Old English}{West Saxon normalization unmarked}!scytil@\iventry{scytil}{}}
`balance beam' (\citeproc{ref-Pheifer1974}{Pheifer 1974, 34}). Both
spellings preserve the vowel before final \emph{l}; the Corpus spelling
independently supplies \emph{e}.

The glossary meaning requires care. Pheifer interprets the earlier item
as the moving beam or lever of a balance, not its tongue or pointer
(\citeproc{ref-Pheifer1974}{Pheifer 1974, 101}, note 632). The
comparative dictionary also identifies Old English \emph{scytel} `dart'
beside Old Norse \emph{skutill} `harpoon, bolt'
(\citeproc{ref-Orel2003}{Orel 2003, 346--47}). CAPR does not claim that
the Latin glossary lemma means `dart'. The selected spelling belongs to
this moving-object noun; the lexical headword retains the dictionary's
projectile sense. Christiansen traces the dictionary confusion over
\emph{momentum} and confirms the moving-beam interpretation
(\citeproc{ref-Christiansen2018}{Christiansen 2018, 583--84}, note 3).

The normalized comparator
{\emph{sċytel}}\index[iv]{02oe@\ivlangheader{Old English}{West Saxon normalization unmarked}!scytel@\iventry{sċytel}{}}
`dart; balance beam' adds only the ordinary editorial palatal-consonant
marking to the attested Corpus spelling. No reconstructed \emph{ie},
altered root vowel or unattested medial vowel is substituted. Orel
explicitly identifies the reconstructed noun as masculine. CAPR selects
the ordinary strong masculine nominative singular; the uninflected
glossary token itself cannot distinguish nominative from accusative.
This citation-cell choice is not an assertion that the Latin lemma
establishes an Old English syntactic nominative.

The case evidence is nevertheless asymmetric. Christiansen quotes
\emph{nim heortes scytel} `take a stag's penis' and \emph{nim fearres
scytel} `take a bull's penis' from the \emph{Medicina de quadrupedibus}:
in each, \emph{scytel} is accusative singular, the object of \emph{nim}
`take' (\citeproc{ref-Christiansen2018}{Christiansen 2018, 589--90}).
The manuscript witnesses transmit a shared Old English translation, not
independent translations establishing separate lexical histories
(\citeproc{ref-Christiansen2018}{Christiansen 2018, 587--88}). Thus the
accusative is positively and more securely attested in syntax, whereas
nominative attestation remains ambiguous in the bare glossary forms.
CAPR nevertheless retains its default nominative citation input
\emph{*skútilaz}: nominative and accusative converge in Old English, and
the existing full-cascade experiments regularly produce \emph{sċytel}
from both that nominative and the inherited accusative \emph{*skútilą}.
The stronger accusative evidence is acknowledged, but it supplies no
substantive reason to depart from the nominative default here. No
selected-input revision follows from the new source.

\subsection{Comparative formation}\label{comparative-formation}

Orel prints the complete
\Recon{skutilaz}\index[iv]{03pgmc@\ivlangheader{Proto-Germanic}{}!skutilaz@\iventry{*skutilaz}{}}
`dart', masculine, and derives it from \emph{*skutan} `shot, shooting'
(\citeproc{ref-Orel2003}{Orel 2003, 346--47}). The Old Norse and Old
English witnesses support the inherited formation independently of the
computational result. The selected
\Recon{skútilaz}\index[iv]{03pgmc@\ivlangheader{Proto-Germanic}{}!skutilaz@\iventry{*skútilaz}{}}
`dart' marks ordinary initial root stress explicitly; no quantity,
suffix or nominative ending has been changed.

Christiansen independently reconstructs \emph{*skutilaz} `projectile',
comparing Old Norse \emph{skutill} `harpoon' and Old High German
\emph{scozila} `harpoon, missile', and discusses the instrumental
formation in \emph{*-ila-} (\citeproc{ref-Christiansen2018}{Christiansen
2018, 593, 595}). Her account separates the etymological stressed
\emph{y} of \emph{scytel} from the vowel of the \emph{scit-} `excrement'
family; that root-vowel issue is not the unstressed medial \emph{i/e}
alternation (\citeproc{ref-Christiansen2018}{Christiansen 2018,
593--98}).

Semantically, she treats medical \emph{scytel} `penis' as a metaphorical
extension of inherited `dart, missile'. The frequent
\emph{scyttel/scyttels} `bar, bolt' formations present additional
questions about gemination and suffixation; their relationship cannot
simply be assumed from overlapping meanings or spellings
(\citeproc{ref-Christiansen2018}{Christiansen 2018, 599--602}). Neither
those formations nor their proposed histories replace the
single-consonant noun selected here.

The projectile noun and its moving-beam usage need no morphological
transponent. The derivational input is an inherited complete word, not a
reconstruction beginning after the loss of its final ending. The precise
semantic development is not turned into a conditioning factor in any
sound law.

\subsection{Closed-syllable regression
protection}\label{closed-syllable-regression-protection}

Ordinary final-ending losses reduce \emph{*skútilaz} to \emph{*skútil}.
Palatalization and \hyperref[rule-OEIUmlaut]{SC055 OEIUmlaut} give
\emph{*ʃytil}. The medial high vowel survives
\hyperref[rule-OEHighVowelApocope]{SC063 OEHighVowelApocope}. At
\hyperref[rule-OELAdjacentSyncope]{SC066 OELAdjacentSyncope}, the word
has a light initial syllable followed by final closed \emph{*til}. The
vowel is not in the open-medial pre-\emph{l} environment of selected
\emph{nettle}. SC066 therefore correctly leaves it unchanged.

Later \hyperref[rule-OEMedUnstressedILowering1]{SC074
OEMedUnstressedILowering1} lowers medial \emph{*i} to \emph{e}, and
ordinary orthography gives \emph{sċytel} `dart; balance beam'. Removing
SC066's following-vowel restriction instead gives \Pred{sċytl}. That
wrong alternative changes none of the earlier selected outputs; this
noun therefore fills a genuine canonical closed-syllable gap rather than
supplying an arbitrary nonfiring word.

\section{\texorpdfstring{deed --- OE
\emph{dǣd}}{deed --- OE dǣd}}\label{deed-oe-dux1e3d}

\index[iv]{02oe@\ivlangheader{Old English}{West Saxon normalization unmarked}!daed@\iventry{dæd}{}}
\index[iv]{02oe@\ivlangheader{Old English}{West Saxon normalization unmarked}!daed@\iventry{dǣd}{}}
\index[iv]{02oe@\ivlangheader{Old English}{West Saxon normalization unmarked}!ded~04@\iventry{dēd}{Angl.}}
\index[iv]{03pgmc@\ivlangheader{Proto-Germanic}{}!dediz@\iventry{*dēdiz}{}}
\index[iv]{03pgmc@\ivlangheader{Proto-Germanic}{}!dediz@\iventry{*dḗdiz}{}}

Derivation: \emph{*dḗdiz} \textgreater{} \emph{dǣd} (regular).

\subsection{Derivation trace}\label{derivation-trace-15}

Proto input: \emph{*dḗdiz}

\begingroup
\setlength{\fboxsep}{6pt}

\noindent\fbox{%
\begin{minipage}{0.97\linewidth}
\small
\begin{tabular}{@{}>{\raggedright\arraybackslash}p{0.582\linewidth}@{\hspace{0.030\linewidth}}>{\raggedright\arraybackslash}p{0.388\linewidth}@{}}
\begin{minipage}[t]{\linewidth}
\centering\textbf{Earlier Germanic changes}\par
\vspace{0.35em}
\raggedright
\centering\textbf{Northwest and West Germanic}\par
\raggedright
\vspace{0.2em}
\begin{tabularx}{\linewidth}{@{}>{\raggedright\arraybackslash}X@{\hspace{0.55em}}>{\raggedright\arraybackslash}p{0.14\linewidth}@{\hspace{0.25em}}}
\mbox{PNWGmc Long E Lowering} & \emph{*dādiz} \\
\mbox{EAF Final Z Deletion} & \emph{*dādi} \\
\mbox{EAF Long A Fronting} & \emph{*dǣdi} \\
\end{tabularx}
\vspace{0.6em}
\centering\textbf{Early Anglo-Frisian}\par
\raggedright
\vspace{0.2em}
\raggedright [no change]\par
\end{minipage}
&
\begin{minipage}[t]{\linewidth}
\centering\textbf{Old English changes}\par
\vspace{0.35em}
\raggedright
\begin{tabularx}{\linewidth}{@{}>{\raggedright\arraybackslash}X@{\hspace{0.55em}}>{\raggedright\arraybackslash}p{0.14\linewidth}@{\hspace{0.25em}}}
\mbox{OE High Vowel Apocope} & \emph{*dǣd} \\
\end{tabularx}
\end{minipage}
\\
\end{tabular}
\end{minipage}%
} \endgroup

Old English form: \emph{dǣd}

\subsection{Reconstruction and comparative
evidence}\label{reconstruction-and-comparative-evidence-14}

Orel reconstructs the noun as \Recon{dēdiz} `deed', and Ringe and Taylor
derive the same inherited i-stem from Proto-Germanic \Recon{dédiz}
`deed' through northwest Germanic \Recon{dadiz} `deed'
(\citeproc{ref-Orel2003}{Orel 2003};
\citeproc{ref-RingeTaylor2014}{Ringe and Taylor 2014}). The acute accent
in \Recon{dḗdiz} `deed' marks stress on the same reconstructed long
vowel.

\subsection{Old English evidence}\label{old-english-evidence-14}

Campbell states that Primitive Germanic \emph{ē} appears as West Saxon
\emph{ǣ} but in other Old English dialects mostly as \emph{ē}, and
Brunner gives the contrast explicitly as West Saxon \emph{dǣd} `deed'
beside non-West-Saxon \emph{dēd} `deed'
(\citeproc{ref-Campbell1959}{Campbell 1959};
\citeproc{ref-SieversBrunner1965}{Brunner 1965}). Clark Hall likewise
lists \emph{dæd} `deed' and cross-refers Anglian \emph{dēd} to it
(\citeproc{ref-ClarkHall1960}{Clark Hall 1960}). West Saxon \emph{dǣd}
is therefore the relevant Old English form here, with Anglian \emph{dēd}
as a dialectal doublet.

\subsection{Development to Old
English}\label{development-to-old-english-14}

From inherited \Recon{dēdiz} `deed', loss of final \emph{-z} and the
West Saxon lowering of stressed long \emph{ē} yield \emph{dǣd} `deed';
Anglian \emph{dēd} `deed' preserves the non-West-Saxon outcome
(\citeproc{ref-Campbell1959}{Campbell 1959};
\citeproc{ref-SieversBrunner1965}{Brunner 1965}). The development
treated here is therefore the regular West Saxon line.

\section{\texorpdfstring{dew --- OE
\emph{dēaw}}{dew --- OE dēaw}}\label{dew-oe-dux113aw}

\index[iv]{02oe@\ivlangheader{Old English}{West Saxon normalization unmarked}!deaw@\iventry{dēaw}{}}
\index[iv]{03pgmc@\ivlangheader{Proto-Germanic}{}!dawwo@\iventry{*dáwwō}{}}

Derivation: \emph{*dáwwō} \textgreater{} \emph{dēaw} (regular).

\subsection{Derivation trace}\label{derivation-trace-16}

Proto input: \emph{*dáwwō}

\begingroup
\setlength{\fboxsep}{6pt}

\noindent\fbox{%
\begin{minipage}{0.97\linewidth}
\small
\begin{tabular}{@{}>{\raggedright\arraybackslash}p{0.485\linewidth}@{\hspace{0.030\linewidth}}>{\raggedright\arraybackslash}p{0.485\linewidth}@{}}
\begin{minipage}[t]{\linewidth}
\centering\textbf{Earlier Germanic changes}\par
\vspace{0.35em}
\raggedright
\centering\textbf{Northwest and West Germanic}\par
\raggedright
\vspace{0.2em}
\begin{tabularx}{\linewidth}{@{}>{\raggedright\arraybackslash}X@{\hspace{0.55em}}>{\raggedright\arraybackslash}p{0.14\linewidth}@{\hspace{0.25em}}}
\mbox{OE WW Simplification} & \emph{*dáuwō} \\
\mbox{PNWGmc Final Long O Raising} & \emph{*dáuwu} \\
\end{tabularx}
\vspace{0.6em}
\centering\textbf{Early Anglo-Frisian}\par
\raggedright
\vspace{0.2em}
\raggedright [no change]\par
\end{minipage}
&
\begin{minipage}[t]{\linewidth}
\centering\textbf{Old English changes}\par
\vspace{0.35em}
\raggedright
\begin{tabularx}{\linewidth}{@{}>{\raggedright\arraybackslash}X@{\hspace{0.55em}}>{\raggedright\arraybackslash}p{0.14\linewidth}@{\hspace{0.25em}}}
\mbox{OE Au Brightening} & \emph{*dáeuwu} \\
\mbox{OE Diphthong Leveling} & \emph{*dēawu} \\
\mbox{OE High Vowel Apocope} & \emph{*dēaw} \\
\end{tabularx}
\end{minipage}
\\
\end{tabular}
\end{minipage}%
} \endgroup

Old English form: \emph{dēaw}

\subsection{Reconstruction and
development}\label{reconstruction-and-development}

The selected PGmc \emph{*dáwwō} gives OE \emph{dēaw} `dew'. Its
inherited geminate is reanalyzed early: short \emph{*a} plus \emph{*ww}
becomes \emph{*au} with retained consonantal \emph{*w}. The nucleus
subsequently undergoes English fronting and offglide realization. Early
reanalysis must therefore be distinguished from the later long English
surface diphthong (\citeproc{ref-RingeTaylor2014}{Ringe and Taylor 2014,
65--66, 171--75}; \citeproc{ref-Campbell1959}{Campbell 1959, 45--47}).

The current path is \hyperref[rule-OEWWSimplification]{SC031
OEWWSimplification}, then \hyperref[rule-OEAuBrightening]{SC030
OEAuBrightening}, then \hyperref[rule-OEDiphthongLeveling]{SC032
OEDiphthongLeveling}. The unchanged reconstruction and final form no
longer depend on a late unconditional ww deletion followed by
singleton-aw lengthening. The selected evaluation uses the explicitly
declared strong-final citation context; lexical accent does not infer
sentence stress (\citeproc{ref-RingeTaylor2014}{Ringe and Taylor 2014,
55, 57--58}).

\section{\texorpdfstring{door --- OE
\emph{dor}}{door --- OE dor}}\label{door-oe-dor}

\index[iv]{02oe@\ivlangheader{Old English}{West Saxon normalization unmarked}!dor@\iventry{dor}{}}
\index[iv]{03pgmc@\ivlangheader{Proto-Germanic}{}!dura@\iventry{*dúrą}{}}
\index[iv]{10ohg@\ivlangheader{Old High German}{}!tura@\iventry{tura}{}}
\index[iv]{11ofris@\ivlangheader{Old Frisian}{}!dore@\iventry{dore}{}}

Derivation: \emph{*dúrą} \textgreater{} \emph{dor} (regular).

\subsection{Derivation trace}\label{derivation-trace-17}

Proto input: \emph{*dúrą}

\begingroup
\setlength{\fboxsep}{6pt}

\noindent\fbox{%
\begin{minipage}{0.97\linewidth}
\small
\begin{tabular}{@{}>{\raggedright\arraybackslash}p{0.359\linewidth}@{\hspace{0.030\linewidth}}>{\raggedright\arraybackslash}p{0.611\linewidth}@{}}
\begin{minipage}[t]{\linewidth}
\centering\textbf{Earlier Germanic changes}\par
\vspace{0.35em}
\raggedright
\centering\textbf{Northwest and West Germanic}\par
\raggedright
\vspace{0.2em}
\begin{tabularx}{\linewidth}{@{}>{\raggedright\arraybackslash}X@{\hspace{0.55em}}>{\raggedright\arraybackslash}p{0.14\linewidth}@{\hspace{0.25em}}}
\mbox{PNWGmc U Lowering} & \emph{*dórą} \\
\end{tabularx}
\vspace{0.6em}
\centering\textbf{Early Anglo-Frisian}\par
\raggedright
\vspace{0.2em}
\raggedright [no change]\par
\end{minipage}
&
\begin{minipage}[t]{\linewidth}
\centering\textbf{Old English changes}\par
\vspace{0.35em}
\raggedright
\begin{tabularx}{\linewidth}{@{}>{\raggedright\arraybackslash}X@{\hspace{0.55em}}>{\raggedright\arraybackslash}p{0.14\linewidth}@{\hspace{0.25em}}}
OE Heavy Syllable Nasal Apocope & \emph{*dór} \\
\end{tabularx}
\end{minipage}
\\
\end{tabular}
\end{minipage}%
} \endgroup

Old English form: \emph{dor}

\subsection{Reconstruction and comparative
evidence}\label{reconstruction-and-comparative-evidence-15}

Kroonen reconstructs a neuter \emph{*dura-} `gate, (single) door' and
cites Old English
{\emph{dor}}\index[iv]{02oe@\ivlangheader{Old English}{West Saxon normalization unmarked}!dor@\iventry{dor}{}}
`door' among its reflexes. In the same entry he separates Old English
{\emph{duru}}\index[iv]{02oe@\ivlangheader{Old English}{West Saxon normalization unmarked}!duru@\iventry{duru}{}}
`door', Old Frisian \emph{dore}, and Old High German
{\emph{tura}}\index[iv]{10ohg@\ivlangheader{Old High German}{}!tura@\iventry{tura}{}}
`door' as reflexes of \emph{*durō-} instead
(\citeproc{ref-Kroonen2013}{Kroonen 2013}).

\subsection{Old English evidence}\label{old-english-evidence-15}

Clark Hall records
{\emph{dor}}\index[iv]{02oe@\ivlangheader{Old English}{West Saxon normalization unmarked}!dor@\iventry{dor}{}}
`door' as a neuter noun and separately records feminine
{\emph{duru}}\index[iv]{02oe@\ivlangheader{Old English}{West Saxon normalization unmarked}!duru@\iventry{duru}{}}
`door' with its own inflection (\citeproc{ref-ClarkHall1960}{Clark Hall
1960}). Ringe and Taylor likewise treat
{\emph{duru}}\index[iv]{02oe@\ivlangheader{Old English}{West Saxon normalization unmarked}!duru@\iventry{duru}{}}
`door' as an early Old English u-stem, originally a root noun shifted
into that class (\citeproc{ref-RingeTaylor2014}{Ringe and Taylor 2014}).
The Old English form here is therefore the attested neuter
{\emph{dor}}\index[iv]{02oe@\ivlangheader{Old English}{West Saxon normalization unmarked}!dor@\iventry{dor}{}}
`door', while
{\emph{duru}}\index[iv]{02oe@\ivlangheader{Old English}{West Saxon normalization unmarked}!duru@\iventry{duru}{}}
`door' remains a parallel Old English reflex from another stem history.

\subsection{Development to Old
English}\label{development-to-old-english-15}

From \Recon{dúrą} `door', Northwest Germanic u-lowering gives
\Recon{dórą} `door', and heavy-syllable nasal apocope then yields
{\emph{dor}}\index[iv]{02oe@\ivlangheader{Old English}{West Saxon normalization unmarked}!dor@\iventry{dor}{}}
`door'. The regular development treated in this entry is therefore
\emph{*dúrą} \textgreater{} \emph{dor}; the feminine
{\emph{duru}}\index[iv]{02oe@\ivlangheader{Old English}{West Saxon normalization unmarked}!duru@\iventry{duru}{}}
`door' belongs to the separate line identified by Kroonen and
Ringe-Taylor (\citeproc{ref-Kroonen2013}{Kroonen 2013};
\citeproc{ref-RingeTaylor2014}{Ringe and Taylor 2014}).

\section{\texorpdfstring{fare --- OE
\emph{faran}}{fare --- OE faran}}\label{fare-oe-faran}

\index[iv]{02oe@\ivlangheader{Old English}{West Saxon normalization unmarked}!faeran@\iventry{færan}{}}
\index[iv]{02oe@\ivlangheader{Old English}{West Saxon normalization unmarked}!faran@\iventry{faran}{}}
\index[iv]{03pgmc@\ivlangheader{Proto-Germanic}{}!farana@\iventry{*fáraną}{}}
\index[iv]{03pgmc@\ivlangheader{Proto-Germanic}{}!faranan@\iventry{*faranan}{}}

Derivation: \emph{*fáraną} \textgreater{} \emph{faran} (regular).

\subsection{Derivation trace}\label{derivation-trace-18}

Proto input: \emph{*fáraną}

\begingroup
\setlength{\fboxsep}{6pt}

\noindent\fbox{%
\begin{minipage}{0.97\linewidth}
\small
\begin{tabular}{@{}>{\raggedright\arraybackslash}p{0.359\linewidth}@{\hspace{0.030\linewidth}}>{\raggedright\arraybackslash}p{0.611\linewidth}@{}}
\begin{minipage}[t]{\linewidth}
\centering\textbf{Earlier Germanic changes}\par
\vspace{0.35em}
\raggedright
\centering\textbf{Northwest and West Germanic}\par
\raggedright
\vspace{0.2em}
\raggedright [no change]\par
\vspace{0.6em}
\centering\textbf{Early Anglo-Frisian}\par
\raggedright
\vspace{0.2em}
\raggedright [no change]\par
\end{minipage}
&
\begin{minipage}[t]{\linewidth}
\centering\textbf{Old English changes}\par
\vspace{0.35em}
\raggedright
\begin{tabularx}{\linewidth}{@{}>{\raggedright\arraybackslash}X@{\hspace{0.55em}}>{\raggedright\arraybackslash}p{0.14\linewidth}@{\hspace{0.25em}}}
\mbox{EAF Brightening} & \emph{*færaną} \\
\mbox{OE A Restoration} & \emph{*faraną} \\
OE Heavy Syllable Nasal Apocope & \emph{*faran} \\
\mbox{OE Secondary Nasalization} & \emph{*farąn} \\
\mbox{OE Weak Tail Reduction} & \emph{*faran} \\
\end{tabularx}
\end{minipage}
\\
\end{tabular}
\end{minipage}%
} \endgroup

Old English form: \emph{faran}

\subsection{Reconstruction and comparative
evidence}\label{reconstruction-and-comparative-evidence-16}

Kroonen gives the inherited strong verb as \emph{*faran-}, and Orel
gives the same lexeme as \Recon{faranan} `fare', both with Old English
\emph{faran} `fare' among the reflexes
(\citeproc{ref-Kroonen2013}{Kroonen 2013}; \citeproc{ref-Orel2003}{Orel
2003, 132}). Campbell also uses \emph{faran} `fare' as a standard
example of Old English A-restoration
(\citeproc{ref-Campbell1959}{Campbell 1959, 61}).

\subsection{Old English evidence}\label{old-english-evidence-16}

Clark Hall lemmatizes the strong verb as \emph{faran} `fare' and
separately records weak \emph{færan} `to frighten';
{\emph{fære}}\index[iv]{02oe@\ivlangheader{Old English}{West Saxon normalization unmarked}!faere@\iventry{fære}{}}
`fare',
{\emph{færst}}\index[iv]{02oe@\ivlangheader{Old English}{West Saxon normalization unmarked}!faerst@\iventry{færst}{}}
`fare', and
{\emph{færð}}\index[iv]{02oe@\ivlangheader{Old English}{West Saxon normalization unmarked}!faerd@\iventry{færð}{}}
`fare' belong to present-tense forms of \emph{faran} rather than to the
infinitive itself (\citeproc{ref-ClarkHall1960}{Clark Hall 1960}).
Bosworth-Toller preserves the same distinction
(\citeproc{ref-BosworthToller1898}{Bosworth and Toller 1898, 108}). The
Old English form here is therefore the attested citation infinitive
\emph{faran} `fare'.

\subsection{Development to Old
English}\label{development-to-old-english-16}

From \Recon{fáraną} `fare', Anglo-Frisian brightening first gives
\Recon{færaną} `fare', but A-restoration before single \emph{r} returns
\Recon{faraną} `fare'; later apocope and weak-tail reduction yield
\emph{faran} `fare' (\citeproc{ref-Campbell1959}{Campbell 1959, 61}).
Fulk's contrast with participial faren- \textless{} \emph{*faræn-}
\textless{} \emph{*faran-} shows why fronting elsewhere in the paradigm
does not alter the infinitive headword (\citeproc{ref-Fulk2018}{Fulk
2018}).

\section{\texorpdfstring{fell --- OE
\emph{fell}}{fell --- OE fell}}\label{fell-oe-fell}

\index[iv]{02oe@\ivlangheader{Old English}{West Saxon normalization unmarked}!fel@\iventry{fel}{}}
\index[iv]{02oe@\ivlangheader{Old English}{West Saxon normalization unmarked}!fell@\iventry{fell}{}}
\index[iv]{02oe@\ivlangheader{Old English}{West Saxon normalization unmarked}!fellan@\iventry{fellan}{}}
\index[iv]{02oe@\ivlangheader{Old English}{West Saxon normalization unmarked}!fellum@\iventry{fellum}{}}
\index[iv]{02oe@\ivlangheader{Old English}{West Saxon normalization unmarked}!feoll@\iventry{feoll}{}}
\index[iv]{03pgmc@\ivlangheader{Proto-Germanic}{}!fella@\iventry{*féllą}{}}

Derivation: \emph{*féllą} \textgreater{} \emph{fell} (regular).

\subsection{Derivation trace}\label{derivation-trace-19}

Proto input: \emph{*féllą}

\begingroup
\setlength{\fboxsep}{6pt}

\noindent\fbox{%
\begin{minipage}{0.97\linewidth}
\small
\begin{tabular}{@{}>{\raggedright\arraybackslash}p{0.359\linewidth}@{\hspace{0.030\linewidth}}>{\raggedright\arraybackslash}p{0.611\linewidth}@{}}
\begin{minipage}[t]{\linewidth}
\centering\textbf{Earlier Germanic changes}\par
\vspace{0.35em}
\raggedright
\centering\textbf{Northwest and West Germanic}\par
\raggedright
\vspace{0.2em}
\raggedright [no change]\par
\vspace{0.6em}
\centering\textbf{Early Anglo-Frisian}\par
\raggedright
\vspace{0.2em}
\raggedright [no change]\par
\end{minipage}
&
\begin{minipage}[t]{\linewidth}
\centering\textbf{Old English changes}\par
\vspace{0.35em}
\raggedright
\begin{tabularx}{\linewidth}{@{}>{\raggedright\arraybackslash}X@{\hspace{0.55em}}>{\raggedright\arraybackslash}p{0.14\linewidth}@{\hspace{0.25em}}}
OE Heavy Syllable Nasal Apocope & \emph{*féll} \\
\end{tabularx}
\end{minipage}
\\
\end{tabular}
\end{minipage}%
} \endgroup

Old English form: \emph{fell}

\subsection{Reconstruction and comparative
evidence}\label{reconstruction-and-comparative-evidence-17}

Kroonen reconstructs the noun as \emph{*fella-} `membrane, skin, hide'
and cites Old English
{\emph{fell}}\index[iv]{02oe@\ivlangheader{Old English}{West Saxon normalization unmarked}!fell@\iventry{fell}{}}
`skin, fell' beside Dutch
{\emph{vel}}\index[iv]{16dutch@\ivlangheader{Dutch}{}!vel@\iventry{vel}{}}
`skin, fell' and German \emph{Fell} (\citeproc{ref-Kroonen2013}{Kroonen
2013}). The derivational input \Recon{féllą} `fell' is the derivable
singular form of that same inherited noun.

\subsection{Old English evidence}\label{old-english-evidence-17}

Clark Hall records
{\emph{fell}}\index[iv]{02oe@\ivlangheader{Old English}{West Saxon normalization unmarked}!fell@\iventry{fell}{}}
`skin, fell' as the noun `fell, skin, hide', and Bright's glossary
likewise gives
{\emph{fell}}\index[iv]{02oe@\ivlangheader{Old English}{West Saxon normalization unmarked}!fell@\iventry{fell}{}}
`skin, fell' with inflected forms such as accusative singular \emph{fel}
and dative plural \emph{fellum} (\citeproc{ref-ClarkHall1960}{Clark Hall
1960}; \citeproc{ref-BrightCassidyRingler1971}{Bright 1971}). The target
is therefore the attested noun
{\emph{fell}}\index[iv]{02oe@\ivlangheader{Old English}{West Saxon normalization unmarked}!fell@\iventry{fell}{}}
`skin, fell', not the verb \emph{fellan} or the preterite \emph{feoll}.

\subsection{Development to Old
English}\label{development-to-old-english-17}

With \Recon{féllą} `fell', no special earlier reshaping is needed:
heavy-syllable nasal apocope yields \Recon{féll} `fell', surfacing as
{\emph{fell}}\index[iv]{02oe@\ivlangheader{Old English}{West Saxon normalization unmarked}!fell@\iventry{fell}{}}
`skin, fell'. The regular development treated here is therefore
\emph{*féllą} \textgreater{} \emph{fell}.

\section{\texorpdfstring{fern --- OE
\emph{fearn}}{fern --- OE fearn}}\label{fern-oe-fearn}

\index[iv]{02oe@\ivlangheader{Old English}{West Saxon normalization unmarked}!fearn@\iventry{fearn}{}}
\index[iv]{02oe@\ivlangheader{Old English}{West Saxon normalization unmarked}!fearna@\iventry{fearna}{}}
\index[iv]{02oe@\ivlangheader{Old English}{West Saxon normalization unmarked}!fearne@\iventry{fearne}{}}
\index[iv]{02oe@\ivlangheader{Old English}{West Saxon normalization unmarked}!fearnes@\iventry{fearnes}{}}
\index[iv]{03pgmc@\ivlangheader{Proto-Germanic}{}!farnan@\iventry{*farnan}{}}
\index[iv]{03pgmc@\ivlangheader{Proto-Germanic}{}!farnaz@\iventry{*fárnaz}{}}

Derivation: \emph{*fárnaz} \textgreater{} \emph{fearn} (regular).

\subsection{Derivation trace}\label{derivation-trace-20}

Proto input: \emph{*fárnaz}

\begingroup
\setlength{\fboxsep}{6pt}

\noindent\fbox{%
\begin{minipage}{0.97\linewidth}
\small
\begin{tabular}{@{}>{\raggedright\arraybackslash}p{0.485\linewidth}@{\hspace{0.030\linewidth}}>{\raggedright\arraybackslash}p{0.485\linewidth}@{}}
\begin{minipage}[t]{\linewidth}
\centering\textbf{Earlier Germanic changes}\par
\vspace{0.35em}
\raggedright
\centering\textbf{Northwest and West Germanic}\par
\raggedright
\vspace{0.2em}
\begin{tabularx}{\linewidth}{@{}>{\raggedright\arraybackslash}X@{\hspace{0.55em}}>{\raggedright\arraybackslash}p{0.14\linewidth}@{\hspace{0.25em}}}
\mbox{EAF Final Z Deletion} & \emph{*fárna} \\
\end{tabularx}
\vspace{0.6em}
\centering\textbf{Early Anglo-Frisian}\par
\raggedright
\vspace{0.2em}
\raggedright [no change]\par
\end{minipage}
&
\begin{minipage}[t]{\linewidth}
\centering\textbf{Old English changes}\par
\vspace{0.35em}
\raggedright
\begin{tabularx}{\linewidth}{@{}>{\raggedright\arraybackslash}X@{\hspace{0.55em}}>{\raggedright\arraybackslash}p{0.14\linewidth}@{\hspace{0.25em}}}
\mbox{PWGmc Final Bare A Loss} & \emph{*fárn} \\
\mbox{EAF Brightening} & \emph{*færn} \\
\mbox{OE Breaking} & \emph{*fearn} \\
\end{tabularx}
\end{minipage}
\\
\end{tabular}
\end{minipage}%
} \endgroup

Old English form: \emph{fearn}

\subsection{Reconstruction and comparative
evidence}\label{reconstruction-and-comparative-evidence-18}

Kroonen cites the noun as masculine \emph{*farna-} and gives Old English
\emph{fearn, fern}, while Orel gives the same lexeme as neuter
\Recon{farnan} `fern' with Old English
{\emph{fearn}}\index[iv]{02oe@\ivlangheader{Old English}{West Saxon normalization unmarked}!fearn@\iventry{fearn}{}}
`fern' (\citeproc{ref-Kroonen2013}{Kroonen 2013};
\citeproc{ref-Orel2003}{Orel 2003, 133}). Those are comparative headword
conventions rather than competing Old English outcomes; the modeled
input here is the nominative-style \Recon{fárnaz} `fern'.

\subsection{Old English evidence}\label{old-english-evidence-18}

Clark Hall gives
{\emph{fearn}}\index[iv]{02oe@\ivlangheader{Old English}{West Saxon normalization unmarked}!fearn@\iventry{fearn}{}}
`fern' as an Old English noun, and Bosworth-Toller records
{\emph{fearn}}\index[iv]{02oe@\ivlangheader{Old English}{West Saxon normalization unmarked}!fearn@\iventry{fearn}{}}
`fern' with inflected forms such as \emph{fearnes}, \emph{fearna}, and
\emph{fearne} (\citeproc{ref-ClarkHall1960}{Clark Hall 1960};
\citeproc{ref-BosworthToller1898}{Bosworth and Toller 1898, 219}).
Kroonen also records
{\emph{fern}}\index[iv]{02oe@\ivlangheader{Old English}{West Saxon normalization unmarked}!fern@\iventry{fern}{}}
`fern', but the local lexical sources give stronger support to
{\emph{fearn}}\index[iv]{02oe@\ivlangheader{Old English}{West Saxon normalization unmarked}!fearn@\iventry{fearn}{}}
`fern' as the citation target (\citeproc{ref-Kroonen2013}{Kroonen
2013}).

\subsection{Development to Old
English}\label{development-to-old-english-18}

From \Recon{fárnaz} `fern', loss of final \emph{-z} and final \emph{-a}
gives \Recon{fárn} `fern'; Anglo-Frisian brightening then yields
\Recon{færn} `fern', and breaking before \emph{r} plus consonant gives
{\emph{fearn}}\index[iv]{02oe@\ivlangheader{Old English}{West Saxon normalization unmarked}!fearn@\iventry{fearn}{}}
`fern' (\citeproc{ref-Campbell1959}{Campbell 1959};
\citeproc{ref-RingeTaylor2014}{Ringe and Taylor 2014}). The development
treated here is therefore the regular \emph{rC}-breaking line.

\section{\texorpdfstring{field --- OE
\emph{feld}}{field --- OE feld}}\label{field-oe-feld}

\index[iv]{02oe@\ivlangheader{Old English}{West Saxon normalization unmarked}!feld@\iventry{feld}{}}
\index[iv]{02oe@\ivlangheader{Old English}{West Saxon normalization unmarked}!felda@\iventry{felda}{}}
\index[iv]{02oe@\ivlangheader{Old English}{West Saxon normalization unmarked}!felde@\iventry{felde}{}}
\index[iv]{03pgmc@\ivlangheader{Proto-Germanic}{}!felthuz@\iventry{*félθuz}{}}

Derivation: \emph{*félθuz} \textgreater{} \emph{feld} (regular).

\subsection{Derivation trace}\label{derivation-trace-21}

Proto input: \emph{*félθuz}

\begingroup
\setlength{\fboxsep}{6pt}

\noindent\fbox{%
\begin{minipage}{0.97\linewidth}
\small
\begin{tabular}{@{}>{\raggedright\arraybackslash}p{0.485\linewidth}@{\hspace{0.030\linewidth}}>{\raggedright\arraybackslash}p{0.485\linewidth}@{}}
\begin{minipage}[t]{\linewidth}
\centering\textbf{Earlier Germanic changes}\par
\vspace{0.35em}
\raggedright
\centering\textbf{Northwest and West Germanic}\par
\raggedright
\vspace{0.2em}
\begin{tabularx}{\linewidth}{@{}>{\raggedright\arraybackslash}X@{\hspace{0.55em}}>{\raggedright\arraybackslash}p{0.14\linewidth}@{\hspace{0.25em}}}
\mbox{EAFL Th Voicing} & \emph{*félduz} \\
\mbox{EAF Final Z Deletion} & \emph{*féldu} \\
\end{tabularx}
\vspace{0.6em}
\centering\textbf{Early Anglo-Frisian}\par
\raggedright
\vspace{0.2em}
\raggedright [no change]\par
\end{minipage}
&
\begin{minipage}[t]{\linewidth}
\centering\textbf{Old English changes}\par
\vspace{0.35em}
\raggedright
\begin{tabularx}{\linewidth}{@{}>{\raggedright\arraybackslash}X@{\hspace{0.55em}}>{\raggedright\arraybackslash}p{0.14\linewidth}@{\hspace{0.25em}}}
\mbox{OE High Vowel Apocope} & \emph{*féld} \\
\end{tabularx}
\end{minipage}
\\
\end{tabular}
\end{minipage}%
} \endgroup

Old English form: \emph{feld}

\subsection{Reconstruction and comparative
evidence}\label{reconstruction-and-comparative-evidence-19}

Ringe and Taylor treat Old English
{\emph{feld}}\index[iv]{02oe@\ivlangheader{Old English}{West Saxon normalization unmarked}!feld@\iventry{feld}{}}
`field' as one of the cases where earlier \emph{*felþu-}
\textasciitilde{} \emph{*feldu-} may reflect either inherited \emph{þ}
\textasciitilde{} \emph{*d} alternation or the regular West Germanic
development \emph{*lþ} \textgreater{} \emph{ld}
(\citeproc{ref-RingeTaylor2014}{Ringe and Taylor 2014, 170}). The
broader proto label \Recon{félθuz} `field' can remain as comparative
background, while that narrower historical ambiguity does not affect the
regular classification.

\subsection{Old English evidence}\label{old-english-evidence-19}

Clark Hall records
{\emph{feld}}\index[iv]{02oe@\ivlangheader{Old English}{West Saxon normalization unmarked}!feld@\iventry{feld}{}}
`field' with oblique forms such as \emph{felda} and \emph{felde}, and
Campbell notes early place-name spellings in \emph{-felth} beside the
later standard form (\citeproc{ref-ClarkHall1960}{Clark Hall 1960, 114};
\citeproc{ref-Campbell1959}{Campbell 1959, 169}). The Old English form
here is therefore the attested citation noun
{\emph{feld}}\index[iv]{02oe@\ivlangheader{Old English}{West Saxon normalization unmarked}!feld@\iventry{feld}{}}
`field', with the older \emph{-felth} spellings as historical support
rather than as rival targets.

\subsection{Development to Old
English}\label{development-to-old-english-19}

In the modeled pathway, medial \emph{*lþ} becomes \emph{ld}, final
\emph{-z} is lost, and high-vowel apocope then yields
{\emph{feld}}\index[iv]{02oe@\ivlangheader{Old English}{West Saxon normalization unmarked}!feld@\iventry{feld}{}}
`field'. Whether the voiced dental ultimately reflects inherited
alternation or the regular \emph{*lþ} \textgreater{} \emph{ld}
development, both accounts converge on the same Old English form
(\citeproc{ref-RingeTaylor2014}{Ringe and Taylor 2014, 170}).

\section{\texorpdfstring{filth --- OE
\emph{fȳlþe}}{filth --- OE fȳlþe}}\label{filth-oe-fux233luxfee}

\index[iv]{02oe@\ivlangheader{Old English}{West Saxon normalization unmarked}!fylthe@\iventry{fȳlþe}{}}
\index[iv]{03pgmc@\ivlangheader{Proto-Germanic}{}!fulitho@\iventry{*fūliþō}{}}
\index[iv]{03pgmc@\ivlangheader{Proto-Germanic}{}!fulitho@\iventry{*fūliþǭ}{}}

Derivation: citation reconstruction \emph{*fūliþō}; form followed here
\emph{*fūliþǭ} \textgreater{} \emph{fȳlþe} (regular).

\subsection{Derivation trace}\label{derivation-trace-22}

Proto input: \emph{*fūliθǭ}

\begingroup
\setlength{\fboxsep}{6pt}

\noindent\fbox{%
\begin{minipage}{0.97\linewidth}
\small
\begin{tabular}{@{}>{\raggedright\arraybackslash}p{0.359\linewidth}@{\hspace{0.030\linewidth}}>{\raggedright\arraybackslash}p{0.611\linewidth}@{}}
\begin{minipage}[t]{\linewidth}
\centering\textbf{Earlier Germanic changes}\par
\vspace{0.35em}
\raggedright
\centering\textbf{Northwest and West Germanic}\par
\raggedright
\vspace{0.2em}
\raggedright [no change]\par
\vspace{0.6em}
\centering\textbf{Early Anglo-Frisian}\par
\raggedright
\vspace{0.2em}
\raggedright [no change]\par
\end{minipage}
&
\begin{minipage}[t]{\linewidth}
\centering\textbf{Old English changes}\par
\vspace{0.35em}
\raggedright
\begin{tabularx}{\linewidth}{@{}>{\raggedright\arraybackslash}X@{\hspace{0.55em}}>{\raggedright\arraybackslash}p{0.14\linewidth}@{\hspace{0.25em}}}
\mbox{OE I Umlaut} & \emph{*fȳliθǭ} \\
\mbox{OE High Vowel Apocope} & \emph{*fȳlθǭ} \\
OE Unstressed Long Vowel Shortening & \emph{*fȳlθæ} \\
\mbox{OE Unstressed AE Merger} & \emph{*fȳlθe} \\
\end{tabularx}
\end{minipage}
\\
\end{tabular}
\end{minipage}%
} \endgroup

Old English form: \emph{fȳlþe}

\subsection{Old English evidence}\label{old-english-evidence-20}

Old English
{\emph{fȳlþ}}\index[iv]{02oe@\ivlangheader{Old English}{West Saxon normalization unmarked}!fylth@\iventry{fȳlþ}{}}
`filth, foulness, impurity' is a feminine abstract noun on the adjective
\emph{fūl} `foul'. The Toller Supplement lemmatizes
{\emph{fylþ}}\index[iv]{02oe@\ivlangheader{Old English}{West Saxon normalization unmarked}!fylth@\iventry{fylþ}{}}
`filth' without marking vowel quantity, in both physical and moral
senses: it glosses Latin \emph{putor} in the Aldhelm glosses (An. Ox.
3323) and \emph{putredo} in the \emph{Liber scintillarum} (Scint. 38,
19), and the entry also records the inflected form
{\emph{fylþu}}\index[iv]{02oe@\ivlangheader{Old English}{West Saxon normalization unmarked}!fylthu@\iventry{fylþu}{}}
`filth' (An. Ox. 1790) (\citeproc{ref-Toller1921}{Toller 1921, 276}).
Clark Hall lemmatizes
{\emph{fýlð}}\index[iv]{02oe@\ivlangheader{Old English}{West Saxon normalization unmarked}!fylth@\iventry{fýlð}{}}
`filth' as feminine and marks the long vowel
(\citeproc{ref-ClarkHall1960}{Clark Hall 1960, 125}). The former Old
English form here \emph{fȳlþ} was the endingless nominative singular,
normalized with marked length after Clark Hall and with thorn for the
dental fricative; the substitution of thorn for eth changes no segment.
The inflected \emph{fylþu} beside endingless \emph{fylþ} is the familiar
alternation of the \emph{-iþō} abstracts, seen also in \emph{strengþu}
`strength' beside \emph{hǣlþ} `health', \emph{þīefþ} `theft' and
\emph{wiermþu} `warmth' (\citeproc{ref-Campbell1959}{Campbell 1959,
236}, §589.6).

The form compared here is now accusative singular
{\emph{fȳlþe}}\index[iv]{02oe@\ivlangheader{Old English}{West Saxon normalization unmarked}!fylthe@\iventry{fȳlþe}{}},
independently attested in \emph{Ðā fylde ādōn of þām Godes temple},
`remove the filth from God's temple'. The noun is the object of
\emph{ādōn}; the feminine accusative ending is not inferred from the
transducer output (\citeproc{ref-Toller1921}{Toller 1921, 276, Hml.S25,
378}). Length and dental spelling are normalized as above. The
nominatives remain documented comparative forms, not alternative
selected cells.

\subsection{Reconstruction and comparative
evidence}\label{reconstruction-and-comparative-evidence-20}

Kroonen reconstructs the adjective base
\Recon{fūla-}\index[iv]{03pgmc@\ivlangheader{Proto-Germanic}{}!fula@\iventry{*fūla-}{}}
`foul, rotten' and cites among its reflexes an exact West Germanic
abstract, Old Saxon
{fulitha}\index[iv]{13os@\ivlangheader{Old Saxon}{}!fulitha@\iventry{fulitha}{}}
`filth' (\citeproc{ref-Kroonen2013}{Kroonen 2013, 158}). Orel likewise
prints the base adjective \emph{*fūlaz}, with OE \emph{fūl} `foul,
dirty, rotten' among its reflexes, but lists no derived abstract
(\citeproc{ref-Orel2003}{Orel 2003, 121}). The Toller entry itself cites
the Old Saxon noun together with Old High German
{fülida}\index[iv]{10ohg@\ivlangheader{Old High German}{}!fulida@\iventry{fülida}{}}
`filth' (\citeproc{ref-Toller1921}{Toller 1921, 276}). Bammesberger
documents the suffix type: the feminine abstracts in \emph{*-iþō} show
the vowel \emph{*-i-} before the dental throughout, as in \emph{*mēriþō}
\textgreater{} \emph{mǣrþu} `fame' and \emph{*hauhiþō} \textgreater{}
\emph{hīehþu} `height' (\citeproc{ref-Bammesberger1990}{Bammesberger
1990, 117--18}, §4.5).

The citation reconstruction
\Recon{fūliþō}\index[iv]{03pgmc@\ivlangheader{Proto-Germanic}{}!fulitho@\iventry{*fūliþō}{}}
`filth' is therefore an assembled morphological reconstruction: the
attested base \emph{*fūla-} composed with the attested suffix
\emph{*-iþō}. No single authority prints the composite verbatim; its two
members are each independently supported, and the Old Saxon and Old High
German nouns attest the derivative itself on the continent. Those
cognates retain the medial vowel, so the medial loss seen in Old English
is an English development, not a common West Germanic one.

The selected accusative
\Recon{fūliþǭ}\index[iv]{03pgmc@\ivlangheader{Proto-Germanic}{}!fulithoacc@\iventry{*fūliþǭ}{}}
combines that supported noun formation with ordinary strong feminine
accusative \emph{*-ǭ}. Citation and form compared here are both
Proto-Germanic cells, not different historical stages. This is the same
inflectional premise used in the selected \emph{sāwle} `soul' history; a
dictionary need not print every regularly inflected reconstructed form.

\subsection{Development to Old
English}\label{development-to-old-english-20}

The former selected derivation was treated as regular at the time; the
following sequence records that superseded analysis, not the current
law. \hyperref[rule-PNWGmcFinalLongORaising]{SC019, final long-\emph{ō}
raising} gives \emph{*fūliþō} \textgreater{} \emph{*fūliþu}.
\hyperref[rule-OEIUmlaut]{SC055, i-umlaut} fronts the root vowel over
the medial syllable, \emph{*fūliþu} \textgreater{} \emph{*fȳliþu}.
\hyperref[rule-OEHighVowelApocope]{SC063, high-vowel apocope} removes
the final high vowel after the heavy sequence, \emph{*fȳliþu}
\textgreater{} \emph{*fȳliþ}, and \hyperref[rule-OEMedialSyncope]{SC065,
medial syncope} deletes the now-medial \emph{*i} before the dental,
\emph{*fȳliþ} \textgreater{} \emph{fȳlþ}.

The selected accusative remains a live chronological witness. Had
deletion preceded i-umlaut, the \emph{*i} that conditions the umlaut
would have been lost first, giving unattested \emph{*fūlþe} rather than
\emph{fȳlþe}; umlaut must therefore precede this deletion. The same
mutation-before-loss relation is visible across the \emph{-iþō} class,
whose members uniformly surface with umlauted root vowels beside
syncopated suffixes (\citeproc{ref-Campbell1959}{Campbell 1959, 236},
§589.6; \citeproc{ref-Bammesberger1990}{Bammesberger 1990, 117--18},
§4.5).

Campbell derives the endingless nominative singular of these abstracts
analogically, after the \emph{lār} type, with the medial syncope taken
as given for the class (\citeproc{ref-Campbell1959}{Campbell 1959, 236},
§589.6). The former derivation reaches the same surface form by the
regular phonological route, through raising, apocope and syncope,
without an analogical step. That route is no longer adopted: its apocope
erroneously removed the final high vowel after a heavy syllable followed
by a light syllable. The corrected conventional history would retain
final \emph{*u}, giving \emph{fȳlþu}, not the endingless target. The
adopted quantitative account instead retains both weak vowels in HLL,
giving \emph{*fȳleþu}; this too is a nonselected nominative comparison.
Fulk additionally supplies positive paradigm evidence for analogical
syncope in abstract nominatives; his conservative light-final analysis
is a substantive competing account (\citeproc{ref-Fulk2010}{Fulk 2010,
140--42}).

The selected accusative avoids those disputed nominative histories.
\emph{*fūliþǭ} undergoes mutation to \emph{*fȳliþǭ}; the retained
nonhigh quantity-bearing ending is not a target of high-vowel deletion.
The single quantitative event under
\hyperref[rule-OEHighVowelApocope]{SC063 OEHighVowelApocope} gives
\emph{*fȳlþǭ}, and the ordinary inherited ending developments yield
\emph{fȳlþe}. \hyperref[rule-OEMedialSyncope]{SC065 OEMedialSyncope} is
an identity anchor, not a second productive deletion. Removing medial
\emph{*i} before mutation would again prevent root fronting. Thus the
independently attested cell compared here preserves the scientific
ordering witness without claiming that an endingless nominative is
regular merely because the former cascade produced it.

\section{\texorpdfstring{flea --- OE
\emph{flēah}}{flea --- OE flēah}}\label{flea-oe-flux113ah}

\index[iv]{02oe@\ivlangheader{Old English}{West Saxon normalization unmarked}!fleah@\iventry{flēah}{}}
\index[iv]{03pgmc@\ivlangheader{Proto-Germanic}{}!flauxz@\iventry{*fláuxz}{}}

Derivation: \emph{*fláuxz} \textgreater{} \emph{flēah} (regular).

\subsection{Derivation trace}\label{derivation-trace-23}

Proto input: \emph{*fláuxz}

\begingroup
\setlength{\fboxsep}{6pt}

\noindent\fbox{%
\begin{minipage}{0.97\linewidth}
\small
\begin{tabular}{@{}>{\raggedright\arraybackslash}p{0.388\linewidth}@{\hspace{0.030\linewidth}}>{\raggedright\arraybackslash}p{0.582\linewidth}@{}}
\begin{minipage}[t]{\linewidth}
\centering\textbf{Earlier Germanic changes}\par
\vspace{0.35em}
\raggedright
\centering\textbf{Northwest and West Germanic}\par
\raggedright
\vspace{0.2em}
\raggedright [no change]\par
\vspace{0.6em}
\centering\textbf{Early Anglo-Frisian}\par
\raggedright
\vspace{0.2em}
\raggedright [no change]\par
\end{minipage}
&
\begin{minipage}[t]{\linewidth}
\centering\textbf{Old English changes}\par
\vspace{0.35em}
\raggedright
\begin{tabularx}{\linewidth}{@{}>{\raggedright\arraybackslash}X@{\hspace{0.55em}}>{\raggedright\arraybackslash}p{0.14\linewidth}@{\hspace{0.25em}}}
\mbox{OE Au Brightening} & \emph{*fláeux} \\
\mbox{OE Diphthong Leveling} & \emph{*flēax} \\
\end{tabularx}
\end{minipage}
\\
\end{tabular}
\end{minipage}%
} \endgroup

Old English form: \emph{flēah}

\subsection{Reconstruction and comparative
evidence}\label{reconstruction-and-comparative-evidence-21}

This word carries a genuine stem-class dispute, not merely a notational
one. Orel reconstructs a root noun with his class-wide explicit
nominative marker, \emph{*flauxz} `flea' (\citeproc{ref-Orel2003}{Orel
2003, 105}). Kroonen instead reconstructs an ō-stem
{\emph{*flauhō-}}\index[iv]{03pgmc@\ivlangheader{Proto-Germanic}{}!flauho@\iventry{*flauhō-}{}}
`flea', with no \emph{-z} at any layer of the analysis
(\citeproc{ref-Kroonen2013}{Kroonen 2013, 145}). Kluge and Seebold offer
yet a third analysis, a masculine a-stem
{\emph{*flauha-}}\index[iv]{03pgmc@\ivlangheader{Proto-Germanic}{}!flauha@\iventry{*flauha-}{}}
`flea' (\citeproc{ref-KlugeSeebold2011}{Kluge 2011, 304}), whose
nominative singular \emph{*flauhaz} would have carried its \emph{-z} in
an unstressed syllable and lost it under the West Germanic unstressed
loss rather than with the root nouns. The attested forms do not decide
the question: Old English \emph{flēah} `flea' and Old Norse \emph{fló}
`flea' are compatible with all three reconstructions. This is a case
where reputable sources disagree about the morphological class itself,
and the disagreement is documented here rather than silently normalized
away.

The derivational input
\Recon{fláuxz}\index[iv]{03pgmc@\ivlangheader{Proto-Germanic}{}!flauxz@\iventry{*fláuxz}{}}
`flea' adopts Orel's root-noun analysis with its morphologically
explicit nominative marker. Under Kroonen's ō-stem analysis the
nominative would instead have ended in \emph{*-ō} and developed like the
other ō-stems; the choice of Orel's form keeps the word in the root-noun
class alongside `book', `goose', and `louse', whose endingless West
Germanic nominatives it shares (\citeproc{ref-RingeTaylor2014}{Ringe and
Taylor 2014, 118}).

\subsection{Old English evidence}\label{old-english-evidence-21}

Old English
{\emph{flēah}}\index[iv]{02oe@\ivlangheader{Old English}{West Saxon normalization unmarked}!fleah@\iventry{flēah}{}}
`flea' is the endingless nominative-accusative singular. The root-noun
class to which Orel assigns it shows no nominative ending anywhere in
West Germanic (\citeproc{ref-RingeTaylor2014}{Ringe and Taylor 2014,
118}); on Kroonen's ō-stem analysis the endingless form would instead
reflect the regular fate of the ō-stem nominative in this shape. Either
way the attested Old English form is the expected one, which is
precisely why the class question remains open.

\subsection{Development to Old
English}\label{development-to-old-english-21}

On the selected root-noun analysis, the nominative marker was eliminated
before Proto-West Germanic (\citeproc{ref-RingeTaylor2014}{Ringe and
Taylor 2014, 118}), giving \Recon{fláux} `flea' from the derivational
input \Recon{fláuxz} `flea'. The diphthong is then fronted and levelled
in Old English, yielding \Recon{flēax} `flea', spelled \emph{flēah}
`flea' with final \emph{h} for the fricative. The early loss of the
root-noun marker is historically distinct from the later West Germanic
loss of final \emph{*z} in unstressed syllables and from the northern
monosyllabic loss with compensatory lengthening; the derivational input
is a stressed monosyllable, but its \emph{*-z} stands after a consonant
and was gone long before the northern change, whose domain is
vowel-final monosyllables.

\section{\texorpdfstring{fly --- OE
\emph{flēogan}}{fly --- OE flēogan}}\label{fly-oe-flux113ogan}

\index[iv]{02oe@\ivlangheader{Old English}{West Saxon normalization unmarked}!flegan~04@\iventry{flégan}{Angl.}}
\index[iv]{02oe@\ivlangheader{Old English}{West Saxon normalization unmarked}!fleogan@\iventry{fléogan}{}}
\index[iv]{02oe@\ivlangheader{Old English}{West Saxon normalization unmarked}!fleogan@\iventry{flēogan}{}}
\index[iv]{02oe@\ivlangheader{Old English}{West Saxon normalization unmarked}!fleoged@\iventry{fleogeð}{}}
\index[iv]{03pgmc@\ivlangheader{Proto-Germanic}{}!fleugana@\iventry{*fléuganą}{}}

Derivation: \emph{*fléuganą} \textgreater{} \emph{flēogan} (regular).

\subsection{Derivation trace}\label{derivation-trace-24}

Proto input: \emph{*fléuganą}

\begingroup
\setlength{\fboxsep}{6pt}

\noindent\fbox{%
\begin{minipage}{0.97\linewidth}
\small
\begin{tabular}{@{}>{\raggedright\arraybackslash}p{0.359\linewidth}@{\hspace{0.030\linewidth}}>{\raggedright\arraybackslash}p{0.611\linewidth}@{}}
\begin{minipage}[t]{\linewidth}
\centering\textbf{Earlier Germanic changes}\par
\vspace{0.35em}
\raggedright
\centering\textbf{Northwest and West Germanic}\par
\raggedright
\vspace{0.2em}
\raggedright [no change]\par
\vspace{0.6em}
\centering\textbf{Early Anglo-Frisian}\par
\raggedright
\vspace{0.2em}
\raggedright [no change]\par
\end{minipage}
&
\begin{minipage}[t]{\linewidth}
\centering\textbf{Old English changes}\par
\vspace{0.35em}
\raggedright
\begin{tabularx}{\linewidth}{@{}>{\raggedright\arraybackslash}X@{\hspace{0.55em}}>{\raggedright\arraybackslash}p{0.20\linewidth}@{\hspace{0.25em}}}
\mbox{OE Diphthong Leveling} & \emph{*flēoganą} \\
OE Heavy Syllable Nasal Apocope & \emph{*flēogan} \\
\mbox{OE Secondary Nasalization} & \emph{*flēogąn} \\
\mbox{OE Weak Tail Reduction} & \emph{*flēogan} \\
\end{tabularx}
\end{minipage}
\\
\end{tabular}
\end{minipage}%
} \endgroup

Old English form: \emph{flēogan}

\subsection{Reconstruction and comparative
evidence}\label{reconstruction-and-comparative-evidence-22}

Ringe and Taylor derive the verb as \emph{*fleugana} `fly'
\textgreater{} OE \emph{fléogan} `fly' and elsewhere contrast West Saxon
\emph{fléogan} `fly' with Anglian \emph{flégan} `fly', alongside related
forms noted in the form note below (\citeproc{ref-RingeTaylor2014}{Ringe
and Taylor 2014}). The derivational input \Recon{fléuganą} `fly'
represents that inherited strong verb in the notation used here.

\subsection{Old English evidence}\label{old-english-evidence-22}

Clark Hall and Bosworth-Toller record \emph{flēogan} `fly' as the
ordinary Old English strong verb, and Bright gives the familiar paradigm
\emph{flēag} `flew', \emph{flugon} `flew', \emph{flogen} `flown' with
present \emph{fleogeð} `flies' (\citeproc{ref-ClarkHall1960}{Clark Hall
1960}; \citeproc{ref-BosworthToller1898}{Bosworth and Toller 1898};
\citeproc{ref-BrightCassidyRingler1971}{Bright 1971}). The target in
this entry is therefore the attested verbal infinitive itself.

\subsection{Form note}\label{form-note-1}

Ringe and Taylor also list related \emph{fléoge} `fly' / \emph{flége}
`fly' and Anglian \emph{flégan} `fly', which belong to the same family
but do not replace the infinitive \emph{flēogan} `fly' treated here
(\citeproc{ref-RingeTaylor2014}{Ringe and Taylor 2014}).

\subsection{Development to Old
English}\label{development-to-old-english-22}

From \Recon{fléuganą} `fly', Old English diphthong leveling gives
\Recon{flēoganą} `fly'; heavy-syllable nasal apocope and weak-tail
reduction then yield \emph{flēogan} `fly'
(\citeproc{ref-RingeTaylor2014}{Ringe and Taylor 2014}). The development
is therefore regular: \emph{*fléuganą} \textgreater{} \emph{flēogan}
`fly'.

\section{\texorpdfstring{forlorn --- OE
\emph{lēosan}}{forlorn --- OE lēosan}}\label{forlorn-oe-lux113osan}

\index[iv]{02oe@\ivlangheader{Old English}{West Saxon normalization unmarked}!forleosan@\iventry{for-leósan}{}}
\index[iv]{02oe@\ivlangheader{Old English}{West Saxon normalization unmarked}!forleosan@\iventry{for-lēosan}{}}
\index[iv]{02oe@\ivlangheader{Old English}{West Saxon normalization unmarked}!forleosan@\iventry{forlēosan}{}}
\index[iv]{02oe@\ivlangheader{Old English}{West Saxon normalization unmarked}!forloren@\iventry{forloren}{}}
\index[iv]{02oe@\ivlangheader{Old English}{West Saxon normalization unmarked}!forlorn@\iventry{forlorn}{}}
\index[iv]{02oe@\ivlangheader{Old English}{West Saxon normalization unmarked}!leosan@\iventry{lēosan}{}}
\index[iv]{03pgmc@\ivlangheader{Proto-Germanic}{}!leusana@\iventry{*léusaną}{}}
\index[iv]{12goth@\ivlangheader{Gothic}{}!fraliusan@\iventry{fra-liusan}{}}

Derivation: \emph{*léusaną} \textgreater{} \emph{lēosan} (regular).

\subsection{Derivation trace}\label{derivation-trace-25}

Proto input: \emph{*léusaną}

\begingroup
\setlength{\fboxsep}{6pt}

\noindent\fbox{%
\begin{minipage}{0.97\linewidth}
\small
\begin{tabular}{@{}>{\raggedright\arraybackslash}p{0.359\linewidth}@{\hspace{0.030\linewidth}}>{\raggedright\arraybackslash}p{0.611\linewidth}@{}}
\begin{minipage}[t]{\linewidth}
\centering\textbf{Earlier Germanic changes}\par
\vspace{0.35em}
\raggedright
\centering\textbf{Northwest and West Germanic}\par
\raggedright
\vspace{0.2em}
\raggedright [no change]\par
\vspace{0.6em}
\centering\textbf{Early Anglo-Frisian}\par
\raggedright
\vspace{0.2em}
\raggedright [no change]\par
\end{minipage}
&
\begin{minipage}[t]{\linewidth}
\centering\textbf{Old English changes}\par
\vspace{0.35em}
\raggedright
\begin{tabularx}{\linewidth}{@{}>{\raggedright\arraybackslash}X@{\hspace{0.55em}}>{\raggedright\arraybackslash}p{0.18\linewidth}@{\hspace{0.25em}}}
\mbox{OE Diphthong Leveling} & \emph{*lēosaną} \\
OE Heavy Syllable Nasal Apocope & \emph{*lēosan} \\
\mbox{OE Secondary Nasalization} & \emph{*lēosąn} \\
\mbox{OE Weak Tail Reduction} & \emph{*lēosan} \\
\end{tabularx}
\end{minipage}
\\
\end{tabular}
\end{minipage}%
} \endgroup

Old English form: \emph{lēosan}

\subsection{Reconstruction and comparative
evidence}\label{reconstruction-and-comparative-evidence-23}

Kroonen reconstructs the verb under \emph{*leusan-} and cites prefixed
daughters such as Gothic \emph{fra-liusan} `lose' and Old English
\emph{for-lēosan} `lose, forlorn'; Orel likewise gives Old English
\emph{for-leósan} `lose' (\citeproc{ref-Kroonen2013}{Kroonen 2013};
\citeproc{ref-Orel2003}{Orel 2003}). The inherited verbal base is
therefore clear, though the daughter set often appears with the prefix.

\subsection{Old English evidence}\label{old-english-evidence-23}

The direct Old English evidence behind English \emph{forlorn} lies in
the prefixed verb \emph{forlēosan} `lose, forlorn' and especially in the
participle \emph{forloren} `forlorn', recorded by Ringe and Taylor and
in the dictionaries (\citeproc{ref-RingeTaylor2014}{Ringe and Taylor
2014}; \citeproc{ref-ClarkHall1960}{Clark Hall 1960};
\citeproc{ref-BosworthToller1898}{Bosworth and Toller 1898}). The
simplex infinitive \emph{lēosan} `lose' represents the verbal base
itself.

\subsection{Form note}\label{form-note-2}

As a base-form comparison, the simplex infinitive is \emph{lēosan}
`lose', while the English adjective continues the prefixed Old English
family \emph{forlēosan} `lose' / \emph{forloren} `forlorn'
(\citeproc{ref-RingeTaylor2014}{Ringe and Taylor 2014}).

\subsection{Development to Old
English}\label{development-to-old-english-23}

From \Recon{léusaną} `lose', Old English diphthong leveling gives
\Recon{lēosaną} `lose', and later nasal apocope and weak-tail reduction
yield \emph{lēosan} `lose' (\citeproc{ref-RingeTaylor2014}{Ringe and
Taylor 2014}). The prefixed forms follow the same verbal base with added
\emph{for-}.

\section{\texorpdfstring{four --- OE
\emph{fēower}}{four --- OE fēower}}\label{four-oe-fux113ower}

\index[iv]{02oe@\ivlangheader{Old English}{West Saxon normalization unmarked}!feower@\iventry{fēower}{}}
\index[iv]{03pgmc@\ivlangheader{Proto-Germanic}{}!fedwor@\iventry{*fédwōr}{}}

Derivation: \emph{*fédwōr} \textgreater{} \emph{fēower} (regular).

\subsection{Derivation trace}\label{derivation-trace-26}

Proto input: \emph{*fédwōr}

\begingroup
\setlength{\fboxsep}{6pt}

\noindent\fbox{%
\begin{minipage}{0.97\linewidth}
\small
\begin{tabular}{@{}>{\raggedright\arraybackslash}p{0.485\linewidth}@{\hspace{0.030\linewidth}}>{\raggedright\arraybackslash}p{0.485\linewidth}@{}}
\begin{minipage}[t]{\linewidth}
\centering\textbf{Earlier Germanic changes}\par
\vspace{0.35em}
\raggedright
\centering\textbf{Northwest and West Germanic}\par
\raggedright
\vspace{0.2em}
\begin{tabularx}{\linewidth}{@{}>{\raggedright\arraybackslash}X@{\hspace{0.55em}}>{\raggedright\arraybackslash}p{0.14\linewidth}@{\hspace{0.25em}}}
\mbox{PWGmc Final Or Lowering} & \emph{*fédwar} \\
\mbox{PWGmc Coronal W Assimilation} & \emph{*féwwar} \\
\mbox{OE WW Simplification} & \emph{*féuwar} \\
\end{tabularx}
\vspace{0.6em}
\centering\textbf{Early Anglo-Frisian}\par
\raggedright
\vspace{0.2em}
\raggedright [no change]\par
\end{minipage}
&
\begin{minipage}[t]{\linewidth}
\centering\textbf{Old English changes}\par
\vspace{0.35em}
\raggedright
\begin{tabularx}{\linewidth}{@{}>{\raggedright\arraybackslash}X@{\hspace{0.55em}}>{\raggedright\arraybackslash}p{0.14\linewidth}@{\hspace{0.25em}}}
\mbox{OE Diphthong Leveling} & \emph{*fēowar} \\
\mbox{EAF Brightening Unstressed} & \emph{*fēowær} \\
\mbox{OE Unstressed AE Merger} & \emph{*fēower} \\
\end{tabularx}
\end{minipage}
\\
\end{tabular}
\end{minipage}%
} \endgroup

Old English form: \emph{fēower}

\subsection{Reconstruction and
development}\label{reconstruction-and-development-1}

The selected PGmc \emph{*fédwōr} gives OE \emph{fēower} `four'.
Coronal-w assimilation creates the geminate; short \emph{*e} plus
\emph{*ww} is then reanalyzed as an \emph{*eu} nucleus with retained
consonantal \emph{*w}. Ringe and Taylor explicitly use this history
alongside the oblique plural pronoun \emph{ēow} `you'
(\citeproc{ref-RingeTaylor2014}{Ringe and Taylor 2014, 41--42}).

Thus \hyperref[rule-PWGmcCoronalWAssimilation]{SC008
PWGmcCoronalWAssimilation} feeds
\hyperref[rule-OEWWSimplification]{SC031 OEWWSimplification}. The
earlier product reaches English long-diphthong realization through
\hyperref[rule-OEDiphthongLeveling]{SC032 OEDiphthongLeveling}, rather
than receiving an English surface vowel in the West Germanic reanalysis
(\citeproc{ref-RingeTaylor2014}{Ringe and Taylor 2014, 65--66,
171--75}). The selected context is the ordinary strong-final citation
convention, not the weak-final selection of the pronoun. No
reconstruction or final target has changed.

\section{\texorpdfstring{gang --- OE
\emph{gang}}{gang --- OE gang}}\label{gang-oe-gang}

\index[iv]{02oe@\ivlangheader{Old English}{West Saxon normalization unmarked}!gang@\iventry{gang}{}}
\index[iv]{02oe@\ivlangheader{Old English}{West Saxon normalization unmarked}!gangan@\iventry{gangan}{}}
\index[iv]{03pgmc@\ivlangheader{Proto-Germanic}{}!gangaz@\iventry{*gangaz}{}}
\index[iv]{03pgmc@\ivlangheader{Proto-Germanic}{}!gangaz@\iventry{*gángaz}{}}
\index[iv]{09on@\ivlangheader{Old Norse}{}!gangr@\iventry{gangr}{}}

Derivation: \emph{*gángaz} \textgreater{} \emph{gang} (regular).

\subsection{Derivation trace}\label{derivation-trace-27}

Proto input: \emph{*gángaz}

\begingroup
\setlength{\fboxsep}{6pt}

\noindent\fbox{%
\begin{minipage}{0.97\linewidth}
\small
\begin{tabular}{@{}>{\raggedright\arraybackslash}p{0.485\linewidth}@{\hspace{0.030\linewidth}}>{\raggedright\arraybackslash}p{0.485\linewidth}@{}}
\begin{minipage}[t]{\linewidth}
\centering\textbf{Earlier Germanic changes}\par
\vspace{0.35em}
\raggedright
\centering\textbf{Northwest and West Germanic}\par
\raggedright
\vspace{0.2em}
\begin{tabularx}{\linewidth}{@{}>{\raggedright\arraybackslash}X@{\hspace{0.55em}}>{\raggedright\arraybackslash}p{0.14\linewidth}@{\hspace{0.25em}}}
\mbox{EAF Final Z Deletion} & \emph{*gánga} \\
\end{tabularx}
\vspace{0.6em}
\centering\textbf{Early Anglo-Frisian}\par
\raggedright
\vspace{0.2em}
\raggedright [no change]\par
\end{minipage}
&
\begin{minipage}[t]{\linewidth}
\centering\textbf{Old English changes}\par
\vspace{0.35em}
\raggedright
\begin{tabularx}{\linewidth}{@{}>{\raggedright\arraybackslash}X@{\hspace{0.55em}}>{\raggedright\arraybackslash}p{0.14\linewidth}@{\hspace{0.25em}}}
\mbox{PWGmc Final Bare A Loss} & \emph{*gáng} \\
\end{tabularx}
\end{minipage}
\\
\end{tabular}
\end{minipage}%
} \endgroup

Old English form: \emph{gang}

\subsection{Reconstruction and comparative
evidence}\label{reconstruction-and-comparative-evidence-24}

Orel reconstructs the noun as \Recon{gangaz} `gang' and cites Old
English
{\emph{gang}}\index[iv]{02oe@\ivlangheader{Old English}{West Saxon normalization unmarked}!gang@\iventry{gang}{}}
`going, way' beside Old Norse \emph{gangr}, Old Frisian \emph{gang} /
\emph{gong}, Old Saxon
{\emph{gang}}\index[iv]{02oe@\ivlangheader{Old English}{West Saxon normalization unmarked}!gang@\iventry{gang}{}}
`going, way', and Old High German
{\emph{gang}}\index[iv]{02oe@\ivlangheader{Old English}{West Saxon normalization unmarked}!gang@\iventry{gang}{}}
`going, way' (\citeproc{ref-Orel2003}{Orel 2003}). The derivational
input \Recon{gángaz} `gang' is the same lexeme in the accent notation
used here.

\subsection{Old English evidence}\label{old-english-evidence-24}

Clark Hall and Bosworth-Toller both record
{\emph{gang}}\index[iv]{02oe@\ivlangheader{Old English}{West Saxon normalization unmarked}!gang@\iventry{gang}{}}
`going, way' as the noun `going, journey, way', and Bright's glossary
gives \emph{gong (gang), m., path, course}
(\citeproc{ref-ClarkHall1960}{Clark Hall 1960};
\citeproc{ref-BosworthToller1898}{Bosworth and Toller 1898, 159};
\citeproc{ref-BrightCassidyRingler1971}{Bright 1971, 392}). The target
is therefore the attested noun headword itself.

\subsection{Form note}\label{form-note-3}

This entry concerns the noun
{\emph{gang}}\index[iv]{02oe@\ivlangheader{Old English}{West Saxon normalization unmarked}!gang@\iventry{gang}{}}
`going, way', not the separate verb \emph{gangan}
(\citeproc{ref-ClarkHall1960}{Clark Hall 1960};
\citeproc{ref-BosworthToller1898}{Bosworth and Toller 1898, 159}).

\subsection{Development to Old
English}\label{development-to-old-english-24}

From \Recon{gángaz} `gang', loss of final \emph{-z} gives \Recon{gánga}
`gang', and later loss of final bare \emph{-a} yields
{\emph{gang}}\index[iv]{02oe@\ivlangheader{Old English}{West Saxon normalization unmarked}!gang@\iventry{gang}{}}
`going, way'. The development is therefore regular: \emph{*gángaz}
\textgreater{} \emph{gang}.

\section{\texorpdfstring{gift --- OE
\emph{ġift}}{gift --- OE ġift}}\label{gift-oe-ux121ift}

\index[iv]{02oe@\ivlangheader{Old English}{West Saxon normalization unmarked}!gift@\iventry{ġift}{}}
\index[iv]{03pgmc@\ivlangheader{Proto-Germanic}{}!giftiz@\iventry{*gíftiz}{}}

Derivation: \emph{*gíftiz} \textgreater{} \emph{ġift} (regular).

\subsection{Derivation trace}\label{derivation-trace-28}

Proto input: \emph{*gíftiz}

\begingroup
\setlength{\fboxsep}{6pt}

\noindent\fbox{%
\begin{minipage}{0.97\linewidth}
\small
\begin{tabular}{@{}>{\raggedright\arraybackslash}p{0.359\linewidth}@{\hspace{0.030\linewidth}}>{\raggedright\arraybackslash}p{0.611\linewidth}@{}}
\begin{minipage}[t]{\linewidth}
\centering\textbf{Earlier Germanic changes}\par
\vspace{0.35em}
\raggedright
\centering\textbf{Northwest and West Germanic}\par
\raggedright
\vspace{0.2em}
\begin{tabularx}{\linewidth}{@{}>{\raggedright\arraybackslash}X@{\hspace{0.55em}}>{\raggedright\arraybackslash}p{0.14\linewidth}@{\hspace{0.25em}}}
\mbox{EAF Final Z Deletion} & \emph{*gífti} \\
\end{tabularx}
\vspace{0.6em}
\centering\textbf{Early Anglo-Frisian}\par
\raggedright
\vspace{0.2em}
\raggedright [no change]\par
\end{minipage}
&
\begin{minipage}[t]{\linewidth}
\centering\textbf{Old English changes}\par
\vspace{0.35em}
\raggedright
\begin{tabularx}{\linewidth}{@{}>{\raggedright\arraybackslash}X@{\hspace{0.55em}}>{\raggedright\arraybackslash}p{0.14\linewidth}@{\hspace{0.25em}}}
\mbox{OE Velar Palatalization} & \emph{*ʝífti} \\
\mbox{OE High Vowel Apocope} & \emph{*ʝíft} \\
\mbox{OE Palatal Fricative Merger} & \emph{*jíft} \\
\end{tabularx}
\end{minipage}
\\
\end{tabular}
\end{minipage}%
} \endgroup

Old English form: \emph{ġift}

\subsection{A. The selected reconstruction and the question it
answers}\label{a.-the-selected-reconstruction-and-the-question-it-answers}

Old English \emph{ġift} `gift' is a small word with unusually important
chronological consequences. A simple high vowel follows an initial
palatal consonant, whereas \emph{ġiefan} `give' shows the diphthongal
treatment associated with that consonant. To explain the difference we
must ask which vowel was present when ordinary palatal diphthongization
operated, not simply which order of rules can reproduce the final
spellings. Campbell lists the short-i gift headword among his
initial-palatal examples (\citeproc{ref-Campbell1959}{Campbell 1959,
173--74}, §427).

We select PGmc
\Recon{giftiz}\index[iv]{03pgmc@\ivlangheader{Proto-Germanic}{}!giftiz@\iventry{*giftiz}{}}
`gift' following Ringe's explicit reconstruction
(\citeproc{ref-Ringe2017}{Ringe 2017, 135}). Our acute root-stress
convention does not add a different vowel or claim that the source
prints an accent. Both the comparative reconstruction and the selected
computational input use the raised vowel. This is a regular inherited
derivation, not a late analogical substitute input. The decision is
nevertheless not a claim that every etymological source prints the same
reconstruction. A substantial e-form tradition must be explained and
preserved rather than suppressed.

\subsection{B. Orel: the full e-vowel noun
reconstruction}\label{b.-orel-the-full-e-vowel-noun-reconstruction}

Orel's noun entry gives an e-vowel reconstruction, normalized here as
\Recon{geftiz}\index[iv]{03pgmc@\ivlangheader{Proto-Germanic}{}!geftiz@\iventry{*geftiz}{}}
`gift', and connects it with the verb `give'
(\citeproc{ref-Orel2003}{Orel 2003, 130}). It has the full iz ending,
not merely a bare root or a stem ending in a hyphen. His evidence
includes the Old English gift headword and the Old Frisian counterpart
\emph{jeft} `gift, privilege', alongside the wider inherited formation.
The e is plainly printed in the original page: it is not an OCR accident
or a vowel introduced by our transcription. Only the initial velar sign
is normalized to g here; no argument about stop versus fricative history
is drawn from that typographic normalization.

Orel thus provides genuine published support for the older e-input used
in our model. However, the entry does not explain why e is reconstructed
despite the high-vowel reflex, identify a specific later raising date,
or argue against the early raising account. Its derivation from `give'
supplies an etymological relationship; that alone does not demonstrate
that the noun's realized root vowel remained identical to the verb's at
every historical stage. The cited bibliography points to morphology and
verbal etymology treatments, including Bammesberger and Seebold,
considered below (\citeproc{ref-Orel2003}{Orel 2003, 130}). Those
references are evidence of the scholarly tradition behind the entry, not
a license to attribute an unstated phonological argument to Orel.

\subsection{C. Kluge--Seebold: an e-vowel feminine
stem}\label{c.-klugeseebold-an-e-vowel-feminine-stem}

Kluge--Seebold gives
\Recon{gefti-}\index[iv]{03pgmc@\ivlangheader{Proto-Germanic}{}!gefti@\iventry{*gefti-}{}}
`gift', explicitly a Germanic feminine stem, and relates the noun to
`give' (\citeproc{ref-KlugeSeebold2011}{Kluge 2011, 359}). This is
independent confirmation that an e-vowel etymological presentation is
established in major reference works. It is not the same printed object
as Orel's full nominative reconstruction. We must not silently append z
and then present the expanded form as a quotation from Kluge--Seebold.

The entry cites the Old English gift headword and the Old Frisian form
\emph{jeft(e)} `gift' within the inherited family
(\citeproc{ref-KlugeSeebold2011}{Kluge 2011, 359}). Its principal
discussion concerns the word's semantic history; it does not supply a
worked chronology of root-vowel raising versus palatal diphthongization.
Consequently the entry establishes the stem and its etymological
affiliation more directly than it establishes the vowel quality at a
narrowly dated phonological node. This does not prove that the stem
notation is purely abstract. It means that such an interpretation, if
adopted, would require an argument beyond the entry.

\subsection{D. Bammesberger: the e-grade derivational
analysis}\label{d.-bammesberger-the-e-grade-derivational-analysis}

Bammesberger places the formation among ti-stems built on an e-grade
verbal root and prints a segmented e-vowel form
(\citeproc{ref-Bammesberger1990}{Bammesberger 1990, 142}, §5.4.2). His
discussion relates the noun to the Germanic root of `give' and treats
the consonantal sequence ft as regular in the ti-abstract formed from a
labial-final root. On the same page the formation also participates in
his explanation of the structural pattern of other abstracts. The
continuation lists further e-grade formations with Old English i-vowel
reflexes (\citeproc{ref-Bammesberger1990}{Bammesberger 1990, 142--43}).

This is a substantive reason for retaining e in the derivational
analysis: the noun belongs to an e-grade formation, not an arbitrarily
invented high-vowel root. But derivational grade and realized surface
quality after a regular conditioned sound change are different
questions. An e-grade formation can have a later i-reflex without
becoming a zero-grade formation. The passage supports the morphological
analysis; it does not itself decide that the inherited e survived until
Old English mutation or state a competing date for raising. We therefore
retain its e-grade evidence in the deeper explanation without using it
alone to override a separately argued phonological reconstruction
(\citeproc{ref-Bammesberger1990}{Bammesberger 1990, 142--43}).

\subsection{E. Seebold: segmented forms and the limits of
notation}\label{e.-seebold-segmented-forms-and-the-limits-of-notation}

Seebold prints a segmented e-vowel formation with the full nominal
ending and supplies more differentiated Old English evidence than a
single normalized citation (\citeproc{ref-Seebold1970}{Seebold 1970,
218}). He distinguishes \emph{giftu} `wedding', explicitly plural,
\emph{gyft} `marriage', and \emph{feoh-gift} `gift of money'. A separate
homonymous noun associated with `food' belongs to a different formation
and must not be recruited as attestation for the noun treated here.
These distinctions matter: the formation's identity is well supported,
but a plural or compound citation is not automatically evidence for a
particular singular manuscript token.

His methodological introduction also makes clear that the reconstructed
notation is not uniformly a direct transcription of one phonological
moment. Earlier states are reconstructed for selected developments,
including the development of ei to long i
(\citeproc{ref-Seebold1970}{Seebold 1970, 35}). The symbols for stem
classes are modelled on common-Germanic forms without claiming exact
phonetic reproduction (\citeproc{ref-Seebold1970}{Seebold 1970,
35--36}). These explicit cautions make it particularly important not to
read every segment mechanically as an input to a later sound-change
cascade. They do not, however, specifically explain the root e in this
noun: the warning about stem-class endings cannot be generalized into
proof that all reconstructed root vowels are abstract morphological
labels. We leave that interpretive limit visible rather than declaring
the e/i discrepancy merely notational.

\subsection{F. Ringe: the realized PGmc i-form and its phonological
rationale}\label{f.-ringe-the-realized-pgmc-i-form-and-its-phonological-rationale}

Ringe explicitly gives PGmc
\Recon{giftiz}\index[iv]{03pgmc@\ivlangheader{Proto-Germanic}{}!giftiz@\iventry{*giftiz}{}}
`gift' as the nominal formation associated with the verb `give'
(\citeproc{ref-Ringe2017}{Ringe 2017, 135}). His discussion places it in
the broader treatment of productive nominal formation and consonantal
developments. The Old English comparison is plural \emph{ġifta}
`wedding', with a change of stem class. We do not turn that plural into
a claim that Ringe attests our normalized singular target.

The vowel has a separate phonological explanation. Ringe states that e
was raised to i when a high front vocalic followed in the same or
immediately succeeding syllable (\citeproc{ref-Ringe2017}{Ringe 2017,
151--52}). The noun's following i supplies precisely such an
environment. This explains how an etymological e-grade formation can be
represented with i in a realized phonological reconstruction. It is not
a lexical exception, a grammatical condition on nouns, or a repair
licensed only because this one word otherwise produces an unwanted
diphthong.

Ringe is more cautious about the exact ancestral dating than a bare
reconstruction might suggest. The merger of stressed e and i in Gothic
prevents direct demonstration of the raising there. His indirect
argument depends on raising preceding j-loss and subsequent contraction,
particularly the history of the numeral `three'. He calls a PGmc date
the most probable hypothesis, not completely certain
(\citeproc{ref-Ringe2017}{Ringe 2017, 152--53}). That qualification
belongs in our account. It concerns the placement of the change in the
deeper ancestry; it must not be misrepresented as doubt that his noun
entry actually prints an i-vowel reconstruction.

Ringe and Taylor independently insist on separating this inherited
raising from Old English front mutation: the older change occurred
hundreds of years earlier, and alleged later repetitions are rare and
doubtful (\citeproc{ref-RingeTaylor2014}{Ringe and Taylor 2014, 220}).
Their discussion therefore supplies an important chronological check
against explaining this noun's inherited high vowel by the general Old
English mutation rule.

\subsection{G. Why we choose i, and what changes in the
derivation}\label{g.-why-we-choose-i-and-what-changes-in-the-derivation}

For the input to a PGmc-to-OE cascade we prefer the explicitly staged
phonological reconstruction, rather than carrying an earlier
etymological e-grade representation forward until a convenient later
rule can correct it. Ringe's noun and regular raising account together
supply that reconstruction (\citeproc{ref-Ringe2017}{Ringe 2017, 135},
pp.~151--153). The competing reference forms remain part of the
evidence: we neither count dictionary spellings as votes nor conclude
that their authors were wrong merely because they did not give a worked
chronology.

With earlier i already present, the noun undergoes the inherited
consonantal and ending developments, then initial palatalization. The
singleton initial is a palatal fricative, distinct from gg/ng stops; the
adopted model keeps it distinct from inherited j during mutation
(\citeproc{ref-Fulk2018}{Fulk 2018, 130--32};
\citeproc{ref-RingeTaylor2014}{Ringe and Taylor 2014, 204}). Ordinary
palatal diphthongization does not change its high vowel, and Old English
mutation need not raise an e that is no longer there. Loss of the final
high vowel and surface spelling yield \emph{ġift} `gift'. Our earlier
implementation instead allowed a selected e-input to reach late mutation
and used that correction to justify putting all palatal diphthongization
afterward. The right final spelling concealed the wrongly conflated
histories.

The independent ordinary chronology is breaking, then palatal
diphthongization, then mutation (\citeproc{ref-Fulk2018}{Fulk 2018, 74};
\citeproc{ref-RingeTaylor2014}{Ringe and Taylor 2014, 215--17, 222}).
The worked contrast is \emph{ġiest} `guest': its low-front vowel first
becomes ea after the palatal, and ea then becomes ie under mutation.
Gift's earlier i is a negative control on ordinary diphthongization, not
an independent proof of that relative order. The later treatment of
\emph{sċēaþ} `sheath' is separately retained after mutation; its
mutation-derived long vowel does not require ordinary diphthongization
as a whole to be late (\citeproc{ref-Campbell1959}{Campbell 1959,
68--69}; \citeproc{ref-RingeTaylor2014}{Ringe and Taylor 2014, 235}).
The complete conditioning of that later treatment remains unresolved.

\subsection{H. Confidence and the remaining
alternatives}\label{h.-confidence-and-the-remaining-alternatives}

Our confidence is strongest in the distinction between inherited raising
and later OE mutation, and in the independently supported ordinary
diphthongization chronology. The choice of Ringe's i-input is well
motivated as an explicit working phonological reconstruction, but is not
advertised as unanimous: Orel's e-form, Kluge--Seebold's e-stem and the
e-grade derivational presentations remain genuine alternatives
(\citeproc{ref-Orel2003}{Orel 2003, 130};
\citeproc{ref-KlugeSeebold2011}{Kluge 2011, 359};
\citeproc{ref-Bammesberger1990}{Bammesberger 1990, 142--43};
\citeproc{ref-Seebold1970}{Seebold 1970, 218}). We have not established
that every one of those forms denotes exactly the same historical
phonological checkpoint as Ringe's input.

Confidence in the precise PGmc date is qualified by Ringe's own indirect
evidence (\citeproc{ref-Ringe2017}{Ringe 2017, 152--53}). Confidence in
our normalized target as a headword is distinct from certification of a
diplomatic paradigm cell (\citeproc{ref-Campbell1959}{Campbell 1959,
173--74}). Variant spellings alone do not establish a different
exceptionless sound law, lexical diffusion, or an analogical history;
those explanations would need their own evidence. We make no such new
claim here.

The computational comparison is deliberately secondary. The adopted
input and ordinary chronology preserve every previously generated final
form, while realizing the source sequence for guest and leaving sheath's
derivation unchanged. That demonstrates compatibility of the selected
account with the present corpus, not its historical truth by output
score. If stronger evidence eventually requires an e-vowel at the input
node, the necessary earlier regular raising must be specified and tested
as such; it must not be silently delegated again to late OE mutation.

\section{\texorpdfstring{give --- OE
\emph{ġiefan}}{give --- OE ġiefan}}\label{give-oe-ux121iefan}

\index[iv]{02oe@\ivlangheader{Old English}{West Saxon normalization unmarked}!forgeofan~05@\iventry{for-geofan}{Merc.}}
\index[iv]{02oe@\ivlangheader{Old English}{West Saxon normalization unmarked}!forgeofan@\iventry{for-geofan}{}}
\index[iv]{02oe@\ivlangheader{Old English}{West Saxon normalization unmarked}!geafa~06@\iventry{geafa}{North.}}
\index[iv]{02oe@\ivlangheader{Old English}{West Saxon normalization unmarked}!geafa@\iventry{geafa}{}}
\index[iv]{02oe@\ivlangheader{Old English}{West Saxon normalization unmarked}!giefan@\iventry{ġiefan}{}}
\index[iv]{03pgmc@\ivlangheader{Proto-Germanic}{}!gebana@\iventry{*gébaną}{}}

Derivation: \emph{*gébaną} \textgreater{} \emph{ġiefan} (regular).

\subsection{Derivation trace}\label{derivation-trace-29}

Proto input: \emph{*gébaną}

\begingroup
\setlength{\fboxsep}{6pt}

\noindent\fbox{%
\begin{minipage}{0.97\linewidth}
\small
\begin{tabular}{@{}>{\raggedright\arraybackslash}p{0.339\linewidth}@{\hspace{0.030\linewidth}}>{\raggedright\arraybackslash}p{0.630\linewidth}@{}}
\begin{minipage}[t]{\linewidth}
\centering\textbf{Earlier Germanic changes}\par
\vspace{0.35em}
\raggedright
\centering\textbf{Northwest and West Germanic}\par
\raggedright
\vspace{0.2em}
\raggedright [no change]\par
\vspace{0.6em}
\centering\textbf{Early Anglo-Frisian}\par
\raggedright
\vspace{0.2em}
\raggedright [no change]\par
\end{minipage}
&
\begin{minipage}[t]{\linewidth}
\centering\textbf{Old English changes}\par
\vspace{0.35em}
\raggedright
\begin{tabularx}{\linewidth}{@{}>{\raggedright\arraybackslash}X@{\hspace{0.55em}}>{\raggedright\arraybackslash}p{0.14\linewidth}@{\hspace{0.25em}}}
OE Heavy Syllable Nasal Apocope & \emph{*géban} \\
\mbox{OE Secondary Nasalization} & \emph{*gébąn} \\
\mbox{OE Velar Palatalization} & \emph{*ʝébąn} \\
OE Ws Palatal Diphthongization & \emph{*ʝíebąn} \\
\mbox{PGmc B Allophony} & \emph{*ʝíeβąn} \\
\mbox{OE Weak Tail Reduction} & \emph{*ʝíeβan} \\
\mbox{OE Palatal Fricative Merger} & \emph{*jíeβan} \\
\end{tabularx}
\end{minipage}
\\
\end{tabular}
\end{minipage}%
} \endgroup

Old English form: \emph{ġiefan}

\subsection{Reconstruction and comparative
evidence}\label{reconstruction-and-comparative-evidence-25}

Kroonen reconstructs the strong verb as \emph{*geban-} and cites Old
English
{\emph{giefan}}\index[iv]{02oe@\ivlangheader{Old English}{West Saxon normalization unmarked}!giefan@\iventry{giefan}{}}
`give' among its reflexes (\citeproc{ref-Kroonen2013}{Kroonen 2013}).
Ringe and Taylor contrast West Saxon \emph{giefan} `give' with Mercian
\emph{for-geofan} `give' and Northumbrian \emph{geafa} `give', showing
that the inherited verb takes different later dialectal shapes
(\citeproc{ref-RingeTaylor2014}{Ringe and Taylor 2014}).

\subsection{Old English evidence}\label{old-english-evidence-25}

Campbell gives \emph{gefan} (W-S \emph{giefan} `give') among examples of
initial palatalization, and Clark Hall records the verb under plain
\emph{giefan} `give' with forms such as \emph{geaf} and \emph{giefen}
(\citeproc{ref-Campbell1959}{Campbell 1959};
\citeproc{ref-ClarkHall1960}{Clark Hall 1960}). I write \emph{ġiefan}
`give' for the palatal initial.

\subsection{Dialect note}\label{dialect-note}

West Saxon \emph{ie} here reflects palatal diphthongization after
initial palatalization; non-West-Saxon forms such as \emph{geafa} `give'
or \emph{for-geofan} `give' continue the same verb without the West
Saxon vocalism (\citeproc{ref-RingeTaylor2014}{Ringe and Taylor 2014}).

\subsection{Development to Old
English}\label{development-to-old-english-25}

The singleton initial belongs to the fricative class, not the gg/ng stop
class. The adopted consonantal account distinguishes its palatal ʝ from
inherited j during mutation and merges them afterwards
(\citeproc{ref-Fulk2018}{Fulk 2018, 130--32};
\citeproc{ref-RingeTaylor2014}{Ringe and Taylor 2014, 204}).

From \Recon{gébaną} `give', initial \emph{g} palatalizes before
\emph{e}; West Saxon palatal diphthongization then yields \emph{ie}, and
later tail reduction gives
{\emph{giefan}}\index[iv]{02oe@\ivlangheader{Old English}{West Saxon normalization unmarked}!giefan@\iventry{giefan}{}}
`give' (\citeproc{ref-Campbell1959}{Campbell 1959};
\citeproc{ref-RingeTaylor2014}{Ringe and Taylor 2014}). The result is
therefore the regular West Saxon infinitive.

\section{\texorpdfstring{gold --- OE
\emph{gold}}{gold --- OE gold}}\label{gold-oe-gold}

\index[iv]{02oe@\ivlangheader{Old English}{West Saxon normalization unmarked}!gold@\iventry{gold}{}}
\index[iv]{03pgmc@\ivlangheader{Proto-Germanic}{}!gultha@\iventry{*gúlθą}{}}

Derivation: \emph{*gúlθą} \textgreater{} \emph{gold} (regular).

\subsection{Derivation trace}\label{derivation-trace-30}

Proto input: \emph{*gúlθą}

\begingroup
\setlength{\fboxsep}{6pt}

\noindent\fbox{%
\begin{minipage}{0.97\linewidth}
\small
\begin{tabular}{@{}>{\raggedright\arraybackslash}p{0.485\linewidth}@{\hspace{0.030\linewidth}}>{\raggedright\arraybackslash}p{0.485\linewidth}@{}}
\begin{minipage}[t]{\linewidth}
\centering\textbf{Earlier Germanic changes}\par
\vspace{0.35em}
\raggedright
\centering\textbf{Northwest and West Germanic}\par
\raggedright
\vspace{0.2em}
\begin{tabularx}{\linewidth}{@{}>{\raggedright\arraybackslash}X@{\hspace{0.55em}}>{\raggedright\arraybackslash}p{0.14\linewidth}@{\hspace{0.25em}}}
\mbox{EAFL Th Voicing} & \emph{*gúldą} \\
\mbox{PNWGmc U Lowering} & \emph{*góldą} \\
\end{tabularx}
\vspace{0.6em}
\centering\textbf{Early Anglo-Frisian}\par
\raggedright
\vspace{0.2em}
\raggedright [no change]\par
\end{minipage}
&
\begin{minipage}[t]{\linewidth}
\centering\textbf{Old English changes}\par
\vspace{0.35em}
\raggedright
\begin{tabularx}{\linewidth}{@{}>{\raggedright\arraybackslash}X@{\hspace{0.55em}}>{\raggedright\arraybackslash}p{0.14\linewidth}@{\hspace{0.25em}}}
OE Heavy Syllable Nasal Apocope & \emph{*góld} \\
\end{tabularx}
\end{minipage}
\\
\end{tabular}
\end{minipage}%
} \endgroup

Old English form: \emph{gold}

\subsection{Reconstruction and comparative
evidence}\label{reconstruction-and-comparative-evidence-26}

Ringe and Taylor cite the noun as \emph{*gulþa-} / \emph{*gulda-}, and
Kroonen gives the same pair (\citeproc{ref-RingeTaylor2014}{Ringe and
Taylor 2014, 42}; \citeproc{ref-Kroonen2013}{Kroonen 2013}). The
derivational input \Recon{gúlθą} `gold' preserves the older consonantal
form while leaving open whether the medial stop reflects inherited
alternation or regular West Germanic development.

\subsection{Old English evidence}\label{old-english-evidence-26}

Bosworth-Toller and Clark Hall both record \emph{gold} as the ordinary
Old English neuter noun (\citeproc{ref-BosworthToller1898}{Bosworth and
Toller 1898, 121}; \citeproc{ref-ClarkHall1960}{Clark Hall 1960, 152}).
The target is therefore the attested citation form itself.

\subsection{Development note}\label{development-note-1}

Ringe and Taylor note that the medial stop can be understood either as
alternation \emph{*gulþa-} / \emph{*gulda-} or as the ordinary West
Germanic change \emph{*lþ} \textgreater{} \emph{ld}; both routes lead to
the same Old English consonantism (\citeproc{ref-RingeTaylor2014}{Ringe
and Taylor 2014, 42}).

\subsection{Development to Old
English}\label{development-to-old-english-26}

From \Recon{gúlθą} `gold', the regular consonant development gives
\Recon{gúldą} `gold'; Northwest Germanic / Old English lowering then
yields \Recon{góldą} `gold', and apocope gives \emph{gold}
(\citeproc{ref-Campbell1959}{Campbell 1959};
\citeproc{ref-RingeTaylor2014}{Ringe and Taylor 2014, 42}).

\section{\texorpdfstring{goose --- OE
\emph{gōs}}{goose --- OE gōs}}\label{goose-oe-gux14ds}

\index[iv]{02oe@\ivlangheader{Old English}{West Saxon normalization unmarked}!gos@\iventry{gōs}{}}
\index[iv]{03pgmc@\ivlangheader{Proto-Germanic}{}!gansz@\iventry{*gánsz}{}}

Derivation: \emph{*gánsz} \textgreater{} \emph{gōs} (regular).

\subsection{Derivation trace}\label{derivation-trace-31}

Proto input: \emph{*gánsz}

\begingroup
\setlength{\fboxsep}{6pt}

\noindent\fbox{%
\begin{minipage}{0.97\linewidth}
\small
\begin{tabular}{@{}>{\raggedright\arraybackslash}p{0.611\linewidth}@{\hspace{0.030\linewidth}}>{\raggedright\arraybackslash}p{0.359\linewidth}@{}}
\begin{minipage}[t]{\linewidth}
\centering\textbf{Earlier Germanic changes}\par
\vspace{0.35em}
\raggedright
\centering\textbf{Northwest and West Germanic}\par
\raggedright
\vspace{0.2em}
\begin{tabularx}{\linewidth}{@{}>{\raggedright\arraybackslash}X@{\hspace{0.55em}}>{\raggedright\arraybackslash}p{0.14\linewidth}@{\hspace{0.25em}}}
EAF Nasal Spirant Lengthening & \emph{*gãns} \\
\mbox{EAF Nasal Spirant Loss} & \emph{*gãs} \\
\mbox{EAF Nasalized Low Rounding} & \emph{*gōs} \\
\end{tabularx}
\vspace{0.6em}
\centering\textbf{Early Anglo-Frisian}\par
\raggedright
\vspace{0.2em}
\raggedright [no change]\par
\end{minipage}
&
\begin{minipage}[t]{\linewidth}
\centering\textbf{Old English changes}\par
\vspace{0.35em}
\raggedright
\raggedright [no change]\par
\end{minipage}
\\
\end{tabular}
\end{minipage}%
} \endgroup

Old English form: \emph{gōs}

\subsection{Reconstruction and comparative
evidence}\label{reconstruction-and-comparative-evidence-27}

The `goose' word is a root noun whose nominative singular has an
unusually layered history. Orel's citation form is \emph{*ǥansz}
`goose', with his class-wide explicit nominative marker \emph{*-z}
stacked on the root-final \emph{s} (\citeproc{ref-Orel2003}{Orel 2003,
126}). Kroonen reconstructs an endingless nominative
{\emph{*gans}}\index[iv]{03pgmc@\ivlangheader{Proto-Germanic}{}!gans@\iventry{*gans}{}}
`goose' with Verner-voiced \emph{z} only in the oblique cells (genitive
{\emph{*gunzaz}}\index[iv]{03pgmc@\ivlangheader{Proto-Germanic}{}!gunzaz@\iventry{*gunzaz}{}}
`goose (gen.)') (\citeproc{ref-Kroonen2013}{Kroonen 2013, 168--69}), and
Kluge/Seebold treat the German cognate under the same root noun
(\citeproc{ref-KlugeSeebold2011}{Kluge 2011, 331}). Bammesberger
supplies the deeper layer: by Szemerényi's law the Proto-Indo-European
nominative \emph{*ǵʰanss} yielded \emph{*ǵʰān}, after which the
Proto-Germanic nominative was rebuilt with a final voiced sibilant,
\emph{*ganz}, reanalyzed within the paradigm
(\citeproc{ref-Bammesberger1990}{Bammesberger 1990, 196}). A mainstream
specialist thus reconstructs a Proto-Germanic `goose' nominative with a
final voiced sibilant; Orel's notation is not idiosyncratic, and
Kroonen's endingless citation reflects a different analytical layer of
the same paradigm rather than a contradictory fact.

The derivational input
\Recon{gánsz}\index[iv]{03pgmc@\ivlangheader{Proto-Germanic}{}!gansz@\iventry{*gánsz}{}}
`goose' follows Orel's morphologically explicit convention, keeping the
nominative marker visible so that its historical elimination is modelled
rather than presupposed.

\subsection{Old English evidence}\label{old-english-evidence-27}

Old English
{\emph{gōs}}\index[iv]{02oe@\ivlangheader{Old English}{West Saxon normalization unmarked}!gos@\iventry{gōs}{}}
`goose' is the endingless nominative-accusative singular, with mutated
plural \emph{ġēs} `geese' (\citeproc{ref-Fulk2018}{Fulk 2018, 165--66}).
The comparative picture matches the rest of the root-noun class: Gothic
replaced the word (its `goose' is the u-stem remodeling behind Spanish
\emph{ganso} `goose' (\citeproc{ref-Orel2003}{Orel 2003, 126})), and no
West Germanic daughter shows any nominative ending
(\citeproc{ref-RingeTaylor2014}{Ringe and Taylor 2014, 118}).

\subsection{Development to Old
English}\label{development-to-old-english-27}

The root-noun nominative marker was eliminated before Proto-West
Germanic (\citeproc{ref-RingeTaylor2014}{Ringe and Taylor 2014, 118}),
giving \Recon{gáns} `goose' from the derivational input \Recon{gánsz}
`goose'. The remaining development is the familiar Ingvaeonic
nasal-spirant sequence: the vowel lengthens and rounds before the nasal
+ voiceless fricative cluster and the nasal is then lost, giving
\Recon{gōs} `goose', which is the attested Old English form. The early
loss of the nominative marker is historically distinct from the later
West Germanic loss of final \emph{*z} in unstressed syllables and from
the northern monosyllabic loss with compensatory lengthening; had the
marker survived to the later monosyllabic change, a lengthened vowel
before \emph{s} would not be the expected outcome of this shape, since
the root-noun ending was already gone.

\section{\texorpdfstring{grave --- OE
\emph{grafan}}{grave --- OE grafan}}\label{grave-oe-grafan}

\index[iv]{02oe@\ivlangheader{Old English}{West Saxon normalization unmarked}!graef@\iventry{græf}{}}
\index[iv]{02oe@\ivlangheader{Old English}{West Saxon normalization unmarked}!graefd@\iventry{græfð}{}}
\index[iv]{02oe@\ivlangheader{Old English}{West Saxon normalization unmarked}!graefen@\iventry{græfen}{}}
\index[iv]{02oe@\ivlangheader{Old English}{West Saxon normalization unmarked}!grafan@\iventry{grafan}{}}
\index[iv]{03pgmc@\ivlangheader{Proto-Germanic}{}!grabana@\iventry{*grábaną}{}}

Derivation: \emph{*grábaną} \textgreater{} \emph{grafan} (regular).

\subsection{Derivation trace}\label{derivation-trace-32}

Proto input: \emph{*grábaną}

\begingroup
\setlength{\fboxsep}{6pt}

\noindent\fbox{%
\begin{minipage}{0.97\linewidth}
\small
\begin{tabular}{@{}>{\raggedright\arraybackslash}p{0.339\linewidth}@{\hspace{0.030\linewidth}}>{\raggedright\arraybackslash}p{0.630\linewidth}@{}}
\begin{minipage}[t]{\linewidth}
\centering\textbf{Earlier Germanic changes}\par
\vspace{0.35em}
\raggedright
\centering\textbf{Northwest and West Germanic}\par
\raggedright
\vspace{0.2em}
\raggedright [no change]\par
\vspace{0.6em}
\centering\textbf{Early Anglo-Frisian}\par
\raggedright
\vspace{0.2em}
\raggedright [no change]\par
\end{minipage}
&
\begin{minipage}[t]{\linewidth}
\centering\textbf{Old English changes}\par
\vspace{0.35em}
\raggedright
\begin{tabularx}{\linewidth}{@{}>{\raggedright\arraybackslash}X@{\hspace{0.55em}}>{\raggedright\arraybackslash}p{0.18\linewidth}@{\hspace{0.25em}}}
\mbox{EAF Brightening} & \emph{*græbaną} \\
\mbox{OE A Restoration} & \emph{*grabaną} \\
OE Heavy Syllable Nasal Apocope & \emph{*graban} \\
\mbox{OE Secondary Nasalization} & \emph{*grabąn} \\
\mbox{PGmc B Allophony} & \emph{*graβąn} \\
\mbox{OE Weak Tail Reduction} & \emph{*graβan} \\
\end{tabularx}
\end{minipage}
\\
\end{tabular}
\end{minipage}%
} \endgroup

Old English form: \emph{grafan}

\subsection{Reconstruction and comparative
evidence}\label{reconstruction-and-comparative-evidence-28}

Campbell gives
{\emph{grafan}}\index[iv]{02oe@\ivlangheader{Old English}{West Saxon normalization unmarked}!grafan@\iventry{grafan}{}}
`dig, grave' among the standard examples of Old English a-restoration
before a single consonant and a following back vowel, and Ringe and
Taylor describe the same development for Class VI infinitives
(\citeproc{ref-Campbell1959}{Campbell 1959, 61};
\citeproc{ref-RingeTaylor2014}{Ringe and Taylor 2014}).

\subsection{Old English evidence}\label{old-english-evidence-28}

Clark Hall records
{\emph{grafan}}\index[iv]{02oe@\ivlangheader{Old English}{West Saxon normalization unmarked}!grafan@\iventry{grafan}{}}
`dig, grave' as the verb `to dig, grave' and separately records noun
\emph{græf} `grave, trench' (\citeproc{ref-ClarkHall1960}{Clark Hall
1960}). The target here is the attested infinitive headword of the verb.

\subsection{Development to Old
English}\label{development-to-old-english-28}

From \Recon{grábaną} `grave', Anglo-Frisian brightening first gives a
fronted stem vowel. A-restoration then returns \emph{a} before single
\emph{b} plus the back-vocalic infinitive ending, and later apocope and
weak-tail reduction yield
{\emph{grafan}}\index[iv]{02oe@\ivlangheader{Old English}{West Saxon normalization unmarked}!grafan@\iventry{grafan}{}}
`dig, grave' (\citeproc{ref-Campbell1959}{Campbell 1959, 61};
\citeproc{ref-RingeTaylor2014}{Ringe and Taylor 2014}).

\subsection{Form note}\label{form-note-4}

Noun \emph{græf} `grave' and verbal forms such as \emph{græfð} `graves'
or past participial \emph{græfen} `graven' belong to other lexical or
paradigm positions and do not replace the infinitive \emph{grafan}
`grave' as the target here (\citeproc{ref-ClarkHall1960}{Clark Hall
1960}).

\section{\texorpdfstring{guest --- OE
\emph{ġiest}}{guest --- OE ġiest}}\label{guest-oe-ux121iest}

\index[iv]{02oe@\ivlangheader{Old English}{West Saxon normalization unmarked}!giest@\iventry{ġiest}{}}
\index[iv]{03pgmc@\ivlangheader{Proto-Germanic}{}!gastiz@\iventry{*gástiz}{}}

Derivation: \emph{*gástiz} \textgreater{} \emph{ġiest} (regular).

\subsection{Derivation trace}\label{derivation-trace-33}

Proto input: \emph{*gástiz}

\begingroup
\setlength{\fboxsep}{6pt}

\noindent\fbox{%
\begin{minipage}{0.97\linewidth}
\small
\begin{tabular}{@{}>{\raggedright\arraybackslash}p{0.339\linewidth}@{\hspace{0.030\linewidth}}>{\raggedright\arraybackslash}p{0.630\linewidth}@{}}
\begin{minipage}[t]{\linewidth}
\centering\textbf{Earlier Germanic changes}\par
\vspace{0.35em}
\raggedright
\centering\textbf{Northwest and West Germanic}\par
\raggedright
\vspace{0.2em}
\begin{tabularx}{\linewidth}{@{}>{\raggedright\arraybackslash}X@{\hspace{0.55em}}>{\raggedright\arraybackslash}p{0.14\linewidth}@{\hspace{0.25em}}}
\mbox{EAF Final Z Deletion} & \emph{*gásti} \\
\end{tabularx}
\vspace{0.6em}
\centering\textbf{Early Anglo-Frisian}\par
\raggedright
\vspace{0.2em}
\raggedright [no change]\par
\end{minipage}
&
\begin{minipage}[t]{\linewidth}
\centering\textbf{Old English changes}\par
\vspace{0.35em}
\raggedright
\begin{tabularx}{\linewidth}{@{}>{\raggedright\arraybackslash}X@{\hspace{0.55em}}>{\raggedright\arraybackslash}p{0.14\linewidth}@{\hspace{0.25em}}}
\mbox{EAF Brightening} & \emph{*gæsti} \\
\mbox{OE Velar Palatalization} & \emph{*ʝæsti} \\
OE Ws Palatal Diphthongization & \emph{*ʝeasti} \\
\mbox{OE I Umlaut} & \emph{*ʝiesti} \\
\mbox{OE High Vowel Apocope} & \emph{*ʝiest} \\
\mbox{OE Palatal Fricative Merger} & \emph{*jiest} \\
\end{tabularx}
\end{minipage}
\\
\end{tabular}
\end{minipage}%
} \endgroup

Old English form: \emph{ġiest}

\subsection{Reconstruction and comparative
evidence}\label{reconstruction-and-comparative-evidence-29}

Campbell and Ringe-Taylor treat the noun as an ordinary i-stem whose
West Saxon development shows palatal diphthongization, while
non-West-Saxon evidence preserves forms of the \emph{gest} `guest' type
(\citeproc{ref-Campbell1959}{Campbell 1959};
\citeproc{ref-RingeTaylor2014}{Ringe and Taylor 2014}).

\subsection{Old English evidence}\label{old-english-evidence-29}

Bosworth-Toller and Clark Hall record the word under forms such as
\emph{gist} `guest', \emph{gest} `guest', \emph{giest} `guest', and
\emph{gyst} `guest' (\citeproc{ref-BosworthToller1898}{Bosworth and
Toller 1898}; \citeproc{ref-ClarkHall1960}{Clark Hall 1960}). The Old
English form here, \emph{ġiest}, `guest' is the normalized West Saxon
form within that attested family.

\subsection{Development to Old
English}\label{development-to-old-english-29}

From \Recon{gástiz} `guest', ordinary English fronting supplies the
low-front vowel. Initial singleton g palatalizes as a fricative, not an
early affricate (\citeproc{ref-Fulk2018}{Fulk 2018, 130--32}). Ordinary
West Saxon palatal diphthongization changes æ to ea before mutation
changes ea to ie. The high-vocalic ending is lost afterwards; the
regular outcome is \emph{ġiest} `guest'
(\citeproc{ref-RingeTaylor2014}{Ringe and Taylor 2014, 215--17, 222,
287}).

\subsection{Dialect note}\label{dialect-note-1}

West Saxon \emph{ġiest} `guest' is the Old English form here. Anglian
\emph{gest} `guest' and related spellings remain real Old English
comparators rather than corrections to that choice
(\citeproc{ref-RingeTaylor2014}{Ringe and Taylor 2014};
\citeproc{ref-BosworthToller1898}{Bosworth and Toller 1898}).

\section{\texorpdfstring{hair --- OE
\emph{hǣr}}{hair --- OE hǣr}}\label{hair-oe-hux1e3r}

\index[iv]{02oe@\ivlangheader{Old English}{West Saxon normalization unmarked}!haer@\iventry{hǣr}{}}
\index[iv]{03pgmc@\ivlangheader{Proto-Germanic}{}!xera@\iventry{*xḗrą}{}}

Derivation: \emph{*xḗrą} \textgreater{} \emph{hǣr} (regular).

\subsection{Derivation trace}\label{derivation-trace-34}

Proto input: \emph{*xḗrą}

\begingroup
\setlength{\fboxsep}{6pt}

\noindent\fbox{%
\begin{minipage}{0.97\linewidth}
\small
\begin{tabular}{@{}>{\raggedright\arraybackslash}p{0.485\linewidth}@{\hspace{0.030\linewidth}}>{\raggedright\arraybackslash}p{0.485\linewidth}@{}}
\begin{minipage}[t]{\linewidth}
\centering\textbf{Earlier Germanic changes}\par
\vspace{0.35em}
\raggedright
\centering\textbf{Northwest and West Germanic}\par
\raggedright
\vspace{0.2em}
\begin{tabularx}{\linewidth}{@{}>{\raggedright\arraybackslash}X@{\hspace{0.55em}}>{\raggedright\arraybackslash}p{0.14\linewidth}@{\hspace{0.25em}}}
\mbox{PNWGmc Long E Lowering} & \emph{*xārą} \\
\mbox{EAF Long A Fronting} & \emph{*xǣrą} \\
\end{tabularx}
\vspace{0.6em}
\centering\textbf{Early Anglo-Frisian}\par
\raggedright
\vspace{0.2em}
\raggedright [no change]\par
\end{minipage}
&
\begin{minipage}[t]{\linewidth}
\centering\textbf{Old English changes}\par
\vspace{0.35em}
\raggedright
\begin{tabularx}{\linewidth}{@{}>{\raggedright\arraybackslash}X@{\hspace{0.55em}}>{\raggedright\arraybackslash}p{0.14\linewidth}@{\hspace{0.25em}}}
OE Velar Fricative Palatalization & \emph{*çǣrą} \\
OE Heavy Syllable Nasal Apocope & \emph{*çǣr} \\
\end{tabularx}
\end{minipage}
\\
\end{tabular}
\end{minipage}%
} \endgroup

Old English form: \emph{hǣr}

\subsection{Reconstruction and comparative
evidence}\label{reconstruction-and-comparative-evidence-30}

Kroonen cites the ordinary Proto-Germanic hair word as \emph{*hēra-}
(\citeproc{ref-Kroonen2013}{Kroonen 2013}). The derivational input
\Recon{xḗrą} `hair' represents that same long-ē stem in the present
derivation.

\subsection{Old English evidence}\label{old-english-evidence-30}

Clark Hall and Bosworth-Toller record \emph{hær} `hair' / \emph{hǣr}
`hair' as the ordinary Old English noun
(\citeproc{ref-ClarkHall1960}{Clark Hall 1960, 158};
\citeproc{ref-BosworthToller1898}{Bosworth and Toller 1898, 510}). The
target is therefore the attested headword itself.

\subsection{Development to Old
English}\label{development-to-old-english-30}

From \Recon{xḗrą} `hair', Northwest Germanic lowering gives a long front
vowel, and later loss of the final nasal leaves the Old English form
\emph{hǣr} `hair'. The development treated here is straightforward and
does not require any special paradigm choice.

\subsection{Form note}\label{form-note-5}

Older references to \Recon{xazwăz} `hair' belong to a different lexeme,
and the separate \emph{haddr} / \emph{heordan} / \emph{hād-} material
does not displace the ordinary simplex \emph{hǣr} `hair' treated here
(\citeproc{ref-Kroonen2013}{Kroonen 2013};
\citeproc{ref-ClarkHall1960}{Clark Hall 1960, 158}).

\section{\texorpdfstring{harvest --- OE
\emph{hierfest}}{harvest --- OE hierfest}}\label{harvest-oe-hierfest}

\index[iv]{02oe@\ivlangheader{Old English}{West Saxon normalization unmarked}!hierfest@\iventry{hierfest}{}}
\index[iv]{03pgmc@\ivlangheader{Proto-Germanic}{}!xarbistuz@\iventry{*xárbistuz}{}}

Derivation: \emph{*xárbistuz} \textgreater{} \emph{hierfest} (regular).

\subsection{Derivation trace}\label{derivation-trace-35}

Proto input: \emph{*xárbistuz}

\begingroup
\setlength{\fboxsep}{6pt}

\noindent\fbox{%
\begin{minipage}{0.97\linewidth}
\footnotesize
\begin{tabular}{@{}>{\raggedright\arraybackslash}p{0.339\linewidth}@{\hspace{0.030\linewidth}}>{\raggedright\arraybackslash}p{0.630\linewidth}@{}}
\begin{minipage}[t]{\linewidth}
\centering\textbf{Earlier Germanic changes}\par
\vspace{0.35em}
\raggedright
\centering\textbf{Northwest and West Germanic}\par
\raggedright
\vspace{0.2em}
\begin{tabularx}{\linewidth}{@{}>{\raggedright\arraybackslash}X@{\hspace{0.55em}}>{\raggedright\arraybackslash}p{0.24\linewidth}@{\hspace{0.25em}}}
\mbox{EAF Final Z Deletion} & \emph{*xárbistu} \\
\end{tabularx}
\vspace{0.6em}
\centering\textbf{Early Anglo-Frisian}\par
\raggedright
\vspace{0.2em}
\raggedright [no change]\par
\end{minipage}
&
\begin{minipage}[t]{\linewidth}
\centering\textbf{Old English changes}\par
\vspace{0.35em}
\raggedright
\begin{tabularx}{\linewidth}{@{}>{\raggedright\arraybackslash}X@{\hspace{0.55em}}>{\raggedright\arraybackslash}p{0.24\linewidth}@{\hspace{0.25em}}}
\mbox{EAF Brightening} & \emph{*xærbistu} \\
\mbox{OE Breaking} & \emph{*xearbistu} \\
OE Velar Fricative Palatalization & \emph{*çearbistu} \\
\mbox{OE I Umlaut} & \emph{*çierbistu} \\
\mbox{PGmc B Allophony} & \emph{*çierβistu} \\
\mbox{OE High Vowel Apocope} & \emph{*çierβist} \\
OE Med Unstressed I Lowering1 & \emph{*çierβest} \\
\end{tabularx}
\end{minipage}
\\
\end{tabular}
\end{minipage}%
} \endgroup

Old English form: \emph{hierfest}

\subsection{Reconstruction and comparative
evidence}\label{reconstruction-and-comparative-evidence-31}

Bammesberger and Ringe-Taylor treat \emph{*harbist-} as the inherited
base and explain that the regular native West Saxon development would be
of the \emph{hierfest} / \emph{hyrfest} type
(\citeproc{ref-Bammesberger1997}{Bammesberger 1997};
\citeproc{ref-RingeTaylor2014}{Ringe and Taylor 2014}).

\subsection{Old English evidence}\label{old-english-evidence-31}

Bosworth-Toller and Clark Hall record \emph{hærfest} `harvest', with
\emph{herfest} `harvest' as a variant in the lexical tradition
(\citeproc{ref-BosworthToller1898}{Bosworth and Toller 1898};
\citeproc{ref-ClarkHall1960}{Clark Hall 1960}). Those attested forms
remain the main dictionary evidence for the noun.

\subsection{Development to Old
English}\label{development-to-old-english-31}

From \Recon{xárbistuz} `harvest', Anglo-Frisian brightening, breaking,
and i-mutation produce a \emph{hierbist-} stage, and later lowering of
unstressed medial \emph{i} to \emph{e} gives \emph{hierfest} `harvest'.
That is the regular West Saxon development treated here
(\citeproc{ref-RingeTaylor2014}{Ringe and Taylor 2014};
\citeproc{ref-Campbell1959}{Campbell 1959}).

\subsection{Source note}\label{source-note-3}

The Old English form here, \emph{hierfest}, represents the regular
native West Saxon outcome discussed by Bammesberger and Ringe-Taylor.
The attested Old English lexical tradition, however, is chiefly
\emph{hærfest} `harvest' / \emph{herfest} `harvest', commonly treated as
non-West-Saxon or Anglian material in West Saxon transmission
(\citeproc{ref-Bammesberger1997}{Bammesberger 1997};
\citeproc{ref-RingeTaylor2014}{Ringe and Taylor 2014}).

\section{\texorpdfstring{hedge --- OE
\emph{heġġ}}{hedge --- OE heġġ}}\label{hedge-oe-heux121ux121}

\index[iv]{02oe@\ivlangheader{Old English}{West Saxon normalization unmarked}!hegg@\iventry{heġġ}{}}
\index[iv]{03pgmc@\ivlangheader{Proto-Germanic}{}!xagjaz@\iventry{*xágjaz}{}}

Derivation: \emph{*xágjaz} \textgreater{} \emph{heġġ} (regular).

\subsection{Derivation trace}\label{derivation-trace-36}

Proto input: \emph{*xágjaz}

\begingroup
\setlength{\fboxsep}{6pt}

\noindent\fbox{%
\begin{minipage}{0.97\linewidth}
\small
\begin{tabular}{@{}>{\raggedright\arraybackslash}p{0.388\linewidth}@{\hspace{0.030\linewidth}}>{\raggedright\arraybackslash}p{0.582\linewidth}@{}}
\begin{minipage}[t]{\linewidth}
\centering\textbf{Earlier Germanic changes}\par
\vspace{0.35em}
\raggedright
\centering\textbf{Northwest and West Germanic}\par
\raggedright
\vspace{0.2em}
\begin{tabularx}{\linewidth}{@{}>{\raggedright\arraybackslash}X@{\hspace{0.55em}}>{\raggedright\arraybackslash}p{0.24\linewidth}@{\hspace{0.25em}}}
\mbox{PWGmc J Gemination} & \emph{*xággjaz} \\
\mbox{EAF Final Z Deletion} & \emph{*xággja} \\
\end{tabularx}
\vspace{0.6em}
\centering\textbf{Early Anglo-Frisian}\par
\raggedright
\vspace{0.2em}
\raggedright [no change]\par
\end{minipage}
&
\begin{minipage}[t]{\linewidth}
\centering\textbf{Old English changes}\par
\vspace{0.35em}
\raggedright
\begin{tabularx}{\linewidth}{@{}>{\raggedright\arraybackslash}X@{\hspace{0.55em}}>{\raggedright\arraybackslash}p{0.14\linewidth}@{\hspace{0.25em}}}
\mbox{PWGmc Final Bare A Loss} & \emph{*xággj} \\
\mbox{EAF Brightening} & \emph{*xæggj} \\
OE Velar Fricative Palatalization & \emph{*çæggj} \\
\mbox{OE Velar Palatalization} & \emph{*çæʤʤj} \\
\mbox{OE I Umlaut} & \emph{*çeʤʤj} \\
\mbox{OE J Loss After Heavy} & \emph{*çeʤʤ} \\
\end{tabularx}
\end{minipage}
\\
\end{tabular}
\end{minipage}%
} \endgroup

Old English form: \emph{heġġ}

\subsection{Reconstruction and comparative
evidence}\label{reconstruction-and-comparative-evidence-32}

The derivational input models a palatal \Recon{-gj-*} `hedge' noun whose
Old English development includes gemination, palatalization, and
i-mutation. The current derivation therefore reaches a palatal-geminate
outcome of the \emph{heġġ} `hedge' type
(\citeproc{ref-Campbell1959}{Campbell 1959}).

\subsection{Old English evidence}\label{old-english-evidence-32}

Bosworth-Toller and Clark Hall record the noun under standard spellings
\emph{hecg} `hedge' / \emph{heċġ} `hedge'
(\citeproc{ref-BosworthToller1898}{Bosworth and Toller 1898};
\citeproc{ref-ClarkHall1960}{Clark Hall 1960}). The lexical item itself
is therefore well attested even though the form compared here is
normalized.

\subsection{Development to Old
English}\label{development-to-old-english-32}

This is the geminate-stop class, distinct from the singleton fricative
in \emph{dæġ} `day'. The model's palatal stop-reflex proxy does not
assert that affrication had already occurred at the articulation
checkpoint (\citeproc{ref-RingeTaylor2014}{Ringe and Taylor 2014,
203--4}; \citeproc{ref-Fulk2018}{Fulk 2018, 131--32}).

From \Recon{xágjaz} `hedge', the current serialized model first applies
West Germanic j-gemination to produce a geminate stop; later Old English
palatalization and loss of final \emph{j} produce \emph{heġġ} `hedge'.
This preserves the selected regular executable derivation, not
independent evidence for the historical position of final low-vowel
loss.

Ringe and Taylor place broader oral/nasal low-vowel apocope before West
Germanic j-gemination and distinguish stranded-glide outcomes from later
paradigmatic leveling (\citeproc{ref-RingeTaylor2014}{Ringe and Taylor
2014, 46--47}). The late oral-a fragment therefore does not establish an
uninterrupted regular history of this selected nominative. Reconciling
its formation and cell with that chronology remains open; no
word-specific analogical repair or change to the selected reconstruction
is asserted here.

\subsection{Form note}\label{form-note-6}

Standard dictionary spelling is \emph{heċġ} `hedge' or \emph{hecg}
`hedge'. Normalized \emph{heġġ} `hedge' is the Old English form here,
while the ordinary lexicographic forms remain the main Old English
citation evidence (\citeproc{ref-BosworthToller1898}{Bosworth and Toller
1898}; \citeproc{ref-ClarkHall1960}{Clark Hall 1960}).

\section{\texorpdfstring{helm --- OE
\emph{helm}}{helm --- OE helm}}\label{helm-oe-helm}

\index[iv]{02oe@\ivlangheader{Old English}{West Saxon normalization unmarked}!helm@\iventry{helm}{}}
\index[iv]{03pgmc@\ivlangheader{Proto-Germanic}{}!xelmaz@\iventry{*xélmaz}{}}

Derivation: \emph{*xélmaz} \textgreater{} \emph{helm} (regular).

\subsection{Derivation trace}\label{derivation-trace-37}

Proto input: \emph{*xélmaz}

\begingroup
\setlength{\fboxsep}{6pt}

\noindent\fbox{%
\begin{minipage}{0.97\linewidth}
\small
\begin{tabular}{@{}>{\raggedright\arraybackslash}p{0.359\linewidth}@{\hspace{0.030\linewidth}}>{\raggedright\arraybackslash}p{0.611\linewidth}@{}}
\begin{minipage}[t]{\linewidth}
\centering\textbf{Earlier Germanic changes}\par
\vspace{0.35em}
\raggedright
\centering\textbf{Northwest and West Germanic}\par
\raggedright
\vspace{0.2em}
\begin{tabularx}{\linewidth}{@{}>{\raggedright\arraybackslash}X@{\hspace{0.55em}}>{\raggedright\arraybackslash}p{0.14\linewidth}@{\hspace{0.25em}}}
\mbox{EAF Final Z Deletion} & \emph{*xélma} \\
\end{tabularx}
\vspace{0.6em}
\centering\textbf{Early Anglo-Frisian}\par
\raggedright
\vspace{0.2em}
\raggedright [no change]\par
\end{minipage}
&
\begin{minipage}[t]{\linewidth}
\centering\textbf{Old English changes}\par
\vspace{0.35em}
\raggedright
\begin{tabularx}{\linewidth}{@{}>{\raggedright\arraybackslash}X@{\hspace{0.55em}}>{\raggedright\arraybackslash}p{0.14\linewidth}@{\hspace{0.25em}}}
\mbox{PWGmc Final Bare A Loss} & \emph{*xélm} \\
OE Velar Fricative Palatalization & \emph{*çélm} \\
\end{tabularx}
\end{minipage}
\\
\end{tabular}
\end{minipage}%
} \endgroup

Old English form: \emph{helm}

\subsection{Reconstruction and comparative
evidence}\label{reconstruction-and-comparative-evidence-33}

Kroonen cites the helmet noun as \emph{*helma-} and separately
distinguishes a different lexeme \emph{*helman-} `rudder'
(\citeproc{ref-Kroonen2013}{Kroonen 2013}). The derivational input
\Recon{xélmaz} `helm' is the nominative-style form used for the helmet
noun itself.

\subsection{Old English evidence}\label{old-english-evidence-33}

Clark Hall and Bosworth-Toller record
{\emph{helm}}\index[iv]{02oe@\ivlangheader{Old English}{West Saxon normalization unmarked}!helm@\iventry{helm}{}}
`helm, helmet' as the ordinary Old English noun for `helmet', while
{\emph{helma}}\index[iv]{02oe@\ivlangheader{Old English}{West Saxon normalization unmarked}!helma@\iventry{helma}{}}
`helm' belongs to a separate rudder lexeme
(\citeproc{ref-ClarkHall1960}{Clark Hall 1960};
\citeproc{ref-BosworthToller1898}{Bosworth and Toller 1898, 542}).

\subsection{Development to Old
English}\label{development-to-old-english-33}

From \Recon{xélmaz} `helm', loss of final \emph{z} and later loss of the
short final vowel yield
{\emph{helm}}\index[iv]{02oe@\ivlangheader{Old English}{West Saxon normalization unmarked}!helm@\iventry{helm}{}}
`helm, helmet'. The development is therefore a straightforward
citation-form match.

\subsection{Form note}\label{form-note-7}

Comparative \emph{*helma-} is headword notation for the helmet cognate
set. It should not be confused with Old English
{\emph{helma}}\index[iv]{02oe@\ivlangheader{Old English}{West Saxon normalization unmarked}!helma@\iventry{helma}{}}
`helm', which is a different noun meaning `helm, rudder'
(\citeproc{ref-Kroonen2013}{Kroonen 2013};
\citeproc{ref-ClarkHall1960}{Clark Hall 1960}).

\section{\texorpdfstring{help --- OE
\emph{helpan}}{help --- OE helpan}}\label{help-oe-helpan}

\index[iv]{02oe@\ivlangheader{Old English}{West Saxon normalization unmarked}!helpan@\iventry{helpan}{}}
\index[iv]{03pgmc@\ivlangheader{Proto-Germanic}{}!xelpana@\iventry{*xélpaną}{}}

Derivation: \emph{*xélpaną} \textgreater{} \emph{helpan} (regular).

\subsection{Derivation trace}\label{derivation-trace-38}

Proto input: \emph{*xélpaną}

\begingroup
\setlength{\fboxsep}{6pt}

\noindent\fbox{%
\begin{minipage}{0.97\linewidth}
\small
\begin{tabular}{@{}>{\raggedright\arraybackslash}p{0.359\linewidth}@{\hspace{0.030\linewidth}}>{\raggedright\arraybackslash}p{0.611\linewidth}@{}}
\begin{minipage}[t]{\linewidth}
\centering\textbf{Earlier Germanic changes}\par
\vspace{0.35em}
\raggedright
\centering\textbf{Northwest and West Germanic}\par
\raggedright
\vspace{0.2em}
\raggedright [no change]\par
\vspace{0.6em}
\centering\textbf{Early Anglo-Frisian}\par
\raggedright
\vspace{0.2em}
\raggedright [no change]\par
\end{minipage}
&
\begin{minipage}[t]{\linewidth}
\centering\textbf{Old English changes}\par
\vspace{0.35em}
\raggedright
\begin{tabularx}{\linewidth}{@{}>{\raggedright\arraybackslash}X@{\hspace{0.55em}}>{\raggedright\arraybackslash}p{0.18\linewidth}@{\hspace{0.25em}}}
OE Velar Fricative Palatalization & \emph{*çélpaną} \\
OE Heavy Syllable Nasal Apocope & \emph{*çélpan} \\
\mbox{OE Secondary Nasalization} & \emph{*çélpąn} \\
\mbox{OE Weak Tail Reduction} & \emph{*çélpan} \\
\end{tabularx}
\end{minipage}
\\
\end{tabular}
\end{minipage}%
} \endgroup

Old English form: \emph{helpan}

\subsection{Reconstruction and comparative
evidence}\label{reconstruction-and-comparative-evidence-34}

The entry treats the strong verb itself rather than the separate noun
{\emph{help}}\index[iv]{02oe@\ivlangheader{Old English}{West Saxon normalization unmarked}!help@\iventry{help}{}}
`help'. Bright's principal parts \emph{helpan, healp, hulpon, holpen}
show the ordinary Old English strong-verb family continued by this input
(\citeproc{ref-BrightCassidyRingler1971}{Bright 1971}).

\subsection{Old English evidence}\label{old-english-evidence-34}

Clark Hall and Bosworth-Toller record
{\emph{helpan}}\index[iv]{02oe@\ivlangheader{Old English}{West Saxon normalization unmarked}!helpan@\iventry{helpan}{}}
`help' as the verbal headword `to help'
(\citeproc{ref-ClarkHall1960}{Clark Hall 1960};
\citeproc{ref-BosworthToller1898}{Bosworth and Toller 1898, 542}). The
target is therefore the attested infinitive citation form.

\subsection{Development to Old
English}\label{development-to-old-english-34}

From \Recon{xélpaną} `help', no special repair is needed beyond the
ordinary reduction of the infinitive ending. The derivation therefore
reaches
{\emph{helpan}}\index[iv]{02oe@\ivlangheader{Old English}{West Saxon normalization unmarked}!helpan@\iventry{helpan}{}}
`help' directly.

\subsection{Form note}\label{form-note-8}

Noun
{\emph{help}}\index[iv]{02oe@\ivlangheader{Old English}{West Saxon normalization unmarked}!help@\iventry{help}{}}
`help' belongs to a separate lexical line and should not replace verbal
{\emph{helpan}}\index[iv]{02oe@\ivlangheader{Old English}{West Saxon normalization unmarked}!helpan@\iventry{helpan}{}}
`help' as the target here (\citeproc{ref-ClarkHall1960}{Clark Hall
1960}; \citeproc{ref-BosworthToller1898}{Bosworth and Toller 1898,
542}).

\section{\texorpdfstring{hew --- OE
\emph{hēawan}}{hew --- OE hēawan}}\label{hew-oe-hux113awan}

\index[iv]{02oe@\ivlangheader{Old English}{West Saxon normalization unmarked}!heawan@\iventry{hēawan}{}}
\index[iv]{03pgmc@\ivlangheader{Proto-Germanic}{}!xawwana@\iventry{*xáwwaną}{}}

Derivation: \emph{*xáwwaną} \textgreater{} \emph{hēawan} (regular).

\subsection{Derivation trace}\label{derivation-trace-39}

Proto input: \emph{*xáwwaną}

\begingroup
\setlength{\fboxsep}{6pt}

\noindent\fbox{%
\begin{minipage}{0.97\linewidth}
\footnotesize
\begin{tabular}{@{}>{\raggedright\arraybackslash}p{0.339\linewidth}@{\hspace{0.030\linewidth}}>{\raggedright\arraybackslash}p{0.630\linewidth}@{}}
\begin{minipage}[t]{\linewidth}
\centering\textbf{Earlier Germanic changes}\par
\vspace{0.35em}
\raggedright
\centering\textbf{Northwest and West Germanic}\par
\raggedright
\vspace{0.2em}
\begin{tabularx}{\linewidth}{@{}>{\raggedright\arraybackslash}X@{\hspace{0.55em}}>{\raggedright\arraybackslash}p{0.24\linewidth}@{\hspace{0.25em}}}
\mbox{OE WW Simplification} & \emph{*xáuwaną} \\
\end{tabularx}
\vspace{0.6em}
\centering\textbf{Early Anglo-Frisian}\par
\raggedright
\vspace{0.2em}
\raggedright [no change]\par
\end{minipage}
&
\begin{minipage}[t]{\linewidth}
\centering\textbf{Old English changes}\par
\vspace{0.35em}
\raggedright
\begin{tabularx}{\linewidth}{@{}>{\raggedright\arraybackslash}X@{\hspace{0.55em}}>{\raggedright\arraybackslash}p{0.22\linewidth}@{\hspace{0.25em}}}
\mbox{OE Au Brightening} & \emph{*xáeuwaną} \\
\mbox{OE Diphthong Leveling} & \emph{*xēawaną} \\
OE Velar Fricative Palatalization & \emph{*çēawaną} \\
OE Heavy Syllable Nasal Apocope & \emph{*çēawan} \\
\mbox{OE Secondary Nasalization} & \emph{*çēawąn} \\
\mbox{OE Weak Tail Reduction} & \emph{*çēawan} \\
\end{tabularx}
\end{minipage}
\\
\end{tabular}
\end{minipage}%
} \endgroup

Old English form: \emph{hēawan}

\subsection{Reconstruction and
development}\label{reconstruction-and-development-2}

PGmc \emph{*xáwwaną} gives OE \emph{hēawan} `hew'. Its inherited
short-vowel-plus-geminate sequence yields \emph{*au} plus retained
consonantal \emph{*w} in the earlier West Germanic reanalysis. English
fronting and realization then produce the long diphthong. The geminate
is not simply deleted at the later English checkpoint
(\citeproc{ref-RingeTaylor2014}{Ringe and Taylor 2014, 65--66, 171--75};
\citeproc{ref-Campbell1959}{Campbell 1959, 45--47}).

Together with dew, this word follows
\hyperref[rule-OEWWSimplification]{SC031 OEWWSimplification},
\hyperref[rule-OEAuBrightening]{SC030 OEAuBrightening}, and
\hyperref[rule-OEDiphthongLeveling]{SC032 OEDiphthongLeveling}. It
contrasts with the j-created geminates of hay and reconstructed West
Saxon strew, which retain their independently governed resolution. The
earlier and later components are separate even though the final spelling
is unchanged (\citeproc{ref-Campbell1959}{Campbell 1959, 45--47};
\citeproc{ref-RingeTaylor2014}{Ringe and Taylor 2014, 53, 173}).

\section{\texorpdfstring{hind --- OE
\emph{hind}}{hind --- OE hind}}\label{hind-oe-hind}

\index[iv]{02oe@\ivlangheader{Old English}{West Saxon normalization unmarked}!hind@\iventry{hind}{}}
\index[iv]{03pgmc@\ivlangheader{Proto-Germanic}{}!xendjo@\iventry{*xéndjō}{}}

Derivation: \emph{*xéndjō} \textgreater{} \emph{hind} (regular).

\subsection{Derivation trace}\label{derivation-trace-40}

Proto input: \emph{*xéndjō}

\begingroup
\setlength{\fboxsep}{6pt}

\noindent\fbox{%
\begin{minipage}{0.97\linewidth}
\small
\begin{tabular}{@{}>{\raggedright\arraybackslash}p{0.485\linewidth}@{\hspace{0.030\linewidth}}>{\raggedright\arraybackslash}p{0.485\linewidth}@{}}
\begin{minipage}[t]{\linewidth}
\centering\textbf{Earlier Germanic changes}\par
\vspace{0.35em}
\raggedright
\centering\textbf{Northwest and West Germanic}\par
\raggedright
\vspace{0.2em}
\begin{tabularx}{\linewidth}{@{}>{\raggedright\arraybackslash}X@{\hspace{0.55em}}>{\raggedright\arraybackslash}p{0.14\linewidth}@{\hspace{0.25em}}}
\mbox{PNWGmc Final Long O Raising} & \emph{*xéndju} \\
\end{tabularx}
\vspace{0.6em}
\centering\textbf{Early Anglo-Frisian}\par
\raggedright
\vspace{0.2em}
\raggedright [no change]\par
\end{minipage}
&
\begin{minipage}[t]{\linewidth}
\centering\textbf{Old English changes}\par
\vspace{0.35em}
\raggedright
\begin{tabularx}{\linewidth}{@{}>{\raggedright\arraybackslash}X@{\hspace{0.55em}}>{\raggedright\arraybackslash}p{0.14\linewidth}@{\hspace{0.25em}}}
OE Velar Fricative Palatalization & \emph{*çéndju} \\
\mbox{OE I Umlaut} & \emph{*çindju} \\
\mbox{OE High Vowel Apocope} & \emph{*çindj} \\
\mbox{OE J Loss After Heavy} & \emph{*çind} \\
\end{tabularx}
\end{minipage}
\\
\end{tabular}
\end{minipage}%
} \endgroup

Old English form: \emph{hind}

\subsection{Reconstruction and comparative
evidence}\label{reconstruction-and-comparative-evidence-35}

Kroonen cites the animal name as \emph{*hindō-} f.~`hind'
(\citeproc{ref-Kroonen2013}{Kroonen 2013}). The derivational input
\Recon{xéndjō} `hind' represents that same noun in the present
derivation.

\subsection{Old English evidence}\label{old-english-evidence-35}

Clark Hall and Bosworth-Toller record
{\emph{hind}}\index[iv]{02oe@\ivlangheader{Old English}{West Saxon normalization unmarked}!hind@\iventry{hind}{}}
`hind, doe' as the noun `hind, female deer'
(\citeproc{ref-ClarkHall1960}{Clark Hall 1960};
\citeproc{ref-BosworthToller1898}{Bosworth and Toller 1898, 554}). The
target is therefore the attested lexical item itself.

\subsection{Development to Old
English}\label{development-to-old-english-35}

From \Recon{xéndjō} `hind', i-mutation produces the front-vocalic Old
English stem, and later apocope plus loss of final \emph{j} yield
\emph{hind}. The outcome is therefore a regular citation-form
derivation.

\subsection{Form note}\label{form-note-9}

\emph{hindan} `from behind, behind' is a different Old English lexeme
and does not belong to the noun history of \emph{hind}
(\citeproc{ref-ClarkHall1960}{Clark Hall 1960};
\citeproc{ref-BosworthToller1898}{Bosworth and Toller 1898, 554}).

\section{\texorpdfstring{hold --- OE
\emph{healdan}}{hold --- OE healdan}}\label{hold-oe-healdan}

\index[iv]{02oe@\ivlangheader{Old English}{West Saxon normalization unmarked}!healdan@\iventry{healdan}{}}
\index[iv]{03pgmc@\ivlangheader{Proto-Germanic}{}!xaldana@\iventry{*xáldaną}{}}

Derivation: \emph{*xáldaną} \textgreater{} \emph{healdan} (regular).

\subsection{Derivation trace}\label{derivation-trace-41}

Proto input: \emph{*xáldaną}

\begingroup
\setlength{\fboxsep}{6pt}

\noindent\fbox{%
\begin{minipage}{0.97\linewidth}
\footnotesize
\begin{tabular}{@{}>{\raggedright\arraybackslash}p{0.339\linewidth}@{\hspace{0.030\linewidth}}>{\raggedright\arraybackslash}p{0.630\linewidth}@{}}
\begin{minipage}[t]{\linewidth}
\centering\textbf{Earlier Germanic changes}\par
\vspace{0.35em}
\raggedright
\centering\textbf{Northwest and West Germanic}\par
\raggedright
\vspace{0.2em}
\raggedright [no change]\par
\vspace{0.6em}
\centering\textbf{Early Anglo-Frisian}\par
\raggedright
\vspace{0.2em}
\raggedright [no change]\par
\end{minipage}
&
\begin{minipage}[t]{\linewidth}
\centering\textbf{Old English changes}\par
\vspace{0.35em}
\raggedright
\begin{tabularx}{\linewidth}{@{}>{\raggedright\arraybackslash}X@{\hspace{0.55em}}>{\raggedright\arraybackslash}p{0.22\linewidth}@{\hspace{0.25em}}}
\mbox{EAF Brightening} & \emph{*xældaną} \\
\mbox{OE Breaking} & \emph{*xealdaną} \\
OE Velar Fricative Palatalization & \emph{*çealdaną} \\
OE Heavy Syllable Nasal Apocope & \emph{*çealdan} \\
\mbox{OE Secondary Nasalization} & \emph{*çealdąn} \\
\mbox{OE Weak Tail Reduction} & \emph{*çealdan} \\
\end{tabularx}
\end{minipage}
\\
\end{tabular}
\end{minipage}%
} \endgroup

Old English form: \emph{healdan}

\subsection{Reconstruction and comparative
evidence}\label{reconstruction-and-comparative-evidence-36}

Campbell and Ringe-Taylor treat the verb as a regular \emph{a + lC}
breaking case, with West Saxon
{\emph{healdan}}\index[iv]{02oe@\ivlangheader{Old English}{West Saxon normalization unmarked}!healdan@\iventry{healdan}{}}
`hold' opposed to Anglian and Mercian
{\emph{haldan}}\index[iv]{02oe@\ivlangheader{Old English}{West Saxon normalization unmarked}!haldan~04@\iventry{haldan}{Angl.}}
`hold' (\citeproc{ref-Campbell1959}{Campbell 1959};
\citeproc{ref-RingeTaylor2014}{Ringe and Taylor 2014}).

\subsection{Old English evidence}\label{old-english-evidence-36}

Bright gives the ordinary strong-verb citation form and principal parts
\emph{healdan, heold, heoldon, healden}
(\citeproc{ref-BrightCassidyRingler1971}{Bright 1971}). The target is
therefore the attested infinitive headword itself.

\subsection{Development to Old
English}\label{development-to-old-english-36}

From \Recon{xáldaną} `hold', Anglo-Frisian brightening first yields a
fronted vowel, and West Saxon breaking then produces \emph{ea} before
\emph{ld}. Later reduction of the infinitive ending gives
{\emph{healdan}}\index[iv]{02oe@\ivlangheader{Old English}{West Saxon normalization unmarked}!healdan@\iventry{healdan}{}}
`hold' (\citeproc{ref-Campbell1959}{Campbell 1959};
\citeproc{ref-RingeTaylor2014}{Ringe and Taylor 2014}).

\subsection{Dialect note}\label{dialect-note-2}

West Saxon
{\emph{healdan}}\index[iv]{02oe@\ivlangheader{Old English}{West Saxon normalization unmarked}!healdan@\iventry{healdan}{}}
`hold' is the Old English form here. Anglian and Mercian
{\emph{haldan}}\index[iv]{02oe@\ivlangheader{Old English}{West Saxon normalization unmarked}!haldan~04@\iventry{haldan}{Angl.}}
`hold' are genuine non-West-Saxon doublets rather than corrections to
that choice (\citeproc{ref-Campbell1959}{Campbell 1959};
\citeproc{ref-RingeTaylor2014}{Ringe and Taylor 2014}).

\section{\texorpdfstring{horn --- OE
\emph{horn}}{horn --- OE horn}}\label{horn-oe-horn}

\index[iv]{02oe@\ivlangheader{Old English}{West Saxon normalization unmarked}!horn@\iventry{horn}{}}
\index[iv]{03pgmc@\ivlangheader{Proto-Germanic}{}!xurna@\iventry{*xúrną}{}}

Derivation: \emph{*xúrną} \textgreater{} \emph{horn} (regular).

\subsection{Derivation trace}\label{derivation-trace-42}

Proto input: \emph{*xúrną}

\begingroup
\setlength{\fboxsep}{6pt}

\noindent\fbox{%
\begin{minipage}{0.97\linewidth}
\small
\begin{tabular}{@{}>{\raggedright\arraybackslash}p{0.359\linewidth}@{\hspace{0.030\linewidth}}>{\raggedright\arraybackslash}p{0.611\linewidth}@{}}
\begin{minipage}[t]{\linewidth}
\centering\textbf{Earlier Germanic changes}\par
\vspace{0.35em}
\raggedright
\centering\textbf{Northwest and West Germanic}\par
\raggedright
\vspace{0.2em}
\begin{tabularx}{\linewidth}{@{}>{\raggedright\arraybackslash}X@{\hspace{0.55em}}>{\raggedright\arraybackslash}p{0.14\linewidth}@{\hspace{0.25em}}}
\mbox{PNWGmc U Lowering} & \emph{*xórną} \\
\end{tabularx}
\vspace{0.6em}
\centering\textbf{Early Anglo-Frisian}\par
\raggedright
\vspace{0.2em}
\raggedright [no change]\par
\end{minipage}
&
\begin{minipage}[t]{\linewidth}
\centering\textbf{Old English changes}\par
\vspace{0.35em}
\raggedright
\begin{tabularx}{\linewidth}{@{}>{\raggedright\arraybackslash}X@{\hspace{0.55em}}>{\raggedright\arraybackslash}p{0.14\linewidth}@{\hspace{0.25em}}}
OE Heavy Syllable Nasal Apocope & \emph{*xórn} \\
\end{tabularx}
\end{minipage}
\\
\end{tabular}
\end{minipage}%
} \endgroup

Old English form: \emph{horn}

\subsection{Reconstruction and comparative
evidence}\label{reconstruction-and-comparative-evidence-37}

Kroonen and Orel cite lemma-style Proto-Germanic headwords of the
\emph{*hurna-} / \Recon{xurnan} `horn' type for this noun
(\citeproc{ref-Kroonen2013}{Kroonen 2013}; \citeproc{ref-Orel2003}{Orel
2003, 234}). The derivational input \Recon{xúrną} `horn' is the
nominative-style form used in the derivation here.

\subsection{Old English evidence}\label{old-english-evidence-37}

Clark Hall, Bosworth-Toller, and Bright all record \emph{horn} as the
ordinary Old English noun (\citeproc{ref-ClarkHall1960}{Clark Hall
1960}; \citeproc{ref-BosworthToller1898}{Bosworth and Toller 1898, 108};
\citeproc{ref-BrightCassidyRingler1971}{Bright 1971}). The target is
therefore the attested citation form.

\subsection{Development to Old
English}\label{development-to-old-english-37}

From \Recon{xúrną} `horn', Northwest Germanic u-lowering gives
\Recon{xórną} `horn', and later loss of the final nasal leaves
\emph{horn}. The development treated here is fully regular.

\subsection{Form note}\label{form-note-10}

The note's oblique \Recon{xurnăn} `horn' belongs to comparative stem
background only. It does not replace the derivational input
\Recon{xúrną} `horn' as the derivational form used here
(\citeproc{ref-Kroonen2013}{Kroonen 2013}; \citeproc{ref-Orel2003}{Orel
2003, 234}).

\section{\texorpdfstring{hue --- OE
\emph{hīew}}{hue --- OE hīew}}\label{hue-oe-hux12bew}

\index[iv]{02oe@\ivlangheader{Old English}{West Saxon normalization unmarked}!hiew@\iventry{hīew}{}}
\index[iv]{03pgmc@\ivlangheader{Proto-Germanic}{}!xiwja@\iventry{*xíwją}{}}

Derivation: \emph{*xíwją} \textgreater{} \emph{hīew} (regular).

\subsection{Derivation trace}\label{derivation-trace-43}

Proto input: \emph{*xíwją}

\begingroup
\setlength{\fboxsep}{6pt}

\noindent\fbox{%
\begin{minipage}{0.97\linewidth}
\small
\begin{tabular}{@{}>{\raggedright\arraybackslash}p{0.339\linewidth}@{\hspace{0.030\linewidth}}>{\raggedright\arraybackslash}p{0.630\linewidth}@{}}
\begin{minipage}[t]{\linewidth}
\centering\textbf{Earlier Germanic changes}\par
\vspace{0.35em}
\raggedright
\centering\textbf{Northwest and West Germanic}\par
\raggedright
\vspace{0.2em}
\begin{tabularx}{\linewidth}{@{}>{\raggedright\arraybackslash}X@{\hspace{0.55em}}>{\raggedright\arraybackslash}p{0.24\linewidth}@{\hspace{0.25em}}}
\mbox{PWGmc J Gemination} & \emph{*xíwwją} \\
\end{tabularx}
\vspace{0.6em}
\centering\textbf{Early Anglo-Frisian}\par
\raggedright
\vspace{0.2em}
\raggedright [no change]\par
\end{minipage}
&
\begin{minipage}[t]{\linewidth}
\centering\textbf{Old English changes}\par
\vspace{0.35em}
\raggedright
\begin{tabularx}{\linewidth}{@{}>{\raggedright\arraybackslash}X@{\hspace{0.55em}}>{\raggedright\arraybackslash}p{0.14\linewidth}@{\hspace{0.25em}}}
\mbox{OE Ew Long Diphthong} & \emph{*xíuwją} \\
\mbox{OE J Glide IO} & \emph{*xīowją} \\
OE Velar Fricative Palatalization & \emph{*çīowją} \\
OE Heavy Syllable Nasal Apocope & \emph{*çīowj} \\
\mbox{OE I Umlaut} & \emph{*çīewj} \\
\mbox{OE J Loss After Heavy} & \emph{*çīew} \\
\end{tabularx}
\end{minipage}
\\
\end{tabular}
\end{minipage}%
} \endgroup

Old English form: \emph{hīew}

\subsection{Reconstruction and comparative
evidence}\label{reconstruction-and-comparative-evidence-38}

The selected reconstruction is realized Proto-Germanic \emph{*xíwją}
`form, appearance', a neuter \emph{ja}-stem. It follows Ringe and
Taylor's explicit PGmc
\Recon{hiwją}\index[iv]{03pgmc@\ivlangheader{Proto-Germanic}{}!hiwja@\iventry{*hiwją}{}}
`form, appearance', normalized with CAPR's dorsal \emph{x} and lexical
acute (\citeproc{ref-RingeTaylor2014}{Ringe and Taylor 2014, 250}). Both
the citation reconstruction and the derivational input use i. This is an
adopted working reconstruction, not a claim that all dictionaries print
it or that their e forms have been shown to belong to a different
historical stage.

\subsection{The published e
reconstructions}\label{the-published-e-reconstructions}

Orel gives
\Recon{xewjan}\index[iv]{03pgmc@\ivlangheader{Proto-Germanic}{}!xewjan@\iventry{*xewjan}{}}
`form, appearance', neuter, with Gothic \emph{hiwi} `surface
appearance', Old Norse \emph{hý} `down on plants' and Old English
\emph{hīw} `form'. He probably derives the noun from \emph{*xawwanan}
and lists older etymological alternatives (\citeproc{ref-Orel2003}{Orel
2003, 171--72}). The e vowel is therefore not an invented CAPR spelling.
The dictionary entry supplies a real competing reconstruction; it does
not license us to call its e an abstract underlying grade or an
automatically pre-PGmc form. His final \emph{-an} must also be
distinguished from the other authors' stem citation or nasal-vowel
notation.

Kroonen gives
\Recon{heuja-}\index[iv]{03pgmc@\ivlangheader{Proto-Germanic}{}!heuja@\iventry{*heuja-}{}},
a neuter stem `visible layer, appearance', on the comparative set Gothic
{hiwi}\index[iv]{12goth@\ivlangheader{Gothic}{}!hiwi@\iventry{hiwi}{}}
`surface appearance, complexion', Old Norse
{hý}\index[iv]{09on@\ivlangheader{Old Norse}{}!hy@\iventry{hý}{}} `down
on plants', Norwegian \emph{hy}, and Old English \emph{hīw}, \emph{hēow}
`shape, form, hue, colour'. He considers an early derivative of
\emph{*kieh₁-u-}, with \emph{*heuja-} probably arising from
\emph{*kēu-ió-} by Dybo's law (\citeproc{ref-Kroonen2013}{Kroonen 2013,
224}). This is not the same derivational argument as Orel's. Their
agreement on e cannot be treated as two identical histories, and their
different endings cannot be silently collapsed into one quoted
whole-word form. Writing prevocalic u as w and h as x explains a
normalization; it does not by itself explain the disagreement between e
and i.

\subsection{The realized-i account and its early
conditioning}\label{the-realized-i-account-and-its-early-conditioning}

Ringe and Taylor explicitly begin with PGmc \emph{*hiwją}. They
distinguish a nongeminated \emph{*hiwi} from the stem with palatalized
geminate w, and propose its development through \emph{*iuw} to pre-OE
long \emph{*īow}, with a still-palatalizing semivowel, then early West
Saxon long \emph{īew}. They call the development surprising and its
phonemic analysis unclear (\citeproc{ref-RingeTaylor2014}{Ringe and
Taylor 2014, 250}). The i choice is consequently not inferred merely
from the final Old English spelling: an author explicitly assigns it to
Proto-Germanic and offers a complete subsequent history for this noun.

Ringe separately describes early e-to-i raising before a high front
vocalic in the same or following syllable. His example \emph{*néwios} to
PGmc \emph{*niwjaz} `new' supplies the relevant w+j class; he qualifies
the evidence for the precise early date (\citeproc{ref-Ringe2017}{Ringe
2017, 151--53}). This passage explains the inherited-i preference but is
not itself a quotation of a whole-word hue reconstruction. Fulk likewise
treats PGmc e-to-i, and eu-to-iu, before following i/j separately from
later front umlaut (\citeproc{ref-Fulk2018}{Fulk 2018, 59}). Under this
working account, the raising precedes the supplied PGmc input. It must
not be replaced by a later one-word e-to-i operation simply because that
would conceal an input disagreement.

\subsection{The glide objection and the strength of the
choice}\label{the-glide-objection-and-the-strength-of-the-choice}

Campbell distinguishes the traditional awj and iwj developments and
prints the hue reflex in the iwj class
(\citeproc{ref-Campbell1959}{Campbell 1959, 46}, §120.2). Ringe and
Taylor's independent gemination comparison uses new, sew and mirth,
including nongeminated beside geminated paradigm stems
(\citeproc{ref-RingeTaylor2014}{Ringe and Taylor 2014, 53}). These
support distinguishing the consonant classes and successive operations;
they do not settle every phonemic detail of the offglide.

Fulk presents a substantive objection to the traditional dismantling
account. He questions how w retained the presumed consonantal status in
these inputs and how the resulting geminate was dismantled. His
discussion considers a different strew input and paradigm regularization
for new, rather than simply endorsing every conventional arrow
(\citeproc{ref-Fulk2018}{Fulk 2018, 71--72}). His general statement of
j-gemination after short vowels, not diphthongs, belongs to that issue
(\citeproc{ref-Fulk2018}{Fulk 2018, 126}, §6.15). He does not supply a
replacement complete hue history in these passages. Neither his
objection nor Ringe and Taylor's uncertainty is removed by our choosing
a computable working account.

We prefer realized PGmc i because the explicit reconstruction,
independently motivated early raising and complete hue chain form a
coherent account of the selected West Saxon reflex. This is not an
author majority, an automatic preference for the most recent book, or an
argument from successful final spelling. Confidence is stronger in the
inherited high-front conditioning and the selected dialectal reflex than
in a uniquely recoverable phonemic analysis of palatalized w. The e
reconstructions remain alternatives with positive published support;
they have not been disproved or silently reclassified.

These differences therefore warrant extended discussion despite a
regular selected result: the authors' formation premises, citation
endpoints and later phonemic uncertainties are distinct questions. No
majority vote or normalization of final spellings resolves them. The
existing source-led \emph{i} choice and full glide analysis are
retained, not replaced by a summary of dictionary vowels.

\subsection{Old English evidence}\label{old-english-evidence-38}

Campbell prints
{\emph{hīew}}\index[iv]{02oe@\ivlangheader{Old English}{West Saxon normalization unmarked}!hiew@\iventry{hīew}{}}
`form' beside \emph{hīow} `form' in the section that derives this stem
type, alongside \emph{nīowe}, \emph{nīewe} `new' and \emph{glīow},
\emph{glīw} `mirth' (\citeproc{ref-Campbell1959}{Campbell 1959, § 120.2}, p.~46). He distinguishes non-West-Saxon io from West Saxon ie.
CAPR selects \emph{hīew} as the source-supported West Saxon counterpart,
not as the only attested spelling. Orel gives \emph{hīw}, and Kroonen
gives \emph{hīw}, \emph{hēow} (\citeproc{ref-Orel2003}{Orel 2003,
171--72}; \citeproc{ref-Kroonen2013}{Kroonen 2013, 224}).

\subsection{Development to Old
English}\label{development-to-old-english-38}

Campbell's decisive class contrast is that awj and iwj share gemination
but thereafter lose different glides: the low-vowel type yields the
palatal consonant, whereas the high-front type retains w and yields
io/ie (\citeproc{ref-Campbell1959}{Campbell 1959, 46}, §120.2). For hue,
the adopted cascade now follows the independently specified Ringe-Taylor
chain rather than using long ēo as an output-compatible proxy for their
long īo (\citeproc{ref-RingeTaylor2014}{Ringe and Taylor 2014, 250}).

{\def\LTcaptype{none} % do not increment counter
\begin{longtable}[]{@{}
  >{\raggedright\arraybackslash}p{(\linewidth - 4\tabcolsep) * \real{0.3333}}
  >{\raggedright\arraybackslash}p{(\linewidth - 4\tabcolsep) * \real{0.3333}}
  >{\raggedright\arraybackslash}p{(\linewidth - 4\tabcolsep) * \real{0.3333}}@{}}
\toprule\noalign{}
\begin{minipage}[b]{\linewidth}\raggedright
Step
\end{minipage} & \begin{minipage}[b]{\linewidth}\raggedright
Selected reconstruction / represented checkpoint
\end{minipage} & \begin{minipage}[b]{\linewidth}\raggedright
Interpretation
\end{minipage} \\
\midrule\noalign{}
\endhead
\bottomrule\noalign{}
\endlastfoot
PGmc input & \emph{*xíwją} & Inherited i, with early raising already
reflected in the reconstruction \\
J-gemination & \emph{*xíwwją} & Palatalized geminate carried as ww+j \\
J-created reanalysis & \emph{*xíuwją} & Offglide formation, retaining
palatalized w as w+j \\
English realization & \emph{*xīowją} & Long īo, distinct from ordinary
ēo \\
Palatal articulation and ending loss & \emph{*çīowj} & Independent
dorsal history and reduction in the existing cascade \\
Mutation & \emph{*çīewj} & Existing long-diphthong mutation conditioned
by the retained palatal carrier \\
Later glide loss and spelling & \emph{hīew} & Retained w and the
selected native orthography \\
\end{longtable}
}

These are CAPR's checkpoints, not diplomatic quotations of all the
authors' intermediate forms. In particular, superscript j on a source w
denotes palatalization. Our separate j symbol carries that phonetic
conditioning until the relevant consumer; it does not decide that a
second j phoneme existed alongside every palatalized w. No new phonetic
marker or lexical exception is needed.

Knee now selects the regular dative \emph{cneowe} `knee (dat.sg.)',
whose short diphthong is produced by \hyperref[rule-OEBreaking]{SC044
OEBreaking}; it no longer shares this long-output component. The source
distinguishes its short pre-ending stem from long endingless \emph{cnēo}
`knee' and restored \emph{w} in \emph{cnēow}
(\citeproc{ref-HoggGrammar2011}{Hogg 2011, 21--22, 86};
\citeproc{ref-RingeTaylor2014}{Ringe and Taylor 2014, 387}).

The j-created geminate is distinct from the earlier inherited-Vww
operation. Its offglide checkpoint belongs to
\hyperref[rule-OEEwLongDiphthong]{SC033 OEEwLongDiphthong}; its long-io
realization is separately represented by
\hyperref[rule-OEJGlideIO]{SC110 OEJGlideIO}, before unchanged ordinary
\hyperref[rule-OEDiphthongLeveling]{SC032 OEDiphthongLeveling}. The
composed \hyperref[rule-OEJWWSimplification]{SC106 OEJWWSimplification}
compatibility operation no longer acts on hue, which already has
singleton w. Its absence of a live application is not universal proof of
redundancy. The earlier inherited-reanalysis placement remains a
representative serialization, not a newly proved strict chronology
(\citeproc{ref-Campbell1959}{Campbell 1959, 45--47};
\citeproc{ref-RingeTaylor2014}{Ringe and Taylor 2014, 53, 65--66,
171--75, 250}).

The contrast with \emph{hīeġ} `hay', from \emph{*xáwją}, remains exact:
the selected inputs differ in one vowel and both geminate, but the
low-vowel word loses w and retains the palatal consonant while hue
retains w. The low-vowel resolution and subsequent fronting do not apply
here (\citeproc{ref-Campbell1959}{Campbell 1959, 46};
\citeproc{ref-RingeTaylor2014}{Ringe and Taylor 2014, 53, 250}).

\subsection{Why this row is in the
corpus}\label{why-this-row-is-in-the-corpus}

The corpus originally lacked a high-front witness separating gemination
of w from the low-vowel resolution illustrated by hay and strew. Hue
remains a positive witness for \hyperref[rule-PWGmcJGemination]{SC010
PWGmcJGemination} and a negative control for
\hyperref[rule-OEAwwjResolution]{SC029 OEAwwjResolution}. The corrected
input and distinct intermediate history now make that contrast
scientifically informative rather than merely spelling-compatible. No
additional lexeme or analogical mechanism is introduced.

\section{\texorpdfstring{lead --- OE
\emph{lǣdan}}{lead --- OE lǣdan}}\label{lead-oe-lux1e3dan}

\index[iv]{02oe@\ivlangheader{Old English}{West Saxon normalization unmarked}!laedan@\iventry{lǣdan}{}}
\index[iv]{03pgmc@\ivlangheader{Proto-Germanic}{}!laidijana@\iventry{*láidijaną}{}}

Derivation: \emph{*láidijaną} \textgreater{} \emph{lǣdan} (regular).

\subsection{Derivation trace}\label{derivation-trace-44}

Proto input: \emph{*láidijaną}

\begingroup
\setlength{\fboxsep}{6pt}

\noindent\fbox{%
\begin{minipage}{0.97\linewidth}
\small
\begin{tabular}{@{}>{\raggedright\arraybackslash}p{0.485\linewidth}@{\hspace{0.030\linewidth}}>{\raggedright\arraybackslash}p{0.485\linewidth}@{}}
\begin{minipage}[t]{\linewidth}
\centering\textbf{Earlier Germanic changes}\par
\vspace{0.35em}
\raggedright
\centering\textbf{Northwest and West Germanic}\par
\raggedright
\vspace{0.2em}
\raggedright [no change]\par
\vspace{0.6em}
\centering\textbf{Early Anglo-Frisian}\par
\raggedright
\vspace{0.2em}
\begin{tabularx}{\linewidth}{@{}>{\raggedright\arraybackslash}X@{\hspace{0.55em}}>{\raggedright\arraybackslash}p{0.20\linewidth}@{\hspace{0.25em}}}
\mbox{EAF Ai Monophthongization} & \emph{*lādijaną} \\
\end{tabularx}
\end{minipage}
&
\begin{minipage}[t]{\linewidth}
\centering\textbf{Old English changes}\par
\vspace{0.35em}
\raggedright
\begin{tabularx}{\linewidth}{@{}>{\raggedright\arraybackslash}X@{\hspace{0.55em}}>{\raggedright\arraybackslash}p{0.18\linewidth}@{\hspace{0.25em}}}
OE Heavy Syllable Nasal Apocope & \emph{*lādijan} \\
\mbox{OE Secondary Nasalization} & \emph{*lādijąn} \\
\mbox{Sievers Law Syncope} & \emph{*lādjąn} \\
\mbox{OE I Umlaut} & \emph{*lǣdjąn} \\
\mbox{OE Weak Tail Reduction} & \emph{*lǣdjan} \\
\mbox{OE J Loss After Heavy} & \emph{*lǣdan} \\
\end{tabularx}
\end{minipage}
\\
\end{tabular}
\end{minipage}%
} \endgroup

Old English form: \emph{lǣdan}

\subsection{Reconstruction and comparative
evidence}\label{reconstruction-and-comparative-evidence-39}

Ringe and Taylor derive Old English \emph{lǣdan} `lead' from
Proto-Germanic \Recon{laidijaną} `lead', and Kroonen likewise cites a
weak verb of the \emph{*laidjan-} type for `lead'
(\citeproc{ref-RingeTaylor2014}{Ringe and Taylor 2014};
\citeproc{ref-Kroonen2013}{Kroonen 2013, 363}).

\subsection{Old English evidence}\label{old-english-evidence-39}

Clark Hall and Bosworth-Toller both record \emph{lædan} `lead' /
\emph{lǣdan} `lead' as the ordinary Old English verb `to lead, guide,
conduct' (\citeproc{ref-ClarkHall1960}{Clark Hall 1960};
\citeproc{ref-BosworthToller1898}{Bosworth and Toller 1898}).

\subsection{Development to Old
English}\label{development-to-old-english-39}

From \Recon{láidijaną} `lead', monophthongization of \emph{*ai} first
gives a \emph{*lād-} stage. Later syncope, i-mutation, weak-tail
reduction, and loss of \emph{j} after a heavy stem yield \emph{lǣdan}
`lead', so the development represented here is fully regular
(\citeproc{ref-RingeTaylor2014}{Ringe and Taylor 2014}).

\section{\texorpdfstring{learn --- OE
\emph{leornian}}{learn --- OE leornian}}\label{learn-oe-leornian}

\index[iv]{02oe@\ivlangheader{Old English}{West Saxon normalization unmarked}!leornian@\iventry{leornian}{}}
\index[iv]{03pgmc@\ivlangheader{Proto-Germanic}{}!liznojana@\iventry{*líznōjaną}{}}

Derivation: \emph{*líznōjaną} \textgreater{} \emph{leornian} (regular).

\subsection{Derivation trace}\label{derivation-trace-45}

Proto input: \emph{*líznōjaną}

\begingroup
\setlength{\fboxsep}{6pt}

\noindent\fbox{%
\begin{minipage}{0.97\linewidth}
\footnotesize
\begin{tabular}{@{}>{\raggedright\arraybackslash}p{0.485\linewidth}@{\hspace{0.030\linewidth}}>{\raggedright\arraybackslash}p{0.485\linewidth}@{}}
\begin{minipage}[t]{\linewidth}
\centering\textbf{Earlier Germanic changes}\par
\vspace{0.35em}
\raggedright
\centering\textbf{Northwest and West Germanic}\par
\raggedright
\vspace{0.2em}
\begin{tabularx}{\linewidth}{@{}>{\raggedright\arraybackslash}X@{\hspace{0.55em}}>{\raggedright\arraybackslash}p{0.22\linewidth}@{\hspace{0.25em}}}
\mbox{PNWGmc I Lowering} & \emph{*léznōjaną} \\
\mbox{EAF Rhotacism} & \emph{*lérnōjaną} \\
\end{tabularx}
\vspace{0.6em}
\centering\textbf{Early Anglo-Frisian}\par
\raggedright
\vspace{0.2em}
\raggedright [no change]\par
\end{minipage}
&
\begin{minipage}[t]{\linewidth}
\centering\textbf{Old English changes}\par
\vspace{0.35em}
\raggedright
\begin{tabularx}{\linewidth}{@{}>{\raggedright\arraybackslash}X@{\hspace{0.55em}}>{\raggedright\arraybackslash}p{0.26\linewidth}@{\hspace{0.25em}}}
\mbox{OE Breaking} & \emph{*léornōjaną} \\
OE Heavy Syllable Nasal Apocope & \emph{*léornōjan} \\
\mbox{OE Secondary Nasalization} & \emph{*léornōjąn} \\
\mbox{OE I Umlaut} & \emph{*léornējąn} \\
\mbox{OE Med Long Short} & \emph{*léornejąn} \\
\mbox{OE Weak Tail Reduction} & \emph{*léornejan} \\
OE Intervocalic J Vocalization & \emph{*léorneian} \\
\mbox{OE Unstressed EI Contraction} & \emph{*léornian} \\
\end{tabularx}
\end{minipage}
\\
\end{tabular}
\end{minipage}%
} \endgroup

Old English form: \emph{leornian}

\subsection{Comparative reconstruction and the inherited
formation}\label{comparative-reconstruction-and-the-inherited-formation}

Old English \emph{leornian} `learn' belongs with Old High German
\emph{lernēn} `learn' beside \emph{lirnēn} `learn', and Old Saxon
\emph{līnon} `learn'. Kroonen reconstructs the weak verb
\emph{*liznōn-}; Orel likewise preserves a high-vowel comparative
formation. The reconstructed \emph{*z} is important both to the nasal
formation and to the later rhotic of English. Neither the dictionary
spelling of the English lemma nor its modern pronunciation establishes
an inherited \emph{*e} (\citeproc{ref-Kroonen2013}{Kroonen 2013, 340};
\citeproc{ref-Orel2003}{Orel 2003, 247}).

The selected infinitive
\Recon{líznōjaną}\index[iv]{03pgmc@\ivlangheader{Proto-Germanic}{}!liznojana@\iventry{*líznōjaną}{}}
`learn' retains that inherited root. Its extended weak-verb shape should
not be confused with a verbatim quotation of Kroonen's citation stem, or
with the distinct imperative and third-singular cells. The family
underwent conjugational remodelling; an English class II paradigm does
not by itself recover every Proto-Germanic ending. Ringe and Taylor
discuss the earlier formation separately from its later English
continuation (\citeproc{ref-RingeTaylor2014}{Ringe and Taylor 2014, 38,
84}).

\subsection{Attested vowels and selected
cells}\label{attested-vowels-and-selected-cells}

The infinitive selected here is independently attested. In Ælfric's
Treatise the clause beginning \emph{Gif þu nelt witan} continues
\emph{þu scealt leornian}: the modal requires an infinitive, so the
grammatical identification does not depend on a dictionary headword. The
edition prints this on p.71. The same volume gives the singular
imperative \emph{leorna} `learn!' in Deuteronomy14:23, followed by a
clause with the singular addressee \emph{þu}. These are textual
occurrences, not cells generated from an instructional paradigm
(\citeproc{ref-Crawford1922}{Crawford 1922, 71, 351}).

For the companion present indicative third singular, Miller's Bede
prints \emph{leornað} `learns' after \emph{swa hwylc swa} `whoever',
coordinated with \emph{rædeð} `reads'. The singular subject
disambiguates the ending, which in isolation could resemble a plural.
CAPR represents final dental-fricative spelling with thorn and therefore
selects \emph{leornaþ} `learns', without changing the root vowel, ending
vowel or grammatical cell (\citeproc{ref-Miller1891}{Miller 1891, 380},
ll. 26--27).

The alternative \emph{liornian} `learn' and finite \emph{liorna}
`learn!' remain important. Campbell's traditional account relates
\emph{io/eo} to breaking, later vowel developments and levelling. It
must not be compressed into a universal equation of \emph{liorn- with
conservative Northumbrian and }leorn- with secondary West Saxon: early
West Saxon also has the high-vowel spelling. A conservative spelling is
not necessarily an unchanged inherited vowel
(\citeproc{ref-Campbell1959}{Campbell 1959, 48, 59, 125}).

\subsection{Lowering, mutation and the weak
paradigm}\label{lowering-mutation-and-the-weak-paradigm}

The account adopted in \hyperref[rule-PNWGmcILowering]{SC015
PNWGmcILowering} lowers \emph{*i} before preconsonantal \emph{*z},
independently of the following vowel. For the selected infinitive, early
\emph{*lizn- > lezn-} is followed by rhotacism and English breaking
before \emph{*rn}. The resulting *leorn- is not obtained by selecting an
inherited mid-vowel reconstruction. Hill explains the competing
high-vowel forms through later i-mutation in extended stems and
subsequent paradigmatic history (\citeproc{ref-Hill2017}{Hill 2017,
153}, n.~25).

That explanation is plausible but not a recovered complete paradigm.
Jasanoff had already connected the family with the history of weak class
III. He considers a lost third-singular \emph{*lizniþ} as a possible
source of the high vowel of OHG \emph{lirnēn}, rather than requiring a
disputed first-singular reconstruction. The proposal matters because it
supplies a particular morphological source for mutation or levelling,
not a general appeal to analogy wherever the desired vowel is missing
(\citeproc{ref-Jasanoff1973}{Jasanoff 1973, 865, 868}).

The original distribution of suffixes, their extension and the direction
of later levelling remain uncertain. Later class II morphology cannot be
projected backwards merely to exclude every earlier high-front trigger;
conversely, one possible trigger does not establish that all three
selected cells once contained it. The modern discussion of the class III
family preserves these limitations (\citeproc{ref-RingeTaylor2014}{Ringe
and Taylor 2014, 38, 84}; \citeproc{ref-TarsiLaker2026}{Tarsi and Laker
2026, 42--43}).

\subsection{Interpretation of the selected
derivation}\label{interpretation-of-the-selected-derivation}

The three entries represent one etymological family: an infinitive, an
imperative singular and a present third singular, not three independent
proofs of a sound law. Their inherited inputs remain distinct, while the
selected \emph{leorn- vocalism receives a regular consonant-conditioned
history. Selection of these textually supported vowels does not disprove
the }liorn- tradition or resolve the whole weak-verb paradigm
(\citeproc{ref-Crawford1922}{Crawford 1922, 71, 351};
\citeproc{ref-Miller1891}{Miller 1891, 380};
\citeproc{ref-Jasanoff1973}{Jasanoff 1973, 865, 868}).

The finite entries retain an analogical family classification for the
conjugational remodelling, not because finite-cell selection is itself
analogy. No additional analogical replacement of the selected root vowel
is required. The distinction between regular phonology of the form
compared here and uncertain earlier paradigm history is essential to
interpreting the derivations (\citeproc{ref-Jasanoff1973}{Jasanoff 1973,
865, 868}; \citeproc{ref-TarsiLaker2026}{Tarsi and Laker 2026, 42--43}).

\section{\texorpdfstring{lid --- OE
\emph{hlid}}{lid --- OE hlid}}\label{lid-oe-hlid}

\index[iv]{02oe@\ivlangheader{Old English}{West Saxon normalization unmarked}!hlid@\iventry{hlid}{}}
\index[iv]{03pgmc@\ivlangheader{Proto-Germanic}{}!xlida@\iventry{*xlídą}{}}
\index[iv]{09on@\ivlangheader{Old Norse}{}!hlitho@\iventry{hliþó}{}}

Derivation: \emph{*xlídą} \textgreater{} \emph{hlid} (regular).

\subsection{Derivation trace}\label{derivation-trace-46}

Proto input: \emph{*xlídą}

\begingroup
\setlength{\fboxsep}{6pt}

\noindent\fbox{%
\begin{minipage}{0.97\linewidth}
\small
\begin{tabular}{@{}>{\raggedright\arraybackslash}p{0.359\linewidth}@{\hspace{0.030\linewidth}}>{\raggedright\arraybackslash}p{0.611\linewidth}@{}}
\begin{minipage}[t]{\linewidth}
\centering\textbf{Earlier Germanic changes}\par
\vspace{0.35em}
\raggedright
\centering\textbf{Northwest and West Germanic}\par
\raggedright
\vspace{0.2em}
\raggedright [no change]\par
\vspace{0.6em}
\centering\textbf{Early Anglo-Frisian}\par
\raggedright
\vspace{0.2em}
\raggedright [no change]\par
\end{minipage}
&
\begin{minipage}[t]{\linewidth}
\centering\textbf{Old English changes}\par
\vspace{0.35em}
\raggedright
\begin{tabularx}{\linewidth}{@{}>{\raggedright\arraybackslash}X@{\hspace{0.55em}}>{\raggedright\arraybackslash}p{0.14\linewidth}@{\hspace{0.25em}}}
OE Heavy Syllable Nasal Apocope & \emph{*xlíd} \\
\end{tabularx}
\end{minipage}
\\
\end{tabular}
\end{minipage}%
} \endgroup

Old English form: \emph{hlid}

\subsection{Reconstruction and comparative
evidence}\label{reconstruction-and-comparative-evidence-40}

Orel cites a neuter lexeme of the \emph{*xliđ-} type with Old English
{\emph{hlid}}\index[iv]{02oe@\ivlangheader{Old English}{West Saxon normalization unmarked}!hlid@\iventry{hlid}{}}
`lid, cover', and Lloyd includes OE \emph{hlid} `lid' beside ON
\emph{hliþó} `lid' and OHG \emph{(h)lit} `lid' among forms that retain
\emph{i} (\citeproc{ref-Orel2003}{Orel 2003};
\citeproc{ref-Lloyd1966}{Lloyd 1966}).

\subsection{Old English evidence}\label{old-english-evidence-40}

Clark Hall and Bosworth-Toller record
{\emph{hlid}}\index[iv]{02oe@\ivlangheader{Old English}{West Saxon normalization unmarked}!hlid@\iventry{hlid}{}}
`lid, cover' as the noun `lid, cover, door, gate'
(\citeproc{ref-ClarkHall1960}{Clark Hall 1960};
\citeproc{ref-BosworthToller1898}{Bosworth and Toller 1898, 563}).

\subsection{Development to Old
English}\label{development-to-old-english-40}

The derivational input already represents the later Germanic
\emph{hliđ-} stage used for the derivation here. From \Recon{xlídą}
`lid', heavy-syllable apocope yields \emph{hlid} `lid', and the form
belongs to the retained-\emph{i} set noted by Lloyd rather than to the
lowered \emph{e} type (\citeproc{ref-Lloyd1966}{Lloyd 1966}).

\subsection{Form note}\label{form-note-11}

An earlier etymological stage \Recon{liþuz} `lid' belongs to comparative
background only. The form represented here is the later \Recon{xlídą}
`lid' \textgreater{} \emph{hlid} `lid' line that matches the attested
Old English noun (\citeproc{ref-Orel2003}{Orel 2003};
\citeproc{ref-Lloyd1966}{Lloyd 1966}).

\section{\texorpdfstring{light --- OE
\emph{līehtan}}{light --- OE līehtan}}\label{light-oe-lux12behtan}

\index[iv]{02oe@\ivlangheader{Old English}{West Saxon normalization unmarked}!liehtan@\iventry{līehtan}{}}
\index[iv]{03pgmc@\ivlangheader{Proto-Germanic}{}!leuxtijana@\iventry{*léuxtijaną}{}}

Derivation: \emph{*léuxtijaną} \textgreater{} \emph{līehtan} (regular).

\subsection{Derivation trace}\label{derivation-trace-47}

Proto input: \emph{*léuxtijaną}

\begingroup
\setlength{\fboxsep}{6pt}

\noindent\fbox{%
\begin{minipage}{0.97\linewidth}
\footnotesize
\begin{tabular}{@{}>{\raggedright\arraybackslash}p{0.339\linewidth}@{\hspace{0.030\linewidth}}>{\raggedright\arraybackslash}p{0.630\linewidth}@{}}
\begin{minipage}[t]{\linewidth}
\centering\textbf{Earlier Germanic changes}\par
\vspace{0.35em}
\raggedright
\centering\textbf{Northwest and West Germanic}\par
\raggedright
\vspace{0.2em}
\raggedright [no change]\par
\vspace{0.6em}
\centering\textbf{Early Anglo-Frisian}\par
\raggedright
\vspace{0.2em}
\raggedright [no change]\par
\end{minipage}
&
\begin{minipage}[t]{\linewidth}
\centering\textbf{Old English changes}\par
\vspace{0.35em}
\raggedright
\begin{tabularx}{\linewidth}{@{}>{\raggedright\arraybackslash}X@{\hspace{0.55em}}>{\raggedright\arraybackslash}p{0.26\linewidth}@{\hspace{0.25em}}}
\mbox{OE Diphthong Leveling} & \emph{*lēoxtijaną} \\
OE Heavy Syllable Nasal Apocope & \emph{*lēoxtijan} \\
\mbox{OE Secondary Nasalization} & \emph{*lēoxtijąn} \\
\mbox{Sievers Law Syncope} & \emph{*lēoxtjąn} \\
\mbox{OE I Umlaut} & \emph{*līextjąn} \\
\mbox{OE Weak Tail Reduction} & \emph{*līextjan} \\
\mbox{OE J Loss After Heavy} & \emph{*līextan} \\
\end{tabularx}
\end{minipage}
\\
\end{tabular}
\end{minipage}%
} \endgroup

Old English form: \emph{līehtan}

\subsection{Reconstruction and comparative
evidence}\label{reconstruction-and-comparative-evidence-41}

Fulk gives Proto-Germanic \Recon{liuxtijanan} `light' with Old English
\emph{līehtan} `illuminate', and Ringe and Taylor likewise derive West
Saxon
{\emph{liehtan}}\index[iv]{02oe@\ivlangheader{Old English}{West Saxon normalization unmarked}!liehtan@\iventry{liehtan}{}}
`light, illuminate' from the same weak-verb formation
(\citeproc{ref-Fulk2018}{Fulk 2018};
\citeproc{ref-RingeTaylor2014}{Ringe and Taylor 2014}).

\subsection{Old English evidence}\label{old-english-evidence-41}

Clark Hall and Bosworth-Toller preserve the verb family under spellings
such as \emph{liehtan} `illuminate',
{\emph{lihtan}}\index[iv]{02oe@\ivlangheader{Old English}{West Saxon normalization unmarked}!lihtan@\iventry{lihtan}{}}
`light, illuminate', and \emph{līhtan} `illuminate', distinct from the
related noun \emph{lēoht} `light' and adjective \emph{leoht} /
\emph{liht} `light' (\citeproc{ref-ClarkHall1960}{Clark Hall 1960};
\citeproc{ref-BosworthToller1898}{Bosworth and Toller 1898}).

\subsection{Development to Old
English}\label{development-to-old-english-41}

From \Recon{léuxtijaną} `light', the regular verbal line preserves
\emph{*xt}, passes through a West Saxon \emph{liehtan} `illuminate'
stage, and is represented here by normalized \emph{līehtan}
`illuminate'. The word treated in this entry is therefore the verb `to
light, illuminate', not the related noun from \Recon{leuxtą} `light'
(\citeproc{ref-Fulk2018}{Fulk 2018};
\citeproc{ref-RingeTaylor2014}{Ringe and Taylor 2014}).

\subsection{Dialect note}\label{dialect-note-3}

Ringe and Taylor and Campbell distinguish West Saxon
{\emph{liehtan}}\index[iv]{02oe@\ivlangheader{Old English}{West Saxon normalization unmarked}!liehtan@\iventry{liehtan}{}}
`light, illuminate' from Anglian
{\emph{lihtan}}\index[iv]{02oe@\ivlangheader{Old English}{West Saxon normalization unmarked}!lihtan~04@\iventry{lihtan}{Angl.}}
`light, illuminate', while later West Saxon also shows
{\emph{lyhtan}}\index[iv]{02oe@\ivlangheader{Old English}{West Saxon normalization unmarked}!lyhtan~03@\iventry{lyhtan}{LWS}}
`light, illuminate' (\citeproc{ref-RingeTaylor2014}{Ringe and Taylor
2014}; \citeproc{ref-Campbell1959}{Campbell 1959}).

\section{\texorpdfstring{linden --- OE
\emph{lind}}{linden --- OE lind}}\label{linden-oe-lind}

\index[iv]{02oe@\ivlangheader{Old English}{West Saxon normalization unmarked}!lind@\iventry{lind}{}}
\index[iv]{03pgmc@\ivlangheader{Proto-Germanic}{}!lindo@\iventry{*líndō}{}}

Derivation: \emph{*líndō} \textgreater{} \emph{lind} (regular).

\subsection{Derivation trace}\label{derivation-trace-48}

Proto input: \emph{*líndō}

\begingroup
\setlength{\fboxsep}{6pt}

\noindent\fbox{%
\begin{minipage}{0.97\linewidth}
\small
\begin{tabular}{@{}>{\raggedright\arraybackslash}p{0.611\linewidth}@{\hspace{0.030\linewidth}}>{\raggedright\arraybackslash}p{0.359\linewidth}@{}}
\begin{minipage}[t]{\linewidth}
\centering\textbf{Earlier Germanic changes}\par
\vspace{0.35em}
\raggedright
\centering\textbf{Northwest and West Germanic}\par
\raggedright
\vspace{0.2em}
\begin{tabularx}{\linewidth}{@{}>{\raggedright\arraybackslash}X@{\hspace{0.55em}}>{\raggedright\arraybackslash}p{0.14\linewidth}@{\hspace{0.25em}}}
\mbox{PNWGmc Final Long O Raising} & \emph{*líndu} \\
\end{tabularx}
\vspace{0.6em}
\centering\textbf{Early Anglo-Frisian}\par
\raggedright
\vspace{0.2em}
\raggedright [no change]\par
\end{minipage}
&
\begin{minipage}[t]{\linewidth}
\centering\textbf{Old English changes}\par
\vspace{0.35em}
\raggedright
\begin{tabularx}{\linewidth}{@{}>{\raggedright\arraybackslash}X@{\hspace{0.55em}}>{\raggedright\arraybackslash}p{0.14\linewidth}@{\hspace{0.25em}}}
\mbox{OE High Vowel Apocope} & \emph{*línd} \\
\end{tabularx}
\end{minipage}
\\
\end{tabular}
\end{minipage}%
} \endgroup

Old English form: \emph{lind}

\subsection{Reconstruction and comparative
evidence}\label{reconstruction-and-comparative-evidence-42}

Kroonen cites \emph{*lindō-} `lime tree' and gives Old English
{\emph{lind}}\index[iv]{02oe@\ivlangheader{Old English}{West Saxon normalization unmarked}!lind@\iventry{lind}{}}
`linden, lime tree' as the relevant reflex
(\citeproc{ref-Kroonen2013}{Kroonen 2013}).

\subsection{Old English evidence}\label{old-english-evidence-42}

Clark Hall and Bosworth-Toller both record
{\emph{lind}}\index[iv]{02oe@\ivlangheader{Old English}{West Saxon normalization unmarked}!lind@\iventry{lind}{}}
`linden, lime tree' as the noun `lime-tree, linden'
(\citeproc{ref-ClarkHall1960}{Clark Hall 1960};
\citeproc{ref-BosworthToller1898}{Bosworth and Toller 1898, 630}).

\subsection{Development to Old
English}\label{development-to-old-english-42}

From \Recon{líndō} `linden', Northwest Germanic final \emph{*ō} raising
first gives a \Recon{líndu} `linden' stage, and later high-vowel apocope
yields
{\emph{lind}}\index[iv]{02oe@\ivlangheader{Old English}{West Saxon normalization unmarked}!lind@\iventry{lind}{}}
`linden, lime tree'. The development is therefore straightforward and
regular.

\subsection{Form note}\label{form-note-12}

The Old English noun represented here is
{\emph{lind}}\index[iv]{02oe@\ivlangheader{Old English}{West Saxon normalization unmarked}!lind@\iventry{lind}{}}
`linden, lime tree'. Clark Hall also has a separate adjectival
\emph{linden} `made of linden-wood', but that is not the noun
counterpart for this entry (\citeproc{ref-ClarkHall1960}{Clark Hall
1960}).

\section{\texorpdfstring{little --- OE
\emph{lȳtelu}}{little --- OE lȳtelu}}\label{little-oe-lux233telu}

\index[iv]{02oe@\ivlangheader{Old English}{West Saxon normalization unmarked}!hlutor@\iventry{hlūtor}{}}
\index[iv]{02oe@\ivlangheader{Old English}{West Saxon normalization unmarked}!lutan@\iventry{lūtan}{}}
\index[iv]{02oe@\ivlangheader{Old English}{West Saxon normalization unmarked}!lytelu@\iventry{lȳtelu}{}}
\index[iv]{03pgmc@\ivlangheader{Proto-Germanic}{}!lutilaz@\iventry{*lūtilaz}{}}
\index[iv]{03pgmc@\ivlangheader{Proto-Germanic}{}!lutilo@\iventry{*lūtilō}{}}

Derivation: citation reconstruction \emph{*lūtilaz}; form followed here
\emph{*lūtilō} \textgreater{} \emph{lȳtelu} (regular).

\subsection{Derivation trace}\label{derivation-trace-49}

Proto input: \emph{*lūtilō}

\begingroup
\setlength{\fboxsep}{6pt}

\noindent\fbox{%
\begin{minipage}{0.97\linewidth}
\small
\begin{tabular}{@{}>{\raggedright\arraybackslash}p{0.485\linewidth}@{\hspace{0.030\linewidth}}>{\raggedright\arraybackslash}p{0.485\linewidth}@{}}
\begin{minipage}[t]{\linewidth}
\centering\textbf{Earlier Germanic changes}\par
\vspace{0.35em}
\raggedright
\centering\textbf{Northwest and West Germanic}\par
\raggedright
\vspace{0.2em}
\begin{tabularx}{\linewidth}{@{}>{\raggedright\arraybackslash}X@{\hspace{0.55em}}>{\raggedright\arraybackslash}p{0.14\linewidth}@{\hspace{0.25em}}}
\mbox{PNWGmc Final Long O Raising} & \emph{*lūtilu} \\
\end{tabularx}
\vspace{0.6em}
\centering\textbf{Early Anglo-Frisian}\par
\raggedright
\vspace{0.2em}
\raggedright [no change]\par
\end{minipage}
&
\begin{minipage}[t]{\linewidth}
\centering\textbf{Old English changes}\par
\vspace{0.35em}
\raggedright
\begin{tabularx}{\linewidth}{@{}>{\raggedright\arraybackslash}X@{\hspace{0.55em}}>{\raggedright\arraybackslash}p{0.14\linewidth}@{\hspace{0.25em}}}
\mbox{OE I Umlaut} & \emph{*lȳtilu} \\
OE Med Unstressed I Lowering1 & \emph{*lȳtelu} \\
\end{tabularx}
\end{minipage}
\\
\end{tabular}
\end{minipage}%
} \endgroup

Old English form: \emph{lȳtelu}

\subsection{Old English evidence}\label{old-english-evidence-43}

The Vespasian Psalter gloss has
{\emph{lytelu}}\index[iv]{02oe@\ivlangheader{Old English}{West Saxon normalization unmarked}!lytelu@\iventry{lytelu}{}}
`little' in \emph{netenu lytelu 7 micel} `small and great animals',
glossing Latin \emph{animalia pusilla et magna} at Psalm 103:25
(\citeproc{ref-Sweet1885}{Sweet 1885, 336}). The adjective agrees with
the neuter plural noun `animals': this is the strong neuter plural, not
the formally identical feminine nominative singular. CAPR selects the
shared nominative--accusative plural cell; the nominative interpretation
belongs to this clause's enumeration of the creatures there.

The normalized comparator
{\emph{lȳtelu}}\index[iv]{02oe@\ivlangheader{Old English}{West Saxon normalization unmarked}!lytelu@\iventry{lȳtelu}{}}
marks quantity that Sweet's diplomatic spelling does not mark. Long
\emph{ȳ} and the grammatical identification are independently supported
by Fulk's West Mercian discussion and Goering's comparison
(\citeproc{ref-Fulk2010}{Fulk 2010, 133--34};
\citeproc{ref-Goering2016}{Goering 2016, 176--77}). Fulk contrasts the
retained medial vowel here with masculine accusative plural
{\emph{lȳtle}}\index[iv]{02oe@\ivlangheader{Old English}{West Saxon normalization unmarked}!lytle@\iventry{lȳtle}{}}
`little'. The forms differ in the quantity of the following syllable at
deletion, not in whether a heavy root happens to permit an exception.

\subsection{Comparative evidence and the reconstruction
problem}\label{comparative-evidence-and-the-reconstruction-problem}

The Germanic evidence does not establish a single, unproblematic vowel
correspondence. Three principal groups must be distinguished:

{\def\LTcaptype{none} % do not increment counter
\begin{longtable}[]{@{}
  >{\raggedright\arraybackslash}p{(\linewidth - 4\tabcolsep) * \real{0.3333}}
  >{\raggedright\arraybackslash}p{(\linewidth - 4\tabcolsep) * \real{0.3333}}
  >{\raggedright\arraybackslash}p{(\linewidth - 4\tabcolsep) * \real{0.3333}}@{}}
\toprule\noalign{}
\begin{minipage}[b]{\linewidth}\raggedright
Evidence
\end{minipage} & \begin{minipage}[b]{\linewidth}\raggedright
Quantity and formation
\end{minipage} & \begin{minipage}[b]{\linewidth}\raggedright
Comparative implication
\end{minipage} \\
\midrule\noalign{}
\endhead
\bottomrule\noalign{}
\endlastfoot
Gothic \emph{leitils}, Old Norse \emph{lítill} `little' & long front
vowel; the Gothic spelling \emph{ei} represents \emph{ī} & supports
\emph{*lītila-}, with nominative masculine \emph{*lītilaz} \\
OE \emph{lȳtel} `little', including the selected \emph{lȳtelu} & long
rounded front vowel, regularly from \emph{*ū} before an umlauting vowel
& supports a back-vowel antecedent \emph{*lūtila-}, not
\emph{*lītila-} \\
OE \emph{lyttel} `little', Old Saxon \emph{luttil}, Old High German
\emph{luzzil} & short root vowel with a geminate or its High German
shifted continuation & requires the history of the short-vowel/geminate
formation to be explained separately \\
\end{longtable}
}

These are comparative dictionary forms, not three identical inflectional
cells or three contemporary manuscript spellings. The Gothic and Norse
forms are masculine citation forms; the directly identified OE plural
has its own evidence above. The grouping follows the dictionaries and
Fulk's explicit contrast, not an assumption that all three groups
descend regularly from one complete word (\citeproc{ref-Orel2003}{Orel
2003, 248}; \citeproc{ref-Kroonen2013}{Kroonen 2013, 339};
\citeproc{ref-Fulk2018}{Fulk 2018, 15, 127}).

The distribution is also less tidy than a division into Gothic--Norse
and West Germanic might suggest. Heidermanns records the rarer OE
\emph{lītel} beside \emph{lȳtel}, Middle Dutch \emph{littel} beside
\emph{luttel}, Old Frisian \emph{lītik} beside \emph{littik}, and Old
Saxon and Middle Dutch \emph{luttik}, with a different suffix. His Dutch
evidence for \emph{littel} and \emph{luttel} begins in Middle Dutch, not
in an attested Proto-Germanic or Old Dutch paradigm. He also records Old
High German \emph{liuzzil} and \emph{liuzilu}, together with the
abstract \emph{liuzila}, and reports Matzel's analysis of these as
mixtures of the front-vowel and geminate stems. These spellings cannot
simply be assigned to an otherwise uniform short-\emph{u} series. Nor is
written High German \emph{zz} an unmediated transcription of Germanic
\emph{*tt}: Heidermanns specifically identifies Notker's \emph{zz} here
as an affricate (\citeproc{ref-Heidermanns1993}{Heidermanns 1993,
371--72, 390--91}). Lühr additionally lists Old Saxon \emph{lītil}, but
identifies no textual witness in that passage; Heidermanns'
corresponding entry does not list this Old Saxon form. Its evidential
status should therefore remain distinct from that of the directly
identified Mercian plural (\citeproc{ref-Luehr2017}{Lühr 2017, 153};
\citeproc{ref-Heidermanns1993}{Heidermanns 1993, 371, 390}).

Orel prints the alternatives
\Recon{lītilaz}\index[iv]{03pgmc@\ivlangheader{Proto-Germanic}{}!litilaz@\iventry{*lītilaz}{}},
\Recon{lūtilaz}\index[iv]{03pgmc@\ivlangheader{Proto-Germanic}{}!lutilaz@\iventry{*lūtilaz}{}}
and
\Recon{luttilaz}\index[iv]{03pgmc@\ivlangheader{Proto-Germanic}{}!luttilaz@\iventry{*luttilaz}{}}
`little' (\citeproc{ref-Orel2003}{Orel 2003, 248}). The long-\emph{ū}
alternative is the comparative citation retained here. It is not a vowel
invented to make the transducer work: the ordinary fronting of inherited
\emph{*ū} before \emph{*i} accounts for the independently established OE
\emph{ȳ}.

The three alternatives register a comparative problem; the tildes do not
supply derivations connecting them. Orel derives the family from his
\emph{*luttiz ~ *luttjaz}, but this cross-reference alone does not
explain how the long front vowel would arise from that base
(\citeproc{ref-Orel2003}{Orel 2003, 248}).

Kroonen instead heads the entry with
\Recon{lītila-}\index[iv]{03pgmc@\ivlangheader{Proto-Germanic}{}!litila@\iventry{*lītila-}{}}
and explicitly leaves the distribution of long-\emph{ī}, long-\emph{ū}
and short-\emph{u}/geminate variants unexplained
(\citeproc{ref-Kroonen2013}{Kroonen 2013, 339}). His entry is not quoted
as a long-\emph{ū} headword. His choice gives a representative entry for
the Germanic adjective family, with a front-vowel reconstruction
directly supported by Gothic and Norse. It is not a demonstrated common
phonological ancestor of every reflex listed beneath it: the
accompanying argument expressly withholds that solution. Indeed, Kroonen
rejects the comparison with `stoop' and Irish `little finger' as a
satisfactory etymology of the whole family precisely because it leaves
the front-vowel variant unaccounted for. Selecting Orel's supported
back-vowel alternative does not settle that objection
(\citeproc{ref-Kroonen2013}{Kroonen 2013, 339}).

\subsection{Antonymic formation and the older
etymologies}\label{antonymic-formation-and-the-older-etymologies}

Meillet already denied that Gothic \emph{leitils} and Old High German
\emph{luzzil} could simply be identified with one another. In a
discussion of the unstable naming of `small', he considered neither
adjective securely etymologized. What he did identify was the influence
of the opposite adjective `great': Gothic \emph{leitils : mikils}, OE
\emph{lȳtel : micel} `little : great', and the corresponding High German
and Norse pairs. His Norse comparison extends to the masculine
accusatives \emph{lítenn} and \emph{mikenn} and to shorter adverbial
formations. The claim concerns characteristic endings and morphological
association, not a sound change between the two root vowels
(\citeproc{ref-Meillet1902}{Meillet 1902, 172--74}, especially p.~173).

The distinction matters for the later literature. Melazzo reports
Thurneysen's analysis of a suffix \emph{-lo-} shared by Gothic
\emph{mikils}, \emph{leitils} and High German \emph{luzzil}, within a
survey of the Greek adjective μέγας, μεγάλη, μέγα `great'. His note
explicitly distinguishes the root vowels of the Germanic `little' forms.
This is useful evidence for the history of suffixal analysis, not a new
etymology of `little', and Thurneysen's original discussion has not been
independently consulted here (\citeproc{ref-Melazzo2014}{Melazzo 2014,
10}, n.~3). Shared suffixes and the existence of shorter formations give
the antonymic comparison a morphological basis; they do not by
themselves establish the direction or every step of the proposed
remodelling.

Krogmann's dedicated study attempts something different: a remote
relationship between two initially distinct root families. He begins
with Germanic \emph{*līt-} and \emph{*leut-}, explicitly not immediately
reconcilable, and associates both `small' and `deceitful' with bending,
crouching and lying concealed. Norse \emph{lúta} and OE \emph{lūtan}
`bend down' anchor the back-vowel comparison. His proposed external
relatives include Baltic words for dejection and Celtic words for
smallness; the transitions from `bent down' to `small' and from `lying
concealed' to `deceiving' supply the semantic argument
(\citeproc{ref-Krogmann1935}{Krogmann 1935, 44--45}).

He does not simply accept every earlier comparison. On pp.~45--46 he
rejects both the Greek `revile/play' comparison and an association with
Latvian `creep', assigning the latter to the family of `slide'. For the
front-vowel group he instead invokes Indo-European \emph{*lei-} `crouch,
diminish', with words for ceasing, weakness and hunger; for the
back-vowel group he uses \emph{*leud-} `bend'. Rejecting an unsupported
common \emph{*lē-} `bend', he ultimately derives the two from the
\emph{i}- and \emph{u}-bases \emph{*elei-} and \emph{*eleu-} of
\emph{*el-} `bend', each with a dental extension. His conclusion is a
distant relationship, not the regular descent of one Germanic vowel from
the other (\citeproc{ref-Krogmann1935}{Krogmann 1935, 46--48}). The
proposed root extensions and the wide semantic grouping carry that
conclusion; the paper does not provide a complete conditioned history of
the three adjective stems, their vowel quantities and their suffixes.

Heidermanns separates those dimensions more explicitly: the root vowels,
single versus geminate dental, and the suffixes represented by the
Norse, Frisian and West Germanic formations. His explanation of
adjective formation again uses `great' as a model, with Gothic
\emph{leita} `a little' and OE \emph{lȳt}, Old Saxon \emph{lūt} `little,
few' as possible shorter bases. He retains the transmitted Gothic adverb
against its conjectural expansion to an adjective. For the
short-vowel/geminate variants he invokes expressive gemination; for the
deeper etymology he offers only a tentative connection with `crouch' and
acknowledges the lack of close comparisons at the same formation level
(\citeproc{ref-Heidermanns1993}{Heidermanns 1993, 372}). That is a
documented proposal, not a specified, exceptionless phonological
environment. It cannot be adopted as a sound law merely because the word
means `small'. His separate treatment of High German \emph{luzzi} also
leaves open whether its geminate was inherited or arose on transfer to
the \emph{ja}-class (\citeproc{ref-Heidermanns1993}{Heidermanns 1993,
390--91}).

\subsection{Two roots, or a more remote
relationship?}\label{two-roots-or-a-more-remote-relationship}

Lühr's alternative is not a new way to turn \emph{ī} into \emph{ū}. She
explicitly denies cognacy between the front- and back-vowel groups. In
her 2017 account the latter belong with `bend, crouch': OE \emph{lūtan},
preterite \emph{lēat}, and Norse \emph{lúta}, preterite \emph{laut}. She
compares Lithuanian \emph{liūdnas} `sad, dejected' with \emph{*lud-no-},
Welsh \emph{llwdn} `young animal' with Celtic \emph{*lud-no-}, and Welsh
\emph{lludd} `weariness' with \emph{*loud-o-}. The Germanic short-vowel
formation is \emph{*lut-ja-la-}, with a glide-bearing stem before the
added \emph{l}-suffix; OE \emph{lȳtel} instead presupposes
\emph{*lūtila-}. These are distinct formations and grades, not
interchangeable spellings (\citeproc{ref-Luehr2017}{Lühr 2017,
152--53}).

For the front-vowel group she proposes Indo-European \emph{*leih₂-}
`cease, diminish', with Gothic dative singular \emph{undarleijin} `one
who is still less than the least' and Greek λῖμος `hunger' as
comparisons. The proposed Germanic base is \emph{*lī-ja-}, not an
inherited whole adjective \emph{*lītila-}. Remodelling after
\emph{*lūtila-} could supply the latter's extended formation; influence
from \emph{*mekila-} `great' is another possibility
(\citeproc{ref-Luehr2017}{Lühr 2017, 153}). A zero-grade root with
\emph{*i} before a laryngeal makes long \emph{ī} phonologically
intelligible. It does not independently establish the membership of the
Gothic comparison or prove the proposed suffixal replacement. Lühr
herself records a competing etymology of \emph{undarleijin}, linking it
instead with High German \emph{liz} `pretext' and the Gothic noun for
`hypocrisy' (\citeproc{ref-Luehr2017}{Lühr 2017, 153}, n.~4).

The earlier Hildebrandlied commentary already contains the two-root
analysis and the proposed remodelling. It is considerably fuller: it
considers an original \emph{i}-stem, transfer to a \emph{ja}-formation
to account for the West Germanic geminate, and subsequent addition of
the \emph{l}-suffix. It explicitly objects that direct gemination before
\emph{l} has not been demonstrated for the required antecedent. For the
front-vowel formation, it also supplies a paradigmatic motive: loss of
the glide before an inflectional \emph{i} could expose a stem
\emph{*lī-}, to which reanalysed \emph{*-ti-} or \emph{*-tila-} could be
attached (\citeproc{ref-Luehr1982}{Lühr 1982}, II, pp.~382--387 in the
author's reissue). The later article therefore restates the central
separation and remodelling hypothesis, while presenting a more compact
antonymic account and an explicitly indexed laryngeal. It is not an
independent second demonstration of the same history.

There are unresolved premises even on this account. The earlier
commentary discusses analogical explanations of the long vowel in the
associated `bend' verb; short \emph{*u} in \emph{*lud-} does not
automatically produce long \emph{*ū}. The Baltic and Celtic comparisons
motivate a semantic family but do not alone establish the complete
long-vowel adjective (\citeproc{ref-Luehr1982}{Lühr 1982}, II, p.~383,
n.~2384 in the author's reissue; \citeproc{ref-Luehr2017}{Lühr 2017,
152--53}). Two roots avoid an unsupported \emph{ī}--\emph{ū} sound
correspondence, but their particular identification and the proposed
morphological convergence still require evidence.

Levitsky offers a different remote unification. His methodological
discussion expressly treats the relation of \emph{i}- and
\emph{u}-extended roots under Persson's account of root extensions and
what he calls anomalous ablaut, rather than ordinary \emph{e/o/zero}
ablaut (\citeproc{ref-Levitsky2003}{Levickij 2003, 37--38}). In the
discussion of `little', he connects \emph{*lei-} and \emph{*leu-} with a
wider \emph{*el-} family. Against Krogmann's `bent' to `small', he
prefers `cut, break into pieces, reduce to small pieces' to `small',
within a broader proposed semantic association of cutting, binding and
bending (\citeproc{ref-Levitsky2003}{Levickij 2003, 58--60}, especially
p.~59).

That offers a genuine alternative to the crouching semantics and a
hypothesis about the prehistory of the root families. It does not supply
a conditioned Germanic \emph{*ī > *ū}, or an independently demonstrated
derivation of all three adjective stems. The parallel root extensions
and semantic transitions would have to be established independently;
grouping words under a sufficiently broad semantic prototype cannot
itself establish the required sound correspondences. The proposal
therefore remains a remote-etymological possibility, not the
phonological basis of the selected OE derivation.

\subsection{West Germanic shortening and
gemination}\label{west-germanic-shortening-and-gemination}

The long-back-vowel versus short-vowel/geminate problem must not be
conflated with the front- versus back-vowel problem. Fulk discusses OE
\emph{hlūtor : hluttor} `clear' beside \emph{lȳtel : lyttel}, Old Saxon
\emph{hluttar}, \emph{luttil} and High German \emph{hluttar},
\emph{luzzil}. Unlike the much disputed gemination after long vowels
before \emph{j}, gemination before liquids has substantial evidence
outside Upper German. His proposed account uses shortening in a closed
syllable, as in \emph{*xlūtr}, and extension of the resulting stem
through the paradigm; the liquids, unlike \emph{j}, were not made
syllabic under Sievers' law (\citeproc{ref-Fulk2018}{Fulk 2018, 127},
§6.15).

The account has phonological conditioning and a comparison outside
`little'; it is not permission to shorten any long vowel and double any
consonant. Fulk's further discussion reconstructs alternation of
\emph{hlūtor} with \emph{*hluttr-}, followed by mixture of stems, and
distinguishes later analogical alternations and new clusters created by
syncope (\citeproc{ref-Fulk2018}{Fulk 2018, 132}, §6.18). Applied to
`little', the account requires independently justified liquid adjacency,
syllabification and subsequent paradigmatic distribution. It competes
with Lühr's glide-bearing formation, whose geminate would arise in a
different environment, and is more specific than Heidermanns' expressive
explanation. None of these accounts, even if it resolves the back-vowel
alternation, explains the Germanic front-vowel series. No new general
shortening or gemination rule is established here.

\subsection{Assessment and selected
cell}\label{assessment-and-selected-cell}

For the executable OE history, the long-back-vowel input remains the
defensible choice: it is explicitly among Orel's reconstructed
alternatives and gives the independently attested long \emph{ȳ} by
ordinary umlaut. Its placement at Proto-Germanic is a comparative
reconstruction adopted for this branch, not proof that every Germanic
language inherited the same adjective. Heidermanns restricts the
back-vowel evidence to West Germanic, and the Gothic--Norse front-vowel
evidence does not independently establish a common long-\emph{ū}
ancestor (\citeproc{ref-Orel2003}{Orel 2003, 248};
\citeproc{ref-Heidermanns1993}{Heidermanns 1993, 372, 390}). The current
citation and input are therefore retained, not silently corrected from a
front-vowel form. `Regular' describes the derivational input-to-OE
history; it does not certify a resolved Indo-European etymology or a
unified history of every member of the family.

The selected
\Recon{lūtilō}\index[iv]{03pgmc@\ivlangheader{Proto-Germanic}{}!lutilo@\iventry{*lūtilō}{}}
combines that independently supported stem with the ordinary strong
neuter nominative--accusative plural ending \emph{*-ō}. Fulk explicitly
reconstructs this ending in the strong-adjective paradigm
(\citeproc{ref-Fulk2018}{Fulk 2018, 209}, §9.2). He also prints the
complete \emph{*lūtilō} as the predecessor of Mercian \emph{lȳtelu}
(\citeproc{ref-Fulk2018}{Fulk 2018, 92}, §5.6). This is a stated
derivation, not a merely illustrative ``form like'' the selected
reconstruction. The independent stem and paradigm evidence remains
important: his printed plural does not resolve Kroonen's uncertainty
about the wider root variants. Citation and selected form are two
Proto-Germanic cells, not different historical stages; no transponent,
analogical replacement or artificial accent is used.

\subsection{Development and discriminatory
value}\label{development-and-discriminatory-value}

\hyperref[rule-PNWGmcFinalLongORaising]{SC019 PNWGmcFinalLongORaising}
gives \emph{*lūtilu}; \hyperref[rule-OEIUmlaut]{SC055 OEIUmlaut} gives
\emph{*lȳtilu}. The heavy first syllable is followed by two light
syllables. Both weak high vowels therefore survive
\hyperref[rule-OEHighVowelApocope]{SC063 OEHighVowelApocope}, the
unified quantitative high-vowel-deletion event. The same HLL retention
is the evidential point of the conservative plurals of `head' and
`livestock' (\citeproc{ref-Fulk2010}{Fulk 2010, 133--38};
\citeproc{ref-Goering2016}{Goering 2016, 176--77, 186--88};
\citeproc{ref-Goering2023}{2023, 60--63, 67--68}).

The following identity trace at \hyperref[rule-OEMedialSyncope]{SC065
OEMedialSyncope} makes no second deletion. The short-root condition of
\hyperref[rule-OELAdjacentSyncope]{SC066 OELAdjacentSyncope} likewise
does not reach this long-root form. Later
\hyperref[rule-OEMedUnstressedILowering1]{SC074
OEMedUnstressedILowering1} gives \emph{*lȳtelu}. The complete cascade
thus reaches the attested comparator without an analogical restoration.

An indiscriminate rule deleting a medial high vowel after every heavy
root instead yields \Pred{lȳtlu}. That wrong alternative changes none of
the previously selected histories, but fails on this plural. An
overgeneralized final apocope in HLL instead gives \Pred{lȳtel},
likewise without changing any previously selected history. Thus the
attested plural independently protects both its medial and final
retained high vowels. The word is therefore a positive historical
witness of retention and a genuine nonapplication control for deletion.
It does not replace the selected singular history of `head', nor does it
imply that every form of `little' should retain its medial vowel.

\section{\texorpdfstring{livestock --- OE
\emph{nīetenu}}{livestock --- OE nīetenu}}\label{livestock-oe-nux12betenu}

\index[iv]{02oe@\ivlangheader{Old English}{West Saxon normalization unmarked}!nietenu@\iventry{nīetenu}{}}
\index[iv]{03pgmc@\ivlangheader{Proto-Germanic}{}!nautina@\iventry{*nautīną}{}}
\index[iv]{03pgmc@\ivlangheader{Proto-Germanic}{}!nautino@\iventry{*náutīnō}{}}

Derivation: citation reconstruction \emph{*nautīną}; form followed here
\emph{*náutīnō} \textgreater{} \emph{nīetenu} (regular).

\subsection{Derivation trace}\label{derivation-trace-50}

Proto input: \emph{*náutīnō}

\begingroup
\setlength{\fboxsep}{6pt}

\noindent\fbox{%
\begin{minipage}{0.97\linewidth}
\small
\begin{tabular}{@{}>{\raggedright\arraybackslash}p{0.485\linewidth}@{\hspace{0.030\linewidth}}>{\raggedright\arraybackslash}p{0.485\linewidth}@{}}
\begin{minipage}[t]{\linewidth}
\centering\textbf{Earlier Germanic changes}\par
\vspace{0.35em}
\raggedright
\centering\textbf{Northwest and West Germanic}\par
\raggedright
\vspace{0.2em}
\begin{tabularx}{\linewidth}{@{}>{\raggedright\arraybackslash}X@{\hspace{0.55em}}>{\raggedright\arraybackslash}p{0.18\linewidth}@{\hspace{0.25em}}}
\mbox{PNWGmc Final Long O Raising} & \emph{*náutīnu} \\
\end{tabularx}
\vspace{0.6em}
\centering\textbf{Early Anglo-Frisian}\par
\raggedright
\vspace{0.2em}
\raggedright [no change]\par
\end{minipage}
&
\begin{minipage}[t]{\linewidth}
\centering\textbf{Old English changes}\par
\vspace{0.35em}
\raggedright
\begin{tabularx}{\linewidth}{@{}>{\raggedright\arraybackslash}X@{\hspace{0.55em}}>{\raggedright\arraybackslash}p{0.20\linewidth}@{\hspace{0.25em}}}
\mbox{OE Au Brightening} & \emph{*náeutīnu} \\
\mbox{OE Diphthong Leveling} & \emph{*nēatīnu} \\
\mbox{OE I Umlaut} & \emph{*nīetīnu} \\
\mbox{OE Med Long Short} & \emph{*nīetinu} \\
OE Med Unstressed I Lowering1 & \emph{*nīetenu} \\
\end{tabularx}
\end{minipage}
\\
\end{tabular}
\end{minipage}%
} \endgroup

Old English form: \emph{nīetenu}

\subsection{Selected cell and Old English
evidence}\label{selected-cell-and-old-english-evidence}

The form compared here is the strong neuter nominative--accusative
plural
{\emph{nīetenu}}\index[iv]{02oe@\ivlangheader{Old English}{West Saxon normalization unmarked}!nietenu@\iventry{nīetenu}{}}
`animals, livestock'. Sweet's edition of the Hatton manuscript of the
\emph{Pastoral Care} prints
{\emph{nietenu}}\index[iv]{02oe@\ivlangheader{Old English}{West Saxon normalization unmarked}!nietenu@\iventry{nietenu}{}}
`animals' repeatedly: \emph{Se monn is on gecynde betera þonne dysig
nietenu} `the human being is by nature better than foolish animals', and
the contrasted subject \emph{ac nietenu} `but animals'
(\citeproc{ref-Sweet1871}{Sweet 1871, 109}, lines 8--11). These are
plural animals, not an endingless singular. The common
nominative--accusative ending is selected without inventing a new
inflectional cell.

Quantity and dialect are not read out of Sweet's unmarked spelling.
Ringe and Taylor identify West Saxon \emph{nīetenu}, against Mercian
{\emph{nētenu}}\index[iv]{02oe@\ivlangheader{Old English}{West Saxon normalization unmarked}!netenu@\iventry{nētenu}{}}
`animals' (\citeproc{ref-RingeTaylor2014}{Ringe and Taylor 2014,
298--99}). Goering likewise identifies these nominative--accusative
plurals and the later West Saxon
{\emph{nytenu}}\index[iv]{02oe@\ivlangheader{Old English}{West Saxon normalization unmarked}!nytenu@\iventry{nytenu}{}}
`animals' (\citeproc{ref-Goering2016}{Goering 2016, 186}). CAPR selects
the independently attested West Saxon endpoint, not a Mercian spelling
imposed on a West Saxon root-vowel history.

\subsection{Reconstruction and morphological
limits}\label{reconstruction-and-morphological-limits}

Fulk prints the complete ancestral plural
\Recon{nautīnō}\index[iv]{03pgmc@\ivlangheader{Proto-Germanic}{}!nautino@\iventry{*nautīnō}{}}
`animals' as the predecessor of Mercian \emph{nētenu} `animals', through
\emph{*nēatinu} (\citeproc{ref-Fulk2018}{Fulk 2018, 92}, §5.6). The
inherited long medial \emph{*ī} and the plain neuter plural \emph{*-ō}
are therefore not values selected from the transducer's prediction.

The comparative base is \emph{*nauta-} `cattle', supported by Old Norse
\emph{naut}, Old English \emph{nēat} and Old High German \emph{nōz}
(\citeproc{ref-Kroonen2013}{Kroonen 2013, 385};
\citeproc{ref-Orel2003}{Orel 2003, 282}). These dictionaries establish
the base, not a quoted complete \emph{*nautīna-} derivative. The
citation \emph{*nautīną} is explicitly reconstructed from Fulk's derived
plural stem and ordinary neuter a-stem singular inflection
(\citeproc{ref-Fulk2018}{Fulk 2018, 145--48}, §§7.7--7.8); it is not
passed off as a dictionary headword. The selected
\Recon{náutīnō}\index[iv]{03pgmc@\ivlangheader{Proto-Germanic}{}!nautino@\iventry{*náutīnō}{}}
`animals' marks the inherited initial root stress explicitly. Fulk's
printed form does not contain that acute accent. This annotation
resolves input tokenization; it does not introduce secondary stress,
change vowel quantity or supply an analogical transponent.

This plural does not depend on the later singular's final \emph{n}. In
the derivational input, \emph{*n} precedes the plural vowel and is not
word-final. Campbell attributes the singular's retained or restored
consonant to inflected forms (\citeproc{ref-Campbell1959}{Campbell 1959,
189}, §473). That account does not fix the restoration's exact date and
is not needed to make the present plural derive regularly.

\subsection{Development and regression
value}\label{development-and-regression-value}

\hyperref[rule-PNWGmcFinalLongORaising]{SC019 PNWGmcFinalLongORaising}
gives \emph{*náutīnu}. Root brightening and diphthong leveling give
\emph{*nēatīnu}; \hyperref[rule-OEIUmlaut]{SC055 OEIUmlaut} gives West
Saxon \emph{*nīetīnu}. The medial nasal remains unchanged under
\hyperref[rule-NWGmcInStemNLoss]{SC064 NWGmcInStemNLoss}: no exception
to final-n loss is introduced.

\hyperref[rule-OEMedLongShort]{SC117 OEMedLongShort}, the early medial
component of the established shortening history, gives \emph{*nīetinu}.
This is HLL, and neither weak high vowel is deleted by
\hyperref[rule-OEHighVowelApocope]{SC063 OEHighVowelApocope}. Later
\hyperref[rule-OEMedUnstressedILowering1]{SC074
OEMedUnstressedILowering1} gives \emph{nīetenu} `animals'. The full
inherited input, not a later-stage surrogate, produces the attested
comparator.

Indiscriminate heavy-root medial deletion instead gives \Pred{nīetnu};
HLL final apocope gives \Pred{nīeten}. More importantly, omitting early
medial shortening also gives \Pred{nīeten}: the still-heavy medial
syllable changes the quantitative deletion environment. The same
omission leaves selected \emph{lȳtelu} `little' and every earlier
selected output unchanged. Livestock therefore protects a real
dependency not supplied by the short-medial little plural or by merely
listing SC117's live firings. The two controls are complementary rather
than duplicate admissions.

\section{\texorpdfstring{louse --- OE
\emph{lūs}}{louse --- OE lūs}}\label{louse-oe-lux16bs}

\index[iv]{02oe@\ivlangheader{Old English}{West Saxon normalization unmarked}!lus@\iventry{lūs}{}}
\index[iv]{03pgmc@\ivlangheader{Proto-Germanic}{}!lusz@\iventry{*lūsz}{}}

Derivation: \emph{*lūsz} \textgreater{} \emph{lūs} (regular).

\subsection{Derivation trace}\label{derivation-trace-51}

Proto input: \emph{*lūsz}

\begingroup
\setlength{\fboxsep}{6pt}

\noindent\fbox{%
\begin{minipage}{0.97\linewidth}
\small
\begin{tabular}{@{}>{\raggedright\arraybackslash}p{0.485\linewidth}@{\hspace{0.030\linewidth}}>{\raggedright\arraybackslash}p{0.485\linewidth}@{}}
\begin{minipage}[t]{\linewidth}
\centering\textbf{Earlier Germanic changes}\par
\vspace{0.35em}
\raggedright
\centering\textbf{Northwest and West Germanic}\par
\raggedright
\vspace{0.2em}
\raggedright [no change]\par
\vspace{0.6em}
\centering\textbf{Early Anglo-Frisian}\par
\raggedright
\vspace{0.2em}
\raggedright [no change]\par
\end{minipage}
&
\begin{minipage}[t]{\linewidth}
\centering\textbf{Old English changes}\par
\vspace{0.35em}
\raggedright
\raggedright [no change]\par
\end{minipage}
\\
\end{tabular}
\end{minipage}%
} \endgroup

Old English form: \emph{lūs}

\subsection{Reconstruction and comparative
evidence}\label{reconstruction-and-comparative-evidence-43}

Orel's citation form is \emph{*lusz} `louse', with the explicit
nominative marker \emph{*-z} he writes across the root-noun class; his
entry points to Bammesberger and Griepentrog for the morphology
(\citeproc{ref-Orel2003}{Orel 2003, 252}). Kroonen has no `louse' entry,
so his dictionary cannot be cited for this word
(\citeproc{ref-Kroonen2013}{Kroonen 2013, 345}). Kluge/Seebold carry the
German cognate as the same root noun
(\citeproc{ref-KlugeSeebold2011}{Kluge 2011, 563}). Bammesberger shows
that the root-final \emph{s} is itself secondary: the root
{\emph{*luw-}}\index[iv]{03pgmc@\ivlangheader{Proto-Germanic}{}!luw@\iventry{*luw-}{}}
`louse' was extended with \emph{s} on the model of the `mouse' word
(\citeproc{ref-Bammesberger1990}{Bammesberger 1990, 195}). The `mouse'
parallel also shows why no extra sibilant is audible in the nominative:
Ringe reconstructs endingless
{\emph{*mūs}}\index[iv]{03pgmc@\ivlangheader{Proto-Germanic}{}!mus@\iventry{*mūs}{}}
`mouse', with the expected double sibilant of the nominative removed by
degemination (\citeproc{ref-Ringe2017}{Ringe 2017, 149}).

The derivational input
\Recon{lūsz}\index[iv]{03pgmc@\ivlangheader{Proto-Germanic}{}!lusz@\iventry{*lūsz}{}}
`louse' follows Orel's morphologically explicit convention. The final
\emph{*-z} is an inflectional marker made visible in the citation form,
not a claim that a phonetic cluster \emph{{[}sz{]}} was ever pronounced;
the notation records the morphology whose elimination the derivation
models.

\subsection{Old English evidence}\label{old-english-evidence-44}

Old English
{\emph{lūs}}\index[iv]{02oe@\ivlangheader{Old English}{West Saxon normalization unmarked}!lus@\iventry{lūs}{}}
`louse' is the endingless nominative-accusative singular of a root noun,
with mutated plural \emph{lȳs} `lice' (\citeproc{ref-Fulk2018}{Fulk
2018, 165--66}). Old Norse \emph{lús} `louse' is likewise endingless, in
keeping with the North Germanic redistribution that drops the reflex of
\emph{*-z} in feminine root nouns (\citeproc{ref-Fulk2018}{Fulk 2018,
167}). No West Germanic daughter shows any ending in this class
(\citeproc{ref-RingeTaylor2014}{Ringe and Taylor 2014, 118}).

\subsection{Development to Old
English}\label{development-to-old-english-43}

The root-noun nominative marker was eliminated before Proto-West
Germanic (\citeproc{ref-RingeTaylor2014}{Ringe and Taylor 2014, 118}).
From the derivational input \Recon{lūsz} `louse' the loss of the marker
gives \Recon{lūs} `louse', which is already the attested Old English
form. Because the root itself ends in \emph{s}, the marker's elimination
here is best read together with the degemination that Ringe describes
for `mouse': at no stage does the reconstruction require an audible
\emph{-sz} sequence. This early morphological development is distinct
from the later West Germanic loss of final \emph{*z} in unstressed
syllables and from the northern monosyllabic loss with compensatory
lengthening; the root-noun ending was gone before either applied, and
the root-final \emph{s} of \emph{lūs} `louse' was never eligible for
either change.

\section{\texorpdfstring{milk --- OE
\emph{meoloc}}{milk --- OE meoloc}}\label{milk-oe-meoloc}

\index[iv]{02oe@\ivlangheader{Old English}{West Saxon normalization unmarked}!meoloc@\iventry{meoloc}{}}
\index[iv]{03pgmc@\ivlangheader{Proto-Germanic}{}!melukz@\iventry{*mélukz}{}}

Derivation: \emph{*mélukz} \textgreater{} \emph{meoloc} (regular).

\subsection{Derivation trace}\label{derivation-trace-52}

Proto input: \emph{*mélukz}

\begingroup
\setlength{\fboxsep}{6pt}

\noindent\fbox{%
\begin{minipage}{0.97\linewidth}
\small
\begin{tabular}{@{}>{\raggedright\arraybackslash}p{0.359\linewidth}@{\hspace{0.030\linewidth}}>{\raggedright\arraybackslash}p{0.611\linewidth}@{}}
\begin{minipage}[t]{\linewidth}
\centering\textbf{Earlier Germanic changes}\par
\vspace{0.35em}
\raggedright
\centering\textbf{Northwest and West Germanic}\par
\raggedright
\vspace{0.2em}
\begin{tabularx}{\linewidth}{@{}>{\raggedright\arraybackslash}X@{\hspace{0.55em}}>{\raggedright\arraybackslash}p{0.14\linewidth}@{\hspace{0.25em}}}
\mbox{EAF Final Z Deletion} & \emph{*méluk} \\
\end{tabularx}
\vspace{0.6em}
\centering\textbf{Early Anglo-Frisian}\par
\raggedright
\vspace{0.2em}
\raggedright [no change]\par
\end{minipage}
&
\begin{minipage}[t]{\linewidth}
\centering\textbf{Old English changes}\par
\vspace{0.35em}
\raggedright
\begin{tabularx}{\linewidth}{@{}>{\raggedright\arraybackslash}X@{\hspace{0.55em}}>{\raggedright\arraybackslash}p{0.14\linewidth}@{\hspace{0.25em}}}
\mbox{OE Back Mutation} & \emph{*méoluk} \\
\mbox{OE Med Unstressed U Lowering} & \emph{*méolok} \\
\end{tabularx}
\end{minipage}
\\
\end{tabular}
\end{minipage}%
} \endgroup

Old English form: \emph{meoloc}

\subsection{Reconstruction and comparative
evidence}\label{reconstruction-and-comparative-evidence-44}

Kroonen and Orel reconstruct the noun with stem notation \emph{*meluk-}
and nominative \Recon{melukz} `milk', and the nominative-style input
used here is \Recon{mélukz} `milk' (\citeproc{ref-Kroonen2013}{Kroonen
2013}; \citeproc{ref-Orel2003}{Orel 2003, 306}).

\subsection{Old English evidence}\label{old-english-evidence-45}

Old English preserves a mixed dossier for this noun. Ringe and Taylor
describe West Saxon \emph{meolc} `milk' \textless{} \emph{meoluc} `milk'
\textless{} \Recon{mélukz} `milk', Campbell likewise discusses
\emph{meoluc} `milk' and
{\emph{meoloc}}\index[iv]{02oe@\ivlangheader{Old English}{West Saxon normalization unmarked}!meoloc@\iventry{meoloc}{}}
`milk', and Anglian shows \emph{milc} `milk'
(\citeproc{ref-RingeTaylor2014}{Ringe and Taylor 2014};
\citeproc{ref-Campbell1959}{Campbell 1959}).

\subsection{Development to Old
English}\label{development-to-old-english-44}

The unsyncopated line from \Recon{mélukz} `milk' loses final \emph{*z},
lowers unstressed \emph{u} to \emph{o}, and with back mutation yields
{\emph{meoloc}}\index[iv]{02oe@\ivlangheader{Old English}{West Saxon normalization unmarked}!meoloc@\iventry{meoloc}{}}
`milk'. That fuller unsyncopated outcome is the form represented here
(\citeproc{ref-RingeTaylor2014}{Ringe and Taylor 2014}).

\subsection{Form comparison}\label{form-comparison-1}

Syncopated \emph{meolc} and Anglian \emph{milc} belong to the competing
leveled tradition associated with oblique forms, whereas \emph{meoloc} /
\emph{meoluc} preserves the fuller nominal shape
(\citeproc{ref-Campbell1959}{Campbell 1959};
\citeproc{ref-RingeTaylor2014}{Ringe and Taylor 2014}).

\section{\texorpdfstring{mother --- OE
\emph{mōder}}{mother --- OE mōder}}\label{mother-oe-mux14dder}

\index[iv]{02oe@\ivlangheader{Old English}{West Saxon normalization unmarked}!moder@\iventry{mōder}{}}
\index[iv]{03pgmc@\ivlangheader{Proto-Germanic}{}!moder@\iventry{*mōdēr}{}}

Derivation: \emph{*mōdēr} \textgreater{} \emph{mōder} (regular).

\subsection{Derivation trace}\label{derivation-trace-53}

Proto input: \emph{*mōdēr}

\begingroup
\setlength{\fboxsep}{6pt}

\noindent\fbox{%
\begin{minipage}{0.97\linewidth}
\small
\begin{tabular}{@{}>{\raggedright\arraybackslash}p{0.359\linewidth}@{\hspace{0.030\linewidth}}>{\raggedright\arraybackslash}p{0.611\linewidth}@{}}
\begin{minipage}[t]{\linewidth}
\centering\textbf{Earlier Germanic changes}\par
\vspace{0.35em}
\raggedright
\centering\textbf{Northwest and West Germanic}\par
\raggedright
\vspace{0.2em}
\raggedright [no change]\par
\vspace{0.6em}
\centering\textbf{Early Anglo-Frisian}\par
\raggedright
\vspace{0.2em}
\raggedright [no change]\par
\end{minipage}
&
\begin{minipage}[t]{\linewidth}
\centering\textbf{Old English changes}\par
\vspace{0.35em}
\raggedright
\begin{tabularx}{\linewidth}{@{}>{\raggedright\arraybackslash}X@{\hspace{0.55em}}>{\raggedright\arraybackslash}p{0.14\linewidth}@{\hspace{0.25em}}}
OE Unstressed Long Vowel Shortening & \emph{*mōder} \\
\end{tabularx}
\end{minipage}
\\
\end{tabular}
\end{minipage}%
} \endgroup

Old English form: \emph{mōder}

\subsection{Reconstruction and comparative
evidence}\label{reconstruction-and-comparative-evidence-45}

Kroonen and Orel cite the Proto-Germanic r-stem kinship noun as
\emph{*mōder-} with citation \Recon{mōdēr} `mother'
(\citeproc{ref-Kroonen2013}{Kroonen 2013}; \citeproc{ref-Orel2003}{Orel
2003}).

\subsection{Old English evidence}\label{old-english-evidence-46}

The transmitted Old English headword tradition is \emph{mōdor} `mother'
/ \emph{modor} `mother', with oblique \emph{mēder} `mother' in the
paradigm. Clark Hall, Campbell, and Ringe and Taylor all preserve that
contrast (\citeproc{ref-ClarkHall1960}{Clark Hall 1960};
\citeproc{ref-Campbell1959}{Campbell 1959};
\citeproc{ref-RingeTaylor2014}{Ringe and Taylor 2014}).

\subsection{Development to Old
English}\label{development-to-old-english-45}

From \Recon{mōdēr} `mother', the regular suffixal development yields
\emph{mōder} `mother'. That regular nominative reflex is the form
represented here, while the more familiar citation form \emph{mōdor}
`mother' reflects later levelling within the r-stem paradigm
(\citeproc{ref-Campbell1959}{Campbell 1959};
\citeproc{ref-RingeTaylor2014}{Ringe and Taylor 2014}).

\subsection{Form comparison}\label{form-comparison-2}

The note therefore concerns inherited vocalism rather than a different
lexeme: \emph{mōder} `mother' is the regularized nominative represented
here, but dictionaries usually print \emph{mōdor} / \emph{modor}
`mother', and the oblique evidence survives in \emph{mēder} `mother'
(\citeproc{ref-ClarkHall1960}{Clark Hall 1960};
\citeproc{ref-RingeTaylor2014}{Ringe and Taylor 2014}).

\section{\texorpdfstring{nest --- OE
\emph{nest}}{nest --- OE nest}}\label{nest-oe-nest}

\index[iv]{02oe@\ivlangheader{Old English}{West Saxon normalization unmarked}!nest@\iventry{nest}{}}
\index[iv]{03pgmc@\ivlangheader{Proto-Germanic}{}!nista@\iventry{*nístą}{}}

Derivation: \emph{*nístą} \textgreater{} \emph{nest} (regular).

\subsection{Derivation trace}\label{derivation-trace-54}

Proto input: \emph{*nístą}

\begingroup
\setlength{\fboxsep}{6pt}

\noindent\fbox{%
\begin{minipage}{0.97\linewidth}
\small
\begin{tabular}{@{}>{\raggedright\arraybackslash}p{0.359\linewidth}@{\hspace{0.030\linewidth}}>{\raggedright\arraybackslash}p{0.611\linewidth}@{}}
\begin{minipage}[t]{\linewidth}
\centering\textbf{Earlier Germanic changes}\par
\vspace{0.35em}
\raggedright
\centering\textbf{Northwest and West Germanic}\par
\raggedright
\vspace{0.2em}
\begin{tabularx}{\linewidth}{@{}>{\raggedright\arraybackslash}X@{\hspace{0.55em}}>{\raggedright\arraybackslash}p{0.14\linewidth}@{\hspace{0.25em}}}
\mbox{PNWGmc I Lowering} & \emph{*néstą} \\
\end{tabularx}
\vspace{0.6em}
\centering\textbf{Early Anglo-Frisian}\par
\raggedright
\vspace{0.2em}
\raggedright [no change]\par
\end{minipage}
&
\begin{minipage}[t]{\linewidth}
\centering\textbf{Old English changes}\par
\vspace{0.35em}
\raggedright
\begin{tabularx}{\linewidth}{@{}>{\raggedright\arraybackslash}X@{\hspace{0.55em}}>{\raggedright\arraybackslash}p{0.14\linewidth}@{\hspace{0.25em}}}
OE Heavy Syllable Nasal Apocope & \emph{*nést} \\
\end{tabularx}
\end{minipage}
\\
\end{tabular}
\end{minipage}%
} \endgroup

Old English form: \emph{nest}

\subsection{Comparative formation and
lowering}\label{comparative-formation-and-lowering}

Old English
{\emph{nest}}\index[iv]{02oe@\ivlangheader{Old English}{West Saxon normalization unmarked}!nest@\iventry{nest}{}}
`nest' continues an inherited neuter formation. Kroonen reconstructs
\emph{*nista-}, Orel the full-word \emph{*nistan}. Their comparanda
include Latin \emph{nīdus} `nest' and Sanskrit \emph{nīḍá-} `nest'. The
traditional analysis combines an element meaning `down' with a
zero-grade formation of the root `sit'. The English word must not be
confused with the separate food noun also spelled \emph{nest}
(\citeproc{ref-Kroonen2013}{Kroonen 2013, 388, 391};
\citeproc{ref-Orel2003}{Orel 2003, 287}).

The selected
\Recon{nístą}\index[iv]{03pgmc@\ivlangheader{Proto-Germanic}{}!nista@\iventry{*nístą}{}}
`nest' retains the inherited high vowel.
\hyperref[rule-PNWGmcILowering]{SC015 PNWGmcILowering} lowers it
immediately before \emph{*st}, rather than because of the next vowel.
This is the principal independent lexical support for Hill's \emph{*st}
condition. The contrast with \emph{*sk} in the fish family is useful,
but it does not turn one inherited nest etymon into numerous independent
witnesses (\citeproc{ref-Hill2017}{Hill 2017, 153--55}).

Later mutation in a derivative with a high-front trigger can conceal an
earlier mid vowel. Such derivatives test chronological compatibility;
they do not establish every detail of a lost paradigm. The general
productivity of lowering before \emph{*st} is less securely demonstrated
than lowering before inherited \emph{*r}, and the adopted unified law
retains that evidential qualification (\citeproc{ref-Hill2017}{Hill
2017, 153--56}).

\section{\texorpdfstring{net --- OE
\emph{nett}}{net --- OE nett}}\label{net-oe-nett}

\index[iv]{02oe@\ivlangheader{Old English}{West Saxon normalization unmarked}!nett@\iventry{nett}{}}
\index[iv]{03pgmc@\ivlangheader{Proto-Germanic}{}!natja@\iventry{*nátją}{}}

Derivation: \emph{*nátją} \textgreater{} \emph{nett} (regular).

\subsection{Derivation trace}\label{derivation-trace-55}

Proto input: \emph{*nátją}

\begingroup
\setlength{\fboxsep}{6pt}

\noindent\fbox{%
\begin{minipage}{0.97\linewidth}
\small
\begin{tabular}{@{}>{\raggedright\arraybackslash}p{0.359\linewidth}@{\hspace{0.030\linewidth}}>{\raggedright\arraybackslash}p{0.611\linewidth}@{}}
\begin{minipage}[t]{\linewidth}
\centering\textbf{Earlier Germanic changes}\par
\vspace{0.35em}
\raggedright
\centering\textbf{Northwest and West Germanic}\par
\raggedright
\vspace{0.2em}
\begin{tabularx}{\linewidth}{@{}>{\raggedright\arraybackslash}X@{\hspace{0.55em}}>{\raggedright\arraybackslash}p{0.24\linewidth}@{\hspace{0.25em}}}
\mbox{PWGmc J Gemination} & \emph{*náttją} \\
\end{tabularx}
\vspace{0.6em}
\centering\textbf{Early Anglo-Frisian}\par
\raggedright
\vspace{0.2em}
\raggedright [no change]\par
\end{minipage}
&
\begin{minipage}[t]{\linewidth}
\centering\textbf{Old English changes}\par
\vspace{0.35em}
\raggedright
\begin{tabularx}{\linewidth}{@{}>{\raggedright\arraybackslash}X@{\hspace{0.55em}}>{\raggedright\arraybackslash}p{0.14\linewidth}@{\hspace{0.25em}}}
\mbox{EAF Brightening} & \emph{*nættją} \\
OE Heavy Syllable Nasal Apocope & \emph{*nættj} \\
\mbox{OE I Umlaut} & \emph{*nettj} \\
\mbox{OE J Loss After Heavy} & \emph{*nett} \\
\end{tabularx}
\end{minipage}
\\
\end{tabular}
\end{minipage}%
} \endgroup

Old English form: \emph{nett}

\subsection{Reconstruction and comparative
evidence}\label{reconstruction-and-comparative-evidence-46}

Orel gives \Recon{natjan} `net' with Old English
{\emph{nett}}\index[iv]{02oe@\ivlangheader{Old English}{West Saxon normalization unmarked}!nett@\iventry{nett}{}}
`net', and Fulk's account of West Germanic gemination before \emph{j}
explains the geminate outcome after a short vowel
(\citeproc{ref-Orel2003}{Orel 2003}; \citeproc{ref-Fulk2018}{Fulk
2018}).

\subsection{Old English evidence}\label{old-english-evidence-47}

Clark Hall and Bosworth-Toller record
{\emph{nett}}\index[iv]{02oe@\ivlangheader{Old English}{West Saxon normalization unmarked}!nett@\iventry{nett}{}}
`net' as the noun, and Campbell notes that final geminates are often
graphically simplified in Old English spelling
(\citeproc{ref-ClarkHall1960}{Clark Hall 1960};
\citeproc{ref-BosworthToller1898}{Bosworth and Toller 1898, 29};
\citeproc{ref-Campbell1959}{Campbell 1959}).

\subsection{Development to Old
English}\label{development-to-old-english-46}

From \Recon{nátją} `net', West Germanic j-gemination first gives
\Recon{náttją} `net'. Later brightening, loss of the weak ending, and
loss of final \emph{j} after a heavy stem yield
{\emph{nett}}\index[iv]{02oe@\ivlangheader{Old English}{West Saxon normalization unmarked}!nett@\iventry{nett}{}}
`net', so the development represented here is regular
(\citeproc{ref-Fulk2018}{Fulk 2018}).

\subsection{Form note}\label{form-note-13}

Spellings in \emph{net} can therefore be graphic simplifications, but
the lexical target supported by the dictionary evidence is
{\emph{nett}}\index[iv]{02oe@\ivlangheader{Old English}{West Saxon normalization unmarked}!nett@\iventry{nett}{}}
`net' (\citeproc{ref-Campbell1959}{Campbell 1959};
\citeproc{ref-Orel2003}{Orel 2003}).

\section{\texorpdfstring{nightmare --- OE
\emph{mære}}{nightmare --- OE mære}}\label{nightmare-oe-muxe6re}

\index[iv]{02oe@\ivlangheader{Old English}{West Saxon normalization unmarked}!maere@\iventry{mære}{}}
\index[iv]{03pgmc@\ivlangheader{Proto-Germanic}{}!maron@\iventry{*márōn}{}}

Derivation: \emph{*márōn} yields \emph{mare} in the current cascade; the
selected Old English form is \emph{mære} (implementation mismatch).

\subsection{Derivation trace}\label{derivation-trace-56}

Proto input: \emph{*márōn}

\begingroup
\setlength{\fboxsep}{6pt}

\noindent\fbox{%
\begin{minipage}{0.97\linewidth}
\small
\begin{tabular}{@{}>{\raggedright\arraybackslash}p{0.359\linewidth}@{\hspace{0.030\linewidth}}>{\raggedright\arraybackslash}p{0.611\linewidth}@{}}
\begin{minipage}[t]{\linewidth}
\centering\textbf{Earlier Germanic changes}\par
\vspace{0.35em}
\raggedright
\centering\textbf{Northwest and West Germanic}\par
\raggedright
\vspace{0.2em}
\begin{tabularx}{\linewidth}{@{}>{\raggedright\arraybackslash}X@{\hspace{0.55em}}>{\raggedright\arraybackslash}p{0.14\linewidth}@{\hspace{0.25em}}}
\mbox{PNWGmc N Stem N Loss} & \emph{*márǭ} \\
\end{tabularx}
\vspace{0.6em}
\centering\textbf{Early Anglo-Frisian}\par
\raggedright
\vspace{0.2em}
\raggedright [no change]\par
\end{minipage}
&
\begin{minipage}[t]{\linewidth}
\centering\textbf{Old English changes}\par
\vspace{0.35em}
\raggedright
\begin{tabularx}{\linewidth}{@{}>{\raggedright\arraybackslash}X@{\hspace{0.55em}}>{\raggedright\arraybackslash}p{0.14\linewidth}@{\hspace{0.25em}}}
\mbox{EAF Brightening} & \emph{*mærǭ} \\
\mbox{OE A Restoration} & \emph{*marǭ} \\
OE Unstressed Long Vowel Shortening & \emph{*maræ} \\
\mbox{OE Unstressed AE Merger} & \Recon{mare} \\
\end{tabularx}
\end{minipage}
\\
\end{tabular}
\end{minipage}%
} \endgroup

Current cascade output: \emph{mare}

Old English form: \emph{mære}

\subsection{Old English evidence and selected
form}\label{old-english-evidence-and-selected-form}

The form compared here is the weak feminine nominative singular
\emph{mære} `nightmare'. The Corpus Glossary gives \emph{maere} under
Latin \emph{incuba}; normalizing the spelling \emph{ae} as \emph{æ}
gives \emph{mære} (\citeproc{ref-Lindsay1921}{Lindsay 1921, 96, I225}).
Hall distinguishes this weak noun from the strong noun \emph{mær}. Both
have a short root vowel and are distinct from the adjective \emph{mǣre}
`famous' (\citeproc{ref-Hall2007}{Hall 2007, 301--2}, §2).

The variant
{\emph{mare}}\index[iv]{02oe@\ivlangheader{Old English}{West Saxon normalization unmarked}!mare@\iventry{mare}{}}
`nightmare' is also attested: \emph{Bald's Leechbook} has the expression
\emph{Gif mon mare ride}. Oblique
{\emph{maran}}\index[iv]{02oe@\ivlangheader{Old English}{West Saxon normalization unmarked}!maran@\iventry{maran}{}}
`nightmare' occurs in \emph{Leechbook III}. Hall's table separates these
contexts from the related glossary copies, which do not constitute
independent attestations of one gloss (\citeproc{ref-Hall2007}{Hall
2007, 313--14}, table 1).

\subsection{Reconstruction and the
variant}\label{reconstruction-and-the-variant}

Ringe and Taylor reconstruct a nominative \emph{*marā} beside an oblique
stem \emph{*marōn-}, followed by the front-root forms \emph{*mærǣ} and
\emph{*mærōn-}. Their printed list includes \emph{mare}, \emph{maran},
and \emph{mære}, not merely \emph{mare} and \emph{mere}
(\citeproc{ref-RingeTaylor2014}{Ringe and Taylor 2014, 192}). Hall
explicitly explains \emph{mare} as levelling of the nominative root
vowel to that of oblique forms such as \emph{maran}, in which retraction
had applied. This is a paradigmatic development, not evidence that the
nominative itself regularly underwent the same restoration
(\citeproc{ref-Hall2007}{Hall 2007, 301}, §2).

\subsection{Current derivation and implementation
gap}\label{current-derivation-and-implementation-gap}

The derivational input remains \emph{*márōn}. The unchanged cascade
produces \emph{mare}, not the selected \emph{mære}. In the current
implementation, the retained nasal ending triggers root-vowel
restoration; its output matches an attested variant but does not
implement the cited front-root nominative history. The unadopted Model B
research experiment yields \emph{mære}. Selecting that attested
nominative does not adopt the experiment or change any production sound
rule.

The English concept label `nightmare' indexes the Old English simplex,
not an attested compound. The regular classification describes the
selected nominative history; the later levelled variant is retained as
comparison, not substituted for the form compared here to obtain a
computational match.

\section{\texorpdfstring{rainbow --- OE
\emph{reġnboga}}{rainbow --- OE reġnboga}}\label{rainbow-oe-reux121nboga}

\index[iv]{02oe@\ivlangheader{Old English}{West Saxon normalization unmarked}!regnboga@\iventry{reġnboga}{}}
\index[iv]{03pgmc@\ivlangheader{Proto-Germanic}{}!regnabugo@\iventry{*régna-bùgô}{}}
\index[iv]{03pgmc@\ivlangheader{Proto-Germanic}{}!regnabugo@\iventry{*régnabùgô}{}}

Derivation: \emph{*régna-bùgô} \textgreater{} \emph{reġnboga} (regular).

\subsection{Derivation trace}\label{derivation-trace-57}

Proto input: \emph{*régnabùgô}

\begingroup
\setlength{\fboxsep}{6pt}

\noindent\fbox{%
\begin{minipage}{0.97\linewidth}
\footnotesize
\begin{tabular}{@{}>{\raggedright\arraybackslash}p{0.485\linewidth}@{\hspace{0.030\linewidth}}>{\raggedright\arraybackslash}p{0.485\linewidth}@{}}
\begin{minipage}[t]{\linewidth}
\centering\textbf{Earlier Germanic changes}\par
\vspace{0.35em}
\raggedright
\centering\textbf{Northwest and West Germanic}\par
\raggedright
\vspace{0.2em}
\begin{tabularx}{\linewidth}{@{}>{\raggedright\arraybackslash}X@{\hspace{0.55em}}>{\raggedright\arraybackslash}p{0.22\linewidth}@{\hspace{0.25em}}}
\mbox{PNWGmc U Lowering} & \emph{*régnabògô} \\
\end{tabularx}
\vspace{0.6em}
\centering\textbf{Early Anglo-Frisian}\par
\raggedright
\vspace{0.2em}
\raggedright [no change]\par
\end{minipage}
&
\begin{minipage}[t]{\linewidth}
\centering\textbf{Old English changes}\par
\vspace{0.35em}
\raggedright
\begin{tabularx}{\linewidth}{@{}>{\raggedright\arraybackslash}X@{\hspace{0.55em}}>{\raggedright\arraybackslash}p{0.24\linewidth}@{\hspace{0.25em}}}
\mbox{OE Strip Secondary Stress} & \emph{*régnabogô} \\
\mbox{EAF Brightening Unstressed} & \emph{*régnæbogô} \\
\mbox{OE A Restoration} & \emph{*régnabogô} \\
\mbox{OE Velar Palatalization} & \emph{*réʝnabogô} \\
\mbox{OE Non High Syn} & \emph{*réʝnbogô} \\
OE Unstressed Long Vowel Shortening & \emph{*réʝnboga} \\
\mbox{OE Palatal Fricative Merger} & \emph{*réjnboga} \\
\end{tabularx}
\end{minipage}
\\
\end{tabular}
\end{minipage}%
} \endgroup

Old English form: \emph{reġnboga}

\subsection{Reconstruction and compound
formation}\label{reconstruction-and-compound-formation}

The compound means `rain-bow'. Orel reconstructs \emph{*reʒna-buʒōn},
explicitly joining the rain noun \emph{*reʒnan ~ *reʒnaz} and the bow
noun \emph{*buʒōn}. The latter belongs to the family of `bend'. Kroonen
independently gives masculine \emph{*bugan-} `bow', an n-stem derivative
of that verb; his rain entry supplies the first constituent, not an
independently printed reconstruction of the complete accented compound
(\citeproc{ref-Orel2003}{Orel 2003, 61, 300};
\citeproc{ref-Kroonen2013}{Kroonen 2013, 61--62, 82, 408}). The bow
noun, the bend verb and the different noun for an arm-ring must not be
treated as interchangeable second members.

The selected \emph{*régna-bùgô} is a complete nominative in CAPR's
quantity and accent notation, not a quotation of Kroonen's bare stem or
a mechanical transcription of Orel's citation ending. Old English
masculine weak \emph{boga} `bow' and oblique \emph{bogan} reflect the
n-stem history. The compound takes the same inflection, with nominative
\emph{-boga} and oblique \emph{-bogan}. The accentuation assigns primary
prominence to the first root and subordinate prominence to the second.
Its exact prehistoric realization is reconstructed; neither the
etymological hyphen nor a dictionary's normalized spelling independently
establishes that accent.

\subsection{The early vowel of the second
root}\label{the-early-vowel-of-the-second-root}

The ordinary short \emph{o} of
{\emph{boga}}\index[iv]{02oe@\ivlangheader{Old English}{West Saxon normalization unmarked}!boga@\iventry{boga}{}}
has an early explanation: inherited short \emph{u} lowered before a
non-high vowel in the next syllable. Campbell, Hogg, Fulk, and Ringe and
Taylor describe the change and its nasal and intervening-\emph{j}
restrictions. Colman expressly applies this explanation to the noun
\emph{boga} (\citeproc{ref-Campbell1959}{Campbell 1959, 43}, §§115--116;
\citeproc{ref-HoggGrammar2011}{Hogg 2011, 54--55}, §§3.10--3.11;
\citeproc{ref-Fulk2018}{Fulk 2018, 56--57}, §4.3;
\citeproc{ref-RingeTaylor2014}{Ringe and Taylor 2014, 27--29};
\citeproc{ref-Colman2011}{Colman 2011, 95--96}).

CAPR infers that the same assimilation applied to the secondarily
stressed root in the compound: \emph{*régna-bùgô > *régna-bògô}. The
following non-high vowel was still present, and neither a nasal
immediately after the target nor an intervening \emph{j} blocks this
particular sequence. The compound extension is a phonological inference,
not a derivation printed with these accents by the cited authorities. It
does not extend early lowering to ordinary unstressed medial \emph{u}. A
regular account is preferable to replacement from the independent noun
unless positive evidence requires that analogy; none has been
established for this selected history.

The name
{\emph{Buga}}\index[iv]{02oe@\ivlangheader{Old English}{West Saxon normalization unmarked}!buga@\iventry{Buga}{}}
beside
{\emph{Boga}}\index[iv]{02oe@\ivlangheader{Old English}{West Saxon normalization unmarked}!boga@\iventry{Boga}{}}
remains relevant. Colman's coins preserve both vowel spellings, and
transmission between moneyer and die-cutter complicates dialect
attribution. The name's retained \emph{u} cannot be explained merely by
saying that names escape sound laws. Nor does this distribution erase
the ordinary noun's independently supported lowering. It remains a
qualification of the broader retention problem, not a reason to
implement an optional change (\citeproc{ref-Colman2011}{Colman 2011,
95--98}).

\subsection{Linking syncope and the regular
development}\label{linking-syncope-and-the-regular-development}

Early lowering retains the subordinate accent. Its technical decoding
gives \emph{*régnabogô}, with the first-member vowel still present.
Ordinary nonhigh medial syncope later gives \emph{*réʝnbogô}. Campbell's
discussion of weakened compound junctures supplies the broader
historical context; the exact ordering of these serialized annotations
is CAPR's explicit representation (\citeproc{ref-Campbell1959}{Campbell
1959, 148--49}, §§356--357, p.~153, §367).

The first velar undergoes palatalization, the final long weak vowel
shortens, and the palatal consonant is represented in the selected
spelling
{\emph{reġnboga}}\index[iv]{02oe@\ivlangheader{Old English}{West Saxon normalization unmarked}!regnboga@\iventry{reġnboga}{}}.
The second velar remains before a back vowel. Thus the adopted chain is
\emph{*régna-bùgô > *régna-bògô > *régnabogô > *réʝnabogô
> *réʝnbogô > *réʝnboga > *réjnboga}, then \emph{reġnboga}. The
intermediate accented and palatal symbols are analytical
representations, not alternative manuscript spellings. The principal
account is \hyperref[rule-PNWGmcULowering]{SC017 PNWGmcULowering}, with
linking deletion at \hyperref[rule-OENonHighSyn]{SC114 OENonHighSyn}.

The inherited \emph{b} is realized positionally after that loss, when it
follows \emph{n} and remains a stop. Hogg describes the stop and
bilabial fricative as early positional realizations of one phoneme,
before later fricative merger (\citeproc{ref-HoggGrammar2011}{Hogg 2011,
29--30}, §§2.53--2.54). Delayed computation of this distribution is not
a new late lenition event or a general postnasal \emph{β > b} law. With
early realization frozen, the same later syncope predicts
\emph{*reġnfoga}. The corrected representation removes the need for an
independently undated early compound-vowel proxy.

The word already contains \emph{o} when
\hyperref[rule-OEMedUnstressedULowering]{SC040 OEMedUnstressedULowering}
operates. It therefore supplies no evidence that compound secondary
stress had disappeared before that late lowering. Its genuine early
vowel dependency is survival of the subordinate accent until root
lowering; later linking loss no longer reads that annotation. This does
not alter the separate weak-vowel chronology of \emph{weorold} `world'
or establish a universal date for compound deaccentuation.

\subsection{Old English documentary evidence and
spelling}\label{old-english-documentary-evidence-and-spelling}

Bosworth--Toller gives \emph{regn-boga} `rainbow' as masculine, with
oblique \emph{-an}. Its quotation ``Hwí wæs se rénboga tó wedde gesette
mancynne?'' has nominative singular
{\emph{rénboga}}\index[iv]{02oe@\ivlangheader{Old English}{West Saxon normalization unmarked}!renboga@\iventry{rénboga}{}},
the subject of the passive question. The following statement has
\emph{ðone rénbogan}, accusative singular. The online entry assigns this
passage to \emph{Boutr. Scrd.} 21, 19--26. Its separate \emph{Genesis}
9:13 quotation, ``Ic sette mínne rénbogan \ldots{} on wolcnum'',
likewise has accusative singular
{\emph{rénbogan}}\index[iv]{02oe@\ivlangheader{Old English}{West Saxon normalization unmarked}!renbogan@\iventry{rénbogan}{}},
not a competing nominative target. The two vocabulary citations,
\emph{Wrt. Voc.} i. 52, 42 and 76, 33, pair the noun with Latin
\emph{iris/yris} and \emph{arcus}
(\citeproc{ref-BosworthToller1898}{Bosworth and Toller 1898},
s.v.~regn-boga).

These are dictionary-mediated documentary witnesses, checked in the
\href{https://bosworthtoller.com/25654}{online entry}, not a new
manuscript collation. The dictionary's \emph{rén-} represents the first
member without written \emph{g} and marks vowel length; its accent is
not CAPR's reconstructed primary-stress mark. The selected \emph{reġn-}
retains the palatal consonant in modern scholarly notation. Hyphenation
and the palatal dot likewise belong to editorial representation. Neither
the contracted first member nor oblique \emph{-bogan} changes the
evidence for second-root \emph{o}. The quotations do not justify
attaching an exact manuscript date or a narrower dialect label to every
normalized spelling. The selected regular span is retained without
claiming to model every attested orthographic or inflectional variant.

\section{\texorpdfstring{sheep --- OE
\emph{sċēap}}{sheep --- OE sċēap}}\label{sheep-oe-sux10bux113ap}

\index[iv]{02oe@\ivlangheader{Old English}{West Saxon normalization unmarked}!sceap@\iventry{sċēap}{}}
\index[iv]{03pgmc@\ivlangheader{Proto-Germanic}{}!skepa@\iventry{*skḗpą}{}}

Derivation: \emph{*skḗpą} \textgreater{} \emph{sċēap} (regular).

\subsection{Derivation trace}\label{derivation-trace-58}

Proto input: \emph{*skḗpą}

\begingroup
\setlength{\fboxsep}{6pt}

\noindent\fbox{%
\begin{minipage}{0.97\linewidth}
\small
\begin{tabular}{@{}>{\raggedright\arraybackslash}p{0.388\linewidth}@{\hspace{0.030\linewidth}}>{\raggedright\arraybackslash}p{0.582\linewidth}@{}}
\begin{minipage}[t]{\linewidth}
\centering\textbf{Earlier Germanic changes}\par
\vspace{0.35em}
\raggedright
\centering\textbf{Northwest and West Germanic}\par
\raggedright
\vspace{0.2em}
\begin{tabularx}{\linewidth}{@{}>{\raggedright\arraybackslash}X@{\hspace{0.55em}}>{\raggedright\arraybackslash}p{0.14\linewidth}@{\hspace{0.25em}}}
\mbox{PNWGmc Long E Lowering} & \emph{*skāpą} \\
\mbox{EAF Long A Fronting} & \emph{*skǣpą} \\
\end{tabularx}
\vspace{0.6em}
\centering\textbf{Early Anglo-Frisian}\par
\raggedright
\vspace{0.2em}
\raggedright [no change]\par
\end{minipage}
&
\begin{minipage}[t]{\linewidth}
\centering\textbf{Old English changes}\par
\vspace{0.35em}
\raggedright
\begin{tabularx}{\linewidth}{@{}>{\raggedright\arraybackslash}X@{\hspace{0.55em}}>{\raggedright\arraybackslash}p{0.14\linewidth}@{\hspace{0.25em}}}
OE Heavy Syllable Nasal Apocope & \emph{*skǣp} \\
\mbox{OE Sk Palatalization} & \emph{*ʃǣp} \\
OE Ws Palatal Diphthongization & \emph{*ʃēap} \\
\end{tabularx}
\end{minipage}
\\
\end{tabular}
\end{minipage}%
} \endgroup

Old English form: \emph{sċēap}

\subsection{Reconstruction and comparative
evidence}\label{reconstruction-and-comparative-evidence-47}

Ringe and Taylor cite a later West Germanic \emph{*skap} `sheep'
\textgreater{} WS \emph{scéap} `sheep', while Orel preserves a
Proto-Germanic noun of the \emph{*skēp-} type for the same lexeme
(\citeproc{ref-RingeTaylor2014}{Ringe and Taylor 2014};
\citeproc{ref-Orel2003}{Orel 2003}).

\subsection{Old English evidence}\label{old-english-evidence-48}

Clark Hall records \emph{scēap} `sheep' with spelling variation, and
Campbell likewise lists West Saxon \emph{scéap} `sheep' among the
palatal-diphthongized forms (\citeproc{ref-ClarkHall1960}{Clark Hall
1960}; \citeproc{ref-Campbell1959}{Campbell 1959}).

\subsection{Development to Old
English}\label{development-to-old-english-47}

From \Recon{skḗpą} `sheep', Northwest Germanic lowering gives
\Recon{skǣpą} `sheep'; after apocope and palatalization the West Saxon
branch diphthongizes to \emph{sċēap} `sheep'. The development
represented here is therefore fully regular.

\subsection{Dialect note}\label{dialect-note-4}

Ringe and Taylor contrast West Saxon \emph{scéap} `sheep' with Mercian
and Kentish \emph{scép} `sheep', and Campbell also notes Northumbrian
\emph{scip} `sheep'. The form represented here is the West Saxon
headword (\citeproc{ref-RingeTaylor2014}{Ringe and Taylor 2014};
\citeproc{ref-Campbell1959}{Campbell 1959}).

\section{\texorpdfstring{shilling --- OE
\emph{sċilling}}{shilling --- OE sċilling}}\label{shilling-oe-sux10billing}

\index[iv]{02oe@\ivlangheader{Old English}{West Saxon normalization unmarked}!scilling@\iventry{sċilling}{}}
\index[iv]{03pgmc@\ivlangheader{Proto-Germanic}{}!skillingaz@\iventry{*skíllingaz}{}}

Derivation: \emph{*skíllingaz} \textgreater{} \emph{sċilling} (regular).

\subsection{Derivation trace}\label{derivation-trace-59}

Proto input: \emph{*skíllingaz}

\begingroup
\setlength{\fboxsep}{6pt}

\noindent\fbox{%
\begin{minipage}{0.97\linewidth}
\small
\begin{tabular}{@{}>{\raggedright\arraybackslash}p{0.359\linewidth}@{\hspace{0.030\linewidth}}>{\raggedright\arraybackslash}p{0.611\linewidth}@{}}
\begin{minipage}[t]{\linewidth}
\centering\textbf{Earlier Germanic changes}\par
\vspace{0.35em}
\raggedright
\centering\textbf{Northwest and West Germanic}\par
\raggedright
\vspace{0.2em}
\begin{tabularx}{\linewidth}{@{}>{\raggedright\arraybackslash}X@{\hspace{0.55em}}>{\raggedright\arraybackslash}p{0.24\linewidth}@{\hspace{0.25em}}}
\mbox{EAF Final Z Deletion} & \emph{*skíllinga} \\
\end{tabularx}
\vspace{0.6em}
\centering\textbf{Early Anglo-Frisian}\par
\raggedright
\vspace{0.2em}
\raggedright [no change]\par
\end{minipage}
&
\begin{minipage}[t]{\linewidth}
\centering\textbf{Old English changes}\par
\vspace{0.35em}
\raggedright
\begin{tabularx}{\linewidth}{@{}>{\raggedright\arraybackslash}X@{\hspace{0.55em}}>{\raggedright\arraybackslash}p{0.20\linewidth}@{\hspace{0.25em}}}
\mbox{PWGmc Final Bare A Loss} & \emph{*skílling} \\
\mbox{OE Sk Palatalization} & \emph{*ʃílling} \\
OE Med Unstressed I Lowering1 & \emph{*ʃílleng} \\
\mbox{OE Med Unstressed I Lowering} & \emph{*ʃílling} \\
\end{tabularx}
\end{minipage}
\\
\end{tabular}
\end{minipage}%
} \endgroup

Old English form: \emph{sċilling}

\subsection{Reconstruction and comparative
evidence}\label{reconstruction-and-comparative-evidence-48}

Kroonen treats the cognate set under \emph{*skellinga-}
\textasciitilde{} \emph{*skillinga-} and connects it with
\emph{*skeld-linga-}, while Orel likewise gives the coin word with OE
{\emph{scilling}}\index[iv]{02oe@\ivlangheader{Old English}{West Saxon normalization unmarked}!scilling@\iventry{scilling}{}}
`shilling' among the reflexes (\citeproc{ref-Kroonen2013}{Kroonen 2013};
\citeproc{ref-Orel2003}{Orel 2003}). The derivational input
\Recon{skíllingaz} `shilling' is the nominative-style form used here to
represent that inherited \emph{*-ing-} derivative.

\subsection{Old English evidence}\label{old-english-evidence-49}

Clark Hall records \emph{scilling} `shilling', and Campbell cites it
among nouns with unstressed \emph{i} in derivational \emph{-ing}
(\citeproc{ref-ClarkHall1960}{Clark Hall 1960};
\citeproc{ref-Campbell1959}{Campbell 1959}). The target is the ordinary
OE citation form, normalized as \emph{sċilling} `shilling'.

\subsection{Development to Old
English}\label{development-to-old-english-48}

From \Recon{skíllingaz} `shilling', loss of final \emph{-az} yields
\Recon{skílling} `shilling'. Old English palatalization of initial
\emph{sk} before front vocalism then gives \emph{sċilling} `shilling'.
The \emph{i} of derivational \emph{-ing-} remains, so the regular
outcome is \emph{sċilling} `shilling', not \Recon{sċilleng} `shilling'
(\citeproc{ref-Campbell1959}{Campbell 1959};
\citeproc{ref-HoggGrammar1992}{Hogg 1992a}).

\subsection{Form note}\label{form-note-14}

Kroonen's \emph{*skellinga-} \textasciitilde{} \emph{*skillinga-} and
his internal analysis \emph{*skeld-linga-} belong to the etymological
background of the cognate set. The derivational input \Recon{skíllingaz}
`shilling' is the specific form used for the derivation represented here
(\citeproc{ref-Kroonen2013}{Kroonen 2013}).

\section{\texorpdfstring{shovel --- OE
\emph{sċofl}}{shovel --- OE sċofl}}\label{shovel-oe-sux10bofl}

\index[iv]{02oe@\ivlangheader{Old English}{West Saxon normalization unmarked}!scofl@\iventry{sċofl}{}}
\index[iv]{03pgmc@\ivlangheader{Proto-Germanic}{}!skuflo@\iventry{*skúflō}{}}

Derivation: \emph{*skúflō} \textgreater{} \emph{sċofl} (regular).

\subsection{Derivation trace}\label{derivation-trace-60}

Proto input: \emph{*skúflō}

\begingroup
\setlength{\fboxsep}{6pt}

\noindent\fbox{%
\begin{minipage}{0.97\linewidth}
\small
\begin{tabular}{@{}>{\raggedright\arraybackslash}p{0.485\linewidth}@{\hspace{0.030\linewidth}}>{\raggedright\arraybackslash}p{0.485\linewidth}@{}}
\begin{minipage}[t]{\linewidth}
\centering\textbf{Earlier Germanic changes}\par
\vspace{0.35em}
\raggedright
\centering\textbf{Northwest and West Germanic}\par
\raggedright
\vspace{0.2em}
\begin{tabularx}{\linewidth}{@{}>{\raggedright\arraybackslash}X@{\hspace{0.55em}}>{\raggedright\arraybackslash}p{0.14\linewidth}@{\hspace{0.25em}}}
\mbox{PNWGmc U Lowering} & \emph{*skóflō} \\
\mbox{PNWGmc Final Long O Raising} & \emph{*skóflu} \\
\end{tabularx}
\vspace{0.6em}
\centering\textbf{Early Anglo-Frisian}\par
\raggedright
\vspace{0.2em}
\raggedright [no change]\par
\end{minipage}
&
\begin{minipage}[t]{\linewidth}
\centering\textbf{Old English changes}\par
\vspace{0.35em}
\raggedright
\begin{tabularx}{\linewidth}{@{}>{\raggedright\arraybackslash}X@{\hspace{0.55em}}>{\raggedright\arraybackslash}p{0.14\linewidth}@{\hspace{0.25em}}}
\mbox{OE Sk Palatalization} & \emph{*ʃóflu} \\
\mbox{OE High Vowel Apocope} & \emph{*ʃófl} \\
\end{tabularx}
\end{minipage}
\\
\end{tabular}
\end{minipage}%
} \endgroup

Old English form: \emph{sċofl}

\subsection{Reconstruction and
development}\label{reconstruction-and-development-3}

The selected short-root-vowel \Recon{skúflō} `shovel' follows the
feminine formation discussed by Orel (\citeproc{ref-Orel2003}{Orel 2003,
347}). With lowering before final oral raising, its root becomes
\emph{*o} before the ending becomes \emph{*u}; subsequent apocope and
consonant development give \emph{sċofl} `shovel'. No later root-vowel
levelling is needed for the form compared here.

This is a conditional witness, not independent proof of the order.
Earlier raising instead predicts a high-root nominative beside a lowered
oblique. Luick's inferred \emph{*scufel} beside \emph{scofel} `shovel'
makes alternation credible but does not directly establish
o-generalization in the early English paradigm
(\citeproc{ref-Luick1914}{Luick 1914-\/-1940, 105--6}, §78). The
assessment of that additional analogy, alongside wool and the other
ō-stems, belongs to \hyperref[rule-PNWGmcFinalLongORaising]{SC019
PNWGmcFinalLongORaising}.

\section{\texorpdfstring{show --- OE
\emph{sċēawian}}{show --- OE sċēawian}}\label{show-oe-sux10bux113awian}

\index[iv]{02oe@\ivlangheader{Old English}{West Saxon normalization unmarked}!sceawian@\iventry{sċēawian}{}}
\index[iv]{03pgmc@\ivlangheader{Proto-Germanic}{}!skawojana@\iventry{*skáwōjaną}{}}

Derivation: \emph{*skáwōjaną} \textgreater{} \emph{sċēawian} (regular).

\subsection{Derivation trace}\label{derivation-trace-61}

Proto input: \emph{*skáwōjaną}

\begingroup
\setlength{\fboxsep}{6pt}

\noindent\fbox{%
\begin{minipage}{0.97\linewidth}
\footnotesize
\begin{tabular}{@{}>{\raggedright\arraybackslash}p{0.339\linewidth}@{\hspace{0.030\linewidth}}>{\raggedright\arraybackslash}p{0.630\linewidth}@{}}
\begin{minipage}[t]{\linewidth}
\centering\textbf{Earlier Germanic changes}\par
\vspace{0.35em}
\raggedright
\centering\textbf{Northwest and West Germanic}\par
\raggedright
\vspace{0.2em}
\raggedright [no change]\par
\vspace{0.6em}
\centering\textbf{Early Anglo-Frisian}\par
\raggedright
\vspace{0.2em}
\raggedright [no change]\par
\end{minipage}
&
\begin{minipage}[t]{\linewidth}
\centering\textbf{Old English changes}\par
\vspace{0.35em}
\raggedright
\begin{tabularx}{\linewidth}{@{}>{\raggedright\arraybackslash}X@{\hspace{0.55em}}>{\raggedright\arraybackslash}p{0.26\linewidth}@{\hspace{0.25em}}}
\mbox{OE Aw Long Diphthong} & \emph{*skḗawōjaną} \\
OE Heavy Syllable Nasal Apocope & \emph{*skḗawōjan} \\
\mbox{OE Secondary Nasalization} & \emph{*skḗawōjąn} \\
\mbox{OE Sk Palatalization} & \emph{*ʃḗawōjąn} \\
\mbox{OE I Umlaut} & \emph{*ʃḗawējąn} \\
\mbox{OE Med Long Short} & \emph{*ʃḗawejąn} \\
\mbox{OE Weak Tail Reduction} & \emph{*ʃḗawejan} \\
OE Intervocalic J Vocalization & \emph{*ʃḗaweian} \\
\mbox{OE Unstressed EI Contraction} & \emph{*ʃḗawian} \\
\end{tabularx}
\end{minipage}
\\
\end{tabular}
\end{minipage}%
} \endgroup

Old English form: \emph{sċēawian}

\subsection{Reconstruction and comparative
evidence}\label{reconstruction-and-comparative-evidence-49}

Orel and Kroonen cite a Class II verb of the type
{\emph{*skawōjan-}}\index[iv]{03pgmc@\ivlangheader{Proto-Germanic}{}!skawojan@\iventry{*skawōjan-}{}},
with OE
{\emph{scēawian}}\index[iv]{02oe@\ivlangheader{Old English}{West Saxon normalization unmarked}!sceawian@\iventry{scēawian}{}}
`show' among the reflexes (\citeproc{ref-Orel2003}{Orel 2003, 337};
\citeproc{ref-Kroonen2013}{Kroonen 2013, 442}). The dictionary singleton
stem is not silently identified with Campbell and Brunner's
inherited-geminate analysis (\citeproc{ref-Campbell1959}{Campbell 1959,
45}; \citeproc{ref-SieversBrunner1965}{Brunner 1965, 49}). Brunner also
records the later West Saxon variant
{\emph{scāwian}}\index[iv]{02oe@\ivlangheader{Old English}{West Saxon normalization unmarked}!scawian@\iventry{scāwian}{}}
`show' in his discussion of diphthong accent shift; it is not evidence
for the same prehistoric singleton quantity
(\citeproc{ref-SieversBrunner1965}{Brunner 1965, 106--7}).

\subsection{Old English evidence}\label{old-english-evidence-50}

The held fourth edition of Bright's reader contains
{\emph{scēawian}}\index[iv]{02oe@\ivlangheader{Old English}{West Saxon normalization unmarked}!sceawian@\iventry{scēawian}{}}
`show' and the related imperative
{\emph{scēawa}}\index[iv]{02oe@\ivlangheader{Old English}{West Saxon normalization unmarked}!sceawa@\iventry{scēawa}{}}
`show' in the Blickling homily (\citeproc{ref-Bright1917}{Bright 1917,
71}). This is not a newly verified passage in the separate
Cassidy/Ringler edition. The sources use
{\emph{scēawian}}\index[iv]{02oe@\ivlangheader{Old English}{West Saxon normalization unmarked}!sceawian@\iventry{scēawian}{}}
`show', while
{\emph{sċēawian}}\index[iv]{02oe@\ivlangheader{Old English}{West Saxon normalization unmarked}!sceawian@\iventry{sċēawian}{}}
`show' supplies CAPR's normalized spelling.

\subsection{Development to Old
English}\label{development-to-old-english-49}

From
\Recon{skáwōjaną}\index[iv]{03pgmc@\ivlangheader{Proto-Germanic}{}!skawojana@\iventry{*skáwōjaną}{}}
`show', the transducer retains \hyperref[rule-OEAwLongDiphthong]{SC034
OEAwLongDiphthong} as a computational proxy mapping to long \emph{ea+w}
before the class-II suffix. Its \emph{*ō} separates root \emph{*w} from
\emph{*j}, excluding direct \emph{aw+j} gemination. Neither that
exclusion nor the successful
{\emph{sċēawian}}\index[iv]{02oe@\ivlangheader{Old English}{West Saxon normalization unmarked}!sceawian@\iventry{sċēawian}{}}
`show' output establishes an exceptionless singleton lengthening law.
The derivational input and \emph{regular} computational classification
remain unchanged, conditional on this retained proxy; their historical
basis is provisional, not a resolution of the root-history premise.
Campbell and Brunner's inherited-geminate account differs from the
dictionary singleton stem (\citeproc{ref-Campbell1959}{Campbell 1959,
45--46}; \citeproc{ref-SieversBrunner1965}{Brunner 1965, 49};
\citeproc{ref-Orel2003}{Orel 2003, 337};
\citeproc{ref-Kroonen2013}{Kroonen 2013, 442}).

The finite companions select an imperative singular and a present third
singular, not alternative dictionary infinitives. Their class-II endings
are independently supported even though successful root outputs cannot
decide the disputed prehistoric quantity
(\citeproc{ref-RingeTaylor2014}{Ringe and Taylor 2014, 80, 314};
\citeproc{ref-Campbell1959}{Campbell 1959, 148}).

\subsection{Form note}\label{form-note-15}

The difference between
{\emph{scēawian}}\index[iv]{02oe@\ivlangheader{Old English}{West Saxon normalization unmarked}!sceawian@\iventry{scēawian}{}}
`show' and
{\emph{sċēawian}}\index[iv]{02oe@\ivlangheader{Old English}{West Saxon normalization unmarked}!sceawian@\iventry{sċēawian}{}}
`show' is CAPR's editorial normalization of initial
\textless{}\emph{sc}\textgreater, not a different lexeme or paradigm
cell. It does not itself settle the consonant's prehistoric
articulation.

\section{\texorpdfstring{sleep --- OE
\emph{slǣpan}}{sleep --- OE slǣpan}}\label{sleep-oe-slux1e3pan}

\index[iv]{02oe@\ivlangheader{Old English}{West Saxon normalization unmarked}!slaepan@\iventry{slǣpan}{}}
\index[iv]{03pgmc@\ivlangheader{Proto-Germanic}{}!slepana@\iventry{*slḗpaną}{}}

Derivation: \emph{*slḗpaną} \textgreater{} \emph{slǣpan} (regular).

\subsection{Derivation trace}\label{derivation-trace-62}

Proto input: \emph{*slḗpaną}

\begingroup
\setlength{\fboxsep}{6pt}

\noindent\fbox{%
\begin{minipage}{0.97\linewidth}
\small
\begin{tabular}{@{}>{\raggedright\arraybackslash}p{0.388\linewidth}@{\hspace{0.030\linewidth}}>{\raggedright\arraybackslash}p{0.582\linewidth}@{}}
\begin{minipage}[t]{\linewidth}
\centering\textbf{Earlier Germanic changes}\par
\vspace{0.35em}
\raggedright
\centering\textbf{Northwest and West Germanic}\par
\raggedright
\vspace{0.2em}
\begin{tabularx}{\linewidth}{@{}>{\raggedright\arraybackslash}X@{\hspace{0.55em}}>{\raggedright\arraybackslash}p{0.24\linewidth}@{\hspace{0.25em}}}
\mbox{PNWGmc Long E Lowering} & \emph{*slāpaną} \\
\mbox{EAF Long A Fronting} & \emph{*slǣpaną} \\
\end{tabularx}
\vspace{0.6em}
\centering\textbf{Early Anglo-Frisian}\par
\raggedright
\vspace{0.2em}
\raggedright [no change]\par
\end{minipage}
&
\begin{minipage}[t]{\linewidth}
\centering\textbf{Old English changes}\par
\vspace{0.35em}
\raggedright
\begin{tabularx}{\linewidth}{@{}>{\raggedright\arraybackslash}X@{\hspace{0.55em}}>{\raggedright\arraybackslash}p{0.14\linewidth}@{\hspace{0.25em}}}
OE Heavy Syllable Nasal Apocope & \emph{*slǣpan} \\
\mbox{OE Secondary Nasalization} & \emph{*slǣpąn} \\
\mbox{OE Weak Tail Reduction} & \emph{*slǣpan} \\
\end{tabularx}
\end{minipage}
\\
\end{tabular}
\end{minipage}%
} \endgroup

Old English form: \emph{slǣpan}

\subsection{Reconstruction and comparative
evidence}\label{reconstruction-and-comparative-evidence-50}

Kroonen preserves the comparative verb as \emph{*slēpan-}, and Fulk
cites the same family under root \emph{*slēb-}
(\citeproc{ref-Kroonen2013}{Kroonen 2013}; \citeproc{ref-Fulk2018}{Fulk
2018, 120}). The derivational input \Recon{slḗpaną} `sleep' is the
infinitive-style form used here for that inherited sleep-verb.

\subsection{Old English evidence}\label{old-english-evidence-51}

Clark Hall gives
{\emph{slæpan}}\index[iv]{02oe@\ivlangheader{Old English}{West Saxon normalization unmarked}!slaepan@\iventry{slæpan}{}}
`sleep' with preterite \emph{slēp} `slept', \emph{slēap} `slept', and
Bright likewise lists
{\emph{slæpan}}\index[iv]{02oe@\ivlangheader{Old English}{West Saxon normalization unmarked}!slaepan@\iventry{slæpan}{}}
`sleep'
({\emph{slāpan}}\index[iv]{02oe@\ivlangheader{Old English}{West Saxon normalization unmarked}!slapan@\iventry{slāpan}{}}
`sleep'), \emph{slēp} `slept', \emph{slēpon} `slept', \emph{slēpen}
`slept' (\citeproc{ref-ClarkHall1960}{Clark Hall 1960};
\citeproc{ref-BrightCassidyRingler1971}{Bright 1971, 435}). The target
represented here is therefore the normalized infinitive \emph{slǣpan}
`sleep', not the preterite forms and not the separate noun \emph{slǣp}
`sleep'.

\subsection{Development to Old
English}\label{development-to-old-english-50}

From \Recon{slḗpaną} `sleep', Northwest Germanic lowering gives
\Recon{slǣpaną} `sleep'. The later OE tail developments then yield
{\emph{slǣpan}}\index[iv]{02oe@\ivlangheader{Old English}{West Saxon normalization unmarked}!slaepan@\iventry{slǣpan}{}}
`sleep' regularly. Brunner and Bülbring show that the OE tradition also
has variant spellings such as West Saxon
{\emph{slāpan}}\index[iv]{02oe@\ivlangheader{Old English}{West Saxon normalization unmarked}!slapan@\iventry{slāpan}{}}
`sleep' /
{\emph{slæpan}}\index[iv]{02oe@\ivlangheader{Old English}{West Saxon normalization unmarked}!slaepan@\iventry{slæpan}{}}
`sleep' and Anglian or Kentish
{\emph{slēpan}}\index[iv]{02oe@\ivlangheader{Old English}{West Saxon normalization unmarked}!slepan@\iventry{slēpan}{}}
`sleep', but those do not displace the infinitive chosen here
(\citeproc{ref-SieversBrunner1965}{Brunner 1965};
\citeproc{ref-Bulbring1902}{Bülbring 1902}).

\subsection{Form note}\label{form-note-16}

The note concerns lemma type rather than a special derivational problem:
this row represents the verb \emph{slǣpan} `sleep', whereas \emph{slǣp}
`sleep' belongs to noun or lookup background and \emph{slēp} `slept' /
\emph{slēap} `slept' are preterite forms
(\citeproc{ref-ClarkHall1960}{Clark Hall 1960};
\citeproc{ref-BrightCassidyRingler1971}{Bright 1971, 435}).

\section{\texorpdfstring{smear --- OE
\emph{smierwan}}{smear --- OE smierwan}}\label{smear-oe-smierwan}

\index[iv]{02oe@\ivlangheader{Old English}{West Saxon normalization unmarked}!smierwan@\iventry{smierwan}{}}
\index[iv]{03pgmc@\ivlangheader{Proto-Germanic}{}!smerwijana@\iventry{*smérwijaną}{}}

Derivation: \emph{*smérwijaną} \textgreater{} \emph{smierwan} (regular).

\subsection{Derivation trace}\label{derivation-trace-63}

Proto input: \emph{*smérwijaną}

\begingroup
\setlength{\fboxsep}{6pt}

\noindent\fbox{%
\begin{minipage}{0.97\linewidth}
\footnotesize
\begin{tabular}{@{}>{\raggedright\arraybackslash}p{0.339\linewidth}@{\hspace{0.030\linewidth}}>{\raggedright\arraybackslash}p{0.630\linewidth}@{}}
\begin{minipage}[t]{\linewidth}
\centering\textbf{Earlier Germanic changes}\par
\vspace{0.35em}
\raggedright
\centering\textbf{Northwest and West Germanic}\par
\raggedright
\vspace{0.2em}
\raggedright [no change]\par
\vspace{0.6em}
\centering\textbf{Early Anglo-Frisian}\par
\raggedright
\vspace{0.2em}
\raggedright [no change]\par
\end{minipage}
&
\begin{minipage}[t]{\linewidth}
\centering\textbf{Old English changes}\par
\vspace{0.35em}
\raggedright
\begin{tabularx}{\linewidth}{@{}>{\raggedright\arraybackslash}X@{\hspace{0.55em}}>{\raggedright\arraybackslash}p{0.26\linewidth}@{\hspace{0.25em}}}
\mbox{OE Breaking} & \emph{*sméorwijaną} \\
OE Heavy Syllable Nasal Apocope & \emph{*sméorwijan} \\
\mbox{OE Secondary Nasalization} & \emph{*sméorwijąn} \\
\mbox{Sievers Law Syncope} & \emph{*sméorwjąn} \\
\mbox{OE I Umlaut} & \emph{*smíerwjąn} \\
\mbox{OE Weak Tail Reduction} & \emph{*smíerwjan} \\
\mbox{OE J Loss After Heavy} & \emph{*smíerwan} \\
\end{tabularx}
\end{minipage}
\\
\end{tabular}
\end{minipage}%
} \endgroup

Old English form: \emph{smierwan}

\subsection{Reconstruction and comparative
evidence}\label{reconstruction-and-comparative-evidence-51}

Kroonen gives the comparative headword as \emph{*smerwjan-}
(\citeproc{ref-Kroonen2013}{Kroonen 2013}). Ringe and Taylor instead
cite a later-stage \Recon{smirwijana} `smear', from which they derive
West Saxon
{\emph{smierwan}}\index[iv]{02oe@\ivlangheader{Old English}{West Saxon normalization unmarked}!smierwan@\iventry{smierwan}{}}
`smear', Mercian
{\emph{smirwan}}\index[iv]{02oe@\ivlangheader{Old English}{West Saxon normalization unmarked}!smirwan~05@\iventry{smirwan}{Merc.}}
`smear', and Northumbrian \emph{smiriga}
(\citeproc{ref-RingeTaylor2014}{Ringe and Taylor 2014}). The
derivational input \Recon{smérwijaną} `smear' therefore represents the
Kroonen-aligned PGmc layer, while the later-stage dialect split belongs
to a different chronological level.

\subsection{Old English evidence}\label{old-english-evidence-52}

The target represented here is the West Saxon citation form
{\emph{smierwan}}\index[iv]{02oe@\ivlangheader{Old English}{West Saxon normalization unmarked}!smierwan@\iventry{smierwan}{}}
`smear'. Campbell's Anglian discussion explains the contrasting
{\emph{smirwan}}\index[iv]{02oe@\ivlangheader{Old English}{West Saxon normalization unmarked}!smirwan@\iventry{smirwan}{}}
`smear', and Clark Hall, Brunner, and Bright show that the same lexical
family later also includes forms such as \emph{smirian}, \emph{smyrian},
and preterite \emph{smyrode} (\citeproc{ref-Campbell1959}{Campbell
1959}; \citeproc{ref-ClarkHall1960}{Clark Hall 1960};
\citeproc{ref-SieversBrunner1965}{Brunner 1965};
\citeproc{ref-BrightCassidyRingler1971}{Bright 1971}).

\subsection{Development to Old
English}\label{development-to-old-english-51}

From \Recon{smérwijaną} `smear', breaking before \emph{r + consonant}
yields \emph{eo}, and later i-umlaut produces \emph{ie}. The result is
West Saxon
{\emph{smierwan}}\index[iv]{02oe@\ivlangheader{Old English}{West Saxon normalization unmarked}!smierwan@\iventry{smierwan}{}}
`smear'. Anglian
{\emph{smirwan}}\index[iv]{02oe@\ivlangheader{Old English}{West Saxon normalization unmarked}!smirwan~04@\iventry{smirwan}{Angl.}}
`smear' reflects the well-known failure of breaking in this environment,
not a different lexeme (\citeproc{ref-Campbell1959}{Campbell 1959};
\citeproc{ref-RingeTaylor2014}{Ringe and Taylor 2014}).

\subsection{Dialect note}\label{dialect-note-5}

The entry therefore represents the West Saxon member of a broader OE
family:
{\emph{smierwan}}\index[iv]{02oe@\ivlangheader{Old English}{West Saxon normalization unmarked}!smierwan@\iventry{smierwan}{}}
`smear' in West Saxon,
{\emph{smirwan}}\index[iv]{02oe@\ivlangheader{Old English}{West Saxon normalization unmarked}!smirwan@\iventry{smirwan}{}}
`smear' in Anglian or Mercian, and related later class-II forms such as
\emph{smirian} or \emph{smyrian} in the same lexical field
(\citeproc{ref-RingeTaylor2014}{Ringe and Taylor 2014};
\citeproc{ref-ClarkHall1960}{Clark Hall 1960}).

\section{\texorpdfstring{sorrow --- OE
\emph{sorg}}{sorrow --- OE sorg}}\label{sorrow-oe-sorg}

\index[iv]{02oe@\ivlangheader{Old English}{West Saxon normalization unmarked}!sorg@\iventry{sorg}{}}
\index[iv]{03pgmc@\ivlangheader{Proto-Germanic}{}!surgo@\iventry{*súrgō}{}}

Derivation: \emph{*súrgō} \textgreater{} \emph{sorg} (regular).

\subsection{Derivation trace}\label{derivation-trace-64}

Proto input: \emph{*súrgō}

\begingroup
\setlength{\fboxsep}{6pt}

\noindent\fbox{%
\begin{minipage}{0.97\linewidth}
\small
\begin{tabular}{@{}>{\raggedright\arraybackslash}p{0.611\linewidth}@{\hspace{0.030\linewidth}}>{\raggedright\arraybackslash}p{0.359\linewidth}@{}}
\begin{minipage}[t]{\linewidth}
\centering\textbf{Earlier Germanic changes}\par
\vspace{0.35em}
\raggedright
\centering\textbf{Northwest and West Germanic}\par
\raggedright
\vspace{0.2em}
\begin{tabularx}{\linewidth}{@{}>{\raggedright\arraybackslash}X@{\hspace{0.55em}}>{\raggedright\arraybackslash}p{0.14\linewidth}@{\hspace{0.25em}}}
\mbox{PNWGmc U Lowering} & \emph{*sórgō} \\
\mbox{PNWGmc Final Long O Raising} & \emph{*sórgu} \\
\end{tabularx}
\vspace{0.6em}
\centering\textbf{Early Anglo-Frisian}\par
\raggedright
\vspace{0.2em}
\raggedright [no change]\par
\end{minipage}
&
\begin{minipage}[t]{\linewidth}
\centering\textbf{Old English changes}\par
\vspace{0.35em}
\raggedright
\begin{tabularx}{\linewidth}{@{}>{\raggedright\arraybackslash}X@{\hspace{0.55em}}>{\raggedright\arraybackslash}p{0.14\linewidth}@{\hspace{0.25em}}}
\mbox{OE High Vowel Apocope} & \emph{*sórg} \\
\end{tabularx}
\end{minipage}
\\
\end{tabular}
\end{minipage}%
} \endgroup

Old English form: \emph{sorg}

\subsection{Reconstruction and
development}\label{reconstruction-and-development-4}

The selected \Recon{súrgō} `sorrow' belongs to the feminine formation
discussed by Kroonen (\citeproc{ref-Kroonen2013}{Kroonen 2013, 493}).
Lowering before final oral raising gives a low root before the ending
becomes \emph{*u}; later apocope yields \emph{sorg} `sorrow'. This
accounts for the selected form without an additional root-vowel
generalization.

Earlier raising would instead protect the nominative root while leaving
the nasal accusative's non-high conditioner available for lowering.
Generalization of that oblique root could recover the low-root paradigm,
but has not been independently established for this particular early
English noun. Generic evidence for such levelling is not specific
attestation of this operation (\citeproc{ref-RingeTaylor2014}{Ringe and
Taylor 2014, 32--34}; \citeproc{ref-Stiles2012}{Stiles 2012, 47--48}).
The comparative economy argument and its limits are developed in
\hyperref[rule-PNWGmcFinalLongORaising]{SC019 PNWGmcFinalLongORaising}.

\section{\texorpdfstring{soul --- OE
\emph{sāwle}}{soul --- OE sāwle}}\label{soul-oe-sux101wle}

\index[iv]{02oe@\ivlangheader{Old English}{West Saxon normalization unmarked}!sawle@\iventry{sāwle}{}}
\index[iv]{03pgmc@\ivlangheader{Proto-Germanic}{}!saiwalo@\iventry{*sáiwalō}{}}
\index[iv]{03pgmc@\ivlangheader{Proto-Germanic}{}!saiwalo@\iventry{*sáiwalǭ}{}}

Derivation: citation reconstruction \emph{*sáiwalō}; form followed here
\emph{*sáiwalǭ} \textgreater{} \emph{sāwle} (regular).

\subsection{Derivation trace}\label{derivation-trace-65}

Proto input: \emph{*sáiwalǭ}

\begingroup
\setlength{\fboxsep}{6pt}

\noindent\fbox{%
\begin{minipage}{0.97\linewidth}
\small
\begin{tabular}{@{}>{\raggedright\arraybackslash}p{0.485\linewidth}@{\hspace{0.030\linewidth}}>{\raggedright\arraybackslash}p{0.485\linewidth}@{}}
\begin{minipage}[t]{\linewidth}
\centering\textbf{Earlier Germanic changes}\par
\vspace{0.35em}
\raggedright
\centering\textbf{Northwest and West Germanic}\par
\raggedright
\vspace{0.2em}
\raggedright [no change]\par
\vspace{0.6em}
\centering\textbf{Early Anglo-Frisian}\par
\raggedright
\vspace{0.2em}
\begin{tabularx}{\linewidth}{@{}>{\raggedright\arraybackslash}X@{\hspace{0.55em}}>{\raggedright\arraybackslash}p{0.14\linewidth}@{\hspace{0.25em}}}
\mbox{EAF Ai Monophthongization} & \emph{*sāwalǭ} \\
\end{tabularx}
\end{minipage}
&
\begin{minipage}[t]{\linewidth}
\centering\textbf{Old English changes}\par
\vspace{0.35em}
\raggedright
\begin{tabularx}{\linewidth}{@{}>{\raggedright\arraybackslash}X@{\hspace{0.55em}}>{\raggedright\arraybackslash}p{0.14\linewidth}@{\hspace{0.25em}}}
\mbox{EAF Brightening Unstressed} & \emph{*sāwælǭ} \\
\mbox{OE A Restoration} & \emph{*sāwalǭ} \\
\mbox{OE Non High Syn} & \emph{*sāwlǭ} \\
OE Unstressed Long Vowel Shortening & \emph{*sāwlæ} \\
\mbox{OE Unstressed AE Merger} & \emph{*sāwle} \\
\end{tabularx}
\end{minipage}
\\
\end{tabular}
\end{minipage}%
} \endgroup

Old English form: \emph{sāwle}

\subsection{Reconstruction and comparative
evidence}\label{reconstruction-and-comparative-evidence-52}

The feminine soul formation is represented by Orel's \emph{*saiwalō} and
Kroonen's \emph{*saiwalō-}. Orel's sea/sacred-lake account and the
alternative associations he reports concern the deeper etymology;
Kroonen's qualifications about its comparative distribution remain
independent of the selected Old English case
(\citeproc{ref-Orel2003}{Orel 2003, 314};
\citeproc{ref-Kroonen2013}{Kroonen 2013, 423}). The comparative citation
\Recon{sáiwalō}\index[iv]{03pgmc@\ivlangheader{Proto-Germanic}{}!saiwalo@\iventry{*sáiwalō}{}}
`soul' is retained, but is not silently equated with the selected
accusative
\Recon{sáiwalǭ}\index[iv]{03pgmc@\ivlangheader{Proto-Germanic}{}!saiwalo@\iventry{*sáiwalǭ}{}}.
The latter is an inherited inflectional reconstruction in CAPR notation,
not a quotation of a dictionary's whole accented accusative.

\subsection{Old English evidence}\label{old-english-evidence-53}

Luke12:20 supplies the selected accusative: \emph{hig feccað þine sawle
fram þe}, `they will demand your soul from you'. Skeat prints it in the
first West Saxon column, the Corpus Christi College Cambridge140 text,
with other manuscripts collated (\citeproc{ref-Skeat1874}{Skeat 1874,
v--vi, 130}). The quantity-normalized target
{\emph{sāwle}}\index[iv]{02oe@\ivlangheader{Old English}{West Saxon normalization unmarked}!sawle@\iventry{sāwle}{}}
`soul' thus has a specific syntactic identification, not merely an
ending compatible with several cases.

The same page preserves dative \emph{minre sawle} at verse19 and
vocative \emph{eala sawel}. That vocative must not be relabelled
nominative. The actual nominative is \emph{seo sawul ys ma} at verse23.
These are real competing paradigm forms; a disyllabic spelling alone is
not proof of spoken syllabicity or its date. Fulk's independent metrical
argument deserves separate weight from such spellings
(\citeproc{ref-Fulk2010}{Fulk 2010, 133--39}).

\subsection{Development to Old
English}\label{development-to-old-english-52}

Monophthongization of atomic stressed \emph{ai} gives long \emph{ā}.
Unstressed brightening and restoration lead to
\Recon{sāwalǭ}\index[iv]{08preoe@\ivlangheader{Pre-Old English / prehistoric English}{}!sawalo@\iventry{*sāwalǭ}{}}.
The short open posttonic vowel precedes a heavy final syllable, so
\hyperref[rule-OENonHighSyn]{SC114 OENonHighSyn} yields
\Recon{sāwlǭ}\index[iv]{08preoe@\ivlangheader{Pre-Old English / prehistoric English}{}!sawlo@\iventry{*sāwlǭ}{}}.
Later decoding/shortening of the inherited ending gives \emph{æ}, then
the unstressed merger gives {\emph{sāwle}}. No medial raising proxy,
epenthesis or analogy occurs within this selected regular span.

The written \emph{ǭ} carries inherited ending identity and quantity
until its existing late decoding. It does not assert that a nasal
rounded vowel physically persisted unchanged to that late date.
Grammatical identity licenses the input; the syncope rule itself refers
only to phonological and independently represented prosodic structure
(\citeproc{ref-HoggGrammar1992}{Hogg 1992a, 220--22}, §§6.14--6.18).

\subsection{Nominatives and wider paradigm
histories}\label{nominatives-and-wider-paradigm-histories}

Ringe and Taylor's nominative line passes through medial nonhigh
syncope, apocope and subsequent epenthesis. Their final high vowel is
restored after an earlier loss; that premise cannot be omitted from the
argument (\citeproc{ref-RingeTaylor2014}{Ringe and Taylor 2014, 265,
330}). Campbell's
{\emph{sāwl}}\index[iv]{02oe@\ivlangheader{Old English}{West Saxon normalization unmarked}!sawl@\iventry{sāwl}{}}
`soul', Luick's expanded
{\emph{sāwul}}\index[iv]{02oe@\ivlangheader{Old English}{West Saxon normalization unmarked}!sawul@\iventry{sāwul}{}}
`soul', and
{\emph{sāwol}}\index[iv]{02oe@\ivlangheader{Old English}{West Saxon normalization unmarked}!sawol@\iventry{sāwol}{}}
`soul' are not interchangeable phonological stages inferred solely from
their spellings (\citeproc{ref-Campbell1959}{Campbell 1959, 143}, §341;
\citeproc{ref-Luick1914}{Luick 1914-\/-1940, 287, 294}, §§309, 317;
\citeproc{ref-SieversBrunner1965}{Brunner 1965, 132}, §159). Fulk's
alternative reconstruction of the early syllabic history remains
relevant; selecting the inherited accusative does not adjudicate every
nominative or the complete remodelling of the paradigm.

The former selected citation
\Recon{sáiwalō}\index[iv]{03pgmc@\ivlangheader{Proto-Germanic}{}!saiwalo@\iventry{*sáiwalō}{}}
no longer reaches \emph{sāwol} by continuous medial raising: that
temporary operation is retired. Its counterfactual current output is
\emph{sāwel}. This result is disclosed, not promoted as a newly
explained nominative. The \emph{regular} classification applies to the
explicitly selected Proto-Germanic-to-Old-English accusative span, not
to an unchanged whole paradigm.

\section{\texorpdfstring{span --- OE
\emph{spannan}}{span --- OE spannan}}\label{span-oe-spannan}

\index[iv]{02oe@\ivlangheader{Old English}{West Saxon normalization unmarked}!spannan@\iventry{spannan}{}}
\index[iv]{03pgmc@\ivlangheader{Proto-Germanic}{}!spannana@\iventry{*spánnaną}{}}

Derivation: \emph{*spánnaną} \textgreater{} \emph{spannan} (regular).

\subsection{Derivation trace}\label{derivation-trace-66}

Proto input: \emph{*spánnaną}

\begingroup
\setlength{\fboxsep}{6pt}

\noindent\fbox{%
\begin{minipage}{0.97\linewidth}
\small
\begin{tabular}{@{}>{\raggedright\arraybackslash}p{0.359\linewidth}@{\hspace{0.030\linewidth}}>{\raggedright\arraybackslash}p{0.611\linewidth}@{}}
\begin{minipage}[t]{\linewidth}
\centering\textbf{Earlier Germanic changes}\par
\vspace{0.35em}
\raggedright
\centering\textbf{Northwest and West Germanic}\par
\raggedright
\vspace{0.2em}
\raggedright [no change]\par
\vspace{0.6em}
\centering\textbf{Early Anglo-Frisian}\par
\raggedright
\vspace{0.2em}
\raggedright [no change]\par
\end{minipage}
&
\begin{minipage}[t]{\linewidth}
\centering\textbf{Old English changes}\par
\vspace{0.35em}
\raggedright
\begin{tabularx}{\linewidth}{@{}>{\raggedright\arraybackslash}X@{\hspace{0.55em}}>{\raggedright\arraybackslash}p{0.18\linewidth}@{\hspace{0.25em}}}
OE Heavy Syllable Nasal Apocope & \emph{*spánnan} \\
\mbox{OE Secondary Nasalization} & \emph{*spánnąn} \\
\mbox{OE Weak Tail Reduction} & \emph{*spánnan} \\
\end{tabularx}
\end{minipage}
\\
\end{tabular}
\end{minipage}%
} \endgroup

Old English form: \emph{spannan}

\subsection{Reconstruction and comparative
evidence}\label{reconstruction-and-comparative-evidence-53}

Kroonen cites the inherited verb as \emph{*spannan-}, with OE
{\emph{spannan}}\index[iv]{02oe@\ivlangheader{Old English}{West Saxon normalization unmarked}!spannan@\iventry{spannan}{}}
`span, fasten' among the reflexes (\citeproc{ref-Kroonen2013}{Kroonen
2013}). The derivational input \Recon{spánnaną} `span' is the
infinitive-style form used here for that same verbal lexeme.

\subsection{Old English evidence}\label{old-english-evidence-54}

Clark Hall lists noun \emph{spann} and verb
{\emph{spannan}}\index[iv]{02oe@\ivlangheader{Old English}{West Saxon normalization unmarked}!spannan@\iventry{spannan}{}}
`span, fasten' as separate headwords, and Brunner likewise records
\emph{sponnan, spannan stv.} (\citeproc{ref-ClarkHall1960}{Clark Hall
1960}; \citeproc{ref-SieversBrunner1965}{Brunner 1965}). The target is
the strong-verb infinitive, not the noun.

\subsection{Development to Old
English}\label{development-to-old-english-53}

From \Recon{spánnaną} `span', the final nasal ending is lost and the
regular OE weak-tail steps surface
{\emph{spannan}}\index[iv]{02oe@\ivlangheader{Old English}{West Saxon normalization unmarked}!spannan@\iventry{spannan}{}}
`span, fasten'. No paradigm-cell substitution is needed: the current
derivation already lands on the infinitive directly.

\subsection{Form note}\label{form-note-17}

English \emph{span} can also reach the noun \emph{spann} in local lookup
material. The form represented here is the verb
{\emph{spannan}}\index[iv]{02oe@\ivlangheader{Old English}{West Saxon normalization unmarked}!spannan@\iventry{spannan}{}}
`span, fasten', with the noun treated elsewhere
(\citeproc{ref-ClarkHall1960}{Clark Hall 1960}).

\section{\texorpdfstring{spar --- OE
\emph{spearra}}{spar --- OE spearra}}\label{spar-oe-spearra}

\index[iv]{02oe@\ivlangheader{Old English}{West Saxon normalization unmarked}!spearra@\iventry{spearra}{}}
\index[iv]{03pgmc@\ivlangheader{Proto-Germanic}{}!sparro@\iventry{*spárrô}{}}

Derivation: \emph{*spárrô} \textgreater{} \emph{spearra} (regular).

\subsection{Derivation trace}\label{derivation-trace-67}

Proto input: \emph{*spárrô}

\begingroup
\setlength{\fboxsep}{6pt}

\noindent\fbox{%
\begin{minipage}{0.97\linewidth}
\small
\begin{tabular}{@{}>{\raggedright\arraybackslash}p{0.359\linewidth}@{\hspace{0.030\linewidth}}>{\raggedright\arraybackslash}p{0.611\linewidth}@{}}
\begin{minipage}[t]{\linewidth}
\centering\textbf{Earlier Germanic changes}\par
\vspace{0.35em}
\raggedright
\centering\textbf{Northwest and West Germanic}\par
\raggedright
\vspace{0.2em}
\raggedright [no change]\par
\vspace{0.6em}
\centering\textbf{Early Anglo-Frisian}\par
\raggedright
\vspace{0.2em}
\raggedright [no change]\par
\end{minipage}
&
\begin{minipage}[t]{\linewidth}
\centering\textbf{Old English changes}\par
\vspace{0.35em}
\raggedright
\begin{tabularx}{\linewidth}{@{}>{\raggedright\arraybackslash}X@{\hspace{0.55em}}>{\raggedright\arraybackslash}p{0.18\linewidth}@{\hspace{0.25em}}}
\mbox{EAF Brightening} & \emph{*spærrô} \\
\mbox{OE Breaking} & \emph{*spearrô} \\
OE Unstressed Long Vowel Shortening & \emph{*spearra} \\
\end{tabularx}
\end{minipage}
\\
\end{tabular}
\end{minipage}%
} \endgroup

Old English form: \emph{spearra}

\subsection{Reconstruction and comparative
evidence}\label{reconstruction-and-comparative-evidence-54}

Kroonen and Orel place this noun in the beam or rafter set
\emph{*spar(r)an-}, with cognates such as Old Saxon and Old High German
{\emph{sparro}}\index[iv]{13os@\ivlangheader{Old Saxon}{}!sparro@\iventry{sparro}{}}
`spar, rafter' (\citeproc{ref-Kroonen2013}{Kroonen 2013};
\citeproc{ref-Orel2003}{Orel 2003}). The derivational input
\Recon{spárrô} `spar' is the OE-facing nominal form used here for that
same lexeme.

\subsection{Old English evidence}\label{old-english-evidence-55}

The noun is
{\emph{spearra}}\index[iv]{02oe@\ivlangheader{Old English}{West Saxon normalization unmarked}!spearra@\iventry{spearra}{}}
`spar, rafter'. English gloss overlap also reaches the unrelated verb
\emph{sperran} `to bar', which does not belong to this row.

\subsection{Development to Old
English}\label{development-to-old-english-54}

From \Recon{spárrô} `spar', Anglo-Frisian brightening gives
\Recon{spærrô} `spar', and OE breaking before geminate \emph{rr} yields
\Recon{spearrô} `spar', later
{\emph{spearra}}\index[iv]{02oe@\ivlangheader{Old English}{West Saxon normalization unmarked}!spearra@\iventry{spearra}{}}
`spar, rafter'. The development is therefore regular for a
breaking-conditioned noun of this type (\citeproc{ref-Luick1914}{Luick
1914-\/-1940}).

\subsection{Form note}\label{form-note-18}

This entry concerns the noun
{\emph{spearra}}\index[iv]{02oe@\ivlangheader{Old English}{West Saxon normalization unmarked}!spearra@\iventry{spearra}{}}
`spar, rafter' only. It should be kept separate from verb
\emph{sperran}, even though the Modern English glosses overlap
(\citeproc{ref-Kroonen2013}{Kroonen 2013}; \citeproc{ref-Orel2003}{Orel
2003}).

\section{\texorpdfstring{spindle --- OE
\emph{spinle}}{spindle --- OE spinle}}\label{spindle-oe-spinle}

\index[iv]{02oe@\ivlangheader{Old English}{West Saxon normalization unmarked}!spinle@\iventry{spinle}{}}
\index[iv]{03pgmc@\ivlangheader{Proto-Germanic}{}!spennilo@\iventry{*spénnilō}{}}
\index[iv]{03pgmc@\ivlangheader{Proto-Germanic}{}!spennilo@\iventry{*spénnilǭ}{}}

Derivation: citation reconstruction \emph{*spénnilō}; form followed here
\emph{*spénnilǭ} \textgreater{} \emph{spinle} (regular).

\subsection{Derivation trace}\label{derivation-trace-68}

Proto input: \emph{*spénnilǭ}

\begingroup
\setlength{\fboxsep}{6pt}

\noindent\fbox{%
\begin{minipage}{0.97\linewidth}
\small
\begin{tabular}{@{}>{\raggedright\arraybackslash}p{0.359\linewidth}@{\hspace{0.030\linewidth}}>{\raggedright\arraybackslash}p{0.611\linewidth}@{}}
\begin{minipage}[t]{\linewidth}
\centering\textbf{Earlier Germanic changes}\par
\vspace{0.35em}
\raggedright
\centering\textbf{Northwest and West Germanic}\par
\raggedright
\vspace{0.2em}
\raggedright [no change]\par
\vspace{0.6em}
\centering\textbf{Early Anglo-Frisian}\par
\raggedright
\vspace{0.2em}
\raggedright [no change]\par
\end{minipage}
&
\begin{minipage}[t]{\linewidth}
\centering\textbf{Old English changes}\par
\vspace{0.35em}
\raggedright
\begin{tabularx}{\linewidth}{@{}>{\raggedright\arraybackslash}X@{\hspace{0.55em}}>{\raggedright\arraybackslash}p{0.20\linewidth}@{\hspace{0.25em}}}
\mbox{OE I Umlaut} & \emph{*spinnilǭ} \\
\mbox{OE High Vowel Apocope} & \emph{*spinnlǭ} \\
OE Preconsonantal Degemination & \emph{*spinlǭ} \\
OE Unstressed Long Vowel Shortening & \emph{*spinlæ} \\
\mbox{OE Unstressed AE Merger} & \emph{*spinle} \\
\end{tabularx}
\end{minipage}
\\
\end{tabular}
\end{minipage}%
} \endgroup

Old English form: \emph{spinle}

\subsection{Old English evidence and selected
cell}\label{old-english-evidence-and-selected-cell-1}

Clark Hall identifies
{\emph{spinel}}\index[iv]{02oe@\ivlangheader{Old English}{West Saxon normalization unmarked}!spinel@\iventry{spinel}{}}
`spindle' as feminine and gives genitive singular
{\emph{spinle}}\index[iv]{02oe@\ivlangheader{Old English}{West Saxon normalization unmarked}!spinle@\iventry{spinle}{}}
(\citeproc{ref-ClarkHall1960}{Clark Hall 1960, 272, s.v}.~spinel). This
positively supports the strong feminine oblique form; the genitive label
is not misquoted as an accusative label. Toller's craft-equipment
passage supplies object usage: \emph{He sceal fela towtōla habban:
flexlīnan, spinle, reōl}, `he must have many tools: flax-lines, a
spindle, a reel'. The listed objects are governed by \emph{habban}
(\citeproc{ref-Toller1921}{Toller 1921, 621, s.v}.~lorh) (Anglia 9,
p.~263, line 11).

The already approved selection is the ordinary strong feminine
accusative singular
{\emph{spinle}}\index[iv]{02oe@\ivlangheader{Old English}{West Saxon normalization unmarked}!spinle@\iventry{spinle}{}}.
It is not a new morphological change made in this coverage audit. Its
vocalic ending is independently licensed by the strong feminine
paradigm, just as in the selected `soul' and `filth' histories. The
spelling
{\emph{spinl}}\index[iv]{02oe@\ivlangheader{Old English}{West Saxon normalization unmarked}!spinl@\iventry{spinl}{}},
also recorded by Clark Hall, remains a nonselected comparison, not the
current endingless target.

\subsection{Comparative premise and its unresolved
limit}\label{comparative-premise-and-its-unresolved-limit}

Orel prints
\Recon{spennilō(n)}\index[iv]{03pgmc@\ivlangheader{Proto-Germanic}{}!spennilo@\iventry{*spennilō(n)}{}}
`spindle', deriving it from the spin family and citing OE \emph{spinel}
(\citeproc{ref-Orel2003}{Orel 2003, 364}). Both the geminate and his
optional n-stem notation are preserved in attribution. CAPR's citation
\Recon{spénnilō}\index[iv]{03pgmc@\ivlangheader{Proto-Germanic}{}!spennilo@\iventry{*spénnilō}{}}
uses its ordinary accent notation and the strong feminine formation; the
selected
\Recon{spénnilǭ}\index[iv]{03pgmc@\ivlangheader{Proto-Germanic}{}!spenniloacc@\iventry{*spénnilǭ}{}}
supplies that formation's accusative ending.

The comparative origin and scope of the geminate remain unresolved.
Kroonen's spin verb and yarn noun are family evidence, not quotations of
this spindle noun (\citeproc{ref-Kroonen2013}{Kroonen 2013, 467}).
Successful output cannot establish where the doubled consonant arose,
settle the relation of the source's ending alternatives, or validate an
otherwise unsupported whole paradigm. The retained input is an explicit
comparative premise; this entry does not initiate a new gemination
investigation.

\subsection{Current derivation and historical
comparisons}\label{current-derivation-and-historical-comparisons}

\hyperref[rule-OEIUmlaut]{SC055 OEIUmlaut} gives \emph{*spinnilǭ}. The
heavy root and following quantity-bearing nonhigh syllable license
medial deletion at \hyperref[rule-OEHighVowelApocope]{SC063
OEHighVowelApocope}, giving \emph{*spinnlǭ}. This is part of the one
quantitative event, not productive syncope at
\hyperref[rule-OEMedialSyncope]{SC065 OEMedialSyncope}, which is now an
identity trace.

\hyperref[rule-OEPreconsonantalDegemination]{SC068
OEPreconsonantalDegemination} then gives \emph{*spinlǭ}; later
weak-ending developments give \emph{spinle}. The ordering controls are
exact: deletion before mutation gives \Pred{spenle}, deletion after
degemination gives \Pred{spinnle}, and omitting quantitative deletion
gives \Pred{spinnele}. These contrasts test the adopted history under
its stated comparative premise, not the independent truth of that
premise (\citeproc{ref-RingeTaylor2014}{Ringe and Taylor 2014, 270--72,
330}).

The former nominative calculation and its old \emph{spinl} target remain
historical records. Under the adopted HLL account, that encoded
nominative instead retains its weak syllables, giving \Pred{spinnelu};
it is not the current regular selected derivation. Conservative vowel
retention and the later paradigmatic alternatives remain relevant to the
larger syncope discussion (\citeproc{ref-Fulk2010}{Fulk 2010, 129--34,
140--42}; \citeproc{ref-Goering2016}{Goering 2016, 176--77, 186--88}).
They do not overturn the independently selected accusative or create a
second productive heavy-root deletion.

\section{\texorpdfstring{still --- OE
\emph{stillan}}{still --- OE stillan}}\label{still-oe-stillan}

\index[iv]{02oe@\ivlangheader{Old English}{West Saxon normalization unmarked}!stillan@\iventry{stillan}{}}
\index[iv]{03pgmc@\ivlangheader{Proto-Germanic}{}!stellijana@\iventry{*stéllijaną}{}}
\index[iv]{17german@\ivlangheader{German}{}!stillen@\iventry{stillen}{}}

Derivation: \emph{*stéllijaną} \textgreater{} \emph{stillan} (regular).

\subsection{Derivation trace}\label{derivation-trace-69}

Proto input: \emph{*stéllijaną}

\begingroup
\setlength{\fboxsep}{6pt}

\noindent\fbox{%
\begin{minipage}{0.97\linewidth}
\footnotesize
\begin{tabular}{@{}>{\raggedright\arraybackslash}p{0.339\linewidth}@{\hspace{0.030\linewidth}}>{\raggedright\arraybackslash}p{0.630\linewidth}@{}}
\begin{minipage}[t]{\linewidth}
\centering\textbf{Earlier Germanic changes}\par
\vspace{0.35em}
\raggedright
\centering\textbf{Northwest and West Germanic}\par
\raggedright
\vspace{0.2em}
\raggedright [no change]\par
\vspace{0.6em}
\centering\textbf{Early Anglo-Frisian}\par
\raggedright
\vspace{0.2em}
\raggedright [no change]\par
\end{minipage}
&
\begin{minipage}[t]{\linewidth}
\centering\textbf{Old English changes}\par
\vspace{0.35em}
\raggedright
\begin{tabularx}{\linewidth}{@{}>{\raggedright\arraybackslash}X@{\hspace{0.55em}}>{\raggedright\arraybackslash}p{0.24\linewidth}@{\hspace{0.25em}}}
OE Heavy Syllable Nasal Apocope & \emph{*stéllijan} \\
\mbox{OE Secondary Nasalization} & \emph{*stéllijąn} \\
\mbox{Sievers Law Syncope} & \emph{*stélljąn} \\
\mbox{OE I Umlaut} & \emph{*stilljąn} \\
\mbox{OE Weak Tail Reduction} & \emph{*stilljan} \\
\mbox{OE J Loss After Heavy} & \emph{*stillan} \\
\end{tabularx}
\end{minipage}
\\
\end{tabular}
\end{minipage}%
} \endgroup

Old English form: \emph{stillan}

\subsection{Reconstruction and comparative
evidence}\label{reconstruction-and-comparative-evidence-55}

The wider West Germanic family includes adjective
{\emph{still}}\index[iv]{02oe@\ivlangheader{Old English}{West Saxon normalization unmarked}!still@\iventry{still}{}}
`still, quiet' and verb \emph{stillen}
(\citeproc{ref-KlugeSeebold2011}{Kluge 2011}). The derivational input
\Recon{stéllijaną} `still' represents the verbal j-formation used for
the OE row.

\subsection{Old English evidence}\label{old-english-evidence-56}

Clark Hall gives \emph{stillan} as the verb and separately
{\emph{stille}}\index[iv]{02oe@\ivlangheader{Old English}{West Saxon normalization unmarked}!stille@\iventry{stille}{}}
`still, quiet' as the adjective (\citeproc{ref-ClarkHall1960}{Clark Hall
1960}). Bosworth-Toller likewise preserves a substantial prefixed verbal
family under \emph{ge-stillan} `still, stop' and related forms
(\citeproc{ref-BosworthToller1898}{Bosworth and Toller 1898, 724}). The
Old English form here is the verb \emph{stillan}, not the adjective.

\subsection{Development to Old
English}\label{development-to-old-english-55}

As a heavy-stem Class I weak verb, \Recon{stéllijaną} `still' undergoes
the expected syncope and i-umlaut, and later loss of \emph{j} after a
heavy stem yields \emph{stillan}. The development represented here is
regular.

\subsection{Form note}\label{form-note-19}

The note concerns lexical framing rather than sound law: \emph{stillan}
is the verb represented here, while
{\emph{stille}}\index[iv]{02oe@\ivlangheader{Old English}{West Saxon normalization unmarked}!stille@\iventry{stille}{}}
`still, quiet' belongs to the related adjectival branch of the family
(\citeproc{ref-ClarkHall1960}{Clark Hall 1960};
\citeproc{ref-KlugeSeebold2011}{Kluge 2011}).

\section{\texorpdfstring{summer --- OE
\emph{sumer}}{summer --- OE sumer}}\label{summer-oe-sumer}

\index[iv]{02oe@\ivlangheader{Old English}{West Saxon normalization unmarked}!sumer@\iventry{sumer}{}}
\index[iv]{03pgmc@\ivlangheader{Proto-Germanic}{}!sumaraz@\iventry{*súmaraz}{}}

Derivation: \emph{*súmaraz} \textgreater{} \emph{sumer} (regular).

\subsection{Derivation trace}\label{derivation-trace-70}

Proto input: \emph{*súmaraz}

\begingroup
\setlength{\fboxsep}{6pt}

\noindent\fbox{%
\begin{minipage}{0.97\linewidth}
\small
\begin{tabular}{@{}>{\raggedright\arraybackslash}p{0.359\linewidth}@{\hspace{0.030\linewidth}}>{\raggedright\arraybackslash}p{0.611\linewidth}@{}}
\begin{minipage}[t]{\linewidth}
\centering\textbf{Earlier Germanic changes}\par
\vspace{0.35em}
\raggedright
\centering\textbf{Northwest and West Germanic}\par
\raggedright
\vspace{0.2em}
\begin{tabularx}{\linewidth}{@{}>{\raggedright\arraybackslash}X@{\hspace{0.55em}}>{\raggedright\arraybackslash}p{0.24\linewidth}@{\hspace{0.25em}}}
\mbox{EAF Final Z Deletion} & \emph{*súmara} \\
\end{tabularx}
\vspace{0.6em}
\centering\textbf{Early Anglo-Frisian}\par
\raggedright
\vspace{0.2em}
\raggedright [no change]\par
\end{minipage}
&
\begin{minipage}[t]{\linewidth}
\centering\textbf{Old English changes}\par
\vspace{0.35em}
\raggedright
\begin{tabularx}{\linewidth}{@{}>{\raggedright\arraybackslash}X@{\hspace{0.55em}}>{\raggedright\arraybackslash}p{0.14\linewidth}@{\hspace{0.25em}}}
\mbox{PWGmc Final Bare A Loss} & \emph{*súmar} \\
\mbox{EAF Brightening Unstressed} & \emph{*súmær} \\
\mbox{OE Unstressed AE Merger} & \emph{*súmer} \\
\end{tabularx}
\end{minipage}
\\
\end{tabular}
\end{minipage}%
} \endgroup

Old English form: \emph{sumer}

\subsection{Reconstruction and comparative
evidence}\label{reconstruction-and-comparative-evidence-56}

Kroonen gives the lexeme as \emph{*sumara-}, and Ringe and Taylor
likewise use \Recon{sumaraz} `summer', while Orel preserves an alternate
\Recon{sumeraz} `summer' (\citeproc{ref-Kroonen2013}{Kroonen 2013};
\citeproc{ref-RingeTaylor2014}{Ringe and Taylor 2014};
\citeproc{ref-Orel2003}{Orel 2003, 425}). The derivational input
\Recon{súmaraz} `summer' follows the \emph{*a} vocalism that underlies
the regular development represented here.

\subsection{Old English evidence}\label{old-english-evidence-57}

Clark Hall gives
{\emph{sumor}}\index[iv]{02oe@\ivlangheader{Old English}{West Saxon normalization unmarked}!sumor@\iventry{sumor}{}}
`summer' m., gs.
{\emph{sumeres}}\index[iv]{02oe@\ivlangheader{Old English}{West Saxon normalization unmarked}!sumeres@\iventry{sumeres}{}}
`summer', ds.
{\emph{sumera}}\index[iv]{02oe@\ivlangheader{Old English}{West Saxon normalization unmarked}!sumera@\iventry{sumera}{}}
`summer' /
{\emph{sumere}}\index[iv]{02oe@\ivlangheader{Old English}{West Saxon normalization unmarked}!sumere@\iventry{sumere}{}}
`summer', and Bright likewise lists \emph{sumor (sumer)} with genitive
{\emph{sumeres}}\index[iv]{02oe@\ivlangheader{Old English}{West Saxon normalization unmarked}!sumeres@\iventry{sumeres}{}}
`summer (gen.)' (\citeproc{ref-ClarkHall1960}{Clark Hall 1960};
\citeproc{ref-BrightCassidyRingler1971}{Bright 1971, 440}). The
tradition therefore preserves both
{\emph{sumor}}\index[iv]{02oe@\ivlangheader{Old English}{West Saxon normalization unmarked}!sumor@\iventry{sumor}{}}
`summer' and
{\emph{sumer}}\index[iv]{02oe@\ivlangheader{Old English}{West Saxon normalization unmarked}!sumer@\iventry{sumer}{}}
`summer', with the oblique forms strongly supporting second-syllable
\emph{e}.

\subsection{Development to Old
English}\label{development-to-old-english-56}

From \Recon{súmaraz} `summer', loss of final \emph{-az} is followed by
fronting and merger in the unstressed second syllable, yielding
{\emph{sumer}}\index[iv]{02oe@\ivlangheader{Old English}{West Saxon normalization unmarked}!sumer@\iventry{sumer}{}}
`summer'. The form compared here is the regularized \emph{e}-form, while
the common citation form
{\emph{sumor}}\index[iv]{02oe@\ivlangheader{Old English}{West Saxon normalization unmarked}!sumor@\iventry{sumor}{}}
`summer' remains part of the attested OE tradition
(\citeproc{ref-RingeTaylor2014}{Ringe and Taylor 2014}).

\subsection{Form note}\label{form-note-20}

The entry does not deny
{\emph{sumor}}\index[iv]{02oe@\ivlangheader{Old English}{West Saxon normalization unmarked}!sumor@\iventry{sumor}{}}
`summer'. It represents
{\emph{sumer}}\index[iv]{02oe@\ivlangheader{Old English}{West Saxon normalization unmarked}!sumer@\iventry{sumer}{}}
`summer' as the regular outcome chosen here, while
{\emph{sumor}}\index[iv]{02oe@\ivlangheader{Old English}{West Saxon normalization unmarked}!sumor@\iventry{sumor}{}}
`summer' remains a common headword spelling and
{\emph{sumeres}}\index[iv]{02oe@\ivlangheader{Old English}{West Saxon normalization unmarked}!sumeres@\iventry{sumeres}{}}
`summer (gen.)' /
{\emph{sumere}}\index[iv]{02oe@\ivlangheader{Old English}{West Saxon normalization unmarked}!sumere@\iventry{sumere}{}}
`summer (dat.)' show that the \emph{e}-vocalism was also real in Old
English (\citeproc{ref-ClarkHall1960}{Clark Hall 1960};
\citeproc{ref-BrightCassidyRingler1971}{Bright 1971, 440}).

\section{\texorpdfstring{sunder --- OE
\emph{sundrian}}{sunder --- OE sundrian}}\label{sunder-oe-sundrian}

\index[iv]{02oe@\ivlangheader{Old English}{West Saxon normalization unmarked}!sundrian@\iventry{sundrian}{}}
\index[iv]{03pgmc@\ivlangheader{Proto-Germanic}{}!sundrojana@\iventry{*súndrōjaną}{}}

Derivation: \emph{*súndrōjaną} \textgreater{} \emph{sundrian} (regular).

\subsection{Derivation trace}\label{derivation-trace-71}

Proto input: \emph{*súndrōjaną}

\begingroup
\setlength{\fboxsep}{6pt}

\noindent\fbox{%
\begin{minipage}{0.97\linewidth}
\footnotesize
\begin{tabular}{@{}>{\raggedright\arraybackslash}p{0.339\linewidth}@{\hspace{0.030\linewidth}}>{\raggedright\arraybackslash}p{0.630\linewidth}@{}}
\begin{minipage}[t]{\linewidth}
\centering\textbf{Earlier Germanic changes}\par
\vspace{0.35em}
\raggedright
\centering\textbf{Northwest and West Germanic}\par
\raggedright
\vspace{0.2em}
\raggedright [no change]\par
\vspace{0.6em}
\centering\textbf{Early Anglo-Frisian}\par
\raggedright
\vspace{0.2em}
\raggedright [no change]\par
\end{minipage}
&
\begin{minipage}[t]{\linewidth}
\centering\textbf{Old English changes}\par
\vspace{0.35em}
\raggedright
\begin{tabularx}{\linewidth}{@{}>{\raggedright\arraybackslash}X@{\hspace{0.55em}}>{\raggedright\arraybackslash}p{0.24\linewidth}@{\hspace{0.25em}}}
OE Heavy Syllable Nasal Apocope & \emph{*súndrōjan} \\
\mbox{OE Secondary Nasalization} & \emph{*súndrōjąn} \\
\mbox{OE I Umlaut} & \emph{*súndrējąn} \\
\mbox{OE Med Long Short} & \emph{*súndrejąn} \\
\mbox{OE Weak Tail Reduction} & \emph{*súndrejan} \\
OE Intervocalic J Vocalization & \emph{*súndreian} \\
\mbox{OE Unstressed EI Contraction} & \emph{*súndrian} \\
\end{tabularx}
\end{minipage}
\\
\end{tabular}
\end{minipage}%
} \endgroup

Old English form: \emph{sundrian}

\subsection{Reconstruction and comparative
evidence}\label{reconstruction-and-comparative-evidence-57}

Orel distinguishes three related formations: adverbial \Recon{sunþraz}
`asunder' \textgreater{}
{\emph{sundor}}\index[iv]{02oe@\ivlangheader{Old English}{West Saxon normalization unmarked}!sundor@\iventry{sundor}{}}
`apart, asunder', Class I verbal \Recon{sunþrjanan} `sunder'
\textgreater{}
{\emph{syndrian}}\index[iv]{02oe@\ivlangheader{Old English}{West Saxon normalization unmarked}!syndrian@\iventry{syndrian}{}}
`separate, sunder', and Class II verbal \Recon{sunþrōjanan} `sunder'
\textgreater{} \emph{sundrian} `sunder' (\citeproc{ref-Orel2003}{Orel
2003}).

Kluge-Seebold aligns the cognate set with German
{\emph{sondern}}\index[iv]{17german@\ivlangheader{German}{}!sondern@\iventry{sondern}{}}
`separate' and OE
{\emph{gesundrian}}\index[iv]{02oe@\ivlangheader{Old English}{West Saxon normalization unmarked}!gesundrian@\iventry{gesundrian}{}}
`separate', so this entry belongs with the Class II verb, not the adverb
(\citeproc{ref-KlugeSeebold2011}{Kluge 2011}).

\subsection{Old English evidence}\label{old-english-evidence-58}

Clark Hall and Bosworth-Toller list \emph{sundrian} `sunder' and
\emph{syndrian} `sunder' separately from adverbial \emph{sundor}
`asunder'; both also record the prefixed verbal family \emph{ā-sundrian}
`sunder apart' (\citeproc{ref-ClarkHall1960}{Clark Hall 1960, 296};
\citeproc{ref-BosworthToller1898}{Bosworth and Toller 1898}). The target
is therefore the weak verb \emph{sundrian} `sunder'.

\subsection{Development to Old
English}\label{development-to-old-english-57}

From \Recon{súndrōjaną} `sunder', the Class II weak-verb suffix yields
regular OE \emph{-ian}, producing \emph{sundrian} `sunder'. Because this
is the \emph{*-ōjan-} verb and not the Class I \emph{*-jan-} formation,
the form represented here does not belong to the umlauted
{\emph{syndrian}}\index[iv]{02oe@\ivlangheader{Old English}{West Saxon normalization unmarked}!syndrian@\iventry{syndrian}{}}
`separate' branch.

\subsection{Form note}\label{form-note-21}

The earlier confusion was lexical, not phonological:
{\emph{sundor}}\index[iv]{02oe@\ivlangheader{Old English}{West Saxon normalization unmarked}!sundor@\iventry{sundor}{}}
`apart' is the separate adverb, and
{\emph{syndrian}}\index[iv]{02oe@\ivlangheader{Old English}{West Saxon normalization unmarked}!syndrian@\iventry{syndrian}{}}
`separate' is a related but different verb. The verb treated here is the
Class II verb \emph{sundrian} (\citeproc{ref-Orel2003}{Orel 2003};
\citeproc{ref-ClarkHall1960}{Clark Hall 1960, 296}).

\section{\texorpdfstring{swallow --- OE
\emph{swealwe}}{swallow --- OE swealwe}}\label{swallow-oe-swealwe}

\index[iv]{02oe@\ivlangheader{Old English}{West Saxon normalization unmarked}!swealwe@\iventry{swealwe}{}}
\index[iv]{03pgmc@\ivlangheader{Proto-Germanic}{}!swalwon@\iventry{*swálwōn}{}}

Derivation: \emph{*swálwōn} \textgreater{} \emph{swealwe} (regular).

\subsection{Derivation trace}\label{derivation-trace-72}

Proto input: \emph{*swálwōn}

\begingroup
\setlength{\fboxsep}{6pt}

\noindent\fbox{%
\begin{minipage}{0.97\linewidth}
\small
\begin{tabular}{@{}>{\raggedright\arraybackslash}p{0.359\linewidth}@{\hspace{0.030\linewidth}}>{\raggedright\arraybackslash}p{0.611\linewidth}@{}}
\begin{minipage}[t]{\linewidth}
\centering\textbf{Earlier Germanic changes}\par
\vspace{0.35em}
\raggedright
\centering\textbf{Northwest and West Germanic}\par
\raggedright
\vspace{0.2em}
\begin{tabularx}{\linewidth}{@{}>{\raggedright\arraybackslash}X@{\hspace{0.55em}}>{\raggedright\arraybackslash}p{0.24\linewidth}@{\hspace{0.25em}}}
\mbox{PNWGmc N Stem N Loss} & \emph{*swálwǭ} \\
\end{tabularx}
\vspace{0.6em}
\centering\textbf{Early Anglo-Frisian}\par
\raggedright
\vspace{0.2em}
\raggedright [no change]\par
\end{minipage}
&
\begin{minipage}[t]{\linewidth}
\centering\textbf{Old English changes}\par
\vspace{0.35em}
\raggedright
\begin{tabularx}{\linewidth}{@{}>{\raggedright\arraybackslash}X@{\hspace{0.55em}}>{\raggedright\arraybackslash}p{0.18\linewidth}@{\hspace{0.25em}}}
\mbox{EAF Brightening} & \emph{*swælwǭ} \\
\mbox{OE Breaking} & \emph{*swealwǭ} \\
OE Unstressed Long Vowel Shortening & \emph{*swealwæ} \\
\mbox{OE Unstressed AE Merger} & \emph{*swealwe} \\
\end{tabularx}
\end{minipage}
\\
\end{tabular}
\end{minipage}%
} \endgroup

Old English form: \emph{swealwe}

\subsection{Reconstruction and comparative
evidence}\label{reconstruction-and-comparative-evidence-58}

Kroonen gives the bird name as \emph{*swalwōn-}, and Ringe and Taylor
cite the later West Germanic stage \Recon{swalwa} `swallow', from which
West Saxon
{\emph{swealwe}}\index[iv]{02oe@\ivlangheader{Old English}{West Saxon normalization unmarked}!swealwe@\iventry{swealwe}{}}
`swallow' and Mercian
{\emph{swalwe}}\index[iv]{02oe@\ivlangheader{Old English}{West Saxon normalization unmarked}!swalwe~05@\iventry{swalwe}{Merc.}}
`swallow' develop (\citeproc{ref-Kroonen2013}{Kroonen 2013, 535};
\citeproc{ref-RingeTaylor2014}{Ringe and Taylor 2014, 200}). The
selected etymological comparison belongs to the swallow-bird family, not
to the verb \emph{swelgan}.

\subsection{Old English evidence}\label{old-english-evidence-59}

Clark Hall records \emph{swealwe (a, o)} as the noun headword
(\citeproc{ref-ClarkHall1960}{Clark Hall 1960}). Campbell and Brunner
also preserve later or oblique-family forms such as \emph{swaluwe},
\emph{swalewan}, and \emph{swealuwe}, but those belong to wider
variation around the noun rather than to the citation form represented
here (\citeproc{ref-Campbell1959}{Campbell 1959};
\citeproc{ref-SieversBrunner1965}{Brunner 1965}).

\subsection{Development to Old
English}\label{development-to-old-english-58}

From \Recon{swálwōn} `swallow', brightening yields \emph{*swælw-}, and
breaking before \emph{lw} gives \emph{*swealw-}. The later noun ending
develops regularly to \emph{swealwe}. The relevant point is that the
bird name has no inherited \emph{*g}: that consonant belongs to the
separate verb \emph{swelgan} (\citeproc{ref-RingeTaylor2014}{Ringe and
Taylor 2014, 200}; \citeproc{ref-Kroonen2013}{Kroonen 2013, 535}).

\section{\texorpdfstring{swine --- OE
\emph{swīn}}{swine --- OE swīn}}\label{swine-oe-swux12bn}

\index[iv]{02oe@\ivlangheader{Old English}{West Saxon normalization unmarked}!swin@\iventry{swīn}{}}
\index[iv]{03pgmc@\ivlangheader{Proto-Germanic}{}!swina@\iventry{*swī́ną}{}}

Derivation: citation reconstruction \emph{*swī́ną}; form followed here
\emph{*swī́ną} \textgreater{} \emph{swīn} (regular).

\subsection{Derivation trace}\label{derivation-trace-73}

Proto input: \emph{*swī́ną}

\begingroup
\setlength{\fboxsep}{6pt}

\noindent\fbox{%
\begin{minipage}{0.97\linewidth}
\small
\begin{tabular}{@{}>{\raggedright\arraybackslash}p{0.359\linewidth}@{\hspace{0.030\linewidth}}>{\raggedright\arraybackslash}p{0.611\linewidth}@{}}
\begin{minipage}[t]{\linewidth}
\centering\textbf{Earlier Germanic changes}\par
\vspace{0.35em}
\raggedright
\centering\textbf{Northwest and West Germanic}\par
\raggedright
\vspace{0.2em}
\raggedright [no change]\par
\vspace{0.6em}
\centering\textbf{Early Anglo-Frisian}\par
\raggedright
\vspace{0.2em}
\raggedright [no change]\par
\end{minipage}
&
\begin{minipage}[t]{\linewidth}
\centering\textbf{Old English changes}\par
\vspace{0.35em}
\raggedright
\begin{tabularx}{\linewidth}{@{}>{\raggedright\arraybackslash}X@{\hspace{0.55em}}>{\raggedright\arraybackslash}p{0.14\linewidth}@{\hspace{0.25em}}}
OE Heavy Syllable Nasal Apocope & \emph{*swī́n} \\
\end{tabularx}
\end{minipage}
\\
\end{tabular}
\end{minipage}%
} \endgroup

Old English form: \emph{swīn}

\subsection{Reconstruction and comparative
evidence}\label{reconstruction-and-comparative-evidence-59}

Kroonen cites the noun as \emph{*swina-} `pig'
(\citeproc{ref-Kroonen2013}{Kroonen 2013}). The selected comparative
form here is the bare neuter citation cell \Recon{swī́ną} `swine', which
matches the singular noun aimed at in Old English rather than an oblique
stem form.

\subsection{Old English evidence}\label{old-english-evidence-60}

Clark Hall records \emph{swin (y)} as the ordinary noun headword
(\citeproc{ref-ClarkHall1960}{Clark Hall 1960}). The target here is that
singular citation form \emph{swīn} `swine', not a plural glossed in
Modern English as \emph{swine}.

\subsection{Development to Old
English}\label{development-to-old-english-59}

From derivational input \Recon{swī́ną} `swine', loss of the final nasal
vowel yields
{\emph{swīn}}\index[iv]{02oe@\ivlangheader{Old English}{West Saxon normalization unmarked}!swin@\iventry{swīn}{}}
`swine, pig'. The outcome is therefore the regular monosyllabic noun
with preserved long root \emph{ī}.

\subsection{Source note}\label{source-note-4}

The derivational input writes stressed long \emph{ī} as \emph{*ī́}, so
comparative \Recon{swī́ną} `swine' and derivational \Recon{swī́ną} `swine'
represent the same lexical form.

\section{\texorpdfstring{there --- OE
\emph{þǣr}}{there --- OE þǣr}}\label{there-oe-uxfeux1e3r}

\index[iv]{02oe@\ivlangheader{Old English}{West Saxon normalization unmarked}!thaer@\iventry{þǣr}{}}
\index[iv]{03pgmc@\ivlangheader{Proto-Germanic}{}!thar@\iventry{*θár}{}}

Derivation: \emph{*θár} \textgreater{} \emph{þǣr} (regular).

\subsection{Derivation trace}\label{derivation-trace-74}

Proto input: \emph{*θár}

\begingroup
\setlength{\fboxsep}{6pt}

\noindent\fbox{%
\begin{minipage}{0.97\linewidth}
\small
\begin{tabular}{@{}>{\raggedright\arraybackslash}p{0.611\linewidth}@{\hspace{0.030\linewidth}}>{\raggedright\arraybackslash}p{0.359\linewidth}@{}}
\begin{minipage}[t]{\linewidth}
\centering\textbf{Earlier Germanic changes}\par
\vspace{0.35em}
\raggedright
\centering\textbf{Northwest and West Germanic}\par
\raggedright
\vspace{0.2em}
\begin{tabularx}{\linewidth}{@{}>{\raggedright\arraybackslash}X@{\hspace{0.55em}}>{\raggedright\arraybackslash}p{0.14\linewidth}@{\hspace{0.25em}}}
PNWGmc Monosyllabic Final R Lengthening & \emph{*θār} \\
\mbox{EAF Long A Fronting} & \emph{*θǣr} \\
\end{tabularx}
\vspace{0.6em}
\centering\textbf{Early Anglo-Frisian}\par
\raggedright
\vspace{0.2em}
\raggedright [no change]\par
\end{minipage}
&
\begin{minipage}[t]{\linewidth}
\centering\textbf{Old English changes}\par
\vspace{0.35em}
\raggedright
\raggedright [no change]\par
\end{minipage}
\\
\end{tabular}
\end{minipage}%
} \endgroup

Old English form: \emph{þǣr}

\subsection{Old English evidence}\label{old-english-evidence-61}

Old English
{\emph{þǣr}}\index[iv]{02oe@\ivlangheader{Old English}{West Saxon normalization unmarked}!thaer@\iventry{þǣr}{}}
`there' is the normalized West Saxon adverb. Clark Hall lemmatizes
{\emph{ðǣr}}\index[iv]{02oe@\ivlangheader{Old English}{West Saxon normalization unmarked}!thaer@\iventry{ðǣr}{}}
`there'; the substitution of thorn for eth does not change the vowel or
the identity of the word (\citeproc{ref-ClarkHall1960}{Clark Hall 1960,
305}). Hill contrasts West Saxon \emph{ðǣr} with Anglian
{\emph{ðēr}}\index[iv]{02oe@\ivlangheader{Old English}{West Saxon normalization unmarked}!ther@\iventry{ðēr}{}}
`there': both have long vowels, with different dialectal qualities
(\citeproc{ref-Hill2017}{Hill 2017, 142--43}). These are normalized
lexical and dialect forms, rather than diplomatic transcriptions of
individual manuscripts.

\subsection{Reconstruction and comparative
evidence}\label{reconstruction-and-comparative-evidence-60}

Kroonen reconstructs
\Recon{þar}\index[iv]{03pgmc@\ivlangheader{Proto-Germanic}{}!thar@\iventry{*þar}{}}
`there', formed from the demonstrative stem \emph{*þa-} with locative
\emph{*-r} (\citeproc{ref-Kroonen2013}{Kroonen 2013, 534}). Orel
likewise identifies the demonstrative formation, but prints the
alternatives \emph{*þar} and \emph{*þēr} (\citeproc{ref-Orel2003}{Orel
2003, 416}). His long-vowel citation must be kept distinct from the
earlier short-vowel antecedent. The selected
\Recon{θár}\index[iv]{03pgmc@\ivlangheader{Proto-Germanic}{}!thar@\iventry{*θár}{}}
`there' represents the short antecedent with an explicit root accent and
a phonological dental symbol; neither addition is a literal quotation of
the dictionaries.

The comparative distribution leaves the date of lengthening open. Gothic
\emph{þar} `there' has a spelling that does not distinguish short
\emph{a} from long \emph{ā}. More informative for conditioning is the
contrast between Old High German \emph{dār} `there' and the short vowels
of \emph{dăra} `thither' and \emph{dărōt} `thither'. Hill uses this
contrast to support lengthening in monosyllables ending in inherited
\emph{r} (\citeproc{ref-Hill2017}{Hill 2017, 142--48}).

\subsection{Development to Old
English}\label{development-to-old-english-60}

The adopted derivation applies
\hyperref[rule-PNWGmcMonosyllabicFinalRLengthening]{SC112, stressed
monosyllabic lengthening before inherited final \emph{r}}, giving
\emph{*þar} \textgreater{} \emph{*þār}. Subsequent
\hyperref[rule-EAFLongAFronting]{SC101, long-\emph{ā} fronting} supplies
\emph{*þǣr}, yielding West Saxon \emph{þǣr} without a change of
formation or an analogical replacement. The vowel becomes long while it
is still back; this is distinct from lengthening an already fronted
short vowel.

Hill argues for a Proto-Germanic date, whereas Ringe and Taylor place
the lengthening in West Germanic and discuss deictic stress and
reanalysis as possible explanations (\citeproc{ref-Hill2017}{Hill 2017,
146--48}; \citeproc{ref-RingeTaylor2014}{Ringe and Taylor 2014, 65}).
The regular short-\emph{a}-to-\emph{ǣ} development is conditional on the
chosen antecedent and early lengthening. Its successful outcome does not
establish that absolute date, and a word without inherited \emph{*i}
cannot independently order SC015 relative to lengthening.

\section{\texorpdfstring{think --- OE
\emph{þenċan}}{think --- OE þenċan}}\label{think-oe-uxfeenux10ban}

\index[iv]{02oe@\ivlangheader{Old English}{West Saxon normalization unmarked}!hycgan@\iventry{hycgan}{}}
\index[iv]{02oe@\ivlangheader{Old English}{West Saxon normalization unmarked}!thencan@\iventry{þenċan}{}}
\index[iv]{03pgmc@\ivlangheader{Proto-Germanic}{}!thankijana@\iventry{*θánkijaną}{}}

Derivation: \emph{*θánkijaną} \textgreater{} \emph{þenċan} (regular).

\subsection{Derivation trace}\label{derivation-trace-75}

Proto input: \emph{*θánkijaną}

\begingroup
\setlength{\fboxsep}{6pt}

\noindent\fbox{%
\begin{minipage}{0.97\linewidth}
\footnotesize
\begin{tabular}{@{}>{\raggedright\arraybackslash}p{0.339\linewidth}@{\hspace{0.030\linewidth}}>{\raggedright\arraybackslash}p{0.630\linewidth}@{}}
\begin{minipage}[t]{\linewidth}
\centering\textbf{Earlier Germanic changes}\par
\vspace{0.35em}
\raggedright
\centering\textbf{Northwest and West Germanic}\par
\raggedright
\vspace{0.2em}
\raggedright [no change]\par
\vspace{0.6em}
\centering\textbf{Early Anglo-Frisian}\par
\raggedright
\vspace{0.2em}
\raggedright [no change]\par
\end{minipage}
&
\begin{minipage}[t]{\linewidth}
\centering\textbf{Old English changes}\par
\vspace{0.35em}
\raggedright
\begin{tabularx}{\linewidth}{@{}>{\raggedright\arraybackslash}X@{\hspace{0.55em}}>{\raggedright\arraybackslash}p{0.22\linewidth}@{\hspace{0.25em}}}
OE Heavy Syllable Nasal Apocope & \emph{*θánkijan} \\
\mbox{OE Secondary Nasalization} & \emph{*θánkijąn} \\
\mbox{Sievers Law Syncope} & \emph{*θánkjąn} \\
\mbox{OE Velar Palatalization} & \emph{*θánʧjąn} \\
\mbox{OE I Umlaut} & \emph{*θenʧjąn} \\
\mbox{OE Weak Tail Reduction} & \emph{*θenʧjan} \\
\mbox{OE J Loss After Heavy} & \emph{*θenʧan} \\
\end{tabularx}
\end{minipage}
\\
\end{tabular}
\end{minipage}%
} \endgroup

Old English form: \emph{þenċan}

\subsection{Reconstruction and comparative
evidence}\label{reconstruction-and-comparative-evidence-61}

Kroonen gives the verb as \emph{*þankjan-} `to think', and Ringe and
Taylor cite fully inflected \Recon{þankijaną} `think' beside OE
\emph{þenċan} `think' (\citeproc{ref-Kroonen2013}{Kroonen 2013};
\citeproc{ref-RingeTaylor2014}{Ringe and Taylor 2014}). The noun
\Recon{þankaz} `think' belongs only to the wider derivational
background.

\subsection{Old English evidence}\label{old-english-evidence-62}

Bosworth-Toller preserves the verb under \emph{þencan} `think' /
\emph{geþencan} `think', and the citation form here is the ordinary
infinitive \emph{þenċan} `think'
(\citeproc{ref-BosworthToller1898}{Bosworth and Toller 1898}).

\subsection{Development to Old
English}\label{development-to-old-english-61}

From \Recon{θánkijaną} `think', palatalization before \emph{*j} and
i-umlaut produce \emph{þenċan} `think'. The infinitive is therefore a
straightforward weak-verb outcome.

\subsection{Lexical note}\label{lexical-note}

Campbell's assibilation discussion uses the same verb \emph{þencan}
`think'; the class-III relic \emph{hycgan} is a different lexeme
(\citeproc{ref-Campbell1959}{Campbell 1959};
\citeproc{ref-HoggGrammar1992}{Hogg 1992a}).

\section{\texorpdfstring{thorn --- OE
\emph{þorn}}{thorn --- OE þorn}}\label{thorn-oe-uxfeorn}

\index[iv]{02oe@\ivlangheader{Old English}{West Saxon normalization unmarked}!thorn@\iventry{þorn}{}}
\index[iv]{03pgmc@\ivlangheader{Proto-Germanic}{}!thurnaz@\iventry{*θúrnaz}{}}

Derivation: \emph{*θúrnaz} \textgreater{} \emph{þorn} (regular).

\subsection{Derivation trace}\label{derivation-trace-76}

Proto input: \emph{*θúrnaz}

\begingroup
\setlength{\fboxsep}{6pt}

\noindent\fbox{%
\begin{minipage}{0.97\linewidth}
\small
\begin{tabular}{@{}>{\raggedright\arraybackslash}p{0.485\linewidth}@{\hspace{0.030\linewidth}}>{\raggedright\arraybackslash}p{0.485\linewidth}@{}}
\begin{minipage}[t]{\linewidth}
\centering\textbf{Earlier Germanic changes}\par
\vspace{0.35em}
\raggedright
\centering\textbf{Northwest and West Germanic}\par
\raggedright
\vspace{0.2em}
\begin{tabularx}{\linewidth}{@{}>{\raggedright\arraybackslash}X@{\hspace{0.55em}}>{\raggedright\arraybackslash}p{0.14\linewidth}@{\hspace{0.25em}}}
\mbox{PNWGmc U Lowering} & \emph{*θórnaz} \\
\mbox{EAF Final Z Deletion} & \emph{*θórna} \\
\end{tabularx}
\vspace{0.6em}
\centering\textbf{Early Anglo-Frisian}\par
\raggedright
\vspace{0.2em}
\raggedright [no change]\par
\end{minipage}
&
\begin{minipage}[t]{\linewidth}
\centering\textbf{Old English changes}\par
\vspace{0.35em}
\raggedright
\begin{tabularx}{\linewidth}{@{}>{\raggedright\arraybackslash}X@{\hspace{0.55em}}>{\raggedright\arraybackslash}p{0.14\linewidth}@{\hspace{0.25em}}}
\mbox{PWGmc Final Bare A Loss} & \emph{*θórn} \\
\end{tabularx}
\end{minipage}
\\
\end{tabular}
\end{minipage}%
} \endgroup

Old English form: \emph{þorn}

\subsection{Reconstruction and comparative
evidence}\label{reconstruction-and-comparative-evidence-62}

Kroonen gives \emph{*þurna-} `thorn, briar', while Orel preserves the
masculine pair \Recon{þurnuz} `thorn' and \Recon{þurnaz} `thorn'
(\citeproc{ref-Kroonen2013}{Kroonen 2013}; \citeproc{ref-Orel2003}{Orel
2003}). The derivational input \Recon{θúrnaz} `thorn' belongs to that
same comparative family.

\subsection{Old English evidence}\label{old-english-evidence-63}

Bright lists \emph{þorn} `thorn' (m.) and Clark Hall likewise treats
\emph{þorn} `thorn' as the ordinary noun headword
(\citeproc{ref-BrightCassidyRingler1971}{Bright 1971};
\citeproc{ref-ClarkHall1960}{Clark Hall 1960}).

\subsection{Development to Old
English}\label{development-to-old-english-62}

The inherited stem shows regular lowering of \emph{u} to \emph{o} before
\emph{r}, and final loss yields \emph{þorn} `thorn'. The noun is
therefore a regular Old English continuation of the Proto-Germanic
thorn-family.

\subsection{Source note}\label{source-note-5}

The comparative sources preserve more than one stem formation, but the
Old English target itself is simply the citation form \emph{þorn}
`thorn'.

\section{\texorpdfstring{thought --- OE
\emph{þōhte}}{thought --- OE þōhte}}\label{thought-oe-uxfeux14dhte}

\index[iv]{02oe@\ivlangheader{Old English}{West Saxon normalization unmarked}!thohte@\iventry{þōhte}{}}
\index[iv]{03pgmc@\ivlangheader{Proto-Germanic}{}!thanxte@\iventry{*θánxtē}{}}

Derivation: \emph{*θánxtē} \textgreater{} \emph{þōhte} (regular).

\subsection{Derivation trace}\label{derivation-trace-77}

Proto input: \emph{*θánxtē}

\begingroup
\setlength{\fboxsep}{6pt}

\noindent\fbox{%
\begin{minipage}{0.97\linewidth}
\small
\begin{tabular}{@{}>{\raggedright\arraybackslash}p{0.485\linewidth}@{\hspace{0.030\linewidth}}>{\raggedright\arraybackslash}p{0.485\linewidth}@{}}
\begin{minipage}[t]{\linewidth}
\centering\textbf{Earlier Germanic changes}\par
\vspace{0.35em}
\raggedright
\centering\textbf{Northwest and West Germanic}\par
\raggedright
\vspace{0.2em}
\begin{tabularx}{\linewidth}{@{}>{\raggedright\arraybackslash}X@{\hspace{0.55em}}>{\raggedright\arraybackslash}p{0.14\linewidth}@{\hspace{0.25em}}}
\mbox{EAF Nasalized Low Rounding} & \emph{*θōxtē} \\
\end{tabularx}
\vspace{0.6em}
\centering\textbf{Early Anglo-Frisian}\par
\raggedright
\vspace{0.2em}
\raggedright [no change]\par
\end{minipage}
&
\begin{minipage}[t]{\linewidth}
\centering\textbf{Old English changes}\par
\vspace{0.35em}
\raggedright
\begin{tabularx}{\linewidth}{@{}>{\raggedright\arraybackslash}X@{\hspace{0.55em}}>{\raggedright\arraybackslash}p{0.14\linewidth}@{\hspace{0.25em}}}
OE Unstressed Long Vowel Shortening & \emph{*θōxte} \\
\end{tabularx}
\end{minipage}
\\
\end{tabular}
\end{minipage}%
} \endgroup

Old English form: \emph{þōhte}

\subsection{Reconstruction and comparative
evidence}\label{reconstruction-and-comparative-evidence-63}

The derivational input is the class I weak preterite third singular of
the verb `think', PGmc
\Recon{θánxtē}\index[iv]{03pgmc@\ivlangheader{Proto-Germanic}{}!thanhte@\iventry{*θánxtē}{}}
`thought' (traditionally *þanhtē). Ringe gives the principal parts in
full as *þankijaną, *þanhtē, *þanhtaz, with the comparative set Gothic
{þagkjan}\index[iv]{12goth@\ivlangheader{Gothic}{}!thagkjan@\iventry{þagkjan}{}}
`think',
{þāhta}\index[iv]{12goth@\ivlangheader{Gothic}{}!thahta@\iventry{þāhta}{}}
`thought', ON
{þekkja}\index[iv]{09on@\ivlangheader{Old Norse}{}!thekkja@\iventry{þekkja}{}}
`think',
{þátti}\index[iv]{09on@\ivlangheader{Old Norse}{}!thatti@\iventry{þátti}{}}
`thought', OE \emph{þenċan}, \emph{þōhte}, \emph{þōht}, OHG
{denken}\index[iv]{10ohg@\ivlangheader{Old High German}{}!denken@\iventry{denken}{}}
`think',
{dāhta}\index[iv]{10ohg@\ivlangheader{Old High German}{}!dahta@\iventry{dāhta}{}}
`thought' (\citeproc{ref-Ringe2017}{Ringe 2017, 281}). The preterite
stem is reconstructed a second time, independently of the paradigm list,
among the past participles that show pre-*t devoicing: post-PIE *tong-
gives PGmc *þank- in *þankijaną beside *þanhtaz `thought'
(\citeproc{ref-Ringe2017}{Ringe 2017, 136}). In both passages the
Proto-Germanic form carries the nasal before the dorsal fricative. The
nasal is therefore present in the derivational input and its loss is
derived by the cascade instead of being built into the reconstruction.

The corpus already holds the present stem of the same verb as
\emph{think}, PGmc
\Recon{θánkijaną}\index[iv]{03pgmc@\ivlangheader{Proto-Germanic}{}!thankijana@\iventry{*θánkijaną}{}}
`think'. The preterite is entered as a separate row because the two
stems have entirely different vowel histories, and only the preterite
passes through the nasal-loss complex.

\subsection{The connectorless
formation}\label{the-connectorless-formation}

Ringe's full principal parts establish an inherited finite word, not a
CAPR assembly from a present stem and a general ending. His set contains
exactly five reconstructable Proto-Germanic class-I verbs with
connectorless pasts; Ringe and Taylor distinguish that inherited set
from later extension (\citeproc{ref-Ringe2017}{Ringe 2017, 281};
\citeproc{ref-RingeTaylor2014}{Ringe and Taylor 2014, 71--75, 97--99}).
The selected third singular is the verb `thought', not the noun.

Its deeper formation is a different question. Fulk discusses
\emph{*þaŋxta} and the cluster \emph{*þaŋk-t-}, reviewing retention of a
present-only suffix, loss of an earlier connector, and analogical
remodelling without finding a consensus (\citeproc{ref-Fulk2018}{Fulk
2018, 299--301}). Those morphological illustrations are not additional
complete selected words. In particular, his counterfactual buy form
\emph{*bu3-dē} predicts the wrong Gothic result; it is neither an
endorsed input nor evidence for thought's finite cell
(\citeproc{ref-Fulk2018}{Fulk 2018, 300}). Successful nasal-loss and
rounding phonology does not explain the connector's absence. The
inherited reconstruction is secure within this account while that
formation-level premise remains qualified.

\subsection{Old English evidence}\label{old-english-evidence-64}

West Saxon
{\emph{þōhte}}\index[iv]{02oe@\ivlangheader{Old English}{West Saxon normalization unmarked}!thohte@\iventry{þōhte}{}}
`thought' is the standard preterite singular of \emph{þenċan}
(\citeproc{ref-ClarkHall1960}{Clark Hall 1960};
\citeproc{ref-BosworthToller1898}{Bosworth and Toller 1898}). Fulk names
this exact form as the Anglo-Frisian witness for the rounding of the
lengthened nasalized low vowel: the lengthened vowels ``may have
remained nasalized for a considerable time, well past the close of the
NWGmc. period, since *ą̄ produced this way developed to ō in
Anglo-Frisian (as in OE pret. sg. þōhte, OFris. thochte `thought') and
did not fall together with OE ā \textless{} ai or OFris. ā \textless{}
ai, au'' (\citeproc{ref-Fulk2018}{Fulk 2018, § 4.1}, p.~55). The Old
Frisian counterpart
{thochte}\index[iv]{11ofris@\ivlangheader{Old Frisian}{}!thochte@\iventry{thochte}{}}
`thought' confirms that the development is Anglo-Frisian and not an
English innovation.

\subsection{Development to Old
English}\label{development-to-old-english-63}

The diagnostic root-vowel history consists of nasal loss and rounding.
Loss of the nasal before the dorsal fricative nasalizes and lengthens
the preceding vowel, so *θánxtē becomes *θãxtē; the Anglo-Frisian
rounding of the nasalized low vowel then gives *θōxtē, which surfaces as
\emph{þōhte} with the fricative written \emph{h}
(\citeproc{ref-Ringe2017}{Ringe 2017, 281}; \citeproc{ref-Fulk2018}{Fulk
2018, § 4.1}, p.~55). The trace also records ordinary final-vowel
shortening; that endpoint adjustment is separate from the diagnostic
root-vowel history.

Three non-firings matter as much as the firings. The Ingvaeonic
nasal-spirant law does not apply, correctly, because the dorsal
fricative is not one of the *f, *þ, *s spirants that law governs.
Preconsonantal loss of the dorsal fricative does not apply either, so
the fricative survives to the surface; this is what makes the preterite
a cleaner diagnostic than \emph{fist} (\citeproc{ref-Ringe2017}{Ringe
2017, 281}). Finally the rounding does not touch ordinary long low
vowels of other origin, as \emph{stone} and \emph{home} show.

\subsection{Why this row is in the
corpus}\label{why-this-row-is-in-the-corpus-1}

The nasal-loss rule has a high-vowel branch and a low-vowel branch.
Before this row only the high-vowel branch had a corpus witness, in
\emph{fist} *fúnxstiz, whose input vowel is *u and which never reaches
the rounding rule at all. The low-vowel pathway *aNx \textgreater{} *ą̄x
\textgreater{} *ōx rested on the handbooks alone. This row supplies an
attested instance of it, and with it the feeding relation between the
nasal loss and the rounding, which was thereby promoted from stage
entailment to an independently demonstrated ordering. See
\hyperref[rule-PGmcNasalLossBeforeX]{SC103 PGmcNasalLossBeforeX} and
\hyperref[rule-EAFNasalizedLowRounding]{SC104 EAFNasalizedLowRounding}.

\section{\texorpdfstring{tide --- OE
\emph{tīd}}{tide --- OE tīd}}\label{tide-oe-tux12bd}

\index[iv]{02oe@\ivlangheader{Old English}{West Saxon normalization unmarked}!tid@\iventry{tīd}{}}
\index[iv]{02oe@\ivlangheader{Old English}{West Saxon normalization unmarked}!tidan@\iventry{tīdan}{}}
\index[iv]{03pgmc@\ivlangheader{Proto-Germanic}{}!tidiz@\iventry{*tī́diz}{}}

Derivation: citation reconstruction \emph{*tī́diz}; form followed here
\emph{*tī́diz} \textgreater{} \emph{tīd} (regular).

\subsection{Derivation trace}\label{derivation-trace-78}

Proto input: \emph{*tī́diz}

\begingroup
\setlength{\fboxsep}{6pt}

\noindent\fbox{%
\begin{minipage}{0.97\linewidth}
\small
\begin{tabular}{@{}>{\raggedright\arraybackslash}p{0.485\linewidth}@{\hspace{0.030\linewidth}}>{\raggedright\arraybackslash}p{0.485\linewidth}@{}}
\begin{minipage}[t]{\linewidth}
\centering\textbf{Earlier Germanic changes}\par
\vspace{0.35em}
\raggedright
\centering\textbf{Northwest and West Germanic}\par
\raggedright
\vspace{0.2em}
\begin{tabularx}{\linewidth}{@{}>{\raggedright\arraybackslash}X@{\hspace{0.55em}}>{\raggedright\arraybackslash}p{0.14\linewidth}@{\hspace{0.25em}}}
\mbox{EAF Final Z Deletion} & \emph{*tī́di} \\
\end{tabularx}
\vspace{0.6em}
\centering\textbf{Early Anglo-Frisian}\par
\raggedright
\vspace{0.2em}
\raggedright [no change]\par
\end{minipage}
&
\begin{minipage}[t]{\linewidth}
\centering\textbf{Old English changes}\par
\vspace{0.35em}
\raggedright
\begin{tabularx}{\linewidth}{@{}>{\raggedright\arraybackslash}X@{\hspace{0.55em}}>{\raggedright\arraybackslash}p{0.14\linewidth}@{\hspace{0.25em}}}
\mbox{OE High Vowel Apocope} & \emph{*tī́d} \\
\end{tabularx}
\end{minipage}
\\
\end{tabular}
\end{minipage}%
} \endgroup

Old English form: \emph{tīd}

\subsection{Reconstruction and comparative
evidence}\label{reconstruction-and-comparative-evidence-64}

Kroonen gives \emph{*tīdi-} `time', and Orel's \Recon{tīđiz} `tide'
points to the same feminine noun (\citeproc{ref-Kroonen2013}{Kroonen
2013}; \citeproc{ref-Orel2003}{Orel 2003}). The related verb
\emph{tīdan} `happen' is separate.

\subsection{Old English evidence}\label{old-english-evidence-65}

Bright records \emph{tīd} `tide' with singular \emph{tīde} `tide' and
plural \emph{tīda} `tides', and Clark Hall treats \emph{tīd} `tide' as
the ordinary noun `time, period, season'
(\citeproc{ref-BrightCassidyRingler1971}{Bright 1971};
\citeproc{ref-ClarkHall1960}{Clark Hall 1960, 309}).

\subsection{Development to Old
English}\label{development-to-old-english-64}

From \Recon{tī́diz} `tide', final \emph{z} is lost and the high final
vowel drops, leaving \emph{tīd} `tide'. The development is
straightforward for a feminine i-stem.

\subsection{Lexical note}\label{lexical-note-1}

English \emph{tide} can also lead to the separate weak verb \emph{tīdan}
`happen'; the noun is \emph{tīd} `tide'.

\section{\texorpdfstring{token --- OE
\emph{tācn}}{token --- OE tācn}}\label{token-oe-tux101cn}

\index[iv]{02oe@\ivlangheader{Old English}{West Saxon normalization unmarked}!tacn@\iventry{tācn}{}}
\index[iv]{03pgmc@\ivlangheader{Proto-Germanic}{}!taikna@\iventry{*táikną}{}}

Derivation: \emph{*táikną} \textgreater{} \emph{tācn} (regular).

\subsection{Derivation trace}\label{derivation-trace-79}

Proto input: \emph{*táikną}

\begingroup
\setlength{\fboxsep}{6pt}

\noindent\fbox{%
\begin{minipage}{0.97\linewidth}
\small
\begin{tabular}{@{}>{\raggedright\arraybackslash}p{0.485\linewidth}@{\hspace{0.030\linewidth}}>{\raggedright\arraybackslash}p{0.485\linewidth}@{}}
\begin{minipage}[t]{\linewidth}
\centering\textbf{Earlier Germanic changes}\par
\vspace{0.35em}
\raggedright
\centering\textbf{Northwest and West Germanic}\par
\raggedright
\vspace{0.2em}
\raggedright [no change]\par
\vspace{0.6em}
\centering\textbf{Early Anglo-Frisian}\par
\raggedright
\vspace{0.2em}
\begin{tabularx}{\linewidth}{@{}>{\raggedright\arraybackslash}X@{\hspace{0.55em}}>{\raggedright\arraybackslash}p{0.14\linewidth}@{\hspace{0.25em}}}
\mbox{EAF Ai Monophthongization} & \emph{*tākną} \\
\end{tabularx}
\end{minipage}
&
\begin{minipage}[t]{\linewidth}
\centering\textbf{Old English changes}\par
\vspace{0.35em}
\raggedright
\begin{tabularx}{\linewidth}{@{}>{\raggedright\arraybackslash}X@{\hspace{0.55em}}>{\raggedright\arraybackslash}p{0.14\linewidth}@{\hspace{0.25em}}}
OE Heavy Syllable Nasal Apocope & \emph{*tākn} \\
\end{tabularx}
\end{minipage}
\\
\end{tabular}
\end{minipage}%
} \endgroup

Old English form: \emph{tācn}

\subsection{Reconstruction and comparative
evidence}\label{reconstruction-and-comparative-evidence-65}

Kroonen cites \emph{*taikna-} and Orel \Recon{taiknan} `token' for the
noun `sign, token' (\citeproc{ref-Kroonen2013}{Kroonen 2013};
\citeproc{ref-Orel2003}{Orel 2003, 438}). The derivational input
\Recon{táikną} `token' is the simple citation-form noun used for the
derivation.

\subsection{Old English evidence}\label{old-english-evidence-66}

Campbell and Sievers-Brunner preserve both unbroken \emph{tācn} `token'
and broken \emph{tācen} `token', with oblique \emph{tācnes} `token'
remaining unbroken (\citeproc{ref-Campbell1959}{Campbell 1959};
\citeproc{ref-SieversBrunner1965}{Brunner 1965}).

\subsection{Development to Old
English}\label{development-to-old-english-65}

Monophthongization of \emph{ai} yields \emph{ā}, and loss of the final
nasal vowel leaves \emph{tācn} `token'. The unbroken form is therefore a
regular Old English outcome.

\subsection{Form note}\label{form-note-22}

\emph{tācn} `token' is the attested unbroken citation form selected
here. Later West Saxon prose often prefers \emph{tācen} `token', but
that does not displace the older unbroken form
(\citeproc{ref-Campbell1959}{Campbell 1959};
\citeproc{ref-SieversBrunner1965}{Brunner 1965}).

\section{\texorpdfstring{town --- OE
\emph{tūn}}{town --- OE tūn}}\label{town-oe-tux16bn}

\index[iv]{02oe@\ivlangheader{Old English}{West Saxon normalization unmarked}!tun@\iventry{tūn}{}}
\index[iv]{03pgmc@\ivlangheader{Proto-Germanic}{}!tuna@\iventry{*tūną}{}}

Derivation: \emph{*tūną} \textgreater{} \emph{tūn} (regular).

\subsection{Derivation trace}\label{derivation-trace-80}

Proto input: \emph{*tūną}

\begingroup
\setlength{\fboxsep}{6pt}

\noindent\fbox{%
\begin{minipage}{0.97\linewidth}
\small
\begin{tabular}{@{}>{\raggedright\arraybackslash}p{0.359\linewidth}@{\hspace{0.030\linewidth}}>{\raggedright\arraybackslash}p{0.611\linewidth}@{}}
\begin{minipage}[t]{\linewidth}
\centering\textbf{Earlier Germanic changes}\par
\vspace{0.35em}
\raggedright
\centering\textbf{Northwest and West Germanic}\par
\raggedright
\vspace{0.2em}
\raggedright [no change]\par
\vspace{0.6em}
\centering\textbf{Early Anglo-Frisian}\par
\raggedright
\vspace{0.2em}
\raggedright [no change]\par
\end{minipage}
&
\begin{minipage}[t]{\linewidth}
\centering\textbf{Old English changes}\par
\vspace{0.35em}
\raggedright
\begin{tabularx}{\linewidth}{@{}>{\raggedright\arraybackslash}X@{\hspace{0.55em}}>{\raggedright\arraybackslash}p{0.14\linewidth}@{\hspace{0.25em}}}
OE Heavy Syllable Nasal Apocope & \emph{*tūn} \\
\end{tabularx}
\end{minipage}
\\
\end{tabular}
\end{minipage}%
} \endgroup

Old English form: \emph{tūn}

\subsection{Reconstruction and comparative
evidence}\label{reconstruction-and-comparative-evidence-66}

Kroonen cites \emph{*tūna-} `fenced area', while Orel gives
\Recon{tūnan} `town' and \Recon{tūnaz} `town'
(\citeproc{ref-Kroonen2013}{Kroonen 2013}; \citeproc{ref-Orel2003}{Orel
2003, 452}). The derivational input \Recon{tūną} `town' is the simple
citation-form noun used in the derivation.

\subsection{Old English evidence}\label{old-english-evidence-67}

Clark Hall records \emph{tūn} `town' as the ordinary headword
`enclosure, yard, village, town' (\citeproc{ref-ClarkHall1960}{Clark
Hall 1960}).

\subsection{Development to Old
English}\label{development-to-old-english-66}

The inherited long \emph{ū} is preserved, and loss of the final nasal
vowel yields \emph{tūn} `town' regularly
(\citeproc{ref-SieversBrunner1965}{Brunner 1965}).

\subsection{Source note}\label{source-note-6}

The comparative headwords vary, but the Old English target here is the
direct citation form \emph{tūn} `town', not an oblique \Recon{tūnăn}
`town'.

\section{\texorpdfstring{wade --- OE
\emph{wadan}}{wade --- OE wadan}}\label{wade-oe-wadan}

\index[iv]{02oe@\ivlangheader{Old English}{West Saxon normalization unmarked}!wadan@\iventry{wadan}{}}
\index[iv]{03pgmc@\ivlangheader{Proto-Germanic}{}!wadana@\iventry{*wádaną}{}}

Derivation: \emph{*wádaną} \textgreater{} \emph{wadan} (regular).

\subsection{Derivation trace}\label{derivation-trace-81}

Proto input: \emph{*wádaną}

\begingroup
\setlength{\fboxsep}{6pt}

\noindent\fbox{%
\begin{minipage}{0.97\linewidth}
\small
\begin{tabular}{@{}>{\raggedright\arraybackslash}p{0.359\linewidth}@{\hspace{0.030\linewidth}}>{\raggedright\arraybackslash}p{0.611\linewidth}@{}}
\begin{minipage}[t]{\linewidth}
\centering\textbf{Earlier Germanic changes}\par
\vspace{0.35em}
\raggedright
\centering\textbf{Northwest and West Germanic}\par
\raggedright
\vspace{0.2em}
\raggedright [no change]\par
\vspace{0.6em}
\centering\textbf{Early Anglo-Frisian}\par
\raggedright
\vspace{0.2em}
\raggedright [no change]\par
\end{minipage}
&
\begin{minipage}[t]{\linewidth}
\centering\textbf{Old English changes}\par
\vspace{0.35em}
\raggedright
\begin{tabularx}{\linewidth}{@{}>{\raggedright\arraybackslash}X@{\hspace{0.55em}}>{\raggedright\arraybackslash}p{0.14\linewidth}@{\hspace{0.25em}}}
\mbox{EAF Brightening} & \emph{*wædaną} \\
\mbox{OE A Restoration} & \emph{*wadaną} \\
OE Heavy Syllable Nasal Apocope & \emph{*wadan} \\
\mbox{OE Secondary Nasalization} & \emph{*wadąn} \\
\mbox{OE Weak Tail Reduction} & \emph{*wadan} \\
\end{tabularx}
\end{minipage}
\\
\end{tabular}
\end{minipage}%
} \endgroup

Old English form: \emph{wadan}

\subsection{Reconstruction and comparative
evidence}\label{reconstruction-and-comparative-evidence-67}

Campbell and Ringe and Taylor describe A-restoration before a following
back vowel, and Luick explicitly includes
{\emph{wadan}}\index[iv]{02oe@\ivlangheader{Old English}{West Saxon normalization unmarked}!wadan@\iventry{wadan}{}}
`wade, go through' among the standard open-syllable examples
(\citeproc{ref-Campbell1959}{Campbell 1959};
\citeproc{ref-RingeTaylor2014}{Ringe and Taylor 2014};
\citeproc{ref-Luick1914}{Luick 1914-\/-1940, 239}).

\subsection{Old English evidence}\label{old-english-evidence-68}

Clark Hall gives
{\emph{wadan}}\index[iv]{02oe@\ivlangheader{Old English}{West Saxon normalization unmarked}!wadan@\iventry{wadan}{}}
`wade, go through' as the verb `to go, move, stride, advance', and
Bright lists the same infinitive in the strong-verb paradigm
(\citeproc{ref-ClarkHall1960}{Clark Hall 1960};
\citeproc{ref-BrightCassidyRingler1971}{Bright 1971}).

\subsection{Development to Old
English}\label{development-to-old-english-67}

From \Recon{wádaną} `wade', Anglo-Frisian brightening first gives
\Recon{wædaną} `wade'. A-restoration before the back-vocalic infinitive
ending then returns \emph{a}, and later reduction yields
{\emph{wadan}}\index[iv]{02oe@\ivlangheader{Old English}{West Saxon normalization unmarked}!wadan@\iventry{wadan}{}}
`wade, go through' (\citeproc{ref-Campbell1959}{Campbell 1959};
\citeproc{ref-RingeTaylor2014}{Ringe and Taylor 2014}).

\subsection{Development note}\label{development-note-2}

The infinitive belongs to the A-restoration class. The citation form is
therefore \emph{wadan} `wade', not a fronted \emph{wæden} `waded'-type
output.

\section{\texorpdfstring{warp --- OE
\emph{weorpan}}{warp --- OE weorpan}}\label{warp-oe-weorpan}

\index[iv]{02oe@\ivlangheader{Old English}{West Saxon normalization unmarked}!weorpan@\iventry{weorpan}{}}
\index[iv]{03pgmc@\ivlangheader{Proto-Germanic}{}!werpana@\iventry{*wérpaną}{}}

Derivation: \emph{*wérpaną} \textgreater{} \emph{weorpan} (regular).

\subsection{Derivation trace}\label{derivation-trace-82}

Proto input: \emph{*wérpaną}

\begingroup
\setlength{\fboxsep}{6pt}

\noindent\fbox{%
\begin{minipage}{0.97\linewidth}
\small
\begin{tabular}{@{}>{\raggedright\arraybackslash}p{0.359\linewidth}@{\hspace{0.030\linewidth}}>{\raggedright\arraybackslash}p{0.611\linewidth}@{}}
\begin{minipage}[t]{\linewidth}
\centering\textbf{Earlier Germanic changes}\par
\vspace{0.35em}
\raggedright
\centering\textbf{Northwest and West Germanic}\par
\raggedright
\vspace{0.2em}
\raggedright [no change]\par
\vspace{0.6em}
\centering\textbf{Early Anglo-Frisian}\par
\raggedright
\vspace{0.2em}
\raggedright [no change]\par
\end{minipage}
&
\begin{minipage}[t]{\linewidth}
\centering\textbf{Old English changes}\par
\vspace{0.35em}
\raggedright
\begin{tabularx}{\linewidth}{@{}>{\raggedright\arraybackslash}X@{\hspace{0.55em}}>{\raggedright\arraybackslash}p{0.20\linewidth}@{\hspace{0.25em}}}
\mbox{OE Breaking} & \emph{*wéorpaną} \\
OE Heavy Syllable Nasal Apocope & \emph{*wéorpan} \\
\mbox{OE Secondary Nasalization} & \emph{*wéorpąn} \\
\mbox{OE Weak Tail Reduction} & \emph{*wéorpan} \\
\end{tabularx}
\end{minipage}
\\
\end{tabular}
\end{minipage}%
} \endgroup

Old English form: \emph{weorpan}

\subsection{Reconstruction and comparative
evidence}\label{reconstruction-and-comparative-evidence-68}

Ringe and Taylor distinguish preterite \Recon{warp} `warp' from
infinitive \Recon{werpana} `warp', and the derivational input here is
the verbal form \Recon{wérpaną} `warp'
(\citeproc{ref-RingeTaylor2014}{Ringe and Taylor 2014}).

\subsection{Old English evidence}\label{old-english-evidence-69}

Clark Hall records \emph{weorpan} as the strong verb headword and
separately lists \emph{wearp} as both noun and preterite. Bright gives
the paradigm \emph{weorpan, wearp, wurpon, worpen}
(\citeproc{ref-ClarkHall1960}{Clark Hall 1960};
\citeproc{ref-BrightCassidyRingler1971}{Bright 1971}).

\subsection{Development to Old
English}\label{development-to-old-english-68}

Breaking before \emph{r + C} yields \emph{weor-}, and the infinitive
develops regularly to \emph{weorpan}
(\citeproc{ref-Campbell1959}{Campbell 1959};
\citeproc{ref-HoggGrammar1992}{Hogg 1992a}).

\subsection{Lexical note}\label{lexical-note-2}

English \emph{warp} also points to related \emph{wearp} material. The
target is the infinitive \emph{weorpan}.

\section{\texorpdfstring{wash --- OE
\emph{wascan}}{wash --- OE wascan}}\label{wash-oe-wascan}

\index[iv]{02oe@\ivlangheader{Old English}{West Saxon normalization unmarked}!wascan@\iventry{wascan}{}}
\index[iv]{03pgmc@\ivlangheader{Proto-Germanic}{}!waskana@\iventry{*wáskaną}{}}

Derivation: \emph{*wáskaną} \textgreater{} \emph{wascan} (regular).

\subsection{Derivation trace}\label{derivation-trace-83}

Proto input: \emph{*wáskaną}

\begingroup
\setlength{\fboxsep}{6pt}

\noindent\fbox{%
\begin{minipage}{0.97\linewidth}
\small
\begin{tabular}{@{}>{\raggedright\arraybackslash}p{0.359\linewidth}@{\hspace{0.030\linewidth}}>{\raggedright\arraybackslash}p{0.611\linewidth}@{}}
\begin{minipage}[t]{\linewidth}
\centering\textbf{Earlier Germanic changes}\par
\vspace{0.35em}
\raggedright
\centering\textbf{Northwest and West Germanic}\par
\raggedright
\vspace{0.2em}
\raggedright [no change]\par
\vspace{0.6em}
\centering\textbf{Early Anglo-Frisian}\par
\raggedright
\vspace{0.2em}
\raggedright [no change]\par
\end{minipage}
&
\begin{minipage}[t]{\linewidth}
\centering\textbf{Old English changes}\par
\vspace{0.35em}
\raggedright
\begin{tabularx}{\linewidth}{@{}>{\raggedright\arraybackslash}X@{\hspace{0.55em}}>{\raggedright\arraybackslash}p{0.18\linewidth}@{\hspace{0.25em}}}
\mbox{EAF Brightening} & \emph{*wæskaną} \\
\mbox{OE A Restoration} & \emph{*waskaną} \\
OE Heavy Syllable Nasal Apocope & \emph{*waskan} \\
\mbox{OE Secondary Nasalization} & \emph{*waskąn} \\
\mbox{OE Weak Tail Reduction} & \emph{*waskan} \\
\end{tabularx}
\end{minipage}
\\
\end{tabular}
\end{minipage}%
} \endgroup

Old English form: \emph{wascan}

\subsection{Reconstruction and comparative
evidence}\label{reconstruction-and-comparative-evidence-69}

Kroonen cites \emph{*waskan-}, Orel \Recon{waskanan} `wash', and Ringe
and Taylor likewise derive Old English \emph{wascan} from the same verb
family (\citeproc{ref-Kroonen2013}{Kroonen 2013};
\citeproc{ref-Orel2003}{Orel 2003, 489};
\citeproc{ref-RingeTaylor2014}{Ringe and Taylor 2014, 142}).

\subsection{Old English evidence}\label{old-english-evidence-70}

Clark Hall heads the verb as \emph{wascan} `wash', while Sievers-Brunner
also notes the variant \emph{wæscan} `wash'
(\citeproc{ref-ClarkHall1960}{Clark Hall 1960};
\citeproc{ref-SieversBrunner1965}{Brunner 1965}).

\subsection{Development to Old
English}\label{development-to-old-english-69}

From \Recon{wáskaną} `wash', brightening gives \Recon{wæskaną} `wash'.
A-restoration before the \emph{sC} cluster restores \emph{a}, and medial
\emph{sc} remains unpalatalized before the following back vowel,
yielding \emph{wascan} (\citeproc{ref-Campbell1959}{Campbell 1959};
\citeproc{ref-RingeTaylor2014}{Ringe and Taylor 2014, 142}).

\subsection{Form note}\label{form-note-23}

The conservative citation form \emph{wascan} `wash' is selected here.
Spellings such as \emph{wæscan} `wash' or \emph{wasċan} `wash' belong to
variant or normalized background rather than to the target of this
entry.

\section{\texorpdfstring{wax --- OE
\emph{weaxan}}{wax --- OE weaxan}}\label{wax-oe-weaxan}

\index[iv]{02oe@\ivlangheader{Old English}{West Saxon normalization unmarked}!weaxan@\iventry{weaxan}{}}
\index[iv]{03pgmc@\ivlangheader{Proto-Germanic}{}!waxsana@\iventry{*wáxsaną}{}}

Derivation: \emph{*wáxsaną} \textgreater{} \emph{weaxan} (regular).

\subsection{Derivation trace}\label{derivation-trace-84}

Proto input: \emph{*wáxsaną}

\begingroup
\setlength{\fboxsep}{6pt}

\noindent\fbox{%
\begin{minipage}{0.97\linewidth}
\footnotesize
\begin{tabular}{@{}>{\raggedright\arraybackslash}p{0.339\linewidth}@{\hspace{0.030\linewidth}}>{\raggedright\arraybackslash}p{0.630\linewidth}@{}}
\begin{minipage}[t]{\linewidth}
\centering\textbf{Earlier Germanic changes}\par
\vspace{0.35em}
\raggedright
\centering\textbf{Northwest and West Germanic}\par
\raggedright
\vspace{0.2em}
\raggedright [no change]\par
\vspace{0.6em}
\centering\textbf{Early Anglo-Frisian}\par
\raggedright
\vspace{0.2em}
\raggedright [no change]\par
\end{minipage}
&
\begin{minipage}[t]{\linewidth}
\centering\textbf{Old English changes}\par
\vspace{0.35em}
\raggedright
\begin{tabularx}{\linewidth}{@{}>{\raggedright\arraybackslash}X@{\hspace{0.55em}}>{\raggedright\arraybackslash}p{0.22\linewidth}@{\hspace{0.25em}}}
\mbox{EAF Brightening} & \emph{*wæxsaną} \\
\mbox{OE Breaking} & \emph{*weaxsaną} \\
OE Heavy Syllable Nasal Apocope & \emph{*weaxsan} \\
\mbox{OE Secondary Nasalization} & \emph{*weaxsąn} \\
\mbox{OE Weak Tail Reduction} & \emph{*weaxsan} \\
\mbox{OE Xs Merge} & \emph{*weaXSan} \\
\end{tabularx}
\end{minipage}
\\
\end{tabular}
\end{minipage}%
} \endgroup

Old English form: \emph{weaxan}

\subsection{Reconstruction and comparative
evidence}\label{reconstruction-and-comparative-evidence-70}

Kroonen cites the verb as \emph{*wahs(j)an-}, Orel as \Recon{waxsanan}
`wax', and Ringe and Taylor discuss the prehistory of Old English
\emph{weaxan} within the same verbal family
(\citeproc{ref-Kroonen2013}{Kroonen 2013}; \citeproc{ref-Orel2003}{Orel
2003, 478}; \citeproc{ref-RingeTaylor2014}{Ringe and Taylor 2014}).

\subsection{Old English evidence}\label{old-english-evidence-71}

Clark Hall gives \emph{weaxan} `wax' as the verb headword and separately
records bare \emph{wax} `wax' as a preterite form; Bright likewise
treats \emph{weaxan} `wax' as the infinitive
(\citeproc{ref-ClarkHall1960}{Clark Hall 1960};
\citeproc{ref-BrightCassidyRingler1971}{Bright 1971}).

\subsection{Development to Old
English}\label{development-to-old-english-70}

From \Recon{wáxsaną} `wax', brightening and breaking yield \emph{weax-},
and the infinitive develops regularly to \emph{weaxan}. The cluster is
preserved here, since \emph{xs} \textgreater{} \emph{s} belongs to forms
where another consonant follows, such as \emph{wæstm} `growth', not to
the infinitive itself (\citeproc{ref-Campbell1959}{Campbell 1959};
\citeproc{ref-SieversBrunner1965}{Brunner 1965}).

\subsection{Lexical note}\label{lexical-note-3}

The target here is the infinitive \emph{weaxan} `wax'. Noun \emph{weax}
`wax' and preterite \emph{wax} / \emph{wēox} `waxed' belong to different
lexical or paradigm slots.

\section{\texorpdfstring{way --- OE
\emph{weġ}}{way --- OE weġ}}\label{way-oe-weux121}

\index[iv]{02oe@\ivlangheader{Old English}{West Saxon normalization unmarked}!weg@\iventry{weġ}{}}
\index[iv]{03pgmc@\ivlangheader{Proto-Germanic}{}!wegaz@\iventry{*wégaz}{}}

Derivation: \emph{*wégaz} \textgreater{} \emph{weġ} (regular).

\subsection{Derivation trace}\label{derivation-trace-85}

Proto input: \emph{*wégaz}

\begingroup
\setlength{\fboxsep}{6pt}

\noindent\fbox{%
\begin{minipage}{0.97\linewidth}
\small
\begin{tabular}{@{}>{\raggedright\arraybackslash}p{0.359\linewidth}@{\hspace{0.030\linewidth}}>{\raggedright\arraybackslash}p{0.611\linewidth}@{}}
\begin{minipage}[t]{\linewidth}
\centering\textbf{Earlier Germanic changes}\par
\vspace{0.35em}
\raggedright
\centering\textbf{Northwest and West Germanic}\par
\raggedright
\vspace{0.2em}
\begin{tabularx}{\linewidth}{@{}>{\raggedright\arraybackslash}X@{\hspace{0.55em}}>{\raggedright\arraybackslash}p{0.14\linewidth}@{\hspace{0.25em}}}
\mbox{EAF Final Z Deletion} & \emph{*wéga} \\
\end{tabularx}
\vspace{0.6em}
\centering\textbf{Early Anglo-Frisian}\par
\raggedright
\vspace{0.2em}
\raggedright [no change]\par
\end{minipage}
&
\begin{minipage}[t]{\linewidth}
\centering\textbf{Old English changes}\par
\vspace{0.35em}
\raggedright
\begin{tabularx}{\linewidth}{@{}>{\raggedright\arraybackslash}X@{\hspace{0.55em}}>{\raggedright\arraybackslash}p{0.14\linewidth}@{\hspace{0.25em}}}
\mbox{PWGmc Final Bare A Loss} & \emph{*wég} \\
\mbox{OE Velar Palatalization} & \emph{*wéʝ} \\
\mbox{OE Palatal Fricative Merger} & \emph{*wéj} \\
\end{tabularx}
\end{minipage}
\\
\end{tabular}
\end{minipage}%
} \endgroup

Old English form: \emph{weġ}

\subsection{Reconstruction and comparative
evidence}\label{reconstruction-and-comparative-evidence-71}

Kroonen cites the noun as \emph{*wega-} `way, road', while the selected
derivational form here is nominative-singular \Recon{wégaz} `way'
(\citeproc{ref-Kroonen2013}{Kroonen 2013}). Campbell, Hogg, and Ringe
and Taylor use the same word as the standard contrast between singular
palatal \emph{weġ} `way' and inflected \emph{wegas} / \emph{wegum}
(\citeproc{ref-Campbell1959}{Campbell 1959};
\citeproc{ref-HoggGrammar1992}{Hogg 1992a};
\citeproc{ref-RingeTaylor2014}{Ringe and Taylor 2014, 341}).

\subsection{Old English evidence}\label{old-english-evidence-72}

The Old English singular is the ordinary noun \emph{weg} `way', here
normalized as \emph{weġ} `way' to show the palatal final. The
contrasting plural and oblique forms \emph{wegas} `ways', \emph{wegum}
`ways' keep a velar stop before the following back vowel
(\citeproc{ref-Campbell1959}{Campbell 1959};
\citeproc{ref-RingeTaylor2014}{Ringe and Taylor 2014, 341}).

\subsection{Development to Old
English}\label{development-to-old-english-71}

From \Recon{wégaz} `way', final \emph{*z} is lost and the weak tail
apocopates, leaving word-final \emph{*g} after a front vowel. In that
environment Old English palatalization yields \emph{weġ} `way', whereas
\emph{wegas} `ways' remains velar because the following \emph{a} blocks
the same outcome (\citeproc{ref-Campbell1959}{Campbell 1959};
\citeproc{ref-HoggGrammar1992}{Hogg 1992a};
\citeproc{ref-RingeTaylor2014}{Ringe and Taylor 2014, 341}).

\subsection{Form note}\label{form-note-24}

Normalized \emph{weġ} `way' and dictionary \emph{weg} `way' represent
the same noun. No Old English form \emph{wē} `we' is supported by the
evidence for `way'.

\section{\texorpdfstring{weapon --- OE
\emph{wǣpn}}{weapon --- OE wǣpn}}\label{weapon-oe-wux1e3pn}

\index[iv]{02oe@\ivlangheader{Old English}{West Saxon normalization unmarked}!waepn@\iventry{wǣpn}{}}
\index[iv]{03pgmc@\ivlangheader{Proto-Germanic}{}!wepna@\iventry{*wḗpną}{}}

Derivation: \emph{*wḗpną} \textgreater{} \emph{wǣpn} (regular).

\subsection{Derivation trace}\label{derivation-trace-86}

Proto input: \emph{*wḗpną}

\begingroup
\setlength{\fboxsep}{6pt}

\noindent\fbox{%
\begin{minipage}{0.97\linewidth}
\small
\begin{tabular}{@{}>{\raggedright\arraybackslash}p{0.485\linewidth}@{\hspace{0.030\linewidth}}>{\raggedright\arraybackslash}p{0.485\linewidth}@{}}
\begin{minipage}[t]{\linewidth}
\centering\textbf{Earlier Germanic changes}\par
\vspace{0.35em}
\raggedright
\centering\textbf{Northwest and West Germanic}\par
\raggedright
\vspace{0.2em}
\begin{tabularx}{\linewidth}{@{}>{\raggedright\arraybackslash}X@{\hspace{0.55em}}>{\raggedright\arraybackslash}p{0.14\linewidth}@{\hspace{0.25em}}}
\mbox{PNWGmc Long E Lowering} & \emph{*wāpną} \\
\mbox{EAF Long A Fronting} & \emph{*wǣpną} \\
\end{tabularx}
\vspace{0.6em}
\centering\textbf{Early Anglo-Frisian}\par
\raggedright
\vspace{0.2em}
\raggedright [no change]\par
\end{minipage}
&
\begin{minipage}[t]{\linewidth}
\centering\textbf{Old English changes}\par
\vspace{0.35em}
\raggedright
\begin{tabularx}{\linewidth}{@{}>{\raggedright\arraybackslash}X@{\hspace{0.55em}}>{\raggedright\arraybackslash}p{0.14\linewidth}@{\hspace{0.25em}}}
OE Heavy Syllable Nasal Apocope & \emph{*wǣpn} \\
\end{tabularx}
\end{minipage}
\\
\end{tabular}
\end{minipage}%
} \endgroup

Old English form: \emph{wǣpn}

\subsection{Reconstruction and comparative
evidence}\label{reconstruction-and-comparative-evidence-72}

Kroonen reconstructs a double-stem noun \emph{*wēbna-} \textasciitilde{}
\emph{*wēpna-} and cites OE \emph{wæpn} `weapon' among its reflexes
(\citeproc{ref-Kroonen2013}{Kroonen 2013, 617}). The derivational input
\Recon{wḗpną} `weapon' represents the unbroken citation-form noun rather
than the later broken simplex.

\subsection{Old English evidence}\label{old-english-evidence-73}

Campbell's cluster-noun discussion preserves unbroken \emph{wépn}
`weapon' beside broken \emph{wépen} `weapon'-type forms
(\citeproc{ref-Campbell1959}{Campbell 1959, 150};
\citeproc{ref-Campbell1959}{1959, 226--27}). Bright contrasts broken
nominative \emph{wǣpen} `weapon' / \emph{wapen} `weapon' with unbroken
oblique \emph{wǣpnes} `weapon', while Clark Hall lemmatizes the noun
under \emph{wapen} `weapon' and also preserves unbroken forms in
compounds and related spellings
(\citeproc{ref-BrightCassidyRingler1971}{Bright 1971, 29};
\citeproc{ref-ClarkHall1960}{Clark Hall 1960, 355}).

\subsection{Development to Old
English}\label{development-to-old-english-72}

Northwest Germanic lowering gives \emph{wǣpn} `weapon', and loss of the
final nasal vowel leaves the unbroken cluster word-finally. The Old
English form here is the attested unbroken form \emph{wǣpn} `weapon'.

\subsection{Form note}\label{form-note-25}

The ordinary late West Saxon simplex headword is \emph{wǣpen} `weapon',
but Campbell's noun-class discussion also preserves unbroken \emph{wépn}
`weapon' beside broken \emph{wépen} `weapon'-type forms
(\citeproc{ref-Campbell1959}{Campbell 1959, 150};
\citeproc{ref-Campbell1959}{1959, 226--27}). \emph{wǣpnes} `weapon'
remains the regular unbroken oblique comparator
(\citeproc{ref-BrightCassidyRingler1971}{Bright 1971, 29};
\citeproc{ref-ClarkHall1960}{Clark Hall 1960, 355}).

\section{\texorpdfstring{west --- OE
\emph{west}}{west --- OE west}}\label{west-oe-west}

\index[iv]{02oe@\ivlangheader{Old English}{West Saxon normalization unmarked}!west@\iventry{west}{}}
\index[iv]{03pgmc@\ivlangheader{Proto-Germanic}{}!westa@\iventry{*wéstą}{}}

Derivation: \emph{*wéstą} \textgreater{} \emph{west} (regular).

\subsection{Derivation trace}\label{derivation-trace-87}

Proto input: \emph{*wéstą}

\begingroup
\setlength{\fboxsep}{6pt}

\noindent\fbox{%
\begin{minipage}{0.97\linewidth}
\small
\begin{tabular}{@{}>{\raggedright\arraybackslash}p{0.359\linewidth}@{\hspace{0.030\linewidth}}>{\raggedright\arraybackslash}p{0.611\linewidth}@{}}
\begin{minipage}[t]{\linewidth}
\centering\textbf{Earlier Germanic changes}\par
\vspace{0.35em}
\raggedright
\centering\textbf{Northwest and West Germanic}\par
\raggedright
\vspace{0.2em}
\raggedright [no change]\par
\vspace{0.6em}
\centering\textbf{Early Anglo-Frisian}\par
\raggedright
\vspace{0.2em}
\raggedright [no change]\par
\end{minipage}
&
\begin{minipage}[t]{\linewidth}
\centering\textbf{Old English changes}\par
\vspace{0.35em}
\raggedright
\begin{tabularx}{\linewidth}{@{}>{\raggedright\arraybackslash}X@{\hspace{0.55em}}>{\raggedright\arraybackslash}p{0.14\linewidth}@{\hspace{0.25em}}}
OE Heavy Syllable Nasal Apocope & \emph{*wést} \\
\end{tabularx}
\end{minipage}
\\
\end{tabular}
\end{minipage}%
} \endgroup

Old English form: \emph{west}

\subsection{Reconstruction and comparative
evidence}\label{reconstruction-and-comparative-evidence-73}

Orel distinguishes a source-labelled neuter \emph{*westan} from the
directional \emph{*westanē̆} formation. They are related words, not
alternative spellings of the same selected ending
(\citeproc{ref-Orel2003}{Orel 2003, 459}). Fulk's account of neuter
\emph{-aⁿ < -an < -om} supports the normalization
\Recon{wéstą}\index[iv]{03pgmc@\ivlangheader{Proto-Germanic}{}!westa@\iventry{*wéstą}{}}
`west' used here (\citeproc{ref-Fulk2018}{Fulk 2018, 146}, §7.8). Stress
and nasal-vowel notation are CAPR's explicit representation, not
additional characters printed in Orel's citation.

Kroonen's \emph{*westera-} belongs to the directional family, with a
different formation. Its deeper associations and the comparisons
reported by Orel remain qualified comparative evidence
(\citeproc{ref-Kroonen2013}{Kroonen 2013, 582--83};
\citeproc{ref-Orel2003}{Orel 2003, 459}). They do not require treating
the neuter base as the former selected from-the-west derivative.

\subsection{Old English evidence}\label{old-english-evidence-74}

{\emph{west}}\index[iv]{02oe@\ivlangheader{Old English}{West Saxon normalization unmarked}!west@\iventry{west}{}}
`west, westwards' is an ordinary attested Old English adverb. Toller
cites the Stilicho passage in \emph{Orosius}6.37, Sweet296.7, including
\emph{þisne her west}, under \emph{west}
(\citeproc{ref-Toller1921}{Toller 1921}, s.v.~west). The selected
inherited formation is neuter; the Old English use is adverbial. These
are compatible descriptions at different historical and grammatical
levels, not a claim that the attested adverb visibly inflects as a
neuter noun.

The genuine
{\emph{westan}}\index[iv]{02oe@\ivlangheader{Old English}{West Saxon normalization unmarked}!westan@\iventry{westan}{}}
`from the west' and
{\emph{westane}}\index[iv]{02oe@\ivlangheader{Old English}{West Saxon normalization unmarked}!westane@\iventry{westane}{}}
`from the west' remain important directional comparanda. Orel's citation
and quantity alternatives must be preserved rather than replaced by the
unsupported selected spelling \emph{westene}
(\citeproc{ref-Orel2003}{Orel 2003, 459}). Withdrawing that former
target is an evidential correction, not a claim that the wider
directional paradigm is solved.

\subsection{Development to Old
English}\label{development-to-old-english-73}

The selected
\Recon{wéstą}\index[iv]{03pgmc@\ivlangheader{Proto-Germanic}{}!westa@\iventry{*wéstą}{}}
has a heavy root syllable. The existing final nasal-vowel apocope
removes the ending, giving {\emph{west}}. There is no medial vowel for
\hyperref[rule-OENonHighSyn]{SC114 OENonHighSyn} to delete. West is
consequently a negative syncope control, not a positive application
manufactured by retaining the longer former input.

The former
\Recon{wéstanē}\index[iv]{03pgmc@\ivlangheader{Proto-Germanic}{}!westane@\iventry{*wéstanē}{}}
now gives counterfactual \emph{westne}, not its old target
\emph{westene}. Neither output establishes the history of the attested
directional variants. The exact inherited ending quantity and its
relative chronology remain open; the regular neuter-base span does not
depend on resolving those wider questions.

\section{\texorpdfstring{wether --- OE
\emph{weþer}}{wether --- OE weþer}}\label{wether-oe-weuxfeer}

\index[iv]{02oe@\ivlangheader{Old English}{West Saxon normalization unmarked}!wether@\iventry{weþer}{}}
\index[iv]{03pgmc@\ivlangheader{Proto-Germanic}{}!wethruz@\iventry{*wéθruz}{}}

Derivation: \emph{*wéθruz} \textgreater{} \emph{weþer} (regular).

\subsection{Derivation trace}\label{derivation-trace-88}

Proto input: \emph{*wéθruz}

\begingroup
\setlength{\fboxsep}{6pt}

\noindent\fbox{%
\begin{minipage}{0.97\linewidth}
\small
\begin{tabular}{@{}>{\raggedright\arraybackslash}p{0.485\linewidth}@{\hspace{0.030\linewidth}}>{\raggedright\arraybackslash}p{0.485\linewidth}@{}}
\begin{minipage}[t]{\linewidth}
\centering\textbf{Earlier Germanic changes}\par
\vspace{0.35em}
\raggedright
\centering\textbf{Northwest and West Germanic}\par
\raggedright
\vspace{0.2em}
\begin{tabularx}{\linewidth}{@{}>{\raggedright\arraybackslash}X@{\hspace{0.55em}}>{\raggedright\arraybackslash}p{0.14\linewidth}@{\hspace{0.25em}}}
\mbox{EAF Final Z Deletion} & \emph{*wéθru} \\
\end{tabularx}
\vspace{0.6em}
\centering\textbf{Early Anglo-Frisian}\par
\raggedright
\vspace{0.2em}
\raggedright [no change]\par
\end{minipage}
&
\begin{minipage}[t]{\linewidth}
\centering\textbf{Old English changes}\par
\vspace{0.35em}
\raggedright
\begin{tabularx}{\linewidth}{@{}>{\raggedright\arraybackslash}X@{\hspace{0.55em}}>{\raggedright\arraybackslash}p{0.14\linewidth}@{\hspace{0.25em}}}
\mbox{OE High Vowel Apocope} & \emph{*wéθr} \\
\mbox{OE Epenthetic Vowel} & \emph{*wéθer} \\
\end{tabularx}
\end{minipage}
\\
\end{tabular}
\end{minipage}%
} \endgroup

Old English form: \emph{weþer}

\subsection{Comparative reconstruction and
development}\label{comparative-reconstruction-and-development}

Old English \emph{weþer} `wether, ram' belongs with Gothic \emph{wiþrus}
`lamb, young ram', Old Norse \emph{veðr} `ram', Old Saxon \emph{withar}
`wether' and Old High German \emph{widar} `ram'. These are masculine
animal names, not the neuter weather noun, despite the similar
spellings. Kroonen reconstructs \emph{*weþru-}; Kluge--Seebold likewise
gives a masculine Germanic \emph{*weþru-}, while allowing a similar stem
formation (\citeproc{ref-Kroonen2013}{Kroonen 2013, 584};
\citeproc{ref-KlugeSeebold2011}{Kluge 2011, 985}).

The inherited root has short \emph{*e}. Gothic \emph{i} does not require
Proto-Germanic \emph{*i}: inherited short \emph{*e} and \emph{*i} merged
as \emph{i} in Gothic outside its breaking environments
(\citeproc{ref-Fulk2018}{Fulk 2018, 59}). The continental \emph{i}-forms
require equal care. Raising before a following \emph{u} is independently
documented in Old Saxon and Old High German, but it is not universal
there; that general development is not a complete, independently dated
history of every cell of this noun (\citeproc{ref-Fulk2018}{Fulk 2018,
69}). The comparison consequently supports inherited \emph{*e} without
requiring every daughter-language vowel to preserve it unchanged.

The principal formation question concerns the inflection, not the root
vowel or masculine gender. Orel records both \emph{*weþruz} and
\emph{*weþraz}, respectively masculine \emph{u}- and \emph{a}-stem
alternatives, and notes traces of \emph{u}-inflection in Norse
(\citeproc{ref-Orel2003}{Orel 2003, 459--60}). Gothic nominative
\emph{wiþrus} supplies particularly useful positive evidence for the
older \emph{u}-stem. CAPR therefore adopts the masculine nominative
singular \emph{*weþruz}, represented as \emph{*wéθruz}, while retaining
\emph{*weþraz} as a legitimate alternative. The latter may represent
early inflectional remodelling; the dictionaries do not establish its
precise date.

The Old English masculine paradigm, including genitive \emph{weþeres},
does not decide the original Germanic stem class. Bosworth--Toller
records \emph{weþer, es; m.} and the senses `a wether, a ram'
(\citeproc{ref-BosworthToller1898}{Bosworth and Toller 1898, 1212}).
Masculine \emph{u}-stems underwent reclassification towards other
declensions in West Germanic, and heavy-stem nominatives lost their
final short \emph{*u} (\citeproc{ref-Fulk2018}{Fulk 2018, 163}). In the
selected nominative, loss of final \emph{*-z} and \emph{*-u} leaves the
consonantal ending resolved as unstressed \emph{-er} in \emph{weþer},
consistent with the vocalization of final postconsonantal sonorants
(\citeproc{ref-Fulk2018}{Fulk 2018, 91--92}). The root vowel is already
\emph{*e}; no early lowering of \emph{*i} is needed for this derivation.

The further etymology is less secure. Orel connects the noun with the
Indo-European word for `year', comparing Hittite, Greek and Albanian
year-words and names for young cattle (\citeproc{ref-Orel2003}{Orel
2003, 459--60}). Kroonen makes that connection conditional: if correct,
the original designation was `yearling, one-year-old animal', with
Sanskrit \emph{vatsá-} `calf' among the comparanda
(\citeproc{ref-Kroonen2013}{Kroonen 2013, 584}). Kluge--Seebold notes
phonological difficulties with the traditional comparison involving
Latin \emph{vitulus} `calf', and also considers derivation from a root
meaning `sheep' (\citeproc{ref-KlugeSeebold2011}{Kluge 2011, 985--86}).
A young-animal designation could develop into `ram' or `wether', but
this attractive semantic account is not as secure as the Germanic
masculine \emph{u}-stem reconstruction.

\section{\texorpdfstring{where --- OE
\emph{hwǣr}}{where --- OE hwǣr}}\label{where-oe-hwux1e3r}

\index[iv]{02oe@\ivlangheader{Old English}{West Saxon normalization unmarked}!hwaer@\iventry{hwǣr}{}}
\index[iv]{03pgmc@\ivlangheader{Proto-Germanic}{}!xwar@\iventry{*xwár}{}}

Derivation: \emph{*xwár} \textgreater{} \emph{hwǣr} (regular).

\subsection{Derivation trace}\label{derivation-trace-89}

Proto input: \emph{*xwár}

\begingroup
\setlength{\fboxsep}{6pt}

\noindent\fbox{%
\begin{minipage}{0.97\linewidth}
\small
\begin{tabular}{@{}>{\raggedright\arraybackslash}p{0.611\linewidth}@{\hspace{0.030\linewidth}}>{\raggedright\arraybackslash}p{0.359\linewidth}@{}}
\begin{minipage}[t]{\linewidth}
\centering\textbf{Earlier Germanic changes}\par
\vspace{0.35em}
\raggedright
\centering\textbf{Northwest and West Germanic}\par
\raggedright
\vspace{0.2em}
\begin{tabularx}{\linewidth}{@{}>{\raggedright\arraybackslash}X@{\hspace{0.55em}}>{\raggedright\arraybackslash}p{0.14\linewidth}@{\hspace{0.25em}}}
PNWGmc Monosyllabic Final R Lengthening & \emph{*xwār} \\
\mbox{EAF Long A Fronting} & \emph{*xwǣr} \\
\end{tabularx}
\vspace{0.6em}
\centering\textbf{Early Anglo-Frisian}\par
\raggedright
\vspace{0.2em}
\raggedright [no change]\par
\end{minipage}
&
\begin{minipage}[t]{\linewidth}
\centering\textbf{Old English changes}\par
\vspace{0.35em}
\raggedright
\raggedright [no change]\par
\end{minipage}
\\
\end{tabular}
\end{minipage}%
} \endgroup

Old English form: \emph{hwǣr}

\subsection{Old English evidence}\label{old-english-evidence-75}

Clark Hall records
{\emph{hwǣr}}\index[iv]{02oe@\ivlangheader{Old English}{West Saxon normalization unmarked}!hwaer@\iventry{hwǣr}{}}
`where', the West Saxon comparison form used here
(\citeproc{ref-ClarkHall1960}{Clark Hall 1960, 170}). Hill sets it
beside Anglian
{\emph{hwēr}}\index[iv]{02oe@\ivlangheader{Old English}{West Saxon normalization unmarked}!hwer@\iventry{hwēr}{}}
`where' (\citeproc{ref-Hill2017}{Hill 2017, 142--43}). The contrast
concerns the quality of a long vowel, not lengthening in West Saxon
alone. The macrons belong to normalized scholarly forms; they do not
imply that every manuscript token explicitly marks quantity.

\subsection{Reconstruction and comparative
evidence}\label{reconstruction-and-comparative-evidence-74}

Kroonen gives
\Recon{hwar}\index[iv]{03pgmc@\ivlangheader{Proto-Germanic}{}!hwar@\iventry{*hwar}{}}
`where', an interrogative formation with locative \emph{*-r}. He
explicitly derives the long-vowel West Germanic forms by monosyllabic
lengthening, rejecting inherited \emph{*hwēr} as their source
(\citeproc{ref-Kroonen2013}{Kroonen 2013, 263}). Orel instead preserves
both
\Recon{xwar}\index[iv]{03pgmc@\ivlangheader{Proto-Germanic}{}!xwar@\iventry{*xwar}{}}
`where' and \emph{*xwēr} as citation alternatives, connecting the
formation with pronominal \emph{*xwa-} (\citeproc{ref-Orel2003}{Orel
2003, 198}). The selected
\Recon{xwár}\index[iv]{03pgmc@\ivlangheader{Proto-Germanic}{}!xwar@\iventry{*xwár}{}}
`where' follows the short-vowel line. Its acute marks the accented root,
while \emph{*xw-} represents the Germanic dorsal-plus-glide sequence
corresponding to the dictionaries' \emph{*hw-} or \emph{*xw-}.

The extended interrogative supplies an important morphological
comparison. Gothic \emph{ƕarjis} `which' and Old Norse \emph{hverr}
`who, which' support an inherited disyllabic \emph{*hwarja-} with short
root vowel. Hill argues that such a formation could precede lengthening
of the monosyllabic adverb; its short vowel therefore need not date
lengthening to West Germanic. Gothic adverbial \emph{ƕar} `where' itself
leaves quantity ambiguous in writing (\citeproc{ref-Hill2017}{Hill 2017,
142, 147--48}).

\subsection{Development to Old
English}\label{development-to-old-english-74}

From short \emph{*hwar},
\hyperref[rule-PNWGmcMonosyllabicFinalRLengthening]{SC112, stressed
monosyllabic lengthening before inherited final \emph{r}} gives
\emph{*hwār}; then \hyperref[rule-EAFLongAFronting]{SC101, long-\emph{ā}
fronting} gives \emph{*hwǣr}. The initial consonant sequence is
represented orthographically as \emph{hw-}. The derivational input with
short \emph{*a} thus reaches \emph{hwǣr} through lengthening followed by
fronting, without inflectional retargeting. The parallel \emph{þǣr}
`there' makes the same vowel development visible in a different
pronominal formation.

The historical interpretation remains narrower than a uniquely dated
reconstruction. Hill's Proto-Germanic account depends in part on his
explanation of the Scandinavian short-vowel forms, while Ringe and
Taylor retain a later West Germanic lengthening
(\citeproc{ref-Hill2017}{Hill 2017, 146--48};
\citeproc{ref-RingeTaylor2014}{Ringe and Taylor 2014, 65}). The early
placement adopted here preserves inherited-\emph{r} conditioning before
rhotacism could obscure consonant ancestry. That analytical choice does
not establish an independently dated ordering against SC015 or the
lowering of inherited \emph{*ē₁} in SC024.

\section{\texorpdfstring{who --- OE
\emph{hwā}}{who --- OE hwā}}\label{who-oe-hwux101}

\index[iv]{02oe@\ivlangheader{Old English}{West Saxon normalization unmarked}!hwa@\iventry{hwā}{}}
\index[iv]{03pgmc@\ivlangheader{Proto-Germanic}{}!xwaz@\iventry{*xwáz}{}}

Derivation: \emph{*xwáz} \textgreater{} \emph{hwā} (regular).

\subsection{Derivation trace}\label{derivation-trace-90}

Proto input: \emph{*xwáz}

\begingroup
\setlength{\fboxsep}{6pt}

\noindent\fbox{%
\begin{minipage}{0.97\linewidth}
\small
\begin{tabular}{@{}>{\raggedright\arraybackslash}p{0.611\linewidth}@{\hspace{0.030\linewidth}}>{\raggedright\arraybackslash}p{0.359\linewidth}@{}}
\begin{minipage}[t]{\linewidth}
\centering\textbf{Earlier Germanic changes}\par
\vspace{0.35em}
\raggedright
\centering\textbf{Northwest and West Germanic}\par
\raggedright
\vspace{0.2em}
\begin{tabularx}{\linewidth}{@{}>{\raggedright\arraybackslash}X@{\hspace{0.55em}}>{\raggedright\arraybackslash}p{0.14\linewidth}@{\hspace{0.25em}}}
\mbox{Monosyllabic Final Z Loss} & \emph{*xwā} \\
\end{tabularx}
\vspace{0.6em}
\centering\textbf{Early Anglo-Frisian}\par
\raggedright
\vspace{0.2em}
\raggedright [no change]\par
\end{minipage}
&
\begin{minipage}[t]{\linewidth}
\centering\textbf{Old English changes}\par
\vspace{0.35em}
\raggedright
\raggedright [no change]\par
\end{minipage}
\\
\end{tabular}
\end{minipage}%
} \endgroup

Old English form: \emph{hwā}

\subsection{Reconstruction and comparative
evidence}\label{reconstruction-and-comparative-evidence-75}

The derivational input is the nominative singular masculine of the
interrogative pronoun, PGmc
\Recon{xwáz}\index[iv]{03pgmc@\ivlangheader{Proto-Germanic}{}!xwaz@\iventry{*xwáz}{}}
`who' (traditionally *hwaz; Gothic
{ƕas}\index[iv]{12goth@\ivlangheader{Gothic}{}!hwas@\iventry{ƕas}{}}
`who'). Ringe and Taylor cite exactly this form among the witnesses for
the northern West Germanic loss of word-final *-z in stressed
monosyllables: ``*hwaz `who' \textgreater{} OE hwā''
(\citeproc{ref-RingeTaylor2014}{Ringe and Taylor 2014, 86}). The
nominative is the natural citation cell of the interrogative and the
direct ancestor of the Old English headword form, so no paradigm-cell
retargeting is involved.

\subsection{Old English evidence}\label{old-english-evidence-76}

West Saxon
{\emph{hwā}}\index[iv]{02oe@\ivlangheader{Old English}{West Saxon normalization unmarked}!hwa@\iventry{hwā}{}}
`who' is the standard headword form
(\citeproc{ref-Campbell1959}{Campbell 1959, § 125}, p.~49). The
vowel is decisive for the derivation: Campbell states that the back
vowel of \emph{hwā} never underwent Anglo-Frisian brightening --- ``*hwǣ
does not exist'' --- and Sievers-Brunner records the same restriction
(\citeproc{ref-Campbell1959}{Campbell 1959, § 125}, p.~49;
\citeproc{ref-SieversBrunner1965}{Brunner 1965, § 137} Anm. 1,
p.~129).

\subsection{Development to Old
English}\label{development-to-old-english-75}

The derivation is two historical steps. First, the later northern West
Germanic loss of word-final *-z in stressed monosyllables removed the
sibilant with compensatory lengthening of the short nucleus:
\Recon{xwáz}\index[iv]{03pgmc@\ivlangheader{Proto-Germanic}{}!xwaz@\iventry{*xwáz}{}}
`who' \textgreater{}
\Recon{hwā}\index[iv]{06nsgmc@\ivlangheader{Northern West Germanic}{}!hwa@\iventry{*hwā}{}}
`who' (\citeproc{ref-RingeTaylor2014}{Ringe and Taylor 2014, 86}). Old
High German instead rhotacized the sibilant (cf.~OHG
{wer}\index[iv]{10ohg@\ivlangheader{Old High German}{}!wer@\iventry{wer}{}}
`who' type inflection), which is why this loss counts among the
diagnostic Ingvaeonic features (\citeproc{ref-Fulk2018}{Fulk 2018, 18}).
Second, the resulting long back *ā stayed back: Anglo-Frisian
brightening of long final *ā requires a preceding nucleus in the word,
and in this monosyllable there is none, so the form surfaces directly as
OE \emph{hwā} (\citeproc{ref-Campbell1959}{Campbell 1959, § 125},
p.~49).

An alternative account derives the endingless pronoun from a generalized
unaccented sentence variant *hwa with later lengthening
(\citeproc{ref-Campbell1959}{Campbell 1959, 166};
\citeproc{ref-SieversBrunner1965}{Brunner 1965, § 182}, p.~160).
Ringe and Taylor reject the unaccented-variant analysis for this witness
set because \emph{mā} `more' and \emph{cū} `cow' are not plausibly
unaccented words (\citeproc{ref-RingeTaylor2014}{Ringe and Taylor 2014,
86}). The project follows Ringe and Taylor; the alternative is recorded,
not adopted.

\subsection{Why this row is in the
corpus}\label{why-this-row-is-in-the-corpus-2}

This row makes SC097, the monosyllabic final \emph{-z loss,
corpus-witnessed: before its addition the change was validated only by
synthetic controls. It also witnesses the ordering SC097 before SC003
rhotacism --- a rhotacized \Pred{hwar} could never yield }hwā* --- and
exercises the narrowed long-final clause of Anglo-Frisian brightening.
See
\texttt{docs/sound\_changes/audits/corpus-maturation-01-candidate-adjudication.md}
§1.

\section{\texorpdfstring{will --- OE
\emph{willa}}{will --- OE willa}}\label{will-oe-willa}

\index[iv]{02oe@\ivlangheader{Old English}{West Saxon normalization unmarked}!willa@\iventry{willa}{}}
\index[iv]{03pgmc@\ivlangheader{Proto-Germanic}{}!weljo@\iventry{*wéljô}{}}

Derivation: \emph{*wéljô} \textgreater{} \emph{willa} (regular).

\subsection{Derivation trace}\label{derivation-trace-91}

Proto input: \emph{*wéljô}

\begingroup
\setlength{\fboxsep}{6pt}

\noindent\fbox{%
\begin{minipage}{0.97\linewidth}
\small
\begin{tabular}{@{}>{\raggedright\arraybackslash}p{0.359\linewidth}@{\hspace{0.030\linewidth}}>{\raggedright\arraybackslash}p{0.611\linewidth}@{}}
\begin{minipage}[t]{\linewidth}
\centering\textbf{Earlier Germanic changes}\par
\vspace{0.35em}
\raggedright
\centering\textbf{Northwest and West Germanic}\par
\raggedright
\vspace{0.2em}
\begin{tabularx}{\linewidth}{@{}>{\raggedright\arraybackslash}X@{\hspace{0.55em}}>{\raggedright\arraybackslash}p{0.24\linewidth}@{\hspace{0.25em}}}
\mbox{PWGmc J Gemination} & \emph{*wélljô} \\
\end{tabularx}
\vspace{0.6em}
\centering\textbf{Early Anglo-Frisian}\par
\raggedright
\vspace{0.2em}
\raggedright [no change]\par
\end{minipage}
&
\begin{minipage}[t]{\linewidth}
\centering\textbf{Old English changes}\par
\vspace{0.35em}
\raggedright
\begin{tabularx}{\linewidth}{@{}>{\raggedright\arraybackslash}X@{\hspace{0.55em}}>{\raggedright\arraybackslash}p{0.14\linewidth}@{\hspace{0.25em}}}
\mbox{OE I Umlaut} & \emph{*willjô} \\
OE Unstressed Long Vowel Shortening & \emph{*willja} \\
\mbox{OE J Loss After Heavy} & \emph{*willa} \\
\end{tabularx}
\end{minipage}
\\
\end{tabular}
\end{minipage}%
} \endgroup

Old English form: \emph{willa}

\subsection{Reconstruction and comparative
evidence}\label{reconstruction-and-comparative-evidence-76}

Kroonen separates noun \Recon{weljan} `will, wish' from verb
\Recon{weljan} `to want', while Orel and Kluge represent the noun as
\Recon{weljōn} `will' (\citeproc{ref-Kroonen2013}{Kroonen 2013};
\citeproc{ref-Orel2003}{Orel 2003};
\citeproc{ref-KlugeSeebold2011}{Kluge 2011}). The selected derivational
form \Recon{wéljô} `will' is the noun-side input used for this row.

\subsection{Old English evidence}\label{old-english-evidence-77}

Clark Hall lemmatizes noun \emph{willa m.} separately from verb
{\emph{willan}}\index[iv]{02oe@\ivlangheader{Old English}{West Saxon normalization unmarked}!willan@\iventry{willan}{}}
`will, want' (\citeproc{ref-ClarkHall1960}{Clark Hall 1960, 368}). The
Old English form here is the noun citation form, not the related verb.

\subsection{Development to Old
English}\label{development-to-old-english-76}

From \Recon{wéljô} `will', j-gemination yields a heavy stem, i-umlaut
gives \emph{will-}, and later shortening plus j-loss produce
{\emph{willa}}\index[iv]{02oe@\ivlangheader{Old English}{West Saxon normalization unmarked}!willa@\iventry{willa}{}}
`will, desire'. The noun is therefore a regular weak masculine outcome.

\subsection{Lexical note}\label{lexical-note-4}

The target here is the noun \emph{willa} `will, wish'. Related verb
{\emph{willan}}\index[iv]{02oe@\ivlangheader{Old English}{West Saxon normalization unmarked}!willan@\iventry{willan}{}}
`will, want' belongs to a separate lexeme and should not be substituted
for the noun row (\citeproc{ref-Kroonen2013}{Kroonen 2013};
\citeproc{ref-ClarkHall1960}{Clark Hall 1960, 368}).

\section{\texorpdfstring{wind --- OE
\emph{windan}}{wind --- OE windan}}\label{wind-oe-windan}

\index[iv]{02oe@\ivlangheader{Old English}{West Saxon normalization unmarked}!windan@\iventry{windan}{}}
\index[iv]{03pgmc@\ivlangheader{Proto-Germanic}{}!windana@\iventry{*wíndaną}{}}

Derivation: \emph{*wíndaną} \textgreater{} \emph{windan} (regular).

\subsection{Derivation trace}\label{derivation-trace-92}

Proto input: \emph{*wíndaną}

\begingroup
\setlength{\fboxsep}{6pt}

\noindent\fbox{%
\begin{minipage}{0.97\linewidth}
\small
\begin{tabular}{@{}>{\raggedright\arraybackslash}p{0.359\linewidth}@{\hspace{0.030\linewidth}}>{\raggedright\arraybackslash}p{0.611\linewidth}@{}}
\begin{minipage}[t]{\linewidth}
\centering\textbf{Earlier Germanic changes}\par
\vspace{0.35em}
\raggedright
\centering\textbf{Northwest and West Germanic}\par
\raggedright
\vspace{0.2em}
\raggedright [no change]\par
\vspace{0.6em}
\centering\textbf{Early Anglo-Frisian}\par
\raggedright
\vspace{0.2em}
\raggedright [no change]\par
\end{minipage}
&
\begin{minipage}[t]{\linewidth}
\centering\textbf{Old English changes}\par
\vspace{0.35em}
\raggedright
\begin{tabularx}{\linewidth}{@{}>{\raggedright\arraybackslash}X@{\hspace{0.55em}}>{\raggedright\arraybackslash}p{0.14\linewidth}@{\hspace{0.25em}}}
OE Heavy Syllable Nasal Apocope & \emph{*wíndan} \\
\mbox{OE Secondary Nasalization} & \emph{*wíndąn} \\
\mbox{OE Weak Tail Reduction} & \emph{*wíndan} \\
\end{tabularx}
\end{minipage}
\\
\end{tabular}
\end{minipage}%
} \endgroup

Old English form: \emph{windan}

\subsection{Reconstruction and comparative
evidence}\label{reconstruction-and-comparative-evidence-77}

Kroonen distinguishes noun \emph{*winda-} from verb \emph{*windan-}, and
the present row belongs to the verb (\citeproc{ref-Kroonen2013}{Kroonen
2013}). Fulk and Ringe and Taylor derive the dental directly from PIE
\emph{*wendh-}, not from a Verner alternant
(\citeproc{ref-Fulk2018}{Fulk 2018};
\citeproc{ref-RingeTaylor2014}{Ringe and Taylor 2014}).

\subsection{Old English evidence}\label{old-english-evidence-78}

Clark Hall and Bosworth-Toller record
{\emph{windan}}\index[iv]{02oe@\ivlangheader{Old English}{West Saxon normalization unmarked}!windan@\iventry{windan}{}}
`wind, coil' as the verb headword (\citeproc{ref-ClarkHall1960}{Clark
Hall 1960}; \citeproc{ref-BosworthToller1898}{Bosworth and Toller 1898,
101}). The Old English form here is the ordinary infinitive of the
strong verb.

\subsection{Development to Old
English}\label{development-to-old-english-77}

The derivational input \Recon{wíndaną} `wind' yields the regular
infinitive \emph{windan} by ordinary heavy-syllable apocope and
weak-tail reduction. The form is therefore a straightforward strong-verb
outcome.

\subsection{Lexical note}\label{lexical-note-5}

English \emph{wind} also names the noun. This row represents the
class-III verb, not the noun (\citeproc{ref-Kroonen2013}{Kroonen 2013};
\citeproc{ref-ClarkHall1960}{Clark Hall 1960}).

\section{\texorpdfstring{wold --- OE
\emph{weald}}{wold --- OE weald}}\label{wold-oe-weald}

\index[iv]{02oe@\ivlangheader{Old English}{West Saxon normalization unmarked}!weald@\iventry{weald}{}}
\index[iv]{03pgmc@\ivlangheader{Proto-Germanic}{}!walthuz@\iventry{*wálθuz}{}}

Derivation: \emph{*wálθuz} \textgreater{} \emph{weald} (regular).

\subsection{Derivation trace}\label{derivation-trace-93}

Proto input: \emph{*wálθuz}

\begingroup
\setlength{\fboxsep}{6pt}

\noindent\fbox{%
\begin{minipage}{0.97\linewidth}
\small
\begin{tabular}{@{}>{\raggedright\arraybackslash}p{0.485\linewidth}@{\hspace{0.030\linewidth}}>{\raggedright\arraybackslash}p{0.485\linewidth}@{}}
\begin{minipage}[t]{\linewidth}
\centering\textbf{Earlier Germanic changes}\par
\vspace{0.35em}
\raggedright
\centering\textbf{Northwest and West Germanic}\par
\raggedright
\vspace{0.2em}
\begin{tabularx}{\linewidth}{@{}>{\raggedright\arraybackslash}X@{\hspace{0.55em}}>{\raggedright\arraybackslash}p{0.14\linewidth}@{\hspace{0.25em}}}
\mbox{EAFL Th Voicing} & \emph{*wálduz} \\
\mbox{EAF Final Z Deletion} & \emph{*wáldu} \\
\end{tabularx}
\vspace{0.6em}
\centering\textbf{Early Anglo-Frisian}\par
\raggedright
\vspace{0.2em}
\raggedright [no change]\par
\end{minipage}
&
\begin{minipage}[t]{\linewidth}
\centering\textbf{Old English changes}\par
\vspace{0.35em}
\raggedright
\begin{tabularx}{\linewidth}{@{}>{\raggedright\arraybackslash}X@{\hspace{0.55em}}>{\raggedright\arraybackslash}p{0.14\linewidth}@{\hspace{0.25em}}}
\mbox{EAF Brightening} & \emph{*wældu} \\
\mbox{OE Breaking} & \emph{*wealdu} \\
\mbox{OE High Vowel Apocope} & \emph{*weald} \\
\end{tabularx}
\end{minipage}
\\
\end{tabular}
\end{minipage}%
} \endgroup

Old English form: \emph{weald}

\subsection{Reconstruction and comparative
evidence}\label{reconstruction-and-comparative-evidence-78}

Kroonen reconstructs the noun as \emph{*walþu-} and gives OE
{\emph{weald}}\index[iv]{02oe@\ivlangheader{Old English}{West Saxon normalization unmarked}!weald@\iventry{weald}{}}
`forest, woodland' beside other West Germanic
{\emph{wald}}\index[iv]{02oe@\ivlangheader{Old English}{West Saxon normalization unmarked}!wald@\iventry{wald}{}}
`forest' forms (\citeproc{ref-Kroonen2013}{Kroonen 2013}). The
derivational input \Recon{wálθuz} `wold' is the nominative singular used
for the derivation.

\subsection{Old English evidence}\label{old-english-evidence-79}

Clark Hall makes
{\emph{weald}}\index[iv]{02oe@\ivlangheader{Old English}{West Saxon normalization unmarked}!weald@\iventry{weald}{}}
`forest, woodland' the main noun headword and cross-refers
{\emph{wald}}\index[iv]{02oe@\ivlangheader{Old English}{West Saxon normalization unmarked}!wald~04@\iventry{wald}{Angl.}}
`forest' and
{\emph{wold}}\index[iv]{02oe@\ivlangheader{Old English}{West Saxon normalization unmarked}!wold@\iventry{wold}{}}
`plain, open country' to it (\citeproc{ref-ClarkHall1960}{Clark Hall
1960}). Clark Hall glosses \emph{wald} as \emph{(N,VPs)} ---
Northumbrian and Vespasian Psalter (Mercian) --- making this an
explicitly dialect-labelled Anglian form.

\subsection{Development to Old
English}\label{development-to-old-english-78}

\emph{*lþ} voices to \emph{ld}, Anglo-Frisian brightening yields
\emph{wæld-}, and breaking before the cluster gives \emph{weald-};
apocope then yields
{\emph{weald}}\index[iv]{02oe@\ivlangheader{Old English}{West Saxon normalization unmarked}!weald@\iventry{weald}{}}
`forest, woodland'. The noun is therefore a regular breaking outcome.

\subsection{Dialect note}\label{dialect-note-6}

{\emph{wald}}\index[iv]{02oe@\ivlangheader{Old English}{West Saxon normalization unmarked}!wald~04@\iventry{wald}{Angl.}}
`forest' survives as an Anglian variant in the same family (Clark Hall
\emph{wald (N,VPs)}: Northumbrian and Vespasian Psalter). The Old
English form here is normalized
{\emph{weald}}\index[iv]{02oe@\ivlangheader{Old English}{West Saxon normalization unmarked}!weald@\iventry{weald}{}}
`forest, woodland', not the variant form
(\citeproc{ref-ClarkHall1960}{Clark Hall 1960};
\citeproc{ref-RingeTaylor2014}{Ringe and Taylor 2014}).

\section{\texorpdfstring{world --- OE
\emph{weorold}}{world --- OE weorold}}\label{world-oe-weorold}

\index[iv]{02oe@\ivlangheader{Old English}{West Saxon normalization unmarked}!weorold@\iventry{weorold}{}}
\index[iv]{03pgmc@\ivlangheader{Proto-Germanic}{}!wiraaldiz@\iventry{*wíra-àldiz}{}}
\index[iv]{03pgmc@\ivlangheader{Proto-Germanic}{}!wiraldiz@\iventry{*wír-àldiz}{}}
\index[iv]{03pgmc@\ivlangheader{Proto-Germanic}{}!wiraldiz@\iventry{*wíràldiz}{}}

Derivation: citation reconstruction \emph{*wíra-àldiz}; form followed
here \emph{*wír-àldiz} \textgreater{} \emph{weorold} (regular).

\subsection{Derivation trace}\label{derivation-trace-94}

Proto input: \emph{*wíràldiz}

\begingroup
\setlength{\fboxsep}{6pt}

\noindent\fbox{%
\begin{minipage}{0.97\linewidth}
\small
\begin{tabular}{@{}>{\raggedright\arraybackslash}p{0.485\linewidth}@{\hspace{0.030\linewidth}}>{\raggedright\arraybackslash}p{0.485\linewidth}@{}}
\begin{minipage}[t]{\linewidth}
\centering\textbf{Earlier Germanic changes}\par
\vspace{0.35em}
\raggedright
\centering\textbf{Northwest and West Germanic}\par
\raggedright
\vspace{0.2em}
\begin{tabularx}{\linewidth}{@{}>{\raggedright\arraybackslash}X@{\hspace{0.55em}}>{\raggedright\arraybackslash}p{0.18\linewidth}@{\hspace{0.25em}}}
\mbox{PWGmc Early I Apocope} & \emph{*wíràldz} \\
\mbox{PNWGmc I Lowering} & \emph{*wéràldz} \\
\mbox{EAF Final Z Deletion} & \emph{*wéràld} \\
\end{tabularx}
\vspace{0.6em}
\centering\textbf{Early Anglo-Frisian}\par
\raggedright
\vspace{0.2em}
\raggedright [no change]\par
\end{minipage}
&
\begin{minipage}[t]{\linewidth}
\centering\textbf{Old English changes}\par
\vspace{0.35em}
\raggedright
\begin{tabularx}{\linewidth}{@{}>{\raggedright\arraybackslash}X@{\hspace{0.55em}}>{\raggedright\arraybackslash}p{0.18\linewidth}@{\hspace{0.25em}}}
\mbox{OE Strip Secondary Stress} & \emph{*wérald} \\
\mbox{OE Inter Stress Raising} & \emph{*wéruld} \\
\mbox{OE Back Mutation} & \emph{*wéoruld} \\
\mbox{OE Med Unstressed U Lowering} & \emph{*wéorold} \\
\end{tabularx}
\end{minipage}
\\
\end{tabular}
\end{minipage}%
} \endgroup

Old English form: \emph{weorold}

\subsection{Reconstruction and comparative
evidence}\label{reconstruction-and-comparative-evidence-79}

The word is the old compound `age of men'. Orel reconstructs
\emph{*wira-alđiz}, with the older first-element vowel, while Ringe and
Taylor discuss the form \Recon{weraldiz} `world' and its pre-Old-English
chain \Recon{weraldu} `world' \textgreater{} \Recon{weruld} `world'
(\citeproc{ref-Orel2003}{Orel 2003, 462};
\citeproc{ref-RingeTaylor2014}{Ringe and Taylor 2014, 326}).
Kluge-Seebold gives the compound \Recon{wira-aldō} `world' and
explicitly includes Old English \emph{weorold} `world'
(\citeproc{ref-KlugeSeebold2011}{Kluge 2011, 981}).

Kroonen's separate entries for \emph{*wira-} `man' and \emph{*aldi-}
`age' support the components, not independently the complete compound or
its accentuation. The first vowel is not beyond dispute: Ringe and
Taylor allow \emph{*weraz} beside \emph{*wiraz} for the simplex, since
Gothic pre-\emph{r} developments and Norse lowering obscure the
comparison. Initial \emph{i} is therefore a supported reconstruction,
not something established by its usefulness in the derivation
(\citeproc{ref-Kroonen2013}{Kroonen 2013, 20, 588};
\citeproc{ref-RingeTaylor2014}{Ringe and Taylor 2014, 34}).

The derivational input
\Recon{wír-àldiz}\index[iv]{03pgmc@\ivlangheader{Proto-Germanic}{}!wiraldiz@\iventry{*wír-àldiz}{}}
`world' adopts the inherited compact i-stem nominative. No early ō-stem
intervention is required in its selected phonological history. The
compound's later class mixing remains relevant, rather than being
silently denied because the conservative nominative suffices
(\citeproc{ref-Orel2003}{Orel 2003, 462};
\citeproc{ref-RingeTaylor2014}{Ringe and Taylor 2014, 202, 326}).

The class contrast is real: Ringe and Taylor oppose i-stem Old Saxon
\emph{werold} `world' and Old High German \emph{weralt} `world' to
ō-stem Old Frisian \emph{warld} `world' and Old Norse \emph{verǫld}
`world'. Their separate qualification that the late Norse word is
probably borrowed prevents counting it as independent evidence for a
common inherited ō-stem innovation (\citeproc{ref-RingeTaylor2014}{Ringe
and Taylor 2014, 202, 326}). Nor did every Old English cell necessarily
abandon the older class. Campbell explicitly records frequent endingless
accusative singular \emph{weorold}, characteristic of feminine i-stems.
This is evidence for a conservative ending, not by itself proof of an
entirely unchanged stem history (\citeproc{ref-Campbell1959}{Campbell
1959, 244}, §610(1) n.2).

The selected inherited nominative is sufficient; a complete
reconstruction of its oblique paradigm is not assumed. Different nominal
endings must still be distinguished, rather than replacing them
indiscriminately with \emph{-u} (\citeproc{ref-Fulk2018}{Fulk 2018,
154--56, 161--62}).

No mixed-stage early analogical transponent is now necessary. The
compact stem nevertheless presupposes resolution or reduction of the
compound junction. The etymological spelling \emph{*wira-aldi-} does not
independently establish a literal hiatus at the chosen stage, and
consonant-separated linking-vowel syncope is not a demonstration of how
adjacent \emph{a+a} would have resolved. An already \emph{*wer-}
compound requires no lowering in its own derivation. Accordingly the
derivational input cannot prove that
\hyperref[rule-PNWGmcILowering]{SC015 PNWGmcILowering} historically
preceded \hyperref[rule-OEInterStressRaising]{SC036
OEInterStressRaising}. The adopted consonantal lowering has no
following-vowel trigger for interstress raising to remove. The two
processes remain compatible in this derivation without that supposed
ordering argument (\citeproc{ref-Orel2003}{Orel 2003, 462};
\citeproc{ref-RingeTaylor2014}{Ringe and Taylor 2014, 326};
\citeproc{ref-Hill2017}{Hill 2017, 153--56}).

\subsection{Old English evidence}\label{old-english-evidence-80}

Old English does not preserve a single isolated form. Ringe and Taylor
give West Saxon
{\emph{weorold}}\index[iv]{02oe@\ivlangheader{Old English}{West Saxon normalization unmarked}!weorold@\iventry{weorold}{}}
`world' \textasciitilde{}
{\emph{worold}}\index[iv]{02oe@\ivlangheader{Old English}{West Saxon normalization unmarked}!worold@\iventry{worold}{}}
`world', Mercian
{\emph{weoruld}}\index[iv]{02oe@\ivlangheader{Old English}{West Saxon normalization unmarked}!weoruld~05@\iventry{weoruld}{Merc.}}
`world', Northumbrian
{\emph{woruld}}\index[iv]{02oe@\ivlangheader{Old English}{West Saxon normalization unmarked}!woruld@\iventry{woruld}{}}
`world', and Kentish
{\emph{wiarald}}\index[iv]{02oe@\ivlangheader{Old English}{West Saxon normalization unmarked}!wiarald~07@\iventry{wiarald}{Kent.}}
`world' (\citeproc{ref-RingeTaylor2014}{Ringe and Taylor 2014, 326}).
Sievers-Brunner and Bright present the same wider set, including the
syncopated
{\emph{world}}\index[iv]{02oe@\ivlangheader{Old English}{West Saxon normalization unmarked}!world@\iventry{world}{}}
`world' and later rounded
{\emph{wurold}}\index[iv]{02oe@\ivlangheader{Old English}{West Saxon normalization unmarked}!wurold@\iventry{wurold}{}}
`world' (\citeproc{ref-SieversBrunner1965}{Brunner 1965, § 113};
\citeproc{ref-BrightCassidyRingler1971}{Bright 1971, 465}).

The Old English form used here is the West Saxon form
{\emph{weorold}}\index[iv]{02oe@\ivlangheader{Old English}{West Saxon normalization unmarked}!weorold@\iventry{weorold}{}}
`world'. It is an attested Old English form within that broader variant
cluster, not the only form the lexeme ever shows.

The two sets of vowels carry different evidence. Initial \emph{eo/o}
concerns back mutation and subsequent rounding after \emph{w}; medial
\emph{a/o/u} concerns weak vocalism. Kentish medial \emph{a} does not
show the breaking expected under stronger prominence. Later \emph{world}
demonstrates loss of that weak vowel, not another raising event. The
variants do not require every medial \emph{o} to have passed through
\emph{u} (\citeproc{ref-Campbell1959}{Campbell 1959, 142}, §338;
\citeproc{ref-HoggGrammar2011}{Hogg 2011, 217--18}, §§6.6--6.7;
\citeproc{ref-RingeTaylor2014}{Ringe and Taylor 2014, 326}).

Čermák places \emph{world} in his class of obscured compounds affected
by loss of an unstressed vowel, comparing \emph{woruld/worold}. This is
evidence for later medial syncope, not another instance of the earlier
raising. His possible reinterpretation of the second element is
explicitly a suggestion; the entry supplies no new primary-text locator
(\citeproc{ref-Cermak2005}{Čermák 2005, 36, 39}). Sauer contrasts Old
English feminine genitive \emph{worulde} with later Middle English
genitive-\emph{s} and manuscript variation. These developments help
distinguish inherited inflection from subsequent restructuring, but
cannot date a prehistoric shift into the ō-stems
(\citeproc{ref-Sauer1992}{Sauer 1992, 89--90, 354--56}).

\subsection{Compound prominence and early Old English
footing}\label{compound-prominence-and-early-old-english-footing}

Goering argues that early Old English \emph{weruld} `world' belonged to
a single initial, trimoraic foot, rather than an initial light foot
followed by a secondary-stressed heavy foot. Apocope would occur under
either analysis and cannot decide between them. The ordinary reduction
and eventual loss of the medial vowel, together with resolution of
inflected \emph{worulde} `world' in \emph{Beowulf} 3068a and 2343a,
support his analysis (\citeproc{ref-Goering2023}{Goering 2023, 79--82},
§4.5.1.2).

This does not rule out earlier compound secondary accent. Kaluza's older
theory allowed subsidiary stress in \emph{weorold} and other
disyllables; Stanley's presentation is evidence for that theory, not
direct written accent evidence (\citeproc{ref-Stanley2022}{Stanley 2022,
34}). The adopted relative history loses independent second-element
prominence before weak rounding. Goering's early Old English system,
perhaps around 600 and compatible with the seventh-/eighth-century
poetic horizon, supplies an approximate upper horizon, not a precise
date for this word's deaccentuation (\citeproc{ref-Goering2023}{Goering
2023, 84--85}, §4.6). Foot structure, lexical prominence and metrical
ictus remain distinct. Hogg's discussion warns against making ``tertiary
stress'' an unexplained name for the desired outcome; Minkova likewise
requires independent linguistic controls on metrical reconstruction
(\citeproc{ref-Hogg1996}{Hogg 1996, 3, 9--11};
\citeproc{ref-Minkova1996}{Minkova 1996, 20--29}). The grave in the
selected reconstruction marks the earlier second-root accent, not an
unstressed linking vowel. It survives the linking-syncope operation and
is then removed before weak raising, not asserted as actual independent
prominence at the latter change. True compound secondary stress can
protect a vowel, while obscured compounds reduce; the history of
transparency cannot be inferred from a hyphen alone
(\citeproc{ref-Fulk2018}{Fulk 2018, 36--38}, §2.2).

\subsection{Competing histories of the weak
vowel}\label{competing-histories-of-the-weak-vowel}

The starting point is weak vocalism, not an obligatory analogy. Campbell
describes retraction before \emph{l/r} plus consonant, rounding to
\emph{o}, and further \emph{u} in some forms. Hogg similarly
distinguishes retraction of a weak vowel from the breaking of a more
prominent one. Luick particularly associates rounding with labials or
following \emph{l}, while expressly leaving the connection between the
\emph{o} and \emph{u} developments uncertain
(\citeproc{ref-Campbell1959}{Campbell 1959, 142}, §338;
\citeproc{ref-HoggGrammar2011}{Hogg 2011, 217--18}, §§6.6--6.7;
\citeproc{ref-Luick1914}{Luick 1914-\/-1940, 125--26}, §113). Bülbring's
assimilation before following \emph{u} supplies one account of
\emph{world}, but Luick's criticism and the wider weak-vowel evidence
prevent treating it as the only possible environment
(\citeproc{ref-Bulbring1902}{Bülbring 1902, 147--49}, §§366--367;
\citeproc{ref-Luick1914}{Luick 1914-\/-1940, 126}, §113;
\citeproc{ref-RingeTaylor2014}{Ringe and Taylor 2014, 202}).

CAPR adopts weak short \emph{a > u} before \emph{l} plus consonant in a
light-initial primary-stressed domain: one intervening consonant
separates the initial short nucleus from the medial weak vowel. This
bounded interpretation joins Luick's environment to Goering's LH
analysis; the exact formal restriction is CAPR's qualified
reconstruction, not either author's quoted law. Heavy-initial
retained-prominence compounds such as \emph{ānfeald} `onefold' argue
against unrestricted raising after removal of accent annotations. Weak
\emph{twifold} `twofold' is a relevant rounded comparison; primary
medial stress is not rewritten (\citeproc{ref-Bulbring1902}{Bülbring
1902, 148--49}, §367; \citeproc{ref-Campbell1959}{Campbell 1959, 142},
§338; \citeproc{ref-HoggGrammar2011}{Hogg 2011, 217--18};
\citeproc{ref-Goering2023}{Goering 2023, 79--82}).

The selected operation precedes early unstressed fronting and breaking.
A post-fronting compression of weak \emph{æ} retraction, rounding and
raising also succeeds and resembles Hogg's sequence. The chosen ordering
does not require that intermediate vowel, without claiming an
independently dated universal law. Versloot's alternative conditional
fronting and later breaking challenge compulsory fronting before
\emph{lC}; his principally stressed-domain analysis does not itself
establish weak \emph{a > u} in \emph{world}
(\citeproc{ref-HoggGrammar2011}{Hogg 2011, 217--18};
\citeproc{ref-Versloot2025}{Versloot 2026, 111--14, 122--23}).

\subsection{Selected regular projection to Old
English}\label{selected-regular-projection-to-old-english}

From selected \emph{*wír-àldiz}, early short-\emph{i} apocope,
\hyperref[rule-PNWGmcILowering]{SC015 PNWGmcILowering}, and
final-\emph{z} loss give \emph{*wér-àld}. After linking-syncope has read
the earlier accent information, grave removal gives \emph{*wérald}.
\hyperref[rule-OEInterStressRaising]{SC036 OEInterStressRaising}
supplies \emph{*wéruld}; \hyperref[rule-OEBackMutation]{SC059
OEBackMutation} gives \emph{*wéoruld}; and
\hyperref[rule-OEMedUnstressedULowering]{SC040 OEMedUnstressedULowering}
gives \emph{weorold}. No final high vowel survives to undergo later
apocope.

This is the longest defensible regular sequence from the selected
compact inherited nominative, not a derivation of the unresolved earlier
compound junction. The \emph{regular} classification concerns that
selected span. The inherited accent, original \emph{i/e}, and exact
weak-vowel chronology retain their qualifications. One raising rewrite
may compress gradual rounding and raising; it does not assert universal
\emph{a > u > o} across the dialect variants
(\citeproc{ref-Luick1914}{Luick 1914-\/-1940, 125--26}, §113;
\citeproc{ref-RingeTaylor2014}{Ringe and Taylor 2014, 202, 326}).

An early rounded \emph{o} alone fails in the existing model at
\hyperref[rule-OEFinalUnstressedOLowering]{SC100
OEFinalUnstressedOLowering}, not because historical rounding is
intrinsically impossible. Furthermore, the present
\hyperref[rule-OEBackMutation]{SC059 OEBackMutation} also admits
\emph{o} as a trigger for initial \emph{e}; its executable behaviour
cannot prove that only \emph{u} feeds the diphthong. The adopted
inherited high-vowel path also succeeds under explicitly restricted
\emph{u}-triggered back mutation. No back-mutation or later merger rule
is changed to rescue it.

\subsection{Stage comparison}\label{stage-comparison}

The comparison below sets the relevant forms side by side. It separates
the comparative headword from the OE-facing stage chosen for the
derivation.

{\def\LTcaptype{none} % do not increment counter
\begin{longtable}[]{@{}
  >{\raggedright\arraybackslash}p{(\linewidth - 6\tabcolsep) * \real{0.2500}}
  >{\raggedright\arraybackslash}p{(\linewidth - 6\tabcolsep) * \real{0.2500}}
  >{\raggedright\arraybackslash}p{(\linewidth - 6\tabcolsep) * \real{0.2500}}
  >{\raggedright\arraybackslash}p{(\linewidth - 6\tabcolsep) * \real{0.2500}}@{}}
\toprule\noalign{}
\begin{minipage}[b]{\linewidth}\raggedright
Stage / interpretation
\end{minipage} & \begin{minipage}[b]{\linewidth}\raggedright
Candidate form
\end{minipage} & \begin{minipage}[b]{\linewidth}\raggedright
Old English outcome or comparison
\end{minipage} & \begin{minipage}[b]{\linewidth}\raggedright
Relevance to this entry
\end{minipage} \\
\midrule\noalign{}
\endhead
\bottomrule\noalign{}
\endlastfoot
comparative compound with older first-element vowel &
{\emph{*wíra-àldiz}}\index[iv]{03pgmc@\ivlangheader{Proto-Germanic}{}!wiraaldiz@\iventry{*wíra-àldiz}{}}
& citation reconstruction / lexeme label & preserves the older
\emph{*wir-} tradition of the compound \\
literature-stage lowered compound after early stem-class shift &
\Recon{weraldiz}\index[iv]{08preoe@\ivlangheader{Pre-Old English / prehistoric English}{}!weraldiz@\iventry{*weraldiz}{}}
\textgreater{}
\Recon{weraldu}\index[iv]{08preoe@\ivlangheader{Pre-Old English / prehistoric English}{}!weraldu@\iventry{*weraldu}{}}
\textgreater{}
\Recon{weruld}\index[iv]{08preoe@\ivlangheader{Pre-Old English / prehistoric English}{}!weruld@\iventry{*weruld}{}}
& Ringe-Taylor background chain to OE
{\emph{weorold}}\index[iv]{02oe@\ivlangheader{Old English}{West Saxon normalization unmarked}!weorold@\iventry{weorold}{}}
`world' \textasciitilde{}
{\emph{worold}}\index[iv]{02oe@\ivlangheader{Old English}{West Saxon normalization unmarked}!worold@\iventry{worold}{}}
`world' & explains the older comparative literature cited for the
word \\
selected inherited compact nominative &
{\emph{*wír-àldiz}}\index[iv]{03pgmc@\ivlangheader{Proto-Germanic}{}!wiraldiz@\iventry{*wír-àldiz}{}}
& regular output: {\emph{weorold}} `world' & exact West Saxon match;
junction pre-encoded, no necessary early analogy \\
broader OE variant cluster & --- &
{\emph{worold}}\index[iv]{02oe@\ivlangheader{Old English}{West Saxon normalization unmarked}!worold@\iventry{worold}{}}
`world',
{\emph{weoruld}}\index[iv]{02oe@\ivlangheader{Old English}{West Saxon normalization unmarked}!weoruld@\iventry{weoruld}{}}
`world',
{\emph{woruld}}\index[iv]{02oe@\ivlangheader{Old English}{West Saxon normalization unmarked}!woruld@\iventry{woruld}{}}
`world',
{\emph{wiarald}}\index[iv]{02oe@\ivlangheader{Old English}{West Saxon normalization unmarked}!wiarald@\iventry{wiarald}{}}
`world',
{\emph{world}}\index[iv]{02oe@\ivlangheader{Old English}{West Saxon normalization unmarked}!world@\iventry{world}{}}
`world' & real attested comparanda that remain outside that West Saxon
line \\
\end{longtable}
}

\section{\texorpdfstring{yarn --- OE
\emph{ġearn}}{yarn --- OE ġearn}}\label{yarn-oe-ux121earn}

\index[iv]{02oe@\ivlangheader{Old English}{West Saxon normalization unmarked}!gearn@\iventry{ġearn}{}}
\index[iv]{03pgmc@\ivlangheader{Proto-Germanic}{}!garna@\iventry{*gárną}{}}

Derivation: \emph{*gárną} \textgreater{} \emph{ġearn} (regular).

\subsection{Derivation trace}\label{derivation-trace-95}

Proto input: \emph{*gárną}

\begingroup
\setlength{\fboxsep}{6pt}

\noindent\fbox{%
\begin{minipage}{0.97\linewidth}
\small
\begin{tabular}{@{}>{\raggedright\arraybackslash}p{0.359\linewidth}@{\hspace{0.030\linewidth}}>{\raggedright\arraybackslash}p{0.611\linewidth}@{}}
\begin{minipage}[t]{\linewidth}
\centering\textbf{Earlier Germanic changes}\par
\vspace{0.35em}
\raggedright
\centering\textbf{Northwest and West Germanic}\par
\raggedright
\vspace{0.2em}
\raggedright [no change]\par
\vspace{0.6em}
\centering\textbf{Early Anglo-Frisian}\par
\raggedright
\vspace{0.2em}
\raggedright [no change]\par
\end{minipage}
&
\begin{minipage}[t]{\linewidth}
\centering\textbf{Old English changes}\par
\vspace{0.35em}
\raggedright
\begin{tabularx}{\linewidth}{@{}>{\raggedright\arraybackslash}X@{\hspace{0.55em}}>{\raggedright\arraybackslash}p{0.14\linewidth}@{\hspace{0.25em}}}
\mbox{EAF Brightening} & \emph{*gærną} \\
\mbox{OE Breaking} & \emph{*gearną} \\
OE Heavy Syllable Nasal Apocope & \emph{*gearn} \\
\mbox{OE Velar Palatalization} & \emph{*ʝearn} \\
\mbox{OE Palatal Fricative Merger} & \emph{*jearn} \\
\end{tabularx}
\end{minipage}
\\
\end{tabular}
\end{minipage}%
} \endgroup

Old English form: \emph{ġearn}

\subsection{Reconstruction and comparative
evidence}\label{reconstruction-and-comparative-evidence-80}

Kroonen cites the noun as \emph{*garna-}, and Ringe and Taylor give the
early chain \Recon{garna} `yarn' \textgreater{} \Recon{geern} `yarn'
\textgreater{} \Recon{gearn} `yarn' \textgreater{} OE
{\emph{gearn}}\index[iv]{02oe@\ivlangheader{Old English}{West Saxon normalization unmarked}!gearn@\iventry{gearn}{}}
`yarn' (\citeproc{ref-Kroonen2013}{Kroonen 2013};
\citeproc{ref-RingeTaylor2014}{Ringe and Taylor 2014, 220}). The
derivational input \Recon{gárną} `yarn' is the nominal citation form
used here, while oblique \Recon{garnăn} `yarn' belongs only to
comparative background.

\subsection{Old English evidence}\label{old-english-evidence-81}

Clark Hall records \emph{gearn (e) n.} `yarn, spun wool', and
Bosworth-Toller glosses
{\emph{gearn}}\index[iv]{02oe@\ivlangheader{Old English}{West Saxon normalization unmarked}!gearn@\iventry{gearn}{}}
`yarn' with Latin \emph{filatum} `spun thread'
(\citeproc{ref-ClarkHall1960}{Clark Hall 1960};
\citeproc{ref-BosworthToller1898}{Bosworth and Toller 1898}).

\subsection{Development to Old
English}\label{development-to-old-english-79}

From \Recon{gárną} `yarn', brightening and breaking before \emph{rn}
yield \emph{gearn} `yarn'; palatalization of initial \emph{g} before the
resulting front-vocalic sequence gives normalized \emph{ġearn} `yarn'.
The derivation is regular.

\subsection{Form note}\label{form-note-26}

Dictionary \emph{gearn} `yarn' and normalized \emph{ġearn} `yarn' refer
to the same noun. The comparative stem \emph{*garna-} and oblique
\Recon{garnăn} `yarn' do not replace the derivational input
\Recon{gárną} `yarn'.

\section{\texorpdfstring{you --- OE
\emph{ēow}}{you --- OE ēow}}\label{you-oe-ux113ow}

\index[iv]{02oe@\ivlangheader{Old English}{West Saxon normalization unmarked}!eow@\iventry{ēow}{}}
\index[iv]{03pgmc@\ivlangheader{Proto-Germanic}{}!izwiz@\iventry{*ízwiz}{}}

Derivation: \emph{*ízwiz} \textgreater{} \emph{ēow} (regular).

\subsection{Derivation trace}\label{derivation-trace-96}

Proto input: \emph{*ízwiz}

\begingroup
\setlength{\fboxsep}{6pt}

\noindent\fbox{%
\begin{minipage}{0.97\linewidth}
\small
\begin{tabular}{@{}>{\raggedright\arraybackslash}p{0.611\linewidth}@{\hspace{0.030\linewidth}}>{\raggedright\arraybackslash}p{0.359\linewidth}@{}}
\begin{minipage}[t]{\linewidth}
\centering\textbf{Earlier Germanic changes}\par
\vspace{0.35em}
\raggedright
\centering\textbf{Northwest and West Germanic}\par
\raggedright
\vspace{0.2em}
\begin{tabularx}{\linewidth}{@{}>{\raggedright\arraybackslash}X@{\hspace{0.55em}}>{\raggedright\arraybackslash}p{0.14\linewidth}@{\hspace{0.25em}}}
\mbox{PWGmc Coronal W Assimilation} & \emph{*íwwiz} \\
\mbox{OE WW Simplification} & \emph{*íuwiz} \\
\mbox{EAF Final Z Deletion} & \emph{*íuwi} \\
PWGmc Unstressed Word Final I Apocope & \emph{*íuw} \\
\end{tabularx}
\vspace{0.6em}
\centering\textbf{Early Anglo-Frisian}\par
\raggedright
\vspace{0.2em}
\raggedright [no change]\par
\end{minipage}
&
\begin{minipage}[t]{\linewidth}
\centering\textbf{Old English changes}\par
\vspace{0.35em}
\raggedright
\begin{tabularx}{\linewidth}{@{}>{\raggedright\arraybackslash}X@{\hspace{0.55em}}>{\raggedright\arraybackslash}p{0.14\linewidth}@{\hspace{0.25em}}}
\mbox{OE Diphthong Leveling} & \emph{*ēow} \\
\end{tabularx}
\end{minipage}
\\
\end{tabular}
\end{minipage}%
} \endgroup

Old English form: \emph{ēow}

\subsection{Reconstruction and comparative
evidence}\label{reconstruction-and-comparative-evidence-81}

The derivational input is the dative(-accusative) plural of the
second-person plural pronoun, PGmc
\Recon{ízwiz}\index[iv]{03pgmc@\ivlangheader{Proto-Germanic}{}!izwiz@\iventry{*ízwiz}{}}
`you (dat. pl.)' (Gothic
{izwis}\index[iv]{12goth@\ivlangheader{Gothic}{}!izwis@\iventry{izwis}{}}
`you'). Ringe and Taylor print the development in full: ``PGmc *izwiz
`you (dat. pl.)' (Goth. izwis) \textgreater{} *iwwi \textgreater{} PWGmc
*iuwi \textasciitilde{} *iuw (see 3.1.4) \textgreater{} OE īow, OF iū,
OS, OHG iu'' (\citeproc{ref-RingeTaylor2014}{Ringe and Taylor 2014,
41--42}), following Stiles's demonstration of the coronal-w assimilation
(\citeproc{ref-Stiles1985}{Stiles 1985, 89--94}). Fulk's discussion of
inherited geminate glides after short vowels supplies the wider
reanalysis context, not a replacement citation for that assimilation
(\citeproc{ref-Fulk2018}{Fulk 2018, 117}). The oblique plural is the
cell that survives as the Old English pronoun \emph{ēow}.

\subsection{Competing origins of the inherited
oblique}\label{competing-origins-of-the-inherited-oblique}

Kroonen derives \emph{*izwiz} by metathesis from \emph{*iw-siz}, with a
suffix abstracted from the first-person plural, and separately
distinguishes older \emph{*iwiz} and the dual formations
(\citeproc{ref-Kroonen2013}{Kroonen 2013, 275--76}). Fulk regards the
reported Krahe--Meid reduplication \emph{*wiz-wiz} and initial-\emph{w}
dissimilation as probably most plausible, while explicitly stating that
there is no consensus (\citeproc{ref-Fulk2018}{Fulk 2018, 187--88}).
This is his assessment of a reported theory, not a newly verified
original argument or a named rebuttal of Kroonen.

Ringe's Katz-based account instead combines case formation and sandhi
history from Indo-European obliques including \emph{*uswé}, with later
prothetic \emph{i}; he also records Weiss's dissimilatory alternative
(\citeproc{ref-Ringe2017}{Ringe 2017, 233--35}). The competing proposals
concern the genealogy and segmentation of the same relevant plural
oblique. Agreement on a later inherited input does not choose between
metathesis and suffix abstraction, reduplication and dissimilation, or
sandhi and case formation. The regular English path and its
independently selected weak-final context cannot settle that earlier
morphology.

\subsection{Old English evidence}\label{old-english-evidence-82}

West Saxon
{\emph{ēow}}\index[iv]{02oe@\ivlangheader{Old English}{West Saxon normalization unmarked}!eow@\iventry{ēow}{}}
`you' beside early West Saxon and Northumbrian
{\emph{īow}}\index[iv]{02oe@\ivlangheader{Old English}{West Saxon normalization unmarked}!iow@\iventry{īow}{}}
`you' shows the normal West Saxon treatment of the diphthong
(\citeproc{ref-Campbell1959}{Campbell 1959, § 702}, p.~283). The
corpus targets the West Saxon citation form.

\subsection{Development to Old
English}\label{development-to-old-english-80}

Coronal-w assimilation first supplies the geminate in \emph{*ízwiz}
\textgreater{} \emph{*íwwiz}. Early inherited-glide reanalysis then
gives an \emph{*iu} nucleus with retained consonantal \emph{*w}, rather
than an already completed English long diphthong. Final-z loss exposes
the last high vowel, and weak-final apocope selects the apocopated
member of Ringe and Taylor's PWGmc \emph{*iuwi} / \emph{*iuw} doublet.
The later English realization gives \emph{ēow} without an i-mutation
trigger (\citeproc{ref-RingeTaylor2014}{Ringe and Taylor 2014, 41--42,
44--45, 57--58, 171--75}).

The selected sentence context is weak and phonologically final at early
apocope. It is recorded independently of the unchanged PGmc
reconstruction and independently of the root accent: the acute is not a
sentence-stress mark. Under a retained-i context the present model
produces predicted \Pred{īei}, because w-loss before i intervenes before
mutation. That is a computational negative control, not an attested
strong Old English spelling. Proclitic \emph{ymbe} `around', weak-final
\emph{and} `and', and stressed \emph{ġiest} `guest' / \emph{fȳr} `fire'
test the boundary and stress contrasts; no pronoun-specific or lexeme-ID
sound-law condition is used (\citeproc{ref-RingeTaylor2014}{Ringe and
Taylor 2014, 55, 57--58}).

\subsection{Why this row is in the
corpus}\label{why-this-row-is-in-the-corpus-3}

The word supplies the assimilation input to
\hyperref[rule-OEWWSimplification]{SC031 OEWWSimplification} and the
selected weak-final context for
\hyperref[rule-PWGmcUnstressedWordFinalIApocope]{SC098
PWGmcUnstressedWordFinalIApocope}. It demonstrates loss of the mutation
trigger, not a late English vocalization-before-degemination order. That
older interpretation belonged to the superseded ww proxy. The
independently justified assimilation feeder and the prosodic condition
now remain separate (\citeproc{ref-RingeTaylor2014}{Ringe and Taylor
2014, 41--42, 55, 57--58}).

\clearpage

\chapter{Attested variants and comparison
forms}\label{attested-variants-and-comparison-forms}

Here the Old English comparator belongs to a documented set of variants.
The variation is part of the evidence and not an inconvenience to be
normalized away.

\section{\texorpdfstring{cud --- OE
\emph{cwedu}}{cud --- OE cwedu}}\label{cud-oe-cwedu}

\index[iv]{02oe@\ivlangheader{Old English}{West Saxon normalization unmarked}!cwedu@\iventry{cwedu}{}}
\index[iv]{03pgmc@\ivlangheader{Proto-Germanic}{}!kweduz@\iventry{*kwéðuz}{}}
\index[iv]{03pgmc@\ivlangheader{Proto-Germanic}{}!kwithuz@\iventry{*kwíθuz}{}}

Derivation: citation reconstruction \emph{*kwíθuz}; form followed here
\emph{*kwéðuz} \textgreater{} \emph{cwedu} (attested variant).

\subsection{Derivation trace}\label{derivation-trace-97}

Proto input: \emph{*kwéðuz}

\begingroup
\setlength{\fboxsep}{6pt}

\noindent\fbox{%
\begin{minipage}{0.97\linewidth}
\small
\begin{tabular}{@{}>{\raggedright\arraybackslash}p{0.582\linewidth}@{\hspace{0.030\linewidth}}>{\raggedright\arraybackslash}p{0.388\linewidth}@{}}
\begin{minipage}[t]{\linewidth}
\centering\textbf{Earlier Germanic changes}\par
\vspace{0.35em}
\raggedright
\centering\textbf{Northwest and West Germanic}\par
\raggedright
\vspace{0.2em}
\begin{tabularx}{\linewidth}{@{}>{\raggedright\arraybackslash}X@{\hspace{0.55em}}>{\raggedright\arraybackslash}p{0.14\linewidth}@{\hspace{0.25em}}}
\mbox{PWGmc Dental Hardening} & \emph{*kwéduz} \\
\mbox{EAF Final Z Deletion} & \emph{*kwédu} \\
\end{tabularx}
\vspace{0.6em}
\centering\textbf{Early Anglo-Frisian}\par
\raggedright
\vspace{0.2em}
\raggedright [no change]\par
\end{minipage}
&
\begin{minipage}[t]{\linewidth}
\centering\textbf{Old English changes}\par
\vspace{0.35em}
\raggedright
\raggedright [no change]\par
\end{minipage}
\\
\end{tabular}
\end{minipage}%
} \endgroup

Old English form: \emph{cwedu}

\subsection{Reconstruction and comparative
evidence}\label{reconstruction-and-comparative-evidence-82}

Kroonen gives the resin word under homonym 2 of \Recon{kwedu-} `cud' and
supplies Old English variants
{\emph{cwidu}}\index[iv]{02oe@\ivlangheader{Old English}{West Saxon normalization unmarked}!cwidu@\iventry{cwidu}{}}
`cud',
{\emph{cweodu}}\index[iv]{02oe@\ivlangheader{Old English}{West Saxon normalization unmarked}!cweodu@\iventry{cweodu}{}}
`cud', and
{\emph{c(w)udu}}\index[iv]{02oe@\ivlangheader{Old English}{West Saxon normalization unmarked}!cwudu@\iventry{c(w)udu}{}}
`cud' (\citeproc{ref-Kroonen2013}{Kroonen 2013, 315}). The homonym
number is not part of the reconstructed phonological stem. Orel also
lists
{\emph{cwidu}}\index[iv]{02oe@\ivlangheader{Old English}{West Saxon normalization unmarked}!cwidu@\iventry{cwidu}{}},
but his feminine headword is a different formation, not another
quotation of Kroonen's masculine u-stem (\citeproc{ref-Orel2003}{Orel
2003, 227}). The derivational input
\Recon{kwéðuz}\index[iv]{03pgmc@\ivlangheader{Proto-Germanic}{}!kweduz@\iventry{*kwéðuz}{}}
`cud' supplies an e-vowel and voiced dental for the chosen comparison
{\emph{cwedu}}\index[iv]{02oe@\ivlangheader{Old English}{West Saxon normalization unmarked}!cwedu@\iventry{cwedu}{}}
`cud'. It is a CAPR whole-word selection, not a diplomatic quotation of
either dictionary's complete form. Its relationship to the different
citation reconstruction above remains under comparative review.

\subsection{Old English evidence}\label{old-english-evidence-83}

The Old English word survives in a wider variant set than one dictionary
headword suggests. Ringe and Taylor discuss
{\emph{cwidu}}\index[iv]{02oe@\ivlangheader{Old English}{West Saxon normalization unmarked}!cwidu@\iventry{cwidu}{}}
`cud' \textgreater{}
{\emph{cwudu}}\index[iv]{02oe@\ivlangheader{Old English}{West Saxon normalization unmarked}!cwudu@\iventry{cwudu}{}}
`cud' \textgreater{}
{\emph{cudu}}\index[iv]{02oe@\ivlangheader{Old English}{West Saxon normalization unmarked}!cudu@\iventry{cudu}{}}
`cud' and also note late West Saxon
{\emph{cweodu}}\index[iv]{02oe@\ivlangheader{Old English}{West Saxon normalization unmarked}!cweodu~03@\iventry{cweodu}{LWS}}
`cud'. Their starting reconstruction is explicitly West Germanic with i,
not an explicitly printed Proto-Germanic form in this passage
(\citeproc{ref-RingeTaylor2014}{Ringe and Taylor 2014, 323}). Clark
Hall's headword
{\emph{cwudu}}\index[iv]{02oe@\ivlangheader{Old English}{West Saxon normalization unmarked}!cwudu@\iventry{cwudu}{}}
`cud' admits eo and i variants, but does not itself attest the selected
bare-e spelling (\citeproc{ref-ClarkHall1960}{Clark Hall 1960, 69}). The
claim that CAPR's comparison is the conservative attested variant needs
further source verification; these passages do not establish it.

\subsection{Development to Old
English}\label{development-to-old-english-81}

The current modeled chain takes
\Recon{kwéðuz}\index[iv]{03pgmc@\ivlangheader{Proto-Germanic}{}!kweduz@\iventry{*kwéðuz}{}}
`cud' to
{\emph{cwedu}}\index[iv]{02oe@\ivlangheader{Old English}{West Saxon normalization unmarked}!cwedu@\iventry{cwedu}{}}
`cud', with the voiced dental hardened along the way. This records the
executable result, not independent proof that the derivational input is
the best source-backed history. Ringe and Taylor support back umlaut and
further reduction within the i-starting variant family; their passage
does not establish the older account's claim that the i-form is leveled
from the bare-e comparison (\citeproc{ref-RingeTaylor2014}{Ringe and
Taylor 2014, 323}). Corpus fields and classification remain unchanged
pending the complete source audit.

\subsection{Variant comparison}\label{variant-comparison}

{\def\LTcaptype{none} % do not increment counter
\begin{longtable}[]{@{}
  >{\raggedright\arraybackslash}p{(\linewidth - 4\tabcolsep) * \real{0.3333}}
  >{\raggedright\arraybackslash}p{(\linewidth - 4\tabcolsep) * \real{0.3333}}
  >{\raggedright\arraybackslash}p{(\linewidth - 4\tabcolsep) * \real{0.3333}}@{}}
\toprule\noalign{}
\begin{minipage}[b]{\linewidth}\raggedright
Variant type
\end{minipage} & \begin{minipage}[b]{\linewidth}\raggedright
Old English form
\end{minipage} & \begin{minipage}[b]{\linewidth}\raggedright
Comment
\end{minipage} \\
\midrule\noalign{}
\endhead
\bottomrule\noalign{}
\endlastfoot
selected comparison &
{\emph{cwedu}}\index[iv]{02oe@\ivlangheader{Old English}{West Saxon normalization unmarked}!cwedu@\iventry{cwedu}{}}
& current target; its attestation and conservative status require
verification \\
i-grade comparison &
{\emph{cwidu}}\index[iv]{02oe@\ivlangheader{Old English}{West Saxon normalization unmarked}!cwidu@\iventry{cwidu}{}}
& source-backed variant; leveling from the Old English form here is not
established by the cited passage \\
back-umlauted forms &
{\emph{cweodu}}\index[iv]{02oe@\ivlangheader{Old English}{West Saxon normalization unmarked}!cweodu@\iventry{cweodu}{}},
{\emph{cwudu}}\index[iv]{02oe@\ivlangheader{Old English}{West Saxon normalization unmarked}!cwudu@\iventry{cwudu}{}}
& later developments within the same OE tradition \\
reduced form &
{\emph{cudu}}\index[iv]{02oe@\ivlangheader{Old English}{West Saxon normalization unmarked}!cudu@\iventry{cudu}{}}
& further reduced member of the same variant set \\
\end{longtable}
}

\section{\texorpdfstring{earthly --- OE
\emph{eorþcunde}}{earthly --- OE eorþcunde}}\label{earthly-oe-eoruxfecunde}

\index[iv]{02oe@\ivlangheader{Old English}{West Saxon normalization unmarked}!eorthcunde@\iventry{eorþcunde}{}}
\index[iv]{03pgmc@\ivlangheader{Proto-Germanic}{}!erthakundaz@\iventry{*érθa-kùnðaz}{}}
\index[iv]{03pgmc@\ivlangheader{Proto-Germanic}{}!erthakundo@\iventry{*érθa-kùnðǭ}{}}
\index[iv]{03pgmc@\ivlangheader{Proto-Germanic}{}!erthakundo@\iventry{*érθakùnðǭ}{}}

Derivation: citation reconstruction \emph{*érθa-kùnðaz}; form followed
here \emph{*érθa-kùnðǭ} \textgreater{} \emph{eorþcunde} (attested
variant).

\subsection{Derivation trace}\label{derivation-trace-98}

Proto input: \emph{*érθakùnðǭ}

\begingroup
\setlength{\fboxsep}{6pt}

\noindent\fbox{%
\begin{minipage}{0.97\linewidth}
\footnotesize
\begin{tabular}{@{}>{\raggedright\arraybackslash}p{0.485\linewidth}@{\hspace{0.030\linewidth}}>{\raggedright\arraybackslash}p{0.485\linewidth}@{}}
\begin{minipage}[t]{\linewidth}
\centering\textbf{Earlier Germanic changes}\par
\vspace{0.35em}
\raggedright
\centering\textbf{Northwest and West Germanic}\par
\raggedright
\vspace{0.2em}
\begin{tabularx}{\linewidth}{@{}>{\raggedright\arraybackslash}X@{\hspace{0.55em}}>{\raggedright\arraybackslash}p{0.22\linewidth}@{\hspace{0.25em}}}
\mbox{PWGmc Dental Hardening} & \emph{*érθakùndǭ} \\
\end{tabularx}
\vspace{0.6em}
\centering\textbf{Early Anglo-Frisian}\par
\raggedright
\vspace{0.2em}
\raggedright [no change]\par
\end{minipage}
&
\begin{minipage}[t]{\linewidth}
\centering\textbf{Old English changes}\par
\vspace{0.35em}
\raggedright
\begin{tabularx}{\linewidth}{@{}>{\raggedright\arraybackslash}X@{\hspace{0.55em}}>{\raggedright\arraybackslash}p{0.26\linewidth}@{\hspace{0.25em}}}
\mbox{EAF Brightening Unstressed} & \emph{*érθækùndǭ} \\
\mbox{OE Breaking} & \emph{*éorθækùndǭ} \\
\mbox{OE A Restoration} & \emph{*éorθakùndǭ} \\
\mbox{OE Non High Syn} & \emph{*éorθkùndǭ} \\
\mbox{OE High Stress Decode} & \emph{*éorθkundǭ} \\
OE Unstressed Long Vowel Shortening & \emph{*éorθkundæ} \\
\mbox{OE Unstressed AE Merger} & \emph{*éorθkunde} \\
\end{tabularx}
\end{minipage}
\\
\end{tabular}
\end{minipage}%
} \endgroup

Old English form: \emph{eorþcunde}

\subsection{Selected form and documentary
status}\label{selected-form-and-documentary-status}

{\def\LTcaptype{none} % do not increment counter
\begin{longtable}[]{@{}
  >{\raggedright\arraybackslash}p{(\linewidth - 2\tabcolsep) * \real{0.5000}}
  >{\raggedright\arraybackslash}p{(\linewidth - 2\tabcolsep) * \real{0.5000}}@{}}
\toprule\noalign{}
\begin{minipage}[b]{\linewidth}\raggedright
Item
\end{minipage} & \begin{minipage}[b]{\linewidth}\raggedright
Value
\end{minipage} \\
\midrule\noalign{}
\endhead
\bottomrule\noalign{}
\endlastfoot
comparative citation in CAPR notation & \emph{*érθa-kùnðaz} \\
selected inherited weak neuter accusative singular &
\emph{*érθa-kùnðǭ} \\
Old English target & \emph{eorþcunde} `earthly' \\
documentary classification & attested\_variant \\
adopted regular span under medial-foot protection & \emph{*érθa-kùnðǭ} →
\emph{eorþcunde} \\
\end{longtable}
}

The Old English form here is an attested inflected adjective, not the
dictionary headword \emph{eorþcund} `earthly'. Its documentary
classification does not assert an irregular sound change, a dialectal
exception or analogical levelling. The citation nominative and the
selected weak neuter accusative are different cells reconstructed at the
same comparative stage.

\subsection{Comparative formation and the limits of the
citation}\label{comparative-formation-and-the-limits-of-the-citation}

Orel gives the whole formation \emph{*erþa-kunđaz}, with Old English
\emph{eorþ-cund} and Gothic \emph{airþa-kunds} `earthly'
(\citeproc{ref-Orel2003}{Orel 2003, 85}). This is comparative evidence
for the compound, not a printed reconstruction of its accented
weak-neuter cell. CAPR's \emph{θ} and \emph{ð} serialize the source's
dental signs; its acute, grave and selected ending are analytical
additions, not diplomatic transcription. The Gothic comparison supports
a formation extending beyond Old English, without independently proving
the date or uninterrupted inheritance of every Old English inflection.

Kroonen independently gives \emph{*kunda-} `born', with Gothic
\emph{-kunds}, Norse \emph{-kundr} and Old English \emph{-cund}, as the
old participle of the root `be born'. His \emph{*erþō-} `earth' supplies
the lexical first constituent, not the compound's linking vowel. The
separate \emph{*kunþa-} `known', associated with `know', is not the same
formation. The bounded consultation recovered component evidence, not an
authored whole earthly compound or a weak-neuter ending
(\citeproc{ref-Kroonen2013}{Kroonen 2013, 118--19, 310--12}).

The second constituent's subsequent suffixal behavior matters. Hogg
includes \emph{-cund} among elements historically derived from compound
constituents, distinguishing older compound stress from later
subordinate stress in disyllabic formations or inflected forms
(\citeproc{ref-HoggGrammar2011}{Hogg 2011, 49--50}, §2.89). Neither a
dictionary hyphen nor the spelling \emph{-cund} independently
establishes prehistoric prominence. No unsourced derivation from the
verb `know' is needed here.

\subsection{Three located passages, not a complete six-token
concordance}\label{three-located-passages-not-a-complete-six-token-concordance}

Ciszek-Kiliszewska's published DOEC survey reports six occurrences of
\emph{eorþcund}, against 761 of \emph{godcund} `divine', and includes
inflections in its wildcard searches. The lexeme-level table is not a
six-item concordance
(\citeproc{ref-CiszekKiliszewska2015}{Ciszek-Kiliszewska 2015,
489--90}). The bounded documentary record locates three distinct
passages: Genesis 1627, An. Ox. 691 and R. Ben. 139, 27--28. It does not
identify the remaining occurrences or explain the discrepancy by
counting unrelated forms. Two dictionary quotations of the Benedict
sentence represent one occurrence, not two.

\subsection{The Genesis object and its weak neuter
agreement}\label{the-genesis-object-and-its-weak-neuter-agreement}

Bosworth--Toller quotes ``Se rinc ageaf eorþ-cunde eád'', with locators
Cd. 79, Th. 98, 8 and Gen.~1627. Its separate \emph{eád} `wealth' entry
explicitly marks the noun as neuter. In context \emph{eád} is the object
of \emph{ageaf} `gave up'; \emph{eorþcunde} agrees with it as a weak
neuter accusative singular (\citeproc{ref-BosworthToller1898}{Bosworth
and Toller 1898}, s.vv.~eorþ-cund, eád).

The passage has also been located in the OCR of Thorpe's 1832 edition on
the page numbered 98. That edition-OCR consultation supports the
sentence context and absence of a determiner. It is not a collation of
Junius 11, Thorpe's page image or Krapp's critical edition. Apparent
eth/thorn differences and hyphenation in OCR must not be described as
verified printed or manuscript spellings. The target \emph{eorþcunde} is
a normalized scholarly representation of the dictionary-located form. No
exact manuscript date or narrower dialect is inferred from it.

\subsection{Weak inflection without a
determiner}\label{weak-inflection-without-a-determiner}

The absence of a demonstrative or possessive does not make the weak
ending an anomaly requiring repair. Campbell describes the freer use of
weak adjectives in poetry, and Fulk independently recognizes the poetic
use without the usual specifier. Their weak-adjective paradigms give
neuter nominative--accusative singular \emph{-e}, including the ordinary
comparison \emph{blinde} `blind' (\citeproc{ref-Campbell1959}{Campbell
1959, 261, 272}, §§638, 656; \citeproc{ref-Fulk2018}{Fulk 2018,
216--17}, §§9.7--9.8).

Applying those general grammatical statements to this sentence is the
report's syntactic inference. The adjective need not be reinterpreted as
a masculine accusative with an inexplicably lost \emph{n}. Nor is weak
inflection here evidence for an otherwise unrecorded levelling event.

\subsection{The inherited weak-neuter
ending}\label{the-inherited-weak-neuter-ending}

Fulk explicitly treats weak adjectives as n-stems. His n-stem discussion
and paradigms supply the relevant neuter singular ending, \emph{*-ōⁿ},
for the West Germanic and Norse comparison (\citeproc{ref-Fulk2018}{Fulk
2018, 168} n.~1, pp.~169--170, §§7.29--7.31). The macron denotes a
bimoric long vowel; superscript \emph{ⁿ} marks nasalization, not a
consonant to be added to the string. CAPR therefore represents this
ending as \emph{*-ǭ}, not trimoric \emph{*-ôⁿ} or literal consonantal
\emph{*-ōn}.

The ending belongs on the unaugmented compound stem. Fulk's neuter
\emph{augō} `eye' paradigm contrasts singular forms without the oblique
\emph{n} formative with plural forms containing it. The class name
``n-stem'' is not permission to insert \emph{*-an-} before the selected
singular ending. CAPR consequently assembles \emph{*erþa-kunđ-ōⁿ},
serialized \emph{*érθa-kùnðǭ}. This whole weak adjective is a
comparative inference from Fulk's class statement and noun paradigm, not
a whole word printed by him. Campbell independently supports the Old
English ending and syntax; he has not been used as an independent check
of that entire Proto-Germanic assembly.

\subsection{Inheritance and the reported origin
alternative}\label{inheritance-and-the-reported-origin-alternative}

Fulk connects the neuter ending with inherited PIE
\emph{*-ōn < *-on-h₂}, citing Jasanoff. He also reports Boutkan's
alternative: analogical extension from the neuter plural
(\citeproc{ref-Fulk2018}{Fulk 2018, 170}, citing Jasanoff 2002, p.~33,
and Boutkan 1995b, pp.~285--286). These are alternative accounts of the
ending's prehistory, not different readings of the OE-relevant macron.
Neither cited work has been independently consulted for this entry.

The report prefers Fulk's regular inherited account, without claiming
that successful execution decides between it and the reported analogy.
Fulk's separate discussion of a perplexing masculine trimoric origin
must not be transferred to this bimoric neuter cell. No
trimoric-to-bimoric repair or additional ending law is adopted.

\subsection{Medial prominence after an initial heavy
foot}\label{medial-prominence-after-an-initial-heavy-foot}

The adopted protection concerns phonological weight and domain, not
membership in \emph{-cund} or the weak-adjective paradigm. In
\emph{eorþ.cun.de}, the initial syllable is heavy and the following
nonfinal closed syllable can head a medial foot. Goering distinguishes
this configuration from an initial light syllable accommodated together
with a following heavy syllable in the initial foot
(\citeproc{ref-Goering2023}{Goering 2023, 79--84}, §§4.5.1.1--4.5.2).

McCully and Hogg's revised quantity-sensitive conditions likewise make
initial-foot accommodation essential, rather than stressing every medial
heavy-looking string (\citeproc{ref-McCullyHogg1990}{McCully and Hogg
1990, 323--33}, especially pp.~330--333). These analyses motivate CAPR's
bounded projection; neither author prints a full accented derivation of
earthly. The exact syllabification and prosodic domain remain explicit
modelling assumptions, not a complete Old English cyclic stress grammar.

\subsection{World, rainbow and noncircular
controls}\label{world-rainbow-and-noncircular-controls}

\emph{Weorold} `world' has the contrasting initial light-plus-heavy
configuration. At the late-lowering checkpoint the inflected diagnostic
\emph{*weo.rul.des} contains one short diphthongal \emph{eo} nucleus in
its initial light syllable. The analytical accented \emph{*éo} is
neither two separately represented vowels nor a long vowel. Initial L+H
is accommodated together in the initial foot rather than giving the
second nucleus an independent medial head. By contrast,
\emph{eorþ.cun.de} has initial H followed by nonfinal H, permitting a
medial head. Its selected weak-vowel history therefore remains
compatible with late lowering (\citeproc{ref-Goering2023}{Goering 2023,
79--82}, §4.5.1.2). \emph{Reġnboga} `rainbow' already has second-root
\emph{o} from \hyperref[rule-PNWGmcULowering]{SC017 PNWGmcULowering}; it
cannot establish that subordinate prominence disappeared before late
lowering. Neither adopted history is reopened.

Hogg warns that ``tertiary stress'' can merely rename a desired
retention. Ordinary \emph{-un \textgreater{} -on} also prevents a
blanket dental-nasal exemption; labial-nasal preservation is a different
environment (\citeproc{ref-Hogg1996}{Hogg 1996, 8--11}, especially p.~11
n.~7; \citeproc{ref-HoggGrammar2011}{Hogg 2011, 239--41}, §§6.55--6.58).
Dictionary-located \emph{godcundne} `divine' and \emph{godcundre}
`divine' are real retention comparisons, but do not by themselves
establish every ending's history. The earlier record's \emph{godcundum}
in \emph{in godcundum mægne} is dative singular in context; its
syncretic \emph{-um} must not be counted as an independently
demonstrated plural. Final uninflected \emph{godcund} remains a
qualification of an ordinary nonfinal-foot account
(\citeproc{ref-Toller1921}{Toller 1921}, s.v.~god-cund). Metrical
prominence is not automatically attested speech stress
(\citeproc{ref-MinkovaZhou2023}{Minkova and Zhou 2023, 86--90}).

\subsection{The selected regular span and its executable
representation}\label{the-selected-regular-span-and-its-executable-representation}

The supported chain is \emph{*érθa-kùnðǭ > *érθa-kùndǭ > *érθakundǭ >
*éorθakundǭ > *éorθkundǭ > *éorθkundæ > *éorθkunde}, followed by
\emph{eorþcunde}. Dental hardening, linking-vowel deletion, breaking and
existing final shortening and merger account for the changes other than
medial protection. The selected bimoric nasal ending is already within
the input notation.

SC038 OEStripSecondaryStress removes the earlier compound annotation at
its inherited position; ordinary nonhigh medial loss later belongs to
\hyperref[rule-OENonHighSyn]{SC114 OENonHighSyn}, without reading that
grave. At \hyperref[rule-OEMedUnstressedULowering]{SC040
OEMedUnstressedULowering}, the independently computed medial prominence
excludes this \emph{u} from unstressed lowering; without that protection
the experimental comparison gives \emph{*eorþconde}.

The adopted calculation retains temporary grave marking for all five
short vowels \emph{a e i o u} in the computed heads and immediately
decodes the marks after lowering. It accommodates an existing pretonic
prefix in the sentence domain and permits consonants preceding the first
syllable boundary. Such consonants do not add a nucleus or change its
weight. No new initial boundary is inserted, and the computation is not
wrapped in the separate pretonic-first-syllable context used by the
lowering operation itself.

Lowering retains its canonical consonant context. Its only additional
preceding-vowel exclusion is temporarily headed \emph{ù}: a headed
\emph{u} must not license lowering of a subsequent \emph{u} merely
because its technical annotation changes. This restores the existing
\emph{u}-exclusion after marking rather than introducing a new
historical condition. Complete-relation checks establish a single-valued
calculation, while complete-corpus controls preserve the existing
derivations. Both the full weak-neuter and genitive probes succeed.
Those results support the selected regular span without establishing an
entire adjective paradigm or removing its documentary qualifications.

The temporary marks are not survival of the original grave, a second
historical deaccentuation, or a new chronology edge. The projection
addresses the short-\emph{u} lowering question; marking five short-vowel
classes does not supply a full stress representation for all nuclei,
boundaries, final-overheavy syllables or transparent compound domains.

\subsection{Other cells and the limits of
adoption}\label{other-cells-and-the-limits-of-adoption}

The Oxford gloss has \emph{eorþcundre} `earthly', feminine genitive
singular, in \emph{ealre eorþcundre gesceafte}, translating Latin
\emph{totius terrestris (i. terrene) creaturae}
(\citeproc{ref-Toller1921}{Toller 1921}, s.v.~eorþ-cund, An. Ox. 691).
The exact edition, compiler and manuscript behind that locator have not
been independently established. Fulk's strong paradigm allows
\emph{*-(a)izōz}; the experimental \emph{*eorþcundere > *eorþcundre}
requires additional pre-\emph{r} syncope research, not an adopted ending
repair (\citeproc{ref-Fulk2018}{Fulk 2018, 209--11};
\citeproc{ref-Campbell1959}{Campbell 1959, 157--59}, §§388--393;
\citeproc{ref-RingeTaylor2014}{Ringe and Taylor 2014, 264--66}).

The Benedict sentence has \emph{eorþcundum} `earthly' agreeing with
masculine dative plural \emph{lustum} `desires', governed by
\emph{filigaþ} `follow'. The dictionary cites it at 139, 27 under
\emph{grædiglice} and 139, 28 under \emph{sceamleáslice}. The underlying
edition and manuscript have not been independently checked
(\citeproc{ref-Toller1921}{Toller 1921}, s.vv.~grædiglice,
sceamleáslice). The conventional strong dative plural \emph{*-aimiz}
currently gives experimental \emph{*-em}, not the attested \emph{-um}.
Secure synchronic inflection does not resolve that whole inherited
derivation (\citeproc{ref-Fulk2018}{Fulk 2018, 209--11}).

The strong genitive \emph{*eorþcundes}, from the supported
\emph{*-as(a)} ending, is an experimentally regular but unattested
fallback, not a documentary substitute for the Genesis cell. Bounded
negative searches are not proof that it never existed. \emph{Eardcundes}
`native' has a different first constituent; \emph{eorþcundlic} `earthly'
is a separately suffixed lexeme. Their dictionary grouping does not make
either an occurrence of the selected adjective. The chosen weak-neuter
span needs neither Oxford r-syncope, a dative ending repair,
final-overheavy protection, nor an additional analogy.

\section{\texorpdfstring{hammer --- OE
\emph{hamer}}{hammer --- OE hamer}}\label{hammer-oe-hamer}

\index[iv]{02oe@\ivlangheader{Old English}{West Saxon normalization unmarked}!hamer@\iventry{hamer}{}}
\index[iv]{03pgmc@\ivlangheader{Proto-Germanic}{}!xamaraz@\iventry{*xámaraz}{}}

Derivation: \emph{*xámaraz} \textgreater{} \emph{hamer} (attested
variant).

\subsection{Derivation trace}\label{derivation-trace-99}

Proto input: \emph{*xámaraz}

\begingroup
\setlength{\fboxsep}{6pt}

\noindent\fbox{%
\begin{minipage}{0.97\linewidth}
\small
\begin{tabular}{@{}>{\raggedright\arraybackslash}p{0.359\linewidth}@{\hspace{0.030\linewidth}}>{\raggedright\arraybackslash}p{0.611\linewidth}@{}}
\begin{minipage}[t]{\linewidth}
\centering\textbf{Earlier Germanic changes}\par
\vspace{0.35em}
\raggedright
\centering\textbf{Northwest and West Germanic}\par
\raggedright
\vspace{0.2em}
\begin{tabularx}{\linewidth}{@{}>{\raggedright\arraybackslash}X@{\hspace{0.55em}}>{\raggedright\arraybackslash}p{0.24\linewidth}@{\hspace{0.25em}}}
\mbox{EAF Final Z Deletion} & \emph{*xámara} \\
\end{tabularx}
\vspace{0.6em}
\centering\textbf{Early Anglo-Frisian}\par
\raggedright
\vspace{0.2em}
\raggedright [no change]\par
\end{minipage}
&
\begin{minipage}[t]{\linewidth}
\centering\textbf{Old English changes}\par
\vspace{0.35em}
\raggedright
\begin{tabularx}{\linewidth}{@{}>{\raggedright\arraybackslash}X@{\hspace{0.55em}}>{\raggedright\arraybackslash}p{0.14\linewidth}@{\hspace{0.25em}}}
\mbox{PWGmc Final Bare A Loss} & \emph{*xámar} \\
\mbox{EAF Brightening Unstressed} & \emph{*xámær} \\
\mbox{OE Unstressed AE Merger} & \emph{*xámer} \\
\end{tabularx}
\end{minipage}
\\
\end{tabular}
\end{minipage}%
} \endgroup

Old English form: \emph{hamer}

\subsection{Reconstruction and comparative
evidence}\label{reconstruction-and-comparative-evidence-83}

The masculine a-stem
\Recon{xámaraz}\index[iv]{03pgmc@\ivlangheader{Proto-Germanic}{}!xamaraz@\iventry{*xámaraz}{}}
`hammer' belongs to the source tradition represented by Kroonen's stem
\emph{*hamara-} and Orel's citation \emph{*xamaraz}. The different
initial notation does not itself constitute an alternative paradigm.
Their deeper metathesis and stone/anvil associations remain comparative
arguments, not consequences of the Old English derivation
(\citeproc{ref-Kroonen2013}{Kroonen 2013, 206--7};
\citeproc{ref-Orel2003}{Orel 2003, 158}). Orel's printed page is 158;
197 is the scan sheet, not the printed folio.

Old English
{\emph{hamer}}\index[iv]{02oe@\ivlangheader{Old English}{West Saxon normalization unmarked}!hamer@\iventry{hamer}{}}
`hammer' and
{\emph{hamor}}\index[iv]{02oe@\ivlangheader{Old English}{West Saxon normalization unmarked}!hamor@\iventry{hamor}{}}
`hammer' are both real citation variants
(\citeproc{ref-ClarkHall1960}{Clark Hall 1960, 160}). Selecting the
former preserves a complete regular inherited span without claiming that
the latter or its obliques have thereby been explained.

\subsection{Old English evidence}\label{old-english-evidence-84}

The Cleopatra glossary prints \emph{Malleus, hamer.} The edition's
heading identifies Cotton Cleopatra A.III, folio 76r, with a
late-tenth-/early-eleventh-century hand. This is not an attribution to
Ælfric (\citeproc{ref-WrightWuelcker1884}{Wright 1884, vol. 1}, pp.~258,
272, line 36). The isolated Old English spelling is
nominative/accusative-syncretic; the Latin citation supports the
nominative selection rather than proving that Old English orthography
distinguished the two cases.

The genitive
{\emph{hameres}}\index[iv]{02oe@\ivlangheader{Old English}{West Saxon normalization unmarked}!hameres@\iventry{hameres}{}}
`of a hammer' is also genuinely attested: \emph{Beátendes hameres
stipnes}, An.Ox.481 (\citeproc{ref-Toller1921}{Toller 1921, 505}). The
earlier attribution to page78 of Bosworth--Toller is withdrawn.
Sievers--Brunner's
{\emph{hamor}}\index[iv]{02oe@\ivlangheader{Old English}{West Saxon normalization unmarked}!hamor@\iventry{hamor}{}}
--
{\emph{hamores}}\index[iv]{02oe@\ivlangheader{Old English}{West Saxon normalization unmarked}!hamores@\iventry{hamores}{}}
paradigm retains a second oblique tradition
(\citeproc{ref-SieversBrunner1965}{Brunner 1965, § 245}). The
coexistence of these forms requires explanation beyond this selected
citation; it is not a reason to discard any of them.

\subsection{Development to Old
English}\label{development-to-old-english-82}

Final \emph{z} disappears, and early loss of bare final \emph{a} reduces
\Recon{xámaraz}\index[iv]{03pgmc@\ivlangheader{Proto-Germanic}{}!xamaraz@\iventry{*xámaraz}{}}
to
\Recon{xámar}\index[iv]{08preoe@\ivlangheader{Pre-Old English / prehistoric English}{}!xamar@\iventry{*xámar}{}}.
Unstressed brightening gives medial \emph{æ}, and its later merger with
\emph{e} yields {\emph{hamer}}. The medial vowel is now in the last
syllable: it is not an open medial target before a following heavy final
syllable. \hyperref[rule-OENonHighSyn]{SC114 OENonHighSyn} consequently
does not apply.

This negative control matters. The older genitive
\Recon{xámaras}\index[iv]{03pgmc@\ivlangheader{Proto-Germanic}{}!xamaras@\iventry{*xámaras}{}}
still has a heavy final syllable when early nonhigh syncope applies. Its
actual output is \emph{hamres}, not attested
{\emph{hameres}}\index[iv]{02oe@\ivlangheader{Old English}{West Saxon normalization unmarked}!hameres@\iventry{hameres}{}}.
The previous claim that brightening alone supplied a complete regular
genitive history is therefore not retained.

\subsection{Paradigm comparison}\label{paradigm-comparison}

{\def\LTcaptype{none} % do not increment counter
\begin{longtable}[]{@{}
  >{\raggedright\arraybackslash}p{(\linewidth - 8\tabcolsep) * \real{0.2000}}
  >{\raggedright\arraybackslash}p{(\linewidth - 8\tabcolsep) * \real{0.2000}}
  >{\raggedright\arraybackslash}p{(\linewidth - 8\tabcolsep) * \real{0.2000}}
  >{\raggedright\arraybackslash}p{(\linewidth - 8\tabcolsep) * \real{0.2000}}
  >{\raggedright\arraybackslash}p{(\linewidth - 8\tabcolsep) * \real{0.2000}}@{}}
\toprule\noalign{}
\begin{minipage}[b]{\linewidth}\raggedright
Cell
\end{minipage} & \begin{minipage}[b]{\linewidth}\raggedright
Inherited input
\end{minipage} & \begin{minipage}[b]{\linewidth}\raggedright
Actual model output
\end{minipage} & \begin{minipage}[b]{\linewidth}\raggedright
Attested comparison
\end{minipage} & \begin{minipage}[b]{\linewidth}\raggedright
Assessment
\end{minipage} \\
\midrule\noalign{}
\endhead
\bottomrule\noalign{}
\endlastfoot
selected nominative singular &
\Recon{xámaraz}\index[iv]{03pgmc@\ivlangheader{Proto-Germanic}{}!xamaraz@\iventry{*xámaraz}{}}
& {\emph{hamer}} &
{\emph{hamer}}\index[iv]{02oe@\ivlangheader{Old English}{West Saxon normalization unmarked}!hamer@\iventry{hamer}{}}
& complete regular span; no nonhigh-syncope application \\
former genitive singular &
\Recon{xámaras}\index[iv]{03pgmc@\ivlangheader{Proto-Germanic}{}!xamaras@\iventry{*xámaras}{}}
& \emph{hamres} &
{\emph{hameres}}\index[iv]{02oe@\ivlangheader{Old English}{West Saxon normalization unmarked}!hameres@\iventry{hameres}{}}
& real attestation, but not a successful current derivation \\
other citation and oblique tradition & same lexical family & not
established here &
{\emph{hamor}}\index[iv]{02oe@\ivlangheader{Old English}{West Saxon normalization unmarked}!hamor@\iventry{hamor}{}},
{\emph{hamores}}\index[iv]{02oe@\ivlangheader{Old English}{West Saxon normalization unmarked}!hamores@\iventry{hamores}{}}
& retained evidence; wider distribution unresolved \\
\end{longtable}
}

The attested-variant classification identifies precisely this situation:
one source-backed regular output is selected from a genuine attested
distribution. It is not an appeal to lexical diffusion or an assertion
that the sound law varies by word.

\section{\texorpdfstring{ten --- OE
\emph{tēon}}{ten --- OE tēon}}\label{ten-oe-tux113on}

\index[iv]{02oe@\ivlangheader{Old English}{West Saxon normalization unmarked}!teon@\iventry{tēon}{}}
\index[iv]{03pgmc@\ivlangheader{Proto-Germanic}{}!texun@\iventry{*téxun}{}}

Derivation: \emph{*téxun} \textgreater{} \emph{tēon} (attested variant).

\subsection{Derivation trace}\label{derivation-trace-100}

Proto input: \emph{*téxun}

\begingroup
\setlength{\fboxsep}{6pt}

\noindent\fbox{%
\begin{minipage}{0.97\linewidth}
\small
\begin{tabular}{@{}>{\raggedright\arraybackslash}p{0.359\linewidth}@{\hspace{0.030\linewidth}}>{\raggedright\arraybackslash}p{0.611\linewidth}@{}}
\begin{minipage}[t]{\linewidth}
\centering\textbf{Earlier Germanic changes}\par
\vspace{0.35em}
\raggedright
\centering\textbf{Northwest and West Germanic}\par
\raggedright
\vspace{0.2em}
\raggedright [no change]\par
\vspace{0.6em}
\centering\textbf{Early Anglo-Frisian}\par
\raggedright
\vspace{0.2em}
\raggedright [no change]\par
\end{minipage}
&
\begin{minipage}[t]{\linewidth}
\centering\textbf{Old English changes}\par
\vspace{0.35em}
\raggedright
\begin{tabularx}{\linewidth}{@{}>{\raggedright\arraybackslash}X@{\hspace{0.55em}}>{\raggedright\arraybackslash}p{0.14\linewidth}@{\hspace{0.25em}}}
\mbox{OE Breaking} & \emph{*téoxun} \\
\mbox{OE Med Unstressed U Lowering} & \emph{*téoxon} \\
\mbox{OE H Loss} & \emph{*téoon} \\
\mbox{OE Contraction} & \emph{*tḗon} \\
\end{tabularx}
\end{minipage}
\\
\end{tabular}
\end{minipage}%
} \endgroup

Old English form: \emph{tēon}

\subsection{Reconstruction and comparative
evidence}\label{reconstruction-and-comparative-evidence-84}

\Recon{tēon} `ten' is the bare cardinal form. Fulk states that Old
English \emph{tien} `ten' shows umlaut from the inflected forms, whereas
the uninflected form without umlaut is reflected in \emph{tēon} `ten'
(\citeproc{ref-Fulk2018}{Fulk 2018, § 10.2}). Brunner gives the same
contrast more broadly: \Recon{tēon} `ten' develops from \Recon{tëhun}
`ten', while West Saxon \emph{tien} `ten', \emph{tȳn} `ten' belong to a
different, umlauted branch of the numeral history
(\citeproc{ref-SieversBrunner1965}{Brunner 1965, §§ 129.2, 129} Anm.
6, 234).

The comparison form \emph{tēon} `ten' therefore represents the bare
cardinal's un-umlauted line, not the later umlauted simplex tradition.

\subsection{Old English evidence}\label{old-english-evidence-85}

The attested simplex forms are varied. Campbell gives \emph{tien} `ten',
north-western West Saxon \emph{tēn} `ten', and late Northumbrian
\emph{tēo} `ten', \emph{tēa} `ten' (\citeproc{ref-Campbell1959}{Campbell
1959, § 682}). Brunner likewise lists West Saxon \emph{tien} `ten',
\emph{tȳn} `ten' beside \emph{tēn} `ten', \emph{tēo} `ten', \emph{tēa}
`ten' in other dialects (\citeproc{ref-SieversBrunner1965}{Brunner 1965,
§ 325}).

Exact simplex \emph{tēon} `ten' is weaker as a directly cited headword
than those spellings. The un-umlauted stem is, however, explicit in
\emph{tēoða} `tenth' and \emph{tēontig} `ten'
(\citeproc{ref-SieversBrunner1965}{Brunner 1965, § 129.2};
\citeproc{ref-Fulk2018}{Fulk 2018, § 10.2}). The comparison form
\emph{tēon} `ten' is therefore a normalized spelling of that un-umlauted
base.

\subsection{Development to Old
English}\label{development-to-old-english-83}

From \Recon{téxun} `ten', lowering of medial unstressed \emph{u} gives
\Recon{téxon} `ten', breaking gives \Recon{téoxon} `ten', loss of
intervocalic \emph{h}/\emph{x} yields \Recon{téoon} `ten', and
contraction produces \Recon{tḗon} `ten', written \emph{tēon} `ten'. This
is the regular bare-cardinal path.

The umlauted forms \emph{tien} / \emph{tīen} `ten' belong to a different
branch, created when the numeral was levelled from inflected forms with
a front-vocalic trigger (\citeproc{ref-Fulk2018}{Fulk 2018, § 10.2};
\citeproc{ref-SieversBrunner1965}{Brunner 1965, § 129} Anm. 6).

\subsection{Variant comparison}\label{variant-comparison-1}

The comparison below sets the relevant forms side by side. It
distinguishes the normalized un-umlauted base from the attested simplex
variants.

{\def\LTcaptype{none} % do not increment counter
\begin{longtable}[]{@{}
  >{\raggedright\arraybackslash}p{(\linewidth - 4\tabcolsep) * \real{0.3333}}
  >{\raggedright\arraybackslash}p{(\linewidth - 4\tabcolsep) * \real{0.3333}}
  >{\raggedright\arraybackslash}p{(\linewidth - 4\tabcolsep) * \real{0.3333}}@{}}
\toprule\noalign{}
\begin{minipage}[b]{\linewidth}\raggedright
Form or branch
\end{minipage} & \begin{minipage}[b]{\linewidth}\raggedright
Status
\end{minipage} & \begin{minipage}[b]{\linewidth}\raggedright
Relevance to this entry
\end{minipage} \\
\midrule\noalign{}
\endhead
\bottomrule\noalign{}
\endlastfoot
\emph{tēon} & normalized un-umlauted comparison form; trace-supported &
Old English form here \\
\emph{tien} / \emph{tīen} & attested West Saxon umlauted simplex forms &
genuine OE variants, but not the bare-cardinal line modeled here \\
\emph{tēn} / \emph{tēo} / \emph{tēa} & attested un-umlauted simplex
variants in other dialects & support the same branch as the comparison
form \\
\end{longtable}
}

\section{\texorpdfstring{three --- OE
\emph{þrīe}}{three --- OE þrīe}}\label{three-oe-uxferux12be}

\index[iv]{02oe@\ivlangheader{Old English}{West Saxon normalization unmarked}!thrie@\iventry{þrīe}{}}
\index[iv]{03pgmc@\ivlangheader{Proto-Germanic}{}!threjez@\iventry{*θréjez}{}}

Derivation: \emph{*θréjez} \textgreater{} \emph{þrīe} (attested
variant).

\subsection{Derivation trace}\label{derivation-trace-101}

Proto input: \emph{*θréjez}

\begingroup
\setlength{\fboxsep}{6pt}

\noindent\fbox{%
\begin{minipage}{0.97\linewidth}
\small
\begin{tabular}{@{}>{\raggedright\arraybackslash}p{0.359\linewidth}@{\hspace{0.030\linewidth}}>{\raggedright\arraybackslash}p{0.611\linewidth}@{}}
\begin{minipage}[t]{\linewidth}
\centering\textbf{Earlier Germanic changes}\par
\vspace{0.35em}
\raggedright
\centering\textbf{Northwest and West Germanic}\par
\raggedright
\vspace{0.2em}
\begin{tabularx}{\linewidth}{@{}>{\raggedright\arraybackslash}X@{\hspace{0.55em}}>{\raggedright\arraybackslash}p{0.14\linewidth}@{\hspace{0.25em}}}
\mbox{EAF Final Z Deletion} & \emph{*θréje} \\
\end{tabularx}
\vspace{0.6em}
\centering\textbf{Early Anglo-Frisian}\par
\raggedright
\vspace{0.2em}
\raggedright [no change]\par
\end{minipage}
&
\begin{minipage}[t]{\linewidth}
\centering\textbf{Old English changes}\par
\vspace{0.35em}
\raggedright
\begin{tabularx}{\linewidth}{@{}>{\raggedright\arraybackslash}X@{\hspace{0.55em}}>{\raggedright\arraybackslash}p{0.14\linewidth}@{\hspace{0.25em}}}
\mbox{OE I Umlaut} & \emph{*θrije} \\
OE Intervocalic J Vocalization & \emph{*θriie} \\
\mbox{OE Contraction} & \emph{*θrīe} \\
\end{tabularx}
\end{minipage}
\\
\end{tabular}
\end{minipage}%
} \endgroup

Old English form: \emph{þrīe}

\subsection{Reconstruction and comparative
evidence}\label{reconstruction-and-comparative-evidence-85}

Kroonen cites the numeral under a broader stem-style reconstruction
rather than under one Old English-ready paradigm cell
(\citeproc{ref-Kroonen2013}{Kroonen 2013, 586}). The input
\Recon{θréjez}\index[iv]{03pgmc@\ivlangheader{Proto-Germanic}{}!threjez@\iventry{*θréjez}{}}
`three' is therefore best understood as the inherited masculine
nominative-accusative singular.

The Old English numeral has no uniform citation form across the
paradigm. The masculine singular line must be kept apart from
feminine-neuter \emph{þrēo} `three' and from later reduced spellings of
the masculine form.

\subsection{Old English evidence}\label{old-english-evidence-86}

Campbell gives masculine nominative-accusative
{\emph{þrīe}}\index[iv]{02oe@\ivlangheader{Old English}{West Saxon normalization unmarked}!thrie@\iventry{þrīe}{}}
`three', feminine and neuter nominative-accusative
{\emph{þrēo}}\index[iv]{02oe@\ivlangheader{Old English}{West Saxon normalization unmarked}!threo@\iventry{þrēo}{}}
`three', genitive
{\emph{þrēora}}\index[iv]{02oe@\ivlangheader{Old English}{West Saxon normalization unmarked}!threora@\iventry{þrēora}{}}
`three', and dative
{\emph{þrim}}\index[iv]{02oe@\ivlangheader{Old English}{West Saxon normalization unmarked}!thrim@\iventry{þrim}{}}
`three', adding that late West Saxon has \emph{þry} `three' / \emph{þri}
`three' for masculine nominative-accusative beside
{\emph{þrīe}}\index[iv]{02oe@\ivlangheader{Old English}{West Saxon normalization unmarked}!thrie@\iventry{þrīe}{}}
`three' (\citeproc{ref-Campbell1959}{Campbell 1959, § 683}). Fulk
presents the same masculine \emph{þrīe} `three' beside the wider numeral
paradigm (\citeproc{ref-Fulk2018}{Fulk 2018, § 10.1}).

The target is therefore an attested Old English paradigm cell.
\emph{þrī} `three' / \emph{þry} `three' belongs to later reduction or
headword-style citation, whereas
{\emph{þrīe}}\index[iv]{02oe@\ivlangheader{Old English}{West Saxon normalization unmarked}!thrie@\iventry{þrīe}{}}
`three' is the conservative masculine nominative-accusative form.

\subsection{Development to Old
English}\label{development-to-old-english-84}

From \Recon{θréjez} `three', loss of final \emph{-z} leaves
\Recon{θréje} `three' as an intermediate form. The following
\Recon{θrije} `three' fronts the stem vowel, then vocalizes between
vowels, and contraction yields \Recon{þrīe} `three'. The compact trace
records the same sequence as θrije \textgreater{} θriie \textgreater{}
\emph{þrīe}.

\subsection{Variant comparison}\label{variant-comparison-2}

The comparison below sets the relevant forms side by side. It separates
the selected masculine cell from the later reduced form and from the
rest of the numeral paradigm.

{\def\LTcaptype{none} % do not increment counter
\begin{longtable}[]{@{}
  >{\raggedright\arraybackslash}p{(\linewidth - 4\tabcolsep) * \real{0.3333}}
  >{\raggedright\arraybackslash}p{(\linewidth - 4\tabcolsep) * \real{0.3333}}
  >{\raggedright\arraybackslash}p{(\linewidth - 4\tabcolsep) * \real{0.3333}}@{}}
\toprule\noalign{}
\begin{minipage}[b]{\linewidth}\raggedright
Form
\end{minipage} & \begin{minipage}[b]{\linewidth}\raggedright
Status
\end{minipage} & \begin{minipage}[b]{\linewidth}\raggedright
Relevance to this entry
\end{minipage} \\
\midrule\noalign{}
\endhead
\bottomrule\noalign{}
\endlastfoot
{\emph{þrīe}}\index[iv]{02oe@\ivlangheader{Old English}{West Saxon normalization unmarked}!thrie@\iventry{þrīe}{}}
& attested masculine nom./acc.; regular output & Old English form
here \\
\emph{þrī} / \emph{þry} & later reduced masculine variant & genuine OE
variant, but not the conservative comparison form \\
{\emph{þrēo}}\index[iv]{02oe@\ivlangheader{Old English}{West Saxon normalization unmarked}!threo@\iventry{þrēo}{}}
& attested feminine-neuter nom./acc. & same numeral, different paradigm
cell \\
{\emph{þrēora}}\index[iv]{02oe@\ivlangheader{Old English}{West Saxon normalization unmarked}!threora@\iventry{þrēora}{}},
{\emph{þrim}}\index[iv]{02oe@\ivlangheader{Old English}{West Saxon normalization unmarked}!thrim@\iventry{þrim}{}}
& attested genitive and dative forms & confirm the wider paradigm, not
the cell compared here \\
\end{longtable}
}

\section{\texorpdfstring{wasp --- OE
\emph{wæfs}}{wasp --- OE wæfs}}\label{wasp-oe-wuxe6fs}

\index[iv]{02oe@\ivlangheader{Old English}{West Saxon normalization unmarked}!waefs@\iventry{wæfs}{}}
\index[iv]{03pgmc@\ivlangheader{Proto-Germanic}{}!wabsaz@\iventry{*wábsaz}{}}

Derivation: \emph{*wábsaz} \textgreater{} \emph{wæfs} (attested
variant).

\subsection{Derivation trace}\label{derivation-trace-102}

Proto input: \emph{*wábsaz}

\begingroup
\setlength{\fboxsep}{6pt}

\noindent\fbox{%
\begin{minipage}{0.97\linewidth}
\small
\begin{tabular}{@{}>{\raggedright\arraybackslash}p{0.388\linewidth}@{\hspace{0.030\linewidth}}>{\raggedright\arraybackslash}p{0.582\linewidth}@{}}
\begin{minipage}[t]{\linewidth}
\centering\textbf{Earlier Germanic changes}\par
\vspace{0.35em}
\raggedright
\centering\textbf{Northwest and West Germanic}\par
\raggedright
\vspace{0.2em}
\begin{tabularx}{\linewidth}{@{}>{\raggedright\arraybackslash}X@{\hspace{0.55em}}>{\raggedright\arraybackslash}p{0.14\linewidth}@{\hspace{0.25em}}}
\mbox{EAF Final Z Deletion} & \emph{*wábsa} \\
\end{tabularx}
\vspace{0.6em}
\centering\textbf{Early Anglo-Frisian}\par
\raggedright
\vspace{0.2em}
\raggedright [no change]\par
\end{minipage}
&
\begin{minipage}[t]{\linewidth}
\centering\textbf{Old English changes}\par
\vspace{0.35em}
\raggedright
\begin{tabularx}{\linewidth}{@{}>{\raggedright\arraybackslash}X@{\hspace{0.55em}}>{\raggedright\arraybackslash}p{0.14\linewidth}@{\hspace{0.25em}}}
\mbox{PWGmc Final Bare A Loss} & \emph{*wábs} \\
\mbox{EAF Brightening} & \emph{*wæbs} \\
\mbox{PGmc B Allophony} & \emph{*wæβs} \\
\end{tabularx}
\end{minipage}
\\
\end{tabular}
\end{minipage}%
} \endgroup

Old English form: \emph{wæfs}

\subsection{Reconstruction and comparative
evidence}\label{reconstruction-and-comparative-evidence-86}

The Proto-Germanic form
\Recon{wábsaz}\index[iv]{03pgmc@\ivlangheader{Proto-Germanic}{}!wabsaz@\iventry{*wábsaz}{}}
`wasp' reaches Old English without any special change of stem or
paradigm cell. The question in this entry is instead which attested Old
English member of the variant set should serve as the comparison form.

Fulk presents the Old English forms together as
{\emph{wæfs}}\index[iv]{02oe@\ivlangheader{Old English}{West Saxon normalization unmarked}!waefs@\iventry{wæfs}{}}
`wasp' with variants
{\emph{wæsp}}\index[iv]{02oe@\ivlangheader{Old English}{West Saxon normalization unmarked}!waesp@\iventry{wæsp}{}}
`wasp' and
{\emph{wæps}}\index[iv]{02oe@\ivlangheader{Old English}{West Saxon normalization unmarked}!waeps@\iventry{wæps}{}}
`wasp' (\citeproc{ref-Fulk2018}{Fulk 2018, § 6.5}). Bülbring and
Brunner then make the chronology more explicit by deriving later
\emph{wæps} and late West Saxon
{\emph{wasp}}\index[iv]{02oe@\ivlangheader{Old English}{West Saxon normalization unmarked}!wasp~03@\iventry{wasp}{LWS}}
`wasp' from earlier \emph{waefs} / \emph{wæfs} through restricted
metatheses (\citeproc{ref-Bulbring1902}{Bülbring 1902, § 484} Anm.
3; \citeproc{ref-SieversBrunner1965}{Brunner 1965, §§ 193, 204}).

\subsection{Old English evidence}\label{old-english-evidence-87}

The earliest directly cited Old English form is
{\emph{wæfs}}\index[iv]{02oe@\ivlangheader{Old English}{West Saxon normalization unmarked}!waefs@\iventry{wæfs}{}}
`wasp', written
{\emph{waefs}}\index[iv]{02oe@\ivlangheader{Old English}{West Saxon normalization unmarked}!waefs@\iventry{waefs}{}}
`wasp' in the Épinal-Corpus material discussed by Bülbring and Brunner
(\citeproc{ref-Bulbring1902}{Bülbring 1902, § 484} Anm. 3;
\citeproc{ref-SieversBrunner1965}{Brunner 1965, § 193}). Later Old
English also shows
{\emph{wæps}}\index[iv]{02oe@\ivlangheader{Old English}{West Saxon normalization unmarked}!waeps@\iventry{wæps}{}}
`wasp' and
{\emph{wæsp}}\index[iv]{02oe@\ivlangheader{Old English}{West Saxon normalization unmarked}!waesp@\iventry{wæsp}{}}
`wasp' /
{\emph{wasp}}\index[iv]{02oe@\ivlangheader{Old English}{West Saxon normalization unmarked}!wasp@\iventry{wasp}{}}
`wasp', and dictionary practice often favors
{\emph{wæps}}\index[iv]{02oe@\ivlangheader{Old English}{West Saxon normalization unmarked}!waeps@\iventry{wæps}{}}
or later spellings as headwords (\citeproc{ref-ClarkHall1960}{Clark Hall
1960, 341}).

This entry therefore distinguishes chronological priority from headword
habit.
{\emph{wæfs}}\index[iv]{02oe@\ivlangheader{Old English}{West Saxon normalization unmarked}!waefs@\iventry{wæfs}{}}
`wasp' is not a convenient reconstruction: it is an attested Old English
form and also the one that matches the regular development most closely.

\subsection{Development to Old
English}\label{development-to-old-english-85}

From
\Recon{wábsaz}\index[iv]{03pgmc@\ivlangheader{Proto-Germanic}{}!wabsaz@\iventry{*wábsaz}{}}
`wasp', the regular Old English path passes through loss of final
\emph{z}, Anglo-Frisian fronting, and the allophonic development of
\emph{b} to a fricative before \emph{s}, yielding {\emph{wæfs}} `wasp'.

The later forms
{\emph{wæps}}\index[iv]{02oe@\ivlangheader{Old English}{West Saxon normalization unmarked}!waeps@\iventry{wæps}{}}
`wasp' and
{\emph{wæsp}}\index[iv]{02oe@\ivlangheader{Old English}{West Saxon normalization unmarked}!waesp@\iventry{wæsp}{}}
`wasp' /
{\emph{wasp}}\index[iv]{02oe@\ivlangheader{Old English}{West Saxon normalization unmarked}!wasp@\iventry{wasp}{}}
`wasp' belong to subsequent, lexically restricted metatheses. They are
genuine Old English forms, but they are later within the variant
history.

\subsection{Variant comparison}\label{variant-comparison-3}

The comparison below sets the relevant forms side by side. It separates
the earliest attested and regular form from the later metathesized
doublets.

{\def\LTcaptype{none} % do not increment counter
\begin{longtable}[]{@{}
  >{\raggedright\arraybackslash}p{(\linewidth - 4\tabcolsep) * \real{0.3333}}
  >{\raggedright\arraybackslash}p{(\linewidth - 4\tabcolsep) * \real{0.3333}}
  >{\raggedright\arraybackslash}p{(\linewidth - 4\tabcolsep) * \real{0.3333}}@{}}
\toprule\noalign{}
\begin{minipage}[b]{\linewidth}\raggedright
Form
\end{minipage} & \begin{minipage}[b]{\linewidth}\raggedright
Status
\end{minipage} & \begin{minipage}[b]{\linewidth}\raggedright
Relevance to this entry
\end{minipage} \\
\midrule\noalign{}
\endhead
\bottomrule\noalign{}
\endlastfoot
{\emph{wæfs}}\index[iv]{02oe@\ivlangheader{Old English}{West Saxon normalization unmarked}!waefs@\iventry{wæfs}{}}
& earliest attested OE form; regular output & Old English form here \\
{\emph{wæps}}\index[iv]{02oe@\ivlangheader{Old English}{West Saxon normalization unmarked}!waeps@\iventry{wæps}{}}
& later attested metathesized variant & genuine OE doublet, but
secondary \\
{\emph{wæsp}}\index[iv]{02oe@\ivlangheader{Old English}{West Saxon normalization unmarked}!waesp@\iventry{wæsp}{}}
/
{\emph{wasp}}\index[iv]{02oe@\ivlangheader{Old English}{West Saxon normalization unmarked}!wasp@\iventry{wasp}{}}
& later West Saxon metathesized variant & genuine OE doublet, but not
the form compared here \\
\end{longtable}
}

\clearpage

\chapter{Early analogy and pre-Old-English input
selection}\label{early-analogy-and-pre-old-english-input-selection}

Here an analogical change already separates the transducer input from
the lexeme's citation reconstruction before the specifically Old English
changes apply. Each entry must therefore justify the remodeled input
independently of the successful output.

\section{\texorpdfstring{bottom --- OE
\emph{botm}}{bottom --- OE botm}}\label{bottom-oe-botm}

\index[iv]{02oe@\ivlangheader{Old English}{West Saxon normalization unmarked}!botm@\iventry{botm}{}}
\index[iv]{03pgmc@\ivlangheader{Proto-Germanic}{}!budmaz@\iventry{*búdmaz}{}}
\index[iv]{03pgmc@\ivlangheader{Proto-Germanic}{}!buttmaz@\iventry{*búttmaz}{}}

Derivation: citation reconstruction \emph{*búdmaz}; form followed here
\emph{*búttmaz} \textgreater{} \emph{botm} (early analogy).

\subsection{Derivation trace}\label{derivation-trace-103}

Proto input: \emph{*búttmaz}

\begingroup
\setlength{\fboxsep}{6pt}

\noindent\fbox{%
\begin{minipage}{0.97\linewidth}
\small
\begin{tabular}{@{}>{\raggedright\arraybackslash}p{0.485\linewidth}@{\hspace{0.030\linewidth}}>{\raggedright\arraybackslash}p{0.485\linewidth}@{}}
\begin{minipage}[t]{\linewidth}
\centering\textbf{Earlier Germanic changes}\par
\vspace{0.35em}
\raggedright
\centering\textbf{Northwest and West Germanic}\par
\raggedright
\vspace{0.2em}
\begin{tabularx}{\linewidth}{@{}>{\raggedright\arraybackslash}X@{\hspace{0.55em}}>{\raggedright\arraybackslash}p{0.18\linewidth}@{\hspace{0.25em}}}
\mbox{PNWGmc U Lowering} & \emph{*bóttmaz} \\
\mbox{EAF Final Z Deletion} & \emph{*bóttma} \\
\end{tabularx}
\vspace{0.6em}
\centering\textbf{Early Anglo-Frisian}\par
\raggedright
\vspace{0.2em}
\raggedright [no change]\par
\end{minipage}
&
\begin{minipage}[t]{\linewidth}
\centering\textbf{Old English changes}\par
\vspace{0.35em}
\raggedright
\begin{tabularx}{\linewidth}{@{}>{\raggedright\arraybackslash}X@{\hspace{0.55em}}>{\raggedright\arraybackslash}p{0.14\linewidth}@{\hspace{0.25em}}}
\mbox{PWGmc Final Bare A Loss} & \emph{*bóttm} \\
OE Preconsonantal Degemination & \emph{*bótm} \\
\end{tabularx}
\end{minipage}
\\
\end{tabular}
\end{minipage}%
} \endgroup

Old English form: \emph{botm}

\subsection{Reconstruction and comparative
evidence}\label{reconstruction-and-comparative-evidence-87}

Kroonen reconstructs the word as a stem complex \Recon{budmō} `bottom'
with genitive \Recon{buttaz} `bottom', summarized as \emph{*budman-}
\textasciitilde{} \emph{*buttman-}, and gives Old English
{\emph{botm}}\index[iv]{02oe@\ivlangheader{Old English}{West Saxon normalization unmarked}!botm@\iventry{botm}{}}
`bottom' as the reflex (\citeproc{ref-Kroonen2013}{Kroonen 2013, 120}).
The comparative label \Recon{búdmaz} `bottom' names the lexeme-level
stem complex, while the derivational input \Recon{búttmaz} `bottom'
represents the pre-Old-English form with oblique \emph{*butt-}
generalized into the nominative formation.

Orel likewise preserves both sides of the comparison under
\Recon{budmaz} `bottom' and \Recon{butmaz} `bottom'
(\citeproc{ref-Orel2003}{Orel 2003, 100}). The derivational input is
thus a historical stem choice, not an arbitrary respelling.

\subsection{Old English evidence}\label{old-english-evidence-88}

The Old English noun itself is secure. Clark Hall gives
{\emph{botm}}\index[iv]{02oe@\ivlangheader{Old English}{West Saxon normalization unmarked}!botm@\iventry{botm}{}}
`bottom' (\citeproc{ref-ClarkHall1960}{Clark Hall 1960, 63}).
Bosworth-Toller cross-references
{\emph{bodan}}\index[iv]{02oe@\ivlangheader{Old English}{West Saxon normalization unmarked}!bodan@\iventry{bodan}{}}
`bottom, floor' to
{\emph{botm}}\index[iv]{02oe@\ivlangheader{Old English}{West Saxon normalization unmarked}!botm@\iventry{botm}{}}
`bottom', showing the wider reflex family without weakening the attested
lemma (\citeproc{ref-BosworthToller1898}{Bosworth and Toller 1898,
112}).

\subsection{Development to Old
English}\label{development-to-old-english-86}

Once the oblique \emph{*butt-} stem has been generalized, the
derivational input \Recon{búttmaz} `bottom' develops regularly to
{\emph{botm}}\index[iv]{02oe@\ivlangheader{Old English}{West Saxon normalization unmarked}!botm@\iventry{botm}{}}
`bottom'. The analogical step is therefore early: it belongs to
pre-Old-English stem formation rather than to a later choice among Old
English paradigm cells.

\section{\texorpdfstring{brand --- OE
\emph{brandes}}{brand --- OE brandes}}\label{brand-oe-brandes}

\index[iv]{02oe@\ivlangheader{Old English}{West Saxon normalization unmarked}!brandes@\iventry{brandes}{}}
\index[iv]{03pgmc@\ivlangheader{Proto-Germanic}{}!brandas@\iventry{*brándas}{}}
\index[iv]{03pgmc@\ivlangheader{Proto-Germanic}{}!brandaz@\iventry{*brándaz}{}}

Derivation: citation reconstruction \emph{*brándaz}; form followed here
\emph{*brándas} \textgreater{} \emph{brandes} (early analogy).

\subsection{Derivation trace}\label{derivation-trace-104}

Proto input: \emph{*brándas}

\begingroup
\setlength{\fboxsep}{6pt}

\noindent\fbox{%
\begin{minipage}{0.97\linewidth}
\small
\begin{tabular}{@{}>{\raggedright\arraybackslash}p{0.359\linewidth}@{\hspace{0.030\linewidth}}>{\raggedright\arraybackslash}p{0.611\linewidth}@{}}
\begin{minipage}[t]{\linewidth}
\centering\textbf{Earlier Germanic changes}\par
\vspace{0.35em}
\raggedright
\centering\textbf{Northwest and West Germanic}\par
\raggedright
\vspace{0.2em}
\raggedright [no change]\par
\vspace{0.6em}
\centering\textbf{Early Anglo-Frisian}\par
\raggedright
\vspace{0.2em}
\raggedright [no change]\par
\end{minipage}
&
\begin{minipage}[t]{\linewidth}
\centering\textbf{Old English changes}\par
\vspace{0.35em}
\raggedright
\begin{tabularx}{\linewidth}{@{}>{\raggedright\arraybackslash}X@{\hspace{0.55em}}>{\raggedright\arraybackslash}p{0.18\linewidth}@{\hspace{0.25em}}}
\mbox{EAF Brightening Unstressed} & \emph{*brándæs} \\
\mbox{OE Unstressed AE Merger} & \emph{*brándes} \\
\end{tabularx}
\end{minipage}
\\
\end{tabular}
\end{minipage}%
} \endgroup

Old English form: \emph{brandes}

\subsection{Reconstruction and comparative
evidence}\label{reconstruction-and-comparative-evidence-88}

The inherited noun is the masculine a-stem
\Recon{brándaz}\index[iv]{03pgmc@\ivlangheader{Proto-Germanic}{}!brandaz@\iventry{*brándaz}{}}
`brand', continued by Old English
{\emph{brand}}\index[iv]{02oe@\ivlangheader{Old English}{West Saxon normalization unmarked}!brand@\iventry{brand}{}}
`brand' and its continental cognates (\citeproc{ref-Orel2003}{Orel 2003,
53}). The selected input
\Recon{brándas}\index[iv]{03pgmc@\ivlangheader{Proto-Germanic}{}!brandas@\iventry{*brándas}{}}
`brand' is not a different lexeme but the genitive singular of that same
a-stem noun.

Both forms belong to the same root and stem class but occupy different
inherited inflectional cells. The derivational input is the oblique
cell.

\subsection{Old English evidence}\label{old-english-evidence-89}

Old English dictionaries lemmatize the noun as
{\emph{brand}}\index[iv]{02oe@\ivlangheader{Old English}{West Saxon normalization unmarked}!brand@\iventry{brand}{}}
`brand' (\citeproc{ref-ClarkHall1960}{Clark Hall 1960, 49};
\citeproc{ref-BosworthToller1898}{Bosworth and Toller 1898, 116}).
Bosworth-Toller also records inflectional forms such as
{\emph{brandas}}\index[iv]{02oe@\ivlangheader{Old English}{West Saxon normalization unmarked}!brandas@\iventry{brandas}{}}
`brand',
{\emph{branda}}\index[iv]{02oe@\ivlangheader{Old English}{West Saxon normalization unmarked}!branda@\iventry{branda}{}}
`brand', and
{\emph{brandum}}\index[iv]{02oe@\ivlangheader{Old English}{West Saxon normalization unmarked}!brandum@\iventry{brandum}{}}
`brand' under the same entry (\citeproc{ref-BosworthToller1898}{Bosworth
and Toller 1898, 116}).

The specific comparison form in this entry,
{\emph{brandes}}\index[iv]{02oe@\ivlangheader{Old English}{West Saxon normalization unmarked}!brandes@\iventry{brandes}{}}
`brand', is the expected genitive singular of that a-stem noun. It is
therefore an inferred Old English paradigm form rather than the ordinary
dictionary headword.

\subsection{Development to Old
English}\label{development-to-old-english-87}

From
\Recon{brándas}\index[iv]{03pgmc@\ivlangheader{Proto-Germanic}{}!brandas@\iventry{*brándas}{}}
`brand', the regular Old English development passes through the usual
unstressed-vowel weakening of the inflectional ending, yielding
{\emph{brandes}}\index[iv]{02oe@\ivlangheader{Old English}{West Saxon normalization unmarked}!brandes@\iventry{brandes}{}}
`brand'. Nothing in the stem itself requires a special repair. The root
consonants and the stressed vowel are the same as in the citation lemma
{\emph{brand}}\index[iv]{02oe@\ivlangheader{Old English}{West Saxon normalization unmarked}!brand@\iventry{brand}{}}
`brand'.

The analytical weight of the entry lies in the ending. By choosing the
oblique singular rather than the nominative citation form, the entry
presents the same lexeme in a different inherited cell.

\subsection{Form comparison}\label{form-comparison-3}

The comparison below sets the relevant forms side by side. It separates
the citation lemma from the selected oblique singular.

{\def\LTcaptype{none} % do not increment counter
\begin{longtable}[]{@{}
  >{\raggedright\arraybackslash}p{(\linewidth - 8\tabcolsep) * \real{0.2000}}
  >{\raggedright\arraybackslash}p{(\linewidth - 8\tabcolsep) * \real{0.2000}}
  >{\raggedright\arraybackslash}p{(\linewidth - 8\tabcolsep) * \real{0.2000}}
  >{\raggedright\arraybackslash}p{(\linewidth - 8\tabcolsep) * \real{0.2000}}
  >{\raggedright\arraybackslash}p{(\linewidth - 8\tabcolsep) * \real{0.2000}}@{}}
\toprule\noalign{}
\begin{minipage}[b]{\linewidth}\raggedright
Form / interpretation
\end{minipage} & \begin{minipage}[b]{\linewidth}\raggedright
Candidate input
\end{minipage} & \begin{minipage}[b]{\linewidth}\raggedright
Expected or documented OE outcome
\end{minipage} & \begin{minipage}[b]{\linewidth}\raggedright
OE comparison form
\end{minipage} & \begin{minipage}[b]{\linewidth}\raggedright
Result
\end{minipage} \\
\midrule\noalign{}
\endhead
\bottomrule\noalign{}
\endlastfoot
citation nominative singular &
{\emph{*brándaz}}\index[iv]{03pgmc@\ivlangheader{Proto-Germanic}{}!brandaz@\iventry{*brándaz}{}}
& expected OE lemma
{\emph{brand}}\index[iv]{02oe@\ivlangheader{Old English}{West Saxon normalization unmarked}!brand@\iventry{brand}{}}
&
{\emph{brand}}\index[iv]{02oe@\ivlangheader{Old English}{West Saxon normalization unmarked}!brand@\iventry{brand}{}}
& regular headword-level outcome \\
genitive singular &
{\emph{*brándas}}\index[iv]{03pgmc@\ivlangheader{Proto-Germanic}{}!brandas@\iventry{*brándas}{}}
& regular output:
{\emph{brandes}}\index[iv]{02oe@\ivlangheader{Old English}{West Saxon normalization unmarked}!brandes@\iventry{brandes}{}}
&
{\emph{brandes}}\index[iv]{02oe@\ivlangheader{Old English}{West Saxon normalization unmarked}!brandes@\iventry{brandes}{}}
& exact match for the oblique cell \\
\end{longtable}
}

The noun itself is straightforwardly inherited. The main point of the
entry is that
{\emph{brandes}}\index[iv]{02oe@\ivlangheader{Old English}{West Saxon normalization unmarked}!brandes@\iventry{brandes}{}}
`brand (gen.)' belongs to the same regular a-stem paradigm as
{\emph{brand}}\index[iv]{02oe@\ivlangheader{Old English}{West Saxon normalization unmarked}!brand@\iventry{brand}{}}
`brand', even though the citation lemma remains the nominative singular.

\section{\texorpdfstring{breast --- OE
\emph{brēost}}{breast --- OE brēost}}\label{breast-oe-brux113ost}

\index[iv]{02oe@\ivlangheader{Old English}{West Saxon normalization unmarked}!breost@\iventry{brēost}{}}
\index[iv]{03pgmc@\ivlangheader{Proto-Germanic}{}!breusta@\iventry{*bréustą}{}}
\index[iv]{03pgmc@\ivlangheader{Proto-Germanic}{}!brustz@\iventry{*brústz}{}}

Derivation: citation reconstruction \emph{*brústz}; form followed here
\emph{*bréustą} \textgreater{} \emph{brēost} (early analogy).

\subsection{Derivation trace}\label{derivation-trace-105}

Proto input: \emph{*bréustą}

\begingroup
\setlength{\fboxsep}{6pt}

\noindent\fbox{%
\begin{minipage}{0.97\linewidth}
\small
\begin{tabular}{@{}>{\raggedright\arraybackslash}p{0.359\linewidth}@{\hspace{0.030\linewidth}}>{\raggedright\arraybackslash}p{0.611\linewidth}@{}}
\begin{minipage}[t]{\linewidth}
\centering\textbf{Earlier Germanic changes}\par
\vspace{0.35em}
\raggedright
\centering\textbf{Northwest and West Germanic}\par
\raggedright
\vspace{0.2em}
\raggedright [no change]\par
\vspace{0.6em}
\centering\textbf{Early Anglo-Frisian}\par
\raggedright
\vspace{0.2em}
\raggedright [no change]\par
\end{minipage}
&
\begin{minipage}[t]{\linewidth}
\centering\textbf{Old English changes}\par
\vspace{0.35em}
\raggedright
\begin{tabularx}{\linewidth}{@{}>{\raggedright\arraybackslash}X@{\hspace{0.55em}}>{\raggedright\arraybackslash}p{0.18\linewidth}@{\hspace{0.25em}}}
\mbox{OE Diphthong Leveling} & \emph{*brēostą} \\
OE Heavy Syllable Nasal Apocope & \emph{*brēost} \\
\end{tabularx}
\end{minipage}
\\
\end{tabular}
\end{minipage}%
} \endgroup

Old English form: \emph{brēost}

\subsection{Reconstruction and comparative
evidence}\label{reconstruction-and-comparative-evidence-89}

The word family shows two related but distinct Proto-Germanic
formations. The root noun
{\emph{*brust-}}\index[iv]{03pgmc@\ivlangheader{Proto-Germanic}{}!brust@\iventry{*brust-}{}}
lies behind forms such as Gothic
{\emph{brusts}}\index[iv]{12goth@\ivlangheader{Gothic}{}!brusts@\iventry{brusts}{}}
`breast', whereas Old English
{\emph{brēost}}\index[iv]{02oe@\ivlangheader{Old English}{West Saxon normalization unmarked}!breost@\iventry{brēost}{}}
`breast' belongs to a thematic formation
{\emph{*breusta-}}\index[iv]{03pgmc@\ivlangheader{Proto-Germanic}{}!breusta@\iventry{*breusta-}{}},
alongside Old Norse
{\emph{brjóst}}\index[iv]{09on@\ivlangheader{Old Norse}{}!brjost@\iventry{brjóst}{}}
`breast' and Old Saxon
{\emph{briost}}\index[iv]{13os@\ivlangheader{Old Saxon}{}!briost@\iventry{briost}{}}
`breast' (\citeproc{ref-Kroonen2013}{Kroonen 2013, 114};
\citeproc{ref-Orel2003}{Orel 2003, 95};
\citeproc{ref-RingeTaylor2014}{Ringe and Taylor 2014, 43}).

The derivational input
\Recon{bréustą}\index[iv]{03pgmc@\ivlangheader{Proto-Germanic}{}!breusta@\iventry{*bréustą}{}}
`breast' therefore differs from the citation label
\Recon{brústz}\index[iv]{03pgmc@\ivlangheader{Proto-Germanic}{}!brustz@\iventry{*brústz}{}}
`breast' because Old English reflects the thematic branch rather than
the root noun. The morphological choice comes before the Old English
sound changes themselves.

\subsection{Old English evidence}\label{old-english-evidence-90}

Clark Hall records the noun as
{\emph{brēost}}\index[iv]{02oe@\ivlangheader{Old English}{West Saxon normalization unmarked}!breost@\iventry{brēost}{}}
`breast' /
{\emph{breóst}}\index[iv]{02oe@\ivlangheader{Old English}{West Saxon normalization unmarked}!breost@\iventry{breóst}{}}
`breast' (\citeproc{ref-ClarkHall1960}{Clark Hall 1960, 65}). The form
is an established Old English lexeme, not a reconstructed target
assembled from comparative evidence alone.

What requires explanation is not the Old English attestation but the
relation between that attested noun and the broader Germanic word
family. The relevant comparison form is therefore the thematic Old
English noun
{\emph{brēost}}\index[iv]{02oe@\ivlangheader{Old English}{West Saxon normalization unmarked}!breost@\iventry{brēost}{}}
`breast'.

\subsection{Development to Old
English}\label{development-to-old-english-88}

From
\Recon{bréustą}\index[iv]{03pgmc@\ivlangheader{Proto-Germanic}{}!breusta@\iventry{*bréustą}{}}
`breast', the regular Old English development gives
{\emph{brēost}}\index[iv]{02oe@\ivlangheader{Old English}{West Saxon normalization unmarked}!breost@\iventry{brēost}{}}
`breast', with the expected \emph{eu} \textgreater{} \emph{ēo} vowel
history (\citeproc{ref-Campbell1959}{Campbell 1959, § 115}). No
special repair is needed once the correct thematic formation is chosen.

The earlier mismatch arose only if the word was forced into the
root-noun line. The Old English noun itself continues the thematic
branch cleanly and directly.

\subsection{Formation comparison}\label{formation-comparison}

The comparison below sets the relevant forms side by side. It separates
the broader root-noun family label from the thematic formation actually
continued in Old English.

{\def\LTcaptype{none} % do not increment counter
\begin{longtable}[]{@{}
  >{\raggedright\arraybackslash}p{(\linewidth - 8\tabcolsep) * \real{0.2000}}
  >{\raggedright\arraybackslash}p{(\linewidth - 8\tabcolsep) * \real{0.2000}}
  >{\raggedright\arraybackslash}p{(\linewidth - 8\tabcolsep) * \real{0.2000}}
  >{\raggedright\arraybackslash}p{(\linewidth - 8\tabcolsep) * \real{0.2000}}
  >{\raggedright\arraybackslash}p{(\linewidth - 8\tabcolsep) * \real{0.2000}}@{}}
\toprule\noalign{}
\begin{minipage}[b]{\linewidth}\raggedright
Formation
\end{minipage} & \begin{minipage}[b]{\linewidth}\raggedright
Candidate input
\end{minipage} & \begin{minipage}[b]{\linewidth}\raggedright
Expected or documented OE outcome
\end{minipage} & \begin{minipage}[b]{\linewidth}\raggedright
OE comparison form
\end{minipage} & \begin{minipage}[b]{\linewidth}\raggedright
Result
\end{minipage} \\
\midrule\noalign{}
\endhead
\bottomrule\noalign{}
\endlastfoot
broader root-noun family &
{\emph{*brústz}}\index[iv]{03pgmc@\ivlangheader{Proto-Germanic}{}!brustz@\iventry{*brústz}{}}
& root-noun type outcomes outside OE & non-OE comparanda & useful family
label, but not the direct source of \emph{brēost} \\
selected thematic formation &
{\emph{*bréustą}}\index[iv]{03pgmc@\ivlangheader{Proto-Germanic}{}!breusta@\iventry{*bréustą}{}}
& regular output: \emph{brēost} &
{\emph{brēost}}\index[iv]{02oe@\ivlangheader{Old English}{West Saxon normalization unmarked}!breost@\iventry{brēost}{}}
& exact match between formation and attested OE noun \\
\end{longtable}
}

The relevant point is the formation split.
{\emph{brēost}}\index[iv]{02oe@\ivlangheader{Old English}{West Saxon normalization unmarked}!breost@\iventry{brēost}{}}
`breast' is the regular Old English outcome of the thematic
{\emph{*breusta-}}\index[iv]{03pgmc@\ivlangheader{Proto-Germanic}{}!breusta@\iventry{*breusta-}{}}
branch, not of the root noun
{\emph{*brust-}}\index[iv]{03pgmc@\ivlangheader{Proto-Germanic}{}!brust@\iventry{*brust-}{}}.

\section{\texorpdfstring{craft --- OE
\emph{cræft}}{craft --- OE cræft}}\label{craft-oe-cruxe6ft}

\index[iv]{02oe@\ivlangheader{Old English}{West Saxon normalization unmarked}!craeft@\iventry{cræft}{}}
\index[iv]{02oe@\ivlangheader{Old English}{West Saxon normalization unmarked}!craft@\iventry{craft}{}}
\index[iv]{02oe@\ivlangheader{Old English}{West Saxon normalization unmarked}!creft@\iventry{creft}{}}
\index[iv]{03pgmc@\ivlangheader{Proto-Germanic}{}!kraftaz@\iventry{*kráftaz}{}}
\index[iv]{03pgmc@\ivlangheader{Proto-Germanic}{}!kraftiz@\iventry{*kráftiz}{}}

Derivation: citation reconstruction \emph{*kráftiz}; form followed here
\emph{*kráftaz} \textgreater{} \emph{cræft} (early analogy).

\subsection{Derivation trace}\label{derivation-trace-106}

Proto input: \emph{*kráftaz}

\begingroup
\setlength{\fboxsep}{6pt}

\noindent\fbox{%
\begin{minipage}{0.97\linewidth}
\small
\begin{tabular}{@{}>{\raggedright\arraybackslash}p{0.485\linewidth}@{\hspace{0.030\linewidth}}>{\raggedright\arraybackslash}p{0.485\linewidth}@{}}
\begin{minipage}[t]{\linewidth}
\centering\textbf{Earlier Germanic changes}\par
\vspace{0.35em}
\raggedright
\centering\textbf{Northwest and West Germanic}\par
\raggedright
\vspace{0.2em}
\begin{tabularx}{\linewidth}{@{}>{\raggedright\arraybackslash}X@{\hspace{0.55em}}>{\raggedright\arraybackslash}p{0.14\linewidth}@{\hspace{0.25em}}}
\mbox{EAF Final Z Deletion} & \emph{*kráfta} \\
\end{tabularx}
\vspace{0.6em}
\centering\textbf{Early Anglo-Frisian}\par
\raggedright
\vspace{0.2em}
\raggedright [no change]\par
\end{minipage}
&
\begin{minipage}[t]{\linewidth}
\centering\textbf{Old English changes}\par
\vspace{0.35em}
\raggedright
\begin{tabularx}{\linewidth}{@{}>{\raggedright\arraybackslash}X@{\hspace{0.55em}}>{\raggedright\arraybackslash}p{0.14\linewidth}@{\hspace{0.25em}}}
\mbox{PWGmc Final Bare A Loss} & \emph{*kráft} \\
\mbox{EAF Brightening} & \emph{*kræft} \\
\end{tabularx}
\end{minipage}
\\
\end{tabular}
\end{minipage}%
} \endgroup

Old English form: \emph{cræft}

\subsection{Reconstruction and comparative
evidence}\label{reconstruction-and-comparative-evidence-90}

Comparative sources disagree about the older stem class. Kroonen gives a
u-stem \emph{*kraftu-}, while Orel prints \Recon{kraftiz} `craft' and
\Recon{kraftuz} `craft' (\citeproc{ref-Kroonen2013}{Kroonen 2013, 340};
\citeproc{ref-Orel2003}{Orel 2003, 259}). The comparative label
\Recon{kráftiz} `craft' remains in view as a lexeme-level shorthand,
while \Recon{kráftaz} `craft' is the pre-Old-English form used for the
Old English derivation.

\subsection{Old English evidence}\label{old-english-evidence-91}

The Old English noun itself is secure. Clark Hall and Bosworth-Toller
both give \emph{cræft} `craft' as the headword
(\citeproc{ref-ClarkHall1960}{Clark Hall 1960, 19};
\citeproc{ref-BosworthToller1898}{Bosworth and Toller 1898, 145}).

\subsection{Development to Old
English}\label{development-to-old-english-89}

The comparison is between possible pre-Old-English inputs. The i-stem
comparator \Recon{kráftiz} `craft' gives \emph{creft}, while the u-stem
comparator \Recon{kráftuz} `craft' gives \emph{craft}. The a-stem-shaped
input \Recon{kráftaz} `craft' yields \emph{cræft} `craft' and is
therefore the form used for the Old English derivation. This does not
require treating the comparative dictionaries as identical; it shows the
narrower point that the Old English derivation needs a pre-Old-English
form without the i-umlaut trigger of \Recon{-iz} `craft' and without the
back-vowel outcome associated with the u-stem comparator.

\subsection{Form comparison}\label{form-comparison-4}

{\def\LTcaptype{none} % do not increment counter
\begin{longtable}[]{@{}lll@{}}
\toprule\noalign{}
Candidate input & OE output & Result \\
\midrule\noalign{}
\endhead
\bottomrule\noalign{}
\endlastfoot
*kráftiz & \emph{creft} & non-match; i-stem comparator \\
*kráftuz & \emph{craft} & non-match; u-stem comparator \\
*kráftaz & \emph{cræft} & exact match; selected pre-OE input \\
\end{longtable}
}

\section{\texorpdfstring{dill --- OE
\emph{dile}}{dill --- OE dile}}\label{dill-oe-dile}

\index[iv]{02oe@\ivlangheader{Old English}{West Saxon normalization unmarked}!dile@\iventry{dile}{}}
\index[iv]{03pgmc@\ivlangheader{Proto-Germanic}{}!deliz@\iventry{*déliz}{}}
\index[iv]{03pgmc@\ivlangheader{Proto-Germanic}{}!deljaz@\iventry{*déljaz}{}}

Derivation: citation reconstruction \emph{*déljaz}; form followed here
\emph{*déliz} \textgreater{} \emph{dile} (early analogy).

\subsection{Derivation trace}\label{derivation-trace-107}

Proto input: \emph{*déliz}

\begingroup
\setlength{\fboxsep}{6pt}

\noindent\fbox{%
\begin{minipage}{0.97\linewidth}
\small
\begin{tabular}{@{}>{\raggedright\arraybackslash}p{0.359\linewidth}@{\hspace{0.030\linewidth}}>{\raggedright\arraybackslash}p{0.611\linewidth}@{}}
\begin{minipage}[t]{\linewidth}
\centering\textbf{Earlier Germanic changes}\par
\vspace{0.35em}
\raggedright
\centering\textbf{Northwest and West Germanic}\par
\raggedright
\vspace{0.2em}
\begin{tabularx}{\linewidth}{@{}>{\raggedright\arraybackslash}X@{\hspace{0.55em}}>{\raggedright\arraybackslash}p{0.14\linewidth}@{\hspace{0.25em}}}
\mbox{EAF Final Z Deletion} & \emph{*déli} \\
\end{tabularx}
\vspace{0.6em}
\centering\textbf{Early Anglo-Frisian}\par
\raggedright
\vspace{0.2em}
\raggedright [no change]\par
\end{minipage}
&
\begin{minipage}[t]{\linewidth}
\centering\textbf{Old English changes}\par
\vspace{0.35em}
\raggedright
\begin{tabularx}{\linewidth}{@{}>{\raggedright\arraybackslash}X@{\hspace{0.55em}}>{\raggedright\arraybackslash}p{0.14\linewidth}@{\hspace{0.25em}}}
\mbox{OE I Umlaut} & \emph{*dili} \\
OE Med Unstressed I Lowering1 & \emph{*dile} \\
\end{tabularx}
\end{minipage}
\\
\end{tabular}
\end{minipage}%
} \endgroup

Old English form: \emph{dile}

\subsection{Reconstruction and comparative
evidence}\label{reconstruction-and-comparative-evidence-91}

Comparative evidence preserves both an i-stem and a ja-stem formation,
with Old English
{\emph{dile}}\index[iv]{02oe@\ivlangheader{Old English}{West Saxon normalization unmarked}!dile@\iventry{dile}{}}
`dill' on one side and continental forms such as Old Saxon
{\emph{dilli}}\index[iv]{13os@\ivlangheader{Old Saxon}{}!dilli@\iventry{dilli}{}}
`dill' and Old High German
{\emph{tilli}}\index[iv]{10ohg@\ivlangheader{Old High German}{}!tilli@\iventry{tilli}{}}
`dill' on the other (\citeproc{ref-Fulk2018}{Fulk 2018, 170}). The
derivational input
\Recon{déliz}\index[iv]{03pgmc@\ivlangheader{Proto-Germanic}{}!deliz@\iventry{*déliz}{}}
`dill' therefore represents the i-stem side of the paradigm, whereas the
citation label
\Recon{déljaz}\index[iv]{03pgmc@\ivlangheader{Proto-Germanic}{}!deljaz@\iventry{*déljaz}{}}
`dill' is a broader comparative headword.

The stem class determines the Old English consonant shape. A ja-stem
with \emph{*-lj-} would be expected to produce gemination, but the Old
English noun shows a single \emph{l}. Fulk's discussion of ja-stems
transferred to the i-stems provides the relevant morphological
background for the OE side (\citeproc{ref-Fulk2018}{Fulk 2018, 170}).

\subsection{Old English evidence}\label{old-english-evidence-92}

Old English dictionaries record the plant name as
{\emph{dile}}\index[iv]{02oe@\ivlangheader{Old English}{West Saxon normalization unmarked}!dile@\iventry{dile}{}}
`dill', alongside the variant
{\emph{dili}}\index[iv]{02oe@\ivlangheader{Old English}{West Saxon normalization unmarked}!dili@\iventry{dili}{}}
`dill' (\citeproc{ref-BosworthToller1898}{Bosworth and Toller 1898,
164}; \citeproc{ref-ClarkHall1960}{Clark Hall 1960, 95}). The form
discussed here is therefore an attested Old English noun with single
\emph{l}.

The Old English evidence is the relevant point. Whatever broader
comparative headword is chosen for the family, the inherited form
reflected in OE is the i-stem type
{\emph{dile}}\index[iv]{02oe@\ivlangheader{Old English}{West Saxon normalization unmarked}!dile@\iventry{dile}{}}
`dill', not a geminated
{\emph{dill}}\index[iv]{02oe@\ivlangheader{Old English}{West Saxon normalization unmarked}!dill@\iventry{dill}{}}
`dill' outcome.

\subsection{Development to Old
English}\label{development-to-old-english-90}

From
\Recon{déliz}\index[iv]{03pgmc@\ivlangheader{Proto-Germanic}{}!deliz@\iventry{*déliz}{}}
`dill', regular loss of final \emph{z} and the later lowering of
unstressed \emph{i} yield
{\emph{dile}}\index[iv]{02oe@\ivlangheader{Old English}{West Saxon normalization unmarked}!dile@\iventry{dile}{}}
`dill'. The stem itself remains ungeminated throughout that path.

The contrast is morphological rather than phonological. If the word were
forced through a ja-stem \emph{*-lj-} pathway, the expected result would
show \emph{ll}. The attested Old English noun instead matches the i-stem
development.

\subsection{Formation comparison}\label{formation-comparison-1}

The comparison below sets the relevant forms side by side. It separates
the broader comparative headword from the stem class actually reflected
in Old English.

{\def\LTcaptype{none} % do not increment counter
\begin{longtable}[]{@{}
  >{\raggedright\arraybackslash}p{(\linewidth - 8\tabcolsep) * \real{0.2000}}
  >{\raggedright\arraybackslash}p{(\linewidth - 8\tabcolsep) * \real{0.2000}}
  >{\raggedright\arraybackslash}p{(\linewidth - 8\tabcolsep) * \real{0.2000}}
  >{\raggedright\arraybackslash}p{(\linewidth - 8\tabcolsep) * \real{0.2000}}
  >{\raggedright\arraybackslash}p{(\linewidth - 8\tabcolsep) * \real{0.2000}}@{}}
\toprule\noalign{}
\begin{minipage}[b]{\linewidth}\raggedright
Formation
\end{minipage} & \begin{minipage}[b]{\linewidth}\raggedright
Candidate input
\end{minipage} & \begin{minipage}[b]{\linewidth}\raggedright
Expected or documented OE outcome
\end{minipage} & \begin{minipage}[b]{\linewidth}\raggedright
OE comparison form
\end{minipage} & \begin{minipage}[b]{\linewidth}\raggedright
Result
\end{minipage} \\
\midrule\noalign{}
\endhead
\bottomrule\noalign{}
\endlastfoot
comparative ja-stem label &
{\emph{*déljaz}}\index[iv]{03pgmc@\ivlangheader{Proto-Germanic}{}!deljaz@\iventry{*déljaz}{}}
& ja-stem type outcome with gemination &
{\emph{dill}}\index[iv]{02oe@\ivlangheader{Old English}{West Saxon normalization unmarked}!dill@\iventry{dill}{}}-type
comparison & useful comparative label, but not the OE form \\
selected i-stem formation &
{\emph{*déliz}}\index[iv]{03pgmc@\ivlangheader{Proto-Germanic}{}!deliz@\iventry{*déliz}{}}
& regular output: {\emph{dile}} &
{\emph{dile}}\index[iv]{02oe@\ivlangheader{Old English}{West Saxon normalization unmarked}!dile@\iventry{dile}{}}
& exact match between formation and attested OE noun \\
\end{longtable}
}

The single \emph{l} is the decisive diagnostic. It identifies
{\emph{dile}}\index[iv]{02oe@\ivlangheader{Old English}{West Saxon normalization unmarked}!dile@\iventry{dile}{}}
`dill' with the i-stem formation rather than with the continental
ja-stem branch.

\section{\texorpdfstring{fast --- OE
\emph{festan}}{fast --- OE festan}}\label{fast-oe-festan}

\index[iv]{02oe@\ivlangheader{Old English}{West Saxon normalization unmarked}!festan@\iventry{festan}{}}
\index[iv]{03pgmc@\ivlangheader{Proto-Germanic}{}!fastena@\iventry{*fastēną}{}}
\index[iv]{03pgmc@\ivlangheader{Proto-Germanic}{}!fastijana@\iventry{*fástijaną}{}}

Derivation: citation reconstruction \emph{*fastēną}; form followed here
\emph{*fástijaną} \textgreater{} \emph{festan} (early analogy).

\subsection{Derivation trace}\label{derivation-trace-108}

Proto input: \emph{*fástijaną}

\begingroup
\setlength{\fboxsep}{6pt}

\noindent\fbox{%
\begin{minipage}{0.97\linewidth}
\footnotesize
\begin{tabular}{@{}>{\raggedright\arraybackslash}p{0.339\linewidth}@{\hspace{0.030\linewidth}}>{\raggedright\arraybackslash}p{0.630\linewidth}@{}}
\begin{minipage}[t]{\linewidth}
\centering\textbf{Earlier Germanic changes}\par
\vspace{0.35em}
\raggedright
\centering\textbf{Northwest and West Germanic}\par
\raggedright
\vspace{0.2em}
\raggedright [no change]\par
\vspace{0.6em}
\centering\textbf{Early Anglo-Frisian}\par
\raggedright
\vspace{0.2em}
\raggedright [no change]\par
\end{minipage}
&
\begin{minipage}[t]{\linewidth}
\centering\textbf{Old English changes}\par
\vspace{0.35em}
\raggedright
\begin{tabularx}{\linewidth}{@{}>{\raggedright\arraybackslash}X@{\hspace{0.55em}}>{\raggedright\arraybackslash}p{0.24\linewidth}@{\hspace{0.25em}}}
\mbox{EAF Brightening} & \emph{*fæstijaną} \\
OE Heavy Syllable Nasal Apocope & \emph{*fæstijan} \\
\mbox{OE Secondary Nasalization} & \emph{*fæstijąn} \\
\mbox{Sievers Law Syncope} & \emph{*fæstjąn} \\
\mbox{OE I Umlaut} & \emph{*festjąn} \\
\mbox{OE Weak Tail Reduction} & \emph{*festjan} \\
\mbox{OE J Loss After Heavy} & \emph{*festan} \\
\end{tabularx}
\end{minipage}
\\
\end{tabular}
\end{minipage}%
} \endgroup

Old English form: \emph{festan}

\subsection{Reconstruction and comparative
evidence}\label{reconstruction-and-comparative-evidence-92}

Kroonen places the verb within the wider
{\emph{*fastu-}}\index[iv]{03pgmc@\ivlangheader{Proto-Germanic}{}!fastu@\iventry{*fastu-}{}}
adjective family and its derived
{\emph{*fasten-}}\index[iv]{03pgmc@\ivlangheader{Proto-Germanic}{}!fasten@\iventry{*fasten-}{}}
verbal line, the comparative background behind Old English `to fast'
(\citeproc{ref-Kroonen2013}{Kroonen 2013, 171}). Ringe and Taylor,
however, distinguish the Old English verb more closely: they treat OE
`to fast' as originally a class-I weak verb that later acquired the
stative meaning through lexical confusion
(\citeproc{ref-RingeTaylor2014}{Ringe and Taylor 2014, 110}).

The derivational input
\Recon{fástijaną}\index[iv]{03pgmc@\ivlangheader{Proto-Germanic}{}!fastijana@\iventry{*fástijaną}{}}
`fast' therefore represents the inherited class-I formation reflected in
Old English, whereas the citation label
\Recon{fastēną}\index[iv]{03pgmc@\ivlangheader{Proto-Germanic}{}!fastena@\iventry{*fastēną}{}}
`fast' belongs to the broader comparative presentation of the lexeme.

\subsection{Old English evidence}\label{old-english-evidence-93}

Old English dictionaries record forms such as
{\emph{festan}}\index[iv]{02oe@\ivlangheader{Old English}{West Saxon normalization unmarked}!festan@\iventry{festan}{}}
`fast', alongside related
{\emph{fæstan}}\index[iv]{02oe@\ivlangheader{Old English}{West Saxon normalization unmarked}!faestan@\iventry{fæstan}{}}
`fast' /
{\emph{fǣstan}}\index[iv]{02oe@\ivlangheader{Old English}{West Saxon normalization unmarked}!faestan@\iventry{fǣstan}{}}
`fast' spellings and meanings
(\citeproc{ref-BosworthToller1898}{Bosworth and Toller 1898, 213}). The
form selected here is
{\emph{festan}}\index[iv]{02oe@\ivlangheader{Old English}{West Saxon normalization unmarked}!festan@\iventry{festan}{}},
which fits the regular class-I phonological development.

The \emph{æ}-forms remain relevant, but they do not control the entry.
In the present analysis they belong to a later analogical reshaping
under the adjective
{\emph{fæst}}\index[iv]{02oe@\ivlangheader{Old English}{West Saxon normalization unmarked}!faest@\iventry{fæst}{}}
`fast', whereas
{\emph{festan}}\index[iv]{02oe@\ivlangheader{Old English}{West Saxon normalization unmarked}!festan@\iventry{festan}{}}
`fast' is the regular inherited class-I comparison form.

\subsection{Development to Old
English}\label{development-to-old-english-91}

From
\Recon{fástijaną}\index[iv]{03pgmc@\ivlangheader{Proto-Germanic}{}!fastijana@\iventry{*fástijaną}{}}
`fast', Anglo-Frisian brightening and subsequent i-umlaut produce the
fronted vowel seen in
{\emph{festan}}\index[iv]{02oe@\ivlangheader{Old English}{West Saxon normalization unmarked}!festan@\iventry{festan}{}}
`fast'. The later weak-tail reductions and loss of \emph{j} after a
heavy syllable complete the regular Old English outcome.

What makes the entry non-regular is not the phonology of
{\emph{festan}}\index[iv]{02oe@\ivlangheader{Old English}{West Saxon normalization unmarked}!festan@\iventry{festan}{}}
`fast' itself, but the choice of formation. Old English continues the
class-I verb, even though the comparative headword is often given under
the parallel
{\emph{*fastēn-}}\index[iv]{03pgmc@\ivlangheader{Proto-Germanic}{}!fasten@\iventry{*fastēn-}{}}
family.

\subsection{Class comparison}\label{class-comparison}

The comparison below sets the relevant forms side by side. It
distinguishes the comparative class-III headword from the class-I
formation actually reflected in Old English.

{\def\LTcaptype{none} % do not increment counter
\begin{longtable}[]{@{}
  >{\raggedright\arraybackslash}p{(\linewidth - 8\tabcolsep) * \real{0.2000}}
  >{\raggedright\arraybackslash}p{(\linewidth - 8\tabcolsep) * \real{0.2000}}
  >{\raggedright\arraybackslash}p{(\linewidth - 8\tabcolsep) * \real{0.2000}}
  >{\raggedright\arraybackslash}p{(\linewidth - 8\tabcolsep) * \real{0.2000}}
  >{\raggedright\arraybackslash}p{(\linewidth - 8\tabcolsep) * \real{0.2000}}@{}}
\toprule\noalign{}
\begin{minipage}[b]{\linewidth}\raggedright
Formation / class
\end{minipage} & \begin{minipage}[b]{\linewidth}\raggedright
Candidate input
\end{minipage} & \begin{minipage}[b]{\linewidth}\raggedright
Expected or documented OE outcome
\end{minipage} & \begin{minipage}[b]{\linewidth}\raggedright
OE comparison form
\end{minipage} & \begin{minipage}[b]{\linewidth}\raggedright
Result
\end{minipage} \\
\midrule\noalign{}
\endhead
\bottomrule\noalign{}
\endlastfoot
comparative class-III headword &
{\emph{*fastēną}}\index[iv]{03pgmc@\ivlangheader{Proto-Germanic}{}!fastena@\iventry{*fastēną}{}}
& class-III type outcome, not
{\emph{festan}}\index[iv]{02oe@\ivlangheader{Old English}{West Saxon normalization unmarked}!festan@\iventry{festan}{}}
& wider family context & useful family label, but not the direct source
of the target \\
selected class-I weak verb &
{\emph{*fástijaną}}\index[iv]{03pgmc@\ivlangheader{Proto-Germanic}{}!fastijana@\iventry{*fástijaną}{}}
& regular output: {\emph{festan}} &
{\emph{festan}}\index[iv]{02oe@\ivlangheader{Old English}{West Saxon normalization unmarked}!festan@\iventry{festan}{}}
& exact match between formation and attested OE verb \\
later analogical reshaping & adjective-driven
{\emph{fæst}}\index[iv]{02oe@\ivlangheader{Old English}{West Saxon normalization unmarked}!faest@\iventry{fæst}{}}
influence &
{\emph{fæstan}}\index[iv]{02oe@\ivlangheader{Old English}{West Saxon normalization unmarked}!faestan@\iventry{fæstan}{}}
/
{\emph{fǣstan}}\index[iv]{02oe@\ivlangheader{Old English}{West Saxon normalization unmarked}!faestan@\iventry{fǣstan}{}}
type spellings & fæstan-type evidence & genuine later OE reshaping, but
secondary to the Old English form here \\
\end{longtable}
}

The relevant point is the class split.
{\emph{festan}}\index[iv]{02oe@\ivlangheader{Old English}{West Saxon normalization unmarked}!festan@\iventry{festan}{}}
`fast' is the regular Old English outcome of the class-I formation,
while the better-known \emph{æ}-forms belong to a later analogical
layer.

\section{\texorpdfstring{fee --- OE
\emph{feoh}}{fee --- OE feoh}}\label{fee-oe-feoh}

\index[iv]{02oe@\ivlangheader{Old English}{West Saxon normalization unmarked}!feoh@\iventry{feoh}{}}
\index[iv]{03pgmc@\ivlangheader{Proto-Germanic}{}!fexu@\iventry{*féxu}{}}
\index[iv]{08preoe@\ivlangheader{Pre-Old English / prehistoric English}{}!fex~01@\iventry{*féx}{transp.}}

Derivation: citation reconstruction \emph{*féxu}; form followed here
\emph{*féx} \textgreater{} \emph{feoh} (early analogy).

\subsection{Derivation trace}\label{derivation-trace-109}

Proto input: \emph{*féx}

\begingroup
\setlength{\fboxsep}{6pt}

\noindent\fbox{%
\begin{minipage}{0.97\linewidth}
\small
\begin{tabular}{@{}>{\raggedright\arraybackslash}p{0.485\linewidth}@{\hspace{0.030\linewidth}}>{\raggedright\arraybackslash}p{0.485\linewidth}@{}}
\begin{minipage}[t]{\linewidth}
\centering\textbf{Earlier Germanic changes}\par
\vspace{0.35em}
\raggedright
\centering\textbf{Northwest and West Germanic}\par
\raggedright
\vspace{0.2em}
\raggedright [no change]\par
\vspace{0.6em}
\centering\textbf{Early Anglo-Frisian}\par
\raggedright
\vspace{0.2em}
\raggedright [no change]\par
\end{minipage}
&
\begin{minipage}[t]{\linewidth}
\centering\textbf{Old English changes}\par
\vspace{0.35em}
\raggedright
\begin{tabularx}{\linewidth}{@{}>{\raggedright\arraybackslash}X@{\hspace{0.55em}}>{\raggedright\arraybackslash}p{0.14\linewidth}@{\hspace{0.25em}}}
\mbox{OE Breaking} & \emph{*féox} \\
\end{tabularx}
\end{minipage}
\\
\end{tabular}
\end{minipage}%
} \endgroup

Old English form: \emph{feoh}

\subsection{Reconstruction and comparative
evidence}\label{reconstruction-and-comparative-evidence-93}

The inherited comparative lexeme is the neuter u-stem
\Recon{féxu}\index[iv]{03pgmc@\ivlangheader{Proto-Germanic}{}!fexu@\iventry{*féxu}{}}
`cattle, wealth'. Kroonen gives the u-stem citation \emph{fehu}, and
Orel gives \emph{fexu}, with a separately attached \emph{fexan}
antecedent for Old English (\citeproc{ref-Kroonen2013}{Kroonen 2013,
134}; \citeproc{ref-Orel2003}{Orel 2003, 97}). The latter notation is
preserved as an alternative source formation, not silently reinterpreted
as proof of an inherited n-stem.

\subsection{Old English evidence}\label{old-english-evidence-94}

Old English
{\emph{feoh}}\index[iv]{02oe@\ivlangheader{Old English}{West Saxon normalization unmarked}!feoh@\iventry{feoh}{}}
`cattle, wealth' follows the a-declension. Campbell's strong paradigm
distinguishes the nominative-accusative from its obliques; Brunner
explicitly classifies the word with the a-stems
(\citeproc{ref-Campbell1959}{Campbell 1959, 225--26}, §574;
\citeproc{ref-SieversBrunner1965}{Brunner 1965, 221}, §275). The cell
compared here remains nominative-accusative singular, not genitive
\emph{feos}.

\subsection{Early transfer and selected
input}\label{early-transfer-and-selected-input}

Ringe and Taylor argue that the u-stem was transferred early to the
a-stems, before the relevant intervocalic h-loss. Their transferred
nominative-accusative is already endingless
\Recon{féx}\index[iv]{08preoe@\ivlangheader{Pre-Old English / prehistoric English}{}!fex~01@\iventry{*féx}{transp.}}
`cattle, wealth', representing their \emph{feh} in the notation used
here (\citeproc{ref-RingeTaylor2014}{Ringe and Taylor 2014, 177}, n.~4).
The inherited u-stem citation and this later transponent therefore have
different historical stages; the a-stem is not reconstructed as
inherited Proto-Germanic.

\subsection{Development to Old
English}\label{development-to-old-english-92}

The transferred \emph{feh} precedes breaking, which yields \emph{feoh}.
Its endingless cell compared here supplies no intervocalic high vowel to
be lost or contracted. Ringe and Taylor contrast that history with the
inherited u-stem's problematic intervening-vowel development
(\citeproc{ref-RingeTaylor2014}{Ringe and Taylor 2014, 177}, n.~4). The
relative date is bounded by breaking and intervocalic h-loss, not an
independently established absolute year.

\subsection{Previous selection and competing
histories}\label{previous-selection-and-competing-histories}

The previous selection used the inherited \emph{féxu} directly. That
input can give \emph{feoh} under a conventional implementation with an
additional post-h deletion, but gives the unaccepted prediction
\emph{feou} under quantitative high-vowel deletion. The source-supported
transfer, not restoration of a match count, motivates the new input. An
a-stem carrier \emph{féxą} also gives \emph{feoh} conditionally, but is
unnecessary: the cited later formation already lacks an ending. The
inherited comparative headword and the attested cell compared here
remain unchanged.

\section{\texorpdfstring{flask --- OE
\emph{flasce}}{flask --- OE flasce}}\label{flask-oe-flasce}

\index[iv]{02oe@\ivlangheader{Old English}{West Saxon normalization unmarked}!flasce@\iventry{flasce}{}}
\index[iv]{03pgmc@\ivlangheader{Proto-Germanic}{}!flasko@\iventry{*flaskō}{}}
\index[iv]{03pgmc@\ivlangheader{Proto-Germanic}{}!flaskon@\iventry{*fláskōn}{}}

Derivation: citation reconstruction \emph{*flaskō}; form followed here
\emph{*fláskōn} \textgreater{} \emph{flasce} (early analogy).

\subsection{Derivation trace}\label{derivation-trace-110}

Proto input: \emph{*fláskōn}

\begingroup
\setlength{\fboxsep}{6pt}

\noindent\fbox{%
\begin{minipage}{0.97\linewidth}
\small
\begin{tabular}{@{}>{\raggedright\arraybackslash}p{0.359\linewidth}@{\hspace{0.030\linewidth}}>{\raggedright\arraybackslash}p{0.611\linewidth}@{}}
\begin{minipage}[t]{\linewidth}
\centering\textbf{Earlier Germanic changes}\par
\vspace{0.35em}
\raggedright
\centering\textbf{Northwest and West Germanic}\par
\raggedright
\vspace{0.2em}
\begin{tabularx}{\linewidth}{@{}>{\raggedright\arraybackslash}X@{\hspace{0.55em}}>{\raggedright\arraybackslash}p{0.24\linewidth}@{\hspace{0.25em}}}
\mbox{PNWGmc N Stem N Loss} & \emph{*fláskǭ} \\
\end{tabularx}
\vspace{0.6em}
\centering\textbf{Early Anglo-Frisian}\par
\raggedright
\vspace{0.2em}
\raggedright [no change]\par
\end{minipage}
&
\begin{minipage}[t]{\linewidth}
\centering\textbf{Old English changes}\par
\vspace{0.35em}
\raggedright
\begin{tabularx}{\linewidth}{@{}>{\raggedright\arraybackslash}X@{\hspace{0.55em}}>{\raggedright\arraybackslash}p{0.14\linewidth}@{\hspace{0.25em}}}
\mbox{EAF Brightening} & \emph{*flæskǭ} \\
\mbox{OE A Restoration} & \emph{*flaskǭ} \\
OE Unstressed Long Vowel Shortening & \emph{*flaskæ} \\
\mbox{OE Unstressed AE Merger} & \emph{*flaske} \\
\end{tabularx}
\end{minipage}
\\
\end{tabular}
\end{minipage}%
} \endgroup

Old English form: \emph{flasce}

\subsection{Reconstruction and comparative
evidence}\label{reconstruction-and-comparative-evidence-94}

The wider Germanic family is often cited under a form such as
\Recon{flaskō}\index[iv]{03pgmc@\ivlangheader{Proto-Germanic}{}!flasko@\iventry{*flaskō}{}}
`flask', but the evidence relevant for Old English points instead to a
weak feminine formation
\Recon{fláskōn}\index[iv]{03pgmc@\ivlangheader{Proto-Germanic}{}!flaskon@\iventry{*fláskōn}{}}
`flask' /
\Recon{flaskǭ}\index[iv]{03pgmc@\ivlangheader{Proto-Germanic}{}!flasko@\iventry{*flaskǭ}{}}
`flask' (\citeproc{ref-Orel2003}{Orel 2003, 104}). That distinction is
crucial for the suffixal history of the noun.

The derivational input therefore differs from the citation label in stem
class. Old English
{\emph{flasce}}\index[iv]{02oe@\ivlangheader{Old English}{West Saxon normalization unmarked}!flasce@\iventry{flasce}{}}
`flask' belongs with the weak feminine line, and the plural or oblique
forms
{\emph{flascan}}\index[iv]{02oe@\ivlangheader{Old English}{West Saxon normalization unmarked}!flascan@\iventry{flascan}{}}
`flask' support that analysis (\citeproc{ref-RingeTaylor2014}{Ringe and
Taylor 2014, 192}).

\subsection{Old English evidence}\label{old-english-evidence-95}

Old English dictionaries record the noun as
{\emph{flasce}}\index[iv]{02oe@\ivlangheader{Old English}{West Saxon normalization unmarked}!flasce@\iventry{flasce}{}}
`flask', with inflectional support from forms such as
{\emph{flascan}}\index[iv]{02oe@\ivlangheader{Old English}{West Saxon normalization unmarked}!flascan@\iventry{flascan}{}}
`flask'; a later West Saxon
{\emph{flaxe}}\index[iv]{02oe@\ivlangheader{Old English}{West Saxon normalization unmarked}!flaxe@\iventry{flaxe}{}}
`flask' is also noted as a secondary variant
(\citeproc{ref-BosworthToller1898}{Bosworth and Toller 1898, 235};
\citeproc{ref-ClarkHall1960}{Clark Hall 1960, 121}).

The relevant comparison form is therefore the weak feminine noun
{\emph{flasce}}\index[iv]{02oe@\ivlangheader{Old English}{West Saxon normalization unmarked}!flasce@\iventry{flasce}{}}
`flask'. The plural and oblique evidence helps explain why the vowel and
ending are preserved as they are in the singular.

\subsection{Development to Old
English}\label{development-to-old-english-93}

From
\Recon{fláskōn}\index[iv]{03pgmc@\ivlangheader{Proto-Germanic}{}!flaskon@\iventry{*fláskōn}{}}
`flask', the weak feminine passes through the expected loss of \emph{n}
and the later Old English development of the unstressed ending, reaching
{\emph{flasce}}\index[iv]{02oe@\ivlangheader{Old English}{West Saxon normalization unmarked}!flasce@\iventry{flasce}{}}
`flask'. Campbell cites restored \emph{a} in exactly this environment,
including
{\emph{flasce}}\index[iv]{02oe@\ivlangheader{Old English}{West Saxon normalization unmarked}!flasce@\iventry{flasce}{}}
`flask' after inflected
{\emph{flascan}}\index[iv]{02oe@\ivlangheader{Old English}{West Saxon normalization unmarked}!flascan@\iventry{flascan}{}}
`flask' (\citeproc{ref-Campbell1959}{Campbell 1959, § 158}). Once
the weak feminine formation is chosen, the noun follows a regular path
to its Old English shape.

The decisive issue is morphological rather than phonological. A simple
strong feminine citation form does not capture the OE weak noun as
cleanly as the selected \Recon{fláskōn} `flask' does.

\subsection{Formation comparison}\label{formation-comparison-2}

The comparison below sets the relevant forms side by side. It separates
the broader comparative headword from the weak feminine formation
actually reflected in Old English.

{\def\LTcaptype{none} % do not increment counter
\begin{longtable}[]{@{}
  >{\raggedright\arraybackslash}p{(\linewidth - 8\tabcolsep) * \real{0.2000}}
  >{\raggedright\arraybackslash}p{(\linewidth - 8\tabcolsep) * \real{0.2000}}
  >{\raggedright\arraybackslash}p{(\linewidth - 8\tabcolsep) * \real{0.2000}}
  >{\raggedright\arraybackslash}p{(\linewidth - 8\tabcolsep) * \real{0.2000}}
  >{\raggedright\arraybackslash}p{(\linewidth - 8\tabcolsep) * \real{0.2000}}@{}}
\toprule\noalign{}
\begin{minipage}[b]{\linewidth}\raggedright
Formation
\end{minipage} & \begin{minipage}[b]{\linewidth}\raggedright
Candidate input
\end{minipage} & \begin{minipage}[b]{\linewidth}\raggedright
Expected or documented OE outcome
\end{minipage} & \begin{minipage}[b]{\linewidth}\raggedright
OE comparison form
\end{minipage} & \begin{minipage}[b]{\linewidth}\raggedright
Result
\end{minipage} \\
\midrule\noalign{}
\endhead
\bottomrule\noalign{}
\endlastfoot
broader comparative headword &
{\emph{*flaskō}}\index[iv]{03pgmc@\ivlangheader{Proto-Germanic}{}!flasko@\iventry{*flaskō}{}}
& broader family label & wider family context & useful lexeme label, but
not the cleanest OE-facing derivation \\
selected weak feminine formation &
{\emph{*fláskōn}}\index[iv]{03pgmc@\ivlangheader{Proto-Germanic}{}!flaskon@\iventry{*fláskōn}{}}
& regular output: {\emph{flasce}} &
{\emph{flasce}}\index[iv]{02oe@\ivlangheader{Old English}{West Saxon normalization unmarked}!flasce@\iventry{flasce}{}}
& exact match between formation and attested OE noun \\
\end{longtable}
}

The weak feminine suffix is the relevant point. It aligns the inherited
form with attested
{\emph{flasce}}\index[iv]{02oe@\ivlangheader{Old English}{West Saxon normalization unmarked}!flasce@\iventry{flasce}{}}
`flask' and its supporting paradigm forms.

\section{\texorpdfstring{follow --- OE
\emph{fylġan}}{follow --- OE fylġan}}\label{follow-oe-fylux121an}

\index[iv]{02oe@\ivlangheader{Old English}{West Saxon normalization unmarked}!fylgan@\iventry{fylġan}{}}
\index[iv]{03pgmc@\ivlangheader{Proto-Germanic}{}!fulgena@\iventry{*fulgēną}{}}
\index[iv]{03pgmc@\ivlangheader{Proto-Germanic}{}!fulgijana@\iventry{*fúlgijaną}{}}

Derivation: citation reconstruction \emph{*fulgēną}; form followed here
\emph{*fúlgijaną} \textgreater{} \emph{fylġan} (early analogy).

\subsection{Derivation trace}\label{derivation-trace-111}

Proto input: \emph{*fúlgijaną}

\begingroup
\setlength{\fboxsep}{6pt}

\noindent\fbox{%
\begin{minipage}{0.97\linewidth}
\footnotesize
\begin{tabular}{@{}>{\raggedright\arraybackslash}p{0.339\linewidth}@{\hspace{0.030\linewidth}}>{\raggedright\arraybackslash}p{0.630\linewidth}@{}}
\begin{minipage}[t]{\linewidth}
\centering\textbf{Earlier Germanic changes}\par
\vspace{0.35em}
\raggedright
\centering\textbf{Northwest and West Germanic}\par
\raggedright
\vspace{0.2em}
\raggedright [no change]\par
\vspace{0.6em}
\centering\textbf{Early Anglo-Frisian}\par
\raggedright
\vspace{0.2em}
\raggedright [no change]\par
\end{minipage}
&
\begin{minipage}[t]{\linewidth}
\centering\textbf{Old English changes}\par
\vspace{0.35em}
\raggedright
\begin{tabularx}{\linewidth}{@{}>{\raggedright\arraybackslash}X@{\hspace{0.55em}}>{\raggedright\arraybackslash}p{0.22\linewidth}@{\hspace{0.25em}}}
OE Heavy Syllable Nasal Apocope & \emph{*fúlgijan} \\
\mbox{OE Secondary Nasalization} & \emph{*fúlgijąn} \\
\mbox{Sievers Law Syncope} & \emph{*fúlgjąn} \\
\mbox{OE Velar Palatalization} & \emph{*fúlʝjąn} \\
\mbox{OE I Umlaut} & \emph{*fylʝjąn} \\
\mbox{OE Weak Tail Reduction} & \emph{*fylʝjan} \\
\mbox{OE J Loss After Heavy} & \emph{*fylʝan} \\
\mbox{OE Palatal Fricative Merger} & \emph{*fyljan} \\
\end{tabularx}
\end{minipage}
\\
\end{tabular}
\end{minipage}%
} \endgroup

Old English form: \emph{fylġan}

\subsection{Reconstruction and comparative
evidence}\label{reconstruction-and-comparative-evidence-95}

Kroonen reconstructs the verb as \emph{*fulgen-} and gives Old English
{\emph{fylgan}}\index[iv]{02oe@\ivlangheader{Old English}{West Saxon normalization unmarked}!fylgan@\iventry{fylgan}{}}
`follow',
{\emph{folgian}}\index[iv]{02oe@\ivlangheader{Old English}{West Saxon normalization unmarked}!folgian@\iventry{folgian}{}}
`follow', adding that Old Norse
{\emph{fylgja}}\index[iv]{09on@\ivlangheader{Old Norse}{}!fylgja@\iventry{fylgja}{}}
`follow' and Old English
{\emph{fylg(e)an}}\index[iv]{02oe@\ivlangheader{Old English}{West Saxon normalization unmarked}!fylgean@\iventry{fylg(e)an}{}}
continue a formation \emph{*fulgjan-}
(\citeproc{ref-Kroonen2013}{Kroonen 2013, 159}). The comparative
headword and the class-I formation are therefore related but not
identical.

Ringe and Taylor distinguish PNWGmc \Recon{fulgija-} `follow'
\textasciitilde{} \Recon{fulgai-} `follow' \textgreater{} OE
{\emph{fylgan}}\index[iv]{02oe@\ivlangheader{Old English}{West Saxon normalization unmarked}!fylgan@\iventry{fylgan}{}}
`follow' \textasciitilde{}
{\emph{folgian}}\index[iv]{02oe@\ivlangheader{Old English}{West Saxon normalization unmarked}!folgian@\iventry{folgian}{}}
`follow' and describe it as a dual formation that probably reflects an
older alternation between j-present and e-stative
(\citeproc{ref-RingeTaylor2014}{Ringe and Taylor 2014, 293--94}). This
is a stem-class choice, not a spelling choice. The derivational input
\Recon{fúlgijaną}\index[iv]{03pgmc@\ivlangheader{Proto-Germanic}{}!fulgijana@\iventry{*fúlgijaną}{}}
`follow' belongs to the class-I \emph{*fulgija-} / \emph{*fulgjan-}
branch; the citation form
\Recon{fulgēną}\index[iv]{03pgmc@\ivlangheader{Proto-Germanic}{}!fulgena@\iventry{*fulgēną}{}}
`follow' belongs to the parallel class-II history behind
{\emph{folgian}}\index[iv]{02oe@\ivlangheader{Old English}{West Saxon normalization unmarked}!folgian@\iventry{folgian}{}}
`follow'.

\subsection{Old English evidence}\label{old-english-evidence-96}

The Old English evidence preserves both formations. Clark Hall lists
{\emph{fylgan}}\index[iv]{02oe@\ivlangheader{Old English}{West Saxon normalization unmarked}!fylgan@\iventry{fylgan}{}}
`follow' with variant spellings
{\emph{fylgian}}\index[iv]{02oe@\ivlangheader{Old English}{West Saxon normalization unmarked}!fylgian@\iventry{fylgian}{}}
`follow' and
{\emph{fyligan}}\index[iv]{02oe@\ivlangheader{Old English}{West Saxon normalization unmarked}!fyligan@\iventry{fyligan}{}}
`follow' (\citeproc{ref-ClarkHall1960}{Clark Hall 1960, 125}).
Bosworth-Toller likewise has a separate
{\emph{fylgean}}\index[iv]{02oe@\ivlangheader{Old English}{West Saxon normalization unmarked}!fylgean@\iventry{fylgean}{}}
`follow' entry (\citeproc{ref-BosworthToller1898}{Bosworth and Toller
1898, 275}).

Bright notes traces of the older conjugation in
{\emph{fylg(e)an}}\index[iv]{02oe@\ivlangheader{Old English}{West Saxon normalization unmarked}!fylgean@\iventry{fylg(e)an}{}}
`follow' (\citeproc{ref-BrightCassidyRingler1971}{Bright 1971, 77}) and
lists
{\emph{folgian}}\index[iv]{02oe@\ivlangheader{Old English}{West Saxon normalization unmarked}!folgian@\iventry{folgian}{}}
`follow'
({\emph{fylgean}}\index[iv]{02oe@\ivlangheader{Old English}{West Saxon normalization unmarked}!fylgean@\iventry{fylgean}{}}
`follow') in the glossary
(\citeproc{ref-BrightCassidyRingler1971}{Bright 1971, 364}). The
relevant comparison form in this entry is therefore the class-I verb
{\emph{fylgan}}\index[iv]{02oe@\ivlangheader{Old English}{West Saxon normalization unmarked}!fylgan@\iventry{fylgan}{}}
`follow' /
{\emph{fylgean}}\index[iv]{02oe@\ivlangheader{Old English}{West Saxon normalization unmarked}!fylgean@\iventry{fylgean}{}},
here normalized as
{\emph{fylġan}}\index[iv]{02oe@\ivlangheader{Old English}{West Saxon normalization unmarked}!fylgan@\iventry{fylġan}{}}
`follow'. The spelling with ġ represents the palatalized velar before a
front-vocalic environment.

\subsection{Development to Old
English}\label{development-to-old-english-94}

\Recon{fúlgijaną}\index[iv]{03pgmc@\ivlangheader{Proto-Germanic}{}!fulgijana@\iventry{*fúlgijaną}{}}
`follow' is a class-I weak-verb formation. In the class-I branch the
\emph{*j} blocks NWGmc lowering of \emph{u} to \emph{o}, since Ringe and
Taylor formulate that lowering for environments in which no \emph{*j}
intervened (\citeproc{ref-RingeTaylor2014}{Ringe and Taylor 2014, 96}).
The same front-vocalic environment then triggers i-umlaut, so \emph{u}
becomes \emph{y} (\citeproc{ref-RingeTaylor2014}{Ringe and Taylor 2014,
§ 6.6.2}).

The subsequent Old English developments are palatalization of the velar,
weak-tail reduction, and loss of \emph{j} after a heavy syllable,
yielding
{\emph{fylġan}}\index[iv]{02oe@\ivlangheader{Old English}{West Saxon normalization unmarked}!fylgan@\iventry{fylġan}{}}
`follow'. This is the regular outcome of the class-I formation. The
class-II form
{\emph{folgian}}\index[iv]{02oe@\ivlangheader{Old English}{West Saxon normalization unmarked}!folgian@\iventry{folgian}{}}
`follow' belongs to the parallel \emph{*-ē-} / \emph{*-ai-} branch and
is not the form modeled here.

\subsection{Class comparison}\label{class-comparison-1}

A class comparison identifies which inherited formation corresponds to
the established Old English form under discussion. The comparison below
is manual; no full automatic class probe is presented here.

{\def\LTcaptype{none} % do not increment counter
\begin{longtable}[]{@{}
  >{\raggedright\arraybackslash}p{(\linewidth - 8\tabcolsep) * \real{0.2000}}
  >{\raggedright\arraybackslash}p{(\linewidth - 8\tabcolsep) * \real{0.2000}}
  >{\raggedright\arraybackslash}p{(\linewidth - 8\tabcolsep) * \real{0.2000}}
  >{\raggedright\arraybackslash}p{(\linewidth - 8\tabcolsep) * \real{0.2000}}
  >{\raggedright\arraybackslash}p{(\linewidth - 8\tabcolsep) * \real{0.2000}}@{}}
\toprule\noalign{}
\begin{minipage}[b]{\linewidth}\raggedright
Formation / class
\end{minipage} & \begin{minipage}[b]{\linewidth}\raggedright
Candidate input
\end{minipage} & \begin{minipage}[b]{\linewidth}\raggedright
Expected or documented OE outcome
\end{minipage} & \begin{minipage}[b]{\linewidth}\raggedright
OE comparison form
\end{minipage} & \begin{minipage}[b]{\linewidth}\raggedright
Result
\end{minipage} \\
\midrule\noalign{}
\endhead
\bottomrule\noalign{}
\endlastfoot
citation class-II formation &
{\emph{*fulgēną}}\index[iv]{03pgmc@\ivlangheader{Proto-Germanic}{}!fulgena@\iventry{*fulgēną}{}}
& probe output: {\emph{folgon}} &
{\emph{folgian}}\index[iv]{02oe@\ivlangheader{Old English}{West Saxon normalization unmarked}!folgian@\iventry{folgian}{}}
`follow' & mismatch: the regular output is not the remodeled infinitive
\emph{folgian} \\
parallel class-II branch & PNWGmc \emph{*fulgai-} & Ringe-Taylor: OE
{\emph{folgian}}\index[iv]{02oe@\ivlangheader{Old English}{West Saxon normalization unmarked}!folgian@\iventry{folgian}{}}
`follow' &
{\emph{folgian}}\index[iv]{02oe@\ivlangheader{Old English}{West Saxon normalization unmarked}!folgian@\iventry{folgian}{}}
`follow' & documents the separate class-II branch, but not the target of
this entry \\
selected class-I formation &
{\emph{*fúlgijaną}}\index[iv]{03pgmc@\ivlangheader{Proto-Germanic}{}!fulgijana@\iventry{*fúlgijaną}{}}
& regular output: {\emph{fylġan}} `follow' &
{\emph{fylġan}}\index[iv]{02oe@\ivlangheader{Old English}{West Saxon normalization unmarked}!fylgan@\iventry{fylġan}{}}
`follow' /
{\emph{fylgan}}\index[iv]{02oe@\ivlangheader{Old English}{West Saxon normalization unmarked}!fylgan@\iventry{fylgan}{}}
`follow' & exact match between input, output, and class \\
\end{longtable}
}

The relevant point is the class split.
{\emph{fylġan}}\index[iv]{02oe@\ivlangheader{Old English}{West Saxon normalization unmarked}!fylgan@\iventry{fylġan}{}}
`follow' is the regular Old English outcome of the class-I
\emph{*fulgija-} / \emph{*fulgjan-} formation, whereas
{\emph{folgian}}\index[iv]{02oe@\ivlangheader{Old English}{West Saxon normalization unmarked}!folgian@\iventry{folgian}{}}
`follow' belongs to the parallel class-II branch.

\section{\texorpdfstring{gall --- OE
\emph{ġealla}}{gall --- OE ġealla}}\label{gall-oe-ux121ealla}

\index[iv]{02oe@\ivlangheader{Old English}{West Saxon normalization unmarked}!gealla@\iventry{ġealla}{}}
\index[iv]{03pgmc@\ivlangheader{Proto-Germanic}{}!galla@\iventry{*gállą}{}}
\index[iv]{03pgmc@\ivlangheader{Proto-Germanic}{}!gallo@\iventry{*gállô}{}}

Derivation: citation reconstruction \emph{*gállą}; form followed here
\emph{*gállô} \textgreater{} \emph{ġealla} (early analogy).

\subsection{Derivation trace}\label{derivation-trace-112}

Proto input: \emph{*gállô}

\begingroup
\setlength{\fboxsep}{6pt}

\noindent\fbox{%
\begin{minipage}{0.97\linewidth}
\small
\begin{tabular}{@{}>{\raggedright\arraybackslash}p{0.359\linewidth}@{\hspace{0.030\linewidth}}>{\raggedright\arraybackslash}p{0.611\linewidth}@{}}
\begin{minipage}[t]{\linewidth}
\centering\textbf{Earlier Germanic changes}\par
\vspace{0.35em}
\raggedright
\centering\textbf{Northwest and West Germanic}\par
\raggedright
\vspace{0.2em}
\raggedright [no change]\par
\vspace{0.6em}
\centering\textbf{Early Anglo-Frisian}\par
\raggedright
\vspace{0.2em}
\raggedright [no change]\par
\end{minipage}
&
\begin{minipage}[t]{\linewidth}
\centering\textbf{Old English changes}\par
\vspace{0.35em}
\raggedright
\begin{tabularx}{\linewidth}{@{}>{\raggedright\arraybackslash}X@{\hspace{0.55em}}>{\raggedright\arraybackslash}p{0.14\linewidth}@{\hspace{0.25em}}}
\mbox{EAF Brightening} & \emph{*gællô} \\
\mbox{OE Breaking} & \emph{*geallô} \\
\mbox{OE Velar Palatalization} & \emph{*ʝeallô} \\
OE Unstressed Long Vowel Shortening & \emph{*ʝealla} \\
\mbox{OE Palatal Fricative Merger} & \emph{*jealla} \\
\end{tabularx}
\end{minipage}
\\
\end{tabular}
\end{minipage}%
} \endgroup

Old English form: \emph{ġealla}

\subsection{Reconstruction and comparative
evidence}\label{reconstruction-and-comparative-evidence-96}

The wider cognate family can be presented under a form such as
\Recon{gállą}\index[iv]{03pgmc@\ivlangheader{Proto-Germanic}{}!galla@\iventry{*gállą}{}}
`gall', but the Old English noun itself belongs with a weak noun
\emph{*gallōn-} `gall', cited here as
\Recon{gállô}\index[iv]{03pgmc@\ivlangheader{Proto-Germanic}{}!gallo@\iventry{*gállô}{}}
`gall' (\citeproc{ref-Kroonen2013}{Kroonen 2013, 165}). The derivational
input therefore differs from the broader comparative headword in stem
class.

The stem class determines the Old English shape. The weak masculine
pathway preserves the ending needed for \emph{ġealla} `gall', whereas a
simple strong-noun headword does not align as closely with the attested
OE noun.

\subsection{Old English evidence}\label{old-english-evidence-97}

Old English dictionaries record the noun as
{\emph{gealla}}\index[iv]{02oe@\ivlangheader{Old English}{West Saxon normalization unmarked}!gealla@\iventry{gealla}{}}
`gall', and Bright also gives the dative
{\emph{geallan}}\index[iv]{02oe@\ivlangheader{Old English}{West Saxon normalization unmarked}!geallan@\iventry{geallan}{}}
`gall', confirming a weak-noun paradigm
(\citeproc{ref-BosworthToller1898}{Bosworth and Toller 1898, 297};
\citeproc{ref-ClarkHall1960}{Clark Hall 1960, 145};
\citeproc{ref-BrightCassidyRingler1971}{Bright 1971, 372}). The
normalized spelling
{\emph{ġealla}}\index[iv]{02oe@\ivlangheader{Old English}{West Saxon normalization unmarked}!gealla@\iventry{ġealla}{}}
`gall' uses ġ for the palatal consonant.

Campbell also notes dialectal variation, contrasting West Saxon or
Kentish
{\emph{gealla}}\index[iv]{02oe@\ivlangheader{Old English}{West Saxon normalization unmarked}!gealla@\iventry{gealla}{}}
`gall' with Anglian
{\emph{galla}}\index[iv]{02oe@\ivlangheader{Old English}{West Saxon normalization unmarked}!galla~04@\iventry{galla}{Angl.}}
`gall' (\citeproc{ref-Campbell1959}{Campbell 1959, § 486}). The
target of this entry is the West Saxon type
{\emph{ġealla}}\index[iv]{02oe@\ivlangheader{Old English}{West Saxon normalization unmarked}!gealla@\iventry{ġealla}{}}
`gall'.

\subsection{Development to Old
English}\label{development-to-old-english-95}

From
\Recon{gállô}\index[iv]{03pgmc@\ivlangheader{Proto-Germanic}{}!gallo@\iventry{*gállô}{}}
`gall', the weak noun develops through the expected Old English history
of the suffix and the regular breaking environment before \emph{ll},
yielding {\emph{ġealla}} `gall' (\citeproc{ref-Campbell1959}{Campbell
1959, § 486}). Once the weak masculine input is chosen, the noun
follows a regular path to its attested Old English form.

The decisive issue is therefore morphological. Old English reflects the
weak noun, while the broader family label belongs to a different way of
presenting the cognate set.

\subsection{Stem comparison}\label{stem-comparison}

The comparison below sets the relevant forms side by side. It separates
the broader comparative headword from the weak noun formation actually
reflected in Old English.

{\def\LTcaptype{none} % do not increment counter
\begin{longtable}[]{@{}
  >{\raggedright\arraybackslash}p{(\linewidth - 8\tabcolsep) * \real{0.2000}}
  >{\raggedright\arraybackslash}p{(\linewidth - 8\tabcolsep) * \real{0.2000}}
  >{\raggedright\arraybackslash}p{(\linewidth - 8\tabcolsep) * \real{0.2000}}
  >{\raggedright\arraybackslash}p{(\linewidth - 8\tabcolsep) * \real{0.2000}}
  >{\raggedright\arraybackslash}p{(\linewidth - 8\tabcolsep) * \real{0.2000}}@{}}
\toprule\noalign{}
\begin{minipage}[b]{\linewidth}\raggedright
Formation
\end{minipage} & \begin{minipage}[b]{\linewidth}\raggedright
Candidate input
\end{minipage} & \begin{minipage}[b]{\linewidth}\raggedright
Expected or documented OE outcome
\end{minipage} & \begin{minipage}[b]{\linewidth}\raggedright
OE comparison form
\end{minipage} & \begin{minipage}[b]{\linewidth}\raggedright
Result
\end{minipage} \\
\midrule\noalign{}
\endhead
\bottomrule\noalign{}
\endlastfoot
broader family label &
{\emph{*gállą}}\index[iv]{03pgmc@\ivlangheader{Proto-Germanic}{}!galla@\iventry{*gállą}{}}
& broader cognate-set headword & wider family context & useful lexeme
label, but not the direct source of \emph{ġealla} \\
selected weak noun &
{\emph{*gállô}}\index[iv]{03pgmc@\ivlangheader{Proto-Germanic}{}!gallo@\iventry{*gállô}{}}
& regular output: {\emph{ġealla}} `gall' &
{\emph{ġealla}}\index[iv]{02oe@\ivlangheader{Old English}{West Saxon normalization unmarked}!gealla@\iventry{ġealla}{}}
`gall' & exact match between formation and attested OE noun \\
dialectal Anglian continuation & weak noun branch & Anglian
{\emph{galla}}\index[iv]{02oe@\ivlangheader{Old English}{West Saxon normalization unmarked}!galla~04@\iventry{galla}{Angl.}}
`gall' type &
{\emph{galla}}\index[iv]{02oe@\ivlangheader{Old English}{West Saxon normalization unmarked}!galla~04@\iventry{galla}{Angl.}}
`gall' & genuine OE variant, but not the West Saxon form used here \\
\end{longtable}
}

The weak-noun stem class is the relevant point. It gives a direct route
to attested \emph{ġealla} `gall', while the broader comparative label
serves only as a family heading.

\section{\texorpdfstring{heaven --- OE
\emph{heofon}}{heaven --- OE heofon}}\label{heaven-oe-heofon}

\index[iv]{02oe@\ivlangheader{Old English}{West Saxon normalization unmarked}!heofon@\iventry{heofon}{}}
\index[iv]{02oe@\ivlangheader{Old English}{West Saxon normalization unmarked}!seofon@\iventry{seofon}{}}
\index[iv]{03pgmc@\ivlangheader{Proto-Germanic}{}!sebun@\iventry{*sébun}{}}
\index[iv]{03pgmc@\ivlangheader{Proto-Germanic}{}!xemenaz@\iventry{*xémenaz}{}}
\index[iv]{06nsgmc@\ivlangheader{Northern West Germanic}{}!xebun@\iventry{*xébun}{}}

Derivation: citation reconstruction \emph{*xémenaz}; form followed here
\emph{*xébun} \textgreater{} \emph{heofon} (early analogy).

\subsection{Derivation trace}\label{derivation-trace-113}

Proto input: \emph{*xébun}

\begingroup
\setlength{\fboxsep}{6pt}

\noindent\fbox{%
\begin{minipage}{0.97\linewidth}
\small
\begin{tabular}{@{}>{\raggedright\arraybackslash}p{0.359\linewidth}@{\hspace{0.030\linewidth}}>{\raggedright\arraybackslash}p{0.611\linewidth}@{}}
\begin{minipage}[t]{\linewidth}
\centering\textbf{Earlier Germanic changes}\par
\vspace{0.35em}
\raggedright
\centering\textbf{Northwest and West Germanic}\par
\raggedright
\vspace{0.2em}
\raggedright [no change]\par
\vspace{0.6em}
\centering\textbf{Early Anglo-Frisian}\par
\raggedright
\vspace{0.2em}
\raggedright [no change]\par
\end{minipage}
&
\begin{minipage}[t]{\linewidth}
\centering\textbf{Old English changes}\par
\vspace{0.35em}
\raggedright
\begin{tabularx}{\linewidth}{@{}>{\raggedright\arraybackslash}X@{\hspace{0.55em}}>{\raggedright\arraybackslash}p{0.14\linewidth}@{\hspace{0.25em}}}
OE Velar Fricative Palatalization & \emph{*çébun} \\
\mbox{OE Back Mutation} & \emph{*çéobun} \\
\mbox{PGmc B Allophony} & \emph{*çéoβun} \\
\mbox{OE Med Unstressed U Lowering} & \emph{*çéoβon} \\
\end{tabularx}
\end{minipage}
\\
\end{tabular}
\end{minipage}%
} \endgroup

Old English form: \emph{heofon}

\subsection{A. Selected CAPR
derivation}\label{a.-selected-capr-derivation}

{\def\LTcaptype{none} % do not increment counter
\begin{longtable}[]{@{}ll@{}}
\toprule\noalign{}
Item & Value \\
\midrule\noalign{}
\endhead
\bottomrule\noalign{}
\endlastfoot
lexical item & heaven \\
citation reconstruction / lexeme label & *xémenaz \\
form followed here & \emph{xébun (northern WGmc / pre-OE }hebun) \\
Old English target & heofon \\
classification & early\_analogy \\
derivational result & *xébun \textgreater{} heofon (multiplicity 1,
exact match) \\
\end{longtable}
}

Compact live trace (corrected cascade, literal adjacent-\emph{mn}
SC022):

{\def\LTcaptype{none} % do not increment counter
\begin{longtable}[]{@{}ll@{}}
\toprule\noalign{}
step & form \\
\midrule\noalign{}
\endhead
\bottomrule\noalign{}
\endlastfoot
proto input & \emph{*xébun} \\
OE Med Unstressed U Lowering & \emph{*xébon} \\
OE Velar Fricative Palatalization & \emph{*çébon} \\
PGmc B Allophony & \emph{*çéβon} \\
OE Back Mutation & \emph{*çéoβon} \\
Old English Orthography & \emph{heofon} \\
\end{longtable}
}

No rule in this span creates the labial: it enters with the input
\emph{b} (→ \emph{β} → \emph{f}). Only vocalic and surface changes are
modelled.

\subsection{B. Why *xébun is not
arbitrary}\label{b.-why-xuxe9bun-is-not-arbitrary}

Ringe and Taylor reconstruct northern West Germanic
\Recon{hebun}\index[iv]{06nsgmc@\ivlangheader{Northern West Germanic}{}!hebun@\iventry{*hebun}{}}
`sky, heaven', gen.
\Recon{hebunas}\index[iv]{06nsgmc@\ivlangheader{Northern West Germanic}{}!hebunas@\iventry{*hebunas}{}}
`heaven', whence West Saxon and Northumbrian
{\emph{heofon}}\index[iv]{02oe@\ivlangheader{Old English}{West Saxon normalization unmarked}!heofon@\iventry{heofon}{}}
`heaven', Mercian
{\emph{heofen}}\index[iv]{02oe@\ivlangheader{Old English}{West Saxon normalization unmarked}!heofen~05@\iventry{heofen}{Merc.}}
`heaven' (\citeproc{ref-RingeTaylor2014}{Ringe and Taylor 2014, 272,
324}). R\&T's \emph{hebun} is the lexeme CAPR writes \emph{*xébun}
(canonical \emph{x} for source \emph{h}), dated to that northern West
Germanic stage, not Proto-Germanic; §D explains its labial.

The word is an inherited mn-stem (\citeproc{ref-Kroonen2013}{Kroonen
2013, 220}): PIE \emph{h₂ék-mōn, gen. }h₂k-mn-ós, remodelled in Germanic
to nom.
\Recon{hemō}\index[iv]{03pgmc@\ivlangheader{Proto-Germanic}{}!hemo@\iventry{*hemō}{}}
`heaven', gen.
\Recon{hemnaz}\index[iv]{03pgmc@\ivlangheader{Proto-Germanic}{}!hemnaz@\iventry{*hemnaz}{}}
`heaven'. In the zero-grade oblique cluster cells the sequence
\emph{-mn- is adjacent, and the change }mn \textgreater{} βn/bn/fn* is
regular there (Fulk §6.14 (\citeproc{ref-Fulk2018}{Fulk 2018, 121})).
Old Norse preserves exactly this paradigmatic distribution ---
{\emph{himinn}}\index[iv]{09on@\ivlangheader{Old Norse}{}!himinn@\iventry{himinn}{}}
`heaven' (nominative, with \emph{m}) beside dative
{\emph{hifni}}\index[iv]{09on@\ivlangheader{Old Norse}{}!hifni@\iventry{hifni}{}}
`heaven' (oblique, with \emph{f} from the syncopated cluster) --- while
Old English and Old Saxon (OS
{\emph{heban}}\index[iv]{13os@\ivlangheader{Old Saxon}{}!heban@\iventry{heban}{}}
`heaven',
{\emph{hebenes}}\index[iv]{13os@\ivlangheader{Old Saxon}{}!hebenes@\iventry{hebenes}{}}
`heaven' (\citeproc{ref-Fulk2018}{Fulk 2018})) generalized the labial
into the vowel-bearing stem. The \emph{b} of \emph{xébun is therefore
analogical }within this nominative input*, but its historical source is
regular and recoverable --- back to the inherited genitive (see §E).

\subsection{\texorpdfstring{C. The unresolved
\emph{u}}{C. The unresolved u}}\label{c.-the-unresolved-u}

The medial \emph{u} is the less secure element and must not be treated
as certain merely because the cascade needs a back vowel to trigger back
mutation. It is the n-stem suffixal vowel, the genuine conditioning
environment for West Saxon back umlaut of \emph{e} (parallel to
\emph{seven}, §D), and it is supported by Ringe--Taylor's reconstruction
and by the Old English back-mutation pattern
(\citeproc{ref-RingeTaylor2014}{Ringe and Taylor 2014, 324}). But its
exact deeper ablaut grade --- \emph{u} here, against the \emph{i} of the
cognate vowel-stem
\Recon{hemina-}\index[iv]{03pgmc@\ivlangheader{Proto-Germanic}{}!hemina@\iventry{*hemina-}{}}
({\emph{himinn}}\index[iv]{09on@\ivlangheader{Old Norse}{}!himinn@\iventry{himinn}{}}
`heaven', Go.
{\emph{himins}}\index[iv]{12goth@\ivlangheader{Gothic}{}!himins@\iventry{himins}{}}
`heaven') --- is not fully recoverable. It is the irreducible residue of
the reconstruction; the \emph{b} is not.

\subsection{D. Regular modelled span}\label{d.-regular-modelled-span}

The northern-WGmc / pre-OE → Old English development \emph{*xébun}
\textgreater{} \emph{heofon} is regular and is the exact structural
parallel of the independently regular numeral \emph{seven},
\emph{*sébun} \textgreater{} \emph{seofon} (row 2174,
\emph{DERIVATION\_CLASS} regular). Both have a medial \emph{u} and a
labial, and both undergo West Saxon back umlaut of \emph{e} before
\emph{u} across that labial (\citeproc{ref-RingeTaylor2014}{Ringe and
Taylor 2014, 324}; \citeproc{ref-HoggGrammar1992}{Hogg 1992a}). The live
traces run in lock-step (medial \emph{u}-lowering, labial
spirantization, back mutation, orthography). Ringe and Taylor explicitly
treat \emph{heaven} and \emph{seven} as parallel back-umlaut examples,
which is why \emph{seven} is retained here as the control.

\subsection{E. Deep oblique controls}\label{e.-deep-oblique-controls}

The obliques go deeper than the selected nominative and derive the
labial themselves. Live outputs under the corrected literal
adjacent-\emph{mn} cascade, with evidential status marked:

{\def\LTcaptype{none} % do not increment counter
\begin{longtable}[]{@{}
  >{\raggedright\arraybackslash}p{(\linewidth - 8\tabcolsep) * \real{0.2000}}
  >{\raggedright\arraybackslash}p{(\linewidth - 8\tabcolsep) * \real{0.2000}}
  >{\raggedright\arraybackslash}p{(\linewidth - 8\tabcolsep) * \real{0.2000}}
  >{\raggedright\arraybackslash}p{(\linewidth - 8\tabcolsep) * \real{0.2000}}
  >{\raggedright\arraybackslash}p{(\linewidth - 8\tabcolsep) * \real{0.2000}}@{}}
\toprule\noalign{}
\begin{minipage}[b]{\linewidth}\raggedright
cell / form
\end{minipage} & \begin{minipage}[b]{\linewidth}\raggedright
type
\end{minipage} & \begin{minipage}[b]{\linewidth}\raggedright
corrected-cascade output
\end{minipage} & \begin{minipage}[b]{\linewidth}\raggedright
status
\end{minipage} & \begin{minipage}[b]{\linewidth}\raggedright
note
\end{minipage} \\
\midrule\noalign{}
\endhead
\bottomrule\noalign{}
\endlastfoot
nom. citation
\Recon{xémenaz}\index[iv]{03pgmc@\ivlangheader{Proto-Germanic}{}!xemenaz@\iventry{*xémenaz}{}}
& reconstructed (citation) &
{\emph{hemen}}\index[iv]{02oe@\ivlangheader{Old English}{West Saxon normalization unmarked}!hemen@\iventry{hemen}{}}
& counterfactual (no labial) & citation a-stem; no cluster, so no
regular \emph{f} \\
gen.sg cluster
\Recon{xémnas}\index[iv]{03pgmc@\ivlangheader{Proto-Germanic}{}!xemnas@\iventry{*xémnas}{}}
& reconstructed (PIE *h₂k-mn-ós) &
{\emph{hefnes}}\index[iv]{02oe@\ivlangheader{Old English}{West Saxon normalization unmarked}!hefnes@\iventry{hefnes}{}}
& reconstructed, \textbf{not attested in OE} & labial \textbf{regular}
here (mn \textgreater{} βn) \\
dat.sg cluster
\Recon{xémni}\index[iv]{03pgmc@\ivlangheader{Proto-Germanic}{}!xemni@\iventry{*xémni}{}}
& reconstructed &
{\emph{hifn}}\index[iv]{02oe@\ivlangheader{Old English}{West Saxon normalization unmarked}!hifn@\iventry{hifn}{}}
& attested only as ON
{\emph{hifni}}\index[iv]{09on@\ivlangheader{Old Norse}{}!hifni@\iventry{hifni}{}}
`heaven' & labial + i-umlaut regular \\
dat.pl cluster
\Recon{xémnum}\index[iv]{03pgmc@\ivlangheader{Proto-Germanic}{}!xemnum@\iventry{*xémnum}{}}
& reconstructed &
{\emph{hefnum}}\index[iv]{02oe@\ivlangheader{Old English}{West Saxon normalization unmarked}!hefnum@\iventry{hefnum}{}}
& reconstructed, \textbf{not attested in OE} & labial regular \\
remodelled stem
\Recon{xébun}\index[iv]{06nsgmc@\ivlangheader{Northern West Germanic}{}!xebun@\iventry{*xébun}{}}
& reconstructed (R\&T *hebun) &
{\emph{heofon}}\index[iv]{02oe@\ivlangheader{Old English}{West Saxon normalization unmarked}!heofon@\iventry{heofon}{}}
& \textbf{attested} (selected row) & labial pre-supplied \\
R\&T genitive
\Recon{hebunas}\index[iv]{06nsgmc@\ivlangheader{Northern West Germanic}{}!hebunas@\iventry{*hebunas}{}}
& reconstructed &
{\emph{heofones}}\index[iv]{02oe@\ivlangheader{Old English}{West Saxon normalization unmarked}!heofones@\iventry{heofones}{}}
& \textbf{attested} & labial pre-supplied \\
OE oblique
{\emph{heofnes}}\index[iv]{02oe@\ivlangheader{Old English}{West Saxon normalization unmarked}!heofnes@\iventry{heofnes}{}}
/
{\emph{heofnum}}\index[iv]{02oe@\ivlangheader{Old English}{West Saxon normalization unmarked}!heofnum@\iventry{heofnum}{}}
& \textbf{attested} & --- & actual OE forms & carry \emph{eo}; from the
re-vowelled stem + late syncope \\
ON
{\emph{himinn}}\index[iv]{09on@\ivlangheader{Old Norse}{}!himinn@\iventry{himinn}{}}
:
{\emph{hifni}}\index[iv]{09on@\ivlangheader{Old Norse}{}!hifni@\iventry{hifni}{}}
`heaven' & \textbf{attested} (North Gmc) & --- & attested cluster line &
keeps the \emph{m} : \emph{f} alternation \\
OS
{\emph{heban}}\index[iv]{13os@\ivlangheader{Old Saxon}{}!heban@\iventry{heban}{}}
/
{\emph{hebenes}}\index[iv]{13os@\ivlangheader{Old Saxon}{}!hebenes@\iventry{hebenes}{}}
& \textbf{attested} (Old Saxon) & --- & labial generalized &
vowel-bearing throughout \\
\end{longtable}
}

Distinctions kept explicit: attested OE forms (\emph{heofon},
\emph{heofones}, \emph{heofnes}, \emph{heofnum}) descend from the
re-vowelled stem; the cluster outputs
\emph{hefnes}/\emph{hifn}/\emph{hefnum} are reconstructed regular
outputs, and for Old English they are counterfactual --- the actual
attested obliques carry \emph{eo} from back mutation on the
retained/\hspace{0pt}restored medial vowel, then late syncope
(\emph{heofonum} \textgreater{} \emph{heofnum}), not from a bare
\emph{-fn-} cluster. The genuine attested witness of the pure cluster
line is Old Norse
{\emph{hifni}}\index[iv]{09on@\ivlangheader{Old Norse}{}!hifni@\iventry{hifni}{}}
`heaven'.

\subsection{F. Why the oblique is not the selected
row}\label{f.-why-the-oblique-is-not-the-selected-row}

The oblique lines are historically essential --- they show that the
labial has a regular origin, all the way back to the inherited genitive
\emph{h₂k-mn-ós --- but they are inferior as the single CAPR
input→target datum. In Old English every oblique cluster form encounters
morphological re-vowelling (restoration of a medial vowel into the
labialized \emph{-βn-} stem) before it reaches an attested Old English
shape. That re-vowelling is a non-regular event; it interrupts any
single deterministic path from a cluster oblique to an attested OE form.
The nominative input }xébun begins after that restructuring and then
proceeds by regular sound change to a secure attested target --- the
longest clean modellable span.

\subsection{G. Attestation warning}\label{g.-attestation-warning}

Bare Old English \emph{hefn} is not securely attested and must not be
reintroduced as an attested comparator. The earliest witnesses (Cædmon's
Hymn) are
{\emph{hefun}}\index[iv]{02oe@\ivlangheader{Old English}{West Saxon normalization unmarked}!hefun~06@\iventry{hefun}{North.}}
`heaven',
{\emph{hefen}}\index[iv]{02oe@\ivlangheader{Old English}{West Saxon normalization unmarked}!hefen~06@\iventry{hefen}{North.}}
`heaven', and (Moore Bede)
{\emph{heben}}\index[iv]{02oe@\ivlangheader{Old English}{West Saxon normalization unmarked}!heben~06@\iventry{heben}{North.}}
`heaven' --- all with a medial vowel. The reconstructed cluster output
\emph{hefn} (§E) is a reconstruction, never a manuscript form.

\subsection{H. Classification
rationale}\label{h.-classification-rationale}

The row is \emph{early\_analogy}, not \emph{late\_analogy} and not
\emph{known\_unmodelled}.

\begin{itemize}
\tightlist
\item
  The selected row models exactly one form,
  \Recon{xébun}\index[iv]{06nsgmc@\ivlangheader{Northern West Germanic}{}!xebun@\iventry{*xébun}{}}
  →
  {\emph{heofon}}\index[iv]{02oe@\ivlangheader{Old English}{West Saxon normalization unmarked}!heofon@\iventry{heofon}{}}.
  The necessary paradigm restructuring (labial levelling + suffix-vowel
  generalization) precedes the derivational input; once that input is
  chosen, the modelled development is regular. This is the definition of
  \emph{early\_analogy} (analogy separates the input from the citation
  reconstruction before the specifically Old English changes apply).
\item
  It is not \emph{late\_analogy}: the analogy is in the input, not a
  later paradigm-cell selection; the later history of \emph{other}
  paradigm cells (oblique syncope, \emph{heo-} levelling) does not touch
  the modelled nominative derivation.
\item
  It is not \emph{known\_unmodelled}: that label (a legitimate,
  conservative alternative discussed in
  \texttt{dossier-heaven-paradigm-history-2026.md}) would discard a
  real, defensible regular span. \emph{*xébun} \textgreater{}
  \emph{heofon} \emph{is} a clean regular derivation to a secure
  attested target, so CAPR models it.
\end{itemize}

The old cross-syllable \emph{mV\ldots n} proxy in SC022 --- which
formerly fabricated the labial from an intervocalic \emph{m} so that the
previous input \emph{*xémonų} could reach \emph{heofon} --- has been
retired in favour of the literal historical adjacent \emph{mn}
\textgreater{} \emph{βn}. See
\texttt{dossier-heaven-paradigm-history-2026.md} §§13--15 and the
implementation audit
\texttt{audits/heaven-sc022-implementation-2026.md}.

\section{\texorpdfstring{knight --- OE
\emph{cniht}}{knight --- OE cniht}}\label{knight-oe-cniht}

\index[iv]{02oe@\ivlangheader{Old English}{West Saxon normalization unmarked}!cniht@\iventry{cniht}{}}
\index[iv]{03pgmc@\ivlangheader{Proto-Germanic}{}!knextaz@\iventry{*knéxtaz}{}}
\index[iv]{03pgmc@\ivlangheader{Proto-Germanic}{}!knixtaz@\iventry{*kníxtaz}{}}

Derivation: citation reconstruction \emph{*kníxtaz}; form followed here
\emph{*knéxtaz} \textgreater{} \emph{cniht} (early analogy).

\subsection{Derivation trace}\label{derivation-trace-114}

Proto input: \emph{*knéxtaz}

\begingroup
\setlength{\fboxsep}{6pt}

\noindent\fbox{%
\begin{minipage}{0.97\linewidth}
\small
\begin{tabular}{@{}>{\raggedright\arraybackslash}p{0.388\linewidth}@{\hspace{0.030\linewidth}}>{\raggedright\arraybackslash}p{0.582\linewidth}@{}}
\begin{minipage}[t]{\linewidth}
\centering\textbf{Earlier Germanic changes}\par
\vspace{0.35em}
\raggedright
\centering\textbf{Northwest and West Germanic}\par
\raggedright
\vspace{0.2em}
\begin{tabularx}{\linewidth}{@{}>{\raggedright\arraybackslash}X@{\hspace{0.55em}}>{\raggedright\arraybackslash}p{0.24\linewidth}@{\hspace{0.25em}}}
\mbox{EAF Final Z Deletion} & \emph{*knéxta} \\
\end{tabularx}
\vspace{0.6em}
\centering\textbf{Early Anglo-Frisian}\par
\raggedright
\vspace{0.2em}
\raggedright [no change]\par
\end{minipage}
&
\begin{minipage}[t]{\linewidth}
\centering\textbf{Old English changes}\par
\vspace{0.35em}
\raggedright
\begin{tabularx}{\linewidth}{@{}>{\raggedright\arraybackslash}X@{\hspace{0.55em}}>{\raggedright\arraybackslash}p{0.14\linewidth}@{\hspace{0.25em}}}
\mbox{PWGmc Final Bare A Loss} & \emph{*knéxt} \\
\mbox{OE Breaking} & \emph{*knéoxt} \\
\mbox{OE Ws Palatal Umlaut} & \emph{*knixt} \\
\end{tabularx}
\end{minipage}
\\
\end{tabular}
\end{minipage}%
} \endgroup

Old English form: \emph{cniht}

\subsection{Reconstruction and comparative
evidence}\label{reconstruction-and-comparative-evidence-97}

The comparative sources align on an \emph{e}-grade reconstruction for
this noun. Ringe and Taylor cite \Recon{kneht} `knight', and Orel gives
\Recon{knextaz}\index[iv]{03pgmc@\ivlangheader{Proto-Germanic}{}!knextaz@\iventry{*knextaz}{}}
`knight' (\citeproc{ref-RingeTaylor2014}{Ringe and Taylor 2014, 142};
\citeproc{ref-Orel2003}{Orel 2003, 256}). Kluge-Seebold likewise points
to \emph{*knehta-} (\citeproc{ref-KlugeSeebold2011}{Kluge 2011, 506}).
The derivational input
\Recon{knéxtaz}\index[iv]{03pgmc@\ivlangheader{Proto-Germanic}{}!knextaz@\iventry{*knéxtaz}{}}
`knight' follows that comparative evidence.

A competing citation reconstruction
\Recon{kníxtaz}\index[iv]{03pgmc@\ivlangheader{Proto-Germanic}{}!knixtaz@\iventry{*kníxtaz}{}}
`knight' remains possible as a label for the word family, but it is not
the reconstruction followed here. The Old English development discussed
below is based on
\Recon{knéxtaz}\index[iv]{03pgmc@\ivlangheader{Proto-Germanic}{}!knextaz@\iventry{*knéxtaz}{}}
`knight'.

\subsection{Old English evidence}\label{old-english-evidence-98}

Old English dictionaries record the noun as
{\emph{cniht}}\index[iv]{02oe@\ivlangheader{Old English}{West Saxon normalization unmarked}!cniht@\iventry{cniht}{}}
`knight' (\citeproc{ref-ClarkHall1960}{Clark Hall 1960, 63};
\citeproc{ref-BosworthToller1898}{Bosworth and Toller 1898, 71}).
Campbell cites plural
{\emph{cneohtas}}\index[iv]{02oe@\ivlangheader{Old English}{West Saxon normalization unmarked}!cneohtas@\iventry{cneohtas}{}}
`knights' among the broken forms, showing the same vowel environment
from another point in the paradigm (\citeproc{ref-Campbell1959}{Campbell
1959, § 146}).

The target is therefore an ordinary attested Old English noun. No
reconstructed OE comparator is needed here.

\subsection{Development to Old
English}\label{development-to-old-english-96}

From
\Recon{knéxtaz}\index[iv]{03pgmc@\ivlangheader{Proto-Germanic}{}!knextaz@\iventry{*knéxtaz}{}}
`knight', the relevant Old English changes include breaking before the
velar cluster and then the later reduction that yields
{\emph{cniht}}\index[iv]{02oe@\ivlangheader{Old English}{West Saxon normalization unmarked}!cniht@\iventry{cniht}{}}
`knight'. Campbell later notes the early West-Saxon alternation
{\emph{cniht}}\index[iv]{02oe@\ivlangheader{Old English}{West Saxon normalization unmarked}!cniht@\iventry{cniht}{}}
`knight' beside plural
{\emph{cneohtas}}\index[iv]{02oe@\ivlangheader{Old English}{West Saxon normalization unmarked}!cneohtas@\iventry{cneohtas}{}}
`knights' (\citeproc{ref-Campbell1959}{Campbell 1959, § 305}).
Sievers-Brunner gives the same contrast as \emph{cniht \ldots{}
cneohtas} (\citeproc{ref-SieversBrunner1965}{Brunner 1965, § 122}).
With that corrected input, the derivation is straightforward.

\subsection{Stem comparison}\label{stem-comparison-1}

The comparison below sets the relevant forms side by side. It separates
the handbook-supported \emph{e}-grade input from a competing citation
reconstruction.

{\def\LTcaptype{none} % do not increment counter
\begin{longtable}[]{@{}
  >{\raggedright\arraybackslash}p{(\linewidth - 8\tabcolsep) * \real{0.2000}}
  >{\raggedright\arraybackslash}p{(\linewidth - 8\tabcolsep) * \real{0.2000}}
  >{\raggedright\arraybackslash}p{(\linewidth - 8\tabcolsep) * \real{0.2000}}
  >{\raggedright\arraybackslash}p{(\linewidth - 8\tabcolsep) * \real{0.2000}}
  >{\raggedright\arraybackslash}p{(\linewidth - 8\tabcolsep) * \real{0.2000}}@{}}
\toprule\noalign{}
\begin{minipage}[b]{\linewidth}\raggedright
Formation / label
\end{minipage} & \begin{minipage}[b]{\linewidth}\raggedright
Candidate input
\end{minipage} & \begin{minipage}[b]{\linewidth}\raggedright
Expected or documented OE outcome
\end{minipage} & \begin{minipage}[b]{\linewidth}\raggedright
OE comparison form
\end{minipage} & \begin{minipage}[b]{\linewidth}\raggedright
Result
\end{minipage} \\
\midrule\noalign{}
\endhead
\bottomrule\noalign{}
\endlastfoot
competing citation reconstruction &
{\emph{*kníxtaz}}\index[iv]{03pgmc@\ivlangheader{Proto-Germanic}{}!knixtaz@\iventry{*kníxtaz}{}}
& not the reconstruction followed here & broader citation tradition &
useful as a competing label, but not the source-based choice used for
the OE derivation \\
handbook-supported reconstruction &
{\emph{*knéxtaz}}\index[iv]{03pgmc@\ivlangheader{Proto-Germanic}{}!knextaz@\iventry{*knéxtaz}{}}
& regular output: {\emph{cniht}} `knight' &
{\emph{cniht}}\index[iv]{02oe@\ivlangheader{Old English}{West Saxon normalization unmarked}!cniht@\iventry{cniht}{}}
`knight' & exact match between comparative reconstruction and attested
OE noun \\
related plural evidence & same stem family & plural
{\emph{cneohtas}}\index[iv]{02oe@\ivlangheader{Old English}{West Saxon normalization unmarked}!cneohtas@\iventry{cneohtas}{}}
`knights' type background &
{\emph{cneohtas}}\index[iv]{02oe@\ivlangheader{Old English}{West Saxon normalization unmarked}!cneohtas@\iventry{cneohtas}{}}
`knights' & supports the vowel environment, but not the Old English form
here cell \\
\end{longtable}
}

\section{\texorpdfstring{lade --- OE
\emph{hladan}}{lade --- OE hladan}}\label{lade-oe-hladan}

\index[iv]{02oe@\ivlangheader{Old English}{West Saxon normalization unmarked}!hladan@\iventry{hladan}{}}
\index[iv]{03pgmc@\ivlangheader{Proto-Germanic}{}!lathojana@\iventry{*laθōjaną}{}}
\index[iv]{03pgmc@\ivlangheader{Proto-Germanic}{}!xladana@\iventry{*xláðaną}{}}

Derivation: citation reconstruction \emph{*laθōjaną}; form followed here
\emph{*xláðaną} \textgreater{} \emph{hladan} (early analogy).

\subsection{Derivation trace}\label{derivation-trace-115}

Proto input: \emph{*xláðaną}

\begingroup
\setlength{\fboxsep}{6pt}

\noindent\fbox{%
\begin{minipage}{0.97\linewidth}
\small
\begin{tabular}{@{}>{\raggedright\arraybackslash}p{0.359\linewidth}@{\hspace{0.030\linewidth}}>{\raggedright\arraybackslash}p{0.611\linewidth}@{}}
\begin{minipage}[t]{\linewidth}
\centering\textbf{Earlier Germanic changes}\par
\vspace{0.35em}
\raggedright
\centering\textbf{Northwest and West Germanic}\par
\raggedright
\vspace{0.2em}
\begin{tabularx}{\linewidth}{@{}>{\raggedright\arraybackslash}X@{\hspace{0.55em}}>{\raggedright\arraybackslash}p{0.24\linewidth}@{\hspace{0.25em}}}
\mbox{PWGmc Dental Hardening} & \emph{*xládaną} \\
\end{tabularx}
\vspace{0.6em}
\centering\textbf{Early Anglo-Frisian}\par
\raggedright
\vspace{0.2em}
\raggedright [no change]\par
\end{minipage}
&
\begin{minipage}[t]{\linewidth}
\centering\textbf{Old English changes}\par
\vspace{0.35em}
\raggedright
\begin{tabularx}{\linewidth}{@{}>{\raggedright\arraybackslash}X@{\hspace{0.55em}}>{\raggedright\arraybackslash}p{0.18\linewidth}@{\hspace{0.25em}}}
\mbox{EAF Brightening} & \emph{*xlædaną} \\
\mbox{OE A Restoration} & \emph{*xladaną} \\
OE Heavy Syllable Nasal Apocope & \emph{*xladan} \\
\mbox{OE Secondary Nasalization} & \emph{*xladąn} \\
\mbox{OE Weak Tail Reduction} & \emph{*xladan} \\
\end{tabularx}
\end{minipage}
\\
\end{tabular}
\end{minipage}%
} \endgroup

Old English form: \emph{hladan}

\subsection{Reconstruction and comparative
evidence}\label{reconstruction-and-comparative-evidence-98}

Ringe and Taylor cite the strong verb
{\emph{hladan}}\index[iv]{02oe@\ivlangheader{Old English}{West Saxon normalization unmarked}!hladan@\iventry{hladan}{}}
`lade, load' directly (\citeproc{ref-RingeTaylor2014}{Ringe and Taylor
2014, 248}). The citation label \Recon{laθōjaną} `lade' is used here
only as a broader family heading, not as the direct source of the OE
derivation.

The derivational input therefore marks an early stem choice. The entry
follows the strong Verner-grade form that reaches Old English
\emph{hladan} directly.

\subsection{Old English evidence}\label{old-english-evidence-99}

Bosworth-Toller records the verb as \emph{hladan} and preserves the
expected strong-verb paradigm material around it
(\citeproc{ref-BosworthToller1898}{Bosworth and Toller 1898, 559}).
Clark Hall likewise records \emph{hladan}
(\citeproc{ref-ClarkHall1960}{Clark Hall 1960, 159}). The target is an
attested infinitive rather than a reconstructed paradigm cell.

For this entry the relevant comparison form is the infinitive
\emph{hladan} itself. The question is how that attested strong verb
relates to the broader comparative family.

\subsection{Development to Old
English}\label{development-to-old-english-97}

From \Recon{xláðaną} `lade', the verb passes through the expected early
voiced stop stage, Anglo-Frisian brightening, and the later
A-restoration that returns the root vowel to \emph{a} before the full
infinitival ending. Campbell treats verbs such as \emph{hladan} as
showing restored \emph{a} in the present system
(\citeproc{ref-Campbell1959}{Campbell 1959, § 744}). Ringe and
Taylor likewise derive \emph{hladan} from the voiced strong-verb line
(\citeproc{ref-RingeTaylor2014}{Ringe and Taylor 2014, 248}). The
resulting Old English infinitive is \emph{hladan}.

\subsection{Class comparison}\label{class-comparison-2}

The comparison below sets the relevant forms side by side. It separates
the wider weak-verb family label from the strong verb actually reflected
in Old English.

{\def\LTcaptype{none} % do not increment counter
\begin{longtable}[]{@{}
  >{\raggedright\arraybackslash}p{(\linewidth - 8\tabcolsep) * \real{0.2000}}
  >{\raggedright\arraybackslash}p{(\linewidth - 8\tabcolsep) * \real{0.2000}}
  >{\raggedright\arraybackslash}p{(\linewidth - 8\tabcolsep) * \real{0.2000}}
  >{\raggedright\arraybackslash}p{(\linewidth - 8\tabcolsep) * \real{0.2000}}
  >{\raggedright\arraybackslash}p{(\linewidth - 8\tabcolsep) * \real{0.2000}}@{}}
\toprule\noalign{}
\begin{minipage}[b]{\linewidth}\raggedright
Formation / class
\end{minipage} & \begin{minipage}[b]{\linewidth}\raggedright
Candidate input
\end{minipage} & \begin{minipage}[b]{\linewidth}\raggedright
Expected or documented OE outcome
\end{minipage} & \begin{minipage}[b]{\linewidth}\raggedright
OE comparison form
\end{minipage} & \begin{minipage}[b]{\linewidth}\raggedright
Result
\end{minipage} \\
\midrule\noalign{}
\endhead
\bottomrule\noalign{}
\endlastfoot
comparative weak-verb label & *laθōjaną & wider family background &
broader family context & useful family label, but not the direct source
of \emph{hladan} \\
selected strong Verner-grade input & *xláðaną & regular output:
\emph{hladan} & \emph{hladan} & exact match between formation and
attested OE infinitive \\
\end{longtable}
}

\section{\texorpdfstring{lap --- OE
\emph{lappa}}{lap --- OE lappa}}\label{lap-oe-lappa}

\index[iv]{02oe@\ivlangheader{Old English}{West Saxon normalization unmarked}!lappa@\iventry{lappa}{}}
\index[iv]{03pgmc@\ivlangheader{Proto-Germanic}{}!labbaz@\iventry{*lábbaz}{}}
\index[iv]{03pgmc@\ivlangheader{Proto-Germanic}{}!lappo@\iventry{*láppô}{}}

Derivation: citation reconstruction \emph{*lábbaz}; form followed here
\emph{*láppô} \textgreater{} \emph{lappa} (early analogy).

\subsection{Derivation trace}\label{derivation-trace-116}

Proto input: \emph{*láppô}

\begingroup
\setlength{\fboxsep}{6pt}

\noindent\fbox{%
\begin{minipage}{0.97\linewidth}
\small
\begin{tabular}{@{}>{\raggedright\arraybackslash}p{0.359\linewidth}@{\hspace{0.030\linewidth}}>{\raggedright\arraybackslash}p{0.611\linewidth}@{}}
\begin{minipage}[t]{\linewidth}
\centering\textbf{Earlier Germanic changes}\par
\vspace{0.35em}
\raggedright
\centering\textbf{Northwest and West Germanic}\par
\raggedright
\vspace{0.2em}
\raggedright [no change]\par
\vspace{0.6em}
\centering\textbf{Early Anglo-Frisian}\par
\raggedright
\vspace{0.2em}
\raggedright [no change]\par
\end{minipage}
&
\begin{minipage}[t]{\linewidth}
\centering\textbf{Old English changes}\par
\vspace{0.35em}
\raggedright
\begin{tabularx}{\linewidth}{@{}>{\raggedright\arraybackslash}X@{\hspace{0.55em}}>{\raggedright\arraybackslash}p{0.14\linewidth}@{\hspace{0.25em}}}
\mbox{EAF Brightening} & \emph{*læppô} \\
\mbox{OE A Restoration} & \emph{*lappô} \\
OE Unstressed Long Vowel Shortening & \emph{*lappa} \\
\end{tabularx}
\end{minipage}
\\
\end{tabular}
\end{minipage}%
} \endgroup

Old English form: \emph{lappa}

\subsection{Reconstruction and comparative
evidence}\label{reconstruction-and-comparative-evidence-99}

The comparative sources point to a weak noun with \emph{pp}: Orel gives
\Recon{lappōn} `lap', and the Old English dictionary tradition preserves
variant \emph{læppa} `lap' (\citeproc{ref-Orel2003}{Orel 2003, 236};
\citeproc{ref-ClarkHall1960}{Clark Hall 1960, 180}). The derivational
input \Recon{láppô} `lap' follows that evidence.

A competing comparative label \Recon{lábbaz} `lap' has also circulated
for the word family, but the cited handbooks do not make it the direct
source of the Old English weak noun. The form relevant to the OE
development is the weak masculine input \Recon{láppô} `lap'.

\subsection{Old English evidence}\label{old-english-evidence-100}

Campbell cites \emph{lappa} `lap' as a case of restored \emph{a}
(\citeproc{ref-Campbell1959}{Campbell 1959, § 158}). Sievers-Brunner
records \emph{lappa} `lap' beside variant \emph{læppa} `lap'
(\citeproc{ref-SieversBrunner1965}{Brunner 1965, § 10}). The
dictionary tradition also preserves \emph{læppa} `lap'
(\citeproc{ref-ClarkHall1960}{Clark Hall 1960, 180};
\citeproc{ref-BosworthToller1898}{Bosworth and Toller 1898, 613}).

The target of this entry is the restored singular \emph{lappa} `lap'.
The variant \emph{læppa} `lap' and the oblique or plural \emph{leappan}
remain part of the Old English record and help frame the noun's vowel
history.

\subsection{Development to Old
English}\label{development-to-old-english-98}

Campbell explicitly lists \emph{lappa} `lap' among the forms with
restored \emph{a} (\citeproc{ref-Campbell1959}{Campbell 1959, § 158}). Sievers-Brunner records \emph{lappa} `lap' beside variant
\emph{læppa} `lap' at the same Old English stage
(\citeproc{ref-SieversBrunner1965}{Brunner 1965, § 10}). With the
weak masculine input chosen, the selected \emph{lappa} `lap' outcome is
therefore the regular Old English comparison.

\subsection{Stem comparison}\label{stem-comparison-2}

The comparison below sets the relevant forms side by side. It separates
the weak masculine formation from a competing voiced comparative label.

{\def\LTcaptype{none} % do not increment counter
\begin{longtable}[]{@{}
  >{\raggedright\arraybackslash}p{(\linewidth - 8\tabcolsep) * \real{0.2000}}
  >{\raggedright\arraybackslash}p{(\linewidth - 8\tabcolsep) * \real{0.2000}}
  >{\raggedright\arraybackslash}p{(\linewidth - 8\tabcolsep) * \real{0.2000}}
  >{\raggedright\arraybackslash}p{(\linewidth - 8\tabcolsep) * \real{0.2000}}
  >{\raggedright\arraybackslash}p{(\linewidth - 8\tabcolsep) * \real{0.2000}}@{}}
\toprule\noalign{}
\begin{minipage}[b]{\linewidth}\raggedright
Formation / label
\end{minipage} & \begin{minipage}[b]{\linewidth}\raggedright
Candidate input
\end{minipage} & \begin{minipage}[b]{\linewidth}\raggedright
Expected or documented OE outcome
\end{minipage} & \begin{minipage}[b]{\linewidth}\raggedright
OE comparison form
\end{minipage} & \begin{minipage}[b]{\linewidth}\raggedright
Result
\end{minipage} \\
\midrule\noalign{}
\endhead
\bottomrule\noalign{}
\endlastfoot
competing voiced comparative label & *lábbaz & not the form followed for
the OE weak-noun derivation & broader comparative background & useful as
a competing label, but not the source-based choice used here \\
selected weak masculine noun & *láppô & regular output: \emph{lappa} &
\emph{lappa} & exact match between formation and attested OE noun \\
attested OE variant line & same noun family & \emph{læppa},
\emph{leappan} & \emph{læppa} / \emph{leappan} & useful control forms
within the same OE tradition \\
\end{longtable}
}

\section{\texorpdfstring{laugh --- OE
\emph{hliehhan}}{laugh --- OE hliehhan}}\label{laugh-oe-hliehhan}

\index[iv]{02oe@\ivlangheader{Old English}{West Saxon normalization unmarked}!hlaehhan@\iventry{hlæhhan}{}}
\index[iv]{02oe@\ivlangheader{Old English}{West Saxon normalization unmarked}!hleahtor@\iventry{hleahtor}{}}
\index[iv]{02oe@\ivlangheader{Old English}{West Saxon normalization unmarked}!hlehhan@\iventry{hlehhan}{}}
\index[iv]{02oe@\ivlangheader{Old English}{West Saxon normalization unmarked}!hliehhan@\iventry{hliehhan}{}}
\index[iv]{03pgmc@\ivlangheader{Proto-Germanic}{}!hlahtraz@\iventry{*hlahtraz}{}}
\index[iv]{03pgmc@\ivlangheader{Proto-Germanic}{}!lakana@\iventry{*lákaną}{}}
\index[iv]{03pgmc@\ivlangheader{Proto-Germanic}{}!xlaxjana@\iventry{*xláxjaną}{}}

Derivation: citation reconstruction \emph{*lákaną}; form followed here
\emph{*xláxjaną} \textgreater{} \emph{hliehhan} (early analogy).

\subsection{Derivation trace}\label{derivation-trace-117}

Proto input: \emph{*xláxjaną}

\begingroup
\setlength{\fboxsep}{6pt}

\noindent\fbox{%
\begin{minipage}{0.97\linewidth}
\footnotesize
\begin{tabular}{@{}>{\raggedright\arraybackslash}p{0.485\linewidth}@{\hspace{0.030\linewidth}}>{\raggedright\arraybackslash}p{0.485\linewidth}@{}}
\begin{minipage}[t]{\linewidth}
\centering\textbf{Earlier Germanic changes}\par
\vspace{0.35em}
\raggedright
\centering\textbf{Northwest and West Germanic}\par
\raggedright
\vspace{0.2em}
\begin{tabularx}{\linewidth}{@{}>{\raggedright\arraybackslash}X@{\hspace{0.55em}}>{\raggedright\arraybackslash}p{0.22\linewidth}@{\hspace{0.25em}}}
\mbox{PWGmc J Gemination} & \emph{*xláxxjaną} \\
\end{tabularx}
\vspace{0.6em}
\centering\textbf{Early Anglo-Frisian}\par
\raggedright
\vspace{0.2em}
\raggedright [no change]\par
\end{minipage}
&
\begin{minipage}[t]{\linewidth}
\centering\textbf{Old English changes}\par
\vspace{0.35em}
\raggedright
\begin{tabularx}{\linewidth}{@{}>{\raggedright\arraybackslash}X@{\hspace{0.55em}}>{\raggedright\arraybackslash}p{0.26\linewidth}@{\hspace{0.25em}}}
\mbox{EAF Brightening} & \emph{*xlæxxjaną} \\
\mbox{OE Breaking} & \emph{*xleaxxjaną} \\
OE Velar Fricative Palatalization & \emph{*xleaxçjaną} \\
OE Heavy Syllable Nasal Apocope & \emph{*xleaxçjan} \\
\mbox{OE Secondary Nasalization} & \emph{*xleaxçjąn} \\
\mbox{OE I Umlaut} & \emph{*xliexçjąn} \\
\mbox{OE Weak Tail Reduction} & \emph{*xliexçjan} \\
\mbox{OE J Loss After Heavy} & \emph{*xliexçan} \\
\end{tabularx}
\end{minipage}
\\
\end{tabular}
\end{minipage}%
} \endgroup

Old English form: \emph{hliehhan}

\subsection{Reconstruction and comparative
evidence}\label{reconstruction-and-comparative-evidence-100}

The wider Germanic family includes a non-j branch represented here by
the citation label \Recon{lákaną} `laugh', while the derivational input
\Recon{xláxjaną} `laugh' reflects the j-present line behind Old English
{\emph{hliehhan}}\index[iv]{02oe@\ivlangheader{Old English}{West Saxon normalization unmarked}!hliehhan@\iventry{hliehhan}{}}
`laugh'.

This branch supplies the geminate fricative and the vowel development
characteristic of the Old English verb. The comparative family label and
the OE-facing input are therefore related but not identical.

\subsection{Old English evidence}\label{old-english-evidence-101}

Bosworth-Toller records
{\emph{hlihhan}}\index[iv]{02oe@\ivlangheader{Old English}{West Saxon normalization unmarked}!hlihhan@\iventry{hlihhan}{}}
`laugh' as the verb `to laugh'
(\citeproc{ref-BosworthToller1898}{Bosworth and Toller 1898, 551}).
Clark Hall cross-references \emph{hlæhan} `laugh', \emph{hlehhan}
`laugh', and \emph{hlihhan} `laugh' to
{\emph{hliehhan}}\index[iv]{02oe@\ivlangheader{Old English}{West Saxon normalization unmarked}!hliehhan@\iventry{hliehhan}{}}
`laugh' (\citeproc{ref-ClarkHall1960}{Clark Hall 1960, 160--61}).
Bright's glossary likewise gives \emph{hlihhan (hliehhan, hlyhhan)}
(\citeproc{ref-BrightCassidyRingler1971}{Bright 1971, 315}). The target
of this entry is the West Saxon
{\emph{hliehhan}}\index[iv]{02oe@\ivlangheader{Old English}{West Saxon normalization unmarked}!hliehhan@\iventry{hliehhan}{}}
`laugh'.

The variants belong to the same background, but the attested lemma
{\emph{hliehhan}}\index[iv]{02oe@\ivlangheader{Old English}{West Saxon normalization unmarked}!hliehhan@\iventry{hliehhan}{}}
`laugh' supplies the evidence followed here.

\subsection{Development to Old
English}\label{development-to-old-english-99}

From \Recon{xláxjaną} `laugh', West Germanic j-gemination yields the
doubled consonant (\citeproc{ref-Campbell1959}{Campbell 1959, § 407}). Ringe and Taylor derive Old English
{\emph{hliehhan}}\index[iv]{02oe@\ivlangheader{Old English}{West Saxon normalization unmarked}!hliehhan@\iventry{hliehhan}{}}
`laugh' from the j-present branch via breaking before the palatalized
geminate (\citeproc{ref-RingeTaylor2014}{Ringe and Taylor 2014, 240}).
They separately compare the related noun \emph{hleahtor} as the outcome
of \Recon{hlahtraz} `laugh' (\citeproc{ref-RingeTaylor2014}{Ringe and
Taylor 2014, 328}).

\subsection{Branch comparison}\label{branch-comparison}

The comparison below sets the relevant forms side by side. It separates
the wider non-j family label from the j-present branch actually
reflected in Old English.

{\def\LTcaptype{none} % do not increment counter
\begin{longtable}[]{@{}
  >{\raggedright\arraybackslash}p{(\linewidth - 8\tabcolsep) * \real{0.2000}}
  >{\raggedright\arraybackslash}p{(\linewidth - 8\tabcolsep) * \real{0.2000}}
  >{\raggedright\arraybackslash}p{(\linewidth - 8\tabcolsep) * \real{0.2000}}
  >{\raggedright\arraybackslash}p{(\linewidth - 8\tabcolsep) * \real{0.2000}}
  >{\raggedright\arraybackslash}p{(\linewidth - 8\tabcolsep) * \real{0.2000}}@{}}
\toprule\noalign{}
\begin{minipage}[b]{\linewidth}\raggedright
Formation / branch
\end{minipage} & \begin{minipage}[b]{\linewidth}\raggedright
Candidate input
\end{minipage} & \begin{minipage}[b]{\linewidth}\raggedright
Expected or documented OE outcome
\end{minipage} & \begin{minipage}[b]{\linewidth}\raggedright
OE comparison form
\end{minipage} & \begin{minipage}[b]{\linewidth}\raggedright
Result
\end{minipage} \\
\midrule\noalign{}
\endhead
\bottomrule\noalign{}
\endlastfoot
wider non-j family & *lákaną & comparative background outside the
selected OE line & wider family context & useful family label, but not
the direct source of \emph{hliehhan} \\
selected j-present branch & *xláxjaną & regular output: \emph{hliehhan}
& \emph{hliehhan} & exact match between branch and attested OE lemma \\
attested OE variants & same OE verb line & \emph{hlæhhan},
\emph{hlehhan} & \emph{hlæhhan} / \emph{hlehhan} & genuine variant
evidence, but secondary to the form compared here \\
\end{longtable}
}

\section{\texorpdfstring{loam --- OE
\emph{lām}}{loam --- OE lām}}\label{loam-oe-lux101m}

\index[iv]{02oe@\ivlangheader{Old English}{West Saxon normalization unmarked}!lam@\iventry{lām}{}}
\index[iv]{03pgmc@\ivlangheader{Proto-Germanic}{}!laimon@\iventry{*laimōn}{}}
\index[iv]{08preoe@\ivlangheader{Pre-Old English / prehistoric English}{}!laima~01@\iventry{*láimą}{transp.}}

Derivation: citation reconstruction \emph{*laimōn}; form followed here
\emph{*láimą} \textgreater{} \emph{lām} (early analogy).

\subsection{Derivation trace}\label{derivation-trace-118}

Proto input: \emph{*láimą}

\begingroup
\setlength{\fboxsep}{6pt}

\noindent\fbox{%
\begin{minipage}{0.97\linewidth}
\small
\begin{tabular}{@{}>{\raggedright\arraybackslash}p{0.485\linewidth}@{\hspace{0.030\linewidth}}>{\raggedright\arraybackslash}p{0.485\linewidth}@{}}
\begin{minipage}[t]{\linewidth}
\centering\textbf{Earlier Germanic changes}\par
\vspace{0.35em}
\raggedright
\centering\textbf{Northwest and West Germanic}\par
\raggedright
\vspace{0.2em}
\raggedright [no change]\par
\vspace{0.6em}
\centering\textbf{Early Anglo-Frisian}\par
\raggedright
\vspace{0.2em}
\begin{tabularx}{\linewidth}{@{}>{\raggedright\arraybackslash}X@{\hspace{0.55em}}>{\raggedright\arraybackslash}p{0.14\linewidth}@{\hspace{0.25em}}}
\mbox{EAF Ai Monophthongization} & \emph{*lāmą} \\
\end{tabularx}
\end{minipage}
&
\begin{minipage}[t]{\linewidth}
\centering\textbf{Old English changes}\par
\vspace{0.35em}
\raggedright
\begin{tabularx}{\linewidth}{@{}>{\raggedright\arraybackslash}X@{\hspace{0.55em}}>{\raggedright\arraybackslash}p{0.14\linewidth}@{\hspace{0.25em}}}
OE Heavy Syllable Nasal Apocope & \emph{*lām} \\
\end{tabularx}
\end{minipage}
\\
\end{tabular}
\end{minipage}%
} \endgroup

Old English form: \emph{lām}

\subsection{Reconstruction and comparative
evidence}\label{reconstruction-and-comparative-evidence-101}

The inherited comparative noun is given as \Recon{laimōn} `loam' or
\emph{*laiman-}, and both Orel and Kroonen identify Old English
\emph{lām} `loam' as a neuter reflex of that family
(\citeproc{ref-Orel2003}{Orel 2003, 272};
\citeproc{ref-Kroonen2013}{Kroonen 2013, 363}). The form followed here,
\Recon{láimą} `loam', differs from the comparative headword because it
represents the stem class that matches the Old English noun most
directly.

This is therefore a class shift within the history of the English branch
rather than a dispute about the OE target itself.

\subsection{Old English evidence}\label{old-english-evidence-102}

Old English dictionaries record the noun as \emph{lām} `loam', a neuter
word for `loam, clay, mud' (\citeproc{ref-BosworthToller1898}{Bosworth
and Toller 1898, 604}; \citeproc{ref-ClarkHall1960}{Clark Hall 1960,
196}). The target is an attested citation form rather than a
reconstructed comparator.

The relevant question is not whether \emph{lām} `loam' is Old English,
but which inherited formation best accounts for that attested neuter
noun.

\subsection{Development to Old
English}\label{development-to-old-english-100}

From \Recon{láimą} `loam', regular monophthongization of \emph{ai} and
the later loss of the final nasal syllable yield \emph{lām} `loam'. With
that OE-facing input, the phonological development is straightforward.

\subsection{Class comparison}\label{class-comparison-3}

The comparison below sets the inherited comparative n-stem label beside
the OE-facing stem class used to derive the attested noun.

{\def\LTcaptype{none} % do not increment counter
\begin{longtable}[]{@{}
  >{\raggedright\arraybackslash}p{(\linewidth - 8\tabcolsep) * \real{0.2000}}
  >{\raggedright\arraybackslash}p{(\linewidth - 8\tabcolsep) * \real{0.2000}}
  >{\raggedright\arraybackslash}p{(\linewidth - 8\tabcolsep) * \real{0.2000}}
  >{\raggedright\arraybackslash}p{(\linewidth - 8\tabcolsep) * \real{0.2000}}
  >{\raggedright\arraybackslash}p{(\linewidth - 8\tabcolsep) * \real{0.2000}}@{}}
\toprule\noalign{}
\begin{minipage}[b]{\linewidth}\raggedright
Formation / class
\end{minipage} & \begin{minipage}[b]{\linewidth}\raggedright
Candidate input
\end{minipage} & \begin{minipage}[b]{\linewidth}\raggedright
Expected or documented OE outcome
\end{minipage} & \begin{minipage}[b]{\linewidth}\raggedright
OE comparison form
\end{minipage} & \begin{minipage}[b]{\linewidth}\raggedright
Result
\end{minipage} \\
\midrule\noalign{}
\endhead
\bottomrule\noalign{}
\endlastfoot
inherited comparative noun & *laimōn & comparative family background &
wider family context & useful headword, but not the direct OE-facing
input \\
OE-facing stem class followed here & *láimą & regular output: \emph{lām}
& \emph{lām} & exact match between the form followed here and the
attested OE noun \\
\end{longtable}
}

\section{\texorpdfstring{lung --- OE
\emph{lungen}}{lung --- OE lungen}}\label{lung-oe-lungen}

\index[iv]{02oe@\ivlangheader{Old English}{West Saxon normalization unmarked}!lungen@\iventry{lungen}{}}
\index[iv]{03pgmc@\ivlangheader{Proto-Germanic}{}!lunganjo@\iventry{*lúnganjō}{}}
\index[iv]{03pgmc@\ivlangheader{Proto-Germanic}{}!lungo@\iventry{*lungō}{}}

Derivation: citation reconstruction \emph{*lungō}; form followed here
\emph{*lúnganjō} \textgreater{} \emph{lungen} (early analogy).

\subsection{Derivation trace}\label{derivation-trace-119}

Proto input: \emph{*lúnganjō}

\begingroup
\setlength{\fboxsep}{6pt}

\noindent\fbox{%
\begin{minipage}{0.97\linewidth}
\small
\begin{tabular}{@{}>{\raggedright\arraybackslash}p{0.485\linewidth}@{\hspace{0.030\linewidth}}>{\raggedright\arraybackslash}p{0.485\linewidth}@{}}
\begin{minipage}[t]{\linewidth}
\centering\textbf{Earlier Germanic changes}\par
\vspace{0.35em}
\raggedright
\centering\textbf{Northwest and West Germanic}\par
\raggedright
\vspace{0.2em}
\begin{tabularx}{\linewidth}{@{}>{\raggedright\arraybackslash}X@{\hspace{0.55em}}>{\raggedright\arraybackslash}p{0.22\linewidth}@{\hspace{0.25em}}}
\mbox{PWGmc J Gemination} & \emph{*lúngannjō} \\
\mbox{PNWGmc Final Long O Raising} & \emph{*lúngannju} \\
\end{tabularx}
\vspace{0.6em}
\centering\textbf{Early Anglo-Frisian}\par
\raggedright
\vspace{0.2em}
\raggedright [no change]\par
\end{minipage}
&
\begin{minipage}[t]{\linewidth}
\centering\textbf{Old English changes}\par
\vspace{0.35em}
\raggedright
\begin{tabularx}{\linewidth}{@{}>{\raggedright\arraybackslash}X@{\hspace{0.55em}}>{\raggedright\arraybackslash}p{0.24\linewidth}@{\hspace{0.25em}}}
\mbox{OE I Umlaut} & \emph{*lúngennju} \\
\mbox{OE High Vowel Apocope} & \emph{*lúngennj} \\
\mbox{OE J Loss After Heavy} & \emph{*lúngenn} \\
OE Final Geminate Simplification & \emph{*lúngen} \\
\end{tabularx}
\end{minipage}
\\
\end{tabular}
\end{minipage}%
} \endgroup

Old English form: \emph{lungen}

\subsection{Reconstruction and comparative
evidence}\label{reconstruction-and-comparative-evidence-102}

Kroonen treats the basic noun as \emph{*lungōn-} and also cites an
OE-facing derivative \emph{*lungunjō-}, continued by Old English
{\emph{lungen}}\index[iv]{02oe@\ivlangheader{Old English}{West Saxon normalization unmarked}!lungen@\iventry{lungen}{}}
`lung' and close West Germanic cognates
(\citeproc{ref-Kroonen2013}{Kroonen 2013, 384}). The derivational input
\Recon{lúnganjō} `lung' models that derived feminine formation rather
than the base noun. The notation differs slightly from Kroonen's
\emph{*lungunjō-}, but both point to the same derived feminine line.

The difference between the citation label and the derivational input is
therefore derivational. Old English
{\emph{lungen}}\index[iv]{02oe@\ivlangheader{Old English}{West Saxon normalization unmarked}!lungen@\iventry{lungen}{}}
`lung' is not a direct reflex of the bare base noun \Recon{lungō}
`lung'; it belongs to an expanded feminine formation.

\subsection{Old English evidence}\label{old-english-evidence-103}

Old English dictionaries record the noun as
{\emph{lungen}}\index[iv]{02oe@\ivlangheader{Old English}{West Saxon normalization unmarked}!lungen@\iventry{lungen}{}}
`lung', with inflected forms such as \emph{lungenne} and \emph{lungena}
(\citeproc{ref-BosworthToller1898}{Bosworth and Toller 1898, 634}).
Clark Hall also preserves a small family of compounds such as
\emph{lungenādl} `lung-disease', \emph{lungensealf}, and
\emph{lungenwyrt} (\citeproc{ref-ClarkHall1960}{Clark Hall 1960, 191}).

The target is an attested Old English lexeme with its own paradigm, not
a rescued inflectional cell.

\subsection{Development to Old
English}\label{development-to-old-english-101}

From the selected derived input, the expected derivational consonant and
vowel adjustments lead to
{\emph{lungen}}\index[iv]{02oe@\ivlangheader{Old English}{West Saxon normalization unmarked}!lungen@\iventry{lungen}{}}
`lung'. Once the expanded feminine formation is chosen, the Old English
outcome is regular.

\subsection{Formation comparison}\label{formation-comparison-3}

The comparison below sets the relevant forms side by side. It separates
the base noun from the derived feminine formation reflected in Old
English.

{\def\LTcaptype{none} % do not increment counter
\begin{longtable}[]{@{}
  >{\raggedright\arraybackslash}p{(\linewidth - 8\tabcolsep) * \real{0.2000}}
  >{\raggedright\arraybackslash}p{(\linewidth - 8\tabcolsep) * \real{0.2000}}
  >{\raggedright\arraybackslash}p{(\linewidth - 8\tabcolsep) * \real{0.2000}}
  >{\raggedright\arraybackslash}p{(\linewidth - 8\tabcolsep) * \real{0.2000}}
  >{\raggedright\arraybackslash}p{(\linewidth - 8\tabcolsep) * \real{0.2000}}@{}}
\toprule\noalign{}
\begin{minipage}[b]{\linewidth}\raggedright
Formation
\end{minipage} & \begin{minipage}[b]{\linewidth}\raggedright
Candidate input
\end{minipage} & \begin{minipage}[b]{\linewidth}\raggedright
Expected or documented OE outcome
\end{minipage} & \begin{minipage}[b]{\linewidth}\raggedright
OE comparison form
\end{minipage} & \begin{minipage}[b]{\linewidth}\raggedright
Result
\end{minipage} \\
\midrule\noalign{}
\endhead
\bottomrule\noalign{}
\endlastfoot
base noun & *lungō & base-noun outcome without the OE derivative suffix
& broader family context & useful headword, but not the direct source of
\emph{lungen} \\
derived OE-facing formation & *lúnganjō & regular output: \emph{lungen}
& \emph{lungen} & exact match between the derived formation and the
attested OE noun \\
Kroonen's cited derivative & \emph{*lungunjō-} & comparative support for
the same OE-facing formation & \emph{lungen} and cognate set & supports
the derived feminine formation, with notation differing from the
normalized input form used here \\
\end{longtable}
}

\section{\texorpdfstring{navel --- OE
\emph{nafola}}{navel --- OE nafola}}\label{navel-oe-nafola}

\index[iv]{02oe@\ivlangheader{Old English}{West Saxon normalization unmarked}!naefla@\iventry{næfla}{}}
\index[iv]{02oe@\ivlangheader{Old English}{West Saxon normalization unmarked}!nafola@\iventry{nafola}{}}
\index[iv]{03pgmc@\ivlangheader{Proto-Germanic}{}!nablo@\iventry{*nablô}{}}
\index[iv]{05pwgmc@\ivlangheader{Proto-West Germanic}{}!nabulo@\iventry{*nábulô}{}}

Derivation: citation reconstruction \emph{*nablô}; form followed here
\emph{*nábulô} \textgreater{} \emph{nafola} (early analogy).

\subsection{Derivation trace}\label{derivation-trace-120}

Proto input: \emph{*nábulô}

\begingroup
\setlength{\fboxsep}{6pt}

\noindent\fbox{%
\begin{minipage}{0.97\linewidth}
\small
\begin{tabular}{@{}>{\raggedright\arraybackslash}p{0.359\linewidth}@{\hspace{0.030\linewidth}}>{\raggedright\arraybackslash}p{0.611\linewidth}@{}}
\begin{minipage}[t]{\linewidth}
\centering\textbf{Earlier Germanic changes}\par
\vspace{0.35em}
\raggedright
\centering\textbf{Northwest and West Germanic}\par
\raggedright
\vspace{0.2em}
\raggedright [no change]\par
\vspace{0.6em}
\centering\textbf{Early Anglo-Frisian}\par
\raggedright
\vspace{0.2em}
\raggedright [no change]\par
\end{minipage}
&
\begin{minipage}[t]{\linewidth}
\centering\textbf{Old English changes}\par
\vspace{0.35em}
\raggedright
\begin{tabularx}{\linewidth}{@{}>{\raggedright\arraybackslash}X@{\hspace{0.55em}}>{\raggedright\arraybackslash}p{0.14\linewidth}@{\hspace{0.25em}}}
\mbox{EAF Brightening} & \emph{*næbulô} \\
\mbox{OE A Restoration} & \emph{*nabulô} \\
\mbox{PGmc B Allophony} & \emph{*naβulô} \\
OE Unstressed Long Vowel Shortening & \emph{*naβula} \\
\mbox{OE Med Unstressed U Lowering} & \emph{*naβola} \\
\end{tabularx}
\end{minipage}
\\
\end{tabular}
\end{minipage}%
} \endgroup

Old English form: \emph{nafola}

\subsection{Reconstruction and comparative
evidence}\label{reconstruction-and-comparative-evidence-103}

Kroonen instead gives a nasal-suffix navel formation with Old English
{\emph{nafela}}\index[iv]{02oe@\ivlangheader{Old English}{West Saxon normalization unmarked}!nafela@\iventry{nafela}{}}
`navel' among its reflexes (\citeproc{ref-Kroonen2013}{Kroonen 2013,
420}), while Ringe and Taylor give the derivational pathway
\Recon{nabulō} `navel' \textgreater{} \Recon{næbula} `navel'
\textgreater{} \emph{nafola} `navel'
(\citeproc{ref-RingeTaylor2014}{Ringe and Taylor 2014, 270}). The
difference is one of stage and notation rather than of lexeme identity:
the derivational input
\Recon{nábulô}\index[iv]{05pwgmc@\ivlangheader{Proto-West Germanic}{}!nabulo@\iventry{*nábulô}{}}
`navel' is the pre-syncope form needed for the Old English development.

For the Old English comparison, the crucial point is simply that the
pre-OE form still contains a medial vowel.

\subsection{Old English evidence}\label{old-english-evidence-104}

Ringe and Taylor note the early West Saxon shift
{\emph{nafola}}\index[iv]{02oe@\ivlangheader{Old English}{West Saxon normalization unmarked}!nafola@\iventry{nafola}{}}
`navel' \textgreater{}
{\emph{nafela}}\index[iv]{02oe@\ivlangheader{Old English}{West Saxon normalization unmarked}!nafela@\iventry{nafela}{}}
`navel' (\citeproc{ref-RingeTaylor2014}{Ringe and Taylor 2014, 336}).
Campbell likewise records \emph{nafela} `navel' beside Corpus
{\emph{nabula}}\index[iv]{02oe@\ivlangheader{Old English}{West Saxon normalization unmarked}!nabula@\iventry{nabula}{}}
`navel' (\citeproc{ref-Campbell1959}{Campbell 1959, § 159}). The
target of this entry is the nominative singular
{\emph{nafola}}\index[iv]{02oe@\ivlangheader{Old English}{West Saxon normalization unmarked}!nafola@\iventry{nafola}{}}
`navel', the form that matches the selected derivational pathway most
directly.

{\emph{nafela}}\index[iv]{02oe@\ivlangheader{Old English}{West Saxon normalization unmarked}!nafela@\iventry{nafela}{}}
`navel' is the better-known later West Saxon spelling, while
{\emph{nabula}}\index[iv]{02oe@\ivlangheader{Old English}{West Saxon normalization unmarked}!nabula@\iventry{nabula}{}}
`navel' preserves a less reduced medial vowel. These forms belong to the
same lexical history, but this entry is centered on
{\emph{nafola}}\index[iv]{02oe@\ivlangheader{Old English}{West Saxon normalization unmarked}!nafola@\iventry{nafola}{}}
`navel'.

\subsection{Development to Old
English}\label{development-to-old-english-102}

Ringe and Taylor give the pre-OE line \Recon{nabulō} `navel'
\textgreater{} \Recon{næbula} `navel' \textgreater{} OE \emph{nafola}
`navel' (\citeproc{ref-RingeTaylor2014}{Ringe and Taylor 2014, 270}).
The trace represents that same development with stress-marked notation
and explicit intermediate weakening. Intervocalic \emph{b} then surfaces
as \emph{f}, and final weak-tail shortening gives \emph{nafola} `navel'.

The medial vowel is still present when A-restoration applies. That is
why the selected pre-syncope input differs from the syncopated
comparative headword.

\subsection{Stage comparison}\label{stage-comparison-1}

The comparison below sets the relevant forms side by side. It separates
the comparative citation form from the pre-syncope input and from the
later OE spellings.

{\def\LTcaptype{none} % do not increment counter
\begin{longtable}[]{@{}
  >{\raggedright\arraybackslash}p{(\linewidth - 8\tabcolsep) * \real{0.2000}}
  >{\raggedright\arraybackslash}p{(\linewidth - 8\tabcolsep) * \real{0.2000}}
  >{\raggedright\arraybackslash}p{(\linewidth - 8\tabcolsep) * \real{0.2000}}
  >{\raggedright\arraybackslash}p{(\linewidth - 8\tabcolsep) * \real{0.2000}}
  >{\raggedright\arraybackslash}p{(\linewidth - 8\tabcolsep) * \real{0.2000}}@{}}
\toprule\noalign{}
\begin{minipage}[b]{\linewidth}\raggedright
Formation / stage
\end{minipage} & \begin{minipage}[b]{\linewidth}\raggedright
Candidate input
\end{minipage} & \begin{minipage}[b]{\linewidth}\raggedright
Expected or documented OE outcome
\end{minipage} & \begin{minipage}[b]{\linewidth}\raggedright
OE comparison form
\end{minipage} & \begin{minipage}[b]{\linewidth}\raggedright
Result
\end{minipage} \\
\midrule\noalign{}
\endhead
\bottomrule\noalign{}
\endlastfoot
syncopated comparative headword &
{\emph{*nablô}}\index[iv]{03pgmc@\ivlangheader{Proto-Germanic}{}!nablo@\iventry{*nablô}{}}
& reduced \emph{næfla}-type outcome rather than
{\emph{nafola}}\index[iv]{02oe@\ivlangheader{Old English}{West Saxon normalization unmarked}!nafola@\iventry{nafola}{}}
`navel' & not the Old English form here & useful citation form, but too
reduced for the pathway modeled here \\
selected pre-syncope input &
{\emph{*nábulô}}\index[iv]{05pwgmc@\ivlangheader{Proto-West Germanic}{}!nabulo@\iventry{*nábulô}{}}
& regular output: {\emph{nafola}} `navel' &
{\emph{nafola}}\index[iv]{02oe@\ivlangheader{Old English}{West Saxon normalization unmarked}!nafola@\iventry{nafola}{}}
`navel' & exact match between derivational input and target \\
later OE reduction stages & same lexical history & attested
{\emph{nafela}}\index[iv]{02oe@\ivlangheader{Old English}{West Saxon normalization unmarked}!nafela@\iventry{nafela}{}}
`navel'; Corpus
{\emph{nabula}}\index[iv]{02oe@\ivlangheader{Old English}{West Saxon normalization unmarked}!nabula@\iventry{nabula}{}}
`navel' &
{\emph{nafela}}\index[iv]{02oe@\ivlangheader{Old English}{West Saxon normalization unmarked}!nafela@\iventry{nafela}{}}
`navel' /
{\emph{nabula}}\index[iv]{02oe@\ivlangheader{Old English}{West Saxon normalization unmarked}!nabula@\iventry{nabula}{}}
`navel' & related OE spellings, but not the chosen comparator \\
\end{longtable}
}

\section{\texorpdfstring{neck --- OE
\emph{hnecca}}{neck --- OE hnecca}}\label{neck-oe-hnecca}

\index[iv]{02oe@\ivlangheader{Old English}{West Saxon normalization unmarked}!hnecca@\iventry{hnecca}{}}
\index[iv]{03pgmc@\ivlangheader{Proto-Germanic}{}!xnakkaz@\iventry{*xnákkaz}{}}
\index[iv]{03pgmc@\ivlangheader{Proto-Germanic}{}!xnekko@\iventry{*xnékkô}{}}

Derivation: citation reconstruction \emph{*xnákkaz}; form followed here
\emph{*xnékkô} \textgreater{} \emph{hnecca} (early analogy).

\subsection{Derivation trace}\label{derivation-trace-121}

Proto input: \emph{*xnékkô}

\begingroup
\setlength{\fboxsep}{6pt}

\noindent\fbox{%
\begin{minipage}{0.97\linewidth}
\small
\begin{tabular}{@{}>{\raggedright\arraybackslash}p{0.359\linewidth}@{\hspace{0.030\linewidth}}>{\raggedright\arraybackslash}p{0.611\linewidth}@{}}
\begin{minipage}[t]{\linewidth}
\centering\textbf{Earlier Germanic changes}\par
\vspace{0.35em}
\raggedright
\centering\textbf{Northwest and West Germanic}\par
\raggedright
\vspace{0.2em}
\raggedright [no change]\par
\vspace{0.6em}
\centering\textbf{Early Anglo-Frisian}\par
\raggedright
\vspace{0.2em}
\raggedright [no change]\par
\end{minipage}
&
\begin{minipage}[t]{\linewidth}
\centering\textbf{Old English changes}\par
\vspace{0.35em}
\raggedright
\begin{tabularx}{\linewidth}{@{}>{\raggedright\arraybackslash}X@{\hspace{0.55em}}>{\raggedright\arraybackslash}p{0.14\linewidth}@{\hspace{0.25em}}}
OE Unstressed Long Vowel Shortening & \emph{*xnékka} \\
\end{tabularx}
\end{minipage}
\\
\end{tabular}
\end{minipage}%
} \endgroup

Old English form: \emph{hnecca}

\subsection{Reconstruction and comparative
evidence}\label{reconstruction-and-comparative-evidence-104}

The noun belongs to an ablauting n-stem family. Kroonen reconstructs a
paradigm with nominative \Recon{hnekkō} `neck', genitive \Recon{hnukkaz}
`neck', and accusative plural \Recon{hnakkuns} `neck', and he places Old
English
{\emph{hnecca}}\index[iv]{02oe@\ivlangheader{Old English}{West Saxon normalization unmarked}!hnecca@\iventry{hnecca}{}}
`neck, nape' among the e-grade descendants
(\citeproc{ref-Kroonen2011}{Kroonen 2011, 167}). Kluge-Seebold likewise
identifies \emph{ae. hnecca} as an ablaut partner of the a-grade
\emph{Nacken} family (\citeproc{ref-KlugeSeebold2011}{Kluge 2011, 347}).

A competing comparative label
\Recon{xnákkaz}\index[iv]{03pgmc@\ivlangheader{Proto-Germanic}{}!xnakkaz@\iventry{*xnákkaz}{}}
`neck' belongs to the wider family, and Orel also gives an a-grade
headword line (\citeproc{ref-Orel2003}{Orel 2003, 218}). The
derivational input
\Recon{xnékkô}\index[iv]{03pgmc@\ivlangheader{Proto-Germanic}{}!xnekko@\iventry{*xnékkô}{}}
`neck', however, is the form that matches the Old English branch.

\subsection{Old English evidence}\label{old-english-evidence-105}

Clark Hall records the weak masculine noun
{\emph{hnecca}}\index[iv]{02oe@\ivlangheader{Old English}{West Saxon normalization unmarked}!hnecca@\iventry{hnecca}{}}
`neck, nape' (\citeproc{ref-ClarkHall1960}{Clark Hall 1960, 162}).
Bosworth-Toller likewise records
{\emph{hnecca}}\index[iv]{02oe@\ivlangheader{Old English}{West Saxon normalization unmarked}!hnecca@\iventry{hnecca}{}}
`neck, nape' (\citeproc{ref-BosworthToller1898}{Bosworth and Toller
1898, 567}). The target is therefore an attested citation form, not an
oblique cell or a reconstructed lemma.

The phonological question is upstream of the Old English evidence. The
attested noun already shows that the branch continued an e-grade form
rather than the a-grade seen in much of the continental material.

\subsection{Development to Old
English}\label{development-to-old-english-103}

From
\Recon{xnékkô}\index[iv]{03pgmc@\ivlangheader{Proto-Germanic}{}!xnekko@\iventry{*xnékkô}{}}
`neck', the derivation is straightforward. The trace shortens the final
long vowel to \Recon{xnékka} `neck', and Old English orthography gives
{\emph{hnecca}}\index[iv]{02oe@\ivlangheader{Old English}{West Saxon normalization unmarked}!hnecca@\iventry{hnecca}{}}
`neck'.

The derivation depends on the earlier selection of the e-grade weak-noun
form continued by Old English.

\subsection{Stem comparison}\label{stem-comparison-3}

The comparison below sets the relevant forms side by side. It separates
the wider a-grade family from the selected e-grade Old English branch.

{\def\LTcaptype{none} % do not increment counter
\begin{longtable}[]{@{}
  >{\raggedright\arraybackslash}p{(\linewidth - 8\tabcolsep) * \real{0.2000}}
  >{\raggedright\arraybackslash}p{(\linewidth - 8\tabcolsep) * \real{0.2000}}
  >{\raggedright\arraybackslash}p{(\linewidth - 8\tabcolsep) * \real{0.2000}}
  >{\raggedright\arraybackslash}p{(\linewidth - 8\tabcolsep) * \real{0.2000}}
  >{\raggedright\arraybackslash}p{(\linewidth - 8\tabcolsep) * \real{0.2000}}@{}}
\toprule\noalign{}
\begin{minipage}[b]{\linewidth}\raggedright
Formation / label
\end{minipage} & \begin{minipage}[b]{\linewidth}\raggedright
Candidate input
\end{minipage} & \begin{minipage}[b]{\linewidth}\raggedright
Expected or documented OE outcome
\end{minipage} & \begin{minipage}[b]{\linewidth}\raggedright
OE comparison form
\end{minipage} & \begin{minipage}[b]{\linewidth}\raggedright
Result
\end{minipage} \\
\midrule\noalign{}
\endhead
\bottomrule\noalign{}
\endlastfoot
competing comparative label &
{\emph{*xnákkaz}}\index[iv]{03pgmc@\ivlangheader{Proto-Germanic}{}!xnakkaz@\iventry{*xnákkaz}{}}
& broader a-grade family rather than the selected OE source &
continental
{\emph{Nacken}}\index[iv]{17german@\ivlangheader{German}{}!nacken@\iventry{Nacken}{}}
line & useful family label, but not the input followed for the Old
English derivation \\
weak noun with a-grade &
{\emph{*xnakkô}}\index[iv]{03pgmc@\ivlangheader{Proto-Germanic}{}!xnakko@\iventry{*xnakkô}{}}
& expected
{\emph{hnacca}}\index[iv]{02oe@\ivlangheader{Old English}{West Saxon normalization unmarked}!hnacca@\iventry{hnacca}{}}
type outcome &
{\emph{hnacca}}\index[iv]{02oe@\ivlangheader{Old English}{West Saxon normalization unmarked}!hnacca@\iventry{hnacca}{}}
& fixes the class, but not the vowel grade \\
selected e-grade nominative &
{\emph{*xnékkô}}\index[iv]{03pgmc@\ivlangheader{Proto-Germanic}{}!xnekko@\iventry{*xnékkô}{}}
& regular output: {\emph{hnecca}} &
{\emph{hnecca}}\index[iv]{02oe@\ivlangheader{Old English}{West Saxon normalization unmarked}!hnecca@\iventry{hnecca}{}}
& exact match between derivational input and attested OE noun \\
oblique paradigm background &
{\emph{*hnukkaz}}\index[iv]{03pgmc@\ivlangheader{Proto-Germanic}{}!hnukkaz@\iventry{*hnukkaz}{}},
{\emph{*hnakkuns}}\index[iv]{03pgmc@\ivlangheader{Proto-Germanic}{}!hnakkuns@\iventry{*hnakkuns}{}}
& ON/OHG/German a-grade continuation &
{\emph{hnakki}}\index[iv]{09on@\ivlangheader{Old Norse}{}!hnakki@\iventry{hnakki}{}}
/
{\emph{Nacken}}\index[iv]{17german@\ivlangheader{German}{}!nacken@\iventry{Nacken}{}}
& shows the wider ablaut family, but not the chosen OE branch \\
\end{longtable}
}

\section{\texorpdfstring{needle --- OE
\emph{nǣdl}}{needle --- OE nǣdl}}\label{needle-oe-nux1e3dl}

\index[iv]{02oe@\ivlangheader{Old English}{West Saxon normalization unmarked}!naedl@\iventry{nǣdl}{}}
\index[iv]{03pgmc@\ivlangheader{Proto-Germanic}{}!nedlo@\iventry{*nḗðlō}{}}
\index[iv]{03pgmc@\ivlangheader{Proto-Germanic}{}!nethlo@\iventry{*nḗθlō}{}}

Derivation: citation reconstruction \emph{*nḗθlō}; form followed here
\emph{*nḗðlō} \textgreater{} \emph{nǣdl} (early analogy).

\subsection{Derivation trace}\label{derivation-trace-122}

Proto input: \emph{*nḗðlō}

\begingroup
\setlength{\fboxsep}{6pt}

\noindent\fbox{%
\begin{minipage}{0.97\linewidth}
\small
\begin{tabular}{@{}>{\raggedright\arraybackslash}p{0.611\linewidth}@{\hspace{0.030\linewidth}}>{\raggedright\arraybackslash}p{0.359\linewidth}@{}}
\begin{minipage}[t]{\linewidth}
\centering\textbf{Earlier Germanic changes}\par
\vspace{0.35em}
\raggedright
\centering\textbf{Northwest and West Germanic}\par
\raggedright
\vspace{0.2em}
\begin{tabularx}{\linewidth}{@{}>{\raggedright\arraybackslash}X@{\hspace{0.55em}}>{\raggedright\arraybackslash}p{0.14\linewidth}@{\hspace{0.25em}}}
\mbox{PNWGmc Long E Lowering} & \emph{*nāðlō} \\
\mbox{PWGmc Dental Hardening} & \emph{*nādlō} \\
\mbox{PNWGmc Final Long O Raising} & \emph{*nādlu} \\
\mbox{EAF Long A Fronting} & \emph{*nǣdlu} \\
\end{tabularx}
\vspace{0.6em}
\centering\textbf{Early Anglo-Frisian}\par
\raggedright
\vspace{0.2em}
\raggedright [no change]\par
\end{minipage}
&
\begin{minipage}[t]{\linewidth}
\centering\textbf{Old English changes}\par
\vspace{0.35em}
\raggedright
\begin{tabularx}{\linewidth}{@{}>{\raggedright\arraybackslash}X@{\hspace{0.55em}}>{\raggedright\arraybackslash}p{0.14\linewidth}@{\hspace{0.25em}}}
\mbox{OE High Vowel Apocope} & \emph{*nǣdl} \\
\end{tabularx}
\end{minipage}
\\
\end{tabular}
\end{minipage}%
} \endgroup

Old English form: \emph{nǣdl}

\subsection{Reconstruction and comparative
evidence}\label{reconstruction-and-comparative-evidence-105}

Ringe and Taylor treat the word as a voiced/voiceless alternant, citing
\Recon{nēþlō} `needle' \textasciitilde{} \Recon{nēdlō} `needle'
\textgreater{} OE \emph{nédl} `needle'
(\citeproc{ref-RingeTaylor2014}{Ringe and Taylor 2014, 329}). The form
followed here,
\Recon{nḗðlō}\index[iv]{03pgmc@\ivlangheader{Proto-Germanic}{}!nedlo@\iventry{*nḗðlō}{}}
`needle', is the voiced Verner-grade form used for the Old English
comparison, while the citation form
\Recon{nḗθlō}\index[iv]{03pgmc@\ivlangheader{Proto-Germanic}{}!nethlo@\iventry{*nḗθlō}{}}
`needle' remains the broader lexeme label.

The development discussed here follows the Ringe-Taylor alternant
framework.

\subsection{Old English evidence}\label{old-english-evidence-106}

Clark Hall records the attested citation form \emph{nǣdl} `needle'
(\citeproc{ref-ClarkHall1960}{Clark Hall 1960, 210}). Campbell lists
\emph{nédl} `needle' among the expected unbroken forms after \emph{t}
and \emph{d} (\citeproc{ref-Campbell1959}{Campbell 1959, § 367}).
Hogg also includes \emph{nidi} / \emph{nǣdl} `needle' in the same
broader cluster history (\citeproc{ref-HoggPhonology1992}{Hogg 1992b}).

The target is therefore an attested citation form. No oblique-cell
substitution is involved in this entry.

\subsection{Development to Old
English}\label{development-to-old-english-104}

Ringe and Taylor give the historical line \Recon{nēþlō} `needle'
\textasciitilde{} \Recon{nēdlō} `needle' \textgreater{} OE \emph{nédl}
`needle' (\citeproc{ref-RingeTaylor2014}{Ringe and Taylor 2014, 329}).
Campbell likewise lists \emph{nédl} `needle' among the expected unbroken
forms after \emph{t} and \emph{d} (\citeproc{ref-Campbell1959}{Campbell
1959, § 367}). The trace expresses the same pathway with the voiced
alternant followed here for the Old English comparison.

The essential choice lies in which Proto-Germanic alternant is taken as
the starting point. Once the voiced form is chosen, the rest of the
pathway is regular.

\subsection{Alternant comparison}\label{alternant-comparison}

The comparison below sets the relevant forms side by side. It separates
the broader citation headword from the voiced alternant used for Old
English.

{\def\LTcaptype{none} % do not increment counter
\begin{longtable}[]{@{}
  >{\raggedright\arraybackslash}p{(\linewidth - 8\tabcolsep) * \real{0.2000}}
  >{\raggedright\arraybackslash}p{(\linewidth - 8\tabcolsep) * \real{0.2000}}
  >{\raggedright\arraybackslash}p{(\linewidth - 8\tabcolsep) * \real{0.2000}}
  >{\raggedright\arraybackslash}p{(\linewidth - 8\tabcolsep) * \real{0.2000}}
  >{\raggedright\arraybackslash}p{(\linewidth - 8\tabcolsep) * \real{0.2000}}@{}}
\toprule\noalign{}
\begin{minipage}[b]{\linewidth}\raggedright
Formation / stage
\end{minipage} & \begin{minipage}[b]{\linewidth}\raggedright
Candidate input
\end{minipage} & \begin{minipage}[b]{\linewidth}\raggedright
Expected or documented OE outcome
\end{minipage} & \begin{minipage}[b]{\linewidth}\raggedright
OE comparison form
\end{minipage} & \begin{minipage}[b]{\linewidth}\raggedright
Result
\end{minipage} \\
\midrule\noalign{}
\endhead
\bottomrule\noalign{}
\endlastfoot
comparative voiceless headword &
{\emph{*nḗθlō}}\index[iv]{03pgmc@\ivlangheader{Proto-Germanic}{}!nethlo@\iventry{*nḗθlō}{}}
& broader word-family label rather than the OE-facing alternant &
\emph{*nēþlō} line & useful citation form, but not the derivational
input for the Old English comparison \\
voiced Verner alternant followed here &
{\emph{*nḗðlō}}\index[iv]{03pgmc@\ivlangheader{Proto-Germanic}{}!nedlo@\iventry{*nḗðlō}{}}
& regular output: {\emph{nǣdl}} &
{\emph{nǣdl}}\index[iv]{02oe@\ivlangheader{Old English}{West Saxon normalization unmarked}!naedl@\iventry{nǣdl}{}}
& exact match between the form followed here and the attested OE noun \\
later hardening stage & *nḗdlō & intermediate pre-OE stage in the same
derivation & \emph{nǣdl} & genuine stage in the pathway, but not the
Proto-Germanic form followed here \\
\end{longtable}
}

\section{\texorpdfstring{nose --- OE
\emph{nosu}}{nose --- OE nosu}}\label{nose-oe-nosu}

\index[iv]{02oe@\ivlangheader{Old English}{West Saxon normalization unmarked}!nasu@\iventry{nasu}{}}
\index[iv]{02oe@\ivlangheader{Old English}{West Saxon normalization unmarked}!nosu@\iventry{nosu}{}}
\index[iv]{03pgmc@\ivlangheader{Proto-Germanic}{}!naso@\iventry{*nasō}{}}
\index[iv]{03pgmc@\ivlangheader{Proto-Germanic}{}!nuso@\iventry{*núsō}{}}

Derivation: citation reconstruction \emph{*nasō}; form followed here
\emph{*núsō} \textgreater{} \emph{nosu} (early analogy).

\subsection{Derivation trace}\label{derivation-trace-123}

Proto input: \emph{*núsō}

\begingroup
\setlength{\fboxsep}{6pt}

\noindent\fbox{%
\begin{minipage}{0.97\linewidth}
\small
\begin{tabular}{@{}>{\raggedright\arraybackslash}p{0.611\linewidth}@{\hspace{0.030\linewidth}}>{\raggedright\arraybackslash}p{0.359\linewidth}@{}}
\begin{minipage}[t]{\linewidth}
\centering\textbf{Earlier Germanic changes}\par
\vspace{0.35em}
\raggedright
\centering\textbf{Northwest and West Germanic}\par
\raggedright
\vspace{0.2em}
\begin{tabularx}{\linewidth}{@{}>{\raggedright\arraybackslash}X@{\hspace{0.55em}}>{\raggedright\arraybackslash}p{0.14\linewidth}@{\hspace{0.25em}}}
\mbox{PNWGmc U Lowering} & \emph{*nósō} \\
\mbox{PNWGmc Final Long O Raising} & \emph{*nósu} \\
\end{tabularx}
\vspace{0.6em}
\centering\textbf{Early Anglo-Frisian}\par
\raggedright
\vspace{0.2em}
\raggedright [no change]\par
\end{minipage}
&
\begin{minipage}[t]{\linewidth}
\centering\textbf{Old English changes}\par
\vspace{0.35em}
\raggedright
\raggedright [no change]\par
\end{minipage}
\\
\end{tabular}
\end{minipage}%
} \endgroup

Old English form: \emph{nosu}

\subsection{Reconstruction and comparative
evidence}\label{reconstruction-and-comparative-evidence-106}

Kroonen reconstructs a Germanic ablaut pair \emph{*nasō-}
\textasciitilde{} \emph{*nusō-} and adds that the root \emph{*nus-} is
likely to have arisen as a secondary zero grade after a remodeling of
the older paradigm (\citeproc{ref-Kroonen2013}{Kroonen 2013, 383}).
Campbell is more specific for Old English, citing \emph{nosu} `nose'
\textless{} \Recon{nusō} `nose' (\citeproc{ref-Campbell1959}{Campbell
1959, 43}).

The citation reconstruction \Recon{nasō} `nose' is therefore best
treated as the full-grade comparative headword, while the derivational
input \Recon{núsō} `nose' represents the remodeled zero-grade line
continued by the Old English form discussed here. Orel's \Recon{nasō}
`nose' alongside OE \emph{nasu} `nose' preserves the competing
full-grade notation and shows that the two lines should not be collapsed
without comment (\citeproc{ref-Orel2003}{Orel 2003, 320}).

\subsection{Old English evidence}\label{old-english-evidence-107}

\emph{Nosu} `nose' is an attested Old English noun. Ringe and Taylor
list it among the few surviving early Old English feminine u-stems
(\citeproc{ref-RingeTaylor2014}{Ringe and Taylor 2014, 385}). Clark Hall
likewise gives \emph{nosu f.} `nose', with genitive-dative singular
\emph{nosa}, and cross-refers
{\emph{nasu}}\index[iv]{02oe@\ivlangheader{Old English}{West Saxon normalization unmarked}!nasu@\iventry{nasu}{}}
`nose' to \emph{nosu} `nose' (\citeproc{ref-ClarkHall1960}{Clark Hall
1960, 211, 217}).

The selected OE target is therefore attested \emph{nosu} `nose', not a
reconstructed placeholder. The lexicographical record also gives
{\emph{nasu}}\index[iv]{02oe@\ivlangheader{Old English}{West Saxon normalization unmarked}!nasu@\iventry{nasu}{}}
`nose' for the full-grade side of the tradition.

\subsection{Development to Old
English}\label{development-to-old-english-105}

From \Recon{núsō} `nose', the regular path is the one documented by the
current trace: \Recon{núsō} `nose' \textgreater{} \Recon{nósō} `nose'
\textgreater{} \Recon{nósu} `nose' \textgreater{} \emph{nosu} `nose'.
The early special step lies in the choice of the zero-grade input, not
in any late Old English repair.

With that input chosen, the OE development is straightforward. The
full-grade line behind \Recon{nasō} `nose' instead points toward
{\emph{nasu}}\index[iv]{02oe@\ivlangheader{Old English}{West Saxon normalization unmarked}!nasu@\iventry{nasu}{}}
`nose', not to the form treated here.

The low root of \emph{nosu} `nose' supports the economy of lowering
before final oral raising, but only conditional on that selected
remodeled stem. Reversing the changes yields a high-root nominative,
whose replacement by a lowered oblique root would be an additional,
individually unproved levelling. The word is therefore not a pristine
inherited nominative proving the chronology. The comparative argument
and its limits are developed in
\hyperref[rule-PNWGmcFinalLongORaising]{SC019 PNWGmcFinalLongORaising}
(\citeproc{ref-Fulk2018}{Fulk 2018, 56};
\citeproc{ref-Kroonen2013}{Kroonen 2013, 383}).

\subsection{Stem comparison}\label{stem-comparison-4}

The comparison below sets the relevant forms side by side. It separates
the full-grade comparative line from the remodeled zero-grade input that
yields the Old English form.

{\def\LTcaptype{none} % do not increment counter
\begin{longtable}[]{@{}
  >{\raggedright\arraybackslash}p{(\linewidth - 8\tabcolsep) * \real{0.2000}}
  >{\raggedright\arraybackslash}p{(\linewidth - 8\tabcolsep) * \real{0.2000}}
  >{\raggedright\arraybackslash}p{(\linewidth - 8\tabcolsep) * \real{0.2000}}
  >{\raggedright\arraybackslash}p{(\linewidth - 8\tabcolsep) * \real{0.2000}}
  >{\raggedright\arraybackslash}p{(\linewidth - 8\tabcolsep) * \real{0.2000}}@{}}
\toprule\noalign{}
\begin{minipage}[b]{\linewidth}\raggedright
Formation / label
\end{minipage} & \begin{minipage}[b]{\linewidth}\raggedright
Candidate input
\end{minipage} & \begin{minipage}[b]{\linewidth}\raggedright
Expected or documented OE outcome
\end{minipage} & \begin{minipage}[b]{\linewidth}\raggedright
OE comparison form
\end{minipage} & \begin{minipage}[b]{\linewidth}\raggedright
Result
\end{minipage} \\
\midrule\noalign{}
\endhead
\bottomrule\noalign{}
\endlastfoot
full-grade comparative line & *nasō & expected full-grade continuation
\emph{nasu} & \emph{nasu} & useful comparative background, but not the
Old English-facing input \\
remodeled zero-grade line & *núsō & regular output: \emph{nosu} &
\emph{nosu} & exact match between derivational input and attested OE
noun \\
\end{longtable}
}

\section{\texorpdfstring{sap --- OE
\emph{sæp}}{sap --- OE sæp}}\label{sap-oe-suxe6p}

\index[iv]{02oe@\ivlangheader{Old English}{West Saxon normalization unmarked}!saep@\iventry{sæp}{}}
\index[iv]{03pgmc@\ivlangheader{Proto-Germanic}{}!sapa@\iventry{*sápą}{}}
\index[iv]{03pgmc@\ivlangheader{Proto-Germanic}{}!sapon@\iventry{*sapōn}{}}

Derivation: citation reconstruction \emph{*sapōn}; form followed here
\emph{*sápą} \textgreater{} \emph{sæp} (early analogy).

\subsection{Derivation trace}\label{derivation-trace-124}

Proto input: \emph{*sápą}

\begingroup
\setlength{\fboxsep}{6pt}

\noindent\fbox{%
\begin{minipage}{0.97\linewidth}
\small
\begin{tabular}{@{}>{\raggedright\arraybackslash}p{0.359\linewidth}@{\hspace{0.030\linewidth}}>{\raggedright\arraybackslash}p{0.611\linewidth}@{}}
\begin{minipage}[t]{\linewidth}
\centering\textbf{Earlier Germanic changes}\par
\vspace{0.35em}
\raggedright
\centering\textbf{Northwest and West Germanic}\par
\raggedright
\vspace{0.2em}
\raggedright [no change]\par
\vspace{0.6em}
\centering\textbf{Early Anglo-Frisian}\par
\raggedright
\vspace{0.2em}
\raggedright [no change]\par
\end{minipage}
&
\begin{minipage}[t]{\linewidth}
\centering\textbf{Old English changes}\par
\vspace{0.35em}
\raggedright
\begin{tabularx}{\linewidth}{@{}>{\raggedright\arraybackslash}X@{\hspace{0.55em}}>{\raggedright\arraybackslash}p{0.14\linewidth}@{\hspace{0.25em}}}
\mbox{EAF Brightening} & \emph{*sæpą} \\
OE Heavy Syllable Nasal Apocope & \emph{*sæp} \\
\end{tabularx}
\end{minipage}
\\
\end{tabular}
\end{minipage}%
} \endgroup

Old English form: \emph{sæp}

\subsection{Reconstruction and comparative
evidence}\label{reconstruction-and-comparative-evidence-107}

The comparative sources do not give one uniform inherited stem. Kroonen
preserves the word family as \Recon{saf} `sap' / \emph{*ppan-}, with Old
English \emph{sæp} `sap' (\citeproc{ref-Kroonen2013}{Kroonen 2013,
420}). Orel preserves the comparative notation \Recon{sapōn} `sap' and
\Recon{sapan} `sap' (\citeproc{ref-Orel2003}{Orel 2003, 319}).

The derivational input
\Recon{sápą}\index[iv]{03pgmc@\ivlangheader{Proto-Germanic}{}!sapa@\iventry{*sápą}{}}
`sap' therefore does not replace those comparative labels. It identifies
the OE-facing stem shape that yields the attested noun treated here.

\subsection{Old English evidence}\label{old-english-evidence-108}

Clark Hall records
{\emph{sæp}}\index[iv]{02oe@\ivlangheader{Old English}{West Saxon normalization unmarked}!saep@\iventry{sæp}{}}
`sap' (e) n. (\citeproc{ref-ClarkHall1960}{Clark Hall 1960, 247}). The
target is therefore an attested neuter Old English noun. Orel's plain
{\emph{sap}}\index[iv]{02oe@\ivlangheader{Old English}{West Saxon normalization unmarked}!sap@\iventry{sap}{}}
`sap, juice' notation belongs to comparative normalization, not to the
spelling adopted here for the Old English form
(\citeproc{ref-Orel2003}{Orel 2003, 319}).

\subsection{Development to Old
English}\label{development-to-old-english-106}

From
\Recon{sápą}\index[iv]{03pgmc@\ivlangheader{Proto-Germanic}{}!sapa@\iventry{*sápą}{}}
`sap', Anglo-Frisian brightening yields \emph{sæ} `sap', and
heavy-syllable nasal apocope then produces {\emph{sæp}} `sap'. That is
the regular path documented by the current trace.

The competing comparative lines do not give the same result. The
inherited n-stem notation
\Recon{sapōn}\index[iv]{03pgmc@\ivlangheader{Proto-Germanic}{}!sapon@\iventry{*sapōn}{}}
`sap' yields
{\emph{sape}}\index[iv]{02oe@\ivlangheader{Old English}{West Saxon normalization unmarked}!sape@\iventry{sape}{}}
`sap', while an i-stem continuation from the
{\emph{*sapi-}}\index[iv]{03pgmc@\ivlangheader{Proto-Germanic}{}!sapi@\iventry{*sapi-}{}}
line leads to
{\emph{sep}}\index[iv]{02oe@\ivlangheader{Old English}{West Saxon normalization unmarked}!sep@\iventry{sep}{}}
`sap' /
{\emph{sepe}}\index[iv]{02oe@\ivlangheader{Old English}{West Saxon normalization unmarked}!sepe@\iventry{sepe}{}}
`sap' rather than to
{\emph{sæp}}\index[iv]{02oe@\ivlangheader{Old English}{West Saxon normalization unmarked}!saep@\iventry{sæp}{}}
`sap'. The special step in this entry is therefore the early stem
choice, not a late OE paradigm-cell selection.

\subsection{Stem comparison}\label{stem-comparison-5}

The comparison below sets the relevant forms side by side. It separates
the competing comparative stem lines from the Old English-facing input.

{\def\LTcaptype{none} % do not increment counter
\begin{longtable}[]{@{}
  >{\raggedright\arraybackslash}p{(\linewidth - 8\tabcolsep) * \real{0.2000}}
  >{\raggedright\arraybackslash}p{(\linewidth - 8\tabcolsep) * \real{0.2000}}
  >{\raggedright\arraybackslash}p{(\linewidth - 8\tabcolsep) * \real{0.2000}}
  >{\raggedright\arraybackslash}p{(\linewidth - 8\tabcolsep) * \real{0.2000}}
  >{\raggedright\arraybackslash}p{(\linewidth - 8\tabcolsep) * \real{0.2000}}@{}}
\toprule\noalign{}
\begin{minipage}[b]{\linewidth}\raggedright
Formation / label
\end{minipage} & \begin{minipage}[b]{\linewidth}\raggedright
Candidate input
\end{minipage} & \begin{minipage}[b]{\linewidth}\raggedright
Expected or documented OE outcome
\end{minipage} & \begin{minipage}[b]{\linewidth}\raggedright
OE comparison form
\end{minipage} & \begin{minipage}[b]{\linewidth}\raggedright
Result
\end{minipage} \\
\midrule\noalign{}
\endhead
\bottomrule\noalign{}
\endlastfoot
comparative n-stem line &
{\emph{*sapōn}}\index[iv]{03pgmc@\ivlangheader{Proto-Germanic}{}!sapon@\iventry{*sapōn}{}}
& local comparator output:
{\emph{sape}}\index[iv]{02oe@\ivlangheader{Old English}{West Saxon normalization unmarked}!sape@\iventry{sape}{}}
`sap' &
{\emph{sape}}\index[iv]{02oe@\ivlangheader{Old English}{West Saxon normalization unmarked}!sape@\iventry{sape}{}}
`sap' & useful comparative background, but not the source of attested
\emph{sæp} \\
inferred i-stem comparator from \emph{*sapi-} &
{\emph{*sapiz}}\index[iv]{03pgmc@\ivlangheader{Proto-Germanic}{}!sapiz@\iventry{*sapiz}{}}
& local comparator output:
{\emph{sepe}}\index[iv]{02oe@\ivlangheader{Old English}{West Saxon normalization unmarked}!sepe@\iventry{sepe}{}}
`sap' &
{\emph{sepe}}\index[iv]{02oe@\ivlangheader{Old English}{West Saxon normalization unmarked}!sepe@\iventry{sepe}{}}
`sap' & confirms that an i-triggering stem does not reach the target \\
selected a-stem input &
{\emph{*sápą}}\index[iv]{03pgmc@\ivlangheader{Proto-Germanic}{}!sapa@\iventry{*sápą}{}}
& regular output: {\emph{sæp}} `sap' &
{\emph{sæp}}\index[iv]{02oe@\ivlangheader{Old English}{West Saxon normalization unmarked}!saep@\iventry{sæp}{}}
`sap' & exact match between derivational input and attested OE noun \\
\end{longtable}
}

\section{\texorpdfstring{sea --- OE
\emph{sǣ}}{sea --- OE sǣ}}\label{sea-oe-sux1e3}

\index[iv]{02oe@\ivlangheader{Old English}{West Saxon normalization unmarked}!sae@\iventry{sǣ}{}}
\index[iv]{03pgmc@\ivlangheader{Proto-Germanic}{}!sai@\iventry{*sái}{}}
\index[iv]{03pgmc@\ivlangheader{Proto-Germanic}{}!saiwiz@\iventry{*sáiwiz}{}}

Derivation: citation reconstruction \emph{*sái}; form followed here
\emph{*sáiwiz} \textgreater{} \emph{sǣ} (early analogy).

\subsection{Derivation trace}\label{derivation-trace-125}

Proto input: \emph{*sáiwiz}

\begingroup
\setlength{\fboxsep}{6pt}

\noindent\fbox{%
\begin{minipage}{0.97\linewidth}
\small
\begin{tabular}{@{}>{\raggedright\arraybackslash}p{0.485\linewidth}@{\hspace{0.030\linewidth}}>{\raggedright\arraybackslash}p{0.485\linewidth}@{}}
\begin{minipage}[t]{\linewidth}
\centering\textbf{Earlier Germanic changes}\par
\vspace{0.35em}
\raggedright
\centering\textbf{Northwest and West Germanic}\par
\raggedright
\vspace{0.2em}
\begin{tabularx}{\linewidth}{@{}>{\raggedright\arraybackslash}X@{\hspace{0.55em}}>{\raggedright\arraybackslash}p{0.14\linewidth}@{\hspace{0.25em}}}
\mbox{EAF Final Z Deletion} & \emph{*sáiwi} \\
\end{tabularx}
\vspace{0.6em}
\centering\textbf{Early Anglo-Frisian}\par
\raggedright
\vspace{0.2em}
\begin{tabularx}{\linewidth}{@{}>{\raggedright\arraybackslash}X@{\hspace{0.55em}}>{\raggedright\arraybackslash}p{0.14\linewidth}@{\hspace{0.25em}}}
\mbox{EAF Ai Monophthongization} & \emph{*sāwi} \\
\end{tabularx}
\end{minipage}
&
\begin{minipage}[t]{\linewidth}
\centering\textbf{Old English changes}\par
\vspace{0.35em}
\raggedright
\begin{tabularx}{\linewidth}{@{}>{\raggedright\arraybackslash}X@{\hspace{0.55em}}>{\raggedright\arraybackslash}p{0.14\linewidth}@{\hspace{0.25em}}}
\mbox{OE I Umlaut} & \emph{*sǣwi} \\
\mbox{OE W Loss Before I} & \emph{*sǣi} \\
\mbox{OE High Vowel Apocope} & \emph{*sǣ} \\
\end{tabularx}
\end{minipage}
\\
\end{tabular}
\end{minipage}%
} \endgroup

Old English form: \emph{sǣ}

\subsection{Reconstruction and comparative
evidence}\label{reconstruction-and-comparative-evidence-108}

Kroonen gives the noun in stem notation as \emph{*saiwi-}, an i-stem
whose English reflex is cited as OE
{\emph{sæ}}\index[iv]{02oe@\ivlangheader{Old English}{West Saxon normalization unmarked}!sae@\iventry{sæ}{}}
`sea' (\citeproc{ref-Kroonen2013}{Kroonen 2013, 423}). Ringe and Taylor
write the fuller form
\Recon{saiwiz}\index[iv]{03pgmc@\ivlangheader{Proto-Germanic}{}!saiwiz@\iventry{*saiwiz}{}}
`sea' and derive it through \Recon{sawi} `sea' \textgreater{}
\Recon{sei} `sea' \textgreater{} OE \emph{sǣ} `sea'
(\citeproc{ref-RingeTaylor2014}{Ringe and Taylor 2014, § 6.7.1}).
The comparative headword is therefore shorter than the form required for
the English history:
\Recon{sái}\index[iv]{03pgmc@\ivlangheader{Proto-Germanic}{}!sai@\iventry{*sái}{}}
`sea' names the lexeme, but
\Recon{sáiwiz}\index[iv]{03pgmc@\ivlangheader{Proto-Germanic}{}!saiwiz@\iventry{*sáiwiz}{}}
`sea' preserves the medial \emph{*w} and the final high vowel that
control the later development.

\subsection{Old English evidence}\label{old-english-evidence-109}

The Old English noun is the ordinary word for `sea'. Kroonen cites it as
{\emph{sæ}}\index[iv]{02oe@\ivlangheader{Old English}{West Saxon normalization unmarked}!sae@\iventry{sæ}{}}
`sea'; the normalized form here is
{\emph{sǣ}}\index[iv]{02oe@\ivlangheader{Old English}{West Saxon normalization unmarked}!sae@\iventry{sǣ}{}}
`sea' (\citeproc{ref-Kroonen2013}{Kroonen 2013, 423}). Campbell likewise
treats \emph{sea} as continuing the same \Recon{saiui} `sea'
\textgreater{} \Recon{sǣi} `sea' history, with loss of \emph{u}/\emph{w}
before \emph{i} (\citeproc{ref-Campbell1959}{Campbell 1959, § 406}).

\subsection{Development to Old
English}\label{development-to-old-english-107}

Once the fuller i-stem input is chosen, the development is regular.
After Proto-West-Germanic monophthongization \Recon{sáiwiz} `sea'
\textgreater{} \Recon{sāwiz} `sea' and final \emph{*-z} loss
\Recon{sāwiz} `sea' \textgreater{} \Recon{sāwi} `sea', the non-initial
\emph{*w} disappears before unstressed \emph{*i}, and the following high
vowel fronts the root vowel before final apocope. The documented chain
is \Recon{sáiwiz} `sea' \textgreater{} \Recon{sāwiz} `sea'
\textgreater{} \Recon{sāwi} `sea' \textgreater{} \Recon{sāi} `sea'
\textgreater{} \Recon{sǣi} `sea' \textgreater{} \emph{sǣ} `sea'
(\citeproc{ref-RingeTaylor2014}{Ringe and Taylor 2014, § 6.7.1};
\citeproc{ref-Campbell1959}{Campbell 1959, § 406}).

\subsection{Stem and stage comparison}\label{stem-and-stage-comparison}

The comparison below sets the relevant forms side by side. It separates
the abbreviated comparative headword from the fuller i-stem input that
yields the Old English form.

{\def\LTcaptype{none} % do not increment counter
\begin{longtable}[]{@{}
  >{\raggedright\arraybackslash}p{(\linewidth - 8\tabcolsep) * \real{0.2000}}
  >{\raggedright\arraybackslash}p{(\linewidth - 8\tabcolsep) * \real{0.2000}}
  >{\raggedright\arraybackslash}p{(\linewidth - 8\tabcolsep) * \real{0.2000}}
  >{\raggedright\arraybackslash}p{(\linewidth - 8\tabcolsep) * \real{0.2000}}
  >{\raggedright\arraybackslash}p{(\linewidth - 8\tabcolsep) * \real{0.2000}}@{}}
\toprule\noalign{}
\begin{minipage}[b]{\linewidth}\raggedright
Formation / label
\end{minipage} & \begin{minipage}[b]{\linewidth}\raggedright
Candidate input
\end{minipage} & \begin{minipage}[b]{\linewidth}\raggedright
OE output or comparison
\end{minipage} & \begin{minipage}[b]{\linewidth}\raggedright
OE comparison form
\end{minipage} & \begin{minipage}[b]{\linewidth}\raggedright
Result
\end{minipage} \\
\midrule\noalign{}
\endhead
\bottomrule\noalign{}
\endlastfoot
abbreviated comparative headword &
{\emph{*sái}}\index[iv]{03pgmc@\ivlangheader{Proto-Germanic}{}!sai@\iventry{*sái}{}}
& too short to preserve the \emph{*w} \ldots{} \emph{*i} environment
needed for the documented chronology &
{\emph{sǣ}}\index[iv]{02oe@\ivlangheader{Old English}{West Saxon normalization unmarked}!sae@\iventry{sǣ}{}}
& useful comparative label, but not the Old English-facing input \\
selected i-stem input &
{\emph{*sáiwiz}}\index[iv]{03pgmc@\ivlangheader{Proto-Germanic}{}!saiwiz@\iventry{*sáiwiz}{}}
& documented regular output: {\emph{sǣ}} &
{\emph{sǣ}}\index[iv]{02oe@\ivlangheader{Old English}{West Saxon normalization unmarked}!sae@\iventry{sǣ}{}}
& exact match between derivational input and Old English target \\
\end{longtable}
}

\section{\texorpdfstring{sieve --- OE
\emph{sife}}{sieve --- OE sife}}\label{sieve-oe-sife}

\index[iv]{02oe@\ivlangheader{Old English}{West Saxon normalization unmarked}!sife@\iventry{sife}{}}
\index[iv]{02oe@\ivlangheader{Old English}{West Saxon normalization unmarked}!spere@\iventry{spere}{}}
\index[iv]{03pgmc@\ivlangheader{Proto-Germanic}{}!sibaz@\iventry{*síbaz}{}}
\index[iv]{03pgmc@\ivlangheader{Proto-Germanic}{}!sibi@\iventry{*síbi}{}}
\index[iv]{03pgmc@\ivlangheader{Proto-Germanic}{}!sibiz@\iventry{*sib-iz}{}}

Derivation: citation reconstruction \emph{*síbaz}; form followed here
\emph{*síbi} \textgreater{} \emph{sife} (early analogy).

\subsection{Derivation trace}\label{derivation-trace-126}

Proto input: \emph{*síbi}

\begingroup
\setlength{\fboxsep}{6pt}

\noindent\fbox{%
\begin{minipage}{0.97\linewidth}
\small
\begin{tabular}{@{}>{\raggedright\arraybackslash}p{0.359\linewidth}@{\hspace{0.030\linewidth}}>{\raggedright\arraybackslash}p{0.611\linewidth}@{}}
\begin{minipage}[t]{\linewidth}
\centering\textbf{Earlier Germanic changes}\par
\vspace{0.35em}
\raggedright
\centering\textbf{Northwest and West Germanic}\par
\raggedright
\vspace{0.2em}
\raggedright [no change]\par
\vspace{0.6em}
\centering\textbf{Early Anglo-Frisian}\par
\raggedright
\vspace{0.2em}
\raggedright [no change]\par
\end{minipage}
&
\begin{minipage}[t]{\linewidth}
\centering\textbf{Old English changes}\par
\vspace{0.35em}
\raggedright
\begin{tabularx}{\linewidth}{@{}>{\raggedright\arraybackslash}X@{\hspace{0.55em}}>{\raggedright\arraybackslash}p{0.14\linewidth}@{\hspace{0.25em}}}
\mbox{PGmc B Allophony} & \emph{*síβi} \\
OE Med Unstressed I Lowering1 & \emph{*síβe} \\
\end{tabularx}
\end{minipage}
\\
\end{tabular}
\end{minipage}%
} \endgroup

Old English form: \emph{sife}

\subsection{Reconstruction and comparative
evidence}\label{reconstruction-and-comparative-evidence-109}

Kluge-Seebold gives wg. \emph{*sibi-} n.~\ldots{} ae. \emph{sife}, and
Campbell groups
{\emph{sife}}\index[iv]{02oe@\ivlangheader{Old English}{West Saxon normalization unmarked}!sife@\iventry{sife}{}}
`sieve' with short neuter i-stems such as \emph{spere}
(\citeproc{ref-KlugeSeebold2011}{Kluge 2011, 847};
\citeproc{ref-Campbell1959}{Campbell 1959, § 609}). The older
morphological background is the s-stem \Recon{sib-iz} `sieve', but the
derivational input is the normalized i-stem form
\Recon{síbi}\index[iv]{03pgmc@\ivlangheader{Proto-Germanic}{}!sibi@\iventry{*síbi}{}}
`sieve'.

Kroonen's nearby \emph{*sebjō-} entry belongs to the separate kinship
lexeme that yields Old English
{\emph{sibb}}\index[iv]{02oe@\ivlangheader{Old English}{West Saxon normalization unmarked}!sibb@\iventry{sibb}{}}
`sieve', not to the sieve word. Orel's
\Recon{sibaz}\index[iv]{03pgmc@\ivlangheader{Proto-Germanic}{}!sibaz@\iventry{*sibaz}{}}
`sieve' \ldots{} OE
{\emph{sife}}\index[iv]{02oe@\ivlangheader{Old English}{West Saxon normalization unmarked}!sife@\iventry{sife}{}}
`sieve' preserves a broader handbook notation, but that a-stem shape
does not fit the Old English form treated here
(\citeproc{ref-Orel2003}{Orel 2003, 328}).

\subsection{Old English evidence}\label{old-english-evidence-110}

Clark Hall gives
{\emph{sibi}}\index[iv]{02oe@\ivlangheader{Old English}{West Saxon normalization unmarked}!sibi@\iventry{sibi}{}}
`sieve' (GL) \ldots{} =
{\emph{sife}}\index[iv]{02oe@\ivlangheader{Old English}{West Saxon normalization unmarked}!sife@\iventry{sife}{}}
`sieve' and also
{\emph{sife}}\index[iv]{02oe@\ivlangheader{Old English}{West Saxon normalization unmarked}!sife@\iventry{sife}{}}
`sieve' n.~`sieve' (\citeproc{ref-ClarkHall1960}{Clark Hall 1960, 263}).
Campbell likewise cites Corpus Glossary \emph{sibi} `sieve' and treats
\emph{sife} `sieve' as a short neuter i-stem
(\citeproc{ref-Campbell1959}{Campbell 1959, §§ 444, 609}). The
normalized Old English target is therefore \emph{sife} `sieve', while
\emph{sibi} `sieve' is an attested earlier spelling rather than a
separate lexeme.

\subsection{Development to Old
English}\label{development-to-old-english-108}

From
\Recon{síbi}\index[iv]{03pgmc@\ivlangheader{Proto-Germanic}{}!sibi@\iventry{*síbi}{}}
`sieve', the regular derivation gives \Recon{síβi} `sieve'
\textgreater{} \Recon{síβe} `sieve' \textgreater{} {\emph{sife}}
`sieve'. Medial \emph{b} is realized as a spirant and later written
\emph{f}, while the final unstressed \emph{i} lowers to \emph{e}. The
older s-stem background \Recon{sib-iz} `sieve' explains the morphology,
but the derivational input
\Recon{síbi}\index[iv]{03pgmc@\ivlangheader{Proto-Germanic}{}!sibi@\iventry{*síbi}{}}
is the immediate pre-Old-English form.

\subsection{Stem comparison}\label{stem-comparison-6}

The comparison below sets the relevant forms side by side. It
distinguishes the accepted i-stem line from its rejected competitors.

{\def\LTcaptype{none} % do not increment counter
\begin{longtable}[]{@{}
  >{\raggedright\arraybackslash}p{(\linewidth - 8\tabcolsep) * \real{0.2000}}
  >{\raggedright\arraybackslash}p{(\linewidth - 8\tabcolsep) * \real{0.2000}}
  >{\raggedright\arraybackslash}p{(\linewidth - 8\tabcolsep) * \real{0.2000}}
  >{\raggedright\arraybackslash}p{(\linewidth - 8\tabcolsep) * \real{0.2000}}
  >{\raggedright\arraybackslash}p{(\linewidth - 8\tabcolsep) * \real{0.2000}}@{}}
\toprule\noalign{}
\begin{minipage}[b]{\linewidth}\raggedright
Formation / label
\end{minipage} & \begin{minipage}[b]{\linewidth}\raggedright
Candidate input
\end{minipage} & \begin{minipage}[b]{\linewidth}\raggedright
Expected or documented OE outcome
\end{minipage} & \begin{minipage}[b]{\linewidth}\raggedright
OE comparison form
\end{minipage} & \begin{minipage}[b]{\linewidth}\raggedright
Result
\end{minipage} \\
\midrule\noalign{}
\endhead
\bottomrule\noalign{}
\endlastfoot
ja-stem kinship line & Kroonen \emph{*sebjō-} / comparator
{\emph{*sibja}}\index[iv]{03pgmc@\ivlangheader{Proto-Germanic}{}!sibja@\iventry{*sibja}{}}
& OE
{\emph{sibb}}\index[iv]{02oe@\ivlangheader{Old English}{West Saxon normalization unmarked}!sibb@\iventry{sibb}{}}
`sieve' &
{\emph{sibb}}\index[iv]{02oe@\ivlangheader{Old English}{West Saxon normalization unmarked}!sibb@\iventry{sibb}{}}
`sieve' & separate lexeme, not the target treated here \\
a-stem handbook line &
{\emph{*síbaz}}\index[iv]{03pgmc@\ivlangheader{Proto-Germanic}{}!sibaz@\iventry{*síbaz}{}}
& expected
{\emph{sif}}\index[iv]{02oe@\ivlangheader{Old English}{West Saxon normalization unmarked}!sif@\iventry{sif}{}}
&
{\emph{sif}}\index[iv]{02oe@\ivlangheader{Old English}{West Saxon normalization unmarked}!sif@\iventry{sif}{}}
& wrong ending for the attested noun \\
selected i-stem line from older \emph{*sib-iz} &
{\emph{*síbi}}\index[iv]{03pgmc@\ivlangheader{Proto-Germanic}{}!sibi@\iventry{*síbi}{}}
& documented regular output: {\emph{sife}} `sieve' &
{\emph{sife}}\index[iv]{02oe@\ivlangheader{Old English}{West Saxon normalization unmarked}!sife@\iventry{sife}{}}
`sieve'; early spelling
{\emph{sibi}}\index[iv]{02oe@\ivlangheader{Old English}{West Saxon normalization unmarked}!sibi@\iventry{sibi}{}}
`sieve' & exact match between derivational input and Old English
evidence \\
\end{longtable}
}

\section{\texorpdfstring{spare --- OE
\emph{sparian}}{spare --- OE sparian}}\label{spare-oe-sparian}

\index[iv]{02oe@\ivlangheader{Old English}{West Saxon normalization unmarked}!sparian@\iventry{sparian}{}}
\index[iv]{03pgmc@\ivlangheader{Proto-Germanic}{}!sparena@\iventry{*sparēną}{}}
\index[iv]{08preoe@\ivlangheader{Pre-Old English / prehistoric English}{}!sparojana~01@\iventry{*spárōjaną}{transp.}}

Derivation: citation reconstruction \emph{*sparēną}; form followed here
\emph{*spárōjaną} \textgreater{} \emph{sparian} (early analogy).

\subsection{Derivation trace}\label{derivation-trace-127}

Proto input: \emph{*spárōjaną}

\begingroup
\setlength{\fboxsep}{6pt}

\noindent\fbox{%
\begin{minipage}{0.97\linewidth}
\footnotesize
\begin{tabular}{@{}>{\raggedright\arraybackslash}p{0.339\linewidth}@{\hspace{0.030\linewidth}}>{\raggedright\arraybackslash}p{0.630\linewidth}@{}}
\begin{minipage}[t]{\linewidth}
\centering\textbf{Earlier Germanic changes}\par
\vspace{0.35em}
\raggedright
\centering\textbf{Northwest and West Germanic}\par
\raggedright
\vspace{0.2em}
\raggedright [no change]\par
\vspace{0.6em}
\centering\textbf{Early Anglo-Frisian}\par
\raggedright
\vspace{0.2em}
\raggedright [no change]\par
\end{minipage}
&
\begin{minipage}[t]{\linewidth}
\centering\textbf{Old English changes}\par
\vspace{0.35em}
\raggedright
\begin{tabularx}{\linewidth}{@{}>{\raggedright\arraybackslash}X@{\hspace{0.55em}}>{\raggedright\arraybackslash}p{0.24\linewidth}@{\hspace{0.25em}}}
\mbox{EAF Brightening} & \emph{*spærōjaną} \\
\mbox{OE A Restoration} & \emph{*sparōjaną} \\
OE Heavy Syllable Nasal Apocope & \emph{*sparōjan} \\
\mbox{OE Secondary Nasalization} & \emph{*sparōjąn} \\
\mbox{OE I Umlaut} & \emph{*sparējąn} \\
\mbox{OE Med Long Short} & \emph{*sparejąn} \\
\mbox{OE Weak Tail Reduction} & \emph{*sparejan} \\
OE Intervocalic J Vocalization & \emph{*spareian} \\
\mbox{OE Unstressed EI Contraction} & \emph{*sparian} \\
\end{tabularx}
\end{minipage}
\\
\end{tabular}
\end{minipage}%
} \endgroup

Old English form: \emph{sparian}

\subsection{Reconstruction and comparative
evidence}\label{reconstruction-and-comparative-evidence-110}

Kroonen keeps the inherited verb under class-III
{\emph{*sparēn-}}\index[iv]{03pgmc@\ivlangheader{Proto-Germanic}{}!sparen@\iventry{*sparēn-}{}}
(\citeproc{ref-Kroonen2013}{Kroonen 2013, 465}). Orel similarly
preserves
\Recon{sparēnan}\index[iv]{03pgmc@\ivlangheader{Proto-Germanic}{}!sparenan@\iventry{*sparēnan}{}}
`spare' (\citeproc{ref-Orel2003}{Orel 2003, 362}). Ringe and Taylor,
however, reconstruct
\Recon{sparai-}\index[iv]{08preoe@\ivlangheader{Pre-Old English / prehistoric English}{}!sparai@\iventry{*sparai-}{}}
\textasciitilde{}
\Recon{sparja-}\index[iv]{08preoe@\ivlangheader{Pre-Old English / prehistoric English}{}!sparja@\iventry{*sparja-}{}}
for the English branch and derive the citation verb from a class-II line
(\citeproc{ref-RingeTaylor2014}{Ringe and Taylor 2014, 162, 191}). The
derivational input
\Recon{spárōjaną}\index[iv]{08preoe@\ivlangheader{Pre-Old English / prehistoric English}{}!sparojana~01@\iventry{*spárōjaną}{transp.}}
`spare' therefore represents the refashioned class-II formation behind
Old English
{\emph{sparian}}\index[iv]{02oe@\ivlangheader{Old English}{West Saxon normalization unmarked}!sparian@\iventry{sparian}{}}
`spare', while the citation reconstruction
\Recon{sparēną}\index[iv]{03pgmc@\ivlangheader{Proto-Germanic}{}!sparena@\iventry{*sparēną}{}}
`spare' remains the inherited comparative headword.

\subsection{Old English evidence}\label{old-english-evidence-111}

Campbell says that
{\emph{sparian}}\index[iv]{02oe@\ivlangheader{Old English}{West Saxon normalization unmarked}!sparian@\iventry{sparian}{}}
`spare' does not show the ordinary class-III characteristics, but the
Ritual forms, normalized here as
{\emph{spæria}}\index[iv]{02oe@\ivlangheader{Old English}{West Saxon normalization unmarked}!spaeria@\iventry{spæria}{}}
`spare',
{\emph{spær}}\index[iv]{02oe@\ivlangheader{Old English}{West Saxon normalization unmarked}!spaer@\iventry{spær}{}}
`spare', and
{\emph{spærede}}\index[iv]{02oe@\ivlangheader{Old English}{West Saxon normalization unmarked}!spaerede@\iventry{spærede}{}}
`spare', together with Vespasian Psalter
{\emph{spearad}}\index[iv]{02oe@\ivlangheader{Old English}{West Saxon normalization unmarked}!spearad@\iventry{spearad}{}}
`spare', point to primitive Old English forms both with and without back
vowels (\citeproc{ref-Campbell1959}{Campbell 1959, § 764}). Brunner
likewise records Northumbrian
{\emph{spæria}}\index[iv]{02oe@\ivlangheader{Old English}{West Saxon normalization unmarked}!spaeria~06@\iventry{spæria}{North.}}
`spare',
{\emph{spærede}}\index[iv]{02oe@\ivlangheader{Old English}{West Saxon normalization unmarked}!spaerede~06@\iventry{spærede}{North.}}
`spare' beside common Old English
{\emph{sparian}}\index[iv]{02oe@\ivlangheader{Old English}{West Saxon normalization unmarked}!sparian@\iventry{sparian}{}}
`spare' and Vespasian Psalter
{\emph{spearad}}\index[iv]{02oe@\ivlangheader{Old English}{West Saxon normalization unmarked}!spearad@\iventry{spearad}{}}
`spare' (\citeproc{ref-SieversBrunner1965}{Brunner 1965, § 364} Anm.
11). The citation form treated here is
{\emph{sparian}}\index[iv]{02oe@\ivlangheader{Old English}{West Saxon normalization unmarked}!sparian@\iventry{sparian}{}}
`spare'; the Anglian forms are relics of the older formation, not
alternative headwords of equal status.

\subsection{Development to Old
English}\label{development-to-old-english-109}

Once the class-II formation
\Recon{spárōjaną}\index[iv]{08preoe@\ivlangheader{Pre-Old English / prehistoric English}{}!sparojana~01@\iventry{*spárōjaną}{transp.}}
`spare' is chosen, the remaining development is regular. The regular
derivation shows brightening, restoration of \emph{a} before the back
vocalism of the suffix, later i-mutation within the weak ending,
weak-tail reduction, and contraction to
{\emph{sparian}}\index[iv]{02oe@\ivlangheader{Old English}{West Saxon normalization unmarked}!sparian@\iventry{sparian}{}}
`spare'. By contrast, Brunner's rule against further apocope of final
\emph{-e} explains why Ritual
{\emph{spær}}\index[iv]{02oe@\ivlangheader{Old English}{West Saxon normalization unmarked}!spaer@\iventry{spær}{}}
`spare' cannot be the regular continuation of inherited
\Recon{spárē}\index[iv]{08preoe@\ivlangheader{Pre-Old English / prehistoric English}{}!spare@\iventry{*spárē}{}}
`spare' (\citeproc{ref-SieversBrunner1965}{Brunner 1965, § 150}).

\subsection{Formation comparison}\label{formation-comparison-4}

The comparison below sets the relevant forms side by side. It contrasts
the inherited class-III formation with the refashioned class-II one that
yields the citation verb.

{\def\LTcaptype{none} % do not increment counter
\begin{longtable}[]{@{}
  >{\raggedright\arraybackslash}p{(\linewidth - 8\tabcolsep) * \real{0.2000}}
  >{\raggedright\arraybackslash}p{(\linewidth - 8\tabcolsep) * \real{0.2000}}
  >{\raggedright\arraybackslash}p{(\linewidth - 8\tabcolsep) * \real{0.2000}}
  >{\raggedright\arraybackslash}p{(\linewidth - 8\tabcolsep) * \real{0.2000}}
  >{\raggedright\arraybackslash}p{(\linewidth - 8\tabcolsep) * \real{0.2000}}@{}}
\toprule\noalign{}
\begin{minipage}[b]{\linewidth}\raggedright
Formation / comparison
\end{minipage} & \begin{minipage}[b]{\linewidth}\raggedright
Candidate input
\end{minipage} & \begin{minipage}[b]{\linewidth}\raggedright
Expected or documented OE outcome
\end{minipage} & \begin{minipage}[b]{\linewidth}\raggedright
OE comparison form
\end{minipage} & \begin{minipage}[b]{\linewidth}\raggedright
Result
\end{minipage} \\
\midrule\noalign{}
\endhead
\bottomrule\noalign{}
\endlastfoot
inherited class-III infinitive &
{\emph{*spárēną}}\index[iv]{03pgmc@\ivlangheader{Proto-Germanic}{}!sparena@\iventry{*spárēną}{}}
& paradigm comparison / probe output:
{\emph{sparen}}\index[iv]{02oe@\ivlangheader{Old English}{West Saxon normalization unmarked}!sparen@\iventry{sparen}{}}
&
{\emph{sparian}}\index[iv]{02oe@\ivlangheader{Old English}{West Saxon normalization unmarked}!sparian@\iventry{sparian}{}}
`spare' & wrong class and wrong ending for the citation verb \\
inherited class-III imperative singular & \Recon{spárē} & paradigm
comparison / probe output:
{\emph{spære}}\index[iv]{02oe@\ivlangheader{Old English}{West Saxon normalization unmarked}!spaere@\iventry{spære}{}}
& Ritual
{\emph{spær}}\index[iv]{02oe@\ivlangheader{Old English}{West Saxon normalization unmarked}!spaer@\iventry{spær}{}}
`spare' & loss of final \emph{-e} is not regular, so the relic form
cannot control the entry \\
inherited class-III finite present & \Recon{spárēθi} & paradigm
comparison / probe output:
{\emph{spæreþ}}\index[iv]{02oe@\ivlangheader{Old English}{West Saxon normalization unmarked}!spaereth@\iventry{spæreþ}{}}
&
{\emph{spearad}}\index[iv]{02oe@\ivlangheader{Old English}{West Saxon normalization unmarked}!spearad@\iventry{spearad}{}}
`spare' & attested form is mixed, not a direct continuation of the
inherited cell \\
selected class-II formation &
{\emph{*spárōjaną}}\index[iv]{08preoe@\ivlangheader{Pre-Old English / prehistoric English}{}!sparojana~01@\iventry{*spárōjaną}{transp.}}
& documented regular output:
{\emph{sparian}}\index[iv]{02oe@\ivlangheader{Old English}{West Saxon normalization unmarked}!sparian@\iventry{sparian}{}}
`spare' &
{\emph{sparian}}\index[iv]{02oe@\ivlangheader{Old English}{West Saxon normalization unmarked}!sparian@\iventry{sparian}{}}
`spare' & exact match between derivational input and Old English
citation form \\
\end{longtable}
}

\section{\texorpdfstring{staff --- OE
\emph{stæf}}{staff --- OE stæf}}\label{staff-oe-stuxe6f}

\index[iv]{02oe@\ivlangheader{Old English}{West Saxon normalization unmarked}!staef@\iventry{stæf}{}}
\index[iv]{03pgmc@\ivlangheader{Proto-Germanic}{}!stabaz@\iventry{*stábaz}{}}
\index[iv]{03pgmc@\ivlangheader{Proto-Germanic}{}!stabiz@\iventry{*stábiz}{}}

Derivation: citation reconstruction \emph{*stábiz}; form followed here
\emph{*stábaz} \textgreater{} \emph{stæf} (early analogy).

\subsection{Derivation trace}\label{derivation-trace-128}

Proto input: \emph{*stábaz}

\begingroup
\setlength{\fboxsep}{6pt}

\noindent\fbox{%
\begin{minipage}{0.97\linewidth}
\small
\begin{tabular}{@{}>{\raggedright\arraybackslash}p{0.388\linewidth}@{\hspace{0.030\linewidth}}>{\raggedright\arraybackslash}p{0.582\linewidth}@{}}
\begin{minipage}[t]{\linewidth}
\centering\textbf{Earlier Germanic changes}\par
\vspace{0.35em}
\raggedright
\centering\textbf{Northwest and West Germanic}\par
\raggedright
\vspace{0.2em}
\begin{tabularx}{\linewidth}{@{}>{\raggedright\arraybackslash}X@{\hspace{0.55em}}>{\raggedright\arraybackslash}p{0.14\linewidth}@{\hspace{0.25em}}}
\mbox{EAF Final Z Deletion} & \emph{*stába} \\
\end{tabularx}
\vspace{0.6em}
\centering\textbf{Early Anglo-Frisian}\par
\raggedright
\vspace{0.2em}
\raggedright [no change]\par
\end{minipage}
&
\begin{minipage}[t]{\linewidth}
\centering\textbf{Old English changes}\par
\vspace{0.35em}
\raggedright
\begin{tabularx}{\linewidth}{@{}>{\raggedright\arraybackslash}X@{\hspace{0.55em}}>{\raggedright\arraybackslash}p{0.14\linewidth}@{\hspace{0.25em}}}
\mbox{PWGmc Final Bare A Loss} & \emph{*stáb} \\
\mbox{EAF Brightening} & \emph{*stæb} \\
\mbox{PGmc B Allophony} & \emph{*stæβ} \\
\end{tabularx}
\end{minipage}
\\
\end{tabular}
\end{minipage}%
} \endgroup

Old English form: \emph{stæf}

\subsection{Reconstruction and comparative
evidence}\label{reconstruction-and-comparative-evidence-111}

The comparative dictionaries do not give one uniform stem class. Kroonen
reconstructs an a-stem \emph{*staba-}
(\citeproc{ref-Kroonen2013}{Kroonen 2013, 471}). Orel writes
\Recon{stábiz}\index[iv]{03pgmc@\ivlangheader{Proto-Germanic}{}!stabiz@\iventry{*stábiz}{}}
`staff' \textasciitilde{}
\Recon{stábaz}\index[iv]{03pgmc@\ivlangheader{Proto-Germanic}{}!stabaz@\iventry{*stábaz}{}}
`staff' (\citeproc{ref-Orel2003}{Orel 2003, 368}). A direct i-stem input
in \Recon{-iz} `staff' would predict i-mutation in Old English, whereas
the attested noun keeps \emph{æ}.

\subsection{Old English evidence}\label{old-english-evidence-112}

The Old English noun itself is the ordinary citation form
{\emph{stæf}}\index[iv]{02oe@\ivlangheader{Old English}{West Saxon normalization unmarked}!staef@\iventry{stæf}{}}
`staff'. Luick lists
{\emph{stæf}}\index[iv]{02oe@\ivlangheader{Old English}{West Saxon normalization unmarked}!staef@\iventry{stæf}{}}
`staff' among closed monosyllables with \emph{æ}
(\citeproc{ref-Luick1914}{Luick 1914-\/-1940, 176}). Ringe and Taylor
pair singular
{\emph{stæf}}\index[iv]{02oe@\ivlangheader{Old English}{West Saxon normalization unmarked}!staef@\iventry{stæf}{}}
`staff' with plural
{\emph{stafas}}\index[iv]{02oe@\ivlangheader{Old English}{West Saxon normalization unmarked}!stafas@\iventry{stafas}{}}
`staffs' (\citeproc{ref-RingeTaylor2014}{Ringe and Taylor 2014, 193}).
The normalized form here is therefore \emph{stæf}; later English
\emph{staff} with \emph{a} belongs to a later stage of the word's
history.

\subsection{Development to Old
English}\label{development-to-old-english-110}

With the selected a-stem input, the development is regular. Final
\emph{*-z} disappears, bare final \emph{-a} is lost, Anglo-Frisian
brightening gives \emph{æ} in the closed monosyllable, and medial
\emph{b} surfaces as a fricative written \emph{f}. The documented chain
is \Recon{stábaz} `staff' \textgreater{} \Recon{stába} `staff'
\textgreater{} \Recon{stáb} `staff' \textgreater{} \emph{stæb} `staff'
\textgreater{} \emph{stæf} `staff'. A direct continuation of
\Recon{stábiz} `staff', by contrast, would produce i-mutated
{\emph{stefe}}\index[iv]{02oe@\ivlangheader{Old English}{West Saxon normalization unmarked}!stefe@\iventry{stefe}{}}
`stave, staff' rather than the attested singular.

\subsection{Formation comparison}\label{formation-comparison-5}

The comparison below sets the relevant forms side by side. It separates
the rejected i-stem line from the selected a-stem input.

{\def\LTcaptype{none} % do not increment counter
\begin{longtable}[]{@{}
  >{\raggedright\arraybackslash}p{(\linewidth - 8\tabcolsep) * \real{0.2000}}
  >{\raggedright\arraybackslash}p{(\linewidth - 8\tabcolsep) * \real{0.2000}}
  >{\raggedright\arraybackslash}p{(\linewidth - 8\tabcolsep) * \real{0.2000}}
  >{\raggedright\arraybackslash}p{(\linewidth - 8\tabcolsep) * \real{0.2000}}
  >{\raggedright\arraybackslash}p{(\linewidth - 8\tabcolsep) * \real{0.2000}}@{}}
\toprule\noalign{}
\begin{minipage}[b]{\linewidth}\raggedright
Formation / label
\end{minipage} & \begin{minipage}[b]{\linewidth}\raggedright
Candidate input
\end{minipage} & \begin{minipage}[b]{\linewidth}\raggedright
OE output or comparison
\end{minipage} & \begin{minipage}[b]{\linewidth}\raggedright
OE comparison form
\end{minipage} & \begin{minipage}[b]{\linewidth}\raggedright
Result
\end{minipage} \\
\midrule\noalign{}
\endhead
\bottomrule\noalign{}
\endlastfoot
comparative i-stem line &
{\emph{*stábiz}}\index[iv]{03pgmc@\ivlangheader{Proto-Germanic}{}!stabiz@\iventry{*stábiz}{}}
& expected {\emph{stefe}} after i-mutation &
{\emph{stæf}}\index[iv]{02oe@\ivlangheader{Old English}{West Saxon normalization unmarked}!staef@\iventry{stæf}{}}
& wrong vowel for the attested singular \\
mixed comparative notation & Orel \emph{*stabiz} \textasciitilde{}
\emph{*stabaz}; Kluge \emph{*stabi-}/a- & source-level stem-class
uncertainty & \emph{stæf} & useful comparative background, but not a
single OE-facing input \\
selected a-stem input &
{\emph{*stábaz}}\index[iv]{03pgmc@\ivlangheader{Proto-Germanic}{}!stabaz@\iventry{*stábaz}{}}
& documented regular output: {\emph{stæf}} &
{\emph{stæf}}\index[iv]{02oe@\ivlangheader{Old English}{West Saxon normalization unmarked}!staef@\iventry{stæf}{}}
& exact match between derivational input and Old English target \\
\end{longtable}
}

\section{\texorpdfstring{stem --- OE
\emph{stefn}}{stem --- OE stefn}}\label{stem-oe-stefn}

\index[iv]{02oe@\ivlangheader{Old English}{West Saxon normalization unmarked}!stefn@\iventry{stefn}{}}
\index[iv]{03pgmc@\ivlangheader{Proto-Germanic}{}!stamnaz@\iventry{*stámnaz}{}}
\index[iv]{03pgmc@\ivlangheader{Proto-Germanic}{}!stamniz@\iventry{*stámniz}{}}

Derivation: citation reconstruction \emph{*stámnaz}; form followed here
\emph{*stámniz} \textgreater{} \emph{stefn} (early analogy).

\subsection{Derivation trace}\label{derivation-trace-129}

Proto input: \emph{*stámniz}

\begingroup
\setlength{\fboxsep}{6pt}

\noindent\fbox{%
\begin{minipage}{0.97\linewidth}
\small
\begin{tabular}{@{}>{\raggedright\arraybackslash}p{0.485\linewidth}@{\hspace{0.030\linewidth}}>{\raggedright\arraybackslash}p{0.485\linewidth}@{}}
\begin{minipage}[t]{\linewidth}
\centering\textbf{Earlier Germanic changes}\par
\vspace{0.35em}
\raggedright
\centering\textbf{Northwest and West Germanic}\par
\raggedright
\vspace{0.2em}
\begin{tabularx}{\linewidth}{@{}>{\raggedright\arraybackslash}X@{\hspace{0.55em}}>{\raggedright\arraybackslash}p{0.14\linewidth}@{\hspace{0.25em}}}
\mbox{EAF Final Z Deletion} & \emph{*stámni} \\
\mbox{PNWGmc Mn Dissimilation} & \emph{*stáβni} \\
\end{tabularx}
\vspace{0.6em}
\centering\textbf{Early Anglo-Frisian}\par
\raggedright
\vspace{0.2em}
\raggedright [no change]\par
\end{minipage}
&
\begin{minipage}[t]{\linewidth}
\centering\textbf{Old English changes}\par
\vspace{0.35em}
\raggedright
\begin{tabularx}{\linewidth}{@{}>{\raggedright\arraybackslash}X@{\hspace{0.55em}}>{\raggedright\arraybackslash}p{0.14\linewidth}@{\hspace{0.25em}}}
\mbox{EAF Brightening} & \emph{*stæβni} \\
\mbox{OE I Umlaut} & \emph{*steβni} \\
\mbox{OE High Vowel Apocope} & \emph{*steβn} \\
\end{tabularx}
\end{minipage}
\\
\end{tabular}
\end{minipage}%
} \endgroup

Old English form: \emph{stefn}

\subsection{Reconstruction and comparative
evidence}\label{reconstruction-and-comparative-evidence-112}

Orel reconstructs the Proto-Germanic source as
\Recon{stámnaz}\index[iv]{03pgmc@\ivlangheader{Proto-Germanic}{}!stamnaz@\iventry{*stámnaz}{}}
`stem, trunk' (also variant \Recon{stamniz} `stem, trunk'), citing Old
Norse
{\emph{stafn}}\index[iv]{09on@\ivlangheader{Old Norse}{}!stafn@\iventry{stafn}{}}
`stem of a ship', Old Frisian
{\emph{stevene}}\index[iv]{11ofris@\ivlangheader{Old Frisian}{}!stevene@\iventry{stevene}{}}
`stem of a ship' (fem.), and Old Saxon
{\emph{stamn}}\index[iv]{13os@\ivlangheader{Old Saxon}{}!stamn@\iventry{stamn}{}}
`stem' as its main continuants (\citeproc{ref-Orel2003}{Orel 2003,
371}).

This word is etymologically unrelated to the Old English homonym
\emph{stefn} `voice, sound' / \emph{stemn} `voice, sound', which
descends from a distinct Proto-Germanic \Recon{stebnō} `voice, sound' /
\Recon{stemnō} `voice, sound' etymon. The two OE lexemes are
distinguished in Clark Hall, Bosworth-Toller, Brunner, and Luick; they
collide on the surface forms \emph{stefn} and \emph{stemn} through
historically independent developments.

Kroonen relates the `stem' family to the \emph{stam(m)} `stem, trunk'
group attested in Old High German, Dutch, and German, without
reconstructing a separate `prow' proto-form
(\citeproc{ref-Kroonen2013}{Kroonen 2013, 479--80}).

\subsection{Old English evidence}\label{old-english-evidence-113}

Clark Hall records
{\emph{stefn}}\index[iv]{02oe@\ivlangheader{Old English}{West Saxon normalization unmarked}!stefn@\iventry{stefn}{}}
`stem, trunk, root, prow, foundation', the weak n-stem
{\emph{stefna}}\index[iv]{02oe@\ivlangheader{Old English}{West Saxon normalization unmarked}!stefna@\iventry{stefna}{}}
`prow or stern of a ship', and the related
{\emph{stofn}}\index[iv]{02oe@\ivlangheader{Old English}{West Saxon normalization unmarked}!stofn@\iventry{stofn}{}}
`trunk, stem, branch, shoot' (\citeproc{ref-ClarkHall1960}{Clark Hall
1960, 276, 341}).

Luick §211 explicitly distinguishes OE \emph{stefn} `stem, trunk' /
\emph{stemn} `stem, trunk' --- noting cognates with Old Saxon
{\emph{stamn}}\index[iv]{13os@\ivlangheader{Old Saxon}{}!stamn@\iventry{stamn}{}}
`stem' and Middle English
{\emph{stam}}\index[iv]{21me@\ivlangheader{Middle English}{}!stam@\iventry{stam}{}}
`stem' --- from the separate OE voice/sound word (\emph{stefn} `voice,
sound' / \emph{stemn} `voice, sound') treated in §75
(\citeproc{ref-Luick1914}{Luick 1914-\/-1940, § 211}). Brunner §205
lists \emph{stefn, stemn Stamm} among words showing the
\emph{fn}/\emph{mn} alternation
(\citeproc{ref-SieversBrunner1965}{Brunner 1965, § 205}).

The primary comparison form used here is
{\emph{stefn}}\index[iv]{02oe@\ivlangheader{Old English}{West Saxon normalization unmarked}!stefn@\iventry{stefn}{}}
`stem, trunk, root, prow', the form attested most directly in the
semantic range relevant to English \emph{stem}. The n-stem
{\emph{stefna}}\index[iv]{02oe@\ivlangheader{Old English}{West Saxon normalization unmarked}!stefna@\iventry{stefna}{}}
`prow/stern' is a closely related but more narrowly attested nautical
specialization.

\subsection{Development to Old
English}\label{development-to-old-english-111}

The derivation of OE
{\emph{stefn}}\index[iv]{02oe@\ivlangheader{Old English}{West Saxon normalization unmarked}!stefn@\iventry{stefn}{}}
`stem, trunk' from the i-stem input
\Recon{stámniz}\index[iv]{03pgmc@\ivlangheader{Proto-Germanic}{}!stamniz@\iventry{*stámniz}{}}
`stem, trunk' (Orel's attested i-stem variant; the citation
reconstruction remains
\Recon{stámnaz}\index[iv]{03pgmc@\ivlangheader{Proto-Germanic}{}!stamnaz@\iventry{*stámnaz}{}}
`stem, trunk') is now regular and modelled, with multiplicity 1. Live
trace under the corrected literal adjacent-\emph{mn} SC022:

{\def\LTcaptype{none} % do not increment counter
\begin{longtable}[]{@{}ll@{}}
\toprule\noalign{}
step & form \\
\midrule\noalign{}
\endhead
\bottomrule\noalign{}
\endlastfoot
proto input & \emph{*stámniz} \\
EAF Final Z Deletion & \emph{*stámni} \\
PNWGmc Mn Dissimilation (adjacent mn \textgreater{} βn) &
\emph{*stáβni} \\
EAF Brightening (á \textgreater{} æ) & \emph{*stæβni} \\
OE i-Umlaut (æ \textgreater{} e, from the i-stem ending) &
\emph{*steβni} \\
OE High Vowel Apocope & \emph{*steβn} \\
Old English Orthography (β \textgreater{} f) & \emph{stefn} \\
\end{longtable}
}

The consonantal change mn → \emph{βn}/\emph{fn} is attested
comparatively: Old Norse
{\emph{stafn}}\index[iv]{09on@\ivlangheader{Old Norse}{}!stafn@\iventry{stafn}{}}
`stem of a ship' and Old Saxon
{\emph{stamn}}\index[iv]{13os@\ivlangheader{Old Saxon}{}!stamn@\iventry{stamn}{}}
`stem' preserve the \emph{fn}/\emph{mn} variants expected from this
family. The handbooks treat the change as a pre-Old-English development
(Luick §211; Brunner §205), and the FST now encodes it as the historical
adjacent \emph{mn} \textgreater{} \emph{βn} (the rule
\emph{PNWGmcMnDissimilation} / SC022). The \emph{e} of \emph{stefn} is
regular from \emph{á} via brightening (\emph{á} \textgreater{} \emph{æ})
plus i-umlaut triggered by the i-stem ending (\emph{æ} \textgreater{}
\emph{e}), so both the vowel and the consonant follow by regular sound
change once the i-stem input is selected.

Note the contrast with \emph{heaven}. Here SC022
(\emph{PNWGmcMnDissimilation}) fires directly inside the selected
derivation, deriving the labial of \emph{stefn} from the adjacent
\emph{-mn-} cluster of \emph{*stámniz}. In \emph{heaven}, by contrast,
the same historical change \emph{mn} \textgreater{} \emph{βn}/\emph{fn}
operated only in the deeper cluster-bearing oblique prehistory; the
resulting labial was then generalized into the vowel-bearing
\emph{*hebun-} stem, so the selected CAPR path \emph{*xébun}
\textgreater{} \emph{heofon} begins \emph{after} that analogy and does
not itself contain SC022 (see
\texttt{dossier-heaven-paradigm-history-2026.md} §§13--15 and
\texttt{audits/heaven-sc022-implementation-2026.md}). The earlier
cross-syllable \emph{mV\ldots n} proxy --- which had formerly fabricated
a labial from an intervocalic \emph{m} --- has been retired in favour of
this literal adjacent \emph{mn} \textgreater{} \emph{βn}. Selecting the
i-stem input \emph{*stámniz} over the a-stem citation \emph{*stámnaz} is
a pre-OE input selection; once selected, the Old English development is
regular. Hence the classification \emph{early\_analogy} (analogy
separates the input from the citation reconstruction before the
specifically Old English changes apply), not \emph{known\_unmodelled}.

\subsection{Homonym note}\label{homonym-note}

Old English \emph{stefn}/\emph{stemn} is a surface homonym for two
etymologically unrelated words:

\begin{enumerate}
\def\labelenumi{\arabic{enumi}.}
\tightlist
\item
  \emph{stefn}/\emph{stemn} `voice, sound': from PGmc \emph{*stebnō}
  (Ringe-Taylor, Orel); via \emph{bn → fn} (b-allophony), giving
  \emph{stefn} `voice, sound'. This is a different row from 2216.
\item
  \emph{stefn}/\emph{stefna}/\emph{stofn}/\emph{stemn} `stem, trunk,
  prow': from PGmc
  \Recon{stámnaz}\index[iv]{03pgmc@\ivlangheader{Proto-Germanic}{}!stamnaz@\iventry{*stámnaz}{}}
  `stem/trunk' (Orel); via the pre-OE cluster change \emph{mn → fn}.
  This is row 2216.
\end{enumerate}

The form followed here, \emph{*stébnō}, used in earlier versions of this
entry was the wrong homonym's transponent and must not be used here.

\subsection{Source comparison}\label{source-comparison}

{\def\LTcaptype{none} % do not increment counter
\begin{longtable}[]{@{}
  >{\raggedright\arraybackslash}p{(\linewidth - 6\tabcolsep) * \real{0.2500}}
  >{\raggedright\arraybackslash}p{(\linewidth - 6\tabcolsep) * \real{0.2500}}
  >{\raggedright\arraybackslash}p{(\linewidth - 6\tabcolsep) * \real{0.2500}}
  >{\raggedright\arraybackslash}p{(\linewidth - 6\tabcolsep) * \real{0.2500}}@{}}
\toprule\noalign{}
\begin{minipage}[b]{\linewidth}\raggedright
Form or label
\end{minipage} & \begin{minipage}[b]{\linewidth}\raggedright
Status
\end{minipage} & \begin{minipage}[b]{\linewidth}\raggedright
OE relation
\end{minipage} & \begin{minipage}[b]{\linewidth}\raggedright
Result
\end{minipage} \\
\midrule\noalign{}
\endhead
\bottomrule\noalign{}
\endlastfoot
\Recon{stámnaz}\index[iv]{03pgmc@\ivlangheader{Proto-Germanic}{}!stamnaz@\iventry{*stámnaz}{}}
`stem, trunk' & Orel's PGmc a-stem citation for the stem/trunk/prow
family & lexeme-level citation & retained as PROTO (citation) \\
\Recon{stámniz}\index[iv]{03pgmc@\ivlangheader{Proto-Germanic}{}!stamniz@\iventry{*stámniz}{}}
`stem, trunk' & Orel's attested i-stem variant & \textbf{selected
derivational input} & regular OE output \emph{stefn} (mult 1) \\
{\emph{stefn}}\index[iv]{02oe@\ivlangheader{Old English}{West Saxon normalization unmarked}!stefn@\iventry{stefn}{}}
III `stem, trunk, root, prow' & primary OE target (strong masc.) &
stem/trunk sense per Clark Hall & target form \\
{\emph{stefna}}\index[iv]{02oe@\ivlangheader{Old English}{West Saxon normalization unmarked}!stefna@\iventry{stefna}{}}
`prow/stern' & OE n-stem specialization & nautical sense per Clark Hall
& comparison form \\
{\emph{stofn}}\index[iv]{02oe@\ivlangheader{Old English}{West Saxon normalization unmarked}!stofn@\iventry{stofn}{}}
`trunk' & OE \emph{o}-grade variant or earlier stage & trunk/stem sense
& comparison form \\
{\emph{stemn}}\index[iv]{02oe@\ivlangheader{Old English}{West Saxon normalization unmarked}!stemn@\iventry{stemn}{}}
`trunk' & late West Saxon \emph{fn → mn} doublet & secondary form &
comparison form \\
\emph{*stébnō} (voice word) & \textbf{wrong homonym} --- belongs to
voice/sound dossier & no relation to stem/trunk sense & \textbf{must not
be used as row 2216's derivational input} \\
\end{longtable}
}

\section{\texorpdfstring{swan --- OE
\emph{swanes}}{swan --- OE swanes}}\label{swan-oe-swanes}

\index[iv]{02oe@\ivlangheader{Old English}{West Saxon normalization unmarked}!swanes@\iventry{swanes}{}}
\index[iv]{03pgmc@\ivlangheader{Proto-Germanic}{}!swanas@\iventry{*swánas}{}}
\index[iv]{03pgmc@\ivlangheader{Proto-Germanic}{}!swanaz@\iventry{*swánaz}{}}

Derivation: citation reconstruction \emph{*swánaz}; form followed here
\emph{*swánas} \textgreater{} \emph{swanes} (early analogy).

\subsection{Derivation trace}\label{derivation-trace-130}

Proto input: \emph{*swánas}

\begingroup
\setlength{\fboxsep}{6pt}

\noindent\fbox{%
\begin{minipage}{0.97\linewidth}
\small
\begin{tabular}{@{}>{\raggedright\arraybackslash}p{0.359\linewidth}@{\hspace{0.030\linewidth}}>{\raggedright\arraybackslash}p{0.611\linewidth}@{}}
\begin{minipage}[t]{\linewidth}
\centering\textbf{Earlier Germanic changes}\par
\vspace{0.35em}
\raggedright
\centering\textbf{Northwest and West Germanic}\par
\raggedright
\vspace{0.2em}
\raggedright [no change]\par
\vspace{0.6em}
\centering\textbf{Early Anglo-Frisian}\par
\raggedright
\vspace{0.2em}
\raggedright [no change]\par
\end{minipage}
&
\begin{minipage}[t]{\linewidth}
\centering\textbf{Old English changes}\par
\vspace{0.35em}
\raggedright
\begin{tabularx}{\linewidth}{@{}>{\raggedright\arraybackslash}X@{\hspace{0.55em}}>{\raggedright\arraybackslash}p{0.14\linewidth}@{\hspace{0.25em}}}
\mbox{EAF Brightening Unstressed} & \emph{*swánæs} \\
\mbox{OE Unstressed AE Merger} & \emph{*swánes} \\
\end{tabularx}
\end{minipage}
\\
\end{tabular}
\end{minipage}%
} \endgroup

Old English form: \emph{swanes}

\subsection{Reconstruction and comparative
evidence}\label{reconstruction-and-comparative-evidence-113}

The Germanic noun is ordinarily cited as the masculine a-stem
\Recon{swánaz}\index[iv]{03pgmc@\ivlangheader{Proto-Germanic}{}!swanaz@\iventry{*swánaz}{}}
`swan' (\citeproc{ref-Orel2003}{Orel 2003, 367}). The derivational input
\Recon{swánas} `swan' is not a competing lexeme reconstruction. It is
the genitive singular
\Recon{swánas}\index[iv]{03pgmc@\ivlangheader{Proto-Germanic}{}!swanas@\iventry{*swánas}{}}
of the same paradigm.

The question here is therefore one of paradigm cell rather than stem
history. The citation form remains
\Recon{swánaz}\index[iv]{03pgmc@\ivlangheader{Proto-Germanic}{}!swanaz@\iventry{*swánaz}{}}
`swan' \textgreater{}
{\emph{swan}}\index[iv]{02oe@\ivlangheader{Old English}{West Saxon normalization unmarked}!swan@\iventry{swan}{}}
`swan'; the comparison form is the genitive singular
\Recon{swánas}\index[iv]{03pgmc@\ivlangheader{Proto-Germanic}{}!swanas@\iventry{*swánas}{}}
`swan' \textgreater{}
{\emph{swanes}}\index[iv]{02oe@\ivlangheader{Old English}{West Saxon normalization unmarked}!swanes@\iventry{swanes}{}}
`swan'.

\subsection{Old English evidence}\label{old-english-evidence-114}

Bright's glossary records the ordinary noun as
{\emph{swan}}\index[iv]{02oe@\ivlangheader{Old English}{West Saxon normalization unmarked}!swan@\iventry{swan}{}}
`swan', m. and also gives the exact inflected form
{\emph{swanes}}\index[iv]{02oe@\ivlangheader{Old English}{West Saxon normalization unmarked}!swanes@\iventry{swanes}{}}
`swan', citing the phrase \emph{swanes feðre}
(\citeproc{ref-BrightCassidyRingler1971}{Bright 1971, 441}).

The target is therefore an attested Old English genitive singular, not a
reconstruction or the ordinary citation lemma.

\subsection{Development to Old
English}\label{development-to-old-english-112}

From \Recon{swánas} `swan', Anglo-Frisian brightening gives
\Recon{swánæs} `swan', and subsequent merger of unstressed \emph{æ} with
\emph{e} yields \emph{swanes} `swan'. The comparison is straightforward
once the genitive singular is chosen as the relevant cell.

\subsection{Paradigm-cell comparison}\label{paradigm-cell-comparison}

The comparison below sets the relevant forms side by side. It separates
the ordinary citation form from the selected inflected cell.

{\def\LTcaptype{none} % do not increment counter
\begin{longtable}[]{@{}
  >{\raggedright\arraybackslash}p{(\linewidth - 8\tabcolsep) * \real{0.2000}}
  >{\raggedright\arraybackslash}p{(\linewidth - 8\tabcolsep) * \real{0.2000}}
  >{\raggedright\arraybackslash}p{(\linewidth - 8\tabcolsep) * \real{0.2000}}
  >{\raggedright\arraybackslash}p{(\linewidth - 8\tabcolsep) * \real{0.2000}}
  >{\raggedright\arraybackslash}p{(\linewidth - 8\tabcolsep) * \real{0.2000}}@{}}
\toprule\noalign{}
\begin{minipage}[b]{\linewidth}\raggedright
PGmc cell / interpretation
\end{minipage} & \begin{minipage}[b]{\linewidth}\raggedright
Candidate input
\end{minipage} & \begin{minipage}[b]{\linewidth}\raggedright
OE output or comparison
\end{minipage} & \begin{minipage}[b]{\linewidth}\raggedright
OE comparison form
\end{minipage} & \begin{minipage}[b]{\linewidth}\raggedright
Result
\end{minipage} \\
\midrule\noalign{}
\endhead
\bottomrule\noalign{}
\endlastfoot
citation nominative singular &
{\emph{*swánaz}}\index[iv]{03pgmc@\ivlangheader{Proto-Germanic}{}!swanaz@\iventry{*swánaz}{}}
& OE headword
{\emph{swan}}\index[iv]{02oe@\ivlangheader{Old English}{West Saxon normalization unmarked}!swan@\iventry{swan}{}}
`swan' &
{\emph{swan}}\index[iv]{02oe@\ivlangheader{Old English}{West Saxon normalization unmarked}!swan@\iventry{swan}{}}
`swan' & ordinary lexeme line \\
genitive singular &
{\emph{*swánas}}\index[iv]{03pgmc@\ivlangheader{Proto-Germanic}{}!swanas@\iventry{*swánas}{}}
& regular output: {\emph{swanes}} `swan' &
{\emph{swanes}}\index[iv]{02oe@\ivlangheader{Old English}{West Saxon normalization unmarked}!swanes@\iventry{swanes}{}}
`swan' & attested cell \\
\end{longtable}
}

\section{\texorpdfstring{thousand --- OE
\emph{þūsend}}{thousand --- OE þūsend}}\label{thousand-oe-uxfeux16bsend}

\index[iv]{02oe@\ivlangheader{Old English}{West Saxon normalization unmarked}!aerende@\iventry{ærende}{}}
\index[iv]{02oe@\ivlangheader{Old English}{West Saxon normalization unmarked}!thusend@\iventry{þūsend}{}}
\index[iv]{03pgmc@\ivlangheader{Proto-Germanic}{}!thusundi@\iventry{*θūs-undī}{}}
\index[iv]{08preoe@\ivlangheader{Pre-Old English / prehistoric English}{}!thusendi~01@\iventry{*θūs-èndi}{transp.}}
\index[iv]{08preoe@\ivlangheader{Pre-Old English / prehistoric English}{}!thusendi~01@\iventry{*θūsèndi}{transp.}}

Derivation: citation reconstruction \emph{*θūs-undī}; form followed here
\emph{*θūs-èndi} \textgreater{} \emph{þūsend} (early analogy).

\subsection{Derivation trace}\label{derivation-trace-131}

Proto input: \emph{*θūsèndi}

\begingroup
\setlength{\fboxsep}{6pt}

\noindent\fbox{%
\begin{minipage}{0.97\linewidth}
\small
\begin{tabular}{@{}>{\raggedright\arraybackslash}p{0.359\linewidth}@{\hspace{0.030\linewidth}}>{\raggedright\arraybackslash}p{0.611\linewidth}@{}}
\begin{minipage}[t]{\linewidth}
\centering\textbf{Earlier Germanic changes}\par
\vspace{0.35em}
\raggedright
\centering\textbf{Northwest and West Germanic}\par
\raggedright
\vspace{0.2em}
\begin{tabularx}{\linewidth}{@{}>{\raggedright\arraybackslash}X@{\hspace{0.55em}}>{\raggedright\arraybackslash}p{0.24\linewidth}@{\hspace{0.25em}}}
\mbox{PWGmc Early I Apocope} & \emph{*θūsènd} \\
\end{tabularx}
\vspace{0.6em}
\centering\textbf{Early Anglo-Frisian}\par
\raggedright
\vspace{0.2em}
\raggedright [no change]\par
\end{minipage}
&
\begin{minipage}[t]{\linewidth}
\centering\textbf{Old English changes}\par
\vspace{0.35em}
\raggedright
\begin{tabularx}{\linewidth}{@{}>{\raggedright\arraybackslash}X@{\hspace{0.55em}}>{\raggedright\arraybackslash}p{0.14\linewidth}@{\hspace{0.25em}}}
\mbox{OE Strip Secondary Stress} & \emph{*θūsend} \\
\end{tabularx}
\end{minipage}
\\
\end{tabular}
\end{minipage}%
} \endgroup

Old English form: \emph{þūsend}

\subsection{Reconstruction and comparative
evidence}\label{reconstruction-and-comparative-evidence-114}

Kroonen reconstructs the Germanic numeral as
{\emph{*þūsundī-}}\index[iv]{03pgmc@\ivlangheader{Proto-Germanic}{}!thusundi@\iventry{*þūsundī-}{}}
and cites Old English
{\emph{þūsend}}\index[iv]{02oe@\ivlangheader{Old English}{West Saxon normalization unmarked}!thusend@\iventry{þūsend}{}}
`thousand' among its continuations (\citeproc{ref-Kroonen2013}{Kroonen
2013, 554}). The derivational input
\Recon{θūs-èndi}\index[iv]{08preoe@\ivlangheader{Pre-Old English / prehistoric English}{}!thusendi~01@\iventry{*θūs-èndi}{transp.}}
`thousand' is not the same claim. It is an OE-oriented transponent with
the second-member vowel already resolved to \emph{e} and the final high
vowel already shortened for apocope.

The chronology must explain why Old English shows
{\emph{þūsend}}\index[iv]{02oe@\ivlangheader{Old English}{West Saxon normalization unmarked}!thusend@\iventry{þūsend}{}}
`thousand' while related languages such as Old Saxon and Old High German
keep \emph{u} in the second syllable?
(\citeproc{ref-Kroonen2013}{Kroonen 2013, 554}).

\subsection{Old English evidence}\label{old-english-evidence-115}

Old English \emph{þūsend} `thousand' is an ordinary citation form, not a
selected oblique or paradigm cell. Campbell treats it as a neuter noun
with normal case forms (\citeproc{ref-Campbell1959}{Campbell 1959, § 689}). The problem lies in the internal history of the word, not in its
lexical status.

\subsection{Development to Old
English}\label{development-to-old-english-113}

If the old final \emph{-ī} had remained long enough to trigger ordinary
double umlaut, Campbell's rule would point toward a form of the
{\emph{*þȳsend}}\index[iv]{02oe@\ivlangheader{Old English}{West Saxon normalization unmarked}!thysend@\iventry{*þȳsend}{}}
`thousand' type rather than attested
{\emph{þūsend}}\index[iv]{02oe@\ivlangheader{Old English}{West Saxon normalization unmarked}!thusend@\iventry{þūsend}{}}
`thousand' (\citeproc{ref-Campbell1959}{Campbell 1959, § 203}).
Preserved root \emph{ū} therefore argues that the umlaut-triggering
vowel was lost or neutralized before the ordinary OE umlaut could
develop.

That early loss, however, does not by itself explain the medial
\emph{e}. Luick compares the word with
{\emph{ærende}}\index[iv]{02oe@\ivlangheader{Old English}{West Saxon normalization unmarked}!aerende@\iventry{ærende}{}}
`errand/message' and later groups \emph{thousand} with forms reshaped on
that pattern (\citeproc{ref-Luick1914}{Luick 1914-\/-1940, §§ 198,
492}). Viredaz is more cautious, arguing that Old English \emph{e} in
this weak position may simply write schwa and so need not prove a unique
\emph{ærende} `message'-type analogy
(\citeproc{ref-GermanicSlavicBaltic2025}{Viredaz 2025, § 2.1.4}).

The selected transponent
\Recon{θūs-èndi}\index[iv]{08preoe@\ivlangheader{Pre-Old English / prehistoric English}{}!thusendi~01@\iventry{*θūs-èndi}{transp.}}
`thousand' captures the OE-side state from which the regular derivation
reaches {\emph{þūsend}} `thousand'.

\subsection{Stage comparison}\label{stage-comparison-2}

The comparison below sets the relevant forms side by side. It separates
the secure chronology from the more interpretive account of the
second-syllable vowel.

{\def\LTcaptype{none} % do not increment counter
\begin{longtable}[]{@{}
  >{\raggedright\arraybackslash}p{(\linewidth - 6\tabcolsep) * \real{0.2500}}
  >{\raggedright\arraybackslash}p{(\linewidth - 6\tabcolsep) * \real{0.2500}}
  >{\raggedright\arraybackslash}p{(\linewidth - 6\tabcolsep) * \real{0.2500}}
  >{\raggedright\arraybackslash}p{(\linewidth - 6\tabcolsep) * \real{0.2500}}@{}}
\toprule\noalign{}
\begin{minipage}[b]{\linewidth}\raggedright
Stage / interpretation
\end{minipage} & \begin{minipage}[b]{\linewidth}\raggedright
Candidate form
\end{minipage} & \begin{minipage}[b]{\linewidth}\raggedright
OE relation
\end{minipage} & \begin{minipage}[b]{\linewidth}\raggedright
Result
\end{minipage} \\
\midrule\noalign{}
\endhead
\bottomrule\noalign{}
\endlastfoot
surviving \emph{-ī} with ordinary double umlaut &
{\emph{*þūsundī-}}\index[iv]{03pgmc@\ivlangheader{Proto-Germanic}{}!thusundi@\iventry{*þūsundī-}{}}
treated as still umlaut-active in OE & would point toward
{\emph{*þȳsend}}\index[iv]{02oe@\ivlangheader{Old English}{West Saxon normalization unmarked}!thysend@\iventry{*þȳsend}{}}
& excluded by preserved \emph{ū} \\
early loss of the trigger without further reshaping & \emph{*þūsund-}
type & explains \emph{ū}, but not why OE alone has medial \emph{e} &
incomplete account \\
selected OE-oriented transponent &
{\emph{*θūs-èndi}}\index[iv]{08preoe@\ivlangheader{Pre-Old English / prehistoric English}{}!thusendi~01@\iventry{*θūs-èndi}{transp.}}
& regular output: {\emph{þūsend}} `thousand' & selected modeling
input \\
\end{longtable}
}

\section{\texorpdfstring{timber --- OE
\emph{timber}}{timber --- OE timber}}\label{timber-oe-timber}

\index[iv]{02oe@\ivlangheader{Old English}{West Saxon normalization unmarked}!timber@\iventry{timber}{}}
\index[iv]{03pgmc@\ivlangheader{Proto-Germanic}{}!timbra@\iventry{*tímbrą}{}}
\index[iv]{03pgmc@\ivlangheader{Proto-Germanic}{}!timra@\iventry{*tímrą}{}}

Derivation: citation reconstruction \emph{*tímrą}; form followed here
\emph{*tímbrą} \textgreater{} \emph{timber} (early analogy).

\subsection{Derivation trace}\label{derivation-trace-132}

Proto input: \emph{*tímbrą}

\begingroup
\setlength{\fboxsep}{6pt}

\noindent\fbox{%
\begin{minipage}{0.97\linewidth}
\small
\begin{tabular}{@{}>{\raggedright\arraybackslash}p{0.359\linewidth}@{\hspace{0.030\linewidth}}>{\raggedright\arraybackslash}p{0.611\linewidth}@{}}
\begin{minipage}[t]{\linewidth}
\centering\textbf{Earlier Germanic changes}\par
\vspace{0.35em}
\raggedright
\centering\textbf{Northwest and West Germanic}\par
\raggedright
\vspace{0.2em}
\raggedright [no change]\par
\vspace{0.6em}
\centering\textbf{Early Anglo-Frisian}\par
\raggedright
\vspace{0.2em}
\raggedright [no change]\par
\end{minipage}
&
\begin{minipage}[t]{\linewidth}
\centering\textbf{Old English changes}\par
\vspace{0.35em}
\raggedright
\begin{tabularx}{\linewidth}{@{}>{\raggedright\arraybackslash}X@{\hspace{0.55em}}>{\raggedright\arraybackslash}p{0.14\linewidth}@{\hspace{0.25em}}}
OE Heavy Syllable Nasal Apocope & \emph{*tímbr} \\
\mbox{OE Epenthetic Vowel} & \emph{*tímber} \\
\end{tabularx}
\end{minipage}
\\
\end{tabular}
\end{minipage}%
} \endgroup

Old English form: \emph{timber}

\subsection{Reconstruction and comparative
evidence}\label{reconstruction-and-comparative-evidence-115}

Kroonen reconstructs the noun as
{\emph{*timbra-}}\index[iv]{03pgmc@\ivlangheader{Proto-Germanic}{}!timbra@\iventry{*timbra-}{}}
and cites Old English
{\emph{timber}}\index[iv]{02oe@\ivlangheader{Old English}{West Saxon normalization unmarked}!timber@\iventry{timber}{}}
`timber' among its continuations (\citeproc{ref-Kroonen2013}{Kroonen
2013, 517}). Ringe and Taylor instead state the history from
\Recon{timra} `timber' through West Germanic
\Recon{timbr}\index[iv]{05pwgmc@\ivlangheader{Proto-West Germanic}{}!timbr@\iventry{*timbr}{}}
`timber' to Old English
{\emph{timber}}\index[iv]{02oe@\ivlangheader{Old English}{West Saxon normalization unmarked}!timber@\iventry{timber}{}}
`timber' (\citeproc{ref-RingeTaylor2014}{Ringe and Taylor 2014, 327}).

The difference is therefore not over the Old English noun itself. It
concerns whether medial \emph{b} belongs in the comparative citation
form or appears in an early pre-Old-English stage of the cluster.

\subsection{Old English evidence}\label{old-english-evidence-116}

Clark Hall lemmatizes the noun as
{\emph{timber}}\index[iv]{02oe@\ivlangheader{Old English}{West Saxon normalization unmarked}!timber@\iventry{timber}{}}
`timber' and also records
{\emph{timbor}}\index[iv]{02oe@\ivlangheader{Old English}{West Saxon normalization unmarked}!timbor@\iventry{timbor}{}}
`timber' as a variant spelling (\citeproc{ref-ClarkHall1960}{Clark Hall
1960, 294}). The Old English form is thus an ordinary citation noun, not
a selected oblique cell or a reconstructed convenience form.

\subsection{Development to Old
English}\label{development-to-old-english-114}

With the consonantal frame \Recon{timbr-} `timber' in place, the rest of
the development is straightforward. Loss of final \emph{-ą} leaves
\Recon{tímbr} `timber', and epenthetic \emph{e} in the final cluster
yields {\emph{timber}} `timber'. Ringe and Taylor's \emph{timra}
\textgreater{} \emph{timbr} \textgreater{} OE \emph{timber} and the
handbook treatment of this epenthetic vowel point to the same Old
English result (\citeproc{ref-RingeTaylor2014}{Ringe and Taylor 2014,
327}; \citeproc{ref-Campbell1959}{Campbell 1959, §§ 463--464}).

\subsection{Formation comparison}\label{formation-comparison-6}

The comparison below sets the relevant forms side by side. It separates
the comparative headword from the OE-facing consonantal input.

{\def\LTcaptype{none} % do not increment counter
\begin{longtable}[]{@{}
  >{\raggedright\arraybackslash}p{(\linewidth - 6\tabcolsep) * \real{0.2500}}
  >{\raggedright\arraybackslash}p{(\linewidth - 6\tabcolsep) * \real{0.2500}}
  >{\raggedright\arraybackslash}p{(\linewidth - 6\tabcolsep) * \real{0.2500}}
  >{\raggedright\arraybackslash}p{(\linewidth - 6\tabcolsep) * \real{0.2500}}@{}}
\toprule\noalign{}
\begin{minipage}[b]{\linewidth}\raggedright
Formation or notation
\end{minipage} & \begin{minipage}[b]{\linewidth}\raggedright
Candidate form
\end{minipage} & \begin{minipage}[b]{\linewidth}\raggedright
OE relation
\end{minipage} & \begin{minipage}[b]{\linewidth}\raggedright
Result
\end{minipage} \\
\midrule\noalign{}
\endhead
\bottomrule\noalign{}
\endlastfoot
Kroonen's comparative citation &
{\emph{*timbra-}}\index[iv]{03pgmc@\ivlangheader{Proto-Germanic}{}!timbra@\iventry{*timbra-}{}}
& already matches the consonantal frame of OE
{\emph{timber}}\index[iv]{02oe@\ivlangheader{Old English}{West Saxon normalization unmarked}!timber@\iventry{timber}{}}
`timber' & closest comparative support for the derivational input \\
Ringe-Taylor citation line &
{\emph{*timra}}\index[iv]{03pgmc@\ivlangheader{Proto-Germanic}{}!timra@\iventry{*timra}{}}
\textgreater{}
\Recon{timbr}\index[iv]{05pwgmc@\ivlangheader{Proto-West Germanic}{}!timbr@\iventry{*timbr}{}}
`timber' & reaches the same OE noun through early cluster expansion &
compatible comparative background \\
modeled input &
{\emph{*tímbrą}}\index[iv]{03pgmc@\ivlangheader{Proto-Germanic}{}!timbra@\iventry{*tímbrą}{}}
& regular output: {\emph{timber}} `timber' & Old English-facing input \\
\end{longtable}
}

\section{\texorpdfstring{wake --- OE
\emph{wacan}}{wake --- OE wacan}}\label{wake-oe-wacan}

\index[iv]{02oe@\ivlangheader{Old English}{West Saxon normalization unmarked}!wacan@\iventry{wacan}{}}
\index[iv]{03pgmc@\ivlangheader{Proto-Germanic}{}!wakana@\iventry{*wákaną}{}}
\index[iv]{03pgmc@\ivlangheader{Proto-Germanic}{}!wakena@\iventry{*wakēną}{}}

Derivation: citation reconstruction \emph{*wakēną}; form followed here
\emph{*wákaną} \textgreater{} \emph{wacan} (early analogy).

\subsection{Derivation trace}\label{derivation-trace-133}

Proto input: \emph{*wákaną}

\begingroup
\setlength{\fboxsep}{6pt}

\noindent\fbox{%
\begin{minipage}{0.97\linewidth}
\small
\begin{tabular}{@{}>{\raggedright\arraybackslash}p{0.359\linewidth}@{\hspace{0.030\linewidth}}>{\raggedright\arraybackslash}p{0.611\linewidth}@{}}
\begin{minipage}[t]{\linewidth}
\centering\textbf{Earlier Germanic changes}\par
\vspace{0.35em}
\raggedright
\centering\textbf{Northwest and West Germanic}\par
\raggedright
\vspace{0.2em}
\raggedright [no change]\par
\vspace{0.6em}
\centering\textbf{Early Anglo-Frisian}\par
\raggedright
\vspace{0.2em}
\raggedright [no change]\par
\end{minipage}
&
\begin{minipage}[t]{\linewidth}
\centering\textbf{Old English changes}\par
\vspace{0.35em}
\raggedright
\begin{tabularx}{\linewidth}{@{}>{\raggedright\arraybackslash}X@{\hspace{0.55em}}>{\raggedright\arraybackslash}p{0.14\linewidth}@{\hspace{0.25em}}}
\mbox{EAF Brightening} & \emph{*wækaną} \\
\mbox{OE A Restoration} & \emph{*wakaną} \\
OE Heavy Syllable Nasal Apocope & \emph{*wakan} \\
\mbox{OE Secondary Nasalization} & \emph{*wakąn} \\
\mbox{OE Weak Tail Reduction} & \emph{*wakan} \\
\end{tabularx}
\end{minipage}
\\
\end{tabular}
\end{minipage}%
} \endgroup

Old English form: \emph{wacan}

\subsection{Reconstruction and comparative
evidence}\label{reconstruction-and-comparative-evidence-116}

Kroonen gives the strong verb as \emph{*wakan-} with Old English
{\emph{wacan}}\index[iv]{02oe@\ivlangheader{Old English}{West Saxon normalization unmarked}!wacan@\iventry{wacan}{}}
`wake' (\citeproc{ref-Kroonen2013}{Kroonen 2013, 568}). Ringe and Taylor
separately derive Old English
{\emph{wacian}}\index[iv]{02oe@\ivlangheader{Old English}{West Saxon normalization unmarked}!wacian@\iventry{wacian}{}}
`wake' from weak \emph{*wakai-} \textasciitilde{} \emph{*wakja-}
(\citeproc{ref-RingeTaylor2014}{Ringe and Taylor 2014, § 3.3.2}).

The difference is therefore lexical and class-based, not graphic. Strong
{\emph{wacan}}\index[iv]{02oe@\ivlangheader{Old English}{West Saxon normalization unmarked}!wacan@\iventry{wacan}{}}
`wake' `wake up, arise' and weak
{\emph{wacian}}\index[iv]{02oe@\ivlangheader{Old English}{West Saxon normalization unmarked}!wacian@\iventry{wacian}{}}
`wake' `be awake, watch' belong to related but distinct histories.

\subsection{Old English evidence}\label{old-english-evidence-117}

Clark Hall lists
{\emph{wacan}}\index[iv]{02oe@\ivlangheader{Old English}{West Saxon normalization unmarked}!wacan@\iventry{wacan}{}}
`wake' and
{\emph{wacian}}\index[iv]{02oe@\ivlangheader{Old English}{West Saxon normalization unmarked}!wacian@\iventry{wacian}{}}
`wake' as separate headwords (\citeproc{ref-ClarkHall1960}{Clark Hall
1960, 338}). Bosworth-Toller cautions under
{\emph{wacan}}\index[iv]{02oe@\ivlangheader{Old English}{West Saxon normalization unmarked}!wacan@\iventry{wacan}{}}
`wake': the simplex infinitive itself does not occur, its place seeming
to be taken by
{\emph{wæcnan}}\index[iv]{02oe@\ivlangheader{Old English}{West Saxon normalization unmarked}!waecnan@\iventry{wæcnan}{}}
`wake' (\citeproc{ref-BosworthToller1898}{Bosworth and Toller 1898,
226}).

The target
{\emph{wacan}}\index[iv]{02oe@\ivlangheader{Old English}{West Saxon normalization unmarked}!wacan@\iventry{wacan}{}}
`wake' is therefore best understood as a normalized strong headword for
the verb family, not as a directly quoted simplex infinitive. It still
remains the correct Old English comparison form for the strong branch.

\subsection{Development to Old
English}\label{development-to-old-english-115}

With strong
\Recon{wákaną}\index[iv]{03pgmc@\ivlangheader{Proto-Germanic}{}!wakana@\iventry{*wákaną}{}}
`wake', Anglo-Frisian brightening first gives a form of the
\Recon{wækaną} `wake' type. A-restoration then returns \emph{a}, and the
ordinary tail reductions yield {\emph{wacan}} `wake'. The weak verb
{\emph{wacian}}\index[iv]{02oe@\ivlangheader{Old English}{West Saxon normalization unmarked}!wacian@\iventry{wacian}{}}
`wake' belongs to a different prehistory and is not the expected outcome
of this input.

\subsection{Class comparison}\label{class-comparison-4}

The comparison below sets the relevant forms side by side. It separates
the strong and weak verb lines.

{\def\LTcaptype{none} % do not increment counter
\begin{longtable}[]{@{}
  >{\raggedright\arraybackslash}p{(\linewidth - 6\tabcolsep) * \real{0.2500}}
  >{\raggedright\arraybackslash}p{(\linewidth - 6\tabcolsep) * \real{0.2500}}
  >{\raggedright\arraybackslash}p{(\linewidth - 6\tabcolsep) * \real{0.2500}}
  >{\raggedright\arraybackslash}p{(\linewidth - 6\tabcolsep) * \real{0.2500}}@{}}
\toprule\noalign{}
\begin{minipage}[b]{\linewidth}\raggedright
Formation / class
\end{minipage} & \begin{minipage}[b]{\linewidth}\raggedright
Candidate input
\end{minipage} & \begin{minipage}[b]{\linewidth}\raggedright
OE outcome or comparison
\end{minipage} & \begin{minipage}[b]{\linewidth}\raggedright
Result
\end{minipage} \\
\midrule\noalign{}
\endhead
\bottomrule\noalign{}
\endlastfoot
weak class-III / class-II branch &
{\emph{*wakēną}}\index[iv]{03pgmc@\ivlangheader{Proto-Germanic}{}!wakena@\iventry{*wakēną}{}},
\emph{*wakai-} \textasciitilde{} \emph{*wakja-} & OE
{\emph{wacian}}\index[iv]{02oe@\ivlangheader{Old English}{West Saxon normalization unmarked}!wacian@\iventry{wacian}{}}
`wake' and related weak forms & related lexeme, but not the target of
this entry \\
strong class-VI branch &
{\emph{*wákaną}}\index[iv]{03pgmc@\ivlangheader{Proto-Germanic}{}!wakana@\iventry{*wákaną}{}}
& regular output: {\emph{wacan}} `wake' & Old English-facing input \\
strong normalized headword &
{\emph{wacan}}\index[iv]{02oe@\ivlangheader{Old English}{West Saxon normalization unmarked}!wacan@\iventry{wacan}{}}
`wake' & dictionary comparison form beside attested strong-family forms
& correct Old English comparator, though not a directly quoted simplex
infinitive \\
\end{longtable}
}

\section{\texorpdfstring{water --- OE
\emph{wæter}}{water --- OE wæter}}\label{water-oe-wuxe6ter}

\index[iv]{02oe@\ivlangheader{Old English}{West Saxon normalization unmarked}!waeter@\iventry{wæter}{}}
\index[iv]{03pgmc@\ivlangheader{Proto-Germanic}{}!watna@\iventry{*wátną}{}}
\index[iv]{03pgmc@\ivlangheader{Proto-Germanic}{}!wator@\iventry{*wátōr}{}}

Derivation: citation reconstruction \emph{*wátną}; form followed here
\emph{*wátōr} \textgreater{} \emph{wæter} (early analogy).

\subsection{Derivation trace}\label{derivation-trace-134}

Proto input: \emph{*wátōr}

\begingroup
\setlength{\fboxsep}{6pt}

\noindent\fbox{%
\begin{minipage}{0.97\linewidth}
\small
\begin{tabular}{@{}>{\raggedright\arraybackslash}p{0.359\linewidth}@{\hspace{0.030\linewidth}}>{\raggedright\arraybackslash}p{0.611\linewidth}@{}}
\begin{minipage}[t]{\linewidth}
\centering\textbf{Earlier Germanic changes}\par
\vspace{0.35em}
\raggedright
\centering\textbf{Northwest and West Germanic}\par
\raggedright
\vspace{0.2em}
\begin{tabularx}{\linewidth}{@{}>{\raggedright\arraybackslash}X@{\hspace{0.55em}}>{\raggedright\arraybackslash}p{0.14\linewidth}@{\hspace{0.25em}}}
\mbox{PWGmc Final Or Lowering} & \emph{*wátar} \\
\end{tabularx}
\vspace{0.6em}
\centering\textbf{Early Anglo-Frisian}\par
\raggedright
\vspace{0.2em}
\raggedright [no change]\par
\end{minipage}
&
\begin{minipage}[t]{\linewidth}
\centering\textbf{Old English changes}\par
\vspace{0.35em}
\raggedright
\begin{tabularx}{\linewidth}{@{}>{\raggedright\arraybackslash}X@{\hspace{0.55em}}>{\raggedright\arraybackslash}p{0.14\linewidth}@{\hspace{0.25em}}}
\mbox{EAF Brightening Unstressed} & \emph{*wátær} \\
\mbox{EAF Brightening} & \emph{*wætær} \\
\mbox{OE Unstressed AE Merger} & \emph{*wæter} \\
\end{tabularx}
\end{minipage}
\\
\end{tabular}
\end{minipage}%
} \endgroup

Old English form: \emph{wæter}

\subsection{Reconstruction and comparative
evidence}\label{reconstruction-and-comparative-evidence-117}

Kroonen reconstructs a heteroclitic noun \emph{*watar-}
\textasciitilde{} \emph{*watan-} and states that the Proto-Germanic
material points to \Recon{watōr} `water' and \Recon{watenaz} `water'
(\citeproc{ref-Kroonen2013}{Kroonen 2013, 616}). Ringe and Taylor
likewise start from singular \Recon{wator} `water' before the Old
English branch (\citeproc{ref-RingeTaylor2014}{Ringe and Taylor 2014,
§ 3.1.4}).

The generalized comparative label is therefore broader than the singular
form that actually corresponds to Old English
{\emph{wæter}}\index[iv]{02oe@\ivlangheader{Old English}{West Saxon normalization unmarked}!waeter@\iventry{wæter}{}}
`water'. The relevant comparator is the inherited nominative-accusative
singular
\Recon{wátōr}\index[iv]{03pgmc@\ivlangheader{Proto-Germanic}{}!wator@\iventry{*wátōr}{}}
`water'.

\subsection{Old English evidence}\label{old-english-evidence-118}

Bright gives the noun as
{\emph{wæter}}\index[iv]{02oe@\ivlangheader{Old English}{West Saxon normalization unmarked}!waeter@\iventry{wæter}{}}
`water' with the regular paradigm
{\emph{wæteres}}\index[iv]{02oe@\ivlangheader{Old English}{West Saxon normalization unmarked}!waeteres@\iventry{wæteres}{}}
`water (gen.)',
{\emph{wætere}}\index[iv]{02oe@\ivlangheader{Old English}{West Saxon normalization unmarked}!waetere@\iventry{wætere}{}}
`water (dat.)', \emph{wæter(u)},
{\emph{wætera}}\index[iv]{02oe@\ivlangheader{Old English}{West Saxon normalization unmarked}!waetera@\iventry{wætera}{}}
`water (gen.pl.)',
{\emph{wæterum}}\index[iv]{02oe@\ivlangheader{Old English}{West Saxon normalization unmarked}!waeterum@\iventry{wæterum}{}}
`water (dat.pl.)' (\citeproc{ref-BrightCassidyRingler1971}{Bright 1971,
29}). Ringe and Taylor add the dialectal contrast between West Saxon
\emph{weeter} `water' and Mercian \emph{weter} `water'
(\citeproc{ref-RingeTaylor2014}{Ringe and Taylor 2014, § 6.5.2}).

The target is therefore an attested Old English citation form within a
normal paradigm. The complication lies on the comparative side of the
lexeme, not in Old English attestation.

\subsection{Development to Old
English}\label{development-to-old-english-116}

From \Recon{wátōr} `water', pre-final \emph{ō} becomes \emph{a},
yielding \Recon{watar} `water' (\citeproc{ref-RingeTaylor2014}{Ringe and
Taylor 2014, § 3.1.4}). Anglo-Frisian brightening then gives
\Recon{wætær} `water', and merger of unstressed \emph{æ}/\emph{e} yields
{\emph{wæter}} `water'.

\subsection{Stage comparison}\label{stage-comparison-3}

The comparison below sets the relevant forms side by side. It separates
the generalized lexeme label from the singular input that matches the
Old English citation form.

{\def\LTcaptype{none} % do not increment counter
\begin{longtable}[]{@{}
  >{\raggedright\arraybackslash}p{(\linewidth - 6\tabcolsep) * \real{0.2500}}
  >{\raggedright\arraybackslash}p{(\linewidth - 6\tabcolsep) * \real{0.2500}}
  >{\raggedright\arraybackslash}p{(\linewidth - 6\tabcolsep) * \real{0.2500}}
  >{\raggedright\arraybackslash}p{(\linewidth - 6\tabcolsep) * \real{0.2500}}@{}}
\toprule\noalign{}
\begin{minipage}[b]{\linewidth}\raggedright
Stage or notation
\end{minipage} & \begin{minipage}[b]{\linewidth}\raggedright
Candidate form
\end{minipage} & \begin{minipage}[b]{\linewidth}\raggedright
OE relation
\end{minipage} & \begin{minipage}[b]{\linewidth}\raggedright
Result
\end{minipage} \\
\midrule\noalign{}
\endhead
\bottomrule\noalign{}
\endlastfoot
generalized comparative label &
{\emph{*wátną}}\index[iv]{03pgmc@\ivlangheader{Proto-Germanic}{}!watna@\iventry{*wátną}{}}
& broader lexeme shorthand, not the singular that corresponds directly
to
{\emph{wæter}}\index[iv]{02oe@\ivlangheader{Old English}{West Saxon normalization unmarked}!waeter@\iventry{wæter}{}}
`water' & useful background only \\
heteroclitic stem notation & \emph{*watar-} \textasciitilde{}
\emph{*watan-} & source-faithful comparative reconstruction & explains
why a singular comparator is needed \\
inherited singular input &
{\emph{*wátōr}}\index[iv]{03pgmc@\ivlangheader{Proto-Germanic}{}!wator@\iventry{*wátōr}{}}
& regular output: {\emph{wæter}} `water' & Old English-facing input \\
\end{longtable}
}

\section{\texorpdfstring{whale --- OE
\emph{hwæl}}{whale --- OE hwæl}}\label{whale-oe-hwuxe6l}

\index[iv]{02oe@\ivlangheader{Old English}{West Saxon normalization unmarked}!hwael@\iventry{hwæl}{}}
\index[iv]{03pgmc@\ivlangheader{Proto-Germanic}{}!walaz@\iventry{*wálaz}{}}
\index[iv]{03pgmc@\ivlangheader{Proto-Germanic}{}!xwalaz@\iventry{*xwálaz}{}}

Derivation: citation reconstruction \emph{*wálaz}; form followed here
\emph{*xwálaz} \textgreater{} \emph{hwæl} (early analogy).

\subsection{Derivation trace}\label{derivation-trace-135}

Proto input: \emph{*xwálaz}

\begingroup
\setlength{\fboxsep}{6pt}

\noindent\fbox{%
\begin{minipage}{0.97\linewidth}
\small
\begin{tabular}{@{}>{\raggedright\arraybackslash}p{0.485\linewidth}@{\hspace{0.030\linewidth}}>{\raggedright\arraybackslash}p{0.485\linewidth}@{}}
\begin{minipage}[t]{\linewidth}
\centering\textbf{Earlier Germanic changes}\par
\vspace{0.35em}
\raggedright
\centering\textbf{Northwest and West Germanic}\par
\raggedright
\vspace{0.2em}
\begin{tabularx}{\linewidth}{@{}>{\raggedright\arraybackslash}X@{\hspace{0.55em}}>{\raggedright\arraybackslash}p{0.14\linewidth}@{\hspace{0.25em}}}
\mbox{EAF Final Z Deletion} & \emph{*xwála} \\
\end{tabularx}
\vspace{0.6em}
\centering\textbf{Early Anglo-Frisian}\par
\raggedright
\vspace{0.2em}
\raggedright [no change]\par
\end{minipage}
&
\begin{minipage}[t]{\linewidth}
\centering\textbf{Old English changes}\par
\vspace{0.35em}
\raggedright
\begin{tabularx}{\linewidth}{@{}>{\raggedright\arraybackslash}X@{\hspace{0.55em}}>{\raggedright\arraybackslash}p{0.14\linewidth}@{\hspace{0.25em}}}
\mbox{PWGmc Final Bare A Loss} & \emph{*xwál} \\
\mbox{EAF Brightening} & \emph{*xwæl} \\
\end{tabularx}
\end{minipage}
\\
\end{tabular}
\end{minipage}%
} \endgroup

Old English form: \emph{hwæl}

\subsection{Reconstruction and comparative
evidence}\label{reconstruction-and-comparative-evidence-118}

The comparative sources are not uniform. Orel gives
\Recon{xwalaz}\index[iv]{03pgmc@\ivlangheader{Proto-Germanic}{}!xwalaz@\iventry{*xwalaz}{}}
`whale' and notes some mixed
\Recon{xwaliz}\index[iv]{03pgmc@\ivlangheader{Proto-Germanic}{}!xwaliz@\iventry{*xwaliz}{}}
`whale' evidence (\citeproc{ref-Orel2003}{Orel 2003, 197}). Kroonen
instead cites
{\emph{*hwali-}}\index[iv]{03pgmc@\ivlangheader{Proto-Germanic}{}!hwali@\iventry{*hwali-}{}}
(\citeproc{ref-Kroonen2013}{Kroonen 2013, 262}).

Both notations agree on inherited initial \emph{hw-/xw-}, but they
differ in stem label. The a-stem-like input followed here is closer to
Orel's notation than to Kroonen's citation form.

\subsection{Old English evidence}\label{old-english-evidence-119}

Clark Hall lemmatizes the noun as
{\emph{hwal}}\index[iv]{02oe@\ivlangheader{Old English}{West Saxon normalization unmarked}!hwal@\iventry{hwal}{}}
`whale', and Bosworth-Toller preserves plural
{\emph{hwalas}}\index[iv]{02oe@\ivlangheader{Old English}{West Saxon normalization unmarked}!hwalas@\iventry{hwalas}{}}
`whales' (\citeproc{ref-ClarkHall1960}{Clark Hall 1960, 170};
\citeproc{ref-BosworthToller1898}{Bosworth and Toller 1898, 326}). The
comparison form is normalized here as
{\emph{hwæl}}\index[iv]{02oe@\ivlangheader{Old English}{West Saxon normalization unmarked}!hwael@\iventry{hwæl}{}}
`whale' for the singular citation form with Anglo-Frisian fronting.

The plural
{\emph{hwalas}}\index[iv]{02oe@\ivlangheader{Old English}{West Saxon normalization unmarked}!hwalas@\iventry{hwalas}{}}
`whales' supplies control evidence. It shows the same lexeme with
\emph{a} in an open syllable beside singular
{\emph{hwæl}}\index[iv]{02oe@\ivlangheader{Old English}{West Saxon normalization unmarked}!hwael@\iventry{hwæl}{}}
`whale' in a closed monosyllable.

\subsection{Development to Old
English}\label{development-to-old-english-117}

From \Recon{xwálaz} `whale', final \emph{-z} disappears and bare final
\emph{-a} is lost. Anglo-Frisian fronting then yields \emph{æ} in the
closed monosyllable, and Old English orthography writes
{\emph{hwæl}}\index[iv]{02oe@\ivlangheader{Old English}{West Saxon normalization unmarked}!hwael@\iventry{hwæl}{}}
`whale'.

\subsection{Formation comparison}\label{formation-comparison-7}

The comparison below sets the relevant forms side by side. It separates
the competing comparative notations from the normalized Old English
singular.

{\def\LTcaptype{none} % do not increment counter
\begin{longtable}[]{@{}
  >{\raggedright\arraybackslash}p{(\linewidth - 6\tabcolsep) * \real{0.2500}}
  >{\raggedright\arraybackslash}p{(\linewidth - 6\tabcolsep) * \real{0.2500}}
  >{\raggedright\arraybackslash}p{(\linewidth - 6\tabcolsep) * \real{0.2500}}
  >{\raggedright\arraybackslash}p{(\linewidth - 6\tabcolsep) * \real{0.2500}}@{}}
\toprule\noalign{}
\begin{minipage}[b]{\linewidth}\raggedright
Comparative line
\end{minipage} & \begin{minipage}[b]{\linewidth}\raggedright
Candidate form
\end{minipage} & \begin{minipage}[b]{\linewidth}\raggedright
OE relation
\end{minipage} & \begin{minipage}[b]{\linewidth}\raggedright
Result
\end{minipage} \\
\midrule\noalign{}
\endhead
\bottomrule\noalign{}
\endlastfoot
Orel's citation &
{\emph{*xwalaz}}\index[iv]{03pgmc@\ivlangheader{Proto-Germanic}{}!xwalaz@\iventry{*xwalaz}{}}
& same stem notation as the modeled singular line & closest comparative
support for the derivational input \\
Kroonen's citation &
{\emph{*hwali-}}\index[iv]{03pgmc@\ivlangheader{Proto-Germanic}{}!hwali@\iventry{*hwali-}{}}
& same initial cluster, different stem label & important comparative
rival, but not the notation followed here \\
modeled input &
{\emph{*xwálaz}}\index[iv]{03pgmc@\ivlangheader{Proto-Germanic}{}!xwalaz@\iventry{*xwálaz}{}}
& regular output: {\emph{hwæl}} `whale' & Old English-facing input \\
plural control &
{\emph{hwalas}}\index[iv]{02oe@\ivlangheader{Old English}{West Saxon normalization unmarked}!hwalas@\iventry{hwalas}{}}
`whales' & attested open-syllable plural beside singular
{\emph{hwæl}}\index[iv]{02oe@\ivlangheader{Old English}{West Saxon normalization unmarked}!hwael@\iventry{hwæl}{}}
`whale' & confirms that the lexeme also preserves an \emph{a}-vocalism
branch \\
\end{longtable}
}

\section{\texorpdfstring{whine --- OE
\emph{hwīnan}}{whine --- OE hwīnan}}\label{whine-oe-hwux12bnan}

\index[iv]{02oe@\ivlangheader{Old English}{West Saxon normalization unmarked}!hwinan@\iventry{hwīnan}{}}
\index[iv]{03pgmc@\ivlangheader{Proto-Germanic}{}!wainojana@\iventry{*wainōjaną}{}}
\index[iv]{03pgmc@\ivlangheader{Proto-Germanic}{}!xwinana@\iventry{*xwī́naną}{}}

Derivation: citation reconstruction \emph{*wainōjaną}; form followed
here \emph{*xwī́naną} \textgreater{} \emph{hwīnan} (early analogy).

\subsection{Derivation trace}\label{derivation-trace-136}

Proto input: \emph{*xwī́naną}

\begingroup
\setlength{\fboxsep}{6pt}

\noindent\fbox{%
\begin{minipage}{0.97\linewidth}
\small
\begin{tabular}{@{}>{\raggedright\arraybackslash}p{0.359\linewidth}@{\hspace{0.030\linewidth}}>{\raggedright\arraybackslash}p{0.611\linewidth}@{}}
\begin{minipage}[t]{\linewidth}
\centering\textbf{Earlier Germanic changes}\par
\vspace{0.35em}
\raggedright
\centering\textbf{Northwest and West Germanic}\par
\raggedright
\vspace{0.2em}
\raggedright [no change]\par
\vspace{0.6em}
\centering\textbf{Early Anglo-Frisian}\par
\raggedright
\vspace{0.2em}
\raggedright [no change]\par
\end{minipage}
&
\begin{minipage}[t]{\linewidth}
\centering\textbf{Old English changes}\par
\vspace{0.35em}
\raggedright
\begin{tabularx}{\linewidth}{@{}>{\raggedright\arraybackslash}X@{\hspace{0.55em}}>{\raggedright\arraybackslash}p{0.14\linewidth}@{\hspace{0.25em}}}
OE Heavy Syllable Nasal Apocope & \emph{*xwī́nan} \\
\mbox{OE Secondary Nasalization} & \emph{*xwī́nąn} \\
\mbox{OE Weak Tail Reduction} & \emph{*xwī́nan} \\
\end{tabularx}
\end{minipage}
\\
\end{tabular}
\end{minipage}%
} \endgroup

Old English form: \emph{hwīnan}

\subsection{Reconstruction and comparative
evidence}\label{reconstruction-and-comparative-evidence-119}

The citation reconstruction preserved in the header belongs to the
lament-family verb seen in German
{\emph{weinen}}\index[iv]{17german@\ivlangheader{German}{}!weinen@\iventry{weinen}{}}
`whine' and Old English
{\emph{wānian}}\index[iv]{02oe@\ivlangheader{Old English}{West Saxon normalization unmarked}!wanian@\iventry{wānian}{}}
`wane'. Kroonen instead separates Old English
{\emph{hwīnan}}\index[iv]{02oe@\ivlangheader{Old English}{West Saxon normalization unmarked}!hwinan@\iventry{hwīnan}{}}
`whine' under \emph{*hwinan-} (\citeproc{ref-Kroonen2013}{Kroonen 2013,
267}). Orel likewise distinguishes strong
\Recon{xwinanan}\index[iv]{03pgmc@\ivlangheader{Proto-Germanic}{}!xwinanan@\iventry{*xwinanan}{}}
`whine' from weak
\Recon{wainōjanan}\index[iv]{03pgmc@\ivlangheader{Proto-Germanic}{}!wainojanan@\iventry{*wainōjanan}{}}
`whine' (\citeproc{ref-Orel2003}{Orel 2003, 201}). Ringe and Taylor make
the same split at the Northwest Germanic level, linking Old Norse
{\emph{hvina}}\index[iv]{09on@\ivlangheader{Old Norse}{}!hvina@\iventry{hvina}{}}
`whine' and Old English
{\emph{hwinan}}\index[iv]{02oe@\ivlangheader{Old English}{West Saxon normalization unmarked}!hwinan@\iventry{hwinan}{}}
`whine' to the same strong verb (\citeproc{ref-RingeTaylor2014}{Ringe
and Taylor 2014, 130}).

The two families also differ phonologically and morphologically. The
lament family has initial \emph{w-}, diphthongal \emph{ai}, and weak-II
morphology, whereas the verb behind Old English \emph{hwīnan} `whine'
has initial \emph{hw-/xw-}, long \emph{ī}, and strong-verb inflection.
The derivational input
\Recon{xwī́naną}\index[iv]{03pgmc@\ivlangheader{Proto-Germanic}{}!xwinana@\iventry{*xwī́naną}{}}
`whine' therefore represents a competing comparative identification
rather than a hidden cell of \Recon{wainōjaną} `whine'.

\subsection{Old English evidence}\label{old-english-evidence-120}

Clark Hall records
{\emph{hwinan}}\index[iv]{02oe@\ivlangheader{Old English}{West Saxon normalization unmarked}!hwinan@\iventry{hwinan}{}}
`whine' with the gloss `to hiss, whizz, whistle'
(\citeproc{ref-ClarkHall1960}{Clark Hall 1960, 171}). Seebold keeps the
verb among the strong verbs and notes that only a present-tense
attestation is directly preserved (\citeproc{ref-Seebold1970}{Seebold
1970, 280}).

The Old English form is normalized here as
{\emph{hwīnan}}\index[iv]{02oe@\ivlangheader{Old English}{West Saxon normalization unmarked}!hwinan@\iventry{hwīnan}{}}
`whine'. That normalization adds the usual vowel length marking to the
dictionary spelling
{\emph{hwinan}}\index[iv]{02oe@\ivlangheader{Old English}{West Saxon normalization unmarked}!hwinan@\iventry{hwinan}{}}
`whine'; it does not turn an unattested verb into a reconstructed one.

\subsection{Development to Old
English}\label{development-to-old-english-118}

Once the strong-verb input
\Recon{xwī́naną}\index[iv]{03pgmc@\ivlangheader{Proto-Germanic}{}!xwinana@\iventry{*xwī́naną}{}}
`whine' is selected, the path to Old English is straightforward. The
compact trace shows heavy-syllable nasal apocope, secondary
nasalization, and weak-tail reduction, after which the form surfaces as
{\emph{hwīnan}}\index[iv]{02oe@\ivlangheader{Old English}{West Saxon normalization unmarked}!hwinan@\iventry{hwīnan}{}}
`whine'.

No special paradigm maneuver is needed for this verb. The comparison is
between two different Germanic verb families: the Old English form
belongs with the strong verb \emph{*hwīnan-}, not with the weak lament
verb.

\subsection{Verb-family comparison}\label{verb-family-comparison}

The comparison below sets the relevant forms side by side. It separates
the competing comparative labels that stand behind the inherited Old
English forms.

{\def\LTcaptype{none} % do not increment counter
\begin{longtable}[]{@{}
  >{\raggedright\arraybackslash}p{(\linewidth - 8\tabcolsep) * \real{0.2000}}
  >{\raggedright\arraybackslash}p{(\linewidth - 8\tabcolsep) * \real{0.2000}}
  >{\raggedright\arraybackslash}p{(\linewidth - 8\tabcolsep) * \real{0.2000}}
  >{\raggedright\arraybackslash}p{(\linewidth - 8\tabcolsep) * \real{0.2000}}
  >{\raggedright\arraybackslash}p{(\linewidth - 8\tabcolsep) * \real{0.2000}}@{}}
\toprule\noalign{}
\begin{minipage}[b]{\linewidth}\raggedright
Verb family / interpretation
\end{minipage} & \begin{minipage}[b]{\linewidth}\raggedright
Candidate input
\end{minipage} & \begin{minipage}[b]{\linewidth}\raggedright
Old English outcome or comparison
\end{minipage} & \begin{minipage}[b]{\linewidth}\raggedright
OE comparison form
\end{minipage} & \begin{minipage}[b]{\linewidth}\raggedright
Result
\end{minipage} \\
\midrule\noalign{}
\endhead
\bottomrule\noalign{}
\endlastfoot
lament-family weak verb &
{\emph{*wainōjaną}}\index[iv]{03pgmc@\ivlangheader{Proto-Germanic}{}!wainojana@\iventry{*wainōjaną}{}}
& comparative continuation in OE
{\emph{wānian}}\index[iv]{02oe@\ivlangheader{Old English}{West Saxon normalization unmarked}!wanian@\iventry{wānian}{}}
`wane' &
{\emph{wānian}}\index[iv]{02oe@\ivlangheader{Old English}{West Saxon normalization unmarked}!wanian@\iventry{wānian}{}}
`wane' & competing citation reconstruction, but not the source of
\emph{hwīnan} \\
selected strong verb &
{\emph{*xwī́naną}}\index[iv]{03pgmc@\ivlangheader{Proto-Germanic}{}!xwinana@\iventry{*xwī́naną}{}}
& regular output: {\emph{hwīnan}} `whine' &
{\emph{hwīnan}}\index[iv]{02oe@\ivlangheader{Old English}{West Saxon normalization unmarked}!hwinan@\iventry{hwīnan}{}}
`whine' & exact match between derivational input and OE verb \\
comparative North Germanic cognate & Northwest Germanic strong verb
behind ON
{\emph{hvina}}\index[iv]{09on@\ivlangheader{Old Norse}{}!hvina@\iventry{hvina}{}}
`whine' / OE
{\emph{hwinan}}\index[iv]{02oe@\ivlangheader{Old English}{West Saxon normalization unmarked}!hwinan@\iventry{hwinan}{}}
`whine' & ON
{\emph{hvina}}\index[iv]{09on@\ivlangheader{Old Norse}{}!hvina@\iventry{hvina}{}}
`whine' / OE
{\emph{hwinan}}\index[iv]{02oe@\ivlangheader{Old English}{West Saxon normalization unmarked}!hwinan@\iventry{hwinan}{}}
`whine' &
{\emph{hwīnan}}\index[iv]{02oe@\ivlangheader{Old English}{West Saxon normalization unmarked}!hwinan@\iventry{hwīnan}{}}
`whine' & supports the strong-verb identification \\
\end{longtable}
}

\section{\texorpdfstring{withy --- OE
\emph{wīþiġ}}{withy --- OE wīþiġ}}\label{withy-oe-wux12buxfeiux121}

\index[iv]{02oe@\ivlangheader{Old English}{West Saxon normalization unmarked}!hunig@\iventry{hunig}{}}
\index[iv]{02oe@\ivlangheader{Old English}{West Saxon normalization unmarked}!withig@\iventry{wīþiġ}{}}
\index[iv]{03pgmc@\ivlangheader{Proto-Germanic}{}!waithiz@\iventry{*wáiθiz}{}}
\index[iv]{03pgmc@\ivlangheader{Proto-Germanic}{}!withaga@\iventry{*wī́θagą}{}}

Derivation: citation reconstruction \emph{*wáiθiz}; form followed here
\emph{*wī́θagą} \textgreater{} \emph{wīþiġ} (early analogy).

\subsection{Derivation trace}\label{derivation-trace-137}

Proto input: \emph{*wī́θagą}

\begingroup
\setlength{\fboxsep}{6pt}

\noindent\fbox{%
\begin{minipage}{0.97\linewidth}
\small
\begin{tabular}{@{}>{\raggedright\arraybackslash}p{0.339\linewidth}@{\hspace{0.030\linewidth}}>{\raggedright\arraybackslash}p{0.630\linewidth}@{}}
\begin{minipage}[t]{\linewidth}
\centering\textbf{Earlier Germanic changes}\par
\vspace{0.35em}
\raggedright
\centering\textbf{Northwest and West Germanic}\par
\raggedright
\vspace{0.2em}
\raggedright [no change]\par
\vspace{0.6em}
\centering\textbf{Early Anglo-Frisian}\par
\raggedright
\vspace{0.2em}
\raggedright [no change]\par
\end{minipage}
&
\begin{minipage}[t]{\linewidth}
\centering\textbf{Old English changes}\par
\vspace{0.35em}
\raggedright
\begin{tabularx}{\linewidth}{@{}>{\raggedright\arraybackslash}X@{\hspace{0.55em}}>{\raggedright\arraybackslash}p{0.14\linewidth}@{\hspace{0.25em}}}
\mbox{EAF Brightening Unstressed} & \emph{*wī́θægą} \\
OE Heavy Syllable Nasal Apocope & \emph{*wī́θæg} \\
\mbox{OE Velar Palatalization} & \emph{*wī́θæʝ} \\
\mbox{OE Unstressed AE Merger} & \emph{*wī́θeʝ} \\
\mbox{OE Late Unstressed Ag Suffix} & \emph{*wī́θiʝ} \\
\mbox{OE Palatal Fricative Merger} & \emph{*wī́θij} \\
\end{tabularx}
\end{minipage}
\\
\end{tabular}
\end{minipage}%
} \endgroup

Old English form: \emph{wīþiġ}

\subsection{Reconstruction and comparative
evidence}\label{reconstruction-and-comparative-evidence-120}

The comparative evidence groups the word with Germanic forms of the
\Recon{wīþja} `withy' (or \emph{*wīþjō-} variant) or \emph{*wiþ-} type
(\citeproc{ref-Orel2003}{Orel 2003, 503}). These forms establish the
cognate set but do not by themselves explain the Old English suffix of
\emph{wīþiġ} `withy'.

For Old English, the relevant point is the suffix history. Campbell's
account of OE \emph{-ig}, including forms such as \emph{hunig}, supports
an analysis in which the \emph{-iġ} of \emph{wīþiġ} `withy' continues a
derivational \emph{*-ag-} sequence rather than a heavy ja-stem
\emph{*-ij-} (\citeproc{ref-Campbell1959}{Campbell 1959, §§ 275,
376}). The derivational input
\Recon{wī́θagą}\index[iv]{03pgmc@\ivlangheader{Proto-Germanic}{}!withaga@\iventry{*wī́θagą}{}}
`withy' is thus a formation choice rather than a mere respelling of the
comparative headword.

\subsection{Old English evidence}\label{old-english-evidence-121}

Clark Hall records the noun as
{\emph{wiðig}}\index[iv]{02oe@\ivlangheader{Old English}{West Saxon normalization unmarked}!withig@\iventry{wiðig}{}}
`withy' (\citeproc{ref-ClarkHall1960}{Clark Hall 1960, 358}). The form
used here,
{\emph{wīþiġ}}\index[iv]{02oe@\ivlangheader{Old English}{West Saxon normalization unmarked}!withig@\iventry{wīþiġ}{}}
`withy', is a normalized Old English spelling with macrons and palatal ġ
marked explicitly.

The relevant comparison form is therefore not a reconstructed dictionary
convenience but an established Old English noun. What requires
explanation is why the selected Proto-Germanic input is
\Recon{wī́θagą}\index[iv]{03pgmc@\ivlangheader{Proto-Germanic}{}!withaga@\iventry{*wī́θagą}{}}
`withy' rather than a comparative headword of the \Recon{wīþja} `withy'
type.

\subsection{Development to Old
English}\label{development-to-old-english-119}

From
\Recon{wī́θagą}\index[iv]{03pgmc@\ivlangheader{Proto-Germanic}{}!withaga@\iventry{*wī́θagą}{}}
`withy', the retained unstressed-fronting component gives the suffixal
front vowel. Later raising recognizes the singleton palatal fricative,
distinct from gg/ng stops; this late suffix history is not prehistoric
i-mutation (\citeproc{ref-RingeTaylor2014}{Ringe and Taylor 2014,
334--35}; \citeproc{ref-Fulk2018}{Fulk 2018, 130--32}). Postmutation
merger and native realization supply the final \emph{ġ}, reaching
{\emph{wīþiġ}} `withy'.

This derivation is regular for the form compared here. The central claim
of the entry is therefore morphological: Old English \emph{wīþiġ}
`withy' belongs with an \emph{*-ag-} derivative, whereas the comparative
\Recon{wīþja} `withy' label belongs to a different way of presenting the
cognate family.

\subsection{Formation comparison}\label{formation-comparison-8}

The comparison below sets the relevant formations side by side. It
distinguishes the comparative headword from the Old English-facing
formation that actually yields the attested noun.

{\def\LTcaptype{none} % do not increment counter
\begin{longtable}[]{@{}
  >{\raggedright\arraybackslash}p{(\linewidth - 8\tabcolsep) * \real{0.2000}}
  >{\raggedright\arraybackslash}p{(\linewidth - 8\tabcolsep) * \real{0.2000}}
  >{\raggedright\arraybackslash}p{(\linewidth - 8\tabcolsep) * \real{0.2000}}
  >{\raggedright\arraybackslash}p{(\linewidth - 8\tabcolsep) * \real{0.2000}}
  >{\raggedright\arraybackslash}p{(\linewidth - 8\tabcolsep) * \real{0.2000}}@{}}
\toprule\noalign{}
\begin{minipage}[b]{\linewidth}\raggedright
Formation
\end{minipage} & \begin{minipage}[b]{\linewidth}\raggedright
Candidate input
\end{minipage} & \begin{minipage}[b]{\linewidth}\raggedright
Expected or documented OE outcome
\end{minipage} & \begin{minipage}[b]{\linewidth}\raggedright
OE comparison form
\end{minipage} & \begin{minipage}[b]{\linewidth}\raggedright
Result
\end{minipage} \\
\midrule\noalign{}
\endhead
\bottomrule\noalign{}
\endlastfoot
comparative family label &
{\emph{*wáiθiz}}\index[iv]{03pgmc@\ivlangheader{Proto-Germanic}{}!waithiz@\iventry{*wáiθiz}{}}
& broader cognate-set headword & OE family context & useful lexeme
label, but not the direct source of \emph{wīþiġ} \\
heavy ja-stem analysis & \Recon{wīþja} `withy' type &
Campbell/Adamczyk-style heavy ja-stem \emph{-e} / zero outcome &
{\emph{wīþiġ}}\index[iv]{02oe@\ivlangheader{Old English}{West Saxon normalization unmarked}!withig@\iventry{wīþiġ}{}}
`withy' & does not account cleanly for the OE suffix \\
\emph{*-ag-} derivative followed here &
{\emph{*wī́θagą}}\index[iv]{03pgmc@\ivlangheader{Proto-Germanic}{}!withaga@\iventry{*wī́θagą}{}}
& regular output: {\emph{wīþiġ}} `withy' &
{\emph{wīþiġ}}\index[iv]{02oe@\ivlangheader{Old English}{West Saxon normalization unmarked}!withig@\iventry{wīþiġ}{}}
`withy' & exact match between formation and target \\
\end{longtable}
}

\section{\texorpdfstring{youth --- OE
\emph{ġeoguþ}}{youth --- OE ġeoguþ}}\label{youth-oe-ux121eoguuxfe}

\index[iv]{02oe@\ivlangheader{Old English}{West Saxon normalization unmarked}!geoguth@\iventry{ġeoguþ}{}}
\index[iv]{03pgmc@\ivlangheader{Proto-Germanic}{}!jugunthiz@\iventry{*júgunθiz}{}}
\index[iv]{05pwgmc@\ivlangheader{Proto-West Germanic}{}!jugunth@\iventry{*júgunθ}{}}

Derivation: citation reconstruction \emph{*júgunθiz}; form followed here
\emph{*júgunθ} \textgreater{} \emph{ġeoguþ} (early analogy).

\subsection{Derivation trace}\label{derivation-trace-138}

Proto input: \emph{*júgunθ}

\begingroup
\setlength{\fboxsep}{6pt}

\noindent\fbox{%
\begin{minipage}{0.97\linewidth}
\small
\begin{tabular}{@{}>{\raggedright\arraybackslash}p{0.485\linewidth}@{\hspace{0.030\linewidth}}>{\raggedright\arraybackslash}p{0.485\linewidth}@{}}
\begin{minipage}[t]{\linewidth}
\centering\textbf{Earlier Germanic changes}\par
\vspace{0.35em}
\raggedright
\centering\textbf{Northwest and West Germanic}\par
\raggedright
\vspace{0.2em}
\begin{tabularx}{\linewidth}{@{}>{\raggedright\arraybackslash}X@{\hspace{0.55em}}>{\raggedright\arraybackslash}p{0.14\linewidth}@{\hspace{0.25em}}}
EAF Nasal Spirant Lengthening & \emph{*júgūnθ} \\
\mbox{EAF Nasal Spirant Loss} & \emph{*júgūθ} \\
\end{tabularx}
\vspace{0.6em}
\centering\textbf{Early Anglo-Frisian}\par
\raggedright
\vspace{0.2em}
\raggedright [no change]\par
\end{minipage}
&
\begin{minipage}[t]{\linewidth}
\centering\textbf{Old English changes}\par
\vspace{0.35em}
\raggedright
\begin{tabularx}{\linewidth}{@{}>{\raggedright\arraybackslash}X@{\hspace{0.55em}}>{\raggedright\arraybackslash}p{0.14\linewidth}@{\hspace{0.25em}}}
OE Unstressed Long Vowel Shortening & \emph{*júguθ} \\
\end{tabularx}
\end{minipage}
\\
\end{tabular}
\end{minipage}%
} \endgroup

Old English form: \emph{ġeoguþ}

\subsection{Reconstruction and comparative
evidence}\label{reconstruction-and-comparative-evidence-121}

The wider etymological tradition reconstructs an earlier form of the
word as
{\emph{*ju(w)unþi-}}\index[iv]{03pgmc@\ivlangheader{Proto-Germanic}{}!juwunthi@\iventry{*ju(w)unþi-}{}}.
The comparative label
\Recon{júgunθiz}\index[iv]{03pgmc@\ivlangheader{Proto-Germanic}{}!jugunthiz@\iventry{*júgunθiz}{}}
`youth' already stands at a later Germanic stage with \emph{g}, and the
derivational input
\Recon{júgunθ}\index[iv]{05pwgmc@\ivlangheader{Proto-West Germanic}{}!jugunth@\iventry{*júgunθ}{}}
`youth' is later again: it represents the form after final \emph{-i} has
been lost.

Ringe and Taylor explicitly give the sequence \Recon{jugunþi} `youth'
\textgreater{} \Recon{juguþ} `youth' \textgreater{} OE geoguþ
\textasciitilde{} iuguþ (\citeproc{ref-RingeTaylor2014}{Ringe and Taylor
2014, 141}). The derivational input therefore differs from the broader
comparative headword because the Old English development must begin
after early loss of final \emph{-i}.

\subsection{Old English evidence}\label{old-english-evidence-122}

The Old English noun is attested with varying spellings. Ringe and
Taylor cite
{\emph{geoguþ}}\index[iv]{02oe@\ivlangheader{Old English}{West Saxon normalization unmarked}!geoguth@\iventry{geoguþ}{}}
`youth' \textasciitilde{}
{\emph{iuguþ}}\index[iv]{02oe@\ivlangheader{Old English}{West Saxon normalization unmarked}!iuguth@\iventry{iuguþ}{}}
`youth' (\citeproc{ref-RingeTaylor2014}{Ringe and Taylor 2014, 141}).
The form is normalized here as
{\emph{ġeoguþ}}\index[iv]{02oe@\ivlangheader{Old English}{West Saxon normalization unmarked}!geoguth@\iventry{ġeoguþ}{}}
`youth': the initial palatal is written with \emph{ġ}, and the attested
spelling variation is treated as orthographic rather than lexical.

Nothing in the source stack suggests that a different paradigm cell
should be chosen. The relevant Old English comparison form is the noun
{\emph{ġeoguþ}}\index[iv]{02oe@\ivlangheader{Old English}{West Saxon normalization unmarked}!geoguth@\iventry{ġeoguþ}{}}
`youth' itself.

\subsection{Development to Old
English}\label{development-to-old-english-120}

The decisive early step is the loss of final \emph{-i} before the Old
English umlaut stage. If that high vowel remained, the word would
develop an over-umlauted \emph{y}-type vowel instead of the attested
form (\citeproc{ref-RingeTaylor2014}{Ringe and Taylor 2014, 141}).

From the derivational input
\Recon{júgunθ}\index[iv]{05pwgmc@\ivlangheader{Proto-West Germanic}{}!jugunth@\iventry{*júgunθ}{}}
`youth', the later development is regular: nasal-spirant lengthening and
loss give \Recon{júgūθ} `youth' (\citeproc{ref-Fulk2018}{Fulk 2018,
109}); unstressed long-vowel shortening then produces \Recon{júguθ}
`youth'. The root vowel remains phonological \emph{u} through
medial-vowel lowering; late West Saxon spelling, not early palatal
fronting, yields {\emph{ġeoguþ}} `youth'
(\citeproc{ref-RingeTaylor2014}{Ringe and Taylor 2014, 5};
\citeproc{ref-SieversBrunner1965}{Brunner 1965, 64--65}, §92.1).
Campbell preserves \emph{u} after accented \emph{u} in forms such as
{\emph{duguþ}}\index[iv]{02oe@\ivlangheader{Old English}{West Saxon normalization unmarked}!duguth@\iventry{duguþ}{}}
`troop' and
{\emph{munuc}}\index[iv]{02oe@\ivlangheader{Old English}{West Saxon normalization unmarked}!munuc@\iventry{munuc}{}}
`monk' (\citeproc{ref-Campbell1959}{Campbell 1959, 156}, §373). Brunner
likewise cites
{\emph{iuzuð}}\index[iv]{02oe@\ivlangheader{Old English}{West Saxon normalization unmarked}!iuzuth@\iventry{iuzuð}{}}
`youth' and
{\emph{munuc}}\index[iv]{02oe@\ivlangheader{Old English}{West Saxon normalization unmarked}!munuc@\iventry{munuc}{}}
`monk' under the same development
(\citeproc{ref-SieversBrunner1965}{Brunner 1965, § 150.3}).

\subsection{Stage comparison}\label{stage-comparison-4}

The comparison below sets the relevant forms side by side. It separates
the broader comparative headword from the later stages relevant to the
Old English noun.

{\def\LTcaptype{none} % do not increment counter
\begin{longtable}[]{@{}
  >{\raggedright\arraybackslash}p{(\linewidth - 6\tabcolsep) * \real{0.2500}}
  >{\raggedright\arraybackslash}p{(\linewidth - 6\tabcolsep) * \real{0.2500}}
  >{\raggedright\arraybackslash}p{(\linewidth - 6\tabcolsep) * \real{0.2500}}
  >{\raggedright\arraybackslash}p{(\linewidth - 6\tabcolsep) * \real{0.2500}}@{}}
\toprule\noalign{}
\begin{minipage}[b]{\linewidth}\raggedright
Stage / interpretation
\end{minipage} & \begin{minipage}[b]{\linewidth}\raggedright
Candidate form
\end{minipage} & \begin{minipage}[b]{\linewidth}\raggedright
Old English outcome or comparison
\end{minipage} & \begin{minipage}[b]{\linewidth}\raggedright
Relevance to this entry
\end{minipage} \\
\midrule\noalign{}
\endhead
\bottomrule\noalign{}
\endlastfoot
earlier etymological headword &
{\emph{*ju(w)unþi-}}\index[iv]{03pgmc@\ivlangheader{Proto-Germanic}{}!juwunthi@\iventry{*ju(w)unþi-}{}}
& comparative family background & older comparative reconstruction of
the lexeme \\
later g-bearing comparative label &
{\emph{*júgunθiz}}\index[iv]{03pgmc@\ivlangheader{Proto-Germanic}{}!jugunthiz@\iventry{*júgunθiz}{}}
& citation reconstruction / lexeme label & preserves the later Germanic
stage behind the selected entry \\
Old English-facing input &
{\emph{*júgunθ}}\index[iv]{05pwgmc@\ivlangheader{Proto-West Germanic}{}!jugunth@\iventry{*júgunθ}{}}
& regular output: {\emph{ġeoguþ}} & exact match for the Old English form
used here \\
full \emph{-i} stage retained too long &
{\emph{*jugunþi}}\index[iv]{03pgmc@\ivlangheader{Proto-Germanic}{}!jugunthi@\iventry{*jugunþi}{}}
& expected over-umlauted \emph{y}-type result & negative control showing
why early \emph{-i} loss must precede the OE umlaut stage \\
\end{longtable}
}

\clearpage

\chapter{Late analogy and paradigm-cell
selection}\label{late-analogy-and-paradigm-cell-selection}

These Old English forms continue a particular paradigm cell or a later
analogical remodeling, not the unaltered citation form. The phonology
may be regular once the proper morphological history has been
identified.

\section{\texorpdfstring{ban --- OE
\emph{bannes}}{ban --- OE bannes}}\label{ban-oe-bannes}

\index[iv]{02oe@\ivlangheader{Old English}{West Saxon normalization unmarked}!bannes@\iventry{bannes}{}}
\index[iv]{03pgmc@\ivlangheader{Proto-Germanic}{}!banna@\iventry{*bánną}{}}
\index[iv]{03pgmc@\ivlangheader{Proto-Germanic}{}!bannas@\iventry{*bánnas}{}}

Derivation: citation reconstruction \emph{*bánną}; form followed here
\emph{*bánnas} \textgreater{} \emph{bannes} (late analogy).

\subsection{Derivation trace}\label{derivation-trace-139}

Proto input: \emph{*bánnas}

\begingroup
\setlength{\fboxsep}{6pt}

\noindent\fbox{%
\begin{minipage}{0.97\linewidth}
\small
\begin{tabular}{@{}>{\raggedright\arraybackslash}p{0.359\linewidth}@{\hspace{0.030\linewidth}}>{\raggedright\arraybackslash}p{0.611\linewidth}@{}}
\begin{minipage}[t]{\linewidth}
\centering\textbf{Earlier Germanic changes}\par
\vspace{0.35em}
\raggedright
\centering\textbf{Northwest and West Germanic}\par
\raggedright
\vspace{0.2em}
\raggedright [no change]\par
\vspace{0.6em}
\centering\textbf{Early Anglo-Frisian}\par
\raggedright
\vspace{0.2em}
\raggedright [no change]\par
\end{minipage}
&
\begin{minipage}[t]{\linewidth}
\centering\textbf{Old English changes}\par
\vspace{0.35em}
\raggedright
\begin{tabularx}{\linewidth}{@{}>{\raggedright\arraybackslash}X@{\hspace{0.55em}}>{\raggedright\arraybackslash}p{0.14\linewidth}@{\hspace{0.25em}}}
\mbox{EAF Brightening Unstressed} & \emph{*bánnæs} \\
\mbox{OE Unstressed AE Merger} & \emph{*bánnes} \\
\end{tabularx}
\end{minipage}
\\
\end{tabular}
\end{minipage}%
} \endgroup

Old English form: \emph{bannes}

\subsection{Reconstruction and comparative
evidence}\label{reconstruction-and-comparative-evidence-122}

Orel cites a bann-noun under
\Recon{bannan}\index[iv]{03pgmc@\ivlangheader{Proto-Germanic}{}!bannan@\iventry{*bannan}{}}
`ban', while Seebold distinguishes bann-stems of both masculine and
neuter type and gives Old English
{\emph{gebann}}\index[iv]{02oe@\ivlangheader{Old English}{West Saxon normalization unmarked}!gebann@\iventry{gebann}{}}
`ban' as the noun reflex (\citeproc{ref-Orel2003}{Orel 2003, 35};
\citeproc{ref-Seebold1970}{Seebold 1970, 89}). The citation
reconstruction
\Recon{bánną}\index[iv]{03pgmc@\ivlangheader{Proto-Germanic}{}!banna@\iventry{*bánną}{}}
`ban' names the lexeme, but the comparison here turns on the genitive
singular
\Recon{bánnas}\index[iv]{03pgmc@\ivlangheader{Proto-Germanic}{}!bannas@\iventry{*bánnas}{}}
`ban'.

The analysis therefore depends on medial, not final, gemination.

\subsection{Old English evidence}\label{old-english-evidence-123}

Old English lexicographic evidence securely supports the noun itself.
Bosworth-Toller records the noun under nominative
{\emph{ge-bann}}\index[iv]{02oe@\ivlangheader{Old English}{West Saxon normalization unmarked}!gebann@\iventry{ge-bann}{}}
`ban', with oblique usage such as
{\emph{gebanne}}\index[iv]{02oe@\ivlangheader{Old English}{West Saxon normalization unmarked}!gebanne@\iventry{gebanne}{}}
`ban' (\citeproc{ref-BosworthToller1898}{Bosworth and Toller 1898,
303}). The exact unprefixed genitive
{\emph{bannes}}\index[iv]{02oe@\ivlangheader{Old English}{West Saxon normalization unmarked}!bannes@\iventry{bannes}{}}
`ban' is less directly cited in the dictionaries, so it is best treated
here as the regular genitive form used for comparison rather than as a
dictionary headword.

\subsection{Development to Old
English}\label{development-to-old-english-121}

From
\Recon{bánnas}\index[iv]{03pgmc@\ivlangheader{Proto-Germanic}{}!bannas@\iventry{*bánnas}{}}
`ban', the geminate remains medial before the case ending and the
unstressed vowel develops regularly to give
{\emph{bannes}}\index[iv]{02oe@\ivlangheader{Old English}{West Saxon normalization unmarked}!bannes@\iventry{bannes}{}}
`ban'. The paradigm comparison therefore sets the genitive against
nominative
{\emph{ban}}\index[iv]{02oe@\ivlangheader{Old English}{West Saxon normalization unmarked}!ban@\iventry{ban}{}}
`ban', the ordinary nominative form of the same noun, rather than
against a directly cited genitive headword.

\subsection{Paradigm comparison}\label{paradigm-comparison-1}

The comparison below sets the relevant forms side by side. It shows why
the genitive singular is the conservative cell used for the entry.

{\def\LTcaptype{none} % do not increment counter
\begin{longtable}[]{@{}
  >{\raggedright\arraybackslash}p{(\linewidth - 8\tabcolsep) * \real{0.2000}}
  >{\raggedright\arraybackslash}p{(\linewidth - 8\tabcolsep) * \real{0.2000}}
  >{\raggedright\arraybackslash}p{(\linewidth - 8\tabcolsep) * \real{0.2000}}
  >{\raggedright\arraybackslash}p{(\linewidth - 8\tabcolsep) * \real{0.2000}}
  >{\raggedright\arraybackslash}p{(\linewidth - 8\tabcolsep) * \real{0.2000}}@{}}
\toprule\noalign{}
\begin{minipage}[b]{\linewidth}\raggedright
PGmc cell / interpretation
\end{minipage} & \begin{minipage}[b]{\linewidth}\raggedright
Candidate input
\end{minipage} & \begin{minipage}[b]{\linewidth}\raggedright
OE output or comparison
\end{minipage} & \begin{minipage}[b]{\linewidth}\raggedright
OE comparison form
\end{minipage} & \begin{minipage}[b]{\linewidth}\raggedright
Result
\end{minipage} \\
\midrule\noalign{}
\endhead
\bottomrule\noalign{}
\endlastfoot
citation nominative singular &
{\emph{*bánną}}\index[iv]{03pgmc@\ivlangheader{Proto-Germanic}{}!banna@\iventry{*bánną}{}}
& regular output: {\emph{ban}} & ban & regular nominative outcome, but
not the Old English form here \\
genitive singular &
{\emph{*bánnas}}\index[iv]{03pgmc@\ivlangheader{Proto-Germanic}{}!bannas@\iventry{*bánnas}{}}
& regular output: \emph{bannes} &
{\emph{bannes}}\index[iv]{02oe@\ivlangheader{Old English}{West Saxon normalization unmarked}!bannes@\iventry{bannes}{}}
& direct match for the conservative genitive \\
\end{longtable}
}

\section{\texorpdfstring{berry --- OE
\emph{berġes}}{berry --- OE berġes}}\label{berry-oe-berux121es}

\index[iv]{02oe@\ivlangheader{Old English}{West Saxon normalization unmarked}!berges@\iventry{berġes}{}}
\index[iv]{03pgmc@\ivlangheader{Proto-Germanic}{}!bazja@\iventry{*bázją}{}}
\index[iv]{03pgmc@\ivlangheader{Proto-Germanic}{}!bazjas@\iventry{*bázjas}{}}

Derivation: citation reconstruction \emph{*bázją}; form followed here
\emph{*bázjas} \textgreater{} \emph{berġes} (late analogy).

\subsection{Derivation trace}\label{derivation-trace-140}

Proto input: \emph{*bázjas}

\begingroup
\setlength{\fboxsep}{6pt}

\noindent\fbox{%
\begin{minipage}{0.97\linewidth}
\small
\begin{tabular}{@{}>{\raggedright\arraybackslash}p{0.359\linewidth}@{\hspace{0.030\linewidth}}>{\raggedright\arraybackslash}p{0.611\linewidth}@{}}
\begin{minipage}[t]{\linewidth}
\centering\textbf{Earlier Germanic changes}\par
\vspace{0.35em}
\raggedright
\centering\textbf{Northwest and West Germanic}\par
\raggedright
\vspace{0.2em}
\begin{tabularx}{\linewidth}{@{}>{\raggedright\arraybackslash}X@{\hspace{0.55em}}>{\raggedright\arraybackslash}p{0.24\linewidth}@{\hspace{0.25em}}}
\mbox{EAF Rhotacism} & \emph{*bárjas} \\
\end{tabularx}
\vspace{0.6em}
\centering\textbf{Early Anglo-Frisian}\par
\raggedright
\vspace{0.2em}
\raggedright [no change]\par
\end{minipage}
&
\begin{minipage}[t]{\linewidth}
\centering\textbf{Old English changes}\par
\vspace{0.35em}
\raggedright
\begin{tabularx}{\linewidth}{@{}>{\raggedright\arraybackslash}X@{\hspace{0.55em}}>{\raggedright\arraybackslash}p{0.14\linewidth}@{\hspace{0.25em}}}
\mbox{EAF Brightening Unstressed} & \emph{*bárjæs} \\
\mbox{EAF Brightening} & \emph{*bærjæs} \\
\mbox{OE I Umlaut} & \emph{*berjæs} \\
\mbox{OE Unstressed AE Merger} & \emph{*berjes} \\
\end{tabularx}
\end{minipage}
\\
\end{tabular}
\end{minipage}%
} \endgroup

Old English form: \emph{berġes}

\subsection{Reconstruction and comparative
evidence}\label{reconstruction-and-comparative-evidence-123}

Kroonen reconstructs the berry noun as
{\emph{*basja-}}\index[iv]{03pgmc@\ivlangheader{Proto-Germanic}{}!basja@\iventry{*basja-}{}}
\textasciitilde{}
{\emph{*bazja-}}\index[iv]{03pgmc@\ivlangheader{Proto-Germanic}{}!bazja@\iventry{*bazja-}{}}
(\citeproc{ref-Kroonen2013}{Kroonen 2013, 54}). The derivational input
\Recon{bázjas}\index[iv]{03pgmc@\ivlangheader{Proto-Germanic}{}!bazjas@\iventry{*bázjas}{}}
`berry' is therefore not a rival lexeme headword, but a specific
genitive singular cell drawn from that paradigm.

The relevant point is that \emph{*rj} did not geminate in Proto-West
Germanic. Ringe and Taylor's
{\emph{here}}\index[iv]{02oe@\ivlangheader{Old English}{West Saxon normalization unmarked}!here@\iventry{here}{}}
`army',
{\emph{herges}}\index[iv]{02oe@\ivlangheader{Old English}{West Saxon normalization unmarked}!herges@\iventry{herges}{}}
`army's' comparison shows the same \emph{rj} environment in an Old
English paradigm without any hidden gemination repair
(\citeproc{ref-RingeTaylor2014}{Ringe and Taylor 2014, 181}).

\subsection{Old English evidence}\label{old-english-evidence-124}

Campbell cites feminine
{\emph{berige}}\index[iv]{02oe@\ivlangheader{Old English}{West Saxon normalization unmarked}!berige@\iventry{berige}{}}
`berry' and notes that \emph{-j-} is retained after \emph{r} in this
type (\citeproc{ref-Campbell1959}{Campbell 1959, 250}). The reviewed
evidence therefore supports the citation form more directly than the
exact genitive
{\emph{berġes}}\index[iv]{02oe@\ivlangheader{Old English}{West Saxon normalization unmarked}!berges@\iventry{berġes}{}}
`berry', which is best read here as the regular genitive comparison form
rather than as a dictionary headword.

\subsection{Development to Old
English}\label{development-to-old-english-122}

Citation
\Recon{bázją}\index[iv]{03pgmc@\ivlangheader{Proto-Germanic}{}!bazja@\iventry{*bázją}{}}
`berry' gives {\emph{bere}} `berry (variant)', not the Old English form
here. The genitive singular
\Recon{bázjas}\index[iv]{03pgmc@\ivlangheader{Proto-Germanic}{}!bazjas@\iventry{*bázjas}{}}
`berry', however, gives
{\emph{berġes}}\index[iv]{02oe@\ivlangheader{Old English}{West Saxon normalization unmarked}!berges@\iventry{berġes}{}}
`berry', with medial \emph{-rġ-} preserved in the same way that Ringe
and Taylor cite
{\emph{herges}}\index[iv]{02oe@\ivlangheader{Old English}{West Saxon normalization unmarked}!herges@\iventry{herges}{}}
`army's' beside
{\emph{here}}\index[iv]{02oe@\ivlangheader{Old English}{West Saxon normalization unmarked}!here@\iventry{here}{}}
`army' (\citeproc{ref-RingeTaylor2014}{Ringe and Taylor 2014, 181}).
This points to paradigm choice rather than to an extra phonological
rule.

\subsection{Paradigm comparison}\label{paradigm-comparison-2}

The comparison below sets the relevant forms side by side. It shows the
contrast between the citation form and the genitive singular cell.

{\def\LTcaptype{none} % do not increment counter
\begin{longtable}[]{@{}
  >{\raggedright\arraybackslash}p{(\linewidth - 8\tabcolsep) * \real{0.2000}}
  >{\raggedright\arraybackslash}p{(\linewidth - 8\tabcolsep) * \real{0.2000}}
  >{\raggedright\arraybackslash}p{(\linewidth - 8\tabcolsep) * \real{0.2000}}
  >{\raggedright\arraybackslash}p{(\linewidth - 8\tabcolsep) * \real{0.2000}}
  >{\raggedright\arraybackslash}p{(\linewidth - 8\tabcolsep) * \real{0.2000}}@{}}
\toprule\noalign{}
\begin{minipage}[b]{\linewidth}\raggedright
PGmc cell / interpretation
\end{minipage} & \begin{minipage}[b]{\linewidth}\raggedright
Candidate input
\end{minipage} & \begin{minipage}[b]{\linewidth}\raggedright
OE output or comparison
\end{minipage} & \begin{minipage}[b]{\linewidth}\raggedright
OE comparison form
\end{minipage} & \begin{minipage}[b]{\linewidth}\raggedright
Result
\end{minipage} \\
\midrule\noalign{}
\endhead
\bottomrule\noalign{}
\endlastfoot
citation nominative singular &
{\emph{*bázją}}\index[iv]{03pgmc@\ivlangheader{Proto-Germanic}{}!bazja@\iventry{*bázją}{}}
& regular output: {\emph{bere}} &
{\emph{berige}}\index[iv]{02oe@\ivlangheader{Old English}{West Saxon normalization unmarked}!berige@\iventry{berige}{}}
/
{\emph{berġe}}\index[iv]{02oe@\ivlangheader{Old English}{West Saxon normalization unmarked}!berge@\iventry{berġe}{}}
& useful citation-form background, but not the Old English form here \\
genitive singular &
{\emph{*bázjas}}\index[iv]{03pgmc@\ivlangheader{Proto-Germanic}{}!bazjas@\iventry{*bázjas}{}}
& regular output: \emph{berġes} &
{\emph{berġes}}\index[iv]{02oe@\ivlangheader{Old English}{West Saxon normalization unmarked}!berges@\iventry{berġes}{}}
& exact match for the conservative cell \\
\end{longtable}
}

\section{\texorpdfstring{bore (3sg) --- OE
\emph{boraþ}}{bore (3sg) --- OE boraþ}}\label{bore-3sg-oe-borauxfe}

\index[iv]{02oe@\ivlangheader{Old English}{West Saxon normalization unmarked}!borath@\iventry{boraþ}{}}
\index[iv]{03pgmc@\ivlangheader{Proto-Germanic}{}!burona@\iventry{*burōną}{}}
\index[iv]{03pgmc@\ivlangheader{Proto-Germanic}{}!burothi@\iventry{*búrōθi}{}}

Derivation: citation reconstruction \emph{*burōną}; form followed here
\emph{*búrōθi} \textgreater{} \emph{boraþ} (late analogy).

\subsection{Derivation trace}\label{derivation-trace-141}

Proto input: \emph{*búrōθi}

\begingroup
\setlength{\fboxsep}{6pt}

\noindent\fbox{%
\begin{minipage}{0.97\linewidth}
\small
\begin{tabular}{@{}>{\raggedright\arraybackslash}p{0.485\linewidth}@{\hspace{0.030\linewidth}}>{\raggedright\arraybackslash}p{0.485\linewidth}@{}}
\begin{minipage}[t]{\linewidth}
\centering\textbf{Earlier Germanic changes}\par
\vspace{0.35em}
\raggedright
\centering\textbf{Northwest and West Germanic}\par
\raggedright
\vspace{0.2em}
\begin{tabularx}{\linewidth}{@{}>{\raggedright\arraybackslash}X@{\hspace{0.55em}}>{\raggedright\arraybackslash}p{0.14\linewidth}@{\hspace{0.25em}}}
\mbox{PWGmc Early I Apocope} & \emph{*búrōθ} \\
\mbox{PNWGmc U Lowering} & \emph{*bórōθ} \\
\end{tabularx}
\vspace{0.6em}
\centering\textbf{Early Anglo-Frisian}\par
\raggedright
\vspace{0.2em}
\raggedright [no change]\par
\end{minipage}
&
\begin{minipage}[t]{\linewidth}
\centering\textbf{Old English changes}\par
\vspace{0.35em}
\raggedright
\begin{tabularx}{\linewidth}{@{}>{\raggedright\arraybackslash}X@{\hspace{0.55em}}>{\raggedright\arraybackslash}p{0.14\linewidth}@{\hspace{0.25em}}}
\mbox{OE Late O Shortening} & \emph{*bóroθ} \\
OE Final Unstressed O Lowering & \emph{*bóraθ} \\
\end{tabularx}
\end{minipage}
\\
\end{tabular}
\end{minipage}%
} \endgroup

Old English form: \emph{boraþ}

\subsection{Reconstruction and comparative
evidence}\label{reconstruction-and-comparative-evidence-124}

Kroonen reconstructs the bore-family verb as \emph{*burojan-} and cites
Old English
{\emph{borian}}\index[iv]{02oe@\ivlangheader{Old English}{West Saxon normalization unmarked}!borian@\iventry{borian}{}}
`bore' among its reflexes (\citeproc{ref-Kroonen2013}{Kroonen 2013,
85}). The form compared here isolates the finite 3sg cell \Recon{búrōθi}
`bores' rather than the infinitive.

Campbell's class-II pattern lufas, \emph{-aþ} (\textless{} \emph{-ōsi},
\emph{-ōþi)} and Ringe and Taylor's account of stable class-II \emph{a}
in finite forms make this 3sg cell the relevant comparison form for the
ending (\citeproc{ref-Campbell1959}{Campbell 1959, § 356.4};
\citeproc{ref-RingeTaylor2014}{Ringe and Taylor 2014, 80}).

\subsection{Old English evidence}\label{old-english-evidence-125}

Clark Hall lemmatizes the verb as
{\emph{borian}}\index[iv]{02oe@\ivlangheader{Old English}{West Saxon normalization unmarked}!borian@\iventry{borian}{}}
`bore'; the imperative singular \emph{bora} and present-third-singular
{\emph{boraþ}}\index[iv]{02oe@\ivlangheader{Old English}{West Saxon normalization unmarked}!borath@\iventry{boraþ}{}}
`bores' are the relevant comparison forms
(\citeproc{ref-ClarkHall1960}{Clark Hall 1960, 48}).

Campbell's dialect survey allows secondary \emph{-e-} and \emph{-o-}
spellings in 2sg and 3sg class-II forms, but the basic ending remains
\emph{-aþ} (\citeproc{ref-Campbell1959}{Campbell 1959, § 757}). The
regular comparison form for this non-\emph{j} 3sg cell is
{\emph{boraþ}}\index[iv]{02oe@\ivlangheader{Old English}{West Saxon normalization unmarked}!borath@\iventry{boraþ}{}}
`bores'.

\subsection{Development to Old
English}\label{development-to-old-english-123}

From \Recon{búrōθi} `bores', loss of final \emph{-i} first gives
\Recon{búrōθ} `bores'. Northwest Germanic lowering then produces
\Recon{bórōθ} `bores', and late shortening of unstressed \emph{ō} yields
\Recon{bóraθ} `bores'.
{\emph{Boraþ}}\index[iv]{02oe@\ivlangheader{Old English}{West Saxon normalization unmarked}!borath@\iventry{Boraþ}{}}
`bores' is therefore the regular comparison form for this non-\emph{j}
3sg cell.

Campbell's class-II ending evidence and Ringe and Taylor's discussion of
stable \emph{a} in finite non-\emph{j} cells support this sequence
(\citeproc{ref-Campbell1959}{Campbell 1959, § 356.4};
\citeproc{ref-RingeTaylor2014}{Ringe and Taylor 2014, 80}).

\subsection{Paradigm comparison}\label{paradigm-comparison-3}

The comparison below sets the relevant forms side by side. It
distinguishes the selected 3sg cell from the bore-family lemma and from
the companion imperative form.

{\def\LTcaptype{none} % do not increment counter
\begin{longtable}[]{@{}
  >{\raggedright\arraybackslash}p{(\linewidth - 8\tabcolsep) * \real{0.2000}}
  >{\raggedright\arraybackslash}p{(\linewidth - 8\tabcolsep) * \real{0.2000}}
  >{\raggedright\arraybackslash}p{(\linewidth - 8\tabcolsep) * \real{0.2000}}
  >{\raggedright\arraybackslash}p{(\linewidth - 8\tabcolsep) * \real{0.2000}}
  >{\raggedright\arraybackslash}p{(\linewidth - 8\tabcolsep) * \real{0.2000}}@{}}
\toprule\noalign{}
\begin{minipage}[b]{\linewidth}\raggedright
PGmc cell / interpretation
\end{minipage} & \begin{minipage}[b]{\linewidth}\raggedright
Candidate input
\end{minipage} & \begin{minipage}[b]{\linewidth}\raggedright
OE outcome or comparison
\end{minipage} & \begin{minipage}[b]{\linewidth}\raggedright
OE comparison form
\end{minipage} & \begin{minipage}[b]{\linewidth}\raggedright
Result
\end{minipage} \\
\midrule\noalign{}
\endhead
\bottomrule\noalign{}
\endlastfoot
lexeme-level infinitive &
{\emph{*burōjaną}}\index[iv]{03pgmc@\ivlangheader{Proto-Germanic}{}!burojana@\iventry{*burōjaną}{}}
& dictionary headword
{\emph{borian}}\index[iv]{02oe@\ivlangheader{Old English}{West Saxon normalization unmarked}!borian@\iventry{borian}{}}
&
{\emph{borian}}\index[iv]{02oe@\ivlangheader{Old English}{West Saxon normalization unmarked}!borian@\iventry{borian}{}}
& family background, but not the cell compared here \\
3sg present &
{\emph{*búrōθi}}\index[iv]{03pgmc@\ivlangheader{Proto-Germanic}{}!burothi@\iventry{*búrōθi}{}}
& regular output {\emph{boraþ}} `bores' &
{\emph{boraþ}}\index[iv]{02oe@\ivlangheader{Old English}{West Saxon normalization unmarked}!borath@\iventry{boraþ}{}}
`bores' & exact match \\
imperative singular companion &
{\emph{*búrô}}\index[iv]{03pgmc@\ivlangheader{Proto-Germanic}{}!buro@\iventry{*búrô}{}}
& related finite form
{\emph{bora}}\index[iv]{02oe@\ivlangheader{Old English}{West Saxon normalization unmarked}!bora@\iventry{bora}{}}
`bore' &
{\emph{bora}}\index[iv]{02oe@\ivlangheader{Old English}{West Saxon normalization unmarked}!bora@\iventry{bora}{}}
`bore' & useful control, but not the target \\
\end{longtable}
}

\section{\texorpdfstring{bore (iptv.2sg) --- OE
\emph{bora}}{bore (iptv.2sg) --- OE bora}}\label{bore-iptv.2sg-oe-bora}

\index[iv]{02oe@\ivlangheader{Old English}{West Saxon normalization unmarked}!bora@\iventry{bora}{}}
\index[iv]{03pgmc@\ivlangheader{Proto-Germanic}{}!buro@\iventry{*búrô}{}}
\index[iv]{03pgmc@\ivlangheader{Proto-Germanic}{}!burona@\iventry{*burōną}{}}

Derivation: citation reconstruction \emph{*burōną}; form followed here
\emph{*búrô} \textgreater{} \emph{bora} (late analogy).

\subsection{Derivation trace}\label{derivation-trace-142}

Proto input: \emph{*búrô}

\begingroup
\setlength{\fboxsep}{6pt}

\noindent\fbox{%
\begin{minipage}{0.97\linewidth}
\small
\begin{tabular}{@{}>{\raggedright\arraybackslash}p{0.359\linewidth}@{\hspace{0.030\linewidth}}>{\raggedright\arraybackslash}p{0.611\linewidth}@{}}
\begin{minipage}[t]{\linewidth}
\centering\textbf{Earlier Germanic changes}\par
\vspace{0.35em}
\raggedright
\centering\textbf{Northwest and West Germanic}\par
\raggedright
\vspace{0.2em}
\begin{tabularx}{\linewidth}{@{}>{\raggedright\arraybackslash}X@{\hspace{0.55em}}>{\raggedright\arraybackslash}p{0.14\linewidth}@{\hspace{0.25em}}}
\mbox{PNWGmc U Lowering} & \emph{*bórô} \\
\end{tabularx}
\vspace{0.6em}
\centering\textbf{Early Anglo-Frisian}\par
\raggedright
\vspace{0.2em}
\raggedright [no change]\par
\end{minipage}
&
\begin{minipage}[t]{\linewidth}
\centering\textbf{Old English changes}\par
\vspace{0.35em}
\raggedright
\begin{tabularx}{\linewidth}{@{}>{\raggedright\arraybackslash}X@{\hspace{0.55em}}>{\raggedright\arraybackslash}p{0.14\linewidth}@{\hspace{0.25em}}}
OE Unstressed Long Vowel Shortening & \emph{*bóra} \\
\end{tabularx}
\end{minipage}
\\
\end{tabular}
\end{minipage}%
} \endgroup

Old English form: \emph{bora}

\subsection{Reconstruction and comparative
evidence}\label{reconstruction-and-comparative-evidence-125}

Kroonen reconstructs the bore-family verb as \emph{*burojan-} and cites
Old English
{\emph{borian}}\index[iv]{02oe@\ivlangheader{Old English}{West Saxon normalization unmarked}!borian@\iventry{borian}{}}
`bore' among its continuants (\citeproc{ref-Kroonen2013}{Kroonen 2013,
85}). Ringe and Taylor give the class-II weak imperative singular ending
as -a \textless{} \emph{*-ō} (\citeproc{ref-RingeTaylor2014}{Ringe and
Taylor 2014, 314}).

The derivational input \Recon{búrô} `bore (iptv.2sg)' is therefore an
imperative cell of the same family, not the citation form of the verb.

\subsection{Old English evidence}\label{old-english-evidence-126}

Clark Hall lemmatizes the verb as
{\emph{borian}}\index[iv]{02oe@\ivlangheader{Old English}{West Saxon normalization unmarked}!borian@\iventry{borian}{}}
`bore' (\citeproc{ref-ClarkHall1960}{Clark Hall 1960, 48}). The
comparison form here is the normalized imperative singular \emph{bora}
`bore (iptv.2sg)', used beside the headword and the related 3sg form
\emph{boraþ} `bores'.

The imperative is thus a paradigm form rather than a replacement for the
dictionary lemma. It is the most direct Old English comparator for the
non-\emph{j} finite cell represented by \Recon{búrô} `bore (iptv.2sg)'.

\subsection{Development to Old
English}\label{development-to-old-english-124}

Northwest Germanic lowering first gives \Recon{bórô} `bore (iptv.2sg)'
from \Recon{búrô} `bore (iptv.2sg)', and late shortening of the
unstressed long vowel then yields \Recon{bóra} `bore (iptv.2sg)', whence
\emph{bora} `bore'.

Ringe and Taylor's class-II imperative singular -a \textless{}
\emph{*-ō} points to exactly this type of outcome
(\citeproc{ref-RingeTaylor2014}{Ringe and Taylor 2014, 314}). The form
compared here therefore isolates the regular finite-cell development
more cleanly than the remodelled infinitive does.

\subsection{Paradigm comparison}\label{paradigm-comparison-4}

The comparison below sets the relevant forms side by side. It
distinguishes the selected imperative cell from the bore-family lemma
and from the companion 3sg form.

{\def\LTcaptype{none} % do not increment counter
\begin{longtable}[]{@{}
  >{\raggedright\arraybackslash}p{(\linewidth - 8\tabcolsep) * \real{0.2000}}
  >{\raggedright\arraybackslash}p{(\linewidth - 8\tabcolsep) * \real{0.2000}}
  >{\raggedright\arraybackslash}p{(\linewidth - 8\tabcolsep) * \real{0.2000}}
  >{\raggedright\arraybackslash}p{(\linewidth - 8\tabcolsep) * \real{0.2000}}
  >{\raggedright\arraybackslash}p{(\linewidth - 8\tabcolsep) * \real{0.2000}}@{}}
\toprule\noalign{}
\begin{minipage}[b]{\linewidth}\raggedright
PGmc cell / interpretation
\end{minipage} & \begin{minipage}[b]{\linewidth}\raggedright
Candidate input
\end{minipage} & \begin{minipage}[b]{\linewidth}\raggedright
OE outcome or comparison
\end{minipage} & \begin{minipage}[b]{\linewidth}\raggedright
OE comparison form
\end{minipage} & \begin{minipage}[b]{\linewidth}\raggedright
Result
\end{minipage} \\
\midrule\noalign{}
\endhead
\bottomrule\noalign{}
\endlastfoot
lexeme-level infinitive &
{\emph{*burōjaną}}\index[iv]{03pgmc@\ivlangheader{Proto-Germanic}{}!burojana@\iventry{*burōjaną}{}}
& dictionary headword
{\emph{borian}}\index[iv]{02oe@\ivlangheader{Old English}{West Saxon normalization unmarked}!borian@\iventry{borian}{}}
&
{\emph{borian}}\index[iv]{02oe@\ivlangheader{Old English}{West Saxon normalization unmarked}!borian@\iventry{borian}{}}
& family background, but not the cell compared here \\
imperative singular &
{\emph{*búrô}}\index[iv]{03pgmc@\ivlangheader{Proto-Germanic}{}!buro@\iventry{*búrô}{}}
& regular output {\emph{bora}} `bore' &
{\emph{bora}}\index[iv]{02oe@\ivlangheader{Old English}{West Saxon normalization unmarked}!bora@\iventry{bora}{}}
`bore' & exact match \\
3sg present companion &
{\emph{*búrōθi}}\index[iv]{03pgmc@\ivlangheader{Proto-Germanic}{}!burothi@\iventry{*búrōθi}{}}
& related finite form
{\emph{boraþ}}\index[iv]{02oe@\ivlangheader{Old English}{West Saxon normalization unmarked}!borath@\iventry{boraþ}{}}
`bores' &
{\emph{boraþ}}\index[iv]{02oe@\ivlangheader{Old English}{West Saxon normalization unmarked}!borath@\iventry{boraþ}{}}
`bores' & useful control, but not the target \\
\end{longtable}
}

\section{\texorpdfstring{bow --- OE
\emph{bēag}}{bow --- OE bēag}}\label{bow-oe-bux113ag}

\index[iv]{02oe@\ivlangheader{Old English}{West Saxon normalization unmarked}!beag@\iventry{bēag}{}}
\index[iv]{03pgmc@\ivlangheader{Proto-Germanic}{}!baug@\iventry{*báug}{}}
\index[iv]{03pgmc@\ivlangheader{Proto-Germanic}{}!beugana@\iventry{*béuganą}{}}

Derivation: citation reconstruction \emph{*béuganą}; form followed here
\emph{*báug} \textgreater{} \emph{bēag} (late analogy).

\subsection{Derivation trace}\label{derivation-trace-143}

Proto input: \emph{*báug}

\begingroup
\setlength{\fboxsep}{6pt}

\noindent\fbox{%
\begin{minipage}{0.97\linewidth}
\small
\begin{tabular}{@{}>{\raggedright\arraybackslash}p{0.388\linewidth}@{\hspace{0.030\linewidth}}>{\raggedright\arraybackslash}p{0.582\linewidth}@{}}
\begin{minipage}[t]{\linewidth}
\centering\textbf{Earlier Germanic changes}\par
\vspace{0.35em}
\raggedright
\centering\textbf{Northwest and West Germanic}\par
\raggedright
\vspace{0.2em}
\raggedright [no change]\par
\vspace{0.6em}
\centering\textbf{Early Anglo-Frisian}\par
\raggedright
\vspace{0.2em}
\raggedright [no change]\par
\end{minipage}
&
\begin{minipage}[t]{\linewidth}
\centering\textbf{Old English changes}\par
\vspace{0.35em}
\raggedright
\begin{tabularx}{\linewidth}{@{}>{\raggedright\arraybackslash}X@{\hspace{0.55em}}>{\raggedright\arraybackslash}p{0.14\linewidth}@{\hspace{0.25em}}}
\mbox{OE Au Brightening} & \emph{*báeug} \\
\mbox{OE Diphthong Leveling} & \emph{*bēag} \\
\end{tabularx}
\end{minipage}
\\
\end{tabular}
\end{minipage}%
} \endgroup

Old English form: \emph{bēag}

\subsection{Reconstruction and comparative
evidence}\label{reconstruction-and-comparative-evidence-126}

The inherited verb belongs to the class-II strong-verb family
\Recon{béuganą}\index[iv]{03pgmc@\ivlangheader{Proto-Germanic}{}!beugana@\iventry{*béuganą}{}}
`bow' (\citeproc{ref-RingeTaylor2014}{Ringe and Taylor 2014, 55}).
Within that paradigm, however, the infinitive and the singular preterite
continue different ablaut grades. The derivational input
\Recon{báug}\index[iv]{03pgmc@\ivlangheader{Proto-Germanic}{}!baug@\iventry{*báug}{}}
`bow' is the singular preterite cell, whereas the citation form
\Recon{béuganą}\index[iv]{03pgmc@\ivlangheader{Proto-Germanic}{}!beugana@\iventry{*béuganą}{}}
is the infinitive.

Campbell's account of Old English class-II strong verbs treats the
singular preterite \emph{au} \textgreater{} \emph{ēa} development as
regular in this environment (\citeproc{ref-Campbell1959}{Campbell 1959,
53}). That is the phonological path relevant for
{\emph{bēag}}\index[iv]{02oe@\ivlangheader{Old English}{West Saxon normalization unmarked}!beag@\iventry{bēag}{}}
`bowed', whereas the analogical \emph{ū} of the present stem belongs to
the separate history behind
{\emph{būgan}}\index[iv]{02oe@\ivlangheader{Old English}{West Saxon normalization unmarked}!bugan@\iventry{būgan}{}}
`to bow' (\citeproc{ref-RingeTaylor2014}{Ringe and Taylor 2014, 55}).

\subsection{Old English evidence}\label{old-english-evidence-127}

Bosworth-Toller and Clark Hall both record
{\emph{bēag}}\index[iv]{02oe@\ivlangheader{Old English}{West Saxon normalization unmarked}!beag@\iventry{bēag}{}}
`bowed' as a preterite form of
{\emph{būgan}}\index[iv]{02oe@\ivlangheader{Old English}{West Saxon normalization unmarked}!bugan@\iventry{būgan}{}}
`to bow' (\citeproc{ref-BosworthToller1898}{Bosworth and Toller 1898,
122}; \citeproc{ref-ClarkHall1960}{Clark Hall 1960, 45}). The form
discussed here is therefore an attested Old English verbal form, not a
reconstructed substitute for the infinitive.

The ordinary dictionary headword remains
{\emph{būgan}}\index[iv]{02oe@\ivlangheader{Old English}{West Saxon normalization unmarked}!bugan@\iventry{būgan}{}}
`to bow', but the relevant comparison form for this entry is the
singular preterite
{\emph{bēag}}\index[iv]{02oe@\ivlangheader{Old English}{West Saxon normalization unmarked}!beag@\iventry{bēag}{}}
`bow, ring' `bowed'. That is the paradigm cell in which the inherited
\emph{*au} grade is preserved most directly.

\subsection{Development to Old
English}\label{development-to-old-english-125}

From
\Recon{báug}\index[iv]{03pgmc@\ivlangheader{Proto-Germanic}{}!baug@\iventry{*báug}{}}
`bow', Anglo-Frisian fronting and the later leveling of the diphthong
produce
{\emph{bēag}}\index[iv]{02oe@\ivlangheader{Old English}{West Saxon normalization unmarked}!beag@\iventry{bēag}{}}
`bow, ring' (\citeproc{ref-Campbell1959}{Campbell 1959, 53}). No special
analogical repair is needed for that cell. The form is the regular Old
English outcome of the singular-preterite grade.

The analogical element in the wider lexeme belongs instead to the
present stem seen in
{\emph{būgan}}\index[iv]{02oe@\ivlangheader{Old English}{West Saxon normalization unmarked}!bugan@\iventry{būgan}{}}
`bow, bend'. The derivational input differs from the citation form
because the regular inherited pathway survives more transparently in the
preterite than in the infinitive.

\subsection{Paradigm comparison}\label{paradigm-comparison-5}

The comparison below sets the relevant forms side by side. It
distinguishes the regular singular preterite from the more familiar
infinitival citation form.

{\def\LTcaptype{none} % do not increment counter
\begin{longtable}[]{@{}
  >{\raggedright\arraybackslash}p{(\linewidth - 8\tabcolsep) * \real{0.2000}}
  >{\raggedright\arraybackslash}p{(\linewidth - 8\tabcolsep) * \real{0.2000}}
  >{\raggedright\arraybackslash}p{(\linewidth - 8\tabcolsep) * \real{0.2000}}
  >{\raggedright\arraybackslash}p{(\linewidth - 8\tabcolsep) * \real{0.2000}}
  >{\raggedright\arraybackslash}p{(\linewidth - 8\tabcolsep) * \real{0.2000}}@{}}
\toprule\noalign{}
\begin{minipage}[b]{\linewidth}\raggedright
PGmc cell / interpretation
\end{minipage} & \begin{minipage}[b]{\linewidth}\raggedright
Candidate input
\end{minipage} & \begin{minipage}[b]{\linewidth}\raggedright
Expected or documented OE outcome
\end{minipage} & \begin{minipage}[b]{\linewidth}\raggedright
OE comparison form
\end{minipage} & \begin{minipage}[b]{\linewidth}\raggedright
Result
\end{minipage} \\
\midrule\noalign{}
\endhead
\bottomrule\noalign{}
\endlastfoot
citation infinitive &
{\emph{*béuganą}}\index[iv]{03pgmc@\ivlangheader{Proto-Germanic}{}!beugana@\iventry{*béuganą}{}}
& inherited present-stem history behind
{\emph{būgan}}\index[iv]{02oe@\ivlangheader{Old English}{West Saxon normalization unmarked}!bugan@\iventry{būgan}{}}
&
{\emph{būgan}}\index[iv]{02oe@\ivlangheader{Old English}{West Saxon normalization unmarked}!bugan@\iventry{būgan}{}}
& establishes the lexeme, but not the Old English form here \\
singular preterite &
{\emph{*báug}}\index[iv]{03pgmc@\ivlangheader{Proto-Germanic}{}!baug@\iventry{*báug}{}}
& regular output:
{\emph{bēag}}\index[iv]{02oe@\ivlangheader{Old English}{West Saxon normalization unmarked}!beag@\iventry{bēag}{}}
&
{\emph{bēag}}\index[iv]{02oe@\ivlangheader{Old English}{West Saxon normalization unmarked}!beag@\iventry{bēag}{}}
& exact match between input, output, and attested cell \\
past participial branch & participial \emph{*bugan-} type & later
participial outcomes & bogen-type evidence & relevant to the paradigm,
but not the clearest match for this entry \\
\end{longtable}
}

The singular preterite is the relevant comparison form. It gives a
direct lautgesetzlich path to attested
{\emph{bēag}}\index[iv]{02oe@\ivlangheader{Old English}{West Saxon normalization unmarked}!beag@\iventry{bēag}{}}
`bow, ring', while the citation form
{\emph{būgan}}\index[iv]{02oe@\ivlangheader{Old English}{West Saxon normalization unmarked}!bugan@\iventry{būgan}{}}
`bow, bend' belongs to a paradigm whose present stem has already
undergone later leveling.

\section{\texorpdfstring{cow --- OE
\emph{cȳ}}{cow --- OE cȳ}}\label{cow-oe-cux233}

\index[iv]{02oe@\ivlangheader{Old English}{West Saxon normalization unmarked}!cum@\iventry{cūm}{}}
\index[iv]{02oe@\ivlangheader{Old English}{West Saxon normalization unmarked}!cy@\iventry{cȳ}{}}
\index[iv]{03pgmc@\ivlangheader{Proto-Germanic}{}!koz@\iventry{*kōz}{}}
\index[iv]{03pgmc@\ivlangheader{Proto-Germanic}{}!kui@\iventry{*kūi}{}}

Derivation: citation reconstruction \emph{*kōz}; form followed here
\emph{*kūi} \textgreater{} \emph{cȳ} (late analogy).

\subsection{Derivation trace}\label{derivation-trace-144}

Proto input: \emph{*kūi}

\begingroup
\setlength{\fboxsep}{6pt}

\noindent\fbox{%
\begin{minipage}{0.97\linewidth}
\small
\begin{tabular}{@{}>{\raggedright\arraybackslash}p{0.388\linewidth}@{\hspace{0.030\linewidth}}>{\raggedright\arraybackslash}p{0.582\linewidth}@{}}
\begin{minipage}[t]{\linewidth}
\centering\textbf{Earlier Germanic changes}\par
\vspace{0.35em}
\raggedright
\centering\textbf{Northwest and West Germanic}\par
\raggedright
\vspace{0.2em}
\raggedright [no change]\par
\vspace{0.6em}
\centering\textbf{Early Anglo-Frisian}\par
\raggedright
\vspace{0.2em}
\raggedright [no change]\par
\end{minipage}
&
\begin{minipage}[t]{\linewidth}
\centering\textbf{Old English changes}\par
\vspace{0.35em}
\raggedright
\begin{tabularx}{\linewidth}{@{}>{\raggedright\arraybackslash}X@{\hspace{0.55em}}>{\raggedright\arraybackslash}p{0.14\linewidth}@{\hspace{0.25em}}}
\mbox{OE I Umlaut} & \emph{*kȳi} \\
\mbox{OE High Vowel Apocope} & \emph{*kȳ} \\
\end{tabularx}
\end{minipage}
\\
\end{tabular}
\end{minipage}%
} \endgroup

Old English form: \emph{cȳ}

\subsection{Reconstruction and comparative
evidence}\label{reconstruction-and-comparative-evidence-127}

Kroonen reconstructs a root noun with the inherited alternation
\emph{*kō-} \textasciitilde{} \emph{*ku-}, explicitly nom.
\emph{\Recon{kōz}} `cow', obl. \emph{*kū-}
(\citeproc{ref-Kroonen2013}{Kroonen 2013, 299}). The citation form
therefore belongs to the nominative singular, whereas the derivational
input \Recon{kūi} `cow' belongs to the oblique stem.

Ringe and Taylor also posit a later PNWGmc nominative \Recon{kūaz} `cow'
\textgreater{} \Recon{kūz} `cow' behind Old English
{\emph{cū}}\index[iv]{02oe@\ivlangheader{Old English}{West Saxon normalization unmarked}!cu@\iventry{cū}{}}
`cow' (\citeproc{ref-RingeTaylor2014}{Ringe and Taylor 2014, § 3.1.3}). That nominative history accounts for the headword, whereas the
form cȳ `cow' depends on the oblique \emph{*kū-} stem and a following
\emph{*i}.

\subsection{Old English evidence}\label{old-english-evidence-128}

Clark Hall lemmatizes the noun under
{\emph{cū}}\index[iv]{02oe@\ivlangheader{Old English}{West Saxon normalization unmarked}!cu@\iventry{cū}{}}
`cow' and records gen.sg.
{\emph{cū}}\index[iv]{02oe@\ivlangheader{Old English}{West Saxon normalization unmarked}!cu@\iventry{cū}{}}
`cow' (or
{\emph{cū(e)}}\index[iv]{02oe@\ivlangheader{Old English}{West Saxon normalization unmarked}!cue@\iventry{cū(e)}{}}),
{\emph{cȳ}}\index[iv]{02oe@\ivlangheader{Old English}{West Saxon normalization unmarked}!cy@\iventry{cȳ}{}}
`cow',
{\emph{cūs}}\index[iv]{02oe@\ivlangheader{Old English}{West Saxon normalization unmarked}!cus@\iventry{cūs}{}}
`cow (gen.)', dat.sg.
{\emph{cȳ}}\index[iv]{02oe@\ivlangheader{Old English}{West Saxon normalization unmarked}!cy@\iventry{cȳ}{}}
`cow', nom.-acc.pl.
{\emph{cȳ}}\index[iv]{02oe@\ivlangheader{Old English}{West Saxon normalization unmarked}!cy@\iventry{cȳ}{}}
`cow', and dat.pl. \emph{cūm} `cow' (\citeproc{ref-ClarkHall1960}{Clark
Hall 1960, 67}). Ringe and Taylor likewise state that the root-noun
{\emph{cū}}\index[iv]{02oe@\ivlangheader{Old English}{West Saxon normalization unmarked}!cu@\iventry{cū}{}}
`cow' exhibits dat.sg., nom.-acc.pl.
{\emph{cȳ}}\index[iv]{02oe@\ivlangheader{Old English}{West Saxon normalization unmarked}!cy@\iventry{cȳ}{}}
`cow' \textless{} \Recon{cūi} `cow', dat.pl. \emph{cūm} \textless{}
\emph{*cūm}, and apparently gen.sg.
{\emph{cā}}\index[iv]{02oe@\ivlangheader{Old English}{West Saxon normalization unmarked}!ca@\iventry{cā}{}}
`cow (gen.)' \textless{} \Recon{cūiz} `cow (gen.)'
(\citeproc{ref-RingeTaylor2014}{Ringe and Taylor 2014, § 6.6.1}).

The dictionary headword is therefore
{\emph{cū}}\index[iv]{02oe@\ivlangheader{Old English}{West Saxon normalization unmarked}!cu@\iventry{cū}{}}
`cow', but
{\emph{cȳ}}\index[iv]{02oe@\ivlangheader{Old English}{West Saxon normalization unmarked}!cy@\iventry{cȳ}{}}
`cow' is an established Old English paradigm form rather than a
convenient reconstruction. It serves as dative singular and also as
nominative-accusative plural. The genitive singular is less stable, with
{\emph{cā}}\index[iv]{02oe@\ivlangheader{Old English}{West Saxon normalization unmarked}!ca@\iventry{cā}{}}
`cow (gen.sg.)',
{\emph{cȳ}}\index[iv]{02oe@\ivlangheader{Old English}{West Saxon normalization unmarked}!cy@\iventry{cȳ}{}}
`cow', \emph{cū(e)}, and
{\emph{cūs}}\index[iv]{02oe@\ivlangheader{Old English}{West Saxon normalization unmarked}!cus@\iventry{cūs}{}}
`cow (gen.sg.)' all appearing in the local source record.

\subsection{Development to Old
English}\label{development-to-old-english-126}

\Recon{kūi} `cow' is the dative singular of the oblique \emph{*kū-}
stem. The following \emph{*i} triggers i-umlaut, so \emph{ū} becomes
\emph{ȳ}, and loss of the final high vowel leaves
{\emph{cȳ}}\index[iv]{02oe@\ivlangheader{Old English}{West Saxon normalization unmarked}!cy@\iventry{cȳ}{}}
`cow'. The development is \Recon{kūi} `cow' \textgreater{} \Recon{kȳi}
`cow' \textgreater{} \Recon{kȳ} `cow' \textgreater{} cȳ.

This is the regular oblique-cell path recognized by Ringe and Taylor's
paradigm statement (\citeproc{ref-RingeTaylor2014}{Ringe and Taylor
2014, § 6.6.1}). It also explains why
{\emph{cȳ}}\index[iv]{02oe@\ivlangheader{Old English}{West Saxon normalization unmarked}!cy@\iventry{cȳ}{}}
`cow' is the cleanest comparison form for the present entry, even though
the ordinary headword is
{\emph{cū}}\index[iv]{02oe@\ivlangheader{Old English}{West Saxon normalization unmarked}!cu@\iventry{cū}{}}
`cow'.

\subsection{Paradigm comparison}\label{paradigm-comparison-6}

A paradigm comparison identifies the Proto-Germanic inflectional cell
that corresponds to an established Old English paradigm form. The
comparison below sets the relevant forms side by side.

{\def\LTcaptype{none} % do not increment counter
\begin{longtable}[]{@{}
  >{\raggedright\arraybackslash}p{(\linewidth - 8\tabcolsep) * \real{0.2000}}
  >{\raggedright\arraybackslash}p{(\linewidth - 8\tabcolsep) * \real{0.2000}}
  >{\raggedright\arraybackslash}p{(\linewidth - 8\tabcolsep) * \real{0.2000}}
  >{\raggedright\arraybackslash}p{(\linewidth - 8\tabcolsep) * \real{0.2000}}
  >{\raggedright\arraybackslash}p{(\linewidth - 8\tabcolsep) * \real{0.2000}}@{}}
\toprule\noalign{}
\begin{minipage}[b]{\linewidth}\raggedright
PGmc cell / interpretation
\end{minipage} & \begin{minipage}[b]{\linewidth}\raggedright
Candidate input
\end{minipage} & \begin{minipage}[b]{\linewidth}\raggedright
OE output or comparison
\end{minipage} & \begin{minipage}[b]{\linewidth}\raggedright
OE comparison form
\end{minipage} & \begin{minipage}[b]{\linewidth}\raggedright
Result
\end{minipage} \\
\midrule\noalign{}
\endhead
\bottomrule\noalign{}
\endlastfoot
citation nominative singular & *kōz & OE headword
{\emph{cū}}\index[iv]{02oe@\ivlangheader{Old English}{West Saxon normalization unmarked}!cu@\iventry{cū}{}}
`cow' belongs to the nominative history of the lexeme & \emph{cū} &
useful background, but not the chosen comparison for
{\emph{cȳ}}\index[iv]{02oe@\ivlangheader{Old English}{West Saxon normalization unmarked}!cy@\iventry{cȳ}{}}
`cow' \\
later generalized nominative & PNWGmc \emph{kūaz \textgreater{} }kūz &
inferred nominative
{\emph{cū}}\index[iv]{02oe@\ivlangheader{Old English}{West Saxon normalization unmarked}!cu@\iventry{cū}{}}
`cow' & \emph{cū} & explains the leveled headword, not the oblique
target \\
dative singular oblique & *kūi & regular output:
{\emph{cȳ}}\index[iv]{02oe@\ivlangheader{Old English}{West Saxon normalization unmarked}!cy@\iventry{cȳ}{}}
`cow' & \emph{cȳ} & exact match between input, output, and paradigm
cell \\
genitive singular oblique & *kūiz & Ringe-Taylor: apparently
{\emph{cā}}\index[iv]{02oe@\ivlangheader{Old English}{West Saxon normalization unmarked}!ca@\iventry{cā}{}}
`cow (gen.sg.)'; Hall: \emph{cū(e)}, \emph{cȳ}, \emph{cūs} & gen.sg.
variable & too unstable to control the entry \\
\end{longtable}
}

The dative singular is the relevant comparison form. It gives a regular
path to attested
{\emph{cȳ}}\index[iv]{02oe@\ivlangheader{Old English}{West Saxon normalization unmarked}!cy@\iventry{cȳ}{}}
`cow', while the broader Old English paradigm shows how far the oblique
\emph{*kū-} grade spread beyond that one cell.

\section{\texorpdfstring{find --- OE
\emph{fundne}}{find --- OE fundne}}\label{find-oe-fundne}

\index[iv]{02oe@\ivlangheader{Old English}{West Saxon normalization unmarked}!fundne@\iventry{fundne}{}}
\index[iv]{03pgmc@\ivlangheader{Proto-Germanic}{}!finthana@\iventry{*fínθaną}{}}
\index[iv]{03pgmc@\ivlangheader{Proto-Germanic}{}!fundanai@\iventry{*fúnðanai}{}}

Derivation: citation reconstruction \emph{*fínθaną}; form followed here
\emph{*fúnðanai} \textgreater{} \emph{fundne} (late analogy).

\subsection{Derivation trace}\label{derivation-trace-145}

Proto input: \emph{*fúnðanai}

\begingroup
\setlength{\fboxsep}{6pt}

\noindent\fbox{%
\begin{minipage}{0.97\linewidth}
\small
\begin{tabular}{@{}>{\raggedright\arraybackslash}p{0.485\linewidth}@{\hspace{0.030\linewidth}}>{\raggedright\arraybackslash}p{0.485\linewidth}@{}}
\begin{minipage}[t]{\linewidth}
\centering\textbf{Earlier Germanic changes}\par
\vspace{0.35em}
\raggedright
\centering\textbf{Northwest and West Germanic}\par
\raggedright
\vspace{0.2em}
\begin{tabularx}{\linewidth}{@{}>{\raggedright\arraybackslash}X@{\hspace{0.55em}}>{\raggedright\arraybackslash}p{0.18\linewidth}@{\hspace{0.25em}}}
PNWGmc Unstressed Ai Monophthongization & \emph{*fúnðanē} \\
\mbox{PWGmc Dental Hardening} & \emph{*fúndanē} \\
\end{tabularx}
\vspace{0.6em}
\centering\textbf{Early Anglo-Frisian}\par
\raggedright
\vspace{0.2em}
\raggedright [no change]\par
\end{minipage}
&
\begin{minipage}[t]{\linewidth}
\centering\textbf{Old English changes}\par
\vspace{0.35em}
\raggedright
\begin{tabularx}{\linewidth}{@{}>{\raggedright\arraybackslash}X@{\hspace{0.55em}}>{\raggedright\arraybackslash}p{0.14\linewidth}@{\hspace{0.25em}}}
\mbox{OE Non High Syn} & \emph{*fúndnē} \\
OE Unstressed Long Vowel Shortening & \emph{*fúndne} \\
\end{tabularx}
\end{minipage}
\\
\end{tabular}
\end{minipage}%
} \endgroup

Old English form: \emph{fundne}

\subsection{Reconstruction and comparative
evidence}\label{reconstruction-and-comparative-evidence-128}

The inherited strong verb
\Recon{fínθaną}\index[iv]{03pgmc@\ivlangheader{Proto-Germanic}{}!finthana@\iventry{*fínθaną}{}}
`find' underlies Old English
{\emph{findan}}\index[iv]{02oe@\ivlangheader{Old English}{West Saxon normalization unmarked}!findan@\iventry{findan}{}}
`find'. Its infinitive identifies the lexeme, not the participial cell
followed here. The dental alternation and the participle's root grade
must be distinguished from its inflectional ending
(\citeproc{ref-RingeTaylor2014}{Ringe and Taylor 2014, 140, 344};
\citeproc{ref-Fulk2018}{Fulk 2018, 108, 253, 284--88}).

The selected
\Recon{fúnðanai}\index[iv]{03pgmc@\ivlangheader{Proto-Germanic}{}!fundanai@\iventry{*fúnðanai}{}}
is a compositional reconstruction: the participial stem in inherited
\emph{-an-} takes the strong masculine nominative plural adjectival
ending \emph{-ai}. Fulk supplies the relevant inflection and
distinguishes the ordinary participial \emph{-an-} formation from rarer
inherited \emph{-in-}; he does not print this whole accented word
(\citeproc{ref-Fulk2018}{Fulk 2018, 209}, §9.2, p.~284, §12.30). The
choice is therefore not a claim that every Old English participial
\emph{-en-} must result from analogy.

\subsection{Old English evidence}\label{old-english-evidence-129}

The selected
{\emph{fundne}}\index[iv]{02oe@\ivlangheader{Old English}{West Saxon normalization unmarked}!fundne@\iventry{fundne}{}}
`found' is directly attested in \emph{Juliana} 335. Strunk prints the
plural construction \emph{gif we gemete sin on moldwege / oþþe feor oþþe
neah fundne weorþen}; his glossary explicitly labels the form \emph{ptc
npm.} The masculine nominative plural identification is thus supported
by syntax and the editor's analysis, not inferred from a successful
derivation (\citeproc{ref-Strunk1904}{Strunk 1904, 15, 82}). The poem,
manuscript and modern edition have different dates; the edition's date
does not date the phonology.

The longer
{\emph{fundene}}\index[iv]{02oe@\ivlangheader{Old English}{West Saxon normalization unmarked}!fundene@\iventry{fundene}{}}
`found' remains genuine evidence. Toller cites \emph{Beón þá herigeata
swa fundene}, a plural construction, not the masculine accusative
singular previously assigned to it (\citeproc{ref-Toller1921}{Toller
1921, 219}). The familiar dictionary participle
{\emph{funden}}\index[iv]{02oe@\ivlangheader{Old English}{West Saxon normalization unmarked}!funden@\iventry{funden}{}}
`found' and the formation \emph{tō-fundennes} `finding' remain relevant
background (\citeproc{ref-ClarkHall1960}{Clark Hall 1960, 124}). None is
suppressed by selecting the independently attested syncopated plural.

\subsection{Development to Old
English}\label{development-to-old-english-127}

Unstressed final \emph{ai} first becomes long \emph{ē}. Dental hardening
gives
\Recon{fúndanē}\index[iv]{08preoe@\ivlangheader{Pre-Old English / prehistoric English}{}!fundane@\iventry{*fúndanē}{}}.
Its medial \emph{a} is in an open posttonic syllable before the
still-heavy final syllable. \hyperref[rule-OENonHighSyn]{SC114
OENonHighSyn} therefore gives
\Recon{fúndnē}\index[iv]{08preoe@\ivlangheader{Pre-Old English / prehistoric English}{}!fundne@\iventry{*fúndnē}{}};
later unstressed shortening supplies final \emph{e}, yielding
{\emph{fundne}}. No analogy intervenes within this selected span.

The retained late-analogy classification concerns the wider participial
paradigm, not an analogical operation needed to manufacture this output.
The relationship of the competing inherited suffixes, the longer plural
and subsequent levelling remains a separate problem. A regular cell does
not by itself settle that paradigm.

\subsection{Paradigm comparison}\label{paradigm-comparison-7}

{\def\LTcaptype{none} % do not increment counter
\begin{longtable}[]{@{}
  >{\raggedright\arraybackslash}p{(\linewidth - 8\tabcolsep) * \real{0.2000}}
  >{\raggedright\arraybackslash}p{(\linewidth - 8\tabcolsep) * \real{0.2000}}
  >{\raggedright\arraybackslash}p{(\linewidth - 8\tabcolsep) * \real{0.2000}}
  >{\raggedright\arraybackslash}p{(\linewidth - 8\tabcolsep) * \real{0.2000}}
  >{\raggedright\arraybackslash}p{(\linewidth - 8\tabcolsep) * \real{0.2000}}@{}}
\toprule\noalign{}
\begin{minipage}[b]{\linewidth}\raggedright
PGmc cell / interpretation
\end{minipage} & \begin{minipage}[b]{\linewidth}\raggedright
Candidate input
\end{minipage} & \begin{minipage}[b]{\linewidth}\raggedright
Actual model output
\end{minipage} & \begin{minipage}[b]{\linewidth}\raggedright
OE comparison
\end{minipage} & \begin{minipage}[b]{\linewidth}\raggedright
Assessment
\end{minipage} \\
\midrule\noalign{}
\endhead
\bottomrule\noalign{}
\endlastfoot
selected strong masculine nominative plural &
\Recon{fúnðanai}\index[iv]{03pgmc@\ivlangheader{Proto-Germanic}{}!fundanai@\iventry{*fúnðanai}{}}
& {\emph{fundne}} &
{\emph{fundne}}\index[iv]{02oe@\ivlangheader{Old English}{West Saxon normalization unmarked}!fundne@\iventry{fundne}{}},
\emph{Juliana} 335 & complete regular selected span \\
participial nominative singular in \emph{-an-} &
\Recon{fúnðanaz}\index[iv]{03pgmc@\ivlangheader{Proto-Germanic}{}!fundanaz@\iventry{*fúnðanaz}{}}
& \emph{fundan} &
{\emph{funden}}\index[iv]{02oe@\ivlangheader{Old English}{West Saxon normalization unmarked}!funden@\iventry{funden}{}}
& not a complete successful derivation of the attested comparison \\
inherited \emph{-in-} alternative &
\Recon{fúnðinaz}\index[iv]{03pgmc@\ivlangheader{Proto-Germanic}{}!fundinaz@\iventry{*fúnðinaz}{}}
& \emph{fynden} &
{\emph{funden}}\index[iv]{02oe@\ivlangheader{Old English}{West Saxon normalization unmarked}!funden@\iventry{funden}{}}
& umlaut prevents treating this probe as a solution \\
full uncontracted masculine accusative &
\Recon{fúnðananǭ}\index[iv]{03pgmc@\ivlangheader{Proto-Germanic}{}!fundanano@\iventry{*fúnðananǭ}{}}
& \emph{fundenne} &
{\emph{fundene}}\index[iv]{02oe@\ivlangheader{Old English}{West Saxon normalization unmarked}!fundene@\iventry{fundene}{}}
& different ending and output; wider history unresolved \\
\end{longtable}
}

The former shortened input
\Recon{fúnðanǭ}\index[iv]{03pgmc@\ivlangheader{Proto-Germanic}{}!fundano@\iventry{*fúnðanǭ}{}}
did not establish an inherited masculine accusative matching the quoted
plural. It is not retained as evidence that the longer attested form has
been explained. Conversely, these unsuccessful comparisons do not
invalidate the independently supported nominative plural.

\section{\texorpdfstring{fright --- OE
\emph{fyrhte}}{fright --- OE fyrhte}}\label{fright-oe-fyrhte}

\index[iv]{02oe@\ivlangheader{Old English}{West Saxon normalization unmarked}!fyrhte@\iventry{fyrhte}{}}
\index[iv]{03pgmc@\ivlangheader{Proto-Germanic}{}!furxtin@\iventry{*furxtīn}{}}
\index[iv]{03pgmc@\ivlangheader{Proto-Germanic}{}!furxtinaz@\iventry{*fúrxtīnaz}{}}

Derivation: citation reconstruction \emph{*furxtīn}; form followed here
\emph{*fúrxtīnaz} \textgreater{} \emph{fyrhte} (late analogy).

\subsection{Derivation trace}\label{derivation-trace-146}

Proto input: \emph{*fúrxtīnaz}

\begingroup
\setlength{\fboxsep}{6pt}

\noindent\fbox{%
\begin{minipage}{0.97\linewidth}
\small
\begin{tabular}{@{}>{\raggedright\arraybackslash}p{0.359\linewidth}@{\hspace{0.030\linewidth}}>{\raggedright\arraybackslash}p{0.611\linewidth}@{}}
\begin{minipage}[t]{\linewidth}
\centering\textbf{Earlier Germanic changes}\par
\vspace{0.35em}
\raggedright
\centering\textbf{Northwest and West Germanic}\par
\raggedright
\vspace{0.2em}
\begin{tabularx}{\linewidth}{@{}>{\raggedright\arraybackslash}X@{\hspace{0.55em}}>{\raggedright\arraybackslash}p{0.24\linewidth}@{\hspace{0.25em}}}
\mbox{EAF Final Z Deletion} & \emph{*fúrxtīna} \\
\end{tabularx}
\vspace{0.6em}
\centering\textbf{Early Anglo-Frisian}\par
\raggedright
\vspace{0.2em}
\raggedright [no change]\par
\end{minipage}
&
\begin{minipage}[t]{\linewidth}
\centering\textbf{Old English changes}\par
\vspace{0.35em}
\raggedright
\begin{tabularx}{\linewidth}{@{}>{\raggedright\arraybackslash}X@{\hspace{0.55em}}>{\raggedright\arraybackslash}p{0.18\linewidth}@{\hspace{0.25em}}}
\mbox{PWGmc Final Bare A Loss} & \emph{*fúrxtīn} \\
\mbox{NWGmc In Stem N Loss} & \emph{*fúrxtī} \\
\mbox{OE I Umlaut} & \emph{*fyrxtī} \\
OE Unstressed Long Vowel Shortening & \emph{*fyrxti} \\
OE Med Unstressed I Lowering1 & \emph{*fyrxte} \\
\end{tabularx}
\end{minipage}
\\
\end{tabular}
\end{minipage}%
} \endgroup

Old English form: \emph{fyrhte}

\subsection{Reconstruction and comparative
evidence}\label{reconstruction-and-comparative-evidence-129}

The noun belongs to the inherited in-stem abstract
\Recon{furxtīn}\index[iv]{03pgmc@\ivlangheader{Proto-Germanic}{}!furxtin@\iventry{*furxtīn}{}}
`fright', the same family as Gothic
{\emph{faurhtei}}\index[iv]{12goth@\ivlangheader{Gothic}{}!faurhtei@\iventry{faurhtei}{}}
`fear, fright' (\citeproc{ref-Orel2003}{Orel 2003, 120}). The
derivational input
\Recon{fúrxtīnaz}\index[iv]{03pgmc@\ivlangheader{Proto-Germanic}{}!furxtinaz@\iventry{*fúrxtīnaz}{}}
`fright' is not a different lexeme but an oblique singular cell within
that in-stem paradigm.

Ringe and Taylor treat the later nominative forms with \emph{-u} or
\emph{-o} as analogically remodeled
(\citeproc{ref-RingeTaylor2014}{Ringe and Taylor 2014, 380--81}). I
therefore use the oblique in-stem form, which continues the older
formation rather than the remodeled nominative.

\subsection{Old English evidence}\label{old-english-evidence-130}

Bosworth-Toller records
{\emph{fyrhte}}\index[iv]{02oe@\ivlangheader{Old English}{West Saxon normalization unmarked}!fyrhte@\iventry{fyrhte}{}}
`fright' with textual attestation, and it also records nominative forms
such as
{\emph{fyrhtu}}\index[iv]{02oe@\ivlangheader{Old English}{West Saxon normalization unmarked}!fyrhtu@\iventry{fyrhtu}{}}
`fright' and
{\emph{fyrhto}}\index[iv]{02oe@\ivlangheader{Old English}{West Saxon normalization unmarked}!fyrhto@\iventry{fyrhto}{}}
`fright' (\citeproc{ref-BosworthToller1898}{Bosworth and Toller 1898,
160}). Clark Hall lists adjective and verb material separately under
\emph{fyrht} / \emph{fyrhtan} (\citeproc{ref-ClarkHall1960}{Clark Hall
1960, 141}).

The relevant comparison form is therefore the attested oblique
{\emph{fyrhte}}\index[iv]{02oe@\ivlangheader{Old English}{West Saxon normalization unmarked}!fyrhte@\iventry{fyrhte}{}}
`fright'. The nominative lemma forms remain part of the Old English
evidence, but the Old English form here of this entry is the oblique
cell.

\subsection{Development to Old
English}\label{development-to-old-english-128}

From
\Recon{fúrxtīnaz}\index[iv]{03pgmc@\ivlangheader{Proto-Germanic}{}!furxtinaz@\iventry{*fúrxtīnaz}{}}
`fright', the oblique in-stem develops through the loss and weakening of
the final ending, yielding
{\emph{fyrhte}}\index[iv]{02oe@\ivlangheader{Old English}{West Saxon normalization unmarked}!fyrhte@\iventry{fyrhte}{}}
`fright'. The form compared here therefore follows the ordinary Old
English reduction of the abstract ending in this paradigm. Final
\emph{*n} after unstressed \emph{*ī} was lost after Proto-West-Germanic
differentiation, once the inherited ending had become word-final, not as
a productive late Old English deletion. Mutation and subsequent
shortening and lowering then give \emph{fyrhte}
(\citeproc{ref-RingeTaylor2014}{Ringe and Taylor 2014, 87--88};
\citeproc{ref-Campbell1959}{Campbell 1959, 189}, §473, pp.~236--237,
§589.7; \citeproc{ref-SieversBrunner1965}{Brunner 1965, 225--26}, §280;
\citeproc{ref-Fulk2018}{Fulk 2018, 173}, §7.34).

The later nominative forms with \emph{-u} or \emph{-o} belong to a
subsequent analogical reshaping of the paradigm. The Old English form
here is earlier in that sense: it is the attested OE cell in which the
inherited in-stem development remains most transparent.

\subsection{Paradigm comparison}\label{paradigm-comparison-8}

The comparison below sets the relevant forms side by side. It separates
the attested oblique form from the later remodeled nominative line.

{\def\LTcaptype{none} % do not increment counter
\begin{longtable}[]{@{}
  >{\raggedright\arraybackslash}p{(\linewidth - 8\tabcolsep) * \real{0.2000}}
  >{\raggedright\arraybackslash}p{(\linewidth - 8\tabcolsep) * \real{0.2000}}
  >{\raggedright\arraybackslash}p{(\linewidth - 8\tabcolsep) * \real{0.2000}}
  >{\raggedright\arraybackslash}p{(\linewidth - 8\tabcolsep) * \real{0.2000}}
  >{\raggedright\arraybackslash}p{(\linewidth - 8\tabcolsep) * \real{0.2000}}@{}}
\toprule\noalign{}
\begin{minipage}[b]{\linewidth}\raggedright
PGmc cell / interpretation
\end{minipage} & \begin{minipage}[b]{\linewidth}\raggedright
Candidate input
\end{minipage} & \begin{minipage}[b]{\linewidth}\raggedright
Expected or documented OE outcome
\end{minipage} & \begin{minipage}[b]{\linewidth}\raggedright
OE comparison form
\end{minipage} & \begin{minipage}[b]{\linewidth}\raggedright
Result
\end{minipage} \\
\midrule\noalign{}
\endhead
\bottomrule\noalign{}
\endlastfoot
citation in-stem headword &
{\emph{*furxtīn}}\index[iv]{03pgmc@\ivlangheader{Proto-Germanic}{}!furxtin@\iventry{*furxtīn}{}}
& broader noun-class label & wider family context & useful lexeme label,
but not the cell compared here \\
remodeled nominative line & nominative in-stem forms &
{\emph{fyrhtu}}\index[iv]{02oe@\ivlangheader{Old English}{West Saxon normalization unmarked}!fyrhtu@\iventry{fyrhtu}{}}
`fright' /
{\emph{fyrhto}}\index[iv]{02oe@\ivlangheader{Old English}{West Saxon normalization unmarked}!fyrhto@\iventry{fyrhto}{}}
`fright' type lemma forms &
{\emph{fyrhtu}}\index[iv]{02oe@\ivlangheader{Old English}{West Saxon normalization unmarked}!fyrhtu@\iventry{fyrhtu}{}}
`fright' /
{\emph{fyrhto}}\index[iv]{02oe@\ivlangheader{Old English}{West Saxon normalization unmarked}!fyrhto@\iventry{fyrhto}{}}
`fright' & genuine OE evidence, but later remodeled \\
selected oblique singular &
{\emph{*fúrxtīnaz}}\index[iv]{03pgmc@\ivlangheader{Proto-Germanic}{}!furxtinaz@\iventry{*fúrxtīnaz}{}}
& regular output: {\emph{fyrhte}} `fright' &
{\emph{fyrhte}}\index[iv]{02oe@\ivlangheader{Old English}{West Saxon normalization unmarked}!fyrhte@\iventry{fyrhte}{}}
`fright' & exact match between input, output, and attested cell \\
\end{longtable}
}

The oblique in-stem form is the relevant comparison form. It yields
attested
{\emph{fyrhte}}\index[iv]{02oe@\ivlangheader{Old English}{West Saxon normalization unmarked}!fyrhte@\iventry{fyrhte}{}}
`fright' directly, while the more familiar nominative forms belong to a
later analogical layer.

\section{\texorpdfstring{have --- OE
\emph{hæfeþ}}{have --- OE hæfeþ}}\label{have-oe-huxe6feuxfe}

\index[iv]{02oe@\ivlangheader{Old English}{West Saxon normalization unmarked}!haefeth@\iventry{hæfeþ}{}}
\index[iv]{03pgmc@\ivlangheader{Proto-Germanic}{}!xabena@\iventry{*xabēną}{}}
\index[iv]{03pgmc@\ivlangheader{Proto-Germanic}{}!xabethi@\iventry{*xábēθi}{}}

Derivation: citation reconstruction \emph{*xabēną}; form followed here
\emph{*xábēθi} \textgreater{} \emph{hæfeþ} (late analogy).

\subsection{Derivation trace}\label{derivation-trace-147}

Proto input: \emph{*xábēθi}

\begingroup
\setlength{\fboxsep}{6pt}

\noindent\fbox{%
\begin{minipage}{0.97\linewidth}
\small
\begin{tabular}{@{}>{\raggedright\arraybackslash}p{0.359\linewidth}@{\hspace{0.030\linewidth}}>{\raggedright\arraybackslash}p{0.611\linewidth}@{}}
\begin{minipage}[t]{\linewidth}
\centering\textbf{Earlier Germanic changes}\par
\vspace{0.35em}
\raggedright
\centering\textbf{Northwest and West Germanic}\par
\raggedright
\vspace{0.2em}
\begin{tabularx}{\linewidth}{@{}>{\raggedright\arraybackslash}X@{\hspace{0.55em}}>{\raggedright\arraybackslash}p{0.14\linewidth}@{\hspace{0.25em}}}
\mbox{PWGmc Early I Apocope} & \emph{*xábēθ} \\
\end{tabularx}
\vspace{0.6em}
\centering\textbf{Early Anglo-Frisian}\par
\raggedright
\vspace{0.2em}
\raggedright [no change]\par
\end{minipage}
&
\begin{minipage}[t]{\linewidth}
\centering\textbf{Old English changes}\par
\vspace{0.35em}
\raggedright
\begin{tabularx}{\linewidth}{@{}>{\raggedright\arraybackslash}X@{\hspace{0.55em}}>{\raggedright\arraybackslash}p{0.14\linewidth}@{\hspace{0.25em}}}
\mbox{EAF Brightening} & \emph{*xæbēθ} \\
OE Velar Fricative Palatalization & \emph{*çæbēθ} \\
\mbox{PGmc B Allophony} & \emph{*çæβēθ} \\
OE Unstressed Long Vowel Shortening & \emph{*çæβeθ} \\
\end{tabularx}
\end{minipage}
\\
\end{tabular}
\end{minipage}%
} \endgroup

Old English form: \emph{hæfeþ}

\subsection{Reconstruction and comparative
evidence}\label{reconstruction-and-comparative-evidence-130}

The verb belongs to the inherited class-III weak paradigm usually cited
under
\Recon{xabēną}\index[iv]{03pgmc@\ivlangheader{Proto-Germanic}{}!xabena@\iventry{*xabēną}{}}
`have' and Old English
{\emph{habban}}\index[iv]{02oe@\ivlangheader{Old English}{West Saxon normalization unmarked}!habban@\iventry{habban}{}}
`have' (\citeproc{ref-Kroonen2013}{Kroonen 2013, 237};
\citeproc{ref-RingeTaylor2014}{Ringe and Taylor 2014, 93}). Within that
paradigm, however, the infinitive and the singular present indicative do
not continue the same stem. Ringe and Taylor distinguish a \emph{-ja-}
stem in the infinitive from a non-geminating -ai- / \emph{-ē-} stem in
the 2sg and 3sg present forms (\citeproc{ref-RingeTaylor2014}{Ringe and
Taylor 2014, 93}).

The derivational input
\Recon{xábēθi}\index[iv]{03pgmc@\ivlangheader{Proto-Germanic}{}!xabethi@\iventry{*xábēθi}{}}
`have' is therefore the 3sg present cell rather than a rephrasing of the
infinitive. For the present analysis, that finite cell is the cleaner
comparator for the inherited non-geminating stem.

\subsection{Old English evidence}\label{old-english-evidence-131}

The ordinary Old English headword is
{\emph{habban}}\index[iv]{02oe@\ivlangheader{Old English}{West Saxon normalization unmarked}!habban@\iventry{habban}{}}
`have' (\citeproc{ref-ClarkHall1960}{Clark Hall 1960, 157}). Campbell's
Anglian paradigm includes unsyncopated 3sg forms of the
{\emph{hæfed}}\index[iv]{02oe@\ivlangheader{Old English}{West Saxon normalization unmarked}!haefed~04@\iventry{hæfed}{Angl.}}
`has' type, and the present paradigm therefore shows forms of the
\emph{hæf-} type that support the normalized target
{\emph{hæfeþ}}\index[iv]{02oe@\ivlangheader{Old English}{West Saxon normalization unmarked}!haefeth@\iventry{hæfeþ}{}}
`has' (\citeproc{ref-Campbell1959}{Campbell 1959, § 762}).

The target form is therefore a normalized finite cell rather than the
ordinary dictionary lemma. It represents the inherited non-geminating
present stem more directly than
{\emph{habban}}\index[iv]{02oe@\ivlangheader{Old English}{West Saxon normalization unmarked}!habban@\iventry{habban}{}}
`have' does.

\subsection{Development to Old
English}\label{development-to-old-english-129}

From
\Recon{xábēθi}\index[iv]{03pgmc@\ivlangheader{Proto-Germanic}{}!xabethi@\iventry{*xábēθi}{}}
`have', the finite form yields
{\emph{hæfeþ}}\index[iv]{02oe@\ivlangheader{Old English}{West Saxon normalization unmarked}!haefeth@\iventry{hæfeþ}{}}
`has' regularly. Ringe and Taylor discuss this non-geminating present
stem under
{\emph{habban}}\index[iv]{02oe@\ivlangheader{Old English}{West Saxon normalization unmarked}!habban@\iventry{habban}{}}
`have' (\citeproc{ref-RingeTaylor2014}{Ringe and Taylor 2014, 364}).
Campbell's Anglian paradigms include unsyncopated 3sg forms of the
{\emph{hæfeþ}}\index[iv]{02oe@\ivlangheader{Old English}{West Saxon normalization unmarked}!haefeth@\iventry{hæfeþ}{}}
`has' /
{\emph{hæfed}}\index[iv]{02oe@\ivlangheader{Old English}{West Saxon normalization unmarked}!haefed@\iventry{hæfed}{}}
`has' type (\citeproc{ref-Campbell1959}{Campbell 1959, § 762}).

The wider lexeme is less straightforward only because the infinitive
{\emph{habban}}\index[iv]{02oe@\ivlangheader{Old English}{West Saxon normalization unmarked}!habban@\iventry{habban}{}}
`have' shows later leveling. The 3sg present cell is therefore the
closer comparison form.

\subsection{Paradigm comparison}\label{paradigm-comparison-9}

The comparison below sets the relevant forms side by side. It separates
the analogically leveled citation form from the regular 3sg present
line.

{\def\LTcaptype{none} % do not increment counter
\begin{longtable}[]{@{}
  >{\raggedright\arraybackslash}p{(\linewidth - 8\tabcolsep) * \real{0.2000}}
  >{\raggedright\arraybackslash}p{(\linewidth - 8\tabcolsep) * \real{0.2000}}
  >{\raggedright\arraybackslash}p{(\linewidth - 8\tabcolsep) * \real{0.2000}}
  >{\raggedright\arraybackslash}p{(\linewidth - 8\tabcolsep) * \real{0.2000}}
  >{\raggedright\arraybackslash}p{(\linewidth - 8\tabcolsep) * \real{0.2000}}@{}}
\toprule\noalign{}
\begin{minipage}[b]{\linewidth}\raggedright
PGmc cell / interpretation
\end{minipage} & \begin{minipage}[b]{\linewidth}\raggedright
Candidate input
\end{minipage} & \begin{minipage}[b]{\linewidth}\raggedright
Expected or documented OE outcome
\end{minipage} & \begin{minipage}[b]{\linewidth}\raggedright
OE comparison form
\end{minipage} & \begin{minipage}[b]{\linewidth}\raggedright
Result
\end{minipage} \\
\midrule\noalign{}
\endhead
\bottomrule\noalign{}
\endlastfoot
citation infinitive & \emph{-ja-} stem of
{\emph{*xabēną}}\index[iv]{03pgmc@\ivlangheader{Proto-Germanic}{}!xabena@\iventry{*xabēną}{}}
& citation form
{\emph{habban}}\index[iv]{02oe@\ivlangheader{Old English}{West Saxon normalization unmarked}!habban@\iventry{habban}{}}
&
{\emph{habban}}\index[iv]{02oe@\ivlangheader{Old English}{West Saxon normalization unmarked}!habban@\iventry{habban}{}}
& important headword, but shaped by later leveling \\
3sg present &
{\emph{*xábēθi}}\index[iv]{03pgmc@\ivlangheader{Proto-Germanic}{}!xabethi@\iventry{*xábēθi}{}}
& regular output: {\emph{hæfeþ}} &
{\emph{hæfeþ}}\index[iv]{02oe@\ivlangheader{Old English}{West Saxon normalization unmarked}!haefeth@\iventry{hæfeþ}{}}
& exact match between input, output, and finite form compared here \\
syncopated finite tradition & same present stem &
{\emph{hæfþ}}\index[iv]{02oe@\ivlangheader{Old English}{West Saxon normalization unmarked}!haefth@\iventry{hæfþ}{}}
type evidence &
{\emph{hæfþ}}\index[iv]{02oe@\ivlangheader{Old English}{West Saxon normalization unmarked}!haefth@\iventry{hæfþ}{}}
& genuine later OE finite form, but not the normalized target used
here \\
\end{longtable}
}

\section{\texorpdfstring{knee --- OE
\emph{cneowe}}{knee --- OE cneowe}}\label{knee-oe-cneowe}

\index[iv]{02oe@\ivlangheader{Old English}{West Saxon normalization unmarked}!cneo@\iventry{cnēo}{}}
\index[iv]{02oe@\ivlangheader{Old English}{West Saxon normalization unmarked}!cneow@\iventry{cnēow}{}}
\index[iv]{02oe@\ivlangheader{Old English}{West Saxon normalization unmarked}!cneowe@\iventry{cneowe}{}}
\index[iv]{02oe@\ivlangheader{Old English}{West Saxon normalization unmarked}!cneowes@\iventry{cneowes}{}}
\index[iv]{03pgmc@\ivlangheader{Proto-Germanic}{}!knewa@\iventry{*knéwą}{}}
\index[iv]{03pgmc@\ivlangheader{Proto-Germanic}{}!knewai@\iventry{*knéwai}{}}

Derivation: citation reconstruction \emph{*knéwą}; form followed here
\emph{*knéwai} \textgreater{} \emph{cneowe} (late analogy).

\subsection{Derivation trace}\label{derivation-trace-148}

Proto input: \emph{*knéwai}

\begingroup
\setlength{\fboxsep}{6pt}

\noindent\fbox{%
\begin{minipage}{0.97\linewidth}
\small
\begin{tabular}{@{}>{\raggedright\arraybackslash}p{0.485\linewidth}@{\hspace{0.030\linewidth}}>{\raggedright\arraybackslash}p{0.485\linewidth}@{}}
\begin{minipage}[t]{\linewidth}
\centering\textbf{Earlier Germanic changes}\par
\vspace{0.35em}
\raggedright
\centering\textbf{Northwest and West Germanic}\par
\raggedright
\vspace{0.2em}
\begin{tabularx}{\linewidth}{@{}>{\raggedright\arraybackslash}X@{\hspace{0.55em}}>{\raggedright\arraybackslash}p{0.14\linewidth}@{\hspace{0.25em}}}
PNWGmc Unstressed Ai Monophthongization & \emph{*knéwē} \\
\end{tabularx}
\vspace{0.6em}
\centering\textbf{Early Anglo-Frisian}\par
\raggedright
\vspace{0.2em}
\raggedright [no change]\par
\end{minipage}
&
\begin{minipage}[t]{\linewidth}
\centering\textbf{Old English changes}\par
\vspace{0.35em}
\raggedright
\begin{tabularx}{\linewidth}{@{}>{\raggedright\arraybackslash}X@{\hspace{0.55em}}>{\raggedright\arraybackslash}p{0.14\linewidth}@{\hspace{0.25em}}}
\mbox{OE Breaking} & \emph{*knéowē} \\
OE Unstressed Long Vowel Shortening & \emph{*knéowe} \\
\end{tabularx}
\end{minipage}
\\
\end{tabular}
\end{minipage}%
} \endgroup

Old English form: \emph{cneowe}

\subsection{Reconstruction and selection of the
dative}\label{reconstruction-and-selection-of-the-dative}

The citation reconstruction is the neuter noun
\Recon{knéwą}\index[iv]{03pgmc@\ivlangheader{Proto-Germanic}{}!knewa@\iventry{*knéwą}{}}
`knee'. Ringe and Taylor expressly give the singular \Recon{knewą}
`knee' and distinguish its endingless development from the inflected
stem (\citeproc{ref-RingeTaylor2014}{Ringe and Taylor 2014, 387},
§7.2.4). The acute in the selected notation marks root stress; it does
not make the vowel long.

The derivational input
\Recon{knéwai}\index[iv]{03pgmc@\ivlangheader{Proto-Germanic}{}!knewai@\iventry{*knéwai}{}}
`knee (dat.sg.)' belongs to another cell of the same inherited paradigm.
It combines the stem with the conventional dative ending. Fulk regards
Proto-Germanic \emph{-ai} as the most likely source of the West Germanic
dative singular, while considering the competing \emph{-ē} analysis and
the uncertainties of the wider Germanic comparison
(\citeproc{ref-Fulk2018}{Fulk 2018, 147}, §7.8). The whole-word input is
therefore an explicit paradigm construction under that defended working
account, not a verbatim reconstruction quoted from Hogg or Campbell. It
is neither a later Old English transponent nor a nominative from which
the dative has been generated by an invented sound law.

The distinction between citation reconstruction and derivational input
is essential. Selecting the dative does not deny the existence of the
familiar nominative \emph{cnēow} `knee', nor replace the inherited
neuter reconstruction with a different lexeme. It selects the cell in
which the regular pre-ending development can be compared without
importing the nominative's analogically restored consonant.

\subsection{Old English quantity and paradigm
evidence}\label{old-english-quantity-and-paradigm-evidence}

The citation and dative are different cells, not rival reconstructions
of one cell. The source comparison accordingly adds an uncertainty about
the inherited ending, not another dispute over the already documented
short dative. Fulk's preferred \emph{-ai} account licenses the selected
paradigm construction without establishing a uniquely recoverable
complete ancestor (\citeproc{ref-Fulk2018}{Fulk 2018, 147}). The
detailed quantity and levelling evidence below remains independently
necessary.

Hogg explicitly uses
{\emph{cneowe}}\index[iv]{02oe@\ivlangheader{Old English}{West Saxon normalization unmarked}!cneowe@\iventry{cneowe}{}}
`knee (dat.sg.)' as an example of short \emph{eo}. His adjacent
long-diphthong example, \emph{cnēowe} `know (past subjunctive)', is a
verb form, not an alternative knee dative. In the discussion of breaking
he again names the knee dative as a regular example before \emph{w},
while assigning the final \emph{w} of nominative \emph{cnēo(w)} to
analogy {[}Hogg (\citeproc{ref-HoggGrammar2011}{2011}), pp.~21--22,
§2.33; p.~86, §5.22 and note 7{]}. The quantity contrast is an explicit
grammatical analysis, not an inference from an unmarked manuscript
digraph.

Campbell's regular paradigm similarly has genitive \emph{cneowes} `knee
(gen.sg.)' and dative \emph{cneowe} beside long endingless
\emph{cnēo(w)}. He then allows the long diphthong to extend into the
inflected forms, giving the secondary \emph{cnēowes} type
(\citeproc{ref-Campbell1959}{Campbell 1959, 232--33}, §584). Luick
independently gives \emph{cneowe(s)} as the genitive and dative products
of short-vowel diphthongization before \emph{w}, with the
following-\emph{i} exclusion (\citeproc{ref-Luick1914}{Luick
1914-\/-1940, 139}, §134). These accounts support the selected short
dative. They do not establish that every occurrence of an inflected knee
form necessarily retained the original quantity.

Ringe and Taylor give the same basic opposition: long endingless
\emph{cnēo} and short \emph{cneow-} before syllabic endings. They
cautiously consider subsequent extension of long quantity into knee's
inflected stem, noting the difficulty of finding decisive verse evidence
{[}Ringe and Taylor (\citeproc{ref-RingeTaylor2014}{2014}),
pp.~187--188, §6.2.4; p.~387, §7.2.4{]}. Fulk describes the wider
two-way leveling pattern of this stem class: the endingless diphthong
spreads into inflected forms, and their \emph{w} spreads into
uninflected forms. His explicit illustrative example is the servant
paradigm, not an independent metrical attestation for knee
(\citeproc{ref-Fulk2018}{Fulk 2018, 154}, §7.12).

\subsection{The regular selected
development}\label{the-regular-selected-development}

Before a retained syllabic ending, the root vowel remains short. The
singleton \emph{w} conditions short \emph{e} to \emph{eo}, unless
followed by a high front vocalic. Hogg's knee dative and Luick's
corresponding examples belong to this regular pre-ending domain
(\citeproc{ref-HoggGrammar2011}{Hogg 2011, 86};
\citeproc{ref-Luick1914}{Luick 1914-\/-1940, 139};
\citeproc{ref-RingeTaylor2014}{Ringe and Taylor 2014, 187--88}). The
weakening of the selected dative ending leaves Old English \emph{-e};
the medial \emph{w} has not become a word-final offglide and does not
need to be restored.

The modeled short diphthong is supplied by
\hyperref[rule-OEBreaking]{SC044 OEBreaking}, not by assigning long
quantity at \hyperref[rule-OEEwLongDiphthong]{SC033 OEEwLongDiphthong}.
The latter now excludes this singleton path and retains only the
separately disclosed j-created representation discussed with \emph{hīew}
`hue'. The inherited-ending and short-diphthong developments are
independently motivated; arriving at the right final spelling is not
itself their historical justification.

\subsection{Why the nominative is a comparison, not the selected
target}\label{why-the-nominative-is-a-comparison-not-the-selected-target}

The endingless development follows a different route. Ringe and Taylor
derive long \emph{cnēo} from the inherited singular through an
offglide-bearing intermediate, and explain the \emph{w} in \emph{cnēow}
as transfer from the inflected forms
(\citeproc{ref-RingeTaylor2014}{Ringe and Taylor 2014, 174, 387}).
Campbell likewise distinguishes endingless diphthong formation from the
later addition of \emph{w} from the obliques
(\citeproc{ref-Campbell1959}{Campbell 1959, 46--47, 232--33}). Luick's
older account distinguishes the endingless knee/tree \emph{eo} and
\emph{eu} sources and their subsequent leveling
(\citeproc{ref-Luick1914}{Luick 1914-\/-1940, 118}, §101). Agreement on
the short/long paradigm opposition does not make all these detailed
early reconstructions identical.

The nominative is thus not an appropriate obligation for a purely
phonological derivation that retains \emph{w} and lengthens the root
before its ending disappears. Such a derivation reaches a real Old
English form by conflating two histories. The selected dative avoids
that conflation without a knee-specific lengthening rule, an analogical
sound change, or grammatical conditioning of a phonetic law. It also
avoids choosing a long oblique which itself reflects the quantity
leveling under investigation.

\subsection{Paradigm comparison and
confidence}\label{paradigm-comparison-and-confidence}

{\def\LTcaptype{none} % do not increment counter
\begin{longtable}[]{@{}
  >{\raggedright\arraybackslash}p{(\linewidth - 4\tabcolsep) * \real{0.3333}}
  >{\raggedright\arraybackslash}p{(\linewidth - 4\tabcolsep) * \real{0.3333}}
  >{\raggedright\arraybackslash}p{(\linewidth - 4\tabcolsep) * \real{0.3333}}@{}}
\toprule\noalign{}
\begin{minipage}[b]{\linewidth}\raggedright
Cell or interpretation
\end{minipage} & \begin{minipage}[b]{\linewidth}\raggedright
Reconstruction / form
\end{minipage} & \begin{minipage}[b]{\linewidth}\raggedright
Historical role
\end{minipage} \\
\midrule\noalign{}
\endhead
\bottomrule\noalign{}
\endlastfoot
PGmc citation singular & *knéwą & Preserved lexeme-level
reconstruction \\
Selected PGmc dative singular & *knéwai & Explicit inherited-cell
construction under the preferred *-ai analysis \\
OE selected dative singular & cneowe & Regular short pre-ending
diphthong \\
OE regular genitive singular & cneowes & Independent short-stem
comparison \\
OE endingless singular & cnēo & Regular long offglide-derived form \\
OE familiar nominative & cnēow & Endingless form with restored w \\
Later long obliques & cnēowe, cnēowes & Possible quantity leveling; not
the selected regular comparison \\
\end{longtable}
}

Confidence is strong in the regular paradigm distinction: Campbell,
Luick, Hogg and Ringe--Taylor independently identify the short
pre-ending stem. Confidence is more qualified in the exact inherited
dative-ending reconstruction and in how widely later long quantity
spread through knee's attested obliques. Neither uncertainty authorizes
an unrestricted short-\emph{ew} to long-\emph{ēow} sound law.

The authors' exact inputs are not silently harmonized. Hogg's breaking
footnote prints
\Recon{knewaz}\index[iv]{03pgmc@\ivlangheader{Proto-Germanic}{}!knewaz@\iventry{*knewaz}{}}
`knee', whereas Ringe and Taylor give the neuter \emph{knewą}
(\citeproc{ref-HoggGrammar2011}{Hogg 2011, 86}, note 7;
\citeproc{ref-RingeTaylor2014}{Ringe and Taylor 2014, 387}). The latter
supplies the retained citation reconstruction. This discrepancy in the
quoted ending is separate from their agreement on short pre-ending
quantity; it is not evidence for replacing the selected dative with a
nominative or for claiming identical whole-word inputs.

This entry belongs to late analogy and paradigm-cell selection because
the conventional citation paradigm has undergone later reshaping. That
placement does not mean that the selected dative itself requires an
analogical operation. It follows the same principle as the regular
selected obliques used for cow, night and meed: distinguish the history
of the lexeme's paradigm from the sound-law path of the chosen cell.

\section{\texorpdfstring{learn (3sg) --- OE
\emph{leornaþ}}{learn (3sg) --- OE leornaþ}}\label{learn-3sg-oe-leornauxfe}

\index[iv]{02oe@\ivlangheader{Old English}{West Saxon normalization unmarked}!leornath@\iventry{leornaþ}{}}
\index[iv]{03pgmc@\ivlangheader{Proto-Germanic}{}!liznojana@\iventry{*liznōjaną}{}}
\index[iv]{03pgmc@\ivlangheader{Proto-Germanic}{}!liznothi@\iventry{*líznōθi}{}}

Derivation: citation reconstruction \emph{*liznōjaną}; form followed
here \emph{*líznōθi} \textgreater{} \emph{leornaþ} (late analogy).

\subsection{Derivation trace}\label{derivation-trace-149}

Proto input: \emph{*líznōθi}

\begingroup
\setlength{\fboxsep}{6pt}

\noindent\fbox{%
\begin{minipage}{0.97\linewidth}
\small
\begin{tabular}{@{}>{\raggedright\arraybackslash}p{0.485\linewidth}@{\hspace{0.030\linewidth}}>{\raggedright\arraybackslash}p{0.485\linewidth}@{}}
\begin{minipage}[t]{\linewidth}
\centering\textbf{Earlier Germanic changes}\par
\vspace{0.35em}
\raggedright
\centering\textbf{Northwest and West Germanic}\par
\raggedright
\vspace{0.2em}
\begin{tabularx}{\linewidth}{@{}>{\raggedright\arraybackslash}X@{\hspace{0.55em}}>{\raggedright\arraybackslash}p{0.14\linewidth}@{\hspace{0.25em}}}
\mbox{PWGmc Early I Apocope} & \emph{*líznōθ} \\
\mbox{PNWGmc I Lowering} & \emph{*léznōθ} \\
\mbox{EAF Rhotacism} & \emph{*lérnōθ} \\
\end{tabularx}
\vspace{0.6em}
\centering\textbf{Early Anglo-Frisian}\par
\raggedright
\vspace{0.2em}
\raggedright [no change]\par
\end{minipage}
&
\begin{minipage}[t]{\linewidth}
\centering\textbf{Old English changes}\par
\vspace{0.35em}
\raggedright
\begin{tabularx}{\linewidth}{@{}>{\raggedright\arraybackslash}X@{\hspace{0.55em}}>{\raggedright\arraybackslash}p{0.18\linewidth}@{\hspace{0.25em}}}
\mbox{OE Breaking} & \emph{*léornōθ} \\
\mbox{OE Late O Shortening} & \emph{*léornoθ} \\
OE Final Unstressed O Lowering & \emph{*léornaθ} \\
\end{tabularx}
\end{minipage}
\\
\end{tabular}
\end{minipage}%
} \endgroup

Old English form: \emph{leornaþ}

\subsection{A textually attested third
singular}\label{a-textually-attested-third-singular}

The selected Old English
{\emph{leornaþ}}\index[iv]{02oe@\ivlangheader{Old English}{West Saxon normalization unmarked}!leornath@\iventry{leornaþ}{}}
`learns' corresponds to Miller's printed \emph{leornað} `learns'. In the
Bede passage, \emph{swa hwylc swa} `whoever' governs coordinated
\emph{rædeð} `reads' and \emph{leornað}. The singular
indefinite-relative subject establishes the person and number; the
ending alone would not suffice. Final eth is represented by thorn
without changing the vowel or cell (\citeproc{ref-Miller1891}{Miller
1891, 380}, ll. 26--27).

The input
\Recon{líznōθi}\index[iv]{03pgmc@\ivlangheader{Proto-Germanic}{}!liznothi@\iventry{*líznōθi}{}}
`learns' is a finite third-singular selection, not the infinitive
citation label. Its inherited \emph{*lizn-} root remains distinct from
the later lowered intermediate. The wider family has a class III
prehistory and subsequent conjugational remodelling; the English class
II form does not recover that complete earlier paradigm
(\citeproc{ref-Kroonen2013}{Kroonen 2013, 340};
\citeproc{ref-Orel2003}{Orel 2003, 247};
\citeproc{ref-Jasanoff1973}{Jasanoff 1973, 865, 868}).

\hyperref[rule-PNWGmcILowering]{SC015 PNWGmcILowering} lowers the
selected root before \emph{*zn}. Rhotacism, English breaking, ending
reduction and dental spelling then give \emph{leornaþ}. This explanation
leaves the inherited input high-vowelled. It does not select an
\emph{*e}-grade just because the Old English word contains \emph{eo}
(\citeproc{ref-Hill2017}{Hill 2017, 153};
\citeproc{ref-RingeTaylor2014}{Ringe and Taylor 2014, 38, 84}).

The high-vowel *liorn- tradition remains important, but the infinitive's
spelling cannot establish the prehistory of this specific cell.
Jasanoff's proposed lost \emph{*lizniþ} and Hill's later-mutation
account identify possible morphological sources of high vowels; neither
yields a fully certain inherited paradigm. The family classification
accordingly retains conjugational remodelling without claiming that the
selected vowel itself requires analogical replacement
(\citeproc{ref-Jasanoff1973}{Jasanoff 1973, 865, 868};
\citeproc{ref-Hill2017}{Hill 2017, 153}, n.~25;
\citeproc{ref-TarsiLaker2026}{Tarsi and Laker 2026, 42--43}).

\section{\texorpdfstring{learn (iptv.2sg) --- OE
\emph{leorna}}{learn (iptv.2sg) --- OE leorna}}\label{learn-iptv.2sg-oe-leorna}

\index[iv]{02oe@\ivlangheader{Old English}{West Saxon normalization unmarked}!leorna@\iventry{leorna}{}}
\index[iv]{03pgmc@\ivlangheader{Proto-Germanic}{}!lizno@\iventry{*líznô}{}}
\index[iv]{03pgmc@\ivlangheader{Proto-Germanic}{}!liznojana@\iventry{*liznōjaną}{}}

Derivation: citation reconstruction \emph{*liznōjaną}; form followed
here \emph{*líznô} \textgreater{} \emph{leorna} (late analogy).

\subsection{Derivation trace}\label{derivation-trace-150}

Proto input: \emph{*líznô}

\begingroup
\setlength{\fboxsep}{6pt}

\noindent\fbox{%
\begin{minipage}{0.97\linewidth}
\small
\begin{tabular}{@{}>{\raggedright\arraybackslash}p{0.485\linewidth}@{\hspace{0.030\linewidth}}>{\raggedright\arraybackslash}p{0.485\linewidth}@{}}
\begin{minipage}[t]{\linewidth}
\centering\textbf{Earlier Germanic changes}\par
\vspace{0.35em}
\raggedright
\centering\textbf{Northwest and West Germanic}\par
\raggedright
\vspace{0.2em}
\begin{tabularx}{\linewidth}{@{}>{\raggedright\arraybackslash}X@{\hspace{0.55em}}>{\raggedright\arraybackslash}p{0.14\linewidth}@{\hspace{0.25em}}}
\mbox{PNWGmc I Lowering} & \emph{*léznô} \\
\mbox{EAF Rhotacism} & \emph{*lérnô} \\
\end{tabularx}
\vspace{0.6em}
\centering\textbf{Early Anglo-Frisian}\par
\raggedright
\vspace{0.2em}
\raggedright [no change]\par
\end{minipage}
&
\begin{minipage}[t]{\linewidth}
\centering\textbf{Old English changes}\par
\vspace{0.35em}
\raggedright
\begin{tabularx}{\linewidth}{@{}>{\raggedright\arraybackslash}X@{\hspace{0.55em}}>{\raggedright\arraybackslash}p{0.14\linewidth}@{\hspace{0.25em}}}
\mbox{OE Breaking} & \emph{*léornô} \\
OE Unstressed Long Vowel Shortening & \emph{*léorna} \\
\end{tabularx}
\end{minipage}
\\
\end{tabular}
\end{minipage}%
} \endgroup

Old English form: \emph{leorna}

\subsection{The imperative singular}\label{the-imperative-singular}

Old English
{\emph{leorna}}\index[iv]{02oe@\ivlangheader{Old English}{West Saxon normalization unmarked}!leorna@\iventry{leorna}{}}
`learn!' is the second-person singular imperative, not an abbreviated
infinitive. Deuteronomy14:23 prints the command followed by \emph{þæt þu
ondræde Drihten} `that you fear the Lord'. The addressee is singular,
and the form is read directly in the edited biblical text
(\citeproc{ref-Crawford1922}{Crawford 1922, 351}).

The derivational input
\Recon{líznô}\index[iv]{03pgmc@\ivlangheader{Proto-Germanic}{}!lizno@\iventry{*líznô}{}}
`learn!' is the imperative cell, distinct from the lexeme-level
infinitive label. The inherited \emph{*lizn-} family is supported by the
comparative dictionaries; its historical weak class III formation and
later remodelling remain uncertain (\citeproc{ref-Kroonen2013}{Kroonen
2013, 340}; \citeproc{ref-Orel2003}{Orel 2003, 247};
\citeproc{ref-Jasanoff1973}{Jasanoff 1973, 865, 868}).

Under \hyperref[rule-PNWGmcILowering]{SC015 PNWGmcILowering}, lowering
before preconsonantal \emph{*z} produces \emph{*léznô}; rhotacism gives
\emph{*lérnô}, and English breaking and ending reduction lead to
\emph{leorna}. The companion infinitive \emph{leornian} and present
third singular \emph{leornaþ} are different cells of the same family.
Their agreement does not multiply the independent etymological evidence
(\citeproc{ref-Hill2017}{Hill 2017, 153};
\citeproc{ref-Crawford1922}{Crawford 1922, 71, 351};
\citeproc{ref-Miller1891}{Miller 1891, 380}).

The attested alternative \emph{liorna} `learn!' remains relevant to the
*liorn-/leorn- distribution. Mutation in an earlier paradigm, later
levelling and dialect history must be considered rather than identifying
every high vowel as unchanged inheritance. The retained family
classification marks conjugational remodelling, not finite-cell
selection or an extra replacement needed for this imperative
(\citeproc{ref-Campbell1959}{Campbell 1959, 48, 59, 125};
\citeproc{ref-Jasanoff1973}{Jasanoff 1973, 865, 868}).

\section{\texorpdfstring{lick (3sg) --- OE
\emph{liccaþ}}{lick (3sg) --- OE liccaþ}}\label{lick-3sg-oe-liccauxfe}

\index[iv]{02oe@\ivlangheader{Old English}{West Saxon normalization unmarked}!liccath@\iventry{liccaþ}{}}
\index[iv]{03pgmc@\ivlangheader{Proto-Germanic}{}!likkona@\iventry{*likkōną}{}}
\index[iv]{03pgmc@\ivlangheader{Proto-Germanic}{}!likkothi@\iventry{*líkkōθi}{}}

Derivation: citation reconstruction \emph{*likkōną}; form followed here
\emph{*líkkōθi} \textgreater{} \emph{liccaþ} (late analogy).

\subsection{Derivation trace}\label{derivation-trace-151}

Proto input: \emph{*líkkōθi}

\begingroup
\setlength{\fboxsep}{6pt}

\noindent\fbox{%
\begin{minipage}{0.97\linewidth}
\small
\begin{tabular}{@{}>{\raggedright\arraybackslash}p{0.359\linewidth}@{\hspace{0.030\linewidth}}>{\raggedright\arraybackslash}p{0.611\linewidth}@{}}
\begin{minipage}[t]{\linewidth}
\centering\textbf{Earlier Germanic changes}\par
\vspace{0.35em}
\raggedright
\centering\textbf{Northwest and West Germanic}\par
\raggedright
\vspace{0.2em}
\begin{tabularx}{\linewidth}{@{}>{\raggedright\arraybackslash}X@{\hspace{0.55em}}>{\raggedright\arraybackslash}p{0.24\linewidth}@{\hspace{0.25em}}}
\mbox{PWGmc Early I Apocope} & \emph{*líkkōθ} \\
\end{tabularx}
\vspace{0.6em}
\centering\textbf{Early Anglo-Frisian}\par
\raggedright
\vspace{0.2em}
\raggedright [no change]\par
\end{minipage}
&
\begin{minipage}[t]{\linewidth}
\centering\textbf{Old English changes}\par
\vspace{0.35em}
\raggedright
\begin{tabularx}{\linewidth}{@{}>{\raggedright\arraybackslash}X@{\hspace{0.55em}}>{\raggedright\arraybackslash}p{0.14\linewidth}@{\hspace{0.25em}}}
\mbox{OE Late O Shortening} & \emph{*líkkoθ} \\
OE Final Unstressed O Lowering & \emph{*líkkaθ} \\
\end{tabularx}
\end{minipage}
\\
\end{tabular}
\end{minipage}%
} \endgroup

Old English form: \emph{liccaþ}

\subsection{Reconstruction and comparative
evidence}\label{reconstruction-and-comparative-evidence-131}

Ringe and Taylor give PWGmc \Recon{ekkōn} `lick (3sg.)' continuing as
Old English
{\emph{liccian}}\index[iv]{02oe@\ivlangheader{Old English}{West Saxon normalization unmarked}!liccian@\iventry{liccian}{}}
`lick', Old Saxon
{\emph{likkon}}\index[iv]{13os@\ivlangheader{Old Saxon}{}!likkon@\iventry{likkon}{}}
`lick', and Old High German
{\emph{lecchon}}\index[iv]{10ohg@\ivlangheader{Old High German}{}!lecchon@\iventry{lecchon}{}}
`lick' (\citeproc{ref-RingeTaylor2014}{Ringe and Taylor 2014, 50}). Orel
gives the fuller weak-verb reconstruction
\Recon{likkōjanan}\index[iv]{03pgmc@\ivlangheader{Proto-Germanic}{}!likkojanan@\iventry{*likkōjanan}{}}
`lick' with the same Old English continuation
(\citeproc{ref-Orel2003}{Orel 2003, 285}).

The form compared here in this entry is the non-\emph{j} present third
singular \Recon{líkkōθi} `licks', not the remodeled infinitive. Campbell
states the class-II present endings as lufas, \emph{-aþ} (\textless{}
\emph{-ōsi}, \emph{-ōþi)} (\citeproc{ref-Campbell1959}{Campbell 1959,
§ 356.4}). Ringe and Taylor likewise note that class-II weak present
2sg. \emph{-as(t)} and 3sg. \emph{-aþ} have stable \emph{a}
(\citeproc{ref-RingeTaylor2014}{Ringe and Taylor 2014, 80}).

\subsection{Old English evidence}\label{old-english-evidence-132}

Bosworth-Toller lemmatizes the verb as \emph{liccian} `lick'
(\citeproc{ref-BosworthToller1898}{Bosworth and Toller 1898, 614}). The
same consonantal frame appears in Campbell's and Brunner's grammatical
citations of \emph{liccian} `lick' (\citeproc{ref-Campbell1959}{Campbell
1959, § 398.1}; \citeproc{ref-SieversBrunner1965}{Brunner 1965,
§ 45} Anm. 3). The Old English headword is therefore clear even
though the entry here is not about the citation form.

The form treated here is the present third singular
{\emph{liccaþ}}\index[iv]{02oe@\ivlangheader{Old English}{West Saxon normalization unmarked}!liccath@\iventry{liccaþ}{}}
`licks'. It is a selected paradigm form beside the lemma
{\emph{liccian}}\index[iv]{02oe@\ivlangheader{Old English}{West Saxon normalization unmarked}!liccian@\iventry{liccian}{}}
`lick' and the related imperative
{\emph{licca}}\index[iv]{02oe@\ivlangheader{Old English}{West Saxon normalization unmarked}!licca@\iventry{licca}{}}
`lick', not a separately lemmatized headword.

\subsection{Development to Old
English}\label{development-to-old-english-130}

\Recon{líkkōθi}\index[iv]{03pgmc@\ivlangheader{Proto-Germanic}{}!likkothi@\iventry{*líkkōθi}{}}
`licks' first loses final \emph{-i}, giving \Recon{líkkōθ} `licks'.
Campbell's class-II present endings lufas, \emph{-aþ} (\textless{}
\emph{-ōsi}, \emph{-ōþi)} support the regular 3sg outcome \emph{-aþ}
(\citeproc{ref-Campbell1959}{Campbell 1959, § 356.4}). Ringe and
Taylor likewise note stable \emph{a} in the class-II 2sg and 3sg
(\citeproc{ref-RingeTaylor2014}{Ringe and Taylor 2014, 80}). Because
this ending never contains \emph{-j-}, the form does not pass through an
i-umlauted \emph{-eþ} stage.

\subsection{Paradigm comparison}\label{paradigm-comparison-10}

The comparison below sets the relevant forms side by side.

{\def\LTcaptype{none} % do not increment counter
\begin{longtable}[]{@{}
  >{\raggedright\arraybackslash}p{(\linewidth - 8\tabcolsep) * \real{0.2000}}
  >{\raggedright\arraybackslash}p{(\linewidth - 8\tabcolsep) * \real{0.2000}}
  >{\raggedright\arraybackslash}p{(\linewidth - 8\tabcolsep) * \real{0.2000}}
  >{\raggedright\arraybackslash}p{(\linewidth - 8\tabcolsep) * \real{0.2000}}
  >{\raggedright\arraybackslash}p{(\linewidth - 8\tabcolsep) * \real{0.2000}}@{}}
\toprule\noalign{}
\begin{minipage}[b]{\linewidth}\raggedright
PGmc cell / interpretation
\end{minipage} & \begin{minipage}[b]{\linewidth}\raggedright
Candidate input
\end{minipage} & \begin{minipage}[b]{\linewidth}\raggedright
Old English outcome or comparison
\end{minipage} & \begin{minipage}[b]{\linewidth}\raggedright
OE comparison form
\end{minipage} & \begin{minipage}[b]{\linewidth}\raggedright
Result
\end{minipage} \\
\midrule\noalign{}
\endhead
\bottomrule\noalign{}
\endlastfoot
lexeme-level infinitive &
{\emph{*líkkōjaną}}\index[iv]{03pgmc@\ivlangheader{Proto-Germanic}{}!likkojana@\iventry{*líkkōjaną}{}}
& regular output
{\emph{liccian}}\index[iv]{02oe@\ivlangheader{Old English}{West Saxon normalization unmarked}!liccian@\iventry{liccian}{}}
`lick' &
{\emph{liccian}}\index[iv]{02oe@\ivlangheader{Old English}{West Saxon normalization unmarked}!liccian@\iventry{liccian}{}}
`lick' & ordinary dictionary headword of the verb, but not the finite
form compared here \\
imperative singular companion &
{\emph{*líkkô}}\index[iv]{03pgmc@\ivlangheader{Proto-Germanic}{}!likko@\iventry{*líkkô}{}}
& regular output {\emph{licca}} `lick' &
{\emph{licca}}\index[iv]{02oe@\ivlangheader{Old English}{West Saxon normalization unmarked}!licca@\iventry{licca}{}}
`lick' & useful family control, but not the target of this entry \\
present third singular &
{\emph{*líkkōθi}}\index[iv]{03pgmc@\ivlangheader{Proto-Germanic}{}!likkothi@\iventry{*líkkōθi}{}}
& regular output {\emph{liccaþ}} `licks' &
{\emph{liccaþ}}\index[iv]{02oe@\ivlangheader{Old English}{West Saxon normalization unmarked}!liccath@\iventry{liccaþ}{}}
`licks' & exact match between the derivational input and the Old English
form here \\
\end{longtable}
}

\section{\texorpdfstring{lick (iptv.2sg) --- OE
\emph{licca}}{lick (iptv.2sg) --- OE licca}}\label{lick-iptv.2sg-oe-licca}

\index[iv]{02oe@\ivlangheader{Old English}{West Saxon normalization unmarked}!licca@\iventry{licca}{}}
\index[iv]{03pgmc@\ivlangheader{Proto-Germanic}{}!likko@\iventry{*líkkô}{}}
\index[iv]{03pgmc@\ivlangheader{Proto-Germanic}{}!likkona@\iventry{*likkōną}{}}

Derivation: citation reconstruction \emph{*likkōną}; form followed here
\emph{*líkkô} \textgreater{} \emph{licca} (late analogy).

\subsection{Derivation trace}\label{derivation-trace-152}

Proto input: \emph{*líkkô}

\begingroup
\setlength{\fboxsep}{6pt}

\noindent\fbox{%
\begin{minipage}{0.97\linewidth}
\small
\begin{tabular}{@{}>{\raggedright\arraybackslash}p{0.359\linewidth}@{\hspace{0.030\linewidth}}>{\raggedright\arraybackslash}p{0.611\linewidth}@{}}
\begin{minipage}[t]{\linewidth}
\centering\textbf{Earlier Germanic changes}\par
\vspace{0.35em}
\raggedright
\centering\textbf{Northwest and West Germanic}\par
\raggedright
\vspace{0.2em}
\raggedright [no change]\par
\vspace{0.6em}
\centering\textbf{Early Anglo-Frisian}\par
\raggedright
\vspace{0.2em}
\raggedright [no change]\par
\end{minipage}
&
\begin{minipage}[t]{\linewidth}
\centering\textbf{Old English changes}\par
\vspace{0.35em}
\raggedright
\begin{tabularx}{\linewidth}{@{}>{\raggedright\arraybackslash}X@{\hspace{0.55em}}>{\raggedright\arraybackslash}p{0.14\linewidth}@{\hspace{0.25em}}}
OE Unstressed Long Vowel Shortening & \emph{*líkka} \\
\end{tabularx}
\end{minipage}
\\
\end{tabular}
\end{minipage}%
} \endgroup

Old English form: \emph{licca}

\subsection{Reconstruction and comparative
evidence}\label{reconstruction-and-comparative-evidence-132}

Ringe and Taylor give PWGmc \Recon{ekkōn} `lick (3sg.)' continuing as
Old English
{\emph{liccian}}\index[iv]{02oe@\ivlangheader{Old English}{West Saxon normalization unmarked}!liccian@\iventry{liccian}{}}
`lick', Old Saxon
{\emph{likkon}}\index[iv]{13os@\ivlangheader{Old Saxon}{}!likkon@\iventry{likkon}{}}
`lick', and Old High German
{\emph{lecchon}}\index[iv]{10ohg@\ivlangheader{Old High German}{}!lecchon@\iventry{lecchon}{}}
`lick' (\citeproc{ref-RingeTaylor2014}{Ringe and Taylor 2014, 50}). Orel
gives the fuller weak-verb reconstruction
\Recon{likkōjanan}\index[iv]{03pgmc@\ivlangheader{Proto-Germanic}{}!likkojanan@\iventry{*likkōjanan}{}}
`lick' with the same Old English continuation
(\citeproc{ref-Orel2003}{Orel 2003, 285}).

Campbell's weak class-II discussion gives present forms such as lufas,
\emph{-aþ} (\textless{} \emph{-ōsi}, \emph{-ōþi)}
(\citeproc{ref-Campbell1959}{Campbell 1959, § 356.4}). Ringe and
Taylor likewise note that class-II weak present 2sg. \emph{-as(t)} and
3sg. \emph{-aþ} have stable \emph{a}
(\citeproc{ref-RingeTaylor2014}{Ringe and Taylor 2014, 80}). The form
treated here is therefore not that remodeled infinitive but a finite
cell in bare trimoric \emph{*-ō}.

\subsection{Old English evidence}\label{old-english-evidence-133}

Bosworth-Toller lemmatizes the verb as \emph{liccian} `lick'
(\citeproc{ref-BosworthToller1898}{Bosworth and Toller 1898, 614}).
Campbell cites \emph{liccian} `lick' among Old English forms with
preserved geminate \emph{cc} (\citeproc{ref-Campbell1959}{Campbell 1959,
§ 398.1}). Brunner likewise cites \emph{liccian} `lick'
(\citeproc{ref-SieversBrunner1965}{Brunner 1965, § 45} Anm. 3). The
Old English evidence therefore establishes the verbal headword and its
consonantal frame securely.

The Old English form here in this entry is the imperative singular
{\emph{licca}}\index[iv]{02oe@\ivlangheader{Old English}{West Saxon normalization unmarked}!licca@\iventry{licca}{}}
`lick'. It is a paradigm form chosen beside the headword
{\emph{liccian}}\index[iv]{02oe@\ivlangheader{Old English}{West Saxon normalization unmarked}!liccian@\iventry{liccian}{}}
`lick' and the related present
{\emph{liccaþ}}\index[iv]{02oe@\ivlangheader{Old English}{West Saxon normalization unmarked}!liccath@\iventry{liccaþ}{}}
`licks', not a separately lemmatized citation word.

\subsection{Development to Old
English}\label{development-to-old-english-131}

With the stem \emph{licc-} established, the remaining development is
brief. Campbell's class-II present endings lufas, \emph{-aþ}
(\textless{} \emph{-ōsi}, \emph{-ōþi)} support late \emph{-a} in this
finite cell (\citeproc{ref-Campbell1959}{Campbell 1959, § 356.4}).
Ringe and Taylor likewise note stable \emph{a} in the class-II 2sg and
3sg (\citeproc{ref-RingeTaylor2014}{Ringe and Taylor 2014, 80}). The
same stem consonantism that appears in \emph{liccian} `lick' is
preserved here, giving \emph{cc} throughout the finite form.

\subsection{Paradigm comparison}\label{paradigm-comparison-11}

The comparison below sets the relevant forms side by side.

{\def\LTcaptype{none} % do not increment counter
\begin{longtable}[]{@{}
  >{\raggedright\arraybackslash}p{(\linewidth - 8\tabcolsep) * \real{0.2000}}
  >{\raggedright\arraybackslash}p{(\linewidth - 8\tabcolsep) * \real{0.2000}}
  >{\raggedright\arraybackslash}p{(\linewidth - 8\tabcolsep) * \real{0.2000}}
  >{\raggedright\arraybackslash}p{(\linewidth - 8\tabcolsep) * \real{0.2000}}
  >{\raggedright\arraybackslash}p{(\linewidth - 8\tabcolsep) * \real{0.2000}}@{}}
\toprule\noalign{}
\begin{minipage}[b]{\linewidth}\raggedright
PGmc cell / interpretation
\end{minipage} & \begin{minipage}[b]{\linewidth}\raggedright
Candidate input
\end{minipage} & \begin{minipage}[b]{\linewidth}\raggedright
Old English outcome or comparison
\end{minipage} & \begin{minipage}[b]{\linewidth}\raggedright
OE comparison form
\end{minipage} & \begin{minipage}[b]{\linewidth}\raggedright
Result
\end{minipage} \\
\midrule\noalign{}
\endhead
\bottomrule\noalign{}
\endlastfoot
lexeme-level infinitive &
{\emph{*líkkōjaną}}\index[iv]{03pgmc@\ivlangheader{Proto-Germanic}{}!likkojana@\iventry{*líkkōjaną}{}}
& regular output
{\emph{liccian}}\index[iv]{02oe@\ivlangheader{Old English}{West Saxon normalization unmarked}!liccian@\iventry{liccian}{}}
`lick' &
{\emph{liccian}}\index[iv]{02oe@\ivlangheader{Old English}{West Saxon normalization unmarked}!liccian@\iventry{liccian}{}}
`lick' & ordinary dictionary headword of the verb, but not the finite
form compared here \\
imperative singular &
{\emph{*líkkô}}\index[iv]{03pgmc@\ivlangheader{Proto-Germanic}{}!likko@\iventry{*líkkô}{}}
& regular output {\emph{licca}} `lick' &
{\emph{licca}}\index[iv]{02oe@\ivlangheader{Old English}{West Saxon normalization unmarked}!licca@\iventry{licca}{}}
`lick' & exact match between the derivational input and the Old English
form here \\
present third singular companion &
{\emph{*líkkōθi}}\index[iv]{03pgmc@\ivlangheader{Proto-Germanic}{}!likkothi@\iventry{*líkkōθi}{}}
& regular output
{\emph{liccaþ}}\index[iv]{02oe@\ivlangheader{Old English}{West Saxon normalization unmarked}!liccath@\iventry{liccaþ}{}}
&
{\emph{liccaþ}}\index[iv]{02oe@\ivlangheader{Old English}{West Saxon normalization unmarked}!liccath@\iventry{liccaþ}{}}
& useful family control, but not the target of this entry \\
\end{longtable}
}

\section{\texorpdfstring{live --- OE
\emph{lifeþ}}{live --- OE lifeþ}}\label{live-oe-lifeuxfe}

\index[iv]{02oe@\ivlangheader{Old English}{West Saxon normalization unmarked}!lifeth@\iventry{lifeþ}{}}
\index[iv]{03pgmc@\ivlangheader{Proto-Germanic}{}!libena@\iventry{*libēną}{}}
\index[iv]{03pgmc@\ivlangheader{Proto-Germanic}{}!libethi@\iventry{*líbēθi}{}}

Derivation: citation reconstruction \emph{*libēną}; form followed here
\emph{*líbēθi} \textgreater{} \emph{lifeþ} (late analogy).

\subsection{Derivation trace}\label{derivation-trace-153}

Proto input: \emph{*líbēθi}

\begingroup
\setlength{\fboxsep}{6pt}

\noindent\fbox{%
\begin{minipage}{0.97\linewidth}
\small
\begin{tabular}{@{}>{\raggedright\arraybackslash}p{0.359\linewidth}@{\hspace{0.030\linewidth}}>{\raggedright\arraybackslash}p{0.611\linewidth}@{}}
\begin{minipage}[t]{\linewidth}
\centering\textbf{Earlier Germanic changes}\par
\vspace{0.35em}
\raggedright
\centering\textbf{Northwest and West Germanic}\par
\raggedright
\vspace{0.2em}
\begin{tabularx}{\linewidth}{@{}>{\raggedright\arraybackslash}X@{\hspace{0.55em}}>{\raggedright\arraybackslash}p{0.14\linewidth}@{\hspace{0.25em}}}
\mbox{PWGmc Early I Apocope} & \emph{*líbēθ} \\
\end{tabularx}
\vspace{0.6em}
\centering\textbf{Early Anglo-Frisian}\par
\raggedright
\vspace{0.2em}
\raggedright [no change]\par
\end{minipage}
&
\begin{minipage}[t]{\linewidth}
\centering\textbf{Old English changes}\par
\vspace{0.35em}
\raggedright
\begin{tabularx}{\linewidth}{@{}>{\raggedright\arraybackslash}X@{\hspace{0.55em}}>{\raggedright\arraybackslash}p{0.14\linewidth}@{\hspace{0.25em}}}
\mbox{PGmc B Allophony} & \emph{*líβēθ} \\
OE Unstressed Long Vowel Shortening & \emph{*líβeθ} \\
\end{tabularx}
\end{minipage}
\\
\end{tabular}
\end{minipage}%
} \endgroup

Old English form: \emph{lifeþ}

\subsection{Reconstruction and comparative
evidence}\label{reconstruction-and-comparative-evidence-133}

The inherited verb belongs to the class-III weak family cited by Kroonen
under \emph{*libēn-}, reflected in Old English
{\emph{libban}}\index[iv]{02oe@\ivlangheader{Old English}{West Saxon normalization unmarked}!libban@\iventry{libban}{}}
`live' (\citeproc{ref-Kroonen2013}{Kroonen 2013, 336}). Ringe and Taylor
show that the paradigm also contained a separate 3sg present stem,
continued in late Northumbrian
{\emph{lifed}}\index[iv]{02oe@\ivlangheader{Old English}{West Saxon normalization unmarked}!lifed~06@\iventry{lifed}{North.}}
`lives', which they treat as an archaism
(\citeproc{ref-RingeTaylor2014}{Ringe and Taylor 2014, 364}).

The derivational input
\Recon{líbēθi}\index[iv]{03pgmc@\ivlangheader{Proto-Germanic}{}!libethi@\iventry{*líbēθi}{}}
`live' therefore represents a finite present cell rather than the
citation infinitive. The later lemma tradition also includes remodeled
forms such as
{\emph{lifian}}\index[iv]{02oe@\ivlangheader{Old English}{West Saxon normalization unmarked}!lifian@\iventry{lifian}{}}
`live'.

\subsection{Old English evidence}\label{old-english-evidence-134}

The ordinary lemma tradition centers on
{\emph{libban}}\index[iv]{02oe@\ivlangheader{Old English}{West Saxon normalization unmarked}!libban@\iventry{libban}{}}
`live' and, in later remodeling,
{\emph{lifian}}\index[iv]{02oe@\ivlangheader{Old English}{West Saxon normalization unmarked}!lifian@\iventry{lifian}{}}
`live'. For this entry, however, the relevant comparison form is the
archaic 3sg present attested as
{\emph{lifed}}\index[iv]{02oe@\ivlangheader{Old English}{West Saxon normalization unmarked}!lifed@\iventry{lifed}{}}
`lives', here normalized as
{\emph{lifeþ}}\index[iv]{02oe@\ivlangheader{Old English}{West Saxon normalization unmarked}!lifeth@\iventry{lifeþ}{}}
`lives' (\citeproc{ref-RingeTaylor2014}{Ringe and Taylor 2014, 364};
\citeproc{ref-Campbell1959}{Campbell 1959, § 762}).

The target is thus a normalized finite form, not the ordinary dictionary
lemma. It preserves the older present-stem history more clearly than the
remodeled lemma tradition does.

\subsection{Development to Old
English}\label{development-to-old-english-132}

From
\Recon{líbēθi}\index[iv]{03pgmc@\ivlangheader{Proto-Germanic}{}!libethi@\iventry{*líbēθi}{}}
`live', regular reduction of the final syllable and later weakening of
the unstressed vowel yield {\emph{lifeþ}} `lives'. The attested spelling
{\emph{lifed}}\index[iv]{02oe@\ivlangheader{Old English}{West Saxon normalization unmarked}!lifed@\iventry{lifed}{}}
`lives' belongs to the same finite form in late Northumbrian orthography
(\citeproc{ref-Campbell1959}{Campbell 1959, § 762};
\citeproc{ref-RingeTaylor2014}{Ringe and Taylor 2014, 364}).

\subsection{Paradigm comparison}\label{paradigm-comparison-12}

The comparison below sets the relevant forms side by side. It separates
the archaic finite cell from the ordinary infinitival and later
remodeled lemma lines.

{\def\LTcaptype{none} % do not increment counter
\begin{longtable}[]{@{}
  >{\raggedright\arraybackslash}p{(\linewidth - 8\tabcolsep) * \real{0.2000}}
  >{\raggedright\arraybackslash}p{(\linewidth - 8\tabcolsep) * \real{0.2000}}
  >{\raggedright\arraybackslash}p{(\linewidth - 8\tabcolsep) * \real{0.2000}}
  >{\raggedright\arraybackslash}p{(\linewidth - 8\tabcolsep) * \real{0.2000}}
  >{\raggedright\arraybackslash}p{(\linewidth - 8\tabcolsep) * \real{0.2000}}@{}}
\toprule\noalign{}
\begin{minipage}[b]{\linewidth}\raggedright
PGmc cell / interpretation
\end{minipage} & \begin{minipage}[b]{\linewidth}\raggedright
Candidate input
\end{minipage} & \begin{minipage}[b]{\linewidth}\raggedright
Expected or documented OE outcome
\end{minipage} & \begin{minipage}[b]{\linewidth}\raggedright
OE comparison form
\end{minipage} & \begin{minipage}[b]{\linewidth}\raggedright
Result
\end{minipage} \\
\midrule\noalign{}
\endhead
\bottomrule\noalign{}
\endlastfoot
citation infinitive line &
{\emph{*libēną}}\index[iv]{03pgmc@\ivlangheader{Proto-Germanic}{}!libena@\iventry{*libēną}{}}
& OE
{\emph{libban}}\index[iv]{02oe@\ivlangheader{Old English}{West Saxon normalization unmarked}!libban@\iventry{libban}{}}
headword tradition &
{\emph{libban}}\index[iv]{02oe@\ivlangheader{Old English}{West Saxon normalization unmarked}!libban@\iventry{libban}{}}
& establishes the lexeme, but not the Old English form here \\
3sg present &
{\emph{*líbēθi}}\index[iv]{03pgmc@\ivlangheader{Proto-Germanic}{}!libethi@\iventry{*líbēθi}{}}
& regular output: {\emph{lifeþ}}; attested
{\emph{lifed}}\index[iv]{02oe@\ivlangheader{Old English}{West Saxon normalization unmarked}!lifed@\iventry{lifed}{}}
&
{\emph{lifed}}\index[iv]{02oe@\ivlangheader{Old English}{West Saxon normalization unmarked}!lifed@\iventry{lifed}{}},
normalized here as
{\emph{lifeþ}}\index[iv]{02oe@\ivlangheader{Old English}{West Saxon normalization unmarked}!lifeth@\iventry{lifeþ}{}}
& selected archaic finite cell \\
later remodeled present tradition & later class-II-type forms &
{\emph{lifian}}\index[iv]{02oe@\ivlangheader{Old English}{West Saxon normalization unmarked}!lifian@\iventry{lifian}{}}
and related finite remodeling &
{\emph{lifian}}\index[iv]{02oe@\ivlangheader{Old English}{West Saxon normalization unmarked}!lifian@\iventry{lifian}{}}
& genuine OE development, but secondary to the cell compared here \\
\end{longtable}
}

\section{\texorpdfstring{make (3sg) --- OE
\emph{macaþ}}{make (3sg) --- OE macaþ}}\label{make-3sg-oe-macauxfe}

\index[iv]{02oe@\ivlangheader{Old English}{West Saxon normalization unmarked}!macath@\iventry{macaþ}{}}
\index[iv]{03pgmc@\ivlangheader{Proto-Germanic}{}!makona@\iventry{*makōną}{}}
\index[iv]{03pgmc@\ivlangheader{Proto-Germanic}{}!makothi@\iventry{*mákōθi}{}}

Derivation: citation reconstruction \emph{*makōną}; form followed here
\emph{*mákōθi} \textgreater{} \emph{macaþ} (late analogy).

\subsection{Derivation trace}\label{derivation-trace-154}

Proto input: \emph{*mákōθi}

\begingroup
\setlength{\fboxsep}{6pt}

\noindent\fbox{%
\begin{minipage}{0.97\linewidth}
\small
\begin{tabular}{@{}>{\raggedright\arraybackslash}p{0.359\linewidth}@{\hspace{0.030\linewidth}}>{\raggedright\arraybackslash}p{0.611\linewidth}@{}}
\begin{minipage}[t]{\linewidth}
\centering\textbf{Earlier Germanic changes}\par
\vspace{0.35em}
\raggedright
\centering\textbf{Northwest and West Germanic}\par
\raggedright
\vspace{0.2em}
\begin{tabularx}{\linewidth}{@{}>{\raggedright\arraybackslash}X@{\hspace{0.55em}}>{\raggedright\arraybackslash}p{0.14\linewidth}@{\hspace{0.25em}}}
\mbox{PWGmc Early I Apocope} & \emph{*mákōθ} \\
\end{tabularx}
\vspace{0.6em}
\centering\textbf{Early Anglo-Frisian}\par
\raggedright
\vspace{0.2em}
\raggedright [no change]\par
\end{minipage}
&
\begin{minipage}[t]{\linewidth}
\centering\textbf{Old English changes}\par
\vspace{0.35em}
\raggedright
\begin{tabularx}{\linewidth}{@{}>{\raggedright\arraybackslash}X@{\hspace{0.55em}}>{\raggedright\arraybackslash}p{0.14\linewidth}@{\hspace{0.25em}}}
\mbox{EAF Brightening} & \emph{*mækōθ} \\
\mbox{OE A Restoration} & \emph{*makōθ} \\
\mbox{OE Late O Shortening} & \emph{*makoθ} \\
OE Final Unstressed O Lowering & \emph{*makaθ} \\
\end{tabularx}
\end{minipage}
\\
\end{tabular}
\end{minipage}%
} \endgroup

Old English form: \emph{macaþ}

\subsection{Reconstruction and comparative
evidence}\label{reconstruction-and-comparative-evidence-134}

Kroonen derives the Old English verb from
{\emph{*makōjan-}}\index[iv]{03pgmc@\ivlangheader{Proto-Germanic}{}!makojan@\iventry{*makōjan-}{}}
on the make-family base \emph{*maka-}
(\citeproc{ref-Kroonen2013}{Kroonen 2013, 350}). Ringe and Taylor
likewise derive Old English
{\emph{macian}}\index[iv]{02oe@\ivlangheader{Old English}{West Saxon normalization unmarked}!macian@\iventry{macian}{}}
`make' from PWGmc \Recon{makon} `make' through \Recon{mekojan} `make'
(\citeproc{ref-RingeTaylor2014}{Ringe and Taylor 2014, 191}).

The derivational input
\Recon{mákōθi}\index[iv]{03pgmc@\ivlangheader{Proto-Germanic}{}!makothi@\iventry{*mákōθi}{}}
`makes' is therefore a finite 3sg cell of the same family, not the
citation form of the verb.

\subsection{Old English evidence}\label{old-english-evidence-135}

Clark Hall lemmatizes the verb as
{\emph{macian}}\index[iv]{02oe@\ivlangheader{Old English}{West Saxon normalization unmarked}!macian@\iventry{macian}{}}
`make' (\citeproc{ref-ClarkHall1960}{Clark Hall 1960, 193}). The
relevant comparison form here is normalized present-third-singular
{\emph{macaþ}}\index[iv]{02oe@\ivlangheader{Old English}{West Saxon normalization unmarked}!macath@\iventry{macaþ}{}}
`makes', set beside the dictionary headword and the related imperative
singular
{\emph{maca}}\index[iv]{02oe@\ivlangheader{Old English}{West Saxon normalization unmarked}!maca@\iventry{maca}{}}
`make'.

Campbell's class-II paradigm gives ordinary 3sg ending \emph{-aþ}, while
his dialect survey allows secondary \emph{-e-} spellings in some
traditions (\citeproc{ref-Campbell1959}{Campbell 1959, § 356.4};
\citeproc{ref-Campbell1959}{1959, § 757}).
{\emph{Macaþ}}\index[iv]{02oe@\ivlangheader{Old English}{West Saxon normalization unmarked}!macath@\iventry{Macaþ}{}}
`makes' is thus the regular comparison form for the non-\emph{j} 3sg
cell.

\subsection{Development to Old
English}\label{development-to-old-english-133}

After early loss of final \emph{-i}, \Recon{mákōθi} `makes' →
\Recon{mákōθ} `makes' after loss of final -i. Anglo-Frisian brightening
gives \Recon{mækōθ} `makes', but Campbell lists
{\emph{macian}}\index[iv]{02oe@\ivlangheader{Old English}{West Saxon normalization unmarked}!macian@\iventry{macian}{}}
`make' among class-II verbs with restored \emph{a}, so the stem returns
to \emph{mak-} before the ending is reduced
(\citeproc{ref-Campbell1959}{Campbell 1959, § 159}).

The ending then follows ordinary class-II 3sg development: \Recon{makōθ}
`makes' \textgreater{} \Recon{makaθ} `makes' \textgreater{} macaþ\_.
Campbell's\_lufas, -aþ (\textless{} -ōsi, -ōþi)\_and Ringe and Taylor's
discussion of stable class-II\_a\_in finite non-\_j` cells support this
sequence (\citeproc{ref-Campbell1959}{Campbell 1959, § 356.4};
\citeproc{ref-RingeTaylor2014}{Ringe and Taylor 2014, 80}).

\subsection{Paradigm comparison}\label{paradigm-comparison-13}

The comparison below sets the relevant forms side by side. It
distinguishes the selected 3sg cell from the make-family lemma and from
the companion imperative form.

{\def\LTcaptype{none} % do not increment counter
\begin{longtable}[]{@{}
  >{\raggedright\arraybackslash}p{(\linewidth - 8\tabcolsep) * \real{0.2000}}
  >{\raggedright\arraybackslash}p{(\linewidth - 8\tabcolsep) * \real{0.2000}}
  >{\raggedright\arraybackslash}p{(\linewidth - 8\tabcolsep) * \real{0.2000}}
  >{\raggedright\arraybackslash}p{(\linewidth - 8\tabcolsep) * \real{0.2000}}
  >{\raggedright\arraybackslash}p{(\linewidth - 8\tabcolsep) * \real{0.2000}}@{}}
\toprule\noalign{}
\begin{minipage}[b]{\linewidth}\raggedright
PGmc cell / interpretation
\end{minipage} & \begin{minipage}[b]{\linewidth}\raggedright
Candidate input
\end{minipage} & \begin{minipage}[b]{\linewidth}\raggedright
OE outcome or comparison
\end{minipage} & \begin{minipage}[b]{\linewidth}\raggedright
OE comparison form
\end{minipage} & \begin{minipage}[b]{\linewidth}\raggedright
Result
\end{minipage} \\
\midrule\noalign{}
\endhead
\bottomrule\noalign{}
\endlastfoot
lexeme-level infinitive &
{\emph{*makōjaną}}\index[iv]{03pgmc@\ivlangheader{Proto-Germanic}{}!makojana@\iventry{*makōjaną}{}}
& dictionary headword
{\emph{macian}}\index[iv]{02oe@\ivlangheader{Old English}{West Saxon normalization unmarked}!macian@\iventry{macian}{}}
`make' &
{\emph{macian}}\index[iv]{02oe@\ivlangheader{Old English}{West Saxon normalization unmarked}!macian@\iventry{macian}{}}
`make' & family background, but not the cell compared here \\
3sg present &
{\emph{*mákōθi}}\index[iv]{03pgmc@\ivlangheader{Proto-Germanic}{}!makothi@\iventry{*mákōθi}{}}
& regular output {\emph{macaþ}} `makes' &
{\emph{macaþ}}\index[iv]{02oe@\ivlangheader{Old English}{West Saxon normalization unmarked}!macath@\iventry{macaþ}{}}
`makes' & exact match \\
imperative singular companion &
{\emph{*mákô}}\index[iv]{03pgmc@\ivlangheader{Proto-Germanic}{}!mako@\iventry{*mákô}{}}
& related finite form
{\emph{maca}}\index[iv]{02oe@\ivlangheader{Old English}{West Saxon normalization unmarked}!maca@\iventry{maca}{}}
`make' &
{\emph{maca}}\index[iv]{02oe@\ivlangheader{Old English}{West Saxon normalization unmarked}!maca@\iventry{maca}{}}
`make' & useful control, but not the target \\
\end{longtable}
}

\section{\texorpdfstring{make (iptv.2sg) --- OE
\emph{maca}}{make (iptv.2sg) --- OE maca}}\label{make-iptv.2sg-oe-maca}

\index[iv]{02oe@\ivlangheader{Old English}{West Saxon normalization unmarked}!lapian@\iventry{lapian}{}}
\index[iv]{02oe@\ivlangheader{Old English}{West Saxon normalization unmarked}!maca@\iventry{maca}{}}
\index[iv]{03pgmc@\ivlangheader{Proto-Germanic}{}!mako@\iventry{*mákô}{}}
\index[iv]{03pgmc@\ivlangheader{Proto-Germanic}{}!makona@\iventry{*makōną}{}}

Derivation: citation reconstruction \emph{*makōną}; form followed here
\emph{*mákô} \textgreater{} \emph{maca} (late analogy).

\subsection{Derivation trace}\label{derivation-trace-155}

Proto input: \emph{*mákô}

\begingroup
\setlength{\fboxsep}{6pt}

\noindent\fbox{%
\begin{minipage}{0.97\linewidth}
\small
\begin{tabular}{@{}>{\raggedright\arraybackslash}p{0.359\linewidth}@{\hspace{0.030\linewidth}}>{\raggedright\arraybackslash}p{0.611\linewidth}@{}}
\begin{minipage}[t]{\linewidth}
\centering\textbf{Earlier Germanic changes}\par
\vspace{0.35em}
\raggedright
\centering\textbf{Northwest and West Germanic}\par
\raggedright
\vspace{0.2em}
\raggedright [no change]\par
\vspace{0.6em}
\centering\textbf{Early Anglo-Frisian}\par
\raggedright
\vspace{0.2em}
\raggedright [no change]\par
\end{minipage}
&
\begin{minipage}[t]{\linewidth}
\centering\textbf{Old English changes}\par
\vspace{0.35em}
\raggedright
\begin{tabularx}{\linewidth}{@{}>{\raggedright\arraybackslash}X@{\hspace{0.55em}}>{\raggedright\arraybackslash}p{0.14\linewidth}@{\hspace{0.25em}}}
\mbox{EAF Brightening} & \emph{*mækô} \\
\mbox{OE A Restoration} & \emph{*makô} \\
OE Unstressed Long Vowel Shortening & \emph{*maka} \\
\end{tabularx}
\end{minipage}
\\
\end{tabular}
\end{minipage}%
} \endgroup

Old English form: \emph{maca}

\subsection{Reconstruction and comparative
evidence}\label{reconstruction-and-comparative-evidence-135}

The make-family belongs to the Old English class-II weak verbs. Campbell
cites \emph{lapian,}
{\emph{macian}}\index[iv]{02oe@\ivlangheader{Old English}{West Saxon normalization unmarked}!macian@\iventry{macian}{}}
`make' among verbs with restored \emph{a}
(\citeproc{ref-Campbell1959}{Campbell 1959, § 159}). Ringe and
Taylor place the Germanic verb in the same class, comparing West
Germanic continuants such as Old Frisian
{\emph{makia}}\index[iv]{11ofris@\ivlangheader{Old Frisian}{}!makia@\iventry{makia}{}}
`make', Old Saxon
{\emph{makon}}\index[iv]{13os@\ivlangheader{Old Saxon}{}!makon@\iventry{makon}{}}
`make', and Old High German
{\emph{mahhon}}\index[iv]{10ohg@\ivlangheader{Old High German}{}!mahhon@\iventry{mahhon}{}}
`make' (\citeproc{ref-RingeTaylor2014}{Ringe and Taylor 2014, 191}).

The derivational input
\Recon{mákô}\index[iv]{03pgmc@\ivlangheader{Proto-Germanic}{}!mako@\iventry{*mákô}{}}
`make (iptv.2sg)' is not the citation form of the lexeme but a finite
paradigm cell. Ringe and Taylor give the class-II weak imperative
singular as -a \textless{} \emph{*-ō}, which makes this cell the
relevant comparison point for the Old English form treated here
(\citeproc{ref-RingeTaylor2014}{Ringe and Taylor 2014, 314}).

\subsection{Old English evidence}\label{old-english-evidence-136}

The dictionary headword is
{\emph{macian}}\index[iv]{02oe@\ivlangheader{Old English}{West Saxon normalization unmarked}!macian@\iventry{macian}{}}
`make' (\citeproc{ref-ClarkHall1960}{Clark Hall 1960, 193}). The form
compared here in this entry is therefore not the lemma but the
imperative singular
{\emph{maca}}\index[iv]{02oe@\ivlangheader{Old English}{West Saxon normalization unmarked}!maca@\iventry{maca}{}}
`make', chosen as a paradigm form beside the headword
{\emph{macian}}\index[iv]{02oe@\ivlangheader{Old English}{West Saxon normalization unmarked}!macian@\iventry{macian}{}}
`make' and the related finite form
{\emph{macaþ}}\index[iv]{02oe@\ivlangheader{Old English}{West Saxon normalization unmarked}!macath@\iventry{macaþ}{}}
`makes'.

The lexical history is still that of \emph{macian} `make', but the
finite cell isolates the regular outcome of trimoric \emph{*ō} more
cleanly than the citation form does.

\subsection{Development to Old
English}\label{development-to-old-english-134}

From \Recon{mákô} `make (iptv.2sg)', Anglo-Frisian brightening first
gives \Recon{mækô} `make (iptv.2sg)'. Campbell cites \emph{macian}
`make' among the class-II verbs with restored \emph{a}
(\citeproc{ref-Campbell1959}{Campbell 1959, § 159}). Ringe and
Taylor's class-II weak imperative singular -a \textless{} \emph{*-ō}
then supports the later finite ending
(\citeproc{ref-RingeTaylor2014}{Ringe and Taylor 2014, 314}).

The same development explains why earlier fronted forms of the
\emph{mæċa} `make' type do not control the entry. Once trimoric
\emph{*ô} is treated as a back-vocalic trigger for restoration, the
imperative singular falls into line with the broader \emph{macian}
`make' family.

\subsection{Paradigm comparison}\label{paradigm-comparison-14}

A paradigm comparison identifies which finite cell of the make-family
matches the Old English form chosen here. The comparison below sets the
relevant forms side by side.

{\def\LTcaptype{none} % do not increment counter
\begin{longtable}[]{@{}
  >{\raggedright\arraybackslash}p{(\linewidth - 8\tabcolsep) * \real{0.2000}}
  >{\raggedright\arraybackslash}p{(\linewidth - 8\tabcolsep) * \real{0.2000}}
  >{\raggedright\arraybackslash}p{(\linewidth - 8\tabcolsep) * \real{0.2000}}
  >{\raggedright\arraybackslash}p{(\linewidth - 8\tabcolsep) * \real{0.2000}}
  >{\raggedright\arraybackslash}p{(\linewidth - 8\tabcolsep) * \real{0.2000}}@{}}
\toprule\noalign{}
\begin{minipage}[b]{\linewidth}\raggedright
PGmc cell / interpretation
\end{minipage} & \begin{minipage}[b]{\linewidth}\raggedright
Candidate input
\end{minipage} & \begin{minipage}[b]{\linewidth}\raggedright
Old English outcome or comparison
\end{minipage} & \begin{minipage}[b]{\linewidth}\raggedright
OE comparison form
\end{minipage} & \begin{minipage}[b]{\linewidth}\raggedright
Result
\end{minipage} \\
\midrule\noalign{}
\endhead
\bottomrule\noalign{}
\endlastfoot
lexeme-level infinitive &
{\emph{*mákōjaną}}\index[iv]{03pgmc@\ivlangheader{Proto-Germanic}{}!makojana@\iventry{*mákōjaną}{}}
& comparative continuation
{\emph{macian}}\index[iv]{02oe@\ivlangheader{Old English}{West Saxon normalization unmarked}!macian@\iventry{macian}{}}
`make' &
{\emph{macian}}\index[iv]{02oe@\ivlangheader{Old English}{West Saxon normalization unmarked}!macian@\iventry{macian}{}}
`make' & ordinary headword of the verb, but not the finite form compared
here \\
imperative singular &
{\emph{*mákô}}\index[iv]{03pgmc@\ivlangheader{Proto-Germanic}{}!mako@\iventry{*mákô}{}}
& regular output: {\emph{maca}} `make' &
{\emph{maca}}\index[iv]{02oe@\ivlangheader{Old English}{West Saxon normalization unmarked}!maca@\iventry{maca}{}}
`make' & exact match between input, output, and selected paradigm
form \\
present third singular companion &
{\emph{*mákōθi}}\index[iv]{03pgmc@\ivlangheader{Proto-Germanic}{}!makothi@\iventry{*mákōθi}{}}
& comparative companion
{\emph{macaþ}}\index[iv]{02oe@\ivlangheader{Old English}{West Saxon normalization unmarked}!macath@\iventry{macaþ}{}}
&
{\emph{macaþ}}\index[iv]{02oe@\ivlangheader{Old English}{West Saxon normalization unmarked}!macath@\iventry{macaþ}{}}
& useful family control, but not the target of this entry \\
\end{longtable}
}

\section{\texorpdfstring{man --- OE
\emph{mannes}}{man --- OE mannes}}\label{man-oe-mannes}

\index[iv]{02oe@\ivlangheader{Old English}{West Saxon normalization unmarked}!mannes@\iventry{mannes}{}}
\index[iv]{03pgmc@\ivlangheader{Proto-Germanic}{}!mannas@\iventry{*mánnas}{}}
\index[iv]{03pgmc@\ivlangheader{Proto-Germanic}{}!mannaz@\iventry{*mánnaz}{}}

Derivation: citation reconstruction \emph{*mánnaz}; form followed here
\emph{*mánnas} \textgreater{} \emph{mannes} (late analogy).

\subsection{Derivation trace}\label{derivation-trace-156}

Proto input: \emph{*mánnas}

\begingroup
\setlength{\fboxsep}{6pt}

\noindent\fbox{%
\begin{minipage}{0.97\linewidth}
\small
\begin{tabular}{@{}>{\raggedright\arraybackslash}p{0.359\linewidth}@{\hspace{0.030\linewidth}}>{\raggedright\arraybackslash}p{0.611\linewidth}@{}}
\begin{minipage}[t]{\linewidth}
\centering\textbf{Earlier Germanic changes}\par
\vspace{0.35em}
\raggedright
\centering\textbf{Northwest and West Germanic}\par
\raggedright
\vspace{0.2em}
\raggedright [no change]\par
\vspace{0.6em}
\centering\textbf{Early Anglo-Frisian}\par
\raggedright
\vspace{0.2em}
\raggedright [no change]\par
\end{minipage}
&
\begin{minipage}[t]{\linewidth}
\centering\textbf{Old English changes}\par
\vspace{0.35em}
\raggedright
\begin{tabularx}{\linewidth}{@{}>{\raggedright\arraybackslash}X@{\hspace{0.55em}}>{\raggedright\arraybackslash}p{0.14\linewidth}@{\hspace{0.25em}}}
\mbox{EAF Brightening Unstressed} & \emph{*mánnæs} \\
\mbox{OE Unstressed AE Merger} & \emph{*mánnes} \\
\end{tabularx}
\end{minipage}
\\
\end{tabular}
\end{minipage}%
} \endgroup

Old English form: \emph{mannes}

\subsection{Reconstruction and comparative
evidence}\label{reconstruction-and-comparative-evidence-136}

The lexeme-level reconstruction is not uniform. Kroonen cites
\emph{*mannan-}, and Orel has \Recon{mannz} `man'
(\citeproc{ref-Kroonen2013}{Kroonen 2013, 354};
\citeproc{ref-Orel2003}{Orel 2003, 299}). The derivational input
\Recon{mánnas}\index[iv]{03pgmc@\ivlangheader{Proto-Germanic}{}!mannas@\iventry{*mánnas}{}}
`man' belongs to a different level: it is the genitive-singular cell
chosen for the Old English comparison.

The target is not the ordinary citation form. Its medial geminate
precedes the genitive ending.

\subsection{Old English evidence}\label{old-english-evidence-137}

Campbell gives the paradigm
{\emph{mann}}\index[iv]{02oe@\ivlangheader{Old English}{West Saxon normalization unmarked}!mann@\iventry{mann}{}}
`man',
{\emph{man}}\index[iv]{02oe@\ivlangheader{Old English}{West Saxon normalization unmarked}!man@\iventry{man}{}}
`man' /
{\emph{mannes}}\index[iv]{02oe@\ivlangheader{Old English}{West Saxon normalization unmarked}!mannes@\iventry{mannes}{}}
`man's /
{\emph{menn}}\index[iv]{02oe@\ivlangheader{Old English}{West Saxon normalization unmarked}!menn@\iventry{menn}{}}
'men' (\citeproc{ref-Campbell1959}{Campbell 1959, § 621}).
Sievers-Brunner likewise cites
{\emph{man}}\index[iv]{02oe@\ivlangheader{Old English}{West Saxon normalization unmarked}!man@\iventry{man}{}}
`man'
{\emph{mannes}}\index[iv]{02oe@\ivlangheader{Old English}{West Saxon normalization unmarked}!mannes@\iventry{mannes}{}}
`man's (\citeproc{ref-SieversBrunner1965}{Brunner 1965, § 226}). He
also explains that word-final simplification underlies forms such as
{\emph{man}}\index[iv]{02oe@\ivlangheader{Old English}{West Saxon normalization unmarked}!man@\iventry{man}{}}
'man' beside inflected
{\emph{monnes}}\index[iv]{02oe@\ivlangheader{Old English}{West Saxon normalization unmarked}!monnes@\iventry{monnes}{}}
'man's (\citeproc{ref-SieversBrunner1965}{Brunner 1965, § 231}).
Clark Hall keeps the dictionary headword under
{\emph{mann}}\index[iv]{02oe@\ivlangheader{Old English}{West Saxon normalization unmarked}!mann@\iventry{mann}{}}
(\citeproc{ref-ClarkHall1960}{Clark Hall 1960, 197}).

The relevant comparison form is therefore the attested genitive singular
{\emph{mannes}}\index[iv]{02oe@\ivlangheader{Old English}{West Saxon normalization unmarked}!mannes@\iventry{mannes}{}}
`man's, not the citation lemma
{\emph{mann}}\index[iv]{02oe@\ivlangheader{Old English}{West Saxon normalization unmarked}!mann@\iventry{mann}{}}
'man'.

\subsection{Development to Old
English}\label{development-to-old-english-135}

Campbell's paradigm \emph{mann}, \emph{man} / \emph{mannes} /
{\emph{menn}}\index[iv]{02oe@\ivlangheader{Old English}{West Saxon normalization unmarked}!menn@\iventry{menn}{}}
`men' confirms the selected genitive singular \emph{mannes}
(\citeproc{ref-Campbell1959}{Campbell 1959, § 621}). In the present
analysis, \Recon{mánnas} `man' develops through Anglo-Frisian
brightening and later unstressed merger to \emph{mannes}. In this cell
the geminate remains medial before \emph{-es}. The citation form behaves
differently because word-final gemination was simplified in Old English
(\citeproc{ref-SieversBrunner1965}{Brunner 1965, § 231}).

\subsection{Paradigm comparison}\label{paradigm-comparison-15}

The comparison below sets the relevant forms side by side. It separates
the citation-form line from the genitive singular.

{\def\LTcaptype{none} % do not increment counter
\begin{longtable}[]{@{}
  >{\raggedright\arraybackslash}p{(\linewidth - 8\tabcolsep) * \real{0.2000}}
  >{\raggedright\arraybackslash}p{(\linewidth - 8\tabcolsep) * \real{0.2000}}
  >{\raggedright\arraybackslash}p{(\linewidth - 8\tabcolsep) * \real{0.2000}}
  >{\raggedright\arraybackslash}p{(\linewidth - 8\tabcolsep) * \real{0.2000}}
  >{\raggedright\arraybackslash}p{(\linewidth - 8\tabcolsep) * \real{0.2000}}@{}}
\toprule\noalign{}
\begin{minipage}[b]{\linewidth}\raggedright
PGmc cell / interpretation
\end{minipage} & \begin{minipage}[b]{\linewidth}\raggedright
Candidate input
\end{minipage} & \begin{minipage}[b]{\linewidth}\raggedright
Expected or documented OE outcome
\end{minipage} & \begin{minipage}[b]{\linewidth}\raggedright
OE comparison form
\end{minipage} & \begin{minipage}[b]{\linewidth}\raggedright
Result
\end{minipage} \\
\midrule\noalign{}
\endhead
\bottomrule\noalign{}
\endlastfoot
citation nominative singular &
{\emph{*mannăz}}\index[iv]{03pgmc@\ivlangheader{Proto-Germanic}{}!mannaz@\iventry{*mannăz}{}}
& expected citation-form outcome {\emph{man}} &
{\emph{mann}}\index[iv]{02oe@\ivlangheader{Old English}{West Saxon normalization unmarked}!mann@\iventry{mann}{}}
/
{\emph{monn}}\index[iv]{02oe@\ivlangheader{Old English}{West Saxon normalization unmarked}!monn@\iventry{monn}{}}
& establishes the lexeme, but not the Old English form here \\
accusative singular &
{\emph{*manną}}\index[iv]{03pgmc@\ivlangheader{Proto-Germanic}{}!manna@\iventry{*manną}{}}
& expected {\emph{man}} &
{\emph{man}}\index[iv]{02oe@\ivlangheader{Old English}{West Saxon normalization unmarked}!man@\iventry{man}{}}
& same word-final simplification as the nominative \\
dative singular &
{\emph{*mannăi}}\index[iv]{03pgmc@\ivlangheader{Proto-Germanic}{}!mannai@\iventry{*mannăi}{}}
& expected {\emph{manne}} &
{\emph{manne}}\index[iv]{02oe@\ivlangheader{Old English}{West Saxon normalization unmarked}!manne@\iventry{manne}{}}
& preserves medial \emph{nn}, but not the chosen cell \\
genitive singular &
{\emph{*mánnas}}\index[iv]{03pgmc@\ivlangheader{Proto-Germanic}{}!mannas@\iventry{*mánnas}{}}
& regular output: {\emph{mannes}} &
{\emph{mannes}}\index[iv]{02oe@\ivlangheader{Old English}{West Saxon normalization unmarked}!mannes@\iventry{mannes}{}}
& exact match between input, output, and attested comparator \\
\end{longtable}
}

\section{\texorpdfstring{meed --- OE
\emph{meorde}}{meed --- OE meorde}}\label{meed-oe-meorde}

\index[iv]{02oe@\ivlangheader{Old English}{West Saxon normalization unmarked}!meorde@\iventry{meorde}{}}
\index[iv]{03pgmc@\ivlangheader{Proto-Germanic}{}!mizdai@\iventry{*mízdai}{}}
\index[iv]{03pgmc@\ivlangheader{Proto-Germanic}{}!mizdo@\iventry{*mizdō}{}}

Derivation: citation reconstruction \emph{*mizdō}; form followed here
\emph{*mízdai} \textgreater{} \emph{meorde} (late analogy).

\subsection{Derivation trace}\label{derivation-trace-157}

Proto input: \emph{*mízdai}

\begingroup
\setlength{\fboxsep}{6pt}

\noindent\fbox{%
\begin{minipage}{0.97\linewidth}
\small
\begin{tabular}{@{}>{\raggedright\arraybackslash}p{0.485\linewidth}@{\hspace{0.030\linewidth}}>{\raggedright\arraybackslash}p{0.485\linewidth}@{}}
\begin{minipage}[t]{\linewidth}
\centering\textbf{Earlier Germanic changes}\par
\vspace{0.35em}
\raggedright
\centering\textbf{Northwest and West Germanic}\par
\raggedright
\vspace{0.2em}
\begin{tabularx}{\linewidth}{@{}>{\raggedright\arraybackslash}X@{\hspace{0.55em}}>{\raggedright\arraybackslash}p{0.14\linewidth}@{\hspace{0.25em}}}
PNWGmc Unstressed Ai Monophthongization & \emph{*mízdē} \\
\mbox{PNWGmc I Lowering} & \emph{*mézdē} \\
\mbox{EAF Rhotacism} & \emph{*mérdē} \\
\end{tabularx}
\vspace{0.6em}
\centering\textbf{Early Anglo-Frisian}\par
\raggedright
\vspace{0.2em}
\raggedright [no change]\par
\end{minipage}
&
\begin{minipage}[t]{\linewidth}
\centering\textbf{Old English changes}\par
\vspace{0.35em}
\raggedright
\begin{tabularx}{\linewidth}{@{}>{\raggedright\arraybackslash}X@{\hspace{0.55em}}>{\raggedright\arraybackslash}p{0.14\linewidth}@{\hspace{0.25em}}}
\mbox{OE Breaking} & \emph{*méordē} \\
OE Unstressed Long Vowel Shortening & \emph{*méorde} \\
\end{tabularx}
\end{minipage}
\\
\end{tabular}
\end{minipage}%
} \endgroup

Old English form: \emph{meorde}

\subsection{Reconstruction and comparative
evidence}\label{reconstruction-and-comparative-evidence-137}

The lexeme-level reconstruction is
\Recon{mizdō}\index[iv]{03pgmc@\ivlangheader{Proto-Germanic}{}!mizdo@\iventry{*mizdō}{}}
`meed', but the derivational input
\Recon{mízdai}\index[iv]{03pgmc@\ivlangheader{Proto-Germanic}{}!mizdai@\iventry{*mízdai}{}}
`meed' is a dative-singular cell rather than the citation form. The Old
English evidence for the
{\emph{meord}}\index[iv]{02oe@\ivlangheader{Old English}{West Saxon normalization unmarked}!meord@\iventry{meord}{}}
`meed, reward' side is oblique.

The wider history of competing
{\emph{mēd}}\index[iv]{02oe@\ivlangheader{Old English}{West Saxon normalization unmarked}!med@\iventry{mēd}{}}
`meed, reward' remains disputed. Kroonen and Fulk explain it through
some form of \emph{z}-loss and compensatory lengthening
(\citeproc{ref-Kroonen2013}{Kroonen 2013, 370};
\citeproc{ref-Fulk2018}{Fulk 2018, 52}), while Orel keeps a doublet
analysis (\citeproc{ref-Orel2003}{Orel 2003, 272}). The comparison here
concerns the attested oblique line \emph{meorde}.

\subsection{Old English evidence}\label{old-english-evidence-138}

The directly attested forms are obliques:
{\emph{meorde}}\index[iv]{02oe@\ivlangheader{Old English}{West Saxon normalization unmarked}!meorde@\iventry{meorde}{}}
`meed' as a dative singular and
{\emph{meorda}}\index[iv]{02oe@\ivlangheader{Old English}{West Saxon normalization unmarked}!meorda@\iventry{meorda}{}}
`meed' as a genitive plural
(\citeproc{ref-BrightCassidyRingler1971}{Bright 1971, 328};
\citeproc{ref-BosworthToller1898}{Bosworth and Toller 1898, 647}).
Lexicographers reconstruct a bare nominative
{\emph{meord}}\index[iv]{02oe@\ivlangheader{Old English}{West Saxon normalization unmarked}!meord@\iventry{meord}{}}
`meed' from those obliques, while West Saxon prose more commonly shows
the competing doublet \emph{mēd} `meed'
(\citeproc{ref-ClarkHall1960}{Clark Hall 1960, 214};
\citeproc{ref-BosworthToller1898}{Bosworth and Toller 1898, 647}).

The target of this entry is therefore the attested oblique
{\emph{meorde}}\index[iv]{02oe@\ivlangheader{Old English}{West Saxon normalization unmarked}!meorde@\iventry{meorde}{}}
`meed', not the reconstructed lemma
{\emph{meord}}\index[iv]{02oe@\ivlangheader{Old English}{West Saxon normalization unmarked}!meord@\iventry{meord}{}}
`meed' and not the better-known West Saxon citation form \emph{mēd}
`meed'.

\subsection{Development to Old
English}\label{development-to-old-english-136}

Ringe and Taylor give the broader noun history as \Recon{mizdo} `reward'
\textgreater{} \Recon{mizdu} `reward' \textgreater{} OE meord
\textasciitilde{} méd (\citeproc{ref-RingeTaylor2014}{Ringe and Taylor
2014, 84}). The fuller oblique-cell path modeled here spells out the
intermediate monophthongization, lowering, rhotacism, breaking, and
unstressed-shortening steps needed for the selected dative-singular
comparison.

This entry therefore follows the attested oblique line within the
broader development. It does not depend on a full decision about the
history of the competing
{\emph{mēd}}\index[iv]{02oe@\ivlangheader{Old English}{West Saxon normalization unmarked}!med@\iventry{mēd}{}}
`meed, reward' tradition.

The consonantal distinction is substantial. Retained \emph{*z} can
become a rhotic, producing the consonantal group that conditions English
breaking in \emph{meord-. Loss of the sibilant with compensatory
lengthening instead leads toward the long-vowel }mēd* tradition. These
are alternative histories of the consonantal sequence, not two
interchangeable spellings of the selected dative
(\citeproc{ref-Kroonen2013}{Kroonen 2013, 370};
\citeproc{ref-Crist2002}{Crist 2002, 3--6}).

\hyperref[rule-PNWGmcILowering]{SC015 PNWGmcILowering} places lowering
before either continuation: its condition is the original preconsonantal
\emph{*z}, not the later rhotic or a particular following vowel. The
selected \emph{*mízdai} therefore supplies a positive application with a
historically relevant consonant group. It does not independently
establish the entire history of the rival nominative or the absolute
date of every lowering environment (\citeproc{ref-Hill2017}{Hill 2017,
153}; \citeproc{ref-Crist2002}{Crist 2002, 3--6}).

\subsection{Paradigm comparison}\label{paradigm-comparison-16}

The comparison below sets the relevant forms side by side. It
distinguishes the attested oblique target from the broader lemma
history.

{\def\LTcaptype{none} % do not increment counter
\begin{longtable}[]{@{}
  >{\raggedright\arraybackslash}p{(\linewidth - 8\tabcolsep) * \real{0.2000}}
  >{\raggedright\arraybackslash}p{(\linewidth - 8\tabcolsep) * \real{0.2000}}
  >{\raggedright\arraybackslash}p{(\linewidth - 8\tabcolsep) * \real{0.2000}}
  >{\raggedright\arraybackslash}p{(\linewidth - 8\tabcolsep) * \real{0.2000}}
  >{\raggedright\arraybackslash}p{(\linewidth - 8\tabcolsep) * \real{0.2000}}@{}}
\toprule\noalign{}
\begin{minipage}[b]{\linewidth}\raggedright
PGmc cell / interpretation
\end{minipage} & \begin{minipage}[b]{\linewidth}\raggedright
Candidate input
\end{minipage} & \begin{minipage}[b]{\linewidth}\raggedright
Expected or documented OE outcome
\end{minipage} & \begin{minipage}[b]{\linewidth}\raggedright
OE comparison form
\end{minipage} & \begin{minipage}[b]{\linewidth}\raggedright
Result
\end{minipage} \\
\midrule\noalign{}
\endhead
\bottomrule\noalign{}
\endlastfoot
citation nominative singular &
{\emph{*mizdō}}\index[iv]{03pgmc@\ivlangheader{Proto-Germanic}{}!mizdo@\iventry{*mizdō}{}}
& inferred lemma outcome
{\emph{meord}}\index[iv]{02oe@\ivlangheader{Old English}{West Saxon normalization unmarked}!meord@\iventry{meord}{}}
`meed' &
{\emph{meord}}\index[iv]{02oe@\ivlangheader{Old English}{West Saxon normalization unmarked}!meord@\iventry{meord}{}}
`meed' & useful background, but the bare lemma is reconstructed rather
than directly attested \\
selected dative singular &
{\emph{*mízdai}}\index[iv]{03pgmc@\ivlangheader{Proto-Germanic}{}!mizdai@\iventry{*mízdai}{}}
& regular output: {\emph{meorde}} `meed' &
{\emph{meorde}}\index[iv]{02oe@\ivlangheader{Old English}{West Saxon normalization unmarked}!meorde@\iventry{meorde}{}}
`meed' & exact match between derivational input and attested target \\
genitive singular &
{\emph{*mizdōz}}\index[iv]{03pgmc@\ivlangheader{Proto-Germanic}{}!mizdoz@\iventry{*mizdōz}{}}
& regular output: {\emph{meorde}} `meed' &
{\emph{meorde}}\index[iv]{02oe@\ivlangheader{Old English}{West Saxon normalization unmarked}!meorde@\iventry{meorde}{}}
`meed' & converges on the same attested string, but the dat.sg. has the
clearest direct support \\
genitive plural control & plural oblique line & attested
{\emph{meorda}}\index[iv]{02oe@\ivlangheader{Old English}{West Saxon normalization unmarked}!meorda@\iventry{meorda}{}}
`meed' &
{\emph{meorda}}\index[iv]{02oe@\ivlangheader{Old English}{West Saxon normalization unmarked}!meorda@\iventry{meorda}{}}
`meed' & confirms the broader oblique tradition, but not the chosen
singular target \\
\end{longtable}
}

\section{\texorpdfstring{night --- OE
\emph{niht}}{night --- OE niht}}\label{night-oe-niht}

\index[iv]{02oe@\ivlangheader{Old English}{West Saxon normalization unmarked}!niht@\iventry{niht}{}}
\index[iv]{03pgmc@\ivlangheader{Proto-Germanic}{}!naxti@\iventry{*náxti}{}}
\index[iv]{03pgmc@\ivlangheader{Proto-Germanic}{}!naxtz@\iventry{*náxtz}{}}

Derivation: citation reconstruction \emph{*náxtz}; form followed here
\emph{*náxti} \textgreater{} \emph{niht} (late analogy).

\subsection{Derivation trace}\label{derivation-trace-158}

Proto input: \emph{*náxti}

\begingroup
\setlength{\fboxsep}{6pt}

\noindent\fbox{%
\begin{minipage}{0.97\linewidth}
\small
\begin{tabular}{@{}>{\raggedright\arraybackslash}p{0.388\linewidth}@{\hspace{0.030\linewidth}}>{\raggedright\arraybackslash}p{0.582\linewidth}@{}}
\begin{minipage}[t]{\linewidth}
\centering\textbf{Earlier Germanic changes}\par
\vspace{0.35em}
\raggedright
\centering\textbf{Northwest and West Germanic}\par
\raggedright
\vspace{0.2em}
\raggedright [no change]\par
\vspace{0.6em}
\centering\textbf{Early Anglo-Frisian}\par
\raggedright
\vspace{0.2em}
\raggedright [no change]\par
\end{minipage}
&
\begin{minipage}[t]{\linewidth}
\centering\textbf{Old English changes}\par
\vspace{0.35em}
\raggedright
\begin{tabularx}{\linewidth}{@{}>{\raggedright\arraybackslash}X@{\hspace{0.55em}}>{\raggedright\arraybackslash}p{0.14\linewidth}@{\hspace{0.25em}}}
\mbox{EAF Brightening} & \emph{*næxti} \\
\mbox{OE Breaking} & \emph{*neaxti} \\
\mbox{OE I Umlaut} & \emph{*niexti} \\
\mbox{OE Ws Palatal Umlaut} & \emph{*nixti} \\
\mbox{OE High Vowel Apocope} & \emph{*nixt} \\
\end{tabularx}
\end{minipage}
\\
\end{tabular}
\end{minipage}%
} \endgroup

Old English form: \emph{niht}

\subsection{Reconstruction and comparative
evidence}\label{reconstruction-and-comparative-evidence-138}

Ringe and Taylor cite gen.sg. \Recon{nahtiz} `night', dat.sg.
\Recon{nahti} `night', and nom.pl. \Recon{nahtiz} `night' for the
high-vowel side of the paradigm, and derive West Saxon
{\emph{niht}}\index[iv]{02oe@\ivlangheader{Old English}{West Saxon normalization unmarked}!niht@\iventry{niht}{}}
`night' from that side (\citeproc{ref-RingeTaylor2014}{Ringe and Taylor
2014, 240}). The citation reconstruction
\Recon{náxtz}\index[iv]{03pgmc@\ivlangheader{Proto-Germanic}{}!naxtz@\iventry{*náxtz}{}}
`night' therefore belongs to the nominative-like headword, while the
derivational input
\Recon{náxti}\index[iv]{03pgmc@\ivlangheader{Proto-Germanic}{}!naxti@\iventry{*náxti}{}}
`night' represents the dative-singular cell.

The word later became the model for endingless datives. Ringe and Taylor
explicitly explain forms such as
{\emph{dæg}}\index[iv]{02oe@\ivlangheader{Old English}{West Saxon normalization unmarked}!daeg@\iventry{dæg}{}}
`day' by analogy with dat.sg. niht \textless{} \Recon{nahti} `night'
(\citeproc{ref-RingeTaylor2014}{Ringe and Taylor 2014, 380}).

\subsection{Old English evidence}\label{old-english-evidence-139}

Clark Hall lemmatizes
{\emph{niht}}\index[iv]{02oe@\ivlangheader{Old English}{West Saxon normalization unmarked}!niht@\iventry{niht}{}}
`night' and cross-references forms such as
{\emph{neaht}}\index[iv]{02oe@\ivlangheader{Old English}{West Saxon normalization unmarked}!neaht@\iventry{neaht}{}}
`night',
{\emph{neht}}\index[iv]{02oe@\ivlangheader{Old English}{West Saxon normalization unmarked}!neht@\iventry{neht}{}}
`night', and
{\emph{nieht}}\index[iv]{02oe@\ivlangheader{Old English}{West Saxon normalization unmarked}!nieht@\iventry{nieht}{}}
`night' (\citeproc{ref-ClarkHall1960}{Clark Hall 1960, 215}). Campbell
likewise preserves the fluctuation between
{\emph{neaht}}\index[iv]{02oe@\ivlangheader{Old English}{West Saxon normalization unmarked}!neaht@\iventry{neaht}{}}
`night' and
{\emph{niht}}\index[iv]{02oe@\ivlangheader{Old English}{West Saxon normalization unmarked}!niht@\iventry{niht}{}}
`night', giving genitive \emph{nihte, nihtes}, dative \emph{niht,
nihte}, nominative plural \emph{niht} `night', and the contrasting
plural-side forms represented by
{\emph{neahtas}}\index[iv]{02oe@\ivlangheader{Old English}{West Saxon normalization unmarked}!neahtas@\iventry{neahtas}{}}
`night' (\citeproc{ref-Campbell1959}{Campbell 1959, § 628.3}).

The comparison form used here is therefore an attested Old English
\emph{niht} `night', not a reconstructed substitute. The broader lexical
record still preserves the non-umlauted side of the paradigm in
{\emph{neaht}}\index[iv]{02oe@\ivlangheader{Old English}{West Saxon normalization unmarked}!neaht@\iventry{neaht}{}}
`night'-type forms.

\subsection{Development to Old
English}\label{development-to-old-english-137}

Ringe and Taylor derive West Saxon \emph{niht} `night' from
\Recon{nahti} `night' via \Recon{nehti} `night' and \Recon{neahti}
`night' (\citeproc{ref-RingeTaylor2014}{Ringe and Taylor 2014, 240}).
Campbell and Brunner preserve the contrasting non-umlauted
{\emph{neaht}}\index[iv]{02oe@\ivlangheader{Old English}{West Saxon normalization unmarked}!neaht@\iventry{neaht}{}}
`night'-type forms elsewhere in the paradigm
(\citeproc{ref-Campbell1959}{Campbell 1959, § 628.3};
\citeproc{ref-SieversBrunner1965}{Brunner 1965, § 284}).

The modeled path is therefore \Recon{náxti} `night' \textgreater{}
\Recon{næxti} `night' \textgreater{} \Recon{neaxti} `night'
\textgreater{} \Recon{niexti} `night' \textgreater{} \Recon{nixti}
`night' \textgreater{} niht.

\subsection{Paradigm comparison}\label{paradigm-comparison-17}

The comparison below sets the relevant forms side by side. It identifies
the inherited cell that matches the attested Old English form.

{\def\LTcaptype{none} % do not increment counter
\begin{longtable}[]{@{}
  >{\raggedright\arraybackslash}p{(\linewidth - 8\tabcolsep) * \real{0.2000}}
  >{\raggedright\arraybackslash}p{(\linewidth - 8\tabcolsep) * \real{0.2000}}
  >{\raggedright\arraybackslash}p{(\linewidth - 8\tabcolsep) * \real{0.2000}}
  >{\raggedright\arraybackslash}p{(\linewidth - 8\tabcolsep) * \real{0.2000}}
  >{\raggedright\arraybackslash}p{(\linewidth - 8\tabcolsep) * \real{0.2000}}@{}}
\toprule\noalign{}
\begin{minipage}[b]{\linewidth}\raggedright
PGmc cell / interpretation
\end{minipage} & \begin{minipage}[b]{\linewidth}\raggedright
Candidate input
\end{minipage} & \begin{minipage}[b]{\linewidth}\raggedright
OE output or comparison
\end{minipage} & \begin{minipage}[b]{\linewidth}\raggedright
OE comparison form
\end{minipage} & \begin{minipage}[b]{\linewidth}\raggedright
Result
\end{minipage} \\
\midrule\noalign{}
\endhead
\bottomrule\noalign{}
\endlastfoot
citation nominative singular &
{\emph{*náxtz}}\index[iv]{03pgmc@\ivlangheader{Proto-Germanic}{}!naxtz@\iventry{*náxtz}{}}
& expected non-umlauted outcome {\emph{neaht}} &
{\emph{neaht}}\index[iv]{02oe@\ivlangheader{Old English}{West Saxon normalization unmarked}!neaht@\iventry{neaht}{}}
& useful background, but not the comparison used for \emph{niht} \\
selected dative singular &
{\emph{*náxti}}\index[iv]{03pgmc@\ivlangheader{Proto-Germanic}{}!naxti@\iventry{*náxti}{}}
& regular output: {\emph{niht}} `night' &
{\emph{niht}}\index[iv]{02oe@\ivlangheader{Old English}{West Saxon normalization unmarked}!niht@\iventry{niht}{}}
`night' & exact match between input, output, and paradigm cell \\
\end{longtable}
}

\section{\texorpdfstring{rest --- OE
\emph{ræste}}{rest --- OE ræste}}\label{rest-oe-ruxe6ste}

\index[iv]{02oe@\ivlangheader{Old English}{West Saxon normalization unmarked}!raeste@\iventry{ræste}{}}
\index[iv]{02oe@\ivlangheader{Old English}{West Saxon normalization unmarked}!rast@\iventry{rast}{}}
\index[iv]{03pgmc@\ivlangheader{Proto-Germanic}{}!rasto@\iventry{*rastō}{}}
\index[iv]{03pgmc@\ivlangheader{Proto-Germanic}{}!rastoz@\iventry{*rástōz}{}}

Derivation: citation reconstruction \emph{*rastō}; form followed here
\emph{*rástōz} \textgreater{} \emph{ræste} (late analogy).

\subsection{Derivation trace}\label{derivation-trace-159}

Proto input: \emph{*rástōz}

\begingroup
\setlength{\fboxsep}{6pt}

\noindent\fbox{%
\begin{minipage}{0.97\linewidth}
\small
\begin{tabular}{@{}>{\raggedright\arraybackslash}p{0.359\linewidth}@{\hspace{0.030\linewidth}}>{\raggedright\arraybackslash}p{0.611\linewidth}@{}}
\begin{minipage}[t]{\linewidth}
\centering\textbf{Earlier Germanic changes}\par
\vspace{0.35em}
\raggedright
\centering\textbf{Northwest and West Germanic}\par
\raggedright
\vspace{0.2em}
\begin{tabularx}{\linewidth}{@{}>{\raggedright\arraybackslash}X@{\hspace{0.55em}}>{\raggedright\arraybackslash}p{0.14\linewidth}@{\hspace{0.25em}}}
\mbox{EAF Final Z Deletion} & \emph{*rástō} \\
\end{tabularx}
\vspace{0.6em}
\centering\textbf{Early Anglo-Frisian}\par
\raggedright
\vspace{0.2em}
\raggedright [no change]\par
\end{minipage}
&
\begin{minipage}[t]{\linewidth}
\centering\textbf{Old English changes}\par
\vspace{0.35em}
\raggedright
\begin{tabularx}{\linewidth}{@{}>{\raggedright\arraybackslash}X@{\hspace{0.55em}}>{\raggedright\arraybackslash}p{0.14\linewidth}@{\hspace{0.25em}}}
PWGmc Surviving Bimoric O Unrounding & \emph{*rástā} \\
\mbox{EAF Brightening} & \emph{*ræstā} \\
\mbox{EAF Brightening Long Final} & \emph{*ræstǣ} \\
OE Unstressed Long Vowel Shortening & \emph{*ræstæ} \\
\mbox{OE Unstressed AE Merger} & \emph{*ræste} \\
\end{tabularx}
\end{minipage}
\\
\end{tabular}
\end{minipage}%
} \endgroup

Old English form: \emph{ræste}

\subsection{Reconstruction and comparative
evidence}\label{reconstruction-and-comparative-evidence-139}

Kroonen treats the noun as a feminine ō-stem
{\emph{*rastō-}}\index[iv]{03pgmc@\ivlangheader{Proto-Germanic}{}!rasto@\iventry{*rastō-}{}},
continued by Old English
{\emph{ræst}}\index[iv]{02oe@\ivlangheader{Old English}{West Saxon normalization unmarked}!raest@\iventry{ræst}{}}
`rest' (\citeproc{ref-Kroonen2013}{Kroonen 2013, 445}). The derivational
input
\Recon{rástōz}\index[iv]{03pgmc@\ivlangheader{Proto-Germanic}{}!rastoz@\iventry{*rástōz}{}}
`rest' therefore does not replace the lexeme-level headword. It
identifies one oblique singular cell on the side of the paradigm that
yields \emph{ræste} `rest'.

The source tradition used here labels that cell specifically as genitive
singular, but the broader local synthesis of the ō-stem paradigm shows
that the oblique singulars converge on the same front-vocalic
\emph{ræste} `rest' side, in contrast to a nominative singular that
would remain \emph{rast} `rest'.

\subsection{Old English evidence}\label{old-english-evidence-140}

The ordinary Old English citation form is
{\emph{ræst}}\index[iv]{02oe@\ivlangheader{Old English}{West Saxon normalization unmarked}!raest@\iventry{ræst}{}}
`rest' (\citeproc{ref-Kroonen2013}{Kroonen 2013, 445}). Clark Hall
likewise gives
{\emph{ræst}}\index[iv]{02oe@\ivlangheader{Old English}{West Saxon normalization unmarked}!raest@\iventry{ræst}{}}
`rest' (\citeproc{ref-ClarkHall1960}{Clark Hall 1960, 239}).
Bosworth-Toller also preserves oblique uses of
{\emph{ræste}}\index[iv]{02oe@\ivlangheader{Old English}{West Saxon normalization unmarked}!raeste@\iventry{ræste}{}}
`rest', including prepositional examples such as \emph{on ræste} `at
rest' and \emph{tó ræste} `to rest'
(\citeproc{ref-BosworthToller1898}{Bosworth and Toller 1898, 121}).

The comparison form used here is therefore an attested oblique
{\emph{ræste}}\index[iv]{02oe@\ivlangheader{Old English}{West Saxon normalization unmarked}!raeste@\iventry{ræste}{}}
`rest', not a reconstructed surrogate. The dictionary headword
{\emph{ræst}}\index[iv]{02oe@\ivlangheader{Old English}{West Saxon normalization unmarked}!raest@\iventry{ræst}{}}
`rest' remains an equally real part of the Old English record.

\subsection{Development to Old
English}\label{development-to-old-english-138}

Once final \emph{*z} is lost, the derivational input moves through the
front-vocalic oblique side of the paradigm rather than the back-vocalic
nominative side. In the modeled derivation, the surviving final long
vowel is first exposed, unrounded, fronted, shortened, and then reduced
to the final \emph{-e} of \emph{ræste} `rest'.

The key point is the paradigm split. Nominative \Recon{rastō} `rest'
yields a regular \emph{rast} `rest', whereas the selected oblique input
yields \emph{ræste} `rest'. The later citation form \emph{ræst} `rest'
is best explained as leveling from that oblique \emph{ræst-} stem.

\subsection{Paradigm comparison}\label{paradigm-comparison-18}

The comparison below sets the relevant forms side by side. It
distinguishes the nominative citation form from the oblique singular
chosen here.

{\def\LTcaptype{none} % do not increment counter
\begin{longtable}[]{@{}
  >{\raggedright\arraybackslash}p{(\linewidth - 8\tabcolsep) * \real{0.2000}}
  >{\raggedright\arraybackslash}p{(\linewidth - 8\tabcolsep) * \real{0.2000}}
  >{\raggedright\arraybackslash}p{(\linewidth - 8\tabcolsep) * \real{0.2000}}
  >{\raggedright\arraybackslash}p{(\linewidth - 8\tabcolsep) * \real{0.2000}}
  >{\raggedright\arraybackslash}p{(\linewidth - 8\tabcolsep) * \real{0.2000}}@{}}
\toprule\noalign{}
\begin{minipage}[b]{\linewidth}\raggedright
PGmc cell / interpretation
\end{minipage} & \begin{minipage}[b]{\linewidth}\raggedright
Candidate input
\end{minipage} & \begin{minipage}[b]{\linewidth}\raggedright
OE output or comparison
\end{minipage} & \begin{minipage}[b]{\linewidth}\raggedright
OE comparison form
\end{minipage} & \begin{minipage}[b]{\linewidth}\raggedright
Result
\end{minipage} \\
\midrule\noalign{}
\endhead
\bottomrule\noalign{}
\endlastfoot
citation nominative singular &
{\emph{*rastō}}\index[iv]{03pgmc@\ivlangheader{Proto-Germanic}{}!rasto@\iventry{*rastō}{}}
& expected regular outcome \emph{rast} &
{\emph{ræst}}\index[iv]{02oe@\ivlangheader{Old English}{West Saxon normalization unmarked}!raest@\iventry{ræst}{}}
`rest' & useful background, but not the cell that matches attested
oblique \emph{ræste} \\
selected oblique singular &
{\emph{*rástōz}}\index[iv]{03pgmc@\ivlangheader{Proto-Germanic}{}!rastoz@\iventry{*rástōz}{}}
& regular output: {\emph{ræste}} `rest' &
{\emph{ræste}}\index[iv]{02oe@\ivlangheader{Old English}{West Saxon normalization unmarked}!raeste@\iventry{ræste}{}}
`rest' & exact match between derivational input and attested OE oblique
form \\
\end{longtable}
}

\section{\texorpdfstring{shoulder --- OE
\emph{sċuldrum}}{shoulder --- OE sċuldrum}}\label{shoulder-oe-sux10buldrum}

\index[iv]{02oe@\ivlangheader{Old English}{West Saxon normalization unmarked}!sculdrum@\iventry{sċuldrum}{}}
\index[iv]{03pgmc@\ivlangheader{Proto-Germanic}{}!skuldramiz@\iventry{*skúldramiz}{}}
\index[iv]{03pgmc@\ivlangheader{Proto-Germanic}{}!skuldro@\iventry{*skuldrō}{}}

Derivation: citation reconstruction \emph{*skuldrō}; form followed here
\emph{*skúldramiz} \textgreater{} \emph{sċuldrum} (late analogy).

\subsection{Derivation trace}\label{derivation-trace-160}

Proto input: \emph{*skúldramiz}

\begingroup
\setlength{\fboxsep}{6pt}

\noindent\fbox{%
\begin{minipage}{0.97\linewidth}
\small
\begin{tabular}{@{}>{\raggedright\arraybackslash}p{0.630\linewidth}@{\hspace{0.030\linewidth}}>{\raggedright\arraybackslash}p{0.339\linewidth}@{}}
\begin{minipage}[t]{\linewidth}
\centering\textbf{Earlier Germanic changes}\par
\vspace{0.35em}
\raggedright
\centering\textbf{Northwest and West Germanic}\par
\raggedright
\vspace{0.2em}
\begin{tabularx}{\linewidth}{@{}>{\raggedright\arraybackslash}X@{\hspace{0.55em}}>{\raggedright\arraybackslash}p{0.24\linewidth}@{\hspace{0.25em}}}
\mbox{PNWGmc A To U Before M} & \emph{*skúldrumiz} \\
\mbox{PWGmc Early I Apocope} & \emph{*skúldrumz} \\
\mbox{EAF Final Z Deletion} & \emph{*skúldrum} \\
\end{tabularx}
\vspace{0.6em}
\centering\textbf{Early Anglo-Frisian}\par
\raggedright
\vspace{0.2em}
\raggedright [no change]\par
\end{minipage}
&
\begin{minipage}[t]{\linewidth}
\centering\textbf{Old English changes}\par
\vspace{0.35em}
\raggedright
\begin{tabularx}{\linewidth}{@{}>{\raggedright\arraybackslash}X@{\hspace{0.55em}}>{\raggedright\arraybackslash}p{0.24\linewidth}@{\hspace{0.25em}}}
\mbox{OE Sk Palatalization} & \emph{*ʃúldrum} \\
\end{tabularx}
\end{minipage}
\\
\end{tabular}
\end{minipage}%
} \endgroup

Old English form: \emph{sċuldrum}

\subsection{Reconstruction and comparative
evidence}\label{reconstruction-and-comparative-evidence-140}

The handbooks do not agree on the reconstruction of the Germanic word.
Orel gives \emph{*skuldr(j)ō} as a feminine ō-/jō-stem and explicitly
notes that Old English
{\emph{sculdor}}\index[iv]{02oe@\ivlangheader{Old English}{West Saxon normalization unmarked}!sculdor@\iventry{sculdor}{}}
`shoulder' is masculine beside OFrisian
{\emph{skulder}}\index[iv]{11ofris@\ivlangheader{Old Frisian}{}!skulder@\iventry{skulder}{}}
`shoulder', Middle Low German
{schulder}\index[iv]{24mlg@\ivlangheader{Middle Low German}{}!schulder@\iventry{schulder}{}}
`shoulder', and Old High German
{\emph{scultra}}\index[iv]{10ohg@\ivlangheader{Old High German}{}!scultra@\iventry{scultra}{}}
`shoulder',
{\emph{scultirra}}\index[iv]{10ohg@\ivlangheader{Old High German}{}!scultirra@\iventry{scultirra}{}}
`shoulder' (\citeproc{ref-Orel2003}{Orel 2003, 345}). Kroonen
reconstructs
{\emph{*skuldra-}}\index[iv]{03pgmc@\ivlangheader{Proto-Germanic}{}!skuldra@\iventry{*skuldra-}{}},
a masculine a-stem, and derives the Old High German feminine forms from
{\emph{*skuldrjōn-}}\index[iv]{03pgmc@\ivlangheader{Proto-Germanic}{}!skuldrjon@\iventry{*skuldrjōn-}{}}
(\citeproc{ref-Kroonen2013}{Kroonen 2013, 478}). Ringe and Taylor cite
PWGmc
\Recon{skuldru}\index[iv]{05pwgmc@\ivlangheader{Proto-West Germanic}{}!skuldru@\iventry{*skuldru}{}}
`shoulder' for the Old English branch
(\citeproc{ref-RingeTaylor2014}{Ringe and Taylor 2014, 142}).

These forms imply different stem classes and different expectations for
the Old English inflection. The question is which inflectional cell best
aligns with the Old English evidence.

A dative/instrumental plural form
\Recon{skúldramiz}\index[iv]{03pgmc@\ivlangheader{Proto-Germanic}{}!skuldramiz@\iventry{*skúldramiz}{}}
`shoulder' aligns with the inherited plural ending that later yields Old
English \emph{-um}, and it corresponds directly to the attested dative
plural discussed below.

\subsection{Old English evidence}\label{old-english-evidence-141}

The ordinary Old English headword is
{\emph{sculdor}}\index[iv]{02oe@\ivlangheader{Old English}{West Saxon normalization unmarked}!sculdor@\iventry{sculdor}{}}
`shoulder'. Clark Hall lemmatizes
{\emph{sculdor}}\index[iv]{02oe@\ivlangheader{Old English}{West Saxon normalization unmarked}!sculdor@\iventry{sculdor}{}}
`shoulder' as the normal dictionary form
(\citeproc{ref-ClarkHall1960}{Clark Hall 1960, 257}). Bosworth-Toller
also preserves the dative plural
{\emph{sculdrum}}\index[iv]{02oe@\ivlangheader{Old English}{West Saxon normalization unmarked}!sculdrum@\iventry{sculdrum}{}}
`shoulder' (\citeproc{ref-BosworthToller1898}{Bosworth and Toller 1898,
85}).

Bosworth-Toller's Supplement records a weak-feminine
{\emph{sculdra}}\index[iv]{02oe@\ivlangheader{Old English}{West Saxon normalization unmarked}!sculdra@\iventry{sculdra}{}}
`shoulder' (\citeproc{ref-BosworthToller1898}{Bosworth and Toller 1898,
699}), so
{\emph{sculdra}}\index[iv]{02oe@\ivlangheader{Old English}{West Saxon normalization unmarked}!sculdra@\iventry{sculdra}{}}
`shoulder' belongs to the Old English record beside the stronger
masculine paradigm headed by
{\emph{sculdor}}\index[iv]{02oe@\ivlangheader{Old English}{West Saxon normalization unmarked}!sculdor@\iventry{sculdor}{}}
`shoulder'. Brunner and Luick also record later spellings such as
{\emph{sceoldor}}\index[iv]{02oe@\ivlangheader{Old English}{West Saxon normalization unmarked}!sceoldor@\iventry{sceoldor}{}}
`shoulder' and the i-mutated dative plural
{\emph{scyldrum}}\index[iv]{02oe@\ivlangheader{Old English}{West Saxon normalization unmarked}!scyldrum@\iventry{scyldrum}{}}
`shoulder', which reflect secondary phonological and analogical
reshaping within Old English (\citeproc{ref-SieversBrunner1965}{Brunner
1965, § 92.2.a}; \citeproc{ref-Luick1914}{Luick 1914-\/-1940, 230}).

The singular and plural evidence point to different parts of the
paradigm. The relevant comparison form here is the attested dative
plural
{\emph{sċuldrum}}\index[iv]{02oe@\ivlangheader{Old English}{West Saxon normalization unmarked}!sculdrum@\iventry{sċuldrum}{}}
`shoulder (dat.pl.)'. The spelling with \emph{sċ-} is a normalized
representation of the same Old English initial cluster.

\subsection{Development to Old
English}\label{development-to-old-english-139}

Proto-Germanic \Recon{skúldramiz} `shoulder' can be interpreted as a
dative/instrumental plural form. In this environment the post-tonic
\emph{a} before \emph{m} is raised to \emph{u}, giving a form of the
\Recon{skúldrumiz} `shoulder' type. Unstressed \emph{u} is regularly
preserved before \emph{m}, especially in the dative plural ending
\emph{-um}: Campbell states this explicitly, and Hogg formulates the
same condition for the dative plural inflexion
(\citeproc{ref-Campbell1959}{Campbell 1959, § 373};
\citeproc{ref-HoggPhonology1992}{Hogg 1992b, § 3.3.1.3}). Brunner
points in the same direction by excluding \emph{m} from the environments
in which medial \emph{o} became general in West Saxon
(\citeproc{ref-SieversBrunner1965}{Brunner 1965, § 44} Anm. 7).

Subsequent reduction of the ending removes the final \emph{*i} and
\emph{*z}, so that the inflectional ending appears in Old English as
\emph{-um}. The initial cluster is written here as \emph{sċ-}, and the
development is \Recon{skúldramiz} `shoulder' \textgreater{}
\Recon{skúldrumiz} `shoulder' \textgreater{} \Recon{skúldrum} `shoulder'
\textgreater{}
{\emph{sċuldrum}}\index[iv]{02oe@\ivlangheader{Old English}{West Saxon normalization unmarked}!sculdrum@\iventry{sċuldrum}{}}
`shoulder'.

\subsection{Paradigm comparison}\label{paradigm-comparison-19}

A paradigm comparison identifies the Proto-Germanic inflectional cell
that corresponds to an established Old English paradigm form. The
comparison below sets the relevant forms side by side.

{\def\LTcaptype{none} % do not increment counter
\begin{longtable}[]{@{}
  >{\raggedright\arraybackslash}p{(\linewidth - 8\tabcolsep) * \real{0.2000}}
  >{\raggedright\arraybackslash}p{(\linewidth - 8\tabcolsep) * \real{0.2000}}
  >{\raggedright\arraybackslash}p{(\linewidth - 8\tabcolsep) * \real{0.2000}}
  >{\raggedright\arraybackslash}p{(\linewidth - 8\tabcolsep) * \real{0.2000}}
  >{\raggedright\arraybackslash}p{(\linewidth - 8\tabcolsep) * \real{0.2000}}@{}}
\toprule\noalign{}
\begin{minipage}[b]{\linewidth}\raggedright
PGmc cell / interpretation
\end{minipage} & \begin{minipage}[b]{\linewidth}\raggedright
Candidate input
\end{minipage} & \begin{minipage}[b]{\linewidth}\raggedright
OE output or comparison
\end{minipage} & \begin{minipage}[b]{\linewidth}\raggedright
OE comparison form
\end{minipage} & \begin{minipage}[b]{\linewidth}\raggedright
Result
\end{minipage} \\
\midrule\noalign{}
\endhead
\bottomrule\noalign{}
\endlastfoot
singular-oriented citation input &
{\emph{*skúldrō}}\index[iv]{03pgmc@\ivlangheader{Proto-Germanic}{}!skuldro@\iventry{*skúldrō}{}}
& probe output: {\emph{sċoldor}} &
{\emph{sculdor}}\index[iv]{02oe@\ivlangheader{Old English}{West Saxon normalization unmarked}!sculdor@\iventry{sculdor}{}}
`shoulder' & fails: the singular output has root \emph{o}, not the
attested \emph{u} \\
serious plural-based singular alternative &
{\emph{*skúldru}}\index[iv]{05pwgmc@\ivlangheader{Proto-West Germanic}{}!skuldru@\iventry{*skúldru}{}}
& probe output: {\emph{sċuldor}} &
{\emph{sculdor}}\index[iv]{02oe@\ivlangheader{Old English}{West Saxon normalization unmarked}!sculdor@\iventry{sculdor}{}}
`shoulder' & close formally, but it compares a plural-stage input with a
singular form \\
dat./inst.pl. input &
\Recon{skúldramiz}\index[iv]{03pgmc@\ivlangheader{Proto-Germanic}{}!skuldramiz@\iventry{*skúldramiz}{}}
`shoulder' & regular output:
{\emph{sċuldrum}}\index[iv]{02oe@\ivlangheader{Old English}{West Saxon normalization unmarked}!sculdrum@\iventry{sċuldrum}{}}
`shoulder (dat.pl.)' &
{\emph{sculdrum}}\index[iv]{02oe@\ivlangheader{Old English}{West Saxon normalization unmarked}!sculdrum@\iventry{sculdrum}{}}
`shoulder' & matches both the output and the dative plural comparison
form \\
later weak-feminine singular & --- & OE
{\emph{sculdra}}\index[iv]{02oe@\ivlangheader{Old English}{West Saxon normalization unmarked}!sculdra@\iventry{sculdra}{}}
`shoulder' &
{\emph{sculdra}}\index[iv]{02oe@\ivlangheader{Old English}{West Saxon normalization unmarked}!sculdra@\iventry{sculdra}{}}
`shoulder' & secondary doublet, useful as a control rather than the
inherited target \\
\end{longtable}
}

The dative plural line is decisive because it matches both the output
and the paradigm cell of Old English
{\emph{sculdrum}}\index[iv]{02oe@\ivlangheader{Old English}{West Saxon normalization unmarked}!sculdrum@\iventry{sculdrum}{}}
`shoulder'. Singular-oriented candidates either lower the root vowel or
compare unlike cells.

\section{\texorpdfstring{shove --- OE
\emph{sċēaf}}{shove --- OE sċēaf}}\label{shove-oe-sux10bux113af}

\index[iv]{02oe@\ivlangheader{Old English}{West Saxon normalization unmarked}!sceaf@\iventry{sċēaf}{}}
\index[iv]{03pgmc@\ivlangheader{Proto-Germanic}{}!skaub@\iventry{*skáub}{}}
\index[iv]{03pgmc@\ivlangheader{Proto-Germanic}{}!skeubana@\iventry{*skéubaną}{}}

Derivation: citation reconstruction \emph{*skéubaną}; form followed here
\emph{*skáub} \textgreater{} \emph{sċēaf} (late analogy).

\subsection{Derivation trace}\label{derivation-trace-161}

Proto input: \emph{*skáub}

\begingroup
\setlength{\fboxsep}{6pt}

\noindent\fbox{%
\begin{minipage}{0.97\linewidth}
\small
\begin{tabular}{@{}>{\raggedright\arraybackslash}p{0.388\linewidth}@{\hspace{0.030\linewidth}}>{\raggedright\arraybackslash}p{0.582\linewidth}@{}}
\begin{minipage}[t]{\linewidth}
\centering\textbf{Earlier Germanic changes}\par
\vspace{0.35em}
\raggedright
\centering\textbf{Northwest and West Germanic}\par
\raggedright
\vspace{0.2em}
\raggedright [no change]\par
\vspace{0.6em}
\centering\textbf{Early Anglo-Frisian}\par
\raggedright
\vspace{0.2em}
\raggedright [no change]\par
\end{minipage}
&
\begin{minipage}[t]{\linewidth}
\centering\textbf{Old English changes}\par
\vspace{0.35em}
\raggedright
\begin{tabularx}{\linewidth}{@{}>{\raggedright\arraybackslash}X@{\hspace{0.55em}}>{\raggedright\arraybackslash}p{0.14\linewidth}@{\hspace{0.25em}}}
\mbox{OE Au Brightening} & \emph{*skáeub} \\
\mbox{OE Diphthong Leveling} & \emph{*skēab} \\
\mbox{OE Sk Palatalization} & \emph{*ʃēab} \\
\mbox{PGmc B Allophony} & \emph{*ʃēaβ} \\
\end{tabularx}
\end{minipage}
\\
\end{tabular}
\end{minipage}%
} \endgroup

Old English form: \emph{sċēaf}

\subsection{Reconstruction and comparative
evidence}\label{reconstruction-and-comparative-evidence-141}

Kroonen reconstructs the strong verb as \emph{*skeuban-}
\textasciitilde{} \emph{*skūban-} and cites Old English present forms
{\emph{scēofan}}\index[iv]{02oe@\ivlangheader{Old English}{West Saxon normalization unmarked}!sceofan@\iventry{scēofan}{}}
`shove',
{\emph{scūfan}}\index[iv]{02oe@\ivlangheader{Old English}{West Saxon normalization unmarked}!scufan@\iventry{scūfan}{}}
`shove' (\citeproc{ref-Kroonen2013}{Kroonen 2013, 444}). Those
present-system forms belong to the same verb family, but the comparison
here uses the singular preterite
\Recon{skáub}\index[iv]{03pgmc@\ivlangheader{Proto-Germanic}{}!skaub@\iventry{*skáub}{}}
`shove', not the infinitive.

\subsection{Old English evidence}\label{old-english-evidence-142}

The ordinary dictionary verb is
{\emph{scūfan}}\index[iv]{02oe@\ivlangheader{Old English}{West Saxon normalization unmarked}!scufan@\iventry{scūfan}{}}
`shove'/{\emph{scēofan}}\index[iv]{02oe@\ivlangheader{Old English}{West Saxon normalization unmarked}!sceofan@\iventry{scēofan}{}}
`shove', but the preterite itself is well attested. Bright gives the
principal parts
{\emph{scufan}}\index[iv]{02oe@\ivlangheader{Old English}{West Saxon normalization unmarked}!scufan@\iventry{scufan}{}}
`shove',
{\emph{sceaf}}\index[iv]{02oe@\ivlangheader{Old English}{West Saxon normalization unmarked}!sceaf@\iventry{sceaf}{}}
`shove',
{\emph{scufon}}\index[iv]{02oe@\ivlangheader{Old English}{West Saxon normalization unmarked}!scufon@\iventry{scufon}{}}
`shove',
{\emph{scofen}}\index[iv]{02oe@\ivlangheader{Old English}{West Saxon normalization unmarked}!scofen@\iventry{scofen}{}}
`shove' (\citeproc{ref-BrightCassidyRingler1971}{Bright 1971, 347}).
Sweet gives the same paradigm (\citeproc{ref-Sweet1953}{Sweet 1953,
29}). The normalized form here is
{\emph{sċēaf}}\index[iv]{02oe@\ivlangheader{Old English}{West Saxon normalization unmarked}!sceaf@\iventry{sċēaf}{}}
`sheaf', regularizing the attested spellings
{\emph{sceaf}}\index[iv]{02oe@\ivlangheader{Old English}{West Saxon normalization unmarked}!sceaf@\iventry{sceaf}{}}
`shove' and prefixed
{\emph{āsceaf}}\index[iv]{02oe@\ivlangheader{Old English}{West Saxon normalization unmarked}!asceaf@\iventry{āsceaf}{}}
`shove'.

\subsection{Development to Old
English}\label{development-to-old-english-140}

From \Recon{skáub} `shove', the development is straightforward.
\emph{*au} fronts and levels to \emph{ēa}, final \emph{*b} becomes a
fricative and is written \emph{f}, and initial \emph{*sk-} undergoes the
usual Old English palatalized spelling in this environment. The
derivation therefore gives \Recon{skáub} `shove' \textgreater{}
\Recon{skáeub} `shove' \textgreater{} \Recon{skēab} `shove'
\textgreater{} \Recon{skēaβ} `shove' \textgreater{}
{\emph{sċēaf}}\index[iv]{02oe@\ivlangheader{Old English}{West Saxon normalization unmarked}!sceaf@\iventry{sċēaf}{}}
`sheaf'.

\subsection{Paradigm comparison}\label{paradigm-comparison-20}

A paradigm comparison is needed here because the ordinary citation verb
and
{\emph{sċēaf}}\index[iv]{02oe@\ivlangheader{Old English}{West Saxon normalization unmarked}!sceaf@\iventry{sċēaf}{}}
`sheaf' belong to different cells of the same strong paradigm. The
comparison below is manual.

{\def\LTcaptype{none} % do not increment counter
\begin{longtable}[]{@{}
  >{\raggedright\arraybackslash}p{(\linewidth - 8\tabcolsep) * \real{0.2000}}
  >{\raggedright\arraybackslash}p{(\linewidth - 8\tabcolsep) * \real{0.2000}}
  >{\raggedright\arraybackslash}p{(\linewidth - 8\tabcolsep) * \real{0.2000}}
  >{\raggedright\arraybackslash}p{(\linewidth - 8\tabcolsep) * \real{0.2000}}
  >{\raggedright\arraybackslash}p{(\linewidth - 8\tabcolsep) * \real{0.2000}}@{}}
\toprule\noalign{}
\begin{minipage}[b]{\linewidth}\raggedright
PGmc cell / interpretation
\end{minipage} & \begin{minipage}[b]{\linewidth}\raggedright
Candidate input
\end{minipage} & \begin{minipage}[b]{\linewidth}\raggedright
OE output or comparison
\end{minipage} & \begin{minipage}[b]{\linewidth}\raggedright
OE comparison form
\end{minipage} & \begin{minipage}[b]{\linewidth}\raggedright
Result
\end{minipage} \\
\midrule\noalign{}
\endhead
\bottomrule\noalign{}
\endlastfoot
citation infinitive &
{\emph{*skéubaną}}\index[iv]{03pgmc@\ivlangheader{Proto-Germanic}{}!skeubana@\iventry{*skéubaną}{}}
& inherited infinitive line
{\emph{sċēofan}}\index[iv]{02oe@\ivlangheader{Old English}{West Saxon normalization unmarked}!sceofan@\iventry{sċēofan}{}};
present system also leveled
{\emph{scūfan}}\index[iv]{02oe@\ivlangheader{Old English}{West Saxon normalization unmarked}!scufan@\iventry{scūfan}{}}
`shove' &
{\emph{scēofan}}\index[iv]{02oe@\ivlangheader{Old English}{West Saxon normalization unmarked}!sceofan@\iventry{scēofan}{}}
`shove' /
{\emph{scūfan}}\index[iv]{02oe@\ivlangheader{Old English}{West Saxon normalization unmarked}!scufan@\iventry{scūfan}{}}
`shove' & necessary background, but not the comparison used for
{\emph{sċēaf}}\index[iv]{02oe@\ivlangheader{Old English}{West Saxon normalization unmarked}!sceaf@\iventry{sċēaf}{}}
`sheaf' \\
1/3 sg. preterite &
{\emph{*skáub}}\index[iv]{03pgmc@\ivlangheader{Proto-Germanic}{}!skaub@\iventry{*skáub}{}}
& documented regular output: {\emph{sċēaf}} `sheaf' &
{\emph{sċēaf}}\index[iv]{02oe@\ivlangheader{Old English}{West Saxon normalization unmarked}!sceaf@\iventry{sċēaf}{}}
`sheaf' & direct match for the singular preterite \\
preterite plural &
{\emph{*skúbun}}\index[iv]{03pgmc@\ivlangheader{Proto-Germanic}{}!skubun@\iventry{*skúbun}{}}
& later leveled plural
{\emph{scufon}}\index[iv]{02oe@\ivlangheader{Old English}{West Saxon normalization unmarked}!scufon@\iventry{scufon}{}}
`shove' beside expected
{\emph{sċufun}}\index[iv]{02oe@\ivlangheader{Old English}{West Saxon normalization unmarked}!scufun@\iventry{sċufun}{}}
under the corrected cascade &
{\emph{scufon}}\index[iv]{02oe@\ivlangheader{Old English}{West Saxon normalization unmarked}!scufon@\iventry{scufon}{}}
`shove' & poorer comparison for the singular-preterite target \\
past participle &
{\emph{*skúbanaz}}\index[iv]{03pgmc@\ivlangheader{Proto-Germanic}{}!skubanaz@\iventry{*skúbanaz}{}}
& attested participial line
{\emph{scofen}}\index[iv]{02oe@\ivlangheader{Old English}{West Saxon normalization unmarked}!scofen@\iventry{scofen}{}}
`shove' &
{\emph{scofen}}\index[iv]{02oe@\ivlangheader{Old English}{West Saxon normalization unmarked}!scofen@\iventry{scofen}{}}
`shove' & valid alternative cell, but not the form compared here \\
\end{longtable}
}

\section{\texorpdfstring{show (3sg) --- OE
\emph{sċēawaþ}}{show (3sg) --- OE sċēawaþ}}\label{show-3sg-oe-sux10bux113awauxfe}

\index[iv]{02oe@\ivlangheader{Old English}{West Saxon normalization unmarked}!sceawad@\iventry{sceawað}{}}
\index[iv]{02oe@\ivlangheader{Old English}{West Saxon normalization unmarked}!sceawath@\iventry{sċēawaþ}{}}
\index[iv]{03pgmc@\ivlangheader{Proto-Germanic}{}!skawona@\iventry{*skawōną}{}}
\index[iv]{03pgmc@\ivlangheader{Proto-Germanic}{}!skawothi@\iventry{*skáwōθi}{}}

Derivation: citation reconstruction \emph{*skawōną}; form followed here
\emph{*skáwōθi} \textgreater{} \emph{sċēawaþ} (late analogy).

\subsection{Derivation trace}\label{derivation-trace-162}

Proto input: \emph{*skáwōθi}

\begingroup
\setlength{\fboxsep}{6pt}

\noindent\fbox{%
\begin{minipage}{0.97\linewidth}
\small
\begin{tabular}{@{}>{\raggedright\arraybackslash}p{0.359\linewidth}@{\hspace{0.030\linewidth}}>{\raggedright\arraybackslash}p{0.611\linewidth}@{}}
\begin{minipage}[t]{\linewidth}
\centering\textbf{Earlier Germanic changes}\par
\vspace{0.35em}
\raggedright
\centering\textbf{Northwest and West Germanic}\par
\raggedright
\vspace{0.2em}
\begin{tabularx}{\linewidth}{@{}>{\raggedright\arraybackslash}X@{\hspace{0.55em}}>{\raggedright\arraybackslash}p{0.24\linewidth}@{\hspace{0.25em}}}
\mbox{PWGmc Early I Apocope} & \emph{*skáwōθ} \\
\end{tabularx}
\vspace{0.6em}
\centering\textbf{Early Anglo-Frisian}\par
\raggedright
\vspace{0.2em}
\raggedright [no change]\par
\end{minipage}
&
\begin{minipage}[t]{\linewidth}
\centering\textbf{Old English changes}\par
\vspace{0.35em}
\raggedright
\begin{tabularx}{\linewidth}{@{}>{\raggedright\arraybackslash}X@{\hspace{0.55em}}>{\raggedright\arraybackslash}p{0.18\linewidth}@{\hspace{0.25em}}}
\mbox{OE Aw Long Diphthong} & \emph{*skḗawōθ} \\
\mbox{OE Sk Palatalization} & \emph{*ʃḗawōθ} \\
\mbox{OE Late O Shortening} & \emph{*ʃḗawoθ} \\
OE Final Unstressed O Lowering & \emph{*ʃḗawaθ} \\
\end{tabularx}
\end{minipage}
\\
\end{tabular}
\end{minipage}%
} \endgroup

Old English form: \emph{sċēawaþ}

\subsection{Reconstruction and comparative
evidence}\label{reconstruction-and-comparative-evidence-142}

Orel reconstructs the verb as
\Recon{skawōjanan}\index[iv]{03pgmc@\ivlangheader{Proto-Germanic}{}!skawojanan@\iventry{*skawōjanan}{}}
`show' and cites Old English
{\emph{sceáwian}}\index[iv]{02oe@\ivlangheader{Old English}{West Saxon normalization unmarked}!sceawian@\iventry{sceáwian}{}}
`show' beside Old Frisian
{\emph{skawia}}\index[iv]{11ofris@\ivlangheader{Old Frisian}{}!skawia@\iventry{skawia}{}}
`show', Old Saxon
{\emph{skawōn}}\index[iv]{13os@\ivlangheader{Old Saxon}{}!skawon@\iventry{skawōn}{}}
`show', and Old High German
{\emph{scouwōn}}\index[iv]{10ohg@\ivlangheader{Old High German}{}!scouwon@\iventry{scouwōn}{}}
`show' (\citeproc{ref-Orel2003}{Orel 2003, 337}). This is
infinitive-family evidence. CAPR assembles the present third singular
\Recon{skáwōθi}\index[iv]{03pgmc@\ivlangheader{Proto-Germanic}{}!skawothi@\iventry{*skáwōθi}{}}
`show (3sg)', a finite class-II cell with stable \emph{a} in the 3sg
ending (\citeproc{ref-RingeTaylor2014}{Ringe and Taylor 2014, 80}).

Campbell states the class-II present endings as lufas, \emph{-aþ}
(\textless{} \emph{-ōsi}, \emph{-ōþi)} in his late-shortening discussion
(\citeproc{ref-Campbell1959}{Campbell 1959, 148}). Ringe and Taylor
likewise note that class-II weak present 2sg. \emph{-as(t)} and 3sg.
\emph{-aþ} have stable \emph{a} (\citeproc{ref-RingeTaylor2014}{Ringe
and Taylor 2014, 80}). The relevant comparison is therefore the 3sg cell
itself, not an i-umlauted alternative.

\subsection{Old English evidence}\label{old-english-evidence-143}

The held fourth edition of Bright's reader contains the simplex
infinitive
{\emph{scēawian}}\index[iv]{02oe@\ivlangheader{Old English}{West Saxon normalization unmarked}!sceawian@\iventry{scēawian}{}}
`show' and imperative
{\emph{scēawa}}\index[iv]{02oe@\ivlangheader{Old English}{West Saxon normalization unmarked}!sceawa@\iventry{scēawa}{}}
`show' in the Blickling homily (\citeproc{ref-Bright1917}{Bright 1917,
71}). The Wanderer contains prefixed \emph{geondscēawað} `surveys'
(\citeproc{ref-Bright1917}{Bright 1917, 161}). This supports the finite
ending pattern, not a separately quoted unprefixed third singular or a
passage in the separate Cassidy/Ringler edition.

The form written here as
{\emph{sċēawaþ}}\index[iv]{02oe@\ivlangheader{Old English}{West Saxon normalization unmarked}!sceawath@\iventry{sċēawaþ}{}}
`shows' is the normalized simplex comparison form for that weak class-II
pattern. It is therefore not a dictionary headword but a finite
comparison form aligned with the attested \emph{scēaw-} evidence and the
directly cited prefixed \emph{geondscēawað} `surveys' ending pattern.

\subsection{Development to Old
English}\label{development-to-old-english-141}

The infinitive entry and \hyperref[rule-OEAwLongDiphthong]{SC034
OEAwLongDiphthong} explain the dictionary-singleton versus
handbook-geminate root premise. The same computational proxy remains
here without claiming that output fit resolves it
(\citeproc{ref-Orel2003}{Orel 2003, 337};
\citeproc{ref-Kroonen2013}{Kroonen 2013, 442};
\citeproc{ref-Campbell1959}{Campbell 1959, 45};
\citeproc{ref-SieversBrunner1965}{Brunner 1965, 49}).

The complete trace of
\Recon{skáwōθi}\index[iv]{03pgmc@\ivlangheader{Proto-Germanic}{}!skawothi@\iventry{*skáwōθi}{}}
`show (3sg)' has already lost final \emph{i} before that proxy: the
actual checkpoint is \emph{*s*k*á*w*ō*θ}. Thus absence of \emph{j} alone
is not the explanation of the compiled lack of mutation. Campbell's
late-shortening account and Ringe and Taylor's stable-\emph{a}
discussion independently support \emph{-aþ}, not \emph{-eþ}
(\citeproc{ref-Campbell1959}{Campbell 1959, 148};
\citeproc{ref-RingeTaylor2014}{Ringe and Taylor 2014, 80}).

\subsection{Paradigm comparison}\label{paradigm-comparison-21}

The comparison below sets the relevant forms side by side.

{\def\LTcaptype{none} % do not increment counter
\begin{longtable}[]{@{}
  >{\raggedright\arraybackslash}p{(\linewidth - 8\tabcolsep) * \real{0.2000}}
  >{\raggedright\arraybackslash}p{(\linewidth - 8\tabcolsep) * \real{0.2000}}
  >{\raggedright\arraybackslash}p{(\linewidth - 8\tabcolsep) * \real{0.2000}}
  >{\raggedright\arraybackslash}p{(\linewidth - 8\tabcolsep) * \real{0.2000}}
  >{\raggedright\arraybackslash}p{(\linewidth - 8\tabcolsep) * \real{0.2000}}@{}}
\toprule\noalign{}
\begin{minipage}[b]{\linewidth}\raggedright
PGmc cell / interpretation
\end{minipage} & \begin{minipage}[b]{\linewidth}\raggedright
Candidate input
\end{minipage} & \begin{minipage}[b]{\linewidth}\raggedright
Old English outcome or comparison
\end{minipage} & \begin{minipage}[b]{\linewidth}\raggedright
OE comparison form
\end{minipage} & \begin{minipage}[b]{\linewidth}\raggedright
Result
\end{minipage} \\
\midrule\noalign{}
\endhead
\bottomrule\noalign{}
\endlastfoot
lexeme-level infinitive &
{\emph{*skáwōjaną}}\index[iv]{03pgmc@\ivlangheader{Proto-Germanic}{}!skawojana@\iventry{*skáwōjaną}{}}
& regular output {\emph{sċēawian}} &
{\emph{scēawian}}\index[iv]{02oe@\ivlangheader{Old English}{West Saxon normalization unmarked}!sceawian@\iventry{scēawian}{}}
`show' & ordinary dictionary headword of the verb, but not the finite
form compared here \\
imperative singular companion &
{\emph{*skáwô}}\index[iv]{03pgmc@\ivlangheader{Proto-Germanic}{}!skawo@\iventry{*skáwô}{}}
& regular output {\emph{sċēawa}} `show' &
{\emph{scēawa}}\index[iv]{02oe@\ivlangheader{Old English}{West Saxon normalization unmarked}!sceawa@\iventry{scēawa}{}}
`show' & useful family control, but not the target of this entry \\
present third singular &
{\emph{*skáwōθi}}\index[iv]{03pgmc@\ivlangheader{Proto-Germanic}{}!skawothi@\iventry{*skáwōθi}{}}
& regular output
{\emph{sċēawaþ}}\index[iv]{02oe@\ivlangheader{Old English}{West Saxon normalization unmarked}!sceawath@\iventry{sċēawaþ}{}}
`shows' & normalized
{\emph{sċēawaþ}}\index[iv]{02oe@\ivlangheader{Old English}{West Saxon normalization unmarked}!sceawath@\iventry{sċēawaþ}{}}
`shows'; source-side pattern \emph{sceawað} & exact match for the finite
form compared here \\
\end{longtable}
}

\section{\texorpdfstring{show (iptv.2sg) --- OE
\emph{sċēawa}}{show (iptv.2sg) --- OE sċēawa}}\label{show-iptv.2sg-oe-sux10bux113awa}

\index[iv]{02oe@\ivlangheader{Old English}{West Saxon normalization unmarked}!sceawa@\iventry{sċēawa}{}}
\index[iv]{03pgmc@\ivlangheader{Proto-Germanic}{}!skawo@\iventry{*skáwô}{}}
\index[iv]{03pgmc@\ivlangheader{Proto-Germanic}{}!skawona@\iventry{*skawōną}{}}
\index[iv]{10ohg@\ivlangheader{Old High German}{}!scouwon@\iventry{scouwōn}{}}
\index[iv]{11ofris@\ivlangheader{Old Frisian}{}!skawia@\iventry{skawia}{}}
\index[iv]{13os@\ivlangheader{Old Saxon}{}!skawon@\iventry{skawōn}{}}

Derivation: citation reconstruction \emph{*skawōną}; form followed here
\emph{*skáwô} \textgreater{} \emph{sċēawa} (late analogy).

\subsection{Derivation trace}\label{derivation-trace-163}

Proto input: \emph{*skáwô}

\begingroup
\setlength{\fboxsep}{6pt}

\noindent\fbox{%
\begin{minipage}{0.97\linewidth}
\small
\begin{tabular}{@{}>{\raggedright\arraybackslash}p{0.359\linewidth}@{\hspace{0.030\linewidth}}>{\raggedright\arraybackslash}p{0.611\linewidth}@{}}
\begin{minipage}[t]{\linewidth}
\centering\textbf{Earlier Germanic changes}\par
\vspace{0.35em}
\raggedright
\centering\textbf{Northwest and West Germanic}\par
\raggedright
\vspace{0.2em}
\raggedright [no change]\par
\vspace{0.6em}
\centering\textbf{Early Anglo-Frisian}\par
\raggedright
\vspace{0.2em}
\raggedright [no change]\par
\end{minipage}
&
\begin{minipage}[t]{\linewidth}
\centering\textbf{Old English changes}\par
\vspace{0.35em}
\raggedright
\begin{tabularx}{\linewidth}{@{}>{\raggedright\arraybackslash}X@{\hspace{0.55em}}>{\raggedright\arraybackslash}p{0.14\linewidth}@{\hspace{0.25em}}}
\mbox{OE Aw Long Diphthong} & \emph{*skḗawô} \\
\mbox{OE Sk Palatalization} & \emph{*ʃḗawô} \\
OE Unstressed Long Vowel Shortening & \emph{*ʃḗawa} \\
\end{tabularx}
\end{minipage}
\\
\end{tabular}
\end{minipage}%
} \endgroup

Old English form: \emph{sċēawa}

\subsection{Reconstruction and comparative
evidence}\label{reconstruction-and-comparative-evidence-143}

Orel reconstructs the verb as \Recon{skawōjanan} `show' and cites Old
English \emph{sceáwian} `show' beside Old Frisian \emph{skawia} `show',
Old Saxon \emph{skawōn} `show', and Old High German \emph{scouwōn}
`show' (\citeproc{ref-Orel2003}{Orel 2003, 337}). This is
infinitive-family evidence, not a quotation of the selected imperative.
CAPR assembles imperative singular \Recon{skáwô} `show (iptv.2sg)',
using the finite class-II cell with imperative -a \textless{} \emph{*-ō}
(\citeproc{ref-RingeTaylor2014}{Ringe and Taylor 2014, 314}).

The imperative singular provides the direct comparison with the Old
English form. The lexical history still belongs to the \emph{sceáwian}
`show' verb, but the cell compared here isolates the finite \emph{-a}
outcome more clearly than the citation form does.

\subsection{Old English evidence}\label{old-english-evidence-144}

The held fourth edition of Bright's reader contains \emph{scēawian}
`show' and the imperative singular \emph{scēawa} `show' in the Blickling
homily (\citeproc{ref-Bright1917}{Bright 1917, 71}). The form treated
here is therefore an attested finite paradigm form, not a reconstructed
convenience form; it is not attributed to an unverified passage in the
separate Cassidy/Ringler edition.

The spelling used in this entry is normalized \emph{sċēawa} `show',
while Bright's cited text gives source spelling \emph{scēawa} `show'.
The ordinary Old English headword remains \emph{scēawian} `show';
\emph{sċēawa} `show' is the imperative singular chosen beside it.

\subsection{Development to Old
English}\label{development-to-old-english-142}

The infinitive entry and \hyperref[rule-OEAwLongDiphthong]{SC034
OEAwLongDiphthong} distinguish the dictionary singleton stem from
Campbell and Brunner's inherited-geminate analysis. The same
computational proxy is retained here; successful finite output does not
resolve that prehistoric root history (\citeproc{ref-Orel2003}{Orel
2003, 337}; \citeproc{ref-Campbell1959}{Campbell 1959, 45};
\citeproc{ref-SieversBrunner1965}{Brunner 1965, 49}).

Ringe and Taylor's class-II weak imperative singular -a \textless{}
\emph{*-ō} supports the finite ending that yields \emph{sċēawa} `show'
(\citeproc{ref-RingeTaylor2014}{Ringe and Taylor 2014, 314}). The ending
comparison is therefore justified independently of the unresolved root
quantity history. The inherited \emph{late\_analogy} corpus
classification is not changed by this root-only adjudication; it does
not make the proxy a demonstrated analogical sound law.

\subsection{Paradigm comparison}\label{paradigm-comparison-22}

The comparison below sets the relevant forms side by side.

{\def\LTcaptype{none} % do not increment counter
\begin{longtable}[]{@{}
  >{\raggedright\arraybackslash}p{(\linewidth - 8\tabcolsep) * \real{0.2000}}
  >{\raggedright\arraybackslash}p{(\linewidth - 8\tabcolsep) * \real{0.2000}}
  >{\raggedright\arraybackslash}p{(\linewidth - 8\tabcolsep) * \real{0.2000}}
  >{\raggedright\arraybackslash}p{(\linewidth - 8\tabcolsep) * \real{0.2000}}
  >{\raggedright\arraybackslash}p{(\linewidth - 8\tabcolsep) * \real{0.2000}}@{}}
\toprule\noalign{}
\begin{minipage}[b]{\linewidth}\raggedright
PGmc cell / interpretation
\end{minipage} & \begin{minipage}[b]{\linewidth}\raggedright
Candidate input
\end{minipage} & \begin{minipage}[b]{\linewidth}\raggedright
Old English outcome or comparison
\end{minipage} & \begin{minipage}[b]{\linewidth}\raggedright
OE comparison form
\end{minipage} & \begin{minipage}[b]{\linewidth}\raggedright
Result
\end{minipage} \\
\midrule\noalign{}
\endhead
\bottomrule\noalign{}
\endlastfoot
lexeme-level infinitive &
{\emph{*skáwōjaną}}\index[iv]{03pgmc@\ivlangheader{Proto-Germanic}{}!skawojana@\iventry{*skáwōjaną}{}}
& regular output
{\emph{sċēawian}}\index[iv]{02oe@\ivlangheader{Old English}{West Saxon normalization unmarked}!sceawian@\iventry{sċēawian}{}}
&
{\emph{scēawian}}\index[iv]{02oe@\ivlangheader{Old English}{West Saxon normalization unmarked}!sceawian@\iventry{scēawian}{}}
`show' & ordinary dictionary headword of the verb, but not the finite
form compared here \\
imperative singular &
{\emph{*skáwô}}\index[iv]{03pgmc@\ivlangheader{Proto-Germanic}{}!skawo@\iventry{*skáwô}{}}
& regular output
{\emph{sċēawa}}\index[iv]{02oe@\ivlangheader{Old English}{West Saxon normalization unmarked}!sceawa@\iventry{sċēawa}{}}
`show' &
{\emph{scēawa}}\index[iv]{02oe@\ivlangheader{Old English}{West Saxon normalization unmarked}!sceawa@\iventry{scēawa}{}}
`show' / normalized
{\emph{sċēawa}}\index[iv]{02oe@\ivlangheader{Old English}{West Saxon normalization unmarked}!sceawa@\iventry{sċēawa}{}}
`show' & exact match between the derivational input and the Old English
form here \\
present third singular companion &
{\emph{*skáwōθi}}\index[iv]{03pgmc@\ivlangheader{Proto-Germanic}{}!skawothi@\iventry{*skáwōθi}{}}
& regular output
{\emph{sċēawaþ}}\index[iv]{02oe@\ivlangheader{Old English}{West Saxon normalization unmarked}!sceawath@\iventry{sċēawaþ}{}}
`show' &
{\emph{sċēawaþ}}\index[iv]{02oe@\ivlangheader{Old English}{West Saxon normalization unmarked}!sceawath@\iventry{sċēawaþ}{}}
`show' & useful family control, but not the target of this entry \\
\end{longtable}
}

\section{\texorpdfstring{span --- OE
\emph{spanne}}{span --- OE spanne}}\label{span-oe-spanne}

\index[iv]{02oe@\ivlangheader{Old English}{West Saxon normalization unmarked}!spanne@\iventry{spanne}{}}
\index[iv]{03pgmc@\ivlangheader{Proto-Germanic}{}!spannai@\iventry{*spánnai}{}}
\index[iv]{03pgmc@\ivlangheader{Proto-Germanic}{}!spanno@\iventry{*spannō}{}}

Derivation: citation reconstruction \emph{*spannō}; form followed here
\emph{*spánnai} \textgreater{} \emph{spanne} (late analogy).

\subsection{Derivation trace}\label{derivation-trace-164}

Proto input: \emph{*spánnai}

\begingroup
\setlength{\fboxsep}{6pt}

\noindent\fbox{%
\begin{minipage}{0.97\linewidth}
\small
\begin{tabular}{@{}>{\raggedright\arraybackslash}p{0.485\linewidth}@{\hspace{0.030\linewidth}}>{\raggedright\arraybackslash}p{0.485\linewidth}@{}}
\begin{minipage}[t]{\linewidth}
\centering\textbf{Earlier Germanic changes}\par
\vspace{0.35em}
\raggedright
\centering\textbf{Northwest and West Germanic}\par
\raggedright
\vspace{0.2em}
\begin{tabularx}{\linewidth}{@{}>{\raggedright\arraybackslash}X@{\hspace{0.55em}}>{\raggedright\arraybackslash}p{0.14\linewidth}@{\hspace{0.25em}}}
PNWGmc Unstressed Ai Monophthongization & \emph{*spánnē} \\
\end{tabularx}
\vspace{0.6em}
\centering\textbf{Early Anglo-Frisian}\par
\raggedright
\vspace{0.2em}
\raggedright [no change]\par
\end{minipage}
&
\begin{minipage}[t]{\linewidth}
\centering\textbf{Old English changes}\par
\vspace{0.35em}
\raggedright
\begin{tabularx}{\linewidth}{@{}>{\raggedright\arraybackslash}X@{\hspace{0.55em}}>{\raggedright\arraybackslash}p{0.14\linewidth}@{\hspace{0.25em}}}
OE Unstressed Long Vowel Shortening & \emph{*spánne} \\
\end{tabularx}
\end{minipage}
\\
\end{tabular}
\end{minipage}%
} \endgroup

Old English form: \emph{spanne}

\subsection{Reconstruction and comparative
evidence}\label{reconstruction-and-comparative-evidence-144}

Seebold gives Old English
{\emph{spann}}\index[iv]{02oe@\ivlangheader{Old English}{West Saxon normalization unmarked}!spann@\iventry{spann}{}}
`span' under this noun family (\citeproc{ref-Seebold1970}{Seebold 1970,
450}). The form followed here,
\Recon{spánnai}\index[iv]{03pgmc@\ivlangheader{Proto-Germanic}{}!spannai@\iventry{*spánnai}{}}
`span', is therefore not a rival headword, but the specific dative
singular form compared on the model of the feminine ō-stem paradigm
(\citeproc{ref-SieversBrunner1965}{Brunner 1965, § 252};
\citeproc{ref-SieversBrunner1965}{1965, § 255.2}).

\subsection{Old English evidence}\label{old-english-evidence-145}

The reviewed lexicographic evidence more directly supports the citation
noun
{\emph{spann}}\index[iv]{02oe@\ivlangheader{Old English}{West Saxon normalization unmarked}!spann@\iventry{spann}{}}
`span' than the exact form
{\emph{spanne}}\index[iv]{02oe@\ivlangheader{Old English}{West Saxon normalization unmarked}!spanne@\iventry{spanne}{}}
`span'. Clark Hall gives
{\emph{spann}}\index[iv]{02oe@\ivlangheader{Old English}{West Saxon normalization unmarked}!spann@\iventry{spann}{}}
`span' (\citeproc{ref-ClarkHall1960}{Clark Hall 1960, 286}), and
{\emph{spanne}}\index[iv]{02oe@\ivlangheader{Old English}{West Saxon normalization unmarked}!spanne@\iventry{spanne}{}}
`span' is accordingly treated as the regular dative singular form
compared here rather than as a dictionary headword.

\subsection{Development to Old
English}\label{development-to-old-english-143}

Citation
\Recon{spannō}\index[iv]{03pgmc@\ivlangheader{Proto-Germanic}{}!spanno@\iventry{*spannō}{}}
`span' yields {\emph{span}} `span'. The oblique cell
\Recon{spánnai}\index[iv]{03pgmc@\ivlangheader{Proto-Germanic}{}!spannai@\iventry{*spánnai}{}}
`span' therefore supplies the conservative comparison form: it preserves
the medial geminate and yields
{\emph{spanne}}\index[iv]{02oe@\ivlangheader{Old English}{West Saxon normalization unmarked}!spanne@\iventry{spanne}{}}
`span', while citation
\Recon{spannō}\index[iv]{03pgmc@\ivlangheader{Proto-Germanic}{}!spanno@\iventry{*spannō}{}}
gives the nominative background form.

\subsection{Paradigm comparison}\label{paradigm-comparison-23}

The comparison below sets the citation form beside the dative singular
form compared here.

{\def\LTcaptype{none} % do not increment counter
\begin{longtable}[]{@{}
  >{\raggedright\arraybackslash}p{(\linewidth - 8\tabcolsep) * \real{0.2000}}
  >{\raggedright\arraybackslash}p{(\linewidth - 8\tabcolsep) * \real{0.2000}}
  >{\raggedright\arraybackslash}p{(\linewidth - 8\tabcolsep) * \real{0.2000}}
  >{\raggedright\arraybackslash}p{(\linewidth - 8\tabcolsep) * \real{0.2000}}
  >{\raggedright\arraybackslash}p{(\linewidth - 8\tabcolsep) * \real{0.2000}}@{}}
\toprule\noalign{}
\begin{minipage}[b]{\linewidth}\raggedright
PGmc cell / interpretation
\end{minipage} & \begin{minipage}[b]{\linewidth}\raggedright
Candidate input
\end{minipage} & \begin{minipage}[b]{\linewidth}\raggedright
OE output or comparison
\end{minipage} & \begin{minipage}[b]{\linewidth}\raggedright
OE comparison form
\end{minipage} & \begin{minipage}[b]{\linewidth}\raggedright
Result
\end{minipage} \\
\midrule\noalign{}
\endhead
\bottomrule\noalign{}
\endlastfoot
citation nominative singular &
{\emph{*spannō}}\index[iv]{03pgmc@\ivlangheader{Proto-Germanic}{}!spanno@\iventry{*spannō}{}}
& regular output: {\emph{span}} `span' &
{\emph{spann}}\index[iv]{02oe@\ivlangheader{Old English}{West Saxon normalization unmarked}!spann@\iventry{spann}{}}
`span' & useful citation-form background, but not the form compared
here \\
dative singular compared here &
{\emph{*spánnai}}\index[iv]{03pgmc@\ivlangheader{Proto-Germanic}{}!spannai@\iventry{*spánnai}{}}
& regular output: {\emph{spanne}} `span' &
{\emph{spanne}}\index[iv]{02oe@\ivlangheader{Old English}{West Saxon normalization unmarked}!spanne@\iventry{spanne}{}}
`span' & exact match for that conservative form \\
\end{longtable}
}

\section{\texorpdfstring{straw --- OE
\emph{strēaw}}{straw --- OE strēaw}}\label{straw-oe-strux113aw}

\index[iv]{02oe@\ivlangheader{Old English}{West Saxon normalization unmarked}!streaw@\iventry{strēaw}{}}
\index[iv]{03pgmc@\ivlangheader{Proto-Germanic}{}!strawa@\iventry{*stráwą}{}}

Derivation: \emph{*stráwą} yields regular \Recon{strēa}; the Old English
form here is \emph{strēaw} (late analogy).

\subsection{Derivation trace}\label{derivation-trace-165}

Proto input: \emph{*stráwą}

\begingroup
\setlength{\fboxsep}{6pt}

\noindent\fbox{%
\begin{minipage}{0.97\linewidth}
\small
\begin{tabular}{@{}>{\raggedright\arraybackslash}p{0.485\linewidth}@{\hspace{0.030\linewidth}}>{\raggedright\arraybackslash}p{0.485\linewidth}@{}}
\begin{minipage}[t]{\linewidth}
\centering\textbf{Earlier Germanic changes}\par
\vspace{0.35em}
\raggedright
\centering\textbf{Northwest and West Germanic}\par
\raggedright
\vspace{0.2em}
\begin{tabularx}{\linewidth}{@{}>{\raggedright\arraybackslash}X@{\hspace{0.55em}}>{\raggedright\arraybackslash}p{0.14\linewidth}@{\hspace{0.25em}}}
\mbox{PWGmc Endingless Aw} & \emph{*stráu} \\
\end{tabularx}
\vspace{0.6em}
\centering\textbf{Early Anglo-Frisian}\par
\raggedright
\vspace{0.2em}
\raggedright [no change]\par
\end{minipage}
&
\begin{minipage}[t]{\linewidth}
\centering\textbf{Old English changes}\par
\vspace{0.35em}
\raggedright
\begin{tabularx}{\linewidth}{@{}>{\raggedright\arraybackslash}X@{\hspace{0.55em}}>{\raggedright\arraybackslash}p{0.14\linewidth}@{\hspace{0.25em}}}
\mbox{OE Au Brightening} & \emph{*stráeu} \\
\mbox{OE Diphthong Leveling} & \Recon{strēa} \\
\end{tabularx}
\end{minipage}
\\
\end{tabular}
\end{minipage}%
} \endgroup

Regular outcome: \Recon{strēa}

Old English form: \emph{strēaw}

\subsection{Reconstruction and comparative
evidence}\label{reconstruction-and-comparative-evidence-145}

Kroonen cites the neuter stem
\Recon{strawa-}\index[iv]{03pgmc@\ivlangheader{Proto-Germanic}{}!strawa@\iventry{*strawa-}{}}
`straw' and Old English
{\emph{strēa(w)}}\index[iv]{02oe@\ivlangheader{Old English}{West Saxon normalization unmarked}!streaw@\iventry{strēa(w)}{}}
`straw' (\citeproc{ref-Kroonen2013}{Kroonen 2013, 484}). The selected
\Recon{stráwą}\index[iv]{03pgmc@\ivlangheader{Proto-Germanic}{}!strawa@\iventry{*stráwą}{}}
`straw' is CAPR's accented nominative/accusative singular, not Kroonen's
verbatim full-word quotation. His derivation of the noun from the strew
formation concerns genealogical direction; it does not explain the later
restoration of the glide (\citeproc{ref-Kroonen2013}{Kroonen 2013,
484}).

\subsection{Old English evidence}\label{old-english-evidence-146}

The comparison remains attested
{\emph{strēaw}}\index[iv]{02oe@\ivlangheader{Old English}{West Saxon normalization unmarked}!streaw@\iventry{strēaw}{}}
`straw', not reconstructed \Recon{strēa} `straw'. Campbell discusses
restoration of the glide and the compound
{\emph{strēaberige}}\index[iv]{02oe@\ivlangheader{Old English}{West Saxon normalization unmarked}!streaberige@\iventry{strēaberige}{}}
`strawberry' (\citeproc{ref-Campbell1959}{Campbell 1959, 46--47,
232--33}). Brunner explicitly treats the regular free simplex and short
obliques as unattested; his paradigm note limits the West Saxon straw
evidence to nominative/accusative singular and a doubtful plural. The
compound preserves an endingless pattern, not an independently attested
free noun or a complete surviving regular paradigm
(\citeproc{ref-SieversBrunner1965}{Brunner 1965, 49, 110, 204--5};
\citeproc{ref-RingeTaylor2014}{Ringe and Taylor 2014, 173, 387}).

\subsection{Development to Old
English}\label{development-to-old-english-144}

The regular endingless branch loses the final vowel and vocalizes final
\Recon{w}: \Recon{stráwą} `straw' \textgreater{} \Recon{stráu} `straw'.
The narrow fragment \hyperref[rule-PWGmcEndinglessAw]{SC111
PWGmcEndinglessAw} combines those operations; its computational holding
site does not establish their order against unrelated changes. Its
product feeds \hyperref[rule-OEAuBrightening]{SC030 OEAuBrightening},
giving \Recon{stráeu} `straw', and
\hyperref[rule-OEDiphthongLeveling]{SC032 OEDiphthongLeveling}, giving
regular \Recon{strēa} `straw' (\citeproc{ref-Campbell1959}{Campbell
1959, 46--47}; \citeproc{ref-RingeTaylor2014}{Ringe and Taylor 2014,
172--73}).

Attested
{\emph{strēaw}}\index[iv]{02oe@\ivlangheader{Old English}{West Saxon normalization unmarked}!streaw@\iventry{strēaw}{}}
`straw' combines that long diphthong with a glide restored from
inflected forms. Reciprocally, the long diphthong spreads into the
formerly short oblique stem. This mutual levelling explains the
long-plus-glide form without inventing a sound law that preserves final
singleton \Recon{w} in the regular endingless branch
(\citeproc{ref-Campbell1959}{Campbell 1959, 232--33};
\citeproc{ref-SieversBrunner1965}{Brunner 1965, 204--5};
\citeproc{ref-RingeTaylor2014}{Ringe and Taylor 2014, 387}). The cascade
stops before that levelling: its reconstructed output remains an
explicit mismatch, while the attested comparison and derivational input
are unchanged.

\subsection{Paradigm comparison}\label{paradigm-comparison-24}

The table distinguishes source-supported historical comparisons from the
one selected compiled derivation; it is not a set of newly compiled
oblique inputs.

{\def\LTcaptype{none} % do not increment counter
\begin{longtable}[]{@{}
  >{\raggedright\arraybackslash}p{(\linewidth - 4\tabcolsep) * \real{0.3333}}
  >{\raggedright\arraybackslash}p{(\linewidth - 4\tabcolsep) * \real{0.3333}}
  >{\raggedright\arraybackslash}p{(\linewidth - 4\tabcolsep) * \real{0.3333}}@{}}
\toprule\noalign{}
\begin{minipage}[b]{\linewidth}\raggedright
Cell or evidence
\end{minipage} & \begin{minipage}[b]{\linewidth}\raggedright
Regular development
\end{minipage} & \begin{minipage}[b]{\linewidth}\raggedright
Attested comparison and qualification
\end{minipage} \\
\midrule\noalign{}
\endhead
\bottomrule\noalign{}
\endlastfoot
selected nominative/accusative singular & \Recon{stráwą} \textgreater{}
\Recon{strēa} &
{\emph{strēaw}}\index[iv]{02oe@\ivlangheader{Old English}{West Saxon normalization unmarked}!streaw@\iventry{strēaw}{}}
`straw': restored glide; regular free simplex unattested \\
short oblique stem, exemplified by genitive singular &
\Recon{strawes}\index[iv]{02oe@\ivlangheader{Old English}{West Saxon normalization unmarked}!strawes@\iventry{*strawes}{}}
`of straw'; corresponding short dative also reconstructed & Brunner's
regular short obliques are unattested; long quantity is extended from
the endingless cells, not demonstrated by a surviving short straw
paradigm \\
remodelled singular and limited plural evidence & long diphthong plus
restored glide & West Saxon
{\emph{strēaw}}\index[iv]{02oe@\ivlangheader{Old English}{West Saxon normalization unmarked}!streaw@\iventry{strēaw}{}}
`straw'; Brunner records plural
{\emph{strewu}}\index[iv]{02oe@\ivlangheader{Old English}{West Saxon normalization unmarked}!strewu@\iventry{strewu}{}}
`straws' with a question mark, not as a secure complete paradigm \\
compound witness & endingless pattern &
{\emph{strēaberige}}\index[iv]{02oe@\ivlangheader{Old English}{West Saxon normalization unmarked}!streaberige@\iventry{strēaberige}{}}
`strawberry', not an attested free
\Recon{strēa}\index[iv]{02oe@\ivlangheader{Old English}{West Saxon normalization unmarked}!strea@\iventry{*strēa}{}}
`straw' \\
\end{longtable}
}

These qualifications follow Brunner's paradigm and attestation
discussion and Ringe and Taylor's explicit mutual levelling account
(\citeproc{ref-SieversBrunner1965}{Brunner 1965, 49, 110, 204--5};
\citeproc{ref-RingeTaylor2014}{Ringe and Taylor 2014, 173, 387}). The
noun/verb genealogical direction remains a separate qualified question;
this entry adopts neither a new ancestry nor an analogical transducer
rule.

\section{\texorpdfstring{thistle --- OE
\emph{þistles}}{thistle --- OE þistles}}\label{thistle-oe-uxfeistles}

\index[iv]{02oe@\ivlangheader{Old English}{West Saxon normalization unmarked}!hrefn@\iventry{hrefn}{}}
\index[iv]{02oe@\ivlangheader{Old English}{West Saxon normalization unmarked}!thistles@\iventry{þistles}{}}
\index[iv]{03pgmc@\ivlangheader{Proto-Germanic}{}!thestilaz@\iventry{*θéstilaz}{}}
\index[iv]{03pgmc@\ivlangheader{Proto-Germanic}{}!thistilas@\iventry{*θístilas}{}}

Derivation: citation reconstruction \emph{*θéstilaz}; form followed here
\emph{*θístilas} \textgreater{} \emph{þistles} (late analogy).

\subsection{Derivation trace}\label{derivation-trace-166}

Proto input: \emph{*θístilas}

\begingroup
\setlength{\fboxsep}{6pt}

\noindent\fbox{%
\begin{minipage}{0.97\linewidth}
\small
\begin{tabular}{@{}>{\raggedright\arraybackslash}p{0.359\linewidth}@{\hspace{0.030\linewidth}}>{\raggedright\arraybackslash}p{0.611\linewidth}@{}}
\begin{minipage}[t]{\linewidth}
\centering\textbf{Earlier Germanic changes}\par
\vspace{0.35em}
\raggedright
\centering\textbf{Northwest and West Germanic}\par
\raggedright
\vspace{0.2em}
\begin{tabularx}{\linewidth}{@{}>{\raggedright\arraybackslash}X@{\hspace{0.55em}}>{\raggedright\arraybackslash}p{0.24\linewidth}@{\hspace{0.25em}}}
\mbox{PNWGmc I Lowering} & \emph{*θéstilas} \\
\end{tabularx}
\vspace{0.6em}
\centering\textbf{Early Anglo-Frisian}\par
\raggedright
\vspace{0.2em}
\raggedright [no change]\par
\end{minipage}
&
\begin{minipage}[t]{\linewidth}
\centering\textbf{Old English changes}\par
\vspace{0.35em}
\raggedright
\begin{tabularx}{\linewidth}{@{}>{\raggedright\arraybackslash}X@{\hspace{0.55em}}>{\raggedright\arraybackslash}p{0.20\linewidth}@{\hspace{0.25em}}}
\mbox{EAF Brightening Unstressed} & \emph{*θéstilæs} \\
\mbox{OE I Umlaut} & \emph{*θistilæs} \\
\mbox{OE High Vowel Apocope} & \emph{*θistlæs} \\
\mbox{OE Unstressed AE Merger} & \emph{*θistles} \\
\end{tabularx}
\end{minipage}
\\
\end{tabular}
\end{minipage}%
} \endgroup

Old English form: \emph{þistles}

\subsection{Reconstruction and comparative
evidence}\label{reconstruction-and-comparative-evidence-146}

Orel prints \emph{*þe(x)stilaz} `thistle' for the lexeme
(\citeproc{ref-Orel2003}{Orel 2003, 458}). The comparative label
\Recon{θéstilaz}\index[iv]{03pgmc@\ivlangheader{Proto-Germanic}{}!thestilaz@\iventry{*θéstilaz}{}}
`thistle' therefore remains in view as the lexeme-level headword, while
the derivational input
\Recon{θístilas}\index[iv]{03pgmc@\ivlangheader{Proto-Germanic}{}!thistilas@\iventry{*θístilas}{}}
`thistle' is a specific genitive singular cell.

\subsection{Old English evidence}\label{old-english-evidence-147}

The ordinary simplex headword tradition is broken
{\emph{þistel}}\index[iv]{02oe@\ivlangheader{Old English}{West Saxon normalization unmarked}!thistel@\iventry{þistel}{}}
`thistle' /
{\emph{ðistel}}\index[iv]{02oe@\ivlangheader{Old English}{West Saxon normalization unmarked}!thistel@\iventry{ðistel}{}}
`thistle'. Clark Hall gives
{\emph{ðistel}}\index[iv]{02oe@\ivlangheader{Old English}{West Saxon normalization unmarked}!thistel@\iventry{ðistel}{}}
as the noun headword (\citeproc{ref-ClarkHall1960}{Clark Hall 1960,
326}). The Old English form here here is the genitive singular
{\emph{þistles}}\index[iv]{02oe@\ivlangheader{Old English}{West Saxon normalization unmarked}!thistles@\iventry{þistles}{}}
`thistle', which preserves the same stem in an oblique form where the
cluster is medial.

\subsection{Development to Old
English}\label{development-to-old-english-145}

Campbell's discussion of cluster nouns shows the contrast clearly.
Simplex forms often develop a parasite vowel in word-final obstruent +
sonorant clusters, while comparable medial clusters remain unbroken; his
examples include \emph{hrefn},
{\emph{tacn}}\index[iv]{02oe@\ivlangheader{Old English}{West Saxon normalization unmarked}!tacn@\iventry{tacn}{}}
`token, sign',
{\emph{wépn}}\index[iv]{02oe@\ivlangheader{Old English}{West Saxon normalization unmarked}!wépn@\iventry{wépn}{}}
`weapon', and
{\emph{botm}}\index[iv]{02oe@\ivlangheader{Old English}{West Saxon normalization unmarked}!botm@\iventry{botm}{}}
`bottom' beside forms with parasitic vowels elsewhere in the same
lexical class (\citeproc{ref-Campbell1959}{Campbell 1959, 151}). The
genitive singular
\Recon{θístilas}\index[iv]{03pgmc@\ivlangheader{Proto-Germanic}{}!thistilas@\iventry{*θístilas}{}}
`thistle' therefore supplies the conservative comparison form: the
cluster is medial and the regular development yields {\emph{þistles}}
`thistle', while the simplex nominative belongs to the broken headword
tradition
{\emph{þistel}}\index[iv]{02oe@\ivlangheader{Old English}{West Saxon normalization unmarked}!thistel@\iventry{þistel}{}}
`thistle'.

\subsection{Paradigm comparison}\label{paradigm-comparison-25}

The comparison below sets the relevant forms side by side. It shows the
contrast between the citation form and the genitive singular cell.

{\def\LTcaptype{none} % do not increment counter
\begin{longtable}[]{@{}
  >{\raggedright\arraybackslash}p{(\linewidth - 8\tabcolsep) * \real{0.2000}}
  >{\raggedright\arraybackslash}p{(\linewidth - 8\tabcolsep) * \real{0.2000}}
  >{\raggedright\arraybackslash}p{(\linewidth - 8\tabcolsep) * \real{0.2000}}
  >{\raggedright\arraybackslash}p{(\linewidth - 8\tabcolsep) * \real{0.2000}}
  >{\raggedright\arraybackslash}p{(\linewidth - 8\tabcolsep) * \real{0.2000}}@{}}
\toprule\noalign{}
\begin{minipage}[b]{\linewidth}\raggedright
PGmc cell / interpretation
\end{minipage} & \begin{minipage}[b]{\linewidth}\raggedright
Candidate input
\end{minipage} & \begin{minipage}[b]{\linewidth}\raggedright
OE output or comparison
\end{minipage} & \begin{minipage}[b]{\linewidth}\raggedright
OE comparison form
\end{minipage} & \begin{minipage}[b]{\linewidth}\raggedright
Result
\end{minipage} \\
\midrule\noalign{}
\endhead
\bottomrule\noalign{}
\endlastfoot
citation nominative singular &
{\emph{*θéstilaz}}\index[iv]{03pgmc@\ivlangheader{Proto-Germanic}{}!thestilaz@\iventry{*θéstilaz}{}}
& computed output: {\emph{þistl}} &
{\emph{þistel}}\index[iv]{02oe@\ivlangheader{Old English}{West Saxon normalization unmarked}!thistel@\iventry{þistel}{}}
`thistle' & useful citation-form background, but not the Old English
form here \\
genitive singular &
{\emph{*θístilas}}\index[iv]{03pgmc@\ivlangheader{Proto-Germanic}{}!thistilas@\iventry{*θístilas}{}}
& computed output: {\emph{þistles}} `thistle' &
{\emph{þistles}}\index[iv]{02oe@\ivlangheader{Old English}{West Saxon normalization unmarked}!thistles@\iventry{þistles}{}}
`thistle' & exact match for the conservative cell \\
\end{longtable}
}

\clearpage

\chapter{Reconstructed Old English
comparators}\label{reconstructed-old-english-comparators}

Direct attestation does not supply the required comparator in these
entries. The target is an explicitly reconstructed Old English form and
carries the corresponding evidential burden.

\section{\texorpdfstring{knob --- OE
\emph{*cnobba}}{knob --- OE *cnobba}}\label{knob-oe-cnobba}

\index[iv]{02oe@\ivlangheader{Old English}{West Saxon normalization unmarked}!cnobba@\iventry{*cnobba}{}}
\index[iv]{03pgmc@\ivlangheader{Proto-Germanic}{}!knubbo@\iventry{*knúbbô}{}}
\index[iv]{03pgmc@\ivlangheader{Proto-Germanic}{}!knuppaz@\iventry{*knúppaz}{}}

Derivation: citation reconstruction \emph{*knúppaz}; form followed here
\emph{*knúbbô} \textgreater{} \emph{*cnobba} (reconstructed Old English
comparator).

\subsection{Derivation trace}\label{derivation-trace-167}

Proto input: \emph{*knúbbô}

\begingroup
\setlength{\fboxsep}{6pt}

\noindent\fbox{%
\begin{minipage}{0.97\linewidth}
\small
\begin{tabular}{@{}>{\raggedright\arraybackslash}p{0.359\linewidth}@{\hspace{0.030\linewidth}}>{\raggedright\arraybackslash}p{0.611\linewidth}@{}}
\begin{minipage}[t]{\linewidth}
\centering\textbf{Earlier Germanic changes}\par
\vspace{0.35em}
\raggedright
\centering\textbf{Northwest and West Germanic}\par
\raggedright
\vspace{0.2em}
\begin{tabularx}{\linewidth}{@{}>{\raggedright\arraybackslash}X@{\hspace{0.55em}}>{\raggedright\arraybackslash}p{0.24\linewidth}@{\hspace{0.25em}}}
\mbox{PNWGmc U Lowering} & \emph{*knóbbô} \\
\end{tabularx}
\vspace{0.6em}
\centering\textbf{Early Anglo-Frisian}\par
\raggedright
\vspace{0.2em}
\raggedright [no change]\par
\end{minipage}
&
\begin{minipage}[t]{\linewidth}
\centering\textbf{Old English changes}\par
\vspace{0.35em}
\raggedright
\begin{tabularx}{\linewidth}{@{}>{\raggedright\arraybackslash}X@{\hspace{0.55em}}>{\raggedright\arraybackslash}p{0.14\linewidth}@{\hspace{0.25em}}}
OE Unstressed Long Vowel Shortening & \emph{*knóbba} \\
\end{tabularx}
\end{minipage}
\\
\end{tabular}
\end{minipage}%
} \endgroup

Old English form: \emph{*cnobba}

\subsection{Reconstruction and comparative
evidence}\label{reconstruction-and-comparative-evidence-147}

The wider knob-family is not uniform. Kroonen's discussion of the
Germanic n-stems points to related voiced and voiceless branches within
this group (\citeproc{ref-Kroonen2011}{Kroonen 2011, 297}). The citation
reconstruction \Recon{knúppaz} `knob' represents the broader cognate-set
headword, while the form followed here, \Recon{knúbbô} `knob',
represents the voiced weak-noun branch treated here.

The Old English record is uneven. The better attested OE material
belongs to the voiceless branch, but the present entry represents the
reconstructed OE form that would continue the voiced branch behind later
English knob.

\subsection{Old English evidence}\label{old-english-evidence-148}

Clark Hall preserves Old English evidence of the
{\emph{cnoppa}}\index[iv]{02oe@\ivlangheader{Old English}{West Saxon normalization unmarked}!cnoppa@\iventry{cnoppa}{}}
`knob' type (\citeproc{ref-ClarkHall1960}{Clark Hall 1960, 79}). Those
forms are genuine Old English evidence, but they belong to the voiceless
branch of the family.

The target \Recon{cnobba} `knob' is a reconstructed Old English form,
not a directly attested one. I use it for the voiced branch because
attested
{\emph{cnoppa}}\index[iv]{02oe@\ivlangheader{Old English}{West Saxon normalization unmarked}!cnoppa@\iventry{cnoppa}{}}
`knob' continues the voiceless branch and therefore represents a
different prehistory.

\subsection{Development to Old
English}\label{development-to-old-english-146}

From the weak-noun form followed here, \Recon{knúbbô} `knob', the
regular Old English outcome is \Recon{cnobba} `knob', with
Proto-Germanic \emph{kn-} represented in Old English as \emph{cn-} and
with the expected weak-noun ending.

The entry therefore does not claim that \Recon{cnobba} `knob' is
attested. Its claim is different: if the voiced weak-noun branch is the
one to be represented, then \Recon{cnobba} is the regular Old English
form corresponding to that branch.

\subsection{Reconstruction status}\label{reconstruction-status}

{\def\LTcaptype{none} % do not increment counter
\begin{longtable}[]{@{}
  >{\raggedright\arraybackslash}p{(\linewidth - 4\tabcolsep) * \real{0.3333}}
  >{\raggedright\arraybackslash}p{(\linewidth - 4\tabcolsep) * \real{0.3333}}
  >{\raggedright\arraybackslash}p{(\linewidth - 4\tabcolsep) * \real{0.3333}}@{}}
\toprule\noalign{}
\begin{minipage}[b]{\linewidth}\raggedright
Form
\end{minipage} & \begin{minipage}[b]{\linewidth}\raggedright
Status
\end{minipage} & \begin{minipage}[b]{\linewidth}\raggedright
Relevance to this entry
\end{minipage} \\
\midrule\noalign{}
\endhead
\bottomrule\noalign{}
\endlastfoot
{\emph{\emph{*knúbbô} \textgreater{} \emph{*cnobba}}} & reconstructed OE
form; regular derivation & reconstructed Old English form compared
here \\
{\emph{cnopp}}\index[iv]{02oe@\ivlangheader{Old English}{West Saxon normalization unmarked}!cnopp@\iventry{cnopp}{}}
/
{\emph{cnoppa}}\index[iv]{02oe@\ivlangheader{Old English}{West Saxon normalization unmarked}!cnoppa@\iventry{cnoppa}{}}
& attested OE branch & important control form, but belongs to the
voiceless branch \\
{\emph{cnæp}}\index[iv]{02oe@\ivlangheader{Old English}{West Saxon normalization unmarked}!cnaep@\iventry{cnæp}{}}
& attested OE form from another family & not part of the present lexeme
line \\
\end{longtable}
}

This remains the most review-sensitive item here, because the choice
between reconstructed \Recon{cnobba} `knob' and attested
{\emph{cnoppa}}\index[iv]{02oe@\ivlangheader{Old English}{West Saxon normalization unmarked}!cnoppa@\iventry{cnoppa}{}}
`knob' is still a comparator-policy question rather than a settled point
of OE attestation.

\section{\texorpdfstring{reek --- OE
\emph{*rēac}}{reek --- OE *rēac}}\label{reek-oe-rux113ac}

\index[iv]{02oe@\ivlangheader{Old English}{West Saxon normalization unmarked}!reac@\iventry{*rēac}{}}
\index[iv]{03pgmc@\ivlangheader{Proto-Germanic}{}!raukaz@\iventry{*ráukaz}{}}
\index[iv]{03pgmc@\ivlangheader{Proto-Germanic}{}!raukiz@\iventry{*ráukiz}{}}

Derivation: citation reconstruction \emph{*ráukiz}; form followed here
\emph{*ráukaz} \textgreater{} \emph{*rēac} (reconstructed Old English
comparator).

\subsection{Derivation trace}\label{derivation-trace-168}

Proto input: \emph{*ráukaz}

\begingroup
\setlength{\fboxsep}{6pt}

\noindent\fbox{%
\begin{minipage}{0.97\linewidth}
\small
\begin{tabular}{@{}>{\raggedright\arraybackslash}p{0.388\linewidth}@{\hspace{0.030\linewidth}}>{\raggedright\arraybackslash}p{0.582\linewidth}@{}}
\begin{minipage}[t]{\linewidth}
\centering\textbf{Earlier Germanic changes}\par
\vspace{0.35em}
\raggedright
\centering\textbf{Northwest and West Germanic}\par
\raggedright
\vspace{0.2em}
\begin{tabularx}{\linewidth}{@{}>{\raggedright\arraybackslash}X@{\hspace{0.55em}}>{\raggedright\arraybackslash}p{0.14\linewidth}@{\hspace{0.25em}}}
\mbox{EAF Final Z Deletion} & \emph{*ráuka} \\
\end{tabularx}
\vspace{0.6em}
\centering\textbf{Early Anglo-Frisian}\par
\raggedright
\vspace{0.2em}
\raggedright [no change]\par
\end{minipage}
&
\begin{minipage}[t]{\linewidth}
\centering\textbf{Old English changes}\par
\vspace{0.35em}
\raggedright
\begin{tabularx}{\linewidth}{@{}>{\raggedright\arraybackslash}X@{\hspace{0.55em}}>{\raggedright\arraybackslash}p{0.14\linewidth}@{\hspace{0.25em}}}
\mbox{OE Au Brightening} & \emph{*ráeuka} \\
\mbox{OE Diphthong Leveling} & \emph{*rēaka} \\
\mbox{PWGmc Final Bare A Loss} & \emph{*rēak} \\
\end{tabularx}
\end{minipage}
\\
\end{tabular}
\end{minipage}%
} \endgroup

Old English form: \emph{*rēac}

\subsection{Reconstruction and comparative
evidence}\label{reconstruction-and-comparative-evidence-148}

The wider noun family is represented by
\Recon{ráukiz}\index[iv]{03pgmc@\ivlangheader{Proto-Germanic}{}!raukiz@\iventry{*ráukiz}{}}
`reek' /
{\emph{*rauki-}}\index[iv]{03pgmc@\ivlangheader{Proto-Germanic}{}!rauki@\iventry{*rauki-}{}},
with Old English
{\emph{rēc}}\index[iv]{02oe@\ivlangheader{Old English}{West Saxon normalization unmarked}!rec@\iventry{rēc}{}}
`reek' as the attested noun reflex in the comparative dictionaries
(\citeproc{ref-Kroonen2013}{Kroonen 2013, 446};
\citeproc{ref-Orel2003}{Orel 2003, 338}). The derivational input
\Recon{ráukaz}\index[iv]{03pgmc@\ivlangheader{Proto-Germanic}{}!raukaz@\iventry{*ráukaz}{}}
`reek' is therefore not the lexeme-level headword, but the form used
here for the Old English derivation.

\subsection{Old English evidence}\label{old-english-evidence-149}

The attested noun is
{\emph{rēc}}\index[iv]{02oe@\ivlangheader{Old English}{West Saxon normalization unmarked}!rec@\iventry{rēc}{}}
`reek', not
\Recon{rēac}\index[iv]{02oe@\ivlangheader{Old English}{West Saxon normalization unmarked}!reac@\iventry{*rēac}{}}
`reek'. Clark Hall records
{\emph{rēc}}\index[iv]{02oe@\ivlangheader{Old English}{West Saxon normalization unmarked}!rec@\iventry{rēc}{}}
`reek' as the noun and also preserves related forms such as
{\emph{rēcels}}\index[iv]{02oe@\ivlangheader{Old English}{West Saxon normalization unmarked}!recels@\iventry{rēcels}{}}
`incense'; Kroonen likewise gives OE
{\emph{rēc}}\index[iv]{02oe@\ivlangheader{Old English}{West Saxon normalization unmarked}!rec@\iventry{rēc}{}}
`reek' under the noun family (\citeproc{ref-ClarkHall1960}{Clark Hall
1960, 255}; \citeproc{ref-Kroonen2013}{Kroonen 2013, 446}). Clark Hall
and Seebold also record verbal
{\emph{rēac}}\index[iv]{02oe@\ivlangheader{Old English}{West Saxon normalization unmarked}!reac@\iventry{rēac}{}}
`reek' as the preterite of
{\emph{rēocan}}\index[iv]{02oe@\ivlangheader{Old English}{West Saxon normalization unmarked}!reocan@\iventry{rēocan}{}}
`reek', but that verbal form is separate from the noun treated here
(\citeproc{ref-ClarkHall1960}{Clark Hall 1960, 254};
\citeproc{ref-Seebold1970}{Seebold 1970, 380}).

The Old English form here
\Recon{rēac}\index[iv]{02oe@\ivlangheader{Old English}{West Saxon normalization unmarked}!reac@\iventry{*rēac}{}}
`reek' is therefore a reconstructed West Saxon noun form, not a directly
attested manuscript headword.

\subsection{Development to Old
English}\label{development-to-old-english-147}

From
\Recon{ráukaz}\index[iv]{03pgmc@\ivlangheader{Proto-Germanic}{}!raukaz@\iventry{*ráukaz}{}}
`reek', the regular West Saxon development gives
\Recon{rēac}\index[iv]{02oe@\ivlangheader{Old English}{West Saxon normalization unmarked}!reac@\iventry{*rēac}{}}
`reek'. The attested noun
{\emph{rēc}}\index[iv]{02oe@\ivlangheader{Old English}{West Saxon normalization unmarked}!rec@\iventry{rēc}{}}
`reek' belongs to the same lexical family, but reflects a later smoothed
surface form rather than the regular noun target represented here.

\subsection{Form note}\label{form-note-27}

The distinction here is between an attested noun headword
{\emph{rēc}}\index[iv]{02oe@\ivlangheader{Old English}{West Saxon normalization unmarked}!rec@\iventry{rēc}{}}
`reek' and a reconstructed regular West Saxon target
\Recon{rēac}\index[iv]{02oe@\ivlangheader{Old English}{West Saxon normalization unmarked}!reac@\iventry{*rēac}{}}
`reek'. The latter is treated as the modelling target, while the former
remains philological background.

\section{\texorpdfstring{strew --- OE
\emph{*strīeġan}}{strew --- OE *strīeġan}}\label{strew-oe-strux12beux121an}

\index[iv]{02oe@\ivlangheader{Old English}{West Saxon normalization unmarked}!striegan@\iventry{*strīeġan}{}}
\index[iv]{03pgmc@\ivlangheader{Proto-Germanic}{}!strawjana@\iventry{*stráwjaną}{}}

Derivation: \emph{*stráwjaną} \textgreater{} \emph{*strīeġan}
(reconstructed Old English comparator).

\subsection{Derivation trace}\label{derivation-trace-169}

Proto input: \emph{*stráwjaną}

\begingroup
\setlength{\fboxsep}{6pt}

\noindent\fbox{%
\begin{minipage}{0.97\linewidth}
\footnotesize
\begin{tabular}{@{}>{\raggedright\arraybackslash}p{0.485\linewidth}@{\hspace{0.030\linewidth}}>{\raggedright\arraybackslash}p{0.485\linewidth}@{}}
\begin{minipage}[t]{\linewidth}
\centering\textbf{Earlier Germanic changes}\par
\vspace{0.35em}
\raggedright
\centering\textbf{Northwest and West Germanic}\par
\raggedright
\vspace{0.2em}
\begin{tabularx}{\linewidth}{@{}>{\raggedright\arraybackslash}X@{\hspace{0.55em}}>{\raggedright\arraybackslash}p{0.24\linewidth}@{\hspace{0.25em}}}
\mbox{PWGmc J Gemination} & \emph{*stráwwjaną} \\
\end{tabularx}
\vspace{0.6em}
\centering\textbf{Early Anglo-Frisian}\par
\raggedright
\vspace{0.2em}
\raggedright [no change]\par
\end{minipage}
&
\begin{minipage}[t]{\linewidth}
\centering\textbf{Old English changes}\par
\vspace{0.35em}
\raggedright
\begin{tabularx}{\linewidth}{@{}>{\raggedright\arraybackslash}X@{\hspace{0.55em}}>{\raggedright\arraybackslash}p{0.26\linewidth}@{\hspace{0.25em}}}
\mbox{OE Awwj Resolution} & \emph{*stráujaną} \\
\mbox{OE Au Brightening} & \emph{*stráeujaną} \\
\mbox{OE Diphthong Leveling} & \emph{*strēajaną} \\
OE Heavy Syllable Nasal Apocope & \emph{*strēajan} \\
\mbox{OE Secondary Nasalization} & \emph{*strēająn} \\
\mbox{OE I Umlaut} & \emph{*strīejąn} \\
\mbox{OE Weak Tail Reduction} & \emph{*strīejan} \\
OE J Strengthening After Front Diphthong & \emph{*strīeʒan} \\
\end{tabularx}
\end{minipage}
\\
\end{tabular}
\end{minipage}%
} \endgroup

Old English form: \emph{*strīeġan}

\subsection{Reconstruction and comparative
evidence}\label{reconstruction-and-comparative-evidence-149}

Kroonen cites the inherited weak verb as
{\emph{*straujan-}}\index[iv]{03pgmc@\ivlangheader{Proto-Germanic}{}!straujan@\iventry{*straujan-}{}}
and gives Old English
{\emph{streowian}}\index[iv]{02oe@\ivlangheader{Old English}{West Saxon normalization unmarked}!streowian@\iventry{streowian}{}}
`strew' as its dictionary continuation
(\citeproc{ref-Kroonen2013}{Kroonen 2013, 483}). Ringe and Taylor
separate the inherited class-I branch from the remodelled class-II
branch: Anglian
{\emph{strēgan}}\index[iv]{02oe@\ivlangheader{Old English}{West Saxon normalization unmarked}!stregan@\iventry{strēgan}{}}
`strew' continues the inherited verb, while West Saxon
{\emph{streowian}}\index[iv]{02oe@\ivlangheader{Old English}{West Saxon normalization unmarked}!streowian@\iventry{streowian}{}}
`strew' is remodelled (\citeproc{ref-RingeTaylor2014}{Ringe and Taylor
2014, § 6.1} n.~27).

Luick groups \Recon{strauwjan} `strew' with the same set as
\emph{*hauwja-} and \Recon{kauwjan} `chew', yielding Anglian
{\emph{strēzan}}\index[iv]{02oe@\ivlangheader{Old English}{West Saxon normalization unmarked}!strezan~04@\iventry{strēzan}{Angl.}}
`strew' beside West Saxon forms of the
{\emph{hīez}}\index[iv]{02oe@\ivlangheader{Old English}{West Saxon normalization unmarked}!hiez@\iventry{hīez}{}}
`hew',
{\emph{ciezan}}\index[iv]{02oe@\ivlangheader{Old English}{West Saxon normalization unmarked}!ciezan@\iventry{ciezan}{}}
`choose' type (\citeproc{ref-Luick1914}{Luick 1914-\/-1940, § 98}).
Fulk likewise allows an early West Saxon \Recon{striegan} `strew'
directly from Proto-Germanic \Recon{straujana} `strew'
(\citeproc{ref-Fulk2018}{Fulk 2018, § 4.10} n.~1).

\subsection{Old English evidence}\label{old-english-evidence-150}

The attested inherited Old English form is
{\emph{strēgan}}\index[iv]{02oe@\ivlangheader{Old English}{West Saxon normalization unmarked}!stregan~04@\iventry{strēgan}{Angl.}}
`strew' in Anglian. The attested West Saxon citation forms are
{\emph{strewian}}\index[iv]{02oe@\ivlangheader{Old English}{West Saxon normalization unmarked}!strewian@\iventry{strewian}{}}
`strew',
{\emph{streowian}}\index[iv]{02oe@\ivlangheader{Old English}{West Saxon normalization unmarked}!streowian@\iventry{streowian}{}}
`strew', and
{\emph{strēawian}}\index[iv]{02oe@\ivlangheader{Old English}{West Saxon normalization unmarked}!streawian@\iventry{strēawian}{}}
`strew', which belong to the remodelled class-II branch
(\citeproc{ref-RingeTaylor2014}{Ringe and Taylor 2014, § 6.1} n.~27;
\citeproc{ref-Campbell1959}{Campbell 1959, § 753.7}).

The target
\Recon{strīeġan}\index[iv]{02oe@\ivlangheader{Old English}{West Saxon normalization unmarked}!striegan@\iventry{*strīeġan}{}}
`strew' is therefore a reconstructed Old English form, not an attested
manuscript lemma. It represents the inferred West Saxon reflex of the
inherited class-I branch.

\subsection{Development to Old
English}\label{development-to-old-english-148}

From \Recon{stráwjaną} `strew', the inherited West Saxon line passes
through \Recon{straujaną} `strew', then a fronted \emph{*strēajan-}
stage, then \Recon{strīejan} `strew' by i-umlaut, with
retained/strengthened glide after the front diphthong to yield
reconstructed \Recon{strīeġan} `strew'.

This differs from Anglian
{\emph{strēgan}}\index[iv]{02oe@\ivlangheader{Old English}{West Saxon normalization unmarked}!stregan~04@\iventry{strēgan}{Angl.}}
`strew', where smoothing removes the diphthongal sequence, and from West
Saxon
{\emph{strewian}}\index[iv]{02oe@\ivlangheader{Old English}{West Saxon normalization unmarked}!strewian@\iventry{strewian}{}}
`strew' /
{\emph{streowian}}\index[iv]{02oe@\ivlangheader{Old English}{West Saxon normalization unmarked}!streowian@\iventry{streowian}{}}
`strew' /
{\emph{strēawian}}\index[iv]{02oe@\ivlangheader{Old English}{West Saxon normalization unmarked}!streawian@\iventry{strēawian}{}}
`strew', where the verb has already been remodelled into class II
(\citeproc{ref-Fulk2018}{Fulk 2018, § 4.10} n.~1;
\citeproc{ref-Campbell1959}{Campbell 1959, § 753.7}).

\subsection{Reconstruction status}\label{reconstruction-status-1}

{\def\LTcaptype{none} % do not increment counter
\begin{longtable}[]{@{}
  >{\raggedright\arraybackslash}p{(\linewidth - 4\tabcolsep) * \real{0.3333}}
  >{\raggedright\arraybackslash}p{(\linewidth - 4\tabcolsep) * \real{0.3333}}
  >{\raggedright\arraybackslash}p{(\linewidth - 4\tabcolsep) * \real{0.3333}}@{}}
\toprule\noalign{}
\begin{minipage}[b]{\linewidth}\raggedright
Form or branch
\end{minipage} & \begin{minipage}[b]{\linewidth}\raggedright
Status
\end{minipage} & \begin{minipage}[b]{\linewidth}\raggedright
Relevance to this entry
\end{minipage} \\
\midrule\noalign{}
\endhead
\bottomrule\noalign{}
\endlastfoot
{\emph{strēgan}}\index[iv]{02oe@\ivlangheader{Old English}{West Saxon normalization unmarked}!stregan~04@\iventry{strēgan}{Angl.}}
`strew' & attested Anglian inherited class-I form & proves that the
inherited verb survived into Old English \\
\Recon{strīeġan}\index[iv]{02oe@\ivlangheader{Old English}{West Saxon normalization unmarked}!striegan@\iventry{*strīeġan}{}}
`strew' & reconstructed West Saxon inherited class-I form;
trace-supported & Old English form here \\
{\emph{strewian}}\index[iv]{02oe@\ivlangheader{Old English}{West Saxon normalization unmarked}!strewian@\iventry{strewian}{}}
/
{\emph{streowian}}\index[iv]{02oe@\ivlangheader{Old English}{West Saxon normalization unmarked}!streowian@\iventry{streowian}{}}
`strew' /
{\emph{strēawian}}\index[iv]{02oe@\ivlangheader{Old English}{West Saxon normalization unmarked}!streawian@\iventry{strēawian}{}}
& attested remodelled West Saxon class-II forms & genuine OE evidence,
but not the inherited branch modeled here \\
\end{longtable}
}

\clearpage

\chapter{Known but unmodelled
developments}\label{known-but-unmodelled-developments}

The historical development is sufficiently supported by comparative and
philological evidence, but the current transducer cascade does not yet
implement the phonological or morphological process required to derive
the target. Naming the gap explains the status without pretending that
no historical account exists.

\section{\texorpdfstring{fire --- OE
\emph{fȳre}}{fire --- OE fȳre}}\label{fire-oe-fux233re}

\index[iv]{02oe@\ivlangheader{Old English}{West Saxon normalization unmarked}!fyre@\iventry{fȳre}{}}
\index[iv]{03pgmc@\ivlangheader{Proto-Germanic}{}!furi@\iventry{*fūri}{}}

Derivation: \emph{*fūri} yields regular \Recon{fȳr}; the Old English
form here is \emph{fȳre} (known but unmodelled development).

\subsection{Derivation trace}\label{derivation-trace-170}

Proto input: \emph{*fūri}

\begingroup
\setlength{\fboxsep}{6pt}

\noindent\fbox{%
\begin{minipage}{0.97\linewidth}
\small
\begin{tabular}{@{}>{\raggedright\arraybackslash}p{0.388\linewidth}@{\hspace{0.030\linewidth}}>{\raggedright\arraybackslash}p{0.582\linewidth}@{}}
\begin{minipage}[t]{\linewidth}
\centering\textbf{Earlier Germanic changes}\par
\vspace{0.35em}
\raggedright
\centering\textbf{Northwest and West Germanic}\par
\raggedright
\vspace{0.2em}
\raggedright [no change]\par
\vspace{0.6em}
\centering\textbf{Early Anglo-Frisian}\par
\raggedright
\vspace{0.2em}
\raggedright [no change]\par
\end{minipage}
&
\begin{minipage}[t]{\linewidth}
\centering\textbf{Old English changes}\par
\vspace{0.35em}
\raggedright
\begin{tabularx}{\linewidth}{@{}>{\raggedright\arraybackslash}X@{\hspace{0.55em}}>{\raggedright\arraybackslash}p{0.14\linewidth}@{\hspace{0.25em}}}
\mbox{OE I Umlaut} & \emph{*fȳri} \\
\mbox{OE High Vowel Apocope} & \Recon{fȳr} \\
\end{tabularx}
\end{minipage}
\\
\end{tabular}
\end{minipage}%
} \endgroup

Regular outcome: \Recon{fȳr}

Old English form: \emph{fȳre}

\subsection{Reconstruction and comparative
evidence}\label{reconstruction-and-comparative-evidence-150}

Kroonen places the lexeme in a heteroclitic family \Recon{fōr} `fire'
\textasciitilde{} \emph{*fun-} and explains the front-mutated West
Germanic forms from an oblique form of the \emph{*fu(w)eri} type
(\citeproc{ref-Kroonen2013}{Kroonen 2013, 151}). The derivational input
\Recon{fūri}\index[iv]{03pgmc@\ivlangheader{Proto-Germanic}{}!furi@\iventry{*fūri}{}}
`fire' is CAPR's oblique-input proxy supplying the \emph{i} needed for
i-umlaut. The account does not independently certify that exact complete
word, its long \emph{ū}, or its cell.

The Old English target combines a regular inherited form
{\emph{fȳr}}\index[iv]{02oe@\ivlangheader{Old English}{West Saxon normalization unmarked}!fyr@\iventry{fȳr}{}}
`fire' with an attested oblique surface form
{\emph{fȳre}}\index[iv]{02oe@\ivlangheader{Old English}{West Saxon normalization unmarked}!fyre@\iventry{fȳre}{}}
`fire', interpreted by CAPR as analogically restored.

\subsection{Competing ancestry and the legitimacy of the
input}\label{competing-ancestry-and-the-legitimacy-of-the-input}

Kroonen distinguishes the dative \emph{*fu(w)eni}, its
\emph{r}-remodelled \emph{*fu(w)eri}, and metathesized \emph{*feuri}.
His preferred pre-Germanic nominative is \emph{*péh₂-ur}; the
\emph{*ph₂-uōr} proposal reported from Simms is rejected, not endorsed
(\citeproc{ref-Kroonen2013}{Kroonen 2013, 151}). Ringe instead supplies
a reaccented zero-grade collective, printed \emph{*ph2uōŕ}, and
separately considers post-Proto-Germanic oblique dissimilation and
\emph{n} replacement by \emph{r} (\citeproc{ref-Ringe2017}{Ringe 2017,
147, 309}). These are consequential formation and accent alternatives,
not interchangeable quoted dative inputs.

Ringe and Taylor reconstruct apparently disyllabic West Germanic
\emph{*fuir} with \emph{r} levelled into the oblique stem, but doubt
that the endingless dative evidence establishes inherited \emph{*fuiri}.
A front-rounded intermediate in the history of \emph{fȳr} is plausible;
the intervening stages remain uncertain
(\citeproc{ref-RingeTaylor2014}{Ringe and Taylor 2014, 119, 225}).
Neither passage endorses CAPR's exact long-\emph{ū} proxy. Equal
citation and input fields select the Proto-Germanic stage by project
convention, not by a new scholarly quotation. Quantity, cell and
historical input legitimacy require separate scientific assessment; no
input migration is made here.

\subsection{Old English evidence}\label{old-english-evidence-151}

Toller's Supplement records
{\emph{fyr}}\index[iv]{02oe@\ivlangheader{Old English}{West Saxon normalization unmarked}!fyr@\iventry{fyr}{}}
`fire' as a noun in its \emph{fyr. Add} entry and separately preserves
oblique
{\emph{fyre}}\index[iv]{02oe@\ivlangheader{Old English}{West Saxon normalization unmarked}!fyre@\iventry{fyre}{}}
`fire', including \emph{in fyre} (\citeproc{ref-Toller1921}{Toller 1921,
276}). These attestations establish the forms, not the date or mechanism
of restoration. The first is the regular output of the phonological
development from the selected input; the second is interpreted by CAPR
as restoration of final \emph{-e} within the paradigm, rather than
demonstrated as such by the dictionary.

The entry therefore distinguishes a regular path from an adopted input
proxy, the comparative legitimacy of that proxy, and an attested surface
form whose additional inflection CAPR treats as morphological
remodeling.

\subsection{Development to Old
English}\label{development-to-old-english-149}

From
\Recon{fūri}\index[iv]{03pgmc@\ivlangheader{Proto-Germanic}{}!furi@\iventry{*fūri}{}}
`fire', i-umlaut changes \emph{ū} to \emph{ȳ}
(\citeproc{ref-HoggPhonology1992}{Hogg 1992b, § 3.3.3.1}).
Subsequent loss of the final high vowel after a heavy syllable yields
{\emph{fȳr}} `fire' (\citeproc{ref-Campbell1959}{Campbell 1959, § 345}). The selected phonological path is complete at that point.

CAPR interprets the final \emph{-e} of
{\emph{fȳre}}\index[iv]{02oe@\ivlangheader{Old English}{West Saxon normalization unmarked}!fyre@\iventry{fȳre}{}}
`fire' as analogical restoration in the Old English paradigm rather than
retention of the derivational input's ending. The form therefore remains
a known but unmodelled remodelling: the deterministic phonology is
regular, but the attested surface form includes later morphological
rebuilding.

\subsection{Paradigm comparison}\label{paradigm-comparison-26}

The comparison below sets the relevant forms side by side. It
distinguishes the lexeme-level etymological background from the
inherited cell that actually produces the front-mutated form and from
the later analogical surface result.

{\def\LTcaptype{none} % do not increment counter
\begin{longtable}[]{@{}
  >{\raggedright\arraybackslash}p{(\linewidth - 8\tabcolsep) * \real{0.2000}}
  >{\raggedright\arraybackslash}p{(\linewidth - 8\tabcolsep) * \real{0.2000}}
  >{\raggedright\arraybackslash}p{(\linewidth - 8\tabcolsep) * \real{0.2000}}
  >{\raggedright\arraybackslash}p{(\linewidth - 8\tabcolsep) * \real{0.2000}}
  >{\raggedright\arraybackslash}p{(\linewidth - 8\tabcolsep) * \real{0.2000}}@{}}
\toprule\noalign{}
\begin{minipage}[b]{\linewidth}\raggedright
PGmc cell / interpretation
\end{minipage} & \begin{minipage}[b]{\linewidth}\raggedright
Candidate input
\end{minipage} & \begin{minipage}[b]{\linewidth}\raggedright
OE output or comparison
\end{minipage} & \begin{minipage}[b]{\linewidth}\raggedright
OE comparison form
\end{minipage} & \begin{minipage}[b]{\linewidth}\raggedright
Result
\end{minipage} \\
\midrule\noalign{}
\endhead
\bottomrule\noalign{}
\endlastfoot
lexeme-level heteroclitic headword & \Recon{fōr} `fire'
\textasciitilde{} \emph{*fun-} & comparative background only & fire
family & explains the wider lexeme, but not the selected oblique
input \\
selected oblique-input proxy &
{\emph{*fūri}}\index[iv]{03pgmc@\ivlangheader{Proto-Germanic}{}!furi@\iventry{*fūri}{}}
& regular output: {\emph{fȳr}} `fire' &
{\emph{fȳr}}\index[iv]{02oe@\ivlangheader{Old English}{West Saxon normalization unmarked}!fyr@\iventry{fȳr}{}}
`fire' & regular output; exact ancestral cell remains qualified \\
later analogical surface form & --- & attested
{\emph{fȳre}}\index[iv]{02oe@\ivlangheader{Old English}{West Saxon normalization unmarked}!fyre@\iventry{fȳre}{}}
`fire' with restored \emph{-e} &
{\emph{fȳre}}\index[iv]{02oe@\ivlangheader{Old English}{West Saxon normalization unmarked}!fyre@\iventry{fȳre}{}}
`fire' & genuine OE target, but not the direct phonological output \\
\end{longtable}
}

\section{\texorpdfstring{tap --- OE
\emph{tæppa}}{tap --- OE tæppa}}\label{tap-oe-tuxe6ppa}

\index[iv]{02oe@\ivlangheader{Old English}{West Saxon normalization unmarked}!taeppa@\iventry{tæppa}{}}
\index[iv]{03pgmc@\ivlangheader{Proto-Germanic}{}!tappo@\iventry{*táppô}{}}

Derivation: \emph{*táppô} yields regular \Recon{tappa}; the Old English
form here is \emph{tæppa} (known but unmodelled development).

\subsection{Derivation trace}\label{derivation-trace-171}

Proto input: \emph{*táppô}

\begingroup
\setlength{\fboxsep}{6pt}

\noindent\fbox{%
\begin{minipage}{0.97\linewidth}
\small
\begin{tabular}{@{}>{\raggedright\arraybackslash}p{0.359\linewidth}@{\hspace{0.030\linewidth}}>{\raggedright\arraybackslash}p{0.611\linewidth}@{}}
\begin{minipage}[t]{\linewidth}
\centering\textbf{Earlier Germanic changes}\par
\vspace{0.35em}
\raggedright
\centering\textbf{Northwest and West Germanic}\par
\raggedright
\vspace{0.2em}
\raggedright [no change]\par
\vspace{0.6em}
\centering\textbf{Early Anglo-Frisian}\par
\raggedright
\vspace{0.2em}
\raggedright [no change]\par
\end{minipage}
&
\begin{minipage}[t]{\linewidth}
\centering\textbf{Old English changes}\par
\vspace{0.35em}
\raggedright
\begin{tabularx}{\linewidth}{@{}>{\raggedright\arraybackslash}X@{\hspace{0.55em}}>{\raggedright\arraybackslash}p{0.14\linewidth}@{\hspace{0.25em}}}
\mbox{EAF Brightening} & \emph{*tæppô} \\
\mbox{OE A Restoration} & \emph{*tappô} \\
OE Unstressed Long Vowel Shortening & \Recon{tappa} \\
\end{tabularx}
\end{minipage}
\\
\end{tabular}
\end{minipage}%
} \endgroup

Regular outcome: \Recon{tappa}

Old English form: \emph{tæppa}

\subsection{Reconstruction and comparative
evidence}\label{reconstruction-and-comparative-evidence-151}

Orel gives the noun under \Recon{tappòn} `tap' and already connects it
with Old English
{\emph{tæppa}}\index[iv]{02oe@\ivlangheader{Old English}{West Saxon normalization unmarked}!taeppa@\iventry{tæppa}{}}
`tap' (\citeproc{ref-Orel2003}{Orel 2003, 402}). The derivational input
\Recon{táppô}\index[iv]{03pgmc@\ivlangheader{Proto-Germanic}{}!tappo@\iventry{*táppô}{}}
`tap' is therefore the inherited noun itself; the entry does not depend
on a different lexeme-level proto or a different inherited noun cell.

\subsection{Old English evidence}\label{old-english-evidence-152}

The Old English noun family is well attested. Orel gives
{\emph{tæppa}}\index[iv]{02oe@\ivlangheader{Old English}{West Saxon normalization unmarked}!taeppa@\iventry{tæppa}{}}
`tap', and Clark Hall records
{\emph{tæppa}}\index[iv]{02oe@\ivlangheader{Old English}{West Saxon normalization unmarked}!taeppa@\iventry{tæppa}{}}
`tap' together with derivatives \emph{tæppere} `tapster' and
\emph{tæppestre} `tapster' (\citeproc{ref-Orel2003}{Orel 2003, 402};
\citeproc{ref-ClarkHall1960}{Clark Hall 1960, 305}). The target is
therefore a real Old English noun form, not a reconstructed convenience
spelling.

\subsection{Development to Old
English}\label{development-to-old-english-150}

From
\Recon{táppô}\index[iv]{03pgmc@\ivlangheader{Proto-Germanic}{}!tappo@\iventry{*táppô}{}}
`tap', the regular inherited noun path gives {\emph{tappa}} `tap'. The
attested target
{\emph{tæppa}}\index[iv]{02oe@\ivlangheader{Old English}{West Saxon normalization unmarked}!taeppa@\iventry{tæppa}{}}
`tap' therefore stands outside that regular phonological development.

The mismatch is historically intelligible, but it is not solved here by
a new inherited input. A related j-verb pathway would give
{\emph{teppan}}\index[iv]{02oe@\ivlangheader{Old English}{West Saxon normalization unmarked}!teppan@\iventry{teppan}{}}
`tap', not the noun target
{\emph{tæppa}}\index[iv]{02oe@\ivlangheader{Old English}{West Saxon normalization unmarked}!taeppa@\iventry{tæppa}{}}
`tap'. The entry accordingly remains a known but unmodelled case.

\subsection{Form comparison}\label{form-comparison-5}

{\def\LTcaptype{none} % do not increment counter
\begin{longtable}[]{@{}
  >{\raggedright\arraybackslash}p{(\linewidth - 6\tabcolsep) * \real{0.2500}}
  >{\raggedright\arraybackslash}p{(\linewidth - 6\tabcolsep) * \real{0.2500}}
  >{\raggedright\arraybackslash}p{(\linewidth - 6\tabcolsep) * \real{0.2500}}
  >{\raggedright\arraybackslash}p{(\linewidth - 6\tabcolsep) * \real{0.2500}}@{}}
\toprule\noalign{}
\begin{minipage}[b]{\linewidth}\raggedright
Form type
\end{minipage} & \begin{minipage}[b]{\linewidth}\raggedright
Input or form
\end{minipage} & \begin{minipage}[b]{\linewidth}\raggedright
OE output or comparison
\end{minipage} & \begin{minipage}[b]{\linewidth}\raggedright
Result
\end{minipage} \\
\midrule\noalign{}
\endhead
\bottomrule\noalign{}
\endlastfoot
regular inherited noun path &
{\emph{*táppô}}\index[iv]{03pgmc@\ivlangheader{Proto-Germanic}{}!tappo@\iventry{*táppô}{}}
& regular output: {\emph{tappa}} `tap' & regular output, but not the
target \\
attested OE target & --- &
{\emph{tæppa}}\index[iv]{02oe@\ivlangheader{Old English}{West Saxon normalization unmarked}!taeppa@\iventry{tæppa}{}}
`tap' & genuine target form, but analogically remodelled in the present
classification \\
related j-verb background &
{\emph{*táppjaną}}\index[iv]{03pgmc@\ivlangheader{Proto-Germanic}{}!tappjana@\iventry{*táppjaną}{}}
&
{\emph{teppan}}\index[iv]{02oe@\ivlangheader{Old English}{West Saxon normalization unmarked}!teppan@\iventry{teppan}{}}
`tap' & related formation, but not the noun target \\
\end{longtable}
}

\clearpage

\chapter{Unexplained or deliberately unmodelled
exceptions}\label{unexplained-or-deliberately-unmodelled-exceptions}

No sufficiently supported account yet reconciles the regular output with
the Old English form. An ad hoc sound law would conceal rather than
explain the mismatch.

\section{\texorpdfstring{buck --- OE
\emph{bucc}}{buck --- OE bucc}}\label{buck-oe-bucc}

\index[iv]{02oe@\ivlangheader{Old English}{West Saxon normalization unmarked}!bucc@\iventry{bucc}{}}
\index[iv]{03pgmc@\ivlangheader{Proto-Germanic}{}!bukkaz@\iventry{*búkkaz}{}}

Derivation: \emph{*búkkaz} yields regular \Recon{bocc}; the Old English
form here is \emph{bucc} (unexplained exception).

\subsection{Derivation trace}\label{derivation-trace-172}

Proto input: \emph{*búkkaz}

\begingroup
\setlength{\fboxsep}{6pt}

\noindent\fbox{%
\begin{minipage}{0.97\linewidth}
\small
\begin{tabular}{@{}>{\raggedright\arraybackslash}p{0.485\linewidth}@{\hspace{0.030\linewidth}}>{\raggedright\arraybackslash}p{0.485\linewidth}@{}}
\begin{minipage}[t]{\linewidth}
\centering\textbf{Earlier Germanic changes}\par
\vspace{0.35em}
\raggedright
\centering\textbf{Northwest and West Germanic}\par
\raggedright
\vspace{0.2em}
\begin{tabularx}{\linewidth}{@{}>{\raggedright\arraybackslash}X@{\hspace{0.55em}}>{\raggedright\arraybackslash}p{0.14\linewidth}@{\hspace{0.25em}}}
\mbox{PNWGmc U Lowering} & \emph{*bókkaz} \\
\mbox{EAF Final Z Deletion} & \emph{*bókka} \\
\end{tabularx}
\vspace{0.6em}
\centering\textbf{Early Anglo-Frisian}\par
\raggedright
\vspace{0.2em}
\raggedright [no change]\par
\end{minipage}
&
\begin{minipage}[t]{\linewidth}
\centering\textbf{Old English changes}\par
\vspace{0.35em}
\raggedright
\begin{tabularx}{\linewidth}{@{}>{\raggedright\arraybackslash}X@{\hspace{0.55em}}>{\raggedright\arraybackslash}p{0.14\linewidth}@{\hspace{0.25em}}}
\mbox{PWGmc Final Bare A Loss} & \emph{*bókk} \\
\end{tabularx}
\end{minipage}
\\
\end{tabular}
\end{minipage}%
} \endgroup

Regular outcome: \Recon{bocc}

Old English form: \emph{bucc}

\subsection{Reconstruction and comparative
evidence}\label{reconstruction-and-comparative-evidence-152}

Kroonen and Orel both reconstruct the word with a geminate stop,
\Recon{bukkaz}\index[iv]{03pgmc@\ivlangheader{Proto-Germanic}{}!bukkaz@\iventry{*bukkaz}{}}
`buck' (\citeproc{ref-Kroonen2013}{Kroonen 2013, 121};
\citeproc{ref-Orel2003}{Orel 2003, 61}). Orel also preserves parallel
n-stem material behind Old English
{\emph{bucca}}\index[iv]{02oe@\ivlangheader{Old English}{West Saxon normalization unmarked}!bucca@\iventry{bucca}{}}
`buck' (\citeproc{ref-Orel2003}{Orel 2003, 62}). The derivational input
therefore remains identical with the lexeme label: no alternative
inherited cell accounts for the form.

\subsection{Old English evidence}\label{old-english-evidence-153}

Old English preserves a mixed lexical picture. Campbell cites
{\emph{bucca}}\index[iv]{02oe@\ivlangheader{Old English}{West Saxon normalization unmarked}!bucca@\iventry{bucca}{}}
`buck' in the exception set for this phonological environment
(\citeproc{ref-Campbell1959}{Campbell 1959, § 115}). Clark Hall and
Bosworth-Toller show that Old English has both
{\emph{bucca}}\index[iv]{02oe@\ivlangheader{Old English}{West Saxon normalization unmarked}!bucca@\iventry{bucca}{}}
`buck' and
{\emph{bucc}}\index[iv]{02oe@\ivlangheader{Old English}{West Saxon normalization unmarked}!bucc@\iventry{bucc}{}}
`buck' (\citeproc{ref-ClarkHall1960}{Clark Hall 1960, 53};
\citeproc{ref-BosworthToller1898}{Bosworth and Toller 1898, 122}). The
a-stem citation form
{\emph{bucc}}\index[iv]{02oe@\ivlangheader{Old English}{West Saxon normalization unmarked}!bucc@\iventry{bucc}{}}
`buck' is the target treated here;
{\emph{bucca}}\index[iv]{02oe@\ivlangheader{Old English}{West Saxon normalization unmarked}!bucca@\iventry{bucca}{}}
`buck' supplies genuine philological background from the same lexical
family.

\subsection{Development to Old
English}\label{development-to-old-english-151}

From
\Recon{búkkaz}\index[iv]{03pgmc@\ivlangheader{Proto-Germanic}{}!bukkaz@\iventry{*búkkaz}{}}
`buck', the regular inherited path gives {\emph{bocc}} `buck'. That is
the form expected under the ordinary lowering pattern in this
environment.
{\emph{bucc}}\index[iv]{02oe@\ivlangheader{Old English}{West Saxon normalization unmarked}!bucc@\iventry{bucc}{}}
`buck' therefore remains outside the deterministic phonology.

No accepted inherited cell repairs the mismatch. A high-vowel
alternative would introduce i-umlaut and produce a
{\emph{byċċ}}\index[iv]{02oe@\ivlangheader{Old English}{West Saxon normalization unmarked}!bycc@\iventry{byċċ}{}}
`buck'-type form rather than the target.
{\emph{bucc}}\index[iv]{02oe@\ivlangheader{Old English}{West Saxon normalization unmarked}!bucc@\iventry{bucc}{}}
`buck' is therefore best treated as a documented exception, not as a
regular paradigm-cell survival.

\subsection{Form comparison}\label{form-comparison-6}

{\def\LTcaptype{none} % do not increment counter
\begin{longtable}[]{@{}
  >{\raggedright\arraybackslash}p{(\linewidth - 6\tabcolsep) * \real{0.2500}}
  >{\raggedright\arraybackslash}p{(\linewidth - 6\tabcolsep) * \real{0.2500}}
  >{\raggedright\arraybackslash}p{(\linewidth - 6\tabcolsep) * \real{0.2500}}
  >{\raggedright\arraybackslash}p{(\linewidth - 6\tabcolsep) * \real{0.2500}}@{}}
\toprule\noalign{}
\begin{minipage}[b]{\linewidth}\raggedright
Form type
\end{minipage} & \begin{minipage}[b]{\linewidth}\raggedright
Input or form
\end{minipage} & \begin{minipage}[b]{\linewidth}\raggedright
OE output or comparison
\end{minipage} & \begin{minipage}[b]{\linewidth}\raggedright
Result
\end{minipage} \\
\midrule\noalign{}
\endhead
\bottomrule\noalign{}
\endlastfoot
regular inherited noun path &
{\emph{*búkkaz}}\index[iv]{03pgmc@\ivlangheader{Proto-Germanic}{}!bukkaz@\iventry{*búkkaz}{}}
& regular output: {\emph{bocc}} `buck' & regular output, but not the
target \\
attested OE target & --- &
{\emph{bucc}}\index[iv]{02oe@\ivlangheader{Old English}{West Saxon normalization unmarked}!bucc@\iventry{bucc}{}}
`buck' & genuine target form, but unexplained in the present
classification \\
parallel OE lexical background & --- &
{\emph{bucca}}\index[iv]{02oe@\ivlangheader{Old English}{West Saxon normalization unmarked}!bucca@\iventry{bucca}{}}
`buck' & related n-stem form, not the present target \\
\end{longtable}
}

\section{\texorpdfstring{fowl --- OE
\emph{fugol}}{fowl --- OE fugol}}\label{fowl-oe-fugol}

\index[iv]{02oe@\ivlangheader{Old English}{West Saxon normalization unmarked}!fugel@\iventry{fugel}{}}
\index[iv]{02oe@\ivlangheader{Old English}{West Saxon normalization unmarked}!fugol@\iventry{fugol}{}}
\index[iv]{03pgmc@\ivlangheader{Proto-Germanic}{}!fuglaz@\iventry{*fúglaz}{}}
\index[iv]{09on@\ivlangheader{Old Norse}{}!fugl@\iventry{fugl}{}}
\index[iv]{10ohg@\ivlangheader{Old High German}{}!fogal@\iventry{fogal}{}}

Derivation: \emph{*fúglaz} yields regular \Recon{fogol}; the Old English
form here is \emph{fugol} (unexplained exception).

\subsection{Derivation trace}\label{derivation-trace-173}

Proto input: \emph{*fúglaz}

\begingroup
\setlength{\fboxsep}{6pt}

\noindent\fbox{%
\begin{minipage}{0.97\linewidth}
\small
\begin{tabular}{@{}>{\raggedright\arraybackslash}p{0.485\linewidth}@{\hspace{0.030\linewidth}}>{\raggedright\arraybackslash}p{0.485\linewidth}@{}}
\begin{minipage}[t]{\linewidth}
\centering\textbf{Earlier Germanic changes}\par
\vspace{0.35em}
\raggedright
\centering\textbf{Northwest and West Germanic}\par
\raggedright
\vspace{0.2em}
\begin{tabularx}{\linewidth}{@{}>{\raggedright\arraybackslash}X@{\hspace{0.55em}}>{\raggedright\arraybackslash}p{0.14\linewidth}@{\hspace{0.25em}}}
\mbox{PNWGmc U Lowering} & \emph{*fóglaz} \\
\mbox{EAF Final Z Deletion} & \emph{*fógla} \\
\end{tabularx}
\vspace{0.6em}
\centering\textbf{Early Anglo-Frisian}\par
\raggedright
\vspace{0.2em}
\raggedright [no change]\par
\end{minipage}
&
\begin{minipage}[t]{\linewidth}
\centering\textbf{Old English changes}\par
\vspace{0.35em}
\raggedright
\begin{tabularx}{\linewidth}{@{}>{\raggedright\arraybackslash}X@{\hspace{0.55em}}>{\raggedright\arraybackslash}p{0.14\linewidth}@{\hspace{0.25em}}}
\mbox{PWGmc Final Bare A Loss} & \emph{*fógl} \\
\mbox{OE Epenthetic Vowel} & \emph{*fógol} \\
\end{tabularx}
\end{minipage}
\\
\end{tabular}
\end{minipage}%
} \endgroup

Regular outcome: \Recon{fogol}

Old English form: \emph{fugol}

\subsection{Reconstruction and comparative
evidence}\label{reconstruction-and-comparative-evidence-153}

The noun is the ordinary Germanic a-stem \Recon{fúglaz} `fowl',
continued by forms such as Old Norse \emph{fugl} and Old High German
{\emph{fogal}}\index[iv]{10ohg@\ivlangheader{Old High German}{}!fogal@\iventry{fogal}{}}
`bird, fowl' (\citeproc{ref-Kroonen2013}{Kroonen 2013, 197};
\citeproc{ref-Orel2003}{Orel 2003, 155}). There is no stem-class or
paradigm-cell dispute behind this entry. The comparative headword and
the derivational input are the same.

The difficulty lies instead in the stressed root vowel. Under the
regular West Germanic and Old English development, that \emph{u} should
lower before the following non-high vowel, yielding an \emph{o}-vocalism
(\citeproc{ref-RingeTaylor2014}{Ringe and Taylor 2014, 42--43};
\citeproc{ref-Campbell1959}{Campbell 1959, 43}).

\subsection{Old English evidence}\label{old-english-evidence-154}

Old English dictionaries record the noun as \emph{fugol} `fowl', with
variant spelling \emph{fugel} `fowl'
(\citeproc{ref-BosworthToller1898}{Bosworth and Toller 1898, 282};
\citeproc{ref-ClarkHall1960}{Clark Hall 1960, 138}). The target is
therefore an attested ordinary Old English noun, not a reconstructed or
selectively chosen paradigm form.

The attested word already contains the crucial problem. Its medial
\emph{-o-} is the ordinary parasite vowel of Old English cluster
phonology, but the root \emph{fu-} retains \emph{u} where the regular
history predicts \emph{fo-} (\citeproc{ref-Campbell1959}{Campbell 1959,
150}).

\subsection{Development to Old
English}\label{development-to-old-english-152}

From \Recon{fúglaz} `fowl', the regular cascade yields \emph{fogol}
`fowl': the root vowel lowers before the following non-high vowel
(\citeproc{ref-RingeTaylor2014}{Ringe and Taylor 2014, 42--43};
\citeproc{ref-Campbell1959}{Campbell 1959, 43}), final \emph{z} is lost,
and the cluster is resolved by the usual medial vowel
(\citeproc{ref-RingeTaylor2014}{Ringe and Taylor 2014, 345};
\citeproc{ref-Campbell1959}{Campbell 1959, 150}). That is the expected
inherited outcome.

The attested Old English noun is \emph{fugol} `fowl', not \emph{fogol}
`fowl'. Luick and later handbooks treat this preservation of \emph{u} as
a small inherited residue, not as a categorical sound law
(\citeproc{ref-Luick1914}{Luick 1914-\/-1940, 148};
\citeproc{ref-RingeTaylor2014}{Ringe and Taylor 2014, 47}). The item
therefore remains a genuine lexical exception rather than the output of
a recoverable regular mechanism.

\subsection{Expected and attested
forms}\label{expected-and-attested-forms}

The comparison below sets the relevant forms side by side. It
distinguishes the regular prediction from the attested Old English noun.

{\def\LTcaptype{none} % do not increment counter
\begin{longtable}[]{@{}
  >{\raggedright\arraybackslash}p{(\linewidth - 4\tabcolsep) * \real{0.3333}}
  >{\raggedright\arraybackslash}p{(\linewidth - 4\tabcolsep) * \real{0.3333}}
  >{\raggedright\arraybackslash}p{(\linewidth - 4\tabcolsep) * \real{0.3333}}@{}}
\toprule\noalign{}
\begin{minipage}[b]{\linewidth}\raggedright
Form
\end{minipage} & \begin{minipage}[b]{\linewidth}\raggedright
Status
\end{minipage} & \begin{minipage}[b]{\linewidth}\raggedright
Relevance to this entry
\end{minipage} \\
\midrule\noalign{}
\endhead
\bottomrule\noalign{}
\endlastfoot
\emph{fogol} & computed regular output from \emph{*fúglaz} & establishes
the expected inherited outcome \\
\emph{fugol} & attested Old English form & Old English form here;
preserves unexplained root \emph{u} \\
\emph{fugel} & attested variant spelling & secondary spelling variant of
the attested noun \\
\end{longtable}
}

The unresolved point lies only in the root vowel. The medial \emph{-o-}
is regular, but no accepted lautgesetzlich pathway has been found from
\Recon{fúglaz} `fowl' to attested \emph{fugol} `fowl'.

\section{\texorpdfstring{rust --- OE
\emph{rust}}{rust --- OE rust}}\label{rust-oe-rust}

\index[iv]{02oe@\ivlangheader{Old English}{West Saxon normalization unmarked}!rust@\iventry{rust}{}}
\index[iv]{03pgmc@\ivlangheader{Proto-Germanic}{}!rustaz@\iventry{*rústaz}{}}

Derivation: \emph{*rústaz} yields regular \Recon{rost}; the Old English
form here is \emph{rust} (unexplained exception).

\subsection{Derivation trace}\label{derivation-trace-174}

Proto input: \emph{*rústaz}

\begingroup
\setlength{\fboxsep}{6pt}

\noindent\fbox{%
\begin{minipage}{0.97\linewidth}
\small
\begin{tabular}{@{}>{\raggedright\arraybackslash}p{0.485\linewidth}@{\hspace{0.030\linewidth}}>{\raggedright\arraybackslash}p{0.485\linewidth}@{}}
\begin{minipage}[t]{\linewidth}
\centering\textbf{Earlier Germanic changes}\par
\vspace{0.35em}
\raggedright
\centering\textbf{Northwest and West Germanic}\par
\raggedright
\vspace{0.2em}
\begin{tabularx}{\linewidth}{@{}>{\raggedright\arraybackslash}X@{\hspace{0.55em}}>{\raggedright\arraybackslash}p{0.14\linewidth}@{\hspace{0.25em}}}
\mbox{PNWGmc U Lowering} & \emph{*róstaz} \\
\mbox{EAF Final Z Deletion} & \emph{*rósta} \\
\end{tabularx}
\vspace{0.6em}
\centering\textbf{Early Anglo-Frisian}\par
\raggedright
\vspace{0.2em}
\raggedright [no change]\par
\end{minipage}
&
\begin{minipage}[t]{\linewidth}
\centering\textbf{Old English changes}\par
\vspace{0.35em}
\raggedright
\begin{tabularx}{\linewidth}{@{}>{\raggedright\arraybackslash}X@{\hspace{0.55em}}>{\raggedright\arraybackslash}p{0.14\linewidth}@{\hspace{0.25em}}}
\mbox{PWGmc Final Bare A Loss} & \emph{*róst} \\
\end{tabularx}
\end{minipage}
\\
\end{tabular}
\end{minipage}%
} \endgroup

Regular outcome: \Recon{rost}

Old English form: \emph{rust}

\subsection{Reconstruction and comparative
evidence}\label{reconstruction-and-comparative-evidence-154}

Orel gives both \Recon{rustō} `rust' and \Recon{rustaz} `rust', with
masculine/feminine alternatives and Old English
{\emph{rust}}\index[iv]{02oe@\ivlangheader{Old English}{West Saxon normalization unmarked}!rust@\iventry{rust}{}}
`rust', Old Saxon
{rost}\index[iv]{13os@\ivlangheader{Old Saxon}{}!rost@\iventry{rost}{}}
`rust' and Old High German
{rost}\index[iv]{10ohg@\ivlangheader{Old High German}{}!rost@\iventry{rost}{}}
`rust' (\citeproc{ref-Orel2003}{Orel 2003, 309}). The conventional
feminine \emph{ō}-stem and masculine \emph{az}-stem are distinct
morphological alternatives within his account. Ringe independently
reconstructs \Recon{rustaz} `rust'; his discussion of alternative
suffixal and cluster histories for \emph{-st-} concerns the formation,
not the explanation of retained \emph{u} (\citeproc{ref-Ringe2017}{Ringe
2017, 247--48}, especially p.~248).

The masculine Old English noun favors the masculine citation
reconstruction, selected here as \Recon{rústaz} `rust' with the
project's acute stress notation. This is a morphological and comparative
choice, not an exclusive disagreement between Orel and Ringe or a claim
that Orel's \emph{ō}-stem is unsupported. The comparative lowered forms
provide context, but output agreement does not choose a stem class: both
formations regularly predict \Recon{rost} `rust', not the attested Old
English noun.

\subsection{Old English evidence}\label{old-english-evidence-155}

The noun is attested, not reconstructed. Clark Hall gives
{\emph{rūst}}\index[iv]{02oe@\ivlangheader{Old English}{West Saxon normalization unmarked}!rust@\iventry{rūst}{}}
`rust' (m.) (\citeproc{ref-ClarkHall1960}{Clark Hall 1960, 245}), and
Toller's supplement to Bosworth-Toller records
{\emph{rúst}}\index[iv]{02oe@\ivlangheader{Old English}{West Saxon normalization unmarked}!rúst@\iventry{rúst}{}}
`rust', beside \emph{rust}, with a qualification about quantity
(\citeproc{ref-Toller1921}{Toller 1921, 691}). Their lexicographical
signs are retained as source spellings; the selected comparator remains
normalized \emph{rust}, without a newly adopted quantity interpretation.
The evidence establishes the noun and masculine gender, not an
independent explanation of its retained vowel.

\subsection{Development to Old
English}\label{development-to-old-english-153}

\hyperref[rule-PNWGmcULowering]{SC017 PNWGmcULowering} regularly lowers
the root \emph{u} of \Recon{rústaz} `rust' before the following non-high
vowel, giving \Recon{róstaz} `rust'. Final-\emph{z} loss and subsequent
regular developments complete the predicted \Recon{rost} `rust', not the
attested comparator \emph{rust}. Campbell describes the lowering and its
unresolved retained-\emph{u} residues
(\citeproc{ref-Campbell1959}{Campbell 1959, 43}, §115); Ringe and Taylor
discuss the remaining exceptions without establishing a general
consonantal exemption (\citeproc{ref-RingeTaylor2014}{Ringe and Taylor
2014, 27--34}).

The control \Recon{frústą} `frost' regularly gives
{\emph{forst}}\index[iv]{02oe@\ivlangheader{Old English}{West Saxon normalization unmarked}!forst@\iventry{forst}{}}
`frost', so \emph{-st-} cannot be exempted generally. Selecting the
masculine stem is not a phonological repair: \emph{rust} remains an
unexplained lexical retention and the mismatch total does not fall. No
exception is inserted into the sound law.

\subsection{Paradigm comparison}\label{paradigm-comparison-27}

Fresh production and isolated-wrapper diagnostics give the following
results. The oblique-type constructions test conditioning; they are not
newly attested Proto-Germanic paradigm cells.

{\def\LTcaptype{none} % do not increment counter
\begin{longtable}[]{@{}
  >{\raggedright\arraybackslash}p{(\linewidth - 4\tabcolsep) * \real{0.3333}}
  >{\raggedright\arraybackslash}p{(\linewidth - 4\tabcolsep) * \real{0.3333}}
  >{\raggedright\arraybackslash}p{(\linewidth - 4\tabcolsep) * \real{0.3333}}@{}}
\toprule\noalign{}
\begin{minipage}[b]{\linewidth}\raggedright
Form
\end{minipage} & \begin{minipage}[b]{\linewidth}\raggedright
Role
\end{minipage} & \begin{minipage}[b]{\linewidth}\raggedright
Result
\end{minipage} \\
\midrule\noalign{}
\endhead
\bottomrule\noalign{}
\endlastfoot
\Recon{rústaz}\index[iv]{03pgmc@\ivlangheader{Proto-Germanic}{}!rustaz@\iventry{*rústaz}{}}
& Selected masculine citation form & \Recon{rost}, mismatch \\
\Recon{rústō}\index[iv]{03pgmc@\ivlangheader{Proto-Germanic}{}!rusto@\iventry{*rústō}{}}
& Orel's alternative \emph{ō}-stem & \Recon{rost}, mismatch \\
\Recon{rústas} & Genitive-singular-type probe & \Recon{rostes} \\
\Recon{rústai} & Dative-singular-type probe & \Recon{roste} \\
\Recon{rústi} & High-front-vowel diagnostic & \Recon{ryst} \\
\Recon{rústu} & Marginal high-vowel diagnostic construction &
\Recon{rust} \\
\end{longtable}
}

The hypothetical instrumental-type \Recon{rústu} `rust' is a diagnostic
construction, not an attested Proto-Germanic citation form or an adopted
historical solution. Its matching output does not establish a source for
the ordinary citation noun.

\section{\texorpdfstring{wolf --- OE
\emph{wulf}}{wolf --- OE wulf}}\label{wolf-oe-wulf}

\index[iv]{02oe@\ivlangheader{Old English}{West Saxon normalization unmarked}!bolt@\iventry{bolt}{}}
\index[iv]{02oe@\ivlangheader{Old English}{West Saxon normalization unmarked}!wulf@\iventry{wulf}{}}
\index[iv]{03pgmc@\ivlangheader{Proto-Germanic}{}!wulfaz@\iventry{*wúlfaz}{}}
\index[iv]{10ohg@\ivlangheader{Old High German}{}!wolf@\iventry{wolf}{}}

Derivation: \emph{*wúlfaz} yields regular \Recon{wolf}; the Old English
form here is \emph{wulf} (unexplained exception).

\subsection{Derivation trace}\label{derivation-trace-175}

Proto input: \emph{*wúlfaz}

\begingroup
\setlength{\fboxsep}{6pt}

\noindent\fbox{%
\begin{minipage}{0.97\linewidth}
\small
\begin{tabular}{@{}>{\raggedright\arraybackslash}p{0.485\linewidth}@{\hspace{0.030\linewidth}}>{\raggedright\arraybackslash}p{0.485\linewidth}@{}}
\begin{minipage}[t]{\linewidth}
\centering\textbf{Earlier Germanic changes}\par
\vspace{0.35em}
\raggedright
\centering\textbf{Northwest and West Germanic}\par
\raggedright
\vspace{0.2em}
\begin{tabularx}{\linewidth}{@{}>{\raggedright\arraybackslash}X@{\hspace{0.55em}}>{\raggedright\arraybackslash}p{0.14\linewidth}@{\hspace{0.25em}}}
\mbox{PNWGmc U Lowering} & \emph{*wólfaz} \\
\mbox{EAF Final Z Deletion} & \emph{*wólfa} \\
\end{tabularx}
\vspace{0.6em}
\centering\textbf{Early Anglo-Frisian}\par
\raggedright
\vspace{0.2em}
\raggedright [no change]\par
\end{minipage}
&
\begin{minipage}[t]{\linewidth}
\centering\textbf{Old English changes}\par
\vspace{0.35em}
\raggedright
\begin{tabularx}{\linewidth}{@{}>{\raggedright\arraybackslash}X@{\hspace{0.55em}}>{\raggedright\arraybackslash}p{0.14\linewidth}@{\hspace{0.25em}}}
\mbox{PWGmc Final Bare A Loss} & \emph{*wólf} \\
\end{tabularx}
\end{minipage}
\\
\end{tabular}
\end{minipage}%
} \endgroup

Regular outcome: \Recon{wolf}

Old English form: \emph{wulf}

\subsection{Reconstruction and comparative
evidence}\label{reconstruction-and-comparative-evidence-155}

The inherited noun is an a-stem: Kroonen gives \emph{*wulfa-}, and Ringe
and Taylor list PGmc \Recon{wulfaz} `wolf' among the inherited words
that preserve \emph{u} in Old English beside Old High German \emph{wolf}
`wolf' (\citeproc{ref-Kroonen2013}{Kroonen 2013, 638};
\citeproc{ref-RingeTaylor2014}{Ringe and Taylor 2014, 47}). Campbell
accordingly names Old English \emph{wulf} as an exception to the regular
lowering of stressed \emph{u} before a following non-high vowel
(\citeproc{ref-Campbell1959}{Campbell 1959, § 115}).

The older literature often notices that the exceptional words cluster
near labials. Bülbring lists \emph{full}, \emph{wulle}, and \emph{wulf}
together, but he also says that the ordinary rule still gives \emph{o}
in comparable forms such as \emph{folc} and \emph{bolt}
(\citeproc{ref-Bulbring1902}{Bülbring 1902, § 116}). Luick rejects a
categorical labial blocker on the same grounds and prefers a lexical or
analogical account instead (\citeproc{ref-Luick1914}{Luick 1914-\/-1940,
148}).

\subsection{Old English evidence}\label{old-english-evidence-156}

Campbell treats \emph{wulf} `wolf' as part of the exceptional \emph{u}
set (\citeproc{ref-Campbell1959}{Campbell 1959, § 115}).
Sievers-Brunner notes that oblique
{\emph{wulfe}}\index[iv]{02oe@\ivlangheader{Old English}{West Saxon normalization unmarked}!wulfe@\iventry{wulfe}{}}
`wolf (dat.)' continues
{\emph{wulfi}}\index[iv]{02oe@\ivlangheader{Old English}{West Saxon normalization unmarked}!wulfi@\iventry{wulfi}{}}
`wolf (dat.)' or older
{\emph{wulfai}}\index[iv]{02oe@\ivlangheader{Old English}{West Saxon normalization unmarked}!wulfai@\iventry{wulfai}{}}
`wolf (dat.)' (\citeproc{ref-SieversBrunner1965}{Brunner 1965, § 160}).

The surviving oblique forms do not supply a regular route back to bare
\emph{wulf} `wolf'. They belong to the same lexeme, but they do not
remove the explanatory problem presented by the citation form.

\subsection{Development to Old
English}\label{development-to-old-english-154}

Under the ordinary Northwest Germanic lowering of stressed \emph{u}
before a following non-high vowel, the citation-form input would point
toward \emph{o}-vocalism (\citeproc{ref-RingeTaylor2014}{Ringe and
Taylor 2014, 42--44}). The compact trace shows exactly that path:
\Recon{wúlfaz} `wolf' \textgreater{} \Recon{wólfaz} `wolf'
\textgreater{} \Recon{wólfa} `wolf' \textgreater{} wolf.

A high-vowel oblique input would behave differently. There the following
high vowel would block the lowering of \emph{u}, but the same
environment would also trigger i-umlaut, so the regular control result
would be
{\emph{wylf}}\index[iv]{02oe@\ivlangheader{Old English}{West Saxon normalization unmarked}!wylf@\iventry{wylf}{}}
`wolf' or \emph{wylfe}, not bare \emph{wulf}. The attested noun
therefore remains unexplained at the citation-form level.

\subsection{Expected and attested
forms}\label{expected-and-attested-forms-1}

The comparison below sets the regular inherited outcomes beside the
attested Old English noun.

{\def\LTcaptype{none} % do not increment counter
\begin{longtable}[]{@{}
  >{\raggedright\arraybackslash}p{(\linewidth - 4\tabcolsep) * \real{0.3333}}
  >{\raggedright\arraybackslash}p{(\linewidth - 4\tabcolsep) * \real{0.3333}}
  >{\raggedright\arraybackslash}p{(\linewidth - 4\tabcolsep) * \real{0.3333}}@{}}
\toprule\noalign{}
\begin{minipage}[b]{\linewidth}\raggedright
Form / interpretation
\end{minipage} & \begin{minipage}[b]{\linewidth}\raggedright
Status
\end{minipage} & \begin{minipage}[b]{\linewidth}\raggedright
Relevance to this entry
\end{minipage} \\
\midrule\noalign{}
\endhead
\bottomrule\noalign{}
\endlastfoot
\emph{*wúlfaz} \textgreater{} \emph{wolf} & computed regular output from
the citation form & shows the regular development expected from the
inherited a-stem \\
\emph{OHG wolf} & comparative regular cognate & confirms that the
\emph{o}-vocalism is the ordinary outcome \\
\emph{*wúlfi} / \emph{*wúlfis} \textgreater{} \emph{wylf} / \emph{wylfe}
& expected high-vowel control forms & shows why oblique high-vowel cells
do not solve the noun's vowel history \\
\emph{wulf} & attested Old English noun & attested Old English form; the
preservation of \emph{u} remains unexplained \\
\end{longtable}
}

\section{\texorpdfstring{wool --- OE
\emph{wull}}{wool --- OE wull}}\label{wool-oe-wull}

\index[iv]{02oe@\ivlangheader{Old English}{West Saxon normalization unmarked}!wull@\iventry{wull}{}}
\index[iv]{03pgmc@\ivlangheader{Proto-Germanic}{}!wullo@\iventry{*wúllō}{}}
\index[iv]{03pgmc@\ivlangheader{Proto-Germanic}{}!wullǭ@\iventry{*wúllǭ}{}}

Derivation: \emph{*wúllō} yields regular \Recon{woll}; the Old English
form here is \emph{wull} (unexplained exception).

\subsection{Derivation trace}\label{derivation-trace-176}

Proto input: \emph{*wúllō}

\begingroup
\setlength{\fboxsep}{6pt}

\noindent\fbox{%
\begin{minipage}{0.97\linewidth}
\small
\begin{tabular}{@{}>{\raggedright\arraybackslash}p{0.611\linewidth}@{\hspace{0.030\linewidth}}>{\raggedright\arraybackslash}p{0.359\linewidth}@{}}
\begin{minipage}[t]{\linewidth}
\centering\textbf{Earlier Germanic changes}\par
\vspace{0.35em}
\raggedright
\centering\textbf{Northwest and West Germanic}\par
\raggedright
\vspace{0.2em}
\begin{tabularx}{\linewidth}{@{}>{\raggedright\arraybackslash}X@{\hspace{0.55em}}>{\raggedright\arraybackslash}p{0.14\linewidth}@{\hspace{0.25em}}}
\mbox{PNWGmc U Lowering} & \emph{*wóllō} \\
\mbox{PNWGmc Final Long O Raising} & \emph{*wóllu} \\
\end{tabularx}
\vspace{0.6em}
\centering\textbf{Early Anglo-Frisian}\par
\raggedright
\vspace{0.2em}
\raggedright [no change]\par
\end{minipage}
&
\begin{minipage}[t]{\linewidth}
\centering\textbf{Old English changes}\par
\vspace{0.35em}
\raggedright
\begin{tabularx}{\linewidth}{@{}>{\raggedright\arraybackslash}X@{\hspace{0.55em}}>{\raggedright\arraybackslash}p{0.14\linewidth}@{\hspace{0.25em}}}
\mbox{OE High Vowel Apocope} & \emph{*wóll} \\
\end{tabularx}
\end{minipage}
\\
\end{tabular}
\end{minipage}%
} \endgroup

Regular outcome: \Recon{woll}

Old English form: \emph{wull}

\subsection{The inherited paradigm}\label{the-inherited-paradigm}

The selected \Recon{wúllō} `wool' represents the feminine nominative
singular. It must be distinguished from accusative \emph{*wúllǭ}, whose
final vowel was nasalized. Ringe explicitly reconstructs both cells;
Orel and Kroonen discuss the inherited formation and its comparative
history (\citeproc{ref-Ringe2017}{Ringe 2017, 171--73, 241--44};
\citeproc{ref-Orel2003}{Orel 2003, 474};
\citeproc{ref-Kroonen2013}{Kroonen 2013, 598--99}). Attested
{\emph{wull}}\index[iv]{02oe@\ivlangheader{Old English}{West Saxon normalization unmarked}!wull@\iventry{wull}{}}
`wool' and
{\emph{wulle}}\index[iv]{02oe@\ivlangheader{Old English}{West Saxon normalization unmarked}!wulle@\iventry{wulle}{}}
`wool' preserve a high root vowel. The latter is not, merely by its
shape, evidence that the whole noun was originally a weak n-stem
(\citeproc{ref-Campbell1959}{Campbell 1959, 43}, §115;
\citeproc{ref-Bulbring1902}{Bülbring 1902, 45--46}, §116).

\subsection{Regular development and the alternative
explanation}\label{regular-development-and-the-alternative-explanation}

CAPR provisionally follows Fulk's lowering-before-final-raising
chronology (\citeproc{ref-Fulk2018}{Fulk 2018, 56}). Both oral \emph{*ō}
and nasal \emph{*ǭ} condition root lowering before their distinct
subsequent developments:

{\def\LTcaptype{none} % do not increment counter
\begin{longtable}[]{@{}
  >{\raggedright\arraybackslash}p{(\linewidth - 4\tabcolsep) * \real{0.3333}}
  >{\raggedright\arraybackslash}p{(\linewidth - 4\tabcolsep) * \real{0.3333}}
  >{\raggedright\arraybackslash}p{(\linewidth - 4\tabcolsep) * \real{0.3333}}@{}}
\toprule\noalign{}
\begin{minipage}[b]{\linewidth}\raggedright
Inherited cell
\end{minipage} & \begin{minipage}[b]{\linewidth}\raggedright
Preferred chronology
\end{minipage} & \begin{minipage}[b]{\linewidth}\raggedright
Raising before lowering
\end{minipage} \\
\midrule\noalign{}
\endhead
\bottomrule\noalign{}
\endlastfoot
nominative \emph{*wúllō} & \emph{*wóllō} → \emph{*wóllu} → regular
\emph{*woll} & \emph{*wúllu} → regular \emph{*wull} \\
accusative \emph{*wúllǭ} & \emph{*wóllǭ} → regular \emph{*wolle} &
\emph{*wóllǭ} → regular \emph{*wolle} \\
\end{longtable}
}

These are sound-law predictions, not additional attestations. Earlier
raising supplies a genuine nominative \emph{*wullu} with high root vowel
and an independently conditioned low-root oblique. Generalization of the
nominative vowel could then produce the attested high-root paradigm:
this is Ringe and Taylor's plausible but unproved explanation, not an
impossible escape cell (\citeproc{ref-RingeTaylor2014}{Ringe and Taylor
2014, 32}). Old Frisian
{\emph{wolle}}\index[iv]{11ofris@\ivlangheader{Old Frisian}{}!wolle@\iventry{wolle}{}}
`wool' and Old High German
{\emph{wolla}}\index[iv]{10ohg@\ivlangheader{Old High German}{}!wolla@\iventry{wolla}{}}
`wool' support the comparative reality of a lowered root, but do not
attest that particular early English accusative
(\citeproc{ref-Kroonen2013}{Kroonen 2013, 598--99}).

The main comparison of the two chronologies, the evidence from other
ō-stems and the costs of the proposed analogies is in
\hyperref[rule-PNWGmcFinalLongORaising]{SC019 PNWGmcFinalLongORaising}.
The preferred ordering does not regularly derive the attested high root
here. The derivational input, target and classification therefore remain
unchanged: regular \Recon{woll} `wool', not attested \emph{wull}, is
what the current sound laws produce. No lexical exception or analogical
levelling has been inserted into those laws.

Luick discusses the alternations in §§77--79, especially §78 (printed
pp.~105--106), not §§105--106. He rejects a proposed labial protection
rule using regular low-root counterexamples
(\citeproc{ref-Luick1914}{Luick 1914-\/-1940, 104--6}, §§77--79). The
separate trimoric genitive \emph{*-ôz} must not be substituted for
either cell above: its analysis differs from the bimoric genitive and
its complete computational pathway remains unmodeled
(\citeproc{ref-Fulk2018}{Fulk 2018, 155};
\citeproc{ref-RingeTaylor2014}{Ringe and Taylor 2014, 59, 299}).

\backmatter

\chapter{References}\label{references}

\protect\phantomsection\label{refs}
\begin{CSLReferences}{1}{1}
\bibitem[\citeproctext]{ref-Adamczyk2001}
Adamczyk, Elżbieta. 2001. {`Old {English} Reflexes of {Sievers'} Law'}.
\emph{Studia Anglica Posnaniensia: An International Review of English
Studies} 36: 61--72. \url{http://hdl.handle.net/10593/18335}.

\bibitem[\citeproctext]{ref-Bammesberger1981}
Bammesberger, Alfred. 1981. {`Die Betonung Der Nominalen *Ga-Komposita
Im Urgermanischen'}. \emph{Beitr{ä}ge Zur Geschichte Der Deutschen
Sprache Und Literatur} 103: 377--91.
\url{https://doi.org/10.1515/bgsl.1981.1981.103.377}.

\bibitem[\citeproctext]{ref-Bammesberger1990}
Bammesberger, Alfred. 1990. \emph{Die Morphologie Des Urgermanischen
Nomens}. Vol. 2. Untersuchungen Zur Vergleichenden Grammatik Der
Germanischen Sprachen. Carl Winter.

\bibitem[\citeproctext]{ref-Bammesberger1992}
Bammesberger, Alfred. 1992. {`The Place of {English} in {Germanic} and
{Indo-European}'}. In \emph{The {Cambridge} History of the {English}
Language. Volume {I}: The Beginnings to 1066}, edited by Richard M.
Hogg. Cambridge University Press.

\bibitem[\citeproctext]{ref-Bammesberger1997}
Bammesberger, Alfred. 1997. {`Die Vorform von Altenglisch {h{æ}rfest}'}.
\emph{Anglia: Zeitschrift f{ü}r Englische Philologie} 115 (2): 223--30.
\url{https://doi.org/10.1515/angl.1997.115.2.223}.

\bibitem[\citeproctext]{ref-Barrack1975}
Barrack, Charles Michael. 1975. \emph{A Diachronic Phonology from
Proto-Germanic to Old English Stressing West-Saxon Conditions}. Vol.
144. Janua Linguarum. Series Practica. Mouton.
\url{https://doi.org/10.1515/9783110873511}.

\bibitem[\citeproctext]{ref-Bennett1950}
Bennett, William H. 1950. {`The {Germanic} Development of
{Indo-European} {ē}'}. \emph{Language} 26 (2): 232--35.
\url{https://www.jstor.org/stable/410059}.

\bibitem[\citeproctext]{ref-Berg2014}
Berg, Joakim P. 2014. {`Old {English} Unstressed Vowels: Dialects and
Diachrony'}. Master's thesis, University of Oslo.
\url{https://hdl.handle.net/10852/39907}.

\bibitem[\citeproctext]{ref-BosworthToller1898}
Bosworth, Joseph, and T. Northcote Toller. 1898. \emph{An {Anglo-Saxon}
Dictionary}. Clarendon Press.

\bibitem[\citeproctext]{ref-Bremmer2008}
Bremmer, Jr., Rolf H. 2008. {`North-{Sea} {Germanic} at the Cross-Roads:
The Emergence of {Frisian} and {Hollandish}'}. \emph{NOWELE:
North-Western European Language Evolution} 54/55: 279--308.
\url{https://doi.org/10.1075/nowele.54-55.09bre}.

\bibitem[\citeproctext]{ref-Bremmer2009}
Bremmer, Jr., Rolf H. 2009. \emph{An Introduction to {Old Frisian}:
History, Grammar, Reader, Glossary}. John Benjamins.
\url{https://doi.org/10.1075/z.147}.

\bibitem[\citeproctext]{ref-Bright1917}
Bright, James W. 1917. \emph{An Anglo-Saxon Reader}. 4th edn. Henry
Holt; Company.

\bibitem[\citeproctext]{ref-BrightCassidyRingler1971}
Bright, James W. 1971. \emph{Bright's {Old English} Grammar and Reader}.
3rd edn. Edited by Frederic G. Cassidy and Richard N. Ringler. Holt,
Rinehart; Winston.

\bibitem[\citeproctext]{ref-SieversBrunner1965}
Brunner, Karl. 1965. \emph{{Altenglische Grammatik nach der
angels{ä}chsischen Grammatik von Eduard Sievers}}. 3rd edn. Max
Niemeyer.

\bibitem[\citeproctext]{ref-Bulbring1902}
Bülbring, Karl D. 1902. \emph{Altenglisches Elementarbuch. 1:
Lautlehre}. Carl Winter.

\bibitem[\citeproctext]{ref-Campbell1939}
Campbell, Alistair. 1939. {`Some {Old} {Frisian} Sound-Changes'}.
\emph{Transactions of the Philological Society} 38 (1): 78--107.
\url{https://doi.org/10.1111/j.1467-968X.1939.tb00203.x}.

\bibitem[\citeproctext]{ref-Campbell1959}
Campbell, Alistair. 1959. \emph{Old {English} Grammar}. Clarendon Press.

\bibitem[\citeproctext]{ref-Cercignani1979}
Cercignani, Fausto. 1979. {`{Proto-Germanic} *{i} and *{e} Revisited'}.
\emph{Journal of English and Germanic Philology} 78 (1): 72--82.

\bibitem[\citeproctext]{ref-Cercignani1980}
Cercignani, Fausto. 1980. {`Early {``Umlaut''} Phenomena in the
{Germanic} Languages'}. \emph{Language} 56 (1): 126--36.

\bibitem[\citeproctext]{ref-Cermak2005}
Čermák, Jan. 2005. {`Notes on Obscured Compounds in {Old English}'}. In
\emph{Patterns: A Festschrift for {Libuše Dušková}}, edited by Jan
Čermák, Aleš Klégr, Markéta Malá, and Pavlína Šaldová. Karolinum.

\bibitem[\citeproctext]{ref-Christiansen2018}
Christiansen, Bethany. 2018. {`{Scytel}: A New {Old English} Word for
{``Penis''}'}. \emph{Anglia} 136 (4): 581--610.
\url{https://doi.org/10.1515/ang-2018-0060}.

\bibitem[\citeproctext]{ref-CiszekKiliszewska2015}
Ciszek-Kiliszewska, Ewa. 2015. {`Old {English} Suffix {-Cund} and Its
Fate'}. \emph{Kwartalnik Neofilologiczny} 62 (3): 487--504.

\bibitem[\citeproctext]{ref-ClarkHall1960}
Clark Hall, J. R. 1960. \emph{A Concise {Anglo-Saxon} Dictionary}. 4th
edn. Cambridge University Press.

\bibitem[\citeproctext]{ref-Colman2011}
Colman, Fran. 2011. {`On the Moneyers' Names {Buga} and {Boia} on
{Anglo-Saxon} Coins'}. \emph{Nomina} 34: 91--120.
\url{https://www.snsbi.org.uk/Nomina_articles/Nomina_34_Colman.pdf}.

\bibitem[\citeproctext]{ref-Crawford1922}
Crawford, S. J., ed. 1922. \emph{The {Old English} Version of the
Heptateuch, {{Æ}lfric}'s Treatise on the Old and New Testament and His
Preface to Genesis}. Humphrey Milford, Oxford University Press, for the
Early English Text Society.

\bibitem[\citeproctext]{ref-Crist2001}
Crist, Sean J. 2001. {`Conspiracy in Historical Phonology'}. Ph.\{D\}.
dissertation, University of Pennsylvania.

\bibitem[\citeproctext]{ref-Crist2002}
Crist, Sean J. 2002. \emph{An Analysis of *{z} Loss in {West Germanic}}.
Paper presented at the Linguistic Society of America Annual Meeting.

\bibitem[\citeproctext]{ref-Fulk2010}
Fulk, R. D. 2010. {`The Roles of Phonology and Analogy in {Old English}
High Vowel Deletion'}. \emph{Transactions of the Philological Society}
108 (2): 126--44.
\url{https://doi.org/10.1111/j.1467-968X.2010.01233.x}.

\bibitem[\citeproctext]{ref-Fulk2018}
Fulk, R. D. 2018. \emph{A Comparative Grammar of the Early {Germanic}
Languages}. Vol. 3. Studies in Germanic Linguistics. John Benjamins.

\bibitem[\citeproctext]{ref-Fulk1998}
Fulk, Robert D. 1998. {`The Chronology of {Anglo-Frisian} Sound
Changes'}. In \emph{Approaches to {Old Frisian} Philology}, edited by
Rolf H. Bremmer Jr., Thomas S. B. Johnston, and Oebele Vries, vol. 49.
Amsterdamer Beitr{ä}ge Zur {ä}lteren Germanistik. Rodopi.

\bibitem[\citeproctext]{ref-Goblirsch1991}
Goblirsch, Kurt Gustav. 1991. {`{Germanic} \emph{Ai} and \emph{Au} in
{Anglo-Frisian}'}. \emph{Amsterdamer Beitr{ä}ge Zur {ä}lteren
Germanistik} 33: 17--23.
\url{https://doi.org/10.1163/18756719-033-01-90000004}.

\bibitem[\citeproctext]{ref-Goering2016}
Goering, Nelson. 2016. {`Early {Old English} Foot Structure'}.
\emph{Transactions of the Philological Society} 114 (2): 171--97.
\url{https://doi.org/10.1111/1467-968X.12079}.

\bibitem[\citeproctext]{ref-Goering2023}
Goering, Nelson. 2023. {`The {Hēafudu}-Problem: Early {Old English} Foot
Structure'}. In \emph{Prosody in Medieval {English} and {Norse}}.
British Academy.
\url{https://doi.org/10.5871/bacad/9780197267462.003.0004}.

\bibitem[\citeproctext]{ref-Hall2007}
Hall, Alaric. 2007. {`The Evidence for {maran}, the {Anglo-Saxon}
'Nightmares''}. \emph{Neophilologus} 91 (2): 299--317.
\url{https://doi.org/10.1007/s11061-005-4256-8}.

\bibitem[\citeproctext]{ref-Heidermanns1993}
Heidermanns, Frank. 1993. \emph{{Etymologisches Wörterbuch der
germanischen Primäradjektive}}. Vol. 33. Studia Linguistica Germanica.
Walter de Gruyter. \url{https://doi.org/10.1515/9783110871616}.

\bibitem[\citeproctext]{ref-HerrenEtAl2026}
Herren, Michael W., Cameron Laird, David W. Porter, Hans Sauer, and
Gernot Wieland, eds. 2026. \emph{The {É}pinal-{Erfurt} Glossary}.
Dictionary of Old English Project, University of Toronto.
\url{https://epinal-erfurt.artsci.utoronto.ca/index.php/new-revised-edition/}.

\bibitem[\citeproctext]{ref-HerrenSauer2016}
Herren, Michael W., and Hans Sauer. 2016. {`Towards a New Edition of the
{É}pinal-{Erfurt} Glossary: A Sample'}. \emph{The Journal of Medieval
Latin} 26: 125--98. \url{https://doi.org/10.1484/J.JML.5.112082}.

\bibitem[\citeproctext]{ref-Hill2017}
Hill, Eugen. 2017. {`{West Germanic} Monosyllabic Lengthening and
{Gothic} Breaking as Partially {Proto-Germanic} Developments: The
Evidence of Pronominal Place Adverbs {``Here''}, {``Where''} and
{``There''}'}. \emph{NOWELE. North-Western European Language Evolution}
70 (2): 135--70. \url{https://doi.org/10.1075/nowele.70.2.02hil}.

\bibitem[\citeproctext]{ref-Hirt1931}
Hirt, Hermann. 1931-\/-1934. \emph{Handbuch Des {Urgermanischen}}. Carl
Winter.

\bibitem[\citeproctext]{ref-Hock1973}
Hock, Hans Henrich. 1973. {`On the Phonemic Status of {Germanic} {e} and
{i}'}. In \emph{Issues in Linguistics: Papers in Honor of {Henry} and
{Ren{é}e Kahane}}, edited by B. B. Kachru, R. B. Lees, Y. Malkiel, A.
Pietrangeli, and S. Saporta. University of Illinois Press.

\bibitem[\citeproctext]{ref-Hogg1979}
Hogg, Richard M. 1979. {`{Old English} Palatalization'}.
\emph{Transactions of the Philological Society} 77 (1): 89--113.
\url{https://doi.org/10.1111/j.1467-968X.1979.tb00855.x}.

\bibitem[\citeproctext]{ref-HoggGrammar1992}
Hogg, Richard M. 1992a. \emph{A Grammar of Old {English}. Volume 1:
Phonology}. Blackwell.

\bibitem[\citeproctext]{ref-HoggPhonology1992}
Hogg, Richard M. 1992b. {`Phonology and Morphology'}. In \emph{The
{Cambridge} History of the {English} Language. Volume {I}: The
Beginnings to 1066}, edited by Richard M. Hogg. Cambridge University
Press.

\bibitem[\citeproctext]{ref-Hogg1996}
Hogg, Richard M. 1996. {`Tertiary Stress in {Old English}: Some
Reflections on Explanatory Inadequacy'}. In \emph{{English Historical
Linguistics 1994}}, edited by Derek Britton, vol. 135. Current Issues in
Linguistic Theory. John Benjamins.
\url{https://doi.org/10.1075/cilt.135.03hog}.

\bibitem[\citeproctext]{ref-HoggGrammar2011}
Hogg, Richard M. 2011. \emph{A Grammar of Old {English}. Volume 1:
Phonology}. Wiley-Blackwell.

\bibitem[\citeproctext]{ref-Howell2024}
Howell, Robert B. 2024. {`What Do We Really Know about {Gothic}
Breaking? {On} the Problem of Consonantally Conditioned Vowel Mutations
in {Germanic}'}. Ch. 4 in \emph{Studies in {Gothic}}, edited by Jared S.
Klein and Artūras Ratkus. Oxford University Press.
\url{https://doi.org/10.1093/oso/9780198896692.003.0004}.

\bibitem[\citeproctext]{ref-Hulden2009}
Hulden, Mans. 2009. {`Foma: A Finite-State Compiler and Library'}.
\emph{Proceedings of the Demonstrations Session at EACL 2009} (Athens),
29--32.

\bibitem[\citeproctext]{ref-Jasanoff1973}
Jasanoff, Jay H. 1973. {`The {Germanic} Third Weak Class'}.
\emph{Language} 49 (4): 850--70. \url{https://doi.org/10.2307/412065}.

\bibitem[\citeproctext]{ref-KaplanKay1994}
Kaplan, Ronald M., and Martin Kay. 1994. {`Regular Models of
Phonological Rule Systems'}. \emph{Computational Linguistics} 20 (3):
331--78.

\bibitem[\citeproctext]{ref-Kastovsky1992}
Kastovsky, Dieter. 1992. {`Semantics and Vocabulary'}. In \emph{The
{Cambridge} History of the {English} Language. Volume {I}: The
Beginnings to 1066}, edited by Richard M. Hogg. Cambridge University
Press.

\bibitem[\citeproctext]{ref-Kilday2024}
Kilday, Douglas G. 2024. {`{Crist's} Law, {Smith's} Law, and {English}
\emph{Wizen}'}. Unpublished manuscript.

\bibitem[\citeproctext]{ref-KlugeSeebold2011}
Kluge, Friedrich. 2011. \emph{{Etymologisches W{ö}rterbuch der deutschen
Sprache}}. 25th edn. Edited by Elmar Seebold. De Gruyter.

\bibitem[\citeproctext]{ref-Krogmann1935}
Krogmann, Willy. 1935. {`{Die Sippen germ. *līt- und *leut- „klein``}'}.
\emph{Indogermanische Forschungen} 53 (1): 44--48.
\url{https://doi.org/10.1515/if-1935-0108}.

\bibitem[\citeproctext]{ref-Kroonen2011}
Kroonen, Guus. 2011. \emph{The {Proto-Germanic} n-Stems: A Study in
Diachronic Morphophonology}. Vol. 18. Leiden Studies in Indo-European.
Rodopi.

\bibitem[\citeproctext]{ref-Kroonen2013}
Kroonen, Guus. 2013. \emph{Etymological Dictionary of {Proto-Germanic}}.
Vol. 11. Leiden Indo-European Etymological Dictionary Series. Brill.

\bibitem[\citeproctext]{ref-Laker2007}
Laker, Stephen. 2007. {`Palatalization of Velars: A Major Link of {Old}
{English} and {Old} {Frisian}'}. \emph{Amsterdamer Beiträge Zur Älteren
Germanistik} 64: 165--84.

\bibitem[\citeproctext]{ref-Levitsky2003}
Levickij, V. V. 2003. {`{Indoevropejskie korni *sep(h)ə-/sp(h)ē-,
*erə-/rē-, *elə-/lē- «rezatʹ, svjazyvatʹ, gnutʹ» i ix derivaty v
germanskix jazykax}'}. \emph{Voprosy Jazykoznanija}, no. 1: 37--62.

\bibitem[\citeproctext]{ref-Lid1952}
Lid, Nils. 1952. {`Den Nordiske Nominativ Singularis Av Maskuline
\emph{an}-Stammer'}. \emph{Norsk Tidsskrift for Sprogvidenskap} (Oslo)
16: 237--40.

\bibitem[\citeproctext]{ref-Lindsay1921}
Lindsay, W. M., ed. 1921. \emph{The {Corpus} Glossary}. Cambridge
University Press.
\url{https://archive.org/details/corpusglossary00linduoft}.

\bibitem[\citeproctext]{ref-Lloyd1966}
Lloyd, Albert L. 1966. {`Is There an {a-}Umlaut of {i} in {Germanic}?'}
\emph{Language} 42 (4): 738--45. \url{https://doi.org/10.2307/411829}.

\bibitem[\citeproctext]{ref-LloydLuhrSpringer1998}
Lloyd, Albert L., Rosemarie Lühr, and Otto Springer. 1998.
\emph{Etymologisches w{ö}rterbuch Des Althochdeutschen. Band {II}: {bı̂
-- ezzo}}. Vandenhoeck \& Ruprecht.

\bibitem[\citeproctext]{ref-Lucas2016}
Lucas, Peter J. 2016. {`Cambridge, {Corpus Christi College} 144: The
{Corpus Glossary}'}. \emph{Anglo-Saxon Manuscripts in Microfiche
Facsimile} 25: 11--18.
\url{https://journals.lib.sfu.ca/index.php/asmmf/article/view/5964}.

\bibitem[\citeproctext]{ref-LucasLucas2012}
Lucas, Peter J., and Angela M. Lucas. 2012. {`{É}pinal, Biblioth{è}que
Municipale 72 (Olim 66:7): Homilies; {É}pinal Glossary'}.
\emph{Anglo-Saxon Manuscripts in Microfiche Facsimile} 18: 67--83.
\url{https://journals.lib.sfu.ca/index.php/asmmf/article/view/5810}.

\bibitem[\citeproctext]{ref-Luehr1982}
Lühr, Rosemarie. 1982. \emph{{Studien zur Sprache des Hildebrandliedes.
Teil II: Kommentar}}. Peter Lang.
\url{https://dwee.saw-leipzig.de/Rosemarie_Luehr/userfiles/downloads/Hildebrandlied.pdf}.

\bibitem[\citeproctext]{ref-Luehr2017}
Lühr, Rosemarie. 2017. {`{Semantische Gegensätze}'}. \emph{Jahrbuch Für
Germanistische Sprachgeschichte} 8: 147--63.
\url{https://doi.org/10.1515/jbgsg-2017-0011}.

\bibitem[\citeproctext]{ref-Luick1914}
Luick, Karl. 1914-\/-1940. \emph{{Historische Grammatik der englischen
Sprache}}. Tauchnitz.

\bibitem[\citeproctext]{ref-McCullyHogg1990}
McCully, C. B., and R. M. Hogg. 1990. {`An Account of {Old English}
Stress'}. \emph{Journal of Linguistics} 26 (2): 315--39.
\url{https://doi.org/10.1017/S0022226700014699}.

\bibitem[\citeproctext]{ref-Meillet1902}
Meillet, Antoine. 1902. \emph{{Études sur l'étymologie et le vocabulaire
du vieux slave}}. Vol. 1. Bibliothèque de l'école Des Hautes Études,
Sciences Historiques Et Philologiques 139. Émile Bouillon.

\bibitem[\citeproctext]{ref-Melazzo2014}
Melazzo, Lucio. 2014. {`{Revisiting the Etymology of Greek μέγας,
μεγάλη, μέγα}'}. \emph{Studi e Saggi Linguistici} 52 (1): 9--17.
\url{https://www.studiesaggilinguistici.it/index.php/ssl/article/view/92}.

\bibitem[\citeproctext]{ref-Miller1891}
Miller, Thomas, ed. 1891. \emph{The {Old English} Version of {Bede}'s
Ecclesiastical History of the English People}. I, second half. For the
Early English Text Society.

\bibitem[\citeproctext]{ref-Minkova1996}
Minkova, Donka. 1996. {`Verse Structure as Evidence for Prosodic
Reconstruction in {Old English}'}. In \emph{{English Historical
Linguistics 1994}}, edited by Derek Britton, vol. 135. Current Issues in
Linguistic Theory. John Benjamins.
\url{https://doi.org/10.1075/cilt.135.04min}.

\bibitem[\citeproctext]{ref-Minkova2008}
Minkova, Donka. 2008. {`Prefixation and Stress in {Old English}'}.
\emph{Word Structure} 1 (1): 21--52.
\url{https://doi.org/10.3366/E1750124508000056}.

\bibitem[\citeproctext]{ref-MinkovaZhou2022}
Minkova, Donka, and Z. L. Zhou. 2022. {`Early Metrical and
Lexicographical Evidence for Functional Stress-Shifts'}. \emph{English
Language \& Linguistics} 26 (3): 533--58.
\url{https://doi.org/10.1017/S1360674322000144}.

\bibitem[\citeproctext]{ref-MinkovaZhou2023}
Minkova, Donka, and Z. L. Zhou. 2023. {`Nominal Compounds in {Old
English} Meter and Prosody'}. \emph{Journal of Germanic Linguistics} 35
(1): 69--95. \url{https://doi.org/10.1017/S1470542722000083}.

\bibitem[\citeproctext]{ref-Nielsen2001}
Nielsen, Hans Frede. 2001. {`{Frisian} and the Grouping of the Older
{Germanic} Languages'}. In \emph{Handbuch Des Friesischen / Handbook of
Frisian Studies}, edited by Horst Haider Munske. Max Niemeyer.

\bibitem[\citeproctext]{ref-Noreen1923}
Noreen, Adolf. 1923. \emph{Altisl{ä}ndische Und Altnorwegische Grammatik
(Laut- Und Flexionslehre)}. 4th edn. Niemeyer.

\bibitem[\citeproctext]{ref-Orel2003}
Orel, Vladimir. 2003. \emph{A Handbook of {Germanic} Etymology}. Brill.

\bibitem[\citeproctext]{ref-OsthoffBrugmann1881}
Osthoff, Hermann, and Karl Brugmann. 1878-\/-1910. \emph{{Morphologische
Untersuchungen auf dem Gebiete der indogermanischen Sprachen}}. 6 vols.
S. Hirzel.

\bibitem[\citeproctext]{ref-Pheifer1974}
Pheifer, J. D., ed. 1974. \emph{Old {English} Glosses in the
{É}pinal-{Erfurt} Glossary}. Clarendon Press.

\bibitem[\citeproctext]{ref-Pheifer1987}
Pheifer, J. D. 1987. {`Early {Anglo-Saxon} Glossaries and the School of
{Canterbury}'}. \emph{Anglo-Saxon England} 16: 17--44.
\url{https://doi.org/10.1017/S0263675100003847}.

\bibitem[\citeproctext]{ref-Quinlin1991}
Quinlin, Daniel P. 1991. {`The Accentuation and Development of {PGmc.}
*/Ga-/'}. \emph{American Journal of Germanic Linguistics and
Literatures} 3 (2): 145--59.
\url{https://doi.org/10.1017/S1040820700000688}.

\bibitem[\citeproctext]{ref-RieggerBogelWalkden2024}
Riegger, Chiara, Tina Bögel, and George Walkden. 2024. {`The Prosody of
the Verbal Prefix {ge-}: Historical and Experimental Evidence'}.
\emph{Proceedings of Interspeech 2024}, 2065--69.
\url{https://doi.org/10.21437/Interspeech.2024-289}.

\bibitem[\citeproctext]{ref-Ringe2017}
Ringe, Don. 2017. \emph{From {Proto-Indo-European} to {Proto-Germanic}}.
2nd edn. Vol. 1. A Linguistic History of English. Oxford University
Press.

\bibitem[\citeproctext]{ref-RingeTaylor2014}
Ringe, Don, and Ann Taylor. 2014. \emph{The Development of {Old
English}}. Vol. 2. A Linguistic History of English. Oxford University
Press.

\bibitem[\citeproctext]{ref-RuneSPforzen}
Runische Schriftlichkeit in den germanischen Sprachen (RuneS). n.d.
\emph{Pforzen, Buckle: Find 201}. Accessed 25 October 2026.
\url{https://www.runesdb.de/find/201}.

\bibitem[\citeproctext]{ref-Russom2001}
Russom, Geoffrey. 2001. {`Metrical Evidence for Subordinate Stress in
Old English'}. \emph{Journal of Germanic Linguistics} 13 (1): 39--64.
\url{https://doi.org/10.1017/S147054270100318X}.

\bibitem[\citeproctext]{ref-Samuels1949}
Samuels, M. L. 1949. {`The Ge-Prefix in the {Old English} Gloss to the
{Lindisfarne Gospels}'}. \emph{Transactions of the Philological Society}
48 (1): 62--116.
\url{https://doi.org/10.1111/j.1467-968X.1949.tb00696.x}.

\bibitem[\citeproctext]{ref-Sauer1992}
Sauer, Hans. 1992. \emph{Nominalkomposita Im Frühmittelenglischen: Mit
Ausblicken Auf Die Geschichte Der Englischen Nominalkomposition}. Vol.
30. Buchreihe Der Anglia. Max Niemeyer.

\bibitem[\citeproctext]{ref-Seebold1970}
Seebold, Elmar. 1970. \emph{{Vergleichendes und etymologisches
W{ö}rterbuch der germanischen starken Verben}}. Mouton.

\bibitem[\citeproctext]{ref-Seiler2017}
Seiler, Annina. 2017. {`Untangling the Strands: The Spelling of the
{É}pinal Glossary'}. In \emph{Exzerpieren -- Kompilieren -- Tradieren:
Transformationen Des Wissens Zwischen Sp{ä}tantike Und
Fr{ü}hmittelalter}, edited by Stephan Dusil, Gerald Schwedler, and
Raphael Schwitter, vol. 64. Millennium-Studien. De Gruyter.
\url{https://doi.org/10.1515/9783110516340-010}.

\bibitem[\citeproctext]{ref-Seiler2019}
Seiler, Annina. 2019. {`The Spelling of the {Erfurt} Glossary'}. In
\emph{Raum Und Sprache: Festschrift f{ü}r Elvira Glaser Zum 65.
Geburtstag}, edited by Andreas Nievergelt and Ludwig Rübekeil, vol. 66.
Germanistische Bibliothek. Universit{ä}tsverlag Winter.
\url{https://www.zora.uzh.ch/id/eprint/168922/}.

\bibitem[\citeproctext]{ref-SimsWilliams2018}
Sims-Williams, Patrick. 2018. {`Mechanising Historical Phonology'}.
\emph{Transactions of the Philological Society} 116 (3): 555--73.
\url{https://doi.org/10.1111/1467-968X.12138}.

\bibitem[\citeproctext]{ref-Skeat1874}
Skeat, Walter W., ed. 1874. \emph{The Gospel According to Saint Luke, in
Anglo-Saxon and Northumbrian Versions Synoptically Arranged}. Cambridge
University Press.
\url{https://archive.org/details/gospelaccording00skeagoog}.

\bibitem[\citeproctext]{ref-Stanley2022}
Stanley, Eric. 2022. \emph{Metricalities: Studies in Old English Verse}.
Edited by Andy Orchard. CLASP Ancillary Publications 2. CLASP.
\url{https://clasp.ell.ox.ac.uk/metricalities.pdf}.

\bibitem[\citeproctext]{ref-StauslandJohnsen2015}
Stausland Johnsen, Sverre. 2015. {`Vowel Reduction in Verbs in {King}
{Alfred's} \emph{Pastoral Care}'}. In \emph{Historical Linguistics 2013:
Selected Papers from the 21st International Conference on Historical
Linguistics, Oslo, 5--9 August 2013}, edited by Dag T. T. Haug. Current
Issues in Linguistic Theory 334. John Benjamins.
\url{https://doi.org/10.1075/cilt.334.02jon}.

\bibitem[\citeproctext]{ref-Stevenson1840}
Stevenson, Joseph, ed. 1840. \emph{Rituale Ecclesi{æ} Dunelmensis}.
Publications of the Surtees Society 10. J. B. Nichols; Son.
\url{https://archive.org/details/ritualeecclesi00cathuoft}.

\bibitem[\citeproctext]{ref-Stiles1985}
Stiles, Patrick V. 1985. {`The Fate of the Numeral {``4''} in Germanic
(1)'}. \emph{NOWELE: North-Western European Language Evolution} 6:
81--104. \url{https://doi.org/10.1075/nowele.6.06sti}.

\bibitem[\citeproctext]{ref-Stiles2012}
Stiles, Patrick V. 2012. {`Older {Runic} Evidence for {North-West
Germanic} {a-}Umlaut of {u} (and {``the Converse of {Polivanov's}
Law''})'}. In \emph{Laws and Rules in Indo-European}, edited by Philomen
Probert and Andreas Willi. Oxford University Press.
\url{https://doi.org/10.1093/acprof:oso/9780199609925.003.0004}.

\bibitem[\citeproctext]{ref-Stiles2017}
Stiles, Patrick V. 2017. {`On the Interpretation of Older Runic
\emph{Swestar} on the {Opedal} Stone'}. Unpublished manuscript.

\bibitem[\citeproctext]{ref-Streitberg1896}
Streitberg, Wilhelm. 1896. \emph{Urgermanische Grammatik}. Carl Winter.

\bibitem[\citeproctext]{ref-Strunk1904}
Strunk, William, Jr., ed. 1904. \emph{Juliana}. D. C. Heath.

\bibitem[\citeproctext]{ref-Sweet1871}
Sweet, Henry, ed. 1871. \emph{King {Alfred}'s {West-Saxon} Version of
{Gregory}'s Pastoral Care}. I. Early English Text Society, Original
Series 45. N. Tr{ü}bner.
\url{https://archive.org/details/kingalfredswests00greg}.

\bibitem[\citeproctext]{ref-Sweet1885}
Sweet, Henry, ed. 1885. \emph{The Oldest {English} Texts}. Early English
Text Society, Original Series 83. N. Tr{ü}bner.
\url{https://archive.org/details/oldestenglishtex83sweeuoft}.

\bibitem[\citeproctext]{ref-Sweet1953}
Sweet, Henry. 1953. \emph{An {Anglo-Saxon} Primer}. 9th edn. Clarendon
Press.

\bibitem[\citeproctext]{ref-TarsiLaker2026}
Tarsi, Matteo, and Stephen Laker. 2026. {`{``To Know''}, {``to Teach''},
and {``to Learn''} in {Germanic}'}. \emph{Journal of Germanic
Linguistics} 38: 40--58.
\url{https://doi.org/10.1017/S1470542726100178}.

\bibitem[\citeproctext]{ref-Toller1921}
Toller, T. Northcote. 1921. \emph{An {Anglo-Saxon} Dictionary:
Supplement}. Oxford University Press.

\bibitem[\citeproctext]{ref-deVaan2016}
Vaan, Michiel de. 2016. \emph{Etymologie: Tweern, Twijn}.
Neerlandistiek. Online tijdschrift voor taal- en letterkunde.
\url{https://neerlandistiek.nl/2016/10/etymologie-tweern-twijn/}.

\bibitem[\citeproctext]{ref-Versloot2017}
Versloot, Arjen P. 2017. {`Proto-{Germanic} *Ai in {North} and {West}
{Germanic}'}. \emph{Folia Linguistica Historica} 38: 281--324.
\url{https://doi.org/10.1515/flih-2017-0010}.

\bibitem[\citeproctext]{ref-Versloot2025}
Versloot, Arjen P. 2026. {`Reconstructing the Historical Phonology of
{Old} {English}'}. \emph{English Language and Linguistics} 30 (1):
103--41. \url{https://doi.org/10.1017/S1360674324000674}.

\bibitem[\citeproctext]{ref-GermanicSlavicBaltic2025}
Viredaz, Rémy. 2025. \emph{{Germanic}, {Slavic} and {Baltic}
{`Thousand'} Once More}. Unpublished manuscript (unfinished), Geneva, 20
October 2025.
\url{https://www.academia.edu/144462167/2025_Germanic_Slavic_and_Baltic_thousand_once_more_full_text_unfinished_}.

\bibitem[\citeproctext]{ref-Waxenberger2019}
Waxenberger, Gaby. 2019. {`Absolute Chronology of Early Sound Changes
Reflected in {Pre}-{Old} {English} Runic Inscriptions'}. \emph{NOWELE:
North-Western European Language Evolution} 72 (1): 60--77.
\url{https://doi.org/10.1075/nowele.00020.wax}.

\bibitem[\citeproctext]{ref-WrightWuelcker1884}
Wright, Thomas. 1884. \emph{Anglo-Saxon and Old English Vocabularies}.
2nd edn. Edited by Richard Paul Wülcker. Vol. 1. Tr{ü}bner.

\bibitem[\citeproctext]{ref-Thorhallsdottir1993}
Þórhallsdóttir, Guðrún. 1993. {`The Development of Intervocalic *j in
{P}roto-{G}ermanic'}. Ph.\{D\}. dissertation, Cornell University.

\end{CSLReferences}

\part*{Index verborum}

\addcontentsline{toc}{part}{Index verborum}

\printindex[iv]

\backmatter
\end{document}
